\documentclass[11pt,a4paper]{article}
\usepackage[T1]{fontenc}
\usepackage[utf8]{inputenc}
\usepackage{lmodern}
\usepackage[a4paper,margin=27mm]{geometry}
\usepackage{amsmath,amssymb,amsthm,mathtools,mathrsfs}
\usepackage{booktabs,longtable,array,microtype}
\newcolumntype{L}[1]{>{\raggedright\arraybackslash}p{#1}}
\usepackage{xcolor}
\usepackage[colorlinks=true,linkcolor=blue!45!black,citecolor=blue!45!black,urlcolor=blue!45!black,breaklinks=true]{hyperref}
\usepackage{xurl}
\usepackage{enumitem}
\usepackage{etoolbox}
\numberwithin{equation}{section}
\newtheorem{theorem}{Theorem}[section]
\newtheorem{lemma}[theorem]{Lemma}
\newtheorem{proposition}[theorem]{Proposition}
\newtheorem{corollary}[theorem]{Corollary}
\theoremstyle{remark}
\newtheorem{remark}[theorem]{Remark}
\newtheorem{example}[theorem]{Example}
\setlength{\parskip}{3pt}
\setlength{\emergencystretch}{3em}
\setcounter{tocdepth}{1}
\makeatletter
\renewcommand{\@pnumwidth}{3em}
\renewcommand{\@tocrmarg}{4em}
\patchcmd{\l@section}{1.5em}{2.2em}{}{\PackageError{unified}{TOC number-width patch failed}{}}
\patchcmd{\l@section}{1.0em}{0.9em}{}{\PackageError{unified}{TOC spacing patch failed}{}}
\makeatother
\hypersetup{pdftitle={Arithmetic and Observation: Classifications, Recovery, and Transcendence Barriers},pdfauthor={Keio}}

\begin{document}
\title{Arithmetic and Observation:\\Classifications, Recovery, and Transcendence Barriers\\[6pt]\large Unified Edition of Twelve Successor Manuscripts}
\author{KEIO\\\small AI-assisted research manuscript}
\date{17 October 2026}
\maketitle
\begin{abstract}
This unified manuscript brings together twelve successive studies of
arithmetic observation, optimization, inverse recovery, and the transfer
of analytic constructions to questions of irrationality and transcendence.
The development begins with observation boundaries and optimal values,
passes through structural recovery and effective logical classification,
and ends with arithmetic certificates, entropy thresholds, primitive
integer forms, and exact quadratic exceptions. The twelve parts retain
their substantive arguments, hypotheses, proofs, examples, target registers,
and limitations. A common introduction, continuous numbering, linked
cross-references, and one consolidated bibliography provide a single
reading structure. An editorial roadmap identifies later strengthenings
and distinguishes model classifications from unresolved transfers to named
constants. The earlier 223-page treatise and the subsequent 447-page
compilation are not reproduced. This edition consolidates the twelve
manuscripts written after those volumes; it does not assert a new
resolution of their outstanding named-constant conjectures.
\end{abstract}

\noindent\textbf{Scope of this edition.}
The source sequence begins with \emph{Arithmetic of Observation Boundaries}
and ends with \emph{Arithmetic Rigidity, Quadratic Exceptions, and Constant
Transfer}. All twelve successor manuscripts are represented. The two
earlier volumes remain cited where they supply mathematical inputs, but
neither is included as a part of this document. The TEX source embeds the
entire present text and bibliography and needs no external chapter files.

\clearpage
\begingroup\small
\setlength{\parskip}{0pt}
\tableofcontents
\endgroup
\clearpage

\section*{Editorial synthesis and reading map}
\addcontentsline{toc}{section}{Editorial synthesis and reading map}
\label{merge:synthesis}

The common question is how much mathematical structure can be recovered
from observations, and which additional information is required before
that recovery decides an arithmetic property. The first parts develop
arithmetic classifications for boundaries and optimal values. The middle
parts test how those mechanisms interact with inverse problems, logic,
and exact finite certificates. The last parts separate analytic
convergence, arithmetic cancellation, and genuine value classifications.
These developments are interdependent, so the original chronological
order is retained. Moving isolated proofs into new thematic chapters
would obscure their local definitions and inherited assumptions.

\subsection*{Twelve parts, one argument sequence}
\begingroup\small
\renewcommand{\arraystretch}{1.12}
\begin{longtable}{@{}L{0.08\textwidth}L{0.35\textwidth}L{0.51\textwidth}@{}}
\toprule
Part & Main subject & Role in the combined development\\
\midrule\endfirsthead
\toprule Part & Main subject & Role in the combined development\\\midrule\endhead
\bottomrule\endfoot
\ref{merge:part01} & Observation boundaries & Arithmetic tests for extraction, stopping, and spectral boundaries; the first kernel and height obstructions.\\
\ref{merge:part02} & Optimal values & Transfer from boundary locations to objective values; mixed-sign spectra and stronger positive-kernel obstructions.\\
\ref{merge:part03} & Independence and entropy optima & Cubic spectra, families of stopping values, integer extraction rates, entropy optimization, and genuinely reduced denominators.\\
\ref{merge:part04} & Rigidity of optimization & The full planar stopping family, entropy-game ranges, differential complexity, and finite algebraic covers.\\
\ref{merge:part05} & Structural and adaptive recovery & Unrestricted mixture identification, finite spectral complexity, robust certificates, and exact adaptive thresholds.\\
\ref{merge:part06} & Logic and effective transfer & The fifty-seven requested logical problems, with distinctions between solved subcases, external announcements, and failed reductions.\\
\ref{merge:part07} & Arithmetic beyond convergence & Effective logarithmic-root and quantile classification, finite cyclotomic certificates, and limits of arithmetic transfer.\\
\ref{merge:part08} & Arithmetic certificates & Inverse realization, moment matching, entropy tests, Mahler relation ideals, and growing arithmetic kernels.\\
\ref{merge:part09} & Thresholds and primitive forms & Sharp alphabet and degree costs, exact Laguerre content, Wilson approximants, and formal Gamma isolation.\\
\ref{merge:part10} & Sparse certificates & Information loss, a four-atom entropy example, rare cancellation, curvature counts, and nonlinear Lambert independence.\\
\ref{merge:part11} & Spectral entropy and transfer & Totally-real entropy thresholds, a third Lambert jet, full-degree positive pencils, and parameter laws.\\
\ref{merge:part12} & Quadratic exceptions & Exact relative and conditional entropy classifications, the cubic ternary obstruction, actual beta-period separation, fixed-parameter content bounds, and repaired Catalan convergence.\\
\end{longtable}
\endgroup

\subsection*{How later parts change the earlier frontier}
Several earlier conclusions remain valid while losing their status as the
strongest statement available in the sequence. The following crosswalk
records the main changes; it does not count each intermediate formulation
as a separate discovery.

\paragraph{Spectra, stopping, and extraction.}
The one-sided rank-two spectral classification in Part~\ref{merge:part01}
is extended to all sign patterns in Part~\ref{merge:part02}; the real-cubic
branch becomes unconditional in Part~\ref{merge:part03}. The unrestricted
higher-rank conditional framework retains its hypotheses. The stopping
branch progresses from boundary arithmetic to objective values and then
to independence of the full planar even-payoff family in
Part~\ref{merge:part04}. An optimizer and its objective value are different
numbers. Likewise, the extraction branch's geometric results for real
shape parameters do not automatically extend arithmetic conclusions
proved only for rational or integer parameters.

\paragraph{Entropy thresholds answer different questions.}
In this roadmap, natural entropy is
$H(p)=-\sum_i p_i\log p_i$, bit entropy is $H_2(p)=H(p)/\log2$,
and effective rank is $E(p)=e^{H(p)}$. Their arithmetic classifications
must not be interchanged; retained parts specify their local notation.
Part~\ref{merge:part09} proves a three-symbol minimum for algebraic
irrational bit entropy among positive algebraic probability vectors.
Part~\ref{merge:part10}
constructs four irrational algebraic probabilities with rational bit
entropy, and Part~\ref{merge:part11} makes four the exact minimum when
all entries are irrational and totally real. The unrestricted ternary question is still open
within these arguments. Part~\ref{merge:part12} excludes every cubic
probability field, leaving degree at least four; when the entropy
denominator is odd, the remaining degree is odd and at least five.
Its separate three-symbol threshold concerns algebraic exponential
relative entropy for nonrational totally positive algebraic laws,
with a positive rational reference and a nonzero real algebraic exponent.
It is not a revision of the totally-real threshold four.

\paragraph{Positivity of conjugates matters.}
Totally real means that every conjugate is real. Totally positive means
that every conjugate is positive. A displayed positive probability need
not be totally positive. The relative-entropy classification in
Part~\ref{merge:part12} uses positivity of the conjugate laws and has
genuine quadratic exceptions. This is compatible with the effective-rank
classification in Part~\ref{merge:part11}; the two exponential expressions
have different factors and permit different cancellations.

\paragraph{Primitive content and rare indices.}
Fixed-kernel estimates, standard denominator estimates, and estimates
after actual primitive reduction are different assertions. The primitive
normalization discussion in Part~\ref{merge:part03} prevents an overbroad
reading of earlier positivity obstructions. For Laguerre forms, the
$O(\log X)$ exceptional count in Part~\ref{merge:part10} concerns
$\log g_n\ge\varepsilon n\log n$. The later
$O_{A,\varepsilon}((\log X)^4)$ count in Part~\ref{merge:part12}
concerns the smaller threshold $\log g_n(A)\ge\varepsilon n$ at fixed
$A$. Both statements are retained. The parameter-residue law of
Part~\ref{merge:part11} and the fixed-parameter index law of
Part~\ref{merge:part12} average different variables. The latter
distribution is for a prime-support upper bound, not for the primitive
content itself. None of these density or counting statements removes
every possible infinite exceptional subsequence.

\paragraph{Catalan approximation spaces.}
The determinant constructions develop through indefinite mixed pencils,
positive full-degree pencils, polynomial-invariant subspaces, and full
paired polynomial spaces. Their conclusions are not interchangeable.
Part~\ref{merge:part12} proves a uniform root gap for the restricted
invariant space and recovers convergence by restoring the omitted sector.
The repaired convergence still lacks the primitive-height estimate needed
for an arithmetic conclusion. Earlier positive interval constructions
are distinct families and are not contradicted by the invariant-space gap.

\paragraph{Lambert functions and Gamma periods.}
Part~\ref{merge:part11} extends the earlier Lambert independence results
to three jets together with a prime series, under its stated derivative
condition. It does not prove independence for every finite jet family or
specialize functional independence to a particular numerical value.
The final determinant-budget optimization concerns a specified sufficient
clearing bound, not the actual arithmetic degree of a Lambert constant.
The Gamma branch moves from a formal identity calculus to an actual
fourteen-dimensional beta-period span in Part~\ref{merge:part12} using
classical transcendence theory. The twelve displayed transcendental
ratios are classical-theorem consequences, not twelve newly solved
open conjectures. Linear independence of that span is not algebraic
independence of its entries.

\subsection*{Editorial conventions and retained limitations}
The mathematical text is preserved at the source-part level. Individual
title pages, abstracts, tables of contents, and bibliographies have been
replaced by common frontmatter and a consolidated bibliography. Labels
and citations have unique internal namespaces, and verification appendices
continue the global section numbering. Repeated definitions and short
supporting arguments are retained when they keep a proof intelligible.
Only explicitly recorded typographic corrections and a terminology
clarification are made; this is not a new full proof audit.

Statements such as ``new'', ``previous'', and ``unresolved here'' inside
a source part refer to its place in the development. The forward notes
identify subsequent progress. Target counts in different registers overlap
and must not be added as counts of discoveries. External announcements
retain their source dates and qualifications; this compilation does not
independently update every literature-status statement or replay external
formal proof projects. Earlier-volume inputs and other cited theorems
remain mathematical dependencies even though the present TEX is a single
independently compilable file.

The final part records zero newly closed previously open named-constant
conjectures among its sixteen targets. The exact model classifications,
inverse constructions, and obstruction theorems elsewhere in the volume
have their stated scopes. Historical priority and journal-level
significance are separate questions from assembling these arguments.

\clearpage
\part{Arithmetic of Observation Boundaries}
\label{merge:part01}
\begingroup

% BEGIN SOURCE: KEIO_Arithmetic_of_Observation_Boundaries.tex

\noindent\textbf{Editorial integration note.} 
This is the first of the twelve included successors. The rank-two spectral theorem is strengthened in Part~\ref{merge:part02}. The positivity obstructions here retain their stated denominator normalization; the primitive-reduction distinction is developed in Part~\ref{merge:part03}.


\section{Scope, provenance, and the arithmetic change of viewpoint}
The starting point is the pair of working manuscripts \emph{Observation and Recovery in Arithmetic and Probability: Prime Flows, Rare Events, and Infinite Mixtures} (223 pages) \cite{p01:root:original} and \emph{Observation and Recovery: Geometry, Mixtures, and Spectral Inversion} (447 pages) \cite{p01:bec17:Previous}. Their recovery theorems identify parameters from observations, while their boundary theorems define distinguished constants by positivity, contact, or optimization. The present manuscript asks for the arithmetic nature of those constants. It does not concatenate the earlier manuscripts, and it does not treat their recovery statements as transcendence theorems.

There are three logically different tasks. An inverse theorem proves that observations determine a number. A numerical method approximates that number with certified error. A Diophantine argument proves that it cannot be rational or algebraic. The third task needs arithmetic information, such as denominator growth, a differential specialization theorem, or a transcendence theorem for exponentiation. Neither uniqueness nor fast numerical convergence supplies that information by itself.

\subsection{What is proved, what is inherited, and what is not claimed}
The proofs below establish transfer statements under their displayed hypotheses. The deep transcendence inputs are classical and are cited explicitly. A new application to a boundary from the preceding working manuscripts is described as a \emph{transfer result}, not as a resolution of a historically documented conjecture. A theorem about a newly formulated class of positive kernels or spectra is described as a \emph{class theorem}. A previously known classification is used only as a calibration case. A numerical value is illustrative and never used as an arithmetic proof.

In particular, proving the irrationality of an entire family of newly studied capacities is different from proving the irrationality of a famous pre-existing constant. Both can be mathematically useful, but they should not be counted together as independent famous problems closed. Rational-affine variants of a single boundary constant are also counted as one family.

The statements here are supported by written proofs and cross-checks within this research workflow. They are not presented as externally refereed results, and no claim of Annals-level novelty follows from the length of this manuscript or the number of examples.

\subsection{The two transfer mechanisms}
The first mechanism starts with a boundary equation
\[
  U(\xi)f(\xi)+V(\xi)f'(\xi)=0,
\]
where $f$ is an $E$-function and $U,V$ are algebraic rational functions. If $f$ and $f'$ are linearly independent over $\overline{\mathbb Q}(z)$, Beukers' theorem excludes every nonzero algebraic ordinary point $\xi$ satisfying this equation. The arithmetic of the \emph{parameters} is crucial: the standard Kummer $E$-function input used here requires rational hypergeometric parameters. This explains why some contact locations are proved transcendental, while an implicitly varying capacity parameter is proved only irrational.

The second mechanism starts with a critical value of a finite positive spectrum. The critical-point and critical-value equations give two algebraic equations on a two-dimensional torus. A one-sided sign condition rules out repeated-factor degeneracy. An irrational-weight derivation has no invariant nonconstant Laurent irreducible factors, so the two equations force algebraic torus coordinates. Gelfond--Schneider then excludes all positive points except the explicitly identifiable point $x=1$.

\subsection{Notation and foundational inputs}
We write $\overline{\mathbb Q}$ for the field of algebraic numbers, $\operatorname{trdeg}$ for transcendence degree, and $M(a;b;z)={}_1F_1(a;b;z)$ for Kummer's function. All real powers of positive numbers use the real logarithm. Algebraic independence of an infinite family means algebraic independence of every finite subfamily. Boundary existence and uniqueness are proved below where needed, or tied explicitly to a displayed analytic characterization from the earlier manuscripts. A literature status is dated to the searches completed for this manuscript; an unverified recent proof claim is not counted as a resolution.


% Integrated source: bessel_results.tex
\section{An arithmetic obstruction on the gamma-extraction boundary}
\label{p01:bec17:section}

The gamma-extraction curve in the preceding manuscript produces a
family of implicitly defined constants.  Its arithmetic nature was not
determined there.  We prove that every rational point of its parameter
axis strictly above $1/2$ gives an irrational boundary constant.  The
argument uses a classical specialization theorem for $E$-functions,
not a new transcendence theorem.  The new application is to the
probabilistic boundary and its exact contact relation.  In particular,
the result does not settle the irrationality of any famous classical
constant whose status was previously open.

\subsection{The inherited boundary and the precise arithmetic question}

For $d>1/2$, the previous manuscript \cite{p01:bec17:Previous} constructed
a unique number
\begin{equation}\label{p01:bec17:range}
 A(d)\in\left(\max(0,1-d),\frac1{4d}\right).
\end{equation}
Put $\eta=a+d-1$ and, for $\eta>0$, define
\begin{equation}\label{p01:bec17:B}
 B_{a,d}(z)={}_1F_1(a-1;\eta;z)
       +\frac{z}{2\eta}{}_1F_1(a;\eta+1;z).
\end{equation}
The defining equation, with uniqueness in \eqref{p01:bec17:range}, is
\begin{equation}\label{p01:bec17:contact}
                B_{A(d),d}\bigl(2(d-A(d))\bigr)=0.
\end{equation}
This equation comes from the inverse Laplace transform of
$(u+1)u^{-a}(u+2)^{-d}$.  At the positivity boundary its density has
a unique double zero at
\begin{equation}\label{p01:bec17:xstar}
                         x_*(d)=d-A(d).
\end{equation}
The exact formula for the contact point is essential to the arithmetic
argument.  A generic unspecified zero of an $E$-function would not
supply it.

On the first excess-power McKay line with rational parameter
$\kappa>1/2$, the maximum endpoint-rate gamma shape is
$h_*(\kappa)=\kappa-A(\kappa)$.  Thus $h_*(d)$ and $x_*(d)$ are
the same number in this normalization.  The ID-remainder capacity is
$d-1/(4d)$, and its deficit from the unrestricted capacity is
$1/(4d)-A(d)$.  These rational-affine variants must not be counted
as independent arithmetic discoveries.

\subsection{The classical specialization input, with its hypotheses checked}

We use the following precise consequence of Beukers' refinement of the
Siegel--Shidlovskii theorem \cite[Corollary~1.4]{p01:bec17:Beukers}.
If a vector of $E$-functions $F$ satisfies $F'=M(z)F$, with
$M\in\overline{\mathbb Q}(z)^{r\times r}$, and its coordinates are
linearly independent over $\overline{\mathbb Q}(z)$, then their
values at a nonzero algebraic point $\xi$ are linearly independent
over $\overline{\mathbb Q}$ whenever $\xi$ is not a pole of $M$.
In particular this is a linear-independence theorem over the full
field of algebraic numbers, not just over the rationals.

\begin{lemma}[Rational Kummer parameters give strict $E$-functions]
\label{p01:bec17:E}
Let $\alpha,\beta\in\mathbb Q$, with
$\beta\notin\{0,-1,-2,\ldots\}$.  The function
$y(z)={}_1F_1(\alpha;\beta;z)$ and its derivative are
$E$-functions.  Their first-order system is
\begin{equation}\label{p01:bec17:system}
 \frac{d}{dz}\begin{pmatrix}y\\y'\end{pmatrix}
 =\begin{pmatrix}0&1\\ \alpha/z&1-\beta/z\end{pmatrix}
                    \begin{pmatrix}y\\y'\end{pmatrix}.
\end{equation}
Its only finite possible singularity is zero.
\end{lemma}
\begin{proof}
The polynomial case is immediate, so suppose that no numerator
Pochhammer symbol vanishes.  Write
\[
 y(z)=\sum_{n\geq0}u_n\frac{z^n}{n!},\qquad
 u_n=\frac{(\alpha)_n}{(\beta)_n}.
\]
The recurrence for $u_n$ gives $|u_n|\leq C^{n+1}$ for a fixed
$C$, so the series is entire of exponential type.  We give the
denominator check explicitly.  Write $\alpha=U/V$, $\beta=W/Z$
in lowest terms with positive $V,Z$.  For a prime $p$ not dividing
$VZ$, and for each power $p^k$, the counts of indices $0\leq j<n$
for which $p^k$ divides $U+Vj$ and $W+Zj$ differ by at most one:
each is the count of one residue class modulo $p^k$ in an interval
of length $n$.  All relevant nonzero factors have absolute value at
most $C_1(n+1)$.  Therefore the denominator valuation of every
$u_j$, $j\leq n$, at such a prime is at most
$\lfloor\log_p(C_1(n+1))\rfloor$.

For the finitely many primes dividing $VZ$, the denominator valuation
is $O(n)$: the explicit factor $(Z/V)^j$ contributes linearly,
and summing the residue-class counts of divisibility over $p^k$
also contributes $O(n)$.  Consequently one common denominator of
$u_0,\ldots,u_n$ divides a fixed-prime factor of size $\exp(O(n))$
times $\operatorname{lcm}(1,\ldots,\lfloor C_1(n+1)\rfloor)$.
The elementary Chebyshev bound for this least common multiple is
$\exp(O(n))$.  Thus the coefficients satisfy the strict exponential
height and denominator bounds in the definition of an $E$-function.
The coefficients of $y'$ are $u_{n+1}$ and satisfy the same bounds.
Finally Kummer's equation
$zy''+(\beta-z)y'-\alpha y=0$ proves
\eqref{p01:bec17:system}; see also \cite[Section~13.2]{p01:bec17:DLMF}.
\end{proof}

The next check avoids any irreducibility assumption about the whole
Kummer differential equation.  That distinction matters when $d$ is
a positive integer.

\begin{lemma}[The analytic solution is not hyperexponential]
\label{p01:bec17:independence}
Suppose $\beta\notin\{0,-1,-2,\ldots\}$ and neither
$\alpha$ nor $\beta-\alpha$ is a nonpositive integer.  Then
${}_1F_1(\alpha;\beta;z)$ and its derivative are linearly independent
over $\mathbb C(z)$.
\end{lemma}
\begin{proof}
Otherwise $y'/y$ would be rational.  Since $y$ is entire and
$y(0)=1$, its finite zeros would be exactly the finitely many poles
of this rational logarithmic derivative, with positive integer
residues.  Let $P$ be the polynomial containing those zeros with
their multiplicities.  The logarithmic derivative of $y/P$ is an
entire rational function, hence a polynomial.  It follows that
$y=e^{Q(z)}P(z)$ for a polynomial $Q$, after absorbing a nonzero
constant into $P$.

The coefficient ratio gives exponential-type growth of $y$, also
for the complex parameters of this lemma.  Thus $Q$ has degree at
most one, and we may write $y=e^{cz}P$ with $n=\deg P\geq0$.
Substitute this expression into Kummer's equation.  The coefficient
of degree $n+1$, after removing $e^{cz}$, is proportional to
$c^2-c$, so $c=0$ or $c=1$.  If $c=0$, the coefficient of degree
$n$ forces $\alpha=-n$.  If $c=1$, the transformed equation is
\[
 zP''+(\beta+z)P'+(\beta-\alpha)P=0,
\]
whose leading coefficient forces $\beta-\alpha=-n$.
Both alternatives contradict the assumptions.
\end{proof}

\begin{proposition}[A standard $E$-function consequence for Robin zeros]
\label{p01:bec17:Robin}
Let $\alpha,\beta\in\mathbb Q$ satisfy the hypotheses of
Lemma~\ref{p01:bec17:independence}.  At every nonzero algebraic $\xi$,
the two numbers
\[
 {}_1F_1(\alpha;\beta;\xi),\qquad
 \left.\frac{d}{dz}{}_1F_1(\alpha;\beta;z)\right|_{z=\xi}
\]
are linearly independent over $\overline{\mathbb Q}$.
Their quotient is therefore transcendental.  If
$P,Q\in\overline{\mathbb Q}[z]$ are relatively prime and not both
zero, every nonzero zero of $P(z)y(z)+Q(z)y'(z)$ is transcendental.
This proposition is a classical-theorem application, not a claim of
a new general $E$-function theorem.
\end{proposition}
\begin{proof}
Apply the stated Beukers theorem to \eqref{p01:bec17:system}, using the
two preceding lemmas.  Every nonzero algebraic point is ordinary
for this system.  Linear independence also ensures that neither
value is zero.  An algebraic quotient would give a forbidden linear
relation.  At an algebraic zero of $Py+Qy'$, coprimality ensures
that $P(\xi),Q(\xi)$ are not both zero, giving the same contradiction.
\end{proof}

\subsection{Irrationality of every rational-parameter curved boundary}

\begin{theorem}[The gamma boundary contains no rational point]
\label{p01:bec17:main}
For every rational $d>1/2$, the number $A(d)$ is irrational.
Consequently the maximum endpoint-rate gamma shape
$h_*(d)=x_*(d)=d-A(d)$ and the capacity gain
$1/(4d)-A(d)$ are irrational as well.  More precisely,
\[
 \mathbb Q(A(d))=\mathbb Q(h_*(d))
              =\mathbb Q\bigl(1/(4d)-A(d)\bigr).
\]
The curved boundary $a=A(d)$ has no point with both coordinates
rational.  In contrast, the linear boundary $A(d)=1-d$ for
$d\leq1/2$ has rational value at every rational $d$.
\end{theorem}
\begin{proof}
Suppose $a=A(d)$ were rational.  By \eqref{p01:bec17:range},
$0<a<1/2$, $\eta=a+d-1>0$, and $d-a>0$.  Set
\[
 \alpha=a-1,\qquad \beta=\eta,\qquad p=1-a.
\]
Then $\alpha\in(-1,-1/2)$ is not a nonpositive integer,
$\beta>0$, and $\beta-\alpha=d>0$ is not a nonpositive integer.
The derivative identity for Kummer's function gives
\begin{equation}\label{p01:bec17:Robin-B}
 B_{a,d}(z)=y(z)-\frac{z}{2p}y'(z).
\end{equation}
The coefficients in this linear combination are relatively prime
polynomials with rational coefficients.  The contact relation
\eqref{p01:bec17:contact} says that it vanishes at
$z_*=2(d-a)>0$, a rational number.  This contradicts
Proposition~\ref{p01:bec17:Robin}.  The remaining assertions are
immediate rational-affine changes of this one irrational number,
and the inherited formula on the linear face.
\end{proof}

For example, the theorem settles the following three members of the
newly defined family, with exact inherited enclosures:
\[
\begin{array}{c|c|c}
d&A(d)&h_*(d)=x_*(d)\\ \hline
1&\text{irrational in }(0,1/4)&\text{irrational in }(3/4,1)\\
3/2&\text{irrational in }(0,1/6)&\text{irrational in }(4/3,3/2)\\
2&\text{irrational in }(0,1/8)&\text{irrational in }(15/8,2)
\end{array}
\]
These rows illustrate one uniform theorem; they are not three
independent open-problem closures.

For integer $d=m\geq1$ the Kummer equation can indeed be reducible.
Its second Frobenius solution is
\[
 z^{1-\eta}{}_1F_1(1-m;2-\eta;z),
\]
a power times a polynomial of degree $m-1$.  When $a$ is rational
in the interval used above, $\eta=a+m-1$ is not an integer, so
this solution is not the analytic normalized solution $y$.
Lemma~\ref{p01:bec17:independence} concerns that particular analytic
solution and remains valid.  Thus $d=1$ and $d=2$ are not omitted
exceptional cases.  The excluded endpoint $a=0$ also matters:
for $d>1$, \eqref{p01:bec17:B} becomes
$B_{0,d}(z)=1-z/\{2(d-1)\}$, whose positive zero is algebraic
when $d$ is rational.  The strict interior assumptions cannot be
discarded.

\subsection{An exact rational-input decision procedure}

The arithmetic obstruction has an effective consequence for the
factorization problem.  This is grouped with Theorem~\ref{p01:bec17:main},
not counted as a second independent irrationality theorem.

\begin{corollary}[Finite rational arithmetic decides rational factor pairs]
\label{p01:bec17:decision}
For rational $\kappa>1/2$ and nonnegative rational shapes
$h_-,h_+$, simultaneous gamma extraction at the two original rates
can be decided by a terminating algorithm using rational arithmetic.
No equality oracle for a special-function value is needed.  The
algorithm makes no uniform complexity or irrationality-measure claim.
\end{corollary}
\begin{proof}
Put $a=\kappa-h_-$ and $d=\kappa-h_+$.  The previous exact diagram
says that feasibility is $a\geq A(d)$.  Reject if $a\leq0$ or
$a+d<1$.  If $d\leq1/2$, accept exactly when $a+d\geq1$.
For $d>1/2$, reject $a+d=1$ and accept $a\geq1/(4d)$.
The only remaining case has
\[
 \max(0,1-d)<a<1/(4d)<1/2.
\]
Set $p=1-a\in(1/2,1)$, $z=2(d-a)>0$, and $\eta=a+d-1>0$.
The one-point criterion reduces feasibility to $B_{a,d}(z)\geq0$.
Here all parameters are rational, and
\begin{equation}\label{p01:bec17:series}
 B_{a,d}(z)=\sum_{n\geq0}b_n,\qquad
 b_n=\left(1-\frac n{2p}\right)
             \frac{(-p)_n}{(\eta)_n}\frac{z^n}{n!}.
\end{equation}
All terms are rational; $b_n>0$ for every $n\geq2$.
For $n\geq3$,
\[
 \frac{b_{n+1}}{b_n}
  =\frac{n+1-2p}{n-2p}\frac{n-p}{n+\eta}\frac z{n+1}
  \leq\frac{2z}{n+1}.
\]
Therefore, for $N\geq\max(3,\lceil4z\rceil)$,
\begin{equation}\label{p01:bec17:certificate}
 \sum_{n=0}^N b_n
 \ \leq B_{a,d}(z)\ \leq
 \sum_{n=0}^N b_n+2b_{N+1}.
\end{equation}
These rational intervals shrink to the value.  Proposition
\ref{p01:bec17:Robin} ensures that the value is not zero.  Increasing
$N$ therefore eventually certifies its sign.  A positive sign
accepts the factor pair and a negative sign rejects it.  This proves
termination and exact correctness, including the cases of negative
$d$ and atomic remainders already handled by rational inequalities.
\end{proof}

The same sign certificates give arbitrarily tight rational brackets
for $A(d)$ when $d>1/2$ is rational.  Start with the rational upper
bound $1/(4d)$ and enumerate rational points decreasing to
$\max(0,1-d)$ until a negative one-point value is certified; then
bisect.  None of the rational midpoints is a boundary point, so every
sign test terminates.  This is an exact computability statement;
it gives no effective bound on how close $A(d)$ can be to rationals
of a given denominator.

\paragraph{What this does not prove.}
Transcendence of $A(1)$, $A(3/2)$, $A(2)$, or of the general
rational-parameter boundary remains unresolved here.  If one merely
assumes $A(d)$ algebraic, the Kummer parameters need not be rational;
Lemma~\ref{p01:bec17:E} no longer applies.  No assertion that such a
function is an $E$-function has been justified.  Also $A(d)$ is a
root in a parameter of a family, not a value of one fixed rational-
parameter $E$-function at a rational point, so irrationality-exponent
theorems for those values cannot simply be applied to $A(d)$.
The same obstruction prevents importing this proof to the entire
intermediate-rate capacity curve, whose density contact location is
not known to be algebraic under a rational-shape assumption.
No result here concerns Euler's constant, Catalan's constant, or an
individual odd zeta value beyond the results already known for them.


% Integrated source: gaussian_results.tex
\section{Arithmetic of stopping boundaries and Bessel obstruction constants}
\label{p01:ga17:section}

This section investigates three targets.  The first is the arithmetic
classification of the power-payoff stopping boundary for a Bessel bridge,
including the reflected Brownian bridge and Brownian excursion.  The
second concerns the locations and values of the Bessel-ratio extrema
that occurred in the earlier HCM obstruction analysis.  Both admit
complete classifications at rational parameters.  The third is the
signed Brownian-bridge stopping constant; the method does not decide
its arithmetic nature.  The stopping problems themselves are established
in the literature \cite{p01:ga17:EW,p01:ga17:HL}.  We have not located a primary
source posing the arithmetic classifications below as open conjectures.
They are therefore recorded as arithmetic transfers to existing
constants, not as closures of named historical conjectures.  The
transcendence input is classical; no new version of the
Siegel--Shidlovskii theorem is claimed.

\subsection{The arithmetic input and a logarithmic-derivative mechanism}

We use the ordinary-point linear specialization theorem stated in
Section~\ref{p01:bec17:section}, in the precise form of Beukers'
Corollary~1.4 \cite{p01:bec17:Beukers}. In particular, there is no
unspecified finite exceptional set at ordinary algebraic points.

Here an $E$-function is an entire function
$f(z)=\sum_{n\geq0}a_nz^n/n!$ with algebraic coefficients, satisfying
a linear differential equation over $\overline{\mathbb Q}(z)$, whose
coefficient heights and common denominators grow at most exponentially.
All functions used below have rational coefficients.  For completeness,
the elementary denominator estimate needed here is supplied next.

\begin{lemma}[Rational hypergeometric coefficients]
\label{p01:ga17:coeff}
For positive rational $a,b$, the numbers
$r_n=(a)_n/(b)_n$ have at most exponential absolute growth and a common
denominator for $r_0,\ldots,r_n$ bounded by $C^{n+1}$, with $C$
depending only on $a,b$.  Consequently
\[
 M(a;b;z)=\sum_{n=0}^{\infty}\frac{(a)_n}{(b)_n}\frac{z^n}{n!}
\]
and its derivative are $E$-functions.  For rational $\mu>-1$, the
normalized modified Bessel function
\[
 H_\mu(z)=\sum_{n=0}^{\infty}
       \frac{(z^2/4)^n}{n!(\mu+1)_n}
\]
and its derivative are also $E$-functions.
\end{lemma}
\begin{proof}
The Kummer assertions and their denominator bounds are
Lemma~\ref{p01:bec17:E}. Their equation is
\begin{equation}\label{p01:ga17:kummerode}
 zf''+(b-z)f'-af=0.
\end{equation}
For the Bessel function the odd exponential coefficients vanish,
whereas the even ones are
\[
 \frac{(2n)!}{4^nn!(\mu+1)_n}
       =\frac{(1/2)_n}{(\mu+1)_n}.
\]
The same rational Pochhammer-ratio estimate applies, and
\begin{equation}\label{p01:ga17:besselode}
 H_\mu''+\frac{2\mu+1}{z}H_\mu'-H_\mu=0
\end{equation}
supplies the differential equation. Derivatives shift the
exponential coefficients and preserve the bounds.
\end{proof}

\begin{lemma}[A classical specialization mechanism]
\label{p01:ga17:mechanism}
Suppose $f,f'$ are $E$-functions satisfying a two-dimensional
first-order system over $\overline{\mathbb Q}(z)$, and $f'/f$ is
not rational.  At every nonzero algebraic point $\xi$ where the
system is regular, $f(\xi)\ne0$ and $f'(\xi)/f(\xi)$ is
transcendental.  Hence if
\[
 Q\bigl(x,f'(x)/f(x)\bigr)=0,
 \qquad Q\in\overline{\mathbb Q}[X,Y],
\]
and $Q(x,Y)$ is not the zero polynomial, then such a nonzero
nonsingular real $x$ cannot be algebraic.
\end{lemma}
\begin{proof}
A rational linear relation between $f$ and $f'$ would make $f'/f$
rational.  Thus Beukers' corollary applies to the pair.  A zero of
$f(\xi)$ or an algebraic value of the quotient would give a
nontrivial algebraic linear relation between the two values.  In
the last assertion, an algebraic $x$ would make the quotient a root
of a nonzero polynomial with algebraic coefficients, a contradiction.
\end{proof}

We will repeatedly use an elementary analytic observation.  If an
entire function of exponential type has rational logarithmic
derivative, it has the form $e^{\lambda z}P(z)$ for a polynomial
$P$.  To verify this, the poles of the logarithmic derivative are
exactly the finitely many zeros, with positive integral residues.
After removing these polynomial factors, integration of the remaining
polynomial logarithmic derivative gives $e^{Q(z)}$.  Exponential
type forces $\deg Q\leq1$.  This observation concerns the particular
entire solution, and does not assume irreducibility of its differential
equation.

\subsection{Target I: an exact arithmetic classification of bridge boundaries}

For $d,q>0$ put $a=q/2$, $b=d/2$ and define $z_{d,q}>0$ by
\begin{equation}\label{p01:ga17:boundary}
 z_{d,q}M'(a;b;z_{d,q})=aM(a;b;z_{d,q}),
 \qquad C_{d,q}=\sqrt{2z_{d,q}}.
\end{equation}
The next proof includes existence and uniqueness, rather than treating
a numerical solution as a defined constant.

\begin{theorem}[Bridge-boundary classification]
\label{p01:ga17:bridge}
For every $d,q>0$, equation~\eqref{p01:ga17:boundary} has exactly one
positive solution.  For rational $d,q>0$,
\begin{equation}\label{p01:ga17:bridgeclass}
 C_{d,q}\text{ is algebraic}
 \quad\Longleftrightarrow\quad
 q-d\in2\mathbb Z_{\geq0}.
\end{equation}
In every other rational-parameter case $C_{d,q}$ is transcendental.
If $q-d=2N$ with $N\geq0$, set
$P_N(z)=M(-N;d/2;-z)$.  The algebraic boundary is specified without
special functions by the unique positive root of
\begin{equation}\label{p01:ga17:boundarypoly}
 zP_N'(z)+(z-q/2)P_N(z)=0.
\end{equation}
\end{theorem}
\begin{proof}
Write $f(z)=\sum_{n\geq0}c_nz^n$ with $c_n>0$.  For $z>0$,
the weights $c_nz^n/f(z)$ define a probability distribution on
the nonnegative integers.  Its mean is $L(z)=zf'(z)/f(z)$ and
\[
 \frac{dL}{d\log z}=\operatorname{Var}_z(n)>0.
\]
The variance is strictly positive since at least two weights are
positive.  Also $L(z)\to0$ as $z\downarrow0$.  For every fixed
$m$, the probability of $n\leq m$ tends to zero as $z\to\infty$:
the numerator has degree at most $m$, while the denominator is at
least its positive term of degree $m+1$.  Thus $L(z)\to\infty$.
This proves existence and uniqueness of $L(z)=a$.

We classify when $f'/f$ can be rational.  The coefficient bounds
in Lemma~\ref{p01:ga17:coeff} imply exponential type.  If its logarithmic
derivative is rational, write $f=e^{\lambda z}P(z)$ and let
$N=\deg P$.  Substitution into~\eqref{p01:ga17:kummerode} first forces
$\lambda(\lambda-1)=0$ by its highest-degree term.  If $\lambda=0$,
the next highest-degree term gives $a=-N$, impossible for $a>0$.
If $\lambda=1$, substitution gives
\[
 zP''+(z+b)P'+(b-a)P=0,
\]
whose highest coefficient forces $a-b=N$.  Conversely, when
$a-b=N\in\mathbb Z_{\geq0}$, Kummer's transformation gives
$f(z)=e^zM(-N;b;-z)$.  The latter polynomial has rational
coefficients when $a,b$ are rational, and is positive for $z>0$.

Outside this exceptional family, the pair $(f,f')$ is linearly
independent and solves
\[
 \binom{f}{f'}'
 =\begin{pmatrix}0&1\\ a/z&1-b/z\end{pmatrix}\binom{f}{f'}.
\]
Its only possible finite singular point is zero.  If $z_{d,q}$
were algebraic, equation~\eqref{p01:ga17:boundary} would contradict
Lemma~\ref{p01:ga17:mechanism}.  In the exceptional family the boundary
equation is precisely~\eqref{p01:ga17:boundarypoly}.  This is a nonzero
polynomial, since its highest term has degree $N+1$, and hence its
positive root is algebraic.  Finally $C_{d,q}$ is algebraic if and
only if $z_{d,q}=C_{d,q}^2/2$ is algebraic.
\end{proof}

The identification with stochastic boundaries is as follows.  The
squared Bessel bridge $Q$ of dimension $d>0$ has local generator
\[
 2u\partial_{uu}+\left(d-\frac{2u}{1-t}\right)\partial_u.
\]
For payoff $Q_\tau^{q/2}$, the ansatz
$V(t,u)=(1-t)^{q/2}g(u/(1-t))$ gives
\[
 4y g''+2(d-y)g'-qg=0.
\]
The solution regular at zero is a multiple of
$M(q/2;d/2;y/2)$.  Value matching and smooth fit at
$y=C_{d,q}^2$ give exactly~\eqref{p01:ga17:boundary}.  Hobson and
Liu's verification theorem establishes that the resulting boundary
$\sqrt{Q_t}=C_{d,q}\sqrt{1-t}$ is optimal; that stochastic
verification is an external input, not a new claim here
\cite[Theorem~3.1]{p01:ga17:HL}.  The same boundary for $d=1$ and payoff
$|X_\tau|^q$ is the reflected Brownian-bridge result of Ekstr\"om
and Wanntorp \cite[Theorem~3.2]{p01:ga17:EW}.

Three examples make the exceptional set concrete:
\begin{itemize}
\item For the reflected Brownian bridge, rational $q>0$ gives an
algebraic boundary exactly when $q=1,3,5,\ldots$.  In particular,
$C_{1,1}=1$ and $C_{1,3}=3^{1/4}$.  These elementary cases are
not being presented as new constants.
\item The quadratic-payoff boundary $C_{1,2}$ is transcendental.
\item The Brownian-excursion linear-payoff boundary
$C_{3,1}\simeq1.50339538$ is transcendental.  In this case the
regular solution is proportional to
$x^{-1}\int_0^x e^{t^2/2}\,dt$, so its equation can also be written
\[
 C_{3,1}e^{C_{3,1}^2/2}
       =2\int_0^{C_{3,1}}e^{t^2/2}\,dt.
\]
\end{itemize}
These examples constitute one classification theorem, not three
independent conjecture closures.  The distinction between reflected
and signed Brownian payoffs is essential.

\subsection{Target II: arithmetic of the Bessel obstruction extrema}

For $\mu>-1$, define
\[
 R_\mu(y)=\frac{I_{\mu+1}(y)}{I_\mu(y)},\qquad
 D_\mu(y)=R_\mu(y)+yR_\mu'(y),\qquad y>0.
\]
These are the real obstruction functions in the earlier HCM analysis.
We reproduce their shape argument so that the constants whose
arithmetic is asserted have fully specified existence and uniqueness.

\begin{lemma}[Shape, with both half-integer endpoints]
\label{p01:ga17:shape}
If $\mu\geq1/2$, $D_\mu$ increases strictly from zero to one;
its supremum is the unattained algebraic number $1$.  If
$-1<\mu<1/2$, it has a unique global maximum at $y_\mu>0$, with
$M_\mu=D_\mu(y_\mu)>1$.  For $-1/2\leq\mu<1/2$ this is its
only stationary point.  For $-1<\mu<-1/2$ there is exactly one
subsequent stationary point, a minimum with value in $(0,1)$.
At every stationary point $y$,
\begin{equation}\label{p01:ga17:curvature}
 D_\mu''(y)=-\frac{2(D_\mu(y)^2-1)}{y}.
\end{equation}
\end{lemma}
\begin{proof}
Put $\alpha=\mu+1/2$, $v=R_\mu+\alpha/y$, and $D=D_\mu$.
The Bessel equation and recurrence give
\begin{equation}\label{p01:ga17:phase}
 R'=1-\frac{2\alpha}{y}R-R^2,\qquad
 v'=1-v^2+\frac{\alpha(\alpha-1)}{y^2}
       =\frac{D-v}{y},\qquad D'=2(1-vD).
\end{equation}
For clarity, positivity needed in this argument does not follow just
from the Riccati equation.  Write $H_\mu(y)=\sum_{n\geq0}c_ny^{2n}$,
where $c_n>0$ for $\mu>-1$.  Under the probability weights
$c_ny^{2n}/H_\mu(y)$, one has $yR(y)=2\mathbb E_y n$ and hence
\[
 D(y)=(yR(y))'=\frac4y\operatorname{Var}_y(n)>0.
\]
All differentiations are justified by convergence of the entire
series.  The standard differentiated large-argument expansions used
below follow from \cite[\S10.40(i)]{p01:bec17:DLMF}.

At a zero of $v'$ one has
$v''=-2\alpha(\alpha-1)/y^3$.  At a zero of $D'$, positivity
gives $v=1/D>0$, and differentiation gives~\eqref{p01:ga17:curvature}.
Such a stationary point cannot have $D=1$.  Otherwise $v=1$ and
$v'=0$, so $\alpha(\alpha-1)=0$.  At these two exceptional
parameters $v=\tanh y$ or $v=\coth y$, respectively, neither of
which equals one at a finite positive point.

The endpoint expansions, also valid after differentiation, are
\begin{align*}
 D(y)&=\frac{y}{\mu+1}+O(y^3) &&(y\downarrow0),\\
 v(y)&=1+\frac{\alpha(\alpha-1)}{2y^2}+O(y^{-3})
       &&(y\to\infty),\\
 D(y)&=1-\frac{\alpha(\alpha-1)}{2y^2}+O(y^{-3})
       &&(y\to\infty).
\end{align*}
If $0<\alpha<1$, $v'$ starts negative, ends positive, and
every zero of $v'$ is a strict upward crossing.  Thus $v$ has
one minimum.  Before that minimum $D<v$, so any stationary point
of $D$ would have $D<1$ and be a strict minimum.  This cannot be
the first zero of $D'$, which starts positive.  After that minimum
$D>v$, so every stationary point of $D$ is a strict maximum
above one.  Since $D'$ ends negative, there is exactly one.

If $\alpha<0$, $v'$ starts positive, ends negative, and each
zero is a strict downward crossing.  Before its unique maximum,
every stationary point of $D$ must be a maximum above one; after
it, every stationary point must be a minimum below one.  There
must be a value $D>1$: otherwise
$[y(1-R(y))]'=1-D\geq0$, whereas $y(1-R(y))$ starts at zero
and tends to $\alpha<0$.  Since $D$ eventually lies below one
and increases to one, the maximum and the later minimum both
exist and are unique.  For $\alpha=0$,
$R=\tanh y$, and $D'=2\operatorname{sech}^2y(1-y\tanh y)$;
$y\tanh y$ increases from zero to infinity, giving one maximum.

For $\alpha>1$, $v'$ starts negative and any zero would have
to be a downward crossing.  There are no such zeros.  Therefore
every possible stationary point of $D$ would be a minimum below
one, which again cannot be the first zero of $D'$.  Thus $D'>0$.
For $\alpha=1$ the conclusion follows from $v=\coth y$ and the
same argument, or direct differentiation.  The endpoint limits
finish the proof.
\end{proof}

\begin{theorem}[Transcendental extremum locations and values]
\label{p01:ga17:extrema}
Let $\mu>-1$ be rational.  Every positive stationary point of
$D_\mu$, and its value, is transcendental.  In particular,
for rational $-1<\mu<1/2$, both $y_\mu$ and $M_\mu$ in
Lemma~\ref{p01:ga17:shape} are transcendental.  For rational
$-1<\mu<-1/2$, the later minimum has transcendental location
and value as well.  The endpoint $\mu=-1/2$ is included:
$y_{-1/2}\tanh y_{-1/2}=1$ and $M_{-1/2}=y_{-1/2}$.
\end{theorem}
\begin{proof}
The entire even function $H_\mu$ of Lemma~\ref{p01:ga17:coeff}
is of exponential type, and its logarithmic derivative is $R_\mu$.
Suppose this derivative were rational.  The analytic observation
above would give $H_\mu=e^{\lambda z}P(z)$.  Evenness implies
$e^{2\lambda z}=P(-z)/P(z)$, forcing $\lambda=0$; otherwise
the exponential has an essential singularity at infinity, whereas
the rational function does not.  Thus $H_\mu$ would be a
polynomial, contrary to its strictly positive infinitely many
even Taylor coefficients.  The system associated with
\eqref{p01:ga17:besselode} has no finite singular point except possibly
zero.  Lemma~\ref{p01:ga17:mechanism} shows that $R_\mu(y)$ is
transcendental for every nonzero algebraic $y$.

At a stationary point, substitute
$D=y(1-R^2)-2\mu R$ into $D'=0$ in~\eqref{p01:ga17:phase}.
Writing $\alpha=\mu+1/2$ gives the polynomial equation
\begin{equation}\label{p01:ga17:criticalpoly}
 y-(yR+\alpha)\{y(1-R^2)-2\mu R\}=0.
\end{equation}
For a fixed nonzero algebraic $y$ this is a nonzero polynomial
in $R$, since its cubic coefficient is $y^2$.  It would force
$R$ to be algebraic.  Therefore a stationary location $y$ is
transcendental.

The value requires a further elimination; transcendence of the
location alone does not suffice.  Write $m=D(y)$.  At a stationary
point, $v=1/m$ and
\[
 D=v+y(1-v^2)+\frac{\alpha(\alpha-1)}{y}.
\]
Eliminating $v$ yields
\begin{equation}\label{p01:ga17:extremalcurve}
 (m^2-1)y^2-m(m^2-1)y
       +(\mu^2-1/4)m^2=0.
\end{equation}
Lemma~\ref{p01:ga17:shape} showed $m>0$ and $m\ne1$.
If $m$ were algebraic, this would be a nonzero quadratic equation
over $\overline{\mathbb Q}$ for $y$, contradicting the preceding
paragraph.  Thus $m$ is transcendental.  The statements about
which extrema exist follow from Lemma~\ref{p01:ga17:shape}.
\end{proof}

\begin{corollary}[The normalized critical scale]
\label{p01:ga17:scale}
For rational $-1<\mu<1/2$, the positive number
\[
 S_\mu=\sqrt{\frac{y_\mu(M_\mu^2-1)}{3}}
\]
is transcendental.  Thus the coefficient $\pi S_\mu$ in the
earlier near-threshold HCM asymptotic has a transcendental ratio
to $\pi$.  This statement does not assert that $\pi S_\mu$
itself is transcendental.
\end{corollary}
\begin{proof}
Set $K=3S_\mu^2=y_\mu(M_\mu^2-1)>0$ and
$A=\mu^2-1/4$.  Equation~\eqref{p01:ga17:extremalcurve} gives
$K(M_\mu-y_\mu)=AM_\mu^2$.  If $K$ were algebraic, substitution
$y_\mu=M_\mu-AM_\mu^2/K$ into its definition would give
\[
 (KM_\mu-AM_\mu^2)(M_\mu^2-1)-K^2=0.
\]
This is a nonzero polynomial over $\overline{\mathbb Q}$
(its constant term is $-K^2$), contradicting
Theorem~\ref{p01:ga17:extrema}.  Hence $K$, and therefore $S_\mu$,
is transcendental.
\end{proof}

The new step relative to the earlier analytic obstruction analysis
is arithmetic: the Riccati stationary relation passes transcendence
from a logarithmic derivative to the extremum location, and the
finite-fiber algebraic relation~\eqref{p01:ga17:extremalcurve} passes it
to the extremal value.  The analytic shape lemma is repeated for
self-containment and is not counted as a new result.  The maximum,
minimum, and normalized scale are grouped as one target family.

\subsection{Target III: the signed Brownian boundary remains undecided}

For signed payoff $X_\tau$ the Brownian-bridge boundary is
$B\sqrt{1-t}$ with $B\simeq0.8399236756$, not the reflected
boundary $C_{1,1}=1$.  Its defining equation is
\begin{equation}\label{p01:ga17:shepp}
 (1-B^2)F(B)=B,\qquad
 F(x)=\int_0^\infty e^{xu-u^2/2}\,du
     =e^{x^2/2}\left(\sqrt{\pi/2}
                    +\int_0^x e^{-u^2/2}\,du\right).
\end{equation}
This is the established stopping constant in
\cite[Theorem~2.1]{p01:ga17:EW}.  Its existence and uniqueness on
$(0,1)$ follow directly: $F'=1+xF$, so
$h(x)=(1-x^2)F(x)-x$ has
$h'(x)=-x(1+x^2)F(x)-x^2<0$ for $x>0$, while
$h(0)>0$ and $h(1)=-1$.

The naive $E$-function argument fails because $F(0)=\sqrt{\pi/2}$
is not algebraic.  Dividing by this value does not repair the issue:
the new derivative at zero is then $\sqrt{2/\pi}$.  More precisely,
if $B$ were algebraic and $z=B^2/2$, equation~\eqref{p01:ga17:shepp}
would yield
\begin{equation}\label{p01:ga17:sheppobstruction}
 \sqrt{\pi/2}
 =\frac{B}{1-B^2}e^{-z}
       -B M(1/2;3/2;-z).
\end{equation}
The right side combines $E$-function values at an algebraic
argument, but the specialization theorem used here concerns
relations with algebraic coefficients; it gives no contradiction
to~\eqref{p01:ga17:sheppobstruction}.  No irrationality or
transcendence assertion for $B$ is proved in this section.  We
also do not claim that the literature search establishes its
global open status.  This is an explicitly recorded failed route,
preventing a normalization constant from being silently treated
as algebraic.

\begin{proposition}[A conditional link to the $E$--$G$ intersection problem]
\label{p01:ga17:sheppconditional}
Let $\mathbf E$ be the ring of values of $E$-functions at algebraic
points.  If $\pi\notin\mathbf E$, the signed stopping constant
$B$ is transcendental.  In particular this follows from the
conjecture $\mathbf E\cap\mathbf G=\overline{\mathbb Q}$,
where $\mathbf G$ is the ring of values of analytic continuations
of $G$-functions at algebraic points, discussed in
\cite[Introduction]{p01:ga17:FR}.
\end{proposition}
\begin{proof}
If $B$ were algebraic, the right side of
\eqref{p01:ga17:sheppobstruction} would belong to $\mathbf E$.
The ring is closed under sums, products, and multiplication by
algebraic numbers; these are the usual closure properties of
$E$-functions, with an algebraic rescaling of the argument if
needed.  Squaring gives $\pi\in\mathbf E$.  On the other hand,
$\pi=6\arcsin(1/2)$ belongs to $\mathbf G$: the series
\[
 \arcsin z=\sum_{n\geq0}
    \frac{\binom{2n}{n}}{4^n(2n+1)}z^{2n+1}
\]
is a $G$-function, with $(1-z^2)g''-zg'=0$ and exponential
coefficient and common-denominator bounds.  Its value at $1/2$
is inside its disk of convergence, so no branch choice is needed.
The intersection conjecture would then force $\pi$ to be algebraic,
a contradiction.  This proves the conditional assertion only.
\end{proof}

\paragraph{Accounting.}
There are two unconditionally proved arithmetic target families and
one target with a conditional reduction but no unconditional decision.
The bridge classification uses established
optimal-stopping formulas; the Bessel classification adds arithmetic
information to the earlier manuscript's obstruction constants.
Neither family has been verified as a historically stated open
arithmetic conjecture.  Numerical approximations in the examples
serve only to identify the constants; every arithmetic conclusion
rests on an exact equation and a proof.


% Integrated source: spectral_results.tex
\section{Arithmetic of critical spectral thresholds and frozen edges}
\label{p01:sc:sec}

The inverse-observation theory supplies constants as minima and as
right boundaries of zero sets.  Their uniqueness does not determine
their arithmetic nature.  This section classifies a substantial class
of finite spectral minima, then contrasts it with finite-prime and
infinite-prime boundary equations.  The main reduction uses the
Gelfond--Schneider theorem: a positive algebraic number other than one,
raised to an algebraic irrational power, is transcendental.  This is
a classical theorem, not a new transcendence principle; a modern
primary proof source is \cite{p01:sc:GS}.  The elementary frozen-edge
consequences of Lindemann--Weierstrass are explicitly separated from
the critical-value argument.

\subsection{A complete classification for a one-sided rank-two spectrum}

Let
\begin{equation}\label{p01:sc:eq:threshold}
 S(x)=\sum_{j=1}^m w_jx^{r_j},\qquad
 T(C)=\min_{x>0}\frac{1+Cx}{S(x)},\qquad
 0<r_j<1,\quad w_j>0,\quad C>0.
\end{equation}
Weights need not sum to one.  Put
\[
             W=\sum_jw_j,\qquad M=\sum_jw_jr_j.
\]
All positive real powers in this section use the real logarithm.

\begin{theorem}[The unique algebraic critical value]
\label{p01:sc:thm:ranktwo}
Suppose the weights and $C$ are algebraic.  Suppose there are an
algebraic irrational number $\theta$, a positive integer $q$, and
integers $m_j,n_j$ such that
\begin{equation}\label{p01:sc:eq:ranktwo}
 r_j=\frac{m_j+n_j\theta}{q},\qquad
 n_j\ge0\ \hbox{for every }j,\qquad
 n_j>0\ \hbox{for at least one }j.
\end{equation}
The same conclusion holds with all $n_j\le0$ and at least one
strict inequality.  Then the following conditions are equivalent:
\begin{enumerate}
\item $T(C)$ is algebraic;
\item the minimizer in \eqref{p01:sc:eq:threshold} is $x=1$;
\item $C=M/(W-M)$.
\end{enumerate}
At this exceptional setting,
\begin{equation}\label{p01:sc:eq:exception}
                T(C)=\frac{1+C}{W}=\frac1{W-M}.
\end{equation}
At every other positive algebraic setting $C$, the critical value
is transcendental.

If all $w_j$ and $C$ are positive rational numbers, the exceptional
setting is impossible, and every such critical value is
transcendental.
\end{theorem}

The one-sided condition is a condition on the coefficients of the
irrational basis element, not on the signs of the physical exponents:
the latter always lie in $(0,1)$.  It permits any finite number of
distinct irrational powers and arbitrarily many rational powers.
It is not merely a two-term formula.  Its role is to prevent a
repeated algebraic component at the critical point; this issue cannot
be omitted from a general algebraic-elimination argument.

\begin{lemma}[Contact geometry]\label{p01:sc:lem:contact}
The minimum in \eqref{p01:sc:eq:threshold} has a unique minimizer $x_C$.
It is characterized by
\begin{equation}\label{p01:sc:eq:contact}
       C=\frac{S'(x_C)}{S(x_C)-x_CS'(x_C)}.
\end{equation}
In particular $x_C=1$ if and only if $C=M/(W-M)$.
\end{lemma}
\begin{proof}
The objective tends to infinity as $x$ tends to zero or infinity,
because all exponents lie strictly between zero and one.  Write
$H=S-xS'>0$.  Strict concavity gives
\[
 \frac{d}{dx}\frac{S'(x)}{H(x)}
                     =\frac{S''(x)S(x)}{H(x)^2}<0.
\]
Moreover $xS'/H$ is bounded above and below by positive constants,
using the smallest and largest exponents.  Thus $S'/H$ decreases
continuously from infinity to zero.  The derivative of the objective
vanishes exactly at \eqref{p01:sc:eq:contact}, which proves existence
and uniqueness.  Evaluation at one gives the last assertion.
\end{proof}

The following elementary algebraic lemma supplies the step that a
formal resultant calculation can otherwise hide.  Let
$\overline{\mathbb Q}$ denote the algebraic numbers.

\begin{lemma}[An irrational diagonal derivation has no invariant curve]
\label{p01:sc:lem:derivation}
On $\overline{\mathbb Q}[u^{\pm1},v^{\pm1}]$, put
\[
                        D=u\partial_u+\theta v\partial_v,
                         \qquad \theta\in\overline{\mathbb Q}
                                             \setminus\mathbb Q.
\]
If $H$ is a nonunit irreducible Laurent polynomial, then $H$ and
$DH$ are coprime.  Every common zero of $H$ and $DH$ in
$(\mathbb C^*)^2$ consequently has both coordinates algebraic.
\end{lemma}
\begin{proof}
If $DH=0$, the coefficient of each monomial $u^av^b$ is multiplied
by $a+b\theta$.  Irrationality forces every nonconstant coefficient
to vanish, so $H$ would be a constant, contrary to the hypothesis.
Suppose instead that $DH=AH$ for a Laurent polynomial $A$.
For a nonzero Laurent polynomial $P$, denote the convex hull of its
exponent vectors by $\mathcal N(P)$.  Products satisfy
\[
                 \mathcal N(PQ)=\mathcal N(P)+\mathcal N(Q).
\]
One direct proof is to maximize a generic real linear functional:
each factor has a unique maximizing exponent, so the corresponding
coefficient in the product cannot cancel.  Continuity gives equality
of support functions in every direction.  Differentiation deletes
monomials but introduces none.  Hence
\[
 \mathcal N(H)+\mathcal N(A)=\mathcal N(DH)\subseteq\mathcal N(H).
\]
Taking support functions in all directions forces
$\mathcal N(A)=\{(0,0)\}$, so $A$ is constant.  The equation
$DH=cH$ then says that all nonzero monomials of $H$ have the same
weight $a+b\theta$.  These weights are distinct for distinct
integer pairs, so $H$ is a single monomial, a unit.  This
contradiction proves coprimality.

Here are details of the algebraicity assertion.  Multiplying $H$ by a
Laurent monomial does not change the common torus zeros of $H$ and
$DH$, by the product rule.  We therefore normalize $H$ to an ordinary
polynomial divisible by neither $u$ nor $v$; it remains irreducible.
If $H$ depends
only on $u$, irreducibility over the algebraically closed field
$\overline{\mathbb Q}$ makes it an associate of $u-a$ for a
nonzero algebraic $a$; it has no common torus zero with $DH$.
Otherwise view the two polynomial equations as equations in $v$ over
$\overline{\mathbb Q}(u)$.  Their coprimality gives a resultant
which is a nonzero polynomial in $u$, vanishing at any
common torus zero.  Thus its $u$ coordinate is algebraic.  The
specialization of $H$ at this coordinate cannot vanish identically
as a polynomial in $v$: otherwise $u-a$ would divide $H$, against
irreducibility and dependence on $v$.  The $v$ coordinate is
therefore algebraic as well.
\end{proof}

\begin{proof}[Proof of Theorem~\ref{p01:sc:thm:ranktwo}]
Only the implication from algebraicity to $x_C=1$ needs proof.
Suppose $\tau=T(C)$ is algebraic and put
\[
 u_*=x_C^{1/q},\qquad v_*=x_C^{\theta/q},\qquad
 P(u,v)=\sum_jw_j u^{m_j}v^{n_j},\qquad
 F(u,v)=P(u,v)-\frac{1+Cu^q}{\tau}.
\]
The value and stationary equations give
\begin{equation}\label{p01:sc:eq:criticalsystem}
                  F(u_*,v_*)=0,\qquad DF(u_*,v_*)=0.
\end{equation}
Indeed, along $(u,v)=(x^{1/q},x^{\theta/q})$, logarithmic
differentiation is $D/q$.  Under \eqref{p01:sc:eq:ranktwo},
\[
          F_v(u_*,v_*)=\sum_jw_jn_j u_*^{m_j}v_*^{n_j-1}>0.
\]
For nonpositive $n_j$ the same expression is strictly negative.
Factor $F$ over the algebraic numbers.  Since its $v$ derivative
does not vanish at the point, exactly one irreducible factor $H$
vanishes there, and its multiplicity in $F$ is one.  Thus
$F=HR$ with $R(u_*,v_*)\ne0$.  Equation
\eqref{p01:sc:eq:criticalsystem} now gives $DH(u_*,v_*)=0$.
Lemma~\ref{p01:sc:lem:derivation} proves that $u_*,v_*$ are algebraic.
But $v_*=u_*^\theta$.  Gelfond--Schneider forces $u_*=1$,
because $u_*>0$ and the exponent is algebraic irrational.
Hence $x_C=1$.

The contact lemma and direct evaluation give the other implications
and \eqref{p01:sc:eq:exception}.  Finally, if the weights are rational,
then $W\in\mathbb Q$ and
\[
 M=\frac1q\sum_jw_jm_j+\frac\theta q\sum_jw_jn_j
\]
is irrational, since the second coefficient is nonzero.  If
$C=M/(W-M)$ were rational, then $M=CW/(1+C)$ would be rational,
a contradiction.
\end{proof}

\begin{remark}[Why one cannot simply assume a squarefree level set]
Without the sign condition, an algebraic level polynomial $F$ can
in principle have a repeated factor through the contact point.
Then $F=DF=0$ may hold on a whole curve, and a resultant can
vanish identically.  The preceding proof establishes a nonvanishing
partial derivative before selecting its factor; it does not assume
away that difficulty.  The same proof works for any rank-two
spectrum whose algebraic critical level is nonsingular at the
contact point.  No unconditional theorem for all mixed-sign
irrational coefficients is claimed here.
\end{remark}

\subsection{Critical locations and three specific threshold constants}

\begin{proposition}[Rational spectra and a single irrational coset]
\label{p01:sc:prop:locations}
Assume the weights and $C$ are positive algebraic.
\begin{enumerate}
\item If all $r_j$ are rational, both $x_C$ and $T(C)$ are algebraic.
\item If $r_j=\alpha+k_j/q$ for a fixed algebraic irrational
      $\alpha$, integers $k_j$, and a positive integer $q$, then
      $x_C$ is algebraic.  If $x_C\ne1$, both $T(C)$ and
      $\log x_C$ are transcendental.
\end{enumerate}
\end{proposition}
\begin{proof}
For the first assertion put $u=x^{1/q}$ with $q$ a common
denominator.  The stationary equation becomes a nonzero polynomial
equation over the algebraic numbers after clearing powers of $u$.
Its positive solution is algebraic, and the objective at that
solution is algebraic.  The polynomial is nonzero, for otherwise
the objective would be constant, contrary to its divergence at
both ends.

For the second assertion put $u=x^{1/q}$ and
$P(u)=\sum_jw_ju^{k_j}$, a positive Laurent polynomial on
$(0,\infty)$.  Then $S(x)=u^{q\alpha}P(u)$, and the
stationary equation is
\begin{equation}\label{p01:sc:eq:cosetstationary}
 qCu^qP(u)-(1+Cu^q)\{q\alpha P(u)+uP'(u)\}=0.
\end{equation}
It is nonzero: for the largest index $k$ occurring in $P$, its
highest coefficient is a positive multiple of
$q-q\alpha-k=q(1-r_{\max})>0$.  Consequently the positive
root $u_C$, and hence $x_C$, is algebraic.  If $u_C\ne1$,
Gelfond--Schneider makes $u_C^{q\alpha}$ transcendental.
The remaining factors in
$(1+Cu_C^q)/(u_C^{q\alpha}P(u_C))$ are nonzero algebraic,
so $T(C)$ is transcendental.  Finally $\log x_C$ cannot be
nonzero algebraic, by the Hermite--Lindemann consequence of
Lindemann--Weierstrass \cite{p01:sc:LW}.
\end{proof}

The following are three concrete constants arising from the same
critical observation mechanism.  Their proofs are not three
independent transcendence breakthroughs; the first is a direct
classical evaluation, while the second and third illustrate the
critical-value classification.

\begin{enumerate}
\item For $S(x)=x^{1/\sqrt2}$ and $C=1$, the contact is
      $x_*=1+\sqrt2$, and
      \begin{equation}\label{p01:sc:eq:constantone}
       \tau_1=\min_{x>0}\frac{1+x}{x^{1/\sqrt2}}
          =\frac{2+\sqrt2}{(1+\sqrt2)^{1/\sqrt2}}
      \end{equation}
      is transcendental by Gelfond--Schneider.  This explicit
      corollary is included as a calibration example, not as a
      newly solved historical problem.

\item The mixed rational--irrational threshold
      \begin{equation}\label{p01:sc:eq:constanttwo}
       \tau_2=\min_{x>0}
                \frac{2(1+x)}{x^{1/3}+x^{1/\sqrt2}}
      \end{equation}
      is transcendental.  Use $q=6$, $\theta=\sqrt2$ and
      exponent pairs $(2,0),(0,3)$ in
      Theorem~\ref{p01:sc:thm:ranktwo}.  Both weights and the setting
      are rational, so no exceptional contact is possible.  This
      proves a statement about the minimum itself without first
      finding an explicit closed form for its minimizer.

\item The two-irrational-frequency threshold
      \begin{equation}\label{p01:sc:eq:constantthree}
       \tau_3=\min_{x>0}
         \frac{2(1+x)}{x^{(\sqrt2-1)/2}+x^{\sqrt2/2}}
      \end{equation}
      is also transcendental, and its contact is algebraic.
      More explicitly, put $\theta=\sqrt2-1$.  If $u_*$ is
      the unique positive root of
      \begin{equation}\label{p01:sc:eq:cubic}
       (1-\theta)u^3+(2-\theta)u^2-(1+\theta)u-\theta=0,
      \end{equation}
      then
      \[
       x_*=u_*^2,\qquad
       \tau_3=\frac{2(1+u_*^2)}{u_*^\theta(1+u_*)}.
      \]
      Equation \eqref{p01:sc:eq:cubic} is the stationary equation
      for $2(1+u^2)/(u^\theta(1+u))$; uniqueness follows from
      Lemma~\ref{p01:sc:lem:contact}.  Substitution of $u=1$
      gives $2-4\theta\ne0$, so Gelfond--Schneider applies.
\end{enumerate}

There is a necessary exception even with an algebraic irrational
frequency.  For any such $r\in(0,1)$, the one-atom spectrum and
the algebraic setting $C=r/(1-r)$ have $x_C=1$ and
$T(C)=1/(1-r)$ algebraic.  Even rational $C$ does not suffice
if the weights are only assumed algebraic: take
\[
 r_1=\frac14,\quad r_2=\frac1{\sqrt2},\quad
 w_1=\frac{r_2-1/2}{r_2-1/4},\quad
 w_2=\frac{1/4}{r_2-1/4}.
\]
The weights are positive algebraic, sum to one, and have weighted
mean $1/2$.  At $C=1$ the threshold is exactly $2$.
These examples rule out an unqualified claim that an irrational
frequency automatically makes every critical value transcendental.

\subsection{Frozen edges: what algebraic frequencies do prove}

For positive distinct frequencies and magnitudes define the frozen
edge by
\begin{equation}\label{p01:sc:eq:frozen}
             \sum_{j=1}^m a_j e^{-\lambda_j S}=1,
                   \qquad a_j,\lambda_j>0.
\end{equation}
There is a unique real solution, since the left side is strictly
decreasing from infinity to zero.  In independent finite-sum
models this is the rightmost zero boundary; the arithmetic
statements below concern the explicitly defined constant $S$.

\begin{proposition}[Classical consequences for algebraic frequencies]
\label{p01:sc:prop:frozen}
If every $a_j$ and $\lambda_j$ is algebraic, then $S=0$ if
$\sum a_j=1$, and otherwise $S$ is transcendental.  More generally,
every nonzero complex zero of
\[
                  1-\sum_j c_j e^{-\lambda_jz},
      \qquad c_j\in\overline{\mathbb Q}\setminus\{0\},
\]
is transcendental when the distinct positive $\lambda_j$ are
algebraic.

In particular let $S_{1,\sqrt2}$ solve
\begin{equation}\label{p01:sc:eq:fourthconstant}
                  e^{-S_{1,\sqrt2}}
                         +e^{-\sqrt2S_{1,\sqrt2}}=1.
\end{equation}
Both $S_{1,\sqrt2}$ and $e^{-S_{1,\sqrt2}}$ are transcendental.
\end{proposition}
\begin{proof}
The zero statement follows immediately from Lindemann--Weierstrass:
if $z\ne0$ were algebraic, the distinct algebraic exponents
$0,-\lambda_1z,\ldots,-\lambda_mz$ would give a nontrivial
algebraic linear relation among their exponentials.  The frozen
edge is zero exactly when $\sum a_j=1$, proving the first claim.

For \eqref{p01:sc:eq:fourthconstant}, the edge is positive and hence
transcendental by the first claim.  Set $X=e^{-S_{1,\sqrt2}}$,
so $0<X<1$ and $X+X^{\sqrt2}=1$.  If $X$ were algebraic,
then $X^{\sqrt2}=1-X$ would be algebraic, contradicting
Gelfond--Schneider.
\end{proof}

These are transfers of classical transcendence theorems, not new
solutions of long-standing logarithm problems.  The zero statement
also does not classify the real and imaginary parts separately:
a transcendental complex number can have an algebraic real part.
The condition on complex coefficients is algebraicity of $c_j$,
not algebraicity of a chosen angle in $c_j=a_je^{i\theta_j}$.

\subsection{Finite prime edges and the barrier at infinitely many primes}

Prime frequencies are logarithms of algebraic numbers, not algebraic
numbers themselves.  Therefore Proposition~\ref{p01:sc:prop:frozen}
does not apply to them.  One can nevertheless obtain a different,
elementary classification at the level of rationality.

\begin{proposition}[Every finite unit-weight prime edge is irrational]
\label{p01:sc:prop:primeirrational}
Let $\mathcal P$ be a finite set of at least two primes.  The unique
positive solution of
\begin{equation}\label{p01:sc:eq:finiteprime}
                         \sum_{p\in\mathcal P}p^{-\sigma}=1
\end{equation}
is irrational.
\end{proposition}
\begin{proof}
Existence and positivity follow from strict decrease, the value
$|\mathcal P|>1$ at zero, and the limit zero at infinity.
Suppose $\sigma=a/b$ with positive integers $a,b$.  Work in a
number field containing all positive radicals $p^{1/b}$ and
choose a nonarchimedean valuation $v$ above any one prime
$p_0\in\mathcal P$.  Normalize it so that $v(p_0)>0$.
For every other rational prime $p$, $v(p)=0$.  Thus
\[
 v(p_0^{-a/b})=-\frac abv(p_0)<0,
       \qquad v(p^{-a/b})=0\quad(p\ne p_0).
\]
In a sum with a unique term of strictly smallest valuation, the
valuation of the sum equals that smallest valuation.  This follows
from the ultrametric inequality applied in both directions to that
term and the sum of the others.  The left side of
\eqref{p01:sc:eq:finiteprime} therefore has negative valuation,
whereas $v(1)=0$, a contradiction.
\end{proof}

For example the natural two-prime constant
\begin{equation}\label{p01:sc:eq:twoprime}
                         2^{-\sigma_{2,3}}+3^{-\sigma_{2,3}}=1
\end{equation}
is irrational.  The corresponding exponential polynomial is a
standard nonlattice self-similar-string model studied in
\cite{p01:sc:fractal}.  The valuation proof is an elementary finite
radical argument; no historical priority is claimed for this
irrationality observation.  We do not prove that
$\sigma_{2,3}$ is transcendental.  If it were algebraic irrational,
Gelfond--Schneider would make each summand transcendental, but that
does not prevent their sum from being one.  Baker's theorem concerns
algebraic linear combinations of finitely many logarithms
\cite{p01:sc:Baker}; the displayed additive relation between two
unknown powers is not such a relation.

For comparison, the commensurable edge
$2^{-S}+4^{-S}=1$ has
\[
              S=\frac{\log((1+\sqrt5)/2)}{\log2}.
\]
It is transcendental: if $S=a/b$ were rational, then
$((1+\sqrt5)/2)^b=2^a$, impossible because the golden ratio
is an algebraic unit whereas a nonzero power of two is not.
If $S$ were algebraic irrational, Gelfond--Schneider would
contradict the algebraicity of $2^S$.  This familiar logarithm
corollary is a useful comparison, not a new difficult constant.

Two infinite constants resist the preceding arguments:
\begin{equation}\label{p01:sc:eq:infiniteconstants}
 P(\sigma_P)=1,\qquad P(s)=\sum_{p\ \mathrm{prime}}p^{-s},
 \qquad\qquad \zeta(\rho)=2.
\end{equation}
The second constant is the exponent in the classical ordered
factorization problem; see \cite{p01:sc:factorization} for its role.
Both equations have unique real solutions greater than one,
by positivity and strict monotonicity.  For $P$, divergence at
one follows from divergence of the sum of prime reciprocals;
for $\zeta$ it is the usual pole.  Both functions tend to zero
and one, respectively, at positive infinity.

No irrationality or transcendence conclusion for $\sigma_P$ or
$\rho$ is obtained here.  The obstruction is specific:
\begin{enumerate}
\item The critical-value theorem is a statement about two Laurent
      variables and a finite algebraic level polynomial.  Neither
      infinite equation supplies such a polynomial or the associated
      critical-point equation.
\item The finite-prime valuation proof concerns equality of a
      \emph{finite} sum in a number field.  Real convergence of an
      infinite radical series does not preserve that valuation
      argument.  Limits of irrational, or even transcendental,
      numbers may be rational.
\item Lindemann--Weierstrass requires algebraic exponents.  The
      products $s\log p$ are not algebraic merely because $s$
      is assumed algebraic.  The usual $E$-function value theorems
      likewise require their algebraic-coefficient differential
      setting; writing a prime Dirichlet series does not provide it.
\item Recovering the full spectrum or locating its boundary
      analytically does not provide the Diophantine estimates needed
      to exclude rational or algebraic boundary locations.
\end{enumerate}
The checked primary sources supply no theorem that resolves these
two arithmetic classifications by the reductions used here.  This
is a limitation of the present work and search, not an assertion of
an exhaustive priority survey.  In particular nothing in this
section classifies the ordinates of the nontrivial Riemann-zeta
zeros or resolves the Riemann hypothesis.

The substantive finite-spectrum result is
Theorem~\ref{p01:sc:thm:ranktwo}: algebraicity of a variational critical
value forces the unique normalized contact $x=1$, and positive
rational data exclude that exception.  The explicit constants,
frozen-edge corollaries, and finite-prime irrationality proof are
kept separate from the unresolved infinite-prime targets.


% Integrated source: mle_results.tex
\section{Exponential-integral values beyond the algebraic exponential field}
\label{p01:mc17:section}

This section records unconditional consequences of classical
transcendence theory. Its main tool is the Siegel--Shidlovskii theorem,
not an inference from statistical identifiability. The consequences
include the transcendence of $\log\operatorname{Ein}(1)$, a simultaneous
statement for logarithmic moments of truncated exponential laws, and
transcendence of an extremum location. We make no priority claim for
these consequences and no claim to resolve the arithmetic nature of
Euler's constant or the classical elementary constants discussed below.

Define the entire exponential integral by
\[
 \operatorname{Ein}(z)=\int_0^z\frac{1-e^{-t}}t\,dt
       =\sum_{n=1}^{\infty}\frac{(-1)^{n-1}z^n}{n\,n!},
\]
and let
\[
 \mathcal F_{\exp}=\overline{\mathbb Q}
               \bigl(e^b:b\in\overline{\mathbb Q}\bigr).
\]
This is a field generated by exponentials of algebraic numbers; it is
not a field closed under taking exponentials of all its own elements.

\begin{theorem}[A simultaneous classical-theory consequence]
\label{p01:mc17:Ein-field}
For every finite collection of distinct nonzero algebraic numbers
$a_1,\ldots,a_s$, the numbers
\[
          \operatorname{Ein}(a_1),\ldots,\operatorname{Ein}(a_s)
\]
are algebraically independent over $\mathcal F_{\exp}$.
In particular, each is transcendental over that entire field, not only
over $\overline{\mathbb Q}(e)$.
\end{theorem}

The arithmetic properties of exponential-integral and polyexponential
values are classical applications of $E$-function theory; see
\cite{p01:mc17:Rivoal,p01:bec17:Beukers}. Related probabilistic interpretations
and higher moment constants also occur in \cite{p01:mc17:Powers}.
We give the finite-field reduction and differential argument explicitly
because independence from every finite collection of algebraic
exponentials is what the logarithmic consequence needs. Our literature
check does not certify novelty of this formulation.

\subsection{The differential and arithmetic steps}

We use the following classical antiderivative criterion
\cite{p01:mc17:Ostrowski,p01:mc17:Srinivasan}. If $K\subset L$ are differential
fields of characteristic zero with the same algebraically closed
constant field $C$, and $f_1',\ldots,f_s'\in K$, then algebraic
dependence of the $f_j$ over $K$ implies
\[
                  \sum_{j=1}^s c_jf_j\in K
\]
for some $c_j\in C$ not all zero. Here the fields will be subfields of
meromorphic functions on $\mathbb C$, containing all complex constants,
so the common constant field is exactly $\mathbb C$.

\begin{lemma}[The formal simple-pole obstruction]
\label{p01:mc17:pole}
Let $b_1,\ldots,b_r\in\mathbb C$ be linearly independent over
$\mathbb Q$, and put
$K=\mathbb C(z,e^{b_1z},\ldots,e^{b_rz})$. Suppose that distinct
nonzero $a_1,\ldots,a_s$ are integer linear combinations of the $b_i$.
Then $\operatorname{Ein}(a_1z),\ldots,\operatorname{Ein}(a_sz)$ are
algebraically independent over $K$.
\end{lemma}
\begin{proof}
The functions $w_i=e^{b_iz}$ are algebraically independent over
$\mathbb C(z)$. Indeed, a polynomial relation becomes a finite
exponential sum $\sum_\lambda p_\lambda(z)e^{\lambda z}=0$ with distinct
$\lambda$ after clearing rational denominators. On a ray on which the
real parts of all $\lambda z$ are distinct, its largest exponential
cannot be canceled by the others with polynomial coefficients.
Thus $K$ is the rational function field $\mathbb C(w_1,\ldots,w_r)(z)$,
with derivation
\[
       D=\frac\partial{\partial z}
               +\sum_{i=1}^r b_iw_i\frac\partial{\partial w_i}.
\]

No rational function $R\in K$ has derivative $H(w)/z$ for a nonzero
$H\in\mathbb C(w_1,\ldots,w_r)$. To prove this, take the formal
Laurent expansion of $R$ in $z$ with coefficient field
$\mathbb C(w_1,\ldots,w_r)$. If $R$ has a pole of order $m\ge1$,
its derivative has a pole of order $m+1$: the derivative of the
leading $A(w)z^{-m}$ contributes $-mA(w)z^{-m-1}$, while the
$w$-derivations have pole order at most $m$. If $R$ is regular at
this formal divisor, then so is $DR$. A simple pole is impossible
in either case. This is a valuation in the abstract rational function
field. It does not treat the analytic values $w_i(0)=1$ as independent
coordinates at a point, and it does not contradict analyticity of
$\operatorname{Ein}$ at the complex number zero.

Write $f_j(z)=\operatorname{Ein}(a_jz)$. We have
\[
 f_j'(z)=\frac{1-e^{-a_jz}}z.
\]
A nonzero constant linear combination of these derivatives is $H(w)/z$
with $H\ne0$: the distinct nonzero $a_j$ yield distinct nonconstant
Laurent monomials $e^{-a_jz}$, which cannot cancel each other or the
constant monomial. The simple-pole obstruction shows that no such
combination has a primitive in $K$. The antiderivative criterion
therefore proves the lemma.
\end{proof}

\begin{proof}[Proof of Theorem~\ref{p01:mc17:Ein-field}]
A putative polynomial relation over $\mathcal F_{\exp}$ involves only
finitely many exponentials $e^{c_1},\ldots,e^{c_u}$ with algebraic
$c_i$. Take a $\mathbb Q$-basis for the span of the $c_i$ and $a_j$,
and divide its elements by a common positive integer so that all
$c_i,a_j$ are integer linear combinations of the resulting basis
$b_1,\ldots,b_r$.

Apply the Siegel--Shidlovskii transcendence-degree theorem in the form
of \cite[Theorem~1.1]{p01:bec17:Beukers} to the vector consisting of
\[
 1,\quad e^{b_iz}\ (1\le i\le r),\quad
 e^{-a_jz},\ \operatorname{Ein}(a_jz)\ (1\le j\le s).
\]
These are $E$-functions over a fixed number field. For the integral
functions the coefficients in exponential-series notation are
$(-1)^{n-1}a_j^n/n$. Their conjugates and common denominators grow
at most exponentially: one uses a fixed denominator for $a_j$ and
$\operatorname{lcm}(1,\ldots,n)\le C^n$. The displayed vector
satisfies a homogeneous first-order linear differential system over
$\overline{\mathbb Q}(z)$, since the derivatives of exponentials are
constant multiples of themselves and
$f_j'=(1-e^{-a_jz})/z$. Its only possible finite singularity is zero.
The inclusion of $1$ makes the system homogeneous.

Lemma~\ref{p01:mc17:pole} shows that the generated function field has
transcendence degree $r+s$ over $\mathbb C(z)$, and hence also over
$\overline{\mathbb Q}(z)$. The additional exponentials $e^{-a_jz}$
are Laurent monomials in the $e^{b_iz}$ and do not increase this degree.
The specialization theorem at the nonsingular algebraic point $z=1$
therefore gives
\[
 \operatorname{trdeg}_{\overline{\mathbb Q}}
 \overline{\mathbb Q}\bigl(e^{b_1},\ldots,e^{b_r},
       \operatorname{Ein}(a_1),\ldots,\operatorname{Ein}(a_s)\bigr)
                                                        =r+s.
\]
There are exactly $r+s$ displayed generators. They are consequently
algebraically independent. This rules out the proposed relation,
because its coefficients lie in
$\overline{\mathbb Q}(e^{b_1},\ldots,e^{b_r})$.
\end{proof}

\subsection{Logarithms and logarithmic statistical moments}

\begin{corollary}[A logarithmic transference consequence]
\label{p01:mc17:log-Ein}
Let $a_1,\ldots,a_s$ be distinct nonzero algebraic numbers and choose
any complex logarithm $L_j$ of $\operatorname{Ein}(a_j)$.
For every nonzero vector $(q_1,\ldots,q_s)\in\mathbb Q^s$,
\[
                         \sum_{j=1}^s q_jL_j
\]
is transcendental. In particular, the real number
$\log\operatorname{Ein}(1)$ is transcendental.
\end{corollary}
\begin{proof}
Theorem~\ref{p01:mc17:Ein-field} ensures the values are nonzero. If the
sum were an algebraic number $\beta$, clear the denominators of the
$q_j$ to obtain integers $n_j$, not all zero, and a positive integer
$d$ with
\[
       \prod_j\operatorname{Ein}(a_j)^{n_j}=e^{d\beta}
                                                    \in\mathcal F_{\exp}.
\]
Multiplying out negative powers gives a nonzero polynomial relation
over $\mathcal F_{\exp}$, a contradiction. For $a=1$ the defining
integral is positive, so its real logarithm is defined.
\end{proof}

This proves transcendence of every nonzero rational linear combination;
it does not assert algebraic independence of the logarithms or their
linear independence over all algebraic coefficients.

\begin{corollary}[Logarithmic moments of truncated exponential laws]
\label{p01:mc17:truncated}
For $a>0$, let $X_a$ have density
\[
       p_a(x)=\frac{a e^{-ax}}{1-e^{-a}},\qquad0<x<1,
\]
and put $m(a)=\mathbb E[-\log X_a]$. If $a_1,\ldots,a_s$ are
distinct positive algebraic numbers, then $m(a_1),\ldots,m(a_s)$
are algebraically independent over $\mathcal F_{\exp}$. Every
nonzero rational linear combination of their real logarithms is
transcendental as well.
\end{corollary}
\begin{proof}
Integration by parts, with primitive $(1-e^{-ax})/a$ for $e^{-ax}$,
gives
\[
 \int_0^1(-\log x)e^{-ax}\,dx
       =\frac1a\int_0^1\frac{1-e^{-ax}}x\,dx
       =\frac{\operatorname{Ein}(a)}a.
\]
The boundary term at zero vanishes because $x\log x\to0$.
Consequently
\[
                    m(a)=\frac{\operatorname{Ein}(a)}{1-e^{-a}}.
\]
The denominators are nonzero elements of $\mathcal F_{\exp}$, so
Theorem~\ref{p01:mc17:Ein-field} proves the independence claim. The same
integer-power argument as in Corollary~\ref{p01:mc17:log-Ein} proves
the logarithmic statement.
\end{proof}

This is a transfer to a precisely specified statistical functional,
not a statement that ordinary exponential-family identifiability proves
transcendence. The arithmetic input is the $E$-function theorem.

\subsection{A transcendental extremum location}

Let $N_x$ be Poisson with mean $x>0$, and conditional on $N_x$ take
that many independent exponential random variables of mean one. Let
$M_x$ be their maximum, with $M_x=0$ when $N_x=0$. For $y\ge0$,
\[
 \mathbb P(M_x\le y)=\exp(-xe^{-y}),\qquad
                    \mathbb E M_x=\operatorname{Ein}(x).
\]
The second identity follows by integrating the survival function and
substituting $t=xe^{-y}$. This connects $\operatorname{Ein}$ with a
Poisson sample-size extreme-value statistic rather than introducing an
unrelated named constant.

\begin{proposition}[Unique power-cost optimum]
\label{p01:mc17:optimizer}
For every algebraic $0<\theta<1$, the function
\[
              x\longmapsto\frac{\mathbb E M_x}{x^\theta}
                      =\frac{\operatorname{Ein}(x)}{x^\theta}
\]
has a unique maximizer $x_\theta\in(0,\infty)$, and $x_\theta$ is
transcendental. It is characterized by
\begin{equation}\label{p01:mc17:stationarity}
             1-e^{-x_\theta}=\theta\operatorname{Ein}(x_\theta).
\end{equation}
\end{proposition}
\begin{proof}
Write $f(x)=1-e^{-x}$, $A(x)=\operatorname{Ein}(x)$, and
$h(x)=x/(e^x-1)$. The function $h$ strictly decreases from one to
zero. Moreover
\[
 \frac{f(x)}{A(x)}
  =\frac{\int_0^x h(t)f(t)\,dt/t}{\int_0^x f(t)\,dt/t}>h(x).
\]
For $\eta(x)=f(x)/A(x)$, differentiation therefore gives
\[
                \eta'(x)=\frac{\eta(x)}x\{h(x)-\eta(x)\}<0.
\]
Since $A(x)\sim x$ at zero and $A(x)\to\infty$ at infinity,
$\eta$ decreases continuously from one to zero. The derivative of
$x^{-\theta}A(x)$ has the sign of $\eta(x)-\theta$, proving a unique
maximum and \eqref{p01:mc17:stationarity}. The endpoint limits are zero:
at zero use $A(x)\sim x$, and at infinity use $A(x)\le1+\log x$
for $x\ge1$.

If $x_\theta$ were algebraic, \eqref{p01:mc17:stationarity} would put
$\operatorname{Ein}(x_\theta)$ in $\mathcal F_{\exp}$, contradicting
Theorem~\ref{p01:mc17:Ein-field}. No arithmetic assertion about the maximal
objective value itself is implied by this argument.
\end{proof}

\subsection{The barrier at the classical elementary constants}

The $E$-function theorem specializes at algebraic points and requires
arithmetic control of Taylor coefficients. Neither hypothesis can be
replaced by observable injectivity or by functional independence alone.
For instance
\[
        F(z)=e^z-e,\qquad H(z)=e^{\sqrt2z}-e^{\sqrt2}
\]
are algebraically independent over $\mathbb C(z)$, by independence of
the two exponentials, but both vanish at the algebraic point $z=1$.
Their constant Taylor coefficients are not algebraic, so the arithmetic
specialization theorem does not apply. Likewise, the observable family
$f_c(z)=e^z+c$ identifies an arbitrary complex parameter $c$ exactly
from $f_c(0)$, regardless of whether $c$ is rational, algebraic
irrational, or transcendental.

For the classical targets $e+\pi$, $e\pi$, $e^e$, and $\pi^\pi$,
no unconditional irrationality proof is obtained here. These are among
the standard open arithmetic targets discussed in
\cite{p01:mc17:Waldschmidt}. The product $e\pi$ must not be confused with
$e^\pi$, whose transcendence is classical. One unconditional
observation is that $e+\pi$ and $e\pi$ cannot both be algebraic:
otherwise $e$ and $\pi$ would be roots of a quadratic polynomial with
algebraic coefficients.

For precision, the usual conditional implications are as follows.
Schanuel's conjecture asserts
\[
 \operatorname{trdeg}_{\mathbb Q}
 \mathbb Q(z_1,\ldots,z_n,e^{z_1},\ldots,e^{z_n})\ge n
\]
when the $z_i$ are linearly independent over $\mathbb Q$.
Applying it to $1,i\pi$ would make $e,\pi$ algebraically independent,
and hence also make $e+\pi,e\pi$ algebraically independent, since
$\mathbb Q(e,\pi)$ is algebraic over
$\mathbb Q(e+\pi,e\pi)$. Applying it to $1,e$ would make $e,e^e$
algebraically independent. Finally $i\pi,\log\pi,\pi\log\pi$ are
$\mathbb Q$-linearly independent: imaginary parts eliminate the first
coefficient, and irrationality of $\pi$ eliminates the other two.
Schanuel's conjecture would therefore make
$\pi,\log\pi,\pi^\pi$ algebraically independent. These are
conditional deductions, not proofs of the conjecture or of the open
individual targets. No result above controls $E$-function values at
arbitrary transcendental arguments, and no such extension is assumed.


% Integrated source: arithmetic_attempts.tex
\section{Arithmetic kernels for Euler, Gompertz, Catalan, and zeta constants}
\label{p01:ac17:section}

This contribution tests four named arithmetic targets: the irrationality of
Euler's constant $\gamma$, the irrationality of the Euler--Gompertz constant
$\delta$, the irrationality of Catalan's constant $G$, and the irrationality of
$\zeta(5)$.  The positive result concerns a related, explicitly specified
family of exponential integrals.  It is a consequence of classical
$E$-function theory and does not identify $\delta$ itself as transcendental.
The negative results below are exact statements about specified approximation
constructions; they do not assert that all approximation methods fail.

\subsection{What the arithmetic transfer must preserve}

Put
\[
 \gamma=-\int_0^\infty e^{-t}\log t\,dt,\qquad
 \delta_x=\int_0^\infty\frac{e^{-t}}{t+x}\,dt\quad(x>0),\qquad
 \delta=\delta_1.
\]
For integers $s\ge2$, write
\[
 H_n^{(s)}=\sum_{j=1}^n j^{-s},\quad H_n=\sum_{j=1}^n j^{-1},
 \quad D_n=\operatorname{lcm}(1,\ldots,n),\quad H_0^{(s)}=H_0=0.
\]
A recoverable integral is not automatically an arithmetic certificate.
The elementary certificate needed here is the following: if
$A_n,B_n\in\mathbb Z$ and
$0<|B_n\xi-A_n|\longrightarrow0$, then $\xi$ is irrational.
Indeed, $\xi=a/b$ would make every nonzero form have absolute value at least
$1/b$.  Consequently, a construction must retain a nonzero coefficient of
the intended constant, prove that the form is nonzero, and control its value
\emph{after} clearing denominators.  These three requirements are separate
from the injectivity and stability of an observation transform.

The primary references provide a useful calibration.  Van Assche and Wolfs
construct improved approximants to both $\gamma$ and $\delta$, but their
published denominator/error analysis still does not establish either
irrationality \cite{p01:ac17:VW}.  Rivoal proves disjunctive transcendence for
$\gamma$ and $\delta$ \cite{p01:mc17:Rivoal}.  The relation
$\gamma+\delta/e=\operatorname{Ein}(1)$ and its cancellation mechanism have
also been studied explicitly by Powers \cite{p01:ac17:Powers}.  None of those
facts is counted here as a new individual-constant result.

\subsection{A positive transfer: a geometric family has at most two exceptions}
\label{p01:ac17:dyadic-section}

Define
\[
 \operatorname{Ein}(z)=\int_0^z\frac{1-e^{-u}}u\,du
 =\sum_{n\ge1}\frac{(-1)^{n+1}z^n}{n\,n!},\qquad
 \mathcal E=\overline{\mathbb Q}
       \bigl(e^a:a\in\overline{\mathbb Q}\bigr).
\]
The multi-point $E$-function theorem proved in
Theorem~\ref{p01:mc17:Ein-field} states that distinct nonzero algebraic arguments
$a_1,\ldots,a_m$ give values $\operatorname{Ein}(a_1),\ldots,
\operatorname{Ein}(a_m)$ algebraically independent over $\mathcal E$.
Its external inputs are the Kolchin--Ostrowski theorem and the
Siegel--Shidlovskii specialization theorem, in the form recorded by Beukers
\cite{p01:bec17:Beukers,p01:mc17:Srinivasan}.  The argument below is a transfer of
that result, not another proof or another count of the same independence
theorem.

\begin{lemma}[The exact boundary split]
\label{p01:ac17:split}
For $x>0$,
\begin{equation}\label{p01:ac17:split-eq}
 \operatorname{Ein}(x)=\gamma+\log x+e^{-x}\delta_x.
\end{equation}
\end{lemma}
\begin{proof}
Split the integral of $e^{-t}\log t$ at $x$.  Integration by parts, using
$1-e^{-t}$ on the first interval, gives
\[
 \int_0^x e^{-t}\log t\,dt
 =(1-e^{-x})\log x-\operatorname{Ein}(x).
\]
The endpoint at zero vanishes because $(1-e^{-t})\log t\to0$.
On the second interval,
\[
 \int_x^\infty e^{-t}\log t\,dt
 =e^{-x}\log x+\int_x^\infty\frac{e^{-t}}t\,dt
 =e^{-x}\log x+e^{-x}\delta_x.
\]
Adding and using the definition of $\gamma$ proves the identity.
\end{proof}

\begin{theorem}[A finite defect in a geometric family]
\label{p01:ac17:dyadic}
Fix an integer $b\ge2$.  If $j_1,\ldots,j_m$ are distinct nonnegative
integers, then
\begin{equation}\label{p01:ac17:rank}
 \operatorname{trdeg}_{\mathcal E}
 \mathcal E\bigl(\delta_{b^{j_1}},\ldots,\delta_{b^{j_m}}\bigr)
 \ge m-2.
\end{equation}
Consequently, at most two members of
$\{\delta_{b^j}:j\ge0\}$ are algebraic over $\mathcal E$, and in particular
at most two are algebraic over $\mathbb Q$.
There is a subset of at most two indices whose removal leaves a family
algebraically independent over $\mathcal E$.
\end{theorem}
\begin{proof}
Write $K=\mathcal E(\delta_{b^{j_1}},\ldots,\delta_{b^{j_m}})$.
Since $e^{-b^{j_i}}\in\mathcal E$ and $\log(b^{j_i})=j_i\log b$,
Lemma~\ref{p01:ac17:split} implies
\[
 \mathcal E\bigl(\operatorname{Ein}(b^{j_1}),\ldots,
                  \operatorname{Ein}(b^{j_m})\bigr)
 \subset K(\gamma,\log b).
\]
The field on the left has transcendence degree $m$ over $\mathcal E$ by
Theorem~\ref{p01:mc17:Ein-field}.  Adjoining two elements raises transcendence
degree by at most two.  This proves \eqref{p01:ac17:rank}.  If three members
were algebraic over $\mathcal E$, apply the bound to those three.

For the last assertion, choose a maximal algebraically independent
subfamily $B$ of the countable family, over $\mathcal E$.  Every omitted
member is algebraic over $\mathcal E(B)$ and hence over a field generated
by a finite subset of $B$.  If three members were omitted, take the union
$B_0$ of three finite witness sets.  The finite family consisting of $B_0$
and those three members would have transcendence degree $|B_0|$, contrary
to \eqref{p01:ac17:rank}, which requires at least $|B_0|+1$.
\end{proof}

For example, the twelve explicit integrals
\begin{equation}\label{p01:ac17:twelve}
 \delta_1,\delta_2,\delta_4,\delta_8,\delta_{16},\delta_{32},
 \delta_{64},\delta_{128},\delta_{256},\delta_{512},
 \delta_{1024},\delta_{2048}
\end{equation}
contain at least ten algebraically independent numbers.  The theorem
neither specifies the two possible exceptional indices nor proves that the
first member $\delta_1$ is irrational.  In particular, twelve entries in
\eqref{p01:ac17:twelve} are not twelve individually resolved conjectures.

\paragraph{Provenance and exact gain.}
This is a self-posed family transfer from the multi-point $E$-function
input.  The one-point split and the associated disjunctive results are
classical \cite{p01:mc17:Rivoal}.  The source search did not establish that the
particular geometric-family formulation is new; no priority claim is made.
The substantive gain over a single disjunction is the uniform deficit of
at most two, for every finite subset, including the infinite-family
consequence.  The two lost dimensions are precisely the common terms
$\gamma$ and $\log b$; observing arbitrarily many scales does not assign
those dimensions to individual members.

\subsection{\texorpdfstring{Target $\delta$: an exact optimization over all integer square kernels}{Target delta: an exact optimization over all integer square kernels}}
\label{p01:ac17:gompertz-section}

The integral representation makes positive polynomial weights a natural
first transfer from positive observation kernels.  Here that entire class
can be optimized exactly.

\begin{lemma}[A Laguerre bound with an arithmetic consequence]
\label{p01:ac17:laguerre}
Let $P\in\mathbb R[t]$ have degree $n$ and leading coefficient $a_n\ne0$.
For every $x>0$,
\begin{equation}\label{p01:ac17:laguerre-bound}
 I_x(P):=\int_0^\infty\frac{e^{-t}P(t)^2}{t+x}\,dt
 \ge\frac{a_n^2(n!)^2}{x+2n+1}.
\end{equation}
If $x$ is a positive integer and $P\in\mathbb Z[t]$, then
\begin{equation}\label{p01:ac17:integer-form}
 I_x(P)=P(-x)^2\delta_x+A_x(P),\qquad A_x(P)\in\mathbb Z.
\end{equation}
\end{lemma}
\begin{proof}
Let $\pi_n(t)=(-1)^n n! L_n(t)$ be the monic Laguerre polynomial, where
\[
 L_n(t)=\sum_{k=0}^n(-1)^k\binom nk\frac{t^k}{k!}.
\]
Rodrigues' formula and $n$ integrations by parts give orthogonality to
all polynomials of degree below $n$ in $L^2(e^{-t}dt)$ and
\[
 \int_0^\infty \pi_n(t)^2e^{-t}dt=(n!)^2.
\]
For completeness, Rodrigues' formula is
$e^{-t}\pi_n(t)=(-1)^n(d/dt)^n(e^{-t}t^n)$; the boundary terms vanish,
and $\pi_n^{(n)}=n!$ gives the displayed norm.
The coefficient of $t^{n-1}$ in $\pi_n$ is $-n^2$.
Comparing the two leading coefficients of $t\pi_n$ and $\pi_{n+1}$,
and then using orthogonality, yields
\[
 \int_0^\infty t\pi_n(t)^2e^{-t}dt=(2n+1)(n!)^2.
\]
Also $\int P\pi_n e^{-t}dt=a_n(n!)^2$.  Apply Cauchy--Schwarz to
$P/\sqrt{t+x}$ and $\sqrt{t+x}\pi_n$ in $L^2(e^{-t}dt)$:
\[
 a_n^2(n!)^4\le I_x(P)(x+2n+1)(n!)^2.
\]
This proves \eqref{p01:ac17:laguerre-bound}.
For the arithmetic assertion, division by the monic polynomial $t+x$
gives
$P(t)^2=P(-x)^2+(t+x)Q(t)$ with $Q\in\mathbb Z[t]$.
Every moment $\int_0^\infty t^k e^{-t}dt=k!$ is an integer.
\end{proof}

\begin{theorem}[The exact positive-square obstruction for $\delta$]
\label{p01:ac17:square-optimum}
Over all nonzero integer polynomials of arbitrary degree,
\begin{equation}\label{p01:ac17:optimum}
 \inf_{0\ne P\in\mathbb Z[t]} I_1(P)=4\delta-2>0.
\end{equation}
Equality holds exactly for $P(t)=t-1$ and $P(t)=1-t$.
Thus no sequence of these integral integer linear forms can tend to zero.
\end{theorem}
\begin{proof}
Polynomial division and factorial moments give
\[
 I_1(t-1)=4\delta-2>0,
 \qquad
 \int_0^\infty\frac{e^{-t}t(t-2)^2}{1+t}\,dt=6-9\delta>0.
\]
Both integrands are positive off a finite set, so
$1/2<\delta<2/3$ and $0<4\delta-2<2/3$.
If $\deg P=n\ge2$, Lemma~\ref{p01:ac17:laguerre} gives
$I_1(P)\ge(n!)^2/(2n+2)\ge2/3$; the last expression is increasing for
$n\ge2$.  A nonzero constant integer polynomial gives $I_1(P)\ge\delta$,
which exceeds $4\delta-2$ because $\delta<2/3$.

It remains to consider $P(t)=at+c$, $a,c\in\mathbb Z$, $a\ne0$.
If $|a|\ge2$, the same bound gives $I_1(P)\ge a^2/4\ge1$.
Changing the sign of $P$ reduces the remaining case to $a=1$.  Direct
integration gives
\[
 I_1(t+c)=2c+(c-1)^2\delta,
\quad
 I_1(t+c)-I_1(t-1)=(c+1)\bigl(2+(c-3)\delta\bigr).
\]
For $c\ge0$ both factors are positive, using $\delta<2/3$ when $c=0$.
For $c\le-2$ both factors are negative, using $\delta>1/2$.
Equality therefore occurs only at $c=-1$.
\end{proof}

The result is an exact solution of a restricted, self-posed optimization
problem.  It is not a lower bound for arbitrary integer forms $q\delta-p$.
The restriction matters: signed kernels, rational kernels, and
Hermite--Pad\'e constructions are not covered.  For a rational polynomial
$P=Q/d$, the integer form obtained by the natural square-denominator
clearance is $d^2I_1(P)=I_1(Q)\ge4\delta-2$; merely normalizing the kernel
to have a small integral cannot evade the obstruction.

\subsection{\texorpdfstring{Target $\gamma$: a vanishing observation can erase the constant}{Target gamma: a vanishing observation can erase the constant}}
\label{p01:ac17:gamma-section}

\begin{proposition}[The exact coefficient of $\gamma$]
\label{p01:ac17:gamma-coefficient}
If $P(t)=\sum_{k=0}^d a_kt^k\in\mathbb Z[t]$, then
\begin{equation}\label{p01:ac17:gamma-form}
 \int_0^\infty e^{-t}P(t)\log t\,dt
 =A(P)-B(P)\gamma,
\quad
 A(P)=\sum_{k=0}^d a_k k!H_k\in\mathbb Z,
\quad B(P)=\sum_{k=0}^d a_k k!\in\mathbb Z.
\end{equation}
In particular, for $n\ge1$,
\begin{equation}\label{p01:ac17:gamma-laguerre}
 \int_0^\infty e^{-t}L_n(t)\,dt=0,
 \qquad
 \int_0^\infty e^{-t}L_n(t)\log t\,dt=-\frac1n.
\end{equation}
The small second integral has zero coefficient of $\gamma$.
After multiplying $L_n$ by $n!$ to obtain an integer polynomial, that
integral equals $-(n-1)!$.
\end{proposition}
\begin{proof}
Differentiating $\Gamma(u+1)=u\Gamma(u)$ successively at nonnegative
integers, starting with $\Gamma'(1)=-\gamma$, gives
$\Gamma'(k+1)=k!(H_k-\gamma)$.  Differentiation under the integral is
justified by exponential decay at infinity and integrability of
$t^k|\log t|$ at zero.  Since $k!/j$ is an integer for $1\le j\le k$,
\eqref{p01:ac17:gamma-form} follows.
The coefficient $B(L_n)$ is $\sum_{k=0}^n(-1)^k\binom nk=0$.
Finally,
\[
 \sum_{k=0}^n(-1)^k\binom nkH_k
 =-\int_0^1\frac{(1-u)^n}{1-u}\,du=-\frac1n,
\]
where $H_k=\int_0^1(1-u^k)/(1-u)\,du$.  This proves all assertions.
\end{proof}

This calculation identifies two failures of this particular attempt:
Laguerre orthogonality removes the target coefficient, and the natural
integer normalization reverses the apparent decay.  A successful
replacement must keep $B(P_n)\ne0$ and make
$0<|A(P_n)-B(P_n)\gamma|\to0$ simultaneously.  The observation-theoretic
statement that a family of transforms determines the exponential law does
not provide either inequality.

\subsection{\texorpdfstring{Why adding finitely many scales does not create a canonical
convergent series for $\gamma$}{Why adding finitely many scales does not create a canonical
convergent series for gamma}}
\label{p01:ac17:borel-section}

The following finite-scale version of the canonical cancellation test is
another restricted obstruction.  Its one-scale case is already present
in Powers' analysis \cite{p01:ac17:Powers}; the proof here uses the locations
and residues of the Borel singularities.

For a formal factorial series $\sum_{n\ge0}a_n$, use the Borel convention
$\widehat A(t)=\sum_{n\ge0}a_nt^n/n!$.  Define the two canonical Borel
transforms
\[
 \widehat\gamma(t)=\frac{e^t/(1+t)-1}{t},
 \qquad \widehat\delta_x(t)=\frac1{x+t}.
\]
The apparent singularity of $\widehat\gamma$ at zero is removable.
Their Laplace integrals are $\gamma$ and $\delta_x$:
\[
 \int_0^\infty e^{-t}\widehat\gamma(t)dt
 =\int_0^\infty\frac{(1+t)^{-1}-e^{-t}}t\,dt=\gamma.
\]
For example, the last identity follows by applying the logarithm identity
$\log y=\int_0^\infty(e^{-t}-e^{-yt})dt/t$ to an exponential random
variable $y$ and taking expectation.  Truncation at both endpoints,
followed by dominated convergence, justifies the operation; the resulting
integrand is bounded at zero and $O(t^{-2})$ at infinity.

\begin{proposition}[Finite canonical pole cancellation]
\label{p01:ac17:borel}
Let $x_1,\ldots,x_m>0$ be distinct.  The function
\[
 c\widehat\gamma(t)+\sum_{j=1}^m c_j\widehat\delta_{x_j}(t)
\]
is entire if and only if every coefficient with $x_j\ne1$ vanishes and,
if $x_j=1$ is present, its coefficient equals $c/e$; if $1$ is absent one
must also have $c=0$.  Hence no nonzero linear combination with algebraic
coefficients $c,c_1,\ldots,c_m$ cancels all the canonical poles.
For such a nonzero algebraic combination, the associated factorial series
fails even the necessary term-to-zero condition for ordinary convergence.
\end{proposition}
\begin{proof}
At $t=-1$, the residue of $\widehat\gamma$ is $-e^{-1}$.
The residue of $\widehat\delta_{x_j}$ at $-x_j$ is $1$.
All the poles are simple, and their locations are distinct except for the
possible match $x_j=1$.  This proves the stated cancellation conditions.
After cancellation, the remaining formula has only removable finite
singularities and is entire.  If all coefficients are algebraic and
$c\ne0$, the required equality $c_j=c/e$ is impossible by the
transcendence of $e$; if $c=0$, distinct poles force every $c_j=0$.

To prove the final statement, subtract the principal parts of the finitely
many uncancelled poles.  The remainder is entire.  The uncancelled pole
nearest zero is unique, since all $x_j$ and $1$ are positive.  If its
location is $-r$ and its nonzero residue is $C$, the Taylor coefficient
$b_n$ of the Borel transform satisfies
$b_n=C(-1)^nr^{-n-1}+o(r^{-n})$ by Cauchy's estimate on a larger circle
containing no other uncancelled pole.  Thus the factorial-series term
$n!b_n$ does not tend to zero.
\end{proof}

This theorem concerns the specified canonical representations.  Divergence
of such a representation proves no arithmetic classification of its
Borel sum.  In particular, it cannot upgrade the known transcendence of
$\gamma+\delta/e$ to the transcendence of $\gamma$ or $\delta$ individually.
It does show exactly why algebraic finite-scale coefficients do not
replicate the exceptional cancellation coefficient $1/e$.

\subsection{\texorpdfstring{Target $G$: localization versus the denominator ledger}{Target G: localization versus the denominator ledger}}
\label{p01:ac17:catalan-section}

The exact integral
\[
 G=\sum_{k\ge0}\frac{(-1)^k}{(2k+1)^2}
   =\int_0^1\frac{-\log x}{1+x^2}\,dx
\]
follows by dominated convergence applied to geometric partial sums.
For $P\in\mathbb Z[z]$, polynomial division gives
\begin{equation}\label{p01:ac17:catalan-form}
 \int_0^1\frac{-\log x}{1+x^2}P(x^2)\,dx
 =P(-1)G+\sum_{k=0}^{d-1}\frac{q_k}{(2k+1)^2},
 \qquad q_k\in\mathbb Z,\quad d=\deg P.
\end{equation}
Thus the coefficient of $G$ is an evaluation outside the observation
interval, at $z=-1$.

\begin{proposition}[A complete ledger for beta localization]
\label{p01:ac17:catalan-beta}
Fix positive integers $a,b$.  Set $d_n=(a+b)n$ and
\[
 P_n(z)=z^{an}(1-z)^{bn},\quad
 I_n=\int_0^1\frac{-\log x}{1+x^2}P_n(x^2)dx,\quad
 O_d=\operatorname{lcm}(1,3,\ldots,2d-1).
\]
Then $O_{d_n}^2I_n\in\mathbb ZG+\mathbb Z$, its coefficient of $G$ is
$O_{d_n}^2(-1)^{an}2^{bn}\ne0$, and
\begin{equation}\label{p01:ac17:catalan-growth}
 \lim_{n\to\infty}\frac1n\log(O_{d_n}^2I_n)
 =4(a+b)+a\log a+b\log b-(a+b)\log(a+b)>0.
\end{equation}
Consequently this specified denominator clearance produces positive
forms tending to infinity, not forms tending to zero.
\end{proposition}
\begin{proof}
The arithmetic assertions follow from \eqref{p01:ac17:catalan-form} and
$P_n(-1)=(-1)^{an}2^{bn}$.  Substitute $u=x^2$:
\[
 I_n=\int_0^1 w(u)e^{n\phi(u)}du,
 \quad w(u)=\frac{-\log u}{4\sqrt u(1+u)},
 \quad\phi(u)=a\log u+b\log(1-u).
\]
The unique maximizer is $u_0=a/(a+b)$, with
$-\phi''(u_0)=(a+b)^3/(ab)$.  Laplace's method, whose hypotheses hold
because $w$ is positive and smooth near this interior maximizer and the
endpoints are exponentially negligible, gives
\[
 I_n\sim w(u_0)\sqrt{\frac{2\pi ab}{(a+b)^3 n}}
 \left(\frac{a^ab^b}{(a+b)^{a+b}}\right)^n.
\]
To check endpoint negligibility directly, choose small $\varepsilon>0$
so that $u^a(1-u)^b$ is uniformly below its maximum outside
$(\varepsilon,1-\varepsilon)$; the integrable factor $w$ then gives an
exponentially smaller bound.  The usual Taylor expansion and a Gaussian
dominating bound give the displayed interior asymptotic.

The prime number theorem in the form $\log D_N=N+o(N)$ implies
$\log O_d=2d+o(d)$, since removing the highest power of $2$ changes the
logarithm by only $O(\log d)$.  Combining the two limits proves the
equality in \eqref{p01:ac17:catalan-growth}.  Finally, with
$\theta=a/(a+b)$,
\[
 a\log a+b\log b-(a+b)\log(a+b)
 =(a+b)\{\theta\log\theta+(1-\theta)\log(1-\theta)\}
 \ge-(a+b)\log2,
\]
which proves strict positivity.
\end{proof}

The proposition does not rule out cancellations that substantially reduce
the actual common denominator of the rational part.  It identifies the
missing requirement precisely: a smaller denominator must be proved,
prime by prime, rather than inferred from the analytic decay of $I_n$.
This distinction is especially relevant to the contemporary Catalan
claim: Sun's September 2026 preprint asserts irrationality, while Wachs'
subsequent primary-source analysis reports an incomplete denominator
ledger, including a missing contribution of the prime $2$
\cite{p01:ac17:Sun,p01:ac17:Wachs}.  We use neither that claim nor its negation as a
theorem about $G$; the result established here is
Proposition~\ref{p01:ac17:catalan-beta}.

\subsection{\texorpdfstring{Target $\zeta(5)$: the boundary that carries the constant}{Target zeta(5): the boundary that carries the constant}}
\label{p01:ac17:zeta-section}

\begin{proposition}[The endpoint coefficient]
\label{p01:ac17:zeta-coefficient}
For $s\ge2$ and $P(x)=\sum_{k=0}^d a_kx^k\in\mathbb Z[x]$,
\begin{equation}\label{p01:ac17:zeta-form}
 \frac1{(s-1)!}\int_0^1
 \frac{(-\log x)^{s-1}P(x)}{1-x}\,dx
 =P(1)\zeta(s)-\sum_{k=0}^d a_kH_k^{(s)}.
\end{equation}
The rational part has denominator dividing $D_d^s$.
In particular, the positive localization
$P_n(x)=x^n(1-x)^n$ removes the coefficient of $\zeta(5)$ exactly.
\end{proposition}
\begin{proof}
For a monomial, expand $(1-x)^{-1}$ and use Tonelli's theorem to obtain
\[
 \frac1{(s-1)!}\int_0^1
 \frac{x^k(-\log x)^{s-1}}{1-x}dx
 =\sum_{j\ge0}\frac1{(k+j+1)^s}=\zeta(s)-H_k^{(s)}.
\]
Linearity proves the formula for an arbitrary polynomial, and every
integer $1\le j\le d$ divides $D_d$.  Finally $P_n(1)=0$.
\end{proof}

For example, at $n=1$, the supposedly small $\zeta(5)$ form is simply
\[
 \frac1{24}\int_0^1 x(-\log x)^4dx=\frac1{32};
\]
there is no $\zeta(5)$ coefficient.  Keeping the endpoint by taking
$P_n(x)=x^n$ instead gives
\[
 \zeta(5)-H_n^{(5)}\sim\frac1{4n^4},\qquad
 D_n^5\bigl(\zeta(5)-H_n^{(5)}\bigr)\longrightarrow\infty.
\]
The first asymptotic follows by integral comparison and the second from
$\log D_n=n+o(n)$.  Thus both natural variants are completely accounted
for: one erases the target and the other loses to denominator growth.
Brown and Zudilin's cellular approximants are a genuinely different and
far stronger construction \cite{p01:ac17:BrownZudilin}; these elementary
calculations do not claim to improve them.  No proof of the individual
irrationality of $\zeta(5)$ is obtained in this contribution.

\subsection{The surviving research requirements}

The four target attempts leave different, explicit requirements.
For $\gamma$, the logarithmic observation must retain its factorial-moment
coefficient while producing a small integer form.
For $\delta$, any new positive-kernel construction must leave the
integer-polynomial-square class optimized in
Theorem~\ref{p01:ac17:square-optimum}, or introduce a quantitatively justified
arithmetic normalization outside it.
For $G$, the beta-localization forms require a much smaller denominator
than the stated least-common-multiple clearance.
For $\zeta(5)$, boundary localization must preserve a nonzero endpoint
coefficient and simultaneously beat its denominator growth.
The family result \eqref{p01:ac17:rank} succeeds for a different reason: it
uses algebraic independence of $E$-function values and explicitly pays
for the two shared nuisance constants.  It supplies a proved family
statement without assigning unproved arithmetic properties to any one
of the four named constants.


% Integrated source: information_results.tex
\section{Arithmetic height after exact observation: three constant targets}
\label{p01:ci17:section}

The three transcendence targets considered here are the transcendence of
\[
 P_2=\sum_{p\ {
m prime}}2^{-p},\qquad
 L_{\mathbb P}=\sum_{p\ {
m prime}}\frac1{2^p-1},\qquad
 E=\sum_{n\geq1}\frac1{2^n-1}.
\]
All sums converge absolutely. None of their transcendence questions is
settled in this section. In particular, irrationality of $L_{\mathbb P}$
is not a new target: it was established by Tao and Ter\"av\"ainen
\cite[Theorem~1.3]{p01:ci17:TT}. Irrationality of $E$ is classical
\cite{p01:ci17:Erdos,p01:ci17:Borwein}, and irrationality of $P_2$ has an elementary
proof below. Our contribution to this investigation is a rigorous audit
of two proposed transfers from exact observation to arithmetic:
height-controlled rational truncation and finite Mahler systems.
The first has sharply different costs for these constants; the second
is impossible for the natural generating functions of $P_2$ and $E$.
These are obstruction and calibration results, not new solutions of
historical transcendence problems. The no-Mahler conclusions are
corollaries of existing structural theorems, with proofs supplied here.

For calibration, put
\[
 K=\sum_{j\geq0}2^{-2^j},\qquad
 F=\sum_{j\geq0}\frac1{2^{2^j}+1}.
\]
Kempner's number $K$ is already transcendental
\cite{p01:ci17:Adamczewski}. The reciprocal Fermat constant $F$ is also
already transcendental, and Coons proved the stronger sharp statement
$\mu(F)=2$ \cite[Theorem~1]{p01:ci17:Coons}. Thus neither is counted as a
newly resolved constant. We use them to identify exactly which arithmetic
information a fast observation error does, and does not, provide.

\subsection{Reduced height and the rational-grid correction}

For a reduced rational approximation $A_N/Q_N$ to a real number $\xi$,
with $Q_N\to\infty$, define its approximation exponent along this
particular sequence by
\[
 \beta_N(\xi;A_N/Q_N)
 =\frac{-\log|\xi-A_N/Q_N|}{\log Q_N}.
\]
This is not the irrationality exponent $\mu(\xi)$, which optimizes over
all rational approximations. A limit below two therefore says nothing
by itself about the value of $\mu(\xi)$.

\begin{lemma}[Rational-grid correction]\label{p01:ci17:grid}
Let $A_N\in\mathbb Z$, $Q_N\in\mathbb N$, $Q_N\to\infty$, and let
$c\in\mathbb Q$. Suppose
\[
 0<|Q_N\xi-A_N-c|\longrightarrow0.
\]
Then $\xi$ is irrational. Moreover, if
$|Q_N\xi-A_N-c|=Q_N^{-\sigma+o(1)}$ with $\sigma>0$, the corrected
rational approximations $(A_N+c)/Q_N$ have error
$Q_N^{-1-\sigma+o(1)}$. If their reduced denominators $q_N$ satisfy
$\log q_N/\log Q_N\to1$, their approximation exponents tend to
$1+\sigma$.
\end{lemma}
\begin{proof}
If $\xi=a/b$ and $c=u/v$ are rational, then
$Q_N\xi-A_N-c$ belongs to $(bv)^{-1}\mathbb Z$. Its nonzero absolute
value is at least $1/(bv)$, a contradiction. The remaining assertions
follow by dividing by $Q_N$ and then taking logarithms.
\end{proof}

The nonzero rational center $c$ matters. A normalized truncation error
converging to one can prove irrationality after correction, even though
it does not tend to zero. The following exact calculation records the
height cost of this correction for the reciprocal Fermat constant.

\begin{theorem}[One carry and an exact denominator formula]
\label{p01:ci17:fermat}
Write $f_j=2^{2^j}+1$ and
\[
 F_N=\sum_{j=0}^N\frac1{f_j}=\frac{A_N}{Q_N},\qquad
 Q_N=\prod_{j=0}^N f_j=2^{2^{N+1}}-1.
\]
The displayed fraction is reduced. The raw truncation and its one-unit
numerator correction satisfy
\begin{align}
 Q_N(F-F_N)&\longrightarrow1, &0<Q_N(F-F_N)&<1,
 \label{p01:ci17:fermatraw}\\
 \frac{A_N+1}{Q_N}-F&=Q_N^{-2}(1+o(1)).
 \label{p01:ci17:fermatcorrected}
\end{align}
Furthermore,
\begin{equation}
 \gcd(A_N+1,Q_N)
 =\prod_{\substack{0\leq i\leq N\\2^{i+1}\mid N-i+1}}f_i.
 \label{p01:ci17:fermatgcd}
\end{equation}
Consequently the reduced denominator $q_N$ of $(A_N+1)/Q_N$ obeys
$Q_N/2^{N+1}\leq q_N\leq Q_N$. The raw approximation exponent tends
to one, whereas the corrected exponent tends to two. In particular,
\eqref{p01:ci17:fermatraw} and Lemma~\ref{p01:ci17:grid} give an elementary
proof of the already known irrationality of $F$.
\end{theorem}
\begin{proof}
The product identity follows by repeatedly factoring $x^2-1$.
The $f_j$ are pairwise coprime: if $i<j$, then
$f_j\equiv2\pmod{f_i}$, while $f_i$ is odd. Modulo $f_i$, the numerator
$A_N=\sum_{j=0}^N Q_N/f_j$ is $Q_N/f_i$, which is a unit modulo $f_i$.
This proves reduction.

Set $x=2^{2^{N+1}}=Q_N+1$. The remaining tail is
\[
 R_N=F-F_N=\sum_{j\geq0}\frac1{x^{2^j}+1}.
\]
It is positive and strictly smaller than
$\sum_{m\geq1}x^{-m}=1/(x-1)$, since $2^j$ are distinct positive
integers and each summand is strictly smaller than $x^{-2^j}$.
Also,
\[
 R_N=\frac1{x+1}+\frac1{x^2+1}+O(x^{-4})
     =x^{-1}+O(x^{-3}).
\]
Thus $(x-1)R_N\to1$, and
\[
 \frac1{Q_N}-R_N
 =\frac1{x-1}-R_N=x^{-2}+O(x^{-3})
 =Q_N^{-2}(1+o(1)).
\]
This proves \eqref{p01:ci17:fermatraw}--\eqref{p01:ci17:fermatcorrected}.

For the exact cancellation, reduction modulo $f_i$ gives
\[
 \frac{Q_N}{f_i}
 =\left(\prod_{j<i}f_j\right)\left(\prod_{i<j\leq N}f_j\right)
 \equiv-2\cdot2^{N-i}=-2^{N-i+1}\pmod{f_i}.
\]
(The empty first product for $i=0$ is also $-2$ modulo $f_0=3$.)
Hence $A_N+1\equiv1-2^{N-i+1}\pmod{f_i}$.
For every prime divisor $\ell$ of $f_i$, the order of $2$ modulo
$\ell$ is exactly $2^{i+1}$: it divides $2^{i+1}$, and
$2^{2^i}\equiv-1\not\equiv1\pmod\ell$.
If $2^{i+1}\nmid N-i+1$, no prime divisor of $f_i$ divides
$A_N+1$. If $2^{i+1}\mid N-i+1$, the congruence
$2^{N-i+1}\equiv1\pmod{f_i}$ holds for the whole integer $f_i$.
Pairwise coprimality proves \eqref{p01:ci17:fermatgcd}.

Every index contributing to that product satisfies
$2^{i+1}\leq N+1$. With $L=\lfloor\log_2(N+1)\rfloor$, the product
is at most $\prod_{i=0}^{L-1}f_i=2^{2^L}-1<2^{N+1}$
(the empty product is one). Therefore
$\log q_N=\log Q_N+O(N)$ and $\log q_N/\log Q_N\to1$.
The exponent statements follow.
\end{proof}

This calculation supplies a denominator formula and an explicit
correction mechanism; it does not replace Coons's upper bound
$\mu(F)\leq2$, nor the classical transcendence theorem. Fast convergence
alone had exponent one. Arithmetic control of a normalized remainder
and its cancellation supplied exponent two.

\subsection{Prime supports: exact denominators and a vanishing height exponent}

Let $p_1<p_2<\cdots$ be the primes. For binary observations supported on
an increasing integer sequence $a_j$, denominators have a different
geometry from those of reciprocal Mersenne observations.

\begin{proposition}[Two exact height laws]\label{p01:ci17:heightlaws}
Let $(a_j)_{j\geq1}$ be strictly increasing positive integers and put
$X=\sum_{j\geq1}2^{-a_j}$. Its $N$th truncation $X_N$ has reduced
denominator $2^{a_N}$, and
\begin{equation}
 2^{-a_{N+1}}\leq X-X_N\leq2^{1-a_{N+1}},\qquad
 \beta_N(X;X_N)=\frac{a_{N+1}}{a_N}+O(a_N^{-1}).
 \label{p01:ci17:binaryheight}
\end{equation}
For the prime Lambert constant, the $N$th truncation
\[
 Y_N=\sum_{j=1}^N\frac1{2^{p_j}-1}=\frac{B_N}{D_N}
\]
is reduced with
\begin{equation}
 D_N=\prod_{j=1}^N(2^{p_j}-1),\qquad
 \beta_N(L_{\mathbb P};Y_N)
 =\frac{p_{N+1}+O(1)}{\sum_{j\leq N}p_j+O(1)}.
 \label{p01:ci17:primeheight}
\end{equation}
By the prime number theorem, as $x=p_N\to\infty$,
\begin{align}
 \beta_N(P_2;\textstyle\sum_{j\leq N}2^{-p_j})&\longrightarrow1,
 \label{p01:ci17:primebinarylimit}\\
 \beta_N(L_{\mathbb P};Y_N)&\sim\frac{2\log x}{x}
 \longrightarrow0.
 \label{p01:ci17:lambertlimit}
\end{align}
\end{proposition}
\begin{proof}
The numerator $2^{a_N}X_N$ is odd, since its final summand is one
and all previous summands are even. Summing over every integer exponent
at least $a_{N+1}$ gives the tail bound, and taking logarithms proves
\eqref{p01:ci17:binaryheight}.

The Euclidean algorithm for exponents gives
$\gcd(2^u-1,2^v-1)=2^{\gcd(u,v)}-1$. Thus the denominators
$2^{p_j}-1$ are pairwise coprime. Reduction of $B_N$ modulo any one
of them leaves the unit $D_N/(2^{p_j}-1)$, proving that $B_N/D_N$
is reduced. Also,
\[
 \log_2 D_N=\sum_{j\leq N}p_j
       +\sum_{j\leq N}\log_2(1-2^{-p_j})
       =\sum_{j\leq N}p_j+O(1).
\]
For $m=p_{N+1}$,
\[
 2^{-m}<L_{\mathbb P}-Y_N
 \leq\frac1{1-2^{-m}}\sum_{r=m}^{\infty}2^{-r}
 \leq\frac{2^{1-m}}{1-2^{-m}}.
\]
This proves \eqref{p01:ci17:primeheight}. Finally, the prime number theorem
gives $p_{N+1}/p_N\to1$ and, by partial summation,
$\sum_{p\leq x}p\sim x^2/(2\log x)$.
Substitution gives the two limits.
\end{proof}

\begin{corollary}[Raw truncations are not continued-fraction approximants]
\label{p01:ci17:nonconvergents}
For each of $P_2$, $L_{\mathbb P}$, and $F$, the raw truncations above
are eventually not continued-fraction convergents. For $K$, none of
its raw truncations is a convergent.
\end{corollary}
\begin{proof}
A convergent $a/q$ to an irrational real number satisfies
$|\xi-a/q|<q^{-2}$. For $P_2$, $L_{\mathbb P}$, and $F$, the proved
limits of the raw exponents are less than two. Their irrationality is
known, and for $P_2$ is proved in the next paragraph. For the $N$th
Kempner truncation, $q=2^{2^N}$ and
\[
 K-K_N=\sum_{j\geq1}q^{-2^j}>q^{-2}.
\]
Thus the necessary inequality fails in every case.
\end{proof}

For completeness, the binary expansion of $P_2$ has digit one exactly
at prime positions. It is neither finite nor eventually all ones. If it
were eventually periodic with period $h\geq1$, choose a prime $p$
beyond the preperiod. The position $p+ph=p(1+h)$ would have the same
digit as $p$ but would be composite, a contradiction. Hence $P_2$ is
irrational. This elementary identification of every prime from an exact
binary expansion does not establish transcendence.

Equation~\eqref{p01:ci17:lambertlimit} is particularly restrictive for the
proposed observation transfer. A truncation giving roughly $x$ binary
digits of absolute accuracy has denominator bit length asymptotic to
$x^2/(2\log x)$. The whole prime support is identifiable, yet these
specific rational approximations spend an increasing number of height
bits per recovered value bit. This does not rule out other rational
approximants: the irrationality theorem of Tao--Ter\"av\"ainen already
uses additional arithmetic structure.

\subsection{The precise finite-Mahler obstruction}

A power series $f(z)\in\overline{\mathbb Q}[[z]]$ is $k$-Mahler if
there exist polynomials $r_0,\ldots,r_s\in\overline{\mathbb Q}[z]$,
not all zero, such that
\[
 \sum_{j=0}^s r_j(z)f(z^{k^j})=0,\qquad k\geq2.
\]
An inhomogeneous relation with rational-function right side also implies
a homogeneous one: its successive dilates span a finite-dimensional
space over $\overline{\mathbb Q}(z)$, including the rational function
$1$. We use two established results of Adamczewski, Bell, and Smertnig
\cite[Theorems~1.1--1.2]{p01:ci17:ABS}: a Mahler series with a finite
coefficient alphabet is automatic; and if a Mahler coefficient sequence
has logarithmic Weil heights $h(a_n)=o(\log n)$, then
$h(a_n)=O(\log\log n)$. These are substantial external inputs.

\begin{theorem}[Two natural generating functions cannot enter a finite
linear Mahler system]\label{p01:ci17:nomahler}
For every integer $k\geq2$, neither
\[
 P(z)=\sum_{p\ \mathrm{prime}}z^p
 \quad\hbox{nor}\quad
 D(z)=\sum_{n\geq1}d(n)z^n
     =\sum_{n\geq1}\frac{z^n}{1-z^n}
\]
is $k$-Mahler. Here $d(n)$ counts positive divisors. In particular,
neither function can be a component of a finite vector $\mathbf f(z)$
satisfying
\[
 \mathbf f(z)=A(z)\mathbf f(z^k),\qquad
 A(z)\in\operatorname{GL}_r(\overline{\mathbb Q}(z)).
\]
The assertion excludes finite linear dilation systems for these natural
functions; it is not an assertion about arbitrary representations of
the individual numbers $P(1/2)$ or $D(1/2)$.
\end{theorem}
\begin{proof}
If $P$ were $k$-Mahler, its coefficients, lying in $\{0,1\}$, would
form a $k$-automatic sequence. Thus a finite automaton would recognize
the base-$k$ representations of primes, with no leading zeroes.
Here is a complete pumping argument excluding this.

Let the automaton have $s$ states. Choose a prime $p>k^s$ with $p>k$,
and write its digit word as $uvw$, where $v$ is a nonempty loop among
the first $s$ transitions. Then $uv^t w$ is accepted for every integer
$t\geq0$. Let $a$ be the integer represented by $u$, let $c$ be the
integer represented by $v$, and let $\ell=|v|$. Appending $v$ acts on
a prefix residue modulo $p$ by the invertible affine map
\[
 T(x)=k^{\ell}x+c\quad\hbox{on }\mathbb F_p.
\]
This map has finite order, say $r\geq1$. Therefore
$T^{1+r}(a)=T(a)$ modulo $p$. Appending the unchanged suffix $w$
preserves this congruence, so the integer represented by $uv^{1+r}w$
is divisible by $p$. It is strictly larger than $p$ and still has no
leading zero. The automaton accepts a composite, a contradiction.
The nonregularity of prime representations is classical; see
\cite{p01:ci17:Yuen} for historical attribution and stronger results.

For $D$, we first establish $\log d(n)=o(\log n)$ elementarily.
Fix $\varepsilon>0$. For primes $p\geq2^{1/\varepsilon}$ and exponents
$a\geq1$, $a+1\leq2^a\leq p^{\varepsilon a}$. For the finitely many
smaller primes, each supremum
$\sup_{a\geq0}(a+1)p^{-\varepsilon a}$ is finite. Multiplication gives
$d(n)\leq C_\varepsilon n^\varepsilon$, and the asserted little-oh
bound follows because $\varepsilon$ was arbitrary.
On the other hand, take $n_m=\prod_{j=1}^m p_j$. Then
$d(n_m)=2^m$. Bertrand's postulate gives $p_j\leq2^j$, hence
\[
 \log n_m\leq\tfrac12m(m+1)\log2,
 \qquad \log\log n_m=O(\log m),
 \qquad \log d(n_m)=m\log2.
\]
Thus $\log d(n)$ is not $O(\log\log n)$. Since the logarithmic Weil
height of the positive integer $d(n)$ is precisely $\log d(n)$,
the quoted height-gap theorem excludes a Mahler equation.

Finally, the displayed vector equation implies
$\mathbf f(z^{k^j})=B_j(z)\mathbf f(z)$ with invertible rational-function
matrices $B_j$. For any chosen component, $r+1$ such scalar dilates
lie in a vector space of dimension at most $r$ over
$\overline{\mathbb Q}(z)$. They satisfy a nonzero linear relation;
clearing denominators gives a Mahler relation. The two scalar
obstructions therefore exclude the vector systems as well.
\end{proof}

The Mahler mechanisms for the calibration constants really do have the
missing finite closure. Their natural functions obey
\[
 \mathcal K(z)=z+\mathcal K(z^2),\qquad
 \mathcal F(z)=\frac z{1+z}+\mathcal F(z^2).
\]
By contrast, the no-Mahler theorem rules out obtaining such finite
linear closure for $P$ or $D$ merely by adding finitely many auxiliary
functions. It neither proves their special values algebraic nor proves
them transcendental. For $\sum\omega(n)z^n$, the height-gap argument
used for $d(n)$ does not apply: $\log\omega(n)=O(\log\log n)$, which
is an allowed height range. No no-Mahler assertion about that function
is made here.

\subsection{Infinite inversion remains possible, but does not transfer heights}

The prime support has a natural exact inversion, closely analogous to
recovering an entire frequency family from its observations. Put
\[
 \mathcal L(z)=\sum_p\frac{z^p}{1-z^p},\qquad |z|<1.
\]
Then
\begin{equation}
 \mathcal L(z)=\sum_{j\geq1}P(z^j),\qquad
 P(z)=\sum_{j\geq1}\mu(j)\mathcal L(z^j),
 \label{p01:ci17:mobius}
\end{equation}
where $\mu$ is the M\"obius function. This exact infinite inversion has
an explicit finite observation error.

\begin{proposition}[A controlled infinite observation identity]
\label{p01:ci17:mobiusbound}
For $0<\rho<1$, $|z|\leq\rho$, and $M\geq1$,
\[
 \left|P(z)-\sum_{j=1}^M\mu(j)\mathcal L(z^j)\right|
 \leq
 \frac{\rho^{2(M+1)}}{(1-\rho)(1-\rho^2)^2}.
\]
\end{proposition}
\begin{proof}
For $|w|\leq\rho$,
\[
 |\mathcal L(w)|
 \leq\sum_{n\geq2}\frac{|w|^n}{1-|w|^n}
 \leq\frac{|w|^2}{(1-\rho)(1-\rho^2)}.
\]
Consequently all sums in \eqref{p01:ci17:mobius} converge absolutely.
Expanding its right side as a double sum, the coefficient of $P(z^m)$
is $\sum_{j\mid m}\mu(j)$, which is one for $m=1$ and zero otherwise.
Summing the preceding bound over $j>M$ proves the claim.
\end{proof}

This inversion does not preserve algebraicity under an infinite limit,
and it supplies no bounds on degrees or heights of hypothetical
algebraic values $\mathcal L(z^j)$. The distinction is mathematical,
not merely computational: even a series of rational summands can
converge exponentially or faster to a transcendental number, as the
Kempner and reciprocal Fermat examples already demonstrate. The
no-Mahler theorem explains why replacing this infinite observation
identity by a finite linear dilation system for $P$ cannot work.

\paragraph{What remains for the three targets.}
For $P_2$, the present proof supplies irrationality, a raw height
exponent of one, and exclusion of every finite Mahler system for its
prime-support generating function. It supplies neither an irrationality
exponent greater than two nor a different transcendence criterion.
For $L_{\mathbb P}$, irrationality is already due to
Tao--Ter\"av\"ainen; our exact denominator formula shows that prime
truncations have height exponent tending to zero. Controlling a
suitable corrected remainder or constructing arithmetic auxiliary forms
is still missing. For $E$, the divisor generating function is excluded
from every finite linear Mahler system by the height gap, and its known
irrationality does not become a transcendence proof. The exact gcd formula and rational-grid correction above are proved
statements within this investigation, while historical priority for
those elementary refinements has not been established.


% Integrated source: target_register.tex
\section{Target register and the remaining arithmetic gaps}
\label{p01:root:register}
The following register contains twenty-three explicitly specified targets. Rows are questions, not a count of independent discoveries. In particular, the two bridge examples belong to one theorem, the Bessel extremum coordinates belong to one family, and affine versions of a capacity are not counted again. ``Not resolved here'' is a statement about this manuscript; for less standard boundary constants it is not a claim that an exhaustive literature search has established a historically recognized open problem.

\small
\begin{longtable}{@{}p{0.045\textwidth}>{\raggedright\arraybackslash}p{0.29\textwidth}>{\raggedright\arraybackslash}p{0.60\textwidth}@{}}
\toprule
No.&Target and arithmetic question&Outcome and precise scope\\
\midrule
\endfirsthead
\toprule No.&Target and arithmetic question&Outcome and precise scope\\\midrule
\endhead
1&Euler's constant $\gamma$: irrationality&Not resolved here. A relation involving $\operatorname{Ein}(1)$ supplies a disjunctive statement, not identification of the arithmetic nature of $\gamma$.\\[3pt]
2&Gompertz's constant $\delta$: irrationality&Not resolved here. The positive squared-integer-polynomial integral method has a sharp positive minimum; signed or more general auxiliary forms remain possible.\\[3pt]
3&Catalan's constant $G$: irrationality&Not resolved here. The attempted elementary truncation route fails its denominator test. Recent claimed proofs are not accepted as verified inputs.\\[3pt]
4&$\zeta(5)$: irrationality&Not resolved here. Simultaneous statements about several odd zeta values cannot be converted into an individual conclusion without an additional elimination argument.\\[3pt]
5&$e+\pi$: irrationality or transcendence&Not resolved here. The elementary quadratic argument gives only that this number and $e\pi$ cannot both be algebraic.\\[3pt]
6&$e\pi$: irrationality or transcendence&Not resolved here. This product is different from $e^\pi$. Schanuel's conjecture would give more, but is not proved.\\[3pt]
7&$e^e$: irrationality or transcendence&Not resolved here. The exponent is transcendental; the hypotheses of Gelfond--Schneider and algebraic-point $E$-function specialization do not apply.\\[3pt]
8&$\pi^\pi$: irrationality or transcendence&Not resolved here. The conditional Schanuel deduction is recorded separately from unconditional theorems.\\[3pt]
9&$P_2=\sum_p2^{-p}$: transcendence&Not resolved here. Its known elementary irrationality is proved for calibration. Its natural prime-support generating function is not Mahler for any integer base.\\[3pt]
10&$L_{\mathbb P}=\sum_p(2^p-1)^{-1}$: transcendence&Not resolved here. Irrationality was already proved by Tao--Ter\"av\"ainen. The raw rational truncation exponent tends to zero.\\[3pt]
11&Erd\H{o}s--Borwein constant $E=\sum_{n\ge1}(2^n-1)^{-1}$: transcendence&Not resolved here. Irrationality is classical. The divisor generating function is excluded from every finite linear Mahler system by the coefficient-height gap.\\[3pt]
12&Gamma capacity $A(1)$: irrationality&Proved, uniformly for $A(d)$ at every rational $d>1/2$ (Theorem~\ref{p01:bec17:main}). Transcendence of these parameter roots is not proved.\\[3pt]
13&Excursion boundary $C_{3,1}$: transcendence&Proved as part of the complete rational-parameter bridge classification (Theorem~\ref{p01:ga17:bridge}). No historically stated arithmetic conjecture is claimed closed.\\[3pt]
14&Quadratic reflected-bridge boundary $C_{1,2}$: transcendence&Proved by the same classification. This is a second example of one theorem, not an independent discovery.\\[3pt]
15&Bessel maximum $(y_0,M_0)$: transcendence&Both coordinates are proved transcendental, within the uniform rational-parameter extremum theorem; the normalized critical scale is covered too (Theorem~\ref{p01:ga17:extrema}).\\[3pt]
16&Mixed threshold $\tau_2$ in \eqref{p01:sc:eq:constanttwo}: transcendence&Proved within the one-sided rank-two spectral classification (Theorem~\ref{p01:sc:thm:ranktwo}). The proof does not need an explicit minimizer.\\[3pt]
17&Two-prime edge $\sigma_{2,3}$: irrationality&Proved by a finite-place valuation argument, uniformly for every finite unit-weight prime set of size at least two. Transcendence is not proved; no historical priority claim is made.\\[3pt]
18&Prime-zeta edge $P(\sigma_P)=1$: irrationality&Not resolved here. The finite radical valuation argument does not pass through a limit defined only by real convergence.\\[3pt]
19&Ordered-factorization exponent $\zeta(\rho)=2$: irrationality&Not resolved here. No finite algebraic differential system or sufficiently strong arithmetic approximation is obtained.\\[3pt]
20&$\log\operatorname{Ein}(1)$: transcendence&Proved as a classical-theory consequence. The proof uses independence of exponential-integral values over the entire algebraic exponential field (Theorem~\ref{p01:mc17:Ein-field}).\\[3pt]
21&Poisson extreme-value optimizer $x_{1/2}$: transcendence&Proved, with existence and uniqueness, as a member of the algebraic-power-cost family (Proposition~\ref{p01:mc17:optimizer}). The optimum value itself is not classified.\\[3pt]
22&Signed Brownian-bridge boundary $B$: irrationality&Not resolved unconditionally. The factor $\sqrt{\pi/2}$ blocks the proposed $E$-function argument. Transcendence follows conditionally from the $E$--$G$ value-intersection conjecture (Proposition~\ref{p01:ga17:sheppconditional}).\\[3pt]
23&Dyadic Gompertz family $\delta_1,\delta_2,\ldots,\delta_{2^{m-1}}$: simultaneous transcendence&The family has transcendence degree at least $m-2$. This is a simultaneous transfer from the exponential-integral theorem; it does not select which individual constants are transcendental.\\
\bottomrule
\end{longtable}
\normalsize

\subsection{Which extensions would be substantive next steps?}
For the gamma capacity, proving transcendence would require control of a Kummer family at algebraic, potentially nonrational, parameters or a different arithmetic representation of its parameter root. The present rational-parameter $E$-function proof cannot supply that extension merely by changing the word ``rational'' to ``algebraic.'' A quantitative irrationality measure is a separate problem as well.

For the spectral classification, a natural next question is to remove the one-sided condition on the irrational exponent coordinates. The precise issue is whether an algebraic critical level can be singular at the positive contact point. The existing proof applies whenever nonsingularity is established by another argument. For spectra of rational rank greater than two, two contact equations need not have a zero-dimensional algebraic intersection, so a higher-dimensional theorem needs additional structure.

For the famous individual constants, the missing step is arithmetic isolation. In the Gompertz family, a lower bound on transcendence degree leaves the identity of the exceptional members unspecified. For a single $\delta$, positive squared-polynomial kernels are ruled out by the sharp bound proved above, but signed Pad\'e constructions are not. For $\gamma$, Catalan's constant, and $\zeta(5)$, one needs controlled nonzero integer linear forms or a new specialization theorem with the relevant hypotheses actually verified.

For prime-supported constants, exact recovery of the prime support and uniform tail estimates are already available. The denominator calculations show why this is insufficient: absolute accuracy and arithmetic height can improve at very different rates. The finite-Mahler obstruction rules out one concrete attempt to import the lacunary method, while leaving nonlinear identities, new auxiliary forms, and other arithmetic structures open.

\subsection{Research assessment}
The strongest outcomes here are uniform statements about pre-existing analytic boundaries: a rational-point exclusion on the gamma capacity curve, an exact algebraic-versus-transcendental bridge classification, an arithmetic classification of Bessel extrema, and the rank-two spectral critical-value theorem with its sharp algebraic exception. The remaining positive results are carefully identified consequences of classical transcendence theory, finite radical arithmetic, or exact denominator bookkeeping. No new individual solution of a documented famous constant conjecture has been established in this investigation. The next assessment should focus on correctness, prior overlap, and the scope of these uniform theorems rather than counting their numerical instances.

% Appendix counter reset suppressed in this unified edition.
\section{Verification record and reproducibility}
This file is a standalone \LaTeX{} source: the bibliography and all
proofs are embedded, and there are no external figure, bibliography,
or input-file dependencies. Compilation uses standard pdf\LaTeX{}
packages. The foundational specialization and transcendence theorems
remain cited external inputs; their consequences are proved explicitly.

The gamma-capacity proof and the two probability boundary families were
cross-read independently. The exponential-integral proof was checked
independently for its common-constant hypothesis, its formal pole
argument, and the specialization point. The rank-two spectral proof
was checked for repeated factors, the Newton-polytope argument, and
the algebraic exceptional setting. The positive-square Gompertz
optimization was independently checked, including all polynomial
degrees and the equality cases. These are internal mathematical
reviews, not formal proof-assistant certification or external peer review.

The substantive assumptions that must not be weakened are rational
Kummer parameters, ordinary nonzero algebraic specialization points,
positivity and one-sided irrational coordinates in the spectral
classification, and the integer-square restriction in the sharp
Gompertz obstruction. The conditional signed-bridge result is kept
distinct from the unconditional results.

% END SOURCE: KEIO_Arithmetic_of_Observation_Boundaries.tex

\endgroup
\clearpage
\part{Arithmetic of Optimal Values}
\label{merge:part02}
\begingroup

% BEGIN SOURCE: KEIO_Arithmetic_of_Optimal_Values.tex

\noindent\textbf{Editorial integration note.} 
The stopping-value and extraction results are extended in Parts~\ref{merge:part03} and~\ref{merge:part04}. Keep the distinction between the boundary location, the attained value, and the chosen denominator normalization.


\begin{quote}\small
\noindent\textit{Original part framing.}\par
\noindent\textbf{Status.} Working manuscript with complete proofs for the
statements marked unconditional, cited classical inputs, and explicitly
marked conditional and unresolved statements. Internal independent
cross-reading is not external peer review or formal certification.
\end{quote}
\section{Research question and advances over the preceding paper}
\label{p02:root:scope}

The objective of this investigation is to turn analytic information about
observation and recovery into arithmetic information about individual
constants.  The starting sources are the 223-page foundational manuscript
\cite{p01:root:original}, the 447-page sequel compilation \cite{p02:root:sequels},
and the subsequent 41-page arithmetic manuscript \cite{p02:root:previous}.
This paper extends the last of those sources.  Their lengths are source
identifiers, not evidence for the strength or correctness of a theorem.
Each inherited statement used below is identified as such.

The earlier paper established arithmetic information about several
boundary \emph{locations}, a rank-two spectral classification with a sign
condition, and family-level transcendence statements.  These do not
automatically classify an optimized \emph{value}, remove a geometric
hypothesis, or identify one particular member of a family.  Those three
gaps motivate the present work.

\begin{center}
\small
\begin{tabular}{L{0.21\textwidth}L{0.32\textwidth}L{0.36\textwidth}}
\toprule
Direction & Previously established & Advance in this paper\\
\midrule
Finite spectra & One-sided irrational coordinates in rank two & All sign patterns; exact exceptional setting retained (Theorem~\ref{p02:sc2:thm:ranktwo})\\[4pt]
Stopping payoff & Arithmetic of the optimal boundary & All even-power optimal values at rational dimensions are transcendental (Theorem~\ref{p02:ga18:evenvalue})\\[4pt]
Joint payoff arithmetic & No exact joint classification of objectives & Full relation lattice and an explicit infinite independent family (Theorems~\ref{p02:ga18:joint} and~\ref{p02:ga18:pell})\\[4pt]
Gamma extraction & Rational-input endpoint shape is irrational & Fixed-shape critical rate and contact are transcendental (Theorem~\ref{p02:be18:rate})\\[4pt]
Boundary tables & Large-order asymptotic & No algebraic-coefficient polynomial recurrence on a rational mesh (Theorem~\ref{p02:be18:nonholonomic})\\[4pt]
Classical constants & Selected positive-kernel failures & Bounds for whole nonnegative integer-polynomial classes, with denominator scope stated\\
\bottomrule
\end{tabular}
\end{center}

The most direct individual-constant result is the even-payoff theorem:
transcendence of a stopping boundary alone cannot imply transcendence of
a function evaluated there.  A separate elimination identity supplies
that missing implication.  The joint-value theorem provides a different
kind of information: the ten displayed values are individually specified,
and \emph{all ten} are algebraically independent.  This is stronger than
saying that some unspecified members of a larger list are transcendental.
Nevertheless, their independence is a transfer of classical
Lindemann--Weierstrass to a deliberately selected stochastic family, not
the resolution of ten independent historical conjectures.

The spectral result uses another mechanism.  At an algebraic putative
optimal level, negative curvature along the physical trajectory and
positive curvature in an auxiliary direction exclude a repeated
irreducible component.  Two coprime equations then force algebraicity
in two variables.  This converts a geometric sign issue into an exact
arithmetic classification.  In higher rank the same equations no longer
force algebraicity; the needed additional conjectural input is stated
without being hidden inside the proof.

Throughout, \(\overline{\mathbb Q}\) denotes the algebraic closure of
\(\mathbb Q\) in \(\mathbb C\).  Independence of an infinite collection
means independence of every finite subcollection.  ``Unresolved here''
means that this manuscript has no proof; it does not by itself certify
global historical open status.  Examples obtained by changing a parameter
in one theorem are counted as one target family.  In particular, the
classical-constant investigations do not supply a new proof of
irrationality or transcendence of Euler's constant, the Euler--Gompertz
constant, Catalan's constant, or an individual odd zeta value beyond
previously known results.


% Integrated source: spectral_results.tex
\section{Opposite curvatures remove the spectral sign restriction}
\label{p02:sc2:sec}

The preceding spectral classification required all coefficients of one
irrational basis element to have the same sign.  That restriction ensured
that an algebraic level set was nonsingular at its point of contact.
Here we remove the restriction completely.  A level set may be singular,
but two opposite curvatures prevent the one type of repeated component
that would invalidate the elimination argument.  The result applies to
every finite spectrum in a two-dimensional rational space containing one.
An additional theorem identifies the exact transcendence-degree input
needed in higher dimension.  Its higher-rank classification is explicitly
conditional, and none of these results is presented as the resolution of
a previously named arithmetic conjecture.

\subsection{The unconditional classification}

Let
\begin{equation}\label{p02:sc2:eq:model}
 S(x)=\sum_{j=1}^{m}w_jx^{r_j},\qquad
 T(C)=\min_{x>0}\frac{1+Cx}{S(x)},\qquad
 W=\sum_jw_j,\quad M=\sum_jw_jr_j.
\end{equation}
Throughout, $0<r_j<1$, all weights are positive, and real powers use the
real logarithm.  The weights need not sum to one.

\begin{theorem}[All rank-two spectra, including mixed signs]
\label{p02:sc2:thm:ranktwo}
Suppose $w_1,\ldots,w_m,C$ are positive algebraic numbers and
\begin{equation}\label{p02:sc2:eq:ranktwo}
 r_j=\frac{m_j+n_j\theta}{q},\qquad
 m_j,n_j\in\mathbb Z,\quad q\in\mathbb Z_{>0},\quad
 \theta\in(\overline{\mathbb Q}\cap\mathbb R)\setminus\mathbb Q,
\end{equation}
where at least one $n_j$ is nonzero.  There is no sign restriction on
the $n_j$.  The minimum has a unique contact $x_C$, and
\begin{equation}\label{p02:sc2:eq:classification}
 T(C)\in\overline{\mathbb Q}
 \quad\Longleftrightarrow\quad x_C=1
 \quad\Longleftrightarrow\quad C=C_*:=\frac{M}{W-M}.
\end{equation}
At the unique exceptional setting,
\begin{equation}\label{p02:sc2:eq:exception}
                  T(C_*)=\frac1{W-M}.
\end{equation}
Every other positive algebraic setting gives a transcendental value.
\end{theorem}

The theorem covers, in particular, arbitrary finite subsets of
$\mathbb Q(\sqrt d)\cap(0,1)$ containing an irrational element, for any
positive nonsquare integer $d$.  It concerns the optimized value, whose
contact need not have a known closed form.  Gelfond--Schneider is the
classical arithmetic input: if $a>0$ is algebraic and $a\ne1$, then
$a^\theta$ is transcendental for algebraic irrational $\theta$;
see the primary proof account \cite{p02:sc2:GS}.  The work here is the
reduction of the variational constant to that input in the presence of
singular algebraic levels.

\begin{lemma}[Contact and curvature]\label{p02:sc2:lem:contact}
The contact in \eqref{p02:sc2:eq:model} is unique and satisfies
\begin{equation}\label{p02:sc2:eq:stationary}
        C=\frac{S'(x_C)}{S(x_C)-x_CS'(x_C)}.
\end{equation}
If $\tau=T(C)$, then
\begin{equation}\label{p02:sc2:eq:contactdata}
 S(x_C)=\frac{1+Cx_C}{\tau},\qquad
 S'(x_C)=\frac C\tau,\qquad S''(x_C)<0.
\end{equation}
\end{lemma}
\begin{proof}
The objective diverges at both ends of $(0,\infty)$ because all
exponents are strictly between zero and one.  Put $H=S-xS'>0$.
Then
\[
 \left(\frac{S'}H\right)'=\frac{S''S}{H^2}<0,
 \qquad
 \frac{xS'}H=
 \frac{\sum_jw_jr_jx^{r_j}}{\sum_jw_j(1-r_j)x^{r_j}}.
\]
The last ratio is bounded above and below by positive constants.
Consequently $S'/H$ decreases from infinity to zero, and the unique
stationary point is the unique minimum.  Its value equation and
\eqref{p02:sc2:eq:stationary} imply \eqref{p02:sc2:eq:contactdata}.
\end{proof}

We give the elimination step in a form that also exposes the
higher-dimensional obstruction.  For $z\in(0,\infty)^d$ and a Laurent
polynomial $F$, its logarithmic Hessian is the ordinary real Hessian at
$y=0$ of the function
\[
       y\longmapsto F(z_1e^{y_1},\ldots,z_de^{y_d}).
\]
It is indefinite if its quadratic form takes both positive and negative
values.  This definition applies whether or not the ordinary gradient
of $F$ vanishes.

\begin{lemma}[Coprime equations impose two algebraic constraints]
\label{p02:sc2:lem:coprime}
If two nonzero Laurent polynomials over $\overline{\mathbb Q}$ in
$d\ge2$ variables are coprime, every common torus zero $z$ satisfies
\[
 \operatorname{trdeg}_{\overline{\mathbb Q}}
       \overline{\mathbb Q}(z_1,\ldots,z_d)\le d-2.
\]
\end{lemma}
\begin{proof}
A nonzero equation excludes transcendence degree $d$.  If the degree
were $d-1$, choose $d-1$ coordinates forming a transcendence basis and
relabel them $z_1,\ldots,z_{d-1}$.  Evaluation embeds the rational
function field in these variables into $\mathbb C$.  The remaining
coordinate has a minimal polynomial over that field.  Every polynomial
equation satisfied by the coordinates is divisible by this minimal
polynomial in the one-variable polynomial ring over the rational
function field.  Clearing denominators and taking its primitive part,
Gauss' lemma gives a common nonconstant irreducible factor of the two
original polynomials after their Laurent monomials have been cleared.
This factor is not a coordinate variable, since all coordinates are
nonzero.  It is therefore a common nonunit Laurent factor, a
contradiction.
\end{proof}

\begin{lemma}[No invariant Laurent hypersurface]
\label{p02:sc2:lem:invariant}
Let $b_1,\ldots,b_d$ be rationally independent real algebraic numbers,
and set
\[
                     D=\sum_{i=1}^d b_i z_i\partial_{z_i}.
\]
For every nonunit irreducible Laurent polynomial $H$ over
$\overline{\mathbb Q}$, the polynomials $H$ and $DH$ are coprime.
\end{lemma}
\begin{proof}
Every monomial $z^k$ is an eigenvector of $D$, with eigenvalue $b\cdot k$.
Rational independence makes these eigenvalues distinct, and only the
constant monomial has eigenvalue zero.  Thus $DH=0$ would force $H$ to
be a constant, which is excluded.  If $H$ divides $DH$, write $DH=AH$.
Newton polytopes satisfy
\[
 \mathcal N(A)+\mathcal N(H)=\mathcal N(DH)\subseteq\mathcal N(H).
\]
For completeness, the product identity follows by maximizing a generic
linear functional: the unique maximizing monomial in each factor gives
a noncancelling maximizing product.  Continuity gives the identity for
all support directions.  Taking support functions in the displayed
inclusion forces every exponent of $A$ to be zero, so $A$ is constant.
Distinct eigenvalues then force $H$ to have only one monomial.  Such a
Laurent polynomial is a unit, again a contradiction.
\end{proof}

\begin{theorem}[An indefinite contact has codimension at least two]
\label{p02:sc2:thm:curvature}
Let $F$ be a nonzero Laurent polynomial with real algebraic
coefficients, let $D$ be as in Lemma~\ref{p02:sc2:lem:invariant}, and let
$z\in(0,\infty)^d$.  Suppose
\begin{equation}\label{p02:sc2:eq:tangency}
                         F(z)=DF(z)=0
\end{equation}
and the logarithmic Hessian of $F$ at $z$ is indefinite.  Then
\begin{equation}\label{p02:sc2:eq:codimension}
 \operatorname{trdeg}_{\overline{\mathbb Q}}
       \overline{\mathbb Q}(z_1,\ldots,z_d)\le d-2.
\end{equation}
For $d=2$, both coordinates are algebraic.
\end{theorem}
\begin{proof}
Factor $F$ in the Laurent polynomial ring over $\overline{\mathbb Q}$.
If two distinct irreducible factors vanish at $z$, they are coprime and
Lemma~\ref{p02:sc2:lem:coprime} applies.  Otherwise exactly one
irreducible factor $H$ vanishes.  Write $F=H^eR$ with $R(z)\ne0$.

If $e=1$, the equation $DF(z)=0$ gives $DH(z)=0$.
Lemmas~\ref{p02:sc2:lem:invariant} and~\ref{p02:sc2:lem:coprime} again apply.
It remains to exclude $e\ge2$.  Pass to logarithmic coordinates and
write $h(y)=H(z_1e^{y_1},\ldots,z_de^{y_d})$ and similarly $r(y)$.
Since $h(0)=0$, the Hessian of $h^er$ at zero is zero for $e\ge3$.
For $e=2$, it is
\begin{equation}\label{p02:sc2:eq:rankone}
                   2r(0)\,\nabla h(0)\nabla h(0)^{\mathsf T}.
\end{equation}
This matrix has complex rank at most one.  It is also the Hessian of a
real-valued function on real coordinates, because $F$ has real
coefficients.  Therefore it is a real symmetric matrix of real rank at
most one.  A real symmetric matrix of rank at most one is positive
semidefinite or negative semidefinite.  It cannot be indefinite.
This contradicts the hypothesis in all cases $e\ge2$.
No choice of a real normalization for the complex factor $H$ is needed.
\end{proof}

\begin{proof}[Proof of Theorem~\ref{p02:sc2:thm:ranktwo}]
Suppose $\tau=T(C)$ is algebraic.  Put
\[
 u=x_C^{1/q},\quad v=x_C^{\theta/q},\quad
 P(U,V)=\sum_jw_jU^{m_j}V^{n_j},\quad
 F(U,V)=P(U,V)-\frac{1+CU^q}{\tau}.
\]
Use $D=U\partial_U+\theta V\partial_V$ and $E=V\partial_V$.
Along $(U,V)=(x^{1/q},x^{\theta/q})$, the operator $D$ is
$q x\,d/dx$.  The contact equations yield
\begin{equation}\label{p02:sc2:eq:opposite}
 \begin{aligned}
 F(u,v)&=DF(u,v)=0,\\
 D^2F(u,v)&=q^2x_C^2S''(x_C)<0,\\
 E^2F(u,v)&=\sum_jw_jn_j^2u^{m_j}v^{n_j}>0.
 \end{aligned}
\end{equation}
In the second line, the additional first-derivative term is
$q^2x_C(S'(x_C)-C/\tau)=0$.  The last line is strictly positive
because at least one $n_j\ne0$.  The operators $D$ and $E$ are constant
directions in logarithmic coordinates, so the last two lines imply that
the logarithmic Hessian is indefinite.  In particular $F$ is nonzero.
Theorem~\ref{p02:sc2:thm:curvature} proves that $u$ and $v$ are algebraic.
Since $v=u^\theta$ and $u>0$, Gelfond--Schneider forces $u=1$.
Thus $x_C=1$.

Conversely, evaluating \eqref{p02:sc2:eq:stationary} at one gives
$C=M/(W-M)$, and at this setting the value is
$(1+C)/W=1/(W-M)$, an algebraic number.  This proves every implication.
\end{proof}

The sign removal is essential: the old proof used $F_V\ne0$, whereas
the new proof only needs $E^2F>0$.  At a mixed-sign contact, $F_V$ may
vanish.  The argument above includes that case and all intersections
of distinct irreducible components; it never assumes that the entire
level polynomial is squarefree.

\subsection{Rational settings and genuinely mixed-sign examples}

\begin{corollary}[The rational cancellation criterion]
\label{p02:sc2:cor:rational}
Under Theorem~\ref{p02:sc2:thm:ranktwo}, assume all weights are rational.
Write
\[
              A=\frac1q\sum_jw_jm_j,\qquad
              B=\frac1q\sum_jw_jn_j,
              \qquad M=A+B\theta.
\]
If $B\ne0$, every positive rational setting has transcendental value.
If $B=0$, the unique algebraic value among positive rational settings
occurs at the rational setting $C_*=A/(W-A)$.
\end{corollary}
\begin{proof}
If $B\ne0$, then $M$ is irrational.  A rational exceptional setting
would give $M=C_*W/(1+C_*)\in\mathbb Q$, a contradiction.  If $B=0$,
the exceptional setting is rational, and the main theorem excludes all
other positive rational settings.
\end{proof}

For a concrete family outside the old one-sided hypothesis, put
\begin{equation}\label{p02:sc2:eq:mixedfamily}
 a=\frac1{\sqrt2},\qquad
 \Theta_k=\min_{x>0}
       \frac{2(1+kx)}{x^{a}+x^{1-a}}\quad(k>0).
\end{equation}
Both exponents lie in $(0,1)$, and their coefficients of $\sqrt2$
have opposite signs.  Equal weights give $W=1$ and $M=1/2$.
Consequently
\begin{equation}\label{p02:sc2:eq:explicitmixed}
 \Theta_1=2,\qquad
 \Theta_2,\Theta_3,\Theta_4,\Theta_5,\Theta_6,
 \Theta_7,\Theta_8,\Theta_9,\Theta_{10},\Theta_{11}
 \quad\hbox{are transcendental}.
\end{equation}
These ten constants are distinct.  If $k_2>k_1$ and $x_2$ is the
minimizer for $k_2$, then
\[
 \Theta_{k_2}>
 \frac{2(1+k_1x_2)}{x_2^a+x_2^{1-a}}\ge\Theta_{k_1}.
\]
Their common proof is one classification theorem, not ten independent
solutions of historical conjectures.  The value $\Theta_1$ also supplies
a necessary correction to any attempted extension of the old rational
corollary: rational weights and a rational setting can give an algebraic
threshold when opposite irrational coefficients cancel.

An asymmetric example is
\begin{equation}\label{p02:sc2:eq:asymmetric}
 \Psi=\min_{x>0}
       \frac{3(1+x)}{2x^{1/\sqrt2}+x^{1-1/\sqrt2}}.
\end{equation}
Here $M=(1+1/\sqrt2)/3$ is irrational.  Thus $\Psi$ is
transcendental, as is the corresponding value at every positive
rational setting.  In contrast with the symmetric family, there is no
rational exceptional setting.

\subsection{Higher rational rank: an exact conditional extension}

The curvature theorem remains valid in any dimension.  What fails
unconditionally is the inference from two algebraic constraints to
algebraicity of all coordinates.  The following statement isolates the
missing input without silently invoking a conjecture.

\begin{theorem}[A sufficient exponential-line bound]
\label{p02:sc2:thm:line}
Suppose $w_j,C$ are positive algebraic, all $r_j\in(0,1)$ are algebraic,
and at least one $r_j$ is irrational.  Let
\[
       d=\dim_{\mathbb Q}\operatorname{span}_{\mathbb Q}
                          \{1,r_1,\ldots,r_m\}\ge2.
\]
Choose a real algebraic rational basis $1,\beta_2,\ldots,\beta_d$
and an integer $q>0$ so that
\[
 r_j=\frac1q\left(m_{j1}+\sum_{i=2}^d m_{ji}\beta_i\right),
                       \qquad m_{ji}\in\mathbb Z.
\]
Set $b_1=1/q$ and $b_i=\beta_i/q$ for $i\ge2$.
Suppose the following bound holds at every real $t\ne0$:
\begin{equation}\label{p02:sc2:eq:linebound}
 \operatorname{trdeg}_{\overline{\mathbb Q}}
 \overline{\mathbb Q}(e^{b_1t},\ldots,e^{b_dt})\ge d-1.
\end{equation}
Then the classification \eqref{p02:sc2:eq:classification} and exceptional
value \eqref{p02:sc2:eq:exception} hold for this spectrum.  It is enough to
assume \eqref{p02:sc2:eq:linebound} only at the nonzero contact logarithm,
if such a contact occurs.
\end{theorem}
\begin{proof}
Assume $\tau=T(C)$ is algebraic and put $t=\log x_C$ and
$z_i=e^{b_it}$.  With multi-index notation define
\[
 P(Z)=\sum_jw_j\prod_{i=1}^d Z_i^{m_{ji}},\qquad
 F(Z)=P(Z)-\frac{1+CZ_1^q}{\tau},\qquad
 D=\sum_{i=1}^d b_iZ_i\partial_{Z_i}.
\]
Along the exponential line, $D=x\,d/dx$.  Thus $F(z)=DF(z)=0$
and $D^2F(z)=x_C^2S''(x_C)<0$.  Some index $i\ge2$ has
$m_{ji}\ne0$ for at least one $j$.  For that index,
\[
 (Z_i\partial_{Z_i})^2F(z)
             =\sum_jw_jm_{ji}^2\prod_{\ell=1}^d z_\ell^{m_{j\ell}}>0.
\]
Theorem~\ref{p02:sc2:thm:curvature} gives transcendence degree at most
$d-2$.  If $t\ne0$, this contradicts \eqref{p02:sc2:eq:linebound}.
Therefore $t=0$ and $x_C=1$.  The converse and exceptional value follow
from the contact calculation exactly as before.
\end{proof}

\begin{corollary}[All finite algebraic spectra, conditional on Schanuel]
\label{p02:sc2:cor:schanuel}
Assume Schanuel's conjecture.  For every finite algebraic spectrum
in $(0,1)$ containing an irrational exponent, and every choice of
positive algebraic weights and setting, the equivalence
\eqref{p02:sc2:eq:classification} holds.
\end{corollary}
\begin{proof}
Schanuel's conjecture asserts that rationally independent complex
numbers $s_1,\ldots,s_d$ satisfy
\[
 \operatorname{trdeg}_{\mathbb Q}
       \mathbb Q(s_1,\ldots,s_d,e^{s_1},\ldots,e^{s_d})\ge d;
\]
see \cite{p02:sc2:Schanuel}.  For $t\ne0$, the numbers $b_it$ are
rationally independent.  Extending the base from $\mathbb Q$ to its
algebraic closure does not change transcendence degree.  Since every
$b_i$ is algebraic and $b_1\ne0$,
their field over $\overline{\mathbb Q}$ is generated by $t$.
Adjoining $t$ increases transcendence degree by at most one.
Schanuel therefore implies \eqref{p02:sc2:eq:linebound}, and
Theorem~\ref{p02:sc2:thm:line} applies.
\end{proof}

For $d=2$, bound \eqref{p02:sc2:eq:linebound} follows unconditionally from
Gelfond--Schneider: two algebraic coordinates on the line would be
related by an algebraic irrational power.  In higher dimension we do
not prove the bound.  The conditional corollary is a precise statement
of the remaining arithmetic input, not an unconditional extension.
For example, the value
\begin{equation}\label{p02:sc2:eq:rankthreeexample}
 \min_{x>0}\frac{2(1+x)}{x^{\sqrt2/2}+x^{\sqrt3/3}}
\end{equation}
is transcendental under the corollary: its mean exponent is not $1/2$
and the rational span with one has dimension three.  This manuscript
does not prove its transcendence unconditionally.

\begin{proposition}[Curvature alone does not eliminate dimension three]
\label{p02:sc2:prop:barrier}
There is a polynomial $F$ with real algebraic coefficients in three
variables, irreducible over $\overline{\mathbb Q}$, for which the
conditions $F=DF=0$ and opposite logarithmic curvatures hold on a
positive real curve containing points with transcendental coordinates.
The diagonal weights of $D$ can be rationally independent algebraic
numbers.
\end{proposition}
\begin{proof}
Take
\[
 \begin{aligned}
 D&=u\partial_u+\sqrt2\,v\partial_v+\sqrt3\,z\partial_z,
                  &E&=v\partial_v,\\
 F&=z-1-\sqrt3(u-1)-5(u-1)^2+(v-1)^2.
 \end{aligned}
\]
The polynomial is irreducible: it is monic and linear in $z$, with
coefficient one.  At $p=(1,1,1)$, direct differentiation gives
\[
 F(p)=DF(p)=0,\qquad D^2F(p)=-3-\sqrt3<0,
                         \qquad E^2F(p)=2>0.
\]
Moreover
\[
 \nabla F(p)=(-\sqrt3,0,1),\qquad
 \nabla(DF)(p)=(-\sqrt3-10,2\sqrt2,\sqrt3)
\]
are independent.  The real implicit function theorem gives a smooth
one-dimensional common zero set near $p$.  In a sufficiently small
neighborhood its coordinates remain positive and the two strict
curvature inequalities persist.  This curve has uncountably many
points, whereas algebraic triples form a countable set.  Some of its
points, in fact uncountably many, therefore have a transcendental
coordinate.
\end{proof}

The proposition is an obstruction to an ambient elimination argument.
It does not say that those curve points lie on the specific exponential
line prescribed by a spectral model, and it is not a counterexample to
the conditional classification.  It locates the gap exactly: coprime
equations cut curves in dimension three, while in dimension two they
cut only algebraic points.  No conclusion about prime-zeta edge
constants, $\zeta(\rho)=2$, or Riemann-zeta zero ordinates follows from
this finite algebraic-spectrum result.

\paragraph{Provenance and count.}
The unconditional gain is one classification: removing the sign
restriction from the rank-two spectral threshold theorem.  Its proof
uses a general curvature-and-factorization statement, which also yields
the conditional higher-rank reduction.  The displayed mixed-sign
constants are applications of that classification.  Classical
Gelfond--Schneider and the conjectural Schanuel input are credited
separately.  The literature check did not establish historical priority
for this exact variational statement; no such priority is asserted.


% Integrated source: gaussian_results.tex
\section{From stopping boundaries to optimal objective values}
\label{p02:ga18:section}

The preceding arithmetic analysis classified certain stopping
\emph{locations}; it did not classify the optimal expected payoffs.
This distinction matters because the value of an analytic function at
a transcendental boundary need not be transcendental.  We investigate
three objective-value targets here.  For every rational dimension and
every positive even integral payoff exponent, the optimal value is
transcendental.  On the algebraic-boundary family we determine all
algebraic relations between finite collections of normalized optimal
values, and construct an infinite algebraically independent collection
with integer dimensions and integer payoff exponents.  The linear
Brownian-excursion value remains undecided by these arguments.

The stochastic optimal-stopping formulas are established results of
Ekstr\"om--Wanntorp and Hobson--Liu \cite{p01:ga17:EW,p02:ga18:HL}.
The new work here is their arithmetic analysis, using classical
transcendence theorems and an explicit Wronskian elimination.
No primary source located in our search states these arithmetic
questions as historical open conjectures, and no such closure is
claimed.  The ten examples below illustrate one joint-independence
theorem; they are not ten independent problem closures.

\subsection{The objective and the inherited boundary input}

Let $R$ be a Bessel bridge of dimension $d>0$, starting at zero and
returning to zero at time one.  For $q>0$ define its optimal origin
value by
\begin{equation}\label{p02:ga18:value}
 \mathcal V_{d,q}=\sup_{0\leq\tau\leq1}
                    \mathbb E[R_\tau^q].
\end{equation}
For integer $d$, this is the radial part of a $d$-dimensional
Brownian bridge.  For arbitrary positive $d$ it is understood through
the usual squared Bessel bridge, including its standard behavior at
zero.  Put $a=q/2$, $b=d/2$, and
\[
 f_{a,b}(z)=M(a;b;z)
          =\sum_{n\geq0}\frac{(a)_n}{(b)_n}\frac{z^n}{n!}.
\]
The established stopping formula \cite{p02:ga18:HL} is
\begin{equation}\label{p02:ga18:boundaryvalue}
 z_{d,q}f_{a,b}'(z_{d,q})=a f_{a,b}(z_{d,q}),\qquad
 C_{d,q}=\sqrt{2z_{d,q}},\qquad
 \mathcal V_{d,q}=\frac{(2z_{d,q})^a}{f_{a,b}(z_{d,q})}.
\end{equation}
The optimal rule first crosses $C_{d,q}\sqrt{1-t}$.
For $d=1$ this is the problem with payoff $|X_\tau|^q$ for a
Brownian bridge, not the signed-payoff stopping problem.

We recall the exact arithmetic input from the preceding manuscript
\cite{p02:root:previous}, without counting it as a new result:
\begin{equation}\label{p02:ga18:oldclassification}
 \begin{gathered}
 d,q\in\mathbb Q_{>0}:\qquad
 z_{d,q}\text{ is algebraic}\ \Longleftrightarrow\
 a-b\in\mathbb Z_{\geq0};\\
 \text{otherwise }z_{d,q}\text{ is transcendental}.
 \end{gathered}
\end{equation}
For clarity, both the existence assertion and the arithmetic mechanism
can be recalled briefly.  The positive Taylor coefficients of $f$
give probability weights proportional to $c_nz^n$.  Then
$zf'/f$ is their mean, its derivative with respect to $\log z$
is their strictly positive variance, and this mean increases from
zero to infinity.  There is exactly one positive solution in
\eqref{p02:ga18:boundaryvalue}.

For rational $a,b>0$, $f$ and $f'$ are $E$-functions satisfying
the first-order system
\[
 \binom f{f'}'=\begin{pmatrix}0&1\\a/z&1-b/z\end{pmatrix}\binom f{f'}.
\]
The denominator estimate for $(a)_n/(b)_n$ follows by counting
residues in arithmetic progressions prime by prime; their common
denominators grow exponentially.  The only finite singular point
of the system is zero.  A rational logarithmic derivative of the
entire exponential-type solution forces $f=e^{\lambda z}P(z)$
with $P$ polynomial.  Kummer's equation then forces either
$\lambda=0$, $a=-\deg P$, or $\lambda=1$, $a-b=\deg P$.
The first alternative is impossible.  Outside the second family,
Beukers' linear-specialization theorem
\cite[Corollary~1.4]{p02:be18:Beukers} excludes the boundary relation
at a nonzero algebraic point.  Inside the second family,
\begin{equation}\label{p02:ga18:exceptional}
 f_{a,b}(z)=e^zP_N(z),\qquad
 P_N(z)=M(-N;b;-z),\qquad N=a-b\in\mathbb Z_{\geq0},
\end{equation}
and the boundary equation is a nonzero polynomial equation.
This proves the recalled input and explains why it concerns the
location rather than the objective.

We also use the classical Lindemann--Weierstrass theorem: exponentials
of distinct algebraic numbers are linearly independent over
$\overline{\mathbb Q}$.  It follows here directly from the same
Beukers theorem.  The functions $e^{\lambda_j z}$ form a system
with constant algebraic diagonal coefficients and are linearly
independent over $\mathbb C(z)$ when the $\lambda_j$ are distinct.
Indeed, after clearing denominators, a polynomial-coefficient
relation is impossible along a ray on which the real parts of
$\lambda_jz$ are all distinct: its largest exponential dominates.
Specializing at $z=1$ gives the stated arithmetic conclusion.

\subsection{Target I: every even-payoff optimum is transcendental}

The essential new identity eliminates the special function from the
optimal value, even when the boundary itself is transcendental.

\begin{lemma}[Contiguous reduction and its constant determinant]
\label{p02:ga18:wronskian}
Let $b>0$ and $m\geq1$ be an integer.  Define polynomials
$A_m,B_m\in\mathbb Q[b,z]$ by $A_1=1$, $B_1=0$, and
\begin{align}
 A_{m+1}&=\frac{zA_m'+(z+m+1-b)A_m}{m},
                                                        \label{p02:ga18:Arec}\\
 B_{m+1}&=\frac{(b-1)A_m+zB_m'+mB_m}{m}.                 \label{p02:ga18:Brec}
\end{align}
Thus rational numerical $b$ gives rational-coefficient polynomials.
Put
\begin{align}
 U_m&=zA_m'+(z+1-b-m)A_m,                                \label{p02:ga18:U}\\
 T_m&=(b-1)A_m+zB_m'-mB_m,                              \label{p02:ga18:T}\\
 K_m(b)&=\frac{\prod_{j=1}^m(j-b)}{(m-1)!}.              \label{p02:ga18:K}
\end{align}
Then
\begin{equation}\label{p02:ga18:reduction}
 M(m;b;z)=A_m(z)M(1;b;z)+B_m(z),\qquad
 B_mU_m-A_mT_m=K_m(b).
\end{equation}
Moreover $\deg A_m=m-1$ with leading coefficient $1/(m-1)!$,
and $\deg U_m=m$ with the same leading coefficient.
\end{lemma}
\begin{proof}
Write $\theta=z\,d/dz$ and $Y=M(1;b;z)$.  Comparison of Taylor
coefficients gives
\begin{equation}\label{p02:ga18:basefirstorder}
 \theta Y=(z+1-b)Y+(b-1),\qquad
 M(m+1;b;z)=\frac{(\theta+m)M(m;b;z)}m.
\end{equation}
These are special cases of the standard contiguous relations
\cite[\S13.3]{p02:be18:DLMF}, but the coefficient check proves them
directly.  Substitution of the first identity into the second
gives exactly \eqref{p02:ga18:Arec}--\eqref{p02:ga18:Brec}, proving the
first formula in \eqref{p02:ga18:reduction}.  The degree and leading
coefficient claims follow by induction.  The constant term is
\begin{equation}\label{p02:ga18:Aconstant}
 A_m(0)=\frac{\prod_{j=2}^m(j-b)}{(m-1)!},\qquad
 A_m(0)+B_m(0)=1.
\end{equation}

Here is a complete derivation of the determinant identity, including
the potentially vanishing exceptional constants.  On $z>0$ let
$H(z)=e^zz^{1-b}$.  The equation
\[
 \mathcal L_m f=0,\qquad
 \mathcal L_m=\theta(\theta+b-1)-z(\theta+m),
\]
is Kummer's equation multiplied by $z$.  Direct substitution gives
$\mathcal L_1H=0$, and the operator identity
\[
 \mathcal L_{m+1}(\theta+m)=(\theta+m)\mathcal L_m
\]
follows from $\theta z=z(\theta+1)$.  Thus repeated application
of the contiguous operator shows that both
$F_m=M(m;b;z)=A_mY+B_m$ and $G_m=A_mH$ solve Kummer's
equation with upper parameter $m$.

The Wronskian in the convention $W(F,G)=FG'-F'G$ therefore obeys
$W'=(1-b/z)W$.  It is a constant times $e^zz^{-b}$.
Since $H'/H=1+(1-b)/z$, direct multiplication using
\eqref{p02:ga18:basefirstorder} gives
\begin{align*}
 W(F_m,G_m)
 &=\frac{H}{z}\left[
 z(B_mA_m'-A_mB_m')+(z+1-b)A_mB_m-(b-1)A_m^2\right]\\
 &=e^zz^{-b}(B_mU_m-A_mT_m).
\end{align*}
It follows that the polynomial in the last parentheses is constant.
At zero, \eqref{p02:ga18:Aconstant} evaluates it as
$(1-b)A_m(0)=K_m(b)$.  The calculation remains valid when this
constant is zero; linear independence of $F_m,G_m$ was not assumed.
\end{proof}

\begin{theorem}[Transcendence of all rational-dimensional even-payoff values]
\label{p02:ga18:evenvalue}
For every $d\in\mathbb Q_{>0}$ and every integer $m\geq1$,
$\mathcal V_{d,2m}$ is transcendental.  More explicitly, put
$b=d/2$ and $z=z_{d,2m}$.  If $b\notin\{1,\ldots,m\}$,
\begin{equation}\label{p02:ga18:polynomialvalue}
 \mathcal V_{d,2m}
       =\frac{(2z)^m U_m(z)}{K_m(b)},
\end{equation}
a nonconstant rational-coefficient polynomial of degree $2m$
in the transcendental boundary variable $z$.  If
$b\in\{1,\ldots,m\}$, then $z$ is positive algebraic and
$\mathcal V_{d,2m}$ is a nonzero algebraic multiple of $e^{-z}$.
\end{theorem}
\begin{proof}
At the optimal boundary the smooth-fit equation is
$(\theta-m)F_m=0$.  By the first identity of
\eqref{p02:ga18:reduction} and \eqref{p02:ga18:basefirstorder}, this
is precisely $U_m(z)Y(z)+T_m(z)=0$.  Multiplying
$F_m=A_mY+B_m$ by $U_m$ and using the constant determinant gives
\begin{equation}\label{p02:ga18:atcontact}
                       F_m(z)U_m(z)=K_m(b).
\end{equation}
If $b\notin\{1,\ldots,m\}$, the right side is nonzero.
In particular $U_m(z)\ne0$, so no division by a potentially
vanishing polynomial is hidden in the argument.  Formula
\eqref{p02:ga18:boundaryvalue} now gives
\eqref{p02:ga18:polynomialvalue}.  Its degree is $2m$ by
Lemma~\ref{p02:ga18:wronskian}.  The boundary classification
\eqref{p02:ga18:oldclassification} says that $z$ is transcendental.
Were $\mathcal V_{d,2m}$ algebraic,
\eqref{p02:ga18:polynomialvalue} would be a nonzero polynomial
equation over $\overline{\mathbb Q}$ for $z$, a contradiction.

If $b\in\{1,\ldots,m\}$, set $N=m-b\geq0$.
Equations \eqref{p02:ga18:oldclassification} and
\eqref{p02:ga18:exceptional} give
\[
 \mathcal V_{d,2m}
      =\frac{(2z)^m}{P_N(z)}e^{-z}.
\]
The prefactor is nonzero algebraic because $z>0$ is algebraic
and $P_N(z)>0$.  Lindemann--Weierstrass makes $e^{-z}$
transcendental.  This covers every positive rational $d$.
\end{proof}

For the quadratic payoff the new polynomial identity is especially
short.  Here $U_1=z-b$, $K_1=1-b$, so for $d\ne2$,
\begin{equation}\label{p02:ga18:quadratic}
 \mathcal V_{d,2}=\frac{C_{d,2}^2(C_{d,2}^2-d)}{2-d}.
\end{equation}
Thus the reflected Brownian-bridge value is
$C_{1,2}^4-C_{1,2}^2$, while the three-dimensional quadratic
value is $C_{3,2}^2(3-C_{3,2}^2)$; both are transcendental.
The exceptional value is $\mathcal V_{2,2}=2/e$.
At the next payoff exponent,
\[
 U_2(z)=z^2+2(1-b)z+(b-2)(b+1),\qquad
 K_2=(1-b)(2-b),
\]
which gives an explicit quartic polynomial in $z$ whenever
$b\ne1,2$.  These are instances of one uniform theorem.
The polynomial relation also shows why no algebraic independence
between a nonexceptional even-payoff value and its own boundary
is asserted: they are algebraically dependent.

\subsection{Target II: all algebraic relations on the algebraic-boundary family}

Take any finite list of rational pairs $(d_i,q_i)$ satisfying
$q_i-d_i=2N_i$ with $N_i\in\mathbb Z_{\geq0}$, and let
$z_i=z_{d_i,q_i}$.  Define the positive algebraic number
\begin{equation}\label{p02:ga18:prefactor}
 a_i^*=\frac{(2z_i)^{q_i/2}}
                    {M(-N_i;d_i/2;-z_i)},\qquad
 W_i=\frac{\mathcal V_{d_i,q_i}}{a_i^*}=e^{-z_i}.
\end{equation}
Rational powers here denote their positive real branches, which
are algebraic at positive algebraic arguments.

\begin{theorem}[Exact transcendence degree and relation lattice]
\label{p02:ga18:joint}
For every such finite list,
\begin{equation}\label{p02:ga18:rank}
 \operatorname{trdeg}_{\overline{\mathbb Q}}
 \overline{\mathbb Q}(\mathcal V_{d_1,q_1},\ldots,
                     \mathcal V_{d_s,q_s})
      =\dim_{\mathbb Q}\operatorname{span}_{\mathbb Q}
                         \{z_1,\ldots,z_s\}.
\end{equation}
More precisely, the kernel of the evaluation map from
$\overline{\mathbb Q}[X_1^{\pm1},\ldots,X_s^{\pm1}]$ to
$\mathbb C$, $X_i\mapsto W_i$, is generated by the binomials
\begin{equation}\label{p02:ga18:lattice}
 X_1^{n_1}\cdots X_s^{n_s}-1,
 \qquad n\in\mathbb Z^s,\quad \sum_i n_i z_i=0.
\end{equation}
Thus the theorem determines all algebraic relations, after the
explicit algebraic rescaling \eqref{p02:ga18:prefactor}.
\end{theorem}
\begin{proof}
Expand any Laurent polynomial at $W_i=e^{-z_i}$ and group its
monomials by the algebraic exponent $\sum_i n_i z_i$.
Lindemann--Weierstrass says that the sum vanishes exactly when
the coefficients in each group add to zero.  Within a group,
subtract one fixed monomial from the others; their exponent
differences lie in the integer relation lattice.  This proves
the kernel statement and shows that it is an exact description,
not merely a list of some relations.

Let the rational rank of the $z_i$ be $r$.  Select $r$ of them
as a basis.  Their exponentials are algebraically independent:
different polynomial monomials give distinct algebraic exponents,
again by Lindemann--Weierstrass.  Every remaining exponential has
an integral positive power equal to a Laurent monomial in the
selected exponentials, because its exponent is a rational linear
combination of the basis.  It is therefore algebraic over the
field they generate.  The nonzero algebraic prefactors do not
change transcendence degree, proving \eqref{p02:ga18:rank}.
\end{proof}

The theorem includes dependent examples.  On the diagonal $q=d$
one has $z=d/2$, so every finite collection has transcendence
degree one, regardless of how many rational dimensions are used.
Consequently merely listing many transcendental optimal values
would not demonstrate independence.  The following construction
does produce independence with actual integer dimensions.

\begin{theorem}[Infinitely many algebraically independent integer-dimensional optima]
\label{p02:ga18:pell}
For each prime $p$, choose any positive integer solution
\[
 x_p^2-py_p^2=1,\qquad x_p>1,
\]
and put $d_p=x_p-1$, $q_p=x_p+1$.  Then the collection
\begin{equation}\label{p02:ga18:pellvalues}
 \left\{\mathcal V_{d_p,q_p}:p\text{ prime}\right\}
\end{equation}
is algebraically independent over $\overline{\mathbb Q}$,
meaning that every finite subcollection is algebraically independent.
Every $d_p$ and $q_p$ is a positive integer, so these are optimal
expected radial-power payoffs for ordinary multidimensional
Brownian bridges.  Explicitly,
\begin{equation}\label{p02:ga18:pellexplicit}
 \mathcal V_{d_p,q_p}
 =\frac{(y_p\sqrt p)^{(x_p+1)/2}}
              {1+y_p\sqrt p/(x_p-1)}
                  \exp\left(-\frac{y_p\sqrt p}{2}\right).
\end{equation}
\end{theorem}
\begin{proof}
Here $q_p=d_p+2$, so $N=1$ and
$P_1(z)=1+2z/d_p$.  The boundary polynomial reduces to
$z^2=d_p(d_p+2)/4$.  Hence
\[
 z_p=\frac{\sqrt{d_p(d_p+2)}}2
     =\frac{y_p\sqrt p}{2}>0.
\]
Formula \eqref{p02:ga18:pellexplicit} follows from
\eqref{p02:ga18:prefactor}.  The numbers $\sqrt p$ for distinct
primes are rationally linearly independent.  One way to see this
is to use the sign-changing automorphisms of the multiquadratic
field generated by these square roots: prime factorization makes
their square classes independent, and changing the sign of just
one square root in a putative linear relation isolates its
coefficient.  Nonzero rational factors $y_p/2$ preserve this
independence.  Theorem~\ref{p02:ga18:joint} now proves the claim.

For completeness, the positive integral Pell solutions used here
exist without any arithmetic conjecture.  Dirichlet's elementary
pigeonhole approximation gives infinitely many positive integers
$r,s$ with $|r-s\sqrt p|<1/s$.  Their nonzero integral norms
$r^2-ps^2$ lie in a fixed bounded set.  Infinitely many pairs
therefore have the same nonzero norm $k$, and an infinite
subcollection has the same residues of $r,s$ modulo $|k|$.
For two distinct members $(r_1,s_1),(r_2,s_2)$ of this
subcollection, chosen with $r_1+s_1\sqrt p>r_2+s_2\sqrt p$,
the quotient
\[
 \frac{r_1+s_1\sqrt p}{r_2+s_2\sqrt p}
 =\frac{r_1r_2-ps_1s_2}{k}
       +\frac{s_1r_2-r_1s_2}{k}\sqrt p
 =x+y\sqrt p
\]
has integral coefficients, norm one, and value greater than one.
Its conjugate is its positive reciprocal, so $x>1$ and $y>0$.
This supplies a solution for each prime.
\end{proof}

Here are ten explicit choices, all satisfying the displayed Pell
equation exactly.  They exhibit one joint theorem, not ten unrelated
arithmetic claims.
\begin{center}
\begin{tabular}{r|r|r|r|r|l}
$p$&$x_p$&$y_p$&$d_p$&$q_p$&$z_p$\\\hline
2&3&2&2&4&$\sqrt2$\\
3&2&1&1&3&$\sqrt3/2$\\
5&9&4&8&10&$2\sqrt5$\\
7&8&3&7&9&$3\sqrt7/2$\\
11&10&3&9&11&$3\sqrt{11}/2$\\
13&649&180&648&650&$90\sqrt{13}$\\
17&33&8&32&34&$4\sqrt{17}$\\
19&170&39&169&171&$39\sqrt{19}/2$\\
23&24&5&23&25&$5\sqrt{23}/2$\\
29&9801&1820&9800&9802&$910\sqrt{29}$
\end{tabular}
\end{center}
For example, the optimal values for $(d,q)=(2,4)$ and $(1,3)$
are algebraically independent, not just individually transcendental.
The infinite family is deliberately selected to make its arithmetic
accessible.  Its probabilistic meaning is standard, while the
selection and the arithmetic question are posed here; historical
novelty is not inferred from the absence of a search match.

\subsection{Target III: the linear Brownian-excursion value is still unresolved}

Let $C=C_{3,1}$ be the established Brownian-excursion stopping
boundary and put $z=C^2/2$.  The preceding boundary theorem proves
that $C$ and $z$ are transcendental, but it does not classify
$\mathcal V_{3,1}$.  Its exact reduction is
\begin{equation}\label{p02:ga18:excursionvalue}
 M(1/2;3/2;z)=\int_0^1e^{zt^2}\,dt,\qquad
 2M(1/2;3/2;z)=e^z,\qquad
 \mathcal V_{3,1}=2C e^{-C^2/2}.
\end{equation}
Indeed differentiation of the integral gives
$z f'(z)=(e^z-f(z))/2$, while smooth fit gives $zf'=f/2$.
This proves the middle identity and then the objective formula.
It agrees with the Brownian-excursion value formula in
\cite{p02:ga18:HL}.

If the value were an algebraic number $v>0$, its boundary variable
would have to satisfy the two exact equations
\begin{equation}\label{p02:ga18:excursionbarrier}
 2M(1/2;3/2;z)-e^z=0,\qquad
                     8z-v^2e^{2z}=0.
\end{equation}
The specialization theorem for $E$-functions only controls their
values at algebraic arguments; here $z$ is already known to be
transcendental.  Nor does Lindemann--Weierstrass apply to this
transcendental exponent.  We have not proved a common-zero theorem
that excludes \eqref{p02:ga18:excursionbarrier}.  Thus neither
irrationality nor transcendence of $\mathcal V_{3,1}$ is established
here.  No global historical open-status claim is made solely from
our literature search.  This target is explicitly recorded as
unresolved by the present work, rather than being inferred from
transcendence of the boundary.

\paragraph{Scope and accounting.}
There are two proved objective-value target families: the even-payoff
transcendence theorem, whose main new step is polynomial elimination
at a transcendental boundary, and the exact joint-relation theorem
on algebraic-boundary values, with its infinite independent family.
The latter uses classical Lindemann--Weierstrass and is not presented
as a new transcendence theorem.  The linear excursion value is a
third investigated target without an unconditional arithmetic
conclusion.  General irrational payoff exponents, generic rational
non-even payoffs outside the algebraic-boundary family, and ordinary
exponential discounting are not resolved by the arguments above.


% Integrated source: bessel_results.tex
\section{Transcendental extraction rates and an obstruction to holonomic boundary tables}
\label{p02:be18:section}

The preceding arithmetic paper proved that the endpoint extraction
constant $A(d)$ is irrational for every rational $d>1/2$.  It did
not prove that these numbers are transcendental.  We first examine
$A(1)$ and $A(3/2)$ separately and make that limitation precise:
under the remaining algebraic-irrational hypothesis, the relevant
Kummer function is provably not an $E$-function.  We then prove
transcendence for a different, exactly specified constant of the same
factorization problem: the critical intermediate rate for extracting
an exponential factor.  Finally, the large-order boundary asymptotic
implies that the boundary table on every positive rational arithmetic
progression is not polynomial-recursive over the algebraic numbers.
These are transfers within the developed model; no named historical
constant conjecture is asserted to be closed.

\subsection{The two endpoint targets and a provable obstruction}

We recall only the exact inputs needed from the earlier manuscript
\cite{p02:root:sequels}.  For $d>1/2$ there is a unique
\[
 A(d)\in\bigl(\max(0,1-d),1/(4d)\bigr)
\]
for which, with $a=A(d)$ and $\eta=a+d-1$,
\begin{equation}\label{p02:be18:endpoint}
 y(z)-\frac{z}{2(1-a)}y'(z)=0,
 \qquad y(z)=M(a-1;\eta;z),\qquad z=2(d-a).
\end{equation}
Here $M={}_1F_1$.  The contact coordinate $x=d-a$ is a double zero
of the boundary remainder density.  The previous arithmetic argument
excludes rational $a$ by applying Beukers' theorem to $(y,y')$.
That theorem cannot simply be extended by replacing ``rational''
with ``algebraic'' in the Kummer parameters.

\begin{proposition}[The remaining algebraic hypothesis is outside the $E$ class]
\label{p02:be18:notE}
Let $d>0$ be rational and let $a$ be algebraic irrational.  Then
$M(a-1;a+d-1;z)$ is not an $E$-function.
\end{proposition}
\begin{proof}
Put $\alpha=a-1$ and $\beta=a+d-1$.  Neither is an integer,
and the series is nonpolynomial.  If
\[
 M(\alpha;\beta;z)
 =\sum_{n\geq0}\frac{(\alpha)_n}{(\beta)_n}\frac{z^n}{n!}
\]
were an $E$-function, its formal Borel transform
\[
 \sum_{n\geq0}\frac{(\alpha)_n}{(\beta)_n}z^n
       ={}_2F_1(\alpha,1;\beta;z)
\]
would be a $G$-function.  Indeed the coefficient height and common-denominator bounds are the same, the latter series has positive
radius of convergence, and its hypergeometric differential equation
supplies holonomicity.

Galochkin's characterization, in the form reproved by Rivoal
\cite[Theorem~1]{p02:be18:Rivoal}, says that every irrational numerator
parameter of a nonpolynomial hypergeometric $G$-function must pair
with an irrational denominator parameter so that numerator minus
denominator is a positive integer, after equal parameters have been
cancelled.  Here the only irrational numerator is $\alpha$, the
only irrational denominator is $\beta$, and
$\alpha-\beta=-d<0$.  They do not cancel and cannot form such a
pair.  This contradicts the characterization.
\end{proof}

Thus the obstruction is an actual failure of the arithmetic-function
hypothesis, not merely an omitted denominator estimate.  Differential
Galois calculations concern functional relations; they do not by
themselves provide the missing arithmetic specialization statement.

\paragraph{Target $A(1)$: an incomplete-gamma reduction.}
Write $a=A(1)$, $p=1-a$, so $0<a<1/4$ and $z=2p$.  The elementary
contiguous identity
\[
 z y'(z)-p y(z)=-p e^z,
 \qquad y(z)=M(-p;1-p;z),
\]
follows termwise from its Taylor series.  Equation
\eqref{p02:be18:endpoint} at $z=2p$ is therefore equivalent to
$y(z)=-e^z$.  Using $y'=-(p/a)M(a;a+1;z)$ gives the exact
positive-integral equation
\begin{equation}\label{p02:be18:A1}
 (1-a)\int_0^1 t^{a-1}e^{2(1-a)t}\,dt=e^{2(1-a)}.
\end{equation}
The integral is $M(a;a+1;2(1-a))/a$, or the corresponding
branch-consistent incomplete-gamma expression; see
\cite[13.6.5]{p02:be18:DLMF}.  If $a$ were algebraic irrational,
the order of that incomplete gamma function would still be
algebraic irrational.  Murty and Saha's special-value theorem
\cite[Corollary~1]{p02:be18:MurtySaha} assumes rational order.  It does
not settle \eqref{p02:be18:A1}.  Transcendence of the exponential on
the right does not contradict equality with an unknown integral on
the left.  We obtain no transcendence proof for $A(1)$.

\paragraph{Target $A(3/2)$: a positive Euler integral.}
An alternative formulation makes the other obstruction explicit.
For positive $a,d$ with $a+d>1$, put $U(z)=M(a;a+d;z)$.  The inverse of
$u^{-a}(u+2)^{-d}$ is
\[
 k(x)=\frac{e^{-2x}x^{a+d-1}}{\Gamma(a+d)}U(2x),
\]
and the endpoint remainder density is $k'+k$.  At
$x=d-a$, its vanishing is equivalent to
\begin{equation}\label{p02:be18:Eulercontact}
 \int_0^1 \{2a-1+2(d-a)t\}
       e^{2(d-a)t}t^{a-1}(1-t)^{d-1}\,dt=0.
\end{equation}
This follows by differentiating the displayed $k$ and applying
Euler's integral for $U$, whose positive beta normalizer cancels.
In particular, for $d=3/2$ it becomes
\begin{equation}\label{p02:be18:Athreehalves}
 \int_0^1 \{2a-1+(3-2a)t\}
       e^{(3-2a)t}t^{a-1}\sqrt{1-t}\,dt=0.
\end{equation}
If $a$ is algebraic irrational, $t^{a-1}$ is not an algebraic
function: continuation about zero multiplies it by
$e^{2\pi i(a-1)}$, which is not a root of unity.  Thus this is
not an exponential integral with an algebraic integrand to which
an algebraic-period assertion could automatically be applied.
The same Kummer non-$E$ obstruction in
Proposition~\ref{p02:be18:notE} applies.  Differential transcendence
of the incomplete gamma function as a function of its order
\cite{p02:be18:Arreche} is not a theorem about its individual values.
No irrational-order special-value theorem closing
\eqref{p02:be18:Athreehalves} has been established in this argument.
The already proved irrationality of $A(3/2)$ remains the strongest
individual arithmetic conclusion here.

\subsection{A different boundary coordinate that is transcendental}

Keep the first excess-power McKay transform
\begin{equation}\label{p02:be18:McKay}
 L(s)=\frac{b+s}{b}
      \left(\frac{b^2-1}{(b+s)^2-1}\right)^\kappa,
 \qquad b>1,\quad \kappa>1.
\end{equation}
For a gamma factor of shape one and rate $\rho=b-c$, with
$0\leq c\leq1$, the candidate remainder transform, after a
positive normalization and a shift of the Laplace variable, is
\[
 F_{c,1}(w)=\frac{w(w+c)}{((w+c)^2-1)^\kappa}.
\]
Define
\begin{equation}\label{p02:be18:ell}
 \nu=\kappa-\frac32,\qquad t=2\kappa-2,\qquad
 \ell_\kappa(x)=\frac{t}{x}+
               \frac{I_{\nu+1}(x)}{I_\nu(x)},\quad x>0.
\end{equation}
The parameter $t$ in this subsection is a positive constant,
not the integration variable of \eqref{p02:be18:Athreehalves}.

We recall the classical arithmetic fact needed for this coordinate.
For rational $\nu>-1$ and nonzero algebraic $x$, the ratio
$I_{\nu+1}(x)/I_\nu(x)$, when interpreted through the normalized
entire functions, is transcendental.  Here is a verification tailored
to our use.  Put
\[
 H_\nu(z)=\sum_{n\geq0}\frac{(z^2/4)^n}{n!(\nu+1)_n}.
\]
Its even exponential Taylor coefficients are
$(1/2)_n/(\nu+1)_n$, and its odd coefficients vanish.  Rational
Pochhammer quotients have exponential coefficient and common-denominator bounds: at primes not dividing the fixed parameter
denominators, the two arithmetic progressions have divisibility
counts differing by at most one at each prime power, giving a
factor dividing $\operatorname{lcm}(1,\ldots,Cn)$; the finitely
many remaining primes contribute $\exp(O(n))$.  Thus $H_\nu$
and $H_\nu'$ are $E$-functions.  They satisfy
\[
 H_\nu''+\frac{2\nu+1}{z}H_\nu'-H_\nu=0.
\]
Their logarithmic derivative cannot be rational.  Otherwise an
entire function of exponential type with finitely many zeros has
form $e^{\lambda z}P(z)$; evenness forces $\lambda=0$, whereas
$H_\nu$ has infinitely many nonzero Taylor coefficients and is
not a polynomial.  Beukers' ordinary-point linear specialization
theorem \cite[Corollary~1.4]{p02:be18:Beukers} now applies to
$(H_\nu,H_\nu')$ at every nonzero algebraic point.  Since
$H_\nu'/H_\nu=I_{\nu+1}/I_\nu$, it proves the asserted
ratio transcendence.  This is a classical input, already used in
the preceding arithmetic paper.

\begin{theorem}[A transcendental intermediate extraction rate]
\label{p02:be18:rate}
Suppose $\kappa>1$ and $\ell_\kappa$ takes a value below one.
Then it has a unique global minimum $c_\kappa\in(0,1)$ at a
point $x_\kappa>0$.  Both satisfy the exact contact relation
\begin{equation}\label{p02:be18:ratecurve}
 (1-c_\kappa^2)x_\kappa^2+t c_\kappa x_\kappa-t=0.
\end{equation}
In the model \eqref{p02:be18:McKay}, extraction of a shape-one gamma
factor at rate $b-c$, $0\leq c\leq1$, is possible exactly when
$c\leq c_\kappa$.  At the boundary its remainder density has
exactly one positive zero, of order two, and the unrestricted
maximal shape is exactly $H_{\rm P}(c_\kappa)=1$.

If $\kappa$ is rational, then $c_\kappa$ and $x_\kappa$ are
both transcendental.  The hypothesis that $\ell_\kappa$ takes
a value below one holds for every real $1<\kappa\leq7/6$.
Consequently this supplies an explicit interval of rational input
parameters producing transcendental critical rates.  When $b>1$
is algebraic, the physical threshold $b-c_\kappa$ is transcendental.
\end{theorem}
\begin{proof}
Let
\[
 f(x)=x^{\kappa-1/2}I_{\kappa-3/2}(x).
\]
Positive constants do not matter in this proof.  The Bessel series
gives $f(x)=C x^t H_\nu(x)$ with $C>0$, so $f>0$ and
$f'>0$, and $f'/f=\ell_\kappa$.  At the two ends,
\[
 \ell_\kappa(x)\longrightarrow\infty\quad(x\downarrow0),
 \qquad \ell_\kappa(x)=1+\frac{\kappa-1}{x}+O(x^{-2})
                      \quad(x\to\infty).
\]
The latter is the standard Bessel expansion
\cite[10.40.1]{p02:be18:DLMF}.  Therefore a value below one implies
an attained positive global minimum below one.

The Bessel equation transforms into
\[
 f''-\frac{t}{x}f'+\left(\frac{t}{x^2}-1\right)f=0,
 \qquad
 \ell'=1+\frac{t\ell}{x}-\frac{t}{x^2}-\ell^2.
\]
At a minimizing point this gives \eqref{p02:be18:ratecurve}.
For a fixed $0<c<1$, the polynomial
$(1-c^2)x^2+tcx-t$ is strictly increasing on the positive axis,
starts negative, and tends to infinity.  It has exactly one
positive root.  All global minimizers have the same value
$c_\kappa$, so there is exactly one.  Differentiating the Riccati
identity at this stationary point gives
\begin{equation}\label{p02:be18:ratecurvature}
 \ell''(x_\kappa)=
       \frac{t(2-c_\kappa x_\kappa)}{x_\kappa^3}>0,
\end{equation}
since \eqref{p02:be18:ratecurve} implies $c_\kappa x_\kappa<1$.

The inverse of $(w+c)/((w+c)^2-1)^\kappa$ is a positive
constant times $e^{-cx}f(x)$.  Since $t>0$, this function tends
to zero at zero.  Multiplication of its transform by $w$ thus
corresponds to differentiation, with no atom or boundary term:
\[
 g_{c,1}(x)=C e^{-cx}f(x)(\ell_\kappa(x)-c),\qquad C>0.
\]
After the prescribed positive exponential tilt it is an integrable
signed density with total mass one.  Uniqueness of Laplace
transforms makes a probability remainder possible exactly when
this density is nonnegative, that is, when $c\leq c_\kappa$.
At equality uniqueness of the minimizer and
\eqref{p02:be18:ratecurvature} give exactly one double zero.

If some larger shape $h>1$ were extractable at this rate, the
shape-one remainder would be the convolution of its probability
remainder $\mu$ with a gamma law of shape $h-1$.  Put
$M=2\kappa-1$.  The quotient Laplace transform of $\mu$ is
asymptotic to $C s^{-(M-h)}$ at positive infinity, for $C>0$.
Necessarily $h\leq M$, since a probability Laplace transform
is bounded.  For $h<M$ this polynomial decay, and for $h=M$ its
positive limit, show that $\mu$ has positive mass arbitrarily
near zero: support bounded below by $\epsilon>0$ would instead
force exponential decay.  Convolving with a positive gamma
density therefore gives a continuous strictly positive lower
bound near every positive point, by restricting $\mu$ to
$[0,x/2]$.  This contradicts the continuous density's zero at
$x_\kappa$.  No shape larger than one is feasible, whereas
shape one is feasible, so $H_{\rm P}(c_\kappa)=1$ without
requiring a separate attainment theorem.

For the arithmetic assertion, first suppose $c_\kappa$ were
algebraic and $\kappa$ rational.  Equation
\eqref{p02:be18:ratecurve}, whose leading coefficient is nonzero,
would force $x_\kappa$ to be algebraic.  Then
\[
 \frac{I_{\nu+1}(x_\kappa)}{I_\nu(x_\kappa)}
       =c_\kappa-\frac{t}{x_\kappa}
\]
would be algebraic, contradicting the classical ratio result.
Conversely, an algebraic $x_\kappa$ makes
\eqref{p02:be18:ratecurve} a nonzero quadratic equation over the
algebraic numbers for $c_\kappa$, giving the same contradiction.
Thus both numbers are transcendental.  Subtracting $c_\kappa$
from an algebraic $b$ preserves transcendence.

It remains to verify the promised explicit range.  Coefficientwise
comparison of $H_{\nu+1}$ and $H_\nu$ gives, for $x>0$,
\[
 0<\frac{I_{\nu+1}(x)}{I_\nu(x)}
 =\frac{x}{2(\nu+1)}\frac{H_{\nu+1}(x)}{H_\nu(x)}
 <\frac{x}{2\kappa-1}.
\]
Taking $x=\sqrt{t(2\kappa-1)}$ therefore yields
\[
 \ell_\kappa(x)<2\sqrt{\frac{2\kappa-2}{2\kappa-1}}\leq1
       \qquad(1<\kappa\leq7/6).
\]
The first inequality is strict also at $\kappa=7/6$, so the
hypothesis holds throughout the stated interval.
\end{proof}

For example, $\kappa=9/8$ gives
$x_\kappa\simeq0.647079685$ and
$c_\kappa\simeq0.856687361$.  These approximations merely
identify the constants.  Their transcendence follows from the
exact argument above.  This theorem changes the boundary coordinate:
its unknown is the rate at fixed shape one, rather than the endpoint
shape $\kappa-A(\kappa)$.  It does not prove transcendence of
$A(1)$, $A(3/2)$, or an arbitrary intermediate-rate shape.

\subsection{The endpoint boundary table cannot obey an algebraic-coefficient recurrence}

The older analytic manuscript established the precise asymptotic
\begin{equation}\label{p02:be18:Aasymptotic}
 A(d)\sim\frac{1}{2\sqrt{2\pi d}}\left(\frac2e\right)^d
       \qquad(d\to+\infty).
\end{equation}
We use this as an inherited analytic input, not as a new estimate.
It has a further exact arithmetic consequence that does not require
knowing whether any individual $A(d)$ is transcendental.

\begin{theorem}[No polynomial recurrence on a rational mesh]
\label{p02:be18:nonholonomic}
Let $u\in\mathbb Q$, $v\in\mathbb Q_{>0}$, and take all
sufficiently large integers $n$ so that $u+vn>1/2$.  There are
no polynomials $p_0,\ldots,p_r\in\overline{\mathbb Q}[X]$,
not all zero, for which
\begin{equation}\label{p02:be18:recurrence}
       \sum_{j=0}^r p_j(n)A(u+v(n+j))=0
\end{equation}
holds for every sufficiently large $n$.
Consequently the ordinary generating series of any such tail is
not $D$-finite over $\overline{\mathbb Q}(z)$.  In particular,
$\sum_{n\geq1}A(n)z^n$ has transcendental radius of convergence
$e/2$ and satisfies no nonzero linear differential equation with
algebraic polynomial coefficients.  The function $A(d)$ itself
is not algebraic over $\mathbb C(d)$.
\end{theorem}
\begin{proof}
Set $q=(2/e)^v$.  This number is transcendental.  Indeed if
$v=r/s>0$ in lowest terms and $q$ were algebraic, then
$q^s=2^r e^{-r}$ would be algebraic, forcing $e^r$ and then
$e$ algebraic, contrary to Hermite's theorem.  From
\eqref{p02:be18:Aasymptotic}, for each fixed $j$,
\[
 \frac{A(u+v(n+j))}{A(u+vn)}\longrightarrow q^j.
\]
Suppose \eqref{p02:be18:recurrence} held.  Let $D$ be the maximum
of the degrees of its nonzero coefficient polynomials, and let
$c_j$ be the coefficient of $X^D$ in $p_j$.  Divide by
$n^D A(u+vn)$ and take the limit.  The result is
$\sum_j c_jq^j=0$, a nonzero polynomial equation with algebraic
coefficients, contradiction.

For completeness, equating Taylor coefficients in a polynomial-coefficient differential equation gives a polynomial recurrence
for the series coefficients: derivatives contribute falling
factorials and multiplication by a monomial shifts the index.
A nonzero differential operator gives a nonzero such recurrence
for all sufficiently large indices.  Thus the excluded recurrence
also excludes $D$-finiteness.  The radius follows directly from
$A(n)^{1/n}\to2/e$.

Finally suppose $A$ satisfied a nonzero polynomial relation over
$\mathbb C(d)$.  Clearing denominators gives
$\sum_{j=0}^m P_j(d)A(d)^j=0$, with $P_j\in\mathbb C[d]$.
Let $j_0$ be the smallest index for which $P_{j_0}$ is nonzero.
After division by $A(d)^{j_0}$, every remaining term with
$j>j_0$ tends to zero exponentially by
\eqref{p02:be18:Aasymptotic}, even after multiplication by its
polynomial.  This would force the nonzero polynomial $P_{j_0}(d)$
to tend to zero on the positive real axis, an impossibility.
\end{proof}

This is an obstruction to exact linear differential or recurrence
compression of the boundary family over algebraic coefficients.
It is not a proof that an individual entry is transcendental.
The conclusion is stated over $\overline{\mathbb Q}$: the proof
does not exclude polynomial recurrences with arbitrary complex
coefficients.  Nor does it establish differential transcendence of
$A(d)$, an irrationality measure for its individual values, or the
arithmetic of every noninteger-shape intermediate boundary.

\paragraph{Result accounting.}
There are two proved target families: the transcendental critical
rate for an exponential factor, and the recurrence obstruction for
the endpoint boundary table.  The separate investigations of
$A(1)$ and $A(3/2)$ remain unsuccessful at proving transcendence.
Their failure is supported by an exact non-$E$ theorem and explicit
integral reductions.  Examples at different rational $\kappa$,
rational meshes, or algebraic base rates are not counted separately.


% Integrated source: arithmetic_results.tex
\section{Positive arithmetic kernels and the two missing boundary values}
\label{p02:ac18:section}

This section makes three further attempts at individual-constant problems:
Euler's constant $\gamma$, the Euler--Gompertz constant $\delta$, and Catalan's
constant $G$; the last argument also treats the same polynomial approach to
$\zeta(5)$ and all integer zeta values.  The new results are exact
obstructions for larger kernel classes and an explicit limitation of
arithmetic inference from finite differences.  None identifies $\gamma$,
$\delta$, $G$, or $\zeta(5)$ as irrational.

We use
\[
 \gamma=-\int_0^\infty e^{-t}\log t\,dt,\qquad
 \delta_x=\int_0^\infty\frac{e^{-t}}{t+x}\,dt,\qquad \delta=\delta_1.
\]
The preceding manuscript \cite{p02:root:previous} solved the square-kernel problem
for $\delta$ and a beta-kernel denominator calculation for $G$.
Here the former is enlarged to \emph{every} integer polynomial nonnegative
on the half-line, and the latter to \emph{every} integer polynomial
nonnegative on the unit interval.  These are restricted approximation
problems, not assertions about all integer linear forms.

\subsection{The full positive polynomial cone for the Gompertz constant}

\begin{lemma}[Two weighted leading-coefficient bounds]
\label{p02:ac18:Laguerre}
For $\alpha\in\{0,1\}$ and a real polynomial $A$ of degree $n$ with
leading coefficient $a$, put
\[
 J_\alpha(A)=\int_0^\infty\frac{t^\alpha e^{-t}A(t)^2}{1+t}\,dt.
\]
Then
\begin{equation}\label{p02:ac18:Laguerre-bound}
 J_\alpha(A)\ge
 \frac{a^2 n!(n+\alpha)!}{2n+\alpha+2}.
\end{equation}
\end{lemma}
\begin{proof}
Let $\pi_n^{(\alpha)}$ be the monic Laguerre polynomial for
$t^\alpha e^{-t}dt$.  Rodrigues' formula
\[
 t^\alpha e^{-t}\pi_n^{(\alpha)}(t)
 =(-1)^n\frac{d^n}{dt^n}(e^{-t}t^{n+\alpha})
\]
shows, by integration by parts, that it is orthogonal to every polynomial
of degree less than $n$, and that its squared norm is
$n!(n+\alpha)!$.  The coefficient of $t^{n-1}$ is $-n(n+\alpha)$.
Comparison with $\pi_{n+1}^{(\alpha)}$ and orthogonality therefore give
\[
 \int_0^\infty t^{\alpha+1}e^{-t}
       \bigl(\pi_n^{(\alpha)}(t)\bigr)^2dt
 =(2n+\alpha+1)n!(n+\alpha)!.
\]
The inner product of $A$ with $\pi_n^{(\alpha)}$ is
$a n!(n+\alpha)!$.  Cauchy--Schwarz, with factors
$A/\sqrt{1+t}$ and $\sqrt{1+t}\,\pi_n^{(\alpha)}$, proves
\eqref{p02:ac18:Laguerre-bound}.
\end{proof}

\begin{theorem}[Exact minimum over the full half-line cone]
\label{p02:ac18:full-cone}
Let $\mathcal P$ be the set of nonzero $P\in\mathbb Z[t]$ such that
$P(t)\ge0$ for $t\ge0$.  Then
\begin{equation}\label{p02:ac18:cone-min}
 \min_{P\in\mathcal P}\int_0^\infty\frac{e^{-t}P(t)}{1+t}\,dt
 =4\delta-2>0.
\end{equation}
The unique minimizing polynomial is $P(t)=(t-1)^2$.
Every integral in \eqref{p02:ac18:cone-min} belongs to
$\mathbb Z\delta+\mathbb Z$.
\end{theorem}
\begin{proof}
We first record rational bounds with certificates.  Direct polynomial
division and the moments $\int_0^\infty t^k e^{-t}dt=k!$ give
\[
 \int_0^\infty\frac{e^{-t}(t-1)^2}{1+t}dt=4\delta-2>0.
\]
With $Q(t)=t^4-20t^3+120t^2-240t+120$, the same calculation gives
\[
 \int_0^\infty\frac{e^{-t}tQ(t)^2}{1+t}dt
 =150300-251001\delta>0.
\]
These strictly positive integrals imply
\begin{equation}\label{p02:ac18:delta-bounds}
 \frac12<\delta<\frac{100}{167}<\frac35,
 \qquad 0<4\delta-2<\frac25.
\end{equation}

Every real polynomial nonnegative on $[0,\infty)$ has a representation
\begin{equation}\label{p02:ac18:half-line-squares}
                  P(t)=A(t)^2+tB(t)^2
\end{equation}
with real polynomials $A,B$.  Here is an elementary factorization proof.
Positive real roots have even multiplicity.  A negative real root gives
a factor $t+c=(\sqrt c)^2+t$, $c>0$, and the factor $t$ has the same
form with zero first square.  An irreducible real quadratic can be written
\[
 t^2-2at+a^2+b^2=(t-r)^2+t\,2(r-a),\qquad
 r=\sqrt{a^2+b^2},\quad b\ne0.
\]
All these factors have the required form, as does a positive constant.
Finally the identity
\[
 (A^2+tB^2)(C^2+tD^2)=(AC-tBD)^2+t(AD+BC)^2
\]
proves \eqref{p02:ac18:half-line-squares}.  The positive highest terms of the
two summands cannot cancel.  Thus if $\deg P=2n$ with leading
coefficient $p$, then $\deg A=n$ and its leading coefficient squared
is $p$; if $\deg P=2n+1$, the corresponding statement holds for $B$.
The integer $p$ is at least one.

For even degree $2n\ge4$, Lemma~\ref{p02:ac18:Laguerre} gives a lower bound
$(n!)^2/(2n+2)\ge2/3$.  For odd degree $2n+1\ge3$, it gives
$n!(n+1)!/(2n+3)\ge2/5$.  Both sequences increase in their stated
ranges.  By \eqref{p02:ac18:delta-bounds} these degrees cannot minimize.
A constant polynomial has integral at least $\delta$.  A polynomial
$at+c$ nonnegative on the half-line has nonnegative integer coefficients,
so its nonzero integral is at least $\min(\delta,1-\delta)>2/5$.

For degree two with leading coefficient $p\ge2$, the even-degree bound
is $p/4\ge1/2$.  It remains to treat $P=t^2+ut+v$ with $u,v\in\mathbb Z$.
If $u\ge0$, then $v\ge0$ and the integral is at least $\delta$.
If $u=-k<0$, nonnegativity implies $v\ge\lceil k^2/4\rceil$, and the
integral is
\[
                       -k+(1+k+v)\delta.
\]
For $k\ge4$, this exceeds $(k-2)^2/8\ge1/2$, using $\delta>1/2$.
For $k=1$ the least value is $3\delta-1>4\delta-2$; for $k=3$ it is
$7\delta-3>4\delta-2$.  For $k=2$, the least value is $4\delta-2$
and equality requires $v=1$.  This proves the minimum and uniqueness.
Finally, division $P(t)=P(-1)+(t+1)R(t)$ with $R\in\mathbb Z[t]$
proves the integer linear-form assertion.
\end{proof}

\paragraph{What this closes, and what it leaves.}
The theorem solves a self-posed lattice optimization over all degrees;
no priority assertion is made.  It strictly includes the preceding
square class because the factors in \eqref{p02:ac18:half-line-squares} need
not have integer coefficients.  Rational normalization cannot help
within this construction: if $P=Q/d$, $Q\in\mathcal P$, clearing that
specified denominator gives the integral of $Q$, bounded below as above.
Signed kernels remain a different problem.  Indeed
$P(t)=q-p(1+t)$ realizes the arbitrary integer form $q\delta-p$,
so dropping positivity would turn the proposed lower bound into an
unproved arithmetic assertion about $\delta$.  The multiple-orthogonal
approximants of Wolfs and Van Assche \cite{p02:ac18:VW} are not excluded by
Theorem~\ref{p02:ac18:full-cone}.

\subsection{A full compact-interval obstruction for Catalan and zeta kernels}

\begin{lemma}[A universal positive integer-polynomial bound]
\label{p02:ac18:compact}
Suppose $w(t)\ge c(1-t)^r$ on $(0,1)$, with $c>0$ and $r\ge0$.
For every nonzero $P\in\mathbb Z[t]$ nonnegative on $[0,1]$, of degree
$d\ge1$,
\begin{equation}\label{p02:ac18:compact-bound}
 \int_0^1 w(t)P(t)dt
 \ge\frac{c}{r+1}\,4^{-d}(8d^3)^{-(r+1)}.
\end{equation}
An infinite integral is allowed on the left.
\end{lemma}
\begin{proof}
Write $M=\max_{[0,1]}P>0$.  The monic Chebyshev extremal inequality
implies $M\ge2^{1-2d}$.  For completeness, the shifted polynomial
$2^{1-2d}T_d(2t-1)$ is monic and alternates between
$\pm2^{1-2d}$ at $d+1$ ordered points.  A monic polynomial of strictly
smaller supremum would differ from it by a degree-at-most-$d-1$
polynomial with at least $d$ sign changes, which is impossible.
Divide the leading coefficient of $P$, whose absolute value is at least
one, to apply this argument.

Expand $P(t)=\sum_{j=0}^d c_jT_j(2t-1)$.  The cosine integral formulas
give $|c_j|\le2M$ for $j\ge1$.  Since
$|T_j'(u)|\le j^2$ on $[-1,1]$, we have
\[
 \|P'\|_\infty\le4M\sum_{j=1}^d j^2\le4Md^3.
\]
The derivative bound for $T_j$ follows from
$T_j(\cos\theta)=\cos(j\theta)$ and
$|\sin(j\theta)|\le j|\sin\theta|$, including endpoints by continuity.
At a maximizing point, therefore, $P\ge M/2$ throughout an interval
in $[0,1]$ of length at least $h=(8d^3)^{-1}$.  Among intervals of
length $h$, the integral of $(1-t)^r$ is smallest at the right endpoint
and equals $h^{r+1}/(r+1)$.  Thus the integral is at least
$cMh^{r+1}/(2(r+1))$.  The lower bound on $M$ proves the lemma.
\end{proof}

\begin{theorem}[Standard denominator clearance fails on both full cones]
\label{p02:ac18:compact-arithmetic}
Let $P\in\mathbb Z[t]\setminus\{0\}$ be nonnegative on $[0,1]$, of
degree $d\ge1$.  Put $D_d=\operatorname{lcm}(1,\ldots,d)$ and
$O_d=\operatorname{lcm}(1,3,\ldots,2d-1)$.  Define
\[
 C(P)=\int_0^1\frac{-\log x}{1+x^2}P(x^2)dx,
 \qquad
 Z_s(P)=\frac1{(s-1)!}\int_0^1
       \frac{(-\log x)^{s-1}P(x)}{1-x}dx\quad(s\ge2).
\]
Then
\begin{align}
 C(P)&\ge\frac{4^{-d}}{1024d^6},\label{p02:ac18:Cbound}\\
 Z_s(P)&\ge
 \frac{4^{-d}}{(s-1)!(s-1)(8d^3)^{s-1}}.
 \label{p02:ac18:Zbound}
\end{align}
Moreover $O_d^2C(P)\in\mathbb ZG+\mathbb Z$ and
$D_d^s Z_s(P)\in\mathbb Z\zeta(s)+\mathbb Z$.  Uniformly over these
polynomials as $d\to\infty$,
\begin{align}
 \log\bigl(O_d^2C(P)\bigr)&\ge(4-\log4)d+o(d),\label{p02:ac18:Cgrowth}\\
 \log\bigl(D_d^sZ_s(P)\bigr)&\ge(s-\log4)d+o(d).
 \label{p02:ac18:Zgrowth}
\end{align}
Thus no sequence in either cone, using these specified clearances, gives
nonzero integer forms tending to zero.
\end{theorem}
\begin{proof}
After $t=x^2$, the Catalan weight is
$(-\log t)/(4\sqrt t(1+t))\ge(1-t)/8$.
Use Lemma~\ref{p02:ac18:compact} with $c=1/8,r=1$.
For zeta, $-\log t\ge1-t$ gives the weight lower bound
$(1-t)^{s-2}/(s-1)!$, so use $c=1/(s-1)!,r=s-2$.

Division of $P(t)$ by $1+t$ gives
\[
 C(P)=P(-1)G+\sum_{k=0}^{d-1}\frac{q_k}{(2k+1)^2},
 \qquad q_k\in\mathbb Z.
\]
For $P(t)=\sum_{k=0}^d a_kt^k$, geometric-series integration gives
\[
 Z_s(P)=P(1)\zeta(s)-\sum_{k=0}^d a_kH_k^{(s)},
 \qquad H_k^{(s)}=\sum_{j=1}^k j^{-s}.
\]
These prove the integrality claims.  The prime number theorem gives
$\log D_d=d+o(d)$ and $\log O_d=2d+o(d)$; exclusion of the power of
$2$ changes the latter logarithm by only $O(\log d)$.
Combining these with \eqref{p02:ac18:Cbound}--\eqref{p02:ac18:Zbound} proves the
logarithmic lower bounds, whose leading constants are positive for
every integer $s\ge2$.
If degrees remain bounded, Lemma~\ref{p02:ac18:compact} gives a positive
lower bound depending only on that bound on the degree.  Constant
polynomials also have a positive lower bound.  This proves the final
assertion for arbitrary sequences of degrees.
\end{proof}

This includes $G$, $\zeta(5)$, and $\zeta(7)$ as one obstruction family,
not as separately solved conjectures.  It improves the beta-only
calculation by allowing arbitrary positive integer-polynomial shapes.
It concerns the stated least-common-multiple multipliers, not the
minimal reduced denominators: a proven arithmetic cancellation could
still change the conclusion for a different clearance.  Also the target
coefficients $P(-1)$ and $P(1)$ must remain nonzero.  The theorem is
compatible with the much richer constructions of Brown--Zudilin
\cite{p01:ac17:BrownZudilin}.

\subsection{Euler's constant: an indefinite lattice criterion and its circular input}

The positive cone behaves differently for the logarithmic weight because
$\log t$ changes sign.  This gives a precise connection to a powerful
homogeneous-dynamics theorem, but also exposes the missing hypothesis.

\begin{proposition}[A ternary criterion for $\gamma$]
\label{p02:ac18:gamma-Oppenheim}
For $v=(a,b,c)\in\mathbb Z^3$, put
\[
 Q_\gamma(v)=\int_0^\infty e^{-t}\log t\,(a+bt+ct^2)^2dt.
\]
Its symmetric matrix is
\begin{equation}\label{p02:ac18:gamma-matrix}
 M_\gamma=\begin{pmatrix}
 -\gamma&1-\gamma&3-2\gamma\\
 1-\gamma&3-2\gamma&11-6\gamma\\
 3-2\gamma&11-6\gamma&50-24\gamma
 \end{pmatrix}.
\end{equation}
This form is nondegenerate of signature $(2,1)$.
The following are equivalent:
\begin{enumerate}
 \item $\gamma$ is irrational;
 \item $Q_\gamma$ is not a real scalar multiple of a rational form;
 \item $Q_\gamma(\mathbb Z^3)$ is dense in $\mathbb R$;
 \item for every $\varepsilon>0$, there is $v\in\mathbb Z^3$ with
       $0<|Q_\gamma(v)|<\varepsilon$.
\end{enumerate}
\end{proposition}
\begin{proof}
The moments are $\int t^ke^{-t}\log t\,dt=k!(H_k-\gamma)$, which
proves \eqref{p02:ac18:gamma-matrix}.  Strict Jensen's inequality for an
exponential random variable gives $\gamma>0$; restricting its negative
logarithmic integral to $(0,1)$ gives
$\gamma<\int_0^1(-\log t)dt=1$.
The determinant is
\[
 \det M_\gamma=-11+3\gamma+10\gamma^2-4\gamma^3<-2,
\]
since the polynomial increases from $-11$ to $-2$ on $[0,1]$.
The first diagonal entry is negative and the last positive.  The
negative determinant therefore forces exactly one negative eigenvalue.

If $\gamma$ is rational, the matrix is rational.  Conversely, if the
matrix is a scalar multiple of a rational matrix, the ratio of its
$(1,2)$ and $(1,1)$ entries is rational.  That ratio is $1-1/\gamma$,
so $\gamma$ is rational.  This proves the equivalence of the first two
conditions.  Margulis' theorem on the Oppenheim conjecture
\cite{p02:ac18:Margulis} applies to a nondegenerate indefinite irrational
ternary form and gives density, proving the implication from the second
to the third.  Density implies the fourth.
Finally, expansion gives
\[
 Q_\gamma(v)=A(v)-B(v)\gamma,\qquad A(v),B(v)\in\mathbb Z,
 \quad B(v)=\int_0^\infty e^{-t}(a+bt+ct^2)^2dt>0
\]
for $v\ne0$.  If $\gamma=p/q$ were rational, every nonzero form would
have absolute value at least $1/q$, contradicting the fourth condition.
\end{proof}

This is a self-posed equivalent criterion, not a proof of any of its
conditions.  Applying Oppenheim by asserting that the displayed matrix
is irrational would assume exactly the desired irrationality of
$\gamma$.  The point of the calculation is that polynomial degree two
already gives a genuine indefinite lattice problem; positivity of the
polynomial square does not imply positivity of this logarithmic
integral.  A new effective search theorem for this specific form would
still have to establish its required Diophantine hypothesis independently.

\subsection{Curvature recovery and the two arithmetic boundary values}

Let $\mathcal F_{\exp}=\overline{\mathbb Q}(e^a:a\in\overline{\mathbb Q})$.
The preceding manuscript proves that the values $\operatorname{Ein}(x)$
at distinct nonzero algebraic $x$ are algebraically independent over
$\mathcal F_{\exp}$ \cite{p02:root:previous}.  This uses classical
Kolchin--Ostrowski and Siegel--Shidlovskii theory; we do not count its
following linear consequences as a new independence theorem.
For an integer $b\ge2$, put
\[
 F_j=\operatorname{Ein}(b^j),\quad
 Y_j=e^{-b^j}\delta_{b^j},\quad
 C_j=Y_j-2Y_{j+1}+Y_{j+2},\qquad j\ge0.
\]
The classical split identity gives
$Y_j=F_j-\gamma-j\log b$.

\begin{proposition}[Exact recovery from independent curvatures]
\label{p02:ac18:curvature}
The sequence $(C_j)_{j\ge0}$ is algebraically independent over
$\mathcal F_{\exp}$, every $C_j$ is positive, and
\begin{equation}\label{p02:ac18:curvature-recovery}
 Y_j=\sum_{r=j}^\infty(r-j+1)C_r.
\end{equation}
For $N\ge j$, the error after $r=N$ lies strictly between zero and
\[
                 (N-j+2)\frac{e^{-b^{N+1}}}{b^{N+1}}.
\]
Moreover
\[
 \mathcal F_{\exp}(Y_j:j\ge0)
 =\mathcal F_{\exp}(C_j:j\ge0)(Y_0,Y_1).
\]
\end{proposition}
\begin{proof}
The common affine terms cancel, so $C_j=F_j-2F_{j+1}+F_{j+2}$.
Any finite collection of these linear forms has independent coefficient
rows: in a nonzero row combination, the row with largest index has a
nonzero coefficient in its last column.  Independent linear forms in
algebraically independent variables are algebraically independent,
proving the first assertion.

Write $Y(x)=\int_1^\infty e^{-xt}dt/t$.  Then
$Y(bx)\le e^{-(b-1)x}Y(x)$, so $Y_{j+1}<Y_j/2$ for $b\ge2$ and $j\ge0$.
Thus $C_j>0$.  Finite summation gives
\[
 \sum_{r=j}^N(r-j+1)C_r
 =Y_j-(N-j+2)Y_{N+1}+(N-j+1)Y_{N+2}.
\]
Since $0<Y_j\le e^{-b^j}/b^j$, passage to the limit proves the recovery
formula and the stated positive error bound.  The recurrence determines
every $Y_j$ from $Y_0,Y_1$ and finitely many $C_r$; conversely each $C_r$
is a linear combination of three $Y$'s.  This proves the field identity.
\end{proof}

The recovery is analytically complete.  Its arithmetic limitation is not
removed by the positivity or the very rapid error bound.  The following
countermodel makes that limitation exact.

\begin{theorem}[Independent curvatures can have two prescribed rational boundaries]
\label{p02:ac18:countermodel}
Let $K\subset\mathbb R$ be a countable field, and let rational numbers
$q_0>2q_1>0$ be given.  Fix $b>1$.  There exists a positive sequence
$(c_j)_{j\ge0}$, algebraically independent over $K$, such that
$c_j=O(e^{-b^j})$ as $j\to\infty$ and
\[
 y_j=\sum_{r=j}^\infty(r-j+1)c_r
\]
satisfies $y_0=q_0$, $y_1=q_1$, and
$y_j-2y_{j+1}+y_{j+2}=c_j$.
\end{theorem}
\begin{proof}
Choose positive numbers $\varepsilon_j\le e^{-b^j}$ for $j\ge2$ such
that $\sum_{j\ge2}2j\varepsilon_j<q_1/2$.
Let the tail coordinates vary in the compact metrizable product
$X=\prod_{j\ge2}[\varepsilon_j,2\varepsilon_j]$.
For every point of $X$, define
\[
 c_1=q_1-\sum_{r\ge2}r c_r,\qquad
 c_0=q_0-2q_1+\sum_{r\ge2}(r-1)c_r.
\]
These are continuous functions on $X$, are positive, and give the two
prescribed sums.  All series converge uniformly on $X$.

For any nonzero polynomial over $K$ in a finite selection of the
coordinates $c_j$, its zero set in $X$ is closed and has empty interior.
To see the latter assertion, take a basic open cylinder.  Besides the
finitely many selected tail coordinates, vary two further coordinates
$c_k,c_l$, with distinct $k,l\ge2$, outside all indices involved in the
polynomial or cylinder.  Their contributions to $(c_0,c_1)$ have matrix
\[
 \begin{pmatrix}k-1&l-1\\-k&-l\end{pmatrix},
 \qquad \det=l-k\ne0.
\]
Together with the selected tail coordinates this gives a linear map of
full row rank onto those coordinates and $(c_0,c_1)$.  Varying them in
open subintervals of their boxes produces an open set in the finite
coordinate space.  A nonzero polynomial cannot vanish on that open set.
This remains valid if one or both of $c_0,c_1$ do not occur in the
polynomial, by projecting away unused coordinates.

There are only countably many such polynomials because $K$ is countable.
The Baire category theorem on the nonempty complete space $X$ supplies a
point outside all their zero sets.  Its full sequence is algebraically
independent over $K$.  The bounds on the boxes give the stated tail decay,
and absolute convergence permits termwise subtraction in the definition
of $y_j$, giving its second difference and the required boundary values.
\end{proof}

Apply the countermodel with
$K=\mathcal F_{\exp}\cap\mathbb R$, which is countable.  Since $\mathcal F_{\exp}$ is closed under complex
conjugation, independence of real numbers over this real subfield is
equivalent to independence over $\mathcal F_{\exp}$: take real and
imaginary parts of a putative polynomial relation.  The countermodel does not
supply a counterexample about the actual Gompertz integrals or the
specific values $F_j$.  It proves that algebraic independence of all
curvatures, positivity, rapid tail decay, and exact boundary reconstruction
alone do not classify even the first two boundary values: in the
countermodel both are rational.  Therefore the preceding defect-two
family statement cannot be turned into an individual classification by
those properties alone.  One must use additional arithmetic information
specific to $\gamma$, $\log b$, or the exponential-integral connection.
The mixed arithmetic-Gevrey problems of Fischler--Rivoal
\cite{p02:ac18:Mixed} locate related missing separation statements; they are
conjectural inputs, not hypotheses silently used here.


% Integrated source: information_results.tex
\section{Prime block complexity, polynomial carries, and modular neighbours}
\label{p02:ci18:section}

We examine three specific transcendence targets,
\[
 P_2=\sum_{p\ {
m prime}}2^{-p},\qquad
 L_{\mathbb P}=\sum_{p\ {
m prime}}\frac1{2^p-1},\qquad
 E=\sum_{n\ge1}\frac1{2^n-1}.
\]
None is proved transcendental here.  The first result is a quantitative
transfer from Maynard's prime-cluster theorem to the block complexity of
the prime indicator.  It excludes a tempting transcendence method for
$P_2$.  The second result makes the additional carry information needed
for $L_{\mathbb P}$ explicit.  The third identifies three algebraically
independent linear combinations of reciprocal-Mersenne moments using
Nesterenko's theorem.  This last transfer is a consequence of established
modular transcendence theory, not a new solution of the problem for $E$.

\subsection{Prime binary digits: superpolynomial complexity with zero entropy}

Write $u_j=1$ if $j$ is prime and $u_j=0$ otherwise, for $j\ge1$.
Let $p_u(n)$ be the number of distinct contiguous length-$n$ words in $u$.
There are infinitely many zero digits, so $u$ is the canonical binary
expansion of $P_2$.  Let $p_u^{\rm rec}(n)$ count only words that occur
infinitely often.

The analytic input is the following consequence of Maynard's
Theorem~3.4 in \cite{p02:ci18:Maynard2016}.  There is an absolute constant
$C_0$ such that for every positive integer $k$ there is an integer
$M=M(k)$, with $k\le M\le\exp(C_0k)$, having this property:
every fixed admissible family of $M$ distinct positive integer linear
forms has infinitely many arguments at which at least $k$ forms are
prime.  Here admissible means that their product has no fixed prime
divisor.  The stated consequence follows by taking a fixed interval
parameter in that theorem and then letting its size tend to infinity;
all coefficients in our application are fixed during that limit.

\begin{theorem}[Prime-word complexity transfer]
\label{p02:ci18:primecomplexity}
Fix $k\ge1$, choose $M=M(k)$ as above, and put
$W=\prod_{p\le M}p$ and $K=\lfloor n/W\rfloor$.
For $K\ge M$,
\begin{equation}\label{p02:ci18:countingbound}
 p_u^{\rm rec}(n)\ge\frac{\binom Kk}{\binom Mk}.
\end{equation}
Consequently there are absolute positive constants $c,C$ such that,
for all sufficiently large $n$,
\begin{equation}\label{p02:ci18:complexityrange}
 \exp\{c\log n\log\log n\}
 \le p_u^{\rm rec}(n)\le p_u(n)
 \le\exp\{Cn/\log\log n\}.
\end{equation}
In particular, the prime word has zero topological entropy, while its
block complexity grows faster than every fixed power of $n$.
\end{theorem}

\begin{proof}
Let $G=\{W,2W,\ldots,KW\}\subseteq\{1,\ldots,n\}$.
For every $M$-element subset $H\subseteq G$, assign distinct primes
$q_i>n$ to the positions $i\in\{1,\ldots,n\}\setminus H$.
The Chinese remainder theorem supplies $r$ modulo
\[
 A=W\prod_{i\notin H}q_i
\]
such that $r\equiv1\pmod W$ and $r\equiv-i\pmod{q_i}$.
Choose a positive representative $r$.  Consider the $M$ forms
$At+r+h$, $h\in H$.

For a prime dividing $W$, their constant terms are all $1$ modulo that
prime.  For $q_i\mid A/W$, the difference $h-i$ is nonzero and has
absolute value less than $q_i$, so no constant term vanishes modulo
$q_i$.  For a prime $p\nmid A$, necessarily $p>M$; the $M$ forms
exclude at most $M$ residues of $t$ modulo $p$.  The family is therefore
admissible.  Maynard's theorem supplies infinitely many $t$ with at
least $k$ prime values among these forms.  At every sufficiently large
such $t$, each position outside $H$ is composite because $At+r+i$ is
a multiple of $q_i$ exceeding $q_i$.

Among the finitely many resulting words, at least one recurs infinitely
often.  Its support $S$ satisfies $S\subseteq H$ and $|S|\ge k$.
Choose one such word for each $H$.  A fixed output support $S$, of size
$s\ge k$, can be chosen by at most
$\binom{K-s}{M-s}\le\binom{K-k}{M-k}$ different sets $H$.
Counting the $\binom KM$ inputs proves
\[
 p_u^{\rm rec}(n)\ge
 \frac{\binom KM}{\binom{K-k}{M-k}}
 =\frac{\binom Kk}{\binom Mk}.
\]

For the uniform lower bound, enlarge $C_0$ if necessary and take
$k=\lfloor(\log\log n)/(4C_0)\rfloor$.  Then
$M\le(\log n)^{1/4}$ and
\[
 \log W\le\log(M!)\le M\log M=o(\log n).
\]
Thus $K/M\ge n^{1/2}$ for large $n$.  Since
$(K-j)/(M-j)\ge K/M$ for $0\le j<k$, the right side of
\eqref{p02:ci18:countingbound} is at least $(K/M)^k$.  This gives the
lower bound in \eqref{p02:ci18:complexityrange}.  There is no interchange
of a prime-counting limit with $k$: for each fixed block length $n$,
all the finitely many forms have already been fixed before Maynard's
theorem is applied.

For the upper bound, let $y\ge3$ and $Q=\prod_{p\le y}p$.
A block starting after position $y$ can have ones only in residue
classes coprime to $Q$.  For each starting residue modulo $Q$, at most
$n\varphi(Q)/Q+Q$ positions are allowed.  There are at most $Q$
starting residues and at most $\lceil y\rceil$ exceptional early
blocks.  Therefore
\begin{equation}\label{p02:ci18:sieveupper}
 p_u(n)\le \lceil y\rceil+
 Q\,2^{n\varphi(Q)/Q+Q}.
\end{equation}
The finite Euler product satisfies
\[
 \frac{Q}{\varphi(Q)}
 =\prod_{p\le y}(1-p^{-1})^{-1}
 \ge\sum_{m\le y}\frac1m\ge\log y,
\]
because every integer $m\le y$ has all its prime factors at most $y$.
Choose $y=\lfloor\log n/(2\log\log n)\rfloor$.
For large $n$, $Q\le y!\le y^y\le n^{1/2}$ and
$\log y\asymp\log\log n$.  Substitution in
\eqref{p02:ci18:sieveupper} proves the upper bound.
\end{proof}

Adamczewski--Bugeaud's criterion states that an irrational algebraic
number must have $p(n)/n\to\infty$ in each integer base
\cite[Theorem~1]{p02:ci18:AB2007}.  Theorem~\ref{p02:ci18:primecomplexity}
shows that this necessary condition is satisfied by $P_2$ with a large
margin.  In particular, obtaining zero entropy or a subexponential
complexity upper bound cannot establish its transcendence by that
criterion.  This is a method exclusion, not evidence that $P_2$ is
algebraic.  The argument is a direct consequence of Maynard's theorem;
we make no global priority claim for this consequence.

\subsection{Prime Lambert series: finite certificates for polynomial carries}

Absolute convergence gives
\[
 L_{\mathbb P}=\sum_{n\ge1}\frac{\omega(n)}{2^n},
\]
where $\omega(n)$ counts distinct prime factors.  Its irrationality is
already a theorem of Tao and Ter\"av\"ainen
\cite[Theorem~1.3]{p02:ci18:TaoTeravainen}.  To seek transcendence one must
exclude every nonconstant polynomial relation, not only linear ones.
The following elementary certificate identifies exactly which
coefficient combinations a carry argument would have to control.

\begin{theorem}[Finite polynomial-carry certificate]
\label{p02:ci18:carrycertificate}
Let $F(z)=\sum_{n\ge1}a_nz^n$, where $a_n$ are integers with
$|a_n|\le n$, and set $\xi=F(1/2)$.
For a nonconstant polynomial
$A(X)=\sum_{j=0}^D c_jX^j\in\mathbb Z[X]$, let
$H=\max_j|c_j|$ and write
\[
 A(F(z))=\sum_{n\ge0}b_nz^n.
\]
For $N\ge0$ and $J\ge1$, define the rational number and explicit bound
\begin{align}
 S_{N,J}&=\sum_{m=1}^J\frac{b_{N+m}}{2^m},\notag\\
 B_{N,J}&=HD(2D-1)!\,2^{2D-1-J}
              (N+J+D+1)^{2D-1}.\label{p02:ci18:carrybound}
\end{align}
If $\operatorname{dist}(S_{N,J},\mathbb Z)>B_{N,J}$, then
$A(\xi)\ne0$.  Furthermore, $\xi$ is transcendental if and only if
every nonconstant $A\in\mathbb Z[X]$ has such a certificate for some
$N,J$.  If $\xi$ is transcendental, one may fix any $N$ in advance and
find a sufficiently large $J$ separately for each $A$.
\end{theorem}

\begin{proof}
Put $s=2D-1$.  Coefficientwise absolute values in $F(z)^j$ are bounded
by those in $(z/(1-z)^2)^j$.  For $n\ge1$ this gives
\[
 |b_n|\le H\sum_{j=1}^D
       \binom{n+j-1}{2j-1}\le HD(n+D)^s,
\]
where a binomial coefficient is zero when $n<j$.
All series below converge absolutely.  Set
$T_N=\sum_{m\ge1}b_{N+m}2^{-m}$.
Writing $m=J+1+t$ and using
$(t+1)^s\le s!\binom{t+s}{s}$ gives
\begin{align*}
 |T_N-S_{N,J}|
 &\le HD\,2^{-J-1}(N+J+D+1)^s
                  \sum_{t\ge0}(t+1)^s2^{-t}\\
 &\le HD\,2^{-J-1}(N+J+D+1)^s s!2^{s+1}
 =B_{N,J}.
\end{align*}
On the other hand,
\begin{equation}\label{p02:ci18:carryintegrality}
 2^N A(\xi)=\sum_{n=0}^N b_n2^{N-n}+T_N.
\end{equation}
The finite sum is an integer.  Thus $A(\xi)=0$ would force $T_N$
to be an integer for every $N$, which is incompatible with the stated
strict distance test.

If $\xi$ is algebraic, take a nonconstant integer polynomial vanishing
at $\xi$; no test for that polynomial can succeed.  If $\xi$ is
transcendental, then $A(\xi)$ is transcendental for every nonconstant
integer polynomial $A$: otherwise an algebraic equation for $A(\xi)$
would compose with $A$ to give an equation for $\xi$.
Consequently $T_N\notin\mathbb Z$ for every fixed $N$.
Since $S_{N,J}\to T_N$ and $B_{N,J}\to0$, the strict test eventually
holds.  This proves both directions.
\end{proof}

For $a_n=\omega(n)$, the coefficients in this test are explicitly
\begin{equation}\label{p02:ci18:omegaconvolution}
 b_n=\sum_{j=1}^D c_j
 \sum_{\substack{n_1+\cdots+n_j=n\\n_i\ge1}}
       \omega(n_1)\cdots\omega(n_j)\qquad(n\ge1).
\end{equation}
For fixed $D,H$, the remainder in \eqref{p02:ci18:carrybound} can be made
smaller than any fixed positive tolerance using $J=O_{D,H}(\log(N+2))$
(with an additional constant depending on the tolerance).  But this
does not prove the strict separation from integers.  Nor is this a
local algorithm using only $\omega(N+1),\ldots,\omega(N+J)$:
the convolution coefficients require information throughout the
prefix up to $N+J$.  The missing arithmetic step is to produce strict
certificates uniformly over all nonconstant polynomials.  The known
irrationality theorem supplies no such conclusion for
\eqref{p02:ci18:omegaconvolution} in degrees $D\ge2$.
Thus the criterion is an explicit reduction and a finite verification
tool; its equivalence must not be counted as a transcendence proof.

\subsection{Erd\H{o}s--Borwein: modular combinations and the digit route}

There is a useful positive comparison with modular transcendence.
For $0<q<1$, put
\[
 A_r(q)=\sum_{n\ge1}\left(\frac{q^n}{1-q^n}\right)^r,
 \qquad r\ge1.
\]
In particular $A_1(1/2)=E$.

\begin{proposition}[Three modular combinations of six moments]
\label{p02:ci18:modularmoments}
For every real algebraic $q$ with $0<q<1$, the following three numbers
are algebraically independent over $\mathbb Q$:
\begin{align*}
 X_1&=A_1(q)+A_2(q),\\
 X_3&=A_1(q)+7A_2(q)+12A_3(q)+6A_4(q),\\
 X_5&=A_1(q)+31A_2(q)+180A_3(q)+390A_4(q)
                   +360A_5(q)+120A_6(q).
\end{align*}
Hence $\operatorname{trdeg}_{\mathbb Q}
\mathbb Q(A_1(q),\ldots,A_6(q))\ge3$.
\end{proposition}

\begin{proof}
For $|x|<1$, differentiate the geometric series, or expand powers in
falling factorials, to obtain
\[
 \sum_{m\ge1}m^j x^m
 =\sum_{r=1}^{j+1}(r-1)!
       \left\{\begin{matrix}j+1\\r\end{matrix}\right\}
       \left(\frac{x}{1-x}\right)^r.
\]
Here the braces are Stirling numbers of the second kind.  For
$j=1,3,5$, summing this identity at $x=q^n$ gives the three displayed
combinations and
\[
 X_j=\sum_{n,m\ge1}m^j q^{nm}
     =\sum_{t\ge1}\sigma_j(t)q^t.
\]
Absolute convergence justifies the rearrangements.
The three Ramanujan series are
$1-24X_1$, $1+240X_3$, and $1-504X_5$.
Nesterenko's theorem \cite{p02:ci18:Nesterenko} says that these three
values, together with $q$, generate a field of transcendence degree
at least three.  Since $q$ is algebraic, the three values are
algebraically independent.  Invertible rational affine changes of
the three coordinates preserve independence, proving the claim.
\end{proof}

This gives genuine transcendental quantities in the same denominator
family as $E$, but it does not determine $A_1(1/2)$ alone.  Even the
first identity only proves that $E+A_2(1/2)$ is transcendental.
Extracting one coordinate from an independent collection of linear
combinations requires additional arithmetic relations; none are
proved here.  Proposition~\ref{p02:ci18:modularmoments} is explicitly
credited as a consequence of Nesterenko's established theorem.

The alternative digit route must also account for recent prior work.
Campbell's 2026 preprint proves infinitely many occurrences of the
binary word $11$ \cite{p02:ci18:Campbell}.  A public September 2026
document package \cite{p02:ci18:Disjunctivity} claims every finite binary
word, together with a quantitative occurrence bound; its accompanying
Lean verification is explicitly conditional on prime-distribution
inputs.  We read the released proof, including its midpoint encoding,
CRT modulus, divisor average and tail comparison.  We found no
obvious gap in those steps, but this is not a claim of specialist peer
review or unconditional machine verification.  We neither claim that
result as new nor use it as a dependency of any theorem above.
If its stated disjunctivity conclusion is taken as established, then
$p_E(n)=2^n$, so the low-complexity criterion cannot prove
transcendence of $E$ either.  Disjunctivity asserts occurrence, not
normality or prescribed limiting digit frequencies, and does not by
itself imply transcendence.


% Integrated source: target_register.tex
\section{Target register and the next missing arithmetic inputs}
\label{p02:root:register}

The register contains eighteen specified targets or target families.
Rows 10--13 share one general obstruction and are not four independent
solutions.  The ten Pell examples belong together in row 3.  A proved
method obstruction is not a proved arithmetic classification of its
target constant.  This distinction is essential to the accounting.

\begingroup\small
\renewcommand{\arraystretch}{1.15}
\begin{longtable}{L{0.04\textwidth}L{0.23\textwidth}L{0.25\textwidth}L{0.35\textwidth}}
\toprule
No.&Target&Result here&Exact qualification\\
\midrule\endfirsthead
\toprule No.&Target&Result here&Exact qualification\\\midrule\endhead
1&Mixed-sign rank-two threshold&Complete classification&One algebraic exceptional setting; all other algebraic settings give transcendental values.\\
2&$\mathcal V_{d,2m}$, rational $d>0$&Transcendental&Every positive integer $m$; polynomial elimination also treats transcendental boundaries.\\
3&Pell-selected optimal values&Algebraically independent&Ten displayed values, and an infinite family; one selected family, not ten historical conjectures.\\
4&$c_{9/8}$ and the rational $\kappa$ family&Transcendental&Fixed-shape critical extraction rate; a different coordinate from the endpoint shape.\\
5&$A(1)$&Transcendence unresolved&Earlier irrationality survives; an algebraic-irrational parameter would put the Kummer function outside the $E$ class.\\
6&$A(3/2)$&Transcendence unresolved&Explicit positive-integral reduction; no suitable specialization theorem proved.\\
7&$\mathcal V_{3,1}$&Arithmetic unresolved&Two explicit common-zero equations; boundary transcendence cannot decide the objective.\\
8&Euler--Gompertz $\delta$&Positive-cone optimum proved&Unique minimizer $(t-1)^2$ over all nonnegative integer polynomials; no individual irrationality proof.\\
9&Euler $\gamma$&Equivalent lattice criterion&The needed irrational-form hypothesis is equivalent to the target; Oppenheim cannot bypass it.\\
10&Catalan $G$&Full polynomial-cone obstruction&Uses the stated standard denominator clearance; smaller reduced denominators are not ruled out.\\
11&$\zeta(5)$&Same obstruction&All positive integer-polynomial shapes covered; no irrationality proof.\\
12&$\zeta(7)$&Same obstruction&An instance of the uniform $s\ge2$ theorem.\\
13&$\zeta(9)$&Same obstruction&An instance of that same theorem, not a separate mechanism.\\
14&Prime-binary $P_2$&Complexity route audited&The block-complexity theorem below does not establish transcendence of the real number.\\
15&Prime Lambert $L_{\mathbb P}(2)$&Transcendence unresolved&Known irrationality is credited to its source; carries prevent transferring prime digit conclusions directly.\\
16&Erd\H{o}s--Borwein $E$&Transcendence unresolved&Existing irrationality and recent digit work are separated from an arithmetic proof.\\
17&Rank-three spectral example&Conditional transcendence&Follows from the displayed exponential-line bound, in particular Schanuel's conjecture.\\
18&Rational-mesh table $A(u+vn)$&No polynomial recurrence&Structural consequence of the inherited asymptotic; not an individual-value classification.\\
\bottomrule
\end{longtable}
\endgroup

\paragraph{Where an additional breakthrough would have to enter.}
For the spectral theorem, a bound of transcendence degree at least
$d-1$ on the specific exponential line would meet the proved geometric
upper bound $d-2$.  This is a concrete higher-rank target.  The argument
does not require the full force of Schanuel if that weaker line bound
can be obtained directly for a chosen spectrum.

For an individual endpoint shape $A(d)$, rational-parameter
$E$-function theory has already done all it can in the present proof:
the remaining algebraic-irrational parameters produce a function
outside that class.  A theorem about those parameter-dependent
special values, or a different arithmetic representation of the contact,
is needed.  The nonrecurrence theorem does not replace such an input.

For the Gompertz family, exact analytic reconstruction from all second
differences still permits rational first values in the countermodel.
An arithmetic connection statement specific to the actual exponential
integral, rather than another limiting or convexity argument alone,
is necessary.  For the positive Catalan and zeta kernels, a successful
replacement must change the denominator analysis, the kernel class,
or the arithmetic framework; merely optimizing a positive polynomial
within the proved class is insufficient.

The present work therefore supplies definite new deductions relative to
the preceding manuscript and several exact failure diagnoses.  It does
not supply a basis for announcing a resolution of a famous individual
constant conjecture, ten historical closures, or an Annals-level
breakthrough.  Historical priority and journal significance require
further comparison and external mathematical assessment.

% Appendix counter reset suppressed in this unified edition.
\section{Verification and source dependencies}

This is a single-file \LaTeX{} manuscript.  The bibliography, tables,
and proof text are embedded; no external input, bibliography database,
or image is needed to compile it.  Classical theorems remain cited
mathematical inputs rather than being reproved in full.

The spectral factorization argument was checked independently for
multiple components, repeated factors, a possibly complex factorization,
and the real logarithmic Hessian.  The stopping-value calculation was
independently checked in both directions, including the degenerate
Wronskian constant, the exceptional algebraic boundaries, and the
relation-lattice proof.  All ten Pell pairs satisfy their norm equations
by exact integer arithmetic.  The critical-rate proof was checked for
uniqueness of the global minimum, the double-zero argument, and both
directions of its algebraic finite-fiber implication.  The positive-cone
proof was independently checked in every degree and in its equality case;
the compact-interval constants and specified denominator clearances were
also checked.  These are internal cross-reviews, not proof-assistant
certificates or external refereeing.

The boundary asymptotic and the exponential-integral independence result
are inherited from the identified project manuscripts.  The former is
the input to Theorem~\ref{p02:be18:nonholonomic}; the latter is the input to
Proposition~\ref{p02:ac18:curvature}.  Their consequences are distinguished
from new foundational transcendence theorems.  The full even-payoff and
spectral proofs do not depend on the conjectural arithmetic of Euler's
constant, Catalan's constant, or an individual odd zeta value.

Numerical approximations, where displayed, identify constants and do not
prove their irrationality.  Absence of a literature search match does
not establish novelty.  Unreviewed contemporary claims are not silently
promoted to theorems.  The explicitly conditional spectral corollary is
not used in any statement marked unconditional.

% END SOURCE: KEIO_Arithmetic_of_Optimal_Values.tex

\endgroup
\clearpage
\part{Independence and Entropy Optima}
\label{merge:part03}
\begingroup

% BEGIN SOURCE: KEIO_Arithmetic_Independence_and_Entropy_Optima.tex

\noindent\textbf{Editorial integration note.} 
The conditional arbitrary-rank spectral framework and the unconditional cubic branch have different scopes. The later planar stopping theorem is in Part~\ref{merge:part04}; the field-closure development continues in Part~\ref{merge:part05}.


\begin{quote}\small
\noindent\textit{Original part framing.}\par
\noindent\textbf{Status.} Working manuscript. Proofs are supplied for
unconditional statements, external theorems are identified, and the
squarefree-factorization implication is explicitly conditional.
Independent internal cross-reading and symbolic checks are not external
peer review or a formal proof certificate.
\end{quote}
\section{The research question and the actual increase in scope}
\label{p03:root3:scope}

This is a sequel to \emph{Arithmetic of Optimal Values}
\cite{p03:root3:previousII}, not a reprint of its results. The earlier
223-page foundation \cite{p03:root3:original}, 447-page sequel compilation
\cite{p03:root3:sequels}, and first arithmetic continuation
\cite{p03:root3:previousI} supply context and explicitly identified analytic
inputs. Page counts identify documents; they do not measure mathematical
strength. The present question is whether those mechanisms can classify
additional individual constants or prove relations substantially stronger
than separate irrationality.

The strongest changes are shown below. Applications within one uniform
theorem are not counted as separate historical conjectures.

\begin{center}
\small
\begin{tabular}{L{0.20\textwidth}L{0.32\textwidth}L{0.36\textwidth}}
\toprule
Direction & Previous boundary & Result proved here\\
\midrule
Spectral values & Higher-rank classification conditional on an exponential-line bound & Unconditional classification throughout real cubic fields, and an identified independent value/contact pair\\[5pt]
Stopping values & Infinite independent family obtained by varying dimension and payoff & Infinite independent family in every fixed even dimension; all prime payoff indices in dimension two\\[5pt]
Gamma extraction & Shape-one interior rate arithmetic & Every positive integer shape; exact double contacts and an explicit finite bound on their number\\[5pt]
Information objectives & Selected arithmetic values and optimizer locations & All rational finite signed-entropy critical values, including nonconvex objectives and boundary supports\\[5pt]
Classical constants & Bounds using standard denominator clearances & Actual reduced-denominator bounds, together with a proof that primitive pair reduction is a different obstruction\\[5pt]
Modular transfer & Modular combinations of moment functions & Three-function algebraic independence over finite dilated modular differential fields; no value-specialization conclusion\\
\bottomrule
\end{tabular}
\end{center}

\subsection{Twenty-two targets and their outcomes}

The register distinguishes arithmetic of a \emph{number} from arithmetic
of a family, and functional independence from independence after
specialization. ``Unresolved here'' describes this manuscript's result,
not a blanket claim about every announcement in the literature.
In particular, the September 24, 2026 OpenAI Math Release preprint
\cite{p03:root3:Catalan} asserts the irrationality of Catalan's constant.
Its full proof is neither audited nor used here. Its signed mixed-moment
determinants fall outside our restricted positive-kernel theorem.

\small
\begin{longtable}{L{0.045\textwidth}L{0.33\textwidth}L{0.51\textwidth}}
\toprule
No. & Target & Outcome and exact limitation\\
\midrule
\endfirsthead
\toprule No. & Target & Outcome and exact limitation\\\midrule
\endhead
\bottomrule\endfoot
1 & Euler's constant $\gamma$: irrationality & Unresolved here. Finite differential cancellation eliminates its coefficient rather than isolates it (Proposition~\ref{p03:ac19:monodromy}).\\[4pt]
2 & Euler--Gompertz $\delta$: irrationality & Unresolved here. The new $E$-value division result does not classify this connection coefficient.\\[4pt]
3 & Catalan $G$: individual arithmetic & No new proof. Actual reduced-denominator obstruction for a specified class (Theorem~\ref{p03:ac19:catalan-reduced}); the external claim above is kept separate.\\[4pt]
4 & $\zeta(5)$: irrationality & Unresolved here. Reduced-denominator cancellation is controlled for an entire support-restricted kernel class (Theorem~\ref{p03:ac19:zeta-reduced}).\\[4pt]
5 & Erd\H{o}s--Borwein $\sum(2^n-1)^{-1}$: transcendence & Unresolved as a number. Its generating function participates in Theorem~\ref{p03:ci19:triple}.\\[4pt]
6 & Prime binary constant $\sum_p2^{-p}$: transcendence & Unresolved as a number. All-order prime asymptotics yield functional independence, not specialization.\\[4pt]
7 & Prime Lambert constant $\sum_p(2^p-1)^{-1}$: transcendence & Same distinction; the triple functional theorem is one result, not three solved arithmetic conjectures.\\[4pt]
8 & Endpoint shape $A(1)$: transcendence & Unresolved. Irrationality was proved previously; integer differential operators do not address the remaining parameter.\\[4pt]
9 & Endpoint shape $A(3/2)$: transcendence & Unresolved for the same precise reason (Section~\ref{p03:be19:section}).\\[4pt]
10 & Linear excursion optimum $\mathcal V_{3,1}$: transcendence & Conditional on the published squarefree-factorization conjecture; an explicit finite-multiple-zero reduction is proved.\\[4pt]
11 & All positive rational $(d,q)$ bridge optima & The same conditional family theorem (Theorem~\ref{p03:ga19:allconditional}); not an additional independent closure.\\[4pt]
12 & Algebraic values for cubic-field spectra & Complete unconditional classification (Theorem~\ref{p03:sc3:thm:cubic}).\\[4pt]
13 & Joint arithmetic of a cubic two-atom value and contact & Algebraic independence, stronger than separate transcendence (Theorem~\ref{p03:sc3:thm:pair}).\\[4pt]
14 & Biquadratic optimized value $Q$ in \eqref{p03:sc3:eq:quartic} & Value unresolved; its contact and elasticity are now individually transcendental.\\[4pt]
15 & Infinitely many independent optima for one fixed bridge dimension & Proved in every positive even dimension (Theorems~\ref{p03:ga19:planar} and~\ref{p03:ga19:evenfixed}), one uniform family.\\[4pt]
16 & Fixed integer-shape gamma critical rates & All interior critical rates and contacts classified as transcendental; exact contact multiplicity and count bound also proved.\\[4pt]
17 & Rational finite-channel capacity & Algebraic optimizer; algebraic exponential of capacity; effective rational/transcendental test in bits (Corollary~\ref{p03:mi3:cor:capacity}).\\[4pt]
18 & Rational finite-source rate--distortion & Same arithmetic classification (Corollary~\ref{p03:mi3:cor:RD}), an application of the common entropy theorem.\\[4pt]
19 & Rational information-bottleneck optimum & Same conclusion despite nonconvexity, including all finite encoder alphabets without extra constraints (Corollary~\ref{p03:mi3:cor:IB}).\\[4pt]
20 & $1/\operatorname{Ein}(1)$ in the ring of $E$-values & Nonmembership proved using a new external morphism theorem; an exact rational-function intersection is described (Theorem~\ref{p03:ac19:rational-division}).\\[4pt]
21 & Entropy of $\operatorname{Poisson}(1)$: irrationality & Unresolved. Infinite logarithmic coefficients are outside the finite entropy proof.\\[4pt]
22 & The earlier optimized exponential-integral payoff & Unresolved. Algebraic-point specialization does not apply at its transcendental optimizer.\\
\end{longtable}
\normalsize

The register is an account of attempts. It does not turn related
corollaries into twenty-two independent discoveries. In particular, the
three information-theoretic classifications share one proof, the three
prime/divisor functions share another, and the general rational stopping
claim remains conditional. The additional reciprocal-moment targets
$A_r(1/2)$ are covered by the same functional analysis and are not used
to inflate the count.

\subsection{Proof architecture and source dependence}

The cubic argument is an intersection of two inequalities: tangency and
indefinite curvature force transcendence degree at most one, whereas
Brownawell's classical theorem forces at least two away from the unit
contact. The fixed-dimension stopping argument instead makes algebraic
boundary degrees carry new prime factors. These factors prevent rational
linear relations between the boundary points; Lindemann--Weierstrass then
excludes polynomial relations between their exponential optimal values.
The entropy argument uses no infinite special-function expansion:
rational logarithmic stationarity becomes a finite semialgebraic system,
and a rational anchor on each face determines the arithmetic of its
critical values.

These mechanisms are genuinely different. Their coexistence gives a
broader research program, not evidence by itself for publication at a
particular level. The precise priority and significance of the transfer
theorems require a specialist literature review beyond the bounded
primary-source comparisons recorded here.


% Integrated source: spectral_results.tex
\section{Cubic spectra: unconditional classification and independent observables}
\label{p03:sc3:sec}

For the optimized spectral boundary
\begin{equation}\label{p03:sc3:eq:model}
 S(x)=\sum_{j=1}^{m}w_jx^{r_j},\qquad
 T(C)=\min_{x>0}\frac{1+Cx}{S(x)},\qquad 0<r_j<1,
\end{equation}
the previous manuscript classified algebraic values when the rational
span of $1,r_1,\ldots,r_m$ has dimension two.  The higher-dimensional
extension was conditional on an exponential-line transcendence-degree
bound.  Here we remove that condition for every spectrum contained in
a real cubic number field.  The decisive classical input is a theorem
of Brownawell that applies even when the base of the powers is
transcendental.  Its use is essential: the contact in
\eqref{p03:sc3:eq:model} is not given to be algebraic.

There are two further conclusions.  For a two-atom spectrum spanning a
cubic field, the optimized value and its contact are algebraically
independent.  For any spectrum in an affine irrational lattice, including
lattices outside cubic fields, the contact and the elasticity each have
an exact algebraic-versus-transcendental classification.  These are
unconditional classifications of the present variational constants.
The classical transcendence theorems are credited explicitly; no
previously named arithmetic conjecture is claimed to be resolved.

\subsection{Contact geometry and the precise cubic input}

Throughout this section, weights and settings are positive real
algebraic numbers, all exponents are real algebraic, and real powers
use the real logarithm.  Write
\begin{equation}\label{p03:sc3:eq:exception}
 W=\sum_jw_j,\qquad M=\sum_jw_jr_j,\qquad
 C_* =\frac{M}{W-M}.
\end{equation}
Partial derivatives with respect to $C$ or a weight are taken in the
open positive parameter domain, holding the exponents fixed.

\begin{lemma}[Contact identities]\label{p03:sc3:lem:contact}
The minimum in \eqref{p03:sc3:eq:model} has a unique contact $x_C>0$.
The contact and value depend real analytically on the positive parameters.
With $x=x_C$, $\tau=T(C)$ and $H=S-xS'$, one has
\begin{align}
 C&=\frac{S'(x)}{H(x)},&
 S(x)&=\frac{1+Cx}{\tau},& S'(x)&=\frac C\tau,
 \label{p03:sc3:eq:contact}\\
 T_C&=\frac{x}{S(x)}=\frac{xT}{1+Cx},&
 B:=T-CT_C&=\frac1{S(x)}>0,&
 T_{w_j}&=-\frac{T x^{r_j}}{S(x)}.
 \label{p03:sc3:eq:envelope}
\end{align}
Moreover $T_{CC}<0$, and
\begin{equation}\label{p03:sc3:eq:unit}
 x_C=1\quad\Longleftrightarrow\quad C=C_*,\qquad
 T(C_*)=\frac1{W-M}.
\end{equation}
\end{lemma}
\begin{proof}
The objective tends to infinity at both ends of $(0,\infty)$, so a
minimum exists.  Each summand of $H$ is $w_j(1-r_j)x^{r_j}>0$, while
$S''(x)<0$.  Direct differentiation gives
\[
 \left(\frac{S'}H\right)'=\frac{S''S}{H^2}<0.
\]
The critical equation is $C=S'/H$, so there is at most one critical
point and hence exactly one minimum.  Its derivative with respect to
$C$ is the reciprocal of the displayed negative derivative.  The
analytic implicit-function theorem gives the asserted parameter
dependence.  The two value identities follow from stationarity.
Differentiating the objective at its contact gives the envelope formulas,
and $T_{CC}=H(x)x_C'/S(x)^2<0$.  Finally $S(1)=W$ and $S'(1)=M$ yield
\eqref{p03:sc3:eq:unit}.
\end{proof}

\begin{lemma}[Classical cubic exponential-line bound]
\label{p03:sc3:lem:brownawell}
If $a>0$, $a\ne1$, and $\beta$ is real algebraic of degree three, then
\begin{equation}\label{p03:sc3:eq:brownawell}
 \operatorname{trdeg}_{\overline{\mathbb Q}}
 \overline{\mathbb Q}(a,a^\beta,a^{\beta^2})\ge2.
\end{equation}
There is no algebraicity assumption on $a$.
\end{lemma}
\begin{proof}
Brownawell's Theorem~2 \cite{p03:sc3:Brownawell}, with the real logarithm
of $a$, asserts in particular that relatively prime nonzero integer
polynomials $P(X,Y)$ and $Q(X,Y,Z)$ cannot both vanish at
$(a,a^\beta,a^{\beta^2})$.  We spell out the transcendence-degree
consequence to retain the crucial arbitrary-base hypothesis.

Suppose the three numbers generate a field of transcendence degree at
most one over $\mathbb Q$.  If $a$ is algebraic, take for $P(X)$ its
minimal polynomial.  The other two numbers are algebraically dependent,
so a nonzero $Q(Y,Z)\in\mathbb Z[Y,Z]$ vanishes at them.  These two
polynomials are relatively prime, a contradiction.  If $a$ is
transcendental, both other numbers are algebraic over $\mathbb Q(a)$.
Their minimal polynomials, after clearing denominators and taking
primitive parts, give irreducible polynomials $P(X,Y)$ and $Q(X,Z)$
over $\mathbb Q$, of positive degree in $Y$ and $Z$, respectively.
They remain irreducible upon adjoining the unused variable and cannot
be associates, so they are relatively prime.  Again Brownawell's
theorem gives a contradiction.  Extending the base from $\mathbb Q$
to its algebraic closure does not change transcendence degree, proving
\eqref{p03:sc3:eq:brownawell}.
\end{proof}

The input in Lemma~\ref{p03:sc3:lem:brownawell} is entirely classical.
It is stronger for the present purpose than the algebraic-base cubic
power theorem alone.  We use only its qualitative nonvanishing
consequence, not its quantitative measure.

\subsection{The algebraic tangency obstruction}

We recall the curvature argument from the preceding manuscript, with
its proof, because it identifies exactly how the cubic input enters.
For a Laurent polynomial $F$ and a positive point $z$, its
\emph{logarithmic Hessian} means the Hessian at $y=0$ of
$F(z_1e^{y_1},\ldots,z_de^{y_d})$.

\begin{lemma}[Curvature forces codimension two]
\label{p03:sc3:lem:curvature}
Let $d\ge2$, let $b_1,\ldots,b_d$ be rationally independent real
algebraic numbers, and put
\[
 D=\sum_{i=1}^{d}b_i Z_i\frac{\partial}{\partial Z_i}.
\]
Suppose a Laurent polynomial $F$ with real algebraic coefficients and
a positive point $z$ satisfy $F(z)=DF(z)=0$.  If the logarithmic
Hessian of $F$ at $z$ is indefinite, then
\begin{equation}\label{p03:sc3:eq:codimension}
 \operatorname{trdeg}_{\overline{\mathbb Q}}
       \overline{\mathbb Q}(z_1,\ldots,z_d)\le d-2.
\end{equation}
\end{lemma}
\begin{proof}
Work in the Laurent unique-factorization domain
$R=\overline{\mathbb Q}[Z_1^{\pm1},\ldots,Z_d^{\pm1}]$.
For every nonunit irreducible $H\in R$, the polynomials $H,DH$ are
relatively prime.  Indeed, $DH=0$ would force $H$ to be constant:
the weights $b\cdot k$, $k\in\mathbb Z^d$, are distinct, and only
the constant monomial has weight zero.  If $DH=AH$ with $A\ne0$,
Newton polytopes give
\[
 \operatorname{Newt}(A)+\operatorname{Newt}(H)
 =\operatorname{Newt}(DH)\subseteq\operatorname{Newt}(H).
\]
Taking support functions in every direction forces
$\operatorname{Newt}(A)\subseteq\{0\}$, so $A$ is constant.
The equation $DH=AH$ then makes $H$ a single monomial, again a unit.

We also record why two relatively prime Laurent polynomials vanishing
at $z$ imply \eqref{p03:sc3:eq:codimension}.  Degree $d$ is excluded by
either equation.  If the degree were $d-1$, choose $d-1$ coordinates
forming a transcendence basis from among the coordinates of $z$.
The remaining coordinate has a minimal polynomial over their rational
function field.  Clearing denominators and taking its primitive part
gives a nonunit Laurent polynomial dividing every Laurent polynomial
that vanishes at $z$, by Gauss's lemma.  The two given polynomials
would have a common factor, a contradiction.

Factor $F$ in $R$.  If two distinct irreducible factors vanish at
$z$, the preceding paragraph applies.  Otherwise there is one
vanishing irreducible factor $H$, and write $F=H^eR_0$, with
$R_0(z)\ne0$.  For $e=1$, the equation $DF(z)=0$ gives $DH(z)=0$,
so the first two paragraphs apply.  It remains to exclude $e\ge2$.
In logarithmic coordinates write $h(y)=H(ze^y)$ and
$r(y)=R_0(ze^y)$.  Since $h(0)=0$, the Hessian of $h^e r$ is zero
if $e\ge3$, and for $e=2$ it is
\[
               2r(0)\,\nabla h(0)\nabla h(0)^{\mathsf T}.
\]
This matrix has complex rank at most one.  It is also real symmetric
because $F$ has real coefficients.  Its real rank is therefore at
most one, so it cannot be indefinite.  This contradiction completes
all cases, including nonreal irreducible factors and singular levels.
\end{proof}

\begin{theorem}[Unconditional cubic-field classification]
\label{p03:sc3:thm:cubic}
Let $K\subset\mathbb R$ be a number field of degree three.  Suppose
the exponents in \eqref{p03:sc3:eq:model} belong to $K$, and at least
one of them is irrational.  Then, for every positive algebraic
setting $C$ and positive algebraic weights,
\begin{equation}\label{p03:sc3:eq:classification}
 T(C)\in\overline{\mathbb Q}
 \quad\Longleftrightarrow\quad x_C=1
 \quad\Longleftrightarrow\quad C=C_*.
\end{equation}
At the exceptional setting the value is $1/(W-M)$.  There is no
restriction on the signs of the coordinates of the exponents in a
rational basis of $K$.
\end{theorem}
\begin{proof}
Choose a real cubic generator $\beta$ of $K$.  For a common positive
integer $q$, write
\[
 r_j=\frac{m_{j0}+m_{j1}\beta+m_{j2}\beta^2}{q},
 \qquad m_{ji}\in\mathbb Z.
\]
At a contact $x=x_C$, put $u=x^{1/q}$, $v=u^\beta$, $z=u^{\beta^2}$.
Suppose $\tau=T(C)$ is algebraic and form
\begin{align*}
 P(U,V,Z)&=\sum_jw_jU^{m_{j0}}V^{m_{j1}}Z^{m_{j2}},\\
 F(U,V,Z)&=P(U,V,Z)-\frac{1+CU^q}{\tau},\\
 D&=U\partial_U+\beta V\partial_V+\beta^2 Z\partial_Z.
\end{align*}
The contact equations give $F(u,v,z)=DF(u,v,z)=0$.  Along the
exponential line, $D$ acts as $qx\,d/dx$.  Consequently
\begin{equation}\label{p03:sc3:eq:opposite}
 D^2F(u,v,z)=q^2x^2S''(x)<0.
\end{equation}
The additional first-derivative term cancels by $S'(x)=C/\tau$.
Some $m_{ji}$ with $i\in\{1,2\}$ is nonzero, since one exponent
is irrational.  The corresponding coordinate field $E_i$ is
$V\partial_V$ or $Z\partial_Z$ and satisfies
\[
 E_i^2F(u,v,z)=\sum_jw_jm_{ji}^{2}x^{r_j}>0.
\]
Thus the logarithmic Hessian is indefinite.  The rational independence
of $1,\beta,\beta^2$ and Lemma~\ref{p03:sc3:lem:curvature} give
$\operatorname{trdeg}_{\overline{\mathbb Q}}
\overline{\mathbb Q}(u,v,z)\le1$.
If $x\ne1$, then $u\ne1$ and
Lemma~\ref{p03:sc3:lem:brownawell} gives the contradictory lower bound two.
Therefore $x=1$.  The converse and the exceptional value follow from
Lemma~\ref{p03:sc3:lem:contact}.
\end{proof}

For spectra of rational rank three, Theorem~\ref{p03:sc3:thm:cubic}
removes a genuine conditional hypothesis from the preceding result.
For spectra already contained in a two-dimensional rational space it
is consistent with, and does not recount, the previous classification.

\subsection{Two atoms identify an independent pair}

\begin{theorem}[The value and contact are algebraically independent]
\label{p03:sc3:thm:pair}
Suppose $m=2$ and
\begin{equation}\label{p03:sc3:eq:fullrank}
 \operatorname{span}_{\mathbb Q}\{1,r_1,r_2\}=K,
 \qquad [K:\mathbb Q]=3.
\end{equation}
At every positive algebraic setting $C\ne C_*$ and for positive
algebraic weights, the pair $(T(C),x_C)$ is algebraically independent
over $\overline{\mathbb Q}$.  Each of the pairs
\[
       (T,T_C),\qquad (x_C,T_C),\qquad
       \left(T,\eta\right),\quad \eta=\frac{CT_C}{T},
\]
is also algebraically independent.  The triple $(T,x_C,T_C)$ has
transcendence degree exactly two and satisfies
\begin{equation}\label{p03:sc3:eq:triple}
                         (1+Cx_C)T_C=x_CT.
\end{equation}
At $C=C_*$ all these numbers are algebraic.
\end{theorem}
\begin{proof}
Use the same $\beta,q,u,v,z$ as in the preceding proof, and let
$L=\overline{\mathbb Q}(u,v,z)$.  Since $x_C\ne1$, Brownawell's
bound gives $\operatorname{trdeg}L\ge2$.  Stationarity gives the
Laurent relation
\begin{equation}\label{p03:sc3:eq:stationaryLaurent}
 \sum_{j=1}^{2}w_j
       \bigl[r_j+(r_j-1)CU^q\bigr]
       U^{m_{j0}}V^{m_{j1}}Z^{m_{j2}}=0
 \quad\hbox{at }(u,v,z).
\end{equation}
This Laurent polynomial is nonzero.  Its first group of monomials
has $D$-weights $qr_j\in(0,q)$, and its second group has weights
$q(r_j+1)\in(q,2q)$; the groups cannot cancel.  Within either
group the two exponents are distinct and the coefficients are nonzero.
Thus $\operatorname{trdeg}L\le2$, and equality holds.

Write $A_j=w_jx^{r_j}$ and, after interchanging the labels if necessary,
assume $r_1<r_2$.  The equations $A_1+A_2=(1+Cx)/T$ and
$r_1A_1+r_2A_2=Cx/T$ yield
\begin{align}
 A_1&=\frac{r_2(1+Cx)-Cx}{T(r_2-r_1)},&
 A_2&=\frac{Cx-r_1(1+Cx)}{T(r_2-r_1)}.
 \label{p03:sc3:eq:recoveratoms}
\end{align}
Hence
\[
 \overline{\mathbb Q}(T,x)
 =\overline{\mathbb Q}(x,x^{r_1},x^{r_2}).
\]
The inclusion from left to right uses $T=(1+Cx)/S(x)$, and the
reverse inclusion uses \eqref{p03:sc3:eq:recoveratoms} and nonzero weights.
By \eqref{p03:sc3:eq:fullrank}, each of $1/q,\beta/q,\beta^2/q$ is a
rational linear combination of $1,r_1,r_2$.  Clearing denominators
shows that $u,v,z$ are algebraic over this field: a positive integer
power of each is a Laurent monomial in $x,x^{r_1},x^{r_2}$.
Conversely those three powers already belong to $L$.
Therefore $L$ is algebraic over $\overline{\mathbb Q}(T,x)$, and
the latter field has transcendence degree two.

The envelope identities and their inverses give
\[
 T_C=\frac{xT}{1+Cx},\qquad
 x=\frac{T_C}{T-CT_C},\qquad
 T=\frac{(1+Cx)T_C}{x},\qquad
 \eta=\frac{Cx}{1+Cx},\qquad
 x=\frac{\eta}{C(1-\eta)}.
\]
Every denominator is nonzero.  These rational field identities prove
all the pair assertions and \eqref{p03:sc3:eq:triple}.  At the exceptional
setting $x=1$ and the same formulas give algebraic numbers.
\end{proof}

\begin{corollary}[The field of first-order observations]
\label{p03:sc3:cor:gradient}
Suppose an arbitrary finite spectrum satisfies
$\operatorname{span}_{\mathbb Q}\{1,r_1,\ldots,r_m\}=K$, where
$K$ is a real cubic field.  At a nonexceptional positive algebraic
parameter point,
\begin{equation}\label{p03:sc3:eq:gradientfield}
 \operatorname{trdeg}_{\overline{\mathbb Q}}
 \overline{\mathbb Q}(T,T_C,T_{w_1},\ldots,T_{w_m})=2.
\end{equation}
In particular at least one of the first partial derivatives is
algebraically independent of $T$.
\end{corollary}
\begin{proof}
The envelope formulas give the exact field identity
\begin{equation}\label{p03:sc3:eq:fieldidentity}
 \overline{\mathbb Q}(T,T_C,T_{w_1},\ldots,T_{w_m})
 =\overline{\mathbb Q}(x,x^{r_1},\ldots,x^{r_m}),
\end{equation}
because $B=T-CT_C=1/S(x)$, $x=T_C/B$ and
$x^{r_j}=-T_{w_j}/(TB)$.  The full-span assumption makes the
cubic power field $L$ algebraic over the right side, as above.
Brownawell supplies degree at least two.  The nonzero stationary
Laurent relation \eqref{p03:sc3:eq:stationaryLaurent}, with the sum now
over all $j$, supplies degree at most two.  Repeated exponents may
be combined; positive coefficients prevent cancellation within each
of the two disjoint weight intervals.  This proves
\eqref{p03:sc3:eq:gradientfield}.  Theorem~\ref{p03:sc3:thm:cubic} makes $T$
transcendental.  If every derivative were algebraic over
$\overline{\mathbb Q}(T)$, their finite joint extension would have
transcendence degree one, a contradiction.
\end{proof}

\subsection{An individual observable beyond cubic fields}

\begin{theorem}[Affine irrational lattices]
\label{p03:sc3:thm:affine}
Suppose that, after combining repeated exponents, the spectrum has
at least two elements and
\begin{equation}\label{p03:sc3:eq:affine}
 r_j=\alpha+\frac{k_j\theta}{q},\qquad
 \alpha\in\overline{\mathbb Q}\cap\mathbb R,\quad
 \theta\in(\overline{\mathbb Q}\cap\mathbb R)\setminus\mathbb Q,
 \quad q\in\mathbb Z_{>0},\quad k_j\in\mathbb Z.
\end{equation}
The integers $k_j$ are distinct, and all $r_j$ lie in $(0,1)$.
There is no requirement that $\alpha\in\mathbb Q+\mathbb Q\theta$.
Then
\begin{equation}\label{p03:sc3:eq:affineclassification}
 x_C\in\overline{\mathbb Q}
 \quad\Longleftrightarrow\quad
 \eta(C):=\frac{CT_C}{T}\in\overline{\mathbb Q}
 \quad\Longleftrightarrow\quad C=C_*.
\end{equation}
Thus both the contact and elasticity are individually transcendental
at every nonexceptional positive algebraic parameter point.
\end{theorem}
\begin{proof}
Suppose $x=x_C$ is algebraic and $x\ne1$.  By the classical
Gelfond--Schneider theorem \cite{p03:sc3:GS},
$y=x^{\theta/q}$ is transcendental.  On factoring $x^\alpha$ out
of the stationary equation, however, one obtains
\begin{equation}\label{p03:sc3:eq:affinepoly}
 \sum_j w_j\bigl[(1+Cx)r_j-Cx\bigr]y^{k_j}=0.
\end{equation}
Its coefficients are algebraic and its exponents are distinct integers.
It is not the zero Laurent polynomial: otherwise every $r_j$ would
equal $Cx/(1+Cx)$, contrary to distinctness.  Multiplication by a
suitable power of $y$ now makes $y$ algebraic, a contradiction.
Consequently an algebraic contact must equal one.  Conversely the
exceptional contact is one by Lemma~\ref{p03:sc3:lem:contact}.
Finally $\eta=Cx/(1+Cx)$ is a nonconstant fractional linear
function of $x$ with algebraic coefficients and inverse
$x=\eta/[C(1-\eta)]$.  This proves the other equivalence.
\end{proof}

Every pair of algebraic exponents with irrational difference satisfies
\eqref{p03:sc3:eq:affine}: take $\alpha=r_1$, $\theta=r_2-r_1$,
$q=1$ and $k_1=0,k_2=1$.  The theorem therefore classifies an
individual observed number even when neither the cubic-field theorem
nor the earlier rank-two theorem classifies the optimized value.
The assumption of at least two distinct exponents matters.  A single
atom has $x_C=r_1/[C(1-r_1)]$, an algebraic contact for all the
parameters under consideration.

\subsection{A finite jet for an unrestricted finite spectrum}

The next consequence requires more observations but no field-degree
restriction.  It is stated separately from the cubic value theorem:
algebraicity of an entire jet is a substantially stronger hypothesis
than algebraicity of $T$ alone.

\begin{proposition}[Finite-jet obstruction]\label{p03:sc3:prop:jet}
Let the $m\ge2$ exponents be distinct, with at least one irrational.
At a positive algebraic parameter point, the $m$ numbers
\begin{equation}\label{p03:sc3:eq:jet}
                       T,T_C,\ldots,\partial_C^{m-1}T
\end{equation}
are all algebraic if and only if $C=C_*$.
At $C=C_*$ all derivatives of $T$ of every finite order with respect
to $C$ are algebraic.
\end{proposition}
\begin{proof}
Assume first that \eqref{p03:sc3:eq:jet} is algebraic.  The identities
$B=T-CT_C>0$, $x=T_C/B$, $S(x)=1/B$ and $S'(x)=C/T$
give algebraic $x,S(x),S'(x)$.  If $m\ge3$, differentiation gives
\begin{equation}\label{p03:sc3:eq:jetrecursion}
 x_C'=\frac{TT_{CC}}{B^2}\ne0,\qquad
 S''(x)=\frac{B^3}{T^3T_{CC}},\qquad
 S^{(\ell+1)}(x)=\frac{\partial_C[S^{(\ell)}(x_C)]}{x_C'}.
\end{equation}
Induction shows that $S^{(\ell)}(x)$ is a rational expression in
$C,T,\ldots,\partial_C^\ell T$ for every $2\le\ell\le m-1$;
its denominators are nonzero by $T>0$, $B>0$ and $T_{CC}<0$.
Thus all $S^{(\ell)}(x)$ for $0\le\ell\le m-1$ are algebraic.
The identities
\[
 x^\ell S^{(\ell)}(x)
   =\sum_{j=1}^{m}w_j
       r_j(r_j-1)\cdots(r_j-\ell+1)x^{r_j},
 \qquad 0\le\ell\le m-1,
\]
form an invertible linear system for the $w_jx^{r_j}$.  Its matrix
is the falling-factorial Vandermonde matrix; its determinant is the
ordinary Vandermonde product in the distinct $r_j$, since the
falling-factorial polynomials are monic of successive degrees.
Hence every $x^{r_j}$ is algebraic.  An irrational $r_j$ and
Gelfond--Schneider force the positive algebraic $x$ to equal one.

Conversely, at $x=1$, every derivative of $S$ is algebraic.  The
analytic function $\phi(x)=S'(x)/(S(x)-xS'(x))$ has algebraic
Taylor coefficients there and $\phi'(1)\ne0$.  Its local inverse
$x_C$ therefore has algebraic Taylor coefficients at the algebraic
point $C_*$: recursive comparison of coefficients divides only by
the nonzero algebraic number $\phi'(1)$.  Composing with
$(1+Cx_C)/S(x_C)$ proves the final assertion and the converse.
\end{proof}

\subsection{Explicit constants and the remaining boundary}

For an integer $n\ge2$ that is not a perfect cube, put
$\beta=\sqrt[3]{n}$ and choose an integer $q>n^{2/3}$.  Define
\begin{equation}\label{p03:sc3:eq:family}
 A_{n,q}(C)=\min_{x>0}
       \frac{2(1+Cx)}{x^{\beta/q}+x^{\beta^2/q}}.
\end{equation}
Here $1,\beta/q,\beta^2/q$ span a cubic field.  The equal weights
are $1/2$, and their mean $M=(\beta+\beta^2)/(2q)$ is irrational.
Hence $C_*=M/(1-M)$ is irrational.  Every positive rational $C$
is nonexceptional.  Theorems~\ref{p03:sc3:thm:cubic} and
\ref{p03:sc3:thm:pair} prove that $A_{n,q}(C)$ and its contact are
algebraically independent, as are its value and elasticity.
In particular this applies to the three concrete constants
\begin{align}
 A_{2,2}(1)&=\min_{x>0}
       \frac{2(1+x)}{x^{\sqrt[3]{2}/2}+x^{\sqrt[3]{4}/2}},
       \label{p03:sc3:eq:A2}\\
 A_{3,3}(1)&=\min_{x>0}
       \frac{2(1+x)}{x^{\sqrt[3]{3}/3}+x^{\sqrt[3]{9}/3}},
       \label{p03:sc3:eq:A3}\\
 A_{5,3}(1)&=\min_{x>0}
       \frac{2(1+x)}{x^{\sqrt[3]{5}/3}+x^{\sqrt[3]{25}/3}}.
       \label{p03:sc3:eq:A5}
\end{align}
These are three applications of one theorem, not three independent
resolutions of historical conjectures.

For comparison, retain the earlier higher-rank example
\begin{equation}\label{p03:sc3:eq:quartic}
 Q=\min_{x>0}\frac{2(1+x)}{x^{\sqrt2/2}+x^{\sqrt3/3}}.
\end{equation}
The exponents belong to a biquadratic field and have irrational
difference.  Their sum is not one, so the setting $C=1$ is
nonexceptional.  Theorem~\ref{p03:sc3:thm:affine} now proves that its
contact and elasticity are each transcendental.  Proposition~\ref{p03:sc3:prop:jet}
also excludes simultaneous algebraicity of its value and first
$C$-derivative.  The present arguments do \emph{not} prove the
transcendence of $Q$ itself.  Brownawell's cubic bound does not
apply merely because the rational span of a spectrum has dimension
three: its embedding into a cubic field is an essential hypothesis.

The new unconditional gain is therefore precise.  Cubic-field
thresholds are classified completely, two full-rank cubic atoms yield
an identified independent pair, and affine-lattice contacts and
elasticities are classified individually.  Spectra outside these
arithmetic structures, and infinite prime-frequency boundaries,
retain their previously stated barriers.


% Integrated source: gaussian_results.tex
\section{Arithmetic of a single bridge and a finite-multiple-zero reduction}
\label{p03:ga19:section}

Three questions are investigated here: the arithmetic of the linear
Brownian-excursion optimum, the classification of optimal values for
general positive rational payoff powers, and joint independence of
different payoff optima for one fixed dimension.  The third question
has an unconditional answer in every positive even integral dimension.
In dimension two, all prime-indexed even-payoff optima, together with
$e$, are algebraically independent.  This conclusion uses a new
Eisenstein calculation for the boundary polynomials; the earlier
construction varied the dimension and used Pell equations.

For the first two questions we obtain a different kind of progress.
We prove that an explicit $E$-function associated with a hypothetical
algebraic optimum has only finitely many multiple zeros, including a
known transcendental double zero.  The squarefree-factorization
conjecture of Fischler--Rivoal would exclude this configuration.
Consequently that published conjecture implies transcendence of
\emph{every} rational-parameter bridge optimum, including the linear
excursion value.  The factorization conjecture is not proved here,
and these individual conclusions remain conditional.

\subsection{Notation and the already established inputs}

For a Bessel bridge of dimension $d>0$ from zero to zero in unit time,
write $\mathcal V_{d,q}$ for the optimal expected payoff $R_\tau^q$,
where $q>0$.  Set $a=q/2$, $b=d/2$, and
\[
 f(z)=M(a;b;z)=\sum_{n\geq0}\frac{(a)_n}{(b)_n}\frac{z^n}{n!}.
\]
The established stopping formula of Hobson--Liu \cite{p02:ga18:HL} is
\begin{equation}\label{p03:ga19:setup}
 z_* f'(z_*)=a f(z_*),\qquad
 C_{d,q}=\sqrt{2z_*},\qquad
 \mathcal V_{d,q}=\frac{(2z_*)^a}{f(z_*)},
\end{equation}
where $z_*>0$ is unique.  For integer $d$, this is the radial problem
for an ordinary $d$-dimensional Brownian bridge; in particular the
fixed dimensions used below do not require a fractional-dimensional
interpretation.

Here is a short verification of the uniqueness and nondegeneracy
facts needed later.  With $c_n=(a)_n/((b)_n n!)>0$, give the
nonnegative integers probabilities $c_nz^n/f(z)$ for real $z>0$.
Then $u(z)=zf'(z)/f(z)$ is their mean and
\begin{equation}\label{p03:ga19:variance}
                 z u'(z)=\operatorname{Var}_z(n)>0.
\end{equation}
The mean starts at zero and tends to infinity.  To see the latter
without differentiating an asymptotic expansion, for each fixed $K$
the total weight on $n\leq K$ tends to zero, by comparison with the
single term of index $K+1$.  Thus $u(z)=a$ has one solution.  Moreover
$h(z)=f(z)/z^a$ has a strict, nondegenerate minimum at $z_*$, since
\begin{equation}\label{p03:ga19:hsecond}
                 h''(z_*)=h(z_*)\frac{u'(z_*)}{z_*}>0.
\end{equation}

We recall, rather than recount as new, the earlier exact boundary
classification \cite{p03:root3:previousI,p03:root3:previousII}:
\begin{equation}\label{p03:ga19:old}
 a,b\in\mathbb Q_{>0}:\quad
 z_*\in\overline{\mathbb Q}
 \quad\Longleftrightarrow\quad a-b\in\mathbb Z_{\geq0}.
\end{equation}
In the exceptional case $a=b+N$, Kummer's identity gives
\begin{equation}\label{p03:ga19:polynomialcase}
 f(z)=e^z P_{N,b}(z),\qquad
 P_{N,b}(z)=M(-N;b;-z)
           =\sum_{k=0}^N\binom Nk\frac{z^k}{(b)_k}.
\end{equation}
Its coefficients are positive.  Then $z_*$ is algebraic and
$\mathcal V_{d,q}$ is a nonzero algebraic multiple of $e^{-z_*}$,
so is transcendental by Lindemann--Weierstrass.

For clarity, the nonexceptional input in \eqref{p03:ga19:old} uses
Beukers' specialization theorem \cite[Corollary~1.4]{p03:ga19:Beukers}.
The rational-parameter functions $f,f'$ are $E$-functions with
a rational differential system singular only at zero.  They are
linearly independent over $\overline{\mathbb Q}(z)$ unless
$f=e^zP$ as in \eqref{p03:ga19:polynomialcase}.  Indeed a rational
logarithmic derivative of an entire exponential-type solution gives
$f=e^{\lambda z}P$; substitution in Kummer's equation gives either
$a=-\deg P$ or $\lambda=1$, $a-b=\deg P$.  The first alternative
is impossible for $a>0$.  Beukers' theorem then excludes the
boundary relation at a nonzero algebraic point.  We also use the
classical Lindemann--Weierstrass theorem in the form that exponentials
of distinct algebraic numbers are linearly independent over
$\overline{\mathbb Q}$.

\subsection{Target I: independent optimal values in one fixed dimension}

The new ingredient is an irreducibility statement for the actual
stopping-boundary polynomial, not an irreducibility assertion for a
Laguerre polynomial substituted without checking its derivative term.

\begin{lemma}[Prime-parameter boundary polynomials]\label{p03:ga19:eisenstein}
Let $b\geq1$ be an integer and let $p\geq b$ be prime.  Set
$N=p-b$, $D=N+1$, and
\begin{equation}\label{p03:ga19:B}
 B_{p,b}(Z)=(b)_N\bigl[ZP_{N,b}'(Z)+(Z-p)P_{N,b}(Z)\bigr].
\end{equation}
Then $B_{p,b}$ is a monic integral polynomial of degree $D$,
Eisenstein at $p$.  In particular
\begin{equation}\label{p03:ga19:degree}
       [\mathbb Q(z_{2b,2p}):\mathbb Q]=p-b+1.
\end{equation}
In dimension two it has the explicit form
\begin{equation}\label{p03:ga19:planarpoly}
 B_{p,1}(Z)=\sum_{k=0}^p
       (2k-p)\frac{p!(p-1)!}{(k!)^2(p-k)!}Z^k.
\end{equation}
\end{lemma}
\begin{proof}
The scaled polynomial $(b)_NP_{N,b}$ is monic and integral,
because its coefficient of $Z^k$ is
$\binom Nk(b+k)_{N-k}$.  Thus \eqref{p03:ga19:B} is monic and
integral.  Its constant term is
\[
       -p(b)_N=-p\frac{(p-1)!}{(b-1)!},
\]
which has $p$-adic valuation exactly one.  For $1\leq k\leq N$,
the coefficient of $Z^k$ equals
\begin{align}
 &(b+k)_{N-k}
   \left[(k-p)\binom Nk+(b+k-1)\binom N{k-1}\right]\\
 &\hspace{12mm}=
 p\,\frac{2k-N-1}{N-k+1}\binom Nk(b+k)_{N-k}.       \label{p03:ga19:coeff}
\end{align}
The identity follows by using
$\binom N{k-1}/\binom Nk=k/(N-k+1)$ and $p=b+N$.
The denominator $N-k+1$ is a positive integer below $p$.
Consequently each integral coefficient in \eqref{p03:ga19:coeff}
is divisible by $p$.  All nonleading coefficients are divisible
by $p$ and the constant is not divisible by $p^2$; Eisenstein's
criterion proves irreducibility.  The case $N=0$ is included:
then $B_{p,p}(Z)=Z-p$.  The stopping equation
\eqref{p03:ga19:setup} is precisely $B_{p,b}(z_*)=0$, which proves
\eqref{p03:ga19:degree}.  Substituting $b=1$, $N=p-1$ in the
coefficient formula, including the leading and constant terms,
gives \eqref{p03:ga19:planarpoly}.
\end{proof}

The first planar polynomials, in increasing degree, are
\begin{align*}
 B_{2,1}(Z)&=Z^2-2,\\
 B_{3,1}(Z)&=Z^3+3Z^2-6Z-6,\\
 B_{5,1}(Z)&=Z^5+15Z^4+40Z^3-120Z^2-360Z-120.
\end{align*}
Only their unique positive roots are stopping boundaries.  Their
other roots nevertheless matter in the field argument below.

\begin{lemma}[Separation from earlier splitting fields]\label{p03:ga19:field}
Let $\alpha_1,\ldots,\alpha_n$ be algebraic numbers of degrees
$D_1,\ldots,D_n$, respectively.  Suppose that for each $j$ there
is a prime $\ell_j\mid D_j$ exceeding every preceding degree
$D_i$ with $i<j$, and $D_1>1$.  Then
$1,\alpha_1,\ldots,\alpha_n$ are linearly independent over
$\mathbb Q$.
\end{lemma}
\begin{proof}
The compositum $K$ of the splitting fields of the minimal
polynomials of $\alpha_1,\ldots,\alpha_{j-1}$ is Galois over
$\mathbb Q$.  Its Galois group injects into the product of the
corresponding symmetric groups.  Thus $[K:\mathbb Q]$ divides
$\prod_{i<j}D_i!$, and is not divisible by $\ell_j$.
If $\alpha_j$ belonged to $K$, its degree $D_j$ would divide
$[K:\mathbb Q]$, a contradiction.  In a proposed rational linear
relation, take the largest index with a nonzero coefficient.  It
would put that $\alpha_j$ in the preceding compositum.  For
$j=1$ it would make $\alpha_1$ rational, also impossible.
\end{proof}

\begin{theorem}[One planar bridge supplies infinitely many independent optima]
\label{p03:ga19:planar}
The collection
\begin{equation}\label{p03:ga19:planarvalues}
          \{e\}\ \cup\
          \{\mathcal V_{2,2p}:p\text{ prime}\}
\end{equation}
is algebraically independent over $\overline{\mathbb Q}$.
Equivalently one can replace $e$ by $\mathcal V_{2,2}=2/e$.
Thus all the constants come from different payoff powers for a
single two-dimensional Brownian bridge.
\end{theorem}
\begin{proof}
Write $z_p=z_{2,2p}$.  Its degree is $p$ by
Lemma~\ref{p03:ga19:eisenstein}.  In increasing prime order these
degrees meet Lemma~\ref{p03:ga19:field}, so every finite subset of
$1,z_2,z_3,z_5,\ldots$ is rationally linearly independent.
The stopping formula gives
\[
 \mathcal V_{2,2p}=a_p e^{-z_p},\qquad
 a_p=\frac{(2z_p)^p}{P_{p-1,1}(z_p)}
                       \in\overline{\mathbb Q}_{>0}.
\]
Distinct monomials in $e,e^{-z_{p_1}},\ldots,e^{-z_{p_n}}$
have distinct algebraic exponents by the stated linear independence.
Lindemann--Weierstrass excludes every nonzero polynomial relation
between them.  Multiplication by the nonzero algebraic $a_p$
preserves independence.  Inverting and scaling the first variable
preserves it as well, proving the equivalent formulation.
\end{proof}

The fixed dimension can be any positive even integer, at the cost
of selecting a subsequence of primes.  No assertion about prime
pairs at a prescribed difference is needed.

\begin{theorem}[Independent optima in every fixed even dimension]
\label{p03:ga19:evenfixed}
For each integer $b\geq1$ there is a strictly increasing sequence
of primes $p_j>b$ such that
\begin{equation}\label{p03:ga19:evenfixedset}
       \{e\}\cup\{\mathcal V_{2b,2p_j}:j\geq1\}
\end{equation}
is algebraically independent over $\overline{\mathbb Q}$.
For $b=1$ all primes can be used.  For $b>1$, a recursive
selection is: after choosing the earlier $p_i$, choose a prime
$\ell_j>\max(b,p_1,\ldots,p_{j-1})$ and then a larger prime
\begin{equation}\label{p03:ga19:dirichletchoice}
                     p_j\equiv b-1\pmod{\ell_j}.
\end{equation}
\end{theorem}
\begin{proof}
For $b>1$, the residue $b-1$ is nonzero modulo $\ell_j$.
Dirichlet's theorem on primes in arithmetic progressions
\cite{p03:ga19:Dirichlet} supplies arbitrarily large primes in
\eqref{p03:ga19:dirichletchoice}.  The degree
$D_j=p_j-b+1$ is divisible by $\ell_j$, and $\ell_j$ exceeds
all earlier $D_i$.  Lemmas~\ref{p03:ga19:eisenstein} and
\ref{p03:ga19:field} show that $1,z_{2b,2p_1},z_{2b,2p_2},\ldots$
are rationally linearly independent.  Each objective is again a
nonzero algebraic multiple of the exponential of its negative
boundary.  The same Lindemann--Weierstrass argument proves the
claim.  The case $b=1$ is Theorem~\ref{p03:ga19:planar}.
\end{proof}

The two theorems are counted as one fixed-dimension family.  They
do not prove independence of all even-payoff optima in every
dimension, nor of the nonexceptional optima at transcendental
boundaries.  The distinction from the prior Pell construction is
that dimension is held fixed, polynomial degrees grow, and the
field separation is obtained from their prime factors.

\subsection{Target II: general rational payoffs and finite multiple zeros}

The following analytic observation is unconditional.  It concerns
complex critical points, although the stopping boundary itself is
positive real.  Controlling all complex critical points is needed
to justify the arithmetic reduction; positivity on the real axis
alone would not suffice.

\begin{lemma}[Escape of the complex critical values]\label{p03:ga19:escape}
Let $a,b>0$, $f=M(a;b;z)$, and suppose
$b-a\notin\{0,-1,-2,\ldots\}$.  At zeros $z$ of
$D(z)=zf'(z)-af(z)$ tending to infinity,
\begin{equation}\label{p03:ga19:criticalasymptotic}
 |f(z)|=\frac{\Gamma(b)}{|\Gamma(b-a)|}|z|^{-a}
                   \bigl(1+O(|z|^{-1})\bigr).
\end{equation}
In particular the magnitude of any branch of the corresponding
critical value $(2z)^a/f(z)$ satisfies
\begin{equation}\label{p03:ga19:valueescape}
 \left|\frac{(2z)^a}{f(z)}\right|
 =\frac{2^a|\Gamma(b-a)|}{\Gamma(b)}|z|^{2a}
                   \bigl(1+O(|z|^{-1})\bigr)\longrightarrow\infty.
\end{equation}
\end{lemma}
\begin{proof}
We use the standard two-term Kummer expansion with separate
remainders, as in \cite[13.7.2 and \S13.7(ii)]{p03:ga19:DLMF}.
It is uniform on either closed half-plane: on the upper half-plane
take $0\leq\arg z\leq\pi$ and the plus sign, and on the lower
one take $-\pi\leq\arg z\leq0$ and the minus sign.  These are
closed subsectors of the respective validity sectors.  Put
\[
 A=\frac{\Gamma(b)}{\Gamma(a)},\qquad
 B_\pm=\frac{\Gamma(b)e^{\pm\pi i a}}{\Gamma(b-a)}.
\]
Both are nonzero, and the expansion reads
\begin{equation}\label{p03:ga19:twoexp}
 f(z)=A e^z z^{a-b}\bigl(1+O(|z|^{-1})\bigr)
       +B_\pm z^{-a}\bigl(1+O(|z|^{-1})\bigr).
\end{equation}
The two error terms are bounded relative to their own displayed
terms; this statement can also be obtained from the exact
connection formula with $U$ and its two standard remainder
estimates.  It is valid when one term is exponentially smaller
than the other and does not divide by their possibly cancelling
sum.

To avoid differentiating an unqualified asymptotic remainder,
apply the same expansion to the exact identity
$f'=(a/b)M(a+1;b+1;z)$.  The exponential leading coefficient
remains $A$, while the algebraic one is $-aB_\pm$.  Therefore
\begin{equation}\label{p03:ga19:Dexp}
 D(z)=A e^z z^{a-b+1}\bigl(1+O(|z|^{-1})\bigr)
       -2aB_\pm z^{-a}\bigl(1+O(|z|^{-1})\bigr).
\end{equation}
At a sufficiently large zero of $D$, the factors $1+O(|z|^{-1})$
are nonzero.  Equation \eqref{p03:ga19:Dexp} then gives
\[
 A e^z z^{a-b}
        =2aB_\pm z^{-a-1}\bigl(1+O(|z|^{-1})\bigr).
\]
Substitution into \eqref{p03:ga19:twoexp} yields
$f(z)=B_\pm z^{-a}(1+O(|z|^{-1}))$.  Since $a$ is real,
$|z^{-a}|=|z|^{-a}$ for either chosen branch, and
$|B_\pm|=\Gamma(b)/|\Gamma(b-a)|$.  This proves
\eqref{p03:ga19:criticalasymptotic} and \eqref{p03:ga19:valueescape}.
\end{proof}

\begin{proposition}[A finite-multiple-zero family]\label{p03:ga19:finite}
Let $a=r/s>0$ be rational in lowest terms, with $r,s$ positive
integers, let $b>0$, and let $c\in\mathbb C^*$.  Then
\begin{equation}\label{p03:ga19:G}
                    G_c(z)=M(a;b;z)^s-cz^r
\end{equation}
has only finitely many multiple zeros in the complex plane.
If $b\in\mathbb Q_{>0}$ and $c\in\overline{\mathbb Q}^*$, then
$G_c$ is an $E$-function.
\end{proposition}
\begin{proof}
The function is not identically zero since $G_c(0)=1$.
At a multiple zero $z\ne0$, $f(z)^s=cz^r\ne0$ and
$sf(z)^{s-1}f'(z)=crz^{r-1}$.  Dividing these nonzero
quantities proves $zf'(z)/f(z)=r/s=a$.  If $b-a$ is not a
nonpositive integer, Lemma~\ref{p03:ga19:escape} gives at large
such points
\[
 \left|\frac{f(z)^s}{z^r}\right|
  =\left(\frac{\Gamma(b)}{|\Gamma(b-a)|}\right)^s
      |z|^{-2r}\bigl(1+O(|z|^{-1})\bigr)\longrightarrow0.
\]
This cannot equal the fixed positive number $|c|$.
Consequently all multiple zeros lie in a bounded set.  They
are finite because the nonzero entire function $G_c$ has no
finite accumulation of zeros.

In the exceptional case, $f=e^zP$ with $P$ polynomial as in
\eqref{p03:ga19:polynomialcase}.  Then
$D=e^z[zP'+(z-a)P]$, and the bracket is a nonzero polynomial
of degree $\deg P+1$.  Every multiple zero of $G_c$ is one of
its finitely many zeros.  Finally rational-parameter Kummer
functions are $E$-functions, and this class is closed under
products and addition of algebraic-coefficient polynomials.
This proves the last assertion when $b\in\mathbb Q_{>0}$
and $c\in\overline{\mathbb Q}^*$.
\end{proof}

The arithmetic step needed to exploit this analytic result is
exactly identifiable.  Fischler and Rivoal's published
Conjecture~1.4 \cite{p03:ga19:FR} asserts a squarefree-factorization
property for $E$-functions: a nonzero $E$-function is a finite
product $\prod_j g_j^j$, where each $g_j$ is an $E$-function
with simple zeros and different $g_j$ have disjoint zero sets.
We call this statement \textnormal{(SF)}.  It is an external
conjecture, not an assumption silently included in a theorem
claimed to be unconditional.

\begin{lemma}[What \textnormal{(SF)} implies for finitely many multiple zeros]
\label{p03:ga19:SFlemma}
Assume \textnormal{(SF)}.  If a nonzero $E$-function has only
finitely many multiple zeros, all those multiple zeros are algebraic.
\end{lemma}
\begin{proof}
First, any $E$-function $g$ with finitely many zeros is a
polynomial over $\overline{\mathbb Q}$ times $e^{\lambda z}$,
with $\lambda\in\overline{\mathbb Q}$.  Indeed the growth
bound for an $E$-function and Hadamard factorization give
$g=P e^{\lambda z}$, initially with complex coefficients and
complex $\lambda$, absorbing a constant factor into $P$.
Take a nonzero differential equation
$\sum_j p_j(z)g^{(j)}(z)=0$ with algebraic polynomial $p_j$.
If $m$ is the greatest degree of the $p_j$ and $L$ is the
leading coefficient of $P$, the coefficient of
$z^{m+\deg P}$ after division by $e^{\lambda z}$ is
\[
         L\sum_{\deg p_j=m}\operatorname{lc}(p_j)\lambda^j.
\]
The polynomial in $\lambda$ is nonzero, so $\lambda$ is
algebraic.  All Taylor coefficients of $g e^{-\lambda z}=P$
are then algebraic, proving the assertion about $P$ and its zeros.

Now factor the original function according to \textnormal{(SF)}.
For $j\geq2$, each zero of $g_j$ is a multiple zero of the
original function, so $g_j$ has finitely many zeros.  The preceding
argument makes all its zeros algebraic.  Every multiple zero
belongs to one of these factors, proving the lemma.
\end{proof}

\begin{theorem}[Conditional classification of every rational-parameter optimum]
\label{p03:ga19:allconditional}
Assume Fischler--Rivoal's squarefree-factorization conjecture
\textnormal{(SF)}.  Then
\begin{equation}\label{p03:ga19:allvalue}
          \mathcal V_{d,q}\text{ is transcendental}
                   \qquad(d,q\in\mathbb Q_{>0}).
\end{equation}
More precisely, outside $q-d\in2\mathbb Z_{\geq0}$, an
algebraic optimum $v$ would give an explicit counterexample to
\textnormal{(SF)}:
\begin{equation}\label{p03:ga19:counterexample}
 M(q/2;d/2;z)^s-\frac{2^r}{v^s}z^r,
                  \qquad q/2=r/s\text{ in lowest terms}.
\end{equation}
This $E$-function would have finitely many multiple zeros and
a transcendental zero of multiplicity exactly two at $z_*$.
\end{theorem}
\begin{proof}
The exceptional family already has an unconditional transcendence
proof, recalled after \eqref{p03:ga19:polynomialcase}.  In all other
cases \eqref{p03:ga19:old} gives a transcendental $z_*$.  If
$v=\mathcal V_{d,q}$ were algebraic, raising
\eqref{p03:ga19:setup} to the $s$th power shows that
\eqref{p03:ga19:counterexample} vanishes at $z_*$.  Differentiating
and using $z_*f'(z_*)=af(z_*)$ shows that its derivative also
vanishes there.

This zero has order exactly two.  On the positive axis put
$h=f/z^a$.  Locally the function in
\eqref{p03:ga19:counterexample} is
$z^r[h(z)^s-h(z_*)^s]$; its second derivative at $z_*$ is
$s z_*^r h(z_*)^{s-1}h''(z_*)>0$ by
\eqref{p03:ga19:hsecond}.  Proposition~\ref{p03:ga19:finite} gives
only finitely many multiple zeros.  Lemma~\ref{p03:ga19:SFlemma}
would make each algebraic under \textnormal{(SF)}, contradicting
the transcendence of $z_*$.  This proves both the implication
and the explicit counterexample statement.
\end{proof}

The analytic finite-zero statement is unconditional; the use of
\textnormal{(SF)} is not.  The known partial factorization theorem
in \cite[Theorem~1.5]{p03:ga19:FR} assumes that \emph{all} zeros
have the same multiplicity (and additional hypotheses).  Our
functions are only known to have finitely many multiple zeros,
so that theorem cannot replace the conjecture.  Nor does the
exponential-polynomial result under Schanuel's conjecture in
that paper automatically apply to general Kummer functions.

\subsection{Target III: the linear Brownian-excursion optimum}

This target is singled out because it is a concrete established
stopping constant, not merely an arbitrarily selected parameter.
For $d=3$, $q=1$, set
\[
 f(z)=M(1/2;3/2;z)=\int_0^1e^{zt^2}\,dt.
\]
Integration of $(t e^{zt^2})'$ gives
$zf'(z)=(e^z-f(z))/2$.  Thus its unique positive stopping
boundary satisfies
\begin{equation}\label{p03:ga19:excursion}
 2f(z_*)=e^{z_*},\qquad C=\sqrt{2z_*},\qquad
              V:=\mathcal V_{3,1}=2C e^{-C^2/2}.
\end{equation}
Its boundary is already known to be transcendental.  Theorem
\ref{p03:ga19:allconditional} proves the following precise new
conditional arithmetic statement:
\begin{equation}\label{p03:ga19:excursionconditional}
 \textnormal{(SF)}\quad\Longrightarrow\quad
                   V\text{ is transcendental}.
\end{equation}
Indeed if $V=v$ were algebraic, then
\begin{equation}\label{p03:ga19:excursionG}
       M(1/2;3/2;z)^2-\frac{2}{v^2}z
\end{equation}
would be an $E$-function with finitely many multiple zeros,
one of them the transcendental double zero $z_*$.  The critical
value escape in this particular case is
\[
   \left|\frac{\sqrt{2z}}{M(1/2;3/2;z)}\right|
       \sim\frac{2\sqrt2}{\sqrt\pi}|z|
       \quad\text{at unbounded complex contact points}.
\]
The general proof therefore supplies the missing global control
that a real-axis tangency calculation alone could not provide.

Unconditionally, the arithmetic of $V$ is still undecided by
this work.  The two equations
$2M(1/2;3/2;z_*)=e^{z_*}$ and $8z_*=v^2e^{2z_*}$ do not
contradict Lindemann--Weierstrass, because $z_*$ is transcendental.
The same issue prevents using the algebraic-point specialization
theorem for $E$-functions.  The finite-multiple-zero construction
locates a sufficient published conjecture rather than treating
those algebraic-argument theorems as statements at arbitrary
transcendental points.

\paragraph{Scope and result accounting.}
One unconditional arithmetic family is proved: infinitely many
jointly independent actual stopping optima in every fixed even
dimension, including a complete prime-indexed planar family.
One unconditional analytic reduction is proved: the associated
$E$-functions have only finitely many multiple zeros.  This gives
one conditional arithmetic family covering all positive rational
dimensions and payoff exponents.  The linear excursion optimum
is an instance of that conditional family and is not counted
again as an unconditional closure.  The three targeted questions
are therefore distinguished from the number of proved families.
None is asserted to close a historically named arithmetic
conjecture, and the published squarefree-factorization conjecture
itself remains open here.


% Integrated source: bessel_results.tex
\section{All integer-shape extraction rates and the multiplicity of their contacts}
\label{p03:be19:section}

The preceding paper proved transcendence of an intermediate critical
rate for extracting a gamma factor of shape one.  We now treat every
positive integer shape.  The obstacle is not merely a larger polynomial
calculation: after repeated differentiation, a rational coefficient
row can vanish, so a naive determinant argument can lose the very
special-function relation needed for transcendence.  A third-order
intertwining equation excludes simultaneous loss of both contact rows.
It also proves that every interior zero of a boundary remainder in
this integer-shape family is exactly double.  No generic-parameter
qualification or omitted half-integer Bessel exception is needed.

The endpoint constants $A(1)$ and $A(3/2)$ remain unresolved at the
level of transcendence.  The argument below uses finite differential
operators and keeps the Bessel order rational; it does not remove
the nonrational-parameter obstruction recorded in the preceding paper.

\subsection{Notation and the inherited probability diagram}

Fix $b>1$ and $\kappa>1/2$.  We use the first excess-power McKay law
with Laplace transform
\begin{equation}\label{p03:be19:LT}
 L(s)=\frac{b+s}{b}
       \left(\frac{b^2-1}{(b+s)^2-1}\right)^\kappa.
\end{equation}
Write $M=2\kappa-1$.  A gamma factor has shape $h$ and rate
$\rho=b-c$, with $0\leq c\leq1$.  The previously proved maximal
shape for an arbitrary probability remainder is $H_\kappa(c)$.
We use the following established analytic facts
\cite{p03:root3:sequels}:
\begin{enumerate}
\item $H_\kappa$ is continuous and strictly decreasing on $[0,1]$,
with $H_\kappa(0)=M$ and
$H_\kappa(1)=\kappa-A(\kappa)$.
\item $0<A(\kappa)<1/(4\kappa)$, with the stronger lower bound
$A(\kappa)>\max(0,1-\kappa)$.
\item At an interior maximal shape the remainder has an analytic
density on $(0,\infty)$, strictly positive near zero and infinity,
and has a nonempty finite set of interior zeros, each of even order.
\end{enumerate}
The earlier arithmetic result additionally says that $A(\kappa)$
is irrational for rational $\kappa>1/2$
\cite{p03:root3:previousI,p03:root3:previousII}.  These are inputs, not new conclusions of
this section.

Let $D=d/dx$ and set
\begin{equation}\label{p03:be19:base}
 t=2\kappa-2,\quad
 j(x)=x^{\kappa-1/2}I_{\kappa-1/2}(x),\quad
 f(x)=j'(x)=x^{\kappa-1/2}I_{\kappa-3/2}(x).
\end{equation}
The derivative identity follows from the standard Bessel recurrence
\cite[10.29.2]{p03:be19:DLMF}.  Direct substitution gives
\begin{align}
 xj''-t j'-xj&=0,                                      \label{p03:be19:jode}\\
 f''-\frac{t}{x}f'+\left(\frac{t}{x^2}-1\right)f&=0.  \label{p03:be19:fode}
\end{align}
For an integer $m\geq1$ define
\begin{equation}\label{p03:be19:gm}
 v_m=(D-c)^m j,\qquad g_m=(D-c)^m f=v_m'.
\end{equation}
When $m<M$, multiplication of the shifted Laplace transform by
$w^m$ is ordinary differentiation without an initial-value term:
$f(x)$ has order $x^{M-1}$ at zero and its first $m-1$
derivatives vanish there.  Up to a positive normalizer, the actual
shape-$m$ remainder density is
\begin{equation}\label{p03:be19:density}
                       e^{-b x}g_m(x).
\end{equation}
Indeed the intermediate inverse is $D^m(e^{-cx}f)=e^{-cx}g_m$,
and its probability tilt is $e^{-(b-c)x}$.  In particular,
nonnegativity is exactly the condition $g_m\geq0$.

\subsection{The third-order equation removes every interior degeneracy}

\begin{lemma}[Intertwining and nonvanishing of the two contact rows]
\label{p03:be19:intertwine}
Let $m\geq1$, let $t,c$ be arbitrary complex constants with
$c^2\ne1$, and work locally away from $x=0$.  For every solution
$j$ of \eqref{p03:be19:jode}, the function $v=(D-c)^m j$ satisfies
\begin{equation}\label{p03:be19:third}
 x v'''+(m+1-t-cx)v''+(ct-x)v'+(cx-m-1)v=0.
\end{equation}
The map $j\mapsto v$ is injective on the two-dimensional solution
space.  Consequently it is impossible that both transformed basis
functions $g_i=((D-c)^m j_i)'$ satisfy
\begin{equation}\label{p03:be19:forbidden}
             g_i(x_0)=g_i'(x_0)=0\qquad(i=1,2)
\end{equation}
at the same nonzero point $x_0$.
\end{lemma}
\begin{proof}
Put $L=x(D^2-1)-tD$ and $T=D-c$.  The commutator is
$[T,L]=D^2-1$, and $D^2-1$ commutes with $T$.  Hence
\[
 T^{m+1}L=LT^{m+1}+(m+1)(D^2-1)T^m.
\]
Applying this to $j$ yields
$LTv+(m+1)(D^2-1)v=0$, which expands to
\eqref{p03:be19:third}.

If $T^m j=0$, then $j=e^{cx}P(x)$ for a polynomial $P$ of
degree at most $m-1$.  Substitution in $Lj=0$ gives
\[
 xP''+(2cx-t)P'+\{(c^2-1)x-tc\}P=0.
\]
For a nonzero polynomial $P$, its highest-degree term multiplied
by $(c^2-1)x$ cannot cancel, because $c^2\ne1$.  Thus $P=0$,
proving injectivity.

Equation \eqref{p03:be19:third} is an ordinary third-order equation
at $x_0\ne0$.  Its solutions with $v'(x_0)=v''(x_0)=0$
form a one-dimensional space, determined by the remaining initial
value $v(x_0)$.  Two independent images $v_1,v_2$ cannot both
belong to that space.  This proves the last assertion, since
$g_i=v_i'$.
\end{proof}

To turn the preceding argument into a finite algebraic contact
relation, define rational functions $U_m,V_m$ by
\[
 (D-c)^m f=U_m f+V_m f'
\]
for every solution of \eqref{p03:be19:fode}.  They are obtained without
any special-function evaluation from $U_0=1$, $V_0=0$ and
\begin{align}
 U_{m+1}&=U_m'-cU_m+(1-t/x^2)V_m,                       \label{p03:be19:Urec}\\
 V_{m+1}&=U_m+V_m'+(t/x-c)V_m.                         \label{p03:be19:Vrec}
\end{align}
Put
\begin{equation}\label{p03:be19:determinant}
 \Delta_m=\det\begin{pmatrix}
 U_m&V_m\\
 U_m'+(1-t/x^2)V_m&U_m+V_m'+tV_m/x
 \end{pmatrix},\qquad
 P_m(t;x,c)=x^{2m}\Delta_m(t;x,c).
\end{equation}

\begin{lemma}[The exact contact polynomial and its fibers]
\label{p03:be19:polynomial}
For every $m\geq1$, $P_m$ belongs to $\mathbb Q[t,x,c]$.
It has degree at most $2m$ in each of $x,c$, with leading
coefficients
\begin{equation}\label{p03:be19:leading}
 [x^{2m}]P_m=(c^2-1)^m,
 \qquad [c^{2m}]P_m=x^{2m}.
\end{equation}
For fixed $t,c$ with $c^2\ne1$, its polynomial in $x$ is
nonzero.  For fixed $t,x\ne0$, its polynomial in $c$ is nonzero.
At every nonzero $x$ and every $c^2\ne1$, the matrix in
\eqref{p03:be19:determinant} is not the zero matrix.
\end{lemma}
\begin{proof}
Induction in \eqref{p03:be19:Urec}--\eqref{p03:be19:Vrec} gives
\[
 U_m=x^{-m}u_m(t;x,c),\qquad
 V_m=x^{1-m}w_m(t;x,c),
\]
where $u_m,w_m$ are polynomials of degree at most $m,m-1$
in $x$.  The determinant therefore has denominator dividing
$x^{2m}$, with numerator of degree at most $2m$ in $x$.
The binomial expansion of $(D-c)^m$ shows that the first row
has leading $c$-term $(-c)^m(1,0)$ and the second row has
leading term $(-c)^m(0,1)$.  The determinant is consequently
monic of degree $2m$ in $c$.  This proves the second identity
in \eqref{p03:be19:leading} and the degree bound in $c$.

For the first identity, take the constant term at infinity in
the rational coefficient matrix.  Equation \eqref{p03:be19:fode}
there reduces to $f''=f$.  The transformation $D-c$ on the
state vector $(f,f')$ is then
\[
 K=\begin{pmatrix}-c&1\\1&-c\end{pmatrix}.
\]
The constant term of the transformed state matrix is $K^m$,
with determinant $(c^2-1)^m$.  This is the coefficient of
$x^{2m}$ after multiplication by $x^{2m}$.

Finally differentiation identifies the solution spaces in
\eqref{p03:be19:jode} and \eqref{p03:be19:fode}.  An explicit inverse
is $j=f'-tf/x$, for which $j'=f$ by \eqref{p03:be19:fode}.
If the displayed matrix vanished at a nonzero point, every
basis pair would satisfy \eqref{p03:be19:forbidden}, contradicting
Lemma~\ref{p03:be19:intertwine}.
\end{proof}

The last assertion is stronger than saying that the determinant is
not identically zero.  For example, when $m=2$, the first row can
vanish at particular positive parameter values.  It is the
simultaneous vanishing of both rows that would invalidate the
arithmetic argument, and the third-order equation rules it out for
all orders.  No search through finitely many small $m$ proves that
assertion; the intertwining identity does.

\subsection{Arithmetic of the entire multiple-contact locus}

We use the classical rational-order Bessel consequence of Beukers'
linear-specialization theorem \cite[Corollary~1.4]{p03:ga19:Beukers}:
for rational $\nu>-1$ and positive algebraic $x$,
\begin{equation}\label{p03:be19:ratio}
                         I_{\nu+1}(x)/I_\nu(x)
                         \quad\hbox{is transcendental}.
\end{equation}
For completeness, the normalized even entire function
\[
 H_\nu(z)=\sum_{n\geq0}\frac{(z^2/4)^n}{n!(\nu+1)_n}
\]
has even exponential Taylor coefficients
$(1/2)_n/(\nu+1)_n$ and odd coefficients zero.  These coefficients
and their common denominators have exponential bounds for rational
$\nu$; the bound follows by counting the two rational arithmetic
progressions modulo prime powers, with the finitely many parameter-denominator primes treated separately.  Thus $H_\nu,H_\nu'$
are $E$-functions.  They solve a system whose only finite possible
singularity is zero.  Since $H_\nu$ is entire of exponential type, a rational logarithmic
derivative would make $H_\nu=e^{az}Q(z)$ for a polynomial $Q$;
evenness forces $a=0$,
contrary to its infinitely many nonzero Taylor coefficients.
The pair is therefore linearly independent over
$\overline{\mathbb Q}(z)$.  Beukers' theorem at a nonzero
algebraic point implies that $H_\nu'/H_\nu$ is transcendental,
and this quotient is \eqref{p03:be19:ratio}.  This input is classical
and was already used at shape one.

\begin{theorem}[Every multiple contact has transcendental coordinates]
\label{p03:be19:contact-arithmetic}
Let $\kappa>1/2$ be rational, $m\geq1$ an integer, and
$0<c<1$.  If the function $g_m$ of \eqref{p03:be19:gm} has a
multiple zero at $x>0$, then both $c$ and $x$ are transcendental.
They obey the explicit polynomial relation
\begin{equation}\label{p03:be19:contact-curve}
                            P_m(2\kappa-2;x,c)=0.
\end{equation}
Each coordinate is algebraic of degree at most $2m$ over the
field generated over $\overline{\mathbb Q}$ by the other.
In particular their joint transcendence degree is exactly one.
At an algebraic rate coordinate $c\in(0,1)$ all positive zeros
of $g_m$ are simple.
\end{theorem}
\begin{proof}
The two equations $g_m(x)=g_m'(x)=0$ are the matrix equation
in \eqref{p03:be19:determinant} applied to $(f(x),f'(x))$.
Since $f(x)>0$ on the positive axis, its determinant vanishes.
This proves \eqref{p03:be19:contact-curve}.

Suppose first that $c$ were algebraic.  The first identity in
\eqref{p03:be19:leading} makes \eqref{p03:be19:contact-curve} a nonzero
polynomial equation with algebraic coefficients for $x$.
Thus $x$ is algebraic.  All entries of the contact matrix are
then algebraic.  By Lemma~\ref{p03:be19:polynomial} at least one
row is nonzero.  Applied to $(f(x),f'(x))$, this row would force
$f'(x)/f(x)$ algebraic; a row with zero second coefficient
would instead contradict $f(x)>0$ immediately.  But
\[
 \frac{f'(x)}{f(x)}=\frac{2\kappa-2}{x}
       +\frac{I_{\kappa-1/2}(x)}{I_{\kappa-3/2}(x)}
\]
is transcendental by \eqref{p03:be19:ratio}.  This is a contradiction.
If $x$ were algebraic to begin with, the second identity in
\eqref{p03:be19:leading} would first force $c$ algebraic, producing
the same contradiction.  Thus both coordinates are transcendental.

The two nonzero leading coefficients also give the two degree
bounds over the corresponding one-coordinate fields.  Since both
coordinates are transcendental and are algebraically dependent,
their joint transcendence degree is exactly one.  The assertion at
an algebraic $c$ is the contrapositive of the first conclusion.
\end{proof}

\subsection{All boundary contacts are exactly double}

\begin{theorem}[Multiplicity and contact bound at an integer-shape boundary]
\label{p03:be19:multiplicity}
Let $\kappa>(m+1)/2$, $m\geq1$, and $0<c<1$.  Suppose the
shape-$m$ remainder density is nonnegative and has a positive zero.
Then every positive zero of that density is exactly double and
lies in
\begin{equation}\label{p03:be19:bound}
                            0<x<\frac{m+1}{c}.
\end{equation}
There are at most $2m$ distinct positive zeros.
These conclusions do not require rational $\kappa$.
\end{theorem}
\begin{proof}
Nonnegativity of the density is $g_m=v_m'\geq0$.
Near zero, the Bessel series gives
\[
 j(x)\sim Cx^{2\kappa-1},\qquad
 v_m(x)\sim C(2\kappa-1)(2\kappa-2)\cdots(2\kappa-m)
                    x^{2\kappa-1-m},\qquad C>0.
\]
All factors and the exponent are positive under the stated strict
inequality.  Hence $v_m(0+)=0$ and $v_m>0$ near zero.
Since $v_m'=g_m\geq0$, it follows that $v_m(x)>0$ for all $x>0$.

At a positive zero of $g_m$, nonnegativity gives $g_m'=0$.
Equation \eqref{p03:be19:third} becomes
\begin{equation}\label{p03:be19:secondcontact}
                   x g_m''=(m+1-cx)v_m.
\end{equation}
At a local minimum the left side is nonnegative, so $cx\leq m+1$.
If equality held, then $g_m''=0$.  Differentiate
\eqref{p03:be19:third}; using $v_m'=v_m''=v_m'''=0$ at that
point gives
\[
                         xg_m'''=-c v_m<0.
\]
The zero would have exact odd order three, contradicting
nonnegativity on both sides.  Thus $cx<m+1$, and
\eqref{p03:be19:secondcontact} gives $g_m''>0$.  This proves
exact multiplicity two and \eqref{p03:be19:bound}.

Each zero satisfies the contact polynomial equation by the same
matrix argument as before; no arithmetic assumption is needed
for this step.  For fixed $t,c$, this is a nonzero polynomial
of degree $2m$ in $x$, so there can be at most $2m$ distinct
positive zeros.
\end{proof}

The multiplicity statement strengthens the inherited assertion
``finitely many zeros of even order.''  It does not assert that
there is only one contact for every integer shape.  Multiple
contacts, if present, all lie in the same algebraic extension of
the field generated by their common rate coordinate.

\subsection{Complete arithmetic classification of the critical integer rate}

\begin{theorem}[The integer-shape phase diagram]
\label{p03:be19:phase}
Fix rational $\kappa>1/2$, an integer $m\geq1$, and restrict
$c$ to $[0,1]$.  Exactly one of the following cases occurs.
\begin{enumerate}
\item If $m>2\kappa-1$, no rate in this interval permits a
shape-$m$ factor.
\item If $m=2\kappa-1$, the only allowable rate coordinate is
$c=0$.  This is an endpoint with an atom in the remainder.
\item If $m<2\kappa-1$ and $m<\kappa-A(\kappa)$, every
$c\in[0,1]$ permits the factor.  The largest rate coordinate
within the prescribed interval is the algebraic endpoint $1$.
\item If $m<2\kappa-1$ and $m>\kappa-A(\kappa)$, there is
one critical coordinate $c_{\kappa,m}\in(0,1)$, and extraction
is possible exactly for $0\leq c\leq c_{\kappa,m}$.
This critical coordinate is transcendental.  Its remainder has
between one and $2m$ positive zeros; every one is exactly double,
transcendental, and smaller than $(m+1)/c_{\kappa,m}$.
The rate and all its contacts together generate a field of
transcendence degree exactly one over $\overline{\mathbb Q}$.
\end{enumerate}
There is no equality case $m=\kappa-A(\kappa)$ at rational
$\kappa$.  In particular, for each integer $m\geq2$, every
rational
\begin{equation}\label{p03:be19:explicit-interval}
                         \frac{m+1}{2}<\kappa\leq m
\end{equation}
belongs to case 4.  This includes integer Bessel parameters
$\kappa=m$ without an exceptional reduction.  For algebraic
$b>1$, the physical critical rate $b-c_{\kappa,m}$ is also
transcendental.
\end{theorem}
\begin{proof}
The cases follow from the inherited strict monotonicity and endpoint
values of $H_\kappa$.  The equality
$m=\kappa-A(\kappa)$ would make $A(\kappa)=\kappa-m$
rational, contradicting its established irrationality.  In case 4,
continuity and strict monotonicity give a unique interior solution
$H_\kappa(c_{\kappa,m})=m$ and the asserted closed feasible
interval.  Its analytic remainder has at least one positive zero
by the inherited boundary theorem.  Since $m<M$, the derivative
representation \eqref{p03:be19:density} applies.  Theorems
\ref{p03:be19:contact-arithmetic} and \ref{p03:be19:multiplicity}
therefore give all the claimed arithmetic, multiplicity, and
counting conclusions.  Every contact is algebraic over the field
of the common rate, so finitely many contacts together do not
increase its transcendence degree.

At $m=M$ and $c=0$, the shifted transform is
$(1-w^{-2})^{-\kappa}$.  Its expansion has a positive constant
term and nonnegative remaining inverse terms, giving the central
endpoint remainder and its atom.  Every positive $c$ has smaller
capacity.  This also explains why the strict inequality $m<M$
was needed to use ordinary differentiation without an origin term.

Finally, \eqref{p03:be19:explicit-interval} ensures $M>m$ and
$\kappa-A(\kappa)<\kappa\leq m$, strictly.  Hence case 4
holds throughout that explicit interval.  Translation by algebraic
$b$ proves the final assertion.
\end{proof}

The two integer values $c=\pm1$ are the spectral endpoints in the
intertwining argument and are deliberately excluded from its local
rank lemma.  This is not an unclassified interior exception.
For example, at integer $\kappa=k\geq1$, equation
\eqref{p03:be19:jode} has an exponential-polynomial solution
$e^xQ_{k-1}(x)$; at $c=1$ the operator $(D-1)^m$ kills it
when $m\geq k$.  For $0<c<1$, the highest-degree argument in
Lemma~\ref{p03:be19:intertwine} excludes such a kernel even when
both Bessel solutions reduce to exponentials times polynomials.
The phase theorem classifies the physical endpoints separately.

\subsection{The endpoint transcendence targets still do not close}

We also reexamined the two specific endpoint targets $A(1)$ and
$A(3/2)$.  Their previous irrationality and the exact integral
reductions remain unchanged.  The new operator identity does not
apply to a hypothetical algebraic-irrational endpoint shape
$h=\kappa-A(\kappa)$: $(D-c)^h$ is then not a finite polynomial
in $D$, and the displayed commutator induction no longer supplies
a third-order equation.  Conversely, forcing an integer shape is
precisely what moves the new unknown to an interior rate.

The previous non-$E$ obstruction is exact: for algebraic irrational
$a$ and positive rational $d$, the Borel transform
${}_2F_1(a-1,1;a+d-1;z)$ violates Galochkin's irrational-parameter
pairing criterion, since its only possible irrational difference
is $-d<0$ \cite[Theorem~1]{p02:be18:Rivoal}.  Thus the old Kummer
function is not an $E$-function under the remaining algebraic
hypothesis.  We have not proved a replacement arithmetic
specialization theorem.  Transcendence of $A(1)$ and $A(3/2)$,
and the general noninteger fixed-shape rate problem, remain
unresolved here.

\paragraph{Accounting.}
The arithmetic theorem is one uniform extension from shape one to
all positive integer shapes, not one new theorem per integer.
The exact double-contact and finite-contact bound resolve a separate
analytic multiplicity question within that same family.  The
endpoint investigations are recorded as unsuccessful.  None of
these internally posed questions is claimed to be a historical
named conjecture or an Annals-level result on the basis of its label.


% Integrated source: mle_results.tex
\section{Arithmetic of finite entropy optimization}
\label{p03:mi3:sec}

The previous exponential-integral arguments classified selected values
at algebraic parameters and certain optimizer locations.  Here the
object is the \emph{optimized information quantity itself}.  Three
targets are finite-channel capacity, finite-source rate--distortion,
and the nonconvex information-bottleneck objective.  A single theorem
handles all three: for rational data, every face-stationary value is
the logarithm of a positive algebraic number, and an algebraic optimizer
exists.  The theorem includes boundary supports and positive-dimensional
critical sets; it assumes neither a unique optimizer nor convexity.

The mechanism is related to established algebraic statistics.
Dukkipati already expressed maximum-entropy and minimum-divergence
estimation through polynomial equations \cite{p03:mi3:Dukkipati}.
Sturmfels, Telen, Vialard, and von Renesse study the toric geometry
of entropy regularization \cite{p03:mi3:toric}.  We do not claim that
exponentiating rational logarithmic stationarity equations is new.
Blakaj and Wolf already prove fixed-law entropy arithmetic
dichotomies: algebraic probabilities give zero or transcendental
entropy in nats, while rational probabilities give rational or
transcendental entropy in an algebraic base
\cite[Theorem 2]{p03:mi3:BlakajWolf}.
Our formulation records the consequences for all signed entropy
combinations, all support faces, and the arithmetic of the attained
values.  Its historical priority has not been established.

All logarithms in this section are natural, and $0\log0=0$.
An algebraic vector means a real vector with algebraic coordinates.
A rational polytope is a bounded polyhedron given by rational linear
equalities and inequalities.

\subsection{All critical values, including nonisolated critical sets}

Let $P\subset\mathbb R^d$ be a nonempty rational polytope.  For
$1\le a\le A$, let
\[
 L_a(x)=(\ell_{a1}(x),\ldots,\ell_{a n_a}(x))
\]
be an affine map with rational coefficients that takes $P$ into a
probability simplex.  Thus each coordinate is nonnegative on $P$ and
$\sum_j\ell_{aj}(x)=1$ there.  Given $c_a\in\mathbb Q$, set
\begin{equation}\label{p03:mi3:eq:objective}
 F(x)=\sum_{a=1}^A c_a H(L_a(x)),\qquad
 H(p)=-\sum_jp_j\log p_j.
\end{equation}
For a nonempty face $Q$ of $P$, a \emph{face-stationary point} is a
point in $\operatorname{relint}Q$ at which the derivative of $F$
vanishes in every direction parallel to $Q$.  This includes every
vertex, where there are no tangent conditions.

\begin{theorem}[Finite arithmetic description of entropy optima]
\label{p03:mi3:thm:entropy}
For the rational data in \eqref{p03:mi3:eq:objective}:
\begin{enumerate}
\item The set of face-stationary values over all faces is finite.
      Each such value is $\log\alpha$ for a positive algebraic
      number $\alpha$, and is attained at an algebraic
      face-stationary point.
\item Both global extrema of $F$ on $P$ are attained at algebraic
      points, and their exponentials are algebraic.
\item From a finite exact rational description of the data, a
      terminating algorithm produces a finite list of algebraic
      points and algebraic positive numbers that represents every
      face-stationary value.  It determines the global extrema by
      comparing those algebraic numbers.
\end{enumerate}
In particular every nonzero extremal value in nats is transcendental.
No polynomial-time complexity bound is asserted.
\end{theorem}

\begin{proof}
Fix a nonempty face $Q$.  An affine function that is nonnegative on
$Q$ and zero at a relative interior point is identically zero on
$Q$: otherwise a sufficiently short line segment through that point
in the direction of a point with positive value would have a negative
value at its opposite end.  Consequently each $\ell_{aj}$ is either
identically zero on $Q$ or strictly positive on
$\operatorname{relint}Q$.  Delete the identically zero coordinates.
The remaining expression for $F$ is smooth on that relative interior.

Choose a rational basis $v_1,\ldots,v_k$ of the linear space parallel
to $Q$.  Since each probability vector has total mass one,
\[
 \sum_j\partial_{v_i}\ell_{aj}=0.
\]
The constant terms in the derivative of $t\log t$ therefore cancel,
and stationarity is exactly
\begin{equation}\label{p03:mi3:eq:logstationary}
 \sum_{a,j}c_a(\partial_{v_i}\ell_{aj})
                  \log\ell_{aj}(x)=0,
                  \qquad 1\le i\le k.
\end{equation}
All coefficients of these logarithms are rational.  Clear their
denominators, obtaining integers $b_{iaj}$.  Because the arguments
are positive, \eqref{p03:mi3:eq:logstationary} is equivalent to the
polynomial equation
\begin{equation}\label{p03:mi3:eq:binomial}
 \prod_{b_{iaj}>0}\ell_{aj}(x)^{b_{iaj}}
       =\prod_{b_{iaj}<0}\ell_{aj}(x)^{-b_{iaj}}.
\end{equation}
An empty product is one.  Thus the set $S_Q$ of face-stationary
points is a semialgebraic set defined over $\mathbb Q$.

We use the standard real-algebraic decomposition theorem: a
semialgebraic set defined over the real algebraic numbers has a
finite effective cell decomposition with algebraic sample points;
each cell is connected by piecewise $C^1$ semialgebraic paths.
These properties follow from cylindrical algebraic decomposition
and semialgebraic path decomposition; see \cite{p03:mi3:BPR}.
Along any such path in $S_Q$, its tangent vector is parallel to
$Q$, and \eqref{p03:mi3:eq:logstationary} makes the derivative of $F$
zero on each smooth piece.  The mean-value theorem and continuity
at the finitely many joins make $F$ constant on the path.  Hence
$F$ is constant on each cell of a finite decomposition of $S_Q$.
In particular it has finitely many values there, and each is
attained at an algebraic sample point $x_*$.  This step does not
assume that the critical equations have finitely many solutions.

Choose a rational point $x_0\in\operatorname{relint}Q$, which exists
because $Q$ is a nonempty rational polytope.  Pair the derivative
at $x_*$ with $x_*-x_0$.  Affineness and stationarity give
\[
 0=-\sum_{a,j}c_a
      \bigl(\ell_{aj}(x_*)-\ell_{aj}(x_0)\bigr)
                                  \log\ell_{aj}(x_*).
\]
Therefore
\begin{equation}\label{p03:mi3:eq:rationalvalue}
 F(x_*)=-\sum_{a,j}c_a\ell_{aj}(x_0)
                                      \log\ell_{aj}(x_*).
\end{equation}
The coefficients on the right are rational, while all logarithm
arguments are positive algebraic numbers.  Taking the exponential
expresses $e^{F(x_*)}$ as a finite product of rational powers of
positive algebraic numbers, so it is algebraic.

There are finitely many faces.  Every global extremizer belongs to
the relative interior of its unique minimal face and is stationary
there.  Continuity of $F$ and compactness of $P$ guarantee existence.
The first two conclusions follow.  For effectiveness, enumerate the
faces, their rational tangent bases, the equations
\eqref{p03:mi3:eq:binomial}, and algebraic sample points of a cell
decomposition of each $S_Q$.  Formula \eqref{p03:mi3:eq:rationalvalue}
computes each positive algebraic exponential exactly, with its
positive real branches.  Algebraic-number order comparison selects
the smallest and largest.  Finally, a nonzero algebraic value of
$F$ would have a transcendental exponential by Hermite--Lindemann,
contradicting what was proved \cite{p03:mi3:Baker}.
\end{proof}

\begin{remark}[What the algorithm outputs]
The output for a value is the symbolic expression $\log\alpha$,
with $\alpha$ specified by a polynomial and an isolating interval;
it is not a claim that a transcendental real has a finite decimal
expansion.  The input consists of exact finite rational data.  This
does not supply a continuous or computable selector from arbitrary
real channels given by approximation oracles.  In particular it
does not contradict the optimizer noncomputability results for
computable-channel representations in \cite{p03:mi3:Boche}.
Likewise, only the rational-face \emph{stationary sets} are being
algebraized.  No semialgebraic description of arbitrary entropy
\emph{level sets} is asserted; the latter can be nowhere
semialgebraic, as proved in \cite{p03:mi3:BlakajWolf}.
\end{remark}

\subsection{An exact test for rationality in bits}

\begin{theorem}[Norm test for rational versus transcendental values]
\label{p03:mi3:thm:bits}
Let $\alpha>0$ be algebraic, let
$d=[\mathbb Q(\alpha):\mathbb Q]$, and put
$N=|N_{\mathbb Q(\alpha)/\mathbb Q}(\alpha)|\in\mathbb Q_{>0}$.
Then $\log_2\alpha$ is rational if and only if there exists an
integer $k$ such that
\begin{equation}\label{p03:mi3:eq:normtest}
                     N=2^k,\qquad \alpha^d=2^k.
\end{equation}
In that case $\log_2\alpha=k/d$.  Otherwise it is transcendental.
Both conditions are decidable by exact rational and algebraic
arithmetic.  Thus every extremum in Theorem~\ref{p03:mi3:thm:entropy}
admits an exact rational-or-transcendental classification in bits.
\end{theorem}
\begin{proof}
Suppose $\log_2\alpha=a/b$ with $a\in\mathbb Z$, $b\ge1$.
Then $\alpha^b=2^a$, and taking absolute norms yields
$N^b=2^{ad}$.  A positive rational number with an integral power
equal to a power of two must be $2^k$ for an integer $k$, by
unique prime factorization in $\mathbb Q$.  Hence $kb=ad$,
so $a/b=k/d$ and $\alpha^d=2^k$.  The converse follows by taking
real logarithms.  If the real number $\log_2\alpha$ were algebraic
irrational, Gelfond--Schneider would make
$2^{\log_2\alpha}=\alpha$ transcendental, a contradiction
\cite{p03:mi3:Baker}.  Factor the numerator and denominator of $N$
to test its first condition, then test the second in the number field.
\end{proof}

\subsection{Three information-theoretic targets}

The optimizations below are established information-theory problems,
not newly named constants.  Channel capacity and rate--distortion
are the two classical settings of Blahut's algorithms
\cite{p03:mi3:Blahut}; the third is the information bottleneck of
Tishby, Pereira, and Bialek \cite{p03:mi3:IB}.

\begin{corollary}[Rational finite-channel capacities, with rational constraints]
\label{p03:mi3:cor:capacity}
Let $W(y\mid x)$ be a finite stochastic matrix with rational entries,
and let $P$ be a nonempty rational polytope of input distributions.
The maximum
\[
                         C_P(W)=\max_{p\in P}I_p(X;Y)
\]
is attained by an algebraic input distribution, and $e^{C_P(W)}$
is algebraic.  A positive value in nats is transcendental; in bits
the value is rational or transcendental, with an exact decision
procedure.  This includes finitely many rational input-cost constraints.
\end{corollary}
\begin{proof}
The input law $p$, joint law $(p_xW(y\mid x))_{x,y}$, and output
law $(\sum_xp_xW(y\mid x))_y$ are rational affine probability
maps.  Apply Theorem~\ref{p03:mi3:thm:entropy} to
$I(X;Y)=H(X)+H(Y)-H(X,Y)$, then Theorem~\ref{p03:mi3:thm:bits}.
\end{proof}

The binary-symmetric rational-crossover special case is already
explicitly discussed and proved in Appendix~B of
\cite{p03:mi3:Boche}, and is not counted as a new result here.
Liang studies algebraic solvability of general capacity formulas
\cite{p03:mi3:Liang}.  Polynomial descriptions for each fixed rational
instance are not formulas in radicals, nor formulas algebraic in
arbitrary real channel parameters.  These distinctions prevent an
incorrect claim to resolve that separate closed-form problem.

\begin{corollary}[Rational finite-source rate--distortion values]
\label{p03:mi3:cor:RD}
Let $p_X$ be a rational distribution on a finite alphabet, let the
reproduction alphabet be finite, let the distortion entries
$d(x,y)\ge0$ and the threshold $D\ge0$ be rational, and suppose
the distortion constraint is feasible.  Then
\[
 R(D)=\min\left\{I(X;Y): P_X=p_X,
                         \ \sum_{x,y}P_{XY}(x,y)d(x,y)\le D\right\}
\]
has an algebraic optimizing joint law, and $e^{R(D)}$ is algebraic.
Every positive $R(D)$ in nats is transcendental, and the value in
bits is rational or transcendental with an exact decision procedure.
\end{corollary}
\begin{proof}
The feasible joint laws form a compact rational polytope.  Their
joint and two marginal distributions are rational affine maps.
Apply the two theorems to the same three-entropy expression for
mutual information.  Zero source probabilities simply give zero
rows; no positive-support hypothesis is needed.  For positive
source probabilities, dividing the algebraic joint law by the
rational source probabilities also produces an algebraic test channel.
\end{proof}

\begin{corollary}[Finite information-bottleneck optima]
\label{p03:mi3:cor:IB}
Fix a rational joint law $p_{XY}$ on finite alphabets, a positive
integer $m$, and a rational $\beta\ge0$.  Over stochastic encoders
$q(t\mid x)$ with $1\le t\le m$, set
\[
 B_{\beta,m}=\min_{q(t\mid x)}
                   \{I(X;T)-\beta I(Y;T)\},\qquad T-X-Y.
\]
There is an algebraic optimizing encoder and $e^{B_{\beta,m}}$
is algebraic.  Every nonzero optimum in nats is transcendental,
and its value in bits is rational or transcendental, decidably.
The same assertions hold after imposing rational linear constraints
on the encoder, provided the feasible set is nonempty.  Convexity of this objective is not required.
Without additional encoder constraints, the same assertions hold
for the infimum over \emph{all finite encoder alphabets}: an optimum
is attained with at most $|\operatorname{supp}p_X|$ symbols.
\end{corollary}
\begin{proof}
The encoder space, with any stated rational linear constraints,
is a rational polytope.  The joint probabilities
$p_{XY}(x,y)q(t\mid x)$ and all their marginals are rational
affine functions of the encoder.  Expand each mutual information
as three entropies and apply the two theorems.  Here $m$ is fixed;
no limiting argument over growing alphabets is needed.

For completeness, the unrestricted cardinality assertion follows
by a direct support reduction.  Delete zero-probability source
letters and write $n=|\operatorname{supp}p_X|$.  Any finite encoder
has output weights $w_t>0$ and posterior vectors
$r_t=P_{X\mid T=t}\in\mathbb R^n$, with
$\sum_tw_tr_t=p_X$.  If $V$ denotes the fixed channel $X\to Y$,
the objective is
\[
 H(X)-\beta H(Y)+\sum_tw_t\{\beta H(r_tV)-H(r_t)\}.
\]
When there are more than $n$ positive weights, the vectors $r_t$
are linearly dependent.  Choose nonzero real $a_t$ with
$\sum_ta_tr_t=0$; summing coordinates gives $\sum_ta_t=0$.
The feasible interval of $s$ for the weights $w_t+sa_t\ge0$
has endpoints of both signs, and at each endpoint some weight
vanishes.  The displayed objective is affine in $s$.  At one
endpoint it is no greater than at $s=0$.  Keep that endpoint and
repeat, preserving the mean posterior and reducing the output
support until at most $n$ symbols remain.  The new posteriors and
weights define an encoder by $q(t\mid x)=w_tr_t(x)/p_X(x)$.
Consequently the infimum over all finite alphabets equals the
minimum on the compact $n$-symbol encoder polytope.  The preceding
arithmetic conclusions apply to that polytope.
\end{proof}

Ordinary constrained maximum entropy is another immediate special
case.  Its polynomial stationarity mechanism is already present in
the cited toric-statistics literature and is not counted as a fourth
independent target.  The signed nonconvex objective in the bottleneck
corollary explains why a proof based solely on strict concavity would
not have sufficed for the general theorem.

\subsection{A globally optimized bottleneck value in closed form}

To exhibit an actual nonzero optimized constant, let $X$ be a fair
bit and $Y=X\oplus Z$, with $Z$ independent and
$\mathbb P(Z=1)=1/4$.  Consider all finite-valued encoders $T-X-Y$.

\begin{proposition}[An explicit nonconvex optimized information constant]
\label{p03:mi3:prop:explicitIB}
For this source,
\begin{equation}\label{p03:mi3:eq:explicitIB}
 \inf_{T-X-Y}\{I(X;T)-8I(Y;T)\}
          =\log\frac{11+5\sqrt5}{32}.
\end{equation}
The infimum is attained by the binary symmetric encoder
$T=X\oplus E$, with independent $E$ and
\[
 \mathbb P(E=1)=\frac{1-2\sqrt{\sqrt5-2}}2.
\]
The optimum is transcendental both in nats and in bits.  Numerically
it is approximately $-0.3665295969$ nats, or $-0.5287904318$ bits.
The underlying binary bottleneck reduction is classical
\cite{p03:mi3:WynerZiv,p03:mi3:Hsu}; the claim here is its exact evaluation
and arithmetic verification at the stated parameters, without a
historical priority claim.
\end{proposition}
\begin{proof}
Write $h(p)=-p\log p-(1-p)\log(1-p)$.  We first recall and prove
the one-bit convexity underlying Mrs.~Gerber's lemma.  For
$0<c<1$, the function
\[
 u=h((1-s)/2)\longmapsto h((1-cs)/2),\qquad 0\le s\le1,
\]
is convex.  For $0<s<1$ its derivative with respect to $u$ is
$c\operatorname{arctanh}(cs)/\operatorname{arctanh}(s)$.
The ratio $\operatorname{arctanh}(cs)/\operatorname{arctanh}(s)$
decreases with $s$: expand both series in $s^2$ after cancelling
$s$, and observe that the ratios of their positive coefficients
are $c^{2n+1}$, strictly decreasing in $n$.  Explicitly, the
numerator of the derivative of the ratio of these series is a sum
over $r>n$ of terms proportional to
$(r-n)(c^{2r+1}-c^{2n+1})(s^2)^{r+n-1}\le0$.
Since $u$ decreases with $s$, the asserted convexity follows;
continuity gives the endpoints.

Apply this with $c=1/2$ to each posterior law of $X$ given $T=t$.
Replacing its parameter by its complement if necessary does not
change either binary entropy.  Jensen's inequality yields
\[
 H(Y\mid T)\ge h(1/4+\varepsilon/2),\qquad
 H(X\mid T)=h(\varepsilon),\quad 0\le\varepsilon\le1/2.
\]
For every such $\varepsilon$, equality is attained by the binary
symmetric encoder with crossover $\varepsilon$.  As both $X$ and
$Y$ are fair, the optimization over all finite $T$ is therefore
exactly the minimum over $0\le\varepsilon\le1/2$ of
\[
        8h(1/4+\varepsilon/2)-h(\varepsilon)-7\log2.
\]
Set $s=1-2\varepsilon$.  Its derivative with respect to $s$ is
\[
             g(s)=\operatorname{arctanh}s
                              -4\operatorname{arctanh}(s/2).
\]
For $0<s<1$, the equation $g(s)=0$, after applying the strictly
increasing hyperbolic tangent and using its quadruple-angle
identity, is equivalent to
\[
                         s^4+16s^2-16=0.
\]
There is exactly one positive solution in $(0,1)$,
$s_*=2\sqrt{\sqrt5-2}$.  Moreover $g'(0)=-1$ and
$g(s)\to+\infty$ as $s\uparrow1$.  Thus the derivative is
negative before $s_*$ and positive after it, proving that this
point is the global minimum, including the two endpoints.

Put $\varepsilon_*=(1-s_*)/2$ and
$r_*=1/4+\varepsilon_*/2$.  Stationarity, or direct cancellation
of the log terms proportional to $\varepsilon_*-1/2$, rewrites
the minimum as
\[
 -4\log(r_*(1-r_*))
       +\tfrac12\log(\varepsilon_*(1-\varepsilon_*))-7\log2.
\]
Since
\[
 \sqrt{\varepsilon_*(1-\varepsilon_*)}
      =\frac{\sqrt5-2}{2},\qquad
 r_*(1-r_*)=\frac{3-\sqrt5}{4},
\]
its exponential equals
\[
       \frac{\sqrt5-2}{(3-\sqrt5)^4}
                              =\frac{11+5\sqrt5}{32}=:\alpha.
\]
Here $1/2<\alpha<1$, so the natural logarithm is nonzero and
transcendental.  The degree is two and the absolute norm is
$1/256=2^{-8}$.  A rational $\log_2\alpha$ would therefore have
to be $-4$ by Theorem~\ref{p03:mi3:thm:bits}, which contradicts
$1/2<\alpha<1$.  The value in bits is consequently transcendental.
\end{proof}

\subsection{Scope, sharp hypotheses, and unsuccessful extensions}

The three targets above concern finite distributions with exact
rational data.  Their common proof is an arithmetic consequence of
semialgebraic critical equations, not a new general transcendence
theorem and not a claim to solve a named historical conjecture.
The following distinctions are essential.
\begin{enumerate}
\item \emph{An algebraic optimizer alone is insufficient.}
For $a=1/\sqrt2$, the exponential of the entropy of the fixed
law $(a,1-a)$ is
\[
 e^{h(a)}=(1-a)^{-1}\left(\frac{1-a}{a}\right)^a,
\]
which is transcendental by Gelfond--Schneider.  In the theorem,
the rational point $x_0$ on the same stationary face changes
algebraic log coefficients into rational ones.  A general
algebraic singleton polytope does not have this property.
\item \emph{Rational inputs cannot be replaced indiscriminately by
algebraic inputs.}  A binary erasure channel with algebraic
irrational erasure probability $\epsilon\in(0,1)$ has capacity
$1-\epsilon$ bits, which is algebraic irrational.  Indeed the
erasure flag is independent of the input and an unerased symbol
reveals it, so $I(X;Y)=(1-\epsilon)H(X)$; the fair input is optimal.
This is an explicit failure of the rational-or-transcendental
classification outside the stated hypotheses.
\item \emph{No arbitrary added rational linear cost is allowed.}
Even adding the constant one to an identically zero entropy
objective makes its optimized exponential equal to $e$.
Rational costs in the capacity and rate--distortion statements
are constraints defining the feasible polytope, not additive
terms in \eqref{p03:mi3:eq:objective}.
\item \emph{Finite alphabets matter.}  For a Poisson law with
algebraic positive mean $\lambda$, its Shannon entropy contains
$e^{-\lambda}\sum_{n\ge0}\lambda^n\log(n!)/n!$.
The logarithmic coefficients do not meet the algebraic-coefficient
$E$-function framework used for $\operatorname{Ein}$.  The finite
critical-cell proof gives no arithmetic classification of this
infinite series.  No irrationality of the entropy of
$\operatorname{Poisson}(1)$ is claimed.
\item \emph{The previous optimized exponential-integral payoff
remains unresolved here.}  At the earlier transcendental critical
point of $\operatorname{Ein}(x)/x^\theta$, the algebraic-point
Siegel--Shidlovskii theorem cannot be invoked.  The present finite
entropy theorem does not repair that missing specialization.
\end{enumerate}

The primary-source comparison is bounded.  Maximum-entropy
algebraization, toric stationarity, the binary entropy transcendence
example, and the binary bottleneck tradeoff are credited above.
The checked sources do not establish historical novelty of the
full signed-entropy critical-value formulation or the norm-based
decision corollary.  These are proved research statements with
uncertified priority, rather than announced resolutions of old
open problems.


% Integrated source: arithmetic_results.tex
\section{Exponential-integral division and genuinely reduced denominators}
\label{p03:ac19:section}

The individual targets in this section are Euler's constant $\gamma$,
the Euler--Gompertz constant $\delta$, Catalan's constant $G$, and
$\zeta(5)$.  None is classified here as irrational by a new proof.
There are two substantive advances over the preceding polynomial-kernel
analysis: an explicit application of the new arithmetic morphism of
Fischler--Rivoal to rational functions of an exponential-integral value,
and denominator obstructions using the \emph{actual reduced rational
remainder}, not a convenient common multiple.  A constructive saturation
theorem then explains exactly why further common-factor division is a
different issue.  These are self-posed results and dependent applications;
no priority assertion is made for them.

We write
\[
 \delta_x=\int_0^\infty\frac{e^{-t}}{t+x}\,dt\quad(x>0),
 \qquad \delta=\delta_1,
 \qquad \operatorname{Ein}(z)=\int_0^z\frac{1-e^{-t}}t\,dt.
\]
The familiar connection formula is
\begin{equation}\label{p03:ac19:connection}
 \operatorname{Ein}(x)=\gamma+\log x+e^{-x}\delta_x\quad(x>0).
\end{equation}
It follows by differentiation and the defining limit of $\gamma$.
We use its analytic continuation with the principal logarithm on the
plane cut along the nonpositive real axis.

\subsection{A new division obstruction for the exponential integral}

Let $\mathcal E$ be the ring of values of $E$-functions at algebraic
points.  The distinction between this ring and its fraction field is
essential.  We use the following external theorem of Fischler--Rivoal
\cite[Proposition 3.1]{p03:ac19:FRring}: there is an injective ring morphism
\[
 \Phi_0:\mathcal E\longrightarrow
 \mathcal R:=\left\{\sum_{\rho\in\overline{\mathbb Q}}
                       c_\rho e^{\rho z}:c_\rho\in\mathbb C,
                       \text{finite support}\right\}.
\]
It fixes algebraic numbers.  If $f(z)=\sum_\rho e^{\rho z}f_\rho(z)$
is the canonical large-sector asymptotic decomposition of an
$E$-function, then
$\Phi_0(f(1))=\sum_\rho f_\rho(1)e^{\rho z}$, with the prescribed
one-sided convention at an exceptional direction.  This external result
is not reproved or counted as our discovery.

\begin{proposition}[The two coefficients of an exponential-integral value]
\label{p03:ac19:PhiEin}
Put $u=\operatorname{Ein}(1)$.  Then
\begin{equation}\label{p03:ac19:Phi-formula}
                   \Phi_0(u)(z)=\gamma+\delta e^{-z}.
\end{equation}
In particular $u$ is not a unit of $\mathcal E$, so that
$1/\operatorname{Ein}(1)\notin\mathcal E$.
More generally, for every positive algebraic $a\ge1$,
$1/\operatorname{Ein}(a)\notin\mathcal E$.
\end{proposition}
\begin{proof}
The coefficients of
$\operatorname{Ein}(z)=\sum_{n\ge1}(-1)^{n+1}z^n/(n\,n!)$
have the required exponential denominator growth, and the function
satisfies a linear differential equation over $\mathbb Q(z)$;
thus it is an $E$-function.  For complex $z$ off the negative real axis,
\[
 \delta_z=\int_0^\infty\frac{e^{-t}}{t+z}\,dt
   \sim\sum_{k\ge0}\frac{(-1)^k k!}{z^{k+1}}.
\]
For completeness, on each closed sector $|\arg z|\le\pi-\varepsilon$
one has $|t+z|\ge c_\varepsilon(t+|z|)$.  Expanding $(t+z)^{-1}$
through $N$ terms and integrating the remainder gives an error at most
$c_\varepsilon^{-1}N!|z|^{-N-1}$.  This is the requisite factorial
asymptotic expansion on sectors of opening greater than $\pi$.
Formula~\eqref{p03:ac19:connection} therefore gives the canonical
decomposition with phases $0,-1$ and coefficient functions
$\gamma+\log z,\delta_z$.  The external morphism theorem gives
\eqref{p03:ac19:Phi-formula}.

Both $\gamma$ and $\delta$ are positive.  Hence
$\gamma+\delta e^{-z}$ has a complex zero.  If $u$ had an inverse
in $\mathcal E$, its image would be an entire exponential polynomial
whose product with this function were $1$, which is impossible at that
zero.  Replacing $z$ by $az$ gives
\[
 \Phi_0(\operatorname{Ein}(a))(z)
          =\gamma+\log a+\delta_a e^{-az}.
\]
For $a\ge1$ both coefficients are positive, and the same argument applies.
\end{proof}

There is an exact, more informative rational-function consequence.
The preceding proposition also shows that $u$ is transcendental: every
nonzero algebraic number and its inverse belong to $\mathcal E$.
Consequently evaluation identifies $\overline{\mathbb Q}(T)$ with
$\overline{\mathbb Q}(u)$.

\begin{theorem}[The only possible exceptional denominator]
\label{p03:ac19:rational-division}
Let $R\in\overline{\mathbb Q}(T)$ be nonpolynomial.  If
$R(u)\in\mathcal E$, then $\gamma$ is algebraic, every finite pole of
$R$ is $\gamma$, and $\delta$ is a unit of $\mathcal E$.
In fact
\begin{equation}\label{p03:ac19:intersection}
 \mathcal E\cap\overline{\mathbb Q}(u)=
 \begin{cases}
 \overline{\mathbb Q}[u,(u-\gamma)^{-1}],
       &\gamma\in\overline{\mathbb Q}\text{ and }\delta\in\mathcal E^\times,\\
 \overline{\mathbb Q}[u],&\text{otherwise}.
 \end{cases}
\end{equation}
Thus any rational function having a finite pole at an algebraic number
outside $(0,1)$ takes at $u$ a value outside $\mathcal E$.
In particular $1/(u-k)\notin\mathcal E$ for every integer $k$.
\end{theorem}
\begin{proof}
Write $R=A/B$ in lowest terms.  If $v=R(u)\in\mathcal E$, apply
$\Phi_0$ to $B(u)v=A(u)$.  On the open set where its denominator is
nonzero this yields
\[
       \Phi_0(v)(z)=R(\gamma+\delta e^{-z}).
\]
The left side is entire.  If $\alpha\ne\gamma$ were a finite pole of
$R$, the equation $e^{-z}=(\alpha-\gamma)/\delta$ would have a complex
solution.  The derivative of $\gamma+\delta e^{-z}$ at that solution
is $-(\alpha-\gamma)\ne0$, so the pole of $R$ would remain a pole
of the composite, a contradiction.  All finite poles are therefore
$\gamma$.  A nonpolynomial rational function has a finite pole, and
its poles are algebraic, proving $\gamma\in\overline{\mathbb Q}$.

Now $B(T)$ is a nonzero constant times $(T-\gamma)^m$ for some $m\ge1$.
Bezout's identity for the coprime polynomials $A,B$ gives
$S A+T_0B=1$ with $S,T_0\in\overline{\mathbb Q}[T]$.
Evaluation shows $1/B(u)=S(u)v+T_0(u)\in\mathcal E$.
Multiplication by a suitable power of $u-\gamma$ gives
$(u-\gamma)^{-1}\in\mathcal E$.  Since
$\delta=e(u-\gamma)$ and $e,e^{-1}\in\mathcal E$, the number
$\delta$ is a unit.  Conversely, if $\gamma$ is algebraic and $\delta$
is a unit, the localization displayed in the first line of
\eqref{p03:ac19:intersection} lies in $\mathcal E$.
The necessary pole condition proves that it contains the whole
intersection.  If either hypothesis fails, the first part of the theorem
excludes every nonpolynomial rational function and gives the second line.
Finally the elementary bounds $0<\gamma<1$ exclude the stated poles.
\end{proof}

The theorem isolates the precise unresolved exception instead of
assuming that every transcendental $E$-value has an inverse in the ring.
It does not decide whether $\gamma$ is algebraic, nor whether $\delta$
belongs to $\mathcal E$.  The mixed-Gevrey separation statements in
\cite{p03:ac19:FRmixed} remain conjectural inputs; the new ring theorem does
not supply them.  The explicit nonmembership above is a dependent
application of \cite{p03:ac19:FRring}, rather than a new proof of individual
transcendence for either connection coefficient.

\subsection{Why a common deformation does not furnish a second equation}

\begin{proposition}[Finite differential cancellation removes the Euler term]
\label{p03:ac19:monodromy}
Let $a_1,\ldots,a_m$ be positive algebraic numbers, and let
$L_j=\sum_{k=0}^{N_j}r_{jk}(z)(d/dz)^k$ with
$r_{jk}\in\overline{\mathbb Q}(z)$.  Put
\[
 Y_j(z)=e^{-a_jz}\delta_{a_jz},\qquad U(z)=\sum_{j=1}^mL_jY_j(z).
\]
If $U$ is single valued on a punctured neighborhood of $0$, then
$\sum_j r_{j0}(z)=0$.  In the connection-formula expansion of $U$,
the coefficient of $\gamma$ consequently vanishes identically.
Adding any single-valued meromorphic expression formed from
$E$-functions and algebraic rational functions does not change this
conclusion.  If all the $r_{jk}$ are regular at $1$, then under the
single-valued hypothesis
\[
                  U(1)\in\mathcal E+
                       \sum_{j=1}^m\overline{\mathbb Q}\log a_j.
\]
\end{proposition}
\begin{proof}
One has $Y_j'(z)=-e^{-a_jz}/z$, so every positive derivative of $Y_j$
is an exponential times an algebraic rational function.  The only
multivalued terms come from its zeroth derivative:
\[
 Y_j(z)=\operatorname{Ein}(a_jz)-\gamma-\log a_j-\log z.
\]
The monodromy of $U$ around $0$ is therefore
$-2\pi i\sum_jr_{j0}(z)$.  Single valuedness forces this rational
function to vanish.  The same sum is the negative coefficient of
$\gamma$.  The assertion at $1$ now follows directly from the remaining
terms and closure of $\mathcal E$ under algebraic scalar multiplication
and addition.
\end{proof}

For the individual Gompertz constant, a common-parameter derivative has
no independent arithmetic content either.  Differentiation under the
integral and integration by parts give
\[
 \delta_x'=\delta_x-x^{-1},\qquad
 \delta_x^{(k)}=\delta_x-
         \sum_{\ell=0}^{k-1}(-1)^\ell\ell!x^{-\ell-1}\quad(k\ge1).
\]
The second identity follows by induction.  At $x=1$ all these
derivatives lie in $\delta+\mathbb Z$.  Differentiating
\eqref{p03:ac19:connection} cancels these occurrences of $\delta$ and yields
the already known derivatives of $\operatorname{Ein}$.  This is an exact
failure of this finite differential route, not a general impossibility
theorem about analytic identities involving $\gamma$ or $\delta$.

\subsection{Actual reduced denominators for positive zeta kernels}

For an integer $s\ge2$, define
\[
 \mathcal Z_s(P)=\frac1{(s-1)!}\int_0^1
              \frac{(-\log x)^{s-1}}{1-x}P(x)\,dx.
\]
If $P(x)=\sum_{k=0}^Na_kx^k\in\mathbb Z[x]$, integration term by term
gives the \emph{formal} coefficient identity
\begin{equation}\label{p03:ac19:zeta-coefficients}
 \mathcal Z_s(P)=P(1)\zeta(s)-R_s(P),\qquad
 R_s(P)=\sum_{k=0}^Na_kH_k^{(s)}\in\mathbb Q.
\end{equation}
Here $H_0^{(s)}=0$ and $H_k^{(s)}=\sum_{j=1}^kj^{-s}$.
No irrationality hypothesis is used in defining these coefficients.
Write $b_s(P)$ for the denominator of $R_s(P)$ in lowest terms,
with denominator $1$ for zero.  This is the smallest positive integer
clearing that rational remainder.

\begin{lemma}[Endpoint amplitude controls a positive polynomial integral]
\label{p03:ac19:amplitude}
If $P$ is a nonnegative real polynomial on $[0,1]$, of degree at most
$N\ge1$, and $P(1)=q>0$, then
\[
 \mathcal Z_s(P)\ge
       \frac{q}{2(s-1)!(s-1)(8N^3)^{s-1}}.
\]
\end{lemma}
\begin{proof}
Set $M=\max_{[0,1]}P\ge q$.  Its shifted Chebyshev expansion has
coefficients of absolute value at most $2M$ in positive degree, by the
cosine integral formulas.  Since $|T_j'|\le j^2$ on $[-1,1]$,
$\|P'\|_\infty\le4M\sum_{j=1}^Nj^2\le4MN^3$.
On an interval of length $h=(8N^3)^{-1}$ adjacent to a maximum,
and contained in $[0,1]$, one therefore has $P\ge M/2$.
The weight satisfies
$(-\log x)^{s-1}/(1-x)\ge(1-x)^{s-2}$.
Among intervals of length $h$ its last weight has least integral at
the right endpoint, namely $h^{s-1}/(s-1)$.  Substitution proves the
claimed inequality.
\end{proof}

\begin{theorem}[Prime cancellation is compensated by the target coefficient]
\label{p03:ac19:zeta-reduced}
Let $n\ge2$, $0\le r<n$, and let $Q\in\mathbb Z[x]$ have degree at
most $r$, be nonnegative on $[0,1]$, and satisfy $q=Q(1)>0$.
Put $P=x^nQ$, $N=n+r$, and
\[
             \Pi_{n,r}=\prod_{(n+r)/2<p\le n}p,
\]
the product being over primes.  For every prime in this product,
\begin{equation}\label{p03:ac19:zeta-valuation}
 v_p(b_s(P))=\max\{0,s-v_p(q)\}.
\end{equation}
Consequently $\Pi_{n,r}^s$ divides $b_s(P)q$, and
\begin{equation}\label{p03:ac19:zeta-lower}
 b_s(P)\mathcal Z_s(P)\ge
 \frac{\Pi_{n,r}^s}{2(s-1)!(s-1)(8N^3)^{s-1}}.
\end{equation}
For every fixed $\eta<1$, uniformly for $r\le\eta n$ and for all such
integer polynomials $Q$ without a coefficient-height restriction,
\[
 \log\bigl(b_s(P)\mathcal Z_s(P)\bigr)
          \ge\frac{s(1-\eta)}2n+o(n)\longrightarrow+\infty.
\]
In particular the reduced-denominator loophole does not rescue these
kernels for $\zeta(5)$, or for any fixed integer $s\ge2$.
\end{theorem}
\begin{proof}
Write $Q=\sum_{j=0}^rc_jx^j$.
For a prime $(n+r)/2<p\le n$, the only multiple of $p$ among
$1,\ldots,n+r$ is $p$ itself, and every harmonic sum $H_{n+j}^{(s)}$
contains its term $p^{-s}$.  Thus
\[
 R_s(P)=\frac{\sum_jc_j}{p^s}+U_p
                 =\frac{q}{p^s}+U_p,
 \qquad U_p\in\mathbb Z_{(p)}.
\]
If $v_p(q)<s$, the first summand has negative $p$-adic valuation and
the second has nonnegative valuation.  Its denominator has exactly
$p^{s-v_p(q)}$.  If $v_p(q)\ge s$, both summands are $p$-integral.
This proves \eqref{p03:ac19:zeta-valuation}, including the possibility
that $R_s(P)=0$.  Multiplying the divisibilities over the distinct
primes gives $b_s(P)q\ge\Pi_{n,r}^s$.
Lemma~\ref{p03:ac19:amplitude} applied to $P$ proves
\eqref{p03:ac19:zeta-lower}.

The prime number theorem, $\vartheta(X)=X+o(X)$, gives
\[
 \log\Pi_{n,r}=\vartheta(n)-\vartheta((n+r)/2)
                     =(n-r)/2+o(n).
\]
The error is uniform for $0\le r\le\eta n$, since both arguments
lie in fixed proportional intervals.  All losses in
\eqref{p03:ac19:zeta-lower} are logarithmic in $n$ for fixed $s$.
This proves the last assertion.  Only this classical weak form of the
prime number theorem is needed; stronger unconditional estimates are
available in \cite{p03:ac19:Dusart}.
\end{proof}

This theorem replaces the standard $\operatorname{lcm}(1,\ldots,N)^s$
by the actual reduced denominator.  Its support hypothesis is material:
the order of vanishing $n$ at $0$ exceeds the remaining degree by a
fixed proportion, and $Q(1)\ne0$.  It does not cover every nonnegative
polynomial or multi-integral construction.  A canceled prime is allowed;
the point is that its loss from the denominator is paid for by its
appearance in $q$, which the positive integral also detects.

\subsection{A reduced-denominator obstruction for Catalan kernels}

Set
\[
 \mathcal C(P)=\int_0^1\frac{-\log x}{1+x^2}P(x^2)\,dx
       =\int_0^1w(t)P(t)\,dt,
 \qquad w(t)=\frac{-\log t}{4\sqrt t(1+t)}\ge\frac{1-t}{8}.
\]
For $m\ge0$, polynomial division gives
\begin{equation}\label{p03:ac19:catalan-moment}
 \mathcal C(t^m)=(-1)^mG+
             \sum_{j=0}^{m-1}\frac{(-1)^{m-1-j}}{(2j+1)^2}.
\end{equation}
An empty sum is zero.  Define the rational remainder $R_C(P)$ by
$\mathcal C(P)=P(-1)G+R_C(P)$ using this identity, and let $b_C(P)$
be its actual reduced denominator.

\begin{theorem}[Exterior amplitude and surviving odd primes]
\label{p03:ac19:catalan-reduced}
Let $n\ge2$ and $0\le r\le n/3$.  Let $Q\in\mathbb Z[t]$ have
degree at most $r$, be nonnegative on $[0,1]$, and satisfy
$q=Q(-1)\ne0$.  Put $P=t^nQ$, $R=\max(1,r)$, and
\[
 \Omega_{n,r}=\prod_{\substack{(2(n+r)-1)/3<p\le2n-1\\p\text{ odd prime}}}p.
\]
Then $\Omega_{n,r}^2$ divides $b_C(P)|q|$, and
\begin{equation}\label{p03:ac19:catalan-lower}
 b_C(P)\mathcal C(P)\ge
       \frac{2^{-n}14^{-r}}{8192R^6}\,\Omega_{n,r}^2.
\end{equation}
Uniformly in all such $Q$,
\[
 \log\bigl(b_C(P)\mathcal C(P)\bigr)
 \ge\left(\frac{20}{9}-\log2-\frac{\log14}{3}\right)n+o(n)
 \longrightarrow+\infty.
\]
The coefficient in parentheses is positive.
\end{theorem}
\begin{proof}
For a prime in the stated interval, the only odd multiple of $p$
among $1,3,\ldots,2(n+r)-1$ is $p$.  Moreover $p\le2n-1$, so its
square denominator occurs in every moment of $t^{n+j}$, $0\le j\le r$.
Writing $Q=\sum_jc_jt^j$ and $k=(p-1)/2$, the coefficient of $p^{-2}$
in \eqref{p03:ac19:catalan-moment} summed over these moments is
\[
 \sum_{j=0}^rc_j(-1)^{n+j-1-k}=(-1)^{n-1-k}Q(-1).
\]
All remaining terms are $p$-integral.  Exactly as in
\eqref{p03:ac19:zeta-valuation},
$v_p(b_C(P))=\max\{0,2-v_p(q)\}$, proving the divisibility.

Let $M=\max_{[1/2,1]}Q$.  The exterior Chebyshev inequality on this
interval yields
\begin{equation}\label{p03:ac19:exterior}
                  |Q(-1)|\le T_r(7)M\le14^rM.
\end{equation}
Here the affine coordinate is $u=4t-3$, for which $t=-1$ corresponds
to $u=-7$.  To recall the elementary exterior inequality, if a degree
at most $r$ polynomial bounded by $1$ on $[-1,1]$ exceeded
$|T_r(u_0)|$ at a real $|u_0|>1$, subtract the corresponding multiple
of $T_r$.  At its $r+1$ alternating extrema this difference has $r$
interior sign changes; the additional root at $u_0$ is impossible.
Finally $T_r(7)\le(7+\sqrt{48})^r<14^r$ for $r\ge1$;
degree zero is immediate.

For $r\ge1$ the same cosine-coefficient derivative estimate used
earlier, now on an interval of length $1/2$, gives
$\|Q'\|_{[1/2,1]}\le8Mr^3$.  There is an interval of length
$h=(16R^3)^{-1}$ in $[1/2,1]$ on which $Q\ge M/2$.
For $r=0$ this assertion holds trivially.  On this interval $t^n\ge2^{-n}$;
the least integral of $(1-t)$ over an interval of length $h$ is $h^2/2$.
Consequently
\[
 \mathcal C(t^nQ)\ge
       2^{-n}\frac M2\frac18\frac{h^2}2
       =\frac{2^{-n}M}{8192R^6}
       \ge\frac{2^{-n}14^{-r}|q|}{8192R^6}.
\]
Combining with $b_C(P)|q|\ge\Omega_{n,r}^2$ proves
\eqref{p03:ac19:catalan-lower}.

The prime number theorem gives, uniformly for $r\le n/3$,
\[
       \log\Omega_{n,r}=(4n-2r)/3+o(n).
\]
The resulting logarithmic lower bound is
$(8n-4r)/3-n\log2-r\log14-O(\log n)+o(n)$.
Its minimum over this range of $r$ is the asserted expression.
For example $\log2<0.7$ and $\log14<2.64$ already show that its
coefficient is greater than $0.64$.
\end{proof}

This result covers genuine coefficient cancellation in the rational
remainder.  The restriction on the remaining degree and the requirement
$Q(-1)\ne0$ are part of the theorem, not conclusions about every
approximating construction for $G$.
An OpenAI Math Release preprint dated 24 September 2026 asserts the
irrationality of $G$ \cite{p03:root3:Catalan}.  Its full proof is not
audited or used here.  Its signed mixed-moment determinant construction
is outside our single positive polynomial kernel and support hypotheses;
the present obstruction neither verifies nor refutes that claim.

\subsection{Complete coefficient reduction: positivity has no remaining force}

The smallest denominator of the rational remainder and the primitive
normalization of the whole coefficient pair must not be conflated.
If an integer kernel gives $A\xi+B$, $A\in\mathbb Z$, $B\in\mathbb Q$,
and $b$ is the denominator of $B$, the coefficients of $b(A\xi+B)$
may still share a common divisor.  The next theorem describes the much
larger class obtained by dividing that common divisor as well.

\begin{theorem}[Constructive saturation of three positive kernel cones]
\label{p03:ac19:saturation}
Let $\xi$ be any of $\delta$, $G$, or $\zeta(s)$ for a fixed integer
$s\ge2$.  For every relatively prime pair $p\in\mathbb Z$, $q\in\mathbb Z_{>0}$
with $q\xi-p>0$, there is a nonzero integer polynomial $P$, nonnegative
on the corresponding integration interval, whose integral has formal
coefficients $Mq,-Mp$ for some positive integer $M$:
\[
 \int_0^\infty\frac{e^{-t}P(t)}{1+t}\,dt=M(q\delta-p),
 \quad\text{or}\quad
 \mathcal C(P)=M(qG-p),
 \quad\text{or}\quad
 \mathcal Z_s(P)=M(q\zeta(s)-p).
\]
Thus its rational remainder is already integral, and primitive reduction
of its complete coefficient pair gives exactly $q\xi-p$.
\end{theorem}
\begin{proof}
In each case we construct a rational polynomial $K_n\ge0$ with target
coefficient $1$ and integral $\xi-r_n>0$, where $r_n\in\mathbb Q$
and $r_n\to\xi$.  We also give a positive polynomial $S$ with target
coefficient $0$ and integral $1$.

For $\xi=\delta$, let
\[
 L_n(t)=\sum_{k=0}^n(-1)^k\binom nk\frac{t^k}{k!},
 \qquad K_n(t)=\frac{L_n(t)^2}{L_n(-1)^2},\qquad S(t)=1+t.
\]
The Laguerre orthogonality normalization is
$\int_0^\infty e^{-t}L_n(t)^2dt=1$.  It follows directly from
Rodrigues' formula
$L_n(t)=e^t(n!)^{-1}(d/dt)^n(e^{-t}t^n)$
by $n$ integrations by parts.  Also $L_n(-1)\ge n+1$, so
\[
 0<\int_0^\infty\frac{e^{-t}K_n(t)}{1+t}\,dt
                         \le\frac1{(n+1)^2}\longrightarrow0.
\]
Its target coefficient is $K_n(-1)=1$.  Division by $1+t$ and
the factorial moments give a rational remainder.  The assertions about
$S$ are immediate.

For $\xi=G$, take
\[
               K_n(t)=2^{-n}(1-t)^n,\qquad S(t)=1+t.
\]
Here $K_n(-1)=1$ and
$0<\mathcal C(K_n)\le2^{-n}G\to0$.  Formula
\eqref{p03:ac19:catalan-moment} gives its rational remainder, while
$\mathcal C(1+t)=\int_0^1-\log x\,dx=1$ and $S(-1)=0$.

For $\xi=\zeta(s)$, take $K_n(x)=x^n$ and $S(x)=1-x$.
Then $r_n=H_n^{(s)}\to\zeta(s)$, while $S(1)=0$ and
\[
 \mathcal Z_s(S)=\frac1{(s-1)!}\int_0^1(-\log x)^{s-1}dx=1.
\]

Given the specified $p,q$, choose $n$ with $r_n>p/q$.  The rational
polynomial
\[
                  W=qK_n+(qr_n-p)S
\]
is nonnegative, has target coefficient $q$, and has integral
$q\xi-p$.  Choose a positive integer $M$ clearing all coefficients
of $W$, and put $P=MW$.  The formal coefficient pair is exactly
$(Mq,-Mp)$, whose greatest common divisor is $M$ because $p,q$ are
coprime.  This proves every assertion, without presupposing
irrationality of any of the three target constants.
\end{proof}

The theorem is a constructive statement about the image of a cone under
its arithmetic moment map.  In particular, a uniform positive lower
bound for \emph{all primitive-normalized} positive kernels cannot be
deduced merely from positivity: that would be a bound for every positive
primitive rational-approximation form $q\xi-p$ itself.  The zeta and
Catalan reduced-denominator theorems above concern the displayed integer
kernels before this extra common-factor division.  They do not close
that logically different loophole.  The saturation construction also
does not control $M$; it gives no previously unknown small integer form
for $\gamma$, $\delta$, $G$, or $\zeta(5)$.


% Integrated source: information_results.tex
\section{Three arithmetic generating functions beyond the modular cusp field}
\label{p03:ci19:section}

The three number-theoretic targets in this section are the Erd\H{o}s--Borwein
constant, the prime binary constant, and the prime Lambert constant,
\[
 E=\sum_{n\ge1}\frac1{2^n-1},\qquad
 P_2=\sum_{p\ {\rm prime}}2^{-p},\qquad
 L_{\mathbb P}=\sum_{p\ {\rm prime}}\frac1{2^p-1}.
\]
Their transcendence is not proved here.  Instead, we prove that their
generating functions are algebraically independent over a field containing
every finite collection of dilated Eisenstein series and their derivatives.
This answers a precise functional question behind an attempted modular
transfer: no finite algebraic elimination within that field can isolate
any of these three functions from the other two.  The result extends to
each individual reciprocal-Mersenne moment.  It does not assert joint
independence of all the moments, which would be false over this field.

For $|q|<1$ and an integer $r\ge1$, write
\begin{equation}\label{p03:ci19:functions}
 A_r(q)=\sum_{n\ge1}\left(\frac{q^n}{1-q^n}\right)^r,
 \qquad P(q)=\sum_{p\ {\rm prime}}q^p,
 \qquad B(q)=\sum_{p\ {\rm prime}}\frac{q^p}{1-q^p}.
\end{equation}
All series converge locally uniformly in the open unit disk.  In particular,
$E=A_1(1/2)$, $P_2=P(1/2)$, and $L_{\mathbb P}=B(1/2)$.
Functional independence of $q$-zeta series has a substantial prior
literature; see, for example, Pupyrev \cite{p03:ci19:Pupyrev2005} and
B\'ezivin \cite{p03:ci19:Bezivin2008}.  We make no priority claim for the general
idea of using boundary asymptotics to establish functional independence.
The statement proved below concerns the specified prime-supported functions
and the explicit enlarged coefficient field.

\subsection{A coefficient field that includes finite modular constructions}

Throughout this section $t$ is positive, $t\downarrow0$, $q=e^{-t}$,
and $L=\log(1/t)$.  For $\lambda>0$, define
\begin{equation}\label{p03:ci19:cuspfield}
 \mathscr K_\lambda
 =\left\{F(t,e^{-\lambda/t}):
       F\in\operatorname{Frac}\mathbb C\{t,z\}\right\}.
\end{equation}
Here $\mathbb C\{t,z\}$ denotes convergent power-series germs at $(0,0)$,
and the expressions in \eqref{p03:ci19:cuspfield} are germs on the positive
real axis.  The next lemma verifies, in particular, that the evaluation is
injective and that nonzero denominators cause no ambiguity near $t=0$.

\begin{lemma}[Leading order in the cusp field]\label{p03:ci19:leading}
For every nonzero $h\in\mathscr K_\lambda$ there exist integers $a,m$ and
$c\in\mathbb C\setminus\{0\}$ such that
\begin{equation}\label{p03:ci19:leadingform}
 h(t)=c t^a e^{-m\lambda/t}(1+O(t)).
\end{equation}
The field contains $t$ and $e^{-t}$ and is closed under differentiation
with respect to $t$.
\end{lemma}

\begin{proof}
For a nonzero convergent germ $G(t,z)$, choose the least $m_0\ge0$ for
which the coefficient $g_{m_0}(t)$ of $z^{m_0}$ is not identically zero.
There are an integer $a_0\ge0$ and a nonzero $c_0$ such that
$g_{m_0}(t)=c_0t^{a_0}(1+O(t))$.  Convergence in a polydisk gives
\[
 G(t,z)=z^{m_0}\bigl(g_{m_0}(t)+zH(t,z)\bigr)
\]
with $H$ analytic and bounded in a smaller polydisk.  Substitution of
$z=e^{-\lambda/t}$ shows that the second term inside the parentheses
is $o(t^N)$ for every $N$, and proves \eqref{p03:ci19:leadingform} for $G$.
Thus the evaluated germ is nonzero and has no zeros for all sufficiently
small positive $t$.  Taking quotients proves the assertion for the field
and proves injectivity of evaluation.  Finally, the chain rule gives
\[
 \frac{d}{dt}F(t,e^{-\lambda/t})
 =\left(F_t+\lambda zt^{-2}F_z\right)(t,e^{-\lambda/t}),
\]
and the meromorphic germ on the right belongs to the same fraction field.
\end{proof}

Let $\mathcal E_2,\mathcal E_4,\mathcal E_6$ denote the normalized
Eisenstein series, whose constant terms are $1$.

\begin{proposition}[Finite modular fields lie in a cusp field]
\label{p03:ci19:modularfield}
Let $\mathcal D$ be a finite nonempty set of positive integers and
$D=\operatorname{lcm}\mathcal D$.  The differential field over $\mathbb C$
generated by $q$ and
\[
 \mathcal E_2(q^d),\quad\mathcal E_4(q^d),\quad
 \mathcal E_6(q^d)\qquad(d\in\mathcal D),
\]
where the derivation is $q\,d/dq$, is contained, after $q=e^{-t}$, in
$\mathscr K_{4\pi^2/D}$.
\end{proposition}

\begin{proof}
The modular transformation and its weight-two correction, in the
normalization of \cite[Section~1]{p03:ci19:KanekoZagier1995}, give
\begin{align*}
 \mathcal E_2(e^{-t})
 &=-\frac{4\pi^2}{t^2}\mathcal E_2(e^{-4\pi^2/t})+\frac{12}{t},\\
 \mathcal E_4(e^{-t})
 &=\frac{16\pi^4}{t^4}\mathcal E_4(e^{-4\pi^2/t}),\\
 \mathcal E_6(e^{-t})
 &=-\frac{64\pi^6}{t^6}\mathcal E_6(e^{-4\pi^2/t}).
\end{align*}
Replace $t$ by $dt$ and set $z=e^{-4\pi^2/(Dt)}$.  The transformed
argument is $z^{D/d}$, an integer power of $z$, and each Eisenstein
series is analytic at argument zero.  Thus all the displayed generators
lie in the indicated field.  The derivation is $-d/dt$, so
Lemma~\ref{p03:ci19:leading} completes the proof.
\end{proof}

\subsection{The logarithmic coordinate of every reciprocal-Mersenne moment}

The following estimate retains a power-saving remainder after removal of
the polar terms.  That remainder, rather than only a leading equivalent,
is useful for the subsequent independence argument.

\begin{lemma}[A flat logarithmic coordinate]\label{p03:ci19:momentasymptotic}
Define the real coefficients $\kappa_{r,j}$ by the Laurent expansion
\[
 (e^x-1)^{-r}=\sum_{j=1}^r\kappa_{r,j}x^{-j}+O(1)
 \qquad(x\longrightarrow0).
\]
Then $\kappa_{r,1}=(-1)^{r-1}$.  There is a real constant $C_r$ such that
\begin{equation}\label{p03:ci19:momentexpansion}
 A_r(e^{-t})=
 \sum_{j=2}^r\kappa_{r,j}\zeta(j)t^{-j}
 +\kappa_{r,1}\frac{L}{t}+\frac{C_r}{t}+O(1).
\end{equation}
Consequently the function
\begin{equation}\label{p03:ci19:Xcoordinate}
 X_r(t)=\frac{1}{\kappa_{r,1}}
 \left(tA_r(e^{-t})-
       \sum_{j=2}^r\kappa_{r,j}\zeta(j)t^{1-j}-C_r\right)
\end{equation}
satisfies $X_r(t)=L+O(t)$.
\end{lemma}

\begin{proof}
The coefficient of $x^{-1}$ is the residue at zero.  Under $z=e^x-1$,
\[
 (e^x-1)^{-r}\,dx=\frac{z^{-r}}{1+z}\,dz,
\]
so that residue is $(-1)^{r-1}$.  Put $\kappa=\kappa_{r,1}$ and
\[
 g_r(x)=(e^x-1)^{-r}
       -\sum_{j=2}^r\kappa_{r,j}x^{-j}
       -\kappa\frac{e^{-x}}x.
\]
The singularity of $g_r$ at zero is removable.  The function extends to
a continuously differentiable function on $[0,\infty)$, and both $g_r$
and $g'_r$ are integrable there: at infinity the subtracted powers have
order at most $x^{-2}$, and all the other terms decay exponentially.
Let $C_r=\int_0^\infty g_r(x)\,dx$.

For each interval $[(n-1)t,nt]$, comparison with its right endpoint gives
\[
 \left|t g_r(nt)-\int_{(n-1)t}^{nt}g_r(x)\,dx\right|
 \le t\int_{(n-1)t}^{nt}|g'_r(x)|\,dx.
\]
Summation yields
$\sum_{n\ge1}g_r(nt)=C_r/t+O(1)$.
The removed terms sum exactly:
\[
 \sum_{n\ge1}(nt)^{-j}=t^{-j}\zeta(j),\qquad
 \sum_{n\ge1}\frac{e^{-nt}}{nt}
 =-\frac1t\log(1-e^{-t})=\frac Lt+O(1).
\]
This proves \eqref{p03:ci19:momentexpansion} and
\eqref{p03:ci19:Xcoordinate}.  No unproved estimate for the divisor problem
is required.
\end{proof}

\subsection{The prime coordinate has a divergent inverse-logarithmic series}

\begin{lemma}[All inverse-logarithmic orders]\label{p03:ci19:primeasymptotic}
For every fixed positive integer $N$,
\begin{equation}\label{p03:ci19:primeexpansion}
 Y(t):=tP(e^{-t})
 =\sum_{k=0}^{N-1}\frac{a_k}{L^{k+1}}+O_N(L^{-N-1}),
 \qquad a_k=(-1)^k\Gamma^{(k)}(1).
\end{equation}
The formal power series
$\widehat S(z)=\sum_{k\ge0}a_kz^{k+1}$ is transcendental over
$\mathbb C(z)$.
\end{lemma}

\begin{proof}
Write $\pi(x)$ for the prime-counting function and
$\operatorname{Li}_2(x)=\int_2^xdu/\log u$ for $x\ge2$.
The exponential error estimate in Trudgian's
Theorem~2 \cite{p03:ci19:Trudgian2016} implies, for every fixed $A>0$,
\[
 \pi(x)-\operatorname{Li}_2(x)=O_A(x/\log^A x).
\]
The normalization of the logarithmic integral changes its value only by
a constant.  Stieltjes integration and then integration by parts give
\begin{align*}
 tP(e^{-t})-
 t\int_2^\infty\frac{e^{-tx}}{\log x}\,dx
 &=t^2\int_2^\infty e^{-tx}
       \bigl(\pi(x)-\operatorname{Li}_2(x)\bigr)\,dx.
\end{align*}
On $2\le x\le t^{-1/2}$ the integral on the right is $O(t)$, using
$\pi(x)+|\operatorname{Li}_2(x)|=O(x)$.  On the rest of the interval,
$\log x\ge L/2$, so it is $O_A(L^{-A})$.  Changing variables $u=tx$
therefore gives
\begin{equation}\label{p03:ci19:primeintegral}
 Y(t)=\int_{2t}^\infty\frac{e^{-u}}{L+\log u}\,du
       +O_A(L^{-A}).
\end{equation}

The portion with $2t\le u\le t^{1/2}$ is $O(t^{1/2})$, since the
denominator is at least $\log2$.  The portion with $u\ge t^{-1/2}$ is
smaller than every power of $t$.  On the middle interval,
$|\log u|\le L/2$, and
\[
 \frac1{L+\log u}
 =\sum_{k=0}^{N-1}\frac{(-\log u)^k}{L^{k+1}}
   +\frac{(-\log u)^N}{L^N(L+\log u)}.
\]
The last term, after integration against $e^{-u}$, is
$O_N(L^{-N-1})$, because
$\int_0^\infty e^{-u}|\log u|^Ndu<\infty$.
Extending the coefficient integrals to $(0,\infty)$ introduces errors
smaller than every inverse power of $L$.  Taking $A=N+1$ in
\eqref{p03:ci19:primeintegral} proves \eqref{p03:ci19:primeexpansion}.

The function $\Gamma(1+w)-(1+w)^{-1}$ is analytic in $|w|<2$.
For every fixed $1<\rho<2$, Cauchy's coefficient estimate gives
\begin{equation}\label{p03:ci19:factorial}
 a_k=k!\bigl(1+O_\rho(\rho^{-k})\bigr).
\end{equation}
Thus $\widehat S$ has radius of convergence zero.  We recall why this
excludes algebraicity, rather than assuming any summability assertion.
Every formal complex power series algebraic over $\mathbb C(z)$ is
convergent.  Indeed, let an irreducible polynomial $Q(z,w)$ annihilate
a formal series $S(z)$.  In characteristic zero,
$Q_w(z,S(z))\ne0$; let its order be $m$.  For a sufficiently long
truncation $S_N$, with $N+1>m$, write
$S=S_N+z^{N+1}U$.  The equation
\[
 z^{-(N+1+m)}Q(z,S_N+z^{N+1}U)=0
\]
is analytic in $(z,U)$: its constant-in-$U$ term has no negative
powers, its coefficient of $U$ has a nonzero constant term, and its
higher coefficients vanish at $z=0$.  Its derivative in $U$ at the
formal constant term $U(0)$ is nonzero.  The analytic implicit function
theorem supplies a convergent solution, whose formal series is unique.
Hence $S$ converges.  Applying this fact to \eqref{p03:ci19:factorial}
proves the asserted formal transcendence.
\end{proof}

\begin{lemma}[The prime Lambert coordinate]\label{p03:ci19:primelambert}
As $t\downarrow0$,
\begin{equation}\label{p03:ci19:primelambertexpansion}
 Z(t):=tB(e^{-t})\sim\log L.
\end{equation}
\end{lemma}

\begin{proof}
Put $X=1/t$.  Uniformly for $p\le X$,
$(e^{p/X}-1)^{-1}=X/p+O(1)$.  For $p>X$ the summand is at most a
constant times $e^{-p/X}$, and comparison with all integers bounds
the tail by $O(X)$.  Therefore
\[
 B(e^{-1/X})=X\sum_{p\le X}\frac1p+O(X).
\]
Partial summation from $\pi(x)\sim x/\log x$ gives
$\sum_{p\le X}1/p\sim\log\log X$.  Dividing by $X$ proves the claim.
\end{proof}

\subsection{The independence theorem and its exact scope}

\begin{theorem}[Three-function independence over the cusp field]
\label{p03:ci19:triple}
For every $\lambda>0$ and every fixed integer $r\ge1$, the three germs
\[
 A_r(e^{-t}),\qquad P(e^{-t}),\qquad B(e^{-t})
\]
are algebraically independent over $\mathscr K_\lambda$.
Consequently, for every finite set $\mathcal D$ of positive integer
dilations, the three analytic functions $A_r(q),P(q),B(q)$ are
algebraically independent over the differential field generated by
$\mathbb C(q)$ and the functions
$\mathcal E_2(q^d),\mathcal E_4(q^d),\mathcal E_6(q^d)$,
$d\in\mathcal D$.
\end{theorem}

\begin{proof}
The affine changes \eqref{p03:ci19:Xcoordinate},
$Y=tP(e^{-t})$, and $Z=tB(e^{-t})$ have coefficients in
$\mathscr K_\lambda$ and nonzero slopes.  It is enough to prove
algebraic independence of $X=X_r,Y,Z$.

First consider a nonzero polynomial $R(U,V)\in\mathbb C[U,V]$.
Because $\widehat S$ is transcendental over $\mathbb C(z)$,
\[
 R(z^{-1},\widehat S(z))\ne0
\]
as a formal Laurent series.  It has a first nonzero term $c z^b$,
where $b\in\mathbb Z$ and $c\ne0$.  By taking sufficiently many terms
in \eqref{p03:ci19:primeexpansion} and using $X=L+O(t)$, we obtain
\begin{equation}\label{p03:ci19:polyleading}
 R(X(t),Y(t))\sim cL^{-b}.
\end{equation}
This step only uses a finite truncation for each specified polynomial;
no convergence of $\widehat S$ is asserted.

Now let $Q(U,V,W)\in\mathbb C[U,V,W]$ be nonzero, and write
$Q=\sum_jR_j(U,V)W^j$ with its nonzero coefficients $R_j$.
By \eqref{p03:ci19:polyleading},
$R_j(X,Y)\sim c_jL^{b_j}$ for suitable integers $b_j$ and
nonzero constants $c_j$.  Choose the largest $b_j$, and among the
indices attaining that largest value choose the largest $j$.
Since $Z\sim\log L$, the corresponding term uniquely dominates:
\begin{equation}\label{p03:ci19:Qleading}
 Q(X,Y,Z)\sim c_jL^{b_j}(\log L)^j.
\end{equation}
In particular, no such polynomial can vanish, or be smaller than every
inverse power of $L$.

Suppose, finally, that a nonzero polynomial with coefficients in
$\mathscr K_\lambda$ vanishes at $(X,Y,Z)$.  Write each nonzero
coefficient in the form \eqref{p03:ci19:leadingform}.  Among this finite
set choose the least exponential order $m$, and then the least power
$a$ among coefficients having that exponential order.  Divide the
polynomial identity by $t^ae^{-m\lambda/t}$.  The leading constants
of the coefficients with these two minimal orders give a nonzero
polynomial $Q\in\mathbb C[U,V,W]$.  All other terms, including the
errors in those selected coefficients, are bounded by $O(t)$ times
a fixed power of $L+\log L+1$, plus terms exponentially small in
$1/t$ times a fixed power of $t^{-1}$ and $L+\log L+1$.
Here we used $X=O(L)$, $Y=O(L^{-1})$, and $Z=O(\log L)$.
The identity would therefore force $Q(X,Y,Z)$ to be smaller than
every inverse power of $L$.  This contradicts
\eqref{p03:ci19:Qleading}.  The functional conclusion follows from
Proposition~\ref{p03:ci19:modularfield}; an analytic identity in the disk
would also hold for the corresponding germs.
\end{proof}

\begin{corollary}[Algebraic extensions do not remove the obstruction]
\label{p03:ci19:algebraicextensions}
The independence in Theorem~\ref{p03:ci19:triple} persists after replacing
$\mathscr K_\lambda$ by any algebraic extension contained in a common
field of germs.  Thus adjoining finitely many algebraic branches of
finite modular constructions does not supply a functional algebraic
relation among the three functions.
\end{corollary}

\begin{proof}
If $K'/\mathscr K_\lambda$ is algebraic, then
$K'(X,Y,Z)/\mathscr K_\lambda(X,Y,Z)$ is algebraic.
Transcendence degree is unchanged by algebraic extension, so
\[
 \operatorname{trdeg}_{K'}K'(X,Y,Z)
 =\operatorname{trdeg}_{\mathscr K_\lambda}
    \mathscr K_\lambda(X,Y,Z)=3.
\]
\end{proof}

There is a useful exact check on the proposed $q\mapsto-q$ route.
Directly from the coefficients $d(n)$ of $A_1(q)$,
\begin{equation}\label{p03:ci19:signidentity}
 A_1(q)+A_1(-q)=4A_1(q^2)-2A_1(q^4).
\end{equation}
Indeed $d(2n)=2d(n)-\mathbf1_{2\mid n}d(n/2)$, as is seen by
writing $n=2^am$ with $m$ odd.  Comparing the even coefficients proves
\eqref{p03:ci19:signidentity}.  It introduces further unknown Lambert
values and is not a modular expression for $A_1$.
For completeness, the analogous divisor-sum identity gives
\[
 \mathcal E_{2k}(-q)=
 2(1+2^{2k-1})\mathcal E_{2k}(q^2)
 -2^{2k}\mathcal E_{2k}(q^4)-\mathcal E_{2k}(q),
 \qquad k=1,2,3.
\]
Thus adding these negative-argument Eisenstein series still stays
inside a finite dilation field covered by the theorem.

The result leaves a precise arithmetic gap.  Functional algebraic
independence does not imply algebraic independence, or even individual
transcendence, at $q=1/2$.  For example, multiplying any functionally
independent family over $\mathbb C(q)$ by $q-1/2$ preserves functional
independence but makes all its specialized values zero.  A successful
proof for $E$, $P_2$, or $L_{\mathbb P}$ therefore still needs a
specialization theorem with verified hypotheses, or a direct
arithmetic approximation argument.  Neither is supplied here.
The theorem also does not exclude differential equations involving
derivatives of the target functions themselves; the derivatives allowed
in its coefficient field are derivatives of the modular generators.
Its contribution is a proved, explicit boundary on one proposed
transfer method, not a resolution of the three transcendence problems.


% Integrated source: conclusion.tex
\section{What has been resolved, and what would constitute the next step}
\label{p03:root3:conclusion}

The unconditional gains concern structured variational constants.
They include a full cubic-field classification, an identified independent
pair of observables, infinite independence in a fixed stochastic model,
all integer extraction shapes, and finite arithmetic descriptions of
signed entropy optima. Their proofs use classical transcendence inputs
in explicitly checked settings. They do not establish a new general
transcendence principle for arbitrary special-function values.

The conditional stopping result identifies a narrower target than a
blanket appeal to Schanuel's conjecture. For rational $a=r/s>0$ and
$b>0$, the function $M(a;b;z)^s-cz^r$ has only finitely many multiple
zeros. It would therefore suffice to prove that multiple zeros of this
specific $E$-function family, for algebraic $c\ne0$, are algebraic.
That restricted statement would already classify every rational bridge
payoff. The full squarefree-factorization conjecture is sufficient, but
it is not proved and is not known here to be necessary.

For the spectral program, the first concrete gap is the biquadratic
value in \eqref{p03:sc3:eq:quartic}. Its optimizing location and elasticity
are transcendental, but these facts do not determine its value. The
cubic proof explains the missing requirement: a lower bound of two for
the relevant exponential-line transcendence degree. Rational dimension
three alone does not provide the cubic-field theorem's hypothesis.

For fixed even-dimensional stopping, an additional question is whether
all even-payoff optima, beyond the prime-indexed or selected subsequence,
are algebraically independent, or which explicit relations their boundary
polynomials permit. The splitting-field argument used here makes no such
claim when the degree factor is already present in the earlier fields.
An answer would require more arithmetic of those specific polynomials.

For classical constants, the reduced-denominator and saturation theorems
show why normalization must be specified before an approximation is
called small. Positivity controls actual denominator-cleared forms in
the displayed support classes. After primitive reduction of the entire
coefficient pair, however, these positive cones realize every positive
primitive rational-approximation form. That universality supplies no new
small form and cannot be counted as an irrationality proof.

The prime/divisor functional theorem has a similarly precise boundary:
even algebraic independence over a large modular coefficient field does
not force transcendence at $q=1/2$. A valid specialization theorem or a
direct arithmetic argument is still needed. Replacing that missing step
by numerical evidence or a functional identity would change the result
into an unsupported claim.

\paragraph{Verification and priority.}
The main sections were independently cross-read within this project.
The checks covered the arbitrary-base hypothesis in Brownawell's theorem,
the curvature factorization argument, the Eisenstein coefficients and
splitting-field separation, the complex Kummer remainder sectors, the
third-order intertwiner and contact multiplicity, the entropy face and
cardinality arguments, the exact reduced-denominator prime windows, and
the cusp-field asymptotics. Exact small-parameter polynomial calculations
were used to check formulas, not to replace uniform proofs. The cited
external conjecture remains conditional. Historical priority is not
certified, and no newly solved named individual-constant conjecture is
reported.

% END SOURCE: KEIO_Arithmetic_Independence_and_Entropy_Optima.tex

\endgroup
\clearpage
\part{Rigidity of Optimization}
\label{merge:part04}
\begingroup

% BEGIN SOURCE: KEIO_Arithmetic_Rigidity_of_Optimization.tex

\noindent\textbf{Editorial integration note.} 
The all-planar independence theorem strengthens earlier subfamilies. The differential classifications and exponential-value closure results are further developed in Part~\ref{merge:part05}. Statements about functions do not automatically classify their values at a specified point.


\begin{quote}\small
\noindent\textit{Original part framing.}\par
\noindent\textbf{Status.} Research working manuscript. Proofs and
external inputs are separated, conditional implications are identified,
and internal cross-reading is not external peer review or formal
verification. Historical priority for the particular transfers has not
been established by an exhaustive literature review.
\end{quote}
\section{The question, the gain, and the remaining arithmetic gap}
\label{p04:root4:scope}

This paper continues \emph{Arithmetic Independence of Variational
Constants} \cite{p04:root4:III}, the preceding 45-page manuscript. The
223-page observation-and-recovery foundation and its collected sequels
are earlier inputs, rather than results reprinted or recounted here.
The research objective is to turn their analytic and identification
methods into proofs about the arithmetic of specific constants.
One must distinguish the arithmetic of an optimizer from that of its
value, and independence of functions from independence of values after
specialization. Both distinctions determine the successful transfers
and the remaining failures below.

\subsection{The most substantial increase in scope}

The leading result is Theorem~\ref{p04:ga20:main}. If
\[
 V_n=\sup_{0\le\tau\le1}\mathbb E\|B_\tau\|^{2n}
\]
for one standard planar Brownian bridge, then
\[
             \boxed{V_1,V_2,V_3,\ldots\text{ are algebraically independent}.}
\]
There is no nonzero polynomial with algebraic coefficients relating
any finite selection of these values. This is stronger than saying
that each is transcendental. The preceding manuscript proved the
corresponding statement only for $e$ and prime-indexed $V_p$;
the present theorem includes every composite index as well. Because
$V_1=2/e$, adjoining $e$ separately to the full sequence would create
a relation, and is deliberately not asserted.

The mechanism is also reusable. The boundary polynomial has degree
$n$ and a doubly transitive Galois action. Its centered roots generate
an irreducible rational representation of dimension $n-1$.
Different degrees therefore supply different representations even
when splitting fields intersect. This removes the need to prove
disjointness of those fields. Lindemann--Weierstrass then converts
linear independence of algebraic boundary parameters into algebraic
independence of optimal values.

In information theory, Theorem~\ref{p04:mi4:thm:IB} gives an exact
arithmetic transition. For rational $0<c<1$ and algebraic
$\beta>c^{-2}$, the unique positive optimal bias satisfies
\[
 \operatorname{arctanh}s_\beta
       =\beta c\operatorname{arctanh}(cs_\beta),\qquad
 s_\beta\in\overline{\mathbb Q}\ \Longleftrightarrow\ \beta\in\mathbb Q.
\]
For algebraic irrational $\beta$, every finite optimal encoder has
a transcendental entry. This is an exact classification in a familiar
optimization model; it is not a proof of transcendence of the optimized
objective itself. Its proof uses Gelfond--Schneider, and its analytic
binary bottleneck boundary is credited to the established literature.

The other changes are summarized without counting individual examples
as independent conjectures.

\small
\begin{longtable}{L{0.20\textwidth}L{0.32\textwidth}L{0.37\textwidth}}
\toprule Direction & Earlier result or limitation & Present result\\\midrule
\endfirsthead
\toprule Direction & Earlier result or limitation & Present result\\\midrule
\endhead
\bottomrule\endfoot
Spectral optimization & Arithmetic of selected values; a biquadratic value unresolved & $D$-finite exactly for rational exponent gap; exact nonlinear order and a complete derivative-field barrier\\[5pt]
Gamma extraction & Integer-shape contact arithmetic and double contact & Double contacts for all real shapes, piecewise analytic frontier and slope comparison; exact rank three at noninteger shape\\[5pt]
Contact observables & Contact and critical rate individually transcendental & At integer Bessel parameter, the contact exponential is also transcendental; the three quantities remain algebraically dependent\\[5pt]
Entropy values & Upper inclusion in logarithms of positive algebraic numbers & Exact converse realization by four signed entropy terms in a pure-saddle game\\[5pt]
$E$-value arithmetic & Rational-function intersection and nonunits & Exact rational-power membership; a finite-cover obstruction and an unconditional all-but-one relative-closure theorem\\[5pt]
Prime generating series & Functional algebraic independence & Joint independence of every derivative in the stated linear sense, over an enlarged coefficient field\\
\end{longtable}
\normalsize

\subsection{Twenty-one targets, with different kinds of outcome}

The register is an account of attempts, not a count of famous open
problems solved. ``Unresolved here'' reports this paper's conclusion.
Known irrationality results and results proved in the earlier manuscripts
are not added to the present success count.

\small
\begin{longtable}{L{0.035\textwidth}L{0.30\textwidth}L{0.55\textwidth}}
\toprule No. & Target & Outcome and precise scope\\\midrule
\endfirsthead
\toprule No. & Target & Outcome and precise scope\\\midrule
\endhead
\bottomrule\endfoot
1 & Euler's constant $\gamma$: irrationality & Unresolved. A nontrivial algebraic extension inside the specified $E$-value field would force $\gamma$ algebraic, but no such extension is constructed.\\[4pt]
2 & Euler--Gompertz $\delta$: irrationality & Unresolved. The new morphism and divisor argument does not determine this connection coefficient.\\[4pt]
3 & Erd\H{o}s--Borwein $\sum_{n\ge1}(2^n-1)^{-1}$: transcendence & Unresolved as a number. A functional coefficient-field enlargement is proved; known linear-form estimates do not settle arbitrary algebraic degree.\\[4pt]
4 & Prime binary $\sum_p2^{-p}$: transcendence & Unresolved as a number. The full derivative family is linearly independent over the stated cusp field with added moment and Lambert coefficients.\\[4pt]
5 & Prime Lambert $\sum_p(2^p-1)^{-1}$: transcendence & Unresolved as a number. Joint derivative independence and a prime-forced shift equation are established; external irrationality is credited.\\[4pt]
6 & Endpoint shape $A(1)$: transcendence & Unresolved. The rational noninteger-shape theorem does not permit an algebraic irrational varying Kummer parameter.\\[4pt]
7 & Endpoint shape $A(3/2)$: transcendence & Unresolved for the same parameter obstruction; the previous irrationality result is not recounted.\\[4pt]
8 & Linear excursion optimum $\mathcal V_{3,1}$ & Unconditional arithmetic unresolved. Mixed zero multiplicities and unbounded exponential denominators exclude a specified available partial factorization theorem.\\[4pt]
9 & Biquadratic spectral optimum $Q$ in \eqref{p04:sc4:eq:Q} & Value unresolved. Completing an exponential table adds an uncontrolled fourth power; every setting derivative remains in the first-derivative field.\\[4pt]
10 & $H(\operatorname{Poisson}(1))$: irrationality & Unresolved. Infinite logarithmic coefficients remain outside finite entropy and $E$-function arguments.\\[4pt]
11 & $\max_{x>0}\operatorname{Ein}(x)/x^\theta$, algebraic $0<\theta<1$ & Value unresolved. The maximizing location is transcendental, so algebraic-argument specialization cannot be applied there.\\[4pt]
12 & All planar even-payoff optimal values jointly & Algebraic independence of the entire sequence, Theorem~\ref{p04:ga20:main}. One infinite-family theorem.\\[4pt]
13 & Algebraic-weight binary bottleneck optimizer & Exact rational/algebraic-irrational classification above threshold, Theorem~\ref{p04:mi4:thm:IB}. Includes every finite output alphabet.\\[4pt]
14 & Range of rational pure-saddle entropy values & Exactly $\log\overline{\mathbb Q}_{>0}$, with four signed entropy terms, Theorem~\ref{p04:mi4:thm:range}.\\[4pt]
15 & Rational monomials in $\operatorname{Ein}(a)$ inside the $E$-value field & Necessary and sufficient integral-exponent criterion for the stated separated family, Theorem~\ref{p04:ac20:saturation}.\\[4pt]
16 & Algebraic extensions over $\overline{\mathbb Q}(\operatorname{Ein}(a))$ & Relative algebraic closure for all but at most one algebraic $a\ge1$, Corollary~\ref{p04:ac20:all-but-one}; the exception is not identified.\\[4pt]
17 & Arithmetic at rational noninteger extraction shapes & Normalized values and logarithmic derivatives at nonzero algebraic points are transcendental; zeros and physical stationary locations are transcendental under the stated rate hypotheses.\\[4pt]
18 & Fractional-shape critical rate itself & Unresolved. Exact rank three explains why two contact equations do not reproduce the integer-shape elimination.\\[4pt]
19 & $e^{2x}$ at integer extraction contacts & Transcendental and in $\overline{\mathbb Q}(x,c)$; the triple has transcendence degree one, Proposition~\ref{p04:be20:exponential}.\\[4pt]
20 & Two-atom optimal-value function and contact map & Complete $D$-finite classification and exact minimal nonlinear differential order, Theorems~\ref{p04:sc4:thm:dfinite}--\ref{p04:sc4:thm:order}.\\[4pt]
21 & All-real-shape extraction frontier & Exact double contact, locally finite analytic branches, signed slope formula and strict comparison law, Theorems~\ref{p04:be20:double} and~\ref{p04:be20:frontier}.\\
\end{longtable}
\normalsize

\subsection{Conventions and external inputs}

Throughout, algebraic and transcendental mean over $\mathbb Q$;
joint algebraic independence is equivalently stated over
$\overline{\mathbb Q}$. All real powers use the real logarithm
unless an analytic continuation is explicitly specified.
The notation $D$-finite refers to a homogeneous linear differential
equation with polynomial coefficients; it is different from satisfying
a nonlinear algebraic differential equation.

The stopping formula, Schur irreducibility, Jordan--Marggraff,
Lindemann--Weierstrass, Gelfond--Schneider, Beukers specialization,
and Lind's Perron realization are external inputs. The arithmetic-ring
section additionally uses a recent Fischler--Rivoal morphism theorem,
identified there as an external preprint result. The analytic existence,
continuity and uniform endpoint positivity of the gamma-extraction
capacity are inherited from the earlier project manuscript and stated
explicitly where used. All new deductions are given with proofs;
these dependency statements do not certify the earlier manuscripts or
external papers beyond the specific inputs checked.

For the two endpoint targets in the register, the earlier exact
definition is recalled here to remove a notation dependency.
For $d>1/2$, $A(d)$ is the unique value
\[
 a\in(\max(0,1-d),1/(4d))
\]
for which, with $y(z)=M(a-1;a+d-1;z)$ and $z=2(d-a)$,
\[
                     y(z)-\frac{z}{2(1-a)}y'(z)=0.
\]
Its existence, uniqueness and irrationality at rational $d$ are earlier
results, not proved anew here. The present question for $A(1),A(3/2)$
is their remaining transcendence classification.

\section{All planar even-payoff optima are algebraically independent}
\label{p04:ga20:section}

This section answers the self-posed arithmetic question for the entire
sequence of even-power stopping optima of a single planar Brownian
bridge.  The preceding paper treated prime-indexed powers.  Here every
positive integer index is included: the resulting optimal values are
jointly algebraically independent.  The argument combines Schur's
irreducibility theorem, a Newton-polygon calculation, the classical
Jordan--Marggraff theorem, and a representation-theoretic separation
of centered roots.  In particular, it requires no disjointness assertion
about the splitting fields.

A second target is the unconditional arithmetic of the linear
Brownian-excursion optimum.  That question remains unresolved.  We
prove two concrete obstructions to applying an available partial
$E$-function factorization theorem: the relevant functions have mixed
zero multiplicities and fail its global-boundedness hypothesis.  These
are limitations of a specified route, not a proof of impossibility of
the desired arithmetic conclusion.

\subsection{The stopping constants and the main result}

Let $R_t$ be the norm of a standard two-dimensional Brownian bridge
from zero to zero on $[0,1]$.  For $n\geq1$ set
\begin{equation}\label{p04:ga20:values}
 V_n=\sup_{0\leq\tau\leq1}\mathbb E[R_\tau^{2n}],
\end{equation}
where the supremum is over stopping times for its natural filtration.
The established stopping formula of Hobson--Liu \cite{p02:ga18:HL} gives
\begin{equation}\label{p04:ga20:stopping}
 f_n(z)=M(n;1;z),\qquad
 z_n f_n'(z_n)=n f_n(z_n),\qquad
 V_n=\frac{(2z_n)^n}{f_n(z_n)}.
\end{equation}
Here $z_n>0$ is the unique solution and the stopping boundary for the
radial process is $\sqrt{2z_n(1-t)}$.

For completeness, uniqueness of the positive solution follows directly
from the positive Taylor coefficients $c_j=(n)_j/(j!)^2$.
The probabilities $c_jz^j/f_n(z)$ have mean
$u(z)=zf_n'(z)/f_n(z)$ and satisfy
\[
 z u'(z)=\operatorname{Var}_z(j)>0.
\]
Their mean starts at zero and tends to infinity: for every fixed $K$,
the total weight on $j\leq K$ tends to zero by comparison with the
term of index $K+1$.  Thus $u(z)=n$ has exactly one positive solution.
The probabilistic optimality assertion itself is the cited stopping
theorem, rather than a new verification theorem here.

\begin{theorem}[The full planar sequence]\label{p04:ga20:main}
The numbers
\[
                         V_1,V_2,V_3,\ldots
\]
are algebraically independent over $\overline{\mathbb Q}$.
Equivalently, $e,V_2,V_3,\ldots$ are algebraically independent.
The associated squared boundary parameters satisfy
\begin{equation}\label{p04:ga20:linindep}
                  1,z_2,z_3,\ldots
       \quad\hbox{are linearly independent over }\mathbb Q.
\end{equation}
Every assertion about an infinite family means the corresponding
assertion for every finite subfamily.
\end{theorem}

Since $V_1=2/e$, the first formulation does not append $e$ to a
family already containing $V_1$.  In particular the ten consecutive
optimal values for payoff powers $2,4,6,\ldots,20$ are algebraically
independent.  This is one uniform theorem, not ten separate closures.

\subsection{Irreducibility of every boundary polynomial}

Kummer's transformation gives
\begin{equation}\label{p04:ga20:P}
 f_n(z)=e^zP_{n-1}(z),\qquad
 P_{n-1}(z)=M(1-n;1;-z)
          =\sum_{k=0}^{n-1}\binom{n-1}{k}\frac{z^k}{k!}.
\end{equation}
In particular $P_{n-1}(z)>0$ for $z>0$.  Define
\begin{equation}\label{p04:ga20:B}
 B_n(z)=(n-1)!\{zP_{n-1}'(z)+(z-n)P_{n-1}(z)\}.
\end{equation}
Its unique positive zero is $z_n$.

\begin{lemma}[Schur form of the stopping polynomial]\label{p04:ga20:Schur}
For every $n\geq1$, $B_n$ is a monic irreducible polynomial in
$\mathbb Z[z]$ of degree $n$.  More explicitly,
\begin{align}
 B_n(z)&=n!\sum_{k=0}^n a_{n,k}\frac{z^k}{k!},
       &a_{n,k}&=\binom{n-1}{k-1}-\binom{n-1}{k},
       \label{p04:ga20:coefficients}\\
 a_{n,0}&=-1,& a_{n,n}&=1.\nonumber
\end{align}
Out-of-range binomial coefficients are zero.  The sum of all roots
of $B_n$ is $n(2-n)$.
\end{lemma}
\begin{proof}
Taking coefficients in \eqref{p04:ga20:B} gives
\[
 \frac{(n-1)!}{k!}
 \left[k\binom{n-1}{k-1}+(k-n)\binom{n-1}{k}\right].
\]
The expression in brackets is
$n[\binom{n-1}{k-1}-\binom{n-1}{k}]$, proving
\eqref{p04:ga20:coefficients}.  Its integrality and monicity follow
immediately.  Schur's classical theorem says that
$\sum_{k=0}^n a_kz^k/k!$ is irreducible over $\mathbb Q$ when
$a_k\in\mathbb Z$ and $|a_0|=|a_n|=1$.  This exact formulation is
stated and generalized in Jindal--Khanduja
\cite[Theorem~1.1 and introduction]{p04:ga20:JK}.  It applies to the
coefficients above; multiplying by $n!$ does not affect
irreducibility.  Finally the coefficient of $z^{n-1}$ is
$n(n-2)$, giving the stated trace by Vieta's formula.
\end{proof}

For comparison with classical orthogonal polynomials,
\[
                  B_n(z)=n!\{L_n(-z)-2L_{n-1}(-z)\}.
\]
This identity follows from the same coefficient calculation.
Irreducibility of an individual Laguerre polynomial would not by
itself justify irreducibility of this particular linear combination;
the Schur form above is the needed input.

\subsection{Double transitivity from one local valuation}

\begin{lemma}[A long prime cycle]\label{p04:ga20:cycle}
Let $n\geq3$ and let $p$ be a prime with $n/2<p<n$.
The Galois group of $B_n$ contains a $p$-cycle.
\end{lemma}
\begin{proof}
Write $n=p+r$ with $1\leq r<p$.  Normalize the $p$-adic valuation
by $v_p(p)=1$.  The constant coefficient of $B_n$ is $-n!$,
of valuation one.  Every coefficient of degree $k<p$ has valuation
at least one, since $n!/k!$ is divisible by $p$.

For the coefficient of degree $p$, observe that, modulo $p$,
\[
 (1+X)^{p+r-1}\equiv(1+X^p)(1+X)^{r-1}.
\]
Consequently
$\binom{n-1}{p-1}\equiv0$ and
$\binom{n-1}{p}\equiv1\pmod p$, and hence
$a_{n,p}\equiv-1\pmod p$.  Since $n!/p!$ is a $p$-adic unit,
the coefficient of degree $p$ has valuation zero.  All other
coefficients are integral and the leading coefficient is one.

The lower Newton polygon therefore has a segment from $(0,1)$
to $(p,0)$ and then a horizontal segment to $(n,0)$.  The
Newton-polygon theorem gives roots of valuation $1/p$.  Thus the
ramification index of a splitting field over $\mathbb Q_p$ is
divisible by $p$.  A completion of the global splitting field has
that splitting property, so its local degree, and therefore the
order of the global Galois group, is divisible by $p$.
Cauchy's theorem gives an element of order $p$.  Because $n<2p$,
its permutation on the $n$ roots is a single $p$-cycle.
\end{proof}

\begin{proposition}[Double transitivity in every degree]
\label{p04:ga20:double}
For every $n\geq2$, the Galois group $G_n$ of $B_n$ acts doubly
transitively on its roots.
\end{proposition}
\begin{proof}
For $n=2$, $B_2(z)=z^2-2$, so the assertion is immediate.
For $n=3$ choose $p=2$.  For $n\geq4$, Bertrand's theorem
applied to $m=\lfloor n/2\rfloor\geq2$ gives a prime
$m<p<2m$, and this prime satisfies $n/2<p<n$.
Lemma~\ref{p04:ga20:cycle} supplies a $p$-cycle in $G_n$.
The group is transitive by Lemma~\ref{p04:ga20:Schur}.

It is also primitive.  Indeed, a nontrivial invariant block system
would have $s\geq2$ blocks of common size $t\geq2$, with $st=n$.
Both $s$ and $t$ are at most $n/2<p$.  The element of order $p$
cannot act nontrivially on the $s$ blocks.  Its nontrivial orbit of
size $p$ would therefore have to lie within one block of size
$t<p$, also impossible.

The classical Jordan--Marggraff theorem states that a primitive
permutation group containing a cycle with $k$ fixed points is
$(k+1)$-transitive.  An elementary proof is given by
Levingston--Taylor \cite[Theorem~A]{p04:ga20:LT}; Jones \cite[Section~3]{p04:ga20:Jones}
discusses its original attribution to Jordan.  Here $k=n-p\geq1$,
so $G_n$ is at least doubly transitive.  Only this classical
theorem is used, not a classification of finite simple groups.
\end{proof}

\subsection{Separating roots without separating splitting fields}

The following elementary representation-theoretic observation is
useful because the degrees of the polynomials need not be prime,
coprime, or ordered by ramification.

\begin{lemma}[Distinct degrees and centered root modules]
\label{p04:ga20:modules}
Let $F_1,\ldots,F_s\in\mathbb Q[z]$ be separable irreducible
polynomials of pairwise distinct degrees $d_i\geq2$, whose Galois
groups are doubly transitive.  Choose one root $\alpha_i$ of each.
Then $1,\alpha_1,\ldots,\alpha_s$ are linearly independent over
$\mathbb Q$.
\end{lemma}
\begin{proof}
Let $K$ be the compositum of the splitting fields and
$G=\operatorname{Gal}(K/\mathbb Q)$.  Its action on the $d_i$
roots of $F_i$ is doubly transitive, since restriction surjects
onto the Galois group of the corresponding splitting field.

Consider the rational permutation module $\mathbb Q^{d_i}$ and
its augmentation submodule
\[
 W_i=\{(v_1,\ldots,v_{d_i}):v_1+\cdots+v_{d_i}=0\}.
\]
This is an irreducible $\mathbb QG$-module.  Here is a direct
verification.  Double transitivity means that the commuting
matrices on $\mathbb Q^{d_i}$ are precisely the matrices constant
on the diagonal and constant off the diagonal, a two-dimensional
space.  The decomposition
$\mathbb Q^{d_i}=\mathbb Q\mathbf1\oplus W_i$ has no nonzero
homomorphisms between its two summands, because the only
invariant vectors and invariant linear forms in the permutation
module are constant.  Hence
$\dim_{\mathbb Q}\operatorname{End}_G(W_i)=1$.
Maschke's theorem then makes $W_i$ irreducible: a decomposition
into two nonzero submodules would provide two independent
commuting projections.

List the roots as $\alpha_{i,1},\ldots,\alpha_{i,d_i}$ and put
$\mu_i=d_i^{-1}\sum_j\alpha_{i,j}\in\mathbb Q$.
The map
\[
 (v_j)_j\longmapsto\sum_j v_j\alpha_{i,j}
              \quad\text{from }W_i\text{ into }K
\]
is $G$-equivariant and nonzero, since the roots are distinct.
It is therefore injective.  Its image $U_i$ is an irreducible
submodule of dimension $d_i-1$, and it contains
$\alpha_i-\mu_i$, the image of a suitable vector
$e_j-d_i^{-1}\mathbf1$.

The modules $U_i$ have different dimensions and are thus pairwise
nonisomorphic.  None is the trivial module: $W_i$ has no nonzero
invariant vector.  This also deals with degree two, when
$\dim U_i=1$.  Semisimplicity now gives the direct sum
\[
             \mathbb Q\cdot1\ \oplus\ U_1\ \oplus\cdots\oplus\ U_s
                         \ \subseteq K.
\]
For clarity, directness follows by induction: an irreducible
submodule intersecting a direct sum of nonisomorphic irreducible
submodules would embed into that sum, and some coordinate
projection would make it isomorphic to a summand, a contradiction.

Each $\alpha_i-\mu_i$ is nonzero, since an irreducible polynomial
of degree at least two has no rational root.  A rational linear
relation among $1,\alpha_i$ can be rewritten using these centered
roots.  Directness forces every coefficient to vanish.
\end{proof}

\begin{proof}[Proof of Theorem~\ref{p04:ga20:main}]
Apply Lemma~\ref{p04:ga20:modules} to any finite selection of the
polynomials $B_n$, $n\geq2$, using Proposition~\ref{p04:ga20:double}.
This proves \eqref{p04:ga20:linindep}.  In this application the
centered root is particularly explicit:
\[
                         z_n-\mu_n=z_n+n-2.
\]
For $n=1$, $z_1=1$.  Hence all the $z_n$, $n\geq1$, are
linearly independent over $\mathbb Q$.

Equations \eqref{p04:ga20:stopping}--\eqref{p04:ga20:P} give
\begin{equation}\label{p04:ga20:algebraicfactor}
 V_n=A_n e^{-z_n},\qquad
 A_n=\frac{(2z_n)^n}{P_{n-1}(z_n)}
                \in\overline{\mathbb Q}^{\times}.
\end{equation}
A nonzero polynomial relation over $\overline{\mathbb Q}$ among
finitely many $V_n$ would therefore give a nonzero algebraic
linear combination of exponentials
\[
                     \exp\!\left(-\sum_n k_n z_n\right).
\]
Distinct monomials have distinct algebraic exponents, by the
linear independence just proved.  Lindemann--Weierstrass rules
out this linear combination.  Finally $P_0=1$ and $z_1=1$
give $V_1=2/e$, proving the equivalent formulation.
\end{proof}

The first boundary polynomials are
\[
\begin{split}
 B_1(z)&=z-1,\qquad B_2(z)=z^2-2,\\
 B_3(z)&=z^3+3z^2-6z-6,\\
 B_4(z)&=z^4+8z^3-48z-24.
\end{split}
\]
The new statement includes $V_4,V_6,V_8,V_9,V_{10}$ and all
other composite-indexed optima together with the preceding prime
family.  Merely observing that every $B_n$ is irreducible would
not prove this conclusion; the centered-module step controls
possible relations across their splitting fields.

\subsection{Target II: a precise obstruction for the excursion value}

The linear Brownian-excursion optimum is the established constant
\begin{equation}\label{p04:ga20:excursion}
 V=\mathcal V_{3,1}=2\sqrt{2z_*}\,e^{-z_*},\qquad
 2f(z_*)=e^{z_*},\qquad
 f(z)=M(1/2;3/2;z)=\int_0^1e^{zt^2}\,dt.
\end{equation}
The preceding arithmetic analysis proves that $z_*$ is
transcendental and that a published squarefree-factorization
conjecture would imply transcendence of $V$
\cite{p04:root4:III,p04:ga20:FR}.  We do not count those earlier
results again.  We investigate whether the particularly simple
first-order identity
\begin{equation}\label{p04:ga20:firstorder}
                       2zf'(z)+f(z)=e^z
\end{equation}
allows the known unconditional factorization theorem to apply.

\begin{proposition}[Two failed hypotheses in the excursion family]
\label{p04:ga20:obstruction}
For $c\in\mathbb C^\times$, put
\begin{equation}\label{p04:ga20:G}
                        G_c(z)=f(z)^2-cz.
\end{equation}
Every multiple zero of $G_c$ has multiplicity exactly two.
There are only finitely many such zeros and infinitely many
simple zeros.  If $c\in\overline{\mathbb Q}^{\times}$, then
$G_c$ is an $E$-function but is not globally bounded in the
exponential-coefficient sense used by Fischler--Rivoal.
\end{proposition}
\begin{proof}
Since $G_c(0)=1$, a multiple zero $w$ is nonzero.  Its two
defining equations give
\[
 f(w)^2=cw,\qquad 2wf'(w)=f(w),\qquad e^w=2f(w).
\]
In particular $f(w)\ne0$.  The Kummer equation
\[
        zf''+(3/2-z)f'-\tfrac12 f=0
\]
then gives the exact identity at this point
\begin{equation}\label{p04:ga20:secondderivative}
                       G_c''(w)=\frac{c(2w-1)}{w}.
\end{equation}
The only candidate for a zero of multiplicity at least three
is $w=1/2$.  But $f(1/2)>1$ whereas $e^{1/2}<2$, so
$2f(1/2)=e^{1/2}$ is impossible.  Every multiple zero is
therefore exactly double.

Finiteness of multiple zeros is the special case already proved
in the preceding paper; the short calculation here records its
global analytic content.  On the closed upper and lower
half-planes separately, the two-term Kummer expansion
\cite[\S13.7]{p03:ga19:DLMF}, with the corresponding branch of
$z^{-1/2}$, is
\[
 f(z)=\frac{e^z}{2z}\{1+O(|z|^{-1})\}
             +B_\pm z^{-1/2}\{1+O(|z|^{-1})\},
 \qquad B_\pm=\pm\frac{i\sqrt\pi}{2}.
\]
The error bounds apply separately to the two displayed terms.
At an unbounded sequence of contacts $e^z=2f(z)$ this implies
\[
 e^z=2B_\pm z^{-1/2}\{1+O(|z|^{-1})\},\qquad
 \frac{f(z)^2}{z}=-\frac{\pi}{4z^2}\{1+O(|z|^{-1})\}\longrightarrow0.
\]
The additional condition $f(z)^2/z=c\ne0$ is thus impossible
outside a fixed disk.  Zeros of the nonzero entire function
$G_c$ have no finite accumulation point, proving finiteness.

The integral defining $f$ gives $|f(z)|\leq e^{|z|}$, so $G_c$
has order at most one.  Laplace's endpoint estimate on the
positive real axis gives
\[
 f(x)=\frac{e^x}{2x}\{1+O(x^{-1})\},\qquad
 G_c(x)=\frac{e^{2x}}{4x^2}\{1+O(x^{-1})\}.
\]
If $G_c$ had only finitely many zeros, Hadamard factorization
would give $G_c=P(z)e^{az+b}$ with a nonzero polynomial $P$.
Comparison of logarithmic absolute values along $x>0$ first
forces $\Re a=2$ and then $\deg P=-2$, a contradiction.
There are infinitely many zeros, and the finite-multiple-zero
conclusion makes infinitely many of them simple.

For the arithmetic assertion, write
\[
 f(z)=\sum_{n\geq0}\frac{1}{2n+1}\frac{z^n}{n!},\qquad
 f(z)^2=\sum_{n\geq0} b_n\frac{z^n}{n!},\qquad
 b_n=\sum_{k=0}^n\binom nk
            \frac{1}{(2k+1)(2n-2k+1)}.
\]
This rational-parameter Kummer function is an $E$-function;
the closure of $E$-functions under products and algebraic
polynomials gives the same assertion for $G_c$.
In the terminology of \cite{p04:ga20:FR}, global boundedness would
require an integer $D\geq1$ for which $D^{n+1}$ times every
exponential Taylor coefficient is an algebraic integer.

Take an arbitrary prime $p\geq5$ and $n=(p-1)/2\geq2$.
The endpoint terms $k=0,n$ in $b_n$ have sum $2/p$.
All remaining denominators are prime to $p$, because both
their positive odd factors lie strictly between $1$ and $p$.
Thus
\begin{equation}\label{p04:ga20:denominators}
                         v_p(b_{(p-1)/2})=-1.
\end{equation}
Subtracting $cz$ changes only the coefficient of index one.
For any proposed $D$, choose a prime $p\geq5$ not dividing $D$.
The coefficient in \eqref{p04:ga20:denominators} remains nonintegral
after multiplication by $D^{n+1}$.  A rational number is an
algebraic integer exactly when it is an integer, so this
contradicts global boundedness, even when $c$ is algebraic
but not rational.
\end{proof}

If $V=v$ were algebraic, then $c=2/v^2$ would be algebraic and
$z_*$ would be a double zero of $G_c$.  Proposition
\ref{p04:ga20:obstruction} shows that this function has both simple
and double zeros, and that it is not globally bounded.
Fischler--Rivoal's unconditional Theorem~1.5 assumes global
boundedness and that all zeros have the same multiplicity
$m\geq2$, in addition to an order restriction on a differential
equation.  Hence that theorem fails two independently verified
hypotheses here.  The first-order identity \eqref{p04:ga20:firstorder}
does not repair either of them.

The remaining arithmetic target can be stated without referring
to any factorization terminology.  Transcendence of $V$ is
equivalent to the following assertion: for every
$c\in\overline{\mathbb Q}^{\times}$, the two entire functions
\[
                   e^z-2f(z),\qquad e^{2z}-4cz
\]
have no common positive real zero.  The unique positive zero
of the first function is $z_*$; at this point the second
equation is exactly $V^2=2/c$.  Algebraic-point specialization
theorems for $E$-functions do not settle this common-zero
question, because $z_*$ is already transcendental.

\paragraph{Scope and provenance.}
The main new result relative to the preceding manuscript is
the full, consecutive planar family in Theorem~\ref{p04:ga20:main},
including every composite payoff index.  The root-module lemma
and its application are proved explicitly; Schur irreducibility,
Bertrand's theorem, the Newton-polygon theorem,
Jordan--Marggraff, Maschke, and Lindemann--Weierstrass remain
classical inputs.  A search of primary literature on Laguerre
Galois groups, Schur-type irreducibility, and bridge stopping
did not verify a prior arithmetic statement for this exact
optimal-value sequence.  This search does not establish a
historical priority claim.

The excursion proposition is an unconditional diagnosis of an
unsuccessful route, not an unconditional irrationality or
transcendence theorem for $V$.  General rational nonexceptional
payoff optima likewise remain outside the unconditional value
classification obtained here.  No named historical conjecture
is claimed to have been closed by this section.

\section{The range of entropy-game values and an arithmetic bottleneck transition}
\label{p04:mi4:sec}

Two questions go beyond classifying extrema of rational entropy
programs.  First, which logarithms of algebraic numbers actually occur
as finite entropy-game values?  Second, does an algebraic optimizer
still exist when a rational information-bottleneck weight is replaced
by an algebraic irrational one?  We answer the first question exactly
for games with a pure saddle, and the second exactly for the binary
symmetric model above its nontriviality threshold.  The constructions
use established results of Lind and Gelfond--Schneider.  They do not
resolve the arithmetic of Poisson entropy or the optimized
exponential-integral payoff discussed at the end.

The number-theoretic representation used below is already known in
a stronger form: Dubickas proves that every positive algebraic number
is a quotient of two Pisot numbers in its own number field
\cite[Theorem 3]{p04:mi4:Dubickas}.  Lind's realization of Perron numbers
as spectral radii is also classical \cite[Theorem 1]{p04:mi4:Lind}.
Our exact entropy-game formulation combines these facts with the
previous rational-entropy stationary-value theorem.  We make no
historical priority claim for this formulation.  The binary
information-bottleneck boundary is likewise established
\cite[Section IV-A]{p03:mi3:Hsu}; the arithmetic distinction between
rational and algebraic irrational weights is the additional
consequence proved here.

All logarithms are natural, and $0\log0=0$.  Write
$\overline{\mathbb Q}_{>0}$ for the positive real algebraic numbers.

\subsection{An exact converse for finite entropy games}

A \emph{rational entropy functional} on a nonempty compact rational
polytope $P$ is
\begin{equation}\label{p04:mi4:eq:entropy}
 F(z)=\sum_{a=1}^{A}c_a H(L_a(z)),\qquad
 H(p)=-\sum_j p_j\log p_j,
\end{equation}
where $c_a\in\mathbb Q$ and every $L_a$ is a rational affine map
from $P$ to a finite probability simplex.  A point is
\emph{face-stationary} when it is in the relative interior of a
face and the derivative vanishes in every direction parallel to
that face.  This definition includes vertices.

We recall the stationary-value fact from the preceding work, with
its proof so that the converse below has explicit hypotheses.

\begin{lemma}[Previously established upper inclusion]
\label{p04:mi4:lem:upper}
Every face-stationary value of \eqref{p04:mi4:eq:entropy} belongs to
$\log\overline{\mathbb Q}_{>0}$.  For fixed rational data there are
only finitely many such values, each attained at an algebraic
face-stationary point.
\end{lemma}

\begin{proof}
On the relative interior of a face $Q$, each affine probability
coordinate is either identically zero or positive.  Delete the zero
coordinates and choose a rational tangent basis $v_i$.  Since each
$L_a$ has total mass one, stationarity is exactly
\[
 \sum_{a,j}c_a(\partial_{v_i}\ell_{aj})
                         \log\ell_{aj}(z)=0.
\]
Clearing denominators and exponentiating turns these conditions
into polynomial equations, with positive arguments, over
$\mathbb Q$.  The stationary set is therefore semialgebraic over
$\mathbb Q$.  A finite semialgebraic cell decomposition supplies
algebraic sample points and piecewise $C^1$ paths within each cell
\cite{p03:mi3:BPR}.  The derivative of $F$ along every such path is zero,
so $F$ is constant on each cell.  Every stationary value is thus
attained at an algebraic point $z_*$.

Choose a rational point $z_0$ in the relative interior of $Q$.
Pair stationarity at $z_*$ with $z_*-z_0$.  Affineness yields
\[
 F(z_*)=-\sum_{a,j}c_a\ell_{aj}(z_0)\log\ell_{aj}(z_*).
\]
The coefficients are rational and the arguments positive algebraic.
Exponentiation gives a positive algebraic number.  There are
finitely many cells and finitely many faces, proving the assertion.
\end{proof}

A \emph{rational entropy game} has strategy polytopes $P,Q$ of the
same rational kind and payoff $F(x,y)$ of the form
\eqref{p04:mi4:eq:entropy} on $P\times Q$.  A pure saddle is a pair
$(x_*,y_*)$ such that
\begin{equation}\label{p04:mi4:eq:saddle}
 F(x,y_*)\le F(x_*,y_*)\le F(x_*,y)
 \quad (x\in P,\ y\in Q).
\end{equation}
The maximizing player chooses $x$ and the minimizing player chooses
$y$.  We do not assume that every rational entropy game has such a
pair.

\begin{theorem}[Exact range of pure-saddle entropy values]
\label{p04:mi4:thm:range}
The following three sets of real numbers are equal:
\begin{enumerate}
\item $\log\overline{\mathbb Q}_{>0}$;
\item the values at pure saddles of all rational entropy games;
\item the face-stationary values of all rational entropy functionals.
\end{enumerate}
Moreover, every number in this common set is realized by a separable
concave--convex game with exactly four signed entropy terms,
\begin{equation}\label{p04:mi4:eq:four}
 F(x,y)=H(x)-H(\pi(x))-H(y)+H(\pi(y)),
\end{equation}
on two rational directed-graph flow polytopes.  The constructed game
has an algebraic pure saddle in the relative interior.  The sizes
of the graphs are allowed to depend on the algebraic number.
\end{theorem}

We give the representation argument in two steps.  The first is
weaker than Dubickas's Pisot quotient theorem, and is included only
as a short direct proof of the precise fact required here.

\begin{lemma}[Perron quotient representation; known consequence]
\label{p04:mi4:lem:quotient}
For every $\alpha\in\overline{\mathbb Q}_{>0}$ there are Perron
numbers $\beta,\gamma>1$ in $\mathbb Q(\alpha)$ with
$\alpha=\gamma/\beta$.  Here a Perron number is a positive algebraic
integer strictly larger in absolute value than all its other
conjugates.
\end{lemma}

\begin{proof}
Put $K=\mathbb Q(\alpha)$ and choose a positive integer $D$ such
that $D\alpha$ is an algebraic integer.  In the Minkowski space
$\mathbb R^{r_1}\times\mathbb C^{r_2}$ of $K$, the image of the
ideal $D\mathcal O_K$ is a full lattice.  Take the first real
coordinate to be the given embedding $K\subset\mathbb R$.
A bounded fundamental parallelepiped implies that for every real
$T>0$ there is $\beta_T\in D\mathcal O_K$ with
\[
 \beta_T=T+O(1),\qquad
 |\sigma(\beta_T)|=O(1)\quad(\sigma\ne\mathrm{id}),
\]
where the constants are independent of $T$.  Indeed, express
$(T,0,\ldots,0)$ in a real lattice basis and round each coefficient
to an integer.  The resulting error lies in a fixed bounded
parallelepiped; both real coordinates of each complex embedding
are bounded.

Both $\beta_T$ and $\gamma_T=\alpha\beta_T$ are algebraic
integers: writing $\beta_T=D b_T$ gives
$\gamma_T=(D\alpha)b_T$ with $b_T\in\mathcal O_K$.
At the distinguished embedding they tend to $+\infty$, while
every other embedding of either number remains bounded.  For
sufficiently large $T$, both distinguished values exceed one
and strictly dominate all other embedding values in modulus.
Every conjugate of an element of $K$ is obtained by an embedding
of $K$, so these two numbers are Perron.  Their quotient is
$\alpha$.  This proves the weaker form needed here; the stronger
Pisot conclusion is \cite[Theorem 3]{p04:mi4:Dubickas}.
\end{proof}

Let $A$ be a primitive nonnegative integral matrix and regard
$A_{ij}$ as the number of directed edges from $i$ to $j$, retaining
parallel edges as distinct elements of the edge set $E_A$.
Its stationary flow polytope is
\begin{equation}\label{p04:mi4:eq:flow}
 P_A=\left\{x\ge0:\ \sum_{e\in E_A}x_e=1,\quad
 \sum_{s(e)=i}x_e=\sum_{t(e)=i}x_e\ \text{for every }i\right\}.
\end{equation}
Write $\pi_i(x)=\sum_{s(e)=i}x_e$ and
$h_A(x)=H(x)-H(\pi(x))$.  These are rational affine probability
maps and a rational entropy functional.  The following is the
finite-state entropy variational principle, traditionally associated
with Shannon and Parry; see the historical account in
\cite[Section 1]{p04:mi4:Lind}.  We supply its elementary form.

\begin{lemma}[Graph entropy maximum]
\label{p04:mi4:lem:flow}
For $\lambda=\rho(A)$,
\[
 \max_{x\in P_A}h_A(x)=\log\lambda.
\]
The maximum has a positive algebraic maximizing flow.  The
functional $h_A$ is concave.
\end{lemma}

\begin{proof}
Choose positive right and left Perron eigenvectors $r,l$, so
$Ar=\lambda r$ and $l^{\mathsf T}A=\lambda l^{\mathsf T}$.
They can be chosen algebraic: the simple eigenvalue $\lambda$ is
algebraic and its eigenvectors are solutions of linear equations
over the real algebraic numbers.  Give each edge $e:i\to j$ the
transition probability
\[
 p_e=\frac{r_j}{\lambda r_i}.
\]
The probabilities leaving each vertex sum to one.  For any
$x\in P_A$ and every $i$ with $\pi_i(x)>0$, put
$q_e=x_e/\pi_i(x)$ for outgoing edges.  Nonnegativity of relative
entropy, followed by flow conservation, gives the exact identity
\begin{align*}
 h_A(x)
 &=-\sum_e x_e\log p_e
       -\sum_{i:\pi_i(x)>0}\pi_i(x)D(q_i\Vert p_i)\\
 &=\log\lambda
       -\sum_{i:\pi_i(x)>0}\pi_i(x)D(q_i\Vert p_i)
 \le\log\lambda.
\end{align*}
In the second equality, the incoming and outgoing averages of
$\log r_i$ cancel.  The positive algebraic flow
\[
 x_e^*=\frac{l_i r_j}{\lambda\sum_k l_k r_k}
       \quad(e:i\to j)
\]
has both incoming and outgoing masses
$\pi_i^*=l_i r_i/\sum_k l_k r_k$, has total mass one, and has
$q_i=p_i$.  It attains equality.  Finally, conditional entropy is
concave in the joint edge law: each vertex contribution
$-\sum_{s(e)=i}x_e\log(x_e/\pi_i)$ is the perspective of the
concave entropy function.  Summing proves concavity.
\end{proof}

\begin{proof}[Proof of Theorem~\ref{p04:mi4:thm:range}]
A pure saddle is face-stationary on the product of the minimal
faces containing its two strategies: each of the two partial
derivatives vanishes along that face by \eqref{p04:mi4:eq:saddle}.
Their tangent spaces span the product tangent space.  Thus (2)
is contained in (3), and Lemma~\ref{p04:mi4:lem:upper} gives (3)
contained in (1).

Conversely fix $\alpha>0$ algebraic.  Use
Lemma~\ref{p04:mi4:lem:quotient} to write $\alpha=\gamma/\beta$ with
$\gamma,\beta>1$ Perron.  Lind's realization theorem gives
primitive nonnegative integral matrices $A,B$ with spectral radii
$\gamma,\beta$, respectively \cite[Theorem 1]{p04:mi4:Lind}.
On $P_A\times P_B$ set $F(x,y)=h_A(x)-h_B(y)$.
If $x_*,y_*$ are the maximizing flows of
Lemma~\ref{p04:mi4:lem:flow}, then
\[
 F(x,y_*)\le\log\gamma-\log\beta
            \le F(x_*,y).
\]
This is a pure saddle of value $\log\alpha$.  Its coordinates
are algebraic and positive, so it is in the relative interior of
the product flow polytope.  Formula \eqref{p04:mi4:eq:four} and
concavity of the two conditional entropies prove all the remaining
claims.
\end{proof}

\begin{remark}[What the converse does and does not say]
\label{p04:mi4:rem:range}
The result is an exact range statement, with no additive constants
inserted into the objective.  In particular, every logarithm of a
positive algebraic number occurs using only four entropy terms
with coefficients $1,-1,-1,1$ and rational linear constraints.
Every nonzero value in this class is transcendental by
Hermite--Lindemann.  This is a consequence of established arithmetic
and entropy theory, not a newly solved classical transcendence
conjecture.  We have not shown that every such logarithm is the
global maximum of a \emph{single} rational entropy functional;
the second player supplies a minimum.  No conclusion is asserted
for arbitrary min--max games without a pure saddle.
\end{remark}

\subsection{An exact arithmetic transition for binary information bottleneck}

Let $X$ be fair binary and let $Y=X\oplus Z$, where $Z$ is
independent Bernoulli with rational parameter
$0<\delta<1/2$.  Set $c=1-2\delta\in(0,1)\cap\mathbb Q$.
For $\beta>0$, consider the usual information-bottleneck objective
\begin{equation}\label{p04:mi4:eq:IB}
 V_c(\beta)=\inf_{T-X-Y}\{I(X;T)-\beta I(Y;T)\},
\end{equation}
where $T$ ranges over arbitrary finite alphabets and arbitrary
stochastic encoders.  The model and its binary symmetric optimum
are classical \cite{p03:mi3:IB,p03:mi3:Hsu}.  The next result classifies
the arithmetic of its optimizers, including encoders with more
than two output symbols.

\begin{theorem}[Algebraic weights and transcendental optimal encoders]
\label{p04:mi4:thm:IB}
Let $\beta>0$ be algebraic.
\begin{enumerate}
\item If $\beta\le c^{-2}$, then $V_c(\beta)=0$ and a rational
      encoder attains it.  Every optimal encoder makes $X$ and $T$
      independent.
\item If $\beta>c^{-2}$, there is a unique $s_\beta\in(0,1)$
      satisfying
      \begin{equation}\label{p04:mi4:eq:critical}
       \operatorname{arctanh}s_\beta
          =\beta c\operatorname{arctanh}(cs_\beta).
      \end{equation}
      A binary symmetric encoder with crossover
      $\varepsilon_\beta=(1-s_\beta)/2$ attains the infimum.
      In every finite optimal encoder, each positive-probability
      output has posterior probability
      \[
       \mathbb P(X=1\mid T=t)
        \in\{\varepsilon_\beta,1-\varepsilon_\beta\}.
      \]
      The total output probability assigned to each of these two
      posterior values is $1/2$.
\item In the second case, $s_\beta$ and $\varepsilon_\beta$ are
      algebraic if and only if $\beta$ is rational.  Consequently
      an optimal encoder with all entries algebraic exists if
      and only if $\beta$ is rational.  If $\beta$ is algebraic
      irrational, every finite optimal encoder has at least one
      transcendental entry.
\end{enumerate}
\end{theorem}

\begin{proof}
Write $h(p)=-p\log p-(1-p)\log(1-p)$.  For an arbitrary finite
encoder let $w_t=\mathbb P(T=t)$ and
$p_t=\mathbb P(X=1\mid T=t)$, omitting outputs of zero probability.
The two marginal laws $X,Y$ are fair and
$\mathbb P(Y=1\mid T=t)=\delta+cp_t$.  Hence its objective is
\begin{equation}\label{p04:mi4:eq:posterior}
 (1-\beta)\log2+
       \sum_t w_t\{\beta h(\delta+cp_t)-h(p_t)\}.
\end{equation}
The expression in braces is symmetric under $p_t\mapsto1-p_t$.
Its minimum over $[0,1]$ is therefore attainable under the prior
constraint $\sum_t w_t p_t=1/2$: take its two symmetric minimizing
arguments with weights $1/2$.  This is exactly a binary symmetric
encoder (or an independent encoder if the minimum is at $p=1/2$).
Conversely, an average cannot lie below its pointwise minimum.
It follows that \eqref{p04:mi4:eq:IB} is attained and equals the minimum
over $0\le s\le1$ of
\begin{equation}\label{p04:mi4:eq:scalar}
 f_\beta(s)=\beta h((1-cs)/2)-h((1-s)/2)
                         -(\beta-1)\log2.
\end{equation}
Furthermore, equality in \eqref{p04:mi4:eq:posterior} is possible only
when every $p_t$ is a pointwise minimizer.

For $0<s<1$,
\[
 f_\beta'(s)=\operatorname{arctanh}s
               -\beta c\operatorname{arctanh}(cs).
\]
The function
\[
 b_c(s)=\frac{\operatorname{arctanh}s}
                  {c\operatorname{arctanh}(cs)}
\]
is strictly increasing from $c^{-2}$ to $+\infty$.
Here is a direct proof of strictness.  Put $u=s^2$ and
$a_n=1/(2n+1)$.  Then
\[
 \frac{\operatorname{arctanh}(cs)}{\operatorname{arctanh}s}
 =\frac{\sum_{n\ge0}a_n c^{2n+1}u^n}
        {\sum_{n\ge0}a_n u^n}.
\]
For $0<u<1$, differentiate this quotient.  After multiplying by
the square of its denominator, pair the terms with indices
$m>n$.  Their sum is
\[
 a_m a_n(m-n)(c^{2m+1}-c^{2n+1})u^{m+n-1}<0.
\]
The series and their derivatives converge absolutely on compact
subintervals, so this computation is valid.  The quotient is
strictly decreasing, proving the claim about $b_c$.  Its endpoint
limits follow from $\operatorname{arctanh}s\sim s$ at zero and
its divergence at one.

If $\beta\le c^{-2}$, then $f_\beta'(s)>0$ for every $s>0$,
so the unique minimizing bias is $s=0$ and $f_\beta(0)=0$.
All optimal posteriors therefore equal $1/2$, which is independence.
If $\beta>c^{-2}$, the equation $b_c(s)=\beta$ has exactly one
solution $s_\beta\in(0,1)$.  The derivative of $f_\beta$ is
negative below that point and positive above it.  Thus the only
minimizing posterior values are
$\varepsilon_\beta,1-\varepsilon_\beta$.  The condition that their
weighted average is $1/2$, and their distinctness, force total
weights $1/2$ and $1/2$.

Exponentiating \eqref{p04:mi4:eq:critical} gives
\begin{equation}\label{p04:mi4:eq:power}
 \frac{1+s_\beta}{1-s_\beta}
  =\left(\frac{1+cs_\beta}{1-cs_\beta}\right)^{\beta c},
\end{equation}
where the power has its positive real meaning.  If $\beta$ is
rational, write $\beta c=a/b$ with positive integers $a,b$.
Then $s_\beta$ is a root of
\[
 (1+s)^b(1-cs)^a-(1-s)^b(1+cs)^a=0.
\]
This is a nonzero rational polynomial, since at $s=1$ its value is
$2^b(1-c)^a>0$.  Thus $s_\beta$ and the optimal binary encoder
are algebraic.

If $\beta$ is algebraic irrational and $s_\beta$ were algebraic,
both ratios in \eqref{p04:mi4:eq:power} would be algebraic numbers
greater than one, while $\beta c$ would be algebraic irrational.
The Gelfond--Schneider theorem makes the right side transcendental,
a contradiction \cite{p03:mi3:Baker}.  Consequently $s_\beta$ and
$\varepsilon_\beta$ are transcendental.  In any finite encoder
with all entries algebraic, both $w_t$ and the positive-probability
posteriors are algebraic, because the input law is rational and
Bayes' formula uses only field operations.  The posterior
classification already proved therefore rules out every such
optimal encoder.  This proves the theorem.
\end{proof}

For a concrete instance, take $\delta=1/4$ and $\beta=4\sqrt2$.
The unique positive optimal bias solves
\[
 \operatorname{arctanh}s=2\sqrt2\operatorname{arctanh}(s/2).
\]
It is transcendental, as is the optimal binary crossover
$(1-s)/2$, and no enlargement of the finite output alphabet permits
an entirely algebraic optimal encoder.  The example is an instance
of a classified family, rather than a separately counted constant.

\begin{remark}[Optimizer arithmetic is not value arithmetic]
\label{p04:mi4:rem:IBvalue}
At the nontrivial stationary point, the elementary identity
\[
 h((1-s)/2)=\log2-\tfrac12\log(1-s^2)
                         -s\operatorname{arctanh}s
\]
and \eqref{p04:mi4:eq:critical} give
\begin{equation}\label{p04:mi4:eq:value}
 V_c(\beta)=\tfrac12\log(1-s_\beta^2)
                  -\tfrac\beta2\log(1-c^2s_\beta^2).
\end{equation}
For rational $\beta$, the algebraic optimizer makes the exponential
of this value algebraic.  For algebraic irrational $\beta$ the
arguments in \eqref{p04:mi4:eq:value} are transcendental; the proof above
does not classify the arithmetic of the value itself.  In
particular, transcendence of an optimizer is not a proof of
transcendence of its optimum value.
\end{remark}

\subsection{Two unresolved transfers}

The Shannon entropy of a Poisson law with intensity one is
\[
 H(\operatorname{Poisson}(1))
       =1+e^{-1}\sum_{n=0}^{\infty}\frac{\log(n!)}{n!}.
\]
The sum converges, but its coefficients $\log(n!)$ are not
algebraic coefficients of an $E$-function.  Neither the finite
entropy stationary-value theorem nor the exponential-integral
independence theorem applies to this series.  Finite truncation
followed by convergence supplies no arithmetic closure principle:
limits of algebraic or transcendental values can have either
arithmetic type.  The bounds of Adell, Lekuona, and Yu
\cite{p04:mi4:Poisson} concern analytic approximation and do not supply
the missing transcendence argument.  We obtain no irrationality
or transcendence result for this individual entropy.

Likewise, for $0<\theta<1$ algebraic, the previously studied maximum
of $\operatorname{Ein}(x)/x^\theta$ has a transcendental maximizing
location.  Siegel--Shidlovskii controls suitable $E$-function values
at algebraic arguments; it cannot simply be evaluated at that
transcendental optimizer.  No arithmetic classification of this
optimized payoff follows here.  These two failed transfers are
recorded as remaining questions, not as solved targets.

\section{Radical saturation and the only possible finite branch point}
\label{p04:ac20:section}

This section distinguishes a ring-membership problem from a
transcendence problem.  Put $k=\overline{\mathbb Q}$, let $\mathcal E$
be the ring of values of $E$-functions at algebraic points, and put
$K=\operatorname{Frac}(\mathcal E)$.  We prove two structural results:
an exact classification of rational monomials in a separated family of
exponential-integral values, and a restriction on every finite algebraic
extension of $k(\operatorname{Ein}(1))$ occurring inside $K$.
These are applications of the external morphism theorem of
Fischler--Rivoal \cite[Proposition 3.1]{p03:ac19:FRring}.  The paper's
nonunit criterion and discussion of factorization are the relevant
antecedents; the statements here are not attributed to that paper.
No independent priority claim is made for these applications.

The logarithmic transcendence target is deliberately not counted again.
The earlier theorem on algebraic independence of distinct
$\operatorname{Ein}$-values over the field generated by algebraic
exponentials already implies the transcendence of every nonzero rational
linear combination of their logarithms.  Neither that conclusion nor
its one-value special case is a new result in this section.

\subsection{Input morphism and its divisor interpretation}

We use the notation
\[
 \delta_a=\int_0^\infty\frac{e^{-t}}{t+a}\,dt,
 \qquad u_a=\operatorname{Ein}(a)
       =\int_0^a\frac{1-e^{-t}}t\,dt \quad(a>0).
\]
The connection formula and the morphism computation established in the
preceding division analysis are
\begin{equation}\label{p04:ac20:morphism}
 u_a=\gamma+\log a+e^{-a}\delta_a,
 \qquad
 \Phi_0(u_a)(z)=A_a+B_a e^{-az},
 \qquad A_a=\gamma+\log a,\quad B_a=\delta_a.
\end{equation}
Here $a$ is algebraic, and
\[
 \Phi_0:\mathcal E\hookrightarrow
 \mathcal R=\left\{\sum_{\rho\in k}c_\rho e^{\rho z}:
                   c_\rho\in\mathbb C,\ \text{finite support}\right\}
\]
fixes $k$.  It extends injectively to $K\to\operatorname{Frac}(\mathcal R)$.
The external theorem constructs $\Phi_0$ from the canonical large-sector
asymptotic decomposition.  In the present case this decomposition is
$\operatorname{Ein}(az)=\gamma+\log a+\log z+e^{-az}\delta_{az}$;
indeed $\delta_{az}\sim\sum_{n\ge0}(-1)^n n!/(az)^{n+1}$ on the
required sectors.  Evaluating the coefficient functions at $z=1$ gives
\eqref{p04:ac20:morphism}.  This recalls the input rather than counting its
proof as another result.

For $a\ge1$, $A_a,B_a>0$.  All zeros of $\Phi_0(u_a)$ are simple,
and have real part
\begin{equation}\label{p04:ac20:height}
 c(a)=\frac1a\log\frac{B_a}{A_a}.
\end{equation}
Their imaginary parts are odd integer multiples of $\pi/a$.
The entire-versus-meromorphic distinction is useful: images of elements
of $\mathcal E$ have nonnegative integral orders at finite points,
whereas images of elements of $K$ have arbitrary integral orders.

\begin{lemma}[A uniform separation range]\label{p04:ac20:separation}
The function $c(a)$ is strictly increasing for real $a\ge3$.
Consequently, the zero sets of $\Phi_0(u_a)$ and $\Phi_0(u_b)$ are
disjoint whenever $a,b\ge3$ are distinct.  Independently of this range,
for $a,b\ge1$ an irrational ratio $a/b$ also ensures disjoint zero sets.
\end{lemma}
\begin{proof}
Let $X$ have density $e^{-x}$ on $[0,\infty)$.  The exact identity
\[
 \frac1{a+x}=\frac1{a+1}-\frac{x-1}{(a+1)^2}
           +\frac{(x-1)^2}{(a+1)^2(a+x)}
\]
and Cauchy--Schwarz give
\[
 \delta_a=\frac1{a+1}+
 \frac1{(a+1)^2}\mathbb E\frac{(X-1)^2}{a+X}
 \ge\frac1{a+1}+\frac1{(a+1)^2(a+3)}
 =\frac{(a+2)^2}{(a+1)^2(a+3)}.
\]
Here $\mathbb E(X-1)^2=1$ and
$\mathbb E[(a+X)(X-1)^2]=a+3$; the inequality is the weighted
Cauchy--Schwarz inequality applied to $X-1$.
Also $\delta_a<1/a$ and differentiation followed by integration by parts
gives $\delta_a'=\delta_a-1/a$.
Write $A=\gamma+\log a$ and $F=\log(A/\delta_a)$, so that $c=-F/a$.
Then
\[
 aF'-F=\frac1A+\frac1{\delta_a}-a-\log(A/\delta_a)
 < H(a):=\frac1A+1-\frac{a+1}{(a+2)^2}-\log(aA).
\]
A direct differentiation yields
\[
 H'(a)=\frac{a}{(a+2)^3}
       -\frac1a\left(1+\frac1A+\frac1{A^2}\right)<0.
\]
The elementary bounds $\gamma>1/2$ and $\log3>13/12$ imply
$A|_{a=3}>19/12$.  Since $\log(19/4)>3/2$,
\[
 H(3)<\frac{12}{19}+\frac{21}{25}-\frac32<0.
\]
Thus $c'(a)=(F-aF')/a^2>0$ for $a\ge3$.
For explicit rational certificates for the small numerical bounds,
$e<11/4$ and $e^{1/12}<12/11$ give $e^{13/12}<3$, while
$(11/4)^{3/2}<19/4$.  Furthermore
$\gamma>H_6-\log7$ and
$\sum_{j=0}^{6}(39/20)^j/j!>7$ give $\gamma>1/2$.
Finally, a common zero for $a,b\ge1$ would have imaginary parts
$(2j+1)\pi/a=(2\ell+1)\pi/b$, forcing $a/b$ rational.
\end{proof}

\begin{theorem}[Exact rational-power saturation]\label{p04:ac20:saturation}
Let $a_1,\ldots,a_m\ge1$ be positive algebraic numbers such that the
zero sets of $\Phi_0(u_{a_j})$ are pairwise disjoint.  For arbitrary
$r_1,\ldots,r_m\in\mathbb Q$, take the positive real powers and put
$v=\prod_{j=1}^m u_{a_j}^{r_j}$.  Then
\begin{align}
 v\in K&\quad\Longleftrightarrow\quad r_j\in\mathbb Z
                                      \quad(1\le j\le m),
                                      \label{p04:ac20:field-saturation}\\
 v\in\mathcal E&\quad\Longleftrightarrow\quad
                        r_j\in\mathbb Z_{\ge0}
                                      \quad(1\le j\le m).
                                      \label{p04:ac20:ring-saturation}
\end{align}
In particular, both equivalences hold for arbitrary distinct algebraic
$a_j\ge3$, with no rational-independence assumption on the parameters.
For one parameter they hold for every algebraic $a\ge1$.
\end{theorem}
\begin{proof}
Choose $q\ge1$ with $m_j=qr_j\in\mathbb Z$.  If $v\in K$, the identity
$v^q=\prod_j u_{a_j}^{m_j}$ can be passed through the extended morphism.
At a zero $z_j$ of $\Phi_0(u_{a_j})$, no other factor vanishes.  Taking
orders of meromorphic functions gives
\[
                  q\,\operatorname{ord}_{z_j}\Phi_0(v)=m_j.
\]
The order on the left is an integer multiple of $q$, hence $r_j$ is an
integer.  If $v\in\mathcal E$, its image is entire and this integer is
nonnegative.  Conversely, integral powers belong to the fraction field,
and nonnegative integral powers belong to the ring.  The separation
lemma proves the stated parameter range.  For $m=1$ there is no
pairwise separation condition.
\end{proof}

Thus, for example, $u_3^{1/2}$ is outside even $K$, while $u_3^{-1}$
is in $K$ but outside $\mathcal E$.  Mixed values such as
$u_3^{2/3}u_4^{-5/7}$ are classified by exactly the same criterion.
This is an arithmetic membership theorem for a self-posed family,
not a new irrationality proof for $\gamma$ or $\delta$.
The range $a\ge3$ is a proved sufficient condition, not a claimed
optimal boundary or a classification of every possible collision in
$1\le a<3$.

\subsection{Rational exponential functions and finite algebraic covers}

The next argument is not limited to radicals.  It concerns every
algebraic element of $K$ over $k(u)$, where $u=u_1$.

\begin{lemma}[One-frequency algebraic descent]\label{p04:ac20:descent}
Let $W\in\operatorname{Frac}(\mathcal R)$ be algebraic over
$\mathbb C(e^{-z})$.  There exists an integer $q\ge1$ such that
$W\in\mathbb C(e^{-z/q})$.
\end{lemma}
\begin{proof}
Choose a $\mathbb Q$-basis $1,\beta_2,\ldots,\beta_s$ of the
$\mathbb Q$-span of $1$ and the finitely many phases occurring in a
numerator and denominator of $W$.  After clearing their rational
coordinates with a positive integer $q$, all these exponential
functions are Laurent monomials in
\[
 X=e^{-z/q},\qquad Y_j=e^{\beta_jz/q}\quad(2\le j\le s).
\]
These functions are algebraically independent over $\mathbb C$:
a polynomial relation would be a linear relation among finitely many
exponentials with distinct exponents, contradicted by differentiating
at $0$ and the Vandermonde determinant.  Hence $W$ is a rational
function in the independent variables $X,Y_2,\ldots,Y_s$.
It is algebraic over $\mathbb C(X)$ because it is algebraic over the
smaller field $\mathbb C(X^q)$.

In characteristic zero a field $F$ is relatively algebraically closed
in $F(Y_2,\ldots,Y_s)$.  Here is a short proof in the needed form.
Differentiate a separable minimal polynomial for an algebraic element
$h$ over $F$ with respect to each $Y_j$; its nonvanishing derivative
forces $\partial h/\partial Y_j=0$.  The only rational functions in a
variable with zero derivative are independent of that variable, and
iteration gives $h\in F$.  Apply this with $F=\mathbb C(X)$.
\end{proof}

\begin{theorem}[The only possible finite branch point]
\label{p04:ac20:branch}
Let $v\in K$ be algebraic over $k(u)$, and let
$d=[k(u,v):k(u)]$.  If $d>1$, then $\gamma\in k$.
More precisely, after extension of constants to $\mathbb C$, the
corresponding cover of the $U$-line is cyclic of degree $d$, is totally
ramified at $U=\gamma$ and at infinity, and is unramified elsewhere.
Consequently, under the explicit hypothesis $\gamma\notin k$,
\begin{align}
 \{v\in K:v\text{ is algebraic over }k(u)\}&=k(u),
                                                 \label{p04:ac20:relclosed}\\
 \{v\in\mathcal E:v\text{ is algebraic over }k(u)\}&=k[u].
                                                 \label{p04:ac20:integral}
\end{align}
The hypothesis in these last two equalities is not established here.
\end{theorem}
\begin{proof}
The image $U=\Phi_0(u)=\gamma+\delta e^{-z}$ is nonconstant; since
$\Phi_0$ fixes $k$ and is injective, $u$ is transcendental.  Choose a
primitive irreducible polynomial $P(U,V)\in k[U,V]$ for $v$ over
$k(u)$.  Because $k$ is algebraically closed, $P$ remains irreducible
over $\mathbb C(U)$.  One elementary justification is Gauss's lemma
followed by descent of a hypothetical polynomial factorization:
the normalized coefficients of factors satisfy finitely many
polynomial equations over $k$; a solution over $\mathbb C$ gives a
solution over $k$ by the weak Nullstellensatz.  Nonzero normalizing
coefficients are encoded by adjoining their inverses.

Apply $\Phi_0$ to the relation.  With $W=\Phi_0(v)$, it gives
$P(\gamma+\delta e^{-z},W)=0$.  The descent lemma puts
$W$ in $\mathbb C(X)$ for $X=e^{-z/q}$ and some $q\ge1$, with
$U=\gamma+\delta X^q$.  Irreducibility shows that
\[
 [\mathbb C(U,W):\mathbb C(U)]=d.
\]
The extension $\mathbb C(X)/\mathbb C(U)$ is cyclic Galois, with
$X\mapsto\zeta X$ for $\zeta^q=1$.  Every subgroup is cyclic, and the
fixed field of its subgroup of order $h$ is $\mathbb C(X^h)$.
Thus
\[
 \mathbb C(U,W)=\mathbb C(X^h),\qquad h\mid q,\qquad d=q/h.
\]
Writing $Y=X^h$, the cover is $U=\gamma+\delta Y^d$.
If $d>1$, it is totally ramified at $\gamma$ and infinity and
unramified elsewhere, as follows directly by differentiating or taking
local orders.

All finite branch points of the cover defined by $P$ are algebraic.
Indeed, outside the zeros of the leading coefficient of $P$ in $V$
and its nonzero $V$-discriminant, its roots are locally distinct
holomorphic functions by the implicit function theorem.  These two
polynomials have coefficients in $k$.  Since $\gamma$ is a branch
point when $d>1$, it must be algebraic.  This proves the unconditional
implication and \eqref{p04:ac20:relclosed}.

For \eqref{p04:ac20:integral}, assume $\gamma\notin k$.  We have just
shown that every relevant $v\in\mathcal E$ is $R(u)$ for
$R\in k(U)$.  If $R$ had a finite pole $\alpha\in k$, then
$\alpha\ne\gamma$ and
$R(\gamma+\delta e^{-z})$ would have a pole at each solution of
$e^{-z}=(\alpha-\gamma)/\delta$.  The derivative of the inner function
there is $-(\alpha-\gamma)\ne0$.  This contradicts the entire image
$\Phi_0(v)$.  Hence $R$ is polynomial.  The reverse inclusion is immediate.
\end{proof}

The implication in Theorem~\ref{p04:ac20:branch} is useful even without
knowing whether $\gamma$ is algebraic: an algebraic curve with two
distinct finite branch values, or with a finite branch value outside
$(0,1)$, cannot acquire a nonrational solution at $U=u$ inside $K$.
There is no assertion that an arbitrary complex algebraic extension of
$k(u)$ embeds into this fraction field.

\begin{corollary}[An unconditional all-but-one family]
\label{p04:ac20:all-but-one}
Among all positive algebraic parameters $a\ge1$, at most one can fail
one of the following equalities:
\begin{align*}
 \{v\in K:v\text{ is algebraic over }k(u_a)\}&=k(u_a),\\
 \{v\in\mathcal E:v\text{ is algebraic over }k(u_a)\}&=k[u_a].
\end{align*}
In particular, for any prescribed $m$ distinct such parameters, both
equalities hold for at least $m-1$ of them.  This does not identify a
possible exceptional parameter, nor assert its existence.
\end{corollary}
\begin{proof}
Repeat the proof of Theorem~\ref{p04:ac20:branch} with $e^{-az}$ in place
of $e^{-z}$ and $A_a+B_ae^{-az}$ in place of $\gamma+\delta e^{-z}$.
In the descent lemma start the rational basis with $a$, giving
$X=e^{-az/q}$.  Every step is unchanged.  Hence either equality can
fail only if $A_a=\gamma+\log a$ is algebraic: for the second equality,
the degree-one pole argument also requires a pole at $A_a$.
If distinct $a,b$ both had algebraic centers, then
$\log(a/b)=A_a-A_b$ would be a nonzero algebraic number whose
exponential is algebraic.  This contradicts the
Hermite--Lindemann theorem \cite{p04:ac20:Lindemann}.  Thus there is at
most one possible exception in the entire parameter family.
\end{proof}

\begin{corollary}[Explicit radical exclusions inside the fraction field]
\label{p04:ac20:polynomial-radicals}
Let $R\in k(U)^\times$, let $p\in\mathbb Z$, $q\ge1$ be coprime,
and suppose that $R$ has a zero or pole at some $\alpha\in k$ with
$\alpha\ne\gamma$ and order $m\ne0$.  If a number $v$ satisfies
$v^q=R(u)^p$ and belongs to $K$, then $q$ divides $m$.
In particular, for every $n\ge1$ and every nonintegral $r\in\mathbb Q$,
\begin{equation}\label{p04:ac20:explicit-roots}
                    (u^n+1)^r\notin K.
\end{equation}
All choices of the complex root have the same nonmembership conclusion.
\end{corollary}
\begin{proof}
At a solution $z_0$ of $\gamma+\delta e^{-z_0}=\alpha$, the inner
function is unramified.  The order of $R(\gamma+\delta e^{-z})$ there
is exactly $m$.  Applying the extended morphism to the power identity
and taking orders gives $q\operatorname{ord}_{z_0}\Phi_0(v)=pm$.
Coprimality forces $q\mid m$.  For $R(U)=U^n+1$, all roots are simple
and none equals $\gamma\in(0,1)$, so $q=1$ would be necessary.
Multiplication by a root of unity does not change membership because
all roots of unity belong to $k\subset\mathcal E$.
\end{proof}

\subsection{What these results settle and what remains}

The two new structural targets are the exact saturation
Theorem~\ref{p04:ac20:saturation} and the finite-cover obstruction
Theorem~\ref{p04:ac20:branch}; Corollaries~\ref{p04:ac20:all-but-one} and~\ref{p04:ac20:polynomial-radicals}
are dependent applications, the former removing the hypothesis for all
but at most one parameter in the entire family.  They address division and
algebraic extraction in the arithmetic ring and its fraction field,
which algebraic independence of a preassigned finite family does not
settle.  The explicit examples $\sqrt{u^n+1}$, $1\le n\le10$, are ten
instances of a single theorem, not ten independent famous conjectures.
Their transcendence already follows from that of $u$; the new property
proved here is nonmembership in $K$.

The attempted individual targets $\gamma$ and $\delta$ remain separated
from these claims.  A nontrivial algebraic extension in $K$ would force
$\gamma$ algebraic, but no such extension is constructed, and no
contradiction from $\gamma$ algebraic is proved.  The positive coefficient
$\delta$ controls the exponential coordinate in the cover; its presence
does not determine whether it is rational or algebraic.  Finally, the
old logarithmic target $\log(\gamma+\delta/e)$ is already settled by the
earlier algebraic-independence argument, so it provides no extra advance
in this round.

\section{The differential complexity of a spectral optimum}
\label{p04:sc4:sec}

The preceding continuation classified cubic-field critical values and
left the biquadratic value
\begin{equation}\label{p04:sc4:eq:Q}
 Q=\min_{x>0}\frac{2(1+x)}
                 {x^{1/\sqrt3}+x^{1/\sqrt2}}
\end{equation}
unresolved.  We investigated that value again.  The conclusion remains
open here; we identify precisely why completing a biquadratic
exponential table does not supply the needed contradiction.

The successful direction concerns the entire optimal-value function.
For every two-atom positive spectrum, we classify exactly when that
function satisfies a linear differential equation with polynomial
coefficients.  We also obtain an explicit nonlinear differential
equation, prove its converse, and determine its least possible
algebraic differential order.  These conclusions include
\eqref{p04:sc4:eq:Q} as a specialization without making an unwarranted
inference about its individual arithmetic value.

\subsection{The parametric curve of values and contacts}

In the functional results below, the exponents and weights may be
arbitrary real numbers satisfying
\begin{equation}\label{p04:sc4:eq:twoatom}
 0<a<b<1,\qquad w_1,w_2>0,\qquad
 S(x)=w_1x^a+w_2x^b,\qquad
 T(C)=\min_{x>0}\frac{1+Cx}{S(x)}\quad(C>0).
\end{equation}
Write $d=b-a$, let $x(C)$ be the unique contact, and put
$H(x)=S(x)-xS'(x)$.  Positive real powers use the real logarithm.
The objective tends to infinity at zero and infinity, and
\[
 \left(\frac{S'}H\right)'=\frac{S''S}{H^2}<0.
\]
These facts give a unique minimum and an analytic contact map with
$x'(C)<0$.  Its value curve has the parametrization
\begin{equation}\label{p04:sc4:eq:param}
 C=\Phi(x):=\frac{S'(x)}{H(x)},\qquad
 T=\Psi(x):=\frac1{H(x)}.
\end{equation}
In particular
\begin{align}
 T'&=\frac{x}{S(x)}>0,&
 B:=T-CT'&=\frac1{S(x)}>0,&
 x&=\frac{T'}B,
 \label{p04:sc4:eq:envelope}\\
 T''&<0,&
 x'&=\frac{TT''}{B^2},&
 S''(x)&=\frac{B^3}{T^3T''}.
 \label{p04:sc4:eq:second}
\end{align}
Here and below primes on $T$ denote differentiation in $C$ and
primes on $S$ denote differentiation in $x$.

We use the standard term \emph{$D$-finite} for an analytic germ that
satisfies a nonzero linear differential equation with polynomial
coefficients, as in \cite{p04:sc4:Stanley}.  All function classifications
are over $\mathbb C(C)$.  They therefore exclude such equations even
when arbitrary complex constants are allowed as coefficients.
Algebraic differential order means the smallest $k\ge0$ for which
a nonzero polynomial over $\mathbb C(C)$ vanishes at
$(T,T',\ldots,T^{(k)})$.

\begin{lemma}[Algebraicity of an optimum]\label{p04:sc4:lem:algebraic}
For a finite sum $S(x)=\sum_jw_jx^{r_j}$ with distinct real
$0<r_j<1$ and positive weights, the analytic germ of its optimized
value is algebraic over $\mathbb C(C)$ if and only if every $r_j$
is rational.
\end{lemma}
\begin{proof}
If $T$ is algebraic, so is $T'$, and the envelope formula makes
$x(C)=T'/(T-CT')$ algebraic.  This function is nonconstant.
Its local inverse $C(x)$ is consequently algebraic, and the identity
$S(x)=1/[T(C(x))-C(x)T'(C(x))]$ makes $S$ algebraic over
$\mathbb C(x)$.

An algebraic germ has finitely many branches.  Continuing the displayed
finite power sum around zero a suitable positive number $N$ of times
must therefore return to the original branch.  Thus
\[
 \sum_jw_j\bigl(e^{2\pi iNr_j}-1\bigr)x^{r_j}=0.
\]
Distinct real powers are linearly independent on an interval: after
writing $x=e^t$, the Vandermonde matrix obtained by successive
derivatives proves this assertion.  Hence $Nr_j\in\mathbb Z$ for
every $j$, proving necessity.  Conversely, rational exponents make
$S,H,\Phi,\Psi$ algebraic in $x$.  The nonconstant local inverse
of $\Phi$ and its composition with $\Psi$ are algebraic.  This
proves sufficiency.
\end{proof}

\subsection{A complete linear-differential classification}

\begin{theorem}[Rational gaps are exactly the $D$-finite case]
\label{p04:sc4:thm:dfinite}
For \eqref{p04:sc4:eq:twoatom}, the following statements are equivalent:
\begin{enumerate}
\item the gap $d=b-a$ is rational;
\item the contact germ $x(C)$ is algebraic over $\mathbb C(C)$;
\item the contact germ $x(C)$ is $D$-finite;
\item the optimized-value germ $T(C)$ is $D$-finite.
\end{enumerate}
In addition, $T(C)$ is algebraic if and only if both $a$ and $b$
are rational.  If $d$ is irrational, analytic continuations of the
real germs of both $T$ and $x$ have infinitely many distinct finite
branch points.  At the explicit branch points constructed below,
the contact has exponent $1/2$ and the value has exponent $1/2$
when $a+b\ne1$; for $a+b=1$ these exponents are respectively
$1/3$ and $2/3$.
\end{theorem}
\begin{proof}
First suppose $d=p/q$ for positive integers $p,q$.  Set
$z=x^{1/q}$.  Stationarity becomes the polynomial equation
\begin{equation}\label{p04:sc4:eq:rationalgap}
 Cz^q\bigl((1-a)w_1+(1-b)w_2z^p\bigr)
              =aw_1+bw_2z^p.
\end{equation}
It is nonzero as a polynomial in $C,z$, so $z(C)$ and $x(C)$
are algebraic.  The value has the form
\[
 T=\frac{z^{-qa}}{(1-a)w_1+(1-b)w_2z^p}.
\]
Its logarithmic derivative belongs to the finite algebraic extension
$L=\mathbb C(C,z)$, which is closed under differentiation.
Consequently $T^{(k)}=R_kT$ with $R_k\in L$ for all $k\ge0$.
If $[L:\mathbb C(C)]=e$, then $R_0,\ldots,R_e$ are linearly
dependent over $\mathbb C(C)$.  Clearing denominators gives a
nonzero linear differential equation with polynomial coefficients
for $T$.  The same finite-dimensional argument proves that every
algebraic function, in particular $x(C)$, is $D$-finite.

Now let $d$ be irrational.  Work on the logarithmic parameter plane
$x=e^t$, on which all powers in \eqref{p04:sc4:eq:param} are single
valued and $\Phi(e^t),\Psi(e^t)$ are meromorphic.  The zeros of
$S(e^t)$ are
\begin{equation}\label{p04:sc4:eq:zeros}
 t_n=\frac{\log(w_1/w_2)+(2n+1)\pi i}{d},\qquad
 x_n=e^{t_n},\qquad n\in\mathbb Z.
\end{equation}
At such a point, put $A_n=w_1x_n^a\ne0$, with the branch determined
by $t_n$.  Then
\begin{align}
 S(x_n)&=0,&
 S'(x_n)&=-\frac{dA_n}{x_n},&
 H(x_n)&=dA_n\ne0,
 \label{p04:sc4:eq:zerodata}\\
 \Phi(x_n)&=c_n=-x_n^{-1},&
 S''(x_n)&=\frac{d(1-a-b)A_n}{x_n^2}.
 \label{p04:sc4:eq:criticalvalues}
\end{align}
The values $c_n$ are distinct: equality for $n\ne m$ would require
$(n-m)/d$ to be an integer and thus make $d$ rational.

The parameter derivatives are
\begin{equation}\label{p04:sc4:eq:paramderivatives}
 \Phi'(x)=\frac{S''(x)S(x)}{H(x)^2},\qquad
 \Psi'(x)=\frac{xS''(x)}{H(x)^2}.
\end{equation}
If $a+b\ne1$, the zero of $S$ is simple and $S''(x_n)\ne0$.
Writing $h=x-x_n$ in its local branch gives
\begin{equation}\label{p04:sc4:eq:sqrt}
 \Phi(x)-c_n=\alpha_n h^2+O(h^3),\qquad
 \Psi(x)-\Psi(x_n)=\beta_n h+O(h^2),
 \quad \alpha_n\beta_n\ne0.
\end{equation}
If $a+b=1$, then instead
\[
 S''(x_n)=0,\qquad S'''(x_n)=\frac{ab\,dA_n}{x_n^3}\ne0.
\]
Thus \eqref{p04:sc4:eq:paramderivatives} gives
\begin{equation}\label{p04:sc4:eq:cuberoot}
 \Phi(x)-c_n=\widetilde\alpha_n h^3+O(h^4),\qquad
 \Psi(x)-\Psi(x_n)=\widetilde\beta_n h^2+O(h^3),
 \quad \widetilde\alpha_n\widetilde\beta_n\ne0.
\end{equation}
These are genuine branch points of the locally inverted curve.
For example, if the value were holomorphic in $C$ at $c_n$,
its nonconstant composition with $\Phi$ would have vanishing order
a multiple of two in \eqref{p04:sc4:eq:sqrt}, or a multiple of three
in \eqref{p04:sc4:eq:cuberoot}; the displayed orders are respectively
one and two.  The contact has order one in $h$ in either case,
and is likewise ramified.

It remains to show that these branches are reachable from the
positive real germ; a list of critical values alone would not suffice.
Suppose either germ satisfied a linear equation, and let $\Sigma$
be its finite set of possible finite singularities, the zeros of
the leading polynomial coefficient.  Choose $c_n\notin\Sigma$
and a positive initial setting outside $\Sigma$.  In the $t$-plane
remove the poles of $\Phi(e^t)$ and $\Psi(e^t)$, the zeros of
$d\Phi(e^t)/dt$, and the inverse image of $\Sigma\cup\{c_n\}$.
This is a discrete set; all functions involved are meromorphic and
$\Phi$ is nonconstant.  Its complement is path connected.  A path
from the real initial parameter can therefore approach $t_n$, with
only its endpoint in the removed set.  Along the path, the local
inverse of $\Phi$ provides analytic continuation of the original
contact and value germs.  The image path avoids $\Sigma$.
Near its endpoint the continuations are precisely the ramified
branches in \eqref{p04:sc4:eq:sqrt} or \eqref{p04:sc4:eq:cuberoot}.
But solutions of a linear differential equation are holomorphic at
an ordinary point $c_n\notin\Sigma$, including every branch reached
by continuation away from $\Sigma$.  This contradiction proves
that neither germ is $D$-finite.

We have proved all four equivalences, since an algebraic germ is
$D$-finite.  The last algebraicity assertion follows from
Lemma~\ref{p04:sc4:lem:algebraic}.
\end{proof}

The result distinguishes two different restrictions.  A common
irrational shift of two rationally spaced exponents prevents
algebraicity of the value, but still permits a linear differential
equation.  An irrational gap prevents every such linear equation.
This statement concerns functions and is independent of the
arithmetic nature of a value at one specified setting.

\subsection{An exact nonlinear equation and its least order}

\begin{theorem}[A differential characterization of two atoms]
\label{p04:sc4:thm:ode}
Every optimum in \eqref{p04:sc4:eq:twoatom} satisfies
\begin{equation}\label{p04:sc4:eq:ode}
 T^2\bigl[abT+(1-a-b)CT'\bigr]T''
                 +(T-CT')^2(T')^2=0.
\end{equation}
Conversely, let a real analytic function $T$ on an interval of
positive settings satisfy \eqref{p04:sc4:eq:ode} and
\begin{equation}\label{p04:sc4:eq:admissible}
 T>0,\quad T'>0,\quad B:=T-CT'>0,\quad T''<0,\qquad
                         a<\frac{CT'}T<b.
\end{equation}
Then, on that interval, $T$ is the optimized value for unique
positive constants $w_1,w_2$ in \eqref{p04:sc4:eq:twoatom}.
\end{theorem}
\begin{proof}
The two powers satisfy the Euler equation
\begin{equation}\label{p04:sc4:eq:euler}
 x^2S''+(1-a-b)xS'+abS=0.
\end{equation}
Insert $x=T'/B$, $S=1/B$, $S'=C/T$ and the last identity in
\eqref{p04:sc4:eq:second}.  Multiplication by $BT^3T''$ gives
exactly \eqref{p04:sc4:eq:ode}.  No differentiation of order three
is needed, and all its coefficients are independent of the weights.

For the converse, define $x=T'/B$.  Direct differentiation gives
$x'=TT''/B^2<0$, so it has a local real analytic inverse.
Set $S(x(C))=1/B(C)$.  The chain rule gives $S'=C/T$ and the
formula for $S''$ in \eqref{p04:sc4:eq:second}.  Thus
\eqref{p04:sc4:eq:ode} implies \eqref{p04:sc4:eq:euler}.  Since $a\ne b$,
the general solution is $S(x)=w_1x^a+w_2x^b$ with constants
$w_1,w_2$.  Writing $\eta=CT'/T$ and using $S=1/B$ gives
\begin{equation}\label{p04:sc4:eq:weights}
 w_1x^a=\frac{b-\eta}{dB},\qquad
 w_2x^b=\frac{\eta-a}{dB}.
\end{equation}
Both are positive by \eqref{p04:sc4:eq:admissible}.  Moreover
$S-xS'=1/T$, so $C=S'/(S-xS')$ and $T=1/(S-xS')$.
The unique-contact argument preceding \eqref{p04:sc4:eq:param} now
identifies $T$ with the global minimum for these positive weights.
The Euler representation and \eqref{p04:sc4:eq:weights} give uniqueness.
Local representations agree on overlaps, so the same two constants
work on the entire interval.
\end{proof}

\begin{theorem}[Exact algebraic differential order]
\label{p04:sc4:thm:order}
Let
\[
                  r=\dim_{\mathbb Q}
                        \operatorname{span}_{\mathbb Q}\{1,a,b\}.
\]
The least algebraic differential order of $T(C)$ is exactly $r-1$.
Thus it is zero, one or two.  In particular, if $1,a,b$ are
rationally independent, then $T$ and $T'$ are algebraically
independent over $\mathbb C(C)$, and
\eqref{p04:sc4:eq:ode} has the least possible order.
\end{theorem}
\begin{proof}
All the following fields are fields of meromorphic germs, not fields
of values at a fixed number.  Parametrize by $x=e^t$.  The contact
and envelope formulas give the exact identity
\begin{equation}\label{p04:sc4:eq:functionfield}
 \mathbb C(C,T,T')=\mathbb C(x,x^a,x^b).
\end{equation}
Indeed, the inclusion from left to right follows from
\eqref{p04:sc4:eq:param} and $T'=x/S$.  For the reverse inclusion,
$x=T'/(T-CT')$, while solving the two contact equations gives
\begin{align}
 w_1x^a&=\frac{b(1+Cx)-Cx}{dT},&
 w_2x^b&=\frac{Cx-a(1+Cx)}{dT}.
 \label{p04:sc4:eq:powerrecovery}
\end{align}

The transcendence degree over $\mathbb C$ of the right side of
\eqref{p04:sc4:eq:functionfield} is exactly $r$.  To verify this
directly, choose a rational basis $\lambda_1,\ldots,\lambda_r$
from the exponent span.  The functions $e^{\lambda_i t}$ are
algebraically independent: any polynomial relation would be a
linear relation among finitely many exponentials with distinct
exponents, contradicted by a Vandermonde matrix.  Clearing
denominators in the rational changes of basis makes the two power
fields algebraic over one another, which preserves transcendence
degree.  Since $C$ is a nonconstant analytic function, it is
transcendental over the constant field.  Hence
\begin{equation}\label{p04:sc4:eq:relativefield}
 \operatorname{trdeg}_{\mathbb C(C)}
                  \mathbb C(C,T,T')=r-1.
\end{equation}

If $r=3$, the two functions $T,T'$ are algebraically independent
over $\mathbb C(C)$, excluding every equation of order at most
one.  Equation \eqref{p04:sc4:eq:ode} supplies order two.  If $r=2$,
the same field degree supplies an equation of order at most one.
At least one exponent is irrational, so
Lemma~\ref{p04:sc4:lem:algebraic} excludes order zero.  If $r=1$,
both exponents are rational and that lemma gives order zero.
\end{proof}

Thus a failure of linear differential equations is not differential
transcendence.  Irrational-gap optima are differential algebraic, and
their minimal nonlinear order is known exactly.  For example,
a rationally independent cubic pair and the biquadratic pair in
\eqref{p04:sc4:eq:Q} both have order two and are not $D$-finite, even
though the present arithmetic information about their individual
values differs.

\begin{corollary}[Higher setting derivatives add no arithmetic field]
\label{p04:sc4:cor:alljets}
Suppose $a,b$ and the positive setting $C_0$ are algebraic.  At
that setting the entire derivative jet generates exactly the field
\begin{equation}\label{p04:sc4:eq:alljets}
 \overline{\mathbb Q}\bigl(T(C_0),T'(C_0),T''(C_0),\ldots\bigr)
              =\overline{\mathbb Q}\bigl(T(C_0),T'(C_0)\bigr).
\end{equation}
This assertion does not require the weights to be algebraic.
\end{corollary}
\begin{proof}
The coefficient
$A=abT+(1-a-b)CT'=T[ab+(1-a-b)\eta]$ in
\eqref{p04:sc4:eq:ode} is positive.  Indeed $\eta=CT'/T$ belongs
to $(a,b)$, and the affine expression in brackets takes the
positive endpoint values $a(1-a)$ and $b(1-b)$.
Define the rational function
\[
 R(C,Y,Z)=-\frac{(Y-CZ)^2Z^2}
                   {Y^2[abY+(1-a-b)CZ]}.
\]
Then $T''=R(C,T,T')$, and successively differentiating shows that
every higher derivative is a rational function of $C,T,T'$
over $\mathbb Q(a,b)$.  For instance
$T'''=R_C+T'R_Y+RR_Z$ evaluated at $(C,T,T')$.
Induction introduces only powers of the displayed nonzero
denominators.  Evaluation at the algebraic setting proves one
field inclusion in \eqref{p04:sc4:eq:alljets}; the other is immediate.
\end{proof}

\subsection{The biquadratic value: an exact reduction and a failed transfer}

For this subsection take
\[
 a=\frac1{\sqrt3},\qquad b=\frac1{\sqrt2},\qquad
 w_1=w_2=\frac12,\qquad C=1,
\]
and write $\xi=x(1)$ and $\tau=T(1)=Q$.  The previous
continuation proved that $\xi$ and the elasticity are individually
transcendental, using the irrational gap; those results are not
recounted as new here \cite{p04:root4:III}.
The weighted-mean contact identity gives
$a<\xi/(1+\xi)<b$, so
\begin{equation}\label{p04:sc4:eq:contactinterval}
                 \frac a{1-a}<\xi<\frac b{1-b}.
\end{equation}
In particular $\xi>1$.  The value is less than two because the
objective equals two at one and has negative derivative there.

\begin{proposition}[What an algebraic value would imply]
\label{p04:sc4:prop:biquadratic}
If $\tau$ is algebraic, then
\begin{equation}\label{p04:sc4:eq:rationalcurve}
 \xi^a=\frac{2[b-(1-b)\xi]}{(b-a)\tau},\qquad
 \xi^b=\frac{2[(1-a)\xi-a]}{(b-a)\tau},
\end{equation}
and
\begin{equation}\label{p04:sc4:eq:threefield}
 \overline{\mathbb Q}(\xi,\xi^a,\xi^b)
             =\overline{\mathbb Q}(\xi),\qquad
 \operatorname{trdeg}_{\overline{\mathbb Q}}
       \overline{\mathbb Q}(\xi,\xi^a,\xi^b)=1.
\end{equation}
Under this same hypothesis, the missing fourth power
$\xi^{ab}=\xi^{1/\sqrt6}$ is transcendental over
$\overline{\mathbb Q}(\xi)$.
\end{proposition}
\begin{proof}
The two affine formulas are \eqref{p04:sc4:eq:powerrecovery} with
$C=1$ and weights $1/2$.  They give the field identity.  The
contact $\xi$ is transcendental by the preceding result; alternatively,
assuming it algebraic, the stationary equation would make
$\xi^{b-a}$ algebraic, contrary to Gelfond--Schneider since
$\xi>1$ and $b-a$ is algebraic irrational.  This proves
\eqref{p04:sc4:eq:threefield}.

For clarity we specify the classical input governing the fourth
power.  Corollary~1.2(a), case $t=1$, of
Roy--Waldschmidt \cite{p04:sc4:RW} implies the following qualitative
fact: for two rationally independent families
$u_1,\ldots,u_p$ and $v_1,\ldots,v_q$, the field generated by
the $pq$ numbers $e^{u_iv_j}$ has transcendence degree at least
two if $pq\ge2(p+q)$.

Apply it with $p=q=4$, the first family
$(1,a,b,ab)$, and the second family
$(\log\xi,a\log\xi,b\log\xi,ab\log\xi)$.
Both are rationally independent, since
$\mathbb Q(\sqrt2,\sqrt3)$ has basis $1,a,b,ab$ and
$\log\xi\ne0$.  The equality $16=2(4+4)$ meets the hypothesis.
As $a^2=1/3$ and $b^2=1/2$, all the products $u_iv_j$
are rational multiples of
$\log\xi,a\log\xi,b\log\xi,ab\log\xi$.
Consequently the exponential-table field is algebraic over
$\mathbb Q(\xi,\xi^a,\xi^b,\xi^{ab})$, and it contains that
field.  The latter therefore has transcendence degree at least
two.  Extending the base to $\overline{\mathbb Q}$ does not
change this bound.  In view of \eqref{p04:sc4:eq:threefield},
$\xi^{ab}$ must be transcendental over
$\overline{\mathbb Q}(\xi)$.
\end{proof}

The proposition explains the precise limitation of this attempted
quartic-field transfer.  The full multiplication table does provide
a second transcendence degree, but it may be contributed entirely
by the fourth coordinate, which the two-atom contact equations do
not constrain.  There is no contradiction.  The ordinary conclusion
of a four- or six-exponential argument, that some entry is
transcendental, is already compatible with
\eqref{p04:sc4:eq:threefield}.  Nor can an algebraic-base theorem simply
be applied at the contact: that contact is itself transcendental.
Corollary~\ref{p04:sc4:cor:alljets} makes the derivative limitation
equally exact: if $\tau$ were algebraic, all setting derivatives
would lie in $\overline{\mathbb Q}(\xi)$, whereas the fourth power
would be transcendental over that field.  Taking more derivatives
with respect to $C$ would not supply the missing coordinate.

There is also a simple obstruction to applying the particular
exponential-table criterion used above while keeping only the three
controlled powers.  Put $E=\operatorname{span}_{\mathbb Q}\{1,a,b\}$.
If rationally independent families $(u_i)_{i=1}^p$ and
$(v_j)_{j=1}^q$ have every product $u_iv_j\in E$, then fixing
one nonzero member in either family gives $p,q\le3$.
It follows that $pq<2(p+q)$, so that criterion cannot apply.
This is a limitation of this specified rectangle construction,
not a claim that every transcendence method must fail.

A sufficient missing input for the previous curvature strategy
would be
\begin{equation}\label{p04:sc4:eq:missing}
 \operatorname{trdeg}_{\overline{\mathbb Q}}
   \overline{\mathbb Q}(x,x^{1/\sqrt3},x^{1/\sqrt2})\ge2
 \qquad(x>0,\ x\ne1),
\end{equation}
or a weaker assertion at the particular contact \eqref{p04:sc4:eq:contactinterval}.
We have not proved \eqref{p04:sc4:eq:missing}, and the full-table
argument above does not prove it.  Thus the irrationality and
transcendence of the individual number $Q$ remain unresolved here.

\paragraph{Accounting.}
The positive results classify two new functional questions: the
linear-differential class of the optimum/contact map, and the exact
minimal algebraic differential order together with a nonlinear
characterization.  The individual biquadratic value was investigated
but not closed.  Functional algebraic independence in
\eqref{p04:sc4:eq:relativefield} must not be confused with algebraic
independence after specializing $C$ to an algebraic number.
No historical named conjecture is counted as solved.

\section{Fractional extraction: exact contact geometry and a sharp arithmetic rank barrier}
\label{p04:be20:section}

The integer-shape argument of the preceding continuation used a
two-dimensional Bessel system.  Fractional extraction changes that
dimension, but it does not destroy all of the useful differential
structure.  We prove two results.  First, every interior boundary
remainder, at every real shape, has only double zeros.  The entire
capacity curve is locally a finite minimum of analytic contact curves,
with an exact integral slope formula and a global comparison inequality.
Second, for a noninteger shape the normalized inverse has differential
order exactly three.  At rational parameters it is an $E$-function,
and its normalized values and logarithmic derivatives at algebraic
points are transcendental.  The rank calculation explains precisely
why the preceding two-row elimination does not settle fractional
critical-rate transcendence.

The individual endpoint targets $A(1)$ and $A(3/2)$ are revisited
at the end.  Their transcendence is not proved here.

\subsection{An inverse and a third-order equation for every real shape}

Fix $\kappa>1/2$, $b>1$, and write
\[
 M=2\kappa-1,\qquad t=2\kappa-2.
\]
The original law has Laplace transform
\[
 L(s)=\frac{b+s}{b}
       \left(\frac{b^2-1}{(b+s)^2-1}\right)^\kappa.
\]
For a gamma shape $0<h<M$ and rate $b-c$, where $0<c<1$, put
$\eta=M-h>0$.  Define coefficients and functions by
\begin{align}
 (1-cz)^h(1-z^2)^{-\kappa}&=\sum_{n\ge0}b_nz^n,
                                                        \label{p04:be20:coeff}\\
 v(x)&=\sum_{n\ge0}\frac{b_nx^{\eta+n}}{\Gamma(\eta+n+1)},
 &g(x)&=v'(x),                                           \label{p04:be20:inverse}\\
 H(x)&=\Gamma(\eta+1)x^{-\eta}v(x),
 &G(x)&=\eta H(x)+xH'(x).
                                                        \label{p04:be20:normalized}
\end{align}
Here $H$ and $G$ are entire, whereas $v$ and $g$ use the positive
real branch when $\eta$ is not an integer.  The series converge
locally uniformly, together with all derivatives away from zero;
their parameter dependence is real analytic in the strict domain.
The actual candidate remainder density is a positive constant times
$e^{-bx}g(x)$.  In particular it is a probability density precisely
when $g\ge0$.  The transform identities are
\begin{equation}\label{p04:be20:laplace}
 V(w):=\int_0^\infty e^{-wx}v(x)\,dx
       =\frac{(w-c)^h}{(w^2-1)^\kappa},\qquad
 \mathcal Lg(w)=wV(w),\qquad \Re w>1.
\end{equation}
They follow first for sufficiently large real $w$ by termwise
integration and then by analytic continuation.  The factor
$v(0+)=0$ justifies the second identity.

We use the previously established capacity theorem as an input
\cite{p04:be20:Previous}: the maximal permissible shape
$\mathcal H(c)=H_\kappa(c)$ is continuous and strictly decreasing
on $[0,1]$, lies strictly between $0$ and $M$ in the interior,
and is attained.  An interior boundary remainder is analytic on
$(0,\infty)$, positive near zero and infinity, and has finitely
many nonempty interior contacts.  The endpoint positivity estimates
are uniform in compact parameter neighborhoods with $h<M$.
These earlier analytic facts are not counted again.

\begin{lemma}[Fractional-shape differential identity]
\label{p04:be20:ode}
For every real $0<h<M$ and $0<c<1$, the inverse $v$ satisfies
\begin{equation}\label{p04:be20:third}
 xv'''+(h+1-t-cx)v''+(ct-x)v'+(cx-h-1)v=0
                         \qquad(x>0).
\end{equation}
\end{lemma}
\begin{proof}
The logarithmic derivative of \eqref{p04:be20:laplace} gives
\[
 (w-c)(w^2-1)V'
   =\{h(w^2-1)-2\kappa w(w-c)\}V.
\]
Under Laplace inversion, $V'$ corresponds to $-xv$ and
multiplication by $w$ to differentiation on the positive axis.
Thus, with $D=d/dx$, the equation becomes
\[
 -(D-c)(D^2-1)(xv)
     =\{h(D^2-1)-2\kappa D(D-c)\}v.
\]
Expansion gives \eqref{p04:be20:third}.  This manipulation can be
verified without assuming that $v''$ or $v'''$ is integrable at
zero: insert the convergent series \eqref{p04:be20:inverse}, compare
coefficients in the preceding first-order equation for $V$, and
then differentiate on a compact positive interval.  The resulting
coefficient identities are exactly \eqref{p04:be20:third}.
In particular no rule for commuting a fractional differential
power with a variable coefficient has been assumed.
\end{proof}

\begin{theorem}[Every real-shape contact is exactly double]
\label{p04:be20:double}
Suppose $0<h<M$, $0<c<1$, and the candidate remainder is
nonnegative.  Every positive zero $x$ of its density is exactly
double and satisfies
\begin{equation}\label{p04:be20:location}
                              0<x<\frac{h+1}{c}.
\end{equation}
This includes every interior probability boundary and every
noninteger shape; no rationality assumption is required.
\end{theorem}
\begin{proof}
Since $g=v'\ge0$ and
$v(x)\sim x^\eta/\Gamma(\eta+1)$ at zero, one has $v(x)>0$
for every $x>0$.  At a zero of the nonnegative analytic function
$g$, we have $g'=0$.  Equation \eqref{p04:be20:third} then gives
\begin{equation}\label{p04:be20:curvature}
                      xg''=(h+1-cx)v.
\end{equation}
The left side is nonnegative, so $cx\le h+1$.  In the equality
case $g''=0$.  Differentiating \eqref{p04:be20:third} at that point,
where $v'=v''=v'''=0$, would give
$xg'''=-cv<0$.  This would be an odd zero of exact order three,
contrary to nonnegativity on both sides.  Therefore $cx<h+1$,
and \eqref{p04:be20:curvature} gives $g''>0$.
\end{proof}

\subsection{The full frontier has analytic contact branches and a signed response}

Partial derivatives below hold $\kappa$ fixed.  At a boundary
contact $x$, define the positive finite integrals
\begin{equation}\label{p04:be20:AB}
 A_x=\int_0^x e^{cy}g(x-y)\,dy,\qquad
 B_x=\int_0^x\frac{e^{cy}g(x-y)}{y}\,dy.
\end{equation}
Finiteness near $y=x$ follows from $\eta>0$, and near $y=0$
from the double zero just proved.

\begin{lemma}[Parameter derivatives at a contact]
\label{p04:be20:response}
At every boundary contact,
\begin{equation}\label{p04:be20:partials}
                          g_c(x)=-hA_x<0,
                 \qquad  g_h(x)=-B_x<0.
\end{equation}
\end{lemma}
\begin{proof}
Differentiation of $wV$ with respect to $c$ multiplies its
transform by $-h/(w-c)$.  Convolution therefore gives the first
identity for every positive $x$.

For the second derivative, the multiplier is $\log(w-c)$.
The regularized identity
\[
 \log(w-c)=\int_0^\infty
             \frac{e^{-y}-e^{-(w-c)y}}{y}\,dy
\]
gives, on the positive axis,
\begin{equation}\label{p04:be20:logkernel}
 g_h(x)=-(\gamma+\log x)g(x)
       +\int_0^x\frac{g(x)-e^{cy}g(x-y)}{y}\,dy,
\end{equation}
where $\gamma$ is Euler's constant.  To justify it directly,
truncate the preceding transform integral at $y=\epsilon$,
invert the resulting convolution, and use
$\int_\epsilon^\infty e^{-y}dy/y=-\gamma-\log\epsilon+o(1)$.
The subtracted integrand converges locally at $y=0$ because
$g$ is differentiable at $x>0$, and is integrable at $y=x$
because $g(u)=O(u^{\eta-1})$.  Thus the limit is
\eqref{p04:be20:logkernel}; it also agrees with the differentiated
locally convergent inverse series.  At a contact $g(x)=0$,
this is precisely $-B_x$.  Both integrals are strictly positive,
since $g$ is nonnegative and positive near the origin.
\end{proof}

\begin{theorem}[Regularity and a universal slope inequality]
\label{p04:be20:frontier}
On every compact subinterval of $(0,1)$, the capacity
$\mathcal H(c)$ is piecewise real analytic with only finitely
many possible switching points.  Locally it is the minimum of
finitely many real analytic functions.  It is locally Lipschitz;
its one-sided derivatives exist everywhere in $(0,1)$, and
\begin{equation}\label{p04:be20:slopebound}
 -\frac{\mathcal H(c)(\mathcal H(c)+1)}{c}
          <\mathcal H'_\pm(c)<0.
\end{equation}
At a switching point the derivative can jump only downward:
$\mathcal H'_-(c)\ge\mathcal H'_+(c)$.

More precisely, a contact branch through $(c,h,x)$ has slope
\begin{equation}\label{p04:be20:slope}
                             h'(c)=-h\frac{A_x}{B_x},
                         \qquad -hx<h'(c)<0
\end{equation}
whenever that branch realizes the probability boundary.  The
normalized frontier
\begin{equation}\label{p04:be20:monotone}
                         c\longmapsto
                         \frac{c\mathcal H(c)}{1+\mathcal H(c)}
\end{equation}
is strictly increasing.  In particular, for $0<c_1<c_2<1$,
\begin{equation}\label{p04:be20:comparison}
 \frac{c_1\mathcal H(c_1)}
      {c_2+(c_2-c_1)\mathcal H(c_1)}
       <\mathcal H(c_2)<\mathcal H(c_1).
\end{equation}
\end{theorem}
\begin{proof}
Fix a boundary point $(c_0,h_0)$ and list its finitely many
contacts $x_1,\ldots,x_N$.  The equations $g=0$, $g_x=0$
have, as equations in $(h,x)$, Jacobian determinant
$g_hg_{xx}\ne0$ at each contact, by Theorem~\ref{p04:be20:double}
and Lemma~\ref{p04:be20:response}.  The analytic implicit-function
theorem therefore gives contact curves $(h_j(c),x_j(c))$.

Here these curves describe the full local probability boundary,
not only necessary conditions for it.  Choose disjoint small
intervals around the contacts.  Their strict minima persist,
since $g_{xx}>0$, and their minimum values have negative
$h$-derivative.  On the intervening compact set the original
density has a positive minimum.  Uniform endpoint positivity,
an inherited input stated above, also persists in a parameter
neighborhood.  Consequently, in that neighborhood,
\[
                     g\ge0\quad\Longleftrightarrow\quad
                     h\le\min_{1\le j\le N}h_j(c).
\]
Continuity of the already established capacity ensures that
its graph stays in this neighborhood.  Thus
$\mathcal H(c)=\min_jh_j(c)$ there.

Two analytic branches either agree identically locally or have
isolated intersections.  A finite family of such branches has
only finitely many switching points in a sufficiently small
neighborhood of a fixed interior point.  Compactness then gives
the asserted local piecewise analyticity and local Lipschitz
property.  At a multiple active branch, the left derivative is
the largest active slope and the right derivative is the smallest,
which proves the direction of a possible jump.

Differentiating $g(c,h(c),x(c))=0$ and using $g_x=0$ gives
$h'=-g_c/g_h=-hA_x/B_x$.  The quotient $A_x/B_x$ is the
mean of $y$ under the positive measure with density
$e^{cy}g(x-y)/y$ on $(0,x)$.  It lies strictly between $0$
and $x$.  This proves \eqref{p04:be20:slope}; combining it with
\eqref{p04:be20:location} proves \eqref{p04:be20:slopebound}, including
both one-sided derivatives at a switch.

On each analytic piece, differentiation gives
\[
 \left(\frac{c\mathcal H}{1+\mathcal H}\right)'
  =\frac{\mathcal H}{1+\mathcal H}
      +\frac{c\mathcal H'}{(1+\mathcal H)^2}>0.
\]
The function is continuous at the finitely many switches on a
compact interval, so it is strictly increasing throughout that
interval.  Rearranging this comparison at $c_1,c_2$ yields the
first inequality of \eqref{p04:be20:comparison}; the second is the
inherited strict decrease of the capacity.
\end{proof}

The theorem allows corners and more than one contact.  It does not
assert global analyticity or a unique contact at every rate.
What it excludes is a flat contact, a locally infinite accumulation
of switches in the interior, and an unbounded local slope.

\subsection{Noninteger shapes have exact differential order three}

The following rank result needs no arithmetic assumptions.
It is the obstruction that the integer-shape calculation did not
encounter.

\begin{theorem}[Three singularities force rank three]
\label{p04:be20:rank}
Let $\kappa>1/2$, $0<c<1$, and $0<h<M$ with
$h\notin\mathbb Z$.  Then
\begin{equation}\label{p04:be20:independent}
                           1,H,H',H''
\end{equation}
are linearly independent over $\mathbb C(x)$.  In particular
the minimal homogeneous differential equation of $H$ has order
exactly three.  The same minimal-order conclusion holds for $G$.
\end{theorem}
\begin{proof}
Suppose a rational relation in \eqref{p04:be20:independent} existed.
The substitution $H=\Gamma(\eta+1)x^{-\eta}v$ changes the
three derivative coefficients by an invertible triangular
rational matrix.  Clearing denominators and multiplying by
$x^2$ would therefore give
\begin{equation}\label{p04:be20:putativerelation}
                  P_0(x)v+P_1(x)v'+P_2(x)v''=Q(x)x^\eta,
\end{equation}
where the $P_i,Q$ are polynomials, all divisible by $x^2$,
and at least one $P_i$ is nonzero.  If all $P_i$ were zero,
the original relation would already force $Q=0$.

Write $P_i(x)=\sum_n p_{in}x^n$ and let $N$ be the largest
index for which some $p_{iN}\ne0$.  Thus $N\ge2$ and
\[
                       R_N(w)=\sum_{i=0}^2p_{iN}w^i
\]
is a nonzero polynomial of degree at most two.  Laplace
transformation of \eqref{p04:be20:putativerelation} is legitimate
after the multiplication by $x^2$: every integrand is integrable
at zero, and integration by parts has no boundary term.
The transform of the left side is
\begin{equation}\label{p04:be20:transformedrelation}
              \sum_{i=0}^2\sum_{n\ge2}
                   p_{in}(-D_w)^n\{w^iV(w)\}.
\end{equation}
The right side is a finite linear combination of powers
$w^{-\eta-n-1}$.  It is therefore holomorphic locally at
each of the three nonzero points $c,1,-1$, with compatible
branches chosen by analytic continuation.

At these three points $\lambda$, the transform has the form
\[
 V(w)=(w-\lambda)^{a_\lambda}U_\lambda(w),\qquad
 U_\lambda(\lambda)\ne0,
 \qquad a_c=h,\quad a_1=a_{-1}=-\kappa.
\]
If $R_N(\lambda)\ne0$, the leading local term of
\eqref{p04:be20:transformedrelation} is a nonzero multiple of
\[
 (a_\lambda)_{\underline N}
 R_N(\lambda)U_\lambda(\lambda)
 (w-\lambda)^{a_\lambda-N}.
\]
No lower derivative order has this exponent.  At $\lambda=c$,
the exponent is nonintegral and the falling factorial is
nonzero because $h$ is not an integer.  At $\lambda=\pm1$,
the exponent is negative and the falling factorial is nonzero
because $-\kappa<0$.  In all cases this term prevents a
holomorphic right side: it is either a pole or the first
nonintegral power of the local branch.  Hence
$R_N(c)=R_N(1)=R_N(-1)=0$.  A nonzero polynomial of degree
at most two cannot have these three distinct roots, contradiction.

Equation \eqref{p04:be20:third}, after the rational gauge
$v=x^\eta H/\Gamma(\eta+1)$, supplies an order-three
equation for $H$, so its minimal order is exactly three.
Finally the rational state transformation from $(v,v',v'')$
to $(g,g',g'')$ has determinant
$(h+1-cx)/x$, by \eqref{p04:be20:third}.  This determinant is
not identically zero.  Thus differentiation preserves the
functional rank three, and the further gauge
$g=x^{\eta-1}G/\Gamma(\eta+1)$ proves the assertion for $G$.
\end{proof}

The pole case at $\lambda=\pm1$ is included, so integer
$\kappa$ is not an exception.  In contrast, when $h$ is an
integer the factor $(w-c)^h$ is locally a polynomial and the
third singularity disappears.  That is precisely the reduction
which allowed the preceding two-row contact determinant.

\subsection{Arithmetic consequences for every rational noninteger shape}

\begin{lemma}[The normalized inverse is an $E$-function]
\label{p04:be20:E}
If $\kappa,h$ are rational in the preceding strict domain and
$c$ is algebraic, then $H,H',H''$ are $E$-functions.  They
satisfy a first-order system over $\overline{\mathbb Q}(x)$
whose only finite possible singularity is zero.
\end{lemma}
\begin{proof}
The exponential Taylor coefficients of $H$ are
\[
                         a_n=\frac{n!b_n}{(\eta+1)_n}.
\]
The coefficients of each rational-power binomial in
\eqref{p04:be20:coeff} have denominators dividing $C^n$ for
some integer $C$.  Indeed, for a rational parameter $A/B$
in lowest terms, primes not dividing $B$ cancel the factorial
denominator by counting a complete residue class modulo every
prime power.  At primes dividing $B$, its denominator valuation
is at most $n\{v_p(B)+1/(p-1)\}$.  Clearing a fixed algebraic
denominator of $c$ therefore makes $C^nb_n$ algebraic integral
in a fixed number field.  Cauchy's estimate applied to
\eqref{p04:be20:coeff} under each embedding bounds all conjugates
of $b_n$ geometrically.

It remains to account for $n!/(\eta+1)_n$.  Write
$\eta+1=A/B>0$ in lowest terms.  For a prime not dividing
$B$, counting the terms $A+Bj$, $0\le j<n$, divisible by
$p^k$ differs from $\lfloor n/p^k\rfloor$ by at most one.
No such term exceeds $Cn$, for a fixed integer $C$.  Thus the
denominator valuation is at most
$\lfloor\log(Cn)/\log p\rfloor$; primes dividing $B$ add
no denominator.  Consequently the common denominator for
these ratios through index $N$ divides
$\operatorname{lcm}(1,\ldots,CN)$, which has exponential
growth.  The ratios themselves have polynomial growth, by the
gamma-ratio estimate, so the coefficients $a_n$ and their
common denominators satisfy the exponential bounds in the
definition of an $E$-function.

The rational gauge in \eqref{p04:be20:third} gives a third-order
equation for $H$ with algebraic coefficients and possible poles
only at zero.  This proves the required differential condition.
Derivative closure gives the same properties for $H'$ and $H''$.
\end{proof}

\begin{theorem}[Normalized values and logarithmic derivatives]
\label{p04:be20:arithmetic}
Assume $\kappa,h\in\mathbb Q$, $0<h<2\kappa-1$,
$h\notin\mathbb Z$, and $c\in\overline{\mathbb Q}\cap(0,1)$.
For every nonzero algebraic $\xi$, the four numbers
\begin{equation}\label{p04:be20:values}
                         1,H(\xi),H'(\xi),H''(\xi)
\end{equation}
are linearly independent over $\overline{\mathbb Q}$.
In particular
\begin{equation}\label{p04:be20:transvalues}
                      G(\xi)\quad\hbox{and}\quad
                      \frac{G'(\xi)}{G(\xi)}
                      \quad\hbox{are transcendental}.
\end{equation}
Every positive zero of the candidate remainder density is
transcendental.  If $b>1$ is also algebraic, every positive
stationary point of that density is transcendental as well.
\end{theorem}
\begin{proof}
Adjoin the constant $1$ to the first-order $E$-function system
in Lemma~\ref{p04:be20:E}.  Its only finite possible singularity
is zero, and its four entries are linearly independent by
Theorem~\ref{p04:be20:rank}.  Beukers' linear specialization
theorem \cite[Corollary~1.4]{p02:be18:Beukers} proves
\eqref{p04:be20:values} at every nonzero algebraic point.

The expression $G(\xi)=\eta H(\xi)+\xi H'(\xi)$ is a
nonzero algebraic linear form in the last three entries.
It cannot be algebraic, since that would give a relation with
the first entry $1$.  If $G'(\xi)/G(\xi)=a$ were algebraic,
then
\[
 -a\eta H(\xi)+(\eta+1-a\xi)H'(\xi)+\xi H''(\xi)=0,
\]
another nonzero relation because $\xi\ne0$.  This proves
\eqref{p04:be20:transvalues}.

On the positive axis the candidate physical density is a
nonzero constant times
$e^{-bx}x^{\eta-1}G(x)$.  Its positive zeros are exactly the
zeros of $G$, so no such zero is algebraic.  Where it is
nonzero, its logarithmic derivative is
\[
                         -b+\frac{\eta-1}{x}+
                         \frac{G'(x)}{G(x)}.
\]
At positive algebraic $x$, this is transcendental, and in
particular nonzero, if $b$ is algebraic.  Thus a stationary
point cannot be algebraic either.
\end{proof}

The normalization in \eqref{p04:be20:transvalues} matters.  We have
not asserted transcendence of the physical density value after
multiplication by its gamma and exponential normalizers.  The
logarithmic derivative and the zero or stationary locations do
not have that ambiguity.

\subsection{What the new rank theorem does not prove}

For a fixed rational noninteger $h$, a boundary contact still
imposes only $g(x)=g'(x)=0$.  At an algebraic candidate rate,
Theorem~\ref{p04:be20:arithmetic} forces the contact $x$ to be
transcendental, but does not contradict its existence.  The
rank-three system has a one-dimensional space of initial states
satisfying those two equations.  Hence the two-row determinant
from the integer argument has no rank-three counterpart of the
same form.  This is a proved structural obstruction to that
specific elimination method, not a proof that fractional-rate
transcendence is impossible by other methods.

The two endpoint questions $A(1)$ and $A(3/2)$ remain separate.
If one of these constants were algebraic irrational, its
corresponding maximal shape would be algebraic irrational, and
the normalized Kummer function in the endpoint representation
would fail the $E$-function criterion.  The earlier exact
application of Galochkin's classification, as formulated by
Rivoal \cite[Theorem~1]{p02:be18:Rivoal}, continues to apply:
the only possible irrational numerator--denominator parameter
pair has difference $-d<0$.  The present rational-parameter
argument therefore supplies no arithmetic specialization at
that hypothetical endpoint.  Nor does linear independence of
functions at a fixed order imply transcendence of a varying
algebraic-irrational order.

\subsection{An exponential observable at the integer contacts}

There is nevertheless a further arithmetic question that can be
settled when the original Bessel parameter is an integer.  It
concerns the exponential of the contact, rather than an
exponential of the critical rate.

\begin{proposition}[The contact exponential is also transcendental]
\label{p04:be20:exponential}
Let $\kappa=k\ge2$ and $1\le m<2k-1$ be integers.  Suppose the
shape-$m$ candidate has a positive multiple contact $x$ at
$0<c<1$.  Then $e^{2x}$ is transcendental.  In fact
\begin{equation}\label{p04:be20:expfield}
 e^{2x}\in\overline{\mathbb Q}(x,c),\qquad
 \operatorname{trdeg}_{\overline{\mathbb Q}}
       \overline{\mathbb Q}(x,c,e^{2x})=1.
\end{equation}
Thus the contact, the rate coordinate, and the contact exponential
are individually transcendental but are algebraically dependent.
The exponential $e^{rx}$ is transcendental for every nonzero
rational $r$ as well.
\end{proposition}
\begin{proof}
Up to a nonzero constant factor, the half-integer Bessel
expression $x^{k-1/2}I_{k-3/2}(x)$ has the form
\[
 f(x)=e^xP(x)+e^{-x}Q(x),
\]
where $P,Q\in\overline{\mathbb Q}[x]$ both have degree
$d=k-1$ and nonzero leading coefficients $p,q$.  The removed
constant has no effect on contacts.  This representation follows
from the elementary formulas for $I_{-1/2},I_{1/2}$ and the
usual Bessel recurrence \cite[10.39 and 10.29]{p04:be20:DLMF}.
The two exponential terms separately solve the base Bessel
equation and form an independent solution basis.
Put
\[
 A=(D+1-c)^mP,\qquad B=(D-1-c)^mQ.
\]
At a contact, writing $E=e^{2x}$ gives the polynomial equations
\begin{equation}\label{p04:be20:exppolynomials}
       EA(x,c)+B(x,c)=0,\qquad
       E(A_x+A)(x,c)+(B_x-B)(x,c)=0.
\end{equation}
The nonzero-contact-row lemma of the preceding integer-shape
paper \cite{p04:root4:III} says that these two coefficient rows
cannot both vanish when $x\ne0$ and $c^2\ne1$.  Consequently
one equation has a nonzero coefficient of $E$; otherwise its
constant coefficient would also vanish at the actual contact.
Solving that equation proves the field inclusion in
\eqref{p04:be20:expfield}.

We now show that for every fixed complex $E\ne0$, the two
equations in \eqref{p04:be20:exppolynomials} permit only finitely
many $x$-coordinates.  First their common zeros cannot have
$|x|\to\infty$.  Along such a hypothetical sequence, if $c$
has a bounded subsequence with limit $c_0$, division by $x^d$
gives the two limiting equations
\[
 Ep(1-c_0)^m+q(-1-c_0)^m=0,\qquad
 Ep(1-c_0)^m-q(-1-c_0)^m=0.
\]
They would require simultaneously $c_0=1$ and $c_0=-1$.
If instead $|c|\to\infty$, divide by $x^dc^m$.
The terms involving derivatives of $P,Q$ are smaller by a
factor tending to zero, and the two limits are
$(-1)^m(Ep+q)=0$ and $(-1)^m(Ep-q)=0$, also impossible.

For completeness, the required algebraic elimination follows
from this boundedness without a generic-fiber assumption.
Regard the two polynomials as polynomials in $c$ over
$\mathbb C(x)$.  If their greatest common divisor had positive
$c$-degree, clearing denominators in a common factor would
produce a common zero in $c$ for every sufficiently large $x$
outside a finite exceptional set.  This contradicts the
boundedness just proved.  Therefore their greatest common
divisor over $\mathbb C(x)[c]$ is one.  A Bezout identity over
$\mathbb C(x)$, with denominators cleared, yields a nonzero
polynomial in $x$ belonging to their ideal; all common
$x$-coordinates are among its finitely many roots.  If $E$ is
algebraic, the same computation takes place over
$\overline{\mathbb Q}(x)$ and gives a polynomial with algebraic
coefficients.

Suppose now that the actual $E=e^{2x}$ were algebraic.  The
preceding elimination would force $x$ algebraic.  Since $x>0$,
the Hermite--Lindemann theorem would then make $e^{2x}$
transcendental, contradiction.  Thus $E$ is transcendental.
The preceding integer-shape theorem already proves that $x,c$
are individually transcendental with joint transcendence degree
one; the field inclusion gives the second assertion of
\eqref{p04:be20:expfield}.  Finally, algebraicity of $e^{rx}$ for
nonzero rational $r=a/b$ would make $e^{2ax}=E^a$ algebraic,
and hence $E$ algebraic, another contradiction.
\end{proof}

This proves dependence rather than independence for the contact
and its own exponential.  It gives no conclusion about
$e^c$, and no application of Lindemann--Weierstrass at a
transcendental argument has been assumed.

\paragraph{Accounting.}
The first new family is the all-real-shape contact and frontier
regularity theorem, including its exact response and comparison
law.  The second is the exact noninteger differential-rank
theorem and its rational-parameter arithmetic specialization.
The contact-exponential proposition closes an additional
observable question inside the preceding integer family, and
is recorded as a related consequence rather than a new family.
Individual examples of a rational shape or algebraic sample
point are not separate discoveries.  Transcendence of the
fractional critical rate itself and of $A(1),A(3/2)$ remains
unresolved.  No historical named constant conjecture is claimed
to have been closed.

\section{A linear differential obstruction for the two prime series}
\label{p04:ci20:section}

We revisit the three fixed-value targets
\[
 E=\sum_{n\ge1}(2^n-1)^{-1},\qquad
 P_2=\sum_{p\ {\rm prime}}2^{-p},\qquad
 L_{\mathbb P}=\sum_{p\ {\rm prime}}(2^p-1)^{-1}.
\]
No transcendence assertion for these numbers is proved here.  The new
result concerns a more permissive functional transfer than algebraic
elimination: we allow arbitrarily many derivatives of both prime series,
and coefficients rational in a reciprocal-Mersenne moment and finite
modular constructions.  Even a nontrivial linear differential relation
is impossible.  The proof separates two meromorphic Borel transforms
by their different sets of poles.  This does not establish nonlinear
differential algebraic independence.

Write, for $|q|<1$,
\[
 A_r(q)=\sum_{n\ge1}\left(\frac{q^n}{1-q^n}\right)^r,
 \qquad P(q)=\sum_pq^p,\qquad B(q)=\sum_p\frac{q^p}{1-q^p}.
\]
For the boundary analysis put $q=e^{-t}$, $L=\log(1/t)$ and
$\partial=-t\,d/dt$, so $\partial L=1$.
For $\lambda>0$ let
\[
 \mathscr K_\lambda
 =\{F(t,e^{-\lambda/t}):F\in\operatorname{Frac}\mathbb C\{t,z\}\}.
\]
These are germs along $t\downarrow0$.  A nonzero element has the form
$ct^ae^{-m\lambda/t}(1+O(t))$, with $a,m\in\mathbb Z$ and $c\ne0$:
expand a nonzero convergent numerator and denominator in $z$, take
their first nonzero coefficients, and then take the first nonzero
powers of $t$ in those coefficients.  The remaining $z$ terms are
smaller than every power of $t$.  This also proves that evaluation is
injective.  The field is closed under $d/dt$, by the chain rule.

\subsection{Two formal transforms with different infinite pole sets}

Let
\begin{equation}\label{p04:ci20:formalFG}
 F(w)=\sum_{k\ge0}a_kw^{-k-1},\quad
 a_k=(-1)^k\Gamma^{(k)}(1),
 \qquad
 G(w)=\log w+M+
       \sum_{k\ge1}\frac{(-1)^{k+1}m_k}{k w^k},
\end{equation}
where $M$ is any complex constant and
\begin{equation}\label{p04:ci20:H}
 H(s)=s\Gamma(s+1)\zeta(s+1),\qquad m_k=H^{(k)}(0).
\end{equation}
The singularity of $H$ at zero is removable and $H(0)=1$.
The following statement is about formal series; it assumes no
convergence of \eqref{p04:ci20:formalFG}.

\begin{lemma}[Formal independence of both derivative families]
\label{p04:ci20:formalindependence}
The family
\[
 1,\quad F,F',F'',\ldots,\quad G,G',G'',\ldots
\]
is linearly independent over $\mathbb C(w)$, with derivatives taken
formally with respect to $w$.
\end{lemma}

\begin{proof}
Clear denominators in a putative finite relation so that all its
coefficients are polynomials in $w$.  Only the undifferentiated $G$
contains $\log w$, so its coefficient must vanish.  The remaining
terms are formal Laurent series with finitely many nonnegative powers.
Define their formal Borel transform by
\[
 \mathcal B\left(\sum_{k\ge0}c_kw^{-k-1}
                      +\text{a polynomial in }w\right)(s)
       =\sum_{k\ge0}\frac{c_k}{k!}s^k.
\]
Direct coefficient comparison gives
\[
 \mathcal B(wU)=\frac d{ds}\mathcal B(U),\qquad
 \mathcal B(U')=-s\mathcal B(U).
\]
In particular,
\begin{equation}\label{p04:ci20:Borelpair}
 \mathcal B(F)=\Gamma(1-s),\qquad
 \mathcal B(G')=H(-s)=-s\Gamma(1-s)\zeta(1-s).
\end{equation}
Both right-hand sides converge near zero and have meromorphic
continuations to the complex plane.  Thus the transformed relation is
\begin{equation}\label{p04:ci20:operatorrelation}
 U_1\Gamma(1-s)+U_2H(-s)=0,
\end{equation}
where $U_1,U_2$ are polynomial-coefficient differential operators in
$s$.  Substitution $w\mapsto d/ds$, $d/dw\mapsto-s$ is an invertible
map of Weyl algebras; hence a nonzero original differential operator
gives a nonzero $U_i$.  Equivalently, one may verify this by writing
operators uniquely as finite sums of powers of $d/ds$ with polynomial
coefficients.

The function $\Gamma(1-s)$ has a simple pole at every positive integer.
The function $H(-s)$ is holomorphic at every odd integer $s\ge3$,
because the trivial zero of $\zeta(1-s)$ cancels the gamma pole.
It has a simple pole at every positive even integer, since
$\zeta(1-2n)\ne0$ for $n\ge1$.  These facts also follow directly from
the functional equation of $\zeta$.

If $U_1=\sum_{j=0}^d c_j(s)(d/ds)^j$ is nonzero, take $d$ maximal.
At all but finitely many positive odd integers, $c_d$ is nonzero and
$U_1\Gamma(1-s)$ has a pole of exact order $d+1$: the derivatives of
lower order cannot cancel its leading polar term.  In contrast,
$U_2H(-s)$ is holomorphic there.  This contradicts
\eqref{p04:ci20:operatorrelation}, so $U_1=0$.  The same argument at the
positive even integers forces $U_2=0$.  Consequently all coefficients
of $F^{(j)}$ and $G^{(j)}$, $j\ge1$, in the original relation vanish.
The remaining polynomial must also vanish, proving the lemma.
\end{proof}

\subsection{All-order asymptotics, including every fixed derivative}

Let $M$ now denote the Meissel--Mertens constant, determined by
\[
 M(x):=\sum_{p\le x}\frac1p=\log\log x+M+o(1).
\]
We use a precise convention for the next lemma: an asymptotic formal
expansion means that truncation at any prescribed power of $L^{-1}$
has an error of the next order.  The expansion is asserted separately
for each fixed derivative order, not uniformly in that order.

\begin{lemma}[Differentiated prime expansions]\label{p04:ci20:primeexpansions}
For $Y(t)=tP(e^{-t})$ and $Z(t)=tB(e^{-t})$, one has, to every
inverse-logarithmic order and for every fixed $j\ge0$,
\begin{equation}\label{p04:ci20:asymptoticpair}
 \partial^jY(t)\sim F^{(j)}(L),\qquad
 \partial^jZ(t)\sim G^{(j)}(L),
\end{equation}
where $F,G$ are \eqref{p04:ci20:formalFG} with this value of $M$.
\end{lemma}

\begin{proof}
The prime-counting estimate in Trudgian's Theorem~2
\cite{p04:ci20:Trudgian2016} implies, for every $A>0$,
\[
 \pi(x)-\int_2^x\frac{du}{\log u}=O_A(x/\log^A x).
\]
Partial summation also gives
\begin{equation}\label{p04:ci20:Mertenserror}
 M(x)=\log\log x+M+O_A((\log x)^{-A})
\end{equation}
for every fixed $A$: the tail integral of the prime-counting error
is bounded by a constant times
$\int_x^\infty du/(u\log^{A+1}u)$, after increasing its requested
exponent if necessary.

For clarity we justify differentiation of the expansions by integral
formulas, rather than by differentiating an unqualified remainder.
Define polynomials $Q_0(u)=1$ and
$Q_{j+1}(u)=(1-u)Q_j(u)+uQ'_j(u)$.  Then
\[
 (t\,d/dt)^jY(t)
   =t\int_{2^-}^\infty Q_j(tx)e^{-tx}\,d\pi(x).
\]
Integration by parts against the prime-counting error, splitting
at $x=t^{-1/2}$, shows that the right-hand side equals
\begin{equation}\label{p04:ci20:Qintegral}
 \int_{2t}^\infty\frac{Q_j(u)e^{-u}}{L+\log u}\,du
      +O_{A,j}(L^{-A}).
\end{equation}
On the first part of the split use $\pi(x)=O(x)$ and the boundedness
of the fixed polynomial kernel and its derivative near zero; its
contribution to the error is $O_j(t)$.  On the second part use
$\log x\ge L/2$ and the integrability of polynomial multiples of
$e^{-u}$.

In \eqref{p04:ci20:Qintegral}, discard $u<t^{1/2}$ and
$u>t^{-1/2}$ at an error smaller than every inverse power of $L$.
On the remaining interval expand $(L+\log u)^{-1}$ to any fixed
order.  All logarithmic moments are integrable.  Finally,
\[
 \int_0^\infty Q_j(u)e^{-u}u^s\,du
       =(-s)^j\Gamma(1+s).
\]
This identity follows by induction and integration by parts, first
for $\Re s>0$ and then near zero by analyticity.  It identifies the
resulting coefficients as those of $(-d/dL)^jF(L)$ for
$(t\,d/dt)^jY$, proving the first part of
\eqref{p04:ci20:asymptoticpair}.

For the second part put $h(u)=u/(e^u-1)$ and
\[
 h_j(u)=(u\,d/du)^jh(u),\qquad w_j(u)=-h'_j(u).
\]
Each $w_j$ is bounded near zero and decays exponentially times a
polynomial at infinity.  Since
$Z(t)=\int h(tx)\,dM(x)$, integration by parts gives
\begin{equation}\label{p04:ci20:Zintegral}
 (t\,d/dt)^jZ(t)=\int_0^\infty w_j(u)M(u/t)\,du,
\end{equation}
where $M(x)=0$ for $x<2$.  The boundary terms vanish.  The part
$u<t^{1/2}$ is $O_j(t^{1/2}(1+\log L))$ and the part
$u>t^{-1/2}$ is smaller than every inverse power of $L$.
In between, \eqref{p04:ci20:Mertenserror} replaces $M(u/t)$ by
$\log(L+\log u)+M$ with an error $O_A(L^{-A})$.
Taylor expansion of this logarithm is valid uniformly because
$|\log u|\le L/2$; its integrated remainder has the requested
inverse-logarithmic order.

The moment-generating identity is
\begin{equation}\label{p04:ci20:wjmoments}
 \int_0^\infty w_j(u)u^s\,du=(-s)^j H(s).
\end{equation}
For $j=0$, integration by parts for $\Re s>0$ gives
$s\int_0^\infty u^s/(e^u-1)\,du
=s\Gamma(s+1)\zeta(s+1)$; geometric expansion of the denominator
justifies the latter identity by absolute convergence.  Repeated
integration by parts gives the factor $(-s)^j$.  Both sides are
analytic near zero, so \eqref{p04:ci20:wjmoments} holds there too.
For $j=0$ it yields precisely the series $G$ in
\eqref{p04:ci20:formalFG}; for general $j$ it yields its $j$th derivative
with the sign corresponding to $t\,d/dt=-\partial$.
This proves the second part without differentiating an asymptotic
error term.
\end{proof}

\subsection{No linear differential transfer even with a moment coefficient}

We record the coordinate estimate needed to include $A_r$ in the
coefficient field.  If
$(e^x-1)^{-r}=\sum_{k=1}^r\kappa_{r,k}x^{-k}+O(1)$, then
$\kappa_{r,1}=(-1)^{r-1}$, by substituting $v=e^x-1$ in its residue.
Subtract from $(e^x-1)^{-r}$ the powers with $k\ge2$ and the function
$\kappa_{r,1}e^{-x}/x$.  The residual $g_r$ extends smoothly to zero,
and $g_r,g'_r$ are integrable on $[0,\infty)$.  Comparing the integral
with the right-endpoint Riemann sum gives
$\sum_{n\ge1}g_r(nt)=t^{-1}\int g_r+O(1)$.  Thus, with
$C_r=\int_0^\infty g_r(x)\,dx$,
\begin{equation}\label{p04:ci20:X}
 X_r(t):=\kappa_{r,1}^{-1}\left(tA_r(e^{-t})
 -\sum_{k=2}^r\kappa_{r,k}\zeta(k)t^{1-k}-C_r\right)
       =L+O(t).
\end{equation}
This supporting coordinate estimate is the same one used for the
preceding functional algebraic-independence result; it is not counted
as a new result here.

\begin{theorem}[Independence of the prime derivative families]
\label{p04:ci20:main}
Fix $\lambda>0$ and an integer $r\ge1$.  Over the coefficient field
$\mathscr K_\lambda(A_r(e^{-t}))$, the following germs are linearly
independent:
\begin{equation}\label{p04:ci20:family}
 1,\quad Y,\partial Y,\partial^2Y,\ldots,
       \quad Z,\partial Z,\partial^2Z,\ldots.
\end{equation}
Equivalently, replace $Y,Z$ by $P(e^{-t}),B(e^{-t})$ and
$\partial$ by $d/dt$ in this assertion.  In particular there is no
nontrivial finite inhomogeneous linear differential relation between
the two prime series with coefficients rational in $A_r$ and elements
of this cusp field.
\end{theorem}

\begin{proof}
The fields obtained by adjoining $A_r(e^{-t})$ and $X_r(t)$ are the
same, by \eqref{p04:ci20:X}.  Clear the finitely many denominators of a
putative relation.  Its coefficients can then be written as polynomials
in $X=X_r$ with coefficients in $\mathscr K_\lambda$, not all zero.
Among all the nonzero cusp-field coefficients choose the least
exponential order $m$, then the least power of $t$, say $a$, at that
order.  Divide by $t^ae^{-m\lambda/t}$ and take the leading constants
of those coefficients.  This produces a nontrivial relation expression
\begin{equation}\label{p04:ci20:leadingrelation}
 C(X)+\sum_{j=0}^J U_j(X)\partial^jY
          +\sum_{j=0}^J V_j(X)\partial^jZ,
 \qquad C,U_j,V_j\in\mathbb C[X],
\end{equation}
which the original exact relation would force to be smaller than
every inverse power of $L$.  Indeed the discarded coefficients are
$O(t)$ after normalization, or exponentially small in $1/t$ times
fixed powers of $t^{-1}$; all the functions multiplying them have
at most a fixed power of $L+\log L+1$ as growth.

By \eqref{p04:ci20:X}, replacing $X$ by $L$ changes
\eqref{p04:ci20:leadingrelation} by less than every inverse power of $L$.
Lemma~\ref{p04:ci20:primeexpansions} gives its formal expansion
\[
 C(L)+\sum_jU_j(L)F^{(j)}(L)+\sum_jV_j(L)G^{(j)}(L).
\]
Lemma~\ref{p04:ci20:formalindependence} says this is nonzero.  It therefore
has a first nonzero power of $L$, or a first nonzero power times
$\log L$; any such leading term contradicts the required smallness.
This proves \eqref{p04:ci20:family}.

Multiplication of each target by $t$, and replacement of $d/dt$ by
$-t\,d/dt$, induce invertible triangular changes between any finite
initial derivative families, with coefficients in
$\mathscr K_\lambda$.  This proves the stated equivalence.
\end{proof}

\begin{corollary}[Rational Lambert coefficients]
\label{p04:ci20:Lambertcoefficients}
The family $1,Y,\partial Y,\partial^2Y,\ldots$ is linearly independent
over the field
\[
 \mathscr K_\lambda(A_r(e^{-t}),B(e^{-t})).
\]
\end{corollary}

\begin{proof}
Clear denominators, use $X=X_r$ and $Z=tB(e^{-t})$, and select the
least cusp orders exactly as in Theorem~\ref{p04:ci20:main}.  A putative
relation would force a nonzero polynomial-in-$Z$ expression
\[
 \sum_{h=0}^H Z^h
   \left(C_h(X)+\sum_{j=0}^J U_{hj}(X)\partial^jY\right),
 \qquad C_h,U_{hj}\in\mathbb C[X],
\]
to be smaller than every inverse power of $L$.
For each nonzero coefficient expression, the formal independence of
$1,F,F',\ldots$ in Lemma~\ref{p04:ci20:formalindependence} gives a
nonzero first Laurent term $c_hL^{b_h}$.  By
\eqref{p04:ci20:X} and Lemma~\ref{p04:ci20:primeexpansions}, the actual
coefficient has this same leading equivalent.  Choose the largest
$b_h$ and then the largest $h$ among ties.  Since $Z\sim\log L$,
the corresponding term $c_hL^{b_h}(\log L)^h$ uniquely dominates.
It cannot be smaller than every inverse power of $L$.
\end{proof}

For a finite collection of integer dilations $d$, take
$D=\operatorname{lcm}d$ and $\lambda=4\pi^2/D$.  The usual Eisenstein
transformation laws imply that all
$\mathcal E_2(q^d),\mathcal E_4(q^d),\mathcal E_6(q^d)$ belong to
$\mathscr K_\lambda$ after $q=e^{-t}$: their transformed arguments
are integer powers of $e^{-4\pi^2/(Dt)}$, their prefactors are powers
of $t$, and the weight-two correction is a constant multiple of
$t^{-1}$.  Their derivatives also belong to this field.  The
normalizations and differential closure are given in
\cite[Section~1]{p04:ci20:KanekoZagier1995}.  Consequently
Theorem~\ref{p04:ci20:main} excludes the same linear differential relation
when its coefficients are rational expressions in a single $A_r$,
$q$, these finite modular generators, and their derivatives.
Derivatives of $A_r$ are not included in that coefficient field by
this statement.

\subsection{The arithmetic routes checked, and the steps still missing}

For the Erd\H{o}s--Borwein target, the primary $q$-logarithm
approximation literature already proves considerably more than
irrationality.  For example, Zudilin's Theorems~1--2 and their
corollary \cite{p04:ci20:Zudilin2005} give quantitative rational linear
independence for appropriate triples involving $A_1(1/b)$,
$A_1(1/b^2)$, and the $q$-dilogarithm, for integer $b\ge2$.
Those theorems are not new consequences of our work.  They estimate
linear forms; they do not give the estimates for polynomials of
unbounded degree in $E$ that would prove its transcendence.  Finite
irrationality measure alone does not distinguish an algebraic
irrational from a transcendental number.

There is also an exact distinction between the functional equation
used by the ordinary $q$-logarithm and its prime-restricted version.
For a fixed integer $b\ge2$, define
\[
 F_b(x)=\sum_{n\ge1}\frac{x^n}{b^n-1},\qquad
 F_{b,\mathbb P}(x)=\sum_p\frac{x^p}{b^p-1}.
\]
Both are analytic for $|x|<b$, and for $|x|<1$ direct subtraction gives
\begin{equation}\label{p04:ci20:shift}
 F_b(bx)-F_b(x)=\frac{x}{1-x},\qquad
 F_{b,\mathbb P}(bx)-F_{b,\mathbb P}(x)=P(x).
\end{equation}
Thus the rational forcing term in the ordinary shift equation is
replaced by the prime generating function.  The identity itself is
elementary, but it identifies exactly why one cannot simply insert
prime support into the ordinary $q$-logarithm proof.  The new theorem
shows that this forcing function cannot be eliminated by the stated
finite linear differential constructions, even when $A_r$ and the
prime Lambert function are available.  It does not rule out other
$q$-difference or arithmetic approximation methods.

For the prime Lambert target, the irrationality theorem of Tao and
Ter\"av\"ainen \cite{p04:ci20:TaoTeravainen2026} is already an external
result.  Its connection to this target is exact:
\[
 B(1/2)=\sum_{n\ge1}\omega(n)2^{-n},
\]
obtained by interchanging the absolutely convergent divisor sums.
We do not re-count this settled irrationality statement as a new
closure.  For $P_2$, our previous prime-word complexity analysis and
the present differential obstruction likewise do not prove
transcendence at $q=1/2$.

Finally, the Borel argument above must not be upgraded silently to
nonlinear differential transcendence.  It uses the exact linear rules
$w\mapsto d/ds$ and $d/dw\mapsto-s$.  Products of formal series
produce convolution under this transform, not pointwise products.
The distinct pole sets therefore prove the asserted linear
independence, but do not by themselves exclude nonlinear differential
polynomials.  Both that stronger functional question and the three
fixed-value transcendence questions remain unresolved by this section.

\section{What would constitute the next arithmetic step?}
\label{p04:root4:closing}

The results above close several precisely stated family questions.
They do not establish a new proof of irrationality for Euler's constant,
Euler--Gompertz, or Poisson entropy, nor a new proof of transcendence
for the three prime and reciprocal-Mersenne targets. The following
remaining steps are substantive mathematical requirements, rather than
requests for additional numerical precision.

For the linear excursion optimum, a squarefree-factorization result
beyond the currently applicable partial theorem would be needed for
the route investigated here. The associated function fails global
boundedness and contains infinitely many simple zeros, alongside its
multiple zeros. Neither failure can be removed by checking a larger
finite collection of Taylor coefficients or zeros.

For the biquadratic optimum, the field-degree argument must control
the three powers in \eqref{p04:sc4:eq:missing}, without relying on a fourth
power that is unconstrained by the contact. The full-jet identity
\eqref{p04:sc4:eq:alljets} proves that repeatedly differentiating with
respect to the setting cannot create this missing arithmetic coordinate.
The cubic-field theorem from the preceding paper therefore does not
silently become a theorem for every three-dimensional exponent span.

For fractional extraction rates, a third independent arithmetic
constraint or a different specialization argument is needed. The
rank-three theorem does yield normalized-value and location arithmetic,
but two contact equations leave a one-dimensional state space. The
all-real-shape frontier theorem supplies geometric control, not the
missing transcendence constraint.

For the prime constants, the derivative-family obstruction is a
functional theorem. A new arithmetic specialization or approximation
argument must control values at $q=1/2$. Conversely, the obstruction
does not exclude nonlinear differential relations or all possible
$q$-difference constructions. It excludes only the exactly stated
coefficient fields and linear differential relations.

The most complete positive result in this continuation is the full
planar independence theorem: it removes an infinite set of missing
composite indices and proves a genuinely joint assertion. The binary
bottleneck arithmetic transition is a second exact classification,
with a particularly short transcendence mechanism once its global
optimization structure is established. These are concrete candidates
for specialist review and a more focused publication project.
Their journal level and historical novelty remain questions for that
review; they cannot be inferred from the number of targets or pages.

\paragraph{Verification record.}
Each of the six subject sections was cross-read independently within
this research session. Particular checks covered the Newton polygon
and representation argument, arbitrary finite encoder alphabets,
branch accessibility in the non-$D$-finite proof, normalization in
$E$-function specialization, fixed-exponential elimination, cyclic-cover
ramification, and leading-term extraction with flat coefficients.
Compilation and page rendering check the artifact, not the truth of
its theorems. No external referee decision, machine-checked proof,
or exhaustive novelty certification is claimed.

% END SOURCE: KEIO_Arithmetic_Rigidity_of_Optimization.tex

\endgroup
\clearpage
\part{Structural Rigidity and Adaptive Recovery}
\label{merge:part05}
\begingroup

% BEGIN SOURCE: KEIO_Structural_Rigidity_from_Adaptive_Observation.tex

\noindent\textbf{Editorial integration note.} 
The exact, adaptive, approximate, and noisy observation models used here are distinct. The finite-entropy counterexample explains why mixing does not preserve every earlier pure-saddle classification. The following part tests transfers to logic without treating recovery alone as a decision procedure.


\begin{quote}\small
\noindent\textit{Original part framing.}\par
\noindent\textbf{Status.} Research working manuscript with complete
arguments for the stated deductions and explicit external inputs.
Internal proof cross-reading is not external peer review or formal
verification. The targeted literature comparison does not certify global
priority. Equivalent formulations and routine consequences are not
counted as independent discoveries.
\end{quote}
\section{What changes when arithmetic is no longer the only objective}
\label{p05:rv5:scope}

The 223-page foundation \cite{p05:rv5:223}, its 447-page collected sequels
\cite{p05:rv5:447}, and the subsequent arithmetic manuscripts supply several
different kinds of observations: replicated random variables, optimized
value curves, adaptive scalar measurements, and morphisms of special-value
rings. The preceding manuscript \cite{p05:rv5:IV} concentrated on arithmetic
rigidity. Here we ask whether those observations determine hidden structure
exactly, how much information they require, and which obstructions survive
changes of model.

The strongest transfers remove hypotheses or give necessary and sufficient
conditions for entire families. They are not obtained by listing many
special numerical constants. The six technical sections are independent
enough to permit separate scrutiny; their common theme is to find an
observation coordinate in which an otherwise nonlinear ambiguity becomes
a rigid support, divisor, or fiber constraint.

\subsection{Two principal classification results}

Let $(Q,R)$ have an arbitrary probability law on
$\mathbb R\times(0,\infty)$, and conditionally on it let
\[
                 Y_j=Q+R\log E_j,\qquad j=1,2,
\]
where $E_1,E_2$ are independent unit exponential variables.
Theorem~\ref{p05:ga21:two} proves
\[
 \boxed{\mathcal L(Y_1,Y_2)\ \text{uniquely determines}\ \mathcal L(Q,R).}
\]
The theorem actually proves injectivity for finite complex mixing measures.
There is no compact support, moment, or scale restriction. In the equivalent
transform $\exp\{-B(u^\alpha+v^\alpha)\}$, it removes the earlier lower
bound on $\alpha$ while also permitting arbitrary positive Weibull shapes.
The improvement is substantive: the old triangular argument needed a first
nonzero scale, whereas the new weighted contraction starts even when the
support accumulates at zero. A nonzero Fourier frequency supplies the
exponentially small gamma factor that makes the contraction possible.

This is population identifiability from a joint law, not estimation from
one realized pair of numbers. One replicate fails by stable subordination.
Nor does injectivity imply a uniformly well-conditioned inverse:
Proposition~\ref{p05:ga21:unstable} gives latent laws a fixed distance apart
whose observed laws become arbitrarily close for every fixed group size.

For the second result, let
\[
 S(x)=\sum_{j=1}^Nw_jx^{r_j},\qquad
 T(C)=\min_{x>0}\frac{1+Cx}{S(x)},\qquad
 w_j>0,\quad 0<r_1<\cdots<r_N<1.
\]
Theorems~\ref{p05:sv5:thm:dfinite} and~\ref{p05:sv5:thm:order} give
\begin{align*}
 T\text{ is $D$-finite}
   &\Longleftrightarrow r_i-r_j\in\mathbb Q\quad\text{for all }i,j,\\
 \operatorname{ord}_{\rm alg.diff}(T)
   &=\dim_{\mathbb Q}\operatorname{span}_{\mathbb Q}
                         \{1,r_1,\ldots,r_N\}-1.
\end{align*}
Thus linear differential finiteness depends on exponent differences,
whereas minimal nonlinear differential order depends on a rational rank
that also includes $1$. The previous two-atom theorem did not by itself
handle collisions of complex zeros or repeated zeros for arbitrary $N$.
The present proof does: Skolem--Mahler--Lech controls zero progressions,
and every projected zero produces an accessible pole or branch point of
the optimized value.

\subsection{Exact gains and their limits}

\small
\begin{longtable}{L{0.19\textwidth}L{0.32\textwidth}L{0.38\textwidth}}
\toprule Direction & Previous restriction or question & Present conclusion\\\midrule
\endfirsthead
\toprule Direction & Previous restriction or question & Present conclusion\\\midrule
\endhead
\bottomrule\endfoot
Latent mixtures & Two-replicate inversion with index bounded away from zero & Arbitrary joint mixing without cutoffs or moments; exact threshold two in the stated Weibull and centered stable models.\\[5pt]
Finite spectra & Two-atom differential classification & Necessary and sufficient $D$-finite criterion and exact nonlinear order for every finite number of atoms.\\[5pt]
Observation design & Classical moment certificates; nonlinear optimized observations & Exact jet budget $2N$, or $2N-1$ with known mass away from contact $x=1$; sharp local square-root stability.\\[5pt]
Entropy programs & Rational-weight elimination and special binary arithmetic transitions & A terminating existence test for algebraic optimizers in the promised smooth concave class with algebraic weights.\\[5pt]
Mixed entropy games & Pure-saddle values lie in logarithms of positive algebraic numbers & Rational Shannon games encode finite logarithmic matrices; the pure-saddle range fails under mixing.\\[5pt]
Special-value fields & Rational-expression intersection and one-coordinate closure & Exact Kummer algebraic closure of exponential fields, plus a joint multivariate exceptional-coordinate theorem.\\[5pt]
Gamma extraction & Locally finite contact branches for general real order & No switches for $1/2<\kappa\le1$; finitely many overall for each order; a two-order mixture has an atomic plateau and a genuine corner.\\[5pt]
Adaptive topology & Arbitrary-accuracy recovery in Euclidean dimension & Exact recovery lies between $d$ and $d+1$ on nonempty compact metric spaces of dimension $d$; closed $d$-manifolds require $d+1$ for $d\ge1$.\\[5pt]
Sparse recovery & Fixed linear measurements and stable/noisy models & Exactly $s+1$ adaptive linear queries for all $s$-sparse vectors in $\mathbb R^m$, $s<m$, with a finite final-query menu; $s$ queries in the stated class have error at least $R/2$ on the radius-$R$ set.\\
\end{longtable}
\normalsize

The last two conclusions concern exact real-valued observations and allow
discontinuous dependence of the next query on the preceding transcript.
The compact-space construction nevertheless has a Borel selection rule
and decoder; the sparse construction has a semialgebraic rule.
Neither result asserts numerical stability, finite precision, or a fast
implementation. In fact exact and arbitrarily accurate recovery differ:
on $S^3$, three queries have error infimum zero but cannot attain exact
recovery, which requires four.

\subsection{Eighteen target questions}

The first twelve rows are resolved within the precise models proved below.
They are not twelve externally documented famous conjectures. The last six
record limits encountered during the same investigation. Multiple
corollaries of one mechanism are grouped in one row.

\small
\begin{longtable}{L{0.035\textwidth}L{0.32\textwidth}L{0.53\textwidth}}
\toprule No. & Question & Outcome\\\midrule
\endfirsthead
\toprule No. & Question & Outcome\\\midrule
\endhead
\bottomrule\endfoot
1 & Can two replicates identify arbitrary joint index--scale mixing down to index zero? & Yes, Theorem~\ref{p05:ga21:two} and Corollary~\ref{p05:ga21:threshold}; no atom at index zero is allowed, but accumulation there is unrestricted.\\[4pt]
2 & Which finite positive spectra produce $D$-finite optimized values? & Exactly the rational-gap spectra, Theorem~\ref{p05:sv5:thm:dfinite}.\\[4pt]
3 & What is the least nonlinear algebraic differential order? & Exactly rational rank minus one, Theorem~\ref{p05:sv5:thm:order}.\\[4pt]
4 & How many value derivatives certify finitely many atoms against every positive continuous competitor? & Sharp normalization-dependent orders and local stability, Theorems~\ref{p05:sv5:thm:certificate} and~\ref{p05:sv5:thm:stability}.\\[4pt]
5 & Does the pure-saddle arithmetic range survive arbitrary mixed strategies? & No. A universal finite-logarithmic-game embedding gives two explicit counterexample candidates with an unconditional at-least-one conclusion, Theorems~\ref{p05:ms5:thm:embedding} and~\ref{p05:ms5:thm:mixedboundary}.\\[4pt]
6 & Can algebraic-optimum existence be decided with algebraic tradeoff weights? & Yes in the specified smooth concave rational entropy class, Theorem~\ref{p05:ms5:thm:decision}; no complexity bound is claimed.\\[4pt]
7 & What is the algebraic closure of an exponential field inside $\operatorname{Frac}(\mathcal E)$? & Exactly its rational-exponent Kummer closure, with the ring intersection and integral closure, Theorem~\ref{p05:ar21:expclosure}.\\[4pt]
8 & Can an algebraic extension depend jointly on $\operatorname{Ein}$ coordinates at $\mathbb Q$-linearly independent algebraic parameters $a_i\ge1$? & Every possible nontrivial finite extension is a cyclic radical in at most one exceptional coordinate, Theorem~\ref{p05:ar21:joint}.\\[4pt]
9 & Can low-order extraction contacts switch, and can all-order switches accumulate at endpoints? & No switches for $1/2<\kappa\le1$; only finitely many for each fixed general order, Theorem~\ref{p05:be21:unique} and Proposition~\ref{p05:be21:finite}.\\[4pt]
10 & Can a richer positive mixture create a genuine extraction phase corner? & Yes. An atomic plateau ends at one corner, with an explicit rational-parameter example, Theorem~\ref{p05:be21:mixture} and Corollary~\ref{p05:be21:example}.\\[4pt]
11 & What changes for exact adaptive recovery on compacta and closed manifolds? & $d\le Q(X)\le d+1$ for nonempty compact metric $X$ of finite dimension $d$; $Q(M)=d+1$ for closed $d$-manifolds with $d\ge1$, Theorems~\ref{p05:in21:dimension} and~\ref{p05:in21:cohomology}.\\[4pt]
12 & What is the exact adaptive linear query count for sparse vectors? & $s+1$ for $1\le s<m$, Theorem~\ref{p05:in21:sparse}; with $s-1$ initial linear queries and one continuous final query, radius-$R$ worst-case error is at least $R/2$, Proposition~\ref{p05:in21:sparse-error}.\\[4pt]
13 & Is two-query arbitrary-accuracy recovery possible on all of $\mathbb R^3$? & Unresolved. The exact cohomology obstruction does not yield a positive error lower bound; \eqref{p05:in21:unattained} shows why.\\[4pt]
14 & Do the full $E$-value unit and factorization conjectures follow? & No. A hypothetical unit can remain transcendental over the exponential field; the closure theorem supplies no missing algebraicity hypothesis.\\[4pt]
15 & Do all-index stopping-value independence results extend to every higher even dimension? & Unresolved. The unshifted Schur criterion used for the planar case does not apply to the shifted-factorial coefficients arising in the proposed extension; no replacement irreducibility theorem is established here.\\[4pt]
16 & Which of the two explicit mixed-game values lies outside the logarithmic algebraic span? & Unresolved individually. The at-least-one conclusion uses the established strong six-exponentials theorem and does not prove four exponentials.\\[4pt]
17 & Is the optimal Chebyshev--Bernstein exponent $\log2$? & Unresolved. An apparently stronger contour shortcut is disproved exactly in Section~\ref{p05:rv5:failed}; the older constant-one bound is not counted anew.\\[4pt]
18 & What are the minimal group sizes with unknown stable location or skewness? & Unresolved. Five and six remain sufficient after removing the index cutoff; minimality is not proved.\\
\end{longtable}
\normalsize

\subsection{Dependencies, terminology, and priority}

``Determines'' refers to injectivity on the stated class, not automatically
to a statistically stable or computable inverse. An $L$-jet includes the
value and all derivatives through order $L$. $D$-finiteness means a linear
differential equation with polynomial coefficients; algebraic differential
order permits nonlinear equations and is a different invariant.

The principal external inputs are explicitly credited in their sections:
Skolem--Mahler--Lech and Hadamard factorization; classical moment
certificates and the separated-reference moment comparison of Wu--Yang;
Baker, Gelfond--Schneider and Roy; real algebraic geometry; the
Fischler--Rivoal morphism for $E$-values; and the Hurewicz, sheaf-cohomology,
and van Kampen--Flores theorems. The gamma-extraction section also states
the exact analytic inputs inherited from the earlier project manuscripts.
No external input becomes a new theorem of this paper merely by being
used in a different model.

In particular, the finite-mixture grouped-sample literature differs from
the unrestricted mixing problem in Section~\ref{p05:ga21:section}, and the
fixed-scale Gumbel example in Powers--Xu differs from the arbitrary joint
location--scale pair theorem. Those distinctions motivate further
literature review; they are not themselves proof of priority.

\section{Two observations identify unrestricted Weibull and stable mixtures}
\label{p05:ga21:section}

The preceding identification theorem required the stable index to
remain bounded away from zero.  Here that restriction is removed.
In fact the proof identifies an arbitrary joint location--scale
mixing law for the log-exponential kernel from two conditionally
independent observations.  Its scale may approach both zero and
infinity, and no latent moment is assumed.  The equivalent Weibull
formulation allows all positive shapes and rates; the stable
formulation allows the entire index range $(0,2]$.

The new mechanism is a weighted uniqueness principle for integer
dilations.  Ordinary recursion needs a first nonempty support
interval.  A nonzero Fourier frequency supplies exponential
flatness at the missing endpoint, and a contraction recovers that
initial interval.  The Gamma estimate establishing this flatness
also justifies the infinite-series exchanges at arbitrarily small
shape.  All identification statements concern the population
\emph{joint law} of observations sharing a latent draw, not recovery
of a distribution from a single realized pair.

\subsection{The unrestricted theorem}

Let $E_1,E_2$ be independent mean-one exponential variables.  For a
finite complex Borel measure $\rho$ on
$\mathbb R\times(0,\infty)$ define the induced finite complex
observation measure by mixing the conditional pair
\begin{equation}\label{p05:ga21:logkernel}
               Y_j=Q+R\log E_j,\qquad j=1,2.
\end{equation}
For probability mixing laws, the latent pair $(Q,R)$ is independent
of $(E_1,E_2)$.

\begin{theorem}[Two observations without support or moment restrictions]
\label{p05:ga21:two}
The map from $\rho$ to the joint observation measure in
\eqref{p05:ga21:logkernel} is injective on all finite complex measures
on $\mathbb R\times(0,\infty)$.
Equivalently, the map
\begin{equation}\label{p05:ga21:pairtransform}
 \mu\longmapsto
 \left\{\int e^{-B(u^\alpha+v^\alpha)}\,d\mu(\alpha,B):u,v>0\right\}
\end{equation}
is injective on all finite complex measures on $(0,\infty)^2$.
No support cutoff, moment, finite-component assumption, or latent
independence is required.
\end{theorem}

The change of coordinates between the two formulations is
\begin{equation}\label{p05:ga21:coordinates}
              R=\alpha^{-1},\qquad Q=-\alpha^{-1}\log B.
\end{equation}
Indeed the upper-orthant probabilities of \eqref{p05:ga21:logkernel}
are $\exp\{-B(e^{\alpha x}+e^{\alpha y})\}$.
Upper orthants determine finite complex measures by inclusion and
exclusion and finite-variation limits at unbounded endpoints.
Thus the equivalence is valid for complex measures as well as
probability laws.

\begin{corollary}[Exact replication threshold for Weibull and stable mixing]
\label{p05:ga21:threshold}
For Weibull kernels with survival function
$\mathbb P(W>x\mid\alpha,B)=e^{-Bx^\alpha}$, $x>0$,
two conditionally independent observations identify every Borel
mixing law on $(\alpha,B)\in(0,\infty)^2$.
For symmetric stable kernels with characteristic function
$\exp\{-B|u|^\alpha\}$, the same conclusion holds on the entire
range $(\alpha,B)\in(0,2]\times(0,\infty)$.
One observation does not identify all probability mixing laws in
either class.  Consequently two is the exact minimal number.
\end{corollary}

The two formulations in this corollary are consequences of the same
transform theorem and are counted as one identification family.
The unrestricted log-exponential statement also applies, after
reflection, to ordinary Gumbel location--scale mixtures.

\begin{proposition}[No uniform inverse stability on the unrestricted class]
\label{p05:ga21:unstable}
For any fixed group size $k\geq1$, there are two sequences of
probability mixing laws separated by a fixed bounded-Lipschitz
distance, whose induced $k$-observation laws converge to one
another in total variation.  Thus exact identification does not
provide a uniform inverse modulus on the unrestricted parameter
space.
\end{proposition}
\begin{proof}
Use the mixing point masses $\rho_R=\delta_{(0,R)}$ and
$\widetilde\rho_R=\delta_{(1,R)}$, with $R\to\infty$.
Equip the parameter space with the bounded metric
\[
 d((Q,R),(Q',R'))=
       \min\{1,|Q-Q'|+|\log R-\log R'|\}.
\]
The function $h(Q,R)=\max\{0,\min\{1,Q\}\}$ is bounded by
one and is one-Lipschitz for this metric.  Its integrals under
the two mixing laws differ by one.

The density of $\log E$ is $g(x)=e^{x-e^x}$ and
\[
 \int_{\mathbb R}|g'(x)|\,dx
       =\int_0^\infty|1-u|e^{-u}\,du=\frac2e.
\]
The one-observation laws under the two mixing point masses are
translations by one at scale $R$, so their total-variation
distance is at most $1/(eR)$.  The telescoping product bound
gives distance at most $k/(eR)$ for the joint $k$-observation
laws.  This tends to zero while the latent bounded-Lipschitz
separation remains at least one.
\end{proof}

\subsection{A weighted principle for dilation inversion}

\begin{lemma}[Starting a dilation recursion at an accumulation point]
\label{p05:ga21:weighted}
Let $\eta$ be a locally finite complex measure on $(0,\infty)$,
with finite variation on $(0,L]$ for every $L<\infty$.
Let $c_m$ be Borel functions,
$c_1=1$, and set $D_m(\alpha)=m\alpha$.
Suppose the series
\begin{equation}\label{p05:ga21:M}
                 M=\sum_{m\geq1}(D_m)_*(c_m\eta)
\end{equation}
converges in total variation after multiplication by $e^{-dx}$
for every $d>0$, and $\int e^{-dx}\,dM(x)=0$ for every $d>0$.
Suppose also that there exist $\varepsilon>0$, a positive Borel
weight $w$, and $\kappa<1$ such that
\begin{align}
 &\int_{(0,\varepsilon]}\frac{d|\eta|(\alpha)}{w(\alpha)}<\infty,
                                                        \label{p05:ga21:weightedfinite}\\
 &\sum_{m\geq2}\sup_{0<\alpha\leq\varepsilon/m}
     |c_m(\alpha)|\frac{w(\alpha)}{w(m\alpha)}\leq\kappa.
                                                        \label{p05:ga21:contraction}
\end{align}
Then $\eta=0$.
\end{lemma}
\begin{proof}
First $M=0$ by uniqueness of the Laplace transform.  To see this
without a finiteness assumption on the unweighted measure, fix
$d_0>0$, multiply by $e^{-d_0x}$, and push forward by $y=e^{-x}$.
This is a finite complex measure on $[0,1]$ with all nonnegative
integer moments zero.  Polynomial density gives the zero measure.
The positive tilt is invertible on bounded intervals, proving
$M=0$.

On $(0,\varepsilon]$, divide \eqref{p05:ga21:M} by $w$.
Writing $\sigma=\eta/w$ there, the total variation of the sum
of terms with $m\geq2$ is at most $\kappa\|\sigma\|$ by
\eqref{p05:ga21:contraction}.  That same bound justifies the weighted
series exchange.  Since the term with $m=1$ is $\sigma$, the
equation $M=0$ implies
$\|\sigma\|\leq\kappa\|\sigma\|$, and hence
$\eta=0$ on $(0,\varepsilon]$.

If $\eta=0$ on $(0,L]$, all terms with $m\geq2$ contributing
to $(0,2L]$ originate in $(0,L]$ and vanish.  Thus $M=0$
implies $\eta=0$ on $(0,2L]$.  Iteration covers $(0,\infty)$.
\end{proof}

The endpoint hypothesis is material.  The lemma does not assert
that an arbitrary integer-dilation equation is uniquely solvable
without control on mass accumulating at zero.

\subsection{The Gamma estimate that supplies the endpoint control}

\begin{lemma}[A globally integrable series envelope]
\label{p05:ga21:Gamma}
For every real $t\ne0$ and $d>0$, the function
\begin{equation}\label{p05:ga21:S}
 S_{t,d}(\alpha)=\alpha\sum_{m\geq1}
       \frac{|\Gamma(m+1+it/\alpha)|}{(m-1)!}e^{-m\alpha d}
\end{equation}
is bounded on the entire interval $(0,\infty)$.
Moreover, for fixed $0<\theta_0<\pi/2$,
\begin{equation}\label{p05:ga21:flat}
 |\Gamma(2+it/\alpha)|
       \leq(\cos\theta_0)^{-2}e^{-|t|\theta_0/\alpha}.
\end{equation}
\end{lemma}
\begin{proof}
Rotation of Euler's Gamma integral into a ray of argument
$\operatorname{sgn}(T)\theta$ gives, for an integer $k\geq1$,
\begin{equation}\label{p05:ga21:rotation}
 |\Gamma(k+iT)|\leq
            \Gamma(k)(\cos\theta)^{-k}e^{-|T|\theta},
            \qquad 0<\theta<\pi/2.
\end{equation}
Indeed the rotated integrand has the factor
$e^{i\theta\operatorname{sgn}(T)(k+iT)}$; its modulus contributes
$e^{-|T|\theta}$.  The radial integral is bounded by
$\int_0^\infty e^{-x\cos\theta}x^{k-1}\,dx$.
Contour rotation is justified by exponential decay on the large
arc and the vanishing small-arc contribution for $k>0$.
Taking $k=2$ proves \eqref{p05:ga21:flat}.

If $\alpha d\leq1$, choose
$\theta=\sqrt{\alpha d}/2\leq1/2$.
The elementary inequality $-\log\cos\theta\leq\theta^2$
on this interval gives
\[
 \frac{e^{-\alpha d}}{\cos\theta}\leq e^{-3\alpha d/4}.
\]
Using \eqref{p05:ga21:rotation} with $k=m+1$ and
$\Gamma(m+1)/(m-1)!=m$, we obtain
\begin{align*}
 S_{t,d}(\alpha)
 &\leq C\alpha e^{-|t|\sqrt d/(2\sqrt\alpha)}
                 \sum_{m\geq1}m e^{-3m\alpha d/4}\\
 &\leq C_d\alpha^{-1}e^{-|t|\sqrt d/(2\sqrt\alpha)}.
\end{align*}
This is bounded down to $\alpha=0$.  If $\alpha d\geq1$,
Euler's integral without rotation gives
$|\Gamma(m+1+it/\alpha)|\leq m!$, and hence
\[
 S_{t,d}(\alpha)
       \leq\frac{\alpha e^{-\alpha d}}{(1-e^{-\alpha d})^2}.
\]
This is bounded over the remaining range and tends to zero as
$\alpha\to\infty$.
\end{proof}

The nonzero frequency in this lemma is essential to the displayed
bound at zero.  The proof below recovers the missing zero Fourier
frequency only at the end, by continuity.

\subsection{Proof of unrestricted two-observation identification}

\begin{proof}[Proof of Theorem~\ref{p05:ga21:two}]
Let $\rho$ be a finite complex mixing measure inducing the zero
joint observation measure in \eqref{p05:ga21:logkernel}.  Fix
$t\in\mathbb R\setminus\{0\}$ and define the finite complex
measure
\[
             \nu_t(dr)=\int e^{itq}\,\rho(dq,dr).
\]
Fourier-weight the first observed coordinate by $e^{itY_1}$
and project to $D=Y_2-Y_1$.  The resulting zero measure has density
\begin{equation}\label{p05:ga21:density}
 T_t(d)=\int\frac{\Gamma(2+itr)}r
             \frac{e^{d/r}}{(1+e^{d/r})^{2+itr}}\,d\nu_t(r).
\end{equation}
To derive it, write $E_1=z$, $E_2=ze^{d/r}$ in the two
exponential densities.  The Jacobian is $ze^{d/r}/r$, and
the Fourier weight contributes $e^{itq}z^{itr}$.
Integration in $z>0$ gives \eqref{p05:ga21:density}.

Its absolute integrand is bounded by the unweighted conditional
difference density, because $|\Gamma(2+itr)|\leq1$.
This also justifies the density calculation for finite complex
mixing measures by Fubini.  On a compact subinterval of $d>0$,
the same integrand is bounded uniformly over all $r>0$ by
$\alpha e^{-\alpha d}$, where $\alpha=1/r$.
Thus $T_t$ is continuous for positive $d$ and, being a zero
density, vanishes at every such $d$.

Push $\nu_t$ forward by $\alpha=1/r$, calling the resulting
finite complex measure $\lambda_t$.  The binomial expansion at
$e^{-\alpha d}$ gives
\begin{equation}\label{p05:ga21:series}
 e^{itd}T_t(d)=\sum_{m\geq1}\int
   \alpha\frac{(-1)^{m-1}\Gamma(m+1+it/\alpha)}{(m-1)!}
                e^{-m\alpha d}\,d\lambda_t(\alpha),\qquad d>0.
\end{equation}
The factor $e^{-itd}$ cancels exactly in this expansion.
Lemma~\ref{p05:ga21:Gamma} bounds the sum of absolute integrands
uniformly in $\alpha$, so the exchange of series and integration
is valid without any shape or scale moment.

Define
\begin{align}
 d\eta_t(\alpha)&=\alpha\Gamma(2+it/\alpha)\,d\lambda_t(\alpha),
                                                    \label{p05:ga21:eta}\\
 c_m(\alpha)&=
       \frac{(-1)^{m-1}(2+it/\alpha)_{m-1}}{(m-1)!}.
                                                    \label{p05:ga21:cm}
\end{align}
The measure $\eta_t$ is finite on every $(0,L]$ by
\eqref{p05:ga21:flat}.  Its dilation series \eqref{p05:ga21:M} has finite
exponentially tilted total variation by the same envelope used
for \eqref{p05:ga21:series}, and its Laplace transform vanishes.

To verify the contraction at zero, choose
$0<c<|t|\theta_0$ for some fixed $0<\theta_0<\pi/2$,
and put $w(\alpha)=e^{-c/\alpha}$.
Equation \eqref{p05:ga21:flat} gives
\[
 \frac{\alpha|\Gamma(2+it/\alpha)|}{w(\alpha)}
 \leq(\cos\theta_0)^{-2}\alpha
             e^{-(|t|\theta_0-c)/\alpha}.
\]
This is bounded on $(0,\varepsilon]$, so
\eqref{p05:ga21:weightedfinite} holds for every fixed
$\varepsilon>0$.

For $m\geq2$, put $\beta=m\alpha$.  Each factor in the
Pochhammer symbol in \eqref{p05:ga21:cm} has modulus at most
$m(1+|t|/\beta)$.  Using
$m^{m-1}/(m-1)!\leq e^m$ gives
\begin{equation}\label{p05:ga21:ratio}
 |c_m(\alpha)|\frac{w(\alpha)}{w(m\alpha)}
 \leq e\left[e(1+|t|/\beta)e^{-c/\beta}\right]^{m-1}.
\end{equation}
The quantity
\[
 r_\varepsilon=\sup_{0<\beta\leq\varepsilon}
                         e(1+|t|/\beta)e^{-c/\beta}
\]
tends to zero as $\varepsilon\downarrow0$.
Choose $\varepsilon$ so small that
$e r_\varepsilon/(1-r_\varepsilon)<1$.
Summing \eqref{p05:ga21:ratio} proves \eqref{p05:ga21:contraction}.
Lemma~\ref{p05:ga21:weighted} now gives $\eta_t=0$ on $(0,\infty)$.

The multiplier in \eqref{p05:ga21:eta} never vanishes.  Its inverse
is bounded on compact subintervals of $(0,\infty)$, so
$\lambda_t=0$ and then $\nu_t=0$.
This holds for every nonzero $t$.  For each Borel set $A$ of
scale coordinates, dominated convergence shows that
$t\mapsto\int_{\mathbb R\times A}e^{itq}\,d\rho(q,r)$ is
continuous.  It is therefore zero also at $t=0$.
Fourier uniqueness gives $\rho(C\times A)=0$ for all Borel
$C\subseteq\mathbb R$, $A\subseteq(0,\infty)$, proving
$\rho=0$.  The coordinate change \eqref{p05:ga21:coordinates}
proves the equivalent formulation.
\end{proof}

\begin{proof}[Proof of Corollary~\ref{p05:ga21:threshold}]
The Weibull pair has joint survival
$\int e^{-B(u^\alpha+v^\alpha)}\,d\mu$.
The symmetric stable pair has that same expression as its joint
characteristic function when $u,v>0$.
Theorem~\ref{p05:ga21:two} proves both positive assertions.

For minimality choose $0<a<b$ and a positive stable random
variable $S$ of index $a/b$, whose Laplace transform is
$\mathbb E e^{-sS}=e^{-s^{a/b}}$.  This classical law exists
with L\'evy measure
\[
             \frac{a/b}{\Gamma(1-a/b)}x^{-1-a/b}\,dx;
\]
see also
\cite{p05:ga21:Simon}.  It is strictly positive because its Laplace
transform tends to zero at infinity.  The identity
\[
                  e^{-u^a}=\mathbb E e^{-S u^b},\qquad u>0,
\]
identifies the one-observation laws of the distinct mixing
measures concentrated respectively at $(a,1)$ and at
$(b,S)$.  It proves failure for Weibull mixing; choosing
$a<b\leq2$ proves failure for stable mixing.
\end{proof}

\subsection{Unrestricted stable location and skewness: inherited bounds}

We record the consequence for correlated latent coordinates, while
keeping its observation count separate from the exact threshold
in Corollary~\ref{p05:ga21:threshold}.
Use the stable parametrization
\begin{equation}\label{p05:ga21:stable}
 \widehat K_{A,\alpha,B,\beta}(t)
       =\exp\{itA-B|t|^\alpha+iB\beta\psi_\alpha(t)\},
\end{equation}
where $A\in\mathbb R$, $B>0$, $0<\alpha\leq2$,
$|\beta|\leq1$, and
\[
 \psi_\alpha(t)=
 \begin{cases}
  \tan(\pi\alpha/2)\operatorname{sgn}(t)|t|^\alpha,
                              &\alpha\notin\{1,2\},\\
  -(2/\pi)t\log|t|,          &\alpha=1,\\
  0,                         &\alpha=2.
 \end{cases}
\]
Set $\psi_\alpha(0)=0$ and impose $\beta=0$ at $\alpha=2$
to remove the unobservable Gaussian skewness coordinate.

\begin{corollary}[Finite replication with unrestricted stable parameters]
\label{p05:ga21:anchors}
For arbitrary Borel joint mixing of $(A,\alpha,B)$ in the
symmetric case $\beta=0$, the joint law of five conditionally
independent observations identifies the mixing law.
For arbitrary Borel joint mixing of all four parameters in
\eqref{p05:ga21:stable}, six observations suffice.
No moments, index cutoff, or latent independence is assumed.
These are sufficient bounds, not assertions of minimality.
\end{corollary}
\begin{proof}
The contrast-and-anchor construction is the one in the preceding
paper \cite{p05:ga21:Previous}.  Its only former index restriction
came from the two-observation theorem, now replaced by
Theorem~\ref{p05:ga21:two}.  We give the steps to check that no
new endpoint assumption enters.

In the symmetric case retain $X_1$ as an anchor and use
$D_1=X_2-X_3$, $D_2=X_4-X_5$.
Conditionally on the latent variables, these contrasts are
independent symmetric stable observations with exponent coefficient
$2B$, independent of the anchor.  For a difference $\rho$ of
two candidate mixing laws and fixed real $t$, the anchor weight
$e^{itA-B|t|^\alpha}$ defines a finite complex measure on
$(\alpha,B)$.  The joint transform with $(D_1,D_2)$ places
it in the kernel of \eqref{p05:ga21:pairtransform}; hence it is
zero.  The real multiplier $e^{-B|t|^\alpha}$ can be removed
on bounded $B$ intervals and then by exhaustion.  Fourier
uniqueness in $A$, after restriction to each Borel set of
$(\alpha,B)$, proves the five-observation assertion.

For six observations retain $X_1,X_2$ and use
$D_1=X_3-X_4$, $D_2=X_5-X_6$.  Oddness of $\psi_\alpha$
cancels both location and skewness in the contrasts.  Applying
Theorem~\ref{p05:ga21:two} for each pair of anchor frequencies
$(p,q)$ and removing the positive real multiplier gives,
after disintegration with respect to $(\alpha,B)$,
\begin{equation}\label{p05:ga21:conditional}
 \int e^{i(p+q)A+iB\beta(\psi_\alpha(p)+\psi_\alpha(q))}
                       \,d\rho_{\alpha,B}(A,\beta)=0.
\end{equation}
The identities hold simultaneously outside one null set: first
use countably many rational pairs $(p,q)$, then continuity in
frequencies to obtain all real pairs.

Fix such an $(\alpha,B)$ with $\alpha<2$ and a nonzero sum
$u=p+q$.  It suffices to inspect $u>0$; reflection of both
frequencies gives the same assertion for $u<0$.
As $p$ varies between $0$ and $u$, the value
$B[\psi_\alpha(p)+\psi_\alpha(u-p)]$ ranges over a
nondegenerate interval.  For $\alpha\ne1$ this follows from
the nonvanishing second derivative of
$p^\alpha+(u-p)^\alpha$ on the appropriate same-sign
interval.  For $\alpha=1$ it follows from the strict convexity
of $p\log p+(u-p)\log(u-p)$.
For each fixed $u$, the Fourier transform in $\beta$ of
$e^{iuA}\rho_{\alpha,B}$ is entire, since $|\beta|\leq1$.
By \eqref{p05:ga21:conditional} it vanishes on that interval and
therefore everywhere.  Continuity supplies $u=0$ and Fourier
uniqueness in $(A,\beta)$ gives $\rho_{\alpha,B}=0$.
At $\alpha=2$ the imposed condition $\beta=0$ reduces the
argument directly to Fourier uniqueness in $A$.
\end{proof}

\paragraph{Relation to prior results and remaining questions.}
The compact-index two-observation theorem and the contrast-and-anchor
architecture were proved in the preceding project manuscript
\cite{p05:ga21:Previous}.  The added result here is the removal of
both shape endpoints for log-exponential/Weibull mixing, including
the formerly inaccessible accumulation at stable index zero.
Lemma~\ref{p05:ga21:weighted} isolates the reusable mechanism.
The one-observation subordination identity is classical.

Grouped finite-mixture identification, such as
\cite{p05:ga21:VS,p05:ga21:BJR}, and continuous-mixture transform transfers,
such as \cite{p05:ga21:PX}, provide relevant prior frameworks.  The
present class permits arbitrary continuous mixing of all kernel
coordinates and is not inferred from a bound on the number of
mixture components.  A primary-source search did not verify an
earlier theorem of this exact unrestricted scope; it does not
establish historical priority or the closure of a named external
conjecture.

Minimal replication with unknown stable location or skewness remains
unresolved here.  So do quantitative stability and estimation from
finitely many noisy observations.  Separately investigated
higher-dimensional stopping-polynomial irreducibility also remains
unresolved: its shifted-factorial coefficients do not satisfy the
planar Schur theorem's hypotheses.  These questions are recorded
in the research ledger rather than counted as additional theorems.

\section{Finite spectra: differential complexity and robust certification}
\label{p05:sv5:sec}
Let $\mu$ be a nonzero finite positive Borel measure supported on a fixed
interval $[a,b]\subset(0,1)$, and set
\begin{equation}\label{p05:sv5:eq:model}
 S_\mu(x)=\int_a^b x^r\,d\mu(r),\qquad
 T_\mu(C)=\min_{x>0}\frac{1+Cx}{S_\mu(x)},\qquad C>0.
\end{equation}
For arbitrary finite spectra, we classify linear differential finiteness
and determine the least nonlinear algebraic differential order exactly.
Against arbitrary positive continuous mixtures, we give finite-jet
certification and a sharp local stability bound.

The arithmetic-progression input below is the classical
Skolem--Mahler--Lech theorem \cite{p05:sv5:Lech,p05:sv5:BellLagarias}.
Moment certificates and polynomial interpolation are classical;
Wu and Yang \cite[Lemma~4 and Proposition~5]{p05:sv5:WuYang} give moment
identifiability and, at separated reference atoms, a square-root
comparison against unrestricted measures.
Our assertions concern the nonlinear optimal-value observation,
including its differential classification and the conversion of
its jets into moment certificates. They do not assert a new
general theorem on moment problems.

\subsection{Observation derivatives as spectral moments}
Write $x=x(C)$ for the contact, $H=S-xS'$, and $B=T-CT'$.
Differentiation under the integral gives $S'>0$, $S''<0$, and $H>0$.
The objective diverges at both ends of the positive axis, and
$(S'/H)'=S''S/H^2<0$. Thus the contact is unique and analytic.
The identities needed throughout are
\begin{align}
 C&=\frac{S'}H,& T&=\frac1H,&
 T'&=\frac{x}S,& B&=\frac1S,&
 x&=\frac{T'}B, \label{p05:sv5:eq:contact}\\
 x'&=\frac{TT''}{B^2}<0,&
 S'&=\frac CT,&
 S''&=\frac{B^3}{T^3T''}. \label{p05:sv5:eq:second}
\end{align}
Primes on $S$ and $T$ refer to their respective variables.
Define an operator acting on observed functions by
\begin{equation}\label{p05:sv5:eq:operator}
 {\cal L}_T=\frac{T'B}{TT''}\frac d{dC},\qquad
 m_k={\cal L}_T^k(B^{-1}).
\end{equation}
It is exactly $x\,d/dx$ in the contact coordinate, so
\begin{equation}\label{p05:sv5:eq:moments}
 m_k(C)=\int r^k x(C)^r\,d\mu(r).
\end{equation}

\begin{lemma}[A finite-jet change of coordinates]\label{p05:sv5:lem:jets}
For $L\ge2$, the $L$-jet of $T$ at $C_0$ determines $x_0=x(C_0)$
and the moments $m_0,\ldots,m_L$ by rational formulas.
The first jet already determines $x_0,m_0,m_1$.
Conversely, $x_0$ and $m_0,\ldots,m_L$, with the stationary relation
at $C_0$, determine the same $L$-jet, locally analytically.
All these maps are locally Lipschitz near every positive spectrum.
\end{lemma}
\begin{proof}
The apparent extra derivative in $ {\cal L}_T$ cancels on its first
application:
\[
 m_0=B^{-1},\qquad m_1=\frac{CT'}{TB}.
\]
Applying $ {\cal L}_T$ once more uses only the second jet, and each
further application increases the required order by one.
The denominators $T,B,T''$ are nonzero.

For the converse, the first two moments give $S,S'$ at $x_0$.
Ordinary derivatives of $S$ through order $L$ are invertible triangular
linear combinations of the Euler derivatives $m_0,\ldots,m_L$.
The formulas above recover $T,T',T''$, with
$T''=B^3/(T^3S'')$. Successively differentiating gives the rest:
$d/dC=(H^2/(S''S))d/dx$, and after the cancellation in $T''$,
each application increases the required $S$-order by one.
These are rational expressions with nonzero denominators.
This also proves local Lipschitz continuity.
\end{proof}

\subsection{A classification for every finite positive spectrum}
An exponential sum below has constant coefficients, not polynomial
coefficients.
A germ is called $D$-finite if it satisfies a nonzero linear
differential equation with polynomial coefficients.
\begin{lemma}[Finitely many projected zeros]\label{p05:sv5:lem:projection}
Let $E(t)=\sum_{j=1}^N c_je^{r_jt}$, with distinct real $r_j$
and nonzero complex $c_j$. If $\{e^t:E(t)=0\}$ is finite,
then $r_j-r_1\in\mathbb Q$ for every $j$.
\end{lemma}
\begin{proof}
For $N=1$ the assertion is immediate. No zero for $N\ge2$ has
multiplicity exceeding $N-1$: the first $N$ vanishing derivatives
would give a nonsingular Vandermonde system for $c_je^{r_jt}$.
Choose logarithms $\tau_h$ of the finitely many projected zeros.
All zeros lie among $\tau_h+2\pi in$, $n\in\mathbb Z$.
For each $h$ and $0\le k\le N-2$, the sequence
$E^{(k)}(\tau_h+2\pi in)$ is a linear recurrence over $\mathbb C$.
Skolem--Mahler--Lech makes its zero set a finite union of complete
arithmetic progressions, augmented by a finite set.
The two-sided version follows from
\cite[Theorem~1.6]{p05:sv5:BellLagarias}, or by applying the ordinary
theorem in both directions: a zero progression on a ray extends
to the whole progression by Vandermonde independence of its
distinct characteristic roots.

Take a common period $Q$ for these finitely many progressions.
On each $\tau_h+2\pi i(\ell+Q\mathbb Z)$ the zero multiplicity
has a constant baseline $d_{h\ell}\in\{0,\ldots,N-1\}$, with
only finitely many exceptions, all increases.
The baseline counts the successive initial derivatives that
vanish on the entire progression.

Division of this divisor requires no factorization conjecture.
If a real exponential sum $F$ vanishes on $t_0+2\pi iQ\mathbb Z$,
group its exponents modulo $(1/Q)\mathbb Z$:
\[
 F(t)=\sum_\gamma e^{\gamma t}P_\gamma(e^{t/Q}),
\]
where $P_\gamma$ are Laurent polynomials.
On the progression, $e^{t/Q}=y_0=e^{t_0/Q}$ is constant and the
characteristic numbers $e^{2\pi iQ\gamma}$ are distinct.
Vandermonde independence implies $P_\gamma(y_0)=0$ for every $\gamma$.
Division by $e^{t/Q}-y_0$ therefore leaves a finite real exponential sum.
Repeat this operation for all baseline multiplicities.
The roots $e^{(\tau_h+2\pi i\ell)/Q}$, $0\le\ell<Q$, are distinct
for distinct $(h,\ell)$. We obtain
\[
 E(t)=A(e^{t/Q})R(t),
\]
where $A$ is a polynomial and $R$ a finite real exponential sum
with finitely many zeros.
Hadamard factorization gives $R(t)=e^{\alpha t+\beta}P(t)$ with
a polynomial $P$, since $R$ has order at most one.
If $\rho$ is its largest real exponent, real-axis asymptotics give
$R'/R\to\rho$, hence $\alpha=\rho$.
Also $e^{-\rho t}R(t)$ has a finite nonzero limit as $t\to+\infty$,
forcing $P$ constant.
Thus $E(t)=e^{\rho t+\beta}A(e^{t/Q})$.
Independence of distinct exponentials puts all original exponents
in $\rho+(1/Q)\mathbb Z$, proving the claim.
\end{proof}

\begin{theorem}[Exact linear differential classification]\label{p05:sv5:thm:dfinite}
Suppose
\[
 S(x)=\sum_{j=1}^N w_jx^{r_j},\qquad
 0<r_1<\cdots<r_N<1,\qquad w_j>0.
\]
The optimized-value germ $T(C)$ is $D$-finite over $\mathbb C(C)$
if and only if every gap $r_j-r_1$ is rational.
If some gap is irrational, analytic continuations of the positive
real value germ have infinitely many distinct finite singularities.
In the rational-gap case the contact $x(C)$ is algebraic and the
logarithmic derivative $T'/T$ is algebraic.
\end{theorem}
\begin{proof}
First write $r_j=r_1+k_j/q$ in the rational-gap case, with
nonnegative integers $k_j$ and $q\ge1$. Put
\[
 z=x^{1/q},\quad P(z)=\sum_jw_jz^{k_j},\quad
 R(z)=r_1P(z)+q^{-1}zP'(z),\quad Q(z)=P(z)-R(z).
\]
Stationarity and the value formula become
\begin{equation}\label{p05:sv5:eq:algebraiccontact}
 Cz^qQ(z)-R(z)=0,\qquad T=\frac{z^{-qr_1}}{Q(z)}.
\end{equation}
The polynomial equation is nonzero, so $z(C)$ and $x(C)$ are
algebraic. Logarithmic differentiation puts $T'/T$ in the finite
algebraic field $\mathbb C(C,z)$, which is closed under differentiation.
Writing $T^{(k)}=R_kT$ in that field and using its finite dimension
gives a linear differential equation over $\mathbb C(C)$.
Clearing denominators proves $D$-finiteness.

Suppose now that some gap is irrational.
Lemma~\ref{p05:sv5:lem:projection} gives infinitely many distinct
projected zeros $x_0=e^{t_0}$ of $S(e^t)$.
Near any such $t_0$, $x=e^t$ is a valid coordinate and the
corresponding branch of $S$ is analytic.
Write its zero multiplicity as $m$ and let $\Phi=S'/H$, $\Psi=1/H$.
For $m\ge2$, putting $h=x-x_0$ gives
\[
 \Phi(x)=-\frac1{x_0}
       +\frac{1-1/m}{x_0^2}h+O(h^2).
\]
Its derivative is nonzero, while $H$ vanishes to order $m-1$.
Thus the value has a pole of order $m-1$ at $c_0=-1/x_0$.
For $m=1$, $H(x_0)\ne0$. If $S''$ vanishes to order $k\ge0$
there, the identities
\[
 \Phi'=\frac{S''S}{H^2},\qquad \Psi'=\frac{xS''}{H^2}
\]
show that $\Phi-c_0$ and $\Psi-\Psi(x_0)$ have orders $k+2$
and $k+1$. The order $k$ is finite since $S''$ is not identically
zero. A holomorphic function of $\Phi$ would have order divisible
by $k+2$, so the value is ramified at $c_0$.
Distinct projected zeros give distinct $c_0$.

These branches are accessible from the real germ. If $T$ satisfied
a linear equation, let $\Sigma$ be the finite zero set of its leading
polynomial coefficient. Choose $c_0\notin\Sigma$ and a positive initial
setting outside $\Sigma$. In the $t$-plane remove the poles of
$\Phi(e^t),\Psi(e^t)$, the zeros of $d\Phi(e^t)/dt$, and the
inverse image of $\Sigma\cup\{c_0\}$. This is a locally finite
discrete set. Its complement is path connected, so a path from
the real contact can approach $t_0$ with only its endpoint removed.
Local inversion along the path continues the real value to the
pole or ramified branch just found. Its image avoids $\Sigma$.
A solution of a linear differential equation is holomorphic at
every ordinary point, including all continuations to $c_0$.
This contradiction proves the theorem.
\end{proof}

Repeated zeros are essential here: they give poles rather than the
branch points used in the two-atom calculation.
We do not classify $D$-finiteness of the contact in the irrational-gap
case, because the contact is locally regular at these repeated zeros.

\subsection{The exact nonlinear differential order}
Algebraic differential order is the least $k\ge0$ for which a
nonzero polynomial over $\mathbb C(C)$ vanishes at
$T,T',\ldots,T^{(k)}$.

\begin{theorem}[Rational rank equals differential complexity]
\label{p05:sv5:thm:order}
For the finite spectrum of Theorem~\ref{p05:sv5:thm:dfinite}, put
\[
 R=\dim_{\mathbb Q}\operatorname{span}_{\mathbb Q}
                 \{1,r_1,\ldots,r_N\}.
\]
The least algebraic differential order of $T$ is $R-1$, and
\begin{equation}\label{p05:sv5:eq:field}
 \mathbb C(C,T,T',T'',\ldots)
       =\mathbb C(x,x^{r_1},\ldots,x^{r_N}).
\end{equation}
An explicit equation, of order at most $N$, is
\begin{equation}\label{p05:sv5:eq:euler}
 \prod_{j=1}^N({\cal L}_T-r_j)(B^{-1})=0.
\end{equation}
Conversely, an analytic $T$ on a connected positive interval satisfying
$T,T',B>0$, $T''<0$, and \eqref{p05:sv5:eq:euler} comes from these
exponents with positive fixed weights if and only if the system
\begin{equation}\label{p05:sv5:eq:weights}
 \sum_{j=1}^N A_j r_j^k=m_k,\qquad 0\le k<N,
\end{equation}
has strictly positive entries $A_j$ at one setting.
\end{theorem}
\begin{proof}
Let $K=\mathbb C(x,x^{r_1},\ldots,x^{r_N})$, a field of meromorphic
germs in $t=\log x$. It is stable under $D=d/dt$.
Both $C,T$ belong to $K$, and $d/dC=(DC)^{-1}D$ preserves $K$.
Thus the left field of \eqref{p05:sv5:eq:field} lies in $K$.
Conversely, that field contains $x=T'/B$, $S=B^{-1}$, and
all $D^kS={\cal L}_T^k(B^{-1})$.
The Vandermonde system \eqref{p05:sv5:eq:weights} recovers every
$w_jx^{r_j}$ and proves equality.

The transcendence degree of $K$ over $\mathbb C$ is $R$.
Indeed, exponentials with rationally independent exponents are
algebraically independent as functions: a polynomial relation
would be a linear relation among exponentials with distinct
exponents, contradicted by a Vandermonde matrix.
Clearing denominators in a rational exponent basis changes the
field only by algebraic extensions. Since $C$ is nonconstant,
the relative transcendence degree over $\mathbb C(C)$ is $R-1$.

Let $k$ be the first order at which the jet is algebraically
dependent; such an order exists by finite transcendence degree.
Then $T,\ldots,T^{(k-1)}$ are algebraically independent and
$T^{(k)}$ is algebraic over their field.
Its minimal polynomial has nonzero separant in characteristic zero.
Differentiating that polynomial expresses $T^{(k+1)}$ rationally
in $\mathbb C(C,T,\ldots,T^{(k)})$.
Repetition puts every later derivative in the same field,
whose transcendence degree is $k$. Hence $k=R-1$.
This argument also covers $k=0$.

The Euler equation $\prod_j(D-r_j)S=0$ gives
\eqref{p05:sv5:eq:euler}. Lemma~\ref{p05:sv5:lem:jets} bounds its order by
$N$ after denominators are cleared. It is not a formal zero:
a single power with exponent different from all $r_j$ produces
an admissible curve on which its left side is nonzero.

For the converse, define $x=T'/B$ and $S(x(C))=B^{-1}$.
The inequalities make $x$ locally invertible and give
\eqref{p05:sv5:eq:second}. The equation becomes
$\prod_j(d/dt-r_j)S(e^t)=0$.
Thus $S=\sum_jw_jx^{r_j}$ with real constants $w_j$.
The system \eqref{p05:sv5:eq:weights} gives $A_j=w_jx^{r_j}$,
so positivity at one point is precisely $w_j>0$ for every $j$.
The contact identities then give the unique global minimum.
The constants agree on overlapping intervals.
\end{proof}

Adding atoms need not increase differential order: the controlling
quantity is the rational rank together with $1$.
This statement concerns functions, not algebraic independence of
their values at one specified setting.

\subsection{Exact atomic certification against continuous mixtures}
The competitors in this subsection are all finite positive measures
on $[a,b]$; they need not have a finite support.
An $L$-jet includes the value and derivatives through order $L$.

\begin{theorem}[Exact jet budget, including normalization]\label{p05:sv5:thm:certificate}
Fix $\mu_0=\sum_{j=1}^Nw_j\delta_{r_j}$ with
$a<r_1<\cdots<r_N<b$, $w_j>0$, a setting $C_0>0$, and contact $x_0$.
Among all positive measures in \eqref{p05:sv5:eq:model}, the $2N$-jet
at $C_0$ uniquely determines $\mu_0$, including its atom count.
No smaller jet certifies it if total mass is unrestricted.

If total mass is specified in advance, the exact minimum order
for this certification is
\begin{equation}\label{p05:sv5:eq:budget}
 \begin{cases}
  2N,&x_0=1,\\
  2N-1,&x_0\ne1.
 \end{cases}
\end{equation}
These are pointwise assertions at every reference spectrum satisfying
the stated interior-node conditions.
\end{theorem}
\begin{proof}
For a supplied $2N$-jet, Lemma~\ref{p05:sv5:lem:jets} recovers $x_0$
and the moments of the tilted finite measure
$\rho(dr)=x_0^r\,d\mu(r)$ through order $2N$.
Put $q(r)=\prod_j(r-r_j)^2$. At $\mu_0$ its integral is zero;
any competing measure with the same moments also has
$\int q\,d\rho=0$. Positivity forces support on the same $N$ nodes.
The first $N$ moments give the weights by Vandermonde inversion.
Multiplication by $x_0^{-r}$ recovers $\mu_0$ itself.

This is constructive even when the nodes are initially unknown.
Let $M_s=(m_{i+j})_{0\le i,j\le s}$.
The conditions $M_{N-1}>0$ and $\det M_N=0$ certify exactly
$N$ atoms for an actually positive representing measure.
Indeed, solve for the unique monic degree-$N$ polynomial $p$
orthogonal to $1,r,\ldots,r^{N-1}$ using $M_{N-1}$.
The zero Schur complement gives $\int p^2\,d\rho=0$.
Thus its roots contain the support, while $M_{N-1}>0$ forces
at least $N$ support points. These roots and Vandermonde weights
recover the spectrum.

To prove sharpness without known mass, write the reference tilted
weights as $A_j=w_jx_0^{r_j}$.
The map
\[
 (A_1,\ldots,A_N,r_1,\ldots,r_N)
       \longmapsto \left(\sum_j A_jr_j^k\right)_{k=0}^{2N-1}
\]
has invertible Jacobian. A left null vector would be a polynomial
of degree at most $2N-1$ vanishing with its derivative at all
$N$ distinct nodes, hence the zero polynomial.
Choose a new node $z\in(a,b)$, different from all $r_j$, and give
it tilted weight $\varepsilon>0$. The implicit function theorem
adjusts the old nodes and weights by $O(\varepsilon)$ so that
all moments through order $2N-1$ stay unchanged.
They remain positive and interior. Untilt by $x_0^{-r}$.
The stationary relation at $C_0$ is unchanged, so $x_0$ is still
the contact, and Lemma~\ref{p05:sv5:lem:jets} gives identical
$(2N-1)$-jets. The competitor has $N+1$ distinct atoms.
When $x_0=1$, the zeroth moment is total mass, so this construction
already proves sharpness with mass fixed.

Now assume mass is fixed and $t=\log x_0\ne0$.
Equality of $(2N-1)$-jets supplies equal tilted moments through
degree $2N-1$. Let $P$ be the degree-at-most-$2N-1$ Hermite
interpolant to $f(r)=e^{-tr}$, matching both value and derivative
at the $N$ reference nodes. The real interpolation remainder gives
\begin{equation}\label{p05:sv5:eq:hermite}
 f(r)-P(r)=\frac{t^{2N}e^{-t\xi_r}}{(2N)!}
                  \prod_j(r-r_j)^2
\end{equation}
for some $\xi_r\in[a,b]$ when $r$ is not a node.
It is strictly positive off the nodes.
Equal mass is exactly equality of $\int f\,d\rho$; the polynomial
integrals also agree. The reference remainder integral is zero,
so positivity forces the competitor onto the same nodes.

For sharpness at $2N-2$ when $N\ge2$, apply the added-atom
construction to the $2N$ functions
\[
 1,\ e^{tr},\ re^{tr},\ldots,r^{2N-2}e^{tr}
\]
using the original, untilted weights. Its square Jacobian is
invertible. Otherwise a nonzero function
$c+e^{tr}P_{2N-2}(r)$ would have $N$ double zeros.
After multiplication by $e^{-tr}$, repeated Rolle's theorem
would force its derivative of order $2N-1$ to have a zero.
That derivative is $c(-t)^{2N-1}e^{-tr}$, impossible if $c\ne0$;
if $c=0$, the polynomial degree gives the contradiction.
The new measure has the same mass and the same tilted moments
through $2N-2$, and hence the same jet of that order.

For $N=1$, only the failure of the zeroth jet needs separate
verification. Split its atom into equal weights at $r+s-h$ and
$r+s+h$, preserving the prescribed mass.
The optimum is analytic in $s,h$ near zero.
Its derivative in $s$ at zero equals $-T(C_0)\log x_0\ne0$
by the envelope formula. The implicit function theorem gives
$s(h)=O(h^2)$ preserving the single observed value.
For $h>0$ small this is a distinct positive two-atom competitor.
\end{proof}

The reduction in \eqref{p05:sv5:eq:budget} is an effect of observation
geometry: away from $x_0=1$, a known total mass supplies a
nonpolynomial tilted moment. At $x_0=1$ it merely repeats $m_0$.

\subsection{A sharp stability certificate with no atom bound on competitors}
For a finite nonzero measure, let $\bar\mu=\mu/\mu([a,b])$.
For a fixed $C_0$ and integer $L$, define
\[
 d_L(\mu,\nu)=\max_{0\le k\le L}
             |T_\mu^{(k)}(C_0)-T_\nu^{(k)}(C_0)|.
\]
These are derivative observations; bounded errors in function values
alone do not supply derivative errors without a further argument.

\begin{theorem}[Sharp local square-root stability]\label{p05:sv5:thm:stability}
Fix the reference spectrum and setting of
Theorem~\ref{p05:sv5:thm:certificate}. There are $K,\delta_0>0$,
depending on that reference and $[a,b]$, such that every finite
positive competing measure $\nu$ with
$\delta=d_{2N}(\mu_0,\nu)\le\delta_0$ satisfies
\begin{equation}\label{p05:sv5:eq:stability}
 W_1(\bar\mu_0,\bar\nu)\le K\sqrt{\delta}.
\end{equation}
No support-cardinality bound or regularity condition on $\nu$ is used.
The exponent $1/2$ cannot be improved, even for competitors having
the same total mass as $\mu_0$.

More precisely, let $\pi_\nu(dr)=x_\nu^r\,d\nu(r)/S_\nu(x_\nu)$
be its normalized tilted measure at $C_0$, and let
$R_0=\{r_1,\ldots,r_N\}$. Partition $[a,b]$ into nearest-node
cells $V_j$, resolving ties arbitrarily. Then
\begin{align}
 \int \operatorname{dist}(r,R_0)^2\,d\pi_\nu(r)&\le K\delta,
 \label{p05:sv5:eq:concentration}\\
 |\pi_\nu(V_j)-\pi_{\mu_0}(\{r_j\})|&\le K\delta.
 \label{p05:sv5:eq:clustermass}
\end{align}
\end{theorem}
\begin{proof}
The rational formulas of Lemma~\ref{p05:sv5:lem:jets}, on a compact
neighborhood of the reference jet, give
$|x_\nu-x_0|\le K\delta$ and
\[
 \left|\int r^k\,d(\pi_\nu-\pi_{\mu_0})\right|\le K\delta,
                         \qquad 0\le k\le2N.
\]
The constants here are uniform over all competitors whose jets
belong to that neighborhood.
For $q(r)=\prod_j(r-r_j)^2$, it follows that
$0\le\int q\,d\pi_\nu\le K\delta$.
Distinct fixed nodes on a compact interval give
$q(r)\ge c\,\operatorname{dist}(r,R_0)^2$:
relative to a nearest node, every other factor has distance at
least half the minimum node separation. For $N=1$ the inequality
is an equality. This proves \eqref{p05:sv5:eq:concentration}.

Let $h_j$ be the degree-at-most-$2N-1$ Hermite polynomial with
$h_j(r_i)=\mathbf1_{\{i=j\}}$ and $h_j'(r_i)=0$.
Taylor's theorem on each cell gives
\[
 |h_j(r)-\mathbf1_{V_j}(r)|
       \le K\operatorname{dist}(r,R_0)^2.
\]
Combine this with the moment bound and
\eqref{p05:sv5:eq:concentration} to obtain
\eqref{p05:sv5:eq:clustermass}.
Project $\pi_\nu$ to its nearest nodes.
The transport cost is at most $K\sqrt{\delta}$ by
Cauchy--Schwarz, and the projected weights differ from those
of $\pi_{\mu_0}$ by $O(\delta)$. Hence
$W_1(\pi_\nu,\pi_{\mu_0})\le K\sqrt{\delta}$.

Finally the untilting operation
\[
 U_x(\pi)(dr)=
       \frac{x^{-r}\,d\pi(r)}{\int x^{-u}\,d\pi(u)}
\]
is jointly Lipschitz in $W_1$ and in $x$ on a compact positive
$x$-interval. Indeed, the weights, their derivatives in $r,x$,
and reciprocal denominators are uniformly bounded; apply the
dual Lipschitz formula for $W_1$ to the weighted integrals.
Since $U_{x_\nu}(\pi_\nu)=\bar\nu$, this proves
\eqref{p05:sv5:eq:stability}.

For sharpness, replace one reference atom $w_j\delta_{r_j}$ by
$\tfrac12w_j(\delta_{r_j-h}+\delta_{r_j+h})$.
For small $h$ all nodes remain interior and the mass is unchanged.
Taylor's theorem gives an $O(h^2)$ change in the derivatives of
$S$ of any prescribed finite order on a contact neighborhood.
Analytic stationary inversion therefore gives
$d_{2N}(\mu_0,\nu_h)\le K_0h^2$.
The one-dimensional transport cost of this split is exactly
$W_1(\bar\mu_0,\bar\nu_h)=w_jh/\mu_0([a,b])$ for small $h$.
Taking $h=\sqrt{\delta/K_0}$ proves the matching lower modulus.
\end{proof}

If observed jets have coordinate errors at most $\varepsilon$,
any positive candidate within $\varepsilon$ of the data is within
$2\varepsilon$ of the true jet. The theorem is therefore a
deterministic certificate for every such candidate near the fixed
reference. It is not a claim that arbitrary noisy jets are feasible,
nor a uniformly conditioned inverse when reference nodes collide.

\subsection{Relation to existing open problems}
Two relevant existing problems were checked rather than silently
identified with the present classification. Shapiro's common-zero
conjecture for exponential sums is stated in
\cite[Section~2]{p05:sv5:ExponentialSurvey}. It concerns arbitrary pairs
with infinitely many common zeros. Lemma~\ref{p05:sv5:lem:projection}
assumes that all zeros project to finitely many points and divides
explicit periodic progressions; it does not settle that conjecture.
The effective Skolem problem asks for a decision procedure for a
zero in a specified recurrence. Recent work still identifies
effectivity as the missing ingredient \cite{p05:sv5:Skolem2026}.
Our use of Skolem--Mahler--Lech is non-effective and does not solve
that decision problem.

The completed targets are the finite-spectrum functional
classification, exact differential-order law, and local certification
theory in the stated observation model. The contact's own
$D$-finite classification for arbitrary spectra with irrational
gaps, quantitative locations of the complex singularities, and
uniform collision-sensitive jet conditioning remain outside these
results.

\section{Arithmetic changes under mixing and algebraic tradeoffs}
\label{p05:ms5:section}

All logarithms are natural, and $H(p)=-\sum_i p_i\log p_i$, with
$0\log0=0$. Two structural questions motivate this section. First, does
the logarithmic algebraicity of rational entropy programs with a pure
saddle survive passage to mixed strategies? Second, when rational
entropy objectives are combined with algebraic weights, can the existence
of an algebraic optimizer still be decided? We give a negative answer to
the first question and a positive answer to the second in a specified
smooth concave class. The transcendence inputs are classical theorems of
Baker, Gelfond--Schneider, and the six-exponentials theory; we do not claim
a new transcendence theorem or a solution of the four-exponentials
conjecture.

\subsection{Finite logarithmic games inside rational entropy payoffs}

Write $\Delta_r=\{x\in\mathbb R_{\ge0}^r:\sum_i x_i=1\}$ and
\[
 \mathcal L_{\rm rat}
   =\operatorname{span}_{\mathbb Q}\{\log n:n\in\mathbb Z,\ n\ge2\}.
\]
A \emph{rational Shannon payoff} on a product of rational polytopes is
a finite rational signed linear combination of functions $H(L(x,y))$,
where each $L$ is a rational affine map taking that product into a
finite probability simplex. Constant probability maps are permitted.
This definition does not permit arbitrary real coefficients or an
arbitrary real additive constant. A mixed strategy is any Borel
probability measure on the player's pure-strategy polytope; the players
randomize independently. The first player maximizes and the second
minimizes the expected payoff.

\begin{theorem}[An entropy embedding of finite logarithmic games]
\label{p05:ms5:thm:embedding}
Given a finite $r\times s$ real matrix $M$ with entries in
$\mathcal L_{\rm rat}$, one can effectively construct a rational Shannon
payoff $F$ on $\Delta_r\times\Delta_s$ such that:
\begin{enumerate}[label=(\roman*)]
\item $F(e_i,e_j)=M_{ij}$ at every pair of vertices;
\item $F(\cdot,y)$ is strictly convex for each $y$, and
$F(x,\cdot)$ is strictly concave for each $x$, whenever the corresponding
simplex has positive dimension;
\item the mixed equilibria of $F$ are exactly the mixed equilibria of
$M$, interpreted as measures supported on simplex vertices.
\end{enumerate}
In particular, the full continuous-strategy game has the same mixed
value as $M$. Statement (iii) holds for arbitrary Borel mixed strategies,
not just for finitely supported strategies chosen in advance.
\end{theorem}

\begin{proof}
The cases of a one-point simplex follow by deleting that variable, so
suppose $r,s\ge2$. Put $h(t)=H(t,1-t)$ and, for $n\ge2$, set
\[
 P_n(u,v)=h((u+v)/n),\qquad 0\le u,v\le1.
\]
For a function $K$ on this square write
$\delta K=K(1,1)-K(1,0)-K(0,1)+K(0,0)$. Direct calculation gives
$\delta P_2=-2\log2$ and, for $n\ge3$,
\[
 \delta P_n=-\log n-\frac2n\log2
 -\frac{n-2}{n}\log(n-2)+\frac{2(n-1)}n\log(n-1).
\]
Here the term involving $\log1$ is zero. Define $K_1=0$, $K_2=-P_2/2$,
and recursively
\[
 K_n=-P_n-\frac2n K_2-\frac{n-2}{n}K_{n-2}
                  +\frac{2(n-1)}nK_{n-1}\quad(n\ge3).
\]
Induction gives $\delta K_n=\log n$. Each $K_n$ is a rational linear
combination of $P_2,\ldots,P_n$, is symmetric, and vanishes at $(0,0)$.
Let $k_n=K_n(1,0)=K_n(0,1)$, and set
\[
 C_n(u,v)=K_n(u,v)-(u+v)k_n.
\]
Its four corner values are zero except $C_n(1,1)=\log n$.

The linear correction is itself an allowed rational Shannon expression.
Indeed, $k_n$ belongs to $\mathcal L_{\rm rat}$ because the entropy of a
rational probability vector is a rational combination of logarithms of
positive integers. For $z\in[0,1]$ and an integer $m\ge2$, the identity
\begin{equation}\label{p05:ms5:eq:linearlog}
 z\log m=H(1-z,\underbrace{z/m,\ldots,z/m}_{m\ {
m entries}})-h(z)
\end{equation}
represents every required linear logarithmic term. In particular no
coefficient $k_n$ is inserted as a new, nonrational entropy coefficient.

Choose a finite expression $M_{ij}=\sum_n q_{ijn}\log n$ with
$q_{ijn}\in\mathbb Q$, and put
\[
 F_0(x,y)=\sum_{i,j,n}q_{ijn}C_n(x_i,y_j).
\]
This is an allowed payoff and agrees with $M$ at the vertices. We next
make its separate curvature have the desired signs without changing
those vertex values.

Write $K_n=\sum_{m=2}^n a_{nm}P_m$, with $a_{nm}\in\mathbb Q$, and set
\[
 B_n=\sum_m|a_{nm}|,\qquad
 A=\sum_{i,j,n}|q_{ijn}|B_n.
\]
For $0<u<1$ and $0\le v\le1$,
\begin{equation}\label{p05:ms5:eq:curvature}
 |\partial_u^2P_m(u,v)|
   =\frac1{(u+v)(m-u-v)}\le\frac1{u(1-u)}.
\end{equation}
To check the inequality, minimize the concave quadratic $z(m-z)$ on
$z\in[u,u+1]$: both endpoint products are at least $u(1-u)$ for $m\ge2$.
The same bound holds after interchanging $u$ and $v$.
For $x$ in the relative interior of $\Delta_r$ and a tangent vector
$a$ with $\sum_i a_i=0$, it follows that
\[
 |D_x^2F_0(x,y)[a,a]|
 \le A\sum_i\frac{a_i^2}{x_i(1-x_i)}
 \le rA\sum_i\frac{a_i^2}{x_i}.
\]
For the last inequality, use
$1/(x_i(1-x_i))=1/x_i+1/(1-x_i)$ and Cauchy--Schwarz in the form
\[
 \frac{a_i^2}{1-x_i}
  =\frac{(\sum_{k\ne i}a_k)^2}{\sum_{k\ne i}x_k}
  \le\sum_{k\ne i}\frac{a_k^2}{x_k}.
\]
Choose the rational number $K=1+\max(r,s)A$ and define
\begin{equation}\label{p05:ms5:eq:embeddedpayoff}
 F(x,y)=F_0(x,y)-K H(x)+K H(y).
\end{equation}
Since $D^2H(x)[a,a]=-\sum_i a_i^2/x_i$, the claimed strict convexity
and strict concavity follow in relative interiors. They also hold on
every nonconstant segment in the closed simplices: restrict to the
smallest face containing the segment, apply the same estimates to its
active coordinates on the open segment, and then use continuity at
the endpoints. The face size only decreases the curvature bound.
No derivative at a zero coordinate is used. The added entropies vanish
at vertices, proving (i) and (ii).

For completeness, the reduction of arbitrary mixed strategies follows
from the two pointwise inequalities
\begin{equation}\label{p05:ms5:eq:vertexreplacement}
 F(x,y)\le\sum_i x_iF(e_i,y),\qquad
 F(x,y)\ge\sum_j y_jF(x,e_j).
\end{equation}
The first is strict for a nonvertex $x$, and the second is strict for a
nonvertex $y$. Replacing a measure $\mu$ on $\Delta_r$ by the vertex
measure with masses $\int x_i\,d\mu(x)$ weakly increases its payoff
against every opponent measure. Replacing the second player's measure
by the corresponding vertex measure weakly decreases the payoff
against every first-player measure. All integrations are justified by
bounded continuity on a compact product. The finite matrix minimax
theorem therefore supplies an equilibrium of the full game with the
same value.

If an equilibrium measure assigned positive mass to nonvertices, its
replacement would strictly improve that player's expected payoff:
the appropriate strict gap in \eqref{p05:ms5:eq:vertexreplacement} is a
positive measurable function at every nonvertex and against every
opponent pure strategy. Its integral is strictly positive whenever
the measure of those nonvertices is positive. Thus every equilibrium
is supported on vertices, and its masses must form a matrix-game
equilibrium. This proves (iii).
\end{proof}

\subsection{An unconditional failure of logarithmic algebraicity}

\begin{theorem}[Two rational entropy games beyond the pure-saddle range]
\label{p05:ms5:thm:mixedboundary}
There are two explicitly constructible rational Shannon games, each on
$\Delta_2\times\Delta_2$, with respective mixed values
\begin{equation}\label{p05:ms5:eq:twovalues}
 V_1=\frac{(\log2)^2}{\log7},\qquad
 V_2=\frac{\log2\log5}{\log7},
\end{equation}
such that each has a unique mixed equilibrium, all its nonzero
equilibrium masses are transcendental, and neither has a pure saddle.
At least one of $e^{V_1},e^{V_2}$ is transcendental.

More strongly, at least one of $V_1,V_2$ lies outside
\[
 \widetilde{\mathcal L}
  =\overline{\mathbb Q}+
    \operatorname{span}_{\overline{\mathbb Q}}
      \{z\in\mathbb C:e^z\in\overline{\mathbb Q}^{\times}\}.
\]
Consequently, for at least one of the two games, both its value and
the exponential of its value are transcendental. These assertions do
not identify which of the two games has this property.
\end{theorem}

\begin{proof}
Let $a=\log2$, $b=\log7$, and $d=\log5$. For $0<u,v<b$, consider
\[
 M(u,v;b)=\begin{pmatrix}u&0\\u+v-b&v\end{pmatrix}.
\]
The unique equilibrium assigns first-row probability $1-u/b$ and
first-column probability $v/b$, and its value is $uv/b$. Indeed,
the difference between the two column payoffs is $u-b+pb$ when
the first-row probability is $p$, while the difference between the
two row payoffs is $qb-v$ when the first-column probability is $q$.
The strict inequalities $0<u,v<b$ put both intersection probabilities
in $(0,1)$ and give uniqueness and the absence of a pure saddle.
Apply Theorem~\ref{p05:ms5:thm:embedding} to $M(a,a;b)$ and $M(a,d;b)$.
Their entries are rational combinations of logarithms of integers,
and the construction is fully explicit through $C_2,C_5,C_7$ and the
rational curvature bound in \eqref{p05:ms5:eq:embeddedpayoff}.

The ratios $a/b$ and $d/b$ are irrational by unique prime
factorization. If either ratio were algebraic irrational, the
Gelfond--Schneider theorem applied to the base $7$ would contradict
$7^{a/b}=2$ or $7^{d/b}=5$. Both ratios, and their complements, are
therefore transcendental. The embedding theorem transfers the unique
matrix equilibria, with exactly these masses, to the full games.

For the exponential assertion apply the classical six-exponentials
theorem to the rationally independent pair $(1,a/b)$ and triple
$(b,a,d)$; see \cite[Theorem~3]{p05:ms5:Waldschmidt}. The six products are
\[
 \begin{pmatrix}b&a&d\\a&a^2/b&ad/b\end{pmatrix}.
\]
All entries except the last two already exponentiate to algebraic
numbers. Thus at least one of $e^{V_1},e^{V_2}$ is transcendental.

For the stronger assertion use Roy's strong six-exponentials theorem,
in the formulation recorded in \cite[Theorem~3]{p05:ms5:Diaz} and credited
there to \cite[Corollary~2]{p05:ms5:Roy}: if two complex numbers are
linearly independent over $\overline{\mathbb Q}$ and three others
are likewise independent, one of their six products lies outside
$\widetilde{\mathcal L}$. Here $(1,a/b)$ is independent by the preceding
Gelfond--Schneider argument, and $(b,a,d)$ is independent by Baker's
theorem applied to their rational independence. Four products are
already in $\widetilde{\mathcal L}$, leaving the claimed alternative.
This space contains both all algebraic numbers and all logarithms of
nonzero algebraic numbers, which proves the last conclusion.
\end{proof}

\begin{remark}[What the counterexample does and does not settle]
\label{p05:ms5:rem:four}
The conclusion is existential but unconditional: it disproves the
universal assertion that every mixed rational Shannon game has
logarithmic algebraic value. It does not prove either individual
exponential in \eqref{p05:ms5:eq:twovalues} transcendental. The usual
four-exponentials conjecture would imply both individual assertions,
by applying it to $(1,a/b)$ and respectively $(b,a)$ or $(b,d)$.
These quotient-of-products questions belong to the classical problem
family explicitly discussed in \cite[Section~1.2]{p05:ms5:Diaz}.
The new structural step here is their realization, with exact
equilibrium control, by rational Shannon payoffs; the number-theoretic
alternative is a direct consequence of classical theory.
The word ``entropy game'' here refers to a one-shot Shannon payoff
game, not to the combinatorial growth-rate games studied in
\cite{p05:ms5:AGKS}; no complexity result for the latter is being claimed.
\end{remark}

\subsection{Baker splitting of algebraic tradeoffs}

The next result handles a different obstruction: an algebraic
irrational tradeoff weight cannot in general be absorbed by clearing
denominators in stationarity equations.

\begin{lemma}[Splitting and exact signs for algebraic logarithmic forms]
\label{p05:ms5:lem:split}
Let $\theta_1,\ldots,\theta_d$ be real algebraic numbers linearly
independent over $\mathbb Q$. If each $L_k$ is a rational linear
combination of real logarithms of positive algebraic numbers, then
\[
                 \sum_{k=1}^d\theta_kL_k=0
                 \quad\Longleftrightarrow\quad L_1=\cdots=L_d=0.
\]
Given finitely many real algebraic coefficients $\beta_i$ and positive
algebraic numbers $\alpha_i$, the sign of
$\sum_i\beta_i\log\alpha_i$ is decidable from exact algebraic-number
representations. No time bound is asserted.
\end{lemma}

\begin{proof}
Choose a rational basis $\lambda_1,\ldots,\lambda_m$ for the vector
space generated by the finitely many logarithms that occur. Baker's
linear-independence theorem states that these rationally independent
logarithms are also independent over $\overline{\mathbb Q}$
\cite{p03:mi3:Baker}. Write $L_k=\sum_j r_{kj}\lambda_j$ with rational
$r_{kj}$. A vanishing sum in the assertion implies
$\sum_k\theta_k r_{kj}=0$ for every $j$, hence all $r_{kj}=0$.
The converse is immediate.

For the algorithm one need not compute a basis of the logarithms.
Compute instead a rational basis $\theta_k$ of the span of the
algebraic coefficients $\beta_i$, using exact number-field linear
algebra, and write $\beta_i=\sum_k r_{ki}\theta_k$.
By the first assertion, the form is zero if and only if, for every $k$,
$\sum_i r_{ki}\log\alpha_i=0$. Clearing rational denominators makes
each of these conditions an exact identity
$\prod_i\alpha_i^{n_{ki}}=1$ with integer exponents. Such identities
are decidable by algebraic-number arithmetic. If the zero test fails,
successively refined certified intervals for the algebraic numbers
and their real logarithms eventually exclude zero and determine the
sign. Termination follows from nonvanishing, not from an assumed
separation bound.
\end{proof}

Fix a nonempty compact rational polytope $P$. Call a function on $P$ a
\emph{smooth rational logarithmic entropy expression} if it has the form
\begin{equation}\label{p05:ms5:eq:class}
 G(z)=\sum_a c_aH(L_a(z))+\sum_b\ell_b(z)\log r_b,
\end{equation}
where $c_a\in\mathbb Q$, $r_b\in\mathbb Q_{>0}$, $\ell_b$ are rational
affine functions, and $L_a$ are rational affine probability maps all
of whose coordinates are strictly positive throughout $P$.
Identically zero coordinates may be deleted before applying this
definition. In particular these functions are smooth on a neighborhood
of $P$ in its affine hull. Their derivatives along rational vectors
are rational linear combinations of logarithms of positive rational
affine functions and positive rational constants: the nonlogarithmic
terms cancel because each $L_a$ has total mass one.

\begin{theorem}[Common stationarity and decidable algebraic optimizers]
\label{p05:ms5:thm:decision}
Let $G_1,\ldots,G_m$ be expressions of the form
\eqref{p05:ms5:eq:class}, let $\lambda_1,\ldots,\lambda_m$ be specified
real algebraic numbers, and suppose
$F=\sum_i\lambda_iG_i$ is concave on $P$.
Choose a rational basis $\theta_1,\ldots,\theta_d$ of the span of the
$\lambda_i$, write $\lambda_i=\sum_k r_{ki}\theta_k$, and put
$A_k=\sum_i r_{ki}G_i$.
\begin{enumerate}[label=(\roman*)]
\item At an algebraic point in the relative interior of any face $Q$
of $P$, stationarity of $F$ along $Q$ is equivalent to simultaneous
stationarity of all $A_k$ along $Q$.
\item There is an effective finite list of algebraic points of $P$
which contains a global maximizer whenever any algebraic global
maximizer exists.
\item It is decidable whether $F$ has an algebraic global maximizer.
If one exists the procedure returns one; otherwise it certifies that
every global maximizer has at least one transcendental coordinate.
\end{enumerate}
The hypotheses specify a promised concave input class; the procedure
does not purport to decide concavity of arbitrary entropy expressions.
It has no polynomial-time guarantee.
\end{theorem}

\begin{proof}
For (i), choose a rational tangent basis of the face. At an algebraic
point each derivative of $A_k$ along such a vector is a rational
linear combination of real logarithms of positive algebraic numbers.
Apply Lemma~\ref{p05:ms5:lem:split} separately to every tangent derivative.

For each face $Q$, let $S_Q$ be the subset of its relative interior
where all $A_k$ are stationary along $Q$. Each equation defining
$S_Q$ is of the form
\[
 \sum_j t_j\log u_j(z)=0,
 \qquad t_j\in\mathbb Q,\quad u_j(z)>0,
\]
where the $u_j$ are rational affine functions or positive rational
constants. Exponentiating, clearing rational exponents, and moving
negative integer powers to the other side gives a polynomial equation
over $\mathbb Q$. Positivity makes this conversion reversible.
Together with the relative-interior conditions for $Q$, these give an
effective semialgebraic description over $\mathbb Q$.

Semialgebraic decomposition algorithms supply finitely many connected
cells covering $S_Q$, each with an algebraic sample point
\cite{p03:mi3:BPR}. On each cell every $A_k$ is constant: connect two of
its points by a piecewise $C^1$ semialgebraic path inside the cell and
differentiate along that path; its tangent vectors lie in the tangent
space of $Q$, on which the derivatives vanish. Thus $F$ is constant
on that cell as well. Include all sample points from all faces in
the candidate list, including zero-dimensional faces, for which
stationarity is vacuous. If an algebraic maximizer $z$ exists, it lies
in the relative interior of its minimal face and is stationary there.
Part (i) puts it in $S_Q$. The sample point of its cell has the same
values of all $A_k$, hence the same maximal value of $F$. This proves
(ii).

For each candidate $z$, concavity and differentiability give the exact
global optimality criterion
\begin{equation}\label{p05:ms5:eq:KKTtest}
           dF_z(v-z)\le0\quad\hbox{for every vertex $v$ of $P$}.
\end{equation}
Necessity follows by differentiating a feasible segment from a
maximizer; sufficiency follows from the concave tangent-plane bound
and the fact that vertices generate $P$. Each left-hand side is an
algebraic linear combination of logarithms of positive algebraic
numbers. Although $v-z$ need not be rational, it is algebraic, which
is all that the exact-sign part of Lemma~\ref{p05:ms5:lem:split} needs.
The mass-one cancellation in \eqref{p05:ms5:eq:class} still applies in
these directions. Test the finitely many inequalities exactly for
each candidate. Part (ii) proves that rejecting every candidate is
equivalent to absence of an algebraic maximizer. A real maximizer
exists by compactness and continuity, proving (iii).
\end{proof}

\begin{corollary}[Weighted information objectives with exact input data]
\label{p05:ms5:cor:channels}
Let $W_1,\ldots,W_m$ be finite rational channels with a common input
alphabet and strictly positive transition probabilities. Let $P$ be
a nonempty compact rational polytope of input distributions, and let
$\lambda_i>0$ be specified real algebraic weights. It is decidable
whether
\[
         \max_{p\in P}\sum_i\lambda_i I(p;W_i)
\]
admits an input distribution all of whose coordinates are algebraic.
If it does, one such distribution can be returned exactly.
\end{corollary}

\begin{proof}
Write
\[
 I(p;W_i)=H(W_ip)-\sum_xp_xH(W_i(\cdot\mid x)).
\]
The first term has strictly positive rational affine probability
coordinates, and the second is an affine rational combination of
logarithms of positive rationals. Mutual information is concave in
the input distribution, so the positive weighted sum satisfies all
hypotheses of Theorem~\ref{p05:ms5:thm:decision}.
\end{proof}

\begin{example}[A positive two-channel family with no algebraic optimizer]
\label{p05:ms5:ex:channels}
Let $W$ be the binary-input, binary-output channel with
$W(1\mid0)=1/4$ and $W(1\mid1)=1/2$, and let $W^{\rm rev}$ be obtained
by interchanging the two input symbols. If $p$ is the probability of
input $1$, put
\[
 J(p)=I(p;W)=h(1/4+p/4)-(1-p)h(1/4)-p h(1/2).
\]
For every positive algebraic $\beta$, the strictly concave objective
$J(p)+\beta J(1-p)$ has a unique maximizer $p_\beta\in(0,1)$.
Moreover,
\[
          p_\beta\in\overline{\mathbb Q}
          \quad\Longleftrightarrow\quad\beta\in\mathbb Q.
\]
To prove this, $J''(p)<0$, $J(0)=J(1)=0$, and $J(p)>0$ on $(0,1)$,
so the weighted maximum is interior and unique. Let $p_0$ maximize
$J$. The derivative identity
\[
                   J'(1/2)=\frac14\log(80/81)<0
\]
shows $p_0<1/2$, while the maximizer of $J(1-p)$ is $1-p_0>1/2$.
For algebraic irrational $\beta$, an algebraic stationary point of
$J(p)+\beta J(1-p)$ would, by Baker splitting, be stationary for both
summands separately, which is impossible. For rational $\beta$, the
stationarity equation exponentiates after clearing denominators to
a polynomial equation over $\mathbb Q$. This polynomial is not
identically zero, because the strictly decreasing derivative has
only one zero. Hence the unique optimizer is algebraic. This example
uses ordinary smooth concave information objectives and strictly
positive channels; it does not rely on a singular boundary channel
or a nonconcave information-bottleneck problem.
\end{example}

\paragraph{Scope and comparison with prior work.}
The algebraic stationary splitting is an application of Baker's
classical linear-independence theorem, and the finite sampling step
uses standard real algebraic geometry. The decision statement concerns
exact rational and algebraic input descriptions with a concavity
promise. It does not contradict impossibility results for selecting
capacity-achieving distributions uniformly from computable-real
channel descriptions, a different input model studied in
\cite{p05:ms5:Boche}. Nor does it settle the arithmetic of a transcendental
optimizer's objective value. The mixed-game embedding and the
algebraic-optimizer test are the reusable structural contributions
established here; historical priority for these formulations has not
been certified by the targeted literature check.

\section{Algebraic closure inside the field of \texorpdfstring{$E$}{E}-values}
\label{p05:ar21:section}

Let $k=\overline{\mathbb Q}$, let $\mathcal E$ denote the ring of
values of $E$-functions at algebraic points, and put
$K=\operatorname{Frac}(\mathcal E)$.  This section studies finite
algebraic extensions inside $K$, rather than another list of individual
rational powers.  Its two targets are an exact algebraic-closure theorem
for exponential values and a joint finite-cover theorem for several
exponential-integral values.  Both depend on the injective arithmetic
morphism of Fischler--Rivoal \cite[Proposition 3.1]{p03:ac19:FRring}.
Their Corollary 3.2 already classifies the intersection of $\mathcal E$
with the field of rational expressions in algebraic exponentials; that
rational-expression statement is an antecedent and is not counted again.
The results below extend the question to arbitrary algebraic elements.
They are proved applications, with no independently established priority
claim and no claim to settle the full unit or factorization questions
posed in \cite{p03:ac19:FRring}.

We use the external morphism and its extension to fraction fields:
\begin{equation}\label{p05:ar21:Phi}
 \Phi:\mathcal E\hookrightarrow\mathcal R,
 \qquad K\hookrightarrow\operatorname{Frac}(\mathcal R),
 \qquad
 \mathcal R=\left\{\sum_{\rho\in k}c_\rho e^{\rho z}:
                   c_\rho\in\mathbb C,\ \text{finite support}\right\}.
\end{equation}
It fixes $k$, and $\Phi(e^\alpha)=e^{\alpha z}$ for $\alpha\in k$.
For the exponential integral, the preceding morphism computation gives
\begin{equation}\label{p05:ar21:EinPhi}
 u_a=\operatorname{Ein}(a),\qquad
 \Phi(u_a)=A_a+B_ae^{-az},\qquad
 A_a=\gamma+\log a,\quad
 B_a=\delta_a=\int_0^\infty\frac{e^{-t}}{t+a}\,dt
\end{equation}
for positive algebraic $a\ge1$.  The constants $A_a,B_a$ are positive.
The choices of real logarithm in this formula are fixed throughout.

\subsection{Two descent lemmas}

\begin{lemma}[Finite phase descent]\label{p05:ar21:phase-descent}
Let $\alpha_1,\ldots,\alpha_m\in k$ be linearly independent over
$\mathbb Q$.  If $W\in\operatorname{Frac}(\mathcal R)$ is algebraic
over $\mathbb C(e^{\alpha_1z},\ldots,e^{\alpha_mz})$, then, for some
positive integer $q$,
\[
 W\in\mathbb C(e^{\alpha_1z/q},\ldots,e^{\alpha_mz/q}).
\]
If $W\in\mathcal R$, it belongs to the Laurent polynomial ring in
these same functions, after increasing $q$ if necessary.
\end{lemma}
\begin{proof}
Extend the given $\alpha_i$ to a rational basis of the span of these
numbers and the finitely many phases in a numerator and denominator of
$W$.  Clear all rational coordinates with one positive integer $q$.
Writing $X_i=e^{\alpha_i z/q}$ and $Y_j=e^{\beta_jz/q}$ for the extra
basis vectors gives $W\in\mathbb C(X_1,\ldots,X_m,Y_1,\ldots,Y_s)$.
These variables are algebraically independent over $\mathbb C$:
a polynomial relation would be a linear relation among finitely many
exponentials with distinct exponents, and successive derivatives at
zero give an invertible Vandermonde matrix.

Put $F=\mathbb C(X_1,\ldots,X_m)$.  An element of $F(Y_1,\ldots,Y_s)$
algebraic over $F$ belongs to $F$.  Indeed, apply each derivation
$\partial/\partial Y_j$ to its separable minimal polynomial over $F$.
The derivative of the polynomial at the element is nonzero, so every
such derivative of the element vanishes.  A rational function in
characteristic zero with zero derivative in a variable is independent
of that variable.  This proves the field assertion.  In the ring case
$W$ began as a Laurent polynomial in $X,Y$.  Independence of the $Y_j$
then removes all its nonzero $Y$-monomials, leaving a Laurent polynomial
in the $X_i$.
\end{proof}

\begin{lemma}[Descent of constant coefficients]\label{p05:ar21:constants}
For independent indeterminates $T_1,\ldots,T_m$, the field
$k(T_1,\ldots,T_m)$ is relatively algebraically closed in
$\mathbb C(T_1,\ldots,T_m)$.  Further, a finite function-field
extension over $k(T_1,\ldots,T_m)$ retains its degree, normality,
and Galois group, when applicable, after extension of constants to
$\mathbb C$.
\end{lemma}
\begin{proof}
For the first assertion let $h\in\mathbb C(T)$ satisfy a nonzero
polynomial over $k(T)$.  Every automorphism of $\mathbb C/k$, acting
on coefficients and fixing $T$, sends $h$ to a root of the same
polynomial, so $h$ has a finite orbit.  Express $h=A/B$ with relatively
prime polynomials and normalize one specified nonzero coefficient of
$B$ to one.  This representation is unique, hence every coefficient
of $A$ and $B$ has a finite orbit under these automorphisms.
A complex number transcendental over $k$ has an infinite orbit:
include it in a transcendence basis, send it to itself plus an arbitrary
integer while fixing the other basis elements, and extend these
isomorphisms to the algebraic closure $\mathbb C$.  All the coefficients
therefore lie in $k$.

For the second assertion, take a primitive element of a finite
extension and clear denominators in its irreducible polynomial.
A polynomial irreducible over an algebraically closed coefficient
field stays irreducible after extension of that field.  In this
setting one can see this directly: a putative factorization has
finitely many normalized coefficients satisfying polynomial equations
over $k$; a solution over $\mathbb C$ gives a solution over $k$ by the
weak Nullstellensatz.  Nonzero normalizing coefficients are handled
by adjoining their inverses.  Gauss's lemma transfers the assertion
between polynomial and rational-function coefficients.  Thus the
primitive-element polynomial has the same degree after the constant
extension.  Apply the same argument to the normal closure: its
original automorphisms extend and already account for the preserved
degree.  Its Galois group, and the subgroups corresponding to the
original intermediate fields, are unchanged.  Normality and the stated
Galois-group assertion follow.
\end{proof}

\subsection{Exact algebraic and integral closure for exponential values}

Fix rationally independent $\alpha_1,\ldots,\alpha_m\in k$ and write
\[
 F_\alpha=k(e^{\alpha_1},\ldots,e^{\alpha_m}),\qquad
 R_{\alpha,q}=k[e^{\alpha_1/q},e^{-\alpha_1/q},\ldots,
                  e^{\alpha_m/q},e^{-\alpha_m/q}],\qquad
 F_{\alpha,q}=\operatorname{Frac}(R_{\alpha,q}).
\]
The exponential values are algebraically independent over $k$; this
follows either from Lindemann--Weierstrass or by applying $\Phi$ and
the exponential Vandermonde argument above.  Thus these really are
Laurent polynomial rings and rational function fields in $m$ variables.

\begin{theorem}[Exact exponential algebraic closure]\label{p05:ar21:expclosure}
One has
\begin{align}
 \{v\in K:v\text{ is algebraic over }F_\alpha\}
       &=\bigcup_{q\ge1}F_{\alpha,q},\label{p05:ar21:fieldclosure}\\
 \{v\in\mathcal E:v\text{ is algebraic over }F_\alpha\}
       &=\bigcup_{q\ge1}R_{\alpha,q}.\label{p05:ar21:ringclosure}
\end{align}
Moreover, $\bigcup_{q\ge1}R_{\alpha,q}$ is exactly the integral
closure of $R_{\alpha,1}$ inside $K$.
The field on the right of \eqref{p05:ar21:fieldclosure} is a Galois
extension of $F_\alpha$ with Galois group
$\widehat{\mathbb Z}^{\,m}$.
\end{theorem}
\begin{proof}
Let $v\in K$ be algebraic over $F_\alpha$.  Its image $W=\Phi(v)$
is algebraic over $k(e^{\alpha_1z},\ldots,e^{\alpha_mz})$.
Lemma~\ref{p05:ar21:phase-descent} gives variables $X_i=e^{\alpha_i z/q}$
for some $q$, with $W\in\mathbb C(X_1,\ldots,X_m)$.  Since $W$
is algebraic over $k(X_1,\ldots,X_m)$, the constant-descent lemma gives
\[ W\in k(X_1,\ldots,X_m). \]  Write $W=S(X)$ with $S$ rational over $k$.
The element $S(e^{\alpha_1/q},\ldots,e^{\alpha_m/q})$ belongs to $K$
and has image $W$.  Its denominator is nonzero by algebraic
independence.  Injectivity of $\Phi$ therefore identifies it with $v$.
This proves one inclusion in \eqref{p05:ar21:fieldclosure}; the reverse
inclusion follows from $(e^{\alpha_i/q})^q=e^{\alpha_i}$.

If $v\in\mathcal E$, the ring assertion of the phase-descent lemma
puts $W$ in $\mathbb C[X_1^{\pm1},\ldots,X_m^{\pm1}]$.
Constant descent and uniqueness of Laurent coefficients now give
\[ W\in k[X_1^{\pm1},\ldots,X_m^{\pm1}]. \]
Injectivity gives $v\in R_{\alpha,q}$.  Conversely, all generators of
these Laurent polynomial rings are $E$-values, proving
\eqref{p05:ar21:ringclosure}.

Every $R_{\alpha,q}$ is integral over $R_{\alpha,1}$, including its
negative-power generators, because each satisfies a monic power
polynomial with coefficient in $R_{\alpha,1}$.  Conversely, if $v\in K$
is integral over $R_{\alpha,1}$, the field assertion puts $v$ in
$F_{\alpha,q}$ for some $q$.  It is then also integral over
$R_{\alpha,q}$.  This Laurent polynomial ring is a unique factorization
domain, hence integrally closed in its fraction field; therefore
$v\in R_{\alpha,q}$.

For each $q$, the extension $F_{\alpha,q}/F_\alpha$ has group
$(\mathbb Z/q\mathbb Z)^m$, acting independently by multiplication
of the $q$th-root generators by $q$th roots of unity.  Independence
of the generators proves that its degree is $q^m$.  The directed union
uses divisibility of $q$, and the restriction maps are the usual
reduction maps.  Its inverse-limit Galois group is consequently
$\widehat{\mathbb Z}^{\,m}$.
\end{proof}

\begin{corollary}[Exponential-field dichotomy]\label{p05:ar21:dichotomy}
Put
\[
 F_{\exp}=k(e^\alpha:\alpha\in k),\qquad
 R_{\exp}=\left\{\sum_{j=1}^n c_je^{\alpha_j}:
                              c_j,\alpha_j\in k\right\}.
\]
Then $F_{\exp}$ is relatively algebraically closed in $K$, and
$R_{\exp}$ is integrally closed inside $K$.  In particular, an
$E$-value is either a finite algebraic-coefficient exponential sum
or is transcendental over $F_{\exp}$.
Any counterexample to the unit conjecture
$\mathcal E^\times=k^\times\exp(k)$ in \cite{p03:ac19:FRring}
would have to be transcendental over $F_{\exp}$.
\end{corollary}
\begin{proof}
A polynomial relation over $F_{\exp}$ uses only finitely many
exponential phases in its coefficients.  Choose a rational basis of
their span and clear denominators.  Theorem~\ref{p05:ar21:expclosure}
then puts every element of $K$ algebraic over $F_{\exp}$ back in
$F_{\exp}$; for elements of $\mathcal E$ it puts them in $R_{\exp}$.
The same finite-support reduction applied to a monic relation proves
the relative integral-closure assertion.

If a unit $v\in\mathcal E^\times$ were algebraic over $F_{\exp}$,
both $v$ and $v^{-1}$ would belong to $R_{\exp}$.  Choose one finite
rational basis and denominator for all their phases.  Their images
are mutually inverse Laurent polynomials.  A unit of a Laurent
polynomial ring is a nonzero constant times a Laurent monomial:
for example, the Newton polytope of a product is the Minkowski sum
of those of its factors, so a product equal to one forces both
supports to consist of one point.  Hence $v=ce^\alpha$ with
$c\in k^\times$, $\alpha\in k$.
\end{proof}

An earlier restricted result of Rivoal--Roques
\cite[Theorem~5]{p05:ar21:RivoalRoques} treats reciprocal values of two
$E$-functions that form a common two-dimensional differential system.
Our restriction instead concerns algebraicity over $F_{\exp}$ and
does not assume such a common system.  These are distinct hypotheses.

The exact closure theorem is not a proof of the full unit conjecture.
It excludes algebraic counterexamples over $F_{\exp}$ and leaves open
possible units transcendental over that field.  Similarly, a factor
of an exponential sum need not be algebraic over the exponential
field merely because its product with a second factor lies there.
That missing implication cannot be inserted into a factorization
argument.

\subsection{Joint exponential-integral fields have at most one ramified coordinate}

Let $a_1,\ldots,a_m\ge1$ be positive algebraic numbers linearly
independent over $\mathbb Q$.  Write
\[
 u_i=u_{a_i},\qquad A_i=\gamma+\log a_i,\qquad B_i=\delta_{a_i},
 \qquad F=k(u_1,\ldots,u_m).
\]
The $u_i$ are algebraically independent: their images are the affine
coordinates $A_i+B_iX_i$ in the algebraically independent exponential
variables $X_i=e^{-a_i z}$.  This simple argument needs the stated
rational independence of the parameters; no claim is made here for
arbitrary rationally dependent frequency families.

\begin{theorem}[A joint cyclic-coordinate obstruction]\label{p05:ar21:joint}
Let $v\in K$ be algebraic over $F$, and set $d=[F(v):F]$.
If $d>1$, there exists an index $i$ such that $A_i\in k$, and
\begin{equation}\label{p05:ar21:cycliccoordinate}
             F(v)=F\bigl((u_i-A_i)^{1/d}\bigr).
\end{equation}
In particular this extension is cyclic.  Among the entire parameter
family there is at most one possible such index, independent of $v$.
If none of the $A_i$ is algebraic, then $F$ is relatively algebraically
closed in $K$.  Unconditionally, one can remove at most one of the
parameters so that the field jointly generated by all the remaining
$u_i$ is relatively algebraically closed in $K$.
\end{theorem}
\begin{proof}
Apply $\Phi$ to a minimal polynomial for $v$ over $F$.  Set
$U_i=A_i+B_ie^{-a_i z}$ and $W=\Phi(v)$.  The phase-descent lemma,
with phases $-a_i$, gives an integer $q\ge1$ and
\[
 W\in\mathbb C(X_1,\ldots,X_m),\qquad
 X_i=e^{-a_i z/q},\qquad U_i=A_i+B_iX_i^q.
\]
The degree of the extension is still $d$ over $\mathbb C(U_1,\ldots,U_m)$,
by Lemma~\ref{p05:ar21:constants}.  The extension
\[
 \mathbb C(X_1,\ldots,X_m)/\mathbb C(U_1,\ldots,U_m)
\]
is Galois with group $G=(\mathbb Z/q\mathbb Z)^m$.  Let $H$ be the
subgroup whose fixed field is $\mathbb C(U,W)$.  As $G$ is abelian,
this intermediate extension is Galois with group $G/H$.

For each $i$, let $G_i$ be the coordinate subgroup acting on $X_i$
alone.  At the generic point of the hyperplane
$D_i=\{U_i=A_i\}$ its inertia group in the full extension is exactly
$G_i$.  To check this explicitly, $U_i-A_i=B_iX_i^q$ has order $q$
along $X_i=0$, while the other $X_j$ are nonzero units there.
The transformations of $X_j$ with $j\ne i$ act nontrivially on the
residue field, whereas those of $X_i$ act trivially on it.
Thus the inertia in the intermediate extension is the image of $G_i$
in $G/H$.

If this image is nontrivial, $D_i$ is a branch divisor of the cover
defined over $k$ by the minimal polynomial of $v$.  More concretely,
clear denominators to obtain a polynomial in $k[U_1,\ldots,U_m,V]$.
Its leading coefficient in $V$ times its nonzero $V$-discriminant
vanishes on every branch divisor.  This follows from the implicit
function theorem off the zero set of that product, where its roots
are distinct local holomorphic functions.  Consequently a nonzero
polynomial $D\in k[U_1,\ldots,U_m]$ vanishes identically when
$U_i=A_i$.  Expand $D$ in the other variables.  Any nonzero coefficient
polynomial in $U_i$ then has $A_i$ as a root, proving $A_i\in k$.

At most one $A_i$ is algebraic.  Otherwise, for distinct $i,j$,
$A_i-A_j=\log(a_i/a_j)$ would be a nonzero algebraic number with
algebraic exponential, contrary to Hermite--Lindemann
\cite{p04:ac20:Lindemann}.  If none is algebraic, every $G_i$ lies in
$H$; since the coordinate subgroups generate $G$, this gives $H=G$
and $d=1$.
If an algebraic center exists, denote its index by $i$.  All other
$G_j$ lie in $H$.  The quotient $G/H$ is therefore a cyclic quotient
of $G_i$, of order $d$, and
\begin{equation}\label{p05:ar21:geometricfield}
 \mathbb C(U,W)=
 \mathbb C\bigl(U_1,\ldots,U_m,(U_i-A_i)^{1/d}\bigr).
\end{equation}
Indeed the fixed field of the subgroup of order $q/d$ of the
$i$th coordinate group is generated by $X_i^{q/d}$; its nonzero
constant factor $B_i^{1/d}$ is available over $\mathbb C$.

It remains to descend the exact field description, not just its
geometric degree.  Lemma~\ref{p05:ar21:constants} shows that $F(v)/F$
itself is cyclic of degree $d$.  As $k$ contains the $d$th roots of
unity, elementary Kummer theory gives an element $y$ with
$F(v)=F(y)$ and $y^d=f(u_1,\ldots,u_m)$ for some $f\in k(U)^\times$.
For completeness, if $\sigma$ generates the cyclic group, the
operator $\sum_{j=0}^{d-1}\zeta^{-j}\sigma^j$ is nonzero by linear
independence of field automorphisms.  Its nonzero value $y$ satisfies
$\sigma(y)=\zeta y$ for a primitive $d$th root $\zeta$; therefore
$y^d\in F$ and $y$ has degree $d$.

By \eqref{p05:ar21:geometricfield}, over $\mathbb C(U)$ this eigenvector
is a rational multiple of $(U_i-A_i)^{b/d}$ for some $b$ coprime to
$d$.  Hence the order of $f$ at every irreducible polynomial divisor
other than $U_i-A_i$ is divisible by $d$, and its order $n$ at
$U_i-A_i$ is coprime to $d$.  Unique factorization in $k[U_1,\ldots,U_m]$
therefore gives
\[
           f(U)=c\,(U_i-A_i)^n R(U)^d,
             \qquad c\in k^\times,\quad R\in k(U)^\times.
\]
Divide $y$ by $c^{1/d}R(u)$.  Bezout's identity for $n,d$ then
produces an element $t\in F(v)$ with $t^d=u_i-A_i$.
This element has degree $d$, by the valuation at $U_i-A_i$ (or the
irreducibility of $T^d-(U_i-A_i)$).  Consequently $F(v)=F(t)$,
which proves \eqref{p05:ar21:cycliccoordinate}.
All choices of the root generate the same field, since their ratios
are roots of unity in $k$.

Finally remove the unique index with algebraic center if it exists.
For the remaining family the already proved no-center case establishes
joint relative algebraic closure.  If there was no such index,
no removal is necessary.
\end{proof}

This conclusion is stronger than separate relative-algebraic-closure
statements for each single $k(u_i)$.  Separate statements do not rule
out a new algebraic extension depending simultaneously on several
coordinates.  The joint theorem rules out every noncyclic extension
and every cyclic extension ramified along a genuinely mixed divisor.
For example, for rationally independent $a,b\ge1$, no root of
\begin{equation}\label{p05:ar21:cubic}
                       T^3-u_aT-u_b=0
\end{equation}
belongs to $K$.  This is a consequence of the joint theorem, not an
additional structural target.

\begin{proof}[Verification of the cubic example]
Over $k(U,V)$ the polynomial $T^3-UT-V$ is irreducible.
Indeed it defines the degree-three rational-function extension
$k(U,T)/k(U,T^3-UT)$: over $k(U)$ the polynomial map $T\mapsto T^3-UT$
has degree three, so the function-field degree is three.
Its discriminant is $4U^3-27V^2$, which is not a square in $k(U,V)$;
it is an irreducible polynomial, since $U^3$ is not a square in
$k(U)$ and the polynomial is quadratic in $V$.
An irreducible cubic has transitive Galois group $A_3$ or $S_3$,
and the discriminant is a square precisely in the first case.
Thus its group is $S_3$ and its degree-three root extension is not
Galois.  Algebraic independence of $u_a,u_b$ identifies their field
with $k(U,V)$, preserving these assertions.  Theorem~\ref{p05:ar21:joint}
would make any such root extension in $K$ cyclic, a contradiction.
\end{proof}

\subsection{Limits of the structural transfer}

The first theorem gives an exact unconditional closure computation and
a corresponding restriction on the unit conjecture.  The second gives
an unconditional description of every possible joint finite extension,
with a single exceptional connection coefficient left explicitly
unclassified.  It does not assert that any exceptional extension exists.
The full ring of $E$-values may contain units transcendental over the
exponential field; the present argument neither constructs nor rules
out such units.  It also does not identify all factorizations of a
nonexponential $E$-value.  These are concrete remaining gaps in the
open ring questions, not consequences suppressed by the notation.

\section{A convexity mechanism for gamma extraction and an atomic corner in mixtures}
\label{p05:be21:section}

The preceding continuation left open whether different interior contacts
can exchange control of a gamma-extraction boundary.  We settle the
question throughout the order interval $1/2<\kappa\leq1$: every
candidate inverse, including an inadmissible one, has a unique global
minimum.  Consequently its probability boundary has exactly one contact
and is analytic throughout the intermediate-rate interval.  The argument
is a signed Cauchy-transform calculation, followed by a second-order
strict differential inequality.

A second investigation concerns positive mixtures of two different
orders.  The same convexity survives, but the geometry changes: the
maximal extractable shape can remain constant on an initial interval,
with an atom in the remainder, and then leave that interval through a
genuine corner.  We give a complete implicit phase diagram and prove
that this phenomenon occurs.  These are extensions of the manuscript's
factorization problems, not claimed resolutions of a named conjecture
in the external Bessel-distribution literature \cite{p05:be21:Baricz}.

\subsection{Notation and the inherited boundary facts}

Fix $b>1$.  For $\kappa>1/2$, let $X_\kappa$ have transform
\[
 L_\kappa(s)=\frac{b+s}{b}
       \left(\frac{b^2-1}{(b+s)^2-1}\right)^\kappa.
\]
Write $M=2\kappa-1$.  Extraction of a gamma variable with shape
$h$ and rate $b-c$, $0<c<1$, leads, after a positive normalization
and the tilt $e^{-bx}$, to the inverse
\begin{align}
 F_{\kappa,c,h}(w)&=
       \frac{w(w-c)^h}{(w-1)^\kappa(w+1)^\kappa},
                                                 \label{p05:be21:F}\\
 g_{\kappa,c,h}(x)&=
       \sum_{n\geq0}\frac{a_n(c,h)x^{\eta+n-1}}
                            {\Gamma(\eta+n)},
 &\eta&=M-h,                                     \label{p05:be21:g}\\
 (1-cz)^h(1-z^2)^{-\kappa}&=
                          \sum_{n\geq0}a_n(c,h)z^n.
                                                 \label{p05:be21:coeff}
\end{align}
For $0<h<M$ the series and its derivatives converge locally uniformly
on $x>0$.  They depend real analytically on all strict parameters.
The candidate is a probability remainder exactly when $g\geq0$.

We explicitly use the earlier capacity and contact theorems
\cite{p05:be21:Capacity,p05:be21:Fractional}.  The maximal shape $H_\kappa(c)$
is attained, belongs to $(0,M)$ for $0<c<1$, and is continuous and
strictly decreasing on $[0,1]$.  A boundary inverse has a nonempty
finite contact set; every contact is double and satisfies
$x<(h+1)/c$.  At a contact of a nonnegative inverse,
\begin{equation}\label{p05:be21:inherited-response}
 g_c=-h\int_0^x e^{cy}g(x-y)\,dy<0,\qquad
 g_h=-\int_0^x\frac{e^{cy}g(x-y)}{y}\,dy<0.
\end{equation}
The finite-contact boundary is locally the minimum of finitely many
analytic contact branches.  Positivity estimates at zero and infinity
are uniform in compact strict parameter neighborhoods.  These are
inputs, rather than results counted again below.

\subsection{A signed cut representation becomes strict convexity}

\begin{theorem}[Strict convexity and a unique minimum]
\label{p05:be21:convex}
Suppose $1/2<\kappa\leq1$, $0<c<1$, and $0<h<2\kappa-1$.
Then
\begin{equation}\label{p05:be21:inequality}
             g''_{\kappa,c,h}(x)+g'_{\kappa,c,h}(x)>0
                       \qquad(x>0).
\end{equation}
The function $r\mapsto g_{\kappa,c,h}(-\log r)$ is strictly
convex on $(0,1)$ and tends to $+\infty$ at both ends.
It has exactly one minimum, which is nondegenerate.  In particular,
the inverse is nonnegative on the whole positive axis if and only
if its value at this single minimum is nonnegative.
\end{theorem}
\begin{proof}
First assume $\kappa<1$.  Then
\[
       0<h<M<\kappa<1,\qquad 0<\eta=M-h<1.
\]
Take principal powers in \eqref{p05:be21:F}, with cut $(-\infty,1]$,
and set $W(s)=-\pi^{-1}\operatorname{Im}F(s+i0)$.  Direct boundary
values give
\begin{equation}\label{p05:be21:cut}
 W(s)=\begin{cases}
 \displaystyle\frac{-s\sin(\pi\eta)}{\pi}
 \frac{(c-s)^h}{(1-s)^\kappa(-1-s)^\kappa},&s<-1,\\[6pt]
 \displaystyle\frac{-s\sin(\pi(h-\kappa))}{\pi}
 \frac{(c-s)^h}{(1-s)^\kappa(1+s)^\kappa},&-1<s<c,\\[6pt]
 \displaystyle\frac{s\sin(\pi\kappa)}{\pi}
 \frac{(s-c)^h}{(1-s)^\kappa(1+s)^\kappa},&c<s<1.
 \end{cases}
\end{equation}
Thus $W$ is positive on $(-\infty,-1)$ and $(0,1)$, and negative
on $(-1,0)$.  The integrable singularities at $\pm1$ have exponent
$-\kappa>-1$; the point $c$ causes no singularity.  At negative
infinity $W(s)=O(|s|^{-\eta})$.

For completeness, this signed cut really inverts the transform.
Apply Cauchy's formula to $F(z)/(z-w)$ on a cut disk, with
$w>1$ fixed, and then let its radius tend to infinity.
The outer-circle contribution is $O(R^{-\eta})$, because
$F(z)=O(R^{-\eta})$ there.  The small circles around $\pm1$
contribute $O(\epsilon^{1-\kappa})$, and that around $c$ tends
to zero.  The two sides of the cut consequently give
\[
       F(w)=\int_{-\infty}^1\frac{W(s)}{w-s}\,ds.
\]
The integral is absolutely convergent.  Fubini's theorem, justified
by the same absolute integral, shows that
\begin{equation}\label{p05:be21:cut-inverse}
       g(x)=\int_{-\infty}^1e^{sx}W(s)\,ds\qquad(x>0).
\end{equation}
Indeed the right side has Laplace transform $F(w)$, and uniqueness
agrees with \eqref{p05:be21:g}.  Every number of $x$ derivatives can
be passed under this integral on a compact positive interval.

The quadratic $s(s+1)$ has exactly the same sign as $W(s)$,
apart from endpoints.  Therefore
\[
       g''(x)+g'(x)
       =\int_{-\infty}^1s(s+1)e^{sx}W(s)\,ds>0.
\]
For $r=e^{-x}$, the operator $D_x^2+D_x$ is $r^2D_r^2$,
which proves the asserted convexity.  The endpoint estimates are
\begin{equation}\label{p05:be21:ends}
 g(x)\sim\frac{x^{\eta-1}}{\Gamma(\eta)}\quad(x\downarrow0),
 \qquad
 g(x)\sim\frac{(1-c)^h}{2^\kappa\Gamma(\kappa)}
                    e^x x^{\kappa-1}\quad(x\longrightarrow\infty).
\end{equation}
The first follows from the series, and the second follows from
\eqref{p05:be21:cut-inverse} at its right endpoint (or from the inherited
uniform tail estimate).  Both limits are positive infinity.
Strict convexity therefore gives a unique interior minimum.
At that minimum $g'=0$, and \eqref{p05:be21:inequality} gives $g''>0$.

Finally fix $\kappa=1$, $0<h<1$, and let $\kappa_j\uparrow1$
with $h<2\kappa_j-1$.  The series gives local convergence of $g$
and its first two derivatives, so the limit in
\eqref{p05:be21:inequality} is nonnegative.  It is strictly positive:
choose any compact interval $I\subset(0,c)$.  On $I$ the positive
integrand in the preceding formula converges uniformly to a
strictly positive function, since
$-\sin\pi(h-\kappa_j)\to\sin\pi h>0$.
All the remaining integrals are nonnegative.  The integral over
$I$ supplies a fixed positive lower bound.  The endpoint estimates
\eqref{p05:be21:ends} still hold, with $\eta=1-h\in(0,1)$.
This proves all assertions also at $\kappa=1$ without an
unjustified integrable-cut formula at its pole endpoints.
\end{proof}

\begin{theorem}[No contact switches in the full low-order interval]
\label{p05:be21:unique}
For every $1/2<\kappa\leq1$ and $0<c<1$, the maximal probability
remainder has exactly one positive zero $x_\kappa(c)$.
This zero is double.  Both $H_\kappa(c)$ and $x_\kappa(c)$ are
real analytic throughout $(0,1)$.  Thus no exchange between
distinct controlling contacts occurs in this order interval.
The exact slope is
\[
 H_\kappa'(c)=-H_\kappa(c)
 \frac{\int_0^{x_\kappa(c)}e^{cy}g(x_\kappa(c)-y)\,dy}
 {\int_0^{x_\kappa(c)}e^{cy}g(x_\kappa(c)-y)\,dy/y}.
\]
\end{theorem}
\begin{proof}
The attained boundary has at least one zero.  By
Theorem~\ref{p05:be21:convex} it must be the unique minimum, and
its second derivative is positive.  At this contact $g_h<0$
by \eqref{p05:be21:inherited-response}.  The Jacobian of $(g,g_x)$
with respect to $(h,x)$ is consequently
$g_hg_{xx}\ne0$.  The analytic implicit-function theorem gives
analytic local functions $h(c),x(c)$.  Uniqueness identifies
these local functions on overlaps with the actual capacity
and its contact.  The first equation differentiated in $c$
gives the displayed slope formula.
\end{proof}

\subsection{Endpoint continuation and finitely many switches for every order}

Here the order is no longer restricted to $\kappa\leq1$.
The earlier endpoint theorem supplies the analytic scalar $A(\kappa)$
and the unique double contact
\[
 h_1=H_\kappa(1)=\kappa-A(\kappa),\qquad x_1=h_1,
 \qquad M-h_1>0
\]
\cite{p05:be21:Endpoint}.  The earlier central-rate bifurcation theorem
also proves, for every fixed $\kappa>1/2$, uniqueness of the boundary
contact for all sufficiently small positive $c$, with
$M-H_\kappa(c)$ and $x_\kappa(c)/c$ analytic in $c^2$
\cite{p05:be21:Capacity}.  In particular that theorem establishes
global uniqueness near zero, not merely a local formal expansion.

\begin{proposition}[Endpoint control and global finiteness]
\label{p05:be21:finite}
For every fixed $\kappa>1/2$, the boundary has a unique contact
for all $c<1$ sufficiently close to $1$.  The capacity and contact
extend as analytic contact curves through $c=1$, and
\[
                         -h_1^2<H_\kappa'(1-)<0.
\]
There are only finitely many points in $(0,1)$ at which the
capacity changes its controlling analytic contact branch.
This does not assert absence of switches when $\kappa>1$.
\end{proposition}
\begin{proof}
Let $c_j\uparrow1$ and let $x_j$ be any boundary contacts.
Continuity gives $H_\kappa(c_j)\to h_1$, while the inherited bound
$x_j<(H_\kappa(c_j)+1)/c_j$ bounds them above.  The local series
at zero, uniformly for $\eta$ near $M-h_1>0$, bounds them away
from zero.  Every subsequential limit must therefore be a zero
of the endpoint inverse.  There is only one, namely $x_1$.

At that point $g_{xx}>0$ and $g_h<0$.  The inverse series is
analytic in parameters near this positive point, also for $c$
slightly larger than $1$.  The response formula remains valid
at $c=1$, since its finite integrals require only $\eta>0$.
The implicit-function theorem gives a unique local contact curve.
All boundary contacts for $c$ close to $1$ lie in this neighborhood,
so they coincide with it.  In the slope formula the ratio of
the two integrals is a strict weighted mean of $y\in(0,x_1)$.
It lies in $(0,x_1)$, giving the stated derivative bounds.

There are no switches in a small interval at either endpoint,
by this argument and the inherited central-rate theorem.
The remaining compact interval is covered by finitely many
neighborhoods on which the capacity is a finite minimum of
analytic functions.  Differences of distinct analytic germs
have isolated zeros; a finite collection has only finitely
many crossings on a sufficiently small closed neighborhood.
A finite subcover proves the claim.  The analytic continuation
past $c=1$ is only a continuation of the contact equations;
it is not an assertion of probabilistic feasibility there.
\end{proof}

\subsection{Two orders: an atomic plateau and a genuine corner}

Choose
\begin{equation}\label{p05:be21:mixture-range}
 \frac12<\kappa_0<\kappa_1\leq1,\qquad
 M_i=2\kappa_i-1,\qquad \delta=M_1-M_0\in(0,1),\qquad
 \epsilon>0.
\end{equation}
Consider the probability law with transform
\begin{equation}\label{p05:be21:mixture-source}
 L(s)=\frac{\Phi(b+s)}{\Phi(b)},\qquad
 \Phi(w)=\frac{w}{(w^2-1)^{\kappa_0}}
             +\epsilon\frac{w}{(w^2-1)^{\kappa_1}}.
\end{equation}
It is a positive mixture of the two laws already considered.
Let $K(c)$ be the greatest gamma shape extractable at rate $b-c$.
For $0<h<M_0$ put
\[
 q_{c,h}(x)=g_{\kappa_0,c,h}(x)
                       +\epsilon g_{\kappa_1,c,h}(x).
\]
At the endpoint $h=M_0$, the first summand has an atom of
unnormalized mass one.  Its continuous part is
\begin{equation}\label{p05:be21:atomic-series}
 k_c(x)=\sum_{n\geq1}
       \frac{a_n(c,M_0)x^{n-1}}{\Gamma(n)},
 \qquad k_c(0+)=-M_0c,
\end{equation}
where the coefficients use $\kappa=\kappa_0$.  Define
\begin{equation}\label{p05:be21:qfull}
 Q_c(x)=k_c(x)+\epsilon g_{\kappa_1,c,M_0}(x).
\end{equation}
Thus the full-shape candidate, before normalization and the tilt,
is $\delta_0+Q_c(x)\,dx$.  More precisely the actual probability
candidate is this tilted measure divided by
$(b-c)^{M_0}\Phi(b)$.

\begin{theorem}[Complete implicit diagram for a two-order mixture]
\label{p05:be21:mixture}
Under \eqref{p05:be21:mixture-range}, the following assertions hold.
\begin{enumerate}
\item For $0<c<1$ and $0<h<M_0$, $q_{c,h}$ has a unique
nondegenerate global minimum and tends to $+\infty$ at both
ends.  The same assertions hold for $Q_c$.  Positivity in
either case is therefore an exact one-point minimum test.
\item There is a number $c_*\in(0,1]$ such that
$K(c)=M_0$ for $0\leq c\leq c_*$ within $[0,1)$,
and $0<K(c)<M_0$ for $c_*<c<1$.  For $0<c<c_*$
the full-shape continuous density $Q_c$ is strictly positive.
For $c_*<c<1$ the capacity has a unique double contact and
is real analytic and strictly decreasing.
\item If $c_*<1$, then $Q_{c_*}$ has exactly one zero $x_*$,
which is double.  A unique analytic contact branch $h(c)$
passes through $(c_*,M_0,x_*)$, and locally
\begin{equation}\label{p05:be21:corner}
                   K(c)=\min\{M_0,h(c)\}.
\end{equation}
Its right derivative is strictly negative whereas its left
derivative is zero.  Writing $Q=Q_{c_*}$, its exact slope is
\begin{equation}\label{p05:be21:atomic-slope}
 K'(c_*+)=-M_0
 \frac{e^{c_*x_*}+\int_0^{x_*}e^{c_*y}Q(x_*-y)\,dy}
 {e^{c_*x_*}/x_*+
                 \int_0^{x_*}e^{c_*y}Q(x_*-y)\,dy/y}<0.
\end{equation}
\end{enumerate}
Consequently the mixture has at most one corner, and it is a
transition from an atomic shape ceiling to a moving contact,
rather than a competition between two interior contacts.
\end{theorem}
\begin{proof}
For $h<M_0$, both component inverses satisfy
Theorem~\ref{p05:be21:convex}.  Their positive combination satisfies
$q''+q'>0$ and has positive infinite limits at both ends.
At $h=M_0$, the series gives local convergence
$g_{\kappa_0,c,h}\to k_c$ together with derivatives away
from zero.  Hence $k_c''+k_c'\geq0$.  The second summand has
strictly positive $g''+g'$ since its remaining exponent is
$\delta\in(0,1)$.  Thus $Q_c''+Q_c'>0$.
At zero,
\begin{equation}\label{p05:be21:atomic-origin}
 Q_c(x)=-M_0c+O(x)
       +\frac{\epsilon x^{\delta-1}}{\Gamma(\delta)}(1+O(x)),
\end{equation}
and at infinity the component of order $\kappa_1$ dominates
with the positive leading term in \eqref{p05:be21:ends}.
This proves the first assertion, including nondegeneracy.

No shape larger than $M_0$ can be feasible: its normalized
candidate Laplace transform grows as a positive constant times
$s^{h-M_0}$ as $s\to\infty$, whereas a probability Laplace
transform is bounded by one.  At $h=M_0$ the transform is the
one of $\delta_0+Q_c\,dx$, because subtracting its limit one
at infinity leaves exactly the series \eqref{p05:be21:atomic-series}
and the second component.  The exponential tilt makes its
total variation finite.  The full shape is therefore feasible
exactly when $Q_c\geq0$.

At $c=0$, all nonzero coefficients of the two inverse series
are positive: in the first summand they are the even coefficients
of $(1-z^2)^{-\kappa_0}$, and the second has remaining exponent
$\delta>0$.  Thus $Q_0(x)>0$ for every $x>0$.
This strict positivity persists for small positive $c$.
Indeed \eqref{p05:be21:atomic-origin} is uniform for $c$ in a small
compact interval, so the positive divergent term controls a
common neighborhood of zero.  The positive large-$x$ leading
term is uniform there as well; its coefficient is bounded
away from zero since $c$ stays below one.  The intervening
compact interval is handled by continuity and its positive
minimum.  This proves an initial plateau of positive length.

For fixed shape, feasibility is preserved on decreasing $c$.
To see this without a density comparison, if $c_1<c_2$, the
ratio of the normalized candidate transforms is the transform
of a compound Poisson variable whose finite Levy measure is
\[
       h\frac{e^{-(b-c_2)x}-e^{-(b-c_1)x}}{x}\,dx>0.
\]
Its density away from its atom is strictly positive on the
whole positive axis.  Thus the feasible full-shape parameters
form an initial interval.  When a full-shape remainder at
$c_2$ is feasible, it has a positive atom at zero, so its
convolution with this compound Poisson variable has strictly
positive continuous density at every $x>0$.  It follows that
$Q_{c_1}>0$.  Pointwise continuity shows that the feasible
interval is closed at any endpoint below one.  Its endpoint
is the asserted $c_*>0$, with $c_*=1$ if it fills $[0,1)$.

For each strict $c$, some positive shape is feasible: choose
$h$ smaller than both individual capacities, and combine the
two nonnegative inverses.  Feasibility in shape is an interval,
because decreasing a shape convolves the remainder with a
gamma variable.  It is closed: for a convergent sequence of
feasible shapes the normalized transforms converge pointwise
to a function continuous at $s=0$ with value one; the continuity
theorem for Laplace transforms gives a probability limit.
The explicit inverse identifies this limit, including the
possible atom at $h=M_0$.  Thus $K(c)$ is attained.

If $K(c)<M_0$, its inverse must have a zero.  Otherwise the
uniform positive endpoint estimates and a positive minimum
on the intervening compact set would permit a slightly larger
shape.  Its zero is the unique minimum and is double, by the
first assertion.  The response identities
\eqref{p05:be21:inherited-response} hold for $q$ as well, because
both summands have the same transform multipliers under
$c$ and $h$ differentiation.  They give $q_c<0$ and $q_h<0$
at the contact.  The implicit-function theorem, uniform
endpoint estimates, and the strict unique minimum identify
a locally analytic capacity, with derivative $-q_c/q_h<0$.
Uniqueness patches these branches throughout $(c_*,1)$.

Suppose now $c_*<1$.  If $Q_{c_*}$ were everywhere positive,
the same compact-set and uniform-endpoint argument would
preserve full-shape feasibility beyond $c_*$.  It therefore
has its unique zero $x_*$, with $Q_{c_*}''(x_*)>0$.
For fixed $x>0$, the series \eqref{p05:be21:g} extends analytically
in $h$ through $M_0$: $1/\Gamma(\eta)$ is entire and its zero
at $\eta=0$ removes the $n=0$ term pointwise.  Thus its
pointwise value there is exactly $Q_c(x)$.

The parameter responses at this atomic contact are
\begin{align}
 q_c(x_*)&=-M_0\left(e^{c_*x_*}
                +\int_0^{x_*}e^{c_*y}Q(x_*-y)\,dy\right)<0,
                                                    \label{p05:be21:atomic-c}\\
 q_h(x_*)&=-\frac{e^{c_*x_*}}{x_*}
          -\int_0^{x_*}\frac{e^{c_*y}Q(x_*-y)}{y}\,dy<0.
                                                    \label{p05:be21:atomic-h}
\end{align}
For the first identity the multiplier $-h/(w-c)$ acts on
the measure $\delta_0+Q\,dx$, and its atom gives the first
displayed term.  For the second, use the regularized identity
\[
 \log(w-c)=\int_0^\infty
               \frac{e^{-y}-e^{-(w-c)y}}{y}\,dy.
\]
On the positive axis the atom contributes $-e^{cx}/x$.
On the continuous part its inverse is
\[
 - (\gamma+\log x)Q(x)
       +\int_0^x\frac{Q(x)-e^{cy}Q(x-y)}{y}\,dy.
\]
This follows by cutting the integral off at $y=\rho$, then
letting $\rho\downarrow0$.  At the contact $Q(x_*)=0$ it
reduces to \eqref{p05:be21:atomic-h}.  The integrals converge at
$y=0$ by the double zero, and at $y=x_*$ because
$Q(u)=O(u^{\delta-1})+O(1)$ with $\delta>0$.
This calculation also agrees with the differentiated inverse
series, and does not discard the atom in the logarithmic response.

The contact Jacobian is $q_hQ''\ne0$, so there is a unique
analytic contact curve $h(c),x(c)$ through this point.  The
strict negative slope $h'=-q_c/q_h$ gives
\eqref{p05:be21:atomic-slope}.  Locally on its right, $h(c)<M_0$;
the unique minimum and uniform endpoint positivity identify
it with the capacity.  To justify uniformity while $h\uparrow M_0$,
put $\eta=M_0-h\geq0$.  The first component is a nonnegative
term $x^{\eta-1}/\Gamma(\eta)$ plus a uniformly bounded
remainder for small $x$; the second has the uniformly positive
leading term $\epsilon x^{\delta+\eta-1}/\Gamma(\delta+\eta)$.
Choose $\eta$ small enough that $\delta+\eta<1$.  This term
dominates the bounded remainder uniformly.  The tail estimate
is uniform because $c_*$ is strictly below one.  On the left
the shape ceiling forces $K=M_0$.  This proves
\eqref{p05:be21:corner} and completes the diagram.
\end{proof}

\begin{corollary}[An explicit mixture with a nonzero corner]
\label{p05:be21:example}
The source \eqref{p05:be21:mixture-source} with
\[
             \kappa_0=\frac34,\qquad \kappa_1=1,
             \qquad\epsilon=\frac1{100}
\]
has $0<c_*<1$, and hence has the genuine corner in
\eqref{p05:be21:corner}.  This holds for every $b>1$.
\end{corollary}
\begin{proof}
Here $M_0=\delta=1/2$.  We show that its full-shape inverse
is negative at $c=1$ and $x=1/100$, which by continuity is
also true for some $c<1$.  For either $\kappa=3/4$ or $1$,
the coefficients of $(1-z)^{1/2}(1-z^2)^{-\kappa}$ obey
$|a_n|\leq2^{n+1}$ by Cauchy's estimate on $|z|=1/2$:
the modulus there is at most $\sqrt{3/2}\,(4/3)<2$.
Since $a_1=-1/2$, the first continuous component satisfies
\[
 k_1(x)\leq-\frac12+4(e^{2x}-1)
                  \leq-\frac12+\frac4{49}
                         \quad(x=1/100).
\]
For the second component, $(1/2)_n\geq n!/2^n$ gives
\[
 |g_{1,1,1/2}(x)|
 \leq\frac{2x^{-1/2}}{\sqrt\pi}e^{4x}
 <20\frac{25}{24}=\frac{125}{6}.
\]
We used $e^u\leq(1-u)^{-1}$ for $0\leq u<1$.
Consequently
\[
 Q_1(1/100)<-\frac12+\frac4{49}+\frac5{24}
                       =-\frac{247}{1176}<0.
\]
The initial plateau has positive length by
Theorem~\ref{p05:be21:mixture}, and this strict inequality forces
its endpoint to lie below one.  The value of $b$ changes
only the positive tilt and normalization, not these signs.
\end{proof}


\section{Exact adaptive recovery: dimension, topology, and sparse models}
\label{p05:in21:section}

This section changes the recovery criterion in a deliberate way.  The
quantity $N(m)$ in the earlier observation papers concerns an algorithm
for every prescribed positive error; the algorithm may change with the
error.  Here an \emph{exact} algorithm must recover every input without
error.  We obtain a dimension theorem for this model, a cohomological
obstruction sharp on closed manifolds, and an exact measurement theorem
for sparse vectors with adaptive linear observations.  These are structural
transfers of classical topology and finite-candidate separation.  They do
not settle the unrestricted two-query approximation problem in
$\mathbb R^3$, and their global priority is not asserted.

\subsection{The exact model and its dimension interval}

For a nonempty compact metrizable space $X$, let $Q(X)$ be the least
number of deterministic adaptive continuous scalar queries that recover
every $x\in X$ exactly.  At a transcript $(y_1,\ldots,y_{j-1})$, the
algorithm selects a continuous function $f_{y_1,\ldots,y_{j-1}}:X\to
\mathbb R$ and receives its exact value.  The selection rule and decoder
are initially unrestricted.  Algorithms with fewer than the allowed
number of queries can be padded by constant queries.  Set $Q(X)=\infty$
if no finite number suffices.  Thus $Q(X)=0$ precisely for a singleton.
The feasible set at each realized transcript is compact.  Queries
defined only on such a feasible set may equivalently be extended to $X$
by the Tietze extension theorem.

We use covering dimension.  Two classical dimension theorems are the
inputs: for a continuous map of compact metric spaces,
\begin{equation}\label{p05:in21:Hurewicz}
 \dim X\leq \dim f(X)+\sup_y\dim f^{-1}(y),
\end{equation}
and a compact metric space of dimension at most $d$ admits a continuous
map to $[0,1]^d$ whose fibers have dimension zero.  The first is the
Hurewicz dimension-lowering theorem, recalled in
\cite{p05:in21:Tsukamoto}; the latter light-map characterization is stated
in \cite[Theorem~2.3]{p05:in21:Kato}.  No manifold or local connectedness
hypothesis is imposed on the fibers.

\begin{theorem}[Exact dimension interval]\label{p05:in21:dimension}
If $X$ is a nonempty compact metrizable space of finite dimension $d$,
then
\begin{equation}\label{p05:in21:interval}
 d\leq Q(X)\leq d+1.
\end{equation}
The upper bound can be realized with the first $d$ queries fixed in
advance, a Borel measurable rule selecting the last query in $C(X)$
with its uniform norm, and a Borel decoder.  If $\dim X=\infty$, then
$Q(X)=\infty$.
\end{theorem}

\begin{proof}
We first prove the lower bound by induction on the query budget $n$.
For $n=0$ the space is a singleton.  For $n\geq1$, let $f$ be the first
query.  Its restriction to each compact fiber leaves an exact algorithm
with $n-1$ queries.  The induction hypothesis and
\eqref{p05:in21:Hurewicz} give
\[
 \dim X\leq \dim f(X)+\sup_y\dim f^{-1}(y)\leq1+(n-1)=n.
\]
This also proves the infinite-dimensional assertion.

For the upper bound choose a light map $g:X\to[0,1]^d$ and query its
coordinates.  Each fiber $F_y=g^{-1}(y)$ is a zero-dimensional compact
metric space, so it embeds into the real line.  An embedding of $F_y$
extends continuously to $X$.  Querying that extension separates all
remaining candidates and proves the upper bound with unrestricted
selection.  For $d=0$ this construction is just a scalar embedding of
the zero-dimensional compactum.

Here are details ensuring that no nonmeasurable choice is necessary.
Write $Y=g(X)$ and $E=C(X)$ with its uniform norm, a separable Banach
space.  Fix a compatible metric $\rho$ on $X$.  For each $j\geq1$ let
\[
 U_j=\{(y,h)\in Y\times E:
       h(x)\ne h(x')\text{ whenever }x,x'\in F_y,
                                  \ \rho(x,x')\geq1/j\}.
\]
The set $U_j$ is open.  Indeed, its complement is the projection of the
closed conditions $g(x)=g(x')=y$, $\rho(x,x')\geq1/j$, and
$h(x)=h(x')$; compactness of $X\times X$ makes this projection closed.
Every section $(U_j)_y$ is dense in $E$.  To approximate a given $h$,
partition the zero-dimensional compactum $F_y$ into finitely many
clopen sets of diameter less than $1/j$, on each of which the oscillation
of $h$ is small.  Replace its values there by distinct nearby constants.
The resulting small continuous perturbation on $F_y$ extends with the
same uniform bound to $X$, and belongs to $(U_j)_y$.

For completeness, a Borel selector from $\bigcap_j(U_j)_y$ can be
constructed directly.  Choose countable bases of $Y$ and of open balls
in $E$.  Recursively, after a ball $B_{j-1}(y)$ has been chosen, take
the first pair of basic sets $V,B_j$ satisfying
\[
 y\in V,\qquad \overline{B_j}\subset B_{j-1}(y),\qquad
 \operatorname{diam}B_j\leq2^{-j},\qquad
 V\times\overline{B_j}\subset U_j.
\]
At the first step omit the inclusion in the previous ball.  Openness
of $U_j$ and density of its sections guarantee such a pair.  For each
fixed previous-ball index, the eligible pairs form a fixed countable
list, and choosing the first one with $y\in V$ is Borel.  The centers
of these balls are Borel maps of $y$, and their nested closures have a
unique limit $h_y\in E$.  This limit is Borel and belongs to every
$(U_j)_y$, hence is injective on $F_y$.

The transcript map $x\mapsto(g(x),h_{g(x)}(x))$ is therefore a Borel
injection into $\mathbb R^{d+1}$.  The Lusin--Souslin theorem gives a
Borel image and a Borel inverse on that image
\cite[Corollary~15.2]{p05:in21:Kechris}.  Extend the decoder arbitrarily
at unrealized transcripts.  This proves the stated measurable version.
\end{proof}

\subsection{A cohomological obstruction sharp on closed manifolds}

All cohomology below is \v{C}ech cohomology, with coefficients in an
abelian group $G$.  For compact metrizable spaces it agrees with
cohomology of the constant sheaf.  This choice matters because the
feasible fibers may be wild compacta.

\begin{theorem}[Top cohomology obstructs exact acquisition]
\label{p05:in21:cohomology}
If $X$ has an exact adaptive continuous algorithm with $n\geq1$
queries, then
\begin{equation}\label{p05:in21:vanishing}
 \check H^q(A;G)=0\qquad(q\geq n)
\end{equation}
for every compact subset $A\subset X$ and every abelian group $G$.
Consequently, if $\dim X=d\geq1$ and
$\check H^d(X;G)\ne0$ for some $G$, then $Q(X)=d+1$.
In particular, for every nonempty closed $d$-manifold, $d\geq1$,
\begin{equation}\label{p05:in21:manifold}
 Q(M)=d+1.
\end{equation}
\end{theorem}

\begin{proof}
We prove the assertion for $X$ by induction on $n$; restricting the
algorithm to a compact subset then proves the hereditary formulation.
If $n=1$, the query is a continuous injection into $\mathbb R$ and is
a homeomorphism onto its compact image.  Compact subsets of the real
line have zero positive-degree \v{C}ech cohomology: their arbitrarily
small neighborhoods can be taken to be finite unions of intervals.

Suppose $n\geq2$ and write $f:X\to Y=f(X)\subset\mathbb R$ for the
first query.  The map is proper.  Each fiber $F_y$ is exactly
recoverable with $n-1$ queries and therefore has
$\check H^q(F_y;G)=0$ for $q\geq n-1$.  Proper base change gives
\[
 (R^qf_*\underline G)_y\cong
 H^q(F_y;\underline G|_{F_y})=\check H^q(F_y;G),
\]
so $R^qf_*\underline G=0$ for $q\geq n-1$.  We use here the
topological proper-base-change theorem, which applies to compact
Hausdorff spaces without local contractibility assumptions
\cite[Tag~09V6]{p05:in21:Stacks}.

The compact metric space $Y\subset\mathbb R$ has covering dimension
at most one.  Cohomology of every sheaf of abelian groups on $Y$
vanishes above degree one.  This is the standard covering-dimension
bound for sheaf cohomology; for precise modern statements see
\cite[Theorem~7.2.3.6 and Corollary~7.2.2.30]{p05:in21:Lurie}.
The Leray spectral sequence
\[
 E_2^{p,q}=H^p(Y,R^qf_*\underline G)
        \ \Longrightarrow\ H^{p+q}(X;\underline G)
\]
thus has possible nonzero entries only for $p\leq1$ and $q\leq n-2$.
It follows that $H^k(X;\underline G)=0$ for $k\geq n$, completing
the induction.  The dimension interval now proves the first
consequence.  A nonempty closed $d$-manifold has nonzero top cohomology
with $\mathbb F_2$ coefficients, so the same argument proves
\eqref{p05:in21:manifold}.
\end{proof}

The interval in Theorem~\ref{p05:in21:dimension} is sharp at both ends:
$Q([0,1]^d)=d$ by coordinate observations, whereas closed manifolds
require $d+1$.  Vanishing of top cohomology is not a converse to the
obstruction.  A tripod has dimension one and trivial positive-degree
cohomology, but cannot embed into the real line, so its exact query
count is two.

There is also a precise reason why this theorem does not sharpen the
earlier metric-width bound.  On the unit sphere $S^3\subset\mathbb R^4$,
\begin{equation}\label{p05:in21:unattained}
 Q(S^3)=4,\qquad
 \inf_{\text{three-query algorithms}}\ 
 \sup_{x\in S^3}\|x-\widehat x\|_2=0.
\end{equation}
The second statement follows by restricting the three-query
arbitrary-precision algorithm of Krieg--Novak--Ullrich on
$\mathbb R^4$ \cite[Theorem~1]{p05:in21:KNU}; if outputs are required on
the sphere, radial normalization changes an error $\varepsilon<1$
to at most $2\varepsilon$.  Thus the infimum is genuinely unattained.
No universal positive-error conclusion can follow solely from the
cohomological exact-recovery obstruction.

\subsection{Finite polyhedra: adaptive linear observations beat embedding dimension}

\begin{proposition}[A finite last-query menu]\label{p05:in21:finite-menu}
Let $X\subset\mathbb R^M$ lie in the union of $L$ affine subspaces,
each of dimension at most $d$.  There are $d$ fixed linear queries
followed by one adaptively chosen linear query that recover $X$
exactly.  The last query can be selected from a fixed finite menu of
at most
\begin{equation}\label{p05:in21:menu-size}
 (M-1)\binom L2+1
\end{equation}
linear functionals.  For a finite union of affine subspaces the rule
and decoder can be chosen semialgebraic; for a finite polyhedron the
same statement holds after restricting candidates to its simplices.
\end{proposition}

\begin{proof}
Choose a linear map $A:\mathbb R^M\to\mathbb R^d$ whose restriction
to each of the finitely many direction spaces is injective.  Such
maps form a nonempty open dense set, as witnessed by the nonvanishing
of the relevant rank minors.  Each affine subspace then contributes
at most one candidate to a fiber $A^{-1}(y)$, so there are at most $L$
distinct candidates.

For distinct real numbers $t_1,\ldots,t_T$, consider the functionals
\[
 \ell_t(x)=x_1+t x_2+\cdots+t^{M-1}x_M,
 \qquad T=(M-1)\binom L2+1.
\]
For each pair of distinct candidates, the polynomial
$\ell_t(x-x')$ is nonzero and has degree at most $M-1$.  Hence at
most $(M-1)\binom L2$ menu entries fail to separate some pair.
Select the first entry separating all candidates, and use its exact
value to identify the input.  A singleton fiber requires no
separation and may use any menu entry.

For the measurability claim, the candidate in a particular affine
subspace, when it exists, is an affine function of $y$ on an affine
set.  Candidate validity and the finitely many tests of equality or
inequality of scalar values are semialgebraic.  Choosing the first
successful entry and the first consistent candidate is therefore
semialgebraic.  Membership in a fixed simplex adds only finitely many
linear inequalities.  For a general subset $X$ in the stated union,
one may use all candidates from the containing affine subspaces;
this already recovers every point of $X$.
\end{proof}

Let $K_d$ be the $d$-skeleton of a $(2d+2)$-simplex, $d\geq1$, in
its usual Euclidean realization.  Let $Q_{\rm na}(K_d)$ denote the
minimum number of fixed continuous scalar queries for exact recovery.
Then
\begin{equation}\label{p05:in21:VKF-gap}
 Q(K_d)=d+1,\qquad Q_{\rm na}(K_d)=2d+1.
\end{equation}
Indeed, $K_d$ has nonzero $d$th cohomology: the boundary of any
$(d+1)$-simplex is a nonzero simplicial $d$-cycle, and there are no
$(d+1)$-chains in $K_d$.  The cohomology theorem proves the adaptive
lower bound, and Proposition~\ref{p05:in21:finite-menu} gives an adaptive
\emph{linear} algorithm attaining it.  Fixed continuous exact
queries are equivalent to an embedding in their Euclidean transcript
space.  The van Kampen--Flores theorem rules out an embedding of
$K_d$ in $\mathbb R^{2d}$; see the primary formulation in
\cite{p05:in21:KishimotoMatsushita}.  A generic placement of its vertices
in $\mathbb R^{2d+1}$ gives a linear embedding: any two simplices
involve at most $2d+2$ vertices, which are affinely independent, so
their convex hulls intersect only in their common face.
This proves \eqref{p05:in21:VKF-gap}.  The separation concerns exact
deterministic acquisition, rather than a stable compressed-sensing
guarantee.

\subsection{An exact adaptive linear threshold for sparse vectors}

Write
\[
 \Sigma_s^m=\{x\in\mathbb R^m:|\operatorname{supp}x|\leq s\},
 \qquad 1\leq s<m.
\]
All measurement values in the next theorem are exact real numbers.
Adaptive choices need not vary continuously with previous values.

\begin{theorem}[Sparse exact linear query count]\label{p05:in21:sparse}
The minimum number of deterministic adaptive linear scalar queries
recovering every $x\in\Sigma_s^m$ exactly is $s+1$.  This remains
the minimum for $\Sigma_s^m\cap\overline B_2(0,R)$ for every $R>0$.
An optimal algorithm uses $s$ fixed queries followed by one query
selected semialgebraically from a finite menu.  For the lower bound,
the last query may even be an arbitrary continuous function, provided
the preceding $s-1$ queries are linear.
\end{theorem}

\begin{proof}
For the upper bound apply Proposition~\ref{p05:in21:finite-menu} to the
$\binom ms$ coordinate subspaces of dimension $s$.  Explicitly, take
an $s\times m$ matrix $A$ all of whose $s\times s$ column minors are
nonzero, for example a Vandermonde matrix with distinct nodes.  Every
support set $S$ of size $s$ gives one candidate $x^S(y)$, obtained by
solving $A_Sx_S=y$ and setting the other coordinates to zero.  Every
input with smaller support is included by extending that support to
size $s$.  Duplicate candidates are harmless.  The finite menu in
\eqref{p05:in21:menu-size} separates the distinct candidates after $s$
fixed observations.  The same algorithm recovers a bounded sparse
class without modification.

Suppose, conversely, that $s$ queries suffice with the first $s-1$
linear and the last continuous.  Follow the transcript generated by
$x=0$ for the first $s-1$ queries.  The selected linear functionals
are then fixed.  Restrict their common kernel to a coordinate
$(s+1)$-plane $E\subset\mathbb R^m$; the resulting linear space
$W\subset E$ has dimension at least two.  We claim that
$W\cap\Sigma_s^m$ contains two distinct line segments crossing at
zero.

If $W$ lies in one coordinate hyperplane of $E$, every vector of $W$
is $s$-sparse and any two independent directions work.  Otherwise,
all $s+1$ coordinate functionals restrict nontrivially to $W$.
If their kernels all coincided, these functionals would be
proportional, and their common kernel would have codimension one
in $W$.  But the common kernel of all coordinates of $E$ is zero,
which would force $\dim W=1$, a contradiction.  Thus two of their
kernels are distinct.  Each has positive dimension; choose nonzero,
nonproportional vectors from two such kernels.  Each vector, and
its entire line, has a zero coordinate in $E$ and is $s$-sparse.
Scaling the segments puts their cross inside the ball of any
prescribed radius $R>0$.

All points of this cross produce the same first $s-1$ observations
as zero, by induction along that transcript.  The final selected
continuous scalar query would therefore have to be injective on the
cross.  This is impossible.  A continuous injection of a compact
cross into $\mathbb R$ would be a homeomorphism onto a compact
interval, whereas removing its crossing point gives four connected
components and removing a point from an interval gives at most two.
This contradiction proves the lower bound.
\end{proof}

\begin{proposition}[A uniform sparse approximation obstruction]
\label{p05:in21:sparse-error}
Every $s$-query algorithm on
$\Sigma_s^m\cap\overline B_2(0,R)$, $1\leq s<m$, whose first $s-1$
queries are linear and whose final query is continuous has worst-case
Euclidean error at least $R/2$.  The decoder and the selection of
queries are unrestricted.
\end{proposition}

\begin{proof}
Follow the zero transcript for the first $s-1$ observations, and form
$W$ inside a coordinate $(s+1)$-plane as in the preceding proof.  If
$W$ is contained in a coordinate hyperplane, it contains a Euclidean
two-dimensional disk of radius $R$ consisting entirely of sparse
vectors.  The final scalar query identifies two antipodal boundary
points by the Borsuk--Ulam theorem, giving the stronger error bound
$R$.

Otherwise choose a two-dimensional subspace $V\subset W$ that is
contained in none of the coordinate hyperplanes.  Such a choice
exists: take a vector of $W$ outside their finite union and extend
it to a two-dimensional subspace.  Identify $V$ isometrically with
$\mathbb R^2$.  Its $s+1$ coordinate functionals are nonzero row
vectors $r_i$ satisfying $\sum_i r_i r_i^{\mathsf T}=I_2$.
If $u_i$ is a unit direction perpendicular to $r_i$, and
$w_i=\|r_i\|_2^2>0$, then
\begin{equation}\label{p05:in21:isotropic-rays}
 \sum_i w_i u_i u_i^{\mathsf T}=I_2,\qquad \sum_i w_i=2.
\end{equation}
The line segments $[-R u_i,R u_i]$ belong to the sparse feasible set.

Write $f$ for the final selected query.  If an endpoint has value
$f(0)$, it and zero already form a same-value pair at distance $R$.
If the two endpoints of any segment have values strictly on the same
side of $f(0)$, take the endpoint whose value is closer to $f(0)$
and use the intermediate value theorem on the opposite half-segment.
This again gives a same-value pair at distance at least $R$.

In the remaining case orient each $u_i$ so that $f(Ru_i)>f(0)$;
orientation does not alter \eqref{p05:in21:isotropic-rays}.  At least two
of these unit vectors have nonpositive inner product.  To see this,
suppose all pairwise inner products were positive.  The vector
$v=\sum_iw_i u_i$ has positive inner product with every $u_i$, so
their angular representatives lie in a single open semicircle.
Pairwise positivity then places all of them in an arc of angular
length strictly less than $\pi/2$.  For its bisector $a$ we have
$(a\cdot u_i)^2>1/2$ for every $i$, contradicting
$1=\sum_iw_i(a\cdot u_i)^2>\tfrac12\sum_iw_i=1$.

Choose $u_i,u_j$ with $u_i\cdot u_j\leq0$, and suppose without loss
that $f(Ru_i)\leq f(Ru_j)$.  By the intermediate value theorem there
is $t\in[0,R]$ with $f(tu_j)=f(Ru_i)$.  Their distance satisfies
\[
 \|Ru_i-tu_j\|_2^2
   =R^2+t^2-2Rt(u_i\cdot u_j)\geq R^2.
\]
They have the same entire transcript, so any common reconstruction
errs by at least half their distance on one of them.
\end{proof}

For comparison, the exact nonadaptive linear count on $\Sigma_s^m$
is $\min(2s,m)$.  Indeed, injectivity is equivalent to the kernel
containing no nonzero vector supported on at most $2s$ coordinates:
every such vector is a difference of two $s$-sparse vectors.  Fewer
than $\min(2s,m)$ rows always have such a kernel vector by a
column-dependence argument, while a full-spark matrix with
$\min(2s,m)$ rows attains the bound.  Thus the adaptive result improves
the exact deterministic count from $2s$ to $s+1$ when $m\geq2s$.
This comparison should not be confused with the much stronger
noise-stable and approximately sparse recovery requirements studied,
for example, in \cite{p05:in21:PriceWoodruff}.  Our finite-candidate
algorithm has no uniform noise guarantee, and its candidate
enumeration can be exponentially expensive in $m$.

\subsection{The non-dyadic approximation target remains open}

The primary version of \cite{p05:in21:KNU}, dated September 2026,
explicitly asks whether two adaptive continuous queries approximate
all of $\mathbb R^3$ uniformly to arbitrary prescribed accuracy.
The earlier width theorem in this project gives
$\lceil\log_2(m+1)\rceil\leq N(m)\leq\lceil\log_2m\rceil+1$,
and hence $2\leq N(3)\leq3$.  The present dimension and cohomology
arguments do not remove that gap.  They start from injective final
fibers, whereas an approximate algorithm only forces its final
fibers to have small diameter.  Replacing one condition by the other
would invalidate both inductions.  Equation~\eqref{p05:in21:unattained}
is an explicit counterexample to that replacement as a general
principle.  Accordingly, the non-dyadic approximation problem is
recorded as an attempted but unresolved external target, not as a
consequence of the exact theorems.

\section{Failed shortcuts and the next proof obligations}
\label{p05:rv5:failed}

The successful transfers do not remove the distinction between an exact
classification and a route toward an external open problem. This section
records specific failures that affect how the results can be used.

\subsection{A contour bound cannot simply lose its prefactor}

In the earlier collected manuscript, shifted Chebyshev polynomials have
Bernstein coefficients defined by
\[
 T_m(2x-1)=\sum_{j=0}^t C(t,m,j)\binom tjx^j(1-x)^{t-j}.
\]
The constant-one exponential estimate was already proved there. The
remaining best-constant question is whether $|C(t,m,j)|\le2^{m^2/t}$
for all $m\le t$, with the inherited bounds $\log2\le c_*\le1$
for the optimal exponent constant. It is not a new closure in this paper.

One possible strengthening of the contour method would be
\[
 |C(t,m,j)|\stackrel{?}{\le}\exp\{t\mathsf H(m/t)\},\qquad
 \mathsf H(a)=\tfrac12[(1+a)\log(1+a)+(1-a)\log(1-a)].
\]
This strengthening is false, even though the same expression with a
prefactor occurs in a valid earlier estimate. Indeed, at $(t,m,j)=(6,4,1)$,
\[
 C(6,4,1)=1-\frac{32}{6}=-\frac{13}{3},\qquad
 \exp\{6\mathsf H(2/3)\}=\frac{3125}{729},\qquad
 \frac{13}{3}=\frac{3159}{729}>\frac{3125}{729}.
\]
The coefficient identity follows from $T_4(2x-1)$ having constant
coefficient $1$ and derivative $-32$ at zero, and the identity
$p'(0)=t(b_1-b_0)$ for degree-$t$ Bernstein coefficients.
This is a counterexample to the proposed shortcut, not to the weaker
$2^{m^2/t}$ conjecture.

\subsection{Four boundaries that remain after the new theorems}

\paragraph{Identification versus rates.}
The unrestricted two-replicate theorem leaves a quantitative inverse
problem. Its proof uses a nonzero Fourier frequency and a weight depending
on that frequency; it does not provide a uniform finite-data estimator.
The explicit large-scale instability forbids uniform continuity in the
stated unrestricted latent metric for any fixed group size. A useful next
theorem must therefore specify tail constraints, a weaker metric, or a
local parameter class before claiming a rate.

\paragraph{Differential classification versus special values.}
The finite-spectrum classification concerns functions and their analytic
continuations. It does not decide whether an optimized value at one
algebraic setting is algebraic or transcendental. The exact jet field
identifies the available functional information; it is not a specialization
theorem. Similarly, the moment-based square-root rate is a credited
classical comparison after a new observation-coordinate conversion.

\paragraph{Exact acquisition versus approximation.}
The manifold obstruction applies to a single algorithm attaining zero
error. It does not exclude a sequence of algorithms with errors tending
to zero. The sphere example proves that such an inference would be false.
The problem $N(3)$ therefore needs a genuinely quantitative obstruction
or a new approximation construction; exact top cohomology alone cannot
settle it. The sparse theorem likewise uses exact scalar values and a
discontinuous adaptive choice; it makes no claim against noisy recovery
lower bounds.

\paragraph{Relative algebraic closure versus all units.}
The arithmetic closure theorem deals exactly with elements already
algebraic over a specified exponential field. It excludes an algebraic
route to an exotic unit, but does not prove that every unit must take that
route. Nor does a product lying in that field make either factor algebraic
over it. The joint $\operatorname{Ein}$ theorem similarly classifies all
possible finite extensions without asserting that its one possible
exception actually occurs.

\subsection{Research judgment}

The unrestricted pair-identification theorem and the finite-spectrum
differential classification are the strongest candidates for focused
follow-up papers: each resolves a complete family and has a proof
mechanism not restricted to a single numerical instance. The algebraic
optimizer decision procedure is another reusable result, but its current
input model and smoothness promise must be kept explicit. The extraction
corner provides a concrete analytic transition, and the exact-query
theorems give sharp results in a different observation model.

Whether any of these formulations has publication-level novelty requires
broader specialist comparison and external checking. The present work
does not justify claiming an Annals-level breakthrough, a Fields-level
result, or a newly solved historical named conjecture. It does provide
explicit statements, proofs, counterexamples, and a finite list of next
obligations on which such an assessment can be based.

% END SOURCE: KEIO_Structural_Rigidity_from_Adaptive_Observation.tex

\endgroup
\clearpage
\part{Logic and Effective Transfer}
\label{merge:part06}
\begingroup

% BEGIN SOURCE: KEIO_Logic_After_Observation_57_Problems.tex

\noindent\textbf{Editorial integration note.} 
The fifty-seven-entry register is preserved in full, including version distinctions and qualified external announcements. Its status columns are dated source assessments, not a fresh literature survey performed for this compilation. Solved restricted cases are not counted as solutions of the unrestricted listed problems.


\begin{quote}\small
\noindent\textit{Original part framing.}\par
\noindent\textbf{Reading convention.}
``Proved here'' describes the argument supplied in this document; it does
not certify novelty. ``Recent preprint'' identifies the authors' stated
result and its checked scope, not external peer review or formal validation.
A proof audit of selected steps is not a certification of an entire paper.
All bibliographic status statements refer to sources checked on
16--17 October 2026. The original A/B/C labels are retained as identifiers,
not as a ranking endorsed by this document.
\end{quote}
\section{Scope, evidence, and the strongest transfer}
\label{p06:logic:scope}

The task is to try every item, not to assume that a list of celebrated
conjectures must contain one that the available methods can settle.
We therefore keep the user's numbering fixed. Two entries explicitly
contain two questions: Skolem and positivity (37), and BMTO and MTWO (43).
Both components are addressed. The two Shelah categoricity labels (7 and
29) are retained separately and their potentially different meanings are
explained. Broad labels such as ``main gap'', ``universality spectrum'',
and ``cardinal transfer for omitted types'' require a specified class,
language, or cardinal range before they have a single truth value.

Earlier observation and optimization manuscripts motivate three routes:
arithmetic information in logarithmic identities, spectral information in
finite observations, and Borel coding of objects. The present document
states the hypotheses needed for each transfer explicitly. It does not
infer a theorem about arbitrary models, arbitrary Borel equivalence
relations, or arbitrary proof systems merely from a result about a finite
or identifiable observation model.

\paragraph{An exact boundary rather than an analogy.}
Let $\mathbb A=\overline{\mathbb Q}\cap\mathbb R$. Consider polynomial
predicates over $\mathbb A$ together with predicates
\[
 P_i(\mathbf x)>0\ (1\leq i\leq m),\qquad
 \sum_{i=1}^m \alpha_i\log P_i(\mathbf x)=0,
 \quad \alpha_i\in\mathbb A\text{ fixed}.
\]
Arbitrary first-order formulas formed from these predicates are effectively
equivalent, on $\mathbb A$, to formulas in the language of real closed
fields. In contrast, adding just
\[
 R(x,y)\quad\Longleftrightarrow\quad
 y>0\ \text{ and }\ x\log 2=\log y
\]
defines $\mathbb Q$ by $\exists y\,R(x,y)$ and makes the full first-order
theory undecidable. Theorem~\ref{p06:lmod:boundary} supplies the proof,
including the algebraic-domain restriction and the external transcendence
and definability inputs. This does not decide the theory of the real
exponential field, and it does not assert undecidability of an existential
fragment. Its significance is a precise failure point for transferring
fixed-coefficient arithmetic tests to variable-coefficient logic.

\paragraph{What a successful closure would require.}
For an original named problem, the target statement, hypotheses, and
conclusion must match those in its source; the new argument must discharge
the missing step. A finite subcase, a conditional reduction, a recent
result of another author, or a reformulation of the obstruction is useful
research information but is not such a closure. No original closure of a
previously open named problem is established here. In particular, the
number of separate entries or proved auxiliary statements is not evidence
of Annals-level significance.

\paragraph{How to use the dossier.}
The original-order index below locates all fifty-seven targets. The six
technical sections group related proof attempts so that shared assumptions
and failed transfers can be checked together. Each entry records a concrete
attempt and its remaining obstruction. The final synthesis identifies
which missing conclusions would actually advance the original targets.


\clearpage
\subsection{Original-order index: all fifty-seven entries}
\label{p06:logic:index}
The short status column is a navigation aid; qualifications, versions, sources, attempts, and verdicts appear at the linked pages. ``Unresolved'' does not certify the absence of every unpublished claim.
\begingroup\small\setlength{\tabcolsep}{3pt}
\begin{longtable}{L{13mm}L{64mm}L{63mm}r}
\toprule ID & Requested target & Status or necessary distinction & Page\\\midrule\endfirsthead
\toprule ID & Requested target & Status or necessary distinction & Page\\\midrule\endhead
1 (A) & Partition Principle implies Choice in ZF & Unresolved; countermodel claim withdrawn & \pageref{p06:lset:item1}\\[2pt]
2 (A) & Superpolynomial lower bounds for Frege proofs & Unrestricted bound unresolved & \pageref{p06:lproof:frege}\\[2pt]
3 (B) & Continuum problem after Cohen & CH independent of ZFC; new-axiom question differs & \pageref{p06:lset:item3}\\[2pt]
4 (B) & Vaught's conjecture & General first-order case unresolved & \pageref{p06:lmod:id4}\\[2pt]
5 (B) & Martin's conjecture on Turing degrees & General case unresolved; order-preserving case known & \pageref{p06:ldes:martin}\\[2pt]
6 (B) & Tarski's exponential-function problem & Unresolved over all reals; restricted boundary proved here & \pageref{p06:lmod:id6}\\[2pt]
7 (B) & Shelah's categoricity conjecture & Exact infinitary threshold distinguished from eventual transfer & \pageref{p06:lmod:id7}\\[2pt]
8 (B) & Shelah's PCF conjecture & Two variants; recent relative counterexample claim & \pageref{p06:lset:item8}\\[2pt]
9 (B) & Woodin's Omega conjecture & Completeness direction unresolved & \pageref{p06:lset:item9}\\[2pt]
10 (B) & HOD conjecture & Conditional negative results; consistency premise essential & \pageref{p06:lset:item10}\\[2pt]
11 (B) & Consistency of Reinhardt cardinals & Choice-free case unresolved; 2026 claim withdrawn & \pageref{p06:lset:item11}\\[2pt]
12 (B) & Automorphisms of the Turing degrees & Global rigidity unresolved & \pageref{p06:lcomp:12}\\[2pt]
13 (B) & NIP field conjecture & General case unresolved; finite-rank results differ & \pageref{p06:lmod:id13}\\[2pt]
14 (B) & Weiss's hyperfiniteness conjecture & General Borel case unresolved & \pageref{p06:ldes:weiss}\\[2pt]
15 (B) & Fra\"iss\'e in $\mathsf{ATR}_0$ & Laver theorem known; ATR implication unresolved & \pageref{p06:lproof:fraisse}\\[2pt]
16 (B) & Kolmogorov-Loveland versus Martin-Lof randomness & Equality unresolved; encoded result does not imply it & \pageref{p06:lcomp:16}\\[2pt]
17 (B) & $\Pi^1_1$ conservation for $\mathsf{RT}^2_2$ & Partial theorem; new full claim not verified & \pageref{p06:lproof:ramsey}\\[2pt]
18 (B) & Kreisel's conjecture & Kleene-presentation version unresolved & \pageref{p06:lcomp:18}\\[2pt]
19 (B) & Sacks's embedding conjecture & General embedding problem unresolved & \pageref{p06:lcomp:19}\\[2pt]
20 (B) & Borel Ruziewicz problem & Borel version distinguished from completed-measure theorem & \pageref{p06:ldes:ruziewicz}\\[2pt]
21 (B) & Roitman: $\mathfrak d=\aleph_1\Rightarrow\mathfrak a=\aleph_1$ & Recent positive proof claim; core construction checked & \pageref{p06:lset:item21}\\[2pt]
22 (B) & Consistency strength of PFA & General equiconsistency unresolved & \pageref{p06:lset:item22}\\[2pt]
23 (B) & Topological Vaught conjecture & General Polish-action case unresolved & \pageref{p06:ldes:vaught}\\[2pt]
24 (B) & Reverse-mathematical strength of Hindman's theorem & Exact strength unresolved; new avoidance result partial & \pageref{p06:lproof:hindman}\\[2pt]
25 (B) & Ordinal analysis of $\Pi^1_2$-$\mathsf{CA}_0$ & Exact axiom system needs matching ordinal calibration & \pageref{p06:lproof:ordinal}\\[2pt]
26 (B) & Friedman's grand conjecture concerning EFA & Publication corpus and formalization must be specified & \pageref{p06:lproof:efa}\\[2pt]
27 (B) & Zilber's quasiminimality conjecture & Full complex exponential case unresolved & \pageref{p06:lmod:id27}\\[2pt]
28 (C) & Cherlin-Zilber conjecture & Finite-Morley-rank conjecture unresolved & \pageref{p06:lmod:id28}\\[2pt]
29 (C) & Shelah's eventual categoricity conjecture & Full AEC proof claims; long proofs not certified here & \pageref{p06:lmod:id29}\\[2pt]
30 (C) & Main gap conjecture & Canonical countable first-order case already a theorem & \pageref{p06:lmod:id30}\\[2pt]
31 (C) & Kueker's conjecture & General case unresolved; stable case known & \pageref{p06:lmod:id31}\\[2pt]
32 (C) & Stable field conjecture & General case unresolved; large-field case known & \pageref{p06:lmod:id32}\\[2pt]
33 (C) & Stable forking conjecture & Recent external counterexample preprint & \pageref{p06:lmod:id33}\\[2pt]
34 (C) & Universality of C-free graphs & General classification not verified solved; tree case known & \pageref{p06:lfin:34}\\[2pt]
35 (C) & Universality spectrum problem & Precise spectrum convention selected; reduction proved & \pageref{p06:lfin:35}\\[2pt]
36 (C) & Generalized star-height problem & Recent height-three bound claim; height-one question remains & \pageref{p06:lfin:36}\\[2pt]
37 (C) & Skolem problem and LRS Positivity & Both general problems unresolved; promised class decided here & \pageref{p06:lcomp:37}\\[2pt]
38 (C) & Finite-model property of Henson graphs & Pseudofiniteness, not a finite model of the full theory & \pageref{p06:lfin:38}\\[2pt]
39 (C) & Finiteness of reducts of homogeneous structures & Finite relational version unresolved; arity reformulation proved & \pageref{p06:ldes:reducts}\\[2pt]
40 (C) & O-minimal theories with a superexponential function & Transexponential target; competing indexed claim unverified & \pageref{p06:lfin:40}\\[2pt]
41 (C) & Categoricity transfer in atomic classes & Underspecified; precise high-cardinal benchmark examined & \pageref{p06:ldes:atomic}\\[2pt]
42 (C) & Algebraic closure of characteristic-zero minimal fields & Characteristic-zero case unresolved & \pageref{p06:lfin:42}\\[2pt]
43 (C) & Decidability of BMTO and MTWO & BMTO decidable; full MTWO not verified resolved & \pageref{p06:lcomp:43}\\[2pt]
44 (C) & Decidability of Laurent series fields & Positive-characteristic pure-field target unresolved & \pageref{p06:lcomp:44}\\[2pt]
45 (C) & Compact logics with Beth and Delta interpolation & Literal existence known; interpolation separation differs & \pageref{p06:lset:item45}\\[2pt]
46 (C) & Cardinality transfer for models omitting types & Sharp Hanf bound known; specific transfers need parameters & \pageref{p06:ldes:omit}\\[2pt]
47 (C) & Barendregt-Geuvers-Klop conjecture & Whole-PTS implication unresolved & \pageref{p06:lproof:bgk}\\[2pt]
48 (C) & Decidability of multiplicative-exponential linear logic & General case unresolved; effective bound missing & \pageref{p06:lfin:48}\\[2pt]
49 (C) & Borel boundedness of countable Borel equivalence relations & General boundedness and converse questions unresolved & \pageref{p06:ldes:bounded}\\[2pt]
50 (C) & Are all Polish groups graphic? & Many graphic classes known; general case unresolved & \pageref{p06:ldes:graphic}\\[2pt]
51 (C) & Scott sets as standard systems of PA models & At most aleph-one known; arbitrary size unresolved & \pageref{p06:lcomp:51}\\[2pt]
52 (C) & Increasing unions of hyperfinite countable Borel equivalence relations & General increasing-union problem unresolved & \pageref{p06:ldes:union}\\[2pt]
53 (C) & Word problem for the S combinator & Normalization decidable; word problem unresolved & \pageref{p06:lfin:53}\\[2pt]
54 (C) & Expansion postponement in pure type systems & Normalizing-types case known; unrestricted version unresolved & \pageref{p06:lproof:expansion}\\[2pt]
55 (C) & Proof-theoretic strength of the graph minor theorem & GMT true; exact proof strength not calibrated here & \pageref{p06:lproof:graphminor}\\[2pt]
56 (C) & Decidability of three-pair Post correspondence & Ordinary three-pair problem not verified resolved & \pageref{p06:lfin:56}\\[2pt]
57 (C) & The $\forall\exists$ theory of c.e. degrees & Pure-order two-block theory not verified resolved & \pageref{p06:lcomp:57}\\[2pt]
\bottomrule\end{longtable}\endgroup
\newpage


\section{Model theory: an effective logarithmic boundary and eleven distinct problems}
\label{p06:lmod:section}

The eleven labels considered here do not denote eleven interchangeable
conjectures. Some have several inequivalent formulations, the canonical
countable first-order main gap is already a theorem, and two labels have
very recent primary manuscripts announcing resolutions. We distinguish
those external claims from results proved in this section. The status
review is dated 17 October 2026. Our positive mathematical contribution
is the following precisely delimited transfer of the logarithmic
splitting method; its priority as a formulation is not certified.

\subsection{Fixed versus variable logarithmic coefficients}

Let $A=\overline{\mathbb Q}\cap\mathbb R$, with its usual ordered-field
structure. Algebraic inputs are represented by integer minimal
polynomials and rational isolating intervals. All logarithms below are
ordinary real logarithms, evaluated externally in $\mathbb R$.
For a finite list of specified coefficients $\lambda_i\in A$ and
polynomials $P_i\in A[X_1,\ldots,X_n]$, introduce the relation
\begin{equation}\label{p06:lmod:atom}
 Z_{\lambda,P}(x)\quad\Longleftrightarrow\quad
 \bigwedge_i P_i(x)>0\quad\text{and}\quad
 \sum_i\lambda_i\log P_i(x)=0,
 \qquad x\in A^n.
\end{equation}
In particular, the relation is false outside its positivity guard.
The subscripts are effective finite syntactic data, not quantified
variables. Let $\mathcal A_{\rm fix}$ have all these effectively indexed
relation symbols and names for algebraic constants.

\begin{theorem}[A decidability and independence boundary]\label{p06:lmod:boundary}
The following assertions hold.
\begin{enumerate}
\item Every relation \eqref{p06:lmod:atom} has an effectively computable
semialgebraic definition over $A$. Consequently the complete
first-order theory of $\mathcal A_{\rm fix}$ is decidable, with arbitrary
quantifier alternation over $A$. This is a definitional expansion of a
real closed field; its theory is o-minimal and NIP.
\item Add instead just the binary relation
\begin{equation}\label{p06:lmod:variable}
 P(x,y)\quad\Longleftrightarrow\quad
 y>0\ \text{and}\ x\log2=\log y,
 \qquad x,y\in A.
\end{equation}
In $(A;0,1,+,-,\cdot,<,P)$ the sets $\mathbb Q$ and $\mathbb Z$
are definable without parameters. Its full first-order theory is
undecidable and has the independence property; the structure is not
o-minimal.
\end{enumerate}
These conclusions do not settle decidability of the real exponential
field. No existential undecidability conclusion is asserted.
\end{theorem}

\begin{proof}
We use Baker's theorem that logarithms of nonzero algebraic numbers
which are linearly independent over $\mathbb Q$ are linearly independent
over $\overline{\mathbb Q}$ \cite{p03:mi3:Baker}. In the present proof all
logarithms are real, so branch choices do not arise.

Compute a $\mathbb Q$-basis $\theta_1,\ldots,\theta_d$ of the span of
the finitely many $\lambda_i$, and write
$\lambda_i=\sum_k r_{ki}\theta_k$, $r_{ki}\in\mathbb Q$.
This is effective: embed the input numbers in a finite number field
using exact algebraic arithmetic, express them in a rational basis of
that field, and perform rational row reduction. Zero coefficients and
the case $d=0$ cause no difficulty.

Fix an algebraic tuple $x$ satisfying the positivity guard. Choose a
$\mathbb Q$-basis $\rho_1,\ldots,\rho_s$ of the span of the real numbers
$\log P_i(x)$. Each $\rho_j$ can be chosen from those logarithms.
Write $\log P_i(x)=\sum_j q_{ij}\rho_j$ with rational $q_{ij}$.
Baker's theorem and then the independence of the $\theta_k$ give
\begin{align*}
 \sum_i\lambda_i\log P_i(x)=0
 &\iff \sum_k\theta_k\sum_i r_{ki}q_{ij}=0
       \quad\text{for every }j\\
 &\iff \sum_i r_{ki}q_{ij}=0
       \quad\text{for every }k,j\\
 &\iff \sum_i r_{ki}\log P_i(x)=0
       \quad\text{for every }k.
\end{align*}
If all logarithms vanish, this is the same argument with an empty
basis. Choose positive integers $D_k$ clearing denominators and put
$n_{ki}=D_kr_{ki}\in\mathbb Z$. By positivity, the last conditions are
equivalent to the polynomial equations
\begin{equation}\label{p06:lmod:polynomial}
 \prod_{i:n_{ki}>0}P_i(x)^{n_{ki}}
   =\prod_{i:n_{ki}<0}P_i(x)^{-n_{ki}}
       \qquad(1\leq k\leq d).
\end{equation}
An empty product is one. Importantly, the algorithm computes the
$r_{ki}$, but never needs to find a basis of the logarithms at an input
point. That latter basis was used only to prove correctness.

Replace each occurrence of \eqref{p06:lmod:atom} in a formula by its guard
and \eqref{p06:lmod:polynomial}. Structural induction on formulas preserves
truth, including negation and every alternation of quantifiers, because
the pointwise equivalence holds at every tuple in $A$. The field $A$ is
real closed. Effective real-closed-field quantifier elimination with
exact algebraic coefficients therefore decides every resulting
sentence \cite{p06:lmod:BPR}. Algebraic constant names can also be removed
using their polynomial and isolating-interval definitions. The new
relations introduce no new definable sets; o-minimality and NIP follow
from the corresponding properties of real closed fields. This proves
the first assertion.

For the second assertion set $D(x):=\exists y\,P(x,y)$. If $x=p/q$ is
rational, with $q>0$, the positive real number $2^{p/q}$ is algebraic,
so $D(x)$ holds. If $x\in A\setminus\mathbb Q$, the
Gelfond--Schneider theorem (also treated in \cite{p03:mi3:Baker}) makes
$2^x$ transcendental, so $D(x)$ fails.
Thus $D(A)=\mathbb Q$. Julia Robinson's definition of the integers in
the rational field, with every quantifier relativized to $D$, now
defines $\mathbb Z$ in our structure \cite{p06:lmod:Robinson}. The same
relativization interprets the full first-order theory of the rational
field. Its undecidability, proved by Robinson, gives the claimed
undecidability.

For completeness, let $Z$ denote the resulting integer predicate and
consider the single formula
\[
 \delta(p,n):=Z(p)\wedge Z(n)\wedge
     \exists k\,(Z(k)\wedge n=pk).
\]
For any $r$ choose distinct positive rational primes $p_1,\ldots,p_r$.
For each subset $S\subseteq\{1,\ldots,r\}$ take
$n_S=\prod_{i\in S}p_i$, with $n_\varnothing=1$.
Unique factorization yields
$\delta(p_i,n_S)$ if and only if $i\in S$. This is the independence
property. Finally, the definable set $D(A)=\mathbb Q$ is dense and
codense in the ordered field $A$, so it is not a finite union of points
and intervals. Hence the expansion is not o-minimal.
\end{proof}

\begin{remark}[Why this is not a solution of Tarski's problem]
\label{p06:lmod:domain}
The equivalence \eqref{p06:lmod:polynomial} is proved on algebraic tuples,
not on arbitrary real tuples. In addition, $A$ is not a definable
domain in the o-minimal real exponential field. Restricting an external
relation to $A$ and quantifying over $A$ changes the structure: the
sentence $\forall x\exists y\,P(x,y)$ is true over $\mathbb R$ and false
over $A$, for instance at $x=\sqrt2$. Thus the negative half of the
theorem cannot be transferred to $\mathbb R_{\exp}$ by merely reusing
the displayed formula. The undecidability proof uses full first-order
arithmetic of $\mathbb Q$; it does not settle Hilbert's tenth problem
over $\mathbb Q$.
\end{remark}

\subsection{ID 4: the first-order Vaught conjecture}
\label{p06:lmod:id4}

The precise statement is that a complete theory $T$ in a countable
first-order language, with infinite models, has either at most
countably many or exactly $2^{\aleph_0}$ countable models up to
isomorphism. Pillay--Tanovi\'c's August 2026 revision still treats this
as open \cite{p06:lmod:PillayTanovic}. Results or announcements for a
single $L_{\omega_1,\omega}$ sentence are not automatically results
for this first-order formulation.

The attempted transfer is to construct complete measurable
observational invariants of countable models. Its precise sufficient
condition is classical: if isomorphism on the standard Borel space
$\operatorname{Mod}_{\omega}(T)$ is coanalytic, Silver's dichotomy
gives countably many classes or a perfect set of pairwise
nonisomorphic models \cite{p06:lmod:Silver}. In the latter case there are
$2^{\aleph_0}$ classes; in the former there are at most countably many.
The model space is Borel because satisfaction of each first-order
sentence is Borel and $T$ is countable. This proves the conditional
transfer. In general isomorphism is analytic, and no coanalyticity
argument follows from finite-dimensional recovery. Producing the
required complete invariant or complexity bound is the missing step.
Verdict: no resolution of the general first-order conjecture.

\subsection{ID 6: Tarski's real exponentiation decision problem}
\label{p06:lmod:id6}

The target is an algorithm deciding every sentence in
$\operatorname{Th}(\mathbb R;0,1,+,-,\cdot,<,\exp)$, with quantifiers
over all real numbers. The unconditional problem remains open in the
reviewed primary literature. Macintyre--Wilkie's conditional
decidability theorem under the appropriate real Schanuel conjecture
is reviewed in Krapp's 2024 paper \cite{p06:lmod:Krapp}.
Theorem~\ref{p06:lmod:boundary} supplies an unconditional algorithm for an
explicit logarithmic equality language over $A$, and locates an
undecidable extension of that language. The exact failures of a
transfer to Tarski's problem are the algebraicity hypothesis in
Baker's theorem and the nondefinability of the quantified domain $A$,
as explained in Remark~\ref{p06:lmod:domain}. Neither arbitrary real
exponential zeros nor inequalities with variable logarithmic
coefficients have been eliminated here. Verdict: a proved restricted
decision boundary, not a closure of ID 6.

\subsection{IDs 7 and 29: two categoricity thresholds}
\label{p06:lmod:id7}

For ID 7 we select the infinitary formulation explicitly posed in
Shelah's Problem 6.9 \cite{p06:lmod:ShelahProblems}: a countable-language
$\psi\in L_{\omega_1,\omega}$ categorical in one
$\lambda\geq\beth_{\omega_1}$ should be categorical in every
$\lambda'\geq\beth_{\omega_1}$. If ordinary countable first-order
categoricity is intended instead, its uncountable transfer is Morley's
theorem, not an open conjecture. Vasey's established universal-sentence
theorem proves an eventual transfer at the larger bound
$\beth_{\beth_{\omega_1}}$ \cite{p06:lmod:VaseyII}; it does not give the
displayed threshold for arbitrary infinitary sentences.

\label{p06:lmod:id29}
For ID 29 the general abstract-elementary-class formulation is:
for every infinite cardinal $\lambda$ there is $\mu(\lambda)$ such
that every AEC $K$ with $\operatorname{LS}(K)\leq\lambda$ which is
categorical in one $\kappa\geq\mu(\lambda)$ is categorical in all
$\kappa'\geq\mu(\lambda)$. Here categoricity includes existence, and
amalgamation, tameness, and absence of maximal models are not added
as assumptions.

There are recent full proof claims that a current status audit must
record. OpenAI's 24 September 2026 primary manuscript states precisely
this uniform assertion as Theorem 1.2 \cite{p06:lmod:OpenAIAEC}.
Esp\'indola's 2019 preprint, revised in 2022, also announces a general
proof \cite{p06:lmod:Espindola}. This section checked the statements and
their scopes; it does not independently certify their long proofs.
Accordingly the status recorded here is ``recent full proof claims,
not independently verified in this report'', rather than either an
unqualified open or an independently proved closed status. An
accompanying OpenAI manuscript claims a CH obstruction to a prescribed
AEC threshold \cite{p06:lmod:OpenAIThreshold}. A counterexample to one
specified bound is compatible with qualitative eventual categoricity.

The attempted transfer between these two IDs is rigorous but limited.
Choose a countable fragment $F$ of $L_{\omega_1,\omega}$ containing
$\psi$ and closed under subformulas and the usual fragment operations.
Then $(\operatorname{Mod}\psi,\preccurlyeq_F)$ is an AEC with
L\"owenheim--Skolem number $\aleph_0$. Coherence is inherited from
agreement on formulas; the union-of-chains property follows by
induction on formulas in $F$, with the existential step being the
Tarski--Vaught argument. A countable closure under chosen witnesses
for the countably many formulas gives downward L\"owenheim--Skolem,
and the same construction over a set $B$ has size at most
$|B|+\aleph_0$. Isomorphism invariance is immediate. Therefore a
uniform AEC eventual theorem implies eventual categoricity for
infinitary sentences with some common threshold. It supplies no
inequality $\mu(\aleph_0)\leq\beth_{\omega_1}$ and hence does not by
itself settle the specified ID 7. Our observational recovery arguments
do not supply the missing amalgamation, type locality, or cardinal
bound. Neither external proof claim is counted as our result.

\subsection{IDs 13 and 32: NIP fields and stable fields}
\label{p06:lmod:id13}

The field formulation of Shelah's NIP conjecture used here is:
every NIP field is finite, real closed, separably closed, or admits a
nontrivial henselian valuation. Johnson's 2026 ICM article states it
as Conjecture 4.2 \cite{p06:lmod:JohnsonICM}. His dp-finite theorem
establishes the relevant classification under finite dp-rank
\cite{p06:lmod:JohnsonDP}. NIP provides finite VC dimension separately for
each formula; it does not assert a uniform finite dp-rank for the
field. Thus interpreting finite-observation bounds as a finite-rank
bound would add an unproved hypothesis. The arbitrary NIP case is
not closed here.

\label{p06:lmod:id32}
The stable fields conjecture says that every infinite field with
stable pure field theory is separably closed. Separably closed cannot
be replaced by algebraically closed in positive characteristic.
Johnson--Tran--Walsberg--Ye prove the assertion for large stable
fields \cite{p06:lmod:StableFields}; Kaplan--Matthews--Walsberg's July
2026 work gives a further proof and applications in the large-field
setting \cite{p06:lmod:BuildingTrees}. The general case is not resolved
by either hypothesis-restricted statement.

\begin{proposition}[A classical conditional reduction]\label{p06:lmod:fieldreduction}
The NIP field conjecture in the displayed formulation implies the
stable fields conjecture.
\end{proposition}
\begin{proof}
Let $K$ be an infinite stable field. Stability implies NIP. Under the
NIP field conjecture, $K$ is real closed, separably closed, or
nontrivially henselian. The real-closed case is impossible: the order
is definable by nonzero squares and has the order property. A field
with a nontrivial henselian valuation is large. The established theorem
that large stable fields are separably closed then handles the
remaining case \cite{p06:lmod:StableFields}.
\end{proof}

This reduction makes the attempted transfer and its obstruction
explicit: it remains to obtain henselianity or largeness without
assuming finite dp-rank or largeness at the outset. It is not a new
proof of either field conjecture.

\subsection{ID 27: quasiminimality of the full complex exponential field}
\label{p06:lmod:id27}

The selected conjecture says that every unary definable subset, with
parameters, of $(\mathbb C;0,1,+,-,\cdot,\exp)$ is countable or
cocountable. Bays--Kirby prove it conditional on exponential
algebraic closedness \cite{p06:lmod:BaysKirby}. Gallinaro--Kirby prove the
corresponding complex-powers reduct quasiminimal in 2024
\cite{p06:lmod:ComplexPowers}; that theorem is not the full exponential
conjecture.

There is a direct restricted transfer from analytic zero sets. Every
one-variable term in the ring-exponential language is entire.
Consequently an atomic equality is either identically true or has a
discrete, hence countable, set of solutions: intersecting a discrete
zero set with successively larger compact disks proves countability.
Finite Boolean combinations of countable or cocountable sets are
again countable or cocountable. Thus the quantifier-free unary
fragment satisfies the desired conclusion. This elementary argument
does not survive arbitrary projections of several-variable zero
sets. The missing step is a quantifier-level theorem such as the
required exponential algebraic closedness, not independence of a
finite list of special values. Verdict: the full conjecture is not
closed by our analytic or transcendence arguments.

\subsection{ID 28: the finite-Morley-rank Cherlin--Zilber conjecture}
\label{p06:lmod:id28}

The formulation used is that every infinite simple group of finite
Morley rank is an algebraic group over an algebraically closed field.
It is distinct from broader conjectures about arbitrary
$\omega$-stable groups. Tent's June 2026 primary exposition still
describes the general problem as open and explains its relation to
nonsplit sharply 2-transitive groups \cite{p06:lmod:Tent}.

We tested the tempting counterexample transfer: use constructions of
nonsplit sharply 2-transitive groups to leave the algebraic category.
The missing hypothesis is finite Morley rank. Their existence as
abstract groups gives no bound on the ranks of all definable sets in
their theories; finite-dimensional numerical observations are not
such a bound. Tent explicitly identifies a suitable finite-rank
construction as a potential route, not a completed construction.
Thus neither an abstract-group example nor an analytic rank analogy
is a counterexample to the stated conjecture. Verdict: no closure.

\subsection{ID 30: which main gap?}
\label{p06:lmod:id30}

For a complete countable first-order theory $T$, the canonical
numerical main-gap statement is the following established theorem:
either $I(T,\kappa)=2^\kappa$ for every uncountable cardinal $\kappa$,
or
\[
 I(T,\aleph_\alpha)<\beth_{\omega_1}(|\alpha|+\aleph_0)
       \qquad(\alpha\geq1).
\]
Here the beth expression denotes the beth iteration with the
specified initial infinite cardinal. Shelah's classification theory
and the complete spectrum analysis of Hart--Hrushovski--Laskowski
are existing results \cite{p06:lmod:HHL}; the displayed formulation is
also recalled in the 2026 survey \cite{p06:lmod:PillayTanovic}.

Consequently this interpretation of ID 30 was already closed before
the present work. A different, explicitly recorded research problem
is Shelah's Problem 4.7 on extending the main gap to uncountable
theories \cite{p06:lmod:ShelahProblems}. The bare name ``main gap
conjecture'' does not specify an AEC or an uncountable-theory variant,
and we do not silently assign one a fabricated uniform cardinal
bound. Our attempted structural transfer would classify models by
trees of observable invariants. Its unproved step is constructing
primary/prime models and decomposition systems over nonforking
configurations; finite-parameter inverse theorems supply neither.
Verdict: the precise countable statement is credited as classical;
no new uncountable or non-elementary main-gap theorem is asserted.

\subsection{ID 31: Kueker's conjecture and a failed stability reduction}
\label{p06:lmod:id31}

For a complete countable theory $T$, the conjecture is: if every
uncountable model of $T$ is $\aleph_0$-saturated, then $T$ is
categorical in some infinite cardinality. Saturation here means
realizing all types over finite parameter sets; the conclusion is
not restricted to uncountable categoricity. Hrushovski proves the
stable case \cite{p06:lmod:HrushovskiKueker}; Tanovi\'c's later paper gives
additional structural results \cite{p06:lmod:TanovicKueker}. No general
resolution was located in the reviewed primary sources.

An attractive attempted reduction is to derive stability from the
saturation premise and invoke Hrushovski. That reduction is false.
Consider the theory of the random graph. Quantifier elimination
shows that a type over a finite parameter set prescribes only
equalities and finitely many adjacency choices for any fixed finite
tuple. The extension axioms realize each consistent prescription,
successively for its distinct coordinates, in every model. Thus
every model is $\aleph_0$-saturated. On the other hand adjacency has
the independence property: for any distinct vertices $a_1,\ldots,a_r$
and any subset $S$, the extension axiom gives a vertex adjacent
exactly to those $a_i$ with $i\in S$. The theory is therefore
unstable. It is $\aleph_0$-categorical, so it satisfies Kueker's
conclusion and is not a counterexample to the conjecture. This
explicit example rules out the proposed proof reduction while
leaving the conjecture itself untouched.

\subsection{ID 33: a 2026 counterexample announcement to stable forking}
\label{p06:lmod:id33}

The stable forking conjecture says that, in a simple theory, every
forking type $\operatorname{tp}(a/B)$ over $C\subseteq B$ contains
an instance $\varphi(x,b)$ of a stable formula, with $b$ from $B$,
which itself forks over $C$. Freitag--Mutchnik's primary preprint
posted 31 August 2026 states the negation as Theorem 1.3
\cite{p06:lmod:StableForking}. Their construction uses a module over a
division ring built from a quantum graph algebra; the paper proves
supersimplicity and then excludes stable witnesses for the specified
forking. This is a recent external counterexample claim, not a
counterexample constructed or independently fully checked here.

The attempted positive transfer would replace nonforking by a
linear-disjointness or rank criterion and then treat its witnesses
as stable. The reported example identifies the failure of that
inference: noncommutative incidence can encode a half-graph even
when the independence calculus has the required simple-theory
axioms. A rank criterion alone does not prove stability of a
witnessing formula. The paper's Question 4.7 separately asks about
the finite-SU-rank supersimple case. We therefore retire the
unrestricted positive target in light of the counterexample
announcement and do not conflate it with that remaining finite-rank
question. Verification in this report covers the formulation and
announced mechanism, not a replacement proof of the source's
construction.

\paragraph{Accounting.}
Theorem~\ref{p06:lmod:boundary} is one proved decision/definability
boundary with a full proof from classical inputs. The conditional
Silver transfer, field reduction, quantifier-free quasiminimality,
and random-graph obstruction are explicitly credited elementary or
classical deductions. None is counted as a new solution of a named
historical conjecture. The eleven-ID ledger preserves the distinction
between a theorem already known, an externally announced resolution,
a restricted result proved here, and a target left unresolved.


\section{Set-theoretic and abstract-logical targets}
\label{p06:lset:section}

This section treats items 1, 3, 8--11, 21, 22 and 45 in the user's fixed
numbering.  The evidence cut-off is 17 October 2026, Korean time.
``A proof is supplied below'' and ``a source announces a proof'' are different
statuses.  None of the historical results or contemporary authors' claims is
counted as a new solution by this investigation.  In particular, old analytic
results proved in ZFC cannot be imported as choice-free lemmas merely by
changing the ambient theory to ZF.

\subsection{Item 1: the Partition Principle versus choice}
\label{p06:lset:item1}

\paragraph{Exact question and current evidence.}
The Partition Principle is the ZF assertion
\[
 \mathrm{PP}:\quad
 (\exists f:X\twoheadrightarrow Y)\Longrightarrow
 (\exists i:Y\hookrightarrow X).
\]
The question is whether $\mathrm{ZF}\vdash\mathrm{PP}\Rightarrow\mathrm{AC}$.
The injection need not be a section of the specified surjection.
Gilson's current revision, dated 10 September 2026, explicitly withdraws its
claimed countermodel; its abstract and correction notice identify the stronger
splitting principle, a misuse of forcing-name symmetry and a failed limit
truncation argument \cite{p06:lset:Gilson}.  Thus the earlier search result claiming
a solution is superseded.  The ordinary implication is not resolved here.

\paragraph{Attempt: force local inverses, then preserve the choice failure.}
Here is a self-contained explanation of the splitting obstruction acknowledged
in that correction.  It does not require analysis of the iteration.

\begin{proposition}[Localized splitting already well-orders the seed]
\label{p06:lset:splitting}
In ZF let $A$ be a set, $S=A^\omega$ and $T=\mathcal P(S)$.
Suppose every surjection between two subsets of $T$ has a right inverse.
Then $A$ is well-orderable.
\end{proposition}
\begin{proof}
The empty case is immediate.  For nonempty $B\subseteq A$ let
\[
 C(B)=\{(b,b,b,\ldots):b\in B\}\subseteq S.
\]
For $a\in B$ put
\[
 E(a,B)=\{(a,b,a,a,\ldots):b\in B\}\subseteq S.
\]
The nonempty set $E(a,B)$ determines $a$ from the zeroth coordinate and
$B$ from the first coordinates.  Consequently
\[
 X=\{C(B):\varnothing\ne B\subseteq A\},\qquad
 Y=\{E(a,B):a\in B\subseteq A\}
\]
are subsets of $T$, and $E(a,B)\mapsto C(B)$ is a well-defined surjection
$Y\to X$.  A section determines a function $c$ on all nonempty subsets of
$A$, with $c(B)\in B$.  Transfinite recursion now repeatedly selects
$c(A\setminus\{a_\xi:\xi<\eta\})$ whenever the remainder is nonempty.
This process must terminate before the Hartogs ordinal of $A$, since otherwise
it gives an injection of that ordinal into $A$.  The resulting enumeration
well-orders $A$.  No choices are made in this recursion beyond the given $c$.
\end{proof}

For the genuine PP the same construction gives only an arbitrary injection
$X\to Y$; its value at $C(B)$ can encode an element of an entirely different
subset of $A$.  Recovering a fiber-respecting map is exactly the unjustified
step.  The proposition rules out the strengthened seed-preserving construction,
not a model of ordinary PP without choice.  This is an audited obstruction,
not a new proof of either answer to item 1.

\subsection{Item 3: what remains of the continuum problem}
\label{p06:lset:item3}

\paragraph{Exact question and status.}
The sentence CH is $2^{\aleph_0}=\aleph_1$ in ZFC.  Its independence,
relative to the consistency of ZFC, is a theorem of G\"odel and Cohen;
Cohen's 1963 paper is \cite{p06:lset:Cohen}.  A later machine-checked treatment
is \cite{p06:lset:HanDoorn}, its main independence theorem.  The question of a
preferred additional axiom deciding CH is not the question whether ZFC decides
CH.  Woodin's 2026 discussion still separates the proposed Ultimate-$L$ axiom
from the conjectures needed to justify its compatibility with large cardinals
\cite{p06:lset:Woodin2026}.

\paragraph{Attempt: identify the continuum through old observation theorems.}
The obstruction is not insufficient numerical precision.  The standard forcing
calculation shows why no ZFC theorem about the previous analytic constructions
can settle CH.

\begin{proposition}[The Cohen calculation, recalled]
\label{p06:lset:cohen}
If $M\models\mathrm{ZFC}$ and $M$ has an infinite cardinal $\kappa$ with
$\kappa^{\aleph_0}=\kappa$, forcing over $M$ with
$\operatorname{Fn}(\kappa\times\omega,2,{<}\omega)$ preserves cardinals
and forces $2^{\aleph_0}=\kappa$.
\end{proposition}
\begin{proof}
An uncountable family of finite conditions has an uncountable delta-system
of domains.  Thinning to equal values on the root gives compatible conditions,
so the forcing is ccc and preserves cardinals.  For each $\alpha<\kappa$
the generic union gives a real $c_\alpha$.  For different $\alpha,\beta$,
conditions making their values differ at a fresh integer are dense; hence the
reals are distinct.  Conversely, a nice name for a real consists of countably
many countable antichains, since the forcing is ccc.  There are at most
$\kappa^{\aleph_0}=\kappa$ such names.  These two estimates give the result.
The familiar ccc preservation and nice-name lemmas are used here as standard
forcing lemmas, not reproved as part of this calculation.
\end{proof}

Starting from a model of GCH, take $\kappa=\aleph_2$; its extension satisfies
$\neg\mathrm{CH}$ while the constructible model gives the positive relative
consistency direction.  Every ZFC theorem proved in the earlier papers remains
a theorem in both settings.  More formally, adding already provable sentences
to a theory does not change its deductive consequences: replace each use of
an added axiom by its proof.  A new nonconservative set-theoretic axiom, with
an independent justification, is the missing ingredient.  There is no new
resolution of the post-Cohen axiom-selection problem here.

\subsection{Item 8: distinguish two PCF questions}
\label{p06:lset:item8}

\paragraph{Formalization and evidence.}
For a progressive set $A$ of infinite regular cardinals
($|A|<\min A$), put
\[
 \operatorname{pcf}(A)=
 \{\operatorname{cf}(\textstyle\prod A/U):U\text{ an ultrafilter on }A\}.
\]
One strong PCF conjecture asserts $|\operatorname{pcf}(A)|=|A|$ for every
such $A$.  The numerical assertion
\[
 \aleph_\omega\text{ strong limit}\quad\Longrightarrow\quad
 2^{\aleph_\omega}<\aleph_{\omega_1}
 \tag{PCF$_\omega$}\label{p06:lset:pcfomega}
\]
is a different target.  Shelah's established bound has $\omega_4$ in place
of $\omega_1$ \cite[\S6]{p06:lset:ShelahPCF}.
Gitik's author-hosted manuscript, revised 3 March 2026, reports a countable
$A$ with $\operatorname{otp}(\operatorname{pcf}(A))=\omega_1+1$;
the critical spectrum calculation is Lemma 7.12, with large-cardinal settings
in \S1 \cite{p06:lset:GitikPCF}.  This is a reported relative counterexample to
the unrestricted cardinality conjecture, not a proof of its negation in every
model.  Its \S8, Question 3, expressly leaves moving the construction down to
$\aleph_\omega$ unresolved.  We have checked this scope, not independently
certified its extender-forcing construction or established a publication status.

\paragraph{Attempt: reconstruct the spectrum from finite restrictions.}
The following elementary calculation exposes a failure of the finite-observation
analogy.  It does not improve Shelah's bound.

\begin{proposition}[Finite decomposition does not extend to countable decomposition]
\label{p06:lset:pcf-finite}
In ZFC, if $A=A_0\cup\cdots\cup A_{m-1}$, then
\[
 \operatorname{pcf}(A)=\bigcup_{i<m}\operatorname{pcf}(A_i),
\]
where empty pieces are omitted.  In contrast, if $\langle\kappa_n:n<\omega\rangle$
is strictly increasing and consists of uncountable regular cardinals, and
$\lambda=\sup_n\kappa_n$, every free ultrafilter $U$ on $\omega$ satisfies
\[
 \operatorname{cf}(\textstyle\prod_n\kappa_n/U)>\lambda.
\]
In particular the PCF spectrum is not the union of the spectra of its singleton
pieces.
\end{proposition}
\begin{proof}
Refine a finite cover to a disjoint partition.  An ultrafilter contains exactly
one piece; restriction to that piece identifies the corresponding reduced
product.  Conversely an ultrafilter on a piece extends to an ultrafilter on
$A$ concentrated there.  Monotonicity under inclusion then gives the stated
formula for the original cover.

Fix a family $F\subseteq\prod_n\kappa_n$ of cardinality $\mu<\lambda$.
For all sufficiently large $n$ we have $\mu<\kappa_n$.  Regularity gives
\[
 g(n)=\sup\{f(n)+1:f\in F\}<\kappa_n
\]
on those coordinates; give $g$ arbitrary legal values on the remaining ones.
Then $[g]_U>[f]_U$ for every $f\in F$.  Thus no family of size less than
$\lambda$ is cofinal in the reduced product.  The quotient is a linear order
without a largest element.  The cofinality of such an order is a regular
cardinal: a singular least cofinal size would, by taking upper bounds of
successively larger proper subfamilies, give a smaller cofinal family.
Since $\lambda$ has cofinality $\omega$, the quotient cofinality cannot equal
$\lambda$, and is therefore larger.  Finally the only ultrafilter on
$\{\kappa_n\}$ is principal, giving spectrum $\{\kappa_n\}$.
\end{proof}

The missing step in the proposed reconstruction is control of free ultrafilters
and true cofinalities, invisible to all finite coordinate restrictions.
For \eqref{p06:lset:pcfomega}, a forcing counterexample would additionally need
the difficult collapses at successors of singular cardinals while preserving
the required scales.  Neither step is supplied here.

\subsection[Item 9: Woodin's Omega conjecture]{Item 9: Woodin's $\Omega$ conjecture}
\label{p06:lset:item9}

\paragraph{Exact formulation.}
Assume a proper class of Woodin cardinals.  For a countable theory $T$,
$T\models_\Omega\varphi$ means that every rank segment satisfying $T$ in
every set-forcing extension satisfies $\varphi$.  One standard proof-relation
convention is Larson's Definition 6.1: $T\vdash_\Omega\varphi$ if some
universally Baire set $A\subseteq\omega^\omega$ has the property that every
countable transitive $A$-closed model $M$ of ZFC containing $T$ satisfies
``$T\models_\Omega\varphi$''.  Concretely, $A$-closed means that for each
poset $P\in M$ there is a countable family $D$ of open dense subsets of $P$,
including all such sets belonging to $M$, such that $A\cap M[G]\in M[G]$
whenever a filter $G$ meets every member of $D$.  The members of $D$ need
not belong to $M$.
The conjecture is, for all sentences $\varphi$,
\[
 \varnothing\models_\Omega\varphi\quad\Longrightarrow\quad
 \varnothing\vdash_\Omega\varphi.
\]
Soundness is already a theorem.  The stronger definability conjectures sometimes
also called an $\Omega$ conjecture are not being identified with this statement.

\paragraph{Source and actual reduction.}
Larson \cite[Definition 6.1, Theorems 6.4 and 9.2]{p06:lset:LarsonOmega}
gives soundness and a conditional completeness route: suitable countable
collapses of rank segments with universally Baire iteration strategies yield
the required forcing countermodels to unprovability.  No proof of the
unrestricted completeness direction was verified in the current search.
The unresolved step in this route is obtaining these strategies for the
relevant collapses of $V$, rather than only for canonical iterable inner
models.  Finite-cover uniqueness from the preceding analytic work does not
produce an iteration strategy or certify all its countable iterates.

\paragraph{Attempt: replace iteration strategies by compactness.}
This fails at an exact, elementary boundary.  If a consistent theory $T$
extends ZF, adjoin constants $c_n$ and sentences saying that all $c_n$ are
ordinals and $c_{n+1}\in c_n$ for every $n<\omega$.  Every finite subset
has a model: interpret its finitely many constants as a sufficiently long
descending list of finite ordinals in a model of $T$.  First-order compactness
therefore gives a model of the entire enlarged theory.  Its membership relation
is externally ill-founded, so it cannot be a transitive model.  Thus the
Henkin construction alone cannot provide the transitive $A$-closed objects
required in $\Omega$ logic.  This proves an obstruction to this attempted
method, not a refutation of $\Omega$ completeness.  The exact missing
strategy hypothesis remains unsolved by our methods.

\subsection{Item 10: HOD conjecture and the stationarity gap}
\label{p06:lset:item10}

\paragraph{Formalization and current status.}
HOD is the class of hereditarily ordinal-definable sets.  A regular uncountable
$\theta$ is $\omega$-strongly measurable in HOD if some infinite cardinal $\delta<\theta$
satisfies $(2^\delta)^{\mathrm{HOD}}<\theta$ and HOD has no partition of
$S^\theta_\omega$ into $\delta$ sets stationary in $V$.
The HOD Hypothesis asserts a proper class of regular cardinals not
$\omega$-strongly measurable in HOD.  With an extendible cardinal, the HOD
dichotomy connects it with correct singular-cardinal successors in HOD.
The HOD \emph{conjecture}, in the formulation of
\cite[\S1.2, Theorems 1.1--1.2 and B]{p06:lset:ABL}, is a provability claim
over ZFC plus an extendible cardinal, not merely the sentence $V=\mathrm{HOD}$.
That paper proves a conditional negative result from the joint consistency
of an exacting cardinal above an extendible cardinal.  It does not establish
this joint consistency from ZFC alone.  Woodin's published 2026 lectures
continue to formulate the associated program as conjectural
\cite{p06:lset:Woodin2026}.  Neither an unconditional proof nor an unconditional
consistency refutation is obtained here.

\paragraph{Attempt: use HOD's choice to manufacture the required partitions.}
HOD satisfies AC, but the sets it regards as stationary need not be stationary
in the outer universe.  The following standard club-shooting argument shows
that even agreement about $\omega_1$ cannot repair this inference.

\begin{proposition}[Stationarity is not upward absolute]
\label{p06:lset:stationarity}
Let $M\models\mathrm{ZFC}$ and let $S\subseteq\omega_1^M$ be stationary
and costationary in $M$.  There is an $\omega$-distributive forcing extension
$M[G]$ preserving $\omega_1$ in which $S$ contains a club.  Hence the
ground-model stationary set $\omega_1\setminus S$ becomes nonstationary.
\end{proposition}
\begin{proof}
Use the poset of countable closed bounded subsets of $S$, ordered by
end-extension.  It adds an unbounded closed union, since extending the maximum
beyond any prescribed countable ordinal is dense.  To prove
$\omega$-distributivity, fix a condition $p$ and open dense sets $D_n$.
Choose a countable elementary submodel $N\prec H_\chi$ containing these
parameters with $\delta=N\cap\omega_1\in S$; stationarity supplies $N$.
Choose in the ambient universe an increasing sequence cofinal in $\delta$.
Recursively choose $p_n\in N$ below $p$ and earlier conditions, in $D_n$,
with maxima exceeding successive terms of that sequence.  Elementarity and
unboundedness of $S$ make each choice possible inside $N$.  The set
\[
 q=\bigcup_{n<\omega}p_n\ \cup\ \{\delta\}
\]
is a countable closed subset of $S$: every limit point below $\delta$ already
occurs in some end-extended closed condition, and the only new limit point
is $\delta$.  It is a lower bound lying in every $D_n$.  This proves
$\omega$-distributivity and hence that no countable sequence of ordinals is
added; in particular $\omega_1$ is preserved.  The generic club is disjoint
from the old complement of $S$.
\end{proof}

This classical construction is included to isolate the obstruction, not to
claim that $M$ in the proposition is HOD of its extension.  Our attempted
HOD argument needs a theorem transferring stationarity of its proposed
partition pieces to $V$.  HOD's internal well-order and uniqueness of a
definition supply no such theorem.  The HOD dichotomy and modern exacting
cardinal results show why this missing ingredient is substantive.

\subsection{Item 11: consistency of Reinhardt cardinals}
\label{p06:lset:item11}

\paragraph{Framework and source correction.}
The intended assertion is a nonidentity elementary class embedding
$j:V\to V$, with critical point $\kappa$, in a choice-free class theory,
or in the usual $\mathrm{ZF}(j)$ framework with the stated class-replacement
instances.  It is not the assertion that a set function has domain $V$,
and it is not restricted to parameter-definable embeddings.  These variants
must not be interchanged.  Kunen's inconsistency uses choice; see
\cite{p06:lset:HKP} for the standard stationary-set argument and precise variants.
McCallum's 2026 inconsistency claim is withdrawn in revision 3 of
23 August 2026, with an identified problem concerning continuity of Skolem
functions \cite{p06:lset:McCallum}.  It therefore supplies no inconsistency
theorem.  The unrestricted choice-free consistency problem remains unresolved
by the verified sources and by this investigation.

\paragraph{Attempt: localize Kunen's argument to exactly what choice supplied.}

\begin{proposition}[Conditional stationary-partition obstruction]
\label{p06:lset:kunen}
Work in the class-embedding framework just specified.  Suppose $j:V\to V$
has critical point $\kappa$.  If $\theta>\kappa$ is a regular uncountable
cardinal fixed by $j$, and $S^\theta_\omega$ has a partition
$\langle S_\alpha:\alpha<\kappa\rangle$ into stationary sets, then a
contradiction follows.
\end{proposition}
\begin{proof}
Write $\langle T_\alpha:\alpha<j(\kappa)\rangle=
j(\langle S_\alpha:\alpha<\kappa\rangle)$.
Since $j(\theta)=\theta$ and $j(\omega)=\omega$, this is a partition of
the same $S^\theta_\omega$ into stationary sets.  The set
\[
 C=\{\beta<\theta:j``\beta\subseteq\beta\}
\]
is club.  It is closed by taking suprema; it is unbounded because starting
above any $\eta<\theta$ we can repeatedly close under $j$ and take the
supremum of the resulting countable sequence below the regular $\theta$.
The required sequences are sets by the replacement assumptions on $j$.
Pick $\beta\in C\cap T_\kappa$.  It has cofinality $\omega$.  For an
increasing cofinal sequence $\langle\beta_n:n<\omega\rangle$ in $\beta$,
elementarity and $j``\beta\subseteq\beta$ give
$j(\beta)=\sup_n j(\beta_n)\leq\beta$, while embeddings of ordinals
never decrease them.  Thus $j(\beta)=\beta$.
There is $\alpha<\kappa$ with $\beta\in S_\alpha$, so
$\beta=j(\beta)\in j(S_\alpha)=T_\alpha$, since $j(\alpha)=\alpha$.
This contradicts the disjointness of $T_\alpha$ and $T_\kappa$.
\end{proof}

This is the classical obstruction with the hypotheses exposed.  In ZFC one
gets a fixed regular successor above the supremum of the critical sequence
and the required stationary partition.  In the choice-free setting these
steps are unavailable in this generality.  A construction of local rank
embeddings, or the consistency results at $V_{\lambda+2}$, is not a
construction of $j:V\to V$; conversely our conditional contradiction does
not remove its partition hypothesis.  The missing step is exactly such a
choice-free global construction or an argument dispensing with that hypothesis.

\subsection{Item 21: Roitman's problem and the September 2026 claim}
\label{p06:lset:item21}

\paragraph{Formalization and version-sensitive status.}
Here $\mathfrak d$ is the least size of an eventually dominating subset of
$\omega^\omega$, and $\mathfrak a$ the least size of an infinite maximal
almost-disjoint family in $[\omega]^\omega$.  The question is whether
\[
 \mathrm{ZFC}\vdash\mathfrak d=\aleph_1\Longrightarrow\mathfrak a=\aleph_1.
\]
The purported counterexample of Mildenberger, arXiv:1404.7343, is withdrawn;
the arXiv record identifies its Lemma 3.2 as flawed \cite{p06:lset:Mildenberger}.
Pelayo-G\'omez's 30 September 2026 preprint now claims the positive theorem
and preservation under outer models retaining domination
\cite[Theorems 1.1 and 1.4]{p06:lset:Pelayo}.  We read its full core construction,
Lemmas 3.1--3.3, and maximality argument; two independent checks in this
investigation found no local gap.  Publication, priority and wider independent
validation were not established.  The correct evidence label is therefore
\emph{recent author proof claim, core checked here}, not an unqualified
``still open'', and not our own solution.

\paragraph{Our attempted reduction and its interface with that claim.}
Normalize a size-$\omega_1$ dominating family by closing under finite
modifications, obtaining pointwise domination.  For a bounding function $d$
and an enumeration $c=\langle c_k:k<\omega\rangle$ of earlier almost-disjoint
sets, use the two observable objects
\[
 K_d=\{x\in[\omega]^\omega:e_x\leq d\},\qquad
 G(c,g)=\omega\setminus\bigcup_k(c_k\setminus[0,g(k))).
\]
If $x$ is almost disjoint from every $c_k$, the trace
$t_x(k)=1+\max(x\cap c_k)$, with value zero for the empty intersection,
satisfies $t_x\leq g\Rightarrow x\subseteq G(c,g)$.
Thus domination controls both enumeration and finite intersections, but it
does not alone build a single new almost-disjoint set hitting every member of
the corresponding compact class.  This simultaneous capture is the precise
step missing from the naive recursion.

The cited preprint's protected-interval fusion is designed to supply that
step.  Our audit checked that its successor stage leaves infinitely many
reserved intervals for each active index, and that countable limit stages
take infinite pseudointersections of the decreasing reservations.  The
maximality proof then uses the trace implication above.  These checks do not
replace the full proof by an analogy with inverse observation.  Our outcome is
a technical audit of an external claimed resolution, with no additional
theorem or claimed closure count.

\subsection{Item 22: consistency strength of PFA}
\label{p06:lset:item22}

\paragraph{Formalization and established bounds.}
PFA asserts that for every proper poset and every family of $\omega_1$
dense subsets there is a filter meeting all of them.  The target is the
equiconsistency of ZFC+PFA with ZFC plus a supercompact cardinal, not the
false literal claim that PFA makes $\omega_2$ supercompact in $V$.
The supercompact upper bound is classical.  A verified lower bound is Trang's
theorem giving a transitive inner model containing all reals and ordinals and
satisfying $\mathrm{AD}_{\mathbb R}+``\Theta\text{ is regular}''$
\cite[Theorems 0.2 and 0.4]{p06:lset:Trang}.  This is an established bound,
not a claim that this search determined the latest optimal lower bound.
Viale--Wei\ss\ prove that a \emph{proper standard iteration} from an
inaccessible $\kappa$ forcing PFA with $\kappa=\omega_2$ requires that
$\kappa$ was supercompact \cite[Corollary 6.11]{p06:lset:VialeWeiss}.
Their qualification cannot be deleted to obtain a general lower bound.

\paragraph{Attempt: descend branches to recover a large-cardinal ground.}
The core descent mechanism can be stated without any ambiguity about forcing
presentations.

\begin{lemma}[Approximation transfers an already supplied branch]
\label{p06:lset:approximation}
Let $M\subseteq W$ be transitive models of ZFC with the same ordinals and
the $\kappa$-approximation property: if $b\in W$ is a subset of a set of
$M$ and $b\cap x\in M$ for every $x\in M$ of $M$-size less than
$\kappa$, then $b\in M$.  Let $\lambda\in M$ be an ordinal and
$\mathcal T=\langle T_x:x\in(P_\kappa\lambda)^M\rangle\in M$,
where $T_x\subseteq\mathcal P(x)$.  If $b\subseteq\lambda$ in $W$
satisfies $b\cap x\in T_x$ for every such $x$, then $b\in M$ and is
a branch for the same system in $M$.
\end{lemma}
\begin{proof}
Every member of $T_x\in M$ belongs to $M$ by transitivity.  For any
$y\in M$ of $M$-size less than $\kappa$, set $x=y\cap\lambda$.
Then $b\cap y=b\cap x\in T_x\subseteq M$.  Approximation yields
$b\in M$.  Membership of each restriction in $T_x$ is absolute between
the transitive models, proving the second assertion.
\end{proof}

This explains the attraction of guessing-model and approximation arguments,
but does not construct an appropriate ground $M$ from an arbitrary model of
PFA.  Nor does PFA automatically provide branches for every system whose
levels are indexed only by the old small subsets: new small subsets in $W$
require additional analysis.  The missing recovery theorem is an inner model
or ground with the required extender/approximation structure, not an estimate
on finitely many measurements.  The general equiconsistency question is not
closed here.

\subsection[Item 45: compactness, Beth and Delta interpolation]{Item 45: compactness, Beth and $\Delta$ interpolation}
\label{p06:lset:item45}

\paragraph{Disambiguation.}
Taken literally, existence of a compact logic having Beth definability and
$\Delta$ interpolation has the elementary positive answer $L_{\omega\omega}$.
The nontrivial problem is a regular abstract logic extending first-order logic
that is compact and has ordinary Beth definability and $\Delta$ interpolation
\emph{but fails Craig interpolation}.  Countable compactness, full compactness,
weak Beth and projective Beth are distinct variants.
Here $\mathrm{PC}(L)$ denotes classes of reducts of models of an $L$ sentence;
$\Delta(L)$ comprises classes for which both the class and its complement are
$\mathrm{PC}(L)$.  $\Delta$ interpolation means $\Delta(L)=L$, up to
definability; Craig interpolation separates arbitrary disjoint projective
classes arising from an implication.  Ordinary Beth concerns an implicitly
unique relation, whereas projective Beth permits additional hidden relations.

\paragraph{Verified neighboring results.}
Shelah's 1985 construction is compact and has Beth, but explicitly fails even
$\Delta$ interpolation \cite[Conclusion 6.11(4)]{p06:lset:ShelahBeth}.
Its Claim 6.12 concerns an enlarged closure; it does not assert failure of
Craig interpolation for that enlarged logic.  The current detailed account
\cite[Proposition 14 and the following discussion]{p06:lset:Vaananen2025}
also warns that the $\Delta$ closure of the cofinality logic need not have
ordinary Beth.  Thus neither of these frequently conflated constructions
settles the conjunction requested here.  No verified resolution of that
precise conjunction was found.

\paragraph{Attempt: alternate Beth and $\Delta$ closure.}
Let $\varphi$ and $\psi$ have common vocabulary $\tau$ and satisfy
$\varphi\models\psi$.  Define
\[
 K=\operatorname{red}_\tau\operatorname{Mod}(\varphi),\qquad
 J=\operatorname{red}_\tau\operatorname{Mod}(\neg\psi).
\]
These are disjoint $\mathrm{PC}(L)$ classes.  If one could show that the
complement of $K$ were also $\mathrm{PC}(L)$, $\Delta$ interpolation would
define $K$ by a $\tau$ sentence $\theta$.  Then every model of $\varphi$
satisfies $\theta$, and no model of $\theta$ satisfies $\neg\psi$, giving
Craig interpolation.  This two-line reduction is valid; its premise is the
missing assertion.  Disjointness $K\cap J=\varnothing$ does not imply that
$K\cup J$ is the class of all $\tau$ structures.  Ordinary Beth does not
turn existence of an arbitrary expansion into an implicitly unique expansion.

Alternating the two closures therefore has three unproved requirements:
compactness at the closure limit, persistence of ordinary Beth after adding
projective complementary definitions, and preservation of a concrete failure
of Craig interpolation.  The cited constructions show that each property
must actually be checked, rather than inherited by name.  Under the additional
countable downward L\"owenheim--Skolem property and finite occurrence
(each sentence uses a finite vocabulary), the classical Lindstr\"om theorem
already forces a compact regular extension back to first-order logic; that
restricted route cannot produce a counterexample.  We obtain a precise
failure diagnosis and the literal first-order answer, but no new solution of
the intended separation problem.

\paragraph{Outcome of the nine-target attempt.}
The outputs are explicit proof barriers and reductions, a corrected distinction
between PCF variants, and audits of recent or withdrawn claims.  CH independence
is historical; the current Roitman proof claim belongs to its named author.
No unresolved external problem is recorded as newly solved by this section.


\section{Descriptive and model-theoretic targets: ten precise tests}
\label{p06:ldes:section}

This section addresses targets 5, 14, 20, 23, 39, 41, 46, 49, 50,
and 52.  The literature audit is dated 17 October 2026.  None of the
general open conjectures in this group is claimed to be settled here.
We give explicit reductions, counterexamples to proposed extensions, and
proofs identifying the additional hypotheses that an observation argument
would have to supply.  The elementary propositions below are method audits;
no priority claim is made for their mathematical content.  In particular,
items 41 and 46 are families of questions, not uniquely specified conjectures.

\subsection{A common obstruction: recovering points does not classify classes}
\label{p06:ldes:coding}

Our earlier observation constructions recover individual points by Borel
adaptive choices.  The following distinction is necessary before applying
them to an equivalence relation.  A complete invariant for an equivalence
relation $E$ on $X$ is a map $c$ such that
\[
 c(x)=c(y)\quad\Longleftrightarrow\quad xEy.
\]
The existence of a Borel complete invariant in a standard Borel space is
the definition of smoothness.

\begin{proposition}[Point information and invariant information]
\label{p06:ldes:point-class}
Let $X$ be compact metrizable.  A finite adaptive scheme whose query
selection maps into $C(X)$ are Borel, and whose decoder is Borel, has a
Borel output map.  Such a scheme can output a complete invariant for $E$
only if $E$ is smooth.  Continuous injective point coding does not imply
this condition, even for a hyperfinite countable Borel equivalence relation.
\end{proposition}
\begin{proof}
Equip $C(X)$ with the uniform norm.  Evaluation $(f,x)\mapsto f(x)$ is
continuous.  Induction on the number of queries therefore shows that every
transcript coordinate is Borel; composing with the decoder proves the
first assertion.  The second is the definition just given.

For the last assertion take $X=2^{\mathbb N}$ and eventual equality $E_0$.
The map
\[
 x\longmapsto \sum_{j\geq0}2x(j)3^{-j-1}
\]
is continuous and injective.  Nevertheless $E_0$ has no Borel complete
invariant.  Here is a direct proof.  If $c:X\to Y$ were such an invariant,
choose a countable family $(B_j)$ of Borel sets separating the points of
the standard Borel space $Y$.  Each $c^{-1}(B_j)$ is unchanged by changing
finitely many bits.  Under fair Bernoulli product measure it is thus a
tail event and has measure zero or one.  For completeness, invariance
under all changes of the first $n$ bits makes it independent of the
first $n$ coordinates; passing to the sigma-algebra they generate makes
it independent of itself, giving this zero--one assertion.  Intersecting
the measure-one choices of $c^{-1}(B_j)$ or its complement gives a
measure-one set on which all the separating tests agree.  On that set
$c$ is constant.  Its fiber would be one countable $E_0$-class, which has
measure zero, a contradiction.  Finally, agreement outside the first
$n$ coordinates defines a finite Borel equivalence relation, and these
relations increase to $E_0$.
\end{proof}

The nonsmoothness of $E_0$ is classical.  Notice also that it does not
contradict a perfect-set dichotomy: partition $\mathbb N$ into infinite
sets $I_j$ and repeat the $j$th bit on $I_j$.  This continuous injection
has compact, hence perfect, image and sends distinct sequences to
different $E_0$-classes.  Thus point recovery, smooth classification,
and constructing a perfect antichain are three different tasks.

\subsection{Target 5: Martin's conjecture and countable adaptive branching}
\label{p06:ldes:martin}

Write $x\leq_T y$ for Turing reducibility and let a cone be
$\{x:a\leq_T x\}$.  One standard formulation assumes
$\mathsf{ZF}+\mathsf{AD}+\mathsf{DC}_{\mathbb R}$ and considers all
Turing-invariant functions $f:2^{\mathbb N}\to2^{\mathbb N}$.
Part I asserts that either the degree of $f(x)$ is constant on a cone,
or $x\leq_T f(x)$ on a cone.  Part II asserts the prewellordering,
modulo cones, of functions above the identity, with pointwise Turing
jump increasing rank by one.  The Borel versions over choice-based
set theory are distinct formulations.

Lutz and Siskind prove Part I for order-preserving functions
\cite[Theorem 4.8]{p06:ldes:LutzSiskind}, and for Borel order-preserving
functions without the extra determinacy assumption
\cite[Theorem 4.13]{p06:ldes:LutzSiskind}.  Yu's October 2026 preprint
studies weaker regularity hypotheses and consistency strength for that
order-preserving case \cite[Theorems 1.2 and 1.4]{p06:ldes:Yu}.
These results do not establish the conjecture for arbitrary
Turing-invariant functions.

\begin{proposition}[Countable cone patching]
\label{p06:ldes:cone-patching}
In ZFC, a countable Borel partition $(A_n)$ of $2^{\mathbb N}$ into
Turing-invariant sets has a member containing a cone.  Consequently a
Turing-invariant, Borel, countably valued branch selector is constant
on a cone.  Patching countably many functions on such a partition cannot
create a cone behavior different from that of every branch.
\end{proposition}
\begin{proof}
The Borel cone theorem, a consequence of Borel determinacy, says that
each $A_n$ or its complement contains a cone.  If no $A_n$ contains
a cone, choose bases $a_n$ for cones in their complements.  The real
$a=\bigoplus_n a_n$ computes every $a_n$, so the cone above $a$ misses
every $A_n$, contrary to the partition property.  On the cone contained
in $A_n$, a patched function equals its $n$th branch.  If that branch
has a specified property on another cone, intersect the two cones by
taking the join of their bases.
\end{proof}

The cone theorem is an external input; see
\cite[Theorem 1.1 and Section 2]{p06:ldes:LutzSiskind}.  The obstruction to
the proposed transfer is exact: general real-valued transcripts are not
countable partitions, and an arbitrary Borel selector need not be
Turing-invariant.  Neither Borel measurability nor exact numerical
recovery supplies uniform indices for Turing reductions.  The attempt
therefore excludes a simple patching strategy and leaves both full
parts unresolved.

\subsection{Target 14: the Borel Weiss problem}
\label{p06:ldes:weiss}

The formulation used here asks whether every Borel action of a
countable amenable group on a standard Borel space has hyperfinite
orbit relation.  Hyperfinite means an increasing union of finite
Borel equivalence relations on the whole space.  Removing a null set
is not allowed.  Naryshkin and Vaccaro's 2026 result proves the treeable
case \cite[Corollary 3]{p06:ldes:NaryshkinVaccaro}; their Theorem 2 uses the
stronger property of hyper-$u$-amenability.  Their introduction explicitly
retains the general Weiss question.

The finite-subgroup mechanism can be checked completely.  If
$\Gamma=\bigcup_n\Gamma_n$ with finite subgroups
$\Gamma_n\subseteq\Gamma_{n+1}$, set
\[
 xF_ny\quad\Longleftrightarrow\quad
 (\exists g\in\Gamma_n)\;g\cdot x=y.
\]
Each $F_n$ is Borel, is an equivalence relation, and has classes of size
at most $|\Gamma_n|$.  Their increasing union is the orbit relation.
This is the familiar locally finite special case, not a new solution.

Replacing subgroups by F\o lner sets destroys the argument at a precise
point: an orbit fragment indexed by a finite set need not be an
equivalence class.  On the translation action of $\mathbb Z$ on itself,
the symmetric one-step edges $k\leftrightarrow k+1$ already generate
one infinite class.  Taking transitive closure can therefore destroy
finiteness before any limiting argument begins.  The missing input is
a Borel construction of coherent finite partitions; selecting F\o lner
sets separately at each point does not provide one.  The treeable
theorem is an applicable external partial result, but no arbitrary
amenable action is shown here to satisfy its extra hypothesis.

\subsection{Target 20: the Borel Ruziewicz problem}
\label{p06:ldes:ruziewicz}

Fix $n\geq2$ and normalized spherical measure $\sigma$ on $S^n$.
The Borel problem asks whether every finitely additive isometry-invariant
probability on the Borel subsets of $S^n$ equals $\sigma$.
This differs from the classical uniqueness theorem on the completed
Lebesgue sigma-algebra.  Conley and Marks explicitly distinguish the
two and prove a related exotic-measure example for a different group
action \cite[Introduction and Theorem 1.7]{p06:ldes:ConleyMarks}.
That example is not a sphere counterexample.

\begin{proposition}[The null-charge obstruction]
\label{p06:ldes:null-charge}
Assume the classical completed-measure Ruziewicz theorem in dimension
$n\geq2$.  If a finitely additive invariant Borel probability $\nu$
vanishes on every $\sigma$-null Borel set, then $\nu=\sigma$.
Thus any counterexample to the Borel problem must charge a Borel null set.
\end{proposition}
\begin{proof}
Every set $A$ in the completion differs from a Borel set $B$ by a
subset of a Borel null set.  Define $\widetilde\nu(A)=\nu(B)$.
Two choices of $B$ differ inside a Borel null set, so finite additivity
and positivity of $\nu$ make this well defined.  For disjoint completed
sets, their Borel representatives have null intersection; removing that
intersection verifies finite additivity of $\widetilde\nu$.
Isometries preserve Borel sets and $\sigma$-null sets, so invariance
also extends.  The completed-measure uniqueness theorem now gives
$\widetilde\nu=\sigma$, and restriction proves the claim.
\end{proof}

This is the standard reduction underlying the distinction, not a new
case of the conjecture.  It locates the failure of a proposed spectral
transfer: a functional charging null sets does not descend to
$L^\infty(\sigma)$ modulo almost-everywhere equality.  A spectral-gap
or $L^2$ observation inequality cannot be applied to it without first
proving null-set annihilation.  That step remains unproved here.

\subsection{Target 23: the topological Vaught conjecture}
\label{p06:ldes:vaught}

For a continuous action of a Polish group $G$ on a Polish space $X$,
the chosen statement is: either there are countably many orbits, or
there is a perfect set of pairwise orbit-inequivalent points.  This is
stronger in formulation than merely excluding intermediate quotient
cardinalities.  Becker and Kechris relate the topological and Borel-action
versions \cite{p06:ldes:BeckerKechris}; Melleray's 2026 exposition explicitly
records the general problem as open immediately after Silver's theorem
\cite[Theorem 15.10 and following discussion]{p06:ldes:Melleray}.

The direct reduction attempt is to make the orbit relation coanalytic
and apply Silver's dichotomy.  It works if the orbit relation is Borel:
it is then coanalytic, and Silver yields exactly the required alternative.
But the assertion that each individual orbit is Borel does not say that
the relation is Borel in $X\times X$.  General Polish orbit relations
are analytic and can fail to be Borel.  The needed global complexity
bound is absent.  Proposition~\ref{p06:ldes:point-class} also prevents
substituting point codes for orbit invariants; its explicit perfect
$E_0$-antichain emphasizes that smoothness is not required by the actual
Vaught conclusion.  No reduction from the general case to the coanalytic
case is obtained.

\subsection{Target 39: Thomas's finite-reduct conjecture}
\label{p06:ldes:reducts}

We use the finite relational version: a countably infinite
ultrahomogeneous structure in a finite relational language has finitely
many parameter-free first-order reducts up to interdefinability.
Cameron restates this as an open problem in 2026
\cite[Question 8.1]{p06:ldes:Cameron}.  Dropping the relational hypothesis
changes its status: Bodor, Cameron and Szab\'o construct homogeneous
structures with infinitely many reducts, including the countable atomless
Boolean algebra in its functional signature \cite{p06:ldes:BodorCameronSzabo}.

Finite relational homogeneity implies $\omega$-categoricity: for each
tuple length there are only finitely many atomic diagrams, and equal
diagrams are carried to one another by homogeneity.  The following
reformulation isolates what finite observation would have to control.

\begin{proposition}[Uniform arity is the missing finite budget]
\label{p06:ldes:arity}
Let $M$ be a countable $\omega$-categorical structure in a countable
language, and let $o_k$ be
the number of $\operatorname{Aut}(M)$-orbits on $M^k$.
For fixed $r$, the number of interdefinability classes of reducts
generated by relations of arity at most $r$ is at most
\[
 2^{N_r},\qquad N_r=\sum_{k=1}^r2^{o_k}.
\]
Moreover, $M$ has finitely many reducts if and only if some $r$ works
for every reduct, in the sense that every reduct is interdefinable
with one generated by relations of arity at most $r$.
\end{proposition}
\begin{proof}
Every parameter-free definable $k$-ary relation is a union of
$\operatorname{Aut}(M)$-orbits, so there are at most $2^{o_k}$ choices.
A reduct generated in arities at most $r$ therefore chooses a subset
of a collection of size at most $N_r$.  This proves the bound and the
reverse implication in the equivalence.

For the forward implication assume there are only finitely many reduct
classes.  Enumerate the distinct relations of a reduct $R$ as
$R_1,R_2,\ldots$; there are only countably many parameter-free formulas
in the countable language.  The reducts generated by $R_1,\ldots,R_j$
form an increasing sequence in the partial order of definability
closures.  A monotone sequence in a finite partial order stabilizes.
After stabilization every later $R_j$ is definable from the same finite
initial list.  Thus $R$ has a finite relational generating list.
Choose such a list for a representative of each of the finitely many
classes, and take the maximum of their arities.
\end{proof}

The fixed-$r$ counting statement does not settle Thomas's conjecture:
the required $r$ must work simultaneously for all reducts, including
ones initially specified using arbitrarily large tuple arities.
Homogeneity of the original structure bounds neither that generating
arity nor the number of definability closures by the argument above.
This is a proved reformulation, not an established uniform bound.

\subsection{Target 41: categoricity transfer for atomic classes}
\label{p06:ldes:atomic}

The supplied label does not state a cardinal threshold, set-theoretic
hypothesis, or amalgamation assumption, so it has no single truth value.
Fix a countable complete first-order theory $T$ and write
$\operatorname{At}(T)$ for its atomic models: every finite tuple realizes
a type isolated over the empty set.  As a precise benchmark, one may
ask whether categoricity in some $\lambda\geq\beth_{\omega_1}$ implies
categoricity in every $\mu\geq\beth_{\omega_1}$ for this class.
This is a specialization of the standard $L_{\omega_1,\omega}$
categoricity conjecture stated in
\cite[Conjecture 1.2]{p06:ldes:Leung}, rather than a uniquely determined
interpretation of the original label.  Known transfer theorems for
short AECs impose structural hypotheses such as amalgamation over sets;
these hypotheses cannot simply be deleted.

Here is an exact translation and a counterexample to a much stronger,
naive transfer.  For every $n$, list the $T$-consistent formulas
$\varphi_{n,i}(\bar x)$ that isolate complete $n$-types.  Then
\[
 \Psi_T=\bigwedge T\ \wedge\
 \bigwedge_{n\geq1}\forall\bar x\,
       \bigvee_i\varphi_{n,i}(\bar x)
\]
is an $L_{\omega_1,\omega}$ sentence, and its models are exactly
$\operatorname{At}(T)$.  Indeed the displayed disjunction says precisely
that each tuple has an isolated type; all conjunctions and disjunctions
are countable.  This encoding gives access to infinitary categoricity
questions, but proves no excellence or amalgamation property.

Categoricity in one arbitrary infinite cardinal does not transfer for
atomic classes.  Let $T$ be dense linear order without endpoints.
Quantifier elimination makes the finite order-and-equality pattern of
each tuple an isolating formula, so every model is atomic.  The countable
model is unique.  Put $L=\sum_{\alpha<\omega_1}\mathbb Q$, with lexicographic
order, and $M=\sum_{j<\omega}L$.  Both are dense without endpoints and
have size $\aleph_1$.  The cofinalities are respectively $\omega_1$ and
$\omega$, so they are not isomorphic.  This elementary counterexample
does not refute the stated high-cardinal benchmark.

The attempt from exact observations stops at the same distinction:
identifying every finite tuple's type gives atomicity, but does not
give compatible amalgams of independent diagrams or uniqueness of
large models.  The outcome for this underspecified target is a precise
formalization and a rejected low-cardinal extrapolation, not a decision
of all atomic categoricity-transfer problems.

\subsection{Target 46: cardinal transfer while omitting types}
\label{p06:ldes:omit}

Again a threshold and a class must be specified.  We consider models
of a countable first-order theory omitting a specified countable family
of partial types in finitely many variables over the empty set.
This is not the same question as a two-cardinal transfer for a named
predicate.  The classical Hanf bound for this countable omitting-types
problem is $\beth_{\omega_1}$; below it no smaller universal bound
works, even after restricting to dependent theories, by the examples
of Chernikov and Shelah \cite[Theorem 4.5]{p06:ldes:ChernikovShelah}.
The following calculations test the most direct extension strategy.

\begin{proposition}[The obstruction is extension, not chain union]
\label{p06:ldes:omit-chain}
There is a complete countable theory and a nonprincipal partial type
having a countable omitting model but no uncountable omitting model.
On the other hand, the union of an elementary chain of models omitting
fixed finite-tuple partial types also omits those types.
\end{proposition}
\begin{proof}
Take the equality language with constants $c_n$ for $n\in\mathbb N$,
and let $T$ be the complete theory of a countably infinite set in which
the $c_n$ name all elements without repetition.  Set
$p(x)=\{x\ne c_n:n\in\mathbb N\}$.
The named model omits $p$, and any omitting model has no elements
beyond the interpretations of the constants, hence is countable.
The type is nonprincipal: a consistent formula mentions only finitely
many constants; in the pure equality theory with these finitely many
names, any unnamed element has the same type as a constant not among
those names.  Consequently no such formula can imply all inequalities
in $p$.  Finite subsets of $p$ are plainly satisfiable.

For the second assertion let $(M_i)$ be an elementary chain and
$M=\bigcup_iM_i$.  Any finite tuple $\bar a$ of $M$ lies in one $M_i$.
If it realized a fixed partial type $p$ in $M$, elementarity would make
every formula of $p$ true of $\bar a$ already in $M_i$, a contradiction.
Apply this separately to each specified type.
\end{proof}

Thus ordinary upward L\"owenheim--Skolem or compactness cannot supply
the necessary omitting extensions.  The chain-union step is valid;
constructing proper extensions that continue to omit the types is the
missing step.  The general Hanf theorem and its sharpness are existing
results, not new closures credited to the observation method.  More
specific cardinal-transfer conjectures require their omitted hypotheses
before they can be evaluated.

\subsection{Target 49: Borel boundedness and representative coherence}
\label{p06:ldes:bounded}

A countable Borel equivalence relation $E$ on $X$ is Borel bounded if
for every Borel $\phi:X\to\mathbb N^{\mathbb N}$ there is Borel $\psi$
such that
\[
 \phi(x)\leq^*\psi(x),\qquad
 xEy\ \Longrightarrow\ \psi(x)=^*\psi(y).
\]
The stars mean eventual domination and eventual equality, not exact
equality.  There are two distinct questions: whether every countable
Borel equivalence relation is Borel bounded, and whether Borel
boundedness implies hyperfiniteness.  Both are explicitly listed as
open by Day and Marks in 2026 \cite[Section 3, before Question 3.6]{p06:ldes:DayMarks}.
Coskey and Schneider give the definition and its connection with the
union problem \cite[Definitions 1.2 and 3.3]{p06:ldes:CoskeySchneider}.

The known implication from hyperfiniteness can be seen by an explicit
observation budget.  If $E=\bigcup_nF_n$, with $F_n$ finite Borel
equivalence relations increasing in $n$, put
\[
 \psi(x)(n)=\max\{\phi(y)(n):yF_nx\}.
\]
Finite-section Borel uniformization makes the maximum Borel.  It
dominates $\phi(x)(n)$, and for $xEy$ the $F_n$-classes of $x,y$ agree
for every sufficiently large $n$.  Hence the required eventual equality
holds.  This is the standard Boykin--Jackson observation, not a new
implication.

For a general relation, enumerating each class and taking diagonal
maxima supplies domination but need not supply coherence.  For example,
let $E$ be the one-class relation on $\mathbb Z$, take
$\phi(k)(n)=|k|$, and enumerate from $x$ the offsets
$0,1,-1,2,-2,\ldots$.  The maximum over the first $2n+1$ entries is
$n+|x|$.  The outputs for $x=0$ and $x=1$ are different at every $n$.
The relation itself is smooth and bounded; this is a counterexample
to that proposed construction, not to boundedness.  A Borel selector
of adequate pointwise bounds from our earlier framework would still
have to satisfy the eventual class-coherence condition.  No such
general selection theorem is proved here.

\subsection{Target 50: are all Polish groups graphic?}
\label{p06:ldes:graphic}

A Polish group is graphic if the orbit relation of every Borel action
on a standard Borel space $X$ is the connectedness relation of a Borel
graph on that same $X$.  Graph degrees need not be countable.  This
does not ask whether the group is the automorphism group of a countable
graph.  Arant, Kechris and Lutz prove graphicness for connected groups
and many other classes, and closure under countable products; they
retain the general question and the stronger closed-subgroup question
\cite[Section 4, Question 3]{p06:ldes:ArantKechrisLutz}.  They also prove
that every Polish orbit relation is strongly Borel $3$-hypergraphable
\cite[Theorem 65 in the arXiv v2 numbering]{p06:ldes:ArantKechrisLutz}.

The tempting observation transfer is to encode the third point and
then project each hyperedge to its pairs.  The following exact example
shows why Borelness does not follow from that operation.

\begin{proposition}[A Borel hypergraph can have a non-Borel pair shadow]
\label{p06:ldes:shadow}
There is a Borel $3$-uniform hypergraph whose relation of pairs lying
together in a hyperedge is not Borel.
\end{proposition}
\begin{proof}
Take an analytic non-Borel set $A$ in a Polish space $X$, represented
as the projection of a closed set $C\subseteq X\times\mathbb N^{\mathbb N}$.
Use the disjoint-union vertex space
\[
 Z=(X\times\{0\})\sqcup(X\times\{1\})
       \sqcup(X\times\mathbb N^{\mathbb N}\times\{2\}).
\]
For $(x,y)\in C$ put one hyperedge
$\{(x,0),(x,1),(x,y,2)\}$.  The corresponding symmetric relation of
ordered triples is Borel, since membership is tested by tags, equality
of the $X$-coordinates, and membership in $C$.
The pair $((x,0),(x,1))$ lies in its shadow exactly when $x\in A$.
The continuous map $x\mapsto((x,0),(x,1))$ would therefore reduce
$A$ to a Borel set if the shadow were Borel, a contradiction.
\end{proof}

This example is not a counterexample to graphicness of a Polish group.
It invalidates the specific witness-erasing projection step.  Keeping
the extra witness changes the vertex space, while the definition
requires a graph on the original space.  A different Borel encoding
that also preserves the required orbit components is still needed.

\subsection{Target 52: increasing unions of hyperfinite relations}
\label{p06:ldes:union}

The question is whether $D_0\subseteq D_1\subseteq\cdots$, each a
hyperfinite countable Borel equivalence relation, implies that
$E=\bigcup_kD_k$ is hyperfinite.  Frisch, Shinko and Vidny\'anszky retain
this as Question 0.2 in their September 2026 revision; their Theorem 0.3
shows that a counterexample would make the collection of hyperfinite
relations $\boldsymbol\Sigma^1_2$-complete
\cite{p06:ldes:FrischShinkoVidnyanszky}.  It is not legitimate to interchange
the two increasing limits without a coherence argument.

Write $D_k=\bigcup_mF_{k,m}$, with each row increasing in $m$ and each
$F_{k,m}$ finite Borel.  The following form of the known
Boykin--Jackson reduction identifies exactly the missing choice.

\begin{proposition}[Coherent rates yield a finite exhaustion]
\label{p06:ldes:coherent-rates}
Suppose a Borel $\psi:X\to\mathbb N^{\mathbb N}$ satisfies, for every
$xEy$, both
\[
 \psi(x)=^*\psi(y),\qquad
 xF_{k,\psi(x)(k)}y\quad\hbox{for all sufficiently large }k.
\]
Then $E$ is hyperfinite.  In particular, if $E$ is Borel bounded, these
rates exist and $E$ is hyperfinite.
\end{proposition}
\begin{proof}
Define $H_n$ by
\[
 xH_ny\quad\Longleftrightarrow\quad
 \forall k\geq n\ 
 \bigl(\psi(x)(k)=\psi(y)(k)\ \hbox{and}\
       xF_{k,\psi(x)(k)}y\bigr).
\]
For each fixed $k$ the displayed condition is an equivalence relation:
on each level set of $\psi(\cdot)(k)$ it is the corresponding $F_{k,m}$.
Thus $H_n$ is a Borel equivalence relation, being a countable
intersection of such relations.  Its class at $x$ is contained in
the finite $F_{n,\psi(x)(n)}$-class, and $H_n\subseteq H_{n+1}$.
Also $H_n\subseteq E$.  The two hypotheses ensure that every pair in
$E$ belongs to some $H_n$, proving $E=\bigcup_nH_n$.

For the last assertion enumerate each $E$-class by Borel maps
$f_i(x)$, allowing repetitions, using countable-section uniformization.
Set $\phi(x)(k)$ to the maximum, over $i\leq k$ for which
$xD_kf_i(x)$, of the least $m$ with $xF_{k,m}f_i(x)$; take maximum
$0$ for the empty set.  These least values and finite maxima are
Borel.  Borel boundedness supplies $\psi$ eventually dominating
$\phi$ and eventually equal on $E$-classes.  Fix $xEy$ and choose
$i$ with $y=f_i(x)$.  For all sufficiently large $k$, one has
$i\leq k$, $xD_ky$, and $\psi(x)(k)\geq\phi(x)(k)$, which gives
the second hypothesis by monotonicity of the row $F_{k,m}$.
\end{proof}

The implication in the last sentence is the existing theorem of
Boykin and Jackson, recorded as Theorem 1.3 by Coskey and Schneider
\cite{p06:ldes:CoskeySchneider}; the proof above makes its missing rate
explicit.  Joining finitely many finite equivalence relations is not
an alternative: on $\mathbb Z$, the pair partitions
$\{2j,2j+1\}$ and $\{2j-1,2j\}$ have an infinite join.  Those two
relations are not nested, so this is only a failure of a proposed
construction, not a counterexample to the increasing-union conjecture.
Neither the coherent rates nor Borel boundedness of the arbitrary
union are established here.  Targets 49 and 52 consequently remain
linked open problems rather than two closures of one selection lemma.


\section{Computability and decision problems: nine audited targets}
\label{p06:lcomp:section}

This section concerns catalogue entries 12, 16, 18, 19, 37, 43, 44,
51, and 57, with a literature checkpoint of 17 October 2026.
Entry 37 contains two different decision problems, and entry 43 contains
two different monadic theories.  We distinguish externally established
results, proved restricted transfers, and questions left open here.
No proof below is claimed to settle a previously open unrestricted
problem.  In particular, qualitative identification of an object,
finite-dimensional representation, and a uniform decision algorithm
are different conclusions.

\subsection{12. Automorphisms of the Turing degrees}
\label{p06:lcomp:12}

Let $\mathcal D=(\mathcal P(\mathbb N)/{\equiv_T},\leq_T)$.
The rigidity conjecture asks whether every order automorphism of
$\mathcal D$ is the identity.  This is a question about all Turing
degrees, not just the computably enumerable degrees or degrees below a
particular jump.  Slaman and Woodin formulate rigidity as
Conjecture 7.5.4 and prove its equivalence to their biinterpretability
statement in Theorem 7.5.3 of \cite{p06:lcomp:SW}.  The introduction of
\cite{p06:lcomp:SS} records the finite-automorphism-base theorem while
distinguishing the still unresolved global rigidity problem from its
own results on $\Delta^0_2$ degrees.  No resolution was verified in the
present source search.

Here is the precise limitation of a reconstruction argument.
An \emph{automorphism base} of a structure $\mathcal A$ is a subset
$B$ whose pointwise stabilizer in $\operatorname{Aut}(\mathcal A)$ is
trivial.

\begin{proposition}[Anchoring a reconstruction base]
\label{p06:lcomp:base}
If $B$ is an automorphism base, restriction to $B$ is an injective map
from $\operatorname{Aut}(\mathcal A)$ to the set of maps $B\to A$.
If each element of $B$ is individually definable without parameters,
then $\mathcal A$ is rigid.
\end{proposition}
\begin{proof}
If automorphisms $f,g$ agree on $B$, then $g^{-1}f$ fixes $B$
pointwise and hence is the identity.  Every automorphism preserves a
parameter-free formula with a unique solution.  Under the second
assumption all automorphisms therefore fix $B$ pointwise.
\end{proof}

This elementary proposition isolates an actual missing step.  A finite
base is a collection of identifying coordinates; it need not be a
collection of invariantly named coordinates.  In the two-element
structure with equality only, a singleton is an automorphism base but
the transposition is an automorphism.  Thus neither a finite-base
theorem nor an analogue of the earlier mixture-identification results
proves rigidity.  One would still have to establish suitable
parameter-free anchors, or another argument forcing the automorphisms
to fix a base.  That task is not achieved here.

\subsection{16. Kolmogorov--Loveland versus Martin--L\"of randomness}
\label{p06:lcomp:16}

Work on Cantor space with fair-coin measure.  A nonmonotonic betting
strategy chooses, computably from the bits already inspected, a fresh
coordinate and a fair nonnegative wager on that bit.  Success means
unbounded capital along the resulting play.  In the usual
Kolmogorov--Loveland setting partial computable strategies are allowed;
only completed legal steps contribute to a play.  The equality
question asks whether avoiding success of all such strategies is
equivalent to Martin--L\"of randomness.  Requiring all strategies to
be total is a convention that must be stated separately; see the
definitions and comparison in \cite{p06:lcomp:Merkle}.

\begin{proposition}[The effective maximal-inequality direction]
\label{p06:lcomp:MLKL}
Every Martin--L\"of random sequence defeats every computable
nonmonotonic fair betting strategy, including a partial strategy on
its completed legal play.
\end{proposition}
\begin{proof}
Normalize the initial capital to one.  A finite legal history of $k$
queried bits specifies a cylinder on $k$ distinct coordinates and has
measure $2^{-k}$.  At any history the two successor histories query the
same new coordinate and have equal conditional probability; fairness
makes capital a nonnegative martingale on this binary history tree.
Stopping at the first completed history whose capital exceeds $2^m$
and truncating the tree at depth $N$, conservation of expected capital
gives measure at most $2^{-m}$ for the stopped histories.  Histories at
which a partial computation does not finish can be made terminal;
this does not increase that measure.  Letting $N$ increase proves the
same bound for all first-hit histories.

Completed histories are effectively enumerable.  Their finite
coordinate cylinders are effectively open in the usual product
topology: each is a finite union of initial-segment cylinders.  Thus
the union $U_m$ of histories with capital exceeding $2^m$ is uniformly
effectively open and has measure at most $2^{-m}$.  A successful play
belongs to every $U_m$, contrary to Martin--L\"of randomness.
\end{proof}

The reverse direction is not proved.  An especially recent result
must not be misreported as doing so.  Petrovi\'c's preprint of
9 October 2026, \cite[Theorem~4.1]{p06:lcomp:Petrovic}, constructs an
encoding $E$ and two total computable nonmonotonic strategies which
succeed on the encoded points associated with a non-Martin--L\"of
random input.  The encoding in its Section 2.8 is given through
clopen partitions; the theorem is not a success theorem on the
unchanged input.  Its Section 5 explicitly explains a failure of the
analogous unencoded construction.  Our attempted transfer would need
to remove this encoding while preserving computable betting and
success on the original sequence.  Injectivity or measure preservation
alone supplies no such conversion.  The missing effective conversion,
and the total/partial convention, remain essential.

\subsection{18. Which Kreisel conjecture?}
\label{p06:lcomp:18}

The primary interpretation here is the bounded-proof-steps conjecture
for \emph{PA with Kleene's specified presentation}.  If $A(x)$ has
only the displayed free variable and some fixed $k$ bounds the number
of inference steps in proofs of $A(\overline n)$ for every standard
$n$, must PA prove $\forall x A(x)$?  The numeral convention, language,
axioms, and proof calculus form part of the problem.

Piazza's 2026 paper \cite[Sections~1 and~5]{p06:lcomp:Piazza} states this
version as Friedman's Problem 34 and proves a finite algebraic
\emph{target} theorem, not the conjecture.  Parikh's positive result
for a relational presentation $\mathrm{PA}^*$ and Hrube\v{s}'s
negative result for a conservative definitional extension
$\mathrm{PA}(-)$ are not contradictory: they concern different proof
systems.  See \cite[Introduction, Theorem~3 and Corollary~1]{p06:lcomp:Hrubes}.
A bound by a total recursive function on \emph{proof codes} is yet
another formulation; the adapted formulation in
\cite[Theorem~2.1]{p06:lcomp:SK} is false.

\begin{proposition}[Uniform proof search is weaker than universal proof]
\label{p06:lcomp:selector}
For any recursively axiomatized theory $T$, if
$T\vdash A(\overline n)$ for every standard $n$, a total computable
function can select a $T$-proof of each instance.  Even for consistent
PA this conclusion does not imply $\mathrm{PA}\vdash\forall x A(x)$.
\end{proposition}
\begin{proof}
On input $n$, enumerate $T$-proofs until a proof of
$A(\overline n)$ occurs.  The assumption makes this algorithm total.
For the second assertion let $A(n)$ say that there is no PA-proof of
$0=1$ with code at most $n$, using the standard arithmetization of
proofs.  If PA is consistent, each individual assertion is a true
finite primitive-recursive calculation and is provable in PA.
Their universal closure is the usual consistency statement for PA,
which PA cannot prove by the second incompleteness theorem.
\end{proof}

This does not refute Kreisel's conjecture: the proof-step lengths in
this example have not been uniformly bounded.  The attempted
finite-state transfer is correspondingly precise.  Finitely many
proof skeletons at a fixed bound yield finitely many combinatorial
states, but the states still require arithmetic realizations.
Piazza's Section 5 separates external membership in these realization
sets from internally proving their coverage and transporting that
coverage to $A(x)$.  Neither our earlier finite spectral descriptions
nor Proposition~\ref{p06:lcomp:selector} fills that arithmetic gap.

\subsection{19. Sacks's conjecture on order embeddings}
\label{p06:lcomp:19}

The conjecture is that every locally countable partial order of
cardinality at most $2^{\aleph_0}$ embeds into $\mathcal D$ by an
order embedding, preserving both comparisons and incomparisons.
Local countability means that every principal lower cone is countable.
These are necessary conditions: there are only countably many sets
computable from any fixed oracle, and there are $2^{\aleph_0}$ Turing
degrees.  They are not presently known to be sufficient in ZFC.

Higuchi and Lutz prove the height-two case and a serious obstruction
at height three in \cite[Theorems~1.1 and~1.2]{p06:lcomp:HL}.
Miller and Greenberg obtain further height-three cases in
\cite[Theorem~1.1 and Proposition~1.2]{p06:lcomp:MG}.
Their restriction counts elements of \emph{depth three}, with depth
defined by chains beginning at an element; it is not a claim about
arbitrarily many uppermost elements under a reversed convention.

\begin{proposition}[Obstruction to a continuous bottom-layer transfer]
\label{p06:lcomp:perfect}
A scheme which always embeds a Cantor-indexed bottom layer using a
continuous injective map into $2^{\mathbb N}$ cannot prove the general
Sacks conjecture by subsequently extending that map to the other
layers.
\end{proposition}
\begin{proof}
A continuous injection from the compact space $2^{\mathbb N}$ into
$2^{\mathbb N}$ is a homeomorphism onto a compact subset without
isolated points, hence has perfect image.  Higuchi--Lutz
\cite[Theorem~1.2]{p06:lcomp:HL} construct a locally countable height-three
partial order of size continuum such that an embedding of its bottom
layer whose image contains a perfect set of representatives cannot
extend to an embedding of the whole order.  Apply the proposed scheme
to that order.  Its prescribed bottom-layer image has exactly the
property forbidden by that theorem.
\end{proof}

This is an obstruction to a specific regular construction, not a
counterexample to Sacks's conjecture in ZFC.  To use analytic
reconstruction here one must additionally construct the required
Turing reductions and prevent all unwanted reductions across an
uncountable family.  A perfect-set coding of independent coordinates
does not provide that extension theorem.

\subsection{37. Skolem and Positivity: an effective spectral subcase}
\label{p06:lcomp:37}

For a rational linear recurrence sequence $(u_n)_{n\geq0}$ specified
by its coefficients and initial values, the \emph{Skolem problem}
asks whether some $u_n=0$.  The \emph{Positivity problem} asks whether
every $u_n\geq0$; it is distinct from eventual positivity.  Integer
and rational formulations are effectively interchangeable by an
appropriate positive exponential rescaling.  The 2026 survey
\cite[Theorem~25, Sections~5.3, 7 and~9]{p06:lcomp:SkolemSurvey} records
decidability of Skolem up to order four and leaves order five open.
For Positivity, general order six and simple order ten remain open;
general order at most five and simple order at most nine are known
decidable.  These statements are literature results, not consequences
of this manuscript.

The projected-zero argument in the preceding structural manuscript
used Skolem--Mahler--Lech qualitatively.  It gave arithmetic
progressions plus a finite exceptional set; it did not compute a bound
for that exceptional set.  The following restricted transfer makes
the extra effective hypothesis explicit.  All algebraic inputs are
given by exact defining polynomials and isolating data.

\begin{theorem}[Decision with a common-power positive-spectrum promise]
\label{p06:lcomp:positive-spectrum}
Let $(u_n)$ be a real algebraic linear recurrence sequence of arbitrary
finite order.  Suppose an integer $M\geq1$ is supplied such that
$\lambda^M$ is positive real for every nonzero characteristic root
$\lambda$.  There is an algorithm computing its zero set as a finite
set together with finitely many arithmetic progressions.  The same
algorithm determines the eventual sign on every residue class modulo
$M$ and decides both Skolem and Positivity for this input class.
\end{theorem}
\begin{proof}
Effective algebraic linear algebra computes the generalized
eigenspace expansion of the recurrence.  The zero characteristic
root contributes only a computably bounded finite transient, which
we first check directly.  Choose a computable $K_*$ such that
$Mk+r$ is beyond the transient whenever $k\geq K_*$ and
$0\leq r<M$.  Combining equal bases gives, on each residue,
\begin{equation}\label{p06:lcomp:spectral-expansion}
 u_{Mk+r}=\sum_{j=1}^{s_r} P_{r,j}(k)\rho_{r,j}^{\,k},
 \qquad 0<\rho_{r,1}<\cdots<\rho_{r,s_r},\quad k\geq K_*.
\end{equation}
The coefficients and bases are algebraic and effectively computable.
They are real: conjugating the expansion leaves the real sequence
unchanged, and independence of exponential-polynomial sequences with
distinct positive bases forces each combined polynomial to equal its
conjugate.  Alternatively one can group conjugate terms directly.
Exact algebraic arithmetic decides which polynomials vanish.

If all polynomials vanish, this residue contributes the entire
progression $\{Mk+r:k\geq K_*\}$.  Otherwise delete zero polynomials
and suppress $r$ in the notation.  Let the largest remaining base be
$\rho_s$, and write $P_s(k)=a k^d+\sum_{i<d}a_i k^i$, with $a\ne0$.
Choose rational $a_*>0$ with $a_*<|a|$, and rational
$B>\sum_{i<d}|a_i|$.  For $k\geq1$, the lower-degree sum is at most
$B k^{d-1}$ when $d\geq1$.  Therefore, for an effectively chosen
$K_0\geq K_*$, the polynomial $P_s(k)$ has sign $\operatorname{sgn}(a)$
and
\begin{equation}\label{p06:lcomp:dominant-lower}
 |P_s(k)|\geq c k^d \quad(k\geq K_0)
\end{equation}
for a computable rational $c>0$, for instance $c=a_*/2$.
For $d=0$ this holds directly.

If $s=1$, the conclusion already follows.  Otherwise algebraic
comparison supplies rational $q\in(0,1)$ exceeding every
$\rho_j/\rho_s$ for $j<s$.  Polynomial coefficient bounds give a
rational $C>0$ and integer $D\geq0$ such that, for all $k\geq1$,
\[
 \left|\sum_{j<s}P_j(k)(\rho_j/\rho_s)^k\right|
       \leq C k^Dq^k.
\]
Set $h=\max(D-d,0)$.  We now give an effective tail cutoff rather
than merely using exponential decay.  Choose rational $q'$ with
$q<q'<1$.  Search for $K_1\geq\max(K_0,1)$ such that
$q(1+1/K_1)^h<q'$.  The ratio
\[
 \frac{(k+1)^h q^{k+1}}{k^h q^k}
       =q(1+1/k)^h
\]
then remains below $q'$ for $k\geq K_1$.  Search for $K\geq K_1$
with $C K^h q^K<c/2$.  This search terminates by the same geometric
ratio bound, and all its comparisons are rational.  For every
$k\geq K$ the lower bases contribute less than
$(c/2) k^d$, whereas \eqref{p06:lcomp:dominant-lower} contributes at least
$c k^d$.  Thus \eqref{p06:lcomp:spectral-expansion} is nonzero with the
fixed sign of $a$ on that tail.

Repeat over the finitely many residues.  Check all remaining initial
terms exactly by real algebraic arithmetic.  The zero progressions
and the finite zeros found in these checks are the full zero set.
For Positivity, every fixed tail sign and every checked initial sign
must be nonnegative.  This is a terminating decision procedure.
\end{proof}

The promise is verifiable for the supplied $M$ by algebraic
arithmetic.  The theorem covers arbitrary Jordan multiplicities and
root-of-unity phases of bounded common order.  It is an explicit
version of a classical dominance argument for a promised subcase;
we make no historical novelty claim.  General rational recurrences
need not satisfy the promise.  Several dominant roots with
non-root-of-unity relative phases produce oscillatory cancellation;
the proof's strict positive-base separation then fails.

\begin{proposition}[Dimension one does not replace exact arithmetic]
\label{p06:lcomp:computable-real}
If initial values are supplied as arbitrary uniformly computable real
numbers, there is no uniform Skolem algorithm even for the fixed
recurrence $u_{n+1}=u_n$.  There is likewise no uniform Positivity
algorithm for that input representation.
\end{proposition}
\begin{proof}
Given a program $e$, let
\[
 c_e=\begin{cases}2^{-T},&\text{if $e$ first halts at step $T\geq1$,}\\
                   0,&\text{if $e$ never halts.}\end{cases}
\]
This is a uniformly computable real.  To approximate within $2^{-m}$,
run $e$ for $m$ steps, returning $2^{-T}$ if it has halted and zero
otherwise.  A later halt changes the value by less than $2^{-m}$.
The constant sequence with initial value $c_e$ has a zero exactly
when $e$ does not halt.  A decision algorithm would decide the halting
problem.  For Positivity use the constant sequence with initial value
$-c_e$: all its terms are nonnegative exactly when $e$ does not halt.
\end{proof}

Here every zero set is either empty or all of $\mathbb N$, and every
spectrum consists of the single root $1$.  Thus even this strongest
qualitative zero-set description does not supply an algorithm for
an arbitrary computable-real representation.  This proposition does
\emph{not} prove the classical rational-input Skolem or Positivity
problem undecidable.  It identifies the exact arithmetic assumption
that an attempted transfer must retain.

\subsection{43. BMTO has a new answer; MTWO is a different question}
\label{p06:lcomp:43}

BMTO means the \emph{Borel monadic theory of the real order}: first
order variables range over $\mathbb R$, the relation is $\leq$, and
monadic second-order variables range over Borel subsets of $\mathbb R$.
It is decidable by Manthe \cite[Theorem~1.1]{p06:lcomp:Manthe}, whose
second preprint version is dated 9 March 2026.  This is an external
resolution and must be removed from an unqualified list of unsolved
problems.  We do not claim to have proved Manthe's theorem.

MTWO means the \emph{monadic theory of well-ordering}, with full
subset quantification.  We interpret the historical question as
asking whether it is consistent with ZFC that the set of monadic
sentences valid in all wellorders is recursive.  Satisfiability in
some wellorder is the equivalent complementary decision problem.
It is not the theory of the single order $(\omega,<)$, and it is not
BMTO.  Shelah's proposed sufficient set-theoretic hypothesis appears
as Conjecture 0.3 in the author-posted version of
\cite{p06:lcomp:ShelahOrder}.  Neeman's primary account
\cite[Theorems~11--16 and final discussion]{p06:lcomp:Neeman} separates
independence at $\omega_2$, consistency of decidability at $\omega_3$,
and still further questions about all ordinals.  No later resolution
of the full historical MTWO question was verified in this audit;
its status here is retained as unresolved, with that explicit
verification limit.

\begin{proposition}[The missing effective table]
\label{p06:lcomp:types}
Fix a set-theoretic universe.  Suppose an algorithm, given $r$,
returns the realized members of a fixed effective finite list of
rank-$r$ monadic Hintikka sentence types, and this returned list
is exactly the set of types realized by wellorders in that universe.
Then MTWO in that universe is decidable.
\end{proposition}
\begin{proof}
An effective finite Hintikka list, allowing unrealizable members,
is obtained by induction on rank and number of free individual and
set variables in this finite relational language.  Given a sentence,
one computes its rank and the types on which it holds by the same
syntactic induction.  It is valid in all wellorders if and only if it
holds on every type returned by the hypothesized algorithm.  This
uses no decision procedure for consistency of the candidate types;
computing which candidates are realized is exactly the hypothesis.
\end{proof}

The hypothesis is not established.  Merely knowing that there are
finitely many rank-$r$ types does not compute which are realized.
Neeman's full-limit automata make the obstacle concrete: their
uncountable-limit transitions consult formulas in the almost-all
theory of the run, involving club and stationary-set behavior.
This is additional semantic data, not an ordinary finite transition
table that can simply be executed on a finite input.  The previous
finite-state or finite-spectrum arguments do not compute it.
Manthe's restriction to Borel subsets of the real line does not
eliminate full-subset quantification over arbitrary uncountable
ordinals, so it supplies no asserted solution of MTWO.

\subsection{44. Laurent series: specify both field and language}
\label{p06:lcomp:44}

The principal open target is the full first-order theory of
$\mathbb F_q((t))$ in the pure ring language, for a fixed finite field
$\mathbb F_q$ of positive characteristic $p$.  The notation
$\mathbb Z_p$ in informal lists must not silently identify the finite
field $\mathbb F_p$ with the ring of $p$-adic integers, which is not
itself a field.  Characteristic-zero Laurent fields have a different
Ax--Kochen--Ershov theory.  Naming a uniformizer, adding valuation
predicates, or adding a predicate for its integral powers also changes
the question.

Anscombe and Fehm \cite[Introduction and Theorems~1.1--1.2]{p06:lcomp:AF}
explicitly distinguish the open full positive-characteristic problem
from known existential decidability without a named uniformizer and
from results conditional on resolution assumptions in stronger
languages.  In a substantially stronger expansion, Gitin
\cite[Theorem~9]{p06:lcomp:Gitin} proves existential undecidability when
a positive-valuation element $\alpha$ and the cyclic subgroup
$\alpha^{\mathbb Z}$ are named.  That is not undecidability of the
pure field.

\begin{proposition}[An exact Frobenius fragment]
\label{p06:lcomp:frobenius}
For an explicitly given rational function $a\in\mathbb F_q(t)$,
the following conditions are equivalent and decidable:
\[
 (\exists x\in\mathbb F_q((t)))\ x^p=a,
 \qquad \frac{da}{dt}=0,
 \qquad a\in\mathbb F_q(t^p).
\]
When they hold, a root can already be computed in $\mathbb F_q(t)$.
\end{proposition}
\begin{proof}
In characteristic $p$, a Laurent series raised to the $p$th power has
only exponents divisible by $p$, so its formal derivative vanishes.
The rational-function derivative is the restriction of the
Laurent-series derivative.  Conversely,
$1,t,\ldots,t^{p-1}$ is a basis of $\mathbb F_q(t)$ over
$\mathbb F_q(t^p)$.  Writing
$a=\sum_{i=0}^{p-1}b_i(t^p)t^i$ and differentiating shows
$a'=\sum_{i=1}^{p-1} i b_i(t^p)t^{i-1}$.  Basis independence forces
$b_i=0$ for $i>0$ when $a'=0$.  Finally, the finite coefficient field
is perfect.  If $a=b(t^p)$, applying inverse Frobenius to the
coefficients of a numerator and denominator of $b$ gives a rational
function whose $p$th power is $a$.  Polynomial arithmetic computes
the derivative and tests whether it is identically zero, and the
same basis calculation computes the root.
\end{proof}

This is a restricted parameterized family of existential instances,
not new quantifier elimination.  It demonstrates where exact
coefficient information makes a decision possible.  It does not
address arbitrary alternating quantifiers or supply the missing
positive-characteristic model theory.

\subsection{51. Scott sets and standard systems}
\label{p06:lcomp:51}

A Scott set is a nonempty Turing ideal $S\subseteq\mathcal P(\mathbb N)$
such that every infinite binary tree coded in $S$ has a path in $S$.
For a nonstandard model $M\models\mathrm{PA}$, its standard system
$\operatorname{SSy}(M)$ consists of the traces on the standard natural
numbers of sets coded in $M$.  Scott's problem asks whether every
Scott set is $\operatorname{SSy}(M)$ for some such $M$.
It is known for countable Scott sets and for Scott sets of cardinality
at most $\aleph_1$; consequently CH gives a positive answer for all
Scott sets.  The unrestricted ZFC question remains open in the
primary accounts \cite[Theorems~1.2--1.3, Question~1.4]{p06:lcomp:Wang}
and \cite[Theorems~1.2--1.3, Corollary~1.4]{p06:lcomp:Gitman}.
Wang's manuscript is dated 2020; a later crawl date is not a new
publication or a new resolution.

\begin{proposition}[Where the chain construction needs an extension theorem]
\label{p06:lcomp:scott-chain}
Let $S$ be a Scott set of uncountable cardinality $\kappa$.
Assume a countable nonstandard $M_0\models\mathrm{PA}$ has
$\operatorname{SSy}(M_0)\subseteq S$.  Suppose that whenever
$M\models\mathrm{PA}$ has cardinality less than $\kappa$ and
$\operatorname{SSy}(M)\subseteq S$, each $A\in S$ can be placed in
$\operatorname{SSy}(N)$ for an elementary extension $N\succeq M$
satisfying
\[
 \operatorname{SSy}(N)\subseteq S,
 \qquad |N|\leq\max(|M|,\aleph_0).
\]
Then $S$ is the standard system of a model of PA.
\end{proposition}
\begin{proof}
Enumerate $S=\{A_\alpha:\alpha<\kappa\}$.
At successor stage use the assumed extension property to put
$A_\alpha$ into the standard system.  At limit stages take the
elementary union.  Every parameter coding a subset of $\mathbb N$
in such a union already belongs to an earlier model.  Elementarity
preserves its membership tests at each standard integer.  Hence the
standard system of a chain union is the union of the earlier standard
systems.  By induction
$|M_\alpha|\leq\max(\aleph_0,|\alpha|)$ for infinite stages
$\alpha<\kappa$, so the hypothesis remains applicable even when
$\kappa$ is singular.  The final union has standard system contained
in $S$ and contains every $A_\alpha$, as required.
\end{proof}

This is the standard Ehrenfeucht extension strategy; compare
\cite[Theorem~1.7]{p06:lcomp:Gitman}.  A countable initial model with
standard system inside $S$ is available by first closing a countable
subfamily under the Scott-set operations and applying the countable
case.  For $\kappa=\aleph_1$ the required countable extension lemma
is known; see \cite[Theorem~2.1]{p06:lcomp:Wang}.  For larger $\kappa$,
the uncountable extension property, with control of all new coded
traces, is the unresolved input.  Recovering a probability law or a
finite family of invariants does not provide that elementary-extension
construction.

\subsection{57. The universal--existential theory of the c.e. degrees}
\label{p06:lcomp:57}

Let $\mathcal R$ be the degrees containing computably enumerable
subsets of $\mathbb N$, with the pure order relation $\leq_T$.
The target asks whether the set of true sentences
$\forall\bar x\exists\bar y\,\theta(\bar x,\bar y)$, with
quantifier-free $\theta$ in this language, is decidable.
The one-block theory is decidable; the three-block theory is
undecidable by Lempp--Nies--Slaman
\cite[Introduction and main theorem]{p06:lcomp:LNS}.  Their two-block
question is still recorded as open in the author's later account
\cite[slides~2--4]{p06:lcomp:Soskova}; no subsequent resolution was
verified in this search.  Results for weak truth-table degrees,
all Turing degrees, or languages with additional join, meet, or jump
operations must not be substituted.

\begin{proposition}[Exact finite-menu reformulation]
\label{p06:lcomp:menus}
For any fixed nonempty partial order $D$, its pure-order
$\forall\exists$ theory is decidable if and only if the following
finite-menu extension problem is decidable: given finite labeled
partial orders $P\subseteq Q_1,\ldots,Q_s$, does every embedding
of $P$ into $D$ extend to an embedding of at least one $Q_j$?
Finite conjunctions of such menu instances are allowed without
changing decidability.
\end{proposition}
\begin{proof}
For a sentence $\forall\bar x\exists\bar y\,\theta$, enumerate
all complete atomic diagrams on the finitely many $x$ variables,
including equality and all comparisons and noncomparisons.  Discard
those incompatible with a partial order.  Quotient each remaining
diagram by its equalities to obtain a finite labeled partial order
$P$.  Enumerate all complete diagrams on $(\bar x,\bar y)$ which
extend that diagram and satisfy $\theta$.  Quotienting equalities
gives a finite menu of extensions $Q_j$ of $P$, recording that some
$y$ variables may name old or repeated points.  The sentence is true
if and only if the corresponding menu condition holds for every
initial diagram.  An unrealizable initial diagram contributes a
vacuously true condition.

Conversely let $\delta_P(\bar x)$ and
$\delta_{Q_j}(\bar x,\bar y_j)$ be complete order diagrams, using
distinct variables for distinct labeled points.  The menu condition
is
\[
 \forall\bar x\left(\delta_P(\bar x)\ \longrightarrow\
       \bigvee_{j=1}^s\exists\bar y_j\,
                   \delta_{Q_j}(\bar x,\bar y_j)\right).
\]
Move the finite existential blocks together.  Variables in unused
disjuncts can be filled arbitrarily because $D$ is nonempty.  This
is a $\forall\exists$ sentence.  An empty menu is handled by a
false consequent.  Finite conjunctions of $\forall\exists$
sentences remain in the same fragment after renaming and merging
their quantifier blocks.
\end{proof}

This is a standard structural translation, explicitly used in
\cite[slides~3--4]{p06:lcomp:Soskova}, and is included to expose the
remaining computational task.  A decision procedure for a
\emph{single} prescribed extension $Q$ does not automatically decide
a menu.  In general
\[
 \forall e\,\bigvee_j \operatorname{Ext}(e,Q_j)
 \quad\not\Longleftrightarrow\quad
 \bigvee_j\forall e\,\operatorname{Ext}(e,Q_j).
\]
For two possible embeddings, one extension may work only for the
first and another only for the second.  The left condition then
holds and the right fails.  One needs a uniform structural
classification of the menu property in the actual c.e. degree
structure.  Neither the reduction nor the previous analytic
identification results provide it.

\subsection{Outcome of this branch}
\label{p06:lcomp:outcome}

The source audit corrects an obsolete status for BMTO, an external
solution.  It also prevents three concrete false transfers: the
recent encoded randomness theorem is not KL$=$ML; a finite algebraic
target for proof skeletons is not Kreisel's conjecture; and
Skolem--Mahler--Lech is not a general zero-finding algorithm.
The fully proved algorithm in
Theorem~\ref{p06:lcomp:positive-spectrum} addresses an exact promised
spectral class and both decision questions in entry 37.
Proposition~\ref{p06:lcomp:computable-real} shows that dropping exact
arithmetic can destroy uniform decidability even in dimension one.
The remaining propositions give explicit reductions or construction
obstructions, with their missing hypotheses stated.  None is
counted as a new closure of an unrestricted historical conjecture.


% Proof theory and reverse mathematics: source audit and transfer attempts.
% Body fragment. Requires amsmath, amssymb, amsthm, hyperref.
\section{Proof theory and reverse mathematics: nine precisely delimited targets}
\label{p06:lproof:section}

This section audits catalogue items 2, 15, 17, 24, 25, 26, 47, 54 and 55 as of
17 October 2026.  None of the nine unrestricted targets is claimed to be solved
here.  The positive statements proved below are transfer criteria and explicit
obstructions to proposed shortcuts; they are not counted as nine new solutions.
In particular, an analytic inverse theorem, a uniquely specified optimizer, or
an existence proof over the real numbers supplies neither a short propositional
proof nor a derivation in a weak arithmetic theory without an additional,
quantitative formalization theorem.

We use the usual second-order arithmetic systems $\mathsf{RCA}_0$,
$\mathsf{WKL}_0$, $\mathsf{ACA}_0$, $\mathsf{ATR}_0$, and
$\Pi^1_n\text{-}\mathsf{CA}_0$ with set parameters in the indicated comprehension
schemes.  A statement of conservativity always specifies both its base theory
and its syntactic class.  Statements about models in the arguments below are
metatheorems in ordinary classical mathematics, not implicit formalizations
inside the weak theory under discussion.

\subsection{Item 2: unrestricted Frege lower bounds}
\label{p06:lproof:frege}

\paragraph{Exact target and verified status.}
Fix a standard finite Frege calculus $F$, with proofs measured by total symbol
length, and let $s_F(\tau)$ be the length of a shortest proof of the tautology
$\tau$.  The target is that $F$ is not polynomially bounded:
\[
 \forall d\,\forall C\,\exists\tau\in\mathrm{TAUT}\qquad
 s_F(\tau)>C(1+|\tau|)^d.
\]
An explicit uniformly constructible hard family is a stronger, usual research
goal.  Extended Frege, bounded-depth Frege, and a restriction on individual line
size are different targets.  Filmus--Pitassi--Santhanam's Theorem~2 gives a
subexponential bounded-depth simulation of short Frege proofs
\cite{p06:lproof:FPS}.  A 2026 result of de Rezende--Engstr\"om--Ghannane--Risse proves
superpolynomial length lower bounds for \emph{tree-like semantic systems with
bounded line size}; its hypotheses do not yield the displayed unrestricted
statement \cite{p06:lproof:REGR}.  Our current primary-source search found no
verified unrestricted resolution.

\begin{proposition}[The quantitative hypothesis required for a transfer]
\label{p06:lproof:proofsize}
Suppose a map sends every $F$-proof $\pi$ of $\tau$ to a certificate in a system
$Q$ for an instance $R(\tau)$, with length at most
$C(|\pi|+|\tau|)^d$, where $C>0$ and $d\geq1$ are fixed.  If all $Q$-certificates
for $R(\tau)$ have length at least $L(\tau)$, then
\[
 s_F(\tau)\geq \bigl(L(\tau)/C\bigr)^{1/d}-|\tau|.
\]
Consequently, a superpolynomial lower bound in $|\tau|$ for these certificates
transfers to $F$.
\end{proposition}
\begin{proof}
Apply both bounds to the image of a shortest $F$-proof and solve the inequality
$L(\tau)\leq C(s_F(\tau)+|\tau|)^d$.  If the conclusion had a polynomial upper
bound, this same inequality would give a polynomial upper bound on $L$.
\end{proof}

\paragraph{Attempt and obstruction.}
The attempted transfer encodes a tautology as a finite positivity or
identification assertion and then uses the earlier analytic uniqueness results
to certify it.  Even granting such a sound encoding, we have neither a map from
\emph{every} short Frege proof to a bounded-size analytic certificate nor a
lower bound on the encoded certificates.  An algorithm that happens to produce
large certificates does not lower-bound other proofs.  The proposition locates
the two missing statements exactly.  It is an elementary conditional transfer,
not a new Frege lower bound.  Even a bounded-depth lower bound must have the
strength required by the simulation in \cite{p06:lproof:FPS}; a lower-bound exponent
that degenerates with depth cannot be silently made uniform.

\subsection[Item 15: Fraisse's theorem over ATR0]{Item 15: Fra\"iss\'e's theorem over $\mathsf{ATR}_0$}
\label{p06:lproof:fraisse}

\paragraph{Exact target and verified status.}
Let $\mathsf{FRA}$ state that every sequence $(L_i)_{i\in\mathbb N}$ of countable
linear orders has $i<j$ and an order embedding $L_i\hookrightarrow L_j$.
The question is $\mathsf{ATR}_0\vdash\mathsf{FRA}$, equivalently the remaining
implication in the proposed equivalence over $\mathsf{RCA}_0$.  This is not the
truth of Fra\"iss\'e's conjecture, which Laver proved.  Freund's 2024 paper,
Theorem~1.2, gives the stronger-than-$\mathsf{ATR}_0$ upper theory of pseudo
$\Pi^1_1$-comprehension with the indicated $\Gamma$ closure; its introduction
records the lower implication to $\mathsf{ATR}_0$ and the unresolved converse
\cite{p06:lproof:Freund}.  That upper bound is not a proof in $\mathsf{ATR}_0$.

\begin{lemma}[Bounded witnesses are the missing compactness datum]
\label{p06:lproof:boundedtree}
Over $\mathsf{WKL}_0$, let $T\subseteq\mathbb N^{<\mathbb N}$ be a coded
prefix-closed tree and let $b:\mathbb N\to\mathbb N$ be a function such that
$\sigma(i)\leq b(i)$ for every $\sigma\in T$ and $i<|\sigma|$.
If $T$ has nodes of every length, it has an infinite path.
\end{lemma}
\begin{proof}
Encode the $i$th entry by a binary block of a fixed positive length large enough
to encode $0,\ldots,b(i)$.  Form the binary tree of prefixes of the encoded
members of $T$.  Membership of a binary string is decidable relative to $T,b$:
at a partial last block there are only boundedly many completions to test.
This is an infinite binary tree.  A path given by weak K\"onig's lemma decodes
at block endpoints into a path through $T$.
\end{proof}

For orders with underlying set $\mathbb N$, suppose that for each $n$ the induced
order on $\{0,\ldots,n-1\}$ in $L$ embeds in $M$, with the image of $i$ bounded
by the \emph{same} function $b(i)$ for all $n$.  The tree of such finite
embeddings satisfies the lemma; its path is an embedding $L\hookrightarrow M$.
This is a valid transfer of finite consistency to a global object.

\begin{proposition}[Finite ages do not provide that bound]
\label{p06:lproof:finiteage}
Every finite suborder of $\omega^*$ embeds in $\omega$, but $\omega^*$ does not
embed in $\omega$.  In particular, equality of the finite ages of two countable
linear orders does not imply that either embeds in the other.
\end{proposition}
\begin{proof}
Every finite linear order embeds in both orders.  An embedding of $\omega^*$
in $\omega$ would give an infinite strictly descending sequence of natural
numbers.  Starting with $a_0$, induction gives $a_n\leq a_0-n$ for
$n\leq a_0+1$, a contradiction.  Similarly an embedding of $\omega$ in
$\omega^*$ would give a descending sequence in the usual natural-number order.
The tree of finite descending sequences in $\omega$ has a node at each level
but no infinite path; it has no uniform bound on its first coordinate.
\end{proof}

\paragraph{Attempt and obstruction.}
Finite observation profiles of orders, or finite embeddings recovered from
those profiles, do not produce the coherent bounds in
Lemma~\ref{p06:lproof:boundedtree}.  Nor does that lemma choose a good pair from an
arbitrary sequence of orders.  The explicit pair $\omega^*,\omega$ refutes the
first shortcut before any reverse-mathematical issue is reached.  Obtaining a
good pair and its embedding using only arithmetical transfinite recursion
remains the missing step; neither is furnished by the earlier analytic
identifiability results.

\subsection{Item 17: the exact conservation problem for Ramsey's theorem}
\label{p06:lproof:ramsey}

\paragraph{Exact target and current evidence.}
Write $B=\mathsf{RCA}_0+\mathsf{B}\Sigma^0_2$.  The target is
\[
 B+\mathsf{RT}^2_2\vdash\varphi\quad\Longrightarrow\quad B\vdash\varphi
 \qquad(\varphi\in\Pi^1_1),
\]
where $\mathsf{RT}^2_2$ concerns two-colourings of unordered pairs and
$\Pi^1_1$ permits a universal set quantifier followed by an arbitrary
arithmetical formula.  Le Hou\'erou--Levy Patey--Yokoyama, corrected version~3
of 6 May 2026, Theorem~1.4, prove $\forall\Pi^0_4$ conservativity, even with
$\mathsf{WKL}_0$ added.  Their Theorem~1.2 reports the exact equivalence of the
full target to $\forall\Pi^0_5$ conservativity; Question~1.1 states the full
problem \cite{p06:lproof:LLPY}.  The correction in version~3 is material.

An author-posted manuscript by Hank Yu-Shiang Hwang, dated 3 October 2026,
claims the full result \cite{p06:lproof:HwangClaim}.  We located its full text and
its binary-partition/model-extension mechanism, but have not independently
verified the complete proof.  Thus the status recorded here is
\emph{a verified partial theorem plus a recent unverified full-resolution
claim}, not an assertion that no such claim exists, and not a claimed solution
by this manuscript.

\begin{proposition}[A sufficient model-extension criterion]
\label{p06:lproof:extension}
Let $B\subseteq T$ be countable theories of second-order arithmetic interpreted
as two-sorted first-order theories.  Suppose every countable model
$(M,\mathcal S)\models B$ has an extension
$(M,\mathcal S')\models T$ with $\mathcal S\subseteq\mathcal S'$ and with the
same first-order structure and the same membership relation on old sets.
Then $T$ is $\Pi^1_1$ conservative over $B$.
\end{proposition}
\begin{proof}
Suppose $B\nvdash\forall X\,\theta(X)$, where $\theta$ is arithmetical.
By completeness and the downward L\"owenheim--Skolem theorem there is a
countable $(M,\mathcal S)\models B$ and $A\in\mathcal S$ satisfying
$\neg\theta(A)$.  Extend this model by the hypothesis.  Induction on the
arithmetical formula shows that its truth with parameter $A$ is unchanged:
all number quantifiers have the same domain and all atomic relations are the
same.  Hence the extension is a model of $T+\neg\theta(A)$, so
$T\nvdash\forall X\theta(X)$.  Several universal set variables can be treated
in the same way, or coded into one set.
\end{proof}

\paragraph{Attempt and obstruction.}
Adding homogeneous sets in the standard natural numbers is not the hypothesis
of Proposition~\ref{p06:lproof:extension}: the relevant countable models may have
nonstandard first-order parts.  Their induction and bounding schemes must
survive the extension.  This is a genuine distinction.  If a consistent
recursively axiomatized $B$ is strong enough for the second incompleteness
theorem and has an $\omega$-model, every $\omega$-model satisfies the true
arithmetical sentence $\operatorname{Con}(B)$, while $B$ does not prove that
sentence.  Consequently, a proof valid in all standard-domain models is not a
syntactic conservation proof.  Our analytic constructions supply no
nonstandard-model preservation theorem, and no $\forall\Pi^0_5$ conservation
argument.  Auditing the new claim would specifically require checking the
coded-history, jump-regularity, and bounding-preservation steps rather than
replacing them by standard compactness.

\subsection{Item 24: Hindman's theorem and the iteration obstruction}
\label{p06:lproof:hindman}

\paragraph{Exact target and verified updates.}
For every $k$ and $c:\mathbb N\to k$, full $\mathsf{HT}$ asks for an infinite
$H$ such that all nonempty sums of \emph{distinct} members of $H$ have one colour.
The reverse-mathematical question is its exact strength over
$\mathsf{RCA}_0$, between $\mathsf{ACA}_0$ and $\mathsf{ACA}_0^+$, where the
latter asserts existence of each $\omega$th Turing jump.  Exactly-two-sum,
at-most-two-sum, and iterated versions must not be interchanged.  For example,
the two-colour exactly-two-sum version follows from $\mathsf{RT}^2_2$, whereas full
$\mathsf{HT}$ implies $\mathsf{ACA}_0$ \cite{p06:lproof:Csima}.

The 2026 preprint of Liu--Patey, Main Theorem~1.5, proves that any arithmetic
instance has a solution avoiding a prescribed non-arithmetic Turing cone.
Its introduction explicitly explains why this one-application result does not
separate $\mathsf{HT}$ from $\mathsf{ACA}_0^+$ \cite{p06:lproof:LiuPatey}.
A new negative result of Liao gives a computable colouring with no infinite
$\Pi^0_3$ solution \cite{p06:lproof:Liao}.  Neither result supplies the exact
reverse-mathematical classification.

\begin{proposition}[A bounded finite-sums transfer]
\label{p06:lproof:boundedht}
Over $\mathsf{WKL}_0$, fix $c:\mathbb N\to k$ and $b:\mathbb N\to\mathbb N$.
Suppose that for each $n$ there are
$x_0<\cdots<x_{n-1}$ with $x_i\leq b(i)$ whose nonempty finite sums all have
the same colour.  Then $c$ has an infinite Hindman solution.
\end{proposition}
\begin{proof}
The strictly increasing tuples with monochromatic finite sums form a
prefix-closed tree.  Checking its membership is a finite computation relative
to $c$.  The bounds $b(i)$ and Lemma~\ref{p06:lproof:boundedtree} give an infinite
path $(x_i)$.  Every two nonempty finite sums occur in one sufficiently long
prefix of this path, so have the same colour.  Its range is the required
infinite set, which exists in $\mathsf{RCA}_0$ for a strictly increasing
sequence.
\end{proof}

The missing bound is a single $b$ valid simultaneously for arbitrarily long
solutions; bounds that grow with the target length do not suffice.  Thus
finite-sums compactness does not resolve the infinite theorem in a weak base.
The recent cone-avoidance theorem suggests a sharper attempted route.
Write $C\leq_{\mathrm{ar}} Z$ if $C\leq_T Z^{(n)}$ for some $n<\omega$.

\begin{proposition}[The iteration hypothesis that would separate the upper bound]
\label{p06:lproof:aravoid}
Let $P$ be a second-order combinatorial principle with sets coding its
instances and solutions and with arithmetical instance and solution
predicates, as for Hindman's theorem.  These predicates are absolute to
the $\omega$-models used below.  Assume the following
relativized arithmetic-cone avoidance property:
\[
 \begin{gathered}
 C\not\leq_{\mathrm{ar}}Z,\quad I\leq_T Z,\quad I\text{ a }P\text{-instance}
 \\
 \Longrightarrow\quad
 \exists H\,[H\text{ solves }I\ \land\ C\not\leq_{\mathrm{ar}}Z\oplus H].
 \end{gathered}                                      \tag{AC}
\]
Then for every non-arithmetic $C$ there is an $\omega$-model of
$\mathsf{ACA}_0+P$ not containing $C$.  In particular, if (AC) holds for
$P=\mathsf{HT}$, then $\mathsf{ACA}_0+\mathsf{HT}$ does not imply
$\mathsf{ACA}_0^+$.
\end{proposition}
\begin{proof}
Start with $Z_0=\varnothing$.  Build an increasing Turing sequence $(Z_s)$,
maintaining $C\not\leq_{\mathrm{ar}}Z_s$.  Schedule all pairs $(n,e)$, with
$n\leq s$, so every pair is eventually treated after $Z_n$ has been defined.
If the $e$th computation relative to $Z_n$ is a code for a $P$-instance $I$,
it is computable from $Z_s$.  By (AC) choose a solution $H$ with
$C\not\leq_{\mathrm{ar}}Z_s\oplus H$, and put
$Z_{s+1}=(Z_s\oplus H)'$.  Otherwise put $Z_{s+1}=Z_s'$.
Taking a jump preserves the invariant, since being arithmetic in a finite
jump of an oracle is equivalent to being arithmetic in that oracle.
This construction takes place in the metatheory; deciding whether a
computation codes a valid instance is not asserted to be effective.

Let $\mathcal I=\{X:(\exists s)X\leq_T Z_s\}$.  It is a Turing ideal, is
closed under finite joins and under the Turing jump, and hence
$(\mathbb N,\mathcal I)\models\mathsf{ACA}_0$.  Every instance in
$\mathcal I$ is eventually scheduled and has its chosen solution in
$\mathcal I$, giving $P$.  The invariant excludes $C$ from $\mathcal I$.
Finally choose $C=\varnothing^{(\omega)}$, which is not arithmetic.  Any
$\omega$-model of $\mathsf{ACA}_0^+$ contains this set, so the constructed model
separates the theories as asserted.
\end{proof}

\paragraph{Attempt and obstruction.}
This is the usual ideal-construction argument with its necessary invariant
written explicitly; no novelty is claimed for that method.  Liu--Patey's
conclusion is Turing nonreducibility, not (AC).  From $C\not\leq_T H$ one
cannot infer $C\not\leq_{\mathrm{ar}}H$: take any $H$ and $C=H'$.  Thus
closing the constructed ideal under jumps can immediately destroy the weaker
invariant.  Proving (AC) for full Hindman solutions, or finding another
preserved invariant strong enough for iteration, is the exact missing step in
this attempted separation.  The earlier analytic inverse results do not
supply it.

\subsection[Item 25: ordinal analysis of Pi12-comprehension]{Item 25: what ``ordinal analysis of $\Pi^1_2$-comprehension'' means}
\label{p06:lproof:ordinal}

\paragraph{Formalization and a necessary status correction.}
The exact theory is $T=\Pi^1_2\text{-}\mathsf{CA}_0$, with arbitrary set
parameters and its specified induction axioms.  Its proof-theoretic ordinal
may be defined as the supremum of order types of recursive orders whose
well-foundedness $T$ proves.  An ordinal \emph{analysis} normally demands an
explicit effective notation system, proof reduction/cut elimination giving
an upper bound, and matching well-ordering or reflection results giving the
lower bound.  Merely naming the supremum is not such an analysis.

Rathjen's 1995 paper announced analyses around $\Pi^1_2$-comprehension
\cite{p06:lproof:Rathjen95}; his 2005 complete paper analyzes a
\emph{parameter-free} $\Pi^1_2$ scheme over $\Delta^1_2$-comprehension plus bar
induction, not automatically $T$ \cite{p06:lproof:Rathjen05}.  More recently,
Arai's preprint, version~4 of 12 August 2025, Theorems 1.1--1.3, gives explicit
ordinal analyses for $\Pi_N$-Collection and associated systems including
$\Sigma^1_{N+2}$-dependent choice plus bar induction.  Theorem~1.4 treats
$Z_2$ in the author's convention, which includes dependent choice
\cite{p06:lproof:Arai}.  Accordingly, the broad assertion that analysis has not
reached beyond $\Pi^1_2$ is not an accurate description of the current
literature.  These stronger-theory results do not, just by inclusion, identify
the exact ordinal of the weaker, precisely specified $T$.  We do not certify
that exact calibration here and do not count this status correction as our
solution.

\begin{proposition}[Two limits on an ordinal shortcut]
\label{p06:lproof:ordinalbarrier}
If $T\subseteq S$, then $|T|\leq |S|$ for the preceding definition of ordinal.
An upper analysis of $S$ alone supplies no equality.  Furthermore, let $T$ be a
consistent recursively axiomatized arithmetic theory to which the second
incompleteness theorem applies.  If $T$ proves
$\operatorname{WF}(O)\rightarrow\operatorname{Con}(T)$ for a coded order $O$,
it cannot also prove $\operatorname{WF}(O)$.
\end{proposition}
\begin{proof}
Every recursive order whose well-foundedness is proved by $T$ is also proved
well-founded by $S$, yielding the inclusion of the two sets of order types and
hence the supremum inequality.  It contains no reverse inclusion.  For the
second assertion, the two supposed proofs would, by modus ponens, give a
$T$-proof of $\operatorname{Con}(T)$, contradicting the second incompleteness
theorem.
\end{proof}

\paragraph{Attempt and obstruction.}
A transfinite rank on the earlier inverse or extraction constructions ranks
those constructions, not arbitrary proofs with $\Pi^1_2$ comprehension.
An ordinal analysis would require a translation of the latter proofs and
verification that each reduction decreases the assigned notation, followed
by a cofinal lower-bound argument.  None has been obtained.  Arai's results
provide relevant existing upper technology, but a claim of equality for $T$
requires matching the exact axiom system and the lower-bound argument.

\subsection{Item 26: Friedman's publication-scoped EFA conjecture}
\label{p06:lproof:efa}

\paragraph{Which conjecture?}
Friedman's FOM posting of 16 April 1999 states two ``grand conjectures''.
The EFA one concerns arithmetical \emph{statements} of theorems published in the
\emph{Annals of Mathematics}; the other concerns eliminating categorical
methods.  Here EFA is $I\Delta_0+\exp$.  The former is not the claim that all
true arithmetic is provable in EFA \cite{p06:lproof:Friedman99}.  In his own 2010
clarification Friedman emphasizes ordinary core mathematics and the need for
an actual arithmetical theorem, rather than a strong intermediate lemma in
one proof \cite{p06:lproof:Friedman10}.

For a precise corpus version one must fix a publication cutoff, a list
$\mathcal A$ of included theorem statements, and their chosen arithmetizations
$\widehat\theta$.  The resulting metamathematical assertion is
\[
       \forall\theta\in\mathcal A\qquad
       \mathsf{EFA}\vdash\widehat\theta.
\]
No cutoff or exclusions are specified in the catalogue label.  We use this
publication-scoped interpretation and record its unresolved/underspecified
status.  An overbroad replacement by ``every true arithmetical sentence'' is
false, assuming EFA is consistent, by its own consistency sentence and the
second incompleteness theorem.  That observation does not refute Friedman's
actual corpus claim.

\begin{proposition}[What an elementary witness bound really permits]
\label{p06:lproof:boundedsearch}
Let $R(x,y)$ be bounded in the language of EFA and let $t(x)$ be an elementary
term (allowing a definitional extension by elementary functions).  If
\[
 \mathsf{EFA}\vdash\forall x\,\exists y\leq t(x)\ R(x,y),
\]
then bounded search gives an elementary function $f$ and EFA proves
$R(x,f(x))$ for every $x$.
\end{proposition}
\begin{proof}
Evaluate the bounded predicate successively at $y=0,\ldots,t(x)$ and return
the first successful $y$, using a default value if none exists.  Bounded
quantification and bounded search over elementary bounds preserve elementary
computability: the bounded truth value can be computed by finite products and
sums of $0$--$1$ values.  Bounded induction proves that the recorded first
success is at most $t(x)$ and satisfies $R$, whenever there is a success.
The assumed EFA theorem eliminates the default case.
\end{proof}

\paragraph{Attempt and obstruction.}
Replacing analytic existence by an explicit search in the earlier work would
help only after both the bound and its correctness have been proved in EFA.
An externally true estimate does not give that EFA proof.  For a
$\Pi^0_1$ statement such as Fermat's Last Theorem, bounding individual finite
checks is likewise not a proof of their universal closure.  No identified
Annals theorem has been reduced to Proposition~\ref{p06:lproof:boundedsearch}
in this work, so no new case of the grand conjecture is claimed.

\subsection{Item 47: weak versus strong normalization of pure type systems}
\label{p06:lproof:bgk}

\paragraph{Exact target and verified status.}
A pure type system is specified by sorts, axioms between sorts, and product
formation rules.  Weak normalization means that \emph{every typable term}
has some finite $\beta$-reduction to normal form; strong normalization means
that no typable term has an infinite reduction sequence.  The
Barendregt--Geuvers--Klop conjecture is that weak normalization of a whole PTS
implies strong normalization of that PTS.  It is not the false assertion that
each individually weakly normalizing untyped term is strongly normalizing.
The problem-originator's TLCA entry records the conjecture and the
non-dependent partial result of Barthe--Hatcliff--S\o rensen
\cite{p06:lproof:BGK,p06:lproof:BHS}.  No unrestricted resolution was verified in the
current search.

\begin{proposition}[Confluence and unique normal forms do not close the gap]
\label{p06:lproof:ars}
There is a finitely branching, acyclic, confluent reduction relation in which
every object is weakly normalizing and has the same unique normal form, but
which is not strongly normalizing.
\end{proposition}
\begin{proof}
Take objects $z,a_0,a_1,\ldots$ and rules
\[
              a_n\longrightarrow z,
              \qquad a_n\longrightarrow a_{n+1}\quad(n\geq0).
\]
The only normal form is $z$, and every $a_n$ reaches it in one step.  Every
two reducts of a common object reduce to $z$, which proves confluence.
Each object has at most two immediate reducts, and no path returns to an
earlier object.  Nevertheless $a_0\to a_1\to a_2\to\cdots$ is infinite.
\end{proof}

\paragraph{Attempt and obstruction.}
A unique analytic terminal state or a terminating chosen reduction strategy
cannot establish termination of all reductions: the proposition is an exact
counterexample to that abstract argument.  It is \emph{not} a counterexample
PTS.  To use the non-dependent positive results for an arbitrary PTS one
would need a translation eliminating dependency, preserving the global weak
normalization hypothesis, and reflecting strong normalization.  The
substitution and dependent-product steps of such a translation are the
unproved requirements.  A well-founded rank decreasing along every typed
reduction would also suffice, but the ranks in the earlier analytic work do
not rank these reductions.

\subsection{Item 54: expansion postponement and the solved normalizing case}
\label{p06:lproof:expansion}

\paragraph{Exact statement and verified correction.}
For a PTS, write $\vdash_r$ for the typing relation in which conversion is
restricted to type reduction, with type expansion removed.  Expansion
postponement asks whether
\[
 \Gamma\vdash a:A\quad\Longrightarrow\quad
 \exists A'\,[\Gamma\vdash_r a:A'\ \land\ A\longrightarrow_\beta^*A'].
                                                               \tag{EP}
\]
The final restoration of the type $A$ can then be done by expansion.  Poll's
1998 paper proves this for every PTS whose typable types have normal forms,
which is weaker than requiring strong normalization of all terms.  Precisely,
Theorem~12 gives $\Gamma\vdash_{nf}a:\operatorname{nf}(A)$ and Theorem~8 embeds
$\vdash_{nf}$ in $\vdash_r$; Corollary~13 is expansion postponement
\cite{p06:lproof:Poll}.  Thus the normalizing-types case is already solved and
cannot be counted as a new result.  The unrestricted all-PTS version is the
remaining target; no verified unrestricted resolution was found here.

\paragraph{Attempted extension and exact failed inference.}
A natural route takes a conversion $A=_\beta B$, uses confluence to choose a
common reduct $C$, and tries to replace each conversion in a typing derivation
by reductions to $C$.  The required derivation is not supplied by the term
rewrite diagram: an abstraction premise has a dependent product type, and
substitution changes the codomain.  Reconstructing a reduction-only typing
derivation requires a subject-reduction or substitution statement for the
modified typing system, not just confluence of $\beta$-reduction.
Poll explicitly distinguishes weak subject reduction for $\vdash_r$ from the
ordinary subject-reduction statement, which does not hold for that relation
in general.  His normal-form intermediate relation solves this reconstruction
problem using existence of normal forms.  If types need not normalize,
$\operatorname{nf}(A)$ is undefined and Theorem~12 cannot be applied.
This identifies the failed proof step, without asserting an unproved
counterexample to (EP).

\paragraph{Relation to item 47.}
Even a global proof of weak normalization for a PTS would already put its types
in Poll's solved class; no appeal to strong normalization is needed for that
conclusion.  Conversely, eliminating expansion in typing judgments says
nothing by itself about whether an arbitrary reduction path terminates.
The two catalogue items are distinct and are not two consequences of one
abstract uniqueness theorem.

\subsection{Item 55: full graph minors, bounded treewidth, and finite size}
\label{p06:lproof:graphminor}

\paragraph{Exact target and verified bounds.}
The graph minor theorem (GMT) says that every infinite sequence of finite
graphs has $i<j$ with $G_i$ a minor of $G_j$.  The open target is its exact
proof-theoretic strength, not its truth.  Krombholz--Rathjen's 2019 preprint
(published 2020), Section~1 and Sections~2--3, gives the upper theory
\[
 \Pi^1_1\text{-}\mathsf{CA}_0+\Pi^1_3\text{-}\mathsf{IND}
       +\Pi^1_2\text{-}\mathsf{BI}.
\]
Their Theorem~5 records equivalences for the \emph{uniform bounded-treewidth}
version, including extended Kruskal, planar GMT, and well-ordering of the indicated
collapsing-function notation; Theorem~12 gives an explicit ordinal lower
bound for extended Kruskal \cite{p06:lproof:KR}.  Each fixed treewidth instance and
the universal statement over all bounds must be distinguished.  These upper
and lower results do not match into an exact full-GMT calibration; no verified
later exact calibration was found in our search.

\begin{proposition}[Finite vertex bounds have an elementary good-pair bound]
\label{p06:lproof:finitegraphs}
For finite simple graphs with at most $N$ vertices, every sequence of length
\[
              K_N+1,\qquad
              K_N=\sum_{r=0}^{N}2^{\binom r2},
\]
contains $i<j$ such that $G_i$ is isomorphic to $G_j$, hence is a minor of
$G_j$.  The finite statement and its proof can be formalized in EFA with the
usual finite-graph coding.
\end{proposition}
\begin{proof}
Relabel each graph on $r$ vertices by $\{0,\ldots,r-1\}$, choosing, for example,
the least adjacency-matrix code among all relabellings.  There are at most
$2^{\binom r2}$ codes for this $r$, hence at most $K_N$ in all.  Two of the
$K_N+1$ codes coincide by the finite pigeonhole principle.  Equal codes give
an isomorphism, and every graph is a minor of an isomorphic graph.
All searches, relabellings, and finite sums here have elementary bounds; for
example $K_N\leq(N+1)2^{N^2}$.  Bounded induction proves the finite pigeonhole
principle by induction on the size of the finite codomain, deleting one
value and, if necessary, its unique preimage.  This gives an EFA proof for
coded finite sequences.
\end{proof}

\paragraph{Attempt and obstruction.}
Finite observational recovery of graphs can be combined with this proposition
only when it bounds their number of vertices.  Bounded treewidth does not do
so: the paths have treewidth at most one and arbitrarily many vertices.
Thus a finite list of recovered graph types is not a rank for a general bad
minor sequence.  The strong uniform bounded-treewidth theorem and the full
GMT require well-quasi-order arguments beyond finite pigeonhole.  No reduction
of those arguments to the earlier finite identification bounds has been
constructed.  The elementary finite-size proposition is a classical special
case with an explicit bound, not a new solution of the proof-strength problem.

\subsection{Outcome of the nine-target attack}
\label{p06:lproof:outcome}

The strongest concrete transfer developed in this section is the exact
arithmetic-cone invariant in Proposition~\ref{p06:lproof:aravoid}; the missing
hypothesis is exposed rather than assumed.  The finite-age and abstract
reduction counterexamples decisively invalidate two proposed shortcuts.
The other positive statements isolate the quantitative or formalization
hypotheses needed to use analytic results in logic.  These are useful proof
checks, but they provide no new unrestricted resolution of the nine named
problems.  In particular, neither recent external claims nor already solved
restricted cases are added to a count of newly closed conjectures.


\section{Finite structures, universality, and effective witnesses}
\label{p06:lfin:section}

This section investigates items 34, 35, 36, 38, 40, 42, 48, 53, and 56.
The status cutoff is 17 October 2026. None of the nine unrestricted
problems is claimed solved by this section. The propositions below are
proved reductions, elementary certificates, or obstructions to proposed
transfers; they are not nine new resolutions. Classical constructions and
recent external advances are explicitly separated from our deductions.
In particular, an external 2026 star-height preprint changes the status of
one boundedness variant, but does not settle the height-one question.

\subsection{34: universality with forbidden ordinary subgraphs}
\label{p06:lfin:34}

Fix a finite family $\mathcal C$ of finite simple graphs. A graph is
$\mathcal C$-free when it contains no member of $\mathcal C$ as an ordinary
subgraph, not necessarily an induced subgraph. The question is whether
there exists a countable $\mathcal C$-free graph into which every countable
$\mathcal C$-free graph embeds as an \emph{induced} subgraph. This is strong
universality. Replacing the forbidden relation by induced subgraphs or
minors changes the problem.

Cherlin--Shelah--Shi develop the algebraic-closure approach for finite
families of connected constraints \cite{p06:lfin:css}. A substantial solved
case is Cherlin--Shelah's Tree Theorem: for one finite tree $T$, weak and
strong countable universality are equivalent, and hold precisely when
$T$ is a path or a near-path (a path with one extra pendant edge)
\cite[Theorem 1]{p06:lfin:tree}. This is not a classification of arbitrary
finite families. No such full classification was verified in our search.

The tempting transfer is to build a universal object by unrestricted
free amalgamation. The following exact obstruction shows its scope.

\begin{proposition}[Exact scope of free amalgamation]\label{p06:lfin:free}
Remove redundant members of $\mathcal C$ so that it is an antichain under
ordinary-subgraph containment. Assume its members have at least two
vertices. The class of finite $\mathcal C$-free graphs is closed under
free amalgamation over induced subgraphs if and only if every member of
$\mathcal C$ is complete. The empty family also has this property.
\end{proposition}
\begin{proof}
A free amalgam of $B$ and $D$ over $A$ has no edges between $B\setminus A$
and $D\setminus A$. A complete forbidden graph cannot meet both these
sets, so if all constraints are complete the amalgam remains free of them.
Conversely, let $F\in\mathcal C$ have nonadjacent vertices $u,v$. Set
$A=F\setminus\{u,v\}$, $B=F\setminus\{v\}$, and
$D=F\setminus\{u\}$. Both $B,D$ are $\mathcal C$-free: any constraint
occurring in either would be a subgraph of $F$, contrary to minimality
and the smaller number of vertices. Their free amalgam is exactly $F$.
\end{proof}

For complete constraints, the usual iterative extension construction now
gives a countable universal graph. Start with the empty graph and
successively realize every finite allowed one-point extension of every
finite subgraph already constructed, repeatedly scheduling all requests.
Free amalgamation preserves the constraint. A countable allowed graph
embeds by enumerating its vertices and extending its embedding one vertex
at a time. This proves the claimed restricted construction without
assuming the desired universal object. It does not prove nonuniversality
when free amalgamation fails: the near-path cases in the Tree Theorem
are explicit counterexamples to that inference. The missing ingredient
is control of the nontrivial algebraic closure generated by amalgamation.

\subsection{35: which universality spectrum problem?}
\label{p06:lfin:35}

For a complete theory $T$ in a countable language with an infinite model,
put
\[
 \operatorname{Univ}_{\preceq}(T)=
 \{\kappa\geq\aleph_0:\text{some }M\models T,
 |M|=\kappa,\text{ elementarily contains every }N\models T
 \text{ of size }\kappa\}.
\]
We interpret ``minimum spectrum'' as a \emph{least set under inclusion}
among these spectra, in a fixed set-theoretic universe. Eventual
inclusion modulo initial segments and comparison in forcing extensions
are separate variants. Allowing complete theories with only finite
models would give an uninteresting empty minimum. Shelah formulates the
classification problem in Definition 2.1 and Problem 2.2, distinguishes
the smallest-spectrum and eventual variants in Problem 2.3, and warns
that cardinal arithmetic is essential \cite{p06:lfin:shelah}. We have not
verified a solution of the least-spectrum variant just specified.

\begin{proposition}[Countable intersection reduction]\label{p06:lfin:intersection}
For any sequence $(T_i)_{i<\omega}$ of complete countable theories with
infinite models there is a complete countable theory $T_\oplus$, with
infinite models, such that for every infinite cardinal $\kappa$,
\[
 \operatorname{Univ}_{\preceq}(T_\oplus)
       =\bigcap_{i<\omega}\operatorname{Univ}_{\preceq}(T_i).
\]
Consequently, if a countable family of these spectra is coinitial under
inclusion among all such spectra, a least spectrum exists.
\end{proposition}
\begin{proof}
First use the disjoint many-sorted sum, with one infinite sort for each
$T_i$ and no symbols connecting sorts. Its sentences reduce to Boolean
combinations of sentences in finitely many sorts, so it is complete.
If $U$ is universal of total size $\kappa$, testing it against a model
whose $i$th sort has size $\kappa$ shows that $|U_i|=\kappa$. Testing
arbitrary $\kappa$-sized models of $T_i$ in that sort shows that $U_i$ is
universal. Conversely, choose universal $U_i$ of size $\kappa$ for all
$i$. Any sort of any model of total size $\kappa$ has size at most
$\kappa$ and has an elementary extension of size $\kappa$, by upward
L\"owenheim--Skolem. It therefore embeds elementarily into $U_i$.
The union of these maps is elementary in the disjoint sum, whose total
size is $\aleph_0\kappa=\kappa$.

For completeness, this does not require enlarging the usual one-sorted
class of theories. Relationalize the $T_i$, use disjoint unary predicates
$P_i$, relativize $T_i$ to $P_i$, and require each relation belonging to
$T_i$ to be false off $P_i$. No infinite disjunction asserting that the
$P_i$ cover the universe is imposed. Take the complete theory of such a
sum with an additional infinite set of points outside all $P_i$.
Completeness and elementarity follow from the following direct formula
decomposition: a formula mentions only finitely many $P_i$ and component
languages; partition its quantified variables among those sorts and
their complement. Relations with mixed sorts are false, equalities
between different sorts are false, and on the complement only equality
remains. The complement is infinite, because unused infinite sorts
remain. Induction on formulas reduces quantification to formulas in the
individual components and equality in an infinite set. This also shows
that the points outside every $P_i$ can be mapped by any injection into
$\kappa$ additional such points. Thus the same equivalence holds for the
one-sorted complete theory.

If $(T_i)$ is coinitial, its intersection is contained in every spectrum
and is itself a spectrum by the construction. It is therefore least.
\end{proof}

This reduction identifies a concrete missing statement: a countable
coinitial family, or another countable-language coding of the full
intersection. There are continuum many countable theories, and taking
one named sort for every one of them does not preserve a countable
language. The proposition supplies no such compression and does not
remove forcing dependence.

\subsection{36: generalized star height and the 2026 boundedness claim}
\label{p06:lfin:36}

A generalized regular expression uses letters, $\varnothing$,
$\{\varepsilon\}$, union, concatenation, complement in the same free
monoid $\Sigma^*$, and Kleene star. Height counts nested stars;
complement does not increase it. The classical unresolved question is
whether some regular language has generalized height at least two,
equivalently whether height one always suffices. A uniform finite upper
bound, an algorithm computing the exact height, and this height-one
question are different assertions. Every individual regular language
having \emph{some} finite height is already immediate from ordinary
regular expressions.

OpenAI's \emph{Generalized Star Height at Most Three}, dated 25 September
2026, states the uniform bound $h_\Sigma(L)\leq3$ in Theorem 1.1. Its
Section 1.1 explicitly leaves the height-one question unanswered
\cite{p06:lfin:star2026}. The repository also supplies a Lean scope document
asserting formalization of the same top-line bound
\cite{p06:lfin:starlean}. We checked the statement, the distinction of
variants, and the periodic-template step below; we did not rebuild the
Lean development or independently audit its entire 38-page construction.
Accordingly this is recorded as an attributed 2026 boundedness result
with a specified verification scope, not as our theorem or as a verified
resolution of the classical height-one question.

\begin{proposition}[A fully checked template transfer]\label{p06:lfin:template}
Let $L$ be recognized by a deterministic automaton with $n$ states. For
fixed words $u,v,w$, with $v\ne\varepsilon$, the language
$L\cap uv^*w$ has an effectively constructible ordinary regular
expression of star height at most one. The same holds with a finite
union of such templates in place of $uv^*w$.
\end{proposition}
\begin{proof}
From the state reached after $u$, repeatedly apply the transition
associated to $v$. A state repeats within $n+1$ terms. Thus there are
integers $0\leq\mu<n$ and $1\leq p\leq n$ such that the state after
$uv^j$, and hence acceptance after the additional suffix $w$, is
periodic with period $p$ for $j\geq\mu$. The accepted exponents are a
finite subset below $\mu$, together with selected progressions
$r+p\mathbb N$ with $\mu\leq r<\mu+p$. Express a progression by
$uv^r(v^p)^*w$ and add the finitely many exceptional words. Finite unions
preserve the bound.
\end{proof}

This verifies the ultimate-periodicity mechanism behind
\cite[Lemma 10.1]{p06:lfin:star2026} once its template hypotheses hold.
The nontrivial additional obligation in that construction is to prove
the independently guessed split tests are sound and their entire
component languages satisfy the template hypotheses (in particular
Proposition 7.1 and Proposition 10.2). We have not replaced that obligation
by the much weaker assertion that the tests are finite-state. Nor does
the proposition force height one for arbitrary languages: even
$\{a,b\}^*$ cannot be covered by finitely many $uv^*w$ templates, since
such a union has only boundedly many words of each fixed length. This
last observation is an obstruction to that cover, not a lower bound on
generalized height, because complement is available.

\subsection{38: Henson graphs and finite extension certificates}
\label{p06:lfin:38}

For $n\geq3$, let $H_n$ be the countable homogeneous universal
$K_n$-free graph. The relevant problem is whether $H_n$ is pseudofinite:
every first-order sentence true in $H_n$ has a finite model. It is
\emph{not} whether its full theory has a finite model; the latter
contains sentences asserting arbitrarily many vertices. The triangle-free
case remains explicitly open in Macpherson--Tent (2025)
\cite[Introduction]{p06:lfin:mt}. Letzter--Sahasrabudhe identify the
already difficult finite $4$-extension case and give quantitative
nonexistence bounds depending on graph order \cite{p06:lfin:ls}.

Call a finite $K_n$-free graph $k$-extendible if, for disjoint $U,V$ with
$|U|+|V|\leq k$ and $U$ containing no $K_{n-1}$, there is a vertex
outside $U\cup V$ adjacent to every vertex of $U$ and to no vertex of $V$.

\begin{proposition}[Exact finite-fragment reduction]\label{p06:lfin:henson}
$H_n$ is pseudofinite if and only if a finite $k$-extendible $K_n$-free
graph exists for every $k$. In the triangle-free case, a $2$-extendible
example exists on ten vertices. Every finite $k$-extendible triangle-free
graph has at least $k+2^k$ vertices.
\end{proposition}
\begin{proof}
The graph axioms, the absence of $K_n$, and all one-point extension
axioms axiomatize $\operatorname{Th}(H_n)$. Indeed they construct a
back-and-forth isomorphism between any two countable models, and they
force infinitude; downward L\"owenheim--Skolem then gives completeness.
For fixed $k$ the relevant extension assertions form a finite conjunction.
Pseudofiniteness supplies a finite model of this conjunction together
with the graph and forbidden-clique axioms. Conversely, every finite
subset of the axiom system is implied by the axioms through some depth
$k$. If all these have finite models, any sentence entailed by the
complete axiom system has a finite model, by compactness applied to its
proof from a finite subset.

For the ten-vertex example take the two-element subsets of a five-element
set, adjacent precisely when disjoint (the Petersen graph). Three
pairwise disjoint two-element sets do not exist, so it is triangle-free.
Every vertex has degree three, and two distinct nonadjacent vertices
have exactly one common neighbor: the complement of their three-element
union. This handles $|U|=2$. If $|U|=|V|=1$, adjacent $u,v$ have no
common neighbor and $u$ has two neighbors other than $v$; nonadjacent
$u,v$ have just one common neighbor, leaving two usable neighbors of
$u$. For $U=\varnothing$, the union of the closed neighborhoods of two
vertices has size at most eight, so a suitable vertex exists. Smaller
requests follow from these cases or directly from degree three.

Finally build an independent set $A$ of size $k$ by repeatedly applying
extension with $U=\varnothing$. For every subset $B\subseteq A$, the
request $(U,V)=(B,A\setminus B)$ has a witness outside $A$. Two
different subsets require different witnesses because their adjacency
patterns on $A$ differ. Thus there are at least $2^k$ vertices outside
$A$.
\end{proof}

The exponential lower bound is an elementary check, weaker than the
published quantitative bound in \cite{p06:lfin:ls}; it is not a new best
bound. The exact obstruction to a closure is producing finite witnesses
at every depth, already beyond the checked depth two. An infinite
Fra\"iss\'e construction realizes the requests in new vertices forever
and does not supply finite witnesses or pseudofiniteness.

\subsection{40: transexponential growth and a failed spectral transfer}
\label{p06:lfin:40}

Here ``superexponential'' must mean eventually larger than every fixed
iterate $\exp^{\circ m}$, rather than merely larger than $\exp$.
The latter is already witnessed by $\exp\circ\exp$ in the o-minimal
real exponential field. The intended question asks for an o-minimal
expansion of an ordered exponential field which is not exponentially
bounded. Krapp--Kuhlmann state it as Question 2.3 and prove that their
particular axiomatized transexponential fields are \emph{not}
o-minimal (Theorem 3.6) \cite{p06:lfin:kk}. Their existence theorem for
transexponential fields therefore does not answer this question.
A contemporaneous Palomar index advertises a result titled
\emph{A transexponential o-minimal structure} \cite{p06:lfin:palomar}.
We could retrieve the indexed description but not its underlying proof
project or exact checked statement. This competing claim is recorded
as unverified here; the older open-status references alone do not settle
its validity or establish the unrestricted problem's October 2026 status.

Our preceding power-mixture optimization suggests a precise transfer to
test, including continuous spectra. It fails for a quantitative reason.

\begin{proposition}[Compact positive spectra cannot supply the growth]
\label{p06:lfin:growth}
Let $\mu$ be a nonzero finite positive measure supported on
$[a,b]\subset(0,1)$, let $M=\mu([a,b])$, and put
\[
 S(x)=\int x^r\,d\mu(r),\qquad
 T(C)=\min_{x>0}\frac{1+Cx}{S(x)}\quad(C>0).
\]
The minimizer $x_C$ is unique. With $A=a/(1-a)$ and $B=b/(1-b)$,
\[
 A\leq Cx_C\leq B,
 \qquad
 \frac{C^a}{MB^a}\leq T(C)
       \leq\frac{1+B}{MA^b}C^b\quad(C\geq B).
\]
For finite support, $T$ is definable, with parameters, in
$\mathbb R_{\exp}$.
\end{proposition}
\begin{proof}
Differentiation under the integral is valid locally on $x>0$, and
$S'>0$, $S''<0$. The objective tends to infinity at both endpoints, so
it has a minimum. Its derivative has the sign of
$N(x)=CS(x)-(1+Cx)S'(x)$, and
$N'(x)=-(1+Cx)S''(x)>0$. Hence there is exactly one minimizer.
Writing $m=xS'(x)/S(x)$, a weighted average of the exponents, stationarity
gives $Cx=m/(1-m)$ and the first bounds follow. For $C\geq B$ one has
$x_C\leq1$, so $Mx_C^b\leq S(x_C)\leq Mx_C^a$. Substituting
$A\leq Cx_C\leq B$ gives both displayed inequalities.
For finitely many atoms, $x^r=\exp(r\log x)$ is definable in
$\mathbb R_{\exp}$, as are finite sums and the quantified condition
that $x$ minimizes the objective. Uniqueness makes this a definable
function, and hence makes $T$ definable.
\end{proof}

Thus neither non-D-finiteness nor differential-algebraic complexity of
these optimum maps implies transexponential growth. Even arbitrary
continuous mixtures in this compact positive regime give polynomially
bounded values. Definability of an arbitrary continuous mixture is not
asserted. A successful transfer would have to leave this regime and
still prove o-minimality, the missing structural requirement.

\subsection{42: characteristic-zero minimal fields}
\label{p06:lfin:42}

Podewski's conjecture says that every infinite minimal field is
algebraically closed. Minimality means that every parameter-definable
subset of the field itself is finite or cofinite; it is weaker than
strong minimality of the theory. The positive-characteristic case is
known. Krupi\'nski--Tanovi\'c--Wagner reduce the remaining case to
definable partial orders, and rule out the almost-linear case
\cite[Theorems 1 and 3]{p06:lfin:podewski}. The March 2026 version of
Johnson--Tran--Walsberg--Ye still states the conjecture and explains a
finite-closed-topology obstruction \cite[Question 11.1 and
Conjecture 11.3]{p06:lfin:large}.

The elementary power-map attack stops exactly where arbitrary
polynomial images cease to be subgroups.

\begin{proposition}[Two exact algebraic/topological reductions]
\label{p06:lfin:minimal}
Let $K$ be an infinite minimal field. For every integer $m\geq1$, the power map
$K^\times\to K^\times$, $x\mapsto x^m$, is surjective. Suppose in
addition that $K$ admits a nondiscrete $T_1$ topology in which translations
are homeomorphisms and every nonconstant polynomial image $f(K)$ is
closed. Then $K$ is algebraically closed.
\end{proposition}
\begin{proof}
The power-map image is definable and infinite (each fiber has at most
$m$ points), hence cofinite. It is a subgroup of $K^\times$. If it were
proper, any other coset would be infinite and disjoint from it,
contradicting cofiniteness. This first assertion is the classical
argument of \cite[Lemma 2.1]{p06:lfin:podewski}.

For a nonconstant $f$, the same finite-fiber argument shows that $f(K)$
is infinite and therefore cofinite. Under the additional topological
hypothesis, its complement is finite and open. If nonempty, a $T_1$
space has an open singleton there: remove the other finitely many
closed singletons from that open set. Translations would make every
singleton open, contradicting nondiscreteness. Thus every nonconstant
polynomial is surjective, and in particular has a root at value zero.
\end{proof}

The second part makes the topological step explicit; it does not produce
the required topology. In the finite-closed topology polynomial images
are closed, but excluding discreteness in the perfect characteristic-zero
case is exactly the outstanding Question 11.1 of \cite{p06:lfin:large}.
Replacing this unknown nondiscreteness by ordinary continuity of
polynomials would be invalid: continuity does not make their images
closed. Likewise power-map surjectivity alone has not supplied roots of
arbitrary polynomials in the present argument.

\subsection{48: MELL and the effective branching bound}
\label{p06:lfin:48}

The question is provability of finite sequents in propositional
multiplicative-exponential linear logic, with the standard single pair
of modalities $!,?$ and no additives. Stra\ss burger diagnoses a fatal
gap in an earlier decidability proof; his Theorem 4.4 proves decidability
of \emph{RMELL}, with unrestricted contraction, not MELL
\cite{p06:lfin:mell}. The equivalence with branching vector-addition-system
reachability is the relevant transfer, not the already decidable plain
VASS problem.

Recent work gives a sharper positive boundary. Bizi\`ere and coauthors
prove that a uniform Strahler
bound exists for every BVASS (Corollary 4.4), decide reachability with a
supplied Strahler bound (Section 6), and compute semilinear reachability
sets for two-dimensional BVASS and five-dimensional \emph{one-state}
BVAS (Theorem 7.5). Their Remark 5.7 explains the missing effective
stabilization test; unrestricted reachability remains open
\cite{p06:lfin:bvass2026}. The two-dimensional advance also appears in
\cite{p06:lfin:bvass2025}.

For a finite rooted tree, its Strahler number is zero at a leaf, is
unchanged at a unary node, and at a branching node is the largest child
number, increased by one if that maximum is attained at least twice.

\begin{proposition}[Exact effectivity gap]\label{p06:lfin:strahler}
Using the bounded-reachability algorithm of \cite[Section 6]{p06:lfin:bvass2026},
general BVASS reachability is decidable if and only if there is a total
computable function $b(\mathcal B,c)$ such that every reachable
configuration $c$ has a witnessing run of Strahler number at most
$b(\mathcal B,c)$. A computable uniform bound depending only on
$\mathcal B$ would suffice, but is not asserted by this equivalence.
\end{proposition}
\begin{proof}
Given the bound, apply bounded reachability at that bound. Conversely,
first decide whether $c$ is reachable. If not, output zero. If yes,
enumerate finite labeled run trees until a valid run with target $c$
appears; validity is decidable by checking all local transitions and
nonnegative counters. Output its computable Strahler number. The search
terminates precisely because the decision was positive.
\end{proof}

The distinction between an instance-dependent computable bound and an
existing system-wide finite bound is essential. No quantitative estimate
from our earlier finite spectral reconstruction controls the finite
basis of BVASS runs in the cited well-quasi-order argument. Replacing
that basis by a computed approximation without a stopping certificate
would simply reintroduce the open problem.

\subsection{53: the word problem for the single combinator S}
\label{p06:lfin:53}

Terms here are ground binary application trees with the sole constant
$S$, and the only rewrite rule is
\[
 (((Sx)y)z)\longrightarrow (xz)(yz).
\]
The word problem asks whether two terms are convertible by this rule
and its inverse. It is not the word problem for $SK$, and it is not
merely normalization. Waldmann proves decidability of normalization of
ground $S$-terms \cite{p06:lfin:waldmann}; Manzonetto's 2023 Problem 4
explicitly retains the pure-$S$ word problem as open
\cite[Section 6.4]{p06:lfin:manzonetto}. We found no primary resolution by
the cutoff.

\begin{proposition}[A decidable region and the missing certificate]
\label{p06:lfin:s}
Convertibility is decidable on pairs for which at least one term has a
normal form. More generally, it is semidecidable, and becomes decidable
if a computable bound on the leaf size of a common reduct is supplied
for every convertible input pair.
\end{proposition}
\begin{proof}
The rule is left-linear and has no nonvariable overlap with itself, so
the standard orthogonal-rewriting confluence theorem applies. Thus
convertible terms have a common reduct, and normal forms, when they
exist, are unique. Use Waldmann's decision procedure on both inputs.
If exactly one normalizes they cannot be convertible, since a common
reduct with a normal form forces both inputs to normalize. If both
normalize, enumerate their reductions until their normal forms are
found and compare them.

For arbitrary inputs enumerate pairs of finite reduction paths and
accept when their endpoints coincide. To check a supplied bound $B$,
observe that leaf size never decreases: a redex with subterm leaf sizes
$|x|,|y|,|z|$ changes from $1+|x|+|y|+|z|$ to
$|x|+|y|+2|z|$, a difference $|z|-1\geq0$. Every path to a reduct
of size at most $B$ stays among the finitely many ground binary trees
of that size. Search that finite graph from both inputs and test whether
their reachable sets meet. If a computable valid $B$ is available for
all convertible pairs, this decides convertibility, returning no when
the finite search finds no meeting.
\end{proof}

Conversely a word-problem decision procedure would compute such an
input-dependent bound by accepting negative cases immediately and
enumerating a common reduct on positive cases. Thus the genuine
remaining region consists of pairs of nonnormalizing terms without a
computable common-reduct bound. Decidable normalization or strong
normalization of a different typed calculus does not supply that bound.

\subsection{56: the three-pair Post correspondence problem}
\label{p06:lfin:56}

We use the ordinary finite-word PCP: three indexed pairs
$(u_i,v_i)$ over a finite alphabet, with repetitions of indices allowed,
ask whether some nonempty word $i_1\cdots i_m$ satisfies
$u_{i_1}\cdots u_{i_m}=v_{i_1}\cdots v_{i_m}$. Empty tile words may be
allowed; the witness sequence must be nonempty. The nonerasing variant
is a restriction, not silently identified with it. An injective
fixed-length binary encoding of the alphabet preserves the number of
pairs and all equalities, so alphabet size is not the issue.

Neary's published Theorem 11 proves undecidability for five pairs and
explicitly leaves three and four pairs open \cite{p06:lfin:pcp}. An early
four-pair title found in secondary indexes is not the published theorem.
Recent results for free-group, reverse, or bi-infinite correspondence
problems concern different variants. No primary proof deciding the
ordinary three-pair problem was verified here.

\begin{proposition}[Finite residual certificate]\label{p06:lfin:pcpbound}
For a PCP instance over an alphabet of size $q\geq1$ and $B\geq0$, it
is decidable whether a solution exists whose unmatched suffix after
each tile has length at most $B$. If such a solution exists, one exists
using at most
\[
             N_B=1+2\sum_{j=1}^{B}q^j
\]
tiles. General three-pair PCP is decidable if and only if a total
computable bound $B(I)$ of this kind exists for every positive
three-pair instance $I$.
\end{proposition}
\begin{proof}
After deleting the longest common prefix of two partial concatenations,
either they disagree at their next letters, in which case no extension
can repair them, or one residual is empty. Record the latter case as
$0$, or $(+,r)$, or $(-,r)$, according as both are empty, only the
upper residual is $r$, or only the lower residual is $r$. Keep only
nonempty $r$ of length at most $B$. These are exactly $N_B$ states.
From $(+,r)$ and tile $(u_i,v_i)$ cancel the common prefix of
$(ru_i,v_i)$; from $(-,r)$ cancel that of $(u_i,rv_i)$; from $0$
cancel that of $(u_i,v_i)$. Discard a letter mismatch or a residual
longer than $B$. This is a finite directed labeled graph.

Solutions with the bound are precisely nonempty walks from $0$ back
to $0$. A shortest such walk has no repeated internal vertex and no
intermediate visit to $0$, because removing the corresponding cycle
would shorten it. Its length is at most the number of vertices.
Thus a finite graph search decides the bounded problem and constructs
the stated witness.

A computable bound reduces unrestricted three-pair PCP to this finite
search. Conversely, given a decision procedure, output zero on a
negative instance. On a positive instance enumerate index words until
a solution is found, and output the largest residual length along it.
This terminates and supplies a total computable bound.
\end{proof}

This is an exact bounded-delay algorithm, valid even with empty tile
words, not an algorithm for unrestricted PCP. Every individual solution
has a finite delay, but that observation gives no computable stopping
bound on negative instances. Group cancellation, approximate equality,
or replacing words by commutative letter counts destroys the prefix
information used in the proof and cannot justify such a bound.

\subsection{Outcome of the nine attempts}
\label{p06:lfin:outcome}

The strongest verified new work in this section is the explicit
countable-intersection reduction and the quantified obstruction to
using positive spectral mixtures for transexponential growth; we do
not claim priority for these deductions. The remaining propositions
make precise which finite certificates would be needed and reproduce
restricted cases with complete proofs. The external 2026 star-height
and BVASS advances are attributed to their authors. No unrestricted
problem in this nine-item group is declared closed by our work.


\section{Synthesis: what the attempts establish and what would change the targets}
\label{p06:logic:synthesis}

All fifty-seven requested entries have been treated. This includes separate
answers to both components of entries 37 and 43, and explicit distinctions
between the two categoricity entries. The mathematical output consists of
proved restricted algorithms, definability boundaries, conditional transfers,
and counterexamples to proposed generalizations. The status audit separately
records established results, recent external claims, withdrawn arguments,
independence, and ambiguous formulations. There is no newly proved solution
of a previously open named conjecture to add to an original-closure count.

\subsection{Four transfers that can be checked without accepting a conjecture}

\paragraph{Logarithmic arithmetic becomes a logical boundary.}
Theorem~\ref{p06:lmod:boundary} is the central positive transfer.
Its fixed-coefficient half handles arbitrary first-order quantifier
alternation, not merely the evaluation of a single numerical expression.
Its negative half needs only one additional binary relation to interpret
the rational field. The effective coefficient-basis computation is uniform
in the input syntax; no algorithm for discovering rational relations among
the logarithms at each quantified point is assumed. This precisely states
what the earlier arithmetic tests can contribute to a decision problem.
The mathematical ingredients are classical, and independent novelty of this
combined formulation has not been established.

\paragraph{Spectral representation becomes a decision algorithm only with a cutoff.}
Theorem~\ref{p06:lcomp:positive-spectrum} computes a finite cutoff after which
each residue class has a fixed sign or vanishes identically. The proof gives
rational searches that provably terminate. The computable-real example in
Proposition~\ref{p06:lcomp:computable-real} shows why exact input representation
is essential even in dimension one. These two statements prevent the
qualitative Skolem--Mahler--Lech conclusion from being mistaken for a general
Skolem algorithm.

\paragraph{Borel bounds must agree on equivalence classes.}
Proposition~\ref{p06:ldes:coherent-rates} constructs an actual increasing
finite exhaustion from an eventual class-coherent rate. Its construction
uses intersections with a finite anchor, so it avoids the infinite classes
that can appear when finite partitions are joined. This is a fully proved
form of a known conditional reduction. It does not construct the missing
rates for an arbitrary increasing union, and so leaves targets 49 and 52
open at precisely that step.

\paragraph{Universality spectra admit a genuine countable intersection.}
Proposition~\ref{p06:lfin:intersection} includes the one-sorted coding needed
to keep the output in the original class of countable theories. The result
reduces existence of a least spectrum to a countable coinitial-family
hypothesis. Naming a sort for every countable theory would violate the
language restriction; no countable compression of that entire family is
asserted. Thus the result is an exact reduction, with an identifiable
remaining set-theoretic and model-theoretic problem.

\subsection{The common obstructions are different mathematical requirements}

The following comparisons are consequences of the explicit attempts, not
a claim that all these problems share one hidden solution.

\begin{longtable}{L{28mm}L{50mm}L{64mm}}
\toprule
Targets & What is available & What a closure would additionally require\\
\midrule
\endfirsthead
\toprule Targets & What is available & What a closure would additionally require\\
\midrule\endhead
1, 11 & An injection or a conditional stationary argument & A fiber-respecting section, or a choice-free global replacement for the missing partition hypothesis.\\[4pt]
2, 26 & Semantic truth or an external algorithm & A quantitative proof translation, or a derivation of the needed bound inside the specified weak theory.\\[4pt]
4, 23, 50 & Point codes, analytic orbit relations, or hypergraph witnesses & The required global Borel/coanalytic complexity bound while preserving the original objects.\\[4pt]
7, 29, 41 & Categoricity of one class at one scale & The exact cardinal threshold and structural transfer, with no unproved amalgamation or locality assumption.\\[4pt]
15, 24, 55 & Finite approximations or finite-size bounds & Uniform coherent witnesses, a preserved iteration invariant, or a rank for the unbounded problem.\\[4pt]
17, 25 & Standard-model arguments or stronger-theory upper bounds & Preservation over nonstandard models, or matching analysis for the exact axiom system.\\[4pt]
37, 48, 53, 56 & Finite witnesses or qualitative finite structure & An effective bound or stopping certificate for negative instances in the original input class.\\[4pt]
39, 49, 52 & Finite arity counts or pointwise bounds & One bound uniform across all reducts, or eventual coherence on every equivalence class.\\[4pt]
46, 51 & A valid chain-union operation & Proper elementary extensions preserving all omitted types or all permitted standard traces.\\
\bottomrule
\end{longtable}

These requirements explain why repeatedly transferring a finite inverse
theorem to another field does not, by itself, close the corresponding
unrestricted conjecture. In each row the missing conclusion has different
quantifiers, invariance, effectivity, or foundational assumptions from
the available result.

\subsection{Research priorities suggested by the proofs}

Three directions merit a more focused investigation after this audit.
First, classify useful intermediate logarithmic languages between the
fixed-coefficient equality fragment and the single relation interpreting
$\mathbb Q$. Any such project must specify the domain, allowed coefficients,
inequalities, and quantifier alternation; the theorem here already supplies
two exact endpoints, but does not classify everything between them.
Second, seek a genuinely new effective tail bound for an oscillatory
algebraic recurrence class beyond the supplied positive-spectrum promise.
The proof above identifies the dominant-term estimate that must be replaced.
Third, investigate conditions producing the coherent rates in
Proposition~\ref{p06:ldes:coherent-rates} that are not simply restatements of
hyperfiniteness or Borel boundedness. A new sufficient condition would need
an independently verified class of examples outside known cases.

These are proposed next investigations, not three additional open problems
claimed to be new to the literature. The present dossier provides complete
arguments for its stated deductions and a source-linked account of the
remaining work. Internal cross-reading does not certify global priority,
formal correctness, or the validity of every external recent preprint.

% END SOURCE: KEIO_Logic_After_Observation_57_Problems.tex

\endgroup
\clearpage
\part{Arithmetic Beyond Convergence}
\label{merge:part07}
\begingroup

% BEGIN SOURCE: KEIO_Arithmetic_Beyond_Convergence.tex

\noindent\textbf{Editorial integration note.} 
The logarithmic-root and quantile classifications lead to the inverse constructions of Part~\ref{merge:part08}. The bound $N\le8r^2$ in the tilted-arcsine analysis bounds a cyclotomic search in that family; it is not a universal algebraicity test for arbitrary constants. Later determinant constructions address limitations of the low-rank families considered here.


\begin{quote}\small
\noindent\textit{Original part framing.}\par
\noindent\textbf{Scope and attribution.}
A proof supplied here establishes its stated conclusion; it does not
by itself establish historical novelty. Baker, Gelfond--Schneider,
Hermite--Lindemann, and Mahler value theorems are external inputs,
credited where used. The new manuscript's results are not formally
verified in Lean. Independent checks described below concern the
written proofs and exact computations, not a machine certification.
Bibliographic status was checked on 16--17 October 2026; an external
preprint or repository is identified with its actual verification
scope. Full-text overlap with the closely related 2025 incomplete-beta
paper remains unverified, so priority for the quantile formulation
is not asserted.
\end{quote}
\section{From failed transfers to exact arithmetic questions}
\label{p07:ra:scope}

The preceding logic manuscript examined failures of transferring
finite descriptions to effective or infinite conclusions. Three
parts survive as useful mathematics here: exact logarithmic zero
tests at algebraic inputs, the need for a computable bound before
finite search becomes an algorithm, and the distinction between
recovering an analytic object and determining its arithmetic type.
We restate every hypothesis used; none of the preceding manuscripts'
unverified assertions is used as an axiom.

The strongest transfer is positive. In an arcsine-type primitive,
Baker's theorem forces an algebraic correction to vanish at an
algebraic-probability algebraic quantile. Gelfond--Schneider forces
a rational angle. Every conjugate of the resulting cyclotomic
coordinate remains in $(0,1)$, so the number of available real roots
bounds the cyclotomic degree. The bound $\varphi(N)\le2r$ then turns
an infinite search into $N\le8r^2$. This is the precise arithmetic
information that a numerical inverse calculation alone does not
provide.

\paragraph{Concrete complete classifications.}
Theorem~\ref{p07:rq:tilt} treats every rational polynomial tilt of
the arcsine law. Theorem~\ref{p07:rga:beta-classification} sharpens
the result for all positive half-integral beta shapes. Theorems
\ref{p07:rq:student} and \ref{p07:rq:fisher} give the corresponding
Student and Fisher conclusions. Theorem~\ref{p07:rlog:classification}
decides the arithmetic type and multiplicity of all isolated roots
in its rational-derivative logarithmic class, and
Theorem~\ref{p07:rlog:binary} covers binary entropy at every algebraic
input. Theorem~\ref{p07:rsc:residues} and
Corollary~\ref{p07:rsc:periodic} give exact algebraic exceptions in
two related infinite-series families. These are uniform results,
rather than a count inflated by listing many instances.

\paragraph{What a named-constant closure still requires.}
For irrationality via integer linear forms one must exhibit
$A_n,B_n\in\mathbb Z$ such that
$0<|A_n+B_nC|\to0$. Analytic smallness before denominator
clearance is insufficient. For transcendence, excluding rational
values alone is insufficient; a proof needs a criterion excluding
all algebraic values, with its hypotheses checked. An applicable
classical value theorem can supply this last step for a new
family, while failing entirely for a superficially similar limit.

The fifteen named targets are retained separately in the following
index. Finding a recent external solution corrects the target
list but is not an original solution by this project. No famous
constant closure has been obtained from the present attempts.
The complete family theorems and the quantified obstructions are
the actual mathematical deliverables; their publication level
requires an assessment beyond this investigation.

\subsection{Fifteen targets, explicit attempts, and outcomes}
\label{p07:ra:targets}

The target column distinguishes irrationality from transcendence.
The outcome column reports this investigation's result; it is not
a guarantee that no unindexed manuscript exists elsewhere.
\begingroup\small\setlength{\tabcolsep}{3pt}
\begin{longtable}{rL{28mm}L{49mm}L{66mm}}
\toprule
ID & Constant and target & Attempt & Outcome\\
\midrule\endfirsthead
\toprule ID & Constant and target & Attempt & Outcome\\\midrule\endhead
1 & Euler $\gamma$: irrationality &
Nonlinear common-log elimination; published approximants after
integer clearance &
Not closed. Finite covers remove the common log; factorial
denominators remain. Section~\ref{p07:rgam:section}.\\[3pt]
2 & Euler--Gompertz $\delta$: irrationality &
Positive rational Laplace kernels and twisted residues &
Not closed. Fixed-denominator gap proved; growing pole order
and actual gcd savings remain distinct loopholes.
Section~\ref{p07:rgam:section}.\\[3pt]
3 & Catalan $G$: irrationality and transcendence &
Positive tail Hankel forms; exact normalization &
Recent external irrationality result and formalization reported,
not ours; transcendence not obtained.
Section~\ref{p07:zrec:section}.\\[3pt]
4 & $\zeta(5)$: irrationality and transcendence &
Determinant smallness, rank, and denominator savings &
Recent external irrationality result reported; no new proof here
and no transcendence conclusion. Section~\ref{p07:zrec:section}.\\[3pt]
5 & $\zeta(3)$: transcendence &
Raw tail moments and determinant growth &
Irrationality is known; the present rank-one family does not
supply a transcendence proof. Section~\ref{p07:zrec:section}.\\[3pt]
6 & $\Gamma(1/5)$: irrationality or transcendence &
Reflection, multiplication, and known transcendence degree &
Not closed; explicit constraint models do not distinguish
this entry. Section~\ref{p07:rga:fifths}.\\[3pt]
7 & $\Gamma(2/5)$: irrationality or transcendence &
The separate second gamma entry under the same relations &
Not closed; the rank statement does not identify it.
Section~\ref{p07:rga:fifths}.\\[3pt]
8 & Glaisher--Kinkelin $A$: irrationality or transcendence &
Period representation and hyperfactorial approximants &
Not closed. $A$ differs from $\log A$; transcendence of every
finite approximant does not pass to the limit.
Section~\ref{p07:rga:glaisher}.\\[3pt]
9 & Erd\H{o}s--Borwein $E$: transcendence &
Finite Mahler certificates and full Lambert singularities &
Not closed. Sparse families classified; the full Lambert
function has no finite linear Mahler equation.
Section~\ref{p07:rsc:section}.\\[3pt]
10 & $\sum_{n\ge1}F_n^{-1}$: transcendence &
Periodic sparse weights, Binet, additive functional equations &
Not closed. Periodic doubling subseries classified; full
spectrum does not satisfy the proposed additive equation.
Section~\ref{p07:rsc:section}.\\[3pt]
11 & Prime Mertens $B_1$: irrationality &
Finite prime logarithms, signs, and explicit tails &
Not closed. Finite noncancellation proved; infinite
cancellation with $\gamma$ remains. Section~\ref{p07:rsc:prime}.\\[3pt]
12 & $e+\pi$: irrationality or transcendence &
Algebraic-input logarithmic elimination &
Not closed; transformed logarithmic bases are not known
algebraic. Section~\ref{p07:rlog:section}.\\[3pt]
13 & $e\pi$: irrationality or transcendence &
Separate product equation and algebraic constraints &
Not closed. The elementary disjunction does not identify
this number. Section~\ref{p07:rlog:section}.\\[3pt]
14 & Khinchin $K$: irrationality &
Ergodic recovery, unbounded averages, tail control &
Not closed. Exact integrability condition and a normal
Liouville counterexample isolate the missing input.
Section~\ref{p07:ric:section}.\\[3pt]
15 & Feigenbaum $\delta_F$: irrationality &
Certified fixed points and isolated unstable eigenvalues &
Not closed. Uniform contraction, a spectral gap, and rational
coefficients do not decide arithmetic type.
Section~\ref{p07:ric:section}.\\
\bottomrule\end{longtable}\endgroup
\newpage



\section{Effective arithmetic classification of logarithmic roots}
\label{p07:rlog:section}

The useful residue of the earlier logical obstruction is an exact
zero test on algebraic inputs. In one variable, a rational derivative
turns this local test into a global classification of every real root.
We give the algorithm and its proof, then a separate algebraic-input
classification for binary entropy. All logarithms are natural unless
a base is displayed. The inputs from transcendence theory are classical
theorems of Baker and Hermite--Lindemann; no new transcendence theorem
of that generality is claimed.

\subsection{A terminating classifier, including multiplicities}

Put $A=\overline{\mathbb Q}\cap\mathbb R$. An algebraic number is
specified by a defining polynomial and an isolating interval. A
semialgebraic input is given by a finite Boolean combination of
polynomial sign conditions with coefficients in $A$.
Consider
\begin{equation}\label{p07:rlog:F}
 F(x)=R(x)+\sum_{i=1}^s\lambda_i\log P_i(x),
 \qquad R\in A(x),\quad \lambda_i\in A,\quad P_i\in A[x].
\end{equation}
Let $D\subseteq\mathbb R$ be a semialgebraic set over $A$ on which
$R$ has no poles and all $P_i$ are strictly positive. A positive
algebraic constant is allowed as a constant polynomial $P_i$.

\begin{theorem}[Effective root classification]\label{p07:rlog:classification}
There is a terminating algorithm with the following outputs for
\eqref{p07:rlog:F} on $D$.
\begin{enumerate}
\item It determines every interval component on which $F$ vanishes
identically. All remaining zeros form a finite set, which it isolates
in pairwise disjoint rational intervals and classifies individually
as algebraic or transcendental. Algebraic zeros are returned by exact
algebraic representations. Each isolating interval can be refined to
any prescribed positive rational width.
\item If $F'\not\equiv0$ as a rational function, write $F'=N/Q$ in
reduced form. Every multiple zero is algebraic. At any zero $a$ the
analytic multiplicity is
\[
 1+\operatorname{ord}_aN.
\]
In particular every transcendental zero is simple, and the algorithm
computes the multiplicities of all zeros.
\item If $D$ has $c$ nondegenerate interval components and $s_0$
isolated points, and $F'\not\equiv0$, the total number of zeros in
$D$ is at most $\deg N+c+s_0$.
\end{enumerate}
In the constant case, isolated points of $D$ are tested separately;
analytic multiplicity is not assigned to an identically zero germ.
No polynomial-time complexity estimate is asserted.
\end{theorem}

\begin{proof}
We first record two exact arithmetic facts. Baker's linear independence
theorem says that logarithms of algebraic numbers independent over
$\mathbb Q$ remain independent over $\overline{\mathbb Q}$. His
inhomogeneous consequence says that a nonzero algebraic-coefficient
linear form in logarithms of algebraic numbers is transcendental
\cite{p07:rlog:BakerIII,p07:rlog:BakerBook}.

Compute a rational basis $\theta_1,\ldots,\theta_d$ of the span of
the fixed coefficients $\lambda_i$, and write
$\lambda_i=\sum_k r_{ki}\theta_k$, $r_{ki}\in\mathbb Q$.
Exact arithmetic in a common number field computes this basis.
Choose positive integers clearing the denominators of each row and
write the resulting integer coefficients as $n_{ki}$. Define
\[
 G_k(x)=\prod_{n_{ki}>0}P_i(x)^{n_{ki}}
          -\prod_{n_{ki}<0}P_i(x)^{-n_{ki}}.
\]
Empty products are one. At every algebraic $a\in D$, Baker's theorem
gives
\begin{equation}\label{p07:rlog:certificate}
 F(a)=0\quad\Longleftrightarrow\quad
 R(a)=0\ \text{and}\ G_1(a)=\cdots=G_d(a)=0.
\end{equation}
Indeed, the inhomogeneous theorem first separates the algebraic term
$R(a)$ from the logarithmic sum. For that sum choose any rational
basis $\rho_j$ from the finitely many logarithms $\log P_i(a)$.
Writing them in this basis and applying Baker gives zero algebraic
coefficients of each $\rho_j$; independence of the $\theta_k$ then
gives $\sum_i r_{ki}\log P_i(a)=0$ for every $k$. Exponentiation
on positive real bases gives precisely $G_k(a)=0$. This reasoning
does not require an algorithm to find the basis of logarithms.

Conversely, the equations on the right of \eqref{p07:rlog:certificate}
imply $F(x)=0$ for \emph{every real} $x\in D$ satisfying them.
Thus they define a semialgebraic subset $S$ of the zero set that
contains exactly all algebraic zeros, together with any intervals it
may contain. Clear the denominator of $R$ in forming $S$; its
denominator is already excluded by the domain guard.

The same argument decides equality to zero of every number
$b+\sum_i\lambda_i\log a_i$, where $b\in A$ and $a_i\in A_{>0}$.
If it is nonzero, its sign is computable by certified approximations
until an interval excluding zero is obtained. Such approximations
exist effectively, for example by the convergent series
\[
 \log a=2\sum_{j\geq0}\frac{z^{2j+1}}{2j+1},
 \qquad z=\frac{a-1}{a+1},\qquad |z|<1,
\]
and its explicit geometric tail bound. Termination uses the exact
zero test, not a claim that small numerical values are nonzero.

Now differentiate:
\begin{equation}\label{p07:rlog:derivative}
 F'=R'+\sum_i\lambda_i\frac{P_i'}{P_i}\in A(x).
\end{equation}
Polynomial arithmetic decides whether this rational function is zero.
If it is zero, $F$ is constant on every interval component of $D$;
an algebraic sample point and the exact zero test decide whether that
constant is zero. The finitely many isolated points and included
endpoints of $D$ can be tested in the same way. This completes this
case.

Suppose instead that the reduced numerator $N$ in
\eqref{p07:rlog:derivative} is nonzero. Isolate its real algebraic roots
and partition each interval component of $D$ at those roots. On every
remaining open interval $F'$ has constant nonzero sign, so $F$ is
strictly monotone. All partition points that belong to $D$ are
algebraic and can be tested exactly. The endpoints of these intervals
are algebraic or infinite. We next show that the signs, including
zero, of all one-sided limits are computable.

At a finite endpoint $b$, factor
$P_i(x)=(x-b)^{m_i}U_i(x)$ with $U_i(b)\ne0$. Along a permitted
side of the domain,
\[
 \log P_i(x)=m_i\log|x-b|+\log|U_i(b)|+o(1).
\]
If $R$ has a pole at $b$, its leading Laurent term dominates all
logarithms and determines the infinite limit and its sign. If $R$
has no pole, the coefficient $\sum_i\lambda_i m_i$ determines any
logarithmic divergence. If that coefficient is zero, the finite
limit is $R(b)+\sum_i\lambda_i\log|U_i(b)|$, to which the exact
sign test applies. At either infinity use polynomial division for
$R$ and
\[
 \log P_i(x)=\deg(P_i)\log|x|+\log|\operatorname{lc}(P_i)|+o(1).
\]
A nonconstant polynomial part of $R$ dominates; otherwise the
coefficient of $\log|x|$ determines divergence, or the constant
limit has the same decidable algebraic-plus-logarithmic form.
All orders, leading coefficients, and cancellations are exact
algebraic computations.

A strictly monotone open interval has one zero precisely when zero
lies strictly between its two endpoint limits. A zero limit at an
excluded endpoint is not an interior root. When a root exists,
enumerate rational points in that interval until points of opposite
sign are found, and refine by the sign test and bisection. This
terminates by continuity and strict monotonicity. Roots at partition
points have already been tested separately. Consequently the full
zero set is finite and effectively isolated.

The set $S$ above is now finite. Its points are effectively computable
algebraic numbers: equivalently one may take the common real roots
of the nonzero polynomials defining it and filter by the domain
conditions. At least one such polynomial is nonzero unless $S$ is
empty or $D$ has only isolated points; an interval of solutions would
contradict $N\ne0$. Standard univariate real algebraic algorithms
handle all cases \cite{p06:lmod:BPR}. Exact order comparisons place each
point of $S$ in its monotonicity interval or at its partition point.
An interval containing a zero but no point of $S$ therefore contains
a transcendental zero. Refining finitely many isolators separates
all roots and gives any requested accuracy.

Since $F$ is analytic near every point of $D$, its derivative has
no pole there. Hence $Q(a)\ne0$ at a zero, and integration of the
first nonzero Taylor term of $F'$ proves the multiplicity formula.
A multiple zero must be a root of $N$, and is therefore algebraic.
Finally Rolle's theorem bounds the number of zeros in an interval
component by one plus the number of critical points there. Summing
and adding the isolated domain points gives the stated bound.
\end{proof}

General root-isolation procedures for elementary functions predate
this argument, including Strzebo\'nski's exp--log--arctan algorithm
\cite{p07:rlog:Strzebonski}. Our statement concerns the exact arithmetic
classification and the explicit terminating zero test in the much
smaller rational-derivative class. It does not assert priority over
that algorithmic literature. In particular, it uses neither a general
oracle for elementary constants nor Schanuel's conjecture.

\subsection{Three natural inverse families}

\begin{proposition}[Erlang quantiles]\label{p07:rlog:erlang}
For any integer $n\geq1$ and algebraic $q\in(0,1)$, the $q$-quantile
of the unit-scale gamma distribution with shape $n$ is transcendental.
It is the unique positive zero of
\[
 x-\log\left(\sum_{k=0}^{n-1}\frac{x^k}{k!}\right)+\log(1-q).
\]
The zero is simple and effectively isolatable.
\end{proposition}
\begin{proof}
Integrating the density $x^{n-1}e^{-x}/(n-1)!$ gives survival function
$e^{-x}S_n(x)$, where $S_n(x)=\sum_{k=0}^{n-1}x^k/k!$.
The displayed function has value $\log(1-q)<0$ at zero, tends to
infinity, and has derivative
$x^{n-1}/((n-1)!S_n(x))>0$ for $x>0$. It has exactly one simple
positive zero. At an algebraic positive zero the nonzero rational
term $x$ would have to vanish by \eqref{p07:rlog:certificate}, a
contradiction. Equivalently, this last conclusion follows directly
from Hermite--Lindemann applied to $e^{-x}S_n(x)=1-q$.
\end{proof}

This is a classical-transcendence consequence for an entire quantile
family, not a claim that individual gamma quantiles were a named
open problem. Algebraic positive changes of scale preserve the conclusion.

\begin{proposition}[The inverse logarithmic mean and its exceptional root]
\label{p07:rlog:logmean}
Extend $L(x)=(x-1)/\log x$ continuously by $L(1)=1$ on $(0,\infty)$.
For algebraic $a>0$, $L(x)=a$ has one solution. It is $1$ when
$a=1$, and is transcendental when $a\ne1$.
For the un-divided equation $x-1-a\log x=0$, there are two distinct
simple roots when $a\ne1$: the algebraic root $1$ and the genuine
inverse-mean solution. At $a=1$ they coalesce to a double root at $1$.
\end{proposition}
\begin{proof}
For $x\ne1$, the sign of $L'(x)$ is the sign of
$\log x-1+1/x$, which is positive: this expression has its unique
minimum zero at $x=1$. The expansion at $1$ gives $L'(1)=1/2$.
Also $L(x)\to0$ at zero and $L(x)\to\infty$ at infinity.
Thus there is a unique inverse-mean solution. If it is algebraic and
different from $1$, the equation says that $\log x=(x-1)/a$ is a
nonzero algebraic number, contradicting Hermite--Lindemann.
For the un-divided equation, $F'=1-a/x$ and $F''=a/x^2>0$.
Its minimum is $a-1-a\log a<0$ unless $a=1$, and its two endpoint
limits are positive infinity. These facts prove the two-root assertion
and simplicity. When $a=1$, $F(1)=F'(1)=0$ and $F''(1)=1$.
\end{proof}

\begin{proposition}[A normalized rational-density quantile]
\label{p07:rlog:density}
On $[0,1]$ normalize the positive rational weight
$w(x)=1+(1+x)^{-1}$ to a probability density.
Every interior algebraic-probability quantile is transcendental.
\end{proposition}
\begin{proof}
The cumulative distribution is
$(x+\log(1+x))/(1+\log2)$. For algebraic $q\in(0,1)$ its quantile
is the unique zero in $(0,1)$ of
\[
 F_q(x)=x-q+\log(1+x)-q\log2.
\]
Uniqueness follows from $F_q'>0$ and the opposite endpoint signs.
If the root were algebraic, \eqref{p07:rlog:certificate} would force
$x=q$ and $\log(1+q)=q\log2$. Strict concavity of $\log$ gives
$\log(1+q)>q\log2$ for $0<q<1$, a contradiction.
\end{proof}

These examples illustrate three distinct tasks: an entire distribution
family, an inverse mean with a removable-root trap, and quantiles whose
normalizing integral is itself transcendental. None is obtained by
declaring an arbitrary numerical approximation irrational.

\subsection{Binary entropy: classification for every algebraic input}

Write
\[
 h_2(p)=\frac{-p\log p-(1-p)\log(1-p)}{\log2},\qquad 0<p<1.
\]
Boche--Schaefer--Poor prove that $h_2(p)$ is transcendental for
rational $p\notin\{0,1/2,1\}$ \cite[Appendix B, Theorem 4]{p07:rlog:BSP}.
Blakaj--Wolf establish the general zero-or-transcendental dichotomy
for natural-log entropy at algebraic spectra and the
rational-or-transcendental dichotomy for algebraic-base entropy at
rational spectra \cite[Theorem 2]{p07:rlog:BlakajWolf}.
The next proof covers all real algebraic binary inputs in base $2$.
The logarithmic
coefficients now depend on $p$, so this is a separate argument rather
than an application of Theorem~\ref{p07:rlog:classification}.

\begin{lemma}\label{p07:rlog:powers2}
If $u,v\in\mathbb Q$ and $2^u+2^v=1$, then $u=v=-1$.
\end{lemma}
\begin{proof}
Choose a common positive denominator $n$ and put $t=2^{1/n}$.
The polynomial $X^n-2$ is Eisenstein at $2$, so
$1,t,\ldots,t^{n-1}$ are linearly independent over $\mathbb Q$.
Writing each integer exponent of $t$ as $kn+r$, $0\leq r<n$,
expresses each summand as a positive rational multiple of $t^r$.
Independence and positivity force both remainders to be zero.
Hence $u,v$ are integers. Both are negative, and two negative
integer powers of $2$ sum to one only if both are $1/2$.
\end{proof}

\begin{theorem}[Algebraic binary probabilities have no intermediate
algebraic entropy]\label{p07:rlog:binary}
For $p\in A\cap(0,1)$,
\[
 h_2(p)\in\overline{\mathbb Q}
       \quad\Longleftrightarrow\quad p=\tfrac12.
\]
Consequently, for every algebraic $q\in(0,1)$ the two solutions of
$h_2(p)=q$ are transcendental and effectively isolatable. A binary
symmetric channel with algebraic crossover probability strictly
between $0$ and $1/2$ has transcendental capacity in bits.
\end{theorem}
\begin{proof}
Suppose $q=h_2(p)$ is algebraic. The defining relation is
\begin{equation}\label{p07:rlog:entropyrelation}
 -\log(1-p)+p\log\frac{1-p}{p}-q\log2=0.
\end{equation}
First suppose $p$ is algebraic irrational. If $q$ does not belong
to the rational span of $1,p$, Baker's coefficient splitting applied
to the independent coefficients $1,p,q$ forces $\log2=0$, which
is impossible. Thus $q=u+vp$ with $u,v\in\mathbb Q$. Splitting
\eqref{p07:rlog:entropyrelation} in the independent coefficients $1,p$
then gives
\[
 \log(1-p)=-u\log2,\qquad
 \log\frac{1-p}{p}=v\log2.
\]
It follows that $1-p=2^{-u}$ and $p=2^{-u-v}$.
Lemma~\ref{p07:rlog:powers2} forces $p=1/2$, contradicting irrationality.

Now let $p=m/n$ in lowest terms, with $0<m<n$.
Baker splitting in \eqref{p07:rlog:entropyrelation} first implies
$q\in\mathbb Q$, since otherwise the coefficients $1,q$ would
again force $\log2=0$. Exponentiating gives
\[
 \frac{n^n}{m^m(n-m)^{n-m}}=2^{nq}.
\]
After clearing the denominator of $q$, the prime valuation of the
left-hand side at every odd prime is zero. The three positive integers
$m,n-m,n$ are pairwise coprime. Therefore none has an odd prime
factor: each is a power of $2$. Since $n\geq2$, coprimality forces
$m=n-m=1$, so $p=1/2$ and $q=1$. The converse is immediate.

The function $h_2$ is symmetric, strictly increasing from $0$ to $1$
on $(0,1/2)$, and strictly decreasing on $(1/2,1)$. Its derivative
is $\log((1-p)/p)/\log2$. Thus every $q\in(0,1)$ has precisely two
roots, and the classification proves their transcendence when $q$
is algebraic. For effective isolation, use rational bisection on
$(0,1/2)$ and symmetry. At every rational interior sample point the
difference from $q$ is a nonzero computable real, so certified
logarithmic approximation eventually decides its sign. Finally the
binary symmetric capacity is $1-h_2(p)$; subtracting a transcendental
number from $1$ preserves transcendence.
\end{proof}

The rational-input part is expressly credited to \cite{p07:rlog:BSP}.
The algebraic-input extension is proved here from Baker's theorem;
priority beyond the primary sources checked is not certified. This
statement is special to the binary entropy measured in bits. For
example, with $p=(\sqrt5-1)/2$ and $b=1/p$, the identity $1-p=p^2$
gives entropy to base $b$ equal to $2-p$, an algebraic irrational.
Changing the base therefore changes the arithmetic conclusion.

\paragraph{Explicit isolators.}
The following decimal endpoints are exact rational numbers, not
rounded assertions about a root. They give four concrete instances
of the proved families:
\begin{center}
\begin{tabular}{lll}
\toprule
Root & Rational isolating interval & Arithmetic type\\
\midrule
Shape-$2$ Erlang median & $(1.678346,1.678347)$ & transcendental\\
$L(x)=2$ & $(3.512862,3.512863)$ & transcendental\\
Rational-density median & $(0.464832,0.464833)$ & transcendental\\
$h_2(p)=1/2$, $p<1/2$ & $(0.110027,0.110028)$ & transcendental\\
\bottomrule
\end{tabular}
\end{center}
For an entirely rational check of the signs, truncate the logarithm
series used above after $240$ terms. For $z=(a-1)/(a+1)$ the absolute
remainder is bounded by
$2|z|^{481}/(481(1-|z|^2))$. Rational interval arithmetic in the
displayed defining equations gives a negative value at each left
endpoint and a positive value at each right endpoint. The monotonicity
proved above establishes isolation. These checks locate the roots;
their transcendence follows from the proofs, not from the intervals.

\subsection{Named targets that the classifier does not close}

We tested $e+\pi$ and $e\pi$ separately. The irrationality of these
individual constants is not established by the reviewed results;
Waldschmidt lists the classical gap \cite{p07:rlog:Waldschmidt}, and a
2026 study still treats $e+\pi$ as open \cite{p07:rlog:Yu}.
The useful failure is a specific hypothesis failure, rather than
an inconclusive decimal search.

For the sum, assuming $e+\pi=a\in A$ gives
$\log(a-\pi)=1$. The logarithm's argument is not known algebraic:
it is $e$. Encoding $\pi$ by a complex logarithm of $-1$ makes
this a nested logarithmic equation, outside the fixed algebraic
polynomial-base class. Baker's theorem cannot be applied to that
unknown base. For the product, $e\pi=a\in A$ similarly gives
$\log(a/\pi)=1$ and again introduces a transcendental argument
and a nested period relation. Treating $\pi$ as a coefficient also
violates the algebraic-coefficient input condition. Numerical
isolation or a rational derivative in a different variable does
not repair these arithmetic failures.

The elementary disjunctive conclusion remains valid: at least one
of $e+\pi$ and $e\pi$ is transcendental. If both were algebraic,
$e$ would be a root of
$X^2-(e+\pi)X+e\pi\in\overline{\mathbb Q}[X]$, contradicting
its transcendence. This does not identify either individual number
and is not counted as a new result. Both named targets remain open
in this branch's ledger.


\section{Gamma values, a logarithmic obstruction, and arithmetic of beta quantiles}
\label{p07:rga:section}

The three pre-existing named targets in this section are
$\Gamma(1/5)$, $\Gamma(2/5)$, and the Glaisher--Kinkelin constant
$A$.  We do not prove any of these three numbers irrational or
transcendental.  We identify exactly why the available transfers do
not do so, and prove a separate effective classification for a natural
two-parameter family of probability quantiles.  Related prior work
on incomplete beta values by Dhillon--Saha must be taken into account
\cite{p07:rga:DhillonSaha}.  This classification
is an application of classical transcendence theorems with a
conjugacy argument; priority for its particular formulation is not
asserted.  All status checks refer to 17 October 2026.

\subsection{The two fifth-denominator gamma values}
\label{p07:rga:fifths}

Write $G_1=\Gamma(1/5)$ and $G_2=\Gamma(2/5)$.
The recent geometric transcendence theorem for the gamma function
does not settle its values at these individual inputs.  Sheth's
2026 paper explicitly leaves $\Gamma(1/5)$ open in the discussion
following Conjecture 4.2 \cite{p07:rga:Sheth}.  The classical result of
Grinspan, together with the related work of Vasil'ev, gives
\begin{equation}\label{p07:rga:rank}
 \operatorname{trdeg}_{\mathbb Q}\mathbb Q(\pi,G_1,G_2)\geq2;
\end{equation}
see \cite{p07:rga:Grinspan} and the precise account in
\cite[Section~3.3.2, p.~91]{p07:rga:Waldschmidt}.
This does imply that at least one of $G_1,G_2$ is transcendental
even over $\mathbb Q(\pi)$.  It does not identify which one.
Neither individual irrationality question is resolved here, and no
primary resolution of either individual transcendence question was
found in the present search.

Reflection and the fivefold multiplication formula give
\begin{align}
 G_1\Gamma(4/5)&=\frac{\pi}{\sin(\pi/5)},&
 G_2\Gamma(3/5)&=\frac{\pi}{\sin(2\pi/5)},\label{p07:rga:reflection}\\
 \prod_{j=1}^4\Gamma(j/5)&=\frac{4\pi^2}{\sqrt5}.
 \label{p07:rga:multiplication}
\end{align}
The sine factors are algebraic.  Since
$\sin(\pi/5)\sin(2\pi/5)=\sqrt5/4$, the last displayed equation
adds no independent relation to the first two.

\begin{proposition}[A concrete limit of reflection and rank transfers]
\label{p07:rga:formal-obstruction}
The relations \eqref{p07:rga:reflection}--\eqref{p07:rga:multiplication},
positivity of their four entries, and the rank bound
\eqref{p07:rga:rank} do not, as a system of numerical constraints,
force either specified entry $G_1$ or $G_2$ to be transcendental.
\end{proposition}
\begin{proof}
Fix the actual number $\pi$ and a positive real number $T$
transcendental over $\overline{\mathbb Q}(\pi)$; such numbers exist
because that field and its algebraic closure are countable.
Set $g_1=2$, $g_2=T$, and define
\[
 g_4=\frac{\pi}{2\sin(\pi/5)},\qquad
 g_3=\frac{\pi}{T\sin(2\pi/5)}.
\]
These four positive numbers satisfy all the displayed relations,
and $\mathbb Q(\pi,g_1,g_2)$ has transcendence degree two, but
$g_1$ is rational.  Interchanging the roles of $g_1$ and $g_2$ and
redefining their reflection partners gives the corresponding
example with $g_2$ rational.
\end{proof}

These are auxiliary numerical tuples, not alternative gamma
functions.  The proposition proves only the insufficiency of the
listed constraints.  A genuine value theorem needs extra arithmetic
information beyond those constraints.  For example, complete beta
integrals naturally involve ratios such as
$B(1/5,1/5)=G_1^2/G_2$, rather than isolating either $G_j$.
The gamma function itself is differentially transcendental, so it
cannot simply be inserted as an $E$-function into a
Siegel--Shidlovskii value theorem.  Likewise, a limiting value at
infinity of an auxiliary function is not a value at a regular
algebraic point.  These are missing hypotheses, not cosmetic changes
of representation.

\subsection{Glaisher--Kinkelin: the constant is not its logarithm}
\label{p07:rga:glaisher}

Our normalization is
\begin{equation}\label{p07:rga:A}
 \log A=\frac1{12}-\zeta'(-1),\qquad
 A=\lim_{N\to\infty}
 \frac{\prod_{k=1}^N k^k}{N^{N^2/2+N/2+1/12}e^{-N^2/4}}.
\end{equation}
See the primary formulas and discussion in
\cite{p07:rga:MPM,p07:rga:Pain}.  The first Stirling--Ramanujan constant in
Mu\~noz--P\'erez-Marco is $S_1=\log A$, not $A$.
Their transcendence conjecture concerns $S_1$.
Furthermore, their Section 1 explicitly distinguishes their
transalgebraic periods over $\mathbb Q(t,e^{-t})$ from the more
restrictive exponential periods of Kontsevich--Zagier.
Consequently the phrase ``exponential period'' in that paper is
not a theorem that $A$ is irrational, and it cannot be substituted
into a period theorem having different hypotheses.  No individual
irrationality or transcendence result for $A$ is claimed here.

There are two precise failed transfers.  First, even a proof that
$\log A$ is transcendental would not by itself prove $A$
transcendental: $\log2$ is transcendental whereas its exponential is
algebraic.  Baker's theorem applies to finite linear forms in
logarithms of algebraic numbers.  Writing an unknown constant as an
integral or as an infinite logarithmic series does not supply that
finite representation.

Second, the natural finite approximants already have a completely
known arithmetic type, but this does not pass to their limit.

\begin{proposition}[Transcendental hyperfactorial approximants]
\label{p07:rga:approximants}
For each positive integer $N$, the finite expression
\[
 A_N=\frac{\prod_{k=1}^N k^k}
 {N^{N^2/2+N/2+1/12}}e^{N^2/4}
\]
is transcendental.  This assertion, together with $A_N\to A$,
does not imply that $A$ is irrational.
\end{proposition}
\begin{proof}
The factor multiplying $e^{N^2/4}$ is a nonzero algebraic number,
because all exponents of positive integers in that factor are
rational.  The exponential is transcendental by the
Hermite--Lindemann theorem.  Multiplication by a nonzero algebraic
number preserves transcendence.  As a general implication,
transcendence is not closed under limits: the transcendental
numbers $e^{1/N}$ converge to the rational number $1$.
\end{proof}

Thus the logical failure identified earlier for noneffective
finite descriptions has a direct arithmetic analogue.  An exact
finite formula and a convergence proof do not furnish a
Diophantine lower bound for the limit.  A new approximation argument
would have to control nonzero integer linear forms in $A$ itself,
including their denominators and errors, or establish another
applicable arithmetic value theorem.  We have not done so.

\subsection{A complete effective classification for half-integral beta shapes}
\label{p07:rga:beta}

For $a,b>0$, define
\[
 I_x(a,b)=\frac{1}{B(a,b)}\int_0^x t^{a-1}(1-t)^{b-1}\,dt,
 \qquad 0\leq x\leq1.
\]
This is the standard beta distribution function
\cite[Equations~8.17.1--8.17.4]{p07:rga:DLMF}.
For each $0<q<1$ there is a unique $Q_{a,b}(q)\in(0,1)$ with
$I_{Q_{a,b}(q)}(a,b)=q$: the density is strictly positive in the
open interval, and the distribution function is continuous with
endpoint values zero and one.

For nonnegative integers $m,n$, set $N=m+n$ and
\begin{equation}\label{p07:rga:K}
 K_{m,n}=\frac{\pi}{B(m+1/2,n+1/2)}
 =\frac{4^N m!n!N!}{(2m)!(2n)!}\in\mathbb Q_{>0}.
\end{equation}
If $N>0$, let $P_{m,n}(X)\in\mathbb Q[X]$ be the polynomial
of degree at most $N-1$ satisfying
\begin{equation}\label{p07:rga:P-equation}
 2X(1-X)P'_{m,n}(X)+(1-2X)P_{m,n}(X)
       =2\bigl(K_{m,n}X^m(1-X)^n-1\bigr).
\end{equation}
Existence and uniqueness, as well as an explicit way to compute it,
are proved below.  Set $P_{0,0}=0$.

\begin{theorem}[Beta quantiles: algebraic and transcendental cases]
\label{p07:rga:beta-classification}
Let $a,b\in\tfrac12\mathbb Z_{>0}$, and let
$q\in\overline{\mathbb Q}\cap(0,1)$.
\begin{enumerate}
\item If at least one of $a,b$ is an integer, $Q_{a,b}(q)$ is
algebraic.
\item For $a=b=1/2$, the quantile is algebraic if and only if
$q$ is rational.  In that case it is $\sin^2(\pi q/2)$.
\item Write $a=m+1/2$, $b=n+1/2$ with $m,n\geq0$ and $m+n>0$.
If $mn=0$, every such quantile is transcendental.
\item If $m,n>0$, the quantile is algebraic if and only if one of
the following exact rational tests succeeds:
\[
 \begin{array}{c|c|c}
 q&\text{test}&Q_{m+1/2,n+1/2}(q)\\ \hline
 1/3&P_{m,n}(1/4)=0&1/4\\
 1/2&m=n&1/2\\
 2/3&P_{m,n}(3/4)=0&3/4.
 \end{array}
\]
For a fixed pair $(m,n)$ at most one row can succeed.  All other
algebraic levels have transcendental quantiles.
\end{enumerate}
Thus this is an effective decision of the arithmetic type for
every exactly specified input in this family.  It does not assert
that every quantile is transcendental.
\end{theorem}

\begin{lemma}[The correction polynomial and its real zero]
\label{p07:rga:correction}
For $m,n\geq0$,
\begin{equation}\label{p07:rga:CDF}
 I_x(m+1/2,n+1/2)=
 \frac{2\arcsin\sqrt{x}+\sqrt{x(1-x)}P_{m,n}(x)}{\pi},
 \qquad 0<x<1.
\end{equation}
For $m+n>0$ the polynomial $P_{m,n}$ is nonzero.
If exactly one of $m,n$ is zero, it has no zero in $(0,1)$.
If $m,n>0$, it has exactly one distinct zero in $(0,1)$.
That zero is $1/2$ if and only if $m=n$.
\end{lemma}
\begin{proof}
Put $x=\sin^2\theta$, $0<\theta<\pi/2$.
The change of variables turns the beta differential into
$2\sin^{2m}\theta\cos^{2n}\theta\,d\theta$.
The function
$w(\theta)=\sin^{2m}\theta\cos^{2n}\theta$
is a rational-coefficient polynomial in $\cos(2\theta)$ and
therefore a rational linear combination
$c_0+\sum_{j=1}^{N}c_j\cos(2j\theta)$.
Its integral over $[0,\pi/2]$ is $c_0\pi/2$, so
$c_0=1/K_{m,n}>0$.  Integrating and normalizing gives
\[
 I_{\sin^2\theta}(m+1/2,n+1/2)
 =\frac{2\theta}{\pi}
   +\frac1\pi\sum_{j=1}^N\frac{c_j}{j c_0}\sin(2j\theta).
\]
Use
$\sin(2j\theta)=\sin(2\theta)U_{j-1}(\cos(2\theta))$,
where $U_{j-1}\in\mathbb Z[X]$ is a Chebyshev polynomial.
Since $\sin(2\theta)=2\sqrt{x(1-x)}$ and
$\cos(2\theta)=1-2x$, this proves \eqref{p07:rga:CDF} with
$P_{m,n}\in\mathbb Q[X]$ of degree at most $N-1$.
Differentiating yields \eqref{p07:rga:P-equation}.
For uniqueness, the homogeneous equation implies that
$\sqrt{x(1-x)}P(x)$ is constant on $(0,1)$.  A polynomial
solution has this product tending to zero at the endpoints, so
the constant and the polynomial are zero.  Nonzero $N$ makes
$w$ nonconstant, proving that the correction is nonzero.

Let
$H(\theta)=I_{\sin^2\theta}(m+1/2,n+1/2)-2\theta/\pi$.
It vanishes at both endpoints and has derivative
\[
 H'(\theta)=\frac2\pi\bigl(K_{m,n}w(\theta)-1\bigr).
\]
If $m=0<n$, $w$ is strictly decreasing, and its mean lies
strictly between its endpoint values.  Consequently $H'$ changes
sign once, from positive to negative.  Hence $H>0$ in the
interior.  The case $n=0<m$ is the reflection, with $H<0$.

If $m,n>0$, $w$ is strictly increasing up to its unique maximum
and strictly decreasing thereafter, and vanishes at both endpoints.
Its maximum exceeds its mean.  Thus $H'$ has exactly the sign
pattern negative, positive, negative, with two transition points.
The first local extremum of $H$ is strictly negative, and the
second is strictly positive because the final decreasing segment
ends at zero.  There is exactly one zero on the strictly increasing
middle segment, and no other interior zero.  By
\eqref{p07:rga:CDF}, zeros of $H$ are precisely those of $P_{m,n}$
under $x=\sin^2\theta$.

For $m=n$, symmetry gives $I_{1/2}=1/2$, so that zero is $1/2$.
If $m>n$, then on $0<x<1/2$ the beta density is strictly smaller
than its reflected density, since their ratio is
$(x/(1-x))^{m-n}<1$.  Integration gives $I_{1/2}<1/2$.
For $m<n$ the reverse strict inequality holds.  This proves the
last assertion.
\end{proof}

\begin{proof}[Proof of Theorem~\ref{p07:rga:beta-classification}]
Suppose first that $a=k$ is a positive integer and $b$ is a
positive half-integer or integer.  In the defining integral use
$u=\sqrt{1-t}$.  The differential becomes
\[
 -2(1-u^2)^{k-1}u^{2b-1}\,du,
\]
which is a polynomial differential with rational coefficients.
Its integral between zero and one is a positive rational number.
Thus the equation $I_x(a,b)=q$ is a nonconstant polynomial
equation over $\overline{\mathbb Q}$ for $u=\sqrt{1-x}$.
It forces $u$ and hence $x$ to be algebraic.  Symmetry handles an
integer $b$ instead.

Now both parameters have the form $m+1/2,n+1/2$.
Assume that a quantile $x$ is algebraic, and set
\[
 \theta=\arcsin\sqrt{x},\quad
 z=1-2x+2i\sqrt{x(1-x)},\quad
 R=\sqrt{x(1-x)}P_{m,n}(x).
\]
Here $z$ is algebraic on the upper unit semicircle,
$\log z=2i\theta$ on the principal branch, and $R$ is algebraic.
Equation \eqref{p07:rga:CDF} gives
\begin{equation}\label{p07:rga:Baker-step}
 \log z-q\log(-1)=-iR,
 \qquad \log(-1)=i\pi.
\end{equation}
Baker's theorem states that a nonzero linear form with algebraic
coefficients in logarithms of nonzero algebraic numbers is
transcendental \cite{p07:rga:BakerIII}.  Therefore \eqref{p07:rga:Baker-step}
forces $R=0$ and $2\theta=\pi q$.  If $q$ were algebraic
irrational, the Gelfond--Schneider theorem would make
$z=\exp(q\log(-1))$ transcendental.  Hence $q$ is rational.
Conversely, when $q$ is rational, $x=\sin^2(\pi q/2)$ is
algebraic; it is the desired quantile exactly when
$P_{m,n}(x)=0$.

For $m=n=0$ the correction is identically zero, giving part 2.
The no-zero statement of Lemma~\ref{p07:rga:correction} gives part 3.
It remains to sharpen the rational-angle condition when $m,n>0$.
Since $q$ is rational, $z$ is a root of unity of order greater
than two, and
\[
 x=\frac{2-z-z^{-1}}4.
\]
Every algebraic conjugate of $x$ has the form
$(2-z^j-z^{-j})/4$ with $\gcd(j,\operatorname{ord}(z))=1$.
Each lies strictly between zero and one.  All are zeros of the
rational polynomial $P_{m,n}$.  Lemma~\ref{p07:rga:correction} allows
only one such real zero, so all conjugates are equal and $x$ is
rational.  Now $z+z^{-1}=2-4x$ is a rational algebraic integer
strictly between $-2$ and $2$.  It is therefore $-1,0$, or $1$,
giving exactly $x=3/4,1/2$, or $1/4$ and respectively
$q=2/3,1/2$, or $1/3$.  The midpoint test is equivalent to $m=n$
by the lemma.  Uniqueness of the interior zero also proves that
at most one row can hold.
\end{proof}

The tests in the theorem require only rational arithmetic.
For explicit implementation write
$K_{m,n}X^m(1-X)^n=\sum_{k=0}^{N}w_kX^k$ and
$P_{m,n}(X)=\sum_{k=0}^{N-1}p_kX^k$.
Coefficient comparison gives
\begin{equation}\label{p07:rga:recurrence}
 p_0=2(w_0-1),\qquad
 p_k=\frac{2w_k+2k p_{k-1}}{2k+1}\quad(1\leq k<N).
\end{equation}
The final coefficient equation $-2N p_{N-1}=2w_N$ follows
from existence of the polynomial already proved.  Exact comparison
of an algebraic input $q$ with the three rational candidates is
also effective.

\begin{corollary}[Median classification and a genuine exceptional quantile]
\label{p07:rga:median}
For positive half-integral beta shapes:
if at least one shape is integral, the median is algebraic;
if both shapes are nonintegral half-integers, the median is
algebraic exactly when the shapes are equal, in which case it is
$1/2$.  In every asymmetric nonintegral half-integer case the
median is transcendental.
Moreover,
\[
 Q_{3/2,5/2}(1/3)=1/4,
 \qquad Q_{5/2,3/2}(2/3)=3/4.
\]
All other algebraic levels for these two distributions have
transcendental quantiles.
\end{corollary}
\begin{proof}
The median statements follow from the theorem at $q=1/2$.
For the explicit exception, \eqref{p07:rga:recurrence} gives
\[
 P_{1,2}(X)=-2+\frac{28}{3}X-\frac{16}{3}X^2
           =\frac23(4X-1)(3-2X).
\]
Its only zero in $(0,1)$ is $1/4$.  The theorem proves the first
formula and excludes all other algebraic levels.  Reflection of
the beta distribution proves the second formula and its exclusion
statement.
\end{proof}

Four individually transcendental examples are the beta medians with
shape pairs $(1/2,3/2)$, $(1/2,5/2)$, $(3/2,5/2)$ and $(5/2,7/2)$.
These are examples of one
uniform family, not four independent historical conjectures.
The asymmetric exception above is also a concrete reason not to
replace the theorem by the false claim that symmetry supplies
the only algebraic interior beta quantiles.

\paragraph{Prior-art boundary.}
Dhillon and Saha's 2025 paper studies the arithmetic nature of
incomplete beta values \cite{p07:rga:DhillonSaha}. Their author slides
explicitly treat positive rational noninteger $a$ with
$a+b\in\mathbb N$ by logarithmic linear forms and Baker's theory
\cite[slide~4/4, compiled PDF p.~222]{p07:rga:SahaTalk}.
This includes the nonintegral half-shape parameter range above, so
neither that range nor the underlying logarithmic method is claimed
as new. The slides do not state the normalized inverse classification
or its three rational-angle tests. The paper's full theorem statements
were not available in the legal open sources checked: the institutional
record links back to the publisher DOI. Exact overlap with the present
classification therefore remains unverified. We present a self-contained
proved application without a claim of historical priority.

\subsection{What was transferred and what remains open}
\label{p07:rga:scope}

The useful transfer is from the effective logarithmic splitting
principle to a normalized inverse problem.  Normalization introduces
the period $\pi$, so ordinary algebraic-level inversion alone is
insufficient.  Baker's theorem removes the algebraic correction,
Gelfond--Schneider restricts the level to a rational angle, and
the one-real-zero property forces the entire Galois orbit to
collapse.  That last step reduces an apparently unbounded
cyclotomic search to three exact rational tests.

This classification does not apply to beta shapes $1/5$ and $2/5$:
the required polynomial correction of an arcsine primitive is
absent, and the normalizing gamma ratios are not algebraically
classified.  It likewise gives no finite logarithmic
representation for $A$.  The three named targets therefore stay
open in this investigation.  The beta theorem is a fully proved
connected-family result; it is not presented as an Annals-level
assessment or as a resolution of those named constants.


\section{Finite cyclotomic search and the Student and Fisher distributions}
\label{p07:rq:section}

The beta argument has a wider consequence. Its arithmetic step does
not require a beta density: it requires a rational polynomial
correction to an arcsine primitive. Counting the real roots of that
correction bounds the degree of an entire cyclotomic orbit. This
gives an explicit finite decision procedure for polynomial tilts of
the arcsine distribution.

\subsection{All rational polynomial tilts of the arcsine law}

Let $w(X)=\sum_{j=0}^d w_jX^j\in\mathbb Q[X]$ be nonzero and
nonnegative on $[0,1]$. Put
\[
 c_w=\frac1\pi\int_0^1\frac{w(t)}{\sqrt{t(1-t)}}dt
     =\sum_{j=0}^d w_j\frac{\binom{2j}{j}}{4^j}\in\mathbb Q_{>0},
 \qquad
 F_w(x)=\frac1{\pi c_w}\int_0^x\frac{w(t)}{\sqrt{t(1-t)}}dt.
\]
The moment identity follows from $t=\sin^2\theta$ and the elementary
recurrence for integrals of even powers of sine. A nonzero
nonnegative polynomial is positive apart from finitely many points;
thus $F_w$ is strictly increasing from zero to one. Denote its
unique $q$-quantile by $Q_w(q)$.

\begin{theorem}[An exact finite test for polynomially tilted quantiles]
\label{p07:rq:tilt}
Suppose $w$ is not constant. There is a unique nonzero polynomial
$P_w\in\mathbb Q[X]$ of degree at most $d-1$ such that
\begin{equation}\label{p07:rq:tiltCDF}
 F_w(x)=\frac{2\arcsin\sqrt{x}+\sqrt{x(1-x)}P_w(x)}{\pi}.
\end{equation}
Let $r$ be the number of its distinct real zeros in $(0,1)$.
There are at most $r$ algebraic numbers $q\in(0,1)$ whose
quantiles are algebraic. All of them, with their quantiles, are
obtained by this finite exact procedure:
\begin{enumerate}
\item Enumerate $3\le N\le8r^2$ and integers $1\le j<N/2$
with $\gcd(j,N)=1$ and $\varphi(N)\le2r$.
\item Set $q=2j/N$ and $x=\sin^2(\pi j/N)$, an explicitly
represented real algebraic number.
\item Retain $(q,x)$ if and only if $P_w(x)=0$.
\end{enumerate}
Every other algebraic $q\in(0,1)$ has a transcendental quantile.
If $r=0$ the list is empty. If $w$ is constant, the separate
arcsine rule applies: $Q_w(q)$ is algebraic exactly for rational $q$.
\end{theorem}
\begin{proof}
Expand $w(\sin^2\theta)$ as a rational linear combination of
$1,\cos(2\theta),\ldots,\cos(2d\theta)$. Its constant term
is $c_w$. Integrating, normalizing, and using
$\sin(2j\theta)=\sin(2\theta)U_{j-1}(\cos(2\theta))$
gives \eqref{p07:rq:tiltCDF}. Uniqueness and nonvanishing follow by
differentiation; explicitly
\[
 2X(1-X)P_w'(X)+(1-2X)P_w(X)=2\bigl(w(X)/c_w-1\bigr).
\]
The coefficient recurrence in \eqref{p07:rga:recurrence}, with
$w/c_w$ in place of $K_{m,n}X^m(1-X)^n$, computes $P_w$.
Real algebraic root isolation computes $r$ exactly.

For algebraic $q,x$ with $F_w(x)=q$, put
$z=1-2x+2i\sqrt{x(1-x)}$. As in
\eqref{p07:rga:Baker-step}, Baker's theorem applied to
$\log z-q\log(-1)=-i\sqrt{x(1-x)}P_w(x)$ forces
$P_w(x)=0$ and $z=\exp(i\pi q)$. Gelfond--Schneider forces
$q\in\mathbb Q$. Conversely a rational $q$ gives an algebraic
$x=\sin^2(\pi q/2)$, and it is the desired quantile exactly
when $P_w(x)=0$.

Let $N>2$ be the order of the resulting root of unity $z$.
Every conjugate of $x=(2-z-z^{-1})/4$ lies in $(0,1)$ and
is a zero of $P_w$. Moreover
$[\mathbb Q(x):\mathbb Q]=\varphi(N)/2$:
$z$ satisfies a quadratic equation over the real field
$\mathbb Q(z+z^{-1})$, is nonreal, and has cyclotomic degree
$\varphi(N)$. Consequently $\varphi(N)\le2r$.

For completeness, the elementary bound
$\varphi(N)\ge\sqrt{N/2}$ makes the search finite without
an unspecified inverse-totient step. By prime factorization,
\[
 \frac{\varphi(N)}{\sqrt N}
    =\prod_{p^e\parallel N}p^{e/2-1}(p-1).
\]
Every factor is at least one except possibly $p=2,e=1$,
whose factor is $1/\sqrt2$. Thus $N\le8r^2$.
The enumeration includes every primitive upper-semicircle root
of unity of each allowed order. Each retained pair satisfies
the CDF equation by \eqref{p07:rq:tiltCDF}; there are at most $r$
pairs because distinct levels give distinct roots in $(0,1)$.
This proves both completeness and soundness of the procedure.
\end{proof}

This theorem also applies to any explicitly given rational convex
mixture of the nonintegral half-integer beta laws: each density is
a rational polynomial times the arcsine density. The number and
locations of algebraic exceptions may increase after mixing;
the finite search, rather than the three-candidate beta rule,
is the applicable conclusion. The classical transcendence inputs
are credited explicitly. We assert no historical priority for
this formulation without a wider literature comparison.

\subsection{A complete parity classification for Student quantiles}

Let $T_\nu$ have the standard Student distribution with positive
integer degrees of freedom $\nu$, and let $t_\nu(q)$ be its
unique $q$-quantile. Its conventional density is
\[
 f_\nu(t)=\frac{\Gamma((\nu+1)/2)}
 {\sqrt{\nu\pi}\,\Gamma(\nu/2)}
 \left(1+\frac{t^2}{\nu}\right)^{-(\nu+1)/2};
\]
see \cite{p07:rq:Rstudent} for the normalization and beta-CDF relation.

\begin{theorem}[Arithmetic parity of Student quantiles]
\label{p07:rq:student}
For every algebraic $q\in(0,1)$,
\[
\begin{array}{c|l}
\text{degrees of freedom}&\text{arithmetic type of }t_\nu(q)\\\hline
\nu\text{ even}&\text{algebraic for every such }q,\\
\nu=1&\text{algebraic if and only if }q\in\mathbb Q,\\
\nu\ge3\text{ odd}&\text{algebraic if and only if }q=1/2.
\end{array}
\]
In the last row the sole algebraic quantile is zero; every other
algebraic-probability quantile is transcendental.
\end{theorem}
\begin{proof}
The strictly increasing change of variables
\[
 U=\frac12\left(1+\frac{T_\nu}{\sqrt{\nu+T_\nu^2}}\right),
 \qquad
 T_\nu=\frac{\sqrt\nu(2U-1)}{2\sqrt{U(1-U)}}
\]
has range $(0,1)$ and transforms the Student density into the
$\operatorname{Beta}(\nu/2,\nu/2)$ density. This can also be
checked directly: $1+t^2/\nu=1/(4u(1-u))$ and
$dt/du=\sqrt\nu/(4[u(1-u)]^{3/2})$; the duplication formula
for gamma gives the beta normalization.
The two displayed transformations are algebraic with algebraic
coefficients, so $t$ is algebraic if and only if its image $u$
is algebraic. They preserve the probability level.
Apply Theorem~\ref{p07:rga:beta-classification}. Even $\nu$ gives
integer beta shapes; $\nu=1$ gives the arcsine law; odd
$\nu\ge3$ gives equal nonintegral half-integer shapes with a
single interior correction zero at $u=1/2$. Its probability
level is $1/2$, and its inverse image is $t=0$.
\end{proof}

For example, at $q=39/40$ the standard $t_3$ critical value is
transcendental, while the standard $t_4$ critical value is
algebraic. This distinction is exact; decimal accuracy cannot
detect it. It is a classification of a familiar statistical
family, not a solution of a previously named conjecture about
Euler's constant or a gamma value.

\subsection{Fisher quantiles and the asymmetric exceptional value}

Write $f_{u,v}(q)$ for the standard Fisher $F$-distribution
quantile with positive integer degrees of freedom $u,v$.
If $X$ has that distribution then
$Y=uX/(v+uX)$ has the $\operatorname{Beta}(u/2,v/2)$
distribution, and $X=(v/u)Y/(1-Y)$.
These are mutually inverse rational maps over $\mathbb Q$.

\begin{corollary}[Arithmetic decision for integer-degree Fisher laws]
\label{p07:rq:fisher}
At every algebraic interior probability level, Fisher quantiles
are algebraic if at least one degree of freedom is even.
If both degrees are odd, all algebraic exceptions and all
transcendental cases are exactly those prescribed by
Theorem~\ref{p07:rga:beta-classification} for $(a,b)=(u/2,v/2)$,
transferred by $x\mapsto(v/u)x/(1-x)$. In particular,
\[
 f_{3,5}(1/3)=\frac59,\qquad f_{5,3}(2/3)=\frac95,
\]
and these are the only algebraic-probability algebraic quantiles
for their respective distributions.
\end{corollary}
\begin{proof}
The beta transformation preserves both algebraicity and the
probability level. For the displayed cases use
$Q_{3/2,5/2}(1/3)=1/4$ and its reflected counterpart in
Corollary~\ref{p07:rga:median}. Substitution gives $5/9$ and $9/5$.
The exclusion statements are inherited from the same corollary.
\end{proof}



\section{Lambert constants: finite certificates and the limit obstruction}
\label{p07:rsc:section}

This section investigates three specified targets:
\[
 E=\sum_{n\geq1}\frac1{2^n-1},\qquad
 R_F=\sum_{n\geq1}\frac1{F_n},\qquad
 B_1=\lim_{x\to\infty}\left(\sum_{p\leq x}\frac1p-\log\log x\right),
\]
where \(F_1=F_2=1\).  The questions are transcendence for the first two
numbers and irrationality for the third.  None of these three questions
is resolved here.  The positive conclusions concern explicitly specified
lacunary families, and the negative conclusions identify exact failures
of proposed transfers to the three original constants.

\subsection{Status and scope of the arithmetic input}

Duverney--Tachiya explicitly distinguish the established irrationality
of \(E\) from its unknown transcendence
\cite[Remark 1.1]{p07:rsc:DT}.  Campbell's May 2026 preprint proves a different
assertion, namely infinitely many occurrences of the binary block \(11\)
in \(E\); it does not claim transcendence \cite{p07:rsc:Campbell}.
Waldschmidt's survey lists the full reciprocal Fibonacci sum as
irrational, while leaving its transcendence unresolved
\cite[slide 59]{p07:rsc:Waldschmidt}.  This must not be confused with
transcendence results for even reciprocal powers or odd-index subseries.
The prime constant \(B_1\) has effective analytic representations and
high-precision computations \cite{p07:rsc:LZ,p07:rsc:KL}; no verified
irrationality proof for \(B_1\) was located in the present search,
completed on 17 October 2026.  Those primary sources support the
representations used below, not a claim that a computation proves
irrationality.

The functional transcendence inputs used in this section are classical.
For a rational function \(R\) over a number field, regular at zero with
\(R(0)=0\), the solution of
\[
 f(z)-f(z^d)=R(z),\qquad f(0)=0,\qquad d\geq2,
\]
is either rational or transcendental as a function; at a nonzero
algebraic point of the disk whose forward orbit avoids the poles of
\(R\), a transcendental solution has a transcendental value.  These
are the rational--transcendental dichotomy and the degree-one case of
Mahler's value theorem; see \cite[Theorem 4.1 and its discussion]{p07:rsc:Adam}.
General rational Mahler solving is already algorithmic
\cite{p07:rsc:CDDM}; complete discrete-residue obstructions to rational
summability are due to Arreche--Zhang \cite{p07:rsc:AZ}.
The elementary certificate below is a bounded-pole specialization of
that theory, with all of its finite bounds proved explicitly.

\subsection{An exact finite certificate for a cyclotomic kernel}

\begin{theorem}[Finite pole certificate]\label{p07:rsc:finite}
Let \(K\) be a number field, \(d\geq2\), and \(R\in K(z)\) satisfy
\(R(0)=0\).  Suppose all finite poles of \(R\) are \(N\)-th roots of
unity and have order at most \(m\).  If \(R\) is a polynomial, one may
take \(m=0,N=1\).  Let \(B\geq0\) bound the degree of its polynomial
part, counting a constant or zero polynomial part as degree zero.
Define, for \(|z|<1\),
\[
 f_R(z)=\sum_{k=0}^{\infty}R(z^{d^k}),\qquad
 D(z)=(1-z^N)^m,\qquad M=mN+\lfloor B/d\rfloor.
\]
For every algebraic \(\alpha\) with \(0<|\alpha|<1\), the following
conditions are equivalent:
\begin{enumerate}
\item \(f_R(\alpha)\) is algebraic;
\item \(f_R\) is rational;
\item there is \(A\in K[z]\), with \(A(0)=0,\deg A\leq M\), such that
\begin{equation}\label{p07:rsc:certificate}
 A(z)D(z^d)-A(z^d)D(z)
   =R(z)D(z)D(z^d).
\end{equation}
\end{enumerate}
Equation \eqref{p07:rsc:certificate} is a finite linear system over \(K\);
its right side is polynomial.  Its normalized solution, when present,
is unique and \(f_R=A/D\).  In particular, a failure certificate for
this linear system proves transcendence of every such disk value.
\end{theorem}

\begin{proof}
Since \(R(z)=O(z)\) at zero and has no poles in the open disk, the
series converges locally uniformly and satisfies
\(f_R(z)-f_R(z^d)=R(z)\).  This is the unique analytic solution with
zero constant term: the first nonzero coefficient of a difference of
two such solutions would contradict \(h(z)=h(z^d)\).
The classical inputs just stated prove the equivalence of (1) and
(2); every forward iterate of \(\alpha\) is in the punctured disk,
so the required regularity holds.

Suppose \(f_R=g\) is rational.  A pole of \(g\) outside the set
\(\mu_N=\{\zeta:\zeta^N=1\}\), other than zero, forces poles at all of
its successive \(d\)-th roots.  Indeed a \(d\)-th root of a point
outside \(\mu_N\) is still outside \(\mu_N\), \(R\) is regular there,
and \(g(z)-g(z^d)=R(z)\) forces cancellation by a pole of \(g(z)\).
There are arbitrarily many distinct such roots, which is impossible
for a rational function.  Zero is regular because \(g=f_R\) there.
Thus all finite poles of \(g\) belong to \(\mu_N\).

If \(g\) had a pole of maximal order \(s>m\), the same backward
argument would propagate poles of order \(s\) to every successive
\(d\)-th root: the rational function \(R\), whose pole orders are at
most \(m\), cannot cancel the order-\(s\) principal part.  Again this
would give infinitely many poles.  Hence \(Dg\) is polynomial.
If the polynomial part of \(g\) has positive degree \(b\), then the
polynomial part of \(g(z)-g(z^d)\) has degree exactly \(db\).
Proper rational parts remain proper under substitution, so
\(db\leq B\).  It follows that \(g=A/D\) with \(\deg A\leq M\).
Substitution gives \eqref{p07:rsc:certificate}.

Although this argument initially permits complex coefficients for
\(A\), the finite linear system has coefficients in \(K\).
Its unique normalized solution therefore belongs to \(K\), by
Gaussian elimination.  Conversely any solution of
\eqref{p07:rsc:certificate} yields a rational function analytic in the
disk with zero constant term satisfying the required equation.
Uniqueness identifies it with \(f_R\).
\end{proof}

The explicit degree bound is useful because it converts an existence
statement into a terminating exact procedure.  It does not extend
Mahler's theorem to an arbitrary infinite spectrum.

\begin{theorem}[All algebraic linear relations in a residue family]
\label{p07:rsc:residues}
Fix \(d\geq2\), an algebraic number \(0<|\alpha|<1\), and
\[
 U_j(\alpha)=\sum_{k=0}^{\infty}
     \frac{\alpha^{j d^k}}{1-\alpha^{d^{k+1}}},
       \qquad 1\leq j\leq d-1.
\]
For algebraic \(c_1,\ldots,c_{d-1}\),
\[
 \sum_{j=1}^{d-1}c_jU_j(\alpha)\ \hbox{is algebraic}
 \quad\Longleftrightarrow\quad c_1=\cdots=c_{d-1}.
\]
If their common value is \(c\), the sum is \(c\alpha/(1-\alpha)\).
Thus the vector space over \(\overline{\mathbb Q}\) spanned by
\(1,U_1(\alpha),\ldots,U_{d-1}(\alpha)\) has dimension exactly \(d-1\).
\end{theorem}

\begin{proof}
Apply Theorem~\ref{p07:rsc:finite} to
\[
 R(z)=\frac{\sum_{j=1}^{d-1}c_jz^j}{1-z^d}.
\]
If its normalized sum \(g\) is rational, its poles are simple and
belong to \(\mu_d\), and its polynomial part is constant.  If
\(\zeta\in\mu_d\setminus\{1\}\) were a pole of \(g\), any
\(\beta^d=\zeta\) would lie outside \(\mu_d\).  Then \(g(z^d)\)
would have a pole at \(\beta\), whereas \(g\) and \(R\) are both
regular there, a contradiction.  Therefore \(g\) has at most a
simple pole at \(1\), is bounded at infinity, and vanishes at zero.
Consequently \(g(z)=c z/(1-z)\).  The identity
\[
 \frac{z}{1-z}-\frac{z^d}{1-z^d}
       =\frac{z+z^2+\cdots+z^{d-1}}{1-z^d}
\]
proves both necessity and sufficiency of equal coefficients.
For the dimension statement, any algebraic linear relation involving
the \(U_j\)'s and \(1\) has all \(U_j\)-coefficients equal; the
displayed identity gives the unique relation up to a scalar.
\end{proof}

For an integer \(b\geq2\), the same theorem classifies exactly the
algebraic numbers in the finite-parameter family
\[
 \sum_{k\geq0}\sum_{j=1}^{d-1}
       \frac{c_j b^{(d-j)d^k}}{b^{d^{k+1}}-1}:
\]
the algebraic locus is the diagonal \(c_1=\cdots=c_{d-1}\), with
value \(c/(b-1)\).  This is an explicit application of classical
Mahler theory, not a newly solved historical conjecture.
For comparison, transcendence of the reciprocal Fermat sum
\(\sum_{k\geq0}(2^{2^k}+1)^{-1}\) was already known; Coons also
proved its irrationality exponent is \(2\) \cite{p07:rsc:Coons}.
It is not counted as an additional discovery here.

\begin{corollary}[Periodic weights on doubling Fibonacci indices]
\label{p07:rsc:periodic}
Let \((a_k)_{k\geq0}\) be a periodic sequence of algebraic numbers.
Then
\[
 V(a)=\sum_{k\geq0}\frac{a_k}{F_{2^k}}
\]
is algebraic if and only if \((a_k)\) is constant.  In the constant
case \(a_k=a\),
\[
 V(a)=a\,\frac{7-\sqrt5}{2}.
\]
\end{corollary}

\begin{proof}
For a period \(h\), put \(d=2^h\) and
\[
 H_a(z)=\sum_{k\geq0}a_k
           \frac{z^{2^k}}{1-z^{2^{k+1}}}.
\]
Grouping \(k\) modulo \(h\), its kernel for dilation by \(d\) is
\[
 \sum_{r=0}^{h-1}\frac{a_rz^{2^r}}{1-z^{2^{r+1}}}
   =\frac{\sum_{j=1}^{d-1}a_{v_2(j)}z^j}{1-z^d}.
\]
Indeed each \(1\leq j<2^h\) has a unique representation
\(j=2^r(2s+1)\).  Theorem~\ref{p07:rsc:residues} says that
\(H_a(\alpha)\) is algebraic precisely when
\(a_0=\cdots=a_{h-1}\).

Let \(\varphi=(1+\sqrt5)/2\) and \(q=\varphi^{-1}\).
Binet's formula, treating the odd index \(2^0=1\) separately, gives
\[
 V(a)=a_0+\sqrt5
       \sum_{k\geq1}a_k\frac{q^{2^k}}{1-q^{2^{k+1}}}.
\]
The last sum is \(H_b(q^2)\), where \(b_k=a_{k+1}\).
The shift of a periodic sequence is constant exactly when the
original sequence is constant.  In that case
\(H_b(q^2)=a q^2/(1-q^2)=a q\), proving the formula.
\end{proof}

The constant-weight identity and special nonconstant-weight
transcendence results have a long history; Nguyen records, for
example, Mignotte's result for \(2+(-1)^k\) \cite{p07:rsc:Nguyen}.
The corollary is a uniform proof for this restricted periodic
family, not a priority claim.  Nguyen's much stronger
Theorem 1.2 proves transcendence for arbitrary Fibonacci/Lucas
choices at indices with \(n_{k+1}/n_k\geq c>2\).
The exact algebraic doubling identity shows why the strict
inequality cannot simply be replaced by \(c=2\).

\subsection{Why the complete Lambert spectrum does not fit the certificate}

Let
\[
 L(z)=\sum_{n\geq1}\frac{z^n}{1-z^n}
     =\sum_{n\geq1}d(n)z^n,\qquad |z|<1.
\]
The next obstruction also follows from Bell--Smertnig's 2026
classification of multiplicative Mahler series \cite{p07:rsc:BS}.
We give a direct proof because it identifies the local quantities
that fail to close under a finite dilation system.

\begin{lemma}[Radial logarithmic coefficient]\label{p07:rsc:radial}
If \(\zeta\) is a root of unity of order \(q\), then, as \(t\downarrow0\),
\[
 L(\zeta e^{-t})
       =\frac1{qt}\log\frac1t+O_q(t^{-1}).
\]
\end{lemma}

\begin{proof}
The terms with \(q\mid n\) contribute \(L(e^{-qt})\).
For \(0<u\leq1\), the function
\((e^u-1)^{-1}-u^{-1}\) is bounded.  Splitting at
\(n=\lfloor1/t\rfloor\), and bounding the remaining terms by a
constant times \(e^{-nt}\), gives
\[
 L(e^{-t})=\frac1t\sum_{n\leq1/t}\frac1n+O(t^{-1})
         =t^{-1}\log(1/t)+O(t^{-1}).
\]
For \(q\nmid n\), the denominators
\(\lvert1-\zeta^n e^{-nt}\rvert\) are bounded below by a positive
constant depending only on \(q\).  Their total absolute contribution
is therefore \(O_q(\sum_{n\geq1}e^{-nt})=O_q(t^{-1})\).
\end{proof}

\begin{theorem}[No finite linear Mahler equation for \(L\)]
\label{p07:rsc:noMahler}
For every \(d\geq2\), every \(r\geq0\), and rational functions
\(a_0,\ldots,a_r,b\in\mathbb C(z)\), an identity
\begin{equation}\label{p07:rsc:mahler}
 \sum_{j=0}^{r}a_j(z)L(z^{d^j})=b(z)
\end{equation}
in the disk implies \(a_0=\cdots=a_r=b=0\).
\end{theorem}

\begin{proof}
Fix \(k\in\{0,\ldots,r\}\), and choose infinitely many distinct roots
\(\zeta\) of orders \(d^k\ell\), where \(\ell\) ranges over primes
coprime to \(d\), avoiding the finitely many poles of all the
rational functions in \eqref{p07:rsc:mahler}.
The order of \(\zeta^{d^j}\) is
\(d^{\max(k-j,0)}\ell\).
Apply Lemma~\ref{p07:rsc:radial} to \(z=\zeta e^{-t}\).
Since all coefficients and \(b\) are regular at \(\zeta\),
the coefficient of \(t^{-1}\log(1/t)\) must vanish.  It is
\[
 \frac1\ell\left(
       d^{-k}\sum_{j=0}^{k}a_j(\zeta)
           +\sum_{j=k+1}^{r}d^{-j}a_j(\zeta)\right).
\]
Terms arising from variation of the rational coefficients are only
\(O(\log(1/t))\); the other remainders are \(O(t^{-1})\).
Consequently the rational function
\[
 P_k(z)=d^{-k}\sum_{j=0}^{k}a_j(z)
                +\sum_{j=k+1}^{r}d^{-j}a_j(z)
\]
has infinitely many zeros, and so vanishes identically.
For \(1\leq k\leq r\),
\[
 P_{k-1}-P_k
       =(d^{-(k-1)}-d^{-k})\sum_{j=0}^{k-1}a_j=0.
\]
Successively \(a_0,\ldots,a_{r-1}\) vanish, and \(P_r=0\)
then gives \(a_r=0\).  For \(r=0\), \(P_0=0\) gives the same
conclusion directly.  Finally \(b=0\).
\end{proof}

In particular the full Lambert function cannot be a component of a
finite linear Mahler system over \(\mathbb C(z)\): successive
dilates of any one component of such a system belong to a
finite-dimensional vector space and would satisfy a nontrivial
relation of the forbidden form.  This is an obstruction to this
specific function-theoretic approach, not a theorem that \(E=L(1/2)\)
cannot admit some unrelated Mahler representation at a single point.

For the Fibonacci target define
\[
 \mathcal F(z)=\sum_{n\geq1}
       \frac{z^n}{1-(-1)^nz^{2n}},
       \qquad R_F=\sqrt5\,\mathcal F(\varphi^{-1}).
\]
Its even-index part is \(L(z^2)-L(z^4)\), while for real \(0<z<1\)
its odd-index part is bounded above by \(\sum_{n\ {\rm odd}}z^n\).
It follows rigorously that
\[
 \mathcal F(e^{-t})
       =\frac1{4t}\log\frac1t+O(t^{-1}).
\]
Thus, for every \(d\geq2\),
\[
 \mathcal F(z)-\mathcal F(z^d)\notin\mathbb C(z):
\]
along \(z=e^{-t}\) the difference has leading term
\((1-d^{-1})(4t)^{-1}\log(1/t)\), incompatible with a rational
function.  This rules out the direct additive certificate, but
does not purport to exclude every possible higher-order equation
for \(\mathcal F\).

\subsection{Denominators and finite prime logarithms}

For clarity, the familiar positive-tail argument has the following
precise limitation.  If \(s_N=P_N/D_N\), \(P_N\in\mathbb Z\),
\(D_N\in\mathbb Z_{>0}\),
and \(S-s_N=\varepsilon_N>0\), then
\[
 D_N\varepsilon_N\longrightarrow0
 \quad\Longrightarrow\quad S\notin\mathbb Q.
\]
Indeed, if \(S=u/v\), the positive integer
\(vD_N(S-s_N)=uD_N-vP_N\) would tend to zero.
For the full Erd\H{o}s--Borwein sum, the product denominator
\(D_N=\prod_{n=1}^N(2^n-1)\) gives
\[
 D_N\varepsilon_N
   >2^{N(N-1)/2-N},
\]
because \(2^n-1\geq2^{n-1}\) and
\(\varepsilon_N>\sum_{n>N}2^{-n}=2^{-N}\).
For the Fibonacci sum, \(F_n\geq\varphi^{n-2}\) and
\(F_n\leq\varphi^{n-1}\), so the product denominator gives
\[
 \left(\prod_{n=1}^NF_n\right)
       \sum_{n>N}\frac1{F_n}
 \geq\varphi^{N(N-3)/2-N}.
\]
Both grow without bound.  These calculations reject the stated
product-denominator certificate; they are not lower bounds for
the reduced denominators and do not reject optimized approximants.
Most importantly, even a successful irrationality certificate
would still not prove either requested transcendence statement.

The prime target admits
\begin{equation}\label{p07:rsc:mertens}
 B_1-\gamma=\sum_p\Delta_p,\qquad
 \Delta_p=\log(1-1/p)+1/p.
\end{equation}
This follows by adding the second and logarithmic third Mertens
formulas, with the convention for \(B_1\) fixed above
\cite{p07:rsc:LZ,p07:rsc:KL}.

\begin{proposition}[No finite rational cancellation of prime corrections]
\label{p07:rsc:prime}
Every nonzero finite rational linear combination of the numbers
\(\Delta_p\) is transcendental.  Consequently
\(\{1\}\cup\{\Delta_p:p\text{ prime}\}\) is linearly independent over
\(\mathbb Q\).  For every integer \(X\geq2\), however, the omitted
tail of \eqref{p07:rsc:mertens} satisfies only the explicit bounds
\[
 \frac1{8X^2}
       < -\sum_{p>X}\Delta_p
       \leq\frac1{2X}.
\]
\end{proposition}

\begin{proof}
After clearing rational coefficients, the logarithmic part is
\(\log A\), with \(A=\prod_p((p-1)/p)^{n_p}\in\mathbb Q_{>0}\).
If not all integers \(n_p\) vanish, take the largest prime \(P\)
with \(n_P\neq0\).  No numerator \(p-1\), for \(p\leq P\), is
divisible by \(P\); among the denominators only \(P\) is.  Hence
\(v_P(A)=-n_P\neq0\), and \(A\neq1\).
Hermite--Lindemann gives transcendence of the real logarithm
\(\log A\).  Adding the rational contribution \(\sum n_p/p\)
preserves transcendence, proving the first assertion and the
linear independence.

The convergent expansion
\[
 -\Delta_p=\sum_{j\geq2}\frac1{jp^j}
\]
lies between \(1/(2p^2)\) and \(1/(2p(p-1))\).
Bertrand's postulate supplies a prime \(X<p<2X\), giving the
strict lower bound.  Summing the upper bound over all integers
\(n>X\), rather than only primes, telescopes to \(1/(2X)\).
\end{proof}

This proposition is a direct classical logarithmic-transcendence
transfer, not new information about the irrationality of \(B_1\).
It makes the missing step exact: finite noncancellation is
available, but neither an arithmetic statement about the infinite
sum \(\sum_p\Delta_p\) nor a noncancellation theorem between that
sum and \(\gamma\) follows.  Convergence, explicit tail signs, and
finite linear independence do not supply either statement.

\paragraph{Outcome.}
The three named targets remain unclosed.  The completed work gives
an explicit finite certificate, a full algebraic-exception locus
for a lacunary family, a periodic Fibonacci-weight classification,
and direct obstruction proofs for transfers to the full spectra.
These results reuse classical machinery and contain known special
cases.  They must not be reported as three new resolutions of
the original constant problems.


\section{Euler and Gompertz constants: finite algebraic extraction and fixed-kernel lattices}
\label{p07:rgam:section}

The two external targets are the irrationality or transcendence of
Euler's constant $\gamma$ and of the Euler--Gompertz constant
\[
 \delta=\int_0^\infty e^{-t}\log(1+t)\,dt
       =\int_0^\infty\frac{e^{-t}}{1+t}\,dt.
\]
The equality follows by integration by parts.  Neither individual target
is resolved here.  The published approximation results of
Wolfs--Van Assche \cite[Theorems 1.2--1.3]{p07:rgam:WV}, and the primary
papers \cite{p07:rgam:MY,p07:rgam:Powers}, explicitly distinguish these targets
from the known disjunctive transcendence statement.  Our current-source
check found no independently verified individual resolution.
The progress below consists of two reusable obstruction theorems, with
complete proofs: finite algebraic extraction cannot retain a common
logarithmic coordinate while giving finite monodromy; and every fixed
positive polynomial denominator gives an attained positive gap on the
integer-polynomial numerator cone.  These strengthen earlier linear
cancellation and single-denominator arguments.  Their priority as
standalone formulations is not claimed.

Write
\begin{equation}\label{p07:rgam:connection}
 \delta_a=\int_0^\infty\frac{e^{-t}}{t+a}\,dt,
 \qquad \operatorname{Ein}(z)=\sum_{n\ge1}\frac{(-1)^{n-1}z^n}{n\,n!},
 \qquad
 e^{-a}\delta_a=\operatorname{Ein}(a)-\gamma-\log a
 \quad(a>0).
\end{equation}
The differential identity for the last expression, followed by the
usual defining limit of $\gamma$, proves the connection formula.
In particular $\gamma+\delta/e=\operatorname{Ein}(1)$; it is not
$\gamma+\delta$.  The earlier $E$-function arguments apply to this
combination.  The new ring morphism in \cite[Proposition 3.1]{p07:rgam:FR}
has also been checked against its primary text: its coefficients need
not be algebraic, so it does not separate $\gamma$ and $\delta$ simply
by displaying them as two coefficients.  Neither obstruction theorem
below assumes that ring morphism or any conjecture about its image.

\subsection{Finite covers do not eliminate a logarithm by retaining its constant term}

The failed logical passage from finite consistency to a suitably
well-founded model has an arithmetic analogue here: a finite algebraic
cover is a much stronger constraint than convergence to an analytic
boundary constant.  The following precise statement supplies the
constraint; the analogy itself supplies no proof.

\begin{theorem}[Finite-orbit algebraic elimination]\label{p07:rgam:finiteorbit}
Let $F$ be a field of characteristic zero, let $T_1,\ldots,T_r$ be
algebraically independent over $F$, and let $w$ be algebraic over
$F(T_1,\ldots,T_r)$.  Suppose that, in a common algebraic extension,
for each $i$ there is an automorphism $\tau_i$ fixing $F$, sending
$T_i$ to $T_i+c_i$ for some $c_i\in F^\times$, and fixing the other
$T_j$.  If each orbit $\{\tau_i^n(w):n\in\mathbb Z\}$ is finite,
then $w$ is algebraic over $F$.  If additionally
$w\in F(T_1,\ldots,T_r)$, then $w\in F$.
\end{theorem}
\begin{proof}
First eliminate $T_1$, putting $F_1=F(T_2,\ldots,T_r)$.  Clear the
denominators of an algebraic relation to obtain a nonzero
$P(T_1,w)=0$ with $P(X,V)\in F_1[X,V]$.  Finiteness of the orbit
provides an integer $N\ge1$ with $\tau_1^N(w)=w$.  Applying its
integer powers gives
\[
                 P(T_1+nNc_1,w)=0\qquad(n\in\mathbb Z).
\]
The polynomial $P(T_1+X,w)$ over the extension field has infinitely
many distinct roots $X=nNc_1$, hence is identically zero.  Translation
of the indeterminate is invertible, so $P(X,w)$ is identically zero.
Writing $P(X,V)=\sum_jp_j(V)X^j$, at least one $p_j\in F_1[V]$
is nonzero and satisfies $p_j(w)=0$.  Thus $w$ is algebraic over $F_1$.
Apply the same argument to $T_2$, and continue.  No commutativity of
the automorphisms on the algebraic extension is needed.

Finally $F$ is relatively algebraically closed in a pure rational
function field over $F$.  For completeness, if $h$ is rational and
algebraic over $F$, differentiating its separable minimal polynomial
with respect to each $T_i$ gives $\partial h/\partial T_i=0$.
In characteristic zero a rational function with zero derivative in
one variable is independent of that variable.  Repetition proves
$h\in F$.
\end{proof}

\begin{corollary}[Nonlinear common-logarithm obstruction]
\label{p07:rgam:commonlog}
Let $a_1,\ldots,a_m>0$ be algebraic, let
$k=\overline{\mathbb Q}(\log a_1,\ldots,\log a_m)$ with real
logarithms, and consider the single-valued meromorphic-germ field
\[
 F=k\bigl(z,e^{-a_1z},\ldots,e^{-a_mz},
                \operatorname{Ein}(a_1z),\ldots,
                \operatorname{Ein}(a_mz)\bigr).
\]
On a logarithmic covering of a punctured neighborhood of zero put
\[
 T=\gamma+\log z,\qquad
 Y_j=e^{-a_jz}\delta_{a_jz}
       =\operatorname{Ein}(a_jz)-\log a_j-T.
\]
Suppose that a germ $w$ algebraic over $F(Y_1,\ldots,Y_m)$ admits
continuation through repeated turns around zero, and has a finite
orbit under that continuation.  Then $w$ is algebraic over $F$.
If $w$ is rational in the $Y_j$ over $F$, then $w\in F$.
If in addition $w$ is a constant germ, then respectively $w$ is
algebraic over $k$, or $w\in k$ in the rational case.
All finite derivatives of the displayed data can be included without
changing these assertions.
\end{corollary}
\begin{proof}
Every element of $F$ is single valued and has a Laurent expansion at
zero with coefficients in $k$.  A turn fixes $F$ and sends $T$ to
$T+2\pi i$.  The element $T$ is transcendental over $F$: an alleged
polynomial relation, continued around the origin, would give infinitely
many roots of a nonzero polynomial at any point where its leading
coefficient is nonzero.  Although $2\pi i$ need not belong to $k$,
we may first enlarge $F$ by this constant for this transcendence
argument.  For the elimination itself one does not need $2\pi i\in F$:
the polynomial-in-$X$ proof of Theorem~\ref{p07:rgam:finiteorbit} takes
place in an extension field containing the nonzero shift, and its
coefficient polynomials still lie in $F$.  Thus it applies without
changing the coefficient field in the conclusion.

The displayed identities give $F(Y_1,\ldots,Y_m)=F(T)$.
Analytic continuation gives the required translation and its action
on the algebraic germs.  More directly, applying the proof just given
to an algebraic relation $P(T,w)=0$ and a turn that fixes $w$ yields
$w$ algebraic over $F$.  Relative algebraic closedness gives the
rational assertion.

For a constant $w$ algebraic over $F$, clear an algebraic relation
into $\sum_jf_j(z)w^j=0$, where not all $f_j\in F$ vanish.  Choose
the lowest Laurent power with a nonzero coefficient in some $f_j$.
Its coefficient yields a nonzero polynomial over $k$ vanishing at $w$.
If $w\in F$ is constant, its constant Laurent coefficient lies in $k$.
Finally $F$ is a differential field, and
$Y_j'=-e^{-a_jz}/z\in F$.  All higher derivatives of $Y_j$ therefore
already belong to $F$.
\end{proof}

For $a_1=1$, the coefficient field is simply $\overline{\mathbb Q}$.
The statement covers arbitrary finite polynomial equations and rational
operations, rather than just linear combinations or differentiation.
It does not assert that the numerical value of every single-valued
expression at $z=1$ is algebraic.  Such values can be transcendental
$E$-values.  Nor does it show that $\gamma$ is algebraic: a formula
extracting the constant germ $\gamma$ by the stipulated operations has
not been constructed.  At the positive boundary,
\[
 \gamma=\lim_{x\to\infty}
           \bigl(\operatorname{Ein}(x)-\log x\bigr),
 \qquad
 0<\operatorname{Ein}(x)-\log x-\gamma
      =e^{-x}\delta_x<e^{-x}/x.
\]
This limit is outside the finite algebraic operation hypothesis.
Exponential convergence is analytic information, not a height estimate
for rational linear forms in $\gamma$.

\subsection{A positive gap for every fixed rational Laplace kernel}

The next result turns positivity into an effective finite-dimensional
reduction.  It is not an inference from compactness alone: the factorial
coercivity below supplies the missing uniform degree bound.

\begin{lemma}[Weighted Laguerre lower bound]\label{p07:rgam:Laguerre}
Let $D(t)=\sum_{j=0}^{m}d_jt^j\in\mathbb R[t]$ be positive on
$[0,\infty)$, put $C_D=\sum_j|d_j|$, and let $Q$ have degree $n$
and leading coefficient $q_n$.  For $\alpha\in\{0,1\}$,
\begin{equation}\label{p07:rgam:Lagbound}
 \int_0^\infty \frac{t^\alpha e^{-t}Q(t)^2}{D(t)}\,dt
 \ge
 \frac{q_n^2 n!(n+\alpha)!}
 {C_D[4(n+m+1)+2]^m}.
\end{equation}
\end{lemma}
\begin{proof}
Let $L_n^{(\alpha)}$ denote the usual Laguerre polynomial, with leading
coefficient $(-1)^n/n!$.  Normalize it by
$\ell_n=L_n^{(\alpha)}/\sqrt{h_n}$, where $h_n=(n+\alpha)!/n!$.
Rodrigues' formula and integration by parts
show that the $\ell_n$ are orthonormal for $t^\alpha e^{-t}dt$ and
\[
 \int_0^\infty t^\alpha e^{-t}Q(t)\ell_n(t)dt
                 =(-1)^nq_n\sqrt{n!(n+\alpha)!}.
\]
Weighted Cauchy--Schwarz bounds the square of this expression by the
left side of \eqref{p07:rgam:Lagbound} times
$\int t^\alpha e^{-t}\ell_n(t)^2D(t)dt$.

Multiplication by $t$ has a three-term recurrence.  Its diagonal is
$2j+\alpha+1$; its off-diagonal magnitudes are
$\sqrt{(j+1)(j+\alpha+1)}$.
In its truncation to indices $0,\ldots,K$ every absolute row sum is
at most $4(K+1)+2\alpha$.  The symmetric matrix therefore has operator
norm at most that number.  Take $K=n+m$.  For $0\le j\le m$,
the matrix element calculating
$\int t^{j+\alpha}e^{-t}\ell_n^2dt$ is unchanged by this truncation:
a length-$j$ path starting at index $n$ cannot reach an index above
$n+m$.  Hence that integral is at most
$[4(n+m+1)+2]^j$.  Summing with the coefficients $|d_j|$ bounds the
remaining factor by $C_D[4(n+m+1)+2]^m$, proving the claim.
\end{proof}

\begin{theorem}[Fixed-denominator lattice gap and finite candidate reduction]
\label{p07:rgam:gap}
Let $D\in\mathbb Q[t]$ be fixed and strictly positive on $[0,\infty)$.
For nonzero $P\in\mathbb Z[t]$ nonnegative on that half-line put
\[
                      J_D(P)=\int_0^\infty e^{-t}\frac{P(t)}{D(t)}dt.
\]
Then
\[
                    \inf_P J_D(P)>0,
\]
and this infimum is attained.  A finite list of integer polynomials
containing every minimizer can be computed from $D$.  This means an
exact finite candidate reduction, not a claim to decide equality of
arbitrary transcendental constants.  More precisely, if $\deg P=2n$
then
\[
 J_D(P)\ge \frac{(n!)^2}{C_D[4(n+m+1)+2]^m},
\]
and if $\deg P=2n+1$ the numerator can be replaced by $n!(n+1)!$.
\end{theorem}
\begin{proof}
A real polynomial nonnegative on $[0,\infty)$ can be written
$P=A^2+tB^2$ with real polynomials $A,B$.  Here is a factorization proof.
A positive real root has even multiplicity and contributes a square.
The factors $t$ and $t+a$ ($a>0$) have the required form.  For an
irreducible quadratic $t^2-2ut+u^2+v^2$, write
\[
 (t-\sqrt{u^2+v^2})^2+
                 2(\sqrt{u^2+v^2}-u)t.
\]
Positive constants also have the required form, and multiplication
preserves it by
\[
 (A^2+tB^2)(C^2+tE^2)=(AC-tBE)^2+t(AE+BC)^2.
\]
This proves the representation.  The nonnegative leading terms show
that for degree $2n$, $A$ has degree $n$, $B$ has degree at most $n-1$,
and the square of the leading coefficient of $A$ is the leading
coefficient of $P$.  For degree $2n+1$ the analogous leading term is
that of $tB^2$.  The leading coefficient of a nonzero such integer
polynomial is at least one.  Applying Lemma~\ref{p07:rgam:Laguerre} to
this leading square proves both displayed bounds.

They tend to infinity with $n$.  Their even-degree lower bound is
increasing once $n\ge\max(m,1)$: the ratio of consecutive terms is
\[
 (n+1)^2\left(1-\frac1{n+m+5/2}\right)^m
 \ge (n+1)^2\left(1-\frac{m}{n+m+5/2}\right)>1.
\]
A rational $b_D>0$ with $D(t)\ge b_D$ for all $t\ge0$ is effectively
obtainable by real algebraic root isolation and an elementary leading
term bound outside a bounded interval.  Thus $J_D(1)\le B:=1/b_D$.
The preceding monotonicity gives a computable degree bound $N\ge1$
for every $P$ with $J_D(P)\le B$.

There is also a computable coefficient bound at those degrees.  On
$[0,1]$, $e^{-t}/D(t)\ge1/(3C_D)$, so
$\int_0^1P\le3C_DB$.  Set $M=\max_{[0,1]}P$.  The shifted Chebyshev
expansion, whose coefficients in positive degree have absolute value
at most $2M$, gives $\|P'\|_\infty\le4MN^3$.  On an interval of
length $1/(8N^3)$ next to a maximizer, $P\ge M/2$.  Consequently
$M\le48N^3C_DB$.  The coefficient $\ell^1$ norm of
$T_j(2t-1)$ is at most $7^j$, by its recurrence; hence every monomial
coefficient of $P$ has absolute value at most
\[
                         144N^3C_DB\,7^N.
\]
There are only finitely many integer polynomials satisfying these
bounds.  Their nonnegativity on the half-line is decidable by real
root isolation.  The resulting finite list contains $1$ and every
possible minimizer.  Every nonzero member has strictly positive
integral, so their minimum is positive and attained; any polynomial
outside the list has integral greater than $B\ge J_D(1)$.
\end{proof}

The constants are deliberately coarse; the theorem covers an entire
fixed rational kernel, not just a predetermined degree or a particular
orthogonal-polynomial sequence.  It does not apply uniformly when $D$
varies.  For example $J_{(1+t)^m}(1)\le1/(m-1)$ for $m>1$, by
omitting $e^{-t}$, so the unscaled integrals tend to zero as pole order
grows.  The arithmetic clearance of those small integrals is a separate
matter, addressed next.

\subsection{Exact arithmetic reduction and the remaining Gompertz loophole}

\begin{proposition}[Twisted residues classify the finite kernel quotient]
\label{p07:rgam:residues}
Let $k$ be a characteristic-zero subfield of $\mathbb C$, and let
$R\in k(t)$ have poles only among $-a_1,\ldots,-a_s$, where the
$a_i\in k$ are distinct positive real numbers.  There are unique
$Q\in k(t)$ and $b_i\in k$, with $Q$ having poles only at these points,
such that
\begin{equation}\label{p07:rgam:reduction}
 R=Q'-Q+\sum_i\frac{b_i}{t+a_i}.
\end{equation}
If the principal part at $-a_i$ is
$\sum_{j=1}^{m_i}c_{ij}(t+a_i)^{-j}$, then
\[
 b_i=\sum_{j=1}^{m_i}\frac{(-1)^{j-1}c_{ij}}{(j-1)!},
 \qquad
 \int_0^\infty e^{-t}R(t)dt=-Q(0)+\sum_i b_i\delta_{a_i}.
\]
Thus the quotient by rational twisted derivatives $Q'-Q$ has precisely
one coordinate per allowed pole.  This is a classification of rational
functions modulo that differential operator; it is not a claim that
the numerical constants $1,\delta_{a_1},\ldots,\delta_{a_s}$ are
linearly independent.
\end{proposition}
\begin{proof}
For a pole of order $j\ge2$, the leading term of
$(d/dt-1)(t+a_i)^{-j+1}$ is $-(j-1)(t+a_i)^{-j}$.
Successively remove these terms; only simple poles remain.  On
polynomials the operator is invertible, its inverse at $H$ being
$-\sum_{j\ge0}H^{(j)}$, a finite sum.  This proves existence.
If the difference of two decompositions were nonzero, a pole of $Q$
of maximal order would give a pole of strictly larger order in $Q'-Q$,
which simple poles cannot cancel.  Thus $Q$ would be polynomial;
then its leading term shows $Q'-Q$ cannot be a nonzero sum of simple
poles or zero.  This proves uniqueness.
Multiplication by $e^{-t}$ makes $Q'-Q$ an exact derivative, so its
residue at every finite pole is zero.  Taking the residue of
$e^{-t}R$ at $-a_i$ and dividing by $e^{a_i}$ gives the formula for
$b_i$.  Integration of $(e^{-t}Q)'$ yields the stated boundary term,
since $Q$ is regular at zero and $e^{-t}Q(t)\to0$ at positive infinity.
\end{proof}

For the individual target $\delta$, take $D=(1+t)^m$, $m\ge1$, and
$P\in\mathbb Z[t]$.  Ordinary division and partial fractions give
\[
 \frac{P(t)}{(1+t)^m}=H(t)+\sum_{j=1}^m\frac{c_j}{(1+t)^j},
                    \qquad H\in\mathbb Z[t],\quad c_j\in\mathbb Z.
\]
Put $I_j=\int_0^\infty e^{-t}(1+t)^{-j}dt$.  Integration by parts
shows $(j-1)I_j=1-I_{j-1}$ for $j\ge2$, with $I_1=\delta$.
Consequently there are explicitly computable integers $A_m(P),B_m(P)$
such that
\begin{equation}\label{p07:rgam:integerform}
 (m-1)!J_{(1+t)^m}(P)=A_m(P)+B_m(P)\delta,
 \quad
 B_m(P)=(m-1)!\sum_{j=1}^m\frac{(-1)^{j-1}c_j}{(j-1)!}.
\end{equation}
The coefficient identity is formal and makes no irrationality assumption.
If $P\ge0$ is nonzero, the form is strictly positive, giving the
required nonvanishing certificate.  For any fixed $m$,
Theorem~\ref{p07:rgam:gap} proves that these positive integer forms cannot
tend to zero, regardless of how numerator degrees or coefficients grow.

There are precise ways outside this theorem.  If
$g_m(P)=\gcd(A_m(P),B_m(P))$ is unbounded, division by this actual
common factor replaces the numerator by the rational polynomial
$P/g_m(P)$; the integer lattice hypothesis is lost.  No bound on that
gcd is proved here.  Allowing $m$ to grow or allowing sign-changing
numerators also escapes the theorem.  In the elementary growing-order
example $P=1$, however, \eqref{p07:rgam:integerform} has $B_m=\pm1$, so
there is no gcd saving at all, and
\[
 (m-1)!I_m\ge
       (m-1)!\frac{e^{-1/m}}{m(1+1/m)^m}
       \ge \frac{(m-1)!}{9m}\longrightarrow\infty.
\]
Here the integral was restricted to $0\le t\le1/m$ and the last
bound uses $e<3$.  Thus this explicit exponentially weighted rational
kernel attempt supplies nonzero forms and exact denominators, but its
cleared errors grow instead of shrink.  It does not resolve $\delta$.

\subsection{The separate approximation test for Euler's constant}

For comparison, write $d_n=\operatorname{lcm}(1,\ldots,n)$.
The verified published construction \cite[Theorem 1.2]{p07:rgam:WV} gives
integer approximants
$P_n=n!d_nF_{n;1}^{(I)}(1)$ and
$Q_n=n!d_nF_{n;2}^{(I)}(1)>0$ with
\[
 \left|\gamma-\frac{P_n}{Q_n}\right|
 \le O\!\left(e^{-4n^{3/4}+n^{1/2}+(3/8)n^{1/4}}\right),
\]
while
\[
 Q_n\asymp n!d_n\,n^{-9/8}
       e^{4n^{3/4}-(1/2)n^{1/2}-(3/8)n^{1/4}}.
\]
These estimates yield only
$|Q_n\gamma-P_n|\le O(n!d_n n^{-9/8}e^{n^{1/2}/2})$.
This is a growing upper bound, not a proof that the forms grow, and
not a proof that they approach zero.  Moreover a quoted upper bound
alone is not a nonvanishing certificate.  In reduced terms, with
$g_n=\gcd(P_n,Q_n)$, a successful use of this particular upper estimate
would require $g_n$ to dominate that factorial-scale bound along an
infinite subsequence, together with nonvanishing there.  Neither input
is supplied by finite-cover elimination, by analytic convergence, or by
the prior uniqueness/observation arguments.  The primary paper itself
states that its improvement does not prove irrationality.

Thus the individual $\gamma$ attempt has an exact nonlinear algebraic
barrier and a separate explicit denominator/error deficit.  The
individual $\delta$ attempt has an exact kernel quotient, nonvanishing
positive integer forms, and a proved fixed-kernel gap.  The two new
obstruction theorems are structural progress within specified routes;
they are not replacements for either missing individual arithmetic proof.


% Fragment: recycled arithmetic, zeta/Catalan branch. Prefix zrec:.
\section{A rank obstruction for arithmetic moment determinants}
\label{p07:zrec:section}

The named targets of this section are Catalan's constant $G$, the irrationality
and transcendence of $\zeta(5)$, and the transcendence of $\zeta(3)$.
The principal result below concerns an actual attempted method, rather than
claiming a new resolution of one of these targets.  It proves that the
superexponential smallness of certain positive Hankel determinants is lost
under \emph{every} rational normalization to an integer polynomial.  The
reason is the rank of the dependence on the desired constant.  This gives a
precise limitation on transferring the earlier moment-positivity and
finite-witness arguments to arithmetic.

\subsection{Current status and external results, checked on 17 October 2026}
\label{p07:zrec:status}

There are important changes after the usual older accounts of these constants.
Fauzan's September 2026 preprint proves irrationality of the real number
$\zeta(5)$ by normalized polynomial determinants \cite{p07:zrec:Fauzan}.
Independent Lean developments have since been released.  We inspected the
actual target, definition and axiom files of the Romik repository:
its target is irrationality of $\sum_{v\geq0}(v+1)^{-5}$, and its declared
additional axiom is precisely the prime number theorem in the form
$\vartheta(x)/x\to1$ \cite{p07:zrec:Romik}.  We did not replay its complete Lean
build in this investigation.  Lai's October 2026 research paper explicitly
records Fauzan's result as Theorem~1.5 \cite{p07:zrec:Lai}.  Thus it would be
inaccurate to present irrationality of $\zeta(5)$ as an untouched open target,
or to count it as our own achievement.

For $G$, Sun's September claim \cite{p07:zrec:Sun} must be distinguished from the
later OpenAI release.  Wachs identifies an unaccounted quadratic $2$-adic
contribution in Sun's denominator calculation \cite{p07:zrec:Wachs}; that criticism
does not address the different OpenAI proof.  The official OpenAI family~005
scope states that Catalan irrationality has been formalized, and supplies the
paper and a comparator statement for the exact alternating series
\cite{p07:zrec:OpenAI}.  Lai also records this later result.  A comparator is a
specification, not by itself the proof implementation; no independent kernel
replay is claimed here.  These are externally announced/formalized
irrationality results.  We found no verified individual transcendence result
for $G$, $\zeta(5)$ or $\zeta(3)$ in the primary material examined.

Brown's 2026 Mellin-period criterion and Fauzan's work demonstrate that
positive determinant methods can have genuine arithmetic power
\cite{p07:zrec:Brown,p07:zrec:Fauzan}.  The obstruction below applies to a particular
rank-constrained construction; it does not refute those richer constructions.
The elementary determinant and Christoffel identities used in its proof are
classical.  Our contribution here is their explicit combination to diagnose
this proposed transfer and to calculate its exact loss.

\subsection{Integer coefficients versus the rank of a positive perturbation}
\label{p07:zrec:rank}

\begin{theorem}[Rank-sensitive integer-normalization bound]
\label{p07:zrec:rankbound}
Let $\theta\in\mathbb R$, let $M$ be a symmetric positive definite $n\times n$ real
matrix, and let $U$ be an $n\times r$ real matrix of rank $r\geq1$.  Suppose
\[
         D(X)=\det\bigl(M+(X-\theta)UU^{\mathsf T}\bigr)
                 \in\mathbb Q[X].
\]
Put $K=U^{\mathsf T}M^{-1}U$.  Then $D$ has degree $r$, its leading coefficient
is $\det(M)\det(K)>0$, and every nonzero rational $c$ for which
$P=cD\in\mathbb Z[X]$ satisfies
\begin{equation}
          |P(\theta)|\geq \frac1{\det K}.
\label{p07:zrec:rankineq}
\end{equation}
In particular, if the columns $u_j$ satisfy
$u_j^{\mathsf T}M^{-1}u_j\leq L$, then $|P(\theta)|\geq L^{-r}$.
For $r=0$, a nonzero integral normalization is a nonzero constant integer,
so its absolute value is at least one.
\end{theorem}
\begin{proof}
The matrix determinant identity gives
\[
 D(X)=\det(M)\det\bigl(I_r+(X-\theta)K\bigr).
\]
Since $M$ is positive definite and $U$ has full column rank, $K$ is positive
definite.  Its determinant is positive, proving both the degree and the
leading-coefficient assertion.  The leading coefficient of $P$ is a nonzero
integer.  Its absolute value is therefore at least one.  As
$P(\theta)=c\det(M)$, division by $\det K$ gives
\eqref{p07:zrec:rankineq}.  Hadamard's inequality applied to the positive definite
matrix $K$ yields $\det K\leq\prod_j K_{jj}\leq L^r$.
\end{proof}

For example, if $r=r_n$ and $L\leq C N_n^a$, then
\[
       -\log|P_n(\theta)|\leq r_n(\log C+a\log N_n).
\]
Consequently $r_n\log(2N_n)=o(n^2)$ rules out a bound
$0<|P_n(\theta)|\leq e^{-\delta n^2}$ for any fixed $\delta>0$.
This is an obstruction independent of any conjecture about $\theta$ and
independent of how accurately denominators are estimated.

\begin{corollary}[Polynomial moment spaces and evaluation rank one]
\label{p07:zrec:christoffel}
Let $\mu$ be a finite positive measure with infinite support, and let
$V\subset\mathbb R[x]$ be a finite-dimensional space with basis
$p_1,\ldots,p_n$.  Define
\[
 M_{ij}=\int p_i p_j\,d\mu,
 \qquad u_i=p_i(a),\qquad
 K_V(a,a)=\sup_{0\ne p\in V}\frac{|p(a)|^2}{\int p^2\,d\mu}.
\]
If $u\ne0$, then $K_V(a,a)=u^{\mathsf T}M^{-1}u$ and
\begin{equation}
 \frac{\det M}{[X]\det\bigl(M+(X-\theta)uu^{\mathsf T}\bigr)}
      =\frac1{K_V(a,a)}.
\label{p07:zrec:exactratio}
\end{equation}
Here $[X]$ denotes the coefficient of $X$; the determinant is affine.
If that determinant has rational coefficients, every nonzero rational
normalization to an integer polynomial has value at $\theta$ at least
$1/K_V(a,a)$ in absolute value.  If $V\subseteq\mathbb R[x]_{\leq N}$, the
same lower bound holds with $K_{\leq N}(a,a)$ in place of $K_V(a,a)$.
\end{corollary}
\begin{proof}
A polynomial $p=\sum c_i p_i$ has norm squared $c^{\mathsf T}Mc$ and value
$c^{\mathsf T}u$.  Cauchy--Schwarz after multiplication by $M^{1/2}$ shows
that the supremum is $u^{\mathsf T}M^{-1}u$, attained by $c=M^{-1}u$.
Theorem~\ref{p07:zrec:rankbound} proves the remaining statements, and inclusion of
spaces can only increase the supremum.  Positive definiteness of $M$ follows
because a nonzero polynomial cannot vanish on an infinite support.
\end{proof}

Both results are unchanged by a rational invertible change of basis, or by
left and right rational invertible matrix transformations: their effect on
the determinant is only a nonzero rational scalar.  Thus a clever rational
basis does not evade the conclusion for a fixed moment space.

\subsection{The zeta tail attempt: a polynomial lower bound after every normalization}
\label{p07:zrec:zetatails}

For an integer $s\geq2$, put $\theta_s=\zeta(s)$ and
\[
 H_m^{(s)}=\sum_{k=1}^m k^{-s},\qquad
 m_j^{(s)}=\theta_s-H_j^{(s)}.
\]
The integral identity
\begin{equation}
 m_j^{(s)}=\int_0^1 x^j w_s(x)\,dx,
 \quad w_s(x)=\frac{(-\log x)^{s-1}}{\Gamma(s)(1-x)}
\label{p07:zrec:zetamoment}
\end{equation}
follows by expanding $(1-x)^{-1}$ and integrating its nonnegative terms.
The individual integrals are
\[
 \int_0^1x^k(-\log x)^{s-1}dx=\frac{\Gamma(s)}{(k+1)^s}.
\]
The measure is finite and strictly positive on $(0,1)$.

The proposed determinant is
\[
 D_{n,s}(X)=\det\bigl(X-H_{i+j}^{(s)}\bigr)_{0\leq i,j<n}.
\]
At $X=\theta_s$ it is strictly positive.  Its dependence on $X$ is the rank-one
matrix $\boldsymbol1\boldsymbol1^{\mathsf T}$, corresponding exactly to
polynomial evaluation at $a=1$.  It therefore has degree one for every $n$.
To quantify the resulting obstruction we calculate the endpoint kernel.

\begin{lemma}[Endpoint kernel bound]
\label{p07:zrec:endpoint}
For fixed integer $s\geq2$, there are constants $0<c_s<C_s<\infty$ such that
for every $N\geq0$,
\[
 c_s(N+1)^{2s-2}\leq K_{\leq N}^{w_s}(1,1)
                      \leq C_s(N+1)^{2s-2}.
\]
\end{lemma}
\begin{proof}
Write $b=s-2\geq0$.  The inequality $-\log x\geq1-x$ gives
\[
 w_s(x)\geq \Gamma(s)^{-1}(1-x)^b.
\]
Also $w_s(x)\leq C x^{-1/2}(1-x)^b$ for a finite constant $C$: the ratio is
\[
 \frac{x^{1/2}}{\Gamma(s)}
       \left(\frac{-\log x}{1-x}\right)^{s-1},
\]
which extends continuously to $[0,1]$, with values zero at zero and
$1/\Gamma(s)$ at one.  It suffices to compute the kernel for the two Jacobi
weights $x^a(1-x)^b$, $a=0,-1/2$.

Here is the calculation, including the normalization.  For the parameters
$a=0,-1/2$ and $b=s-2\geq0$ let
\[
 J_k(x)=\frac{(-1)^k}{k!\,x^a(1-x)^b}
          \frac{d^k}{dx^k}\{x^{k+a}(1-x)^{k+b}\}.
\]
Leibniz's rule shows that this is a polynomial.  Integration by parts $k$
times proves its orthogonality to degrees below $k$, since the boundary
terms vanish.  The same calculation with $J_k$ itself, followed by the beta
integral, gives
\[
 J_k(1)=\frac{\Gamma(k+b+1)}{k!\Gamma(b+1)},\qquad
 h_k=\int_0^1J_k^2x^a(1-x)^b dx
 =\frac{\Gamma(k+a+1)\Gamma(k+b+1)}
 {(2k+a+b+1)k!\Gamma(k+a+b+1)}.
\]
For the norm formula one can use the leading coefficient
$\Gamma(2k+a+b+1)/(k!\Gamma(k+a+b+1))$ and integrate
$x^{k+a}(1-x)^{k+b}$.  These normalizations also agree with the standard
Jacobi table \cite{p07:zrec:DLMF}.
Thus
\[
 K_{\leq N}^{a,b}(1,1)=\sum_{k=0}^N\frac{J_k(1)^2}{h_k},
 \qquad
 \frac{J_k(1)^2}{h_k}\sim\frac{2k^{2b+1}}{\Gamma(b+1)^2}.
\]
The last relation follows directly from Stirling's formula; summation gives
$K_{\leq N}^{a,b}(1,1)\asymp(N+1)^{2b+2}$.
Finally, increasing a measure decreases its evaluation kernel.  Applying this
observation to the two weight comparisons proves the assertion.
\end{proof}

\begin{theorem}[No exponential arithmetic gain for the entire polynomial tail family]
\label{p07:zrec:zetanogo}
Let $s\geq2$ be a fixed integer and define on $\mathbb Q[x]$ the rational
polynomial-valued functional
\[
        \Lambda_X(x^j)=X-H_j^{(s)}.
\]
Let $V_N\subseteq\mathbb Q[x]_{\leq N}$ be any nonzero finite-dimensional
space with any rational basis $p_i$.  Put
$D_{V_N}(X)=\det(\Lambda_X(p_i p_j))$.
For every nonzero rational $c$ with $P=cD_{V_N}\in\mathbb Z[X]$,
\begin{equation}
       0<|P(\zeta(s))|\quad\hbox{and}\quad
       |P(\zeta(s))|\geq c'_s(N+1)^{-(2s-2)},
\label{p07:zrec:zetalower}
\end{equation}
where $c'_s>0$ depends only on $s$.  In particular, if $N$ is bounded by a
fixed power of a construction index $n$, no such normalized determinant can
be $e^{-\delta n}$, or $e^{-\delta n^2}$, for all sufficiently large $n$.
For the full space $V=\mathbb Q[x]_{<n}$, writing
$D_{n,s}(X)=A_{n,s}X+B_{n,s}$ gives
\begin{equation}
 A_{n,s}>0,\qquad
 \zeta(s)+\frac{B_{n,s}}{A_{n,s}}
      =\frac1{K_{<n}^{w_s}(1,1)}\asymp n^{-(2s-2)}.
\label{p07:zrec:zetaapprox}
\end{equation}
\end{theorem}
\begin{proof}
For a polynomial $p$, the coefficient of $X$ in $\Lambda_X(p)$ is $p(1)$.
At $X=\zeta(s)$ the functional is integration against $w_s$.
Consequently its Gram matrix is a positive definite matrix plus
$(X-\zeta(s))uu^{\mathsf T}$, where $u_i=p_i(1)$.
If $u=0$, the determinant is a nonzero rational constant and any nonzero
integral normalization has absolute value at least one.
Otherwise Corollary~\ref{p07:zrec:christoffel} and
Lemma~\ref{p07:zrec:endpoint} give \eqref{p07:zrec:zetalower}.
The exact ratio and the two-sided kernel bound prove \eqref{p07:zrec:zetaapprox}.
A fixed power of $N$ cannot decay exponentially in $n$ when $N$ is polynomial
in $n$, proving the last assertion.
\end{proof}

For $s=3$ the exponent is $4$; for $s=5$ it is $8$.
These are statements about this precise family, not about all rational
approximations to the constants.  In particular, they do not contradict
Ap\'ery's or Fauzan's constructions.  They show why replacing their arithmetic
input by positivity of the simplest tail matrices cannot work as proposed.

\subsection{Catalan tails: the exact exterior-evaluation exponential rate}
\label{p07:zrec:catalan}

Write
\[
 S_{j-1}=\sum_{k=0}^{j-1}\frac{(-1)^k}{(2k+1)^2},\qquad
 t_j=(-1)^j(G-S_{j-1}),\qquad S_{-1}=0.
\]
A geometric-series integration and $y=x^2$ give
\begin{equation}
 t_j=\int_0^1\frac{x^{2j}(-\log x)}{1+x^2}dx
    =\int_0^1 y^j w_G(y)dy,
 \qquad w_G(y)=\frac{-\log y}{4\sqrt y(1+y)}.
\label{p07:zrec:gmoment}
\end{equation}
The geometric expansion is justified either first on $[0,r]$, $r<1$, and
then by dominated convergence, or by its uniformly bounded partial sums.
The coefficient of $G$ is $(-1)^j$, so the associated polynomial functional
has evaluation at $a=-1$ as its rank-one coefficient.

\begin{lemma}[Exterior kernel]
\label{p07:zrec:exterior}
Let $\rho=3+2\sqrt2$.  For the weight $w_G$ there are constants $c,C>0$ such
that, for $N\geq0$,
\[
          c\rho^{2N}\leq K_{\leq N}^{w_G}(-1,-1)
                 \leq C(N+1)^2\rho^{2N}.
\]
In particular $N^{-1}\log K_{\leq N}^{w_G}(-1,-1)\to2\log\rho$.
\end{lemma}
\begin{proof}
Since $-\log y\geq1-y$ and $4\sqrt y(1+y)\leq8$, we have
$w_G(y)\geq(1-y)/8$.  The shifted Jacobi polynomial for the latter weight is
$J_k(y)=P_k^{(1,0)}(2y-1)$, with squared norm $1/(2k+2)$.
The elementary identity
\[
       P_k^{(1,0)}(t)=\frac{P_k(t)-P_{k+1}(t)}{1-t}
\]
follows, for example, by comparing orthogonality and leading coefficients;
$P_k$ is the Legendre polynomial.  At $t=-3$ its magnitude is at most
$(\rho^k+\rho^{k+1})/4$.  Indeed, the elementary Laplace representation
\[
       P_k(3)=\frac1\pi\int_0^\pi(3+\sqrt8\cos u)^kdu
\]
gives $|P_k(-3)|=P_k(3)\leq\rho^k$.
Both polynomial identities can also be checked by their three-term
recurrences and initial values.  Summing $J_k(-1)^2/h_k$ and using the weight
comparison proves the upper bound.

For the lower bound test the variational kernel with
$p(y)=T_N(2y-1)$, where $T_N$ is a Chebyshev polynomial.
On $[0,1]$ its absolute value is at most one, so
$\int p^2w_G\leq G$.  At $y=-1$,
$|T_N(-3)|=(\rho^N+\rho^{-N})/2\geq\rho^N/2$.
Hence the ratio in the variational definition is at least
$\rho^{2N}/(4G)$.  The limiting statement follows by taking logarithms.
\end{proof}

\begin{theorem}[Catalan integer-normalization obstruction]
\label{p07:zrec:gnogo}
Define $\Lambda_X^G(y^j)=(-1)^j(X-S_{j-1})$.
For any rational polynomial space $V\subseteq\mathbb Q[y]_{\leq N}$, any
rational basis of $V$, and any nonzero rational normalization
\[
       P(X)=c\det\bigl(\Lambda_X^G(p_i p_j)\bigr)\in\mathbb Z[X],
\]
we have
\[
       0<|P(G)|\quad\hbox{and}\quad
       |P(G)|\geq c_0(N+1)^{-2}\rho^{-2N}
\]
for an absolute $c_0>0$.  Thus bounded linear degree $N=O(n)$ rules out
quadratic exponential decay $e^{-\delta n^2}$ after \emph{any} integer
normalization.  For the full space of degrees below $n$, its affine
polynomial $D_{n,G}(X)=A_{n,G}X+B_{n,G}$ satisfies
\[
       \lim_{n\to\infty}\frac1n
           \log\left(G+\frac{B_{n,G}}{A_{n,G}}\right)
                   =-2\log(3+2\sqrt2).
\]
\end{theorem}
\begin{proof}
The coefficient functional is $p\mapsto p(-1)$, and evaluation at $G$ is the
positive measure \eqref{p07:zrec:gmoment}.  Apply
Corollary~\ref{p07:zrec:christoffel} and Lemma~\ref{p07:zrec:exterior}.  If all
polynomials in $V$ vanish at $-1$, the normalized determinant is a nonzero
integer constant and the conclusion still holds after decreasing $c_0$.
For the full polynomial space use the exact ratio
\eqref{p07:zrec:exactratio} and the kernel limit.
\end{proof}

This argument does not rule out an ordinary exponentially small Catalan
linear form with satisfactory arithmetic.  It rules out the tempting
inference that increasing the size of these raw tail matrices automatically
produces a quadratic-exponential irrationality criterion.  The needed
arithmetic saving must come from a different functional or different family,
not a rescaling of these same determinants.

\subsection{The exact geometric decay and a complete denominator ledger}
\label{p07:zrec:ledgersec}

The obstruction is not caused by an incorrectly estimated analytic decay.
The unnormalized determinants really do decay on a quadratic exponential
scale.  We give that scale and then account for the denominators explicitly.

\begin{proposition}[Sharp geometric scale]
\label{p07:zrec:geometry}
Let $\mu$ be a finite positive measure on $[a,b]$, $a<b$, having a density
that is continuous and strictly positive on every compact subinterval of
$(a,b)$.  Set
\[
 \Delta_n(\mu)=\det\left(\int x^{i+j}\,d\mu(x)\right)_{0\leq i,j<n}.
\]
Then
\[
 \lim_{n\to\infty}\frac{\log\Delta_n(\mu)}{n^2}
                      =\log\frac{b-a}{4}.
\]
In particular, both $D_{n,s}(\zeta(s))$ and $D_{n,G}(G)$ have limit
$-\log4$ after taking their logarithms and dividing by $n^2$.
\end{proposition}
\begin{proof}
An affine change of variable multiplies the determinant by $(b-a)^{n(n-1)}$,
apart from a total-mass normalization contributing only $O(n)$ to its
logarithm.  It suffices to prove the result on $[0,1]$.
Gram--Schmidt expresses $\Delta_n$ as the product of the squared norms of
the monic orthogonal polynomials of degrees $0,\ldots,n-1$.
The monic shifted Chebyshev polynomial of degree $k\geq1$ has supremum
$2^{1-2k}$ on $[0,1]$.  Since the orthogonal polynomial minimizes the norm
among monic polynomials of its degree, if $M=\mu([0,1])$,
\[
       \Delta_n\leq M^n 2^{-2(n-1)^2}.
\]
This gives the upper limit $-\log4$.

For $0<\varepsilon<1/2$, let $m_\varepsilon>0$ be a lower bound for the
density on $[\varepsilon,1-\varepsilon]$.  The Gram matrix dominates
$m_\varepsilon$ times the Lebesgue Gram matrix on this interval.  Positivity
and conjugation by the inverse square root of the latter matrix give the
same inequality for determinants, with factor $m_\varepsilon^n$.
The monic shifted Legendre norms give the exact identity
\[
 \Delta_n(dx|_{[u,v]})=(v-u)^{n^2}
       \prod_{k=0}^{n-1}\frac{1}{(2k+1)\binom{2k}{k}^{2}}.
\]
Indeed the degree-$k$ Legendre polynomial on $[0,1]$ has leading coefficient
$\binom{2k}{k}$ and squared norm $1/(2k+1)$; rescaling gives the displayed
formula.  Stirling's formula now gives a lower limit
$\log((1-2\varepsilon)/4)$.  Letting $\varepsilon\downarrow0$ proves the
claim.  The two weights used above satisfy all the hypotheses.
\end{proof}

Put $d_m=\operatorname{lcm}(1,\ldots,m)$ and
$o_m=\operatorname{lcm}(1,3,\ldots,2m-1)$, with $d_0=o_0=1$.
The ordinary prime number theorem implies
\[
        \log d_m=m+o(m),\qquad \log o_m=2m+o(m).
\]
For the second identity, deleting the powers of $2$ from the least common
multiple through $2m-1$ changes its logarithm by only $O(\log m)$.

\begin{proposition}[Certified denominator cost]
\label{p07:zrec:denominators}
The explicit integers
\[
 C_{n,s}=\prod_{i=0}^{n-1}d_{n-1+i}^{s},\qquad
 C_{n,G}=\prod_{i=0}^{n-1}o_{n-1+i}^{2}
\]
satisfy $C_{n,s}D_{n,s}(X),C_{n,G}D_{n,G}(X)\in\mathbb Z[X]$ and
\[
 \lim_{n\to\infty}\frac{\log(C_{n,s}D_{n,s}(\zeta(s)))}{n^2}
       =\frac{3s}{2}-\log4>0,
 \qquad
 \lim_{n\to\infty}\frac{\log(C_{n,G}D_{n,G}(G))}{n^2}
       =6-\log4>0.
\]
Thus these particular, completely justified integer normalizations grow on
the quadratic exponential scale.
\end{proposition}
\begin{proof}
The denominator of $H_j^{(s)}$ divides $d_j^s$.  Multiplication of row $i$
by $d_{n-1+i}^s$ therefore makes all its polynomial entries integral.
The alternating partial sum $S_{j-1}$ has denominator dividing $o_j^2$,
which proves the Catalan assertion by the same row operation.
Summing $\log d_m=m+o(m)$ over $m=n-1,\ldots,2n-2$ gives
$\log C_{n,s}=(3s/2)n^2+o(n^2)$; the Catalan sum gives
$\log C_{n,G}=6n^2+o(n^2)$.  Proposition~\ref{p07:zrec:geometry} completes the
calculation.  The little-oh errors can be summed uniformly because their
arguments lie between $n-1$ and $2n-2$.
\end{proof}

This also identifies precisely what a hypothetical content improvement would
have to do.  Let $g_n$ be the greatest common divisor of the coefficients of
one of the integer polynomials $C_nD_n$.  If
$\log g_n/n^2\to\gamma$ and $\log C_n/n^2\to\alpha$, then its primitive
normalization has logarithmic value
\[
 \lim_{n\to\infty}n^{-2}\log|(C_n/g_n)D_n(\theta)|
                     =\alpha-\gamma-\log4.
\]
A negative limit would require $\gamma>\alpha-\log4$.
Theorems~\ref{p07:zrec:zetanogo} and \ref{p07:zrec:gnogo} prove that this strict
inequality is impossible for the present raw-tail families.  The argument
therefore rules out even a presently undiscovered content saving of a fixed
positive quadratic margin.  It does not merely exhibit one inefficient
denominator estimate.

\begin{proposition}[Rational dilation does not improve the quadratic balance]
\label{p07:zrec:dilation}
Let $s\geq2$ and $z=a/b\in(0,1]$, with $a,b$ coprime positive integers.
Write
\[
 L=\operatorname{Li}_s(z),\qquad
 D_{n,s,z}(X)=\det\left(X-\sum_{k=1}^{i+j}\frac{z^k}{k^s}\right)_{0\leq i,j<n}.
\]
Then $D_{n,s,z}$ is affine and positive at $L$, and
\[
 b^{n(n-1)}C_{n,s}D_{n,s,z}(X)\in\mathbb Z[X],\qquad
 \lim_{n\to\infty}\frac1{n^2}
  \log\left(b^{n(n-1)}C_{n,s}D_{n,s,z}(L)\right)
                  =\log a+\frac{3s}{2}-\log4>0.
\]
In particular, shrinking a rational argument with numerator one produces
no improvement at all in this certified quadratic-scale balance.
\end{proposition}
\begin{proof}
Expansion of $(1-x)^{-1}$ and the substitution $x=zu$ give
\[
 L-\sum_{k=1}^j\frac{z^k}{k^s}
       =\int_0^z x^j\frac{\log(z/x)^{s-1}}{\Gamma(s)(1-x)}\,dx.
\]
The measure is positive and finite; its support is $[0,z]$.  The perturbation
matrix is again rank one.  Proposition~\ref{p07:zrec:geometry} gives geometric
limit $\log(z/4)$.  Multiply row $i$ by $b^i d_{n-1+i}^s$ and column $j$ by
$b^j$.  Each rational term becomes
$a^k b^{i+j-k}d_{n-1+i}^s/k^s$, an integer because $k\leq i+j\leq n-1+i$.
The coefficient of $X$ is also integral.  Taking logarithms of the
determinant multiplier gives $\log b+3s/2$ on the quadratic scale.
Adding $\log(z/4)$ proves the asserted cancellation of $b$.
\end{proof}

An exact rational calculation, included with this section's supporting code,
checks the affine identity and primitive integer coefficients for
$1\leq n\leq6$.  Some small cases are
\[
\begin{array}{c|c|r|r}
\text{target}&n&[X]P_n&[1]P_n\\\hline
\zeta(3)&2&7&-8\\
\zeta(3)&3&4105&-4887\\
\zeta(5)&2&31&-32\\
\zeta(5)&3&4458601&-4622103\\
G&2&10&-9\\
G&3&61794&-56575
\end{array}
\]
These computations illustrate the exact construction; the theorems above
hold for every size and do not rely on these finite checks.

\subsection{What would actually be needed for transcendence}
\label{p07:zrec:heightsec}

The difference between irrationality and transcendence is not repaired by
producing a small integer polynomial alone.  The height and degree must be
tracked together with its value and its nonvanishing.

\begin{proposition}[A height-aware quadratic-scale criterion]
\label{p07:zrec:height}
Let $\theta$ be algebraic of degree $d$ with primitive minimal polynomial
having leading coefficient $a\ne0$ and conjugates
$\theta=\theta_1,\ldots,\theta_d$.  For
$P\in\mathbb Z[X]$ of degree $m$ and height $H\geq1$, if $P(\theta)\ne0$,
then
\[
 |P(\theta)|\geq |a|^{-m}\bigl((m+1)H\bigr)^{-(d-1)}
                    \prod_{j=2}^d\max(1,|\theta_j|)^{-m}.
\]
Consequently, suppose a sequence satisfies
\[
 m_n=o(n^2),\quad \log H(P_n)\leq\beta n^2+o(n^2),\quad
 0<|P_n(\theta)|\leq\exp(-\alpha n^2+o(n^2)),
\]
where $\alpha>0$ and $\beta\geq0$.  Then $\theta$ cannot have algebraic
degree $d$ whenever $\alpha>(d-1)\beta$.
\end{proposition}
\begin{proof}
The resultant of the minimal polynomial and $P$ is the nonzero integer
$a^m\prod_{j=1}^dP(\theta_j)$, so its absolute value is at least one.
For each other conjugate,
$|P(\theta_j)|\leq(m+1)H\max(1,|\theta_j|)^m$.
Division proves the inequality.  Taking logarithms gives
\[
 \log|P(\theta)|\geq-(d-1)\log H-O_\theta(m+\log(m+1)),
\]
which contradicts the specified upper bound in the indicated range.
\end{proof}

To exclude all algebraic degrees with positive quadratic height cost, the
ratio $\alpha/\beta$ would need to be unbounded across a family of
constructions.  Alternatively, $\beta=0$ and $\alpha>0$ suffice, subject to
the stated degree and nonvanishing hypotheses.  None of our tail determinants
provides either alternative.  At the two zeta targets they do not even
provide exponentially small normalized values, by
Theorem~\ref{p07:zrec:zetanogo}.  At $G$, Theorem~\ref{p07:zrec:gnogo} already rules
out a positive quadratic exponent with degree space $O(n)$.

For an explicit check on this logical distinction, write
$(1+\sqrt2)^k=p_k+q_k\sqrt2$, where $p_k,q_k$ are positive integers.  Its
conjugate is $(1-\sqrt2)^k=p_k-q_k\sqrt2$.  With $k=n^2$ the degree-one
integer polynomials $Q_n(X)=q_{n^2}X-p_{n^2}$ satisfy
\[
 0<|Q_n(\sqrt2)|=(\sqrt2-1)^{n^2},\qquad
 \log H(Q_n)=n^2\log(1+\sqrt2)+O(1).
\]
Thus even quadratic exponential smallness of nonzero integer linear forms
is compatible with algebraic irrationality.  Here $\alpha=\beta$, exactly
at the degree-two threshold of Proposition~\ref{p07:zrec:height}.

Nor does a linear-independence theorem for additional constants settle an
individual transcendence question.  Lai proves that
$1,\zeta(3),\zeta(4)$ are linearly independent over $\mathbb Q$
\cite{p07:zrec:Lai}.  For comparison, $1,\alpha,\pi^4$ are linearly independent
over $\mathbb Q$ for \emph{every} algebraic irrational $\alpha$: a nonzero
coefficient of $\pi^4$ would make $\pi^4$ algebraic, whereas a zero coefficient
reduces to irrationality of $\alpha$.  Since $\zeta(4)=\pi^4/90$, the form
of this independence result is compatible with algebraic irrationality of
the middle entry.  No individual transcendence conclusion follows from it.

\subsection{The actual transfer boundary and the remaining tasks}
\label{p07:zrec:conclusions}

The earlier positivity and finite-witness methods transfer successfully to
strict nonvanishing of our determinants.  They do not transfer the required
denominator savings.  The obstruction here identifies the exact failure:
all polynomial row changes leave the target dependence as one evaluation
functional, and an integer leading coefficient forces the value to exceed
the reciprocal Christoffel kernel.  For zeta tails this reciprocal is only
a power of the polynomial degree.  For Catalan tails it is exponential in
that degree, although the raw determinant is quadratic exponential.

The recent OpenAI Catalan construction illustrates a genuine escape from
these hypotheses, rather than a counterexample.  Its paper explicitly uses
mixed signed moments, cancellation of an auxiliary $\zeta(2)$ contribution,
and a separate arithmetic nonvanishing certificate at prime scales
\cite[Sections 1.2 and 4]{p07:zrec:OpenAI}.  Those are not rational basis changes
of the single positive tail functional studied here.  In particular, our
positive rank-one theorem cannot be applied to that mixed signed matrix.
Fauzan's polynomial determinants likewise must not be replaced by these
affine raw-tail determinants without proving a new identity.

For individual transcendence of $G$, $\zeta(5)$ and $\zeta(3)$, the unresolved
task in this attempted approach is therefore specific: construct a family
outside the rank-constrained tail class, certify that its integer-polynomial
value is nonzero, and simultaneously obtain height, degree and decay bounds
meeting Proposition~\ref{p07:zrec:height} for every algebraic degree.  Neither
functional independence of a generating function nor finite numerical
checks provide those three estimates.  This section supplies a proved
method obstruction and an exact arithmetic ledger; it records no new
resolution of these individual transcendence questions.


\section{Arithmetic information missing from ergodic and analytic recovery}
\label{p07:ric:section}

This section tests two named targets: the Khinchin constant $K$ and the
period-doubling Feigenbaum constant $\delta$.  Neither constant is proved
irrational or transcendental here.  Instead, two precise transfers identify
information that recovery arguments do not supply.  Full finite-cylinder
statistics need not control the unbounded logarithmic observable defining
$K$.  Even a uniformly effective, normalized analytic contraction, with a
uniform spectral gap and a unique solution, need not permit an algorithmic
decision of the arithmetic type of that solution.  The latter assertion is
proved for an explicit family with rational Taylor coefficients; it is not
an undecidability assertion about either named constant.

\subsection{Khinchin's constant: reconstruction of the law versus its logarithmic moment}
\label{p07:ric:khinchin}

Let $\mu$ be the Gauss probability measure on $(0,1)$,
\[
 d\mu(x)=\frac{dx}{(1+x)\log 2},\qquad
 p_k=\mu\left(\frac1{k+1}<x<\frac1k\right)
     =\log_2\frac{(k+1)^2}{k(k+2)}\quad(k\ge1).
\]
Writing the continued fraction of an irrational $x\in(0,1)$ as
$[0;a_1,a_2,\ldots]$, the classical almost-everywhere geometric-mean theorem
identifies
\begin{equation}\label{p07:ric:K-definition}
 L:=\log K=\sum_{k\ge1}p_k\log k,
 \qquad
 \lim_{n\to\infty}\frac1n\sum_{j=1}^n\log a_j=L
 \quad\text{for $\mu$-almost every $x$}.
\end{equation}
Bailey, Borwein and Crandall give the product representation, accelerated
evaluation methods, and a discussion of the unresolved arithmetic question
\cite{p07:ric:BBC}.  Their finite searches for algebraic relations do not
constitute a proof of irrationality.  No such proof is supplied by the
sources verified in this investigation.

For completeness, even the exact one-digit probabilities exhibit a warning
against an arithmetic inference from the recovered probability law.

\begin{proposition}\label{p07:ric:probabilities}
Every $p_k$ is transcendental, whereas $\sum_{k\ge1}p_k=1$.  Moreover, for
every integer $B\ge2$,
\begin{equation}\label{p07:ric:K-tail}
 0\le L-\sum_{k\le B}p_k\log k
 \le\frac{\log B+1}{B\log2}.
\end{equation}
\end{proposition}

\begin{proof}
Put $r_k=(k+1)^2/[k(k+2)]$.  This is rational and lies strictly between
$1$ and $2$.  If $p_k=a/b$ were rational in lowest terms, the equality
$2^a=r_k^b$ and the $2$-adic valuation would imply $b\mid a$, hence
$b=1$, contradicting $0<p_k<1$.  If $p_k$ were algebraic irrational,
the Gelfond--Schneider theorem would make $2^{p_k}=r_k$ transcendental,
also a contradiction.  This application of the classical theorem has no
novelty claim.  Telescoping gives
\[
 \prod_{k=1}^B r_k=\frac{2(B+1)}{B+2}\longrightarrow2,
\]
and therefore the probabilities sum to one.  Finally,
\[
 p_k=\frac{\log(1+1/[k(k+2)])}{\log2}
 \le\frac1{k^2\log2}.
\]
The function $(\log t)/t^2$ decreases for $t\ge2$, so its integral on
$[B,\infty)$ bounds the tail and gives \eqref{p07:ric:K-tail}.
\end{proof}

The missing observation condition can be characterized exactly.  The next
statement is the elementary uniform-integrability principle specialized to
the continued-fraction observable; its role is to expose the required
additional hypothesis, rather than to claim a new general convergence
theorem.

\begin{proposition}[Exact tail requirement]\label{p07:ric:UI}
Let $(a_j)$ be any sequence of positive integers such that
\[
 \frac1n\#\{j\le n:a_j=k\}\longrightarrow p_k
 \quad\text{for every }k\ge1.
\]
Then $n^{-1}\sum_{j\le n}\log a_j\to L$ if and only if
\begin{equation}\label{p07:ric:UI-condition}
 \lim_{B\to\infty}\limsup_{n\to\infty}
 \frac1n\sum_{\substack{j\le n\\a_j>B}}\log a_j=0.
\end{equation}
\end{proposition}

\begin{proof}
For fixed $B$ the truncated average converges to
$L_B=\sum_{k\le B}p_k\log k$.  If the full averages converge to $L$,
their difference from the truncated averages converges to $L-L_B$;
\eqref{p07:ric:K-tail} then proves \eqref{p07:ric:UI-condition}.  Conversely,
nonnegativity gives $\liminf_n n^{-1}\sum_{j\le n}\log a_j\ge L_B$.
The limsup is at most $L_B$ plus the limsup of the tail in
\eqref{p07:ric:UI-condition}.  Let $B$ tend to infinity in these two bounds.
\end{proof}

Observing every finite word, rather than only individual digits, still does
not remove this obstruction.  A Gauss-normal continued fraction means here
that each finite word $w$ of positive integers has limiting occurrence
frequency $\mu(C_w)$, where $C_w$ is its continued-fraction cylinder.

\begin{theorem}[A computable exceptional orbit with the full cylinder law]
\label{p07:ric:normal-Liouville}
There is a computable Liouville number $x\in(0,1)$ whose continued fraction
is Gauss-normal and satisfies
\[
 \limsup_{n\to\infty}\frac1n\sum_{j=1}^n\log a_j=+\infty.
\]
Thus the entire collection of limiting finite-cylinder frequencies does not
determine the logarithmic orbit average in \eqref{p07:ric:K-definition}.
\end{theorem}

\begin{proof}
Start with a computable Gauss-normal digit sequence $(b_i)$; the explicit
construction of Adler, Keane and Smorodinsky supplies one
\cite{p07:ric:AKS}.  Retain all these digits in their order.  Immediately after
the original digit $b_{2^j}$, insert one digit $A_j$.  If the already
constructed prefix before this insertion has length $m_j$ and convergent
denominator $q_{m_j}$, choose an integer
\[
 A_j\ge \max\{2,q_{m_j}^{j+2}\}.
\]
This recursively defines a computable sequence of positive integers.
Its infinite continued fraction is computable: its rational convergents
are effective and their denominators have a uniform Fibonacci lower bound.

There are $O(\log n)$ insertions among the first $n$ digits.  For a fixed
word of length $r$, deleting one inserted digit can change at most $2r$
candidate occurrences near that digit.  After deleting all insertions,
the prefix becomes a prefix of $(b_i)$ of length $n-O(\log n)$.
The discrepancy of its word count from the original prefix count is
$O(r\log n)$, including the terminal boundary.  Dividing by $n$ proves
that every finite word has the same limiting frequency as in $(b_i)$.

At the convergent immediately before $A_j$, the standard
continued-fraction estimate gives
\[
 0<\left|x-\frac{p_{m_j}}{q_{m_j}}\right|
 <\frac1{A_jq_{m_j}^2}
 \le q_{m_j}^{-j-4}.
\]
The denominators tend to infinity.  For each positive integer exponent
these inequalities give infinitely many sufficiently accurate rational
approximations, so $x$ is Liouville and hence transcendental.
Finally the recurrence for convergent denominators gives
$q_m\ge2^{(m-1)/2}$: indeed $q_{m+2}\ge2q_m$, with the usual
initial values.  At the prefix ending in $A_j$ the logarithmic average is
at least
\[
 \frac{\log A_j}{m_j+1}
 \ge \frac{(j+2)(m_j-1)}{2(m_j+1)}\log2
 \longrightarrow+\infty.
\]
This proves the assertion.
\end{proof}

The construction is an explicit application of sparse insertion and
classical continued-fraction estimates; no priority claim for normal
Liouville constructions is intended.  It also does not contradict the
almost-everywhere theorem: the orbit constructed is exceptional.  For the
named constant $K$, the precise failed transfer is now visible.  Identifying
the invariant distribution and proving an effective real tail estimate
does not yield algebraic-height-controlled rational approximations to its
exponentiated logarithmic moment.  In particular, the weights $p_k$ are
themselves transcendental, so Baker's theorem for finite linear forms with
algebraic coefficients cannot simply be applied to
$\sum p_k\log k$.  Neither the transcendence of the individual weights nor
their exact identification proves an arithmetic classification of $K$.

\subsection{Feigenbaum's constant: uniqueness and a spectral gap are insufficient}
\label{p07:ric:Feigenbaum}

The target here is the period-doubling parameter-scaling constant
$\delta=4.669201609\ldots$.  For the principal period-doubling cascade of
the logistic family, let $r_n$ be its superstable parameter of period
$2^n$.  With this convention,
\[
 \delta=\lim_{n\to\infty}
       \frac{r_n-r_{n-1}}{r_{n+1}-r_n}.
\]
The equivalent renormalization description identifies the expanding
scaling eigenvalue at the universal fixed point.  Rigorous fixed-point
methods originate in Lanford's work \cite{p07:ric:Lanford}; high-precision
calculations and digit statistics are discussed in
\cite{p07:ric:Briggs,p07:ric:KaramanosKotsireas,p07:ric:Molteni}.
The verified sources do not establish its irrationality or transcendence.
One must distinguish $\delta$ from the spatial-scaling constant
$\alpha=-2.5029\ldots$: the polynomial-time theorem of Hertling and Spandl
computes the fixed function and $\alpha=1/g(1)$, and is not a stated
polynomial-time theorem for $\delta$ \cite[Theorem~2]{p07:ric:HertlingSpandl}.

Could an arithmetic conclusion follow from effective analytic recovery,
uniqueness and hyperbolicity alone?  The following negative transfer is
considerably more restrictive than encoding an arbitrary computable real
as the constant term of a function.  Both the constant and linear Taylor
coefficients are fixed; all remaining coefficients are nonnegative
rationals with a uniform algorithm and uniform analytic bounds.

\begin{theorem}[Normalized analytic recovery does not decide arithmetic type]
\label{p07:ric:contraction}
There is a uniformly specified family of entire functions
$F_e(z)=\sum_{n\ge0}c_{e,n}z^n$, indexed by finite machine descriptions,
with the following properties.
\begin{enumerate}
\item Every coefficient is a nonnegative rational, and
$c_{e,0}=1/2$, $c_{e,1}=1/8$.  Given $e$ and $n$ in unary, $c_{e,n}$ is
computable exactly in time polynomial in the description length of $e$
and in $n$.
\item $\sup_{|z|\le2}|F_e(z)|<1$,
$F_e([0,1])\subset[1/2,3/4]$, and
$1/8\le F'_e(x)<1/4$ on $[0,1]$.
\item The unique fixed point $x_e\in[0,1]$ is uniformly computable to
error $2^{-p}$ in time polynomial in $p$ and the description length of $e$.
\item Each $x_e$ is either $4/7$ or transcendental.  There is no algorithm
that, given these coefficient-program descriptions with all the preceding
uniform guarantees, always decides which alternative holds.
\end{enumerate}
\end{theorem}

\begin{proof}
Fix a universal machine model and pad the computation so that any first
halting step $N$ satisfies $N\ge2$.  Put
\[
 F_0(z)=\frac12+\frac z8,\qquad \epsilon_N=2^{-2N-6}.
\]
If machine $e$ never halts, define $F_e=F_0$.  If it first halts at step
$N$, define
\begin{equation}\label{p07:ric:F-halt}
 F_e(z)=F_0(z)+\epsilon_Nz^Ne^z.
\end{equation}
These are mathematical definitions, not an algorithm required to decide
whether halting occurs.  To compute the coefficient of degree $n\ge2$,
simulate the machine for just $n$ steps.  If it has not halted, output
zero.  If its first halting step is $N\le n$, output
\[
 c_{e,n}=\frac{2^{-2N-6}}{(n-N)!}.
\]
This is correct even if the machine halts later than $n$.  Its rational
output has $O(n\log(n+2))$ bits.  A universal simulation for $n$ steps,
factorial calculation and exact integer arithmetic all have polynomial
cost in the stated input lengths.  Thus the family has uniform coefficient
algorithms, without using a halting oracle.

For $|z|\le2$, in the halting case,
\[
 |F_e(z)|\le\frac34+\epsilon_N2^Ne^2
 =\frac34+2^{-N-6}e^2<1.
\]
On $[0,1]$, its values lie between $1/2$ and
$5/8+\epsilon_N e<3/4$.  Also
\[
 F'_e(x)=\frac18+\epsilon_Ne^x(Nx^{N-1}+x^N)
 \le\frac18+e(N+1)2^{-2N-6}<\frac14.
\]
For the final bound use $e<3$ and
$(N+1)4^{-N}\le3/16$ for $N\ge2$.  The nonhalting case satisfies the
same bounds directly.  The contraction theorem now gives a unique
$x_e\in[1/2,3/4]$.

We verify the claimed complexity rather than infer it from existence.
Nonnegativity of the coefficients and $F_e(2)<1$ give the uniform
truncation estimate
\[
 \sup_{|x|\le1}\left|F_e(x)-\sum_{n=0}^m c_{e,n}x^n\right|
 \le 2^{-m-1}\sum_{n>m}2^nc_{e,n}<2^{-m-1}.
\]
At a dyadic argument with $O(p)$ bits, the polynomial of degree $p+5$
can be evaluated by rational arithmetic with polynomial bit cost and then
rounded to error at most $2^{-p-4}$.  Including the truncation error, the
total evaluation error is less than $\eta=2^{-p-3}$.  Start at $1/2$ and perform
$p+2$ such approximate iterations, retaining $O(p)$ bits at each step.
The computed iterates remain in $[0,1]$ because the exact image is in
$[1/2,3/4]$ and $\eta\le1/16$ for $p\ge1$.
Their distance from $x_e$ is at most
\[
 4^{-p-2}+\frac{\eta}{1-1/4}<2^{-p}.
\]
Only polynomially many rational operations on polynomially many bits are
used; each coefficient is obtained by the preceding bounded simulation.

If $e$ never halts, its fixed point solves $x=1/2+x/8$ and equals $4/7$.
If $e$ halts at $N$ and its fixed point were algebraic, then
\[
 e^{x_e}=\frac{(7/8)x_e-1/2}{\epsilon_Nx_e^N}
\]
would be algebraic.  This contradicts the Hermite--Lindemann theorem,
since $x_e\ge1/2$ is a nonzero algebraic number.  Thus the fixed point is
transcendental in the halting case.  A decision procedure for the two
arithmetic alternatives, on this uniformly generated family of promised
inputs, would decide the halting problem \cite{p07:ric:Turing}.
\end{proof}

The representation in this theorem matters.  The input is an effective
program for rational Taylor coefficients, with fixed analytic estimates;
it is not a supplied elementary expression announcing in advance the
halting time $N$.  Every particular function happens to have one of the
elementary forms above, but there is no uniform way to select the correct
form by deciding whether the machine will halt.  This is a limitation on a
general arithmetic decision procedure, not on approximation: all of the
fixed points remain efficiently computable real numbers.

\begin{corollary}[Finite jets, convex minimizers and expanding eigenvalues]
\label{p07:ric:three-transfers}
The family of Theorem~\ref{p07:ric:contraction} has the following additional
properties.
\begin{enumerate}
\item For every finite Taylor order and every positive tolerance there
are two functions satisfying its uniform analytic bounds, agreeing to
that order and within the tolerance in the uniform norm on $|z|\le2$,
whose fixed points are respectively rational and transcendental.
\item The entire functions
\[
 V_e(x)=\frac{x^2}{2}-\int_0^xF_e(t)\,dt
\]
have rational Taylor coefficients, $V_e(0)=0$, $V'_e(0)=-1/2$, and
$3/4<V''_e\le7/8$ on $[0,1]$.  Their unique minimizers on this interval
have the same undecidable rational/transcendental dichotomy.
\item The entire maps
\[
 \mathcal R_e(u,v)=\bigl(F_e(u),(4+u)v\bigr)
\]
have a unique fixed point in $[0,1]\times\mathbb R$.  It is hyperbolic,
with one eigenvalue in $[1/8,1/4)$ and a simple expanding eigenvalue
$\lambda_e\in[9/2,19/4]$.  These expanding eigenvalues are uniformly
polynomial-time computable, but their rational/transcendental dichotomy
is undecidable on the same input representation.
\end{enumerate}
\end{corollary}

\begin{proof}
The perturbation $\epsilon_Nz^Ne^z$ has vanishing coefficients below
degree $N$ and satisfies
\[
 \sup_{|z|\le2}|\epsilon_Nz^Ne^z|
 \le e^2 2^{-N-6}\longrightarrow0.
\]
Compare it with $F_0$ and choose $N$ larger than the prescribed order.
Both functions can be realized by machines in the construction.
For the second assertion, $V'_e=x-F_e(x)$ and $V''_e=1-F'_e$.
The derivative changes sign across the unique zero $x_e$, which lies in
the interior; the asserted uniform convexity proves that it is the unique
minimizer.  Integration preserves rational Taylor coefficients.
For the third assertion a fixed point must have $u=x_e$ and
$v=(4+x_e)v$, hence $v=0$.  At that point the derivative matrix is
\[
 D\mathcal R_e(x_e,0)=
 \begin{pmatrix}F'_e(x_e)&0\\0&4+x_e\end{pmatrix}.
\]
Its expanding eigenvalue is $\lambda_e=4+x_e$, so rational translation
preserves the dichotomy and the computational bounds.
\end{proof}

The corollary gives a precise obstruction to the attempted Feigenbaum
transfer: rational coefficient access, analytic control, a unique recovered
object and a simple expanding eigenvalue, even with uniform quantitative
gaps, cannot by themselves provide an arithmetic classification algorithm.
The maps $\mathcal R_e$ are model maps, not Feigenbaum renormalization
operators.  The theorem therefore neither resolves $\delta$ nor rules out
a proof using its specific renormalization identities.  Such a proof would
need additional arithmetic structure---for example an identity governed
by a transcendence theorem or approximants with controlled algebraic
heights---which neither a spectral enclosure nor an arbitrarily long
digit calculation supplies.  These model consequences form one
obstruction result; they are not counted as additional named constants.


\section{What has been recovered, and what would constitute a further advance}
\label{p07:ra:conclusion}

The completed classifications have a common shape: analytic recovery is
supplemented by an arithmetic invariant that is computable or rigid.
In the quantile theorem it is a finite cyclotomic degree bound; in the
logarithmic-root theorem it is the algebraic zero certificate; in the
Mahler families it is a finite rational functional equation. Removing
that invariant destroys the inference to arithmetic type.

\begin{center}\small\setlength{\tabcolsep}{3pt}
\begin{tabular}{L{43mm}L{51mm}L{54mm}}
\toprule
Earlier failed transfer & Recovered ingredient & Exact result here\\
\midrule
Logarithmic logic to arbitrary real inversion &
Baker splitting at algebraic inputs plus a rational derivative &
All isolated roots and their arithmetic types are decided in
Theorem~\ref{p07:rlog:classification}.\\[3pt]
Normalization of an inverse problem by a transcendental period &
An algebraic correction and its full Galois orbit &
Finite order bound $N\le8r^2$ in Theorem~\ref{p07:rq:tilt}.\\[3pt]
Finite algebraic manipulation to a connection constant &
Finite monodromy under translation &
The common logarithmic coordinate is eliminated in
Theorem~\ref{p07:rgam:finiteorbit}.\\[3pt]
Finite-dimensional compactness to a computable search &
Factorial coercivity in numerator degree &
Every fixed positive rational Laplace kernel has a finite
integer candidate reduction, Theorem~\ref{p07:rgam:gap}.\\[3pt]
Small positive moment determinants to small integer forms &
Rank and the evaluation kernel &
Every rational integer normalization obeys
Theorem~\ref{p07:zrec:rankbound}; raw zeta tails lose exponential gain.\\[3pt]
Exact limiting observations to the arithmetic of a limit &
Uniform integrability and arithmetic structure &
Theorems~\ref{p07:ric:normal-Liouville} and \ref{p07:ric:contraction}
supply explicit countermodels to the unrestricted transfer.\\
\bottomrule
\end{tabular}
\end{center}

\paragraph{A narrower question genuinely decided.}
For a given rational polynomial arcsine tilt, the question of which
algebraic probability levels have algebraic quantiles has an exact
terminating answer here. The same is true for every positive
half-integral beta law and every integer-degree Student or Fisher
law. These are full classifications of stated natural families.
They are not merely examples of a number that was defined to be
transcendental, and they do not rely on floating-point tests for
algebraicity. They are also not fifteen separate historical conjectures.

\paragraph{The remaining novelty question.}
The underlying transcendence theorems are classical. Dhillon--Saha's
2025 work explicitly studies incomplete beta values by logarithmic
methods in a parameter range containing our half-integral case;
the available public slides confirm that overlap. The normalized
inverse classification, its real-root argument, and the finite
cyclotomic bound need a line-by-line comparison with the full
published article and the wider inverse-special-function literature
before a priority claim is made. The entropy and Mahler statements
likewise include explicitly credited antecedents. A complete proof
and a new historical contribution are different assessments.

\paragraph{A quantitatively meaningful next target.}
The Hankel obstruction gives a concrete design constraint:
a positive perturbation of small rank cannot create the desired
integer-polynomial decay unless its inverse-Gram evaluation
determinant is correspondingly large. A next construction must
escape that condition by a genuinely different family, such as
mixed signed functionals or a growing arithmetic perturbation rank,
and must preserve nonvanishing. Even then, the height criterion
in Proposition~\ref{p07:zrec:height} has to exclude every algebraic
degree before a transcendence conclusion follows. The externally
reported recent irrationality constructions illustrate some escapes;
they do not automatically furnish the missing transcendence bounds.

For $\gamma$ and $\delta$, growing denominator kernels and
provable common-factor savings are outside the fixed-lattice
obstruction, but no such successful saving is established here.
For $B_1$ and $K$, a new theorem controlling the infinite arithmetic
combination is needed, not just faster evaluation. For the gamma
fifths, rank information must be replaced by a statement identifying
an individual value. For $\delta_F$, the special renormalization
identities must contribute arithmetic information beyond
hyperbolicity. These are the explicit unfinished mathematical
steps in the fifteen attempts.

\subsection{Verification record}

The beta classification and polynomial-tilt extension were independently
checked, including the conjugacy argument, the bound $8r^2$, and the
asymmetric beta and Fisher exceptions. The logarithmic zero and binary
entropy proofs were checked separately from their numerical examples.
The Laguerre degree and coefficient bounds, the Mahler pole bounds and
algebraic loci, the Christoffel normalization bounds, and the computable
contraction and normal-Liouville constructions received independent
written-proof reviews. Exact rational checks verify the logarithmic
isolators and the small Hankel coefficient tables; those checks do not
replace the all-parameter proofs.

This document is a single self-contained \LaTeX\ source with an embedded
bibliography. It uses no external figures, source inclusions, or
bibliography database. Its cross-references and citations are checked
during assembly, and the final PDF is compiled and inspected for layout.
None of these checks establishes historical novelty or formal verification
of a theorem by a proof assistant.

% END SOURCE: KEIO_Arithmetic_Beyond_Convergence.tex

\endgroup
\clearpage
\part{Arithmetic Certificates}
\label{merge:part08}
\begingroup

% BEGIN SOURCE: KEIO_Arithmetic_Certificates.tex

\noindent\textbf{Editorial integration note.} 
The five-atom entropy construction remains valid; subsequent parts give a three-atom example for algebraic irrational bit entropy and a four-atom example for rational bit entropy with irrational probabilities. The full-degree mixed Catalan pencil here is not asserted to be positive: later parts construct and distinguish positive alternatives. Exact moment and degree costs are sharpened in Part~\ref{merge:part09}.


\begin{quote}\small
\noindent\textit{Original part framing.}\par
\noindent\textbf{Reading and verification note.}
This is a standalone research manuscript: its proofs do not require
access to the earlier 223-page or 447-page compilations.
The results added here are claims with supplied mathematical proofs,
not formal machine certificates or an assertion of journal-level
acceptance. Exact computations and separate proof reviews test
specific identities and hypotheses. Bibliographic priority is not
established merely by proving a statement.
The foundational inputs are credited to Baker, Gelfond--Schneider,
W\"ustholz, Nishioka, and Adamczewski--Faverjon. Previously proved
project results are identified as such, rather than counted again.
The current investigation and its source checks were completed on
16--17 October 2026.
\end{quote}
\section{What survives a failed transfer}
\label{p08:intro:scope}

The preceding investigation supplied the forward criterion for
algebraic-level quantiles of polynomially tilted arcsine laws and the
bound $N\le8r^2$ for possible cyclotomic orders. It also identified
several invalid inferences: rapid convergence does not control
integer denominators, linear independence does not exclude nonlinear
relations, and finite observations do not automatically determine
arithmetic type. Rather than repeat these failures, we ask for
the extra structure that makes each transfer valid.

The present manuscript develops four main positive directions.
First, cyclotomic necessity can be reversed: every exception pattern
allowed by Galois conjugacy is constructible, and finite moment
constraints do not obstruct it. Second, a logarithmic value with
algebraic coefficients admits a finite norm certificate, which
classifies entropy in every finite alphabet. Third, additive
functional equations permit polynomial dependence to be reduced to
a rational coboundary, turning a relation ideal into finite linear
algebra. Fourth, an exact differential on an elliptic curve reduces
the quantile problem to finitely many algebraic torsion tests.

\begin{center}
\begin{tabular}{L{0.28\linewidth}L{0.61\linewidth}}
\toprule
Earlier available result & Additional result here\\
\midrule
Forward quantile test and finite order bound &
All admissible finite exception sets, their minimum realization
degree, positive beta mixtures, and arbitrarily many exact moment
constraints (Theorems~\ref{p08:cr:realization}, \ref{p08:cm:main}).\\[3pt]
Binary entropy classification &
An exact classifier for every finite algebraic-weight logarithmic
expression in bits, with explicit nonbinary counterexamples
(Theorems~\ref{p08:cte:algorithm}, \ref{p08:cte:five}).\\[3pt]
Linear Mahler relation locus &
The complete polynomial relation ideal and its exact transcendence
degree, including arbitrary finite collections of integer dilations
(Theorem~\ref{p08:ctm:ideal}, Corollary~\ref{p08:ctm:mixed}).\\[3pt]
Unperturbed elliptic division &
A finite exception classifier after adding a nonzero algebraic
exact-differential correction (Theorem~\ref{p08:cet:finite}).\\[3pt]
Fixed-denominator and rank-one obstructions &
Actual reduced denominators for unrestricted odd Taylor schedules;
growing-pole content; full-rank signed odd-zeta pencils
(Theorems~\ref{p08:rlap:odd}, \ref{p08:rlap:primes}, \ref{p08:dct:roots}).\\
\bottomrule
\end{tabular}
\end{center}

\paragraph{The point of the inverse construction.}
For an admissible set $S$ of $s$ probability levels, the minimum
polynomial tilt degree is exactly $s+1$. If the first $M$ moments
must also stay fixed, our constructive degree bound becomes
$2s+M+2$. This latter bound is not asserted to be optimal.
The realization does not say that an arbitrary subset of rational
levels is possible: complete real cyclotomic conjugacy blocks are
necessary and sufficient.

\paragraph{The point of the complete ideal.}
A linear relation test cannot exclude, for example, a quadratic
relation among numbers. Theorem~\ref{p08:ctm:ideal} proves that every
polynomial relation in its specified additive cyclotomic class
is generated by finitely many computable affine linear ones.
The upgrade therefore supplies information that the previous
linear locus did not contain. The deep value and purity theorems
used for this passage are existing results, not new inputs proved here.

\paragraph{What is not counted as a new result.}
The bound $N\le8r^2$, classical elliptic division, the previously
covered Student progressions, and the earlier transcendence proof
for $\log\operatorname{Ein}(1)$ are not new closures.
The known algebraic independence of sparse Lambert blocks is
attributed to Bundschuh--V\"a\"an\"anen. Recent external work on
$\zeta(5)$ and Catalan's constant is not assigned to this project.
Our fifteen named-target attempts are recorded at the end, including
the exact point where each still stops.



\section{Inverse realization of arithmetic quantile patterns}
\label{p08:cr:section}

The cyclotomic divisibility criterion answers whether a specified
polynomial tilt has an algebraic quantile. Here we solve the inverse
problem: which finite sets of algebraic probability levels can be
exactly the levels with algebraic quantiles? Every set allowed by
Galois conjugacy can be realized, at the smallest possible polynomial
degree, by a strictly positive rational tilt arbitrarily close to the
arcsine law. The realization can additionally be chosen as a positive
rational mixture of familiar half-integral beta distributions.
These are self-posed inverse questions, not resolutions of named
historical conjectures. The logarithmic transcendence input is
classical \cite{p07:rga:BakerIII}; prior incomplete-beta applications are
credited below \cite{p08:cr:DhillonSaha,p08:cr:SahaTalk}.

\subsection{Notation and the necessary Galois constraint}

Let $\mu_0$ be the arcsine probability measure on $(0,1)$, with
\[
 d\mu_0(x)=\frac{dx}{\pi\sqrt{x(1-x)}},\qquad
 F_0(x)=\frac{2}{\pi}\arcsin\sqrt{x}.
\]
For a nonzero rational polynomial $w$ nonnegative on $[0,1]$, put
$c_w=\int w\,d\mu_0\in\mathbb Q_{>0}$ and
$d\mu_w=(w/c_w)d\mu_0$. The rationality of $c_w$ follows from
\[
 \int x^j\,d\mu_0(x)=4^{-j}\binom{2j}{j}.
\]
Write $F_w$ and $Q_w$ for its CDF and quantile function. Since a
nonzero polynomial has only finitely many zeros, $F_w$ is strictly
increasing and every interior quantile is unique. Define the
\emph{algebraic-level exception set}
\begin{equation}\label{p08:cr:E}
 \mathcal E(w)=\{q\in\overline{\mathbb Q}\cap(0,1):
                       Q_w(q)\in\overline{\mathbb Q}\}.
\end{equation}
This definition makes no assertion about transcendental probability
levels.

For each integer $N\ge3$, define the finite primitive-angle block
and its corresponding set of quantile locations by
\begin{align}
 \mathcal C_N&=\{2j/N:1\le j<N/2,\ \gcd(j,N)=1\},\label{p08:cr:CN}\\
 \mathcal X_N&=\{\sin^2(\pi j/N):1\le j<N/2,\ \gcd(j,N)=1\}.
 \label{p08:cr:XN}
\end{align}
Each has $h_N=\varphi(N)/2$ elements. The sets for different $N$
are disjoint: $N$ is the order of $e^{i\pi q}$ for a level $q$,
and the map $q\mapsto\sin^2(\pi q/2)$ is injective on $(0,1)$.
Let
\begin{equation}\label{p08:cr:Psi}
 \Psi_N(X)=\prod_{x\in\mathcal X_N}(X-x)\in\mathbb Q[X].
\end{equation}
This is an irreducible monic polynomial of degree $h_N$.
Indeed, if $\zeta$ is a primitive $N$-th root of unity, then
$x=(2-\zeta-\zeta^{-1})/4$ has exactly these conjugates, and
$[\mathbb Q(\zeta+\zeta^{-1}):\mathbb Q]=\varphi(N)/2$:
$\zeta$ is nonreal and quadratic over this real field. The usual
cyclotomic construction therefore computes $\Psi_N$ exactly.
In particular,
\[
 \Psi_3=X-\tfrac34,\quad
 \Psi_4=X-\tfrac12,\quad
 \Psi_6=X-\tfrac14,\quad
 \Psi_5=X^2-\tfrac54X+\tfrac5{16}.
\]

Use the rational differential operator
\begin{equation}\label{p08:cr:L}
 \mathcal L(P)(X)=\frac{(1-2X)P(X)+2X(1-X)P'(X)}2.
\end{equation}
For every polynomial $P$,
\begin{equation}\label{p08:cr:derivative}
 \frac{\mathcal L(P)(x)}{\sqrt{x(1-x)}}
       =\frac{d}{dx}\bigl(\sqrt{x(1-x)}P(x)\bigr),
 \qquad \int\mathcal L(P)\,d\mu_0=0.
\end{equation}
The endpoint terms vanish in the last integral. If $P\ne0$ has
degree $d$ and leading coefficient $a_d$, then $\mathcal L(P)$
has degree $d+1$ and leading coefficient $-(d+1)a_d$.

\begin{lemma}[Arithmetic necessity and exact cyclotomic blocks]
\label{p08:cr:necessity}
For nonconstant $w$, there is a unique nonzero polynomial
$P_w\in\mathbb Q[X]$ such that
\begin{equation}\label{p08:cr:CDF}
 \frac{w}{c_w}=1+\mathcal L(P_w),\qquad
 F_w(x)=F_0(x)+\frac{\sqrt{x(1-x)}P_w(x)}{\pi}.
\end{equation}
It satisfies $\deg P_w=\deg w-1$. Moreover,
\begin{equation}\label{p08:cr:block-test}
 \mathcal E(w)=\bigcup_{\substack{N\ge3\\\Psi_N\mid P_w}}
                  \mathcal C_N.
\end{equation}
This is a finite disjoint union, and its exact cardinality is
\[
 |\mathcal E(w)|=
 \sum_{\substack{N\ge3\\\Psi_N\mid P_w}}\frac{\varphi(N)}2
 \le \deg P_w.
\]
All algebraic irrational levels have transcendental quantiles.
\end{lemma}
\begin{proof}
Put $x=\sin^2\theta$ and expand $w(\sin^2\theta)$ as a rational
linear combination of $\cos(2j\theta)$ for
$0\le j\le\deg(w)$. Its constant coefficient is $c_w$.
Integrate and use
$\sin(2j\theta)=\sin(2\theta)U_{j-1}(\cos(2\theta))$,
where $U_{j-1}\in\mathbb Z[X]$. This gives the CDF formula
with $P_w\in\mathbb Q[X]$. Differentiation gives the first
formula. The kernel of $\mathcal L$ on polynomials is zero:
\eqref{p08:cr:derivative} makes $\sqrt{x(1-x)}P(x)$ constant,
and its endpoint limits force that constant to be zero.
Nonconstancy of $w$ then gives $P_w\ne0$. The leading-coefficient
calculation following \eqref{p08:cr:derivative} proves the degree claim.

Suppose that $q,x\in\overline{\mathbb Q}\cap(0,1)$ and
$F_w(x)=q$. For the principal logarithms, set
\[
 z=1-2x+2i\sqrt{x(1-x)},\qquad \log(-1)=i\pi.
\]
The point $z$ is algebraic on the upper unit semicircle, and
\eqref{p08:cr:CDF} yields
\[
 \log z-q\log(-1)=-i\sqrt{x(1-x)}P_w(x).
\]
The right side is algebraic. Baker's theorem says a nonzero
linear form with algebraic coefficients in logarithms of nonzero
algebraic numbers is transcendental \cite{p07:rga:BakerIII}. Consequently
$P_w(x)=0$ and $z=e^{i\pi q}$. Gelfond--Schneider forces
$q\in\mathbb Q$: an algebraic irrational exponent of $-1$
cannot have an algebraic value. Thus $x=\sin^2(\pi q/2)$.
Conversely, for rational $q$, this algebraic $x$ is the quantile
exactly when $P_w(x)=0$.

Let $N$ be the order of $e^{i\pi q}$. The last equality holds
if and only if the minimal polynomial $\Psi_N$ divides $P_w$.
Every root of this irreducible factor gives the corresponding
level in $\mathcal C_N$. Distinct factors have disjoint roots,
which proves \eqref{p08:cr:block-test} and the cardinality assertion.
\end{proof}

For direct comparison with the preceding divisibility test, if
$d=\deg P_w$ then
\[
 A_{P_w}(Z)=(4Z)^dP_w\!\left(\frac{2-Z-Z^{-1}}4\right)
 =\sum_{k=0}^d p_k(-1)^k(4Z)^{d-k}(Z-1)^{2k}.
\]
This has degree $2d$ and nonzero constant coefficient
$(-1)^dp_d$. For $N\ge3$, $\Psi_N\mid P_w$ is equivalent
to $\Phi_N\mid A_{P_w}$. The factors $N=1,2$ correspond to
endpoints and are excluded. The inverse problem below uses the
real factors $\Psi_N$ to keep the degree count sharp.

\subsection{Realization, the smallest possible degree, and closeness}

\begin{theorem}[Complete inverse realization with minimal degree]
\label{p08:cr:realization}
A finite set $S\subset(0,1)$ is equal to $\mathcal E(w)$ for a
nonconstant rational polynomial $w\ge0$ on $[0,1]$ if and only if
\begin{equation}\label{p08:cr:S}
 S=\bigcup_{N\in\mathcal N}\mathcal C_N
\end{equation}
for a finite set $\mathcal N\subset\{3,4,\ldots\}$.
The empty set is allowed. Set $s=|S|$ and
\[
 H_S(X)=\prod_{N\in\mathcal N}\Psi_N(X),
 \qquad H_{\varnothing}=1.
\]
Then $\deg H_S=s$. For every sufficiently small nonzero rational
$\varepsilon$, the explicit polynomial
\begin{equation}\label{p08:cr:w-epsilon}
 w_\varepsilon=1+\varepsilon\mathcal L(H_S)
\end{equation}
is strictly positive on $[0,1]$, is already normalized by
$c_{w_\varepsilon}=1$, and satisfies
\[
 \mathcal E(w_\varepsilon)=S,\qquad
 \deg w_\varepsilon=s+1.
\]
Every nonconstant polynomial realization of $S$ has degree at
least $s+1$. Thus the displayed degree is minimal, including
$s=0$. At minimal degree its correction polynomial is necessarily
a nonzero rational multiple of $H_S$.

For every fixed integer $k\ge0$, these realizations can be made
arbitrarily close to the arcsine law in the $C^k[0,1]$ norm of
the density ratio $w_\varepsilon$, in total variation, in uniform
CDF distance, and in uniform quantile distance.
\end{theorem}
\begin{proof}
Necessity is Lemma~\ref{p08:cr:necessity}. Conversely, let $S$ have
\eqref{p08:cr:S}. Since all roots of $H_S$ are exactly the locations
in the chosen disjoint cyclotomic orbits, its degree is $s$.
The polynomial $\mathcal L(H_S)$ is nonzero. Choose a positive
rational bound $B\ge\|\mathcal L(H_S)\|_\infty$; for example,
one plus the sum of the absolute values of its rational
coefficients suffices. Any rational $\varepsilon$ with
$0<|\varepsilon|<1/(2B)$ gives $w_\varepsilon>1/2$.
Equation \eqref{p08:cr:derivative} gives normalization exactly, and
its integrated version gives correction $P=\varepsilon H_S$.
Lemma~\ref{p08:cr:necessity} proves that its exception set is exactly
$S$. The degree calculation gives $s+1$.

For any other nonconstant realization, all the distinct factors
of $H_S$ divide its correction polynomial. Therefore
$\deg P_w\ge s$ and $\deg w=\deg P_w+1\ge s+1$.
At equality the quotient $P_w/H_S$ is a nonzero rational constant.
For $S=\varnothing$, a constant correction gives degree one and
no exceptions. A constant tilt cannot realize a finite $S$,
because it is the arcsine law and has all rational interior levels
as exceptions.

For quantitative closeness, set $L=\mathcal L(H_S)$. Directly,
\begin{align}
 \|w_\varepsilon-1\|_{C^k}
     &=|\varepsilon|\|L\|_{C^k},\label{p08:cr:Ck}\\
 \|\mu_{w_\varepsilon}-\mu_0\|_{\mathrm{TV}}
     &=\frac{|\varepsilon|}{2}\int|L|\,d\mu_0
       \le\frac{|\varepsilon|}{2}\|L\|_\infty,\label{p08:cr:TV}\\
 \|F_{w_\varepsilon}-F_0\|_\infty
     &\le\frac{|\varepsilon|}{2\pi}\|H_S\|_\infty.
 \label{p08:cr:CDF-bound}
\end{align}
Here total variation is the supremum of the differences on
measurable sets. The inverse arcsine CDF is
$Q_0(q)=\sin^2(\pi q/2)$ and is $\pi/2$-Lipschitz on $[0,1]$.
For $x=Q_{w_\varepsilon}(q)$, the CDF estimate gives
$|F_0(x)-q|\le\|F_{w_\varepsilon}-F_0\|_\infty$.
Applying $Q_0$ proves
\begin{equation}\label{p08:cr:quantile-bound}
 \sup_{0\le q\le1}|Q_{w_\varepsilon}(q)-Q_0(q)|
       \le\frac{|\varepsilon|}{4}\|H_S\|_\infty.
\end{equation}
All these distances tend to zero with a sequence of nonzero
rational $\varepsilon\to0$.
\end{proof}

The normalized realizations of minimal degree admit an exact further
description. Put $m_S=\min_{[0,1]}\mathcal L(H_S)$ and
$M_S=\max_{[0,1]}\mathcal L(H_S)$. Then $m_S<0<M_S$, because
this nonzero polynomial has zero arcsine mean. All normalized
minimal-degree realizations are $1+\varepsilon\mathcal L(H_S)$,
with rational $\varepsilon\ne0$. They are strictly positive
exactly when
\[
 -\frac1{M_S}<\varepsilon<-\frac1{m_S};
\]
nonnegativity allows the two endpoints. The extrema and hence the
endpoints are effectively computable real algebraic numbers.
This follows by evaluating the rational polynomial at endpoints
and at the finitely many real roots of its derivative.

The theorem separates a genuine obstruction from an artificial one:
a single level $2/5$ cannot occur alone, since it forces the other
level $4/5$ in $\mathcal C_5$; but the complete pair can occur
arbitrarily close to the arcsine distribution. The level $1/3$
can occur alone because $\mathcal C_6=\{1/3\}$.
More generally, an arbitrary proposed finite list is decided by
computing the orders of $e^{i\pi q}$ and checking closure under
the corresponding blocks. The bound on the number of exceptions
is attained by every realization above: its correction has exactly
$s$ simple interior roots, and all of them give algebraic quantiles.

\subsection{Realization inside positive rational beta mixtures}

\begin{theorem}[A positive mixture realization]
\label{p08:cr:mixture}
Every finite set $S$ in Theorem~\ref{p08:cr:realization} can be
realized by a convex mixture with strictly positive rational weights
of the $s+2$ distributions
\[
 \operatorname{Beta}(k+1/2,s+3/2-k),\qquad 0\le k\le s+1.
\]
The normalized tilt has the minimal degree $s+1$ and can be made
arbitrarily close to the arcsine law in all the senses stated in
Theorem~\ref{p08:cr:realization}. No minimality assertion is made
about the number of mixture components.
\end{theorem}
\begin{proof}
Put $\ell=s+1$. The Bernstein polynomials
$b_{k,\ell}(x)=\binom\ell kx^k(1-x)^{\ell-k}$ form a basis
over $\mathbb Q$ for polynomials of degree at most $\ell$, and
$\sum_{k=0}^{\ell}b_{k,\ell}=1$. Thus the constructed tilt has
an exact expansion
\[
 w_\varepsilon(x)=\sum_{k=0}^{\ell}
                  (1+\varepsilon a_k)b_{k,\ell}(x),
 \qquad a_k\in\mathbb Q.
\]
Make $\varepsilon\ne0$ rational sufficiently small that all
$1+\varepsilon a_k$ are strictly positive. Define
\[
 K_{k,\ell-k}=\frac{\pi}{B(k+1/2,\ell-k+1/2)}
 =\frac{4^\ell k!(\ell-k)!\ell!}{(2k)!(2\ell-2k)!}
 \in\mathbb Q_{>0},
\]
and
\[
 \lambda_k=(1+\varepsilon a_k)
                       \frac{\binom\ell k}{K_{k,\ell-k}}.
\]
These are positive rational numbers. The identity between densities
is exactly
\[
 w_\varepsilon(x)\frac1{\pi\sqrt{x(1-x)}}
 =\sum_{k=0}^{\ell}\lambda_k
   \frac{x^{k-1/2}(1-x)^{\ell-k-1/2}}
        {B(k+1/2,\ell-k+1/2)}.
\]
Integrating and using $c_{w_\varepsilon}=1$ gives
$\sum_k\lambda_k=1$. The degree, exception set, and closeness
were already proved. Beta normalization and the CDF convention
are standard \cite{p08:cr:DLMF}.
\end{proof}

\begin{example}[Algebraic exceptions created by mixing]
\label{p08:cr:mixture-example}
Take the following rational mixture, in the displayed order:
\[
 \frac1{512}(185,41,131,155)
\]
on
\[
 \operatorname{Beta}(1/2,7/2),\quad
 \operatorname{Beta}(3/2,5/2),\quad
 \operatorname{Beta}(5/2,3/2),\quad
 \operatorname{Beta}(7/2,1/2).
\]
Its algebraic-level algebraic quantiles are exactly
\begin{equation}\label{p08:cr:golden-quantiles}
 Q(2/5)=\frac{5-\sqrt5}{8},\qquad
 Q(4/5)=\frac{5+\sqrt5}{8}.
\end{equation}
Each of the four component distributions has a transcendental
quantile at both of these levels.
\end{example}
\begin{proof}
Use
$P=\Psi_5=X^2-5X/4+5/16$ and
\[
 w=1+\mathcal L(P)
   =\frac{37}{32}-\frac{35}{16}X+5X^2-3X^3.
\]
Its degree-three Bernstein coefficients are
$(37/32,41/96,131/96,31/32)$, all positive.
The conversion in Theorem~\ref{p08:cr:mixture} gives exactly the
four weights above. Lemma~\ref{p08:cr:necessity} gives the full
exception set $\mathcal C_5=\{2/5,4/5\}$ and the two roots
of $\Psi_5$ give \eqref{p08:cr:golden-quantiles}.

For completeness, the four component correction polynomials,
in the same order, are
\[
 2+\frac43(1-X)+\frac{16}{15}(1-X)^2,\quad
 \frac23(4X-1)(3-2X),\quad
 -\frac23(3-4X)(1+2X),\quad
 -2-\frac43X-\frac{16}{15}X^2.
\]
The first and last have no interior zeros; the middle two have
only $1/4$ and $3/4$, corresponding respectively to levels
$1/3$ and $2/3$. Thus neither $2/5$ nor $4/5$ is an
algebraic-quantile level of any component, by the same logarithmic
criterion. This verifies the final assertion individually.
\end{proof}

\subsection{An attempted global classification of the beta quarter exception}
\label{p08:cr:quarter}

A separate question is whether the positive integer solutions of
\begin{equation}\label{p08:cr:quarter-equation}
 I_{1/4}(m+1/2,n+1/2)=1/3
\end{equation}
consist only of $(m,n)=(1,2)$. The construction theorems do not
settle this more rigid two-parameter question. We record an exact
integer recurrence which isolates the remaining arithmetic issue;
we do not infer a global classification from a finite search.

Let $N=m+n$ and define
\[
 E_{m,n}=\int_0^{\pi/6}\sin^{2m}\theta\cos^{2n}\theta\,d\theta
       -\frac13\int_0^{\pi/2}\sin^{2m}\theta\cos^{2n}\theta\,d\theta,
 \qquad R_{m,n}=\frac{2N!4^N}{\sqrt3}E_{m,n}.
\]
Then $R_{0,0}=0$, and integration by parts gives the integer
recurrences
\begin{align}
 R_{0,n}&=2(2n-1)R_{0,n-1}+(n-1)!3^{n-1},&&n\ge1,
 \label{p08:cr:R0}\\
 R_{m,n}&=2(2m-1)R_{m-1,n}-(m+n-1)!3^n,&&m\ge1.
 \label{p08:cr:Rmn}
\end{align}
To verify the second, differentiate
$\sin^{2m-1}\theta\cos^{2n+1}\theta$. Its derivative is
\[
 (2m-1)\sin^{2m-2}\theta\cos^{2n}\theta
 -2N\sin^{2m}\theta\cos^{2n}\theta.
\]
The endpoint contribution at $\pi/6$ is
$\sqrt3\,3^n/4^N$; the complete integral has zero endpoint
contribution. This gives \eqref{p08:cr:Rmn}. For \eqref{p08:cr:R0},
differentiate $\sin\theta\cos^{2n-1}\theta$ instead, whose
endpoint contribution is $\sqrt3\,3^{n-1}/4^n$.
Thus every $R_{m,n}$ is an explicitly computed integer and
\eqref{p08:cr:quarter-equation} is equivalent to $R_{m,n}=0$.

The normalized incomplete beta recurrences
\cite[Equations~8.17.20--8.17.21]{p08:cr:DLMF} show that the left
side of \eqref{p08:cr:quarter-equation} is strictly decreasing in
$m$ and strictly increasing in $n$. Therefore there is at most
one solution in each fixed row or column. Direct calculation gives
$R_{1,2}=0$, and consequently this is the only solution with
$m=1$, or with $n=2$. These statements do not exclude solutions
with other pairs of indices. The missing step is a proof that
\eqref{p08:cr:R0}--\eqref{p08:cr:Rmn} never vanishes again for positive
indices, or a counterexample. No such proof or counterexample is
claimed here.

\subsection{Prior art and scope of the additional results}

Baker's logarithmic value theorem and cyclotomic conjugacy are
classical. Dhillon--Saha study incomplete beta values
\cite{p08:cr:DhillonSaha}, and Saha's primary slides explicitly cover
rational noninteger parameters whose sum is a positive integer
using logarithmic linear forms \cite[slide~4/4]{p08:cr:SahaTalk}.
That covers the arithmetic method used for individual half-integral
beta components above. The preceding paper in this project already
proved the exact cyclotomic test for polynomial tilts; it is
recalled here solely to prove the inverse theorems.

The additional statements are the complete converse realization,
the exact minimal degree, arbitrary closeness, and realization
inside positive rational beta mixtures. They form one connected
inverse-realization theory, not a list of separately solved famous
conjectures. A primary-source search did not identify an explicit
prior formulation of this inverse problem, but the full published
Dhillon--Saha theorem statements were not accessible in the legal
open sources checked. No historical-priority claim is made.


\section{Finite moments can conceal any admissible arithmetic pattern}
\label{p08:cm:section}

The inverse realization theorem can be made compatible with any
prescribed finite list of moments of the arcsine law. This is a stronger
statement than small perturbation: the moment equalities below are exact.
It identifies a specific limitation of transferring statistical
identifiability to arithmetic identifiability.

\begin{theorem}[Moment-preserving inverse realization]\label{p08:cm:main}
Let $S$ be any finite union of the blocks $\mathcal C_N$ in
\eqref{p08:cr:CN}, including the empty union, and put $s=|S|$.
For every integer $M\ge0$, every integer $k\ge0$, and every $\eta>0$,
there exists a nonconstant $w\in\mathbb Q[X]$, strictly positive on
$[0,1]$, such that
\begin{align}
 &\int w\,d\mu_0=1,\qquad \mathcal E(w)=S,\label{p08:cm:pattern}\\
 &\int x^j w(x)\,d\mu_0(x)=4^{-j}\binom{2j}{j}
                  &&(1\le j\le M),\label{p08:cm:moments}\\
 &\|w-1\|_{C^k[0,1]}<\eta,\qquad
                  \deg w\le 2s+M+2.\label{p08:cm:degree}
\end{align}
The construction is effective for rational positive $\eta$, and $w\mu_0$
can additionally be expressed as a strictly positive rational mixture
of finitely many half-integral beta laws.
\end{theorem}
\begin{proof}
Let $H=H_S$ be as in Theorem~\ref{p08:cr:realization}. Consider the positive
measure
\[
                    d\nu(x)=H(x)^2\sqrt{x(1-x)}\,dx.
\]
Its moment matrices are positive definite: a nonzero polynomial cannot
vanish on an interval, and the density is positive except at finitely
many points. Its moments divided by $\pi$ are rational, by the
half-integral beta integral. Gram--Schmidt therefore gives monic
orthogonal polynomials $R_j\in\mathbb Q[X]$ of every degree $j$.
Consecutive polynomials $R_M,R_{M+1}$ have no common complex root.
Indeed, their three-term recurrence has nonzero, positive backward
coefficient; a common root would propagate to $R_0=1$.

Put $D=2s+M+1$ and consider
\[
           B_t=H^2(R_{M+1}+tR_M),\qquad t\in\mathbb Q.
\]
This polynomial has degree $D$ for every $t$. We choose $t$ so that
its only factors $\Psi_N$ are those already dividing $H$.
There are only finitely many candidates: a factor requires
$\varphi(N)/2\le D$, which implies $N\le8D^2$. For completeness,
$\varphi(N)^2\ge N/2$ follows by multiplying the contributions
$p^{a-2}(p-1)^2$ of prime powers $p^a\Vert N$; the only factor
less than $1$ can be $1/2$ from $p^a=2$.
For any root $\alpha$ of a candidate $\Psi_N$ not dividing $H$,
the equation $R_{M+1}(\alpha)+tR_M(\alpha)=0$ forbids at most
one value of $t$, since the two orthogonal polynomials have no
common root. Thus finitely many choices are forbidden. Enumerating
rational $t$ and using exact polynomial divisibility tests finds an
allowed choice after finitely many steps.

Choose a sufficiently small nonzero rational $\epsilon$ and set
\[
            P=\epsilon B_t,\qquad w=1+\mathcal L(P).
\]
The derivative identity \eqref{p08:cr:derivative} gives normalization;
Lemma~\ref{p08:cr:necessity} gives the exact exception set $S$.
For $1\le j\le M$, integration by parts gives
\begin{align*}
 \int x^j(w-1)\,d\mu_0
 &=-\frac{j}{\pi}\int_0^1
                x^{j-1}\sqrt{x(1-x)}P(x)\,dx\\
 &=-\frac{j\epsilon}{\pi}\int_0^1
                x^{j-1}(R_{M+1}+tR_M)\,d\nu=0.
\end{align*}
Boundary terms vanish and orthogonality proves the last equality.
Shrinking $|\epsilon|$ ensures positivity and the stated $C^k$ bound;
these can be enforced by rational coefficient bounds. The degree is
$D+1$ by the leading coefficient calculation for $\mathcal L$.
Finally express $w$ in the Bernstein basis of degree $D+1$.
The constant polynomial $1$ has all Bernstein coefficients equal to
$1$, so all coefficients remain positive for small $|\epsilon|$.
Multiplying these basis functions by $d\mu_0$ and normalizing gives
the claimed rational beta mixture, as in Theorem~\ref{p08:cr:mixture}.
\end{proof}

\begin{example}[The mean does not determine median arithmetic]
\label{p08:cm:example}
Set
\[
 P(x)=\frac1{16}\left(x^2-\frac5{16}\right),\qquad
 w(x)=1+\frac1{16}\left(-3x^3+\frac52x^2+
                         \frac5{16}x-\frac5{32}\right).
\]
The sum of the absolute coefficients of $w-1$ is $191/512<1$,
so $w>0$ on $[0,1]$. Its integral against $\mu_0$ is $1$.
Moreover $\int xw\,d\mu_0=1/2$, because
\[
 \frac{\int_0^1 x^2\sqrt{x(1-x)}\,dx}
      {\int_0^1\sqrt{x(1-x)}\,dx}=\frac5{16}.
\]
The positive root $\sqrt5/4$ of $P$ has the negative conjugate
$-\sqrt5/4$. It cannot be a root of any $\Psi_N$, all of whose
conjugates lie in $(0,1)$. Hence $\mathcal E(w)$ is empty.
This law has the same algebraic mean as the arcsine law but a
transcendental median; indeed every algebraic interior probability
has a transcendental quantile.
\end{example}

\begin{corollary}[No finite-moment classifier on this class]
\label{p08:cm:obstruction}
For each fixed $M$, the first $M$ exact moments do not determine
the algebraic-level exception set of a strictly positive rational
polynomial tilt of the arcsine law. This remains true in every
$C^k$ neighborhood of the constant weight $1$ and within positive
rational finite beta mixtures.
\end{corollary}
\begin{proof}
Apply Theorem~\ref{p08:cm:main} to two different admissible sets, for
example $S=\varnothing$ and $S=\{1/2\}=\mathcal C_4$.
The resulting laws have identical specified moments and different
exception sets.
\end{proof}

The degree bound is a constructive upper bound, not a claim of
minimality under the moment constraints. The result does not prevent
classification from the full polynomial density, which is precisely
what the cyclotomic algorithm accomplishes.


\section{Exact arithmetic of finite entropy: a norm algorithm and counterexamples}
\label{p08:cte:section}

The binary entropy phenomenon does not extend unchanged to larger
alphabets. We give an exact algorithm for deciding the algebraicity of
finite information quantities at algebraic data, then exhibit five
algebraic probabilities whose Shannon entropy in bits is an algebraic
irrational number. Even rational probabilities with rational entropy
need not be dyadic. These are precise corrections to tempting extensions,
not resolutions of historically named conjectures.

The relevant prior result is Blakaj--Wolf's Theorem~2
\cite{p08:cte:BlakajWolf}: entropy in nats at algebraic spectra is zero or
transcendental, whereas entropy in an algebraic base at rational spectra
is rational or transcendental. Their second hypothesis is rationality
of the spectrum, not algebraicity. Our example respects both statements.
Boche--Schaefer--Poor prove the binary rational-input exception theorem
\cite[Appendix B, Theorem 4]{p08:cte:BSP}. Our previous binary algebraic-input
classification therefore cannot be extrapolated solely from its
rational-input special case. The results below use classical Baker theory;
priority of the explicit algorithm or examples beyond these checked
primary sources is not certified.

\subsection{A finite exact test for algebraic information values}

Write $\mathbb A=\overline{\mathbb Q}\cap\mathbb R$.
All logarithms in this subsection are real logarithms of positive
numbers. Algebraic inputs are specified by defining polynomials and
isolating data, so exact field operations, minimal polynomials,
number-field norms, and rational linear algebra are effective.

\begin{lemma}[Detecting a rational power of two]\label{p08:cte:norm}
Let $A>0$ be algebraic, let $d=[\mathbb Q(A):\mathbb Q]$, and let
$n_A=|N_{\mathbb Q(A)/\mathbb Q}(A)|\in\mathbb Q_{>0}$.
Then
\[
 A\in 2^{\mathbb Q}
 \quad\Longleftrightarrow\quad
 n_A=2^m\text{ for some }m\in\mathbb Z
 \quad\hbox{and}\quad A^d=2^m.
\]
These conditions are decidable by finite exact algebraic arithmetic.
When they hold, $\log_2 A=m/d$.
\end{lemma}
\begin{proof}
If $A=2^{a/b}$ with $a\in\mathbb Z$ and $b\ge1$, then every conjugate
of $A$ has modulus $A$, since its $b$-th power is $2^a$.
Consequently $n_A=A^d=2^{ad/b}$. A positive rational number that is a
rational power of $2$ is an integral power of $2$: clear the exponent's
denominator and compare prime valuations. Thus $n_A=2^m$ as required.
Conversely, positivity and $A^d=2^m$ give $A=2^{m/d}$.
To implement the test, reduce the rational number $n_A$ to lowest terms,
check that its numerator and denominator are powers of $2$, and check
the stated algebraic equality. This also covers $A=1$.
\end{proof}

We will use the following standard coefficient-splitting consequence of
Baker's theorem \cite{p07:rlog:BakerIII,p03:mi3:Baker}. If algebraic numbers
$\theta_1,\ldots,\theta_s$ are linearly independent over $\mathbb Q$,
and each $L_j$ is a rational linear combination of real logarithms of
positive algebraic numbers, then
\[
                  \sum_j\theta_jL_j=0
               \quad\Longrightarrow\quad L_j=0\quad\hbox{for all }j.
\]
Indeed, expand the $L_j$ in a $\mathbb Q$-basis of the finitely many
logarithms involved. Baker's theorem says that this basis remains
linearly independent over $\overline{\mathbb Q}$; the remaining
coefficient equations split by independence of the $\theta_j$.
This is an argument for correctness and does not require an algorithm
for finding a basis of logarithms.

\begin{theorem}[Exact arithmetic classification]\label{p08:cte:algorithm}
Given finitely many $\lambda_i\in\mathbb A$ and $\alpha_i\in\mathbb A_{>0}$,
there is a terminating algorithm which determines whether
\begin{equation}\label{p08:cte:value}
             V=\frac{\sum_i\lambda_i\log\alpha_i}{\log2}
\end{equation}
is algebraic or transcendental. In the algebraic case it returns $V$
as an exact algebraic number and determines whether it is rational.

More explicitly, choose a $\mathbb Q$-basis $\theta_1,\ldots,\theta_s$
of the span of the $\lambda_i$, write
$\lambda_i=\sum_j r_{ji}\theta_j$ with $r_{ji}\in\mathbb Q$, and choose
$D\ge1$ with $Dr_{ji}\in\mathbb Z$. Put
\begin{equation}\label{p08:cte:A}
              A_j=\prod_i\alpha_i^{Dr_{ji}}.
\end{equation}
Then $V$ is algebraic if and only if every $A_j$ is a rational power
of $2$. If the norm test of Lemma~\ref{p08:cte:norm} returns
$\log_2 A_j=m_j/d_j$, then
\begin{equation}\label{p08:cte:output}
                  V=\frac1D\sum_j\theta_j\frac{m_j}{d_j}.
\end{equation}
If all coefficients are zero, the algorithm simply returns $V=0$.
\end{theorem}
\begin{proof}
The sufficiency and formula follow from
$DV\log2=\sum_j\theta_j\log A_j$.
Suppose conversely that $V$ is algebraic. First $V$ belongs to the
$\mathbb Q$-span of the $\theta_j$. Otherwise these numbers together
with $V$ are independent over $\mathbb Q$, and coefficient splitting
in this equality forces $D\log2=0$, a contradiction.
Write $V=\sum_j s_j\theta_j$, $s_j\in\mathbb Q$. Coefficient splitting
now gives
\[
                    \log A_j-Ds_j\log2=0
                       \qquad(1\le j\le s).
\]
Hence every $A_j\in2^{\mathbb Q}$.

All the $\theta_j$, rational coefficients, and products $A_j$ are
computed in a common number field. The norm test uses only exact
algebraic operations and finitely many checks on rational integers;
no search for unknown logarithmic relations or unbounded approximation
is needed. Formula~\eqref{p08:cte:output} gives exact algebraic output.
Testing whether its minimal polynomial has degree one decides
rationality. If any norm test fails, the proved equivalence certifies
transcendence. No running-time bound is asserted.
\end{proof}

The equivalent multiplicative-lattice criterion is also explicit.
For fixed positive algebraic $\alpha_1,\ldots,\alpha_n$, set
\[
 \Lambda=\{(k_1,\ldots,k_n,k_0)\in\mathbb Z^{n+1}:
               \alpha_1^{k_1}\cdots\alpha_n^{k_n}2^{k_0}=1\}.
\]
Baker's theorem gives, for any proposed algebraic $v$,
\[
 V=v\quad\Longleftrightarrow\quad
 (\lambda_1,\ldots,\lambda_n,-v)
       \in\operatorname{span}_{\overline{\mathbb Q}}\Lambda.
\]
To verify this equivalence, the rational kernel of the row of real
logarithms is exactly $\mathbb Q\Lambda$ after clearing denominators,
and Baker says its algebraic kernel is its scalar extension.
The norm algorithm avoids having to compute the whole lattice when
only the arithmetic nature of the single value~\eqref{p08:cte:value} is sought.

\begin{corollary}[Finite information quantities]\label{p08:cte:information}
At exact algebraic finite probability data, algebraicity or
transcendence in bits is decidable for Shannon entropy, conditional
entropy, mutual information, and finite Kullback--Leibler divergence.
The algebraic value, when present, is exactly recoverable. The same
holds for any specified algebraic linear combination of these quantities.
\end{corollary}
\begin{proof}
Delete zero-mass terms. Positive joint and marginal probabilities are
positive algebraic numbers. Each indicated quantity is a finite sum
of algebraic coefficients times logarithms of positive algebraic
probabilities or their ratios, divided by $\log2$.
For relative entropy, a positive first-distribution mass paired with
zero second-distribution mass is detected exactly and gives $+\infty$;
otherwise the finite-value theorem applies. Thus the support cases,
as well as all hypotheses of Theorem~\ref{p08:cte:algorithm}, are decidable.
\end{proof}

This is a pointwise arithmetic decision procedure with an exact
algebraic input description. It does not make a positive-dimensional
entropy level set semialgebraic, and does not contradict
Blakaj--Wolf's non-semialgebraicity theorem.

\subsection{Rational probabilities: the full valuation criterion}

\begin{proposition}\label{p08:cte:rational}
Let $p_i=a_i/N>0$, where $a_i\in\mathbb Z_{>0}$ and $\sum_i a_i=N$.
No coprimality hypothesis on this representation is needed. Then
$H_2(p)=-\sum_i p_i\log_2p_i$ is algebraic if and only if, for every
odd prime $\ell$,
\begin{equation}\label{p08:cte:odd}
          N v_\ell(N)=\sum_i a_i v_\ell(a_i).
\end{equation}
Only primes dividing $N\prod_i a_i$ need be tested. When these
conditions hold, the entropy is the rational number
\begin{equation}\label{p08:cte:rationalvalue}
       H_2(p)=v_2(N)-\frac1N\sum_i a_i v_2(a_i).
\end{equation}
If one condition fails, the entropy is transcendental.
\end{proposition}
\begin{proof}
The identity
\[
 NH_2(p)\log2=\log\left(\frac{N^N}{\prod_i a_i^{a_i}}\right)
\]
together with Theorem~\ref{p08:cte:algorithm} and
Lemma~\ref{p08:cte:norm}, here applied to a positive rational number,
shows that algebraicity is equivalent to the rational number in
parentheses being an integral power of $2$. Prime valuations give
exactly~\eqref{p08:cte:odd} and~\eqref{p08:cte:rationalvalue}.
This refines the known rational-or-transcendental dichotomy by an
explicit exact test; the dichotomy itself is already in
\cite{p08:cte:BlakajWolf}.
\end{proof}

\begin{example}[Odd-prime cancellation with four outcomes]\label{p08:cte:four}
The rational distribution
\[
                  p=\left(\frac38,\frac14,\frac13,\frac1{24}\right)
\]
has entropy $7/4$ bits, despite its non-dyadic probabilities.
In the denominator-$24$ representation its numerators are $(9,6,8,1)$.
The only odd prime to check is $3$, and
\[
             24v_3(24)=24=9v_3(9)+6v_3(6).
\]
The $2$-adic formula gives
$3-(6+24)/24=7/4$. Thus neither cancellation of odd primes nor
rational entropy requires each probability to be a negative integral
power of $2$. We make no claim that four is the least possible
alphabet size for such a rational counterexample.
\end{example}

\subsection{A five-outcome algebraic irrational entropy}

The construction uses a double root of a positive Laurent polynomial.
This mechanism is useful independently of the particular example.

\begin{proposition}[Double-root construction]\label{p08:cte:construction}
Let $k_i,m_i\in\mathbb Z$ and $u>0$ be algebraic. Suppose
\[
 F(t)=\sum_{i=1}^n2^{-k_i}t^{m_i},\qquad F(u)=1,\qquad F'(u)=0.
\]
Then $p_i=2^{-k_i}u^{m_i}$ form a positive algebraic probability vector,
and
\begin{equation}\label{p08:cte:doublevalue}
                        H_2(p)=\sum_i k_i p_i.
\end{equation}
In particular this entropy is algebraic. If every coefficient of a
positive Laurent polynomial $F$ is dyadic rational, its monomials can
be split into finitely many terms of the displayed form.
\end{proposition}
\begin{proof}
Normalization is $F(u)=1$, and differentiating gives
$\sum_i m_i p_i=uF'(u)=0$. Therefore
\[
 -\sum_i p_i\log_2p_i
       =\sum_i k_i p_i-\log_2u\sum_i m_i p_i
       =\sum_i k_i p_i.
\]
Every positive dyadic rational is a finite sum of integral powers of
$2$, giving the last assertion. Notice that the double-root identity
cancels the coefficient of $\log u$ exactly; a numerical critical point
would not supply this arithmetic certificate.
\end{proof}

\begin{theorem}[Explicit algebraic irrational entropy]\label{p08:cte:five}
Put $u=1+\sqrt3$. The five probabilities
\begin{equation}\label{p08:cte:fivevector}
 p=\left(\frac{u^2}{16},\frac1{2u},\frac2{u^2},
                   \frac1{4u^2},\frac1{u^3}\right)
\end{equation}
are positive algebraic irrational numbers with sum $1$, and
\begin{equation}\label{p08:cte:fiveentropy}
                    H_2(p)=\frac{6\sqrt3-3}{4}.
\end{equation}
Thus the rational-or-transcendental dichotomy fails for finite
algebraic probability vectors in base $2$.
\end{theorem}
\begin{proof}
The polynomial identity
\begin{equation}\label{p08:cte:factorization}
 (t^2-2t-2)^2(t+4)=t^5-16t^3+8t^2+36t+16
\end{equation}
shows that
\[
 F(t)=\frac{t^2}{16}+\frac1{2t}
                  +\frac2{t^2}+\frac1{4t^2}+\frac1{t^3}
\]
satisfies $F(u)=1$ and $F'(u)=0$; indeed
$F(t)-1=(t^2-2t-2)^2(t+4)/(16t^3)$.
Apply Proposition~\ref{p08:cte:construction} with
\[
           (k_i)=(4,1,-1,2,0),\qquad (m_i)=(2,-1,-2,-2,-3).
\]
The resulting entropy is
\[
   4p_1+p_2-p_3+2p_4
       =\frac{u^2}{4}+\frac1{2u}-\frac3{2u^2}
       =\frac{6u-9}{4}=\frac{6\sqrt3-3}{4},
\]
where $u^2=2u+2$ was used. For a separate normalization check,
writing $s=\sqrt3$ gives
\[
 p=\left(\frac{2+s}{8},\frac{s-1}{4},2-s,
                    \frac{2-s}{8},\frac{3s-5}{4}\right).
\]
Every entry is positive and irrational, and their rational and
$s$ coefficients sum to $1$ and $0$, respectively. The entropy in
\eqref{p08:cte:fiveentropy} is algebraic irrational since $s$ is irrational.
\end{proof}

\begin{corollary}[Rational entropy without rational probabilities]
\label{p08:cte:eleven}
Replace the second and third atoms of~\eqref{p08:cte:fivevector} by four
equal subatoms each, leaving the other three atoms intact. The resulting
$11$-outcome probability vector has only algebraic irrational entries
and has rational entropy $11/4$ bits.
\end{corollary}
\begin{proof}
Splitting an atom of mass $a$ into four equal atoms increases the
entropy by $2a$. Therefore the new entropy is
\[
       \frac{6\sqrt3-3}{4}+2\frac{\sqrt3-1}{4}+2(2-\sqrt3)
                          =\frac{11}{4}.
\]
Division by $4$ preserves irrationality of the affected entries.
\end{proof}

These examples answer two distinct extrapolation questions negatively:
algebraic entropy does not force rational entropy, and rational entropy
does not force rational probabilities. Example~\ref{p08:cte:four} separately
shows that rational probabilities with rational entropy need not be
dyadic. No smallest-alphabet assertion is needed for any conclusion.
The same examples are spectra of diagonal algebraic density matrices,
so they also apply to von Neumann entropy in bits.

\subsection{Three named constants: exact outcomes and limits}

\paragraph{$e+\pi$ and $e\pi$.}
These remain unresolved here, including their irrationality.
Waldschmidt's primary notes record the classical gap
\cite{p07:rlog:Waldschmidt}; a 2026 preprint still treats $e+\pi$ as an open
irrationality target \cite{p07:rlog:Yu}. The norm algorithm requires
algebraic logarithm bases and algebraic coefficients. Assuming
$e+\pi=a\in\overline{\mathbb Q}$ gives $\log(a-\pi)=1$, whose base is
$e$, not an algebraic number. Assuming $e\pi=a$ gives
$\log(a/\pi)=1$ with the same defect. Replacing $\pi$ by a branch of
$\log(-1)/i$ introduces nested logarithms and does not repair this
hypothesis. The new entropy construction supplies no theorem about
such expressions.

There is also a precise obstruction to transferring an $E$-function
value theorem: $e=e^1$ is an algebraic-point $E$-value, but a joint
$E$-function system supplying $\pi$ at an algebraic point has not been
provided. Writing $\sin\pi=0$ instead evaluates at the unknown
transcendental point $\pi$, outside the algebraic-point specialization
theorem. The elementary disjunction that at least one of $e+\pi,e\pi$
is transcendental follows from the quadratic with these two
coefficients and is classical; it is not an individual closure.

\paragraph{$\log\operatorname{Ein}(1)$.}
This constant is already proved transcendental in the earlier
$E$-function section of this project, so it is not counted as a new
result of this turn. The applicable input is the stronger established
statement
\[
 \operatorname{Ein}(1)\text{ is transcendental over }
 \mathcal F_{\exp}=\overline{\mathbb Q}
                   (e^a:a\in\overline{\mathbb Q}).
\]
For completeness, its mechanism is the Kolchin--Ostrowski primitive
criterion followed by the Siegel--Shidlovskii specialization theorem
\cite{p08:cte:Ostrowski,p08:cte:Beukers}. For any finite set of algebraic
exponents, choose a rational basis $b_1,\ldots,b_r$ of their span,
rescaled so that the exponents and $1$ are integral combinations.
In the abstract field
$\mathbb C(z,w_1,\ldots,w_r)$, $w_i=e^{b_i z}$, the derivative
$(1-e^{-z})/z$ has a nonzero simple pole at the formal divisor $z=0$.
A rational primitive cannot have such a derivative: a pole of order
$m$ gives derivative order $m+1$, while a regular function has regular
derivative, with the $w$-derivation unable to cancel the leading term.
The primitive criterion gives transcendence of $\operatorname{Ein}(z)$
over that exponential field. The vector consisting of $1$, these
exponentials, $e^{-z}$, and $\operatorname{Ein}(z)$ is an $E$-function
solution of a homogeneous system whose only finite possible
singularity is $0$. Specialization at $1$ therefore preserves the
transcendence degree \cite[Theorem 1.1]{p08:cte:Beukers}, proving the stated
finite-field independence. A relation over $\mathcal F_{\exp}$ uses
only finitely many exponents, so the whole-field claim follows.

Since $\operatorname{Ein}(1)>0$, its real logarithm is defined. If it
were algebraic $b$, then $\operatorname{Ein}(1)=e^b$ would belong to
$\mathcal F_{\exp}$, a contradiction. This proof genuinely applies to
this constant, but neither creates a new proof of algebraic
independence of its logarithms nor solves $e+\pi$ or $e\pi$.


\section{From roots of unity to elliptic torsion: another finite classifier}
\label{p08:cet:section}

The cyclotomic test has a genus-one counterpart. Here a finite
torsion test replaces cyclotomic divisibility. The transcendence input
is the classical analytic subgroup theorem, not a new transcendence
theorem. We give the full reduction and a terminating algorithm for
a specified nonzero correction. Earlier project manuscripts
already treated unperturbed elliptic division; the additional object
here is the exact algebraic correction and its finite exceptional set.

\begin{lemma}[Two consequences of the analytic subgroup theorem]
\label{p08:cet:AST}
Let $E$ be an elliptic curve over $\overline{\mathbb Q}$ and use a
nonzero invariant differential defined over $\overline{\mathbb Q}$
to identify its Lie algebra with $\mathbb C$.
\begin{enumerate}[label=(\roman*)]
\item If $u_1,\ldots,u_m$ are elliptic logarithms of algebraic
points and $c_1,\ldots,c_m$ are algebraic, then
$\sum c_ju_j$ is either zero or transcendental.
\item If $E$ is real, $u,v$ are nonzero real elliptic logarithms
of algebraic points, and $u/v$ is algebraic, then $u/v\in\mathbb Q$.
\end{enumerate}
\end{lemma}
\begin{proof}
We use the logarithmic form of the analytic subgroup theorem:
the smallest $\overline{\mathbb Q}$-linear Lie subspace containing
a logarithm of an algebraic point is the Lie algebra of an algebraic
subgroup \cite{p08:cet:Wustholz,p08:cet:Chang}.
For (i), suppose $a=\sum c_ju_j$ is nonzero and algebraic, and
apply the theorem in $\mathbb G_a\times E^m$ to
$(a,u_1,\ldots,u_m)$, in the hyperplane
$t-\sum c_jz_j=0$. A connected algebraic subgroup of
$\mathbb G_a\times E^m$ that projects nontrivially to
$\mathbb G_a$ contains $\mathbb G_a\times\{0\}$.
Indeed, its affine part maps trivially to the abelian variety,
whereas a homomorphism from an abelian variety to $\mathbb G_a$
is zero. This contradicts containment of its Lie algebra in the
displayed hyperplane.

For (ii), put $q=u/v$. The analytic subgroup theorem in $E^2$
applied to $(u,v)$ places it in the Lie algebra of a one-dimensional
algebraic subgroup lying on $z_1=qz_2$. Both projections are
nonzero isogenies, so their tangent ratio belongs to
$\operatorname{End}^0(E)$. In characteristic zero this algebra is
$\mathbb Q$ or an imaginary quadratic field; its intersection
with $\mathbb R$ is $\mathbb Q$. Thus the real number $q$ is
rational. This is the usual elliptic form of the
Gelfond--Schneider phenomenon.
\end{proof}

Consider the smooth projective curve
\[
                 E:\ y^2=1-x^4,\qquad O=(-1,0),
\]
with the indicated point as group identity. Its real locus is the
oval formed by the upper and lower branches for $-1\le x\le1$.
The differential $dx/y$ is regular and nonzero. Define
\[
 I(x)=\int_{-1}^x\frac{dt}{\sqrt{1-t^4}},\qquad
 \Omega=I(1)>0.
\]
On the upper branch, the elliptic logarithm of $(x,y)$ is $I(x)$.
The full real period is $2\Omega$ and $(1,0)$ is a point of order
two: the involution $(x,y)\mapsto(x,-y)$ is negation with this
choice of origin.

Let $P\in\mathbb Q[X]\setminus\{0\}$ and choose any nonzero
rational $\epsilon$ sufficiently small that
\begin{equation}\label{p08:cet:weight}
 w(x)=1+\epsilon\bigl((1-x^4)P'(x)-2x^3P(x)\bigr)>0
                      \qquad(-1\le x\le1).
\end{equation}
Such a choice is effective by a coefficient bound. The probability
density and its CDF are
\begin{equation}\label{p08:cet:CDF}
 f(x)=\frac{w(x)}{\Omega\sqrt{1-x^4}},\qquad
 F(x)=\frac{I(x)+\epsilon\sqrt{1-x^4}P(x)}{\Omega}.
\end{equation}
Differentiation proves this formula, and the correction vanishes at
both endpoints. Let $Q=F^{-1}$ on $(0,1)$.

\begin{theorem}[Effective elliptic exception classifier]
\label{p08:cet:finite}
For algebraic $q\in(0,1)$, the quantile $Q(q)$ is algebraic if and
only if there is a real root $x\in(-1,1)$ of $P$ for which
\[
 T_x=(x,\sqrt{1-x^4})\in E(\overline{\mathbb Q})
       \text{ is torsion},\qquad q=I(x)/\Omega.
\]
Every such $q$ is rational. The entire list of exceptions, including
their exact rational levels and algebraic quantiles, is computable
by a terminating algorithm. There are at most $\deg P$ exceptions;
for nonzero constant $P$ there are none.
\end{theorem}
\begin{proof}
If $x=Q(q)$ and $x,q$ are algebraic, \eqref{p08:cet:CDF} gives
\[
                I(x)-q\Omega=-\epsilon\sqrt{1-x^4}P(x).
\]
Both $I(x)$ and $\Omega$ are elliptic logarithms of algebraic
points. The right side is algebraic, so Lemma~\ref{p08:cet:AST}(i)
forces it to vanish. Hence $P(x)=0$ and $I(x)=q\Omega$.
Part (ii) of the lemma gives $q\in\mathbb Q$. Since $2\Omega$
is a period, this makes $T_x$ torsion.
Conversely, a torsion point on the upper real arc has
$I(x)/(2\Omega)\in\mathbb Q$. If $P(x)=0$, the CDF formula
gives the stated algebraic quantile. Strict increase of $F$ ensures
that different roots give different probability levels.

It remains to justify effectiveness, rather than merely recognize
torsion by an unbounded search. Isolate the finitely many real roots
of $P$ in $(-1,1)$. For each one, work in the explicit number field
$K=\mathbb Q(x,\sqrt{1-x^4})$ and use any algebraic Weierstrass
model of $(E,O)$. Choose two primes of good reduction of $K$ with
distinct residue characteristics $p$ and $\ell$. They can be found
by enumeration since there are only finitely many primes of bad
reduction. Compute the finite integers
\[
             A=\#E(k_{\mathfrak p}),\qquad
             B=\#E(k_{\mathfrak l}),\qquad D=AB.
\]
If $T_x$ has finite order $n$, its prime-to-$p$ part injects under
the first reduction, and its $p$-primary part injects under the
second reduction. Consequently $n\mid D$.
The reduction fact follows from the formal group: multiplication
by an integer prime to the residue characteristic is invertible
on the reduction kernel. Thus
\[
                 T_x\text{ is torsion}\quad\Longleftrightarrow
                               [D]T_x=O.
\]
This is a finite exact algebraic computation. If it succeeds,
factor $D$ and test its divisors to find the exact order $n$.
The real arc coordinate is then $I(x)/(2\Omega)=j/n$ with
$0<j<n/2$ and $\gcd(j,n)=1$. Certified real approximations to
the convergent integrals separate these finitely many candidates
and recover the unique integer $j$, yielding the exact answer
$q=2j/n$. The approximation is used only to select from an
already proved finite rational list. This completes the algorithm.
\end{proof}

\begin{example}
With $P=1$ and sufficiently small nonzero rational $\epsilon$,
every algebraic-level quantile is transcendental. With $P=x$,
the only algebraic-level algebraic quantile is $Q(1/2)=0$.
Indeed $I(0)=\Omega/2$ by symmetry, and $P$ has no other root.
These statements concern the perturbed laws in \eqref{p08:cet:CDF};
for the unperturbed law $\epsilon=0$, every rational-level
quantile is algebraic by elliptic division.
\end{example}

This transfer removes the special dependence on roots of unity:
the finite candidates are now zeros of a correction intersected
with an elliptic torsion subgroup. It does not decide arbitrary
elliptic constants or arbitrary real numbers. Nor is a polynomial
complexity bound for the algorithm asserted.



\section{Even-power rational densities and arithmetic transfer}
\label{p08:rd:section}

We give a complete arithmetic classification for algebraic-probability
quantiles of the rational densities proportional to $1/(1+x^{2m})$.
The mechanism is a sign-definite logarithmic part of the primitive,
separated from its argument part by Baker's theorem.  This is a different
family from the earlier Student densities.  No global priority claim is
made for this application.  We explicitly record the overlap of a
related fractional Student calculation with the supplied earlier paper,
then give three separate attempts on named constants that remain unresolved.

\subsection{Separating modulus and phase in a rational primitive}

We use Baker's theorem that logarithms of nonzero algebraic numbers
which are linearly independent over $\mathbb Q$ are linearly independent
over $\overline{\mathbb Q}$; arbitrary determinations of the logarithms
are allowed \cite{p08:rd:BakerIII}.

\begin{lemma}[Real-logarithm and phase separation]\label{p08:rd:separation}
Suppose $b_j>0$ are algebraic, $\ell_k\in i\mathbb R$ are logarithms of
nonzero algebraic numbers, and $a_j,c_k$ are real algebraic numbers.
If
\[
 B+A=0,\qquad B=\sum_j a_j\log b_j,
 \qquad A=\frac1i\sum_k c_k\ell_k,
\]
then $B=A=0$.
\end{lemma}
\begin{proof}
Choose a $\mathbb Q$-basis from each of the two finite lists
$\{\log b_j\}$ and $\{\ell_k\}$, omitting zero elements.  The union of
these bases is $\mathbb Q$-linearly independent: its real and imaginary
parts must vanish separately.  All its members are logarithms of
algebraic numbers.  Baker's theorem makes the union
$\overline{\mathbb Q}$-linearly independent as well.  Expressing the
hypothesized relation in that union forces every coefficient to vanish;
in particular its two separate sums vanish.
\end{proof}

\begin{theorem}[Algebraic quantiles of even-power rational laws]
\label{p08:rd:even-power}
For an integer $m\ge2$, let $Q_m(q)$ be the $q$-quantile of the probability
density
\[
 f_m(x)=\frac{m\sin(\pi/(2m))}{\pi}\frac1{1+x^{2m}},
 \qquad x\in\mathbb R.
\]
For every algebraic $q\in(0,1)$,
\[
 Q_m(q)\in\overline{\mathbb Q}
 \quad\Longleftrightarrow\quad q=\frac12.
\]
The algebraic quantile in that case is $0$.  Thus every noncentral
algebraic-probability quantile is transcendental, including all noncentral
rational-probability quantiles.
\end{theorem}
\begin{proof}
Write $Z_m=\int_{\mathbb R}(1+t^{2m})^{-1}dt
=\pi\csc(\pi/(2m))/m$, by the substitution $u=t^{2m}$ and Euler's beta
integral.  Put $J_m(x)=\int_0^x(1+t^{2m})^{-1}dt$.  The quantile equation
is $J_m(x)=(q-1/2)Z_m$.

The simple poles $\zeta$ of the integrand satisfy $\zeta^{2m}=-1$.
Their residues are $-\zeta/(2m)$.  None lies on the real axis.  Define
$L_\zeta(x)$ as the continuous logarithm of
$(x-\zeta)/(-\zeta)$ along the real segment from $0$ to $x$, starting
with value $0$.  Partial fractions, integrated with this normalization,
give
\[
 J_m(x)=-\frac1{2m}\sum_{\zeta^{2m}=-1}\zeta L_\zeta(x).
\]
For an upper-half-plane pole $\zeta=c+is$, $s>0$, put
$D_c(x)=x^2-2cx+1>0$ and write
$L_\zeta(x)=\tfrac12\log D_c(x)+i\theta_\zeta(x)$.
The lower pole is its complex conjugate, with conjugate logarithm.
Their joint contribution is
\[
 -\frac{c}{2m}\log D_c(x)+\frac{s}{m}\theta_\zeta(x).
\]
Consequently $J_m(x)=B_m(x)+A_m(x)$, where $A_m$ is the sum of the
argument terms, and pairing the upper poles with real parts $c$ and $-c$
gives
\begin{equation}\label{p08:rd:positive-log}
 B_m(x)=\sum_{\substack{\zeta^{2m}=-1\\
                  \operatorname{Im}\zeta>0,\ \operatorname{Re}\zeta=c>0}}
 \frac{c}{2m}
 \log\frac{x^2+2cx+1}{x^2-2cx+1}.
\end{equation}
An upper pole of real part zero contributes nothing to this sum.  Since
$m\ge2$, at least one pole in the displayed sum exists.  Every summand
is positive for $x>0$ and negative for $x<0$.  Therefore
$B_m(x)=0$ if and only if $x=0$.

Now suppose that $x$ and $q$ are algebraic.  All the positive numbers
$D_c(x)$ are algebraic.  Furthermore
\[
 \exp(i\theta_\zeta(x))=
 \frac{(x-\zeta)/(-\zeta)}{\sqrt{D_c(x)}}
\]
is algebraic and has modulus one, so $i\theta_\zeta(x)$ is a purely
imaginary logarithm of an algebraic number.  Since
$Z_m/\pi$ is real algebraic and $i\pi$ is a logarithm of $-1$, the
quantile equation is a relation of exactly the type in
Lemma~\ref{p08:rd:separation}: its real-logarithm part is $B_m(x)$, while
$A_m(x)-(q-1/2)Z_m$ is an algebraic linear combination of purely
imaginary logarithms divided by $i$.  The lemma forces $B_m(x)=0$.
Thus $x=0$ and $q=1/2$.  The converse follows from symmetry.
\end{proof}

For $m=1$ the modulus part disappears: this is the Cauchy law, whose
rational-probability quantiles are algebraic.  The strict sign in
\eqref{p08:rd:positive-log} is the additional obstruction when $m\ge2$.
In particular the quartic density $\sqrt2/(\pi(1+x^4))$ has
transcendental lower and upper quartiles.  The theorem is a positive,
uniform classification, not an inference from numerical approximations.

\subsection{A recovered Student classification, not an additional discovery}
\label{p08:rd:Student}
The same investigation recovers the following result already contained
in the supplied arc-length sequel \cite{p08:rd:PriorArc}.  For
$\nu_0\in\{1/3,1/2,2/3\}$ and algebraic $q\in(0,1)$, a Student quantile
with $\nu_0$ degrees of freedom is algebraic exactly when $q$ is rational;
with $\nu_0+2n$ degrees of freedom, $n\ge1$, it is algebraic exactly
when $q=1/2$.  We do not count either assertion as new.

Here is the short mechanism, to identify the overlap precisely.
Set $a=(\nu+1)/2$, $I_a(s)=\int_0^s(1+t^2)^{-a}dt$ and
$L_a=I_a(+\infty)$.  The three starting exponents $2/3,3/4,5/6$ are
first-kind elliptic integrals, by the respective substitutions
\[
 (X,Y)=((1+s^2)^{1/3},s),\quad
 s=\frac{2t}{1-t^2},\quad
 (u,v)=((1+s^2)^{-1/3},s/\sqrt{1+s^2}),
\]
on $Y^2=X^3-1$, $y^2=1-t^4$, and $v^2=1-u^3$.
Real algebraic fractions of these elliptic integrals are rational,
by Schneider's theorem \cite{p08:rd:Schneider,p08:rd:BertolinWaldschmidt}.
For every increment one has
\[
 I_{a+1}(s)=\frac{s}{2a(1+s^2)^a}+\frac{2a-1}{2a}I_a(s).
\]
Iteration gives $I_{a+n}=C_nI_a+D_n$, $L_{a+n}=C_nL_a$, where
$C_n>0$ is rational and $D_n(s)$ is nonzero algebraic of the sign of $s$
at nonzero algebraic $s$.  An algebraic-probability quantile would make
$I_a(s)-(2q-1)L_a=-D_n(s)/C_n$.  A linear form over the algebraic numbers
in elliptic logarithms cannot be nonzero algebraic, by the analytic
subgroup theorem \cite{p08:rd:Wustholz}; hence $s=0$.  This is precisely the
elliptic-division and exact-differential mechanism of the earlier sequel's
elliptic-arc theorem and algebraic-parameter beta theorem (its source
labels \texttt{dd-elliptic-alg} and \texttt{bs-algebraic-ratios}).

\subsection{Three named constants: different transfers and their exact stopping points}
\label{p08:rd:named}

The following attempts do not establish irrationality of Khinchin's
$K$, the period-doubling Feigenbaum $\delta$, or the Glaisher--Kinkelin
$A$.  They test whether the finite-dimensional arithmetic structures just
used are actually present in natural representations of those constants.

\paragraph{Khinchin's constant: an algebraic-integrand test.}
The Gauss product formula \cite{p07:ric:BBC} gives
\[
 (\log2)\log K=
 \sum_{k\ge1}\log k\log\frac{(k+1)^2}{k(k+2)}.
\]
One can turn this into a positive Stieltjes integral, but its density
does not become an algebraic differential.

\begin{proposition}\label{p08:rd:K-Stieltjes}
The meromorphic function
\[
 H(z)=\sum_{k=2}^\infty\frac{\log(k/(k-1))}{k+z}
\]
satisfies
$\int_0^1H(t)\,dt=(\log2)\log K$ and is not algebraic over
$\mathbb C(z)$.
\end{proposition}
\begin{proof}
The series converges locally uniformly away from the negative integers
$-2,-3,\ldots$, since its tail is $O(k^{-2})$ on every compact set.
It has a simple pole at each $-k$, with nonzero residue
$\log(k/(k-1))$.  On $[0,1]$ summation by parts yields
\[
 H(t)=\sum_{k\ge1}\log k
          \left(\frac1{k+t}-\frac1{k+1+t}\right).
\]
The discarded boundary term is $O((\log k)/k)$.
Integration, or Tonelli's theorem on the nonnegative summands, gives the
claimed formula.  An algebraic function over $\mathbb C(z)$ has only
finitely many finite singularities: these lie among the zeros of the
leading coefficient and discriminant of a defining polynomial.
The infinitely many poles of $H$ exclude algebraicity.
\end{proof}

This rules out interpreting this particular Stieltjes density as a
first-kind algebraic differential in the same parameter.  It does not
rule out some other, presently unavailable representation of its definite
integral.  The unresolved step is a finite algebraic-group or
height-controlled identity for the value, not recovery of $H$.

\paragraph{Feigenbaum's constant: bounded heights, but no degree control.}
Let $r_n$ be the principal superstable logistic parameter of period $2^n$,
and put
\[
 d_n=\frac{r_n-r_{n-1}}{r_{n+1}-r_n},\qquad
 \delta=\lim_{n\to\infty}d_n.
\]
The existence and numerical evaluation of this scaling limit are distinct
from an arithmetic classification \cite{p07:ric:Molteni}.
Here $h$ denotes the absolute logarithmic Weil height.

\begin{proposition}\label{p08:rd:delta-height}
The approximants $d_n$ are algebraic and satisfy $h(d_n)\le14\log2$.
The construction supplies the degree bound
\[
 [\mathbb Q(d_n):\mathbb Q]\le2^{7\cdot2^{n-1}}
 \qquad(n\ge1).
\]
If a sequence of distinct algebraic approximants $\alpha_j\to\xi$ has
bounded heights and
\begin{equation}\label{p08:rd:degree-criterion}
 \frac{-\log|\xi-\alpha_j|}
 {[\mathbb Q(\alpha_j):\mathbb Q]}\longrightarrow+\infty,
\end{equation}
then $\xi$ is transcendental.  The bounds above do not verify
\eqref{p08:rd:degree-criterion} for $\delta$.
\end{proposition}
\begin{proof}
The logistic map is affinely conjugate to $z\mapsto z^2+c_n$, where
$c_n=r_n(2-r_n)/4$ and the critical point $0$ has period $2^n$.
The polynomial $f_c^{\circ m}(0)$ is monic of degree $2^{m-1}$ in $c$.
Hence $c_n$ is an algebraic integer of degree at most $2^{2^n-1}$.
Every conjugate of $c_n$ has modulus at most $2$: if $|c|>2$, the orbit
of $0$ escapes, since after its first step successive moduli grow by at
least the factor $|c|-1>1$.  Therefore $h(c_n)\le\log2$.

The relation $r_n=1\pm\sqrt{1-4c_n}$ and the standard height inequalities
$h(ab)\le h(a)+h(b)$, $h(a+b)\le h(a)+h(b)+\log2$, and
$h(a^m)=|m|h(a)$ give $h(r_n)\le\log8$.
Applying the same inequalities to the two differences and their quotient
gives $h(d_n)\le4\log8+2\log2=14\log2$.
Also $[\mathbb Q(r_n):\mathbb Q]\le2^{2^n}$, so the degree of the field
generated by $r_{n-1},r_n,r_{n+1}$ is at most the product of these three
bounds, which is the stated bound for $d_n$.

Finally, if $\xi$ were algebraic of degree $D$ and $h(\alpha_j)\le C$,
the product formula applied to the nonzero algebraic number
$\xi-\alpha_j$ would give
\[
 |\xi-\alpha_j|\ge
 \exp\left(-D[\mathbb Q(\alpha_j):\mathbb Q]
                 (h(\xi)+C+\log2)\right).
\]
Here we used $|z|\ge\exp(-\deg(z)h(z))$ for nonzero algebraic $z$,
which follows by applying the height definition to $z^{-1}$.
This contradicts \eqref{p08:rd:degree-criterion}.
\end{proof}

Thus even these uniformly bounded heights are not the missing arithmetic
estimate by themselves.  One would need approximation relative to the
actual degrees, as in \eqref{p08:rd:degree-criterion}, or a different
arithmetic identity.  A geometric convergence estimate and the available
large degree upper bound do not establish that condition.  No lower bound
for the actual degrees, and no impossibility theorem for a sharper
approach, is being asserted.

\paragraph{Glaisher--Kinkelin's constant: the prime-product transfer.}
Use $\log A=1/12-\zeta'(-1)$, a normalization also used in the recent
automatic-product paper \cite{p08:rd:Campbell}.  Differentiating the zeta
functional equation \cite[25.4.2]{p08:rd:DLMF} at $-1$ gives the exact
reduction
\begin{equation}\label{p08:rd:A-prime}
 12\log A=\log(2\pi)+\gamma-
                       \frac{\zeta'(2)}{\zeta(2)}
 =\log(2\pi)+\gamma+
                       \sum_p\frac{\log p}{p^2-1}.
\end{equation}
Indeed the logarithmic derivative at $-1$ is
$\log(2\pi)-1+\gamma-\zeta'(2)/\zeta(2)$ and
$\zeta(-1)=-1/12$; the absolutely convergent prime sum is the
logarithmic derivative of the Euler product at $2$.

This is a concrete link to prime observation, but not a finite logarithmic
certificate.  To make the limitation quantitative, put
\[
 S_y=\sum_{p\le y}\frac{\log p}{p^2-1},\qquad
 D_y=\operatorname{lcm}_{p\le y}(p^2-1),\qquad
 B_y=\prod_{p\le y}p^{D_y/(p^2-1)}.
\]
Then $e^{S_y}=B_y^{1/D_y}$ is algebraic, has degree at most $D_y$, and
has height exactly $S_y$, since $B_y$ is a positive integer and heights
scale under powers.  These heights remain bounded.  Nevertheless for
$y\ge2$ the immediate effective estimate is only
\[
 0\le \sum_p\frac{\log p}{p^2-1}-S_y
 \le\frac43\frac{\log y+1}{y},
\]
by comparison with $\sum_{m>y}(\log m)/m^2$ (take integer $y$).
The denominator-clearing degrees are not controlled by a fixed number
field.  Moreover \eqref{p08:rd:A-prime} still contains $\gamma$ and $\pi$.
Neither this product identity nor the presence of automatic sequences in
other product evaluations proves that $A$ is algebraic, irrational, or
transcendental.  The missing step is an arithmetic relation or
independence theorem for the actual combined value in
\eqref{p08:rd:A-prime}; the elliptic quantile theorem supplies none.


\section{Complete relation ideals for cyclotomic Mahler sums}
\label{p08:ctm:section}

The earlier linear classification of sparse Lambert values admits a
stronger, finite form: in the additive cyclotomic setting, every
polynomial relation between algebraic-point values follows from a
computable affine linear relation.  This statement also survives mixing
multiplicatively independent dilations.  Its transcendence inputs are
classical Mahler theory and the purity theorem of Adamczewski--Faverjon;
the finite certificate and the applications below are proved explicitly.
These results do not establish transcendence of the Erd\H{o}s--Borwein
constant or of the full reciprocal Fibonacci sum, or irrationality of
the prime Mertens constant.  The last subsection isolates the remaining
obstructions through different decompositions of those three constants.

\subsection{From algebraic dependence to an additive coboundary}

We write \(\overline{\mathbb Q}\) for the algebraic numbers and
\(\mu_N=\{z:z^N=1\}\).  A rational additive coboundary for the dilation
\(z\mapsto z^d\) is a function \(g(z)-g(z^d)\), with \(g\in\mathbb C(z)\).
The following elementary argument is the additive part of the classical
Mahler algebraic-independence mechanism; compare Kubota's theorem and
the applications of Nishioka's theory in \cite{p08:ctm:Kubota,p08:ctm:BV15}.

\begin{lemma}[The additive dependence criterion]\label{p08:ctm:additive}
Let \(d\geq2\), let \(R_1,\ldots,R_s\in\mathbb C(z)\) be regular at
zero, and let analytic germs \(f_i\) satisfy
\[
 f_i(z^d)=f_i(z)-R_i(z).
\]
The functions \(f_1,\ldots,f_s\) are algebraically dependent over
\(\mathbb C(z)\) if and only if there are constants
\(c_1,\ldots,c_s\), not all zero, and \(g\in\mathbb C(z)\) such that
\[
 \sum_i c_iR_i(z)=g(z)-g(z^d).
\]
\end{lemma}

\begin{proof}
Suppose there is a polynomial relation.  In its nonzero ideal in
\(\mathbb C(z)[Y_1,\ldots,Y_s]\), choose a polynomial \(P\) whose
leading monomial is least in graded lexicographic order, and normalize
the leading coefficient to one.  The polynomial
\[
 \mathcal T P(z,Y)=P(z^d,Y_1-R_1(z),\ldots,Y_s-R_s(z))
\]
is another relation with the same leading monomial and coefficient.
Consequently \(\mathcal T P-P=0\); otherwise its leading monomial
would be smaller.  Write \(m=\deg_Y P\), necessarily positive, and
let \(P_j\) be its homogeneous part of degree \(j\).  The degree-
\(m\) identity says that every coefficient of \(P_m\) is fixed by
\(z\mapsto z^d\).  A rational function fixed by this substitution is
constant: comparison of the first nonconstant Laurent coefficient at
zero proves this directly.  Thus \(P_m\in\mathbb C[Y]\).

In degree \(m-1\), the same identity gives
\[
 \sum_i R_i(z)\frac{\partial P_m}{\partial Y_i}(Y)
       =P_{m-1}(z^d,Y)-P_{m-1}(z,Y).
\]
At least one coefficient vector of the gradient on the left is nonzero,
since the characteristic is zero and \(m>0\).  Comparing that monomial
coefficient gives the required constants and a rational coboundary
(with the sign of its rational function reversed).

Conversely, a rational \(g\) in the displayed coboundary cannot have
a pole at zero: a pole of order \(a>0\) would give a pole of order
\(da\) in \(g(z^d)\), which neither other term could cancel.  Thus
\(\sum_i c_if_i-g\) is analytic at zero and invariant under
\(z\mapsto z^d\).  Its first nonconstant Taylor coefficient would
contradict invariance.  It is constant, giving a linear relation over
\(\mathbb C(z)\).
\end{proof}

For clarity, the arithmetic transfer used below is a stated external
input.  For a one-variable linear Mahler system over
\(\overline{\mathbb Q}(z)\), analytic at zero, the transcendence degree
of its values at a regular algebraic \(0<|\alpha|<1\) equals the
functional transcendence degree.  This is Nishioka's theorem, also
\cite[Corollary 2.5 and equation (2.4), preprint numbering]{p08:ctm:AF}.
For several systems whose integer dilations are pairwise
multiplicatively independent, the ideal of all value relations is
generated by the ideals within the individual systems
\cite[Theorem 2.8, preprint numbering]{p08:ctm:AF}.  The older
regular-singular version already suffices here
\cite[Theorem 2.4]{p08:ctm:AF18}.  No unproved generalization of
Loxton--van der Poorten's early multivariate claims is used.

\subsection{One finite linear system computes the entire ideal}

\begin{theorem}[Complete cyclotomic relation certificate]\label{p08:ctm:ideal}
Let \(K\) be a number field and \(R_1,\ldots,R_s\in K(z)\) satisfy
\(R_i(0)=0\).  Suppose their finite poles lie in \(\mu_N\), have
orders at most \(m\), and their polynomial parts have degree at most
\(B\), with a constant or zero part assigned degree zero.  Put
\[
 D(z)=(1-z^N)^m,\qquad M=mN+\lfloor B/d\rfloor,\qquad
 f_i(z)=\sum_{k\geq0}R_i(z^{d^k}),\qquad d\geq2.
\]
If all kernels are polynomials, take \(m=0,N=1\).  Solve the finite
homogeneous linear system over \(K\)
\begin{equation}\label{p08:ctm:matrix}
 A(z)D(z^d)-A(z^d)D(z)
   =D(z)D(z^d)\sum_{i=1}^s c_iR_i(z),\qquad
 A(0)=0,\quad \deg A\leq M,
\end{equation}
in \(c\in K^s\) and \(A\in K[z]\).  Let \(V\subset K^s\) be the
projection of its solution space to \(c\), let \(r=\dim_K V\),
and choose a basis \(c^{(1)},\ldots,c^{(r)}\), with corresponding
polynomials \(A_1,\ldots,A_r\).

For every algebraic \(0<|\alpha|<1\), the entire ideal of polynomial
relations over \(\overline{\mathbb Q}\) between
\(f_1(\alpha),\ldots,f_s(\alpha)\) is
\begin{equation}\label{p08:ctm:idealformula}
 \left(\ \sum_{i=1}^s c_i^{(j)}X_i-\frac{A_j(\alpha)}{D(\alpha)}
                    \ :\ 1\leq j\leq r\ \right).
\end{equation}
In particular its transcendence degree is exactly \(s-r\).  The
normalized \(A\) associated with any \(c\in V\) is unique.
\end{theorem}

\begin{proof}
The defining sums converge locally uniformly in the disk and satisfy
\(f_i(z)-f_i(z^d)=R_i(z)\).  We first verify that
\eqref{p08:ctm:matrix} captures every rational coboundary, including those
initially having complex coefficients.  Normalize its rational
solution \(g\) by \(g(0)=0\); regularity at zero follows from the
last paragraph of Lemma~\ref{p08:ctm:additive}.

A pole of \(g\) outside \(\mu_N\) forces poles at every successive
\(d\)-th root of that point: these roots remain outside \(\mu_N\),
where the right side of the coboundary equation is regular.  Their
number is unbounded, impossible for a rational function.  All finite
poles therefore lie in \(\mu_N\).  If a pole had maximal order
larger than \(m\), the same backward propagation, now of that pole
order, would again give infinitely many poles.  Thus \(Dg\) is
polynomial.  If the polynomial part of \(g\) has positive degree
\(b\), that of \(g-g(z^d)\) has degree exactly \(db\), so
\(b\leq\lfloor B/d\rfloor\).  It follows that \(g=A/D\) with
\(\deg A\leq M\).  This gives exactly \eqref{p08:ctm:matrix}.
Conversely every solution of that identity gives a rational
coboundary.  A normalized solution is unique because a rational
function invariant under \(z\mapsto z^d\) is constant.

The coefficient equations are a finite linear system over \(K\).
Its complex solution space is the scalar extension of its solution
space over \(K\).  In particular the space of all complex coefficient
vectors giving rational coboundaries is \(V\otimes_K\mathbb C\).
Extend the chosen basis of \(V\) to the rows of an invertible matrix
\(S\in\operatorname{GL}_s(K)\), and put \(h=Sf\).  Its first \(r\)
coordinates are the rational functions \(A_j/D\), by normalized
analytic uniqueness.  The remaining \(s-r\) coordinates are
algebraically independent over \(\mathbb C(z)\).  Indeed, a relation
among them would, by Lemma~\ref{p08:ctm:additive}, produce a nonzero
coboundary vector in the span of the complementary rows of \(S\),
contradicting their direct-sum decomposition with \(V\).

Apply the stated one-variable value theorem to these remaining
coordinates together with the constant function \(1\).  Their
system matrix has identity diagonal and the rational kernels in its
first column.  It and its inverse have poles only on the unit circle,
and determinant one.  The forward orbit of \(\alpha\) is entirely
inside the punctured disk, so \(\alpha\) is regular; the
one-variable pair is admissible.  Thus these \(s-r\) values remain
algebraically independent.  In the coordinates \(Y=SX\), all
relations are therefore generated by
\(Y_j-A_j(\alpha)/D(\alpha)\), \(1\leq j\leq r\): division by
these linear polynomials leaves a polynomial in algebraically
independent coordinates.  Changing back to \(X\) proves
\eqref{p08:ctm:idealformula}.
\end{proof}

The theorem makes the arithmetic output exact, rather than merely
giving a lower bound for the number of independent values.  Its
additive-dependence principle and its value theorem are classical;
general rational Mahler solving and discrete residues are also
established subjects \cite{p07:rsc:CDDM,p07:rsc:AZ}.  The specific contribution
here is the explicit bounded matrix that computes the complete
relation ideal for this cyclotomic input class.  This is not a claim
to a new general Mahler transcendence theorem.

\begin{corollary}[Arbitrary finite collections of integer dilations]
\label{p08:ctm:mixed}
For each \(i\), let \(d_i\geq2\) and let \(R_i\) be as in
Theorem~\ref{p08:ctm:ideal}, allowing different cyclotomic pole sets and
bounds.  At a common algebraic \(0<|\alpha|<1\), all algebraic
relations between
\[
 f_i(\alpha)=\sum_{k\geq0}R_i(\alpha^{d_i^k})
\]
are generated by affine linear relations, computable as follows.
Partition the dilations into classes of multiplicative dependence.
For each class choose a common power \(D_0=d_i^{e_i}\), replace
\(R_i\) by \(\sum_{k=0}^{e_i-1}R_i(z^{d_i^k})\), and apply
\eqref{p08:ctm:matrix} with dilation \(D_0\) within that class.
The union of the resulting class relations generates the full ideal;
the transcendence degrees of the classes add.
\end{corollary}

\begin{proof}
Multiplicative dependence of positive integers greater than one is
an equivalence relation: their prime-exponent vectors are rationally
proportional.  Each finite class therefore admits the stated common
power, effectively obtainable from those integer vectors.  Iterating
the functional equation gives precisely the replacement kernel.
Its poles are still roots of unity, with finite explicitly computable
orders, and it vanishes at zero.  Theorem~\ref{p08:ctm:ideal} computes
the ideal for each class.  Common powers for distinct classes are
pairwise multiplicatively independent.  The purity theorem stated
above now gives exactly the sum of these extended ideals.

All its regularity hypotheses hold as in the preceding proof.  Even
if one uses only the older regular-singular purity theorem, our systems
qualify: for a vector \(f\) the analytic triangular matrix
\(Y(z)=\left(\begin{smallmatrix}1&0\\f(z)&I\end{smallmatrix}\right)\)
satisfies \(Y(z^{D_0})=A(z)Y(z)\) and has determinant one, providing
an analytic gauge to the identity system.  Its coefficients are
algebraic.  Finally purity is a statement about the entire ideal,
so no additional nonlinear mixed relations remain.
\end{proof}

\subsection{Residues, Fibonacci contrasts, and infinite independent families}

\begin{theorem}[The complete residue-family ideal]\label{p08:ctm:residue}
For \(d\geq2\), define
\[
 U_{d,j}(z)=\sum_{k\geq0}\frac{z^{j d^k}}{1-z^{d^{k+1}}},
       \qquad 1\leq j\leq d-1.
\]
At every algebraic \(0<|\alpha|<1\), their relation ideal is
\[
 \left(X_1+\cdots+X_{d-1}-\frac{\alpha}{1-\alpha}\right),
\]
and their transcendence degree is \(d-2\).  For pairwise
multiplicatively independent \(d_1,\ldots,d_t\), the union of these
families, evaluated at arbitrary respective algebraic disk points
\(\alpha_1,\ldots,\alpha_t\), has only these \(t\) relations and
transcendence degree \(\sum_{i=1}^t(d_i-2)\).
\end{theorem}

\begin{proof}
For a coefficient vector \(c\), its kernel is
\(R(z)=(\sum_{j=1}^{d-1}c_jz^j)/(1-z^d)\).  The pole argument in
Theorem~\ref{p08:ctm:ideal} shows that a normalized rational coboundary
solution has only simple poles in \(\mu_d\) and is bounded at
infinity.  A pole at \(\zeta\in\mu_d\setminus\{1\}\) is impossible:
every \(d\)-th root of \(\zeta\) is outside \(\mu_d\), where
\(g(z)\) and \(R(z)\) are regular but \(g(z^d)\) would have a
pole.  Hence \(g(z)=c z/(1-z)\).  Since
\[
 \frac{z}{1-z}-\frac{z^d}{1-z^d}
       =\frac{z+\cdots+z^{d-1}}{1-z^d},
\]
the coboundary space consists exactly of constant vectors.  The
complete-ideal theorem proves the first assertion.  Purity proves
the second, with the same regularity verification even when the
evaluation points differ.  The case \(d=2\) is interpreted literally:
the single value is algebraic and contributes zero to the dimension.
\end{proof}

For example, \(U_{3,1}(1/2),U_{4,1}(1/2),U_{4,2}(1/2)\) are
algebraically independent.  Bases \(3\) and \(4\) are
multiplicatively independent although the latter is a power of \(2\).
By contrast one cannot invoke purity between bases \(2\) and \(4\);
their relations must first be computed at a common power.

\begin{theorem}[All periodic Fibonacci residue relations]
\label{p08:ctm:fibonacci}
Let \(F_1=F_2=1\), and for \(h\geq1\) put
\[
 V_{h,r}=\sum_{\substack{k\geq0\\k\equiv r\pmod h}}
                 \frac1{F_{2^k}},\qquad 0\leq r<h.
\]
The entire algebraic relation ideal of this \(h\)-tuple over
\(\overline{\mathbb Q}\) is
\[
 \left(\sum_{r=0}^{h-1}X_r-\frac{7-\sqrt5}{2}\right).
\]
In particular its transcendence degree is \(h-1\).  Moreover the
countable family
\begin{equation}\label{p08:ctm:haar}
 W_{m,r}=V_{2^{m+1},r}-V_{2^{m+1},r+2^m},
       \qquad m\geq0,\quad0\leq r<2^m,
\end{equation}
is algebraically independent over \(\overline{\mathbb Q}\).
\end{theorem}

\begin{proof}
For a vector \(a=(a_0,\ldots,a_{h-1})\), extend its entries
periodically and define
\[
 H_a(z)=\sum_{k\geq0}a_k\frac{z^{2^k}}{1-z^{2^{k+1}}}.
\]
For dilation \(d=2^h\), its kernel is
\[
 \sum_{r=0}^{h-1}\frac{a_rz^{2^r}}{1-z^{2^{r+1}}}
    =\frac{\sum_{j=1}^{d-1}a_{v_2(j)}z^j}{1-z^d}.
\]
Each index \(j\) has the unique form \(2^r(2s+1)\), proving
the equality.  The preceding proof shows that this kernel is a
rational coboundary precisely when all \(a_r\) are equal.
Theorem~\ref{p08:ctm:ideal} applied to the \(h\) coordinate kernels
therefore gives precisely one affine linear relation among their
values and transcendence degree \(h-1\).

Let \(q=(\sqrt5-1)/2\).  For every periodic \(a\), Binet's formula
gives, with \(b_k=a_{k+1}\),
\[
 \sum_{k\geq0}\frac{a_k}{F_{2^k}}
       =a_0+\sqrt5\,H_b(q^2).
\]
This is an invertible affine transformation of the \(h\) coordinate
values: the shift of the residues is a permutation.  The only
relation has constant coefficients \(a_r=1\).  In that case
\(H_b(z)=z/(1-z)\), whence the sum is
\(1+\sqrt5 q=(7-\sqrt5)/2\).  This proves the full ideal assertion.

For the second assertion, fix \(M\geq1\).  Refinement gives
\(V_{h,r}=V_{2h,r}+V_{2h,r+h}\).  At each node this sum, together
with its difference in \eqref{p08:ctm:haar}, recovers the two children
by division by two.  Thus the \(2^M\) values
\(V_{2^M,r}\) are an invertible rational linear transform of their
total sum and the \(2^M-1\) differences with \(m<M\).  Their
transcendence degree is \(2^M-1\), and the total sum is algebraic.
Consequently those \(2^M-1\) differences are algebraically
independent.  Every finite subfamily of \eqref{p08:ctm:haar} lies in
one such level, completing the proof.
\end{proof}

The total-sum identity is classical, and substantial algebraic
independence theorems for sparse Fibonacci/Lucas sums, including
periodic weights, predate this work \cite{p08:ctm:BV14,p08:ctm:BV15,p08:ctm:Nguyen}.
We do not assert priority for sparse Fibonacci transcendence.
The advance over the previous project result is the full polynomial
relation ideal, and the explicit compatible basis of independent
contrasts at all dyadic refinement levels.  These are consequences
of the complete additive certificate, not separate historical
conjecture closures.

\subsection{Three named constants: transfers and the remaining gaps}

\paragraph{Erd\H{o}s--Borwein: independent blocks do not control their sum.}
For \(b\geq2\) and a prime \(p\), let
\[
 E_b=\sum_{n\geq1}\frac1{b^n-1},\qquad
 A_{p,m}(b)=\sum_{k\geq0}\frac1{b^{m p^k}-1}\quad(p\nmid m).
\]
Then
\begin{equation}\label{p08:ctm:EBblocks}
 E_b=\sum_{\substack{m\geq1\\p\nmid m}}A_{p,m}(b),
\end{equation}
and the countable family on the right is algebraically independent.
This independence is already a special case of
\cite[Theorem 1]{p08:ctm:BV15}; it is not a new result.  A short
certificate proof helps explain its relation to our framework.
Write \(G_p(z)=\sum_{k\geq0}z^{p^k}/(1-z^{p^k})\).
For a finite set of indices \(m\) prime to \(p\), the kernels of
\(G_p(z^m)\) are \(z^m/(1-z^m)\).  If a combination is a
rational coboundary, the normalized rational solution has simple
poles in \(\mu_N\), where \(N\) is the least common multiple of
the indices, and is bounded at infinity.  Since \(p\nmid N\),
each point of \(\mu_N\) has \(p-1\) of its \(p\)-th roots
outside \(\mu_N\).  Such a pole would be impossible by backward
propagation.  The solution is therefore constant.  The kernels are
linearly independent: a primitive root of the largest remaining
index is a pole of its kernel alone, and induction removes all
indices (with \(m=1\) treated last).  Theorem~\ref{p08:ctm:ideal}
proves their algebraic independence at \(z=1/b\).
The unique factorization \(n=m p^k\), \(p\nmid m\), proves
\eqref{p08:ctm:EBblocks} by positive rearrangement.

There is a genuine logical obstruction to drawing a conclusion
about \(E_b\) from this independence alone.  There exist positive,
algebraically independent real numbers \(t_1,t_2,\ldots\) whose
sum is \(1\).  To prove this, put \(r_0=1\), and recursively choose
\(t_n\in(r_{n-1}/3,2r_{n-1}/3)\) transcendental over
\(\overline{\mathbb Q}(t_1,\ldots,t_{n-1})\), then set
\(r_n=r_{n-1}-t_n\).  Such a choice exists because the algebraic
closure of this finitely generated countable field is countable.
Independence follows inductively and \(0<r_n<(2/3)^n\), proving
the sum assertion.  Thus finite-block Mahler independence cannot
by itself prove transcendence of \(E_2\).  An arithmetic estimate
controlling the infinite tail relative to algebraic approximation
is still missing.  Duverney--Tachiya explicitly distinguish this
transcendence question from known irrationality; Campbell's 2026
binary-pattern result does not close it \cite{p08:ctm:DT,p07:rsc:Campbell}.

\paragraph{The full reciprocal Fibonacci sum: a theta--Lambert split.}
Set \(L(z)=\sum_{n\geq1}z^n/(1-z^n)\), \(q=(\sqrt5-1)/2\), and
\(\theta_2(Q)=\sum_{n\in\mathbb Z}Q^{(n+1/2)^2}\) for real
\(0<Q<1\).  Then the exact classical decomposition is
\begin{equation}\label{p08:ctm:thetaFib}
 \sum_{n\geq1}\frac1{F_n}
   =\sqrt5\left(L(q^2)-L(q^4)
                       +\frac14\theta_2(q^2)^2\right).
\end{equation}
Indeed, Binet's formula gives \(1/F_n=\sqrt5 q^n/
(1-(-1)^nq^{2n})\).  The even-index contribution divided by
\(\sqrt5\) is \(L(q^2)-L(q^4)\), using
\(t/(1-t)-t^2/(1-t^2)=t/(1-t^2)\).  For the odd indices,
Jacobi's two-square identity yields
\[
 \frac14\theta_2(q^2)^2
       =\sum_{\substack{m\geq1\\m\ \mathrm{odd}}}
                                      \frac{q^m}{1+q^{2m}}.
\]
For completeness, the change of variables
\((a,b)=(u+v+1,u-v)\) is a bijection from \(\mathbb Z^2\)
onto the pairs of opposite parity, and
\(a^2+b^2=2(u+1/2)^2+2(v+1/2)^2\).  Hence the theta square is
\(\sum_{n\ \mathrm{odd}}r_2(n)q^n\).  Jacobi's formula
\(r_2(n)=4\sum_{d\mid n}\chi_4(d)\), where
\(\chi_4(2j)=0\) and \(\chi_4(2j+1)=(-1)^j\), turns it into
\(4\sum_{m\ \mathrm{odd}}\sum_{j\geq0}(-1)^j q^{m(2j+1)}\),
proving the displayed identity by absolute convergence.  See
\cite{p08:ctm:Hirschhorn} for the classical two-square formula.

The odd-index sum is already known to be transcendental
\cite[slide 58]{p08:ctm:Waldschmidt}.  This does not exclude its
algebraic cancellation with the Lambert difference in
\eqref{p08:ctm:thetaFib}.  Precisely, proving that
\[
 L(q^2)-L(q^4)+\tfrac14\theta_2(q^2)^2
                         \notin\overline{\mathbb Q}
\]
would close the original target; the finite additive certificate
does not control this mixed theta--Lambert relation.  The full-sum
transcendence question is explicitly left unresolved in
\cite[slide 59]{p08:ctm:Waldschmidt}.  Thus this route identifies the
missing mixed relation instead of reusing the earlier no-Mahler
argument.

\paragraph{Prime Mertens: a zeta-value representation with an exponential tail.}
Let \(\mu\) be the M\"obius function and \(\gamma\) Euler's constant.
The prime Mertens constant has the absolutely convergent forms
\begin{align}
 B_1-\gamma
   &=\sum_{m\geq2}\frac{\mu(m)}m\log\zeta(m),
                                      \label{p08:ctm:Mertens}\\
 B_1
   &=\gamma-\frac12\log\zeta(2)
     +\sum_{\substack{n\geq3\\n\ \mathrm{odd}}}
       \frac{\mu(n)}n\log\frac{\zeta(n)}{\sqrt{\zeta(2n)}}.
                                      \label{p08:ctm:Mertenspair}
\end{align}
The tail in \eqref{p08:ctm:Mertens} after \(m=M\geq1\) has absolute
value at most \(3\cdot2^{-M}/(M+1)\).  These representations belong
to the classical computational theory of prime constants; compare
\cite{p07:rsc:LZ}.  To verify them directly, write
\(P(s)=\sum_p p^{-s}\).  Euler's product gives
\(\log\zeta(m)=\sum_{k\geq1}P(mk)/k\).  Absolute convergence
allows rearrangement, and the coefficient of \(P(n)\) in the
first series is
\[
 \frac1n\sum_{\substack{m\mid n\\m\geq2}}\mu(m)=-\frac1n
       \quad(n\geq2).
\]
It is therefore
\(\sum_p(\log(1-1/p)+1/p)=B_1-\gamma\).
Pairing the terms \(n\) and \(2n\) for odd \(n\), using
\(\mu(2n)=-\mu(n)\) and \(\mu(4n)=0\), proves
\eqref{p08:ctm:Mertenspair}.  Finally
\[
 0<\log\zeta(m)<\zeta(m)-1
 \leq 2^{-m}+\int_2^\infty x^{-m}\,dx\leq3\cdot2^{-m}
       \quad(m\geq2)
\]
implies the asserted tail bound by summing a geometric series.

Although \(\zeta(2n)\) is an algebraic multiple of \(\pi^{2n}\),
neither the arithmetic nature of the required odd zeta values nor
their infinite interaction with \(\gamma\) is controlled here.
Baker's theorem for a finite combination of logarithms of algebraic
numbers has no such hypotheses.  No verified irrationality proof
for \(B_1\) was located in the primary-source check.  The faster
tail in this representation is consequently a computational gain,
not a proof that \(B_1\) is irrational.


\section{Growing Laplace kernels and actual arithmetic cancellation}
\label{p08:rlap:section}

The named targets here are Euler's constant $\gamma$ and the
Euler--Gompertz constant
$\delta=\int_0^\infty e^{-t}/(1+t)\,dt$.
The published 2025 work of Wolfs--Van Assche improves rational
approximations to both constants without proving their irrationality
\cite{p08:rlap:WV}.  The conditional application in
Murabayashi--Yoshida \cite{p08:rlap:MY} also retains an unproved input.
No verified individual resolution was found in our current primary-source
check, and none is claimed below.  Known transcendence of
$\operatorname{Ein}(1)=\gamma+\delta/e$ is a different statement.

We now leave the fixed-denominator integer-polynomial cone of the
preceding investigation.  We examine growing pole orders, primitive
integer forms, and integer-valued rather than integer-coefficient
polynomials.  The main new internal results are an obstruction for
\emph{every} odd Taylor-order schedule using the actual reduced
denominator, and a modular description of every prime occurring in the
content of the classical Gompertz Pad\'e forms.  The Pad\'e family itself
is classical; no priority claim is made for the modular formulation.
The proofs distinguish a proved exclusion from the higher-valuation
estimate still missing for a positive irrationality argument.

\subsection{Euler's constant: all odd Taylor schedules fail after primitive clearance}

For integers $n\ge1$ and $M\ge1$, put
\[
 S_M(n)=\sum_{k=1}^M\frac{(-1)^{k-1}n^k}{k\,k!}
       =\frac{a_M(n)}{b_M(n)},\qquad
 \gcd(a_M(n),b_M(n))=1,\quad b_M(n)>0.
\]
The denominator here is the actual reduced denominator, not a chosen
factorial or least common multiple.  Write
$\delta_n=\int_0^\infty e^{-t}/(t+n)\,dt$.

\begin{theorem}[Primitive odd-Taylor obstruction]\label{p08:rlap:odd}
For every choice of odd positive integers $M=M(n)$,
\[
 b_M(n)\bigl[S_M(n)-\log n-\gamma\bigr]\longrightarrow+\infty
                         \qquad(n\to\infty).
\]
In particular the primitive integer linear forms
$a_M(n)-b_M(n)\gamma-b_M(n)\log n$ are strictly positive, but never
approach zero along an unbounded sequence of growing arguments and
odd truncation orders.  No growth restriction on $M(n)$ is imposed.
\end{theorem}
\begin{proof}
Taylor's integral remainder for $e^{-t}$, integrated after division by
$t$, gives
\[
 \operatorname{Ein}(n)=S_M(n)+(-1)^M E_M(n),\qquad
 E_M(n)=\frac1{M!}\int_0^n\frac1t
                \int_0^t e^{-s}(t-s)^M\,ds\,dt>0.
\]
Together with
$\operatorname{Ein}(n)=\gamma+\log n+e^{-n}\delta_n$, this proves
for odd $M$ that
\begin{equation}\label{p08:rlap:positive-error}
 S_M(n)-\log n-\gamma=E_M(n)+e^{-n}\delta_n>0.
\end{equation}

First suppose $n\ge4$ and $M\le3n/2$.  Restrict the outer integral to
$n-n/(M+1)\le t\le n$ and the inner one to $0\le s\le1$.
These ranges are compatible since $t\ge n/2\ge2$.  They give
\[
 E_M(n)\ge
 \frac{e^{-1}n^M}{(M+1)M!}
       \left(1-\frac1{M+1}-\frac1n\right)^M
 \ge \frac{e^{-11}n^M}{(M+1)M!}.
\]
Indeed, with $u=1/(M+1)+1/n$, one has $u\le3/4$,
$Mu\le5/2$, and $\log(1-u)\ge-u/(1-u)$.
The last lower bound tends to infinity uniformly along all such
schedules.  To verify this, the arithmetic--geometric mean inequality
$M!\le((M+1)/2)^M$ and $M+1\le7n/4$ imply
\[
 \frac{n^M}{(M+1)M!}\ge\frac{(8/7)^M}{M+1}.
\]
This tends to infinity whenever $M\to\infty$; for $M$ in any fixed
finite set the original expression tends to infinity with $n$.
A subsequence argument combines these alternatives.  Since $b_M(n)\ge1$,
\eqref{p08:rlap:positive-error} proves the desired assertion in this range.

It remains to handle $M>3n/2$.  If $M/2<p\le M$ is prime and
$p\nmid n$, the $k=p$ term of $S_M(n)$ has $p$-adic valuation $-2$.
Every other summand has valuation at least $-1$: there is no other
multiple of $p$ among $1,\ldots,M$, and $v_p(k!)\le1$ there.
The uniquely smallest valuation cannot cancel.  Therefore
\begin{equation}\label{p08:rlap:exactvaluation}
             v_p\bigl(S_M(n)\bigr)=-2,
             \qquad p^2\mid b_M(n).
\end{equation}
Writing $\vartheta(x)=\sum_{p\le x}\log p$, the distinct primes
in this interval that divide $n$ contribute at most $\log n$.  Hence
\[
 \log b_M(n)\ge2\bigl(\vartheta(M)-\vartheta(M/2)\bigr)-2\log n
                     =M+o(M)-2\log n,
\]
by the prime number theorem.  Jensen's inequality applied to an
exponential random variable of mean one gives $\delta_n\ge1/(n+1)$.
Thus
\[
 \log\bigl(b_M(n)e^{-n}\delta_n\bigr)
 \ge M-n+o(M)-2\log n-\log(n+1)
 \ge M/3+o(M)-3\log(n+1)\longrightarrow+\infty.
\]
Equation~\eqref{p08:rlap:positive-error} completes the proof, including
schedules alternating between the two ranges.
\end{proof}

This is an actual reduced-denominator statement.  Its forms have
coefficient gcd one because $a_M(n),b_M(n)$ are coprime.
It does not prove a barrier for arbitrary approximants to $\gamma$:
even truncation orders permit cancellation between the two positive
terms in \eqref{p08:rlap:positive-error}, and combinations at several
arguments might cancel logarithms or introduce new common factors.
No such cancellation estimate is supplied here.  The theorem closes
this specific growing-argument odd-Taylor route, including orders
chosen after inspecting all previous coefficients.

\subsection{Gompertz: a growing-pole Pad\'e family and its exact content support}

Define integers $P_n,Q_n$ by
\begin{equation}\label{p08:rlap:recurrence}
 P_0=0,\quad P_1=1,\qquad Q_0=1,\quad Q_1=2,
 \qquad U_n=2nU_{n-1}-(n-1)^2U_{n-2}\quad(n\ge2).
\end{equation}
The recurrence applies separately to $P$ and $Q$.
Let $g_n=\gcd(P_n,Q_n)>0$ for $n\ge1$.

\begin{proposition}[Exact positive error and growing poles]
\label{p08:rlap:Pade}
For every $n\ge0$,
\begin{align}
 Q_n&=n!L_n(-1)=\sum_{j=0}^n\binom nj^2j!,\label{p08:rlap:Qsum}\\
 Q_n\delta-P_n&=n!I_n>0,\qquad
 I_n=\int_0^\infty\frac{e^{-t}t^n}{(1+t)^{n+1}}\,dt.
 \label{p08:rlap:remainder}
\end{align}
Here $L_n$ is the ordinary Laguerre polynomial.  In particular
$P_n/Q_n$ is a rational approximation with a certified nonzero error,
and its primitive integer form is exactly
$(Q_n/g_n)\delta-P_n/g_n=n!I_n/g_n$.
\end{proposition}
\begin{proof}
The Rodrigues formula
$L_n(t)=e^t(d/dt)^n(e^{-t}t^n)/n!$ and $n$ integrations by parts give
\[
 \int_0^\infty\frac{e^{-t}L_n(t)}{1+t}dt
   =\int_0^\infty\frac{e^{-t}t^n}{(1+t)^{n+1}}dt.
\]
All boundary terms vanish: the derivatives of $e^{-t}t^n$ of order
less than $n$ vanish at zero, and exponential decay handles infinity.
The polynomial $A_n(t)=n!L_n(t)$ has integer coefficients.
Consequently
\[
 Q_n=A_n(-1),\qquad
 P_n=\int_0^\infty e^{-t}\frac{A_n(-1)-A_n(t)}{1+t}dt
\]
are integers, since the quotient is an integer polynomial and
$\int_0^\infty e^{-t}t^kdt=k!$.
They have the stated initial values.  The Laguerre recurrence is
$A_n=(2n-1-t)A_{n-1}-(n-1)^2A_{n-2}$.
At $t=-1$ this gives the recurrence for $Q$; integration after
$1/(1+t)$ gives the same recurrence for the remainders when $n\ge2$,
because $\int e^{-t}A_{n-1}(t)dt=0$.  It follows for $P$ as well.
The finite Laguerre series, with its index reversed, gives
\eqref{p08:rlap:Qsum}.  Positivity of $I_n$ proves nonvanishing.
\end{proof}

The integral and recurrence are classical Stieltjes--Laguerre
approximants, discussed alongside the newer constructions in
\cite{p08:rlap:WV,p08:rlap:MY}; they are not new approximants.
The following arithmetic reduction directly addresses their actual
common factors.

\begin{theorem}[Prime support of the primitive content]
\label{p08:rlap:primes}
For every $n\ge1$ and prime $p$,
\[
 p\mid g_n
 \quad\Longleftrightarrow\quad
 p\le n\ \hbox{ and }\ Q_r\equiv0\pmod p,
                   \qquad r=n\bmod p\in\{0,\ldots,p-1\}.
\]
Consequently $3$ and $5$ never divide $g_n$; $2\mid g_n$ exactly
when $n\ge3$ is odd; and $7\mid g_n$ exactly when
$n\ge7$ and $n\equiv2\pmod7$.
\end{theorem}
\begin{proof}
The recurrence gives the exact determinant
\begin{equation}\label{p08:rlap:determinant}
 P_nQ_{n-1}-P_{n-1}Q_n=((n-1)!)^2\qquad(n\ge1).
\end{equation}
This follows from its value one at $n=1$ and multiplication by
$(n-1)^2$ at each next step.  Thus $p>n$ cannot divide $g_n$.

Modulo $p$, every summand in \eqref{p08:rlap:Qsum} with $j\ge p$
vanishes.  If $n=kp+r$, the identity
$(1+X)^n\equiv(1+X^p)^k(1+X)^r\pmod p$
shows that $\binom nj\equiv\binom rj$ for $0\le j<p$.
It follows that
\begin{equation}\label{p08:rlap:Qmod}
                         Q_n\equiv Q_r\pmod p.
\end{equation}
For $k\ge1$, the recurrence at $kp+1$ gives
$P_{kp+1}\equiv2P_{kp}\pmod p$, because $(kp)^2\equiv0$.
Induction through the rest of the block, using the same recurrence
coefficients as at indices $1,\ldots,p-1$, yields
\[
                     P_{kp+r}\equiv P_{kp}Q_r\pmod p
                                      \qquad(0\le r<p).
\]
Hence, for $p\le n$, $Q_r\equiv0$ implies that both $P_n$ and $Q_n$
vanish modulo $p$, and the reverse implication follows from
\eqref{p08:rlap:Qmod}.  The small-prime assertions follow by evaluating
$Q_0,\ldots,Q_{p-1}$ modulo $p$: modulo $3$ these are $1,2,1$;
modulo $5$ they are $1,2,2,4,4$; and modulo $7$ they are
$1,2,0,6,6,6,6$.  Modulo $2$ the two values are $1,0$.
\end{proof}

This theorem describes prime divisibility, not all prime powers.
For example direct recurrence evaluation gives $g_9=14$, $g_{16}=49$,
and $g_{23}=266$, so even a prime occurring only in one residue class
can occur to higher order.  Replacing the missing valuation analysis
by an observed small gcd would be an unjustified arithmetic step.

\begin{proposition}[Error scale and the exact remaining saving]
\label{p08:rlap:scale}
As $n\to\infty$,
\begin{equation}\label{p08:rlap:Iasymp}
 I_n\sim\sqrt\pi e^{1/2}n^{-1/4}e^{-2\sqrt n},
 \qquad
 \log\left|\delta-\frac{P_n}{Q_n}\right|
                       =-4\sqrt n+O(\log n).
\end{equation}
Therefore along any subsequence the primitive positive forms tend to
zero if and only if
\begin{equation}\label{p08:rlap:gcdthreshold}
 \frac{g_n}{n!\,n^{-1/4}e^{-2\sqrt n}}\longrightarrow+\infty.
\end{equation}
Such a subsequence would prove $\delta$ irrational, but no estimate
establishing \eqref{p08:rlap:gcdthreshold} is proved here.
\end{proposition}
\begin{proof}
Put $\lambda=\sqrt n$ and $t=\lambda u$.  On every compact interval
of positive $u$ the integrand becomes
\[
 \frac1{u+1/\lambda}
 \exp\left[-\lambda(u+u^{-1})+\frac1{2u^2}
                                    +O(\lambda^{-1})\right]du.
\]
The phase $u+u^{-1}$ has its unique minimum $2$ at $u=1$, with
second derivative $2$.  Taylor expansion there, followed by
$v=\sqrt\lambda(u-1)$, gives the Gaussian factor
$e^{1/2}\sqrt{\pi/\lambda}$ multiplying $e^{-2\lambda}$.
Here the usual dominated convergence step can be justified by a
fixed small neighborhood of $1$, where the phase exceeds
$2+c(u-1)^2$, and then by a compact annulus where its excess has a
positive minimum.

To check the noncompact tails rather than assume uniformity of the
expansion, use $\log(1+1/t)\ge1/(t+1)$ and
$f(t)=t+n/(t+1)$.  For large $n$, on $0\le t\le\lambda/2$,
$f(t)\ge f(\lambda/2)\ge(12/5)\lambda$.
On $t\ge2\lambda$, $f'(t)\ge3/4$ and
$f(2\lambda)\ge(12/5)\lambda$.  The two tail integrals are thus
bounded by a polynomial in $\lambda$ times $e^{-(12/5)\lambda}$.
They are negligible, proving the first asymptotic.

For the second, the series
$L_n(-1)=\sum_{k=0}^n\binom nk/k!$ satisfies
\[
 L_n(-1)\le\sum_{k\ge0}\frac{n^k}{(k!)^2}\le e^{2\sqrt n}.
\]
The last inequality follows by retaining diagonal terms in the square
of $\sum_{k\ge0}(\sqrt n)^k/k!$.
The single term with $k=\lfloor\sqrt n\rfloor$, using
$\binom nk\ge n^k e^{-O(k^2/n)}/k!$ and Stirling's formula,
gives $\log L_n(-1)\ge2\sqrt n-O(\log n)$.
Since the rational approximation error equals $I_n/L_n(-1)$,
the stated logarithmic estimate follows.  Finally
\eqref{p08:rlap:remainder} and \eqref{p08:rlap:Iasymp} give the exact
subsequence equivalence.  A nonzero integer form in a rational
number of fixed denominator is bounded away from zero, which proves
the irrationality implication if the missing estimate were available.
\end{proof}

The determinant bound \eqref{p08:rlap:determinant} alone is much too weak
for this purpose.  The prime-support theorem removes certain possible
factors exactly, but does not establish a sufficiently large saving or
rule one out.  Thus the growing-pole attempt has reached a purely
integer-content problem with certified nonvanishing and a sharp error
scale; it has not resolved the named constant.

\subsection{Integer-valued squared-binomial kernels: an explicit escaping attempt}

The next family deliberately violates the earlier integer-coefficient
hypothesis.  Let
\[
 B_n(t)=\binom tn^2=\left(\frac{t(t-1)\cdots(t-n+1)}{n!}\right)^2
                                      \qquad(n\ge1).
\]
It is nonnegative on the real line and integer-valued on every integer,
but usually has nonintegral coefficients.  Put
\[
 J_n=\int_0^\infty\frac{e^{-t}B_n(t)}{1+t}dt,
 \quad M_n=\int_0^\infty e^{-t}B_n(t)dt,
 \quad T_n=\int_0^\infty e^{-t}B_n(t)\log t\,dt.
\]

\begin{proposition}[Exact arithmetic and failure of this integer-valued family]
\label{p08:rlap:binomial}
There are rationals $a_n,A_n$ with
\[
 J_n=\delta+a_n,\qquad T_n=A_n-M_n\gamma,
                  \qquad M_n\in\mathbb Q_{>0}.
\]
The canonical integer forms obtained by multiplying these identities by
$(n!)^2$ have integer coefficients.  Nevertheless $J_n\to+\infty$
and $T_n/M_n\to+\infty$.  Thus neither primitive clearance of the
first family nor reduction of the rational candidate $A_n/M_n$ in
the second gives an approximation proving irrationality of the named
constant.
\end{proposition}
\begin{proof}
Since $B_n(-1)=1$, polynomial division by $1+t$ proves the first
identity with rational remainder; its coefficients have denominators
dividing $(n!)^2$.  If $B_n=\sum c_kt^k$, the moment identity
\[
 \int_0^\infty e^{-t}t^k\log t\,dt=k!(H_k-\gamma)
\]
gives $M_n=\sum c_kk!$ and $A_n=\sum c_kk!H_k$.
Here $H_0=0$ and $k!H_k\in\mathbb Z$ for $k\ge1$, proving
all the integrality assertions.

For $3n\le t\le3n+1$, the falling product is increasing and
$B_n(t)\ge\binom{3n}{n}^2$.  The largest term of the
$\operatorname{Binomial}(3n,1/3)$ distribution occurs at $n$;
since it is at least $1/(3n+1)$,
\[
 \binom{3n}{n}\ge\frac{(27/4)^n}{3n+1}.
\]
Consequently
\begin{equation}\label{p08:rlap:binomialgrowth}
 J_n\ge\frac{e^{-1}}{(3n+2)(3n+1)^2}
                 \left(\frac{729}{16e^3}\right)^n\longrightarrow+\infty.
\end{equation}
The base exceeds one since $e<3$.
If $a_n=u_n/v_n$ is in lowest terms, the primitive form is exactly
$v_n\delta+u_n=v_nJ_n$, whose coefficient gcd is one.  Its positive
value is at least $J_n$, so actual common-factor division cannot rescue
this family.

For $0\le t\le n$, write $k\le t\le k+1$ with
$0\le k\le n-1$.  Bounding each factor in the falling product gives
\[
 \left|\binom tn\right|
 \le\frac{(k+1)!(n-1-k)!}{n!}\le1.
\]
Thus $\int_0^n e^{-t}B_n(t)dt\le1$.
For $0\le t\le1$ the sharper bound
$|\binom tn|\le t/n$ gives a negative logarithmic contribution of
absolute value at most $1/(9n^2)$.  On $t\ge n$ the logarithm is
at least $\log n$, and on $1\le t\le n$ it is nonnegative.  Hence
\[
 \frac{T_n}{M_n}\ge
   (\log n)\left(1-\frac1{M_n}\right)-\frac1{9n^2M_n}.
\]
The same interval used in \eqref{p08:rlap:binomialgrowth} shows
$M_n\to\infty$ exponentially.  The right side tends to infinity.
Since $A_n/M_n=\gamma+T_n/M_n$, reducing its numerator and denominator
does not alter this divergence or convert it into an approximation
to $\gamma$.
\end{proof}

These computations address three genuine escapes from the previous
fixed-kernel theorem.  Growing pole order produces a useful classical
Pad\'e family, whose actual gcd now has an exact prime-support rule.
Growing Taylor arguments still fail at every odd order even after
primitive reduction.  Integer-valued squared-binomial numerators are
allowed, but their mass moves far into the exponential tail and their
primitive forms fail for an explicit reason.  None of these statements
rules out other mixed kernels, cancellations across distinct logarithms,
or a sufficiently strong new valuation theorem for $g_n$.


% Prefix dct:. Original investigation of rational and mixed moment transfers.
\section{Beyond rank one: rational moment pencils and mixed Catalan moments}
\label{p08:dct:section}

We investigate four specified targets: transcendence of $G$, $\zeta(3)$ and
$\zeta(5)$, and irrationality of $\zeta(7)$.  None is resolved here.
The advance over the preceding raw-tail obstruction is constructive: we
give full-rank rational moment pencils for every odd zeta value, prove their
real-root and central-gap structure, and obtain rational two-sided
approximants with a proved exponential error.  We also exhibit an explicit
full-degree mixed Catalan family that cancels its auxiliary period.
Arithmetic estimates needed to turn these constructions into the named
resolutions are stated separately.

The starting point is the actual rational functional in Fauzan's 2026
paper, Sections 2.1--2.4, rather than the rank-one tail functional
\cite{p08:dct:Fauzan}.  Its $\zeta(5)$ determinant has growing degree and a
positive value at the target; positivity of the value does not mean that
its coefficient matrix is positive.  The extension below makes the
alternating signature explicit.  Classical ingredients, including Hermite's
formula, symmetric matrix pencils, and Gaussian quadrature, are identified
as such; we claim no priority for those ingredients.

\subsection{A positive rational functional for every odd weight}

Fix an odd integer $s\geq3$ and put $\theta_s=\zeta(s)$,
$H_j^{(s)}=\sum_{v=1}^jv^{-s}$.  On rational functions with only simple poles
at negative integer squares, define a $\mathbb Q$-linear functional by
\begin{align}
 \mu_{s,X}(t^e)&=(-1)^eB_{2e+2}
       \frac{(2e+s)!}{(s-1)!(2e+2)!},\label{p08:dct:moments}\\
 \mu_{s,X}\left(\frac1{t+j^2}\right)
       &=j^{s-1}(X-H_j^{(s)})-\frac1{s-1}+\frac1{2j}.
       \label{p08:dct:poles}
\end{align}
Polynomial division and simple partial fractions define the functional
uniquely.  The Bernoulli convention is $B_1=-1/2$.

\begin{theorem}[Odd-weight positive representation]
\label{p08:dct:positive}
For every rational function $R$ in this domain,
\[
 \mu_{s,\theta_s}(R)=\int_0^\infty R(y^2)w_s(y)\,dy,
 \qquad
 w_s(y)=\frac{2(2\pi)^{s-1}}{(s-1)!}\,
                 y^s\sum_{\ell\geq1}\ell^{s-1}e^{-2\pi\ell y}>0.
\]
The measure is finite, has all moments, and is strictly positive on
$(0,\infty)$.  This includes $s=3,5,7$, with no arithmetic hypothesis on
their zeta values.
\end{theorem}
\begin{proof}
Let $f(y)=(e^{2\pi y}-1)^{-1}$.  Since $s-1$ is even,
$w_s(y)=2y^s f^{(s-1)}(y)/(s-1)!$.  The expansion at zero shows
$w_s(y)\to1/\pi$, and at infinity it is
$O_s(y^s e^{-2\pi y})$.
Termwise integration, justified by nonnegativity, gives
\[
 m_e:=\int_0^\infty y^{2e}w_s(y)dy
 =\frac{2(2e+s)!}{(s-1)!(2\pi)^{2e+2}}\zeta(2e+2).
\]
Euler's even-zeta formula gives \eqref{p08:dct:moments}.

For $a>0$ one has the elementary identity
\[
 \frac{d^{s-1}}{dy^{s-1}}\frac{y^s}{y^2+a^2}
             =(s-1)!a^{s-1}\operatorname{Im}(a-iy)^{-s}.
\]
Polynomial division leaves a polynomial of degree $s-2$ and remainder
\[
 (-1)^{(s-1)/2}a^{s-1}\frac{y}{y^2+a^2}.
\]
Differentiate its two simple fractions.  Integrating by parts $s-1$ times introduces
no sign.  All boundary terms vanish: at zero each product is $O(y)$,
since the derivative orders sum to $s-2$, and at infinity the derivatives
of $f$ decay exponentially.  Hermite's formula for Hurwitz zeta
\cite[25.11.29]{p08:dct:DLMF} yields
\[
 \int_0^\infty\frac{w_s(y)}{y^2+a^2}dy
             =a^{s-1}\zeta(s,a)-\frac1{2a}-\frac1{s-1}.
\]
At $a=j$ use $\zeta(s,j)=\theta_s-H_j^{(s)}+j^{-s}$ to obtain
\eqref{p08:dct:poles}.  Linearity proves the assertion for every $R$.
\end{proof}

\subsection{Full rank, real roots, and the exact central gap}

Let $0\leq N<K$, $h=K-N$, $r\geq0$ be integers, and set
\[
 D_N(t)=\prod_{j=1}^N(t+j^2),\quad
 E(t)=\prod_{j=N+1}^K(t+j^2),\quad W(t)=D_N(t)^r.
\]
Define the $h\times h$ rational affine pencil and its polynomial by
\[
 A_{s,N,K,r}(X)=
   \left[\mu_{s,X}\left(\frac{W(t)t^{i+j}}{E(t)}\right)\right]_{0\leq i,j<h},
 \qquad P_{s,N,K,r}(X)=\det A_{s,N,K,r}(X).
\]

\begin{theorem}[Signature and real-root structure]
\label{p08:dct:roots}
The polynomial $P_{s,N,K,r}\in\mathbb Q[X]$ has degree exactly $h$, is
positive at $\theta_s$, and has only real roots, counted with multiplicity.
Its leading coefficient is
\[
 (-1)^{h(h-1)/2}\prod_{j=N+1}^Kj^{s-1}D_N(-j^2)^r.
\]
If $rN$ is even, precisely $\lceil h/2\rceil$ roots are below $\theta_s$
and $\lfloor h/2\rfloor$ are above it.  If $rN$ is odd, the counts are
reversed.  No root equals $\theta_s$.

In particular, for $K=40n$, $N=3n$ and $r=s$ the degree is $37n$;
there are exactly $\lfloor37n/2\rfloor$ roots below $\zeta(s)$ and
$\lceil37n/2\rceil$ above.  Thus the target lies in the central gap of
an explicitly defined rational polynomial for every odd $s\geq3$.
For $s=5$ this is the determinant family in Fauzan's construction,
before its scalar normalizations.
\end{theorem}
\begin{proof}
At the target the matrix $M=A(\theta_s)$ is positive definite: its
quadratic form at a nonzero real polynomial $q$ of degree below $h$ is
$\int_0^\infty W(y^2)q(y^2)^2w_s(y)/E(y^2)dy>0$.
Write $A(X)=M+(X-\theta_s)B$.  Partial fractions show
\[
 B=V\operatorname{diag}(\alpha_j)V^{\mathsf T},\qquad
 V_{ij}=(-j^2)^i,\qquad
 \alpha_j=\frac{j^{s-1}W(-j^2)}{E'(-j^2)}.
\]
Every $\alpha_j$ is nonzero and $V$ is square invertible.
The Vandermonde identity gives the stated determinant of $B$.
Moreover
$\operatorname{sgn}(\alpha_j)=(-1)^{rN+j-N-1}$, so Sylvester's law of
inertia gives the stated alternating positive and negative counts for $B$.

The real symmetric matrix $T=M^{-1/2}BM^{-1/2}$ is invertible.  If its
eigenvalues are $\gamma_1,\ldots,\gamma_h$, then
\[
 P(X)=\det M\prod_{i=1}^h(1+(X-\theta_s)\gamma_i).
\]
Its roots are $\theta_s-1/\gamma_i$.  Positive eigenvalues give roots
below the target and negative eigenvalues give roots above it.  Congruence
preserves inertia, proving all assertions.  Finally, $sN$ has the parity
of $n$ and $37n$ has the parity of $n$, giving the central counts.
\end{proof}

\begin{corollary}[Certified algebraic positivity interval]
\label{p08:dct:interval}
If both signatures in Theorem~\ref{p08:dct:roots} occur, let $a<\theta_s<b$
be its two adjacent roots.  Then
\[
       A(X)>0\quad\Longleftrightarrow\quad a<X<b.
\]
Both endpoints are real algebraic numbers that can be isolated using the
explicit rational pencil, without inserting a numerical value of
$\theta_s$.  On this interval $\log\det A(X)$ is strictly concave and has
exactly one critical point.  That point is algebraic.

For principal subspaces of a fixed pencil, positivity intervals decrease
under inclusion.  The eigenvalues of the corresponding relative matrices
$M^{-1/2}BM^{-1/2}$ interlace.  These assertions do not claim nesting when
the denominator $E$ or the numerator $W$ is changed.
\end{corollary}
\begin{proof}
All factors $1+(X-\theta_s)\gamma_i$ must be positive, giving exactly
the stated interval.  Standard exact real-root isolation of rational
polynomials and rational positive-definiteness tests identify it; the
root counts in Theorem~\ref{p08:dct:roots} specify its position.
Twice differentiating gives
\[
 \frac{d^2}{dX^2}\log\det A(X)
     =-\operatorname{tr}\bigl((A(X)^{-1/2}BA(X)^{-1/2})^2\bigr)<0.
\]
The first derivative tends to $+\infty$ at $a$ and $-\infty$ at $b$.
The unique critical point is a root of the rational polynomial $P'$.
Restriction of quadratic forms proves interval nesting.  The relative
operator on a subspace is an orthogonal compression in the $M$ inner
product; the usual min--max argument proves eigenvalue interlacing.
\end{proof}

This is a genuine escape from the rank-one obstruction: the target
coefficient has full rank $37n$, and its two signatures both grow.
It supplies algebraic enclosures, not a transcendence conclusion.

\subsection{Rational two-sided approximants with exponential error}

Let $\nu_s$ be the pushforward of $w_s(y)dy$ under $t=y^2$.  Its moments
$m_e$ are the rational numbers in \eqref{p08:dct:moments}.  Write
$S_s(z)=\int(t+z)^{-1}d\nu_s(t)$ for $z>0$.
Let $q_n$ be the monic degree-$n$ orthogonal polynomial for $\nu_s$ and
$r_n$ that for $t\,d\nu_s$.  Their coefficients are rational, because
the defining positive definite moment systems have rational entries.
Define
\[
 L_n(z)=\frac{1}{q_n(-z)}
       \int\frac{q_n(-z)-q_n(t)}{t+z}\,d\nu_s(t),
 \qquad
 U_n(z)=\frac{m_0-L_n^{\,t\nu_s}(z)}{z}.
\]
The integrands after division are polynomials.  These are rational
functions over $\mathbb Q$, not integrals that still contain an unknown
zeta value.  This is the classical Gaussian/Radau construction, here
applied to the odd-zeta functional.

\begin{theorem}[Explicit rational enclosure family]
\label{p08:dct:pade}
For $n\geq1$, $j\geq1$, put
\[
 b_{s,j}=H_j^{(s)}+\frac{j^{1-s}}{s-1}-\frac1{2j^s},\quad
 \ell_{n,j}=b_{s,j}+j^{1-s}L_n(j^2),\quad
 u_{n,j}=b_{s,j}+j^{1-s}U_n(j^2).
\]
Then $\ell_{n,j},u_{n,j}\in\mathbb Q$ and
\begin{align*}
 0<\theta_s-\ell_{n,j}&\leq m_{2n}j^{-s-4n-1},\\
 0<u_{n,j}-\theta_s&\leq m_{2n+1}j^{-s-4n-3}.
\end{align*}
For each fixed $j$, the lower bounds increase and the upper bounds decrease
with $n$.  If $n(j)=\lfloor\pi j/2\rfloor$, then for all sufficiently large
$j$,
\[
       0<u_{n(j),j}-\ell_{n(j),j}
                         \leq C_s j^{-1/2}e^{-2\pi j}.
\]
All quantities entering the rational endpoints are obtained from finite
linear systems with explicit Bernoulli moments.
\end{theorem}
\begin{proof}
Orthogonal polynomials for a positive measure of infinite support in
$(0,\infty)$ have $n$ simple positive zeros.  For completeness, if fewer
than $n$ zeros in this interval changed sign, multiplying by the product
of their linear factors would give a polynomial of degree below $n$ whose
product with the orthogonal polynomial has one sign and nonzero integral,
a contradiction.  This also excludes other zeros by degree counting.
Hence $|q_n(-z)|,|r_n(-z)|\geq z^n$.
Polynomial division and orthogonality give the exact positive errors
\[
 S_s(z)-L_n(z)=\frac{1}{q_n(-z)^2}
                   \int\frac{q_n(t)^2}{t+z}d\nu_s(t)>0,
\]
\[
 U_n(z)-S_s(z)=\frac{1}{z r_n(-z)^2}
                   \int\frac{t r_n(t)^2}{t+z}d\nu_s(t)>0.
\]
For the first identity, subtract the definition of $L_n$ from $S_s$;
the difference between $q_n(t)$ and $q_n(t)^2/q_n(-z)$, divided by
$t+z$, is $q_n$ times a polynomial of degree below $n$.
The second follows by applying the first to $t\nu_s$ and using
$zS_s(z)=m_0-\int t/(t+z)d\nu_s$.
The minimal-norm property of monic orthogonal polynomials bounds their
squared norms by $m_{2n}$ and $m_{2n+1}$, respectively.  Thus the errors
are at most $m_{2n}z^{-2n-1}$ and $m_{2n+1}z^{-2n-2}$.
Equation~\eqref{p08:dct:poles} proves the stated bounds for $\theta_s$.

For monotonicity one may use the exact variational formula
\[
 L_n(z)=\max_{\deg p<n}\int\bigl(2p(t)-(t+z)p(t)^2\bigr)d\nu_s(t).
\]
Completing the square proves it: the optimizer is
$(1-q_n(t)/q_n(-z))/(t+z)$.  Nested polynomial spaces give monotonicity;
apply the same argument to $t\nu_s$ for the upper bound.

Finally $m_k=2(2k+s)!\zeta(2k+2)/((s-1)!(2\pi)^{2k+2})$.
For $n=\lfloor\pi j/2\rfloor$ the integers $4n+s$ and $4n+s+2$ differ
from $2\pi j$ by bounded amounts depending only on $s$.
Stirling's inequality applied to the two error bounds therefore gives
$O_s(j^{-1/2}e^{-2\pi j})$ for each.  Their sum proves the final assertion.
\end{proof}

\begin{proposition}[A certified height bound, not an irrationality proof]
\label{p08:dct:height}
Write the endpoints in Theorem~\ref{p08:dct:pade} as reduced fractions.
Their positive denominators have logarithms at most
\[
             O_s\bigl(n^2\log(n+j+2)+j\log(j+2)\bigr).
\]
For $n\asymp j$ this is $O_s(j^2\log j)$, whereas the certified error
is only $e^{-2\pi j+O_s(\log j)}$.
\end{proposition}
\begin{proof}
The identity
$B_m=\sum_{k=0}^m(-1)^k k!\,\left\{\begin{smallmatrix}m\\k\end{smallmatrix}\right\}/(k+1)$,
obtained by expanding $z/(e^z-1)$, shows that the denominator of $B_m$
divides $\operatorname{lcm}(1,\ldots,m+1)$.
Thus all $m_e$, $e\leq2n+1$, have a common denominator
$D=(s-1)!\operatorname{lcm}(1,\ldots,4n+5)$.
The factorial formula bounds their logarithmic magnitudes by
$O_s(n\log(n+2))$; the elementary bound $\log D=O_s(n\log(n+2))$ suffices.

The variational formula gives an explicit rational expression
$L_n(z)=b^{\mathsf T}A^{-1}b$, with
$A_{ik}=m_{i+k+1}+zm_{i+k}$ and $b_i=m_i$, $0\leq i,k<n$.
At $z=j^2$, $DA$ and $Db$ are integer arrays of logarithmic size
$O_s(n\log(n+j+2))$.  The adjugate formula and the elementary determinant
bound $|\det C|\leq n!\max|C_{ik}|^n$ yield the asserted quadratic bound
for numerator and denominator.  The same argument applies to $t\nu_s$.
Adding $b_{s,j}$ introduces at most
$s\log(j!)+O_s(\log j)$, and multiplication by powers of $j$ has the
same stated bound.  Reduction of a rational fraction cannot increase its
denominator.
\end{proof}

This proves rapid rational approximation for $s=3,5,7$ by one uniform
construction.  It does not prove irrationality of $\zeta(7)$: the estimates
do not show that denominator times error tends to zero.  They also do not
prove that the actual denominators grow at the upper-bound rate.

\subsection{An explicit full-degree mixed Catalan family}

The different escape mechanism in OpenAI's Catalan paper is cancellation
between two moment kernels \cite[Section 2.1]{p08:dct:OpenAI}.  Put
$f(t)=\sqrt{1-t^2}$ and define
\[
 M(i,j)=\int_{-1}^1\int_0^1\frac{|t|t^i u^j}{f(t)(1-tu)}du\,dt,
 \qquad
 Z(i,j)=\int_0^1\int_0^1\frac{t^i u^j}{1-tu}du\,dt.
\]
For $c_k=4^{-k}\binom{2k}{k}$, interpreted as zero unless $k$ is a
nonnegative integer, the paper's exact moment evaluation is
\[
 M(i,j)=M^0(i,j)+4G c_{(i-j)/2}
                    +\tfrac32\zeta(2)c_{(j-i-1)/2},\qquad
 Z(i,j)=Z^0(i,j)+[i=j]\zeta(2),
\]
where $M^0(i,j),Z^0(i,j)$ are rational.  The following new choice of rows
is much simpler than the paper's arithmetically optimized rows; it is used
only to isolate what cancellation and rank growth accomplish.

\begin{proposition}[Full degree after exact auxiliary-period cancellation]
\label{p08:dct:catalan}
For $n\geq1$ and $0\leq r<n$ set
\[
 P_r(t)=t^r,\qquad
 D_r(t)=\sum_{\substack{k\geq0\\r+1+2k<n}}c_k t^{r+1+2k}.
\]
Extend $M,Z$ linearly in their first argument.  The matrix
$F_n=[M(P_r,j)-\tfrac32 Z(D_r,j)]_{0\leq r,j<n}$ has entries in
$\mathbb Q+\mathbb QG$.  Replacing $G$ by $X$ gives a rational matrix
$F_n(X)$ such that $\det F_n(X)$ has degree exactly $n$ and leading
coefficient $4^n$.

The actual value is not always positive.  In dimension three,
\[
 \det F_3(X)=64(X-7/8)(X^2-27X/32-9/128),
                  \qquad \det F_3(G)<0.
\]
Thus full rank and exact period cancellation hold uniformly, but they
do not supply the positive-Gram nonvanishing mechanism.
\end{proposition}
\begin{proof}
Since $f(t)^{-1}=\sum_{k\geq0}c_kt^{2k}$,
$tP_r/f-D_r$ vanishes to order at least $n$ at zero.  The coefficient of
$\zeta(2)$ in column $j<n$ is
$\tfrac32[t^j](tP_r/f-D_r)=0$.
The $X$ coefficient is $4c_{(r-j)/2}$, a lower triangular matrix with
diagonal entries $4$.  Its determinant is $4^n$, proving exact degree.

For an entirely rational check of the dimension-three formula, the matrix is
\[
 F_3(X)=
 \begin{pmatrix}
 4X-3/2&3/2&5/4\\
 7/8&4X-11/4&7/8\\
 2X+1&2&4X-2
 \end{pmatrix}.
\]
Direct expansion gives the displayed polynomial.
The alternating series gives $7/8<G<23/25$.  The lower inequality follows
already from $1-1/9=8/9$; for the upper inequality use the partial sum
through $j=8$, which is less than $23/25$ by rational arithmetic.
The quadratic $X^2-27X/32-9/128$ is increasing on $[7/8,23/25]$ and its
value at $23/25$ is negative.  The determinant is consequently negative.
\end{proof}

The optimized OpenAI construction includes separate arithmetic
nonvanishing and prime-divisibility arguments; our elementary rows do not
inherit them.  The proposition gives a reproducible high-rank family and
an exact demonstration of the additional issue that remains.

\subsection{Four target verdicts and the quantitative missing inputs}

\paragraph{Transcendence of $\zeta(5)$.}
The new signature theorem and rational enclosures apply.  The external
Fauzan estimates, however, are explicit about their current reach:
for $K\in40\mathbb N$, $\deg Q_K=(37/40)K$,
\[
 0<Q_K(\zeta(5))\leq e^{-(79/1600)K^2},\qquad
 \log H(Q_K)\leq(511067/40000)K^2+o(K^2)
\]
\cite[Appendices C.2--C.3]{p08:dct:Fauzan}.  A resultant lower bound for an
algebraic target of degree $d$ requires a decay-to-height ratio exceeding
$d-1$.  Here that ratio is only $1975/511067<0.004$.
It excludes degree one but supplies no degree-two exclusion.  The
real-root factorization does not change these numerical estimates.

\paragraph{Transcendence of $\zeta(3)$.}
Theorems~\ref{p08:dct:positive}--\ref{p08:dct:pade} apply with $s=3$ and give full
rank and a strict exponential rational enclosure.  Ap\'ery's irrationality
is external and classical.  Neither these bounds nor Lai's 2026 linear
independence of $1,\zeta(3),\zeta(4)$ supplies a degree-uniform
transcendence estimate \cite{p08:dct:Lai}.  In particular, the heights in
Proposition~\ref{p08:dct:height} do not make its rational errors small enough
to exclude algebraic irrationality.

\paragraph{Irrationality of $\zeta(7)$.}
Positivity, exact degree and the central-root counts now follow from the
same proved theorem, rather than from analogy with $s=5$.  The arithmetic
side remains unproved.  A September 2026 claim by Anand must not be
confused with an individual formalized result: the related Lean repository
states that at least one of $\zeta(7),\zeta(9),\ldots,\zeta(21)$ is
irrational, which does not specify $\zeta(7)$ \cite{p08:dct:ZetaWindow}.
A primary issue report identifies an inconsistent outer-prime integral
in Anand's proposed transfer \cite{p08:dct:ZetaIssue}.  We do not use that
preprint's claimed negative margin.  Its last piece can be checked
independently without a numerical integration.  For the printed scalar
normalization, $K<p\leq2h<2K$ implies
\[
 v_p(S_K)=-2\sum_{i=1}^{h-1}\left\lfloor\frac{2i}{p}\right\rfloor
                   =-(2h-p-1).
\]
The other factorials are $p$-units, and the paper explicitly chooses
$\gamma_p^{\rm out}=0$ for $p>K$.  Hence the scaled outer integrand on
$1<y\leq37/20$ is $37/20-y$, not the negative expression in its Table~5.
Its integral is $289/800$, whereas the table's last row integrates to
$-629/50$; the discrepancy is exactly $10353/800$.  This verifies a
specific defect in the proposed transfer, without claiming that this
single correction settles all other estimates in that preprint.
For our own constructions the exact
missing assertion is an integral normalization $Q_n$ with nonzero value,
degree $O(n)$ and $\log|Q_n(\zeta(7))|\leq-\delta n^2$ for some
$\delta>0$.  The positive representation proves nonvanishing of the
unnormalized determinant; it does not supply this denominator estimate.

\paragraph{Transcendence of $G$.}
The mixed family in Proposition~\ref{p08:dct:catalan} cancels $\zeta(2)$ and
has degree $n$.  It supplies neither a uniform nonvanishing theorem nor
height and decay estimates.  The later OpenAI irrationality result is an
external input, not a transcendence theorem.  A successful transfer would
need all three estimates for the same integer polynomials; full degree by
itself is insufficient.

For clarity, the resultant criterion used above follows from one line of
arithmetic.  If $\theta$ has degree $d$, minimal-polynomial leading
coefficient $a$, and conjugates $\theta_j$, then for a nonzero value of an
integer polynomial $Q$ of degree $m$ and height $H$,
\[
 1\leq\left|a^m\prod_{j=1}^dQ(\theta_j)\right|,
 \qquad
 \log|Q(\theta)|\geq-(d-1)\log H-O_\theta(m+\log(m+1)).
\]
Thus a fixed positive height cost cannot prove transcendence merely by
increasing determinant size.  One needs improved decay-to-height ratios
excluding every degree, or a different arithmetic principle.
The new results here concern the structure and approximation power of
full-rank families; none is counted as a new resolution of the four
specified individual questions.


\section{Fifteen named targets and the remaining arithmetic step}
\label{p08:end:targets}

The table records fifteen separately examined targets, not fifteen
new theorems. Fourteen are not resolved by this investigation.
The fifteenth is a previously proved project result, included
because the new transfer had to be checked against it. Entries
concern the precise arithmetic question in the second column;
known irrationality does not settle a transcendence target.

{\small
\begin{longtable}{rL{0.20\linewidth}L{0.37\linewidth}L{0.28\linewidth}}
\toprule
No. & Target & Attempt and certified output & Remaining step\\
\midrule
\endfirsthead
\toprule
No. & Target & Attempt and certified output & Remaining step\\
\midrule
\endhead
1 & $\gamma$ irrationality &
Actual reduced-denominator analysis excludes every odd Taylor-order
schedule; integer-valued square kernels also checked. &
A different family with primitive, nonzero integer forms tending
to zero.\\[4pt]
2 & Gompertz $\delta$ irrationality &
Growing-pole Pad\'e errors and every prime in the gcd content
are determined exactly. &
Higher prime-power valuations and enough total cancellation.\\[4pt]
3 & $e+\pi$ irrationality &
Norm and $E$-function transfers are tested; both fail a concrete
algebraic-input hypothesis. &
Control involving both exponential and period values.\\[4pt]
4 & $e\pi$ irrationality &
The analogous multiplicative substitution introduces the same
uncontrolled logarithm base. &
An individual nonvanishing result, not the classical disjunction.\\[4pt]
5 & Catalan $G$ transcendence &
Auxiliary $\zeta(2)$ is removed in a mixed moment family;
the determinant has full degree. &
Uniform nonvanishing and all-degree value/height control.\\[4pt]
6 & $\zeta(3)$ transcendence &
Positive rational moments, exact real-root counts, and strict
exponentially shrinking rational enclosures. &
Arithmetic height bounds strong enough for every algebraic degree.\\[4pt]
7 & $\zeta(5)$ transcendence &
The actual full-rank pencil is analyzed; the printed
decay/height ratio is computed. &
Its stated bounds do not exclude even algebraic degree two.\\[4pt]
8 & $\zeta(7)$ irrationality &
The positive odd-weight functional and enclosures are proved
directly; an external denominator calculation is checked. &
A valid integral normalization with sufficiently small value.\\[4pt]
9 & Erd\H{o}s--Borwein sum transcendence &
Exact decomposition into countably algebraically independent
sparse blocks, with prior result credited. &
A finite-to-infinite arithmetic transfer with quantitative tails.\\[4pt]
10 & Full reciprocal Fibonacci sum transcendence &
Exact theta--Lambert decomposition; sparse contrast ideals
are completely classified. &
Exclude cancellation between the full theta and Lambert values.\\[4pt]
11 & Prime Mertens constant irrationality &
M\"obius--zeta expansion with an explicit exponential tail. &
Control the infinite combination involving $\gamma$ and odd zeta
values.\\[4pt]
12 & Khinchin constant irrationality &
A positive Stieltjes representation has infinitely many poles,
explaining failure of the finite-kernel transfer. &
An arithmetic theorem for this infinite representation.\\[4pt]
13 & Feigenbaum constant irrationality &
Algebraic period-doubling approximants have bounded absolute
height but increasing uncontrolled degree. &
Degree-sensitive approximation separation for the limit.\\[4pt]
14 & Glaisher constant irrationality &
Functional-equation and prime-logarithm reductions are made
explicit. &
Control an exponential involving uncontrolled constants and an
infinite prime sum.\\[4pt]
15 & $\log\operatorname{Ein}(1)$ transcendence &
Previously proved via independence over finite exponential
fields; the argument is restated. &
No new closure counted; outside the algebraic-input norm algorithm.\\
\bottomrule
\end{longtable}
}

The named-target proofs and formulas are in Sections
\ref{p08:cte:section}, \ref{p08:rlap:section}, \ref{p08:dct:section},
\ref{p08:ctm:section}, and \ref{p08:rd:section}.
No numerical precision, count of examples, or number of pages is
used as evidence of irrationality or transcendence.

\subsection{A quantitative gate for a genuine further advance}

The determinant and approximation branches now give concrete
objects to improve. For a transcendental target candidate $\xi$,
an argument through integer polynomials requires both a nonzero
value and a value bound that defeats algebraic separation at every
fixed degree. For example, suppose nonzero integer polynomials
$P_n$ satisfy
\[
 \deg P_n=O(n),\qquad
 \log H(P_n)\le A n^2+o(n^2),\qquad
 0<|P_n(\xi)|\le\exp(-B n^2+o(n^2)).
\]
For an algebraic $\xi$ of degree $d$, a resultant estimate yields
\[
 \log|P_n(\xi)|\ge -(d-1)\log H(P_n)-O_\xi(n+\log n).
\]
Indeed, the resultant with the minimal polynomial of $\xi$ is a
nonzero integer, and the other conjugate factors are bounded by
$(\deg P_n+1)H(P_n)\max(1,|\xi_j|)^{\deg P_n}$.
Thus $B>(d-1)A$ excludes that degree. A single finite ratio $B/A$
does not exclude every degree. This is why real-rootedness and
fast analytic error alone do not close the transcendence targets.

The analogous linear-form gate is even more concrete:
after full reduction to coprime integer coefficients, one needs a
nonzero form that tends to zero. The content theorem for Gompertz
identifies where prime factors occur, but does not yet provide the
factorial-scale total content needed for that gate.
The odd-Taylor theorem excludes one broad family outright.

\subsection{Proved boundaries and research priorities}

The strongest completed conclusions in this manuscript concern
exact arithmetic structure, not a new famous constant. Finite
moment data can leave the complete algebraic-quantile pattern
undetermined; full algebraic density data can determine it.
Finite algebraic logarithmic data admit a terminating entropy
classifier; nested logarithms and transcendental inputs are outside
its proof. Additive cyclotomic Mahler values have a computable
complete ideal; an infinite sum of independent blocks is outside
that finite certificate. These are proved distinctions rather
than heuristic cautions.

The most direct next mathematical questions are correspondingly
specific: optimize the degree under moment constraints; classify
which algebraic entropy values are realizable at a fixed alphabet
size; extend the finite relation certificate to nonadditive systems
without merely invoking an existing general algorithm; and obtain
a genuinely stronger primitive content or determinant-height
estimate for one named constant. None of these extensions is
silently assumed in the present conclusions.

\paragraph{Validation performed.}
The exact entropy identities were checked independently in
$\mathbb Q(\sqrt3)$; beta-mixture weights and the moment example
were checked with rational arithmetic. The determinant branch
has exact rational Pad\'e and mixed-moment checks.
The Gompertz recurrence and its prime-support consequences were
checked at finite indices in addition to their general proofs.
Separate proof reviews examined the main construction, norm
criterion, analytic subgroup reduction, additive dependence
criterion, and determinant estimates. These checks are not a
formal verification of the manuscript or a substitute for external
expert refereeing.

% END SOURCE: KEIO_Arithmetic_Certificates.tex

\endgroup
\clearpage
\part{Sharp Thresholds and Primitive Forms}
\label{merge:part09}
\begingroup

% BEGIN SOURCE: KEIO_Arithmetic_Thresholds_and_Primitive_Forms.tex

\noindent\textbf{Editorial integration note.} 
The several minimum-alphabet questions have different hypotheses. Parts~\ref{merge:part10}--\ref{merge:part12} strengthen the rational-entropy branch, primitive-content estimates, and formal Gamma analysis while preserving these distinctions. The present prime-grid isolation bound is extended to unrestricted auxiliary denominators in Part~\ref{merge:part10}.


\begin{quote}\small
\noindent\textit{Original part framing.}\par
\noindent\textbf{Scope and verification.}
This standalone manuscript supplies the statements, hypotheses, and
proofs used in its arguments.  It does not require the earlier
223-page or 447-page compilations.  Separate proof reviews and exact
algebraic checks are not a formal machine verification, and a supplied
proof does not establish historical priority.  The depth theorems of
Baker, Gelfond--Schneider, W\"ustholz, and the CM torsion literature
are credited as external inputs.  The Wilson--Ap\'ery connection
already appears in Pr\'evost's 2017 work and is not claimed as a
discovery here.  Source checks and this investigation date from
16--17 October 2026.  Publication level and priority remain matters
for further expert review.
\end{quote}
\section{From examples and convergence to exact costs}
\label{p09:pfi:introduction}

The preceding manuscript established exact arithmetic classifiers
for specified quantile and entropy families, together with finite
Mahler relation certificates.  It also located failures in several
attempts to pass from rapid analytic convergence to arithmetic
conclusions.  The question here is more demanding: which complexity
costs are unavoidable, what is the actual reduced denominator, and
can a failed finite transfer be repaired within its proposed class?

\paragraph{Three and four are sharp entropy thresholds.}
Among positive algebraic probability vectors, the earlier five-atom
example is reduced to the exact minimum three:
\[
 p=\left(\frac{1+\sqrt3}{4},\ 2-\sqrt3,\;
                \frac{3\sqrt3-5}{4}\right),\qquad
 H_2(p)=\frac{3(\sqrt3-1)}2.
\]
Theorem~\ref{p09:pfent:three} proves both existence and minimality.
A different optimization problem has answer four: rational
probabilities with rational entropy can first include a non-dyadic
entry at four atoms (Corollary~\ref{p09:pfent:fourminimal}).
The Galois criterion in Theorem~\ref{p09:pfent:galois} detects
transcendence from a distinct simultaneous positive conjugate,
and Corollary~\ref{p09:pfent:quantum} transfers it to quantum spectra.

\paragraph{Exact moment costs and elliptic degree costs.}
Preserving the first $M$ arcsine moments requires weight degree at
least $M+1$.  At that degree the exception pattern is forced to be
the grid $k/(M+1)$; deleting every exception requires exactly
$M+2$ when $M\ge1$ (Theorem~\ref{p09:pm:sharp}).
Theorem~\ref{p09:pm:algorithm} computes the precise optimum for every
admissible finite pattern.  In the elliptic family, the cost depends
on all Galois conjugates, including nonreal ones.  Real CM cyclotomy
after a quadratic twist gives $n\le8d^2$, and the first noncentral
exception has correction degree four and weight degree seven
(Theorems~\ref{p09:ef:bound} and \ref{p09:ef:threshold}).

\paragraph{Primitive arithmetic, rather than an unreduced denominator.}
Theorem~\ref{p09:pg:content} determines the entire gcd of the classical
Laguerre numerator and denominator, not just its prime support.
The residue tree and its local analytic branches describe all
prime powers.  For Euler's constant, Theorem~\ref{p09:pg:gamma}
reduces the attempted even Taylor escape to at most one crossover
order per large argument, with a necessary exponential cancellation.
For each fixed bound, the arguments admitting any bounded primitive
even-order form have natural density zero (Corollary~\ref{p09:pg:density});
an infinite sparse exceptional subsequence is not excluded.
For the odd-zeta trial family, Theorem~\ref{p09:pz:enclosure} proves
width $\exp(-8\log(1+\sqrt2)n+O_s(\log n))$ together with
$\log\operatorname{den}=O_s(n)$.  The constants still matter:
an order-$n$ denominator bound by itself is not an irrationality
criterion.

\paragraph{Two complete limits on proposed transfers.}
Theorem~\ref{p09:pfm:independence} excludes rational linear identities
among all finite collections of full Lambert dilates and their
derivatives.  The same obstruction survives the theta perturbation
in the reciprocal Fibonacci function.  Separately,
Theorem~\ref{p09:gpf:depth} determines the exact number of beta factors
needed for a specified formal gamma isolation.  It explicitly
concerns the defined identity system; it does not assert that
actual gamma values have no additional relations.

These are uniform results with supplied proofs.  They are not
obtained by counting separate numerical examples as separate
conjecture solutions.  The final target ledger records what remains
unproved for each named constant.


\section{Sharp alphabet thresholds and Galois rigidity of entropy}
\label{p09:pfent:section}

Throughout this section the entropy is measured in bits:
\[
 H_2(p)=-\sum_{i=1}^n p_i\log_2 p_i,
 \qquad p_i>0,\qquad\sum_i p_i=1.
\]
Zero entries, if present, are deleted; alphabet size means the number
of positive entries. The arithmetic dichotomy for rational probability
vectors follows from classical transcendence theory and is explicitly
recorded by Blakaj--Wolf \cite[Theorem 2]{p09:pfent:BW}. Here we prove two
sharp small-alphabet statements, a construction in every real quadratic
field, and a Galois obstruction which transfers to quantum entropy.
These are not claims about the irrationality of a historically named
constant. The priority of the specific refinements below has not been
certified beyond the primary literature checked here.

\subsection{Two elementary consequences of Baker's theorem}

We use the following form of Baker's theorem
\cite{p07:rlog:BakerIII}: real logarithms of positive algebraic numbers which
are linearly independent over $\mathbb Q$ remain linearly independent
over $\overline{\mathbb Q}$. Two consequences will be used repeatedly.

\begin{lemma}[Splitting and transporting logarithmic relations]
\label{p09:pfent:split}
Suppose that $\theta_1,\ldots,\theta_d$ are algebraic numbers independent
over $\mathbb Q$, and each $L_j$ is a rational linear combination of
real logarithms of positive algebraic numbers. Then
\[
                 \sum_j\theta_j L_j=0
                 \quad\Longrightarrow\quad L_j=0\text{ for all }j.
\]
Equivalently, if $\alpha_i>0$ are algebraic, the algebraic linear
relations among $\log\alpha_1,\ldots,\log\alpha_m$ are the
$\overline{\mathbb Q}$-span of their rational linear relations.
Every rational relation, after clearing denominators, is the logarithm
of an identity $\prod_i\alpha_i^{k_i}=1$, with $k_i\in\mathbb Z$.
\end{lemma}
\begin{proof}
Choose a $\mathbb Q$-basis of the finite-dimensional space generated by
the logarithms. Baker's theorem makes this basis independent over
$\overline{\mathbb Q}$. Comparing its coefficients proves both
statements; in the first statement compare also the coefficients of
the independent $\theta_j$. Positivity makes exponentiation and taking
real logarithms mutually inverse in the last assertion.
\end{proof}

\begin{lemma}[Positive rational powers of two]
\label{p09:pfent:powers}
If $r_i\in\mathbb Q$ and $\sum_{i=1}^n2^{r_i}=1$, then every $r_i$
is an integer. In particular, if $2^r+2^s=1$, then $r=s=-1$.
\end{lemma}
\begin{proof}
Take a common positive denominator $d$ and put $v=2^{1/d}$. Its
minimal polynomial $X^d-2$ is irreducible by Eisenstein's criterion,
so $1,v,\ldots,v^{d-1}$ are independent over $\mathbb Q$. Write each
integer $dr_i$ as $d a_i+b_i$ with $0\le b_i<d$. Then
$2^{r_i}=2^{a_i}v^{b_i}$. In the claimed sum every coefficient of a
nonconstant basis vector is nonnegative, and is strictly positive if
some $b_i$ equals that vector's exponent. Independence forces all
$b_i=0$. For two terms, write the negative integral exponents as
$-m,-n$, with $1\le m\le n$. The equality
$2^{n-m}+1=2^n$ forces $m=n=1$.
\end{proof}

\subsection{The least alphabet for algebraic irrational entropy is three}

\begin{proposition}[Binary lower bound]
\label{p09:pfent:binary}
For algebraic $0<p<1$,
\[
        H_2(p,1-p)\in\overline{\mathbb Q}
        \quad\Longleftrightarrow\quad p=\tfrac12.
\]
The algebraic value is then $1$.
\end{proposition}
\begin{proof}
First suppose $p$ is irrational and the entropy $h$ is algebraic. The
identity
\[
 -\log(1-p)-p\log\frac{p}{1-p}=h\log2
\]
and Lemma~\ref{p09:pfent:split} imply $h\in\mathbb Q+\mathbb Qp$:
otherwise $1,p,h$ would be independent and would force $\log2=0$.
Write $h=a+bp$ with $a,b\in\mathbb Q$. Splitting the coefficients of
$1,p$ now gives
\[
        1-p=2^{-a},\qquad \frac p{1-p}=2^{-b}.
\]
Thus both $p$ and $1-p$ are rational powers of two, contradicting
irrationality by Lemma~\ref{p09:pfent:powers}.

For rational $p=a/N$, take positive coprime integers $a,N-a$.
The identity
\[
 NH_2(p,1-p)\log2
     =\log\frac{N^N}{a^a(N-a)^{N-a}}
\]
and Baker's theorem show that algebraicity forces the positive rational
number in the logarithm to be an integral power of two. For each odd
prime $\ell$ this gives
\[
 N v_\ell(N)=a v_\ell(a)+(N-a)v_\ell(N-a).
\]
If $\ell\mid N$, coprimality implies it divides neither numerator,
a contradiction. If $\ell\nmid N$, the equation prohibits it from
dividing either numerator. Hence $N,a,N-a$ are all powers of two.
Coprimality gives $a=N-a=1$ and $N=2$. The converse is immediate.
\end{proof}

\begin{theorem}[Sharp three-atom construction]
\label{p09:pfent:three}
The least number of positive algebraic probabilities whose entropy in
bits is algebraic irrational is exactly three. An explicit vector is
\begin{equation}\label{p09:pfent:threevector}
 p=\left(\frac{1+\sqrt3}{4},\;2-\sqrt3,\;
                         \frac{3\sqrt3-5}{4}\right),
 \qquad H_2(p)=\frac{3(\sqrt3-1)}2.
\end{equation}
Every entry of this vector is irrational.
\end{theorem}
\begin{proof}
Put $u=1+\sqrt3$, so that $u^2-2u-2=0$. The identity
\[
 (t^2-2t-2)^2=t^4-4t^3+8t+4
\]
gives, for $F(t)=t/4+2/t^2+1/t^3$,
\[
       F(t)-1=\frac{(t^2-2t-2)^2}{4t^3},
       \qquad F(u)=1,\qquad F'(u)=0.
\]
Consequently the positive probabilities
$(p_1,p_2,p_3)=(u/4,2/u^2,1/u^3)$ sum to one and satisfy
$p_1-2p_2-3p_3=0$. Expanding their logarithms therefore gives
\[
 H_2(p)=2p_1-p_2-(p_1-2p_2-3p_3)\log_2u
       =\frac u2-\frac2{u^2}
       =\frac{3(\sqrt3-1)}2.
\]
The displayed radical formulas follow by rationalization. Positivity
uses $5/3<\sqrt3<2$; the entropy is algebraic irrational. One atom has
zero entropy, and Proposition~\ref{p09:pfent:binary} rules out two atoms.
\end{proof}

This improves the preceding five-atom construction to the exact
minimum. The conclusion is a sharp arithmetic threshold, not a
numerical search result.

\begin{corollary}[Five irrational atoms with rational entropy]
\label{p09:pfent:fiveirrational}
Put $A=(1+\sqrt3)/4$, $B=2-\sqrt3$ and
$C=(3\sqrt3-5)/4$. The five positive irrational probabilities
\[
                         (A,B/2,B/4,B/4,C)
\]
have entropy $3/2$ bits. Thus rational entropy can occur with every
probability irrational using five atoms. We do not assert that five
is the minimum for this different property.
\end{corollary}
\begin{proof}
Splitting an atom of mass $B$ according to a conditional distribution
of entropy $h$ increases total entropy by $Bh$, as follows directly
by expanding the logarithms. Here the split proportions are
$(1/2,1/4,1/4)$ and their entropy is $3/2$. Hence the new entropy is
\[
       \frac{3(\sqrt3-1)}2+\frac32(2-\sqrt3)=\frac32.
\]
Every displayed atom is a nonzero rational multiple of an irrational
entry of~\eqref{p09:pfent:threevector}.
\end{proof}

\subsection{Complete rational classification for three atoms}

\begin{lemma}[Rational valuation test]
\label{p09:pfent:valuation}
Let $a_1,\ldots,a_n$ be positive integers with sum $N$.
For $p_i=a_i/N$, algebraicity of $H_2(p)$ is equivalent to
\begin{equation}\label{p09:pfent:valuationeq}
       N v_\ell(N)=\sum_i a_i v_\ell(a_i)
                  \quad\hbox{for every odd prime }\ell.
\end{equation}
If it holds, the entropy is rational.
\end{lemma}
\begin{proof}
Use $NH_2(p)\log2=\log R$, where
$R=N^N/\prod_i a_i^{a_i}\in\mathbb Q_{>0}$.
Baker's theorem implies that algebraicity of the ratio
$\log R/\log2$ forces it to be rational. If a positive rational number
is a rational power of two, comparison of prime valuations after
clearing the exponent's denominator shows it is an integral power
of two. Thus the condition is precisely that all odd valuations of
$R$ vanish. Sufficiency is the same identity read backwards.
\end{proof}

\begin{theorem}[Rational ternary classification]
\label{p09:pfent:rationalthree}
For positive rational $p_1,p_2,p_3$ with sum one,
\[
 H_2(p_1,p_2,p_3)\in\overline{\mathbb Q}
 \quad\Longleftrightarrow\quad
 (p_1,p_2,p_3)\text{ is a permutation of }
                  (\tfrac12,\tfrac14,\tfrac14).
\]
The algebraic entropy is $3/2$; every other positive rational ternary
distribution has transcendental entropy in bits.
\end{theorem}
\begin{proof}
Write $p=(a/N,b/N,c/N)$ with $a,b,c$ positive integers,
$a+b+c=N$ and $\gcd(a,b,c)=1$. Assume the valuation criterion
of Lemma~\ref{p09:pfent:valuation}. An odd prime not dividing $N$ cannot
divide any of $a,b,c$, by~\eqref{p09:pfent:valuationeq}.

Fix an odd prime $\ell\mid N$. At least one numerator is not divisible
by $\ell$, by primitivity. At least two are not divisible by $\ell$:
otherwise reducing their sum modulo $\ell$ would contradict
$\ell\mid N$. At least one numerator is divisible by $\ell$, since
the left side of~\eqref{p09:pfent:valuationeq} is positive. Exactly one
numerator therefore contains $\ell$. If that numerator is $a$, put
$t=v_\ell(N)$ and $s=v_\ell(a)$. The valuation identity reads
\[
                  \frac aN=\frac ts,
                  \qquad s>t\ge1.
\]
The reduced denominator of $a/N$ is at most $s$, and
$s\le\log_\ell a<\log_3 N<\sqrt N$. The last strict inequality
holds for every $N>0$ for which the logarithm is positive: indeed
$\log N/\sqrt N\le2/e<\log3$.

If odd primes occurred in two different numerators, the reduced
denominators $d_1,d_2$ of those two probabilities would both be
smaller than $\sqrt N$. Since the third probability is one minus
their sum, its denominator divides $\operatorname{lcm}(d_1,d_2)$.
Primitivity says that $N$ is the least common denominator of all three
probabilities, hence
\[
                 N=\operatorname{lcm}(d_1,d_2)
                   \le d_1d_2<N,
\]
a contradiction. Therefore all odd primes, if any exist, occur in
only one numerator. Relabel it $a$, so that $b=2^u$ and $c=2^v$.

We now exclude this remaining possibility. Write
$a=2^h A$ and $N=2^k B$, where $A,B$ are odd. They are greater
than one in the present case. The odd valuation identities give
$A^a=B^N$. Reduce $a/N=t/s$, with $\gcd(s,t)=1$ and $s>t$.
Then $A^t=B^s$, so unique factorization gives an odd integer $C>1$
with
\[
                         A=C^s,\qquad B=C^t.
\]
The equality $sa=tN$ becomes
\[
                  s C^{s-t}=t2^{k-h}.
\]
The left side is greater than $t$, so $j=k-h>0$.
Coprimality forces $s=2^l$ for some $l\ge1$; $t$ is odd.
Taking the $2$-adic valuation gives $j=l$, leaving
\begin{equation}\label{p09:pfent:oddC}
         t=C^d,\qquad d=s-t=2^l-t>0\text{ odd},
         \qquad C^d+d=2^l.
\end{equation}
Furthermore
\[
             b+c=N-a=2^h C^t d.
\]
If $u=v$, the left side is a power of two, impossible since $C>1$.
Otherwise comparison of odd parts of the sum of two powers of two
gives an integer $w\ge1$ with
\begin{equation}\label{p09:pfent:mod8}
                         C^t d=1+2^w.
\end{equation}
If $l\ge3$, the oddness of $C,d,t$ and~\eqref{p09:pfent:oddC} imply
$t\equiv C\pmod8$ and $d\equiv-C\pmod8$. Hence
$C^t d\equiv-C^2\equiv7\pmod8$, whereas $1+2^w$ is congruent
to $3$, $5$ or $1$ modulo $8$. If $l=2$, equation
$C^d+d=4$ allows only $C=3,d=1$, which would require
$27=1+2^w$, also impossible. If $l=1$, already $C^d+d\ge4>2$.
Thus odd primes do not occur.

It follows that $N,a,b,c$ are powers of two, so all three probabilities
are negative integral powers of two. Order their exponents as
$1\le r_1\le r_2\le r_3$. If $r_1\ge2$, their sum is at most
$3/4$. Thus $r_1=1$, and the remaining two terms sum to $1/2$.
Lemma~\ref{p09:pfent:powers}, or direct subtraction, gives
$r_2=r_3=2$. The converse follows by evaluating the entropy.
\end{proof}

\begin{corollary}[A second sharp alphabet threshold]
\label{p09:pfent:fourminimal}
The least alphabet size for a rational probability vector with
rational entropy and at least one non-dyadic probability is four.
An example is
\[
       p=(\tfrac38,\tfrac14,\tfrac13,\tfrac1{24}),
       \qquad H_2(p)=\tfrac74.
\]
\end{corollary}
\begin{proof}
Proposition~\ref{p09:pfent:binary} and
Theorem~\ref{p09:pfent:rationalthree} exclude smaller alphabets.
For the displayed denominator-$24$ numerators $(9,6,8,1)$, the odd
prime condition is $24=9\cdot2+6\cdot1$ at $3$. The $2$-adic
valuation gives the entropy $3-(6+24)/24=7/4$.
The example was available in the preceding manuscript; its minimality
is the additional conclusion proved here.
\end{proof}

\subsection{Algebraic irrational entropy in every real quadratic field}

\begin{proposition}[Dyadic critical-point construction]
\label{p09:pfent:critical}
Suppose $u>0$ is algebraic and
$F(t)=\sum_i2^{-k_i}t^{m_i}$, with $k_i,m_i\in\mathbb Z$, satisfies
$F(u)=1$ and $F'(u)=0$. Then the positive probabilities
$p_i=2^{-k_i}u^{m_i}$ have algebraic entropy
$H_2(p)=\sum_i k_i p_i$.
The same conclusion holds after splitting each positive dyadic
coefficient of a Laurent polynomial into a finite sum of powers of two.
\end{proposition}
\begin{proof}
The derivative identity is $\sum_i m_i p_i=0$. Substitute
$\log_2p_i=-k_i+m_i\log_2u$ into the entropy. A positive dyadic
rational has a finite binary expansion, allowing the indicated split.
\end{proof}

\begin{theorem}[Every real quadratic field occurs]
\label{p09:pfent:quadratic}
Let $D>1$ be a squarefree integer. There is a finite positive
probability vector with all entries in $\mathbb Q(\sqrt D)$ whose
entropy in bits generates $\mathbb Q(\sqrt D)$ over $\mathbb Q$.
The construction is effective using rational arithmetic and
real-algebraic comparisons.
\end{theorem}
\begin{proof}
Choose a positive dyadic rational $q$ such that $1<q^2D<3$, and set
\[
 b=q^2D-1\in(0,2),\qquad u=1+q\sqrt D.
\]
Thus $u^2-2u-b=0$. The Laurent polynomial
\[
 F(t)=\frac t4+\frac{1-b/2}{t}
                    +\frac b{t^2}+\frac{b^2}{4t^3}
\]
has strictly positive dyadic coefficients and satisfies
\[
                  F(t)-1=\frac{(t^2-2t-b)^2}{4t^3}.
\]
Split its coefficients into powers of two, leaving the first monomial
$t/4$ intact. Proposition~\ref{p09:pfent:critical} gives a positive
probability vector with entropy $h\in\mathbb Q(\sqrt D)$ and an
irrational atom $u/4$. If $h$ is irrational, the construction is done.
If $h$ is rational, split this atom into two equal atoms. Such a split
increases entropy by $u/4$, so the new entropy is irrational and still
belongs to $\mathbb Q(\sqrt D)$. Any irrational element of a quadratic
field generates that field. Choosing $q$, splitting the dyadic
coefficients and testing the rationality of $h$ are all effective.
\end{proof}

\begin{proposition}[A genuine limit on the single-parameter construction]
\label{p09:pfent:positiveconjugate}
Let $F\in\mathbb Q[t,t^{-1}]$ have positive coefficients, with some
nonzero exponent. If an algebraic $u>0$ satisfies $F(u)=1$ and
$F'(u)=0$, then $u$ has no other positive real conjugate over
$\mathbb Q$.
\end{proposition}
\begin{proof}
For $G(x)=F(e^x)$ one has
$G''(x)=\sum_m c_m m^2 e^{mx}>0$. Consequently $G$ has at most one
critical point. Any positive conjugate $v$ of $u$ also satisfies
$F(v)=1$ and $F'(v)=0$, and would give a second critical point
$\log v$ unless $v=u$.
\end{proof}

Thus this construction does not establish that every positive
algebraic number is itself an entropy value, nor does it directly
extend to a prescribed totally positive primitive element of a
higher-degree field.

\subsection{Galois rigidity and quantum transfer}

\begin{theorem}[Positive-conjugate obstruction]
\label{p09:pfent:galois}
Let $p=(p_1,\ldots,p_n)$ be a positive algebraic probability vector.
If $H_2(p)$ is algebraic and a field embedding $\sigma$ of a number
field containing its entries has $\sigma(p_i)>0$ for every $i$, then
\[
                            \sigma(p_i)=p_i\quad\text{for all }i.
\]
Equivalently, a distinct simultaneous positive real conjugate of $p$
certifies transcendence of its entropy in bits.
\end{theorem}
\begin{proof}
Write $h=H_2(p)\in\overline{\mathbb Q}$. The vector
$(p_1,\ldots,p_n,h)$ is an algebraic relation among
$\log p_1,\ldots,\log p_n,\log2$. By
Lemma~\ref{p09:pfent:split}, it is an algebraic linear combination of
integer relation vectors $k$ satisfying
\[
                          \prod_i p_i^{k_i}2^{k_0}=1.
\]
Apply $\sigma$ to the bases in this multiplicative identity. Since
$q_i=\sigma(p_i)>0$ and $\sigma(2)=2$, taking real logarithms gives
the same integer relation among $\log q_1,\ldots,\log q_n,\log2$.
Use the same algebraic linear combination of these relations, without
conjugating its coefficients. It follows that
\[
               \sum_i p_i\log q_i+h\log2=0.
\]
Subtracting the original entropy identity yields
$\sum_i p_i\log(p_i/q_i)=0$. Both $p$ and $q$ are positive
probability vectors, since $\sum q_i=\sigma(1)=1$.
The strict equality case of the elementary relative-entropy
inequality now gives $p=q$. Explicitly,
\[
 \sum_i p_i\log\frac{q_i}{p_i}
       \le \log\left(\sum_i p_i\frac{q_i}{p_i}\right)=0,
\]
and strict concavity of the logarithm makes equality possible only
when every ratio $q_i/p_i$ is the same; normalization makes it one.
\end{proof}

\begin{corollary}[Totally positive probabilities]
\label{p09:pfent:totallypositive}
If every entry of a positive algebraic probability vector is totally
positive over $\mathbb Q$, and at least one entry is irrational, its
entropy in bits is transcendental.
\end{corollary}
\begin{proof}
Use a normal closure of the common number field. Every automorphism
sends every entry to a positive real number. Were the entropy
algebraic, Theorem~\ref{p09:pfent:galois} would make all entries fixed by
the entire Galois group, hence rational.
\end{proof}

\begin{corollary}[Rational characteristic polynomials]
\label{p09:pfent:quantum}
Let $\rho$ be a finite-dimensional density matrix whose characteristic
polynomial has rational coefficients. If any eigenvalue of $\rho$ is
irrational, its von Neumann entropy in bits is transcendental.
Consequently its entropy is algebraic if and only if all nonzero
eigenvalues are rational and satisfy the odd-prime valuation test
in Lemma~\ref{p09:pfent:valuation}; in that case it is rational.
This applies, in particular, to rational density matrices.
\end{corollary}
\begin{proof}
The eigenvalues are nonnegative algebraic numbers and sum to one.
After deleting zeros, every algebraic conjugate of any remaining
eigenvalue is a nonzero root of the rational characteristic polynomial,
hence another positive eigenvalue. Thus every remaining eigenvalue is
totally positive. Apply Corollary~\ref{p09:pfent:totallypositive}, and
then Lemma~\ref{p09:pfent:valuation} in the rational-spectrum case.
\end{proof}

For example, the rational positive definite density matrix
\[
          \rho=\frac18
             \begin{pmatrix}5&1\\1&3\end{pmatrix}
\]
has eigenvalues $(4\pm\sqrt2)/8$, so its entropy in bits is
transcendental. The Galois statement also applies in arbitrary rank
and does not require a formula for the eigenvalues. In contrast,
the diagonal density matrix with spectrum~\eqref{p09:pfent:threevector}
has algebraic irrational entropy; its characteristic polynomial
cannot have rational coefficients. This distinction identifies an
arithmetic hypothesis that survives the passage from classical
probabilities to quantum spectra.


\section{Sharp degree costs under exact moment constraints}
\label{p09:pm:section}

The inverse arcsine construction from the preceding manuscript can be
sharpened from an existence theorem to an exact minimum-degree
algorithm.  Its first two degree thresholds have closed forms.

On \((0,1)\), let \(\rho_0(x)=1/(\pi\sqrt{x(1-x)})\).  For
\(P\in\mathbb Q[x]\), define
\[
 \mathcal L(P)=\frac{(1-2x)P+2x(1-x)P'}2,\qquad
 w_P=1+\mathcal L(P),
\]
and suppose \(P\ne0\) and \(w_P>0\) on \([0,1]\).
Then \(w_P\rho_0\) is a probability density, with CDF
\[
 F_P(x)=\frac{2\arcsin\sqrt x+\sqrt{x(1-x)}P(x)}\pi.
\]
Every fixed nonzero rational polynomial becomes admissible after
multiplication by a sufficiently small nonzero rational scalar.
For a polynomial of degree \(d\),
\begin{equation}\label{p09:pm:degree}
                         \deg w_P=d+1,
\end{equation}
because its leading coefficient \(a\) becomes \(-(d+1)a\)
in degree \(d+1\) under \(\mathcal L\).

We use the established arithmetic classifier from the preceding
manuscript: for algebraic \(q\in(0,1)\),
\begin{equation}\label{p09:pm:classifier}
 F_P^{-1}(q)\in\overline{\mathbb Q}
 \quad\Longleftrightarrow\quad
 q\in\mathbb Q\quad\hbox{and}\quad
 P\!\left(\sin^2\frac{\pi q}2\right)=0.
\end{equation}
Its transcendence input is Baker's theorem for algebraic linear
forms in logarithms, together with Gelfond--Schneider in the
irrational-algebraic exponent case.  In outline, write the arcsine
term as the logarithm of an algebraic number when \(x\) is
algebraic.  The CDF equation makes its difference from an algebraic
multiple of \(\log(-1)\) algebraic.  Baker's theorem forces this
difference to vanish; hence \(P(x)=0\).  Gelfond--Schneider then
forces the algebraic level \(q\) to be rational.  The converse is
immediate.  We recall this established classifier solely to specify
the input to the new degree problem; see also the classical
logarithmic results \cite{p09:pfm:BakerI,p09:pfm:BakerIII}.

\begin{lemma}[Moment constraints are orthogonality constraints]
\label{p09:pm:orthogonality}
For \(M\geq0\), the density \(w_P\rho_0\) has the same first
\(M\) raw moments as \(\rho_0\) if and only if
\begin{equation}\label{p09:pm:orthogonal}
 \int_0^1x^jP(x)\sqrt{x(1-x)}\,dx=0
                         \qquad(0\leq j<M).
\end{equation}
These are homogeneous rational linear conditions on the
coefficients of \(P\).
\end{lemma}
\begin{proof}
The difference of CDFs is
\(\sqrt{x(1-x)}P(x)/\pi\), which vanishes at both endpoints.
Integration by parts therefore gives, for \(k\geq1\),
\[
 \int_0^1x^k(w_P-1)\rho_0\,dx
 =-\frac{k}{\pi}\int_0^1x^{k-1}P(x)\sqrt{x(1-x)}\,dx.
\]
This proves the equivalence.  Each monomial integral is a beta
integral \(B(j+3/2,3/2)\), a rational multiple of \(\pi\).
Thus division by \(\pi\) makes every constraint rational.
\end{proof}

\begin{theorem}[Two sharp thresholds]\label{p09:pm:sharp}
Fix \(M\geq0\).  Among nonconstant positive rational polynomial
tilts preserving the first \(M\) arcsine moments, the minimum
weight degree is exactly \(M+1\).  At this degree the correction
is uniquely determined up to a nonzero scalar:
\[
                    P(x)=c\,U_M(2x-1),
\]
where \(U_M\) is the Chebyshev polynomial of the second kind.
Its full algebraic-quantile exception set is
\begin{equation}\label{p09:pm:grid}
          G_M=\{k/(M+1):1\leq k\leq M\}.
\end{equation}
For \(M\geq1\), requiring the exception set to be empty raises
the exact minimum weight degree by one, to
\begin{equation}\label{p09:pm:empty}
                              \boxed{M+2}.
\end{equation}
Both minima can be attained arbitrarily close to the unperturbed
weight \(1\) in any prescribed finite \(C^k\) norm.
\end{theorem}
\begin{proof}
The positive weight \(\sqrt{x(1-x)}\) makes its polynomial
inner product positive definite.  A nonzero polynomial of degree
less than \(M\) cannot be orthogonal to every polynomial of degree
less than \(M\), since it would be orthogonal to itself.  At degree
\(M\), the orthogonal complement is one-dimensional.  The change
\(u=2x-1\) identifies its generator with \(U_M(u)\), whose roots
are \(\cos(k\pi/(M+1))\), \(1\leq k\leq M\).
Equations \eqref{p09:pm:degree} and \eqref{p09:pm:classifier} prove the
first minimum and the exact grid \eqref{p09:pm:grid}.

For \(M\geq1\) this grid is nonempty, so an empty exception set
requires correction degree at least \(M+1\).  The polynomials
\[
 P_t(x)=U_{M+1}(2x-1)+tU_M(2x-1),\qquad t\in\mathbb Q,
\]
all satisfy \eqref{p09:pm:orthogonal}.  Consecutive Chebyshev
polynomials have no common zero, as follows from their three-term
recurrence and \(U_0=1\).  Therefore each candidate algebraic
root \(\alpha\) excludes at most one parameter \(t\).
Only finitely many candidates need be considered: if a rational
polynomial of degree \(M+1\) vanishes at
\(\sin^2(\pi q/2)\) with rational \(q\), the order \(N\) of
\(e^{i\pi q}\) satisfies \(\varphi(N)\leq2(M+1)\), and hence
\(N\leq8(M+1)^2\).  The last inequality follows from
\(N\leq2\varphi(N)^2\).  Avoid the finitely many forbidden
rational values of \(t\).  The resulting polynomial has no
exception, proving \eqref{p09:pm:empty}.  Finally multiply the
chosen correction by an arbitrarily small nonzero rational
scalar.  Orthogonality, degree, and its root set stay unchanged,
whereas positivity and the asserted norm closeness follow.
\end{proof}

\subsection{An algorithm for the exact optimum of every admissible pattern}
For \(N\geq3\), put
\[
 C_N=\{2j/N:0<j<N/2,\ \gcd(j,N)=1\},\qquad
 \Psi_N(X)=\prod_{q\in C_N}
                 \left(X-\sin^2\frac{\pi q}2\right).
\]
Classical real cyclotomy gives distinct irreducible monic
polynomials \(\Psi_N\in\mathbb Q[X]\), of degree
\(\varphi(N)/2\).  The admissible finite exception sets are
exactly finite unions of the disjoint blocks \(C_N\).
For such a set \(S\), write
\[
 H_S=\prod_{C_N\subset S}\Psi_N,\qquad s=\deg H_S=|S|;
 \qquad H_{\varnothing}=1.
\]

\begin{theorem}[Exact minimum-degree algorithm with prescribed moments]
\label{p09:pm:algorithm}
Given \(M\geq0\) and an admissible finite exception set \(S\),
the following finite algorithm computes the exact minimum weight
degree and constructs a strictly positive rational polynomial
tilt attaining it while preserving the first \(M\) arcsine moments.

For \(d=0,1,\ldots,2s+M+1\), compute by rational linear algebra
\[
 V_d=\{P\in\mathbb Q[x]:\deg P\leq d,\ H_S\mid P,
                   \ P\text{ satisfies }\eqref{p09:pm:orthogonal}\}.
\]
If \(V_d=\{0\}\), reject this degree.  Otherwise form the
monic greatest common divisor \(G_d\) of any rational basis of
\(V_d\).  Accept the first degree for which every real cyclotomic
factor \(\Psi_N\) of \(G_d\) also divides \(H_S\).
The exact minimum weight degree is then \(d+1\).

All tests are finite: at degree \(d\geq1\), it is enough to
test \(3\leq N\leq8d^2\) with \(\varphi(N)\leq2d\).
At \(d=0\) there are no such factors.  An attaining polynomial
is obtained by enumerating rational coefficient combinations of
the basis until none is divisible by an unselected \(\Psi_N\),
and then scaling by a sufficiently small nonzero rational number.
\end{theorem}
\begin{proof}
By \eqref{p09:pm:classifier}, a correction with all exceptions in
\(S\) and containing \(S\) must be a nonzero member of \(V_d\)
and have no real cyclotomic factor outside \(H_S\).
If such a factor divides \(G_d\), it divides every member of
\(V_d\), so degree \(d\) is impossible.

Conversely suppose no unwanted cyclotomic factor divides \(G_d\).
For each of the finitely many unwanted factors of degree at most
\(d\), its divisibility condition cuts out a proper rational
linear subspace of \(V_d\).  A finite-dimensional vector space
over the infinite field \(\mathbb Q\) is not the union of
finitely many proper linear subspaces.  Thus there is a nonzero
\(P\in V_d\) avoiding all these factors.  Rational enumeration
finds one, and the classifier gives its exact exception set.
This proves both the test and its minimality.  Positivity follows
after scaling without changing the constraints or exceptions.

It remains to prove the stated finite stopping bound.  Let
\(R_j\) be the monic orthogonal polynomial of degree \(j\)
for the positive measure
\[
                    H_S(x)^2\sqrt{x(1-x)}\,dx
                              \quad(0<x<1).
\]
Its moments divided by \(\pi\) are rational, so each
\(R_j\in\mathbb Q[x]\).  Positive definiteness gives its usual
three-term recurrence with nonzero subdiagonal coefficients;
consequently \(R_M\) and \(R_{M+1}\) are coprime.
The polynomial
\[
        P_t=H_S^2(R_{M+1}+tR_M),\qquad t\in\mathbb Q,
\]
has degree \(2s+M+1\) and satisfies every constraint in
\eqref{p09:pm:orthogonal}.  Outside the selected factors of \(H_S\),
each possible cyclotomic root forbids at most one \(t\), by
coprimality.  Avoiding finitely many values therefore gives exactly
\(S\).  The search must terminate within the asserted range.
\end{proof}

The result computes an exact optimum, not merely an upper bound.
For instance, if only the first moment is fixed and no algebraic
quantile exception is allowed, the sharp weight degree is three:
\(P=x^2-5/16\) has the required orthogonality, and its only
root in \((0,1)\) is \(\sqrt5/4\).  Its other conjugate is
negative, whereas every conjugate of a real cyclotomic quantile
coordinate lies in \((0,1)\).  It therefore has no cyclotomic
factor.  Scaling this \(P\) makes a positive tilt with mean
exactly \(1/2\) and a transcendental quantile at every algebraic
level, including the median.


\section{Inverse elliptic exceptions, exact degree cost, and a sharp first threshold}
\label{p09:ef:section}

The preceding project paper gave a finite torsion test for a fixed
elliptic correction.  Here we solve the inverse problem, account for
nonreal Galois conjugates in its exact degree cost, and identify the
first degree at which a noncentral exception can occur.  The uniform
torsion bound used below is a consequence of classical real cyclotomy
for CM curves; it is not a new bound for CM torsion.

\subsection{The family and the previously established classifier}
Let
\[
 E:y^2=1-x^4,\quad O=(-1,0),\qquad
 I(x)=\int_{-1}^x\frac{dt}{\sqrt{1-t^4}},\quad\Omega=I(1).
\]
The smooth projective curve is elliptic with identity $O$.  Its upper
real arc has elliptic logarithm $I(x)$, and its real period is $2\Omega$.
For a nonzero rational polynomial $P$ and a sufficiently small nonzero
rational $\epsilon$, set
\begin{align}
 L_E(P)(x)&=(1-x^4)P'(x)-2x^3P(x),\label{p09:ef:operator}\\
 w_P(x)&=1+\epsilon L_E(P)(x)>0\quad(-1\leq x\leq1),\nonumber\\
 F_P(x)&=\frac{I(x)+\epsilon\sqrt{1-x^4}P(x)}{\Omega}.
 \label{p09:ef:cdf}
\end{align}
This is the CDF of the density $w_P(x)/(\Omega\sqrt{1-x^4})$.
For $q\in\mathbb Q\cap(0,1)$, denote by $\xi_q$ the unique point in
$(-1,1)$ satisfying $I(\xi_q)=q\Omega$.  Elliptic division implies
that $\xi_q$ is algebraic.  The order of
$T_q=(\xi_q,\sqrt{1-\xi_q^4})$ is
\begin{equation}\label{p09:ef:order}
 n(q)=\operatorname{den}(q/2).
\end{equation}
Here $\operatorname{den}$ means the denominator in lowest terms.

We recall explicitly the previously proved classification:
\begin{equation}\label{p09:ef:classifier}
 \{q\in\overline{\mathbb Q}\cap(0,1):F_P^{-1}(q)
          \in\overline{\mathbb Q}\}
 =\{q\in\mathbb Q\cap(0,1):P(\xi_q)=0\}.
\end{equation}
Indeed, if $q$ and $x=F_P^{-1}(q)$ are algebraic, then
$I(x)-q\Omega=-\epsilon\sqrt{1-x^4}P(x)$ is algebraic.
The analytic subgroup theorem implies that an algebraic linear
combination of elliptic logarithms of algebraic points is either zero
or transcendental, so $P(x)=0$.  The real elliptic
Gelfond--Schneider theorem then gives $I(x)/\Omega\in\mathbb Q$.
Both consequences follow from the analytic subgroup theorem in
$\mathbb G_a\times E^2$ and $E^2$, respectively
\cite{p08:cet:Wustholz,p09:ef:Chang}.  Conversely $P(\xi_q)=0$ gives
$F_P(\xi_q)=q$.  We use \eqref{p09:ef:classifier} as an established
input, not as an additional discovery of this section.

\subsection{Galois packets and the inverse problem}
Let $m_q\in\mathbb Q[X]$ be the monic minimal polynomial of $\xi_q$.
Define its real packet by
\begin{equation}\label{p09:ef:packet}
 \mathcal C(q)=\{q'\in\mathbb Q\cap(0,1):m_q(\xi_{q'})=0\}.
\end{equation}
These are finite equivalence classes: two packets are equal exactly
when their minimal polynomials agree.  They need not have cardinality
$\deg m_q$, since some conjugates can be nonreal or lie outside the
probability interval in the $x$ coordinate.

\begin{theorem}[Exact inverse elliptic realization]\label{p09:ef:inverse}
A finite set $S\subset\overline{\mathbb Q}\cap(0,1)$ occurs as the full
exception set on the left side of \eqref{p09:ef:classifier} for a nonzero
rational polynomial correction if and only if $S$ is a finite union
of packets $\mathcal C(q)$.  Write
\[
 H_S=\prod_{\mathcal C\subset S}m_{\mathcal C},\qquad
 D(S)=\deg H_S=\sum_{\mathcal C\subset S}\deg m_{\mathcal C},
\]
where one minimal polynomial is included for each distinct packet;
for $S=\varnothing$, take $H_S=1$ and $D(S)=0$.
Then the exact minimum correction degree is $D(S)$, and the exact
minimum weight degree is
\begin{equation}\label{p09:ef:degree}
                  \boxed{\deg w_P=D(S)+3}.
\end{equation}
The minimum is attained by $P=H_S$.  The resulting weight can be made
arbitrarily close to $1$ in every prescribed finite $C^k$ norm by
choosing $\epsilon$ sufficiently small and nonzero.
\end{theorem}
\begin{proof}
If an exception $q$ occurs, $m_q$ divides $P$.  Every real root
$x\in(-1,1)$ of $m_q$ is the coordinate of a torsion point: choose
an embedding carrying $\xi_q$ to $x$ and extend it to the coordinates
of $T_q$.  Its image has coordinates $(x,y)$, with
$y=\pm\sqrt{1-x^4}$, and still has order $n(q)$.  Changing the sign
of $y$ is group inversion.  Thus the upper-arc point is torsion,
its level $I(x)/\Omega$ is rational, and the entire packet occurs.
Distinct minimal polynomials are coprime.  Consequently $H_S\mid P$
and $\deg P\geq D(S)$.

Conversely the roots of $H_S$ in the interval give exactly the selected
packets, so \eqref{p09:ef:classifier} proves exact realization.  The
polynomial $L_E(H_S)$ is bounded on the compact interval; taking
$0<|\epsilon|<1/(1+\max|L_E(H_S)|)$ ensures strict positivity.
For a polynomial $P$ of degree $d$ and leading coefficient $a\ne0$,
$L_E(P)$ has leading term $-(d+2)a x^{d+3}$.
Hence $w_P$ has degree exactly $d+3$, including when $d=0$.
This proves both minimality assertions.  Every derivative of
$w_P-1$ is proportional to $\epsilon$, which proves the last claim.
\end{proof}

\subsection{A degree-only bound: real cyclotomy survives the transfer}
The real cyclotomy theorem of Aoki and Bourdon--Clark--Stankewicz
states that if an elliptic curve over a real number field $F$ has
complex multiplication by $\mathbb Q(i)$ and an $F$-rational point of
order $n$, then
\[
 \mathbb Q(\zeta_n)^+\subset F,
 \qquad \mathbb Q(\zeta_n)\subset F(i).
\]
Here $\mathbb Q(\zeta_n)^+$ is the maximal real cyclotomic subfield.
We use the version in \cite[Theorem 4.12]{p09:ef:BCS}.  For curves with
CM by the maximal order $\mathbb Z[i]$, of discriminant $-4$, the
same paper's square-root SPY bound gives, for odd $n\geq3$,
\begin{equation}\label{p09:ef:spy}
                  \varphi(n)^2\leq4[F:\mathbb Q]
\end{equation}
because the real field does not contain $\mathbb Q(i)$ and
$\gcd(n,-4)=1$ \cite[Theorem 4.19]{p09:ef:BCS}.

\begin{theorem}[Elliptic degree certificate]\label{p09:ef:bound}
For $q\in\mathbb Q\cap(0,1)$, put
$d=[\mathbb Q(\xi_q):\mathbb Q]$ and $n=n(q)$.  Then
\begin{equation}\label{p09:ef:cyclotomy}
 \mathbb Q(\zeta_n)^+\subset\mathbb Q(\xi_q),\qquad
 \varphi(n)\leq2d,\qquad \boxed{n\leq8d^2}.
\end{equation}
For the odd part $n_{\mathrm{odd}}$ of $n$, one also has
$\varphi(n_{\mathrm{odd}})^2\leq4d$ whenever
$n_{\mathrm{odd}}>1$.
Consequently every exceptional level for a correction $P$ of positive
degree $D$ belongs to the finite list
\begin{equation}\label{p09:ef:candidates}
 q=2j/n,\quad 3\leq n\leq8D^2,\quad
 0<j<n/2,\quad\gcd(j,n)=1,\quad\varphi(n)\leq2D.
\end{equation}
Thus elliptic division and exact substitution into $P$ give a second
terminating classifier, with an a priori order cutoff depending only
on the correction degree.
\end{theorem}
\begin{proof}
The explicit change of variables
\begin{equation}\label{p09:ef:birational}
 U=\frac{2(1-x)}{1+x},\qquad
 V=\frac{4y}{(1+x)^2}
\end{equation}
is a rational isomorphism from $(E,O)$ to
$E':V^2=U^3+4U$, carrying $O$ to infinity.
It identifies the coordinate fields $\mathbb Q(x)=\mathbb Q(U)$.
This Weierstrass curve has CM by $\mathbb Z[i]$ via
$(U,V)\mapsto(-U,iV)$.
Let $U_0=2(1-\xi_q)/(1+\xi_q)>0$, set
$F=\mathbb Q(\xi_q)$ and $a=U_0^3+4U_0>0$.
The quadratic twist
\[
                       aZ^2=U^3+4U
\]
is defined over the real number field $F$, has the same geometric CM,
and contains the $F$-rational point $(U_0,1)$ of order $n$.
The order is unchanged by its isomorphism to $E'$ over $F(\sqrt a)$.
Applying real cyclotomy to this twist gives the two field inclusions
and $\varphi(n)\leq[F(i):\mathbb Q]=2d$.

For completeness, $n\leq2\varphi(n)^2$ follows directly from
\[
 \frac{\varphi(n)^2}{n}
   =\prod_{p^a\parallel n}p^{a-2}(p-1)^2\geq\frac12.
\]
Each odd-prime factor is at least one, and the factor at $2$ is at
least $1/2$.  This proves the stated quadratic cutoff.
Multiplying the twisted point by the $2$-part of $n$ produces an
$F$-rational point of order $n_{\mathrm{odd}}$, to which
\eqref{p09:ef:spy} applies.  Finally a root of $P$ has degree at most $D$,
so \eqref{p09:ef:candidates} follows from \eqref{p09:ef:order}.
Division points on an elliptic curve over a number field are computable
by division polynomials and exact algebraic arithmetic.  Selecting
the upper real point with logarithm $2j\Omega/n$ uses certified real
isolation among an already finite algebraic list.  No unrestricted
numerical zero test is required.
\end{proof}

\subsection{A sharp first noncentral exception}
The packet cost is genuinely different from the number of real
exceptions.  The following result determines its first nontrivial case.

\begin{theorem}[Degree seven is the first noncentral weight]\label{p09:ef:threshold}
In the family \eqref{p09:ef:cdf}, every correction of degree at most three
has exception set either $\varnothing$ or $\{1/2\}$.
A noncentral exception first becomes possible at correction degree
four, or equivalently weight degree seven.  More precisely,
\[
 P_+(x)=x^4-2x^3-2x+1
\]
has exactly the exception $q=2/3$, and $P_-(x)=P_+(-x)$ has exactly
the exception $q=1/3$.  Each is degree-optimal for its singleton set.
For $P=P_+$ the explicit choice $\epsilon=1/64$ is valid.
\end{theorem}
\begin{proof}
First determine the real points of order three on $E'$.
Its third division polynomial is
\[
                    \psi_3(U)=3U^4+24U^2-16.
\]
After \eqref{p09:ef:birational},
\[
 (1+x)^4\psi_3\!\left(\frac{2(1-x)}{1+x}\right)
                       =128(x^4-2x^3-2x+1).
\]
Writing $t=x+x^{-1}$ transforms $P_+(x)=0$ into
$t^2-2t-2=0$.  For $t=1+\sqrt3>2$ there are two positive reciprocal
real roots, exactly one in $(0,1)$; for $t=1-\sqrt3\in(-2,2)$ the
two roots are nonreal.  The polynomial is irreducible over $\mathbb Q$:
$t$ has degree two and $t^2-4=2\sqrt3$ is not a square in
$\mathbb Q(\sqrt3)$, since
$(a+b\sqrt3)^2=2\sqrt3$ with $a,b\in\mathbb Q$ would give
$a^2+3b^2=0$ and $2ab=2$.
The unique root in $(-1,1)$ is the upper real point of order three,
whose arc coordinate is $q=2/3$.  Its negative has level $1/3$ by
$I(-x)=\Omega-I(x)$, and $P_-$ is its minimal polynomial.

Now suppose $\xi_q$ has degree $d\leq3$ and its point has order $n$.
The odd-part bound in Theorem~\ref{p09:ef:bound} yields
$\varphi(n_{\mathrm{odd}})^2\leq12$ if the odd part is nontrivial,
so $n_{\mathrm{odd}}=3$.  In that case a multiple of the point has
order three.  Its $x$ coordinate belongs to $\mathbb Q(\xi_q)$:
the $x$ coordinate of multiplication by an integer is a rational
function of $x$ on this degree-two quotient by group inversion.
That coordinate is the degree-four number just found, a contradiction
to $d\leq3$.  Thus $n$ is a power of two.
Real cyclotomy gives $\varphi(n)\leq6$, hence $n\leq8$.
Order one and order two give arc endpoints; order four gives only
$\xi_{1/2}=0$ on the open upper arc.

For order eight, doubling gives one of the two order-four points
$(U,V)=(2,\pm4)$ on $E'$.  The doubling formula
\[
 U([2]T)=\frac{(U(T)^2-4)^2}{4U(T)(U(T)^2+4)}
\]
therefore gives, after the same substitution, the equation
\[
                        x^4+2x^2-1=0.
\]
Its real roots are $\pm\sqrt{\sqrt2-1}$, both of degree four:
$\sqrt2-1$ has norm $-1$ in $\mathbb Q(\sqrt2)$ and hence is not
 a square in that real quadratic field.  This excludes order eight
when $d\leq3$ and proves the lower threshold.

Finally
\[
 L_E(P_+)=-6x^7+10x^6+6x^4+2x^3-6x^2-2.
\]
The sum of the absolute values of its coefficients is $32$, so
$1+L_E(P_+)/64\geq1/2$ on $[-1,1]$.  Classification
\eqref{p09:ef:classifier} proves the claimed singleton exception, and
Theorem~\ref{p09:ef:inverse} proves exact optimality.
\end{proof}

\begin{corollary}[The two quartiles form one forced packet]\label{p09:ef:quartiles}
The set $\{1/4,3/4\}$ is one real Galois packet, with minimal
polynomial $P_8=x^4+2x^2-1$ and minimum weight degree seven.
It can be the full exception set, although its two quantiles are
opposites and the median of the perturbed law is transcendental.
Neither singleton $\{1/4\}$ nor $\{3/4\}$ is realizable in this family.
\end{corollary}
\begin{proof}
The order-eight calculation in the theorem identifies the two
interior real roots as $\xi_{1/4}=-\sqrt{\sqrt2-1}$ and
$\xi_{3/4}=\sqrt{\sqrt2-1}$, with irreducible polynomial $P_8$.
Its other two roots are nonreal, so the real packet is exactly the
stated pair.  The inverse theorem gives the degree and the impossibility
of splitting the packet.  Taking $P=P_8$ gives exactly these exceptions;
since $P_8(0)=-1$, the classifier excludes an algebraic median.
For example $\epsilon=1/48$ is positive and valid: the coefficient
absolute-value sum of
$L_E(P_8)=-6x^7-8x^5+6x^3+4x$ is $24$.
\end{proof}

For comparison, the singleton median has minimal polynomial $x$,
so it costs one correction degree and four weight degrees.  The
singleton $\{2/3\}$ costs four correction degrees and seven weight
degrees because its other Galois conjugates must also be present in
any rational polynomial, even though they do not give real probability
levels.  This distinction is absent from the elementary count of
visible real exceptions.  The results here solve an inverse and
minimal-degree problem for a specified elliptic probability family;
they do not settle the arithmetic of unrelated classical constants.


\section{Primitive Laguerre content and the single remaining Taylor crossover}
\label{p09:pg:section}

This section returns to two named targets: Euler's constant $\gamma$
and the Euler--Gompertz constant
$\delta=\int_0^\infty e^{-t}/(1+t)\,dt$.
The classical Laguerre approximants themselves are not new
\cite{p09:pg:Lagarias}.  Nor is their congruence-preserving property:
for parameter one, it was explicitly recorded by Bala in 2022
\cite{p09:pg:Bala}.  Our improvement within the present investigation is
to determine their \emph{entire primitive content}, including all prime
powers, and to identify the local analytic objects controlling those
powers.  For $\gamma$ we isolate the one even truncation order per
argument at which the previous odd-order obstruction can possibly
be evaded.  Neither result proves either constant irrational.

\subsection{The complete content of a shifted Laguerre family}

For a positive integer $A$, put
\[
 \delta_A=\int_0^\infty\frac{e^{-t}}{1+At}\,dt,
 \qquad Q_n(A)=\sum_{j=0}^n j!\binom nj^2 A^j
             =A^n n!L_n(-1/A).
\]
Define $P_0(A)=0$, $P_1(A)=1$, and apply the recurrence
\begin{equation}\label{p09:pg:rec}
 U_n=(1+(2n-1)A)U_{n-1}-(n-1)^2A^2U_{n-2}
 \qquad(n\ge2)
\end{equation}
to both $P_n$ and $Q_n$.  Write $g_n(A)=\gcd(P_n(A),Q_n(A))$.

\begin{theorem}[Exact primitive content]\label{p09:pg:content}
For every $A,n\ge1$,
\begin{align}
 \gcd(Q_n(A),Q_{n-1}(A))&=1,\label{p09:pg:adjacent}\\
 g_n(A)&=\gcd\bigl(Q_n(A),((n-1)!)^2\bigr).\label{p09:pg:gcd}
\end{align}
In particular, for every prime $p$,
\[
 v_p(g_n(A))=
 \min\{v_p(Q_n(A)),\,2v_p((n-1)!)\}.
\]
There is no unknown numerator cancellation left in this formula.
\end{theorem}
\begin{proof}
If $p\mid A$, the defining sum gives $Q_n(A)\equiv1\pmod p$.
If $p\nmid A$, write $n=kp+r$, $0\le r<p$.  Terms with $j\ge p$
vanish modulo $p$, and Lucas's elementary binomial congruence gives
$Q_n(A)\equiv Q_r(A)\pmod p$.  In particular $Q_{kp}(A)\equiv1$.
If $p$ divided two consecutive terms, their residues would give
two consecutive zero terms among $Q_0(A),\ldots,Q_{p-1}(A)$.
For indices $2\le j<p$ the backward coefficient
$(j-1)^2A^2$ in \eqref{p09:pg:rec} is nonzero modulo $p$.
Backward propagation would therefore give $Q_0(A)\equiv0$, a
contradiction.  This proves \eqref{p09:pg:adjacent}.

The recurrence and the initial determinant one imply
\[
 P_n(A)Q_{n-1}(A)-P_{n-1}(A)Q_n(A)
             =A^{2n-2}((n-1)!)^2.
\]
Reduction modulo $Q_n(A)$, followed by
$\gcd(Q_n(A),Q_{n-1}(A))=\gcd(Q_n(A),A)=1$, gives
\eqref{p09:pg:gcd}.  Taking valuations proves the last assertion.
\end{proof}

The exact positive remainder is
\begin{equation}\label{p09:pg:error}
 Q_n(A)\delta_A-P_n(A)
 =A^{2n}n!\int_0^\infty
        \frac{e^{-t}t^n}{(1+At)^{n+1}}\,dt>0.
\end{equation}
For completeness, take $B_n(t)=A^n n!L_n(t)$.  The polynomial
$[B_n(-1/A)-B_n(t)]/(1+At)$ has integer coefficients: division
of each difference $(-1/A)^j-t^j$ leaves powers $A^{n-j}$ with
nonnegative exponents.  Its integral against $e^{-t}$ is therefore
an integer.  The Laguerre recurrence identifies this integer with
$P_n(A)$.  Rodrigues's formula and $n$ integrations by parts give
\eqref{p09:pg:error}.  This also supplies nonvanishing independently of
the arithmetic nature of $\delta_A$.

\subsection{All prime powers as a finite residue tree}

\begin{lemma}[Congruence preservation]\label{p09:pg:period}
For every positive integer $m$ and every $n\ge0$,
\[
                 Q_{n+m}(A)\equiv Q_n(A)\pmod m.
\]
\end{lemma}
\begin{proof}
It is enough to work modulo $p^a$.  For nonnegative integers $u,v$
with $u\equiv v\pmod{p^a}$, Vandermonde's identity implies
\[
 v_p\left(\binom uj-\binom vj\right)
       \ge a-\lfloor\log_p j\rfloor\qquad(j\ge1),
\]
where a negative right side is harmless.  One first proves this
for a shift by $p^a$, using
$v_p\binom{p^a}{h}=a-v_p(h)$ for $1\le h\le p^a$,
and then telescopes.  Since
$v_p(j!)\ge\lfloor\log_p j\rfloor$, the difference
$j!A^j[\binom uj^2-\binom vj^2]$ is divisible by $p^a$.
Summing the finitely many nonzero terms proves the claim.
\end{proof}

Let
\[
 R_{p,a}(A)=\{0\le r<p^a:Q_r(A)\equiv0\pmod{p^a}\}.
\]
Starting from $R_{p,1}(A)$, construct $R_{p,a+1}(A)$ by testing
the $p$ children $r+bp^a$, $0\le b<p$, of every retained residue.
Congruence preservation proves that no other residue can survive.
Theorem~\ref{p09:pg:content} now gives the exact finite formula
\begin{equation}\label{p09:pg:tree}
 v_p(g_n(A))=
 \sum_{a=1}^{2v_p((n-1)!)}
 \mathbf 1_{\{n\bmod p^a\in R_{p,a}(A)\}}.
\end{equation}
Thus the primitive coefficients can be determined through finitely
many modular computations without computing the numerator $P_n(A)$.
This is an exact algorithm, not an asymptotic upper bound on $g_n(A)$.

\subsection{Simple roots give exact $p$-adic distance formulas}

For an odd prime $p$, define the continuous interpolation
\[
 \mathcal Q_{p,A}(x)=\sum_{j=0}^\infty j!A^j\binom{x}{j}^{\!2}
                    \qquad(x\in\mathbb Z_p).
\]
The series converges uniformly, since $\binom{x}{j}\in\mathbb Z_p$
and $v_p(j!)\to\infty$.  At nonnegative integer arguments it
equals $Q_n(A)$.  If $0\le r<p$, put
\[
 d_r(A)=2\sum_{j=1}^r j!A^j\binom rj^2(H_r-H_{r-j}),
 \qquad H_0=0.
\]
This number is $p$-integral.

\begin{theorem}[A complete branch above each simple residue root]
\label{p09:pg:hensel}
Suppose $p$ is odd, $p\nmid A$, $0\le r<p$, and
\[
                  Q_r(A)\equiv0\pmod p,
                  \qquad d_r(A)\not\equiv0\pmod p.
\]
There is a unique $\alpha_{p,r,A}\in r+p\mathbb Z_p$ with
$\mathcal Q_{p,A}(\alpha_{p,r,A})=0$, and for every
$x\in r+p\mathbb Z_p$,
\begin{equation}\label{p09:pg:distance}
 v_p\bigl(\mathcal Q_{p,A}(x)\bigr)
                         =v_p(x-\alpha_{p,r,A}).
\end{equation}
In particular the contribution of this residue disk to the actual
content is exactly
\[
 \min\{v_p(n-\alpha_{p,r,A}),\,2v_p((n-1)!)\}.
\]
\end{theorem}
\begin{proof}
Work on $x=r+py$.  In the falling factorial $(r+py)_j$, at least
$\lfloor j/p\rfloor$ factors have all coefficients divisible by $p$.
Consequently every coefficient of
$(r+py)_j^2 A^j/j!$ has valuation at least
$2\lfloor j/p\rfloor-v_p(j!)$, which tends to infinity for odd $p$.
The series is therefore a restricted power series in $y$.

We compute it modulo $p^2$, coefficient by coefficient.
For $j<p$, Taylor expansion at $r$ gives
$Q_r(A)+p d_r(A)y$ modulo $p^2$; the terms $r<j<p$ have a
double zero at $r$ and do not contribute.  Terms with $j\ge2p$
have all coefficients divisible by $p^2$, since
$2\lfloor j/p\rfloor-v_p(j!)\ge2$ in that range.
For $j=p+s$, $0\le s\le r$, there is exactly one factor
congruent to zero at $r$ in $(r+py)_j$.  Wilson's theorem gives
\[
 \frac{(r+py)_{p+s}^{2} A^{p+s}}{(p+s)!}
       \equiv-py^2 A\,\frac{(r)_s^2 A^s}{s!}\pmod{p^2}.
\]
If $r<s<p$, there are two such factors and the term vanishes
modulo $p^2$.  Summation yields the restricted-series identity
\begin{equation}\label{p09:pg:local}
 \mathcal Q_{p,A}(r+py)
       \equiv Q_r(A)+p d_r(A)y-pA Q_r(A)y^2\pmod{p^2}.
\end{equation}
The root assumption therefore gives
\[
 \mathcal Q_{p,A}(r+py)/p
             =c+d_r(A)y+pG(y),
 \qquad c=Q_r(A)/p,\quad G\in\mathbb Z_p\langle y\rangle.
\]
Since $d_r(A)$ is a unit, Hensel's lemma gives a unique zero
$\beta\in\mathbb Z_p$.  Moreover $G(y)-G(\beta)$ is divisible
by $y-\beta$ with an integral restricted-series quotient.  Thus
\[
 \mathcal Q_{p,A}(r+py)/p
  =(y-\beta)\bigl[d_r(A)+p\widetilde G(y)\bigr],
\]
and the bracket is a unit.  Taking
$\alpha_{p,r,A}=r+p\beta$ proves \eqref{p09:pg:distance}.
\end{proof}

The simplicity hypothesis is essential to the proved statement;
it has not been established for every prime and every root.
Without it, the finite tree \eqref{p09:pg:tree} remains valid.
For $A=1$, the prime $2$ needs no analytic argument:
periodicity modulo $4$, together with
$(Q_0,Q_1,Q_2,Q_3)\equiv(1,2,3,2)\pmod4$, gives
$v_2(Q_n)=1$ for odd $n$ and $0$ for even $n$.
For $p=7$, the only root is $r=2$ and $d_2=10\equiv3$;
its unique lift modulo $49$ is $16$.
For $p=11$, the roots $4,7$ have derivatives $4,6$ modulo $11$.
Every higher valuation in these three disks is therefore governed
by one explicitly computable $p$-adic zero.

\subsection{A precise bridge to the left-factorial problem}

Write $!p=\sum_{j=0}^{p-1}j!$.  The familiar congruence
$\binom{p-1}{j}\equiv(-1)^j\pmod p$ gives
\[
                        Q_{p-1}(1)\equiv !p\pmod p.
\]
Connections between this Laguerre sequence and Kurepa's hypothesis
already occur in the literature \cite{p09:pg:Petojevic}; we record the
following consequence for the \emph{actual reduced approximants}.

\begin{proposition}\label{p09:pg:kurepa}
Kurepa's assertion $p\nmid !p$ for every odd prime $p$ is equivalent to
\[
 \gcd(g_n(1),n+1)=
 \begin{cases}
  1,&n\text{ even or }n=1,\\
  2,&n\ge3\text{ odd}.
 \end{cases}
\]
\end{proposition}
\begin{proof}
If an odd prime $p$ divides $n+1$ and $p\le n$, then
$n\equiv p-1\pmod p$ and
$p\mid g_n(1)$ if and only if $p\mid !p$, by
Theorem~\ref{p09:pg:content} and the prime-residue computation.
If $p>n$, it cannot divide $g_n(1)$.
This proves the forward implication, including the exact factor
$2$ established above.  Conversely, if $p\mid !p$ for an odd prime,
choose $n=2p-1$.  Then $p\mid g_n(1)$ and $p\mid n+1$,
contradicting the displayed formula.
\end{proof}

This equivalence does not settle Kurepa's assertion, whose continuing
computational and theoretical difficulty is discussed in Harvey's
2026 ICM account \cite{p09:pg:Harvey}.  It is a warning against replacing
a genuine local arithmetic problem by an unproved universal
``no exceptional residue'' assumption.

\subsection{What the exact content does and does not yield}

The saddle point in \eqref{p09:pg:error}, after writing $c=1/A$ and
$t=\sqrt{nc}\,u$, gives
\[
 Q_n(A)\delta_A-P_n(A)
 \sim \sqrt\pi e^{1/(2A)}A^{n-3/4}n!n^{-1/4}
                           e^{-2\sqrt{n/A}}.
\]
The phase $u+u^{-1}$ has a unique nondegenerate minimum at $u=1$;
the same elementary Laplace calculation as in the parameter-one
case proves this formula.  Accordingly the primitive positive form
tends to zero along a subsequence exactly when
\begin{equation}\label{p09:pg:margin}
 \log g_n(A)-\log(n!A^n)+2\sqrt{n/A}+\tfrac14\log n
                                      \longrightarrow+\infty.
\end{equation}
Equations \eqref{p09:pg:tree} and \eqref{p09:pg:distance} express the entire
left side arithmetically, without the unknown constant $\delta_A$.
But neither an upper bound ruling out \eqref{p09:pg:margin} nor a lower
bound establishing it has been proved here.  A fixed-prime Hensel
description is not a bound for the sum of contributions from all
primes up to $n$.  This is the precise global estimate still missing.

\subsection{Euler's constant: only one even crossover can escape}

Put
\[
 S_M(n)=\sum_{k=1}^M\frac{(-1)^{k-1}n^k}{k\,k!}
       =\frac{a_M(n)}{b_M(n)},\qquad
 \gcd(a_M(n),b_M(n))=1,\qquad b_M(n)>0.
\]
Let
\[
 E_M(n)=\frac1{M!}\int_0^n\frac1t
             \int_0^t e^{-s}(t-s)^M\,ds\,dt,
 \quad T_n=e^{-n}\int_0^\infty\frac{e^{-t}}{n+t}\,dt.
\]
The exact remainder identity is
\[
 S_M(n)-\log n-\gamma=T_n+(-1)^{M+1}E_M(n).
\]
Odd orders have two errors of the same sign and were excluded in the
previous investigation.  Even orders require a different argument.
Let $\kappa>e$ be the unique solution of
\[
                    \kappa(\log\kappa-1)=1
                    \qquad(\kappa\approx3.59112).
\]

\begin{theorem}[A unique possible even crossover]\label{p09:pg:gamma}
For every sufficiently large integer $n$, at most one even integer
$M\ge2n$ satisfies
\begin{equation}\label{p09:pg:window}
                         \tfrac12T_n\le E_M(n)\le2T_n.
\end{equation}
As $n\to\infty$, for every sequence of positive even integers $M=M(n)$
outside this possible exceptional choice,
\[
 b_M(n)\,\bigl|S_M(n)-\log n-\gamma\bigr|\longrightarrow\infty.
\]
Any orders satisfying \eqref{p09:pg:window} necessarily obey
\begin{equation}\label{p09:pg:criticalorder}
             M=\kappa n-\frac{\log n}{2\log\kappa}+O(1).
\end{equation}
If an even-order sequence actually has primitive forms tending to zero,
then, in addition,
\begin{equation}\label{p09:pg:cancellation}
 \left|\frac{E_M(n)}{T_n}-1\right|
           \le\exp\bigl(-[\kappa-1+o(1)]n\bigr).
\end{equation}
\end{theorem}
\begin{proof}
The integral defining $E_M$ gives the exact inequality
\[
 \frac{E_{M+2}(n)}{E_M(n)}
          \le\frac{n^2}{(M+1)(M+2)}<\frac14
                                  \qquad(M\ge2n).
\]
Two even orders cannot both satisfy \eqref{p09:pg:window}.
Outside that window, $|T_n-E_M(n)|\ge T_n/2$.

For $M\le2n$, restricting the outer integral to
$n-n/(M+1)\le t\le n$ and the inner integral to $0\le s\le1$
gives, for $n\ge4$,
\[
 E_M(n)\ge e^{-13}\frac{n^M}{(M+1)M!}.
\]
This tends to infinity uniformly over positive orders $M\le2n$.
For $M\le\sqrt n$ use $M!\le M^M$; for
$\sqrt n<M\le2n$ use Stirling's bound and $e n/M\ge e/2>1$.
Since $T_n\to0$ and $b_M(n)\ge1$, these orders are excluded.

For $M\ge2n$, every prime $M/2<p\le M$ with $p\nmid n$ occurs
to exact valuation $-2$ in the sum $S_M(n)$: its $k=p$ term has
valuation $-2$ and every other term has valuation at least $-1$.
Thus the \emph{actual reduced denominator} satisfies
\[
 \log b_M(n)\ge
  2\bigl(\vartheta(M)-\vartheta(M/2)\bigr)-2\log n
                  =M+o(M)-2\log n.
\]
Also $T_n\ge e^{-n}/(n+1)$ by Jensen's inequality.  Hence
$b_M(n)T_n\to\infty$ uniformly for $M\ge2n$, proving the
assertion outside \eqref{p09:pg:window}.

It remains to locate that window.  Uniformly for $M/n$ in a compact
subset of $(0,\infty)$, boundary Laplace asymptotics give
\begin{equation}\label{p09:pg:Easymp}
 E_M(n)\sim\frac{n^M}{M!\,M(1+M/n)}.
\end{equation}
Indeed, set $t=n-v$ in the defining double integral.  After division
by $n^{M-1}/M!$, the integrand tends to
$\exp[-cv-(1+c)s]$ when $M/n\to c>0$; its integral is
$1/[c(1+c)]$.  On $t\ge n/2$, a fixed exponential in $v+s$
dominates uniformly for $c$ in a compact interval, while the
remaining range is exponentially negligible.  This proves the
stated uniform asymptotic.

Since $T_n\sim e^{-n}/n$, Stirling's formula gives, with $c=M/n$,
\[
 \log\frac{E_M(n)}{T_n}
 =n\{1+c(1-\log c)\}-\tfrac12\log n
  -\tfrac32\log c-\log(1+c)-\tfrac12\log(2\pi)+o(1).
\]
First compare orders near $2n$ and $4n$ and use monotonicity to
confine the window to $2<M/n<4$.  The displayed formula then
forces $M/n\to\kappa$.  Expansion at $\kappa$, whose derivative
is $-\log\kappa$, proves \eqref{p09:pg:criticalorder}.
Finally, if $b_M|T_n-E_M|\to0$, divide by $b_MT_n$ and apply the
denominator lower bound and \eqref{p09:pg:criticalorder} to obtain
\eqref{p09:pg:cancellation}.
\end{proof}

\begin{corollary}[Density-one exclusion of bounded even-order forms]
\label{p09:pg:density}
For every fixed $B>0$, the set
\[
 \left\{n\ge1:\ \text{there is an even }M\ge2
 \text{ with } b_M(n)|S_M(n)-\log n-\gamma|\le B\right\}
\]
has natural density zero.  More precisely, the set of arguments $n$
at which the factor-two window \eqref{p09:pg:window} contains an even
order has natural density
\[
                         \frac{\log2}{\log\kappa}.
\]
\end{corollary}
\begin{proof}
Define
\[
 C_\kappa=\tfrac32\log\kappa+\log(1+\kappa)
                  +\tfrac12\log(2\pi),\qquad
 c_n=\kappa n-\frac{\log n}{2\log\kappa}
                  -\frac{C_\kappa}{\log\kappa}.
\]
The expansion in the proof of Theorem~\ref{p09:pg:gamma} yields, uniformly
for bounded $M-c_n$,
\begin{equation}\label{p09:pg:center}
 \log(E_M(n)/T_n)=-\log\kappa\,(M-c_n)+o(1).
\end{equation}
The number $\kappa$ is transcendental: if it were algebraic,
$\log\kappa=1+1/\kappa$ would be a nonzero algebraic number whose
exponential is algebraic, contrary to the Hermite--Lindemann theorem.
In particular $\kappa$ is irrational.

The sequence $c_n$ is equidistributed modulo $2$.  Indeed, for any
nonzero integer $h$, its Weyl sum is a constant of modulus one times
\[
 \sum_{n\le N} e^{\pi i h\kappa n}
                       n^{-\pi i h/(2\log\kappa)}.
\]
The geometric partial sums are bounded because
$e^{\pi i h\kappa}\ne1$.  Partial summation, using a derivative
of magnitude $O_h(1/n)$ for the second factor, bounds the displayed
sum by $O_h(\log N)$.  Weyl's criterion proves the assertion.

Write $a=\log2/\log\kappa<1$.  Equations
\eqref{p09:pg:criticalorder} and \eqref{p09:pg:center} show that the existence
of a window order is equivalent, apart from a boundary of width
$o(1)$, to $\operatorname{dist}(c_n,2\mathbb Z)\le a$.
Equidistribution modulo $2$ gives density $2a/2=a$; the shrinking
boundary contributes density zero.

If the primitive magnitude is at most $B$, the uniform exclusions in
Theorem~\ref{p09:pg:gamma} first place $M$ in the window.  There
$b_M(n)T_n\to\infty$ uniformly, so
$|E_M/T_n-1|\le B/(b_MT_n)\to0$ uniformly over these choices.
Equation~\eqref{p09:pg:center} now forces
$\operatorname{dist}(c_n,2\mathbb Z)\to0$ on the set in question.
For every $\varepsilon>0$, all sufficiently large members therefore
lie in $\operatorname{dist}(c_n,2\mathbb Z)\le\varepsilon$,
a set of density $\varepsilon$ when $0<\varepsilon<1$.
Letting $\varepsilon\downarrow0$ proves density zero.
\end{proof}

This is a genuine narrowing of the failed Taylor attack: arbitrary
even schedules are reduced to at most one crossover per argument,
and that crossover would need exponentially precise cancellation.
The density statement does not rule out an infinite sparse sequence
of successful crossovers.  Further, the forms still contain
$\log n$, so a successful proof for $\gamma$ alone would also have
to eliminate that logarithmic term.  No estimate accomplishing
either remaining task is asserted here.


% All labels and bibliography keys in this fragment begin with pz:.
\section{Wilson trial functions, primitive determinants, and four zeta targets}
\label{p09:pz:section}

The previous construction produced positive odd-zeta moment pencils and
rational enclosures by solving growing moment systems.  The corresponding
quadratic logarithmic height bound was too expensive for arithmetic
applications.  Here we replace those linear systems by one explicit
polynomial family.  This gives a proved improvement of the height estimate
from $O_s(n^2\log n)$ to $O_s(n)$, together with a stronger
exponential enclosure rate.  It does not establish the transcendence of
$G$, $\zeta(3)$ or $\zeta(5)$, or the irrationality of $\zeta(7)$.

There is an important attribution correction.  The positive Stieltjes
representation of the Hurwitz zeta function, including positive real
weights beyond odd integers, is already in Pr\'evost's work
\cite[Theorems 1--3]{p09:pz:Prevost2017}.  Wilson-polynomial remainder Pad\'e
approximation is also developed by Pr\'evost and Rivoal
\cite{p09:pz:PR2021}.  Neither ingredient is claimed as a new discovery here.
Our statement below is an explicit comparison and transfer for the
particular moment-enclosure construction used in this project; historical
priority for this precise formulation is not asserted.

\subsection{An exact Wilson specialization}

For an odd integer $s\geq3$, write $\theta_s=\zeta(s)$ and
\[
 d\nu_s(t)=w_s(y)\,dy,\qquad t=y^2,\qquad
 w_s(y)=\frac{2(2\pi)^{s-1}}{(s-1)!}y^s
                  \sum_{\ell\geq1}\ell^{s-1}e^{-2\pi\ell y}.
\]
Its moments are
\begin{equation}\label{p09:pz:moments}
 m_{s,e}=\int t^e\,d\nu_s(t)
 =(-1)^e B_{2e+2}\frac{(2e+s)!}{(s-1)!(2e+2)!}.
\end{equation}
Put $S_s(z)=\int(t+z)^{-1}d\nu_s(t)$.  Hermite's formula gives, for
integers $j\geq1$,
\begin{equation}\label{p09:pz:target}
 \theta_s=b_{s,j}+j^{1-s}S_s(j^2),\qquad
 b_{s,j}=\sum_{v=1}^jv^{-s}+\frac{j^{1-s}}{s-1}-\frac1{2j^s}.
\end{equation}

Define the rational polynomial
\begin{equation}\label{p09:pz:R}
 R_n(t)=\sum_{k=0}^n
 \frac{(-1)^k}{(k!)^2}\binom nk\binom{n+k+1}{k}
                  \prod_{r=0}^{k-1}(t+r^2),\qquad n\geq0.
\end{equation}
The empty product equals one.

\begin{lemma}[Wilson identity and exact norm]\label{p09:pz:wilson}
The polynomials $R_n$ are orthogonal for $\nu_3$, and
\[
                 \int R_n(t)^2\,d\nu_3(t)=\frac{n+1}{4}.
\]
If $q_n$ is their monic normalization, then
\[
 t q_n=q_{n+1}+\beta_nq_n+\lambda_nq_{n-1},
\]
where
\[
 \beta_n=\frac{(n+1)^2(2n^2+4n+1)}{(2n+1)(2n+3)},\qquad
 \lambda_n=\frac{n^3(n+1)^3}{4(2n+1)^2}\quad(n\geq1).
\]
\end{lemma}
\begin{proof}
The gamma identities
$|\Gamma(iy)|^2=\pi/(y\sinh\pi y)$ and
$|\Gamma(1+iy)|^2=\pi y/\sinh\pi y$ show that
\[
 w_3(y)=\pi^2y^3\frac{\cosh\pi y}{\sinh^3\pi y}
 =\frac1{4\pi}
 \left|\frac{\Gamma(iy)\Gamma(1+iy)^3}{\Gamma(2iy)}\right|^2.
\]
Thus this is the Wilson weight with parameters $(0,1,1,1)$.
The hypergeometric definition of Wilson polynomials gives
$R_n(t)=W_n(t;0,1,1,1)/(n!)^3$
\cite{p09:pz:DLMF,p09:pz:Wilson1980}.
The limiting parameter zero causes no extra mass: take the first
parameter $a\downarrow0$ in the positive-parameter orthogonality.
Near $y=0$ the factor
$|\Gamma(a+iy)/\Gamma(2iy)|^2$ is bounded by a constant times
$y^2/(a^2+y^2)$, hence is uniformly bounded; at infinity a common
exponential majorant times a polynomial applies.  Dominated convergence
therefore proves the asserted limiting orthogonality.

The Wilson recurrence at these parameters gives the displayed
$\beta_n,\lambda_n$.  Since $m_{3,0}=1/4$, the monic squared norms are
\[
 h_n=\frac14\prod_{k=1}^n\lambda_k
 =\frac{(n!)^5((n+1)!)^3}{4((2n+1)!)^2}.
\]
The leading coefficient of $R_n$ is
$(-1)^n(2n+1)!/((n+1)!(n!)^3)$.  Multiplication by its square gives
$(n+1)/4$.
\end{proof}

\begin{lemma}[A uniform transfer of the squared norm]\label{p09:pz:norm}
For fixed odd $s\geq3$, there is an effective constant $C_s>0$ such that
\[
                   0<\int R_n^2\,d\nu_s\leq C_s(n+1)^{s-2}.
\]
For $k=(s-3)/2$, the more explicit estimate
\[
 \int R_n^2\,d\nu_s
 \leq\frac{n+1}{2}
 \bigl((1+4\pi^2)[1+4(n+k+2)^2]\bigr)^k
\]
is valid.
\end{lemma}
\begin{proof}
For $q\in(0,1)$ and integer $m\geq2$,
\[
 \sum_{\ell\geq1}\ell^m q^\ell
 \leq\frac{m!q}{(1-q)^{m+1}},\qquad
 \sum_{\ell\geq1}\ell^2q^\ell
 =\frac{q(1+q)}{(1-q)^3}.
\]
The first inequality follows from
$\ell^m\leq m!\binom{\ell+m-1}{m}$ and its generating function.
Using $y/(1-e^{-2\pi y})\leq y+(2\pi)^{-1}$, it follows that
\[
 \frac{w_s(y)}{w_3(y)}
 \leq2(1+2\pi y)^{s-3}
 \leq2(1+4\pi^2)^k(1+y^2)^k.
\]

In the orthonormal Wilson basis, multiplication by $t$ is the Jacobi
matrix with diagonal $\beta_n$ and off-diagonals $\sqrt{\lambda_n}$.
Its finite section through index $n+k$ has operator norm at most
$4(n+k+2)^2$, by the displayed recurrence and the maximum row-sum bound.
A path contributing to its $(n,n)$ entry in any power at most $k$ never
uses an index above $n+k$.  Consequently
\[
 \int(1+t)^kR_n(t)^2\,d\nu_3(t)
 \leq[1+4(n+k+2)^2]^k\frac{n+1}{4}.
\]
The weight comparison proves the result.  Strict positivity follows
because $R_n$ is nonzero and $w_s>0$.
\end{proof}

\subsection{Explicit rational intervals without moment-matrix inversion}

For $n,j\geq1$ let $a_{n,j}=R_n(-j^2)>0$, and put
\[
 p_{n,j}(t)=\frac{1-R_n(t)/a_{n,j}}{t+j^2}\in\mathbb Q[t].
\]
Positivity of $a_{n,j}$ follows directly from \eqref{p09:pz:R}: every term
with $k\leq j$ is positive and every term with $k>j$ is zero.
For a polynomial $p$, define
\begin{align*}
 \mathcal L_s(p;z)&=\int[2p-(t+z)p^2]d\nu_s(t),\\
 \mathcal U_s(p;z)&=\frac{m_{s,0}-\int[2tp-t(t+z)p^2]d\nu_s(t)}{z}.
\end{align*}
All these integrals are finite rational combinations of the explicit
moments \eqref{p09:pz:moments}.

\begin{lemma}[Integer-valued arithmetic and Bernoulli evaluation]
\label{p09:pz:intvalued}
Let $\operatorname{Int}(\mathbb Z)$ denote the rational polynomials
mapping $\mathbb Z$ into $\mathbb Z$, let
$L_d=\operatorname{lcm}(1,\ldots,d)$ with $L_0=1$, and let
$\mathcal B:\mathbb Q[x]\longrightarrow\mathbb Q$ be defined by
$\mathcal B(x^m)=B_m$ with $B_1=-1/2$.
For $f\in\operatorname{Int}(\mathbb Z)$ of degree at most $d$:
\begin{enumerate}
\item $L_df'$ is integer-valued;
\item if $a\in\mathbb Z$ and $f(a)=0$, then
      $L_df(x)/(x-a)$ is integer-valued;
\item $L_{d+1}\mathcal B(f)\in\mathbb Z$.
\end{enumerate}
Moreover, for every polynomial $F\in\mathbb Q[t]$,
\[
 \int F(t)d\nu_s(t)
   =\frac{\mathcal B\bigl((x^s F(-x^2))^{(s-2)}\bigr)}{(s-1)!}.
\]
\end{lemma}
\begin{proof}
An integer-valued polynomial has integer coefficients in the Newton basis
$\binom{x}{k}$, because these coefficients are its finite differences
at zero.  The identities
\[
 \frac{d}{dx}\binom{x}{k}
   =\sum_{j=1}^k\frac{(-1)^{j-1}}j\binom{x}{k-j},\qquad
 \frac1{x-a}\binom{x-a}{k}
   =\frac1k\binom{x-a-1}{k-1}
\]
prove the first two assertions; for the second, expand at $a$ and use
$f(a)=0$ to remove the constant term.
The formal identity
$\mathcal B((1+z)^x)=\log(1+z)/z$ gives
$\mathcal B(\binom{x}{k})=(-1)^k/(k+1)$, proving the third assertion.
Finally the displayed integral identity follows on each monomial
$F(t)=t^e$ from \eqref{p09:pz:moments}, and then by linearity.
\end{proof}

\begin{theorem}[A quantitative Wilson trial enclosure]\label{p09:pz:enclosure}
Set
\[
 \ell_{s,n,j}=b_{s,j}+j^{1-s}\mathcal L_s(p_{n,j};j^2),\qquad
 u_{s,n,j}=b_{s,j}+j^{1-s}\mathcal U_s(p_{n,j};j^2).
\]
Then both endpoints are rational and
\[
 \ell_{s,n,j}<\zeta(s)<u_{s,n,j},\qquad
 u_{s,n,j}-\ell_{s,n,j}
   =\frac{\displaystyle\int R_n(t)^2d\nu_s(t)}
          {j^{s+1}R_n(-j^2)^2}.
\]
On the diagonal $j=n$ one has
\begin{equation}\label{p09:pz:rate}
 0<u_{s,n,n}-\ell_{s,n,n}
 \leq\exp\bigl(-8\log(1+\sqrt2)n+O_s(\log n)\bigr).
\end{equation}
The positive reduced denominators of both endpoints satisfy
\begin{equation}\label{p09:pz:height}
                 \log\operatorname{den}(\ell_{s,n,n}),\quad
                 \log\operatorname{den}(u_{s,n,n})
                       \leq O_s(n).
\end{equation}
In particular the exponent in the enclosure rate is about $7.05099n$.
The construction involves explicit polynomial arithmetic and moment
summation, with no growing linear system to invert.
\end{theorem}
\begin{proof}
Completing the square gives, for every polynomial $p$,
\[
 S_s(z)-\mathcal L_s(p;z)
 =\int\frac{[1-(t+z)p(t)]^2}{t+z}\,d\nu_s(t),
\]
\[
 \mathcal U_s(p;z)-S_s(z)
 =\frac1z\int\frac{t[1-(t+z)p(t)]^2}{t+z}\,d\nu_s(t).
\]
For the specified $p$, their sum is
$\int R_n^2d\nu_s/(z a_{n,j}^2)$; both terms are strictly positive.
Equation \eqref{p09:pz:target} proves the exact interval and width.

Let $A_n=\sum_{k=0}^n\binom nk^2\binom{n+k}{k}^2$ be the classical
Ap\'ery integer.  Direct factorial simplification in \eqref{p09:pz:R} gives
\[
 R_n(-n^2)=\sum_{k=0}^n
 \frac{n(n+k+1)}{(n+1)(n+k)}
       \binom nk^2\binom{n+k}{k}^2.
\]
Hence $nA_n/(n+1)\leq R_n(-n^2)\leq A_n$.
Uniform elementary factorial estimates show that the logarithm of the
$k$th summand of $A_n$ equals $n f(k/n)+O(\log(n+1))$, where
\[
 f(x)=2[(1+x)\log(1+x)-(1-x)\log(1-x)-2x\log x].
\]
The endpoint terms use $0\log0=0$.  Its maximum occurs at
$x=1/\sqrt2$, where $f(x)=4\log(1+\sqrt2)$.
Bounding the sum by $(n+1)$ times its maximum term, and using the integer
nearest $n/\sqrt2$ for a lower bound, proves
\[
       \log R_n(-n^2)=4\log(1+\sqrt2)n+O(\log n).
\]
Together with Lemma~\ref{p09:pz:norm}, this proves \eqref{p09:pz:rate}.

For the denominator bound, first observe that
\[
 R_n(-x^2)=\sum_{k=0}^n\binom nk\binom{n+k+1}{k}
                          \binom{x}{k}\binom{x+k-1}{k}
                    \in\operatorname{Int}(\mathbb Z).
\]
The identity follows from pairing the factors $x-r$ and $x+r$.
In particular $a_n=R_n(-n^2)$ is a positive integer.
The numerator $a_n-R_n(-x^2)$ vanishes at both $x=n$ and $x=-n$.
Apply Lemma~\ref{p09:pz:intvalued} twice to conclude that
\[
                   a_nL_{2n}^2p_{n,n}(-x^2)
                          \in\operatorname{Int}(\mathbb Z).
\]
Consequently, multiplying either of the polynomials
$2p-(t+n^2)p^2$ and $2tp-t(t+n^2)p^2$ by
$E_n=a_n^2L_{2n}^4$ makes it integer-valued after $t=-x^2$.
Their degrees in $t$ are at most $2n-1$ and $2n$, respectively.
The Bernoulli-evaluation formula in Lemma~\ref{p09:pz:intvalued}, followed
by its derivative and denominator assertions, proves that
\[
        E_n (s-1)!L_{4n+s+1}^{s-1}
\]
clears their two moment integrals.  Here the undifferentiated polynomial
$x^sF(-x^2)$ has degree at most $4n+s$, and there are $s-2$
derivatives followed by one Bernoulli evaluation.
The same integer also clears $m_{s,0}$.  Thus one explicit common
denominator for the two endpoints is
\begin{equation}\label{p09:pz:denominator}
 Q_{s,n}=2(s-1)n^{s+1}L_n^s a_n^2L_{2n}^4
                         (s-1)!L_{4n+s+1}^{s-1}.
\end{equation}
The elementary Chebyshev bound $\log L_m=O(m)$ suffices here.  It follows,
for example, by bounding the product of primes between $m$ and $2m$
by $\binom{2m}{m}\leq4^m$, summing dyadic intervals, and then summing
prime powers.  The proved exponential size of $a_n$ now gives
$\log Q_{s,n}=O_s(n)$.
Reduction cannot increase a denominator, completing the proof.
\end{proof}

For an explicit asymptotic accounting, put $\psi(m)=\log L_m$.
Equation~\eqref{p09:pz:denominator} gives
\[
 \log Q_{s,n}=2\log a_n+4\psi(2n)+(s-1)\psi(4n+s+1)
                    +s\psi(n)+O_s(\log n).
\]
If the prime number theorem $\psi(x)=x+o(x)$ is used, its leading
coefficient is
$8\log(1+\sqrt2)+5s+4$.  This upper bound is larger than the decay
coefficient $8\log(1+\sqrt2)$; it proves no positive irrationality
margin.  It is not a lower bound for the actual reduced denominators.

The old estimate $e^{-2\pi j+O_s(\log j)}$ used moment-polynomial
degree $\lfloor\pi j/2\rfloor$ and a generic linear-system height bound.
The new statement uses degree $j$ and the exact explicit norm comparison.
It is a strict improvement over those proved estimates, not an assertion
that the previous actual reduced denominators had quadratic growth.
For $s=3$ the lower trial is the actual Gaussian Pad\'e approximant;
classical cancellations can improve its arithmetic further.  For
$s\geq5$ it is a variational trial and need not be the Pad\'e optimizer.

\subsection{Why changing a basis cannot change a primitive determinant}

For $0\ne P\in\mathbb Q[X]$, denote by $\operatorname{prim}(P)$ its
unique primitive integer multiple with positive leading coefficient.

\begin{proposition}[Primitive invariance]\label{p09:pz:primitive}
Let $A(X)$ be a square matrix over $\mathbb Q[X]$ with nonzero determinant.
For any constant matrices $C,D\in\operatorname{GL}_h(\mathbb Q)$,
\[
          \operatorname{prim}(\det(CA(X)D))
                         =\operatorname{prim}(\det A(X)).
\]
Consequently every rational change from monomials to an integer-valued,
Newton, orthogonal, or Lagrange basis leaves the primitive integer
polynomial, its height, and its value at a fixed target exactly unchanged.
This applies both to congruences of the odd-zeta moment pencils and to
independent row and column changes in a mixed Catalan matrix.
\end{proposition}
\begin{proof}
The determinants differ by the nonzero rational scalar $\det C\det D$.
Primitive normalization removes precisely such a scalar.  All conclusions
are then identities, not asymptotic comparisons.
\end{proof}

This does not make a good basis useless: it can expose large common
factors and improve a previously crude bound.  It does rule out the
claim that smaller entry denominators, by themselves, have created a
new primitive arithmetic certificate.  A target-dependent orthogonal
basis is outside the proposition and may have coefficients involving
the very constant whose arithmetic nature is being investigated; it
cannot be treated as a rational integerization step.

There is an analogous invariant for the variational approximants above:
expressing the same polynomial $p$ in a different rational basis does
not change either endpoint.  Replacing the optimizer by $p_{n,j}$,
however, changes the approximation problem, which is why the constructive
transfer in Theorem~\ref{p09:pz:enclosure} is not merely a basis relabeling.

\subsection{Four individual target conclusions}

\paragraph{The transcendence of $\zeta(3)$.}
The Wilson identification supplies exact norms, explicit approximants,
and a stronger estimate than our previous generic moment calculation.
Its links with Ap\'ery's sequence and remainder Pad\'e approximation are
classical, not new irrationality or transcendence results
\cite{p09:pz:Prevost2017,p09:pz:PR2021}.  Our present height and width bounds do
not exclude any algebraic degree beyond previously known results.

\paragraph{The transcendence of $\zeta(5)$.}
Theorem~\ref{p09:pz:enclosure} applies with $s=5$ and quantitatively improves
the proved enclosure/height tradeoff in this construction, including
removal of its factorial denominator loss.
Proposition~\ref{p09:pz:primitive} shows that reexpressing the older full-rank
pencil in a rational orthogonal or integer-valued basis cannot itself
improve the primitive polynomial's actual height.  A new common-divisor
estimate can improve its certified bound; none sufficient for a
transcendence conclusion has been proved here.

\paragraph{The irrationality of $\zeta(7)$.}
Theorem~\ref{p09:pz:enclosure} gives an explicit pair of rational bounds at
rate $e^{-7.05099n+O(\log n)}$.
The proved logarithmic denominator bound is now $O(n)$.  The explicit constants in \eqref{p09:pz:denominator} still do not imply the
required smallness after integerization.  The missing step is a substantially
stronger estimate for the actual reduced denominators or a different
certificate whose denominator and error bounds have a strict positive
exponential margin.  No claim of irrationality follows from the faster
numerical convergence alone.

\paragraph{The transcendence of Catalan's constant $G$.}
The previous elementary mixed Catalan determinant cancels its auxiliary
$\zeta(2)$ term and has full degree.  Rational row/column transformations
preserve its primitive polynomial by Proposition~\ref{p09:pz:primitive}.
They therefore cannot remove its missing uniform nonvanishing theorem
or create a better primitive height by scalar bookkeeping.  A genuinely
different kernel or polynomial family is required for this proposed
escape; no such transcendence certificate is supplied here.


\section{Logarithmic boundary signatures and the finite-to-infinite barrier}
\label{p09:pfm:section}

Let
\[
 L(z)=\sum_{n\geq1}\frac{z^n}{1-z^n}
     =\sum_{n\geq1}\tau(n)z^n,\qquad |z|<1,
 \qquad \mathcal D=z\frac{d}{dz}.
\]
Here \(\tau(n)\) counts positive divisors.  The classical Lambert
identity is recorded in \cite{p09:pfm:DLMF}.  The previously obtained
complete relation certificates apply to finite additive Mahler
systems.  They do not automatically extend to the full Lambert
series.  We now give a direct obstruction that survives arbitrary
finite rational differential filters and arbitrary integer dilations.
This is a functional result, not a transcendence theorem for
\(L(1/2)\).  In particular, the already known failure of an ordinary
Mahler equation for \(L\), also a consequence of the recent
classification of Bell--Smertnig \cite{p09:pfm:BS}, is not counted as a
new result.

\subsection{A greatest-common-divisor matrix detects every linear filter}

\begin{lemma}[Radial logarithmic signature]\label{p09:pfm:radial}
If \(\zeta\) is a primitive \(N\)-th root of unity, \(m\geq1\), and
\(r\geq0\), then, as \(t\downarrow0\),
\begin{equation}\label{p09:pfm:signature}
 \left.\mathcal D^r L(z^m)\right|_{z=\zeta e^{-t}}
 =\frac{r!\gcd(N,m)}{Nm}\,
       \frac{\log(1/t)}{t^{r+1}}+O_{N,m,r}(t^{-r-1}).
\end{equation}
\end{lemma}

\begin{proof}
Put \(f(u)=(e^u-1)^{-1}\).  Its Laurent expansion at zero gives
\[
 (-1)^r f^{(r)}(u)=r!u^{-r-1}+O_r(1)\quad(0<u\leq1),
\]
and the left side is \(O_r(e^{-u})\) for \(u\geq1\).
Termwise differentiation, justified on compact subintervals of
\(t>0\), and separation at \(n=\lfloor1/t\rfloor\) therefore give
\[
 \left(-\frac d{dt}\right)^r L(e^{-t})
 =r!t^{-r-1}\sum_{n\leq1/t}\frac1n+O_r(t^{-r-1})
 =r!t^{-r-1}\log(1/t)+O_r(t^{-r-1}).
\]
Indeed the small-argument error sums as \(\sum_{n\leq1/t}n^r\),
and the large-argument tail is bounded by
\(\sum_{n>1/t}n^re^{-nt}=O_r(t^{-r-1})\).

For a primitive \(Q\)-th root \(\xi\), the summands of
\(L(\xi e^{-u})\) whose index is divisible by \(Q\) sum to
\(L(e^{-Qu})\).  In each remaining residue class the function
\(\xi^a e^{-v}/(1-\xi^a e^{-v})\), together with all its
derivatives, is bounded near \(v=0\) and decays exponentially
at infinity.  Its differentiated contribution is thus
\(O_{Q,r}(u^{-r-1})\), by the same estimates.  The preceding
display consequently has leading coefficient \(r!/Q\) at
\(\xi\).  Finally take \(\xi=\zeta^m\),
\(Q=N/\gcd(N,m)\), and \(u=mt\).  Since
\(\mathcal D=-d/dt\) on the chosen radial path, the factor from
the derivatives is \(m^r\), which proves \eqref{p09:pfm:signature}.
\end{proof}

\begin{theorem}[No finite differential--dilation relation]
\label{p09:pfm:independence}
For any finite set \(M\) of distinct positive integers and any
\(R\geq0\), the functions
\[
 1,\qquad \mathcal D^r L(z^m)
       \quad(m\in M,\ 0\leq r\leq R)
\]
are linearly independent over \(\mathbb C(z)\).
More strongly, every combination
\[
 G(z)=\sum_{m\in M}\sum_{r=0}^R
                       a_{m,r}(z)\mathcal D^r L(z^m),
 \qquad a_{m,r}\in\mathbb C(z),
\]
with at least one nonzero coefficient has logarithmic singularities
at a dense subset of the unit circle.  It admits no meromorphic
continuation through any open arc of that circle.
\end{theorem}

\begin{proof}
Take the largest \(R\) for which some \(a_{m,R}\ne0\).
At every root of unity \(\zeta\) of order \(N\) that is not a
pole of a coefficient, Lemma~\ref{p09:pfm:radial} gives
\begin{equation}\label{p09:pfm:top}
 G(\zeta e^{-t})=
 \frac{R!}{N}\left(\sum_{m\in M}
          \frac{\gcd(N,m)}m a_{m,R}(\zeta)\right)
 t^{-R-1}\log(1/t)+O(t^{-R-1}).
\end{equation}
Replacing \(a_{m,R}(\zeta e^{-t})\) by its value at \(\zeta\)
costs only \(O(t^{-R}\log(1/t))\), which is absorbed by the
error term.  Lower derivative orders are absorbed as well.

For each \(d\in M\) define the rational function
\[
 Q_d(z)=\sum_{m\in M}\frac{\gcd(d,m)}m a_{m,R}(z).
\]
The matrix \((\gcd(d,m))_{d,m\in M}\) is nonsingular.  To
verify this without a determinant formula, use
\[
 \gcd(d,m)=\sum_{e\mid d,\ e\mid m}\varphi(e).
\]
This represents it as \(E\,\operatorname{diag}(\varphi(e))E^T\),
where \(E_{d,e}=1_{e\mid d}\) and the columns run through all
divisors of members of \(M\).  The submatrix consisting of the
columns \(e\in M\), ordered increasingly, is lower triangular
with diagonal one.  Thus \(E\) has full row rank, and the matrix
is positive definite.  Some \(Q_d\) is therefore nonzero.

Let \(\ell\) range over primes not dividing any member of \(M\),
and set \(N=d\ell\).  Then
\(\gcd(N,m)=\gcd(d,m)\) for all \(m\in M\).  Apart from the
finitely many zeros and poles of \(Q_d\) and of the coefficients,
every primitive \(d\ell\)-th root has a nonzero logarithmic
coefficient in \eqref{p09:pfm:top}.  Such primitive roots are dense on
the circle as \(\ell\to\infty\): in any fixed arc, the number
of eligible numerators coprime to \(d\) grows proportionally to
\(\ell\), whereas at most \(d\) are multiples of \(\ell\).

A nonzero term \(t^{-R-1}\log(1/t)\), with error
\(O(t^{-R-1})\), cannot be the radial expansion of a meromorphic
function at that point.  This proves the continuation assertion.
In particular \(G\) cannot equal a rational function.  This is
exactly the claimed linear independence, including the function
\(1\).
\end{proof}

\begin{corollary}[Every finite Lambert filter retains the obstruction]
\label{p09:pfm:filter}
Let \(A\subset\mathbb Z_{>0}\) be finite, let
\((c_a)_{a\in A}\subset\mathbb C\) be nonzero, and suppose
\(H\) is analytic in the unit disk and satisfies
\[
 \mathcal D^rH(\zeta e^{-t})=O_{\zeta,r}(t^{-r-1})
 \quad(t\downarrow0)
\]
for every root of unity \(\zeta\) and every \(r\geq0\).  Set
\[
 F(z)=\sum_{a\in A}c_aL(z^a)+H(z).
\]
Then \(1\) and all finitely many selected functions
\(\mathcal D^rF(z^m)\), for distinct positive \(m\) and
nonnegative \(r\), are linearly independent over \(\mathbb C(z)\).
Every nonzero rational differential combination of these dilates
has the same natural-boundary conclusion as in
Theorem~\ref{p09:pfm:independence}.
\end{corollary}

\begin{proof}
At the highest derivative order, the Lambert coefficients become
the finite Dirichlet convolution
\[
 b_n(z)=\sum_{am=n}c_a a_{m,R}(z).
\]
This convolution cannot vanish identically.  If \(a_0\) is the
least index with \(c_{a_0}\ne0\), and \(m_0\) is the least
index with \(a_{m_0,R}\ne0\), then
\(b_{a_0m_0}=c_{a_0}a_{m_0,R}\ne0\): every other contributing
pair would need one index smaller than its respective minimum.
The terms involving \(H(z^m)\) have size \(O(t^{-R-1})\),
so they do not change the logarithmic coefficients.  The proof
of Theorem~\ref{p09:pfm:independence} applies to the nonzero vector
\((b_n)\).
\end{proof}

\subsection{Two familiar constants: what the obstruction now excludes}

For the Erd\H{o}s--Borwein constant \(E_2=L(1/2)\), the new
conclusion is exact: no choice of finitely many integer dilations,
finitely many derivatives, and rational coefficient functions puts
the full Lambert function into a nontrivial linear relation.
Consequently enlarging the earlier additive system by finitely many
such derivatives or dilates cannot repair the attempted transfer.
An ordinary finite-dimensional linear Mahler system containing
\(L\) would already give a rational linear dependence of finitely
many iterates and is therefore excluded.

This obstruction concerns the proposed method, not the arithmetic
nature of \(E_2\).  Known value results are stronger than mere
irrationality in a different direction: Duverney--Tachiya prove that
\(1,L(1/b),\ldots,L(1/b^h)\) are linearly independent over
\(\mathbb Q\), for integer \(|b|>1\)
\cite[Corollary 1.1]{p09:pfm:DT}.  Their result does not give
algebraic independence or transcendence of \(E_2\), and neither
does the functional theorem above.

For the full reciprocal Fibonacci sum, put
\[
 \mathscr F(z)=\sum_{n\geq1}\frac{z^n}{1-(-1)^nz^{2n}},
 \qquad q=\frac{\sqrt5-1}{2}.
\]
Binet's formula gives
\(\sum_{n\geq1}F_n^{-1}=\sqrt5\,\mathscr F(q)\).
The theta--Lambert decomposition is
\begin{equation}\label{p09:pfm:fibonacci}
 \mathscr F(z)=L(z^2)-L(z^4)+T(z),\qquad
 T(z)=\frac14\sum_{u,v\in\mathbb Z}
             z^{2(u+1/2)^2+2(v+1/2)^2}.
\end{equation}
All exponents in \(T\) are positive integers.  The identity follows
by separating even and odd \(n\) and applying Jacobi's two-square
identity to the odd part.  It is the same classical decomposition
used in the preceding manuscript, rather than a new identity.

The theta perturbation satisfies the hypothesis of
Corollary~\ref{p09:pfm:filter}.  Indeed, with
\(A_{u,v}=2(u+1/2)^2+2(v+1/2)^2\), absolute values give
\[
 \left|\mathcal D^rT(\zeta e^{-t})\right|
 \leq\frac14\sum_{u,v\in\mathbb Z}A_{u,v}^r e^{-tA_{u,v}}
 =O_r(t^{-r-1}).
\]
The final estimate follows by comparison with a two-dimensional
Gaussian integral, or by bounding
\(x^re^{-x}\leq C_re^{-x/2}\) and then separating the two
one-dimensional Gaussian sums.  Thus:

\begin{corollary}[The full Fibonacci function has no finite repair]
\label{p09:pfm:fibonacci-cor}
The functions \(1\) and any finite collection
\(\mathcal D^r\mathscr F(z^m)\) with distinct pairs \((m,r)\)
are linearly independent over \(\mathbb C(z)\).  In particular,
no finite rational differential--dilation filter removes the
Lambert part of \eqref{p09:pfm:fibonacci} to leave a function with
only the indicated theta-type radial growth.
\end{corollary}

The remaining numerical target is still
\(\mathscr F(q)\notin\overline{\mathbb Q}\).  Functional linear
independence at a variable does not prove this assertion at the
single algebraic point \(q\); the required finite Mahler system
is precisely what the corollary rules out.  Nonlinear, modular,
or genuinely infinite-dimensional methods are not excluded.

\subsection{Prime Mertens: exact independence of every finite cutoff}

The standard prime correction representation is
\[
 B_1-\gamma=\sum_p c_p,\qquad
 c_p=\log\left(1-\frac1p\right)+\frac1p.
\]
See \cite{p09:pfm:LZ} for the classical constant and its computation.
We test a distinct proposed repair: cancel the finite logarithmic
part by combining several prime cutoffs.  The following is an
explicit application of Baker's classical theorem, not a new
transcendence theorem of comparable generality.

\begin{proposition}[No finite algebraic cancellation of prime corrections]
\label{p09:pfm:mertens}
For every finite set of distinct primes \(S\), the numbers
\(1\) and \(c_p\), \(p\in S\), are linearly independent over
\(\overline{\mathbb Q}\).  Consequently the same holds for
\(1\) and the successive prime cutoffs
\(C_k=\sum_{j=1}^k c_{p_j}\), \(1\leq k\leq n\).
Moreover the reduced denominator of
\(A_k=\sum_{j=1}^k p_j^{-1}\) is exactly
\(P_k=\prod_{j=1}^k p_j\).
\end{proposition}

\begin{proof}
The rational numbers \((p-1)/p\), \(p\in S\), are
multiplicatively independent.  In a product of their integer
powers equal to one, let \(p\) be the largest prime with a
nonzero exponent.  Its valuation in every smaller factor is zero,
whereas its valuation in its own factor is minus that exponent,
a contradiction.  Thus their real logarithms are linearly
independent over \(\mathbb Q\).  Baker's theorem makes them
linearly independent over \(\overline{\mathbb Q}\), and its
inhomogeneous form says that a nonzero algebraic-coefficient
linear form in these logarithms is transcendental
\cite{p09:pfm:BakerI,p09:pfm:BakerIII}.  A relation among \(1,c_p\)
would make such a form algebraic.  It must therefore vanish,
forcing every coefficient to be zero.  Passing to the \(C_k\)
is an invertible triangular rational change of coordinates.

Finally, the numerator of \(A_k\) over denominator \(P_k\)
is \(\sum_{p\leq p_k}P_k/p\).  Modulo any prime \(p\mid P_k\),
only the term \(P_k/p\) is nonzero.  No factor of \(P_k\)
can cancel, proving the denominator assertion.
\end{proof}

For completeness the prime number theorem makes the mismatch of
the unaccelerated cutoff scales explicit.  If
\[
 T(x)=-\sum_{p>x}c_p,
\]
then \(T(x)\sim(2x\log x)^{-1}\), since
\(-c_p=1/(2p^2)+O(p^{-3})\) and partial summation gives
\(\sum_{p>x}p^{-2}\sim(x\log x)^{-1}\).  Meanwhile
\(\log\prod_{p\leq x}p\sim x\).  Hence
\[
 \left(\prod_{p\leq x}p\right)T(x)\longrightarrow\infty.
\]
The logarithms in these approximants do not disappear; the
proposition proves that finite algebraic recombinations cannot
make them disappear.  The display is therefore an audit of this
specific cutoff strategy, not a universal lower bound on every
possible approximation to \(B_1\).  It leaves both the infinite
prime tail and its interaction with \(\gamma\) unresolved, so
it supplies no irrationality claim for \(B_1\).


\section{The exact multiplicative depth of a prime-grid gamma isolation}
\label{p09:gpf:section}

The five individual targets in this section are
\[
 \Gamma(1/5),\quad\Gamma(2/5),\quad
 \Gamma(1/7),\quad\Gamma(2/7),\quad\Gamma(3/7).
\]
None is proved irrational or transcendental here. The earlier project
argument established that reflection identities and the known
transcendence-degree bound for the fifth-denominator values do not
isolate either individual value. We now replace that qualitative
obstruction by a sharp calculation: in a precisely specified
reflection-and-recurrence model, isolating a fifth-denominator gamma
power needs at least three beta factors, and isolating a
seventh-denominator gamma power needs at least five. The bounds are
attained by explicit formulas. The result is a theorem about a
certificate model, not a lower bound for all possible transcendence
proofs.

\subsection{What the classical period theorems do and do not supply}

The analytic subgroup theorem and its applications to Fermat
Jacobians give very strong \emph{linear} information about beta
values. Wolfart--W\"ustholz determine the algebraic-coefficient linear
relations among rational beta values through the classical
Deligne--Koblitz--Ogus relations
\cite{p09:gpf:WolfartWustholz,p09:gpf:ABP}. This is different from determining
all polynomial relations among gamma values. The positive-characteristic
gamma analogue does have substantially stronger known algebraic
independence results; transferring those results to Euler's gamma
function is not legitimate. The distinction is explicit in
\cite[Sections 1.4.3--1.4.4]{p09:gpf:ABP}.

For fifth denominators, the known lower bound
\[
 \operatorname{trdeg}_{\mathbb Q}
 \mathbb Q\bigl(\pi,\Gamma(1/5),\Gamma(2/5)\bigr)\geq2
\]
does not identify which gamma value must supply transcendence;
see \cite[Section 3.3.2]{p09:gpf:Waldschmidt}.
Sheth's 2026 paper explicitly records that the transcendence of
$\Gamma(1/5)$ remains open \cite[discussion after Conjecture 4.2]{p09:gpf:Sheth}.
Our primary-source search did not locate a resolution of the other
four specified individual questions. Thus they remain unresolved
\emph{in this investigation}; the statement is not a claim to have
exhaustively certified the entire current literature.

\subsection{A formal model with a verifiable interpretation}

Let $p\geq5$ be prime and $h=(p-1)/2$. Work in the Laurent polynomial
ring
\[
 \mathcal R_p=\overline{\mathbb Q}
       [T^{\pm1},X_1^{\pm1},\ldots,X_h^{\pm1}].
\]
The variable $T$ represents $\pi$, and $X_j$ represents
$\Gamma(j/p)$. Define formal gamma values for $1\leq a<p$ by
\begin{equation}\label{p09:gpf:formal-gamma}
 \mathcal G_a=
 \begin{cases}
 X_a,&a\leq h,\\
 \displaystyle\frac{T}{\sin(\pi a/p)X_{p-a}},&a>h.
 \end{cases}
\end{equation}
The sine coefficients are nonzero algebraic numbers. Also set
$\mathcal G_p=1$, and, for $p<a<2p$, set
$\mathcal G_a=((a-p)/p)\mathcal G_{a-p}$. Finally put
\begin{equation}\label{p09:gpf:formal-beta}
             \mathcal B_{a,b}
                =\frac{\mathcal G_a\mathcal G_b}{\mathcal G_{a+b}}
                  \qquad(1\leq a,b<p).
\end{equation}
Every identity in this ring becomes a true identity of gamma and
beta values after the indicated substitution. The converse is not
assumed: an arithmetic identity not implied by these formal rules
could exist. Standard gamma and beta identities used here are
\cite[Sections 5.5 and 5.12]{p09:gpf:DLMF}.

For $m\in\mathbb Z$ and $1\leq j\leq h$, define $\lambda_p(j,m)$
to be the smallest value of $\sum_{a,b}|n_{a,b}|$ in a formal
identity
\begin{equation}\label{p09:gpf:certificate}
 X_j^m=cT^s\prod_{1\leq a,b<p}\mathcal B_{a,b}^{\,n_{a,b}},
 \qquad c\in\overline{\mathbb Q}^{\times},\quad
 s,n_{a,b}\in\mathbb Z.
\end{equation}
If no identity exists, set $\lambda_p(j,m)=\infty$.
Thus powers of $\pi$ and algebraic factors have zero cost; beta
factors and their reciprocals each have unit cost, with multiplicity.
This convention is part of the definition. Arbitrary rational
powers and beta values with other denominators are not included.

\begin{theorem}[Sharp gamma-isolation depth]\label{p09:gpf:depth}
For every prime $p\geq5$, $1\leq j\leq(p-1)/2$, and
$m\in\mathbb Z$,
\[
 \boxed{\displaystyle
 \lambda_p(j,m)=
 \begin{cases}
 |m/p|(p-2),&p\mid m,\\
 \infty,&p\nmid m.
 \end{cases}}
\]
In particular, the first nonzero isolatable power is the $p$th
power, and its exact beta-factor cost is $p-2$.
\end{theorem}

\begin{proof}
First discard the power of $T$ and the algebraic coefficient of
each Laurent monomial and retain its exponent vector in
$\mathbb Z^h$. Write $e_1,\ldots,e_h$ for the standard basis and
define, with indices reduced modulo $p$,
\[
 v(0)=0,\qquad
 v(a)=\begin{cases}e_a,&1\leq a\leq h,\\
                  -e_{p-a},&h<a<p.
       \end{cases}
\]
Then the exponent vector of a beta factor is
\begin{equation}\label{p09:gpf:beta-vector}
                 b(a,b)=v(a)+v(b)-v(a+b).
\end{equation}
The homomorphism
\[
 \chi:\mathbb Z^h\longrightarrow\mathbb Z/p\mathbb Z,
 \qquad\chi(e_a)=a\pmod p
\]
satisfies $\chi(v(a))=a$ and consequently annihilates every
$b(a,b)$. If \eqref{p09:gpf:certificate} holds, then $mj=0$ modulo
$p$. Since $p$ is prime and $j\ne0$ modulo $p$, necessarily
$p\mid m$.

For the sharp metric lower bound, let $t$ be the inverse of $j$
modulo $p$, and let $r_t(a)\in\{1,\ldots,p-1\}$ be the least
positive residue of $ta$ for $a\not\equiv0\pmod p$. Define the
real linear functional $L_t$ on $\mathbb R^h$ by
\[
              L_t(e_a)=1-\frac{2r_t(a)}p.
\]
The identity $r_t(p-a)=p-r_t(a)$ gives
\[
 L_t(v(a))=1-\frac{2r_t(a)}p\quad(a\ne0),
                    \qquad L_t(v(0))=0.
\]
If $a+b\ne0$ modulo $p$, the difference
$r_t(a)+r_t(b)-r_t(a+b)$ is either $0$ or $p$.
Hence
\[
             L_t(b(a,b))\in\{-1,1\}.
\]
If $a+b=0$ modulo $p$, the same expression is $0$.
Thus $|L_t(b(a,b))|\leq1$ in every case, whereas
$L_t(e_j)=1-2/p$. Applying $L_t$ to the exponent-vector
identity in \eqref{p09:gpf:certificate} yields
\begin{equation}\label{p09:gpf:lower}
        \sum_{a,b}|n_{a,b}|\geq |m|\frac{p-2}{p}.
\end{equation}

For the matching construction, put
$r_k\in\{1,\ldots,p-1\}$ equal to the least positive residue of
$kj$, for $1\leq k\leq p-1$. Thus $r_1=j$ and
$r_{p-1}=p-j$. For $1\leq k\leq p-2$, write
\[
 r_k+j=r_{k+1}+p\varepsilon_k,\qquad
                         \varepsilon_k\in\{0,1\}.
\]
Recurrence and telescoping give
\[
 \prod_{k=1}^{p-2}\mathcal B_{r_k,j}
 =\frac{\mathcal G_j^{p-1}}{\mathcal G_{p-j}}
       \prod_{k=1}^{p-2}
                  \left(\frac{p}{r_{k+1}}\right)^{\varepsilon_k}.
\]
Reflection then proves the formal and actual identity
\begin{equation}\label{p09:gpf:optimal-identity}
 \boxed{\displaystyle
 \Gamma(j/p)^p
 =\pi C_{p,j}\prod_{k=1}^{p-2}B(r_k/p,j/p),\qquad
 C_{p,j}=\frac1{\sin(\pi j/p)}
             \prod_{k=1}^{p-2}
                \left(\frac{r_{k+1}}p\right)^{\varepsilon_k}.}
\end{equation}
The coefficient $C_{p,j}$ is positive algebraic. This uses exactly
$p-2$ beta factors. Raising the identity to an integer power,
or taking its reciprocal, attains \eqref{p09:gpf:lower} for all
$m$ divisible by $p$. The case $m=0$ is the empty product.
\end{proof}

\begin{remark}[The model boundary is essential]
Theorem~\ref{p09:gpf:depth} does not prove that every possible expression
for $\Gamma(j/p)$ requires $p-2$ periods. It treats integer products
of the specified beta values under the stated formal relations.
It neither assumes nor proves the Rohrlich--Lang conjecture.
A theorem about nonlinear period products, an identity involving
other denominators, or a direct Diophantine approximation argument
could supply genuinely additional information.
\end{remark}

\subsection{A finite-degree non-detection theorem}

The lower bound also describes exactly one way in which low-degree
period tests fail to see whether a specified gamma coordinate has
become algebraic. For an integer $d\geq0$, let
$\mathcal V_d\subset\mathcal R_p$ be the
$\overline{\mathbb Q}[T^{\pm1}]$-module spanned by products of at
most $d$ formal beta factors or reciprocals. It contains constants.
Equivalently, its exponent support in the $X$ variables is contained
in the radius-$d$ word ball for the generators \eqref{p09:gpf:beta-vector}.

\begin{theorem}[Low-degree specialization preserves the formal tests]
\label{p09:gpf:lowdegree}
Fix $j$ and any nonzero algebraic number $\alpha$. If
$2d<p-2$, then substitution $X_j=\alpha$ is injective on
$\mathcal V_d$.
In particular, this is true for linear beta tests when $p\geq5$,
and for tests using at most two beta factors per monomial when
$p\geq7$.
\end{theorem}

\begin{proof}
Two distinct $X$-exponent vectors $u,w$ in this word ball can
coalesce after substitution only if
$u-w=m e_j$ with $m\ne0$. Their difference has a representation
using at most $2d$ beta generators or their negatives. The
congruence and lower-bound arguments in the proof of
Theorem~\ref{p09:gpf:depth} give $p\mid m$ and
$2d\geq|m/p|(p-2)\geq p-2$, a contradiction.
Thus distinct $X$-monomials remain distinct after the substitution.
Multiplication of their coefficients by the nonzero algebraic
powers of $\alpha$ cannot make a nonzero coefficient vanish.
The coefficients remain Laurent polynomials in the independent
variable $T$, proving injectivity.
\end{proof}

For a numerical interpretation one may set $T$ equal to the
actual $\pi$, set $X_j$ to any positive algebraic $\alpha$, and
choose the other $X$ variables positive and algebraically
independent over $\overline{\mathbb Q}(\pi)$. Then every nonzero
formal test in $\mathcal V_d$ remains nonzero. Such tuples exist
by countability. They are not alternative gamma functions and do
not match every analytic property of gamma. They show the
insufficiency of precisely this collection of formal finite-degree
tests. In the fifth case they can also satisfy the known
transcendence-degree lower bound quoted above.

\subsection{The five named targets and their optimal certificates}

For $p=5$, write $B_{a,b}=B(a/5,b/5)$. Formula
\eqref{p09:gpf:optimal-identity} gives
\begin{align}
 \Gamma(1/5)^5
   &=\frac{\pi}{\sin(\pi/5)}B_{1,1}B_{2,1}B_{3,1},
       \label{p09:gpf:fifth-one}\\
 \Gamma(2/5)^5
   &=\frac{\pi}{5\sin(2\pi/5)}B_{2,2}B_{4,2}B_{1,2}.
       \label{p09:gpf:fifth-two}
\end{align}
Each uses three factors, and three is minimal in the model.
For example, reflection also reduces
$B_{3,1}=[\sin(\pi/5)/\sin(2\pi/5)]B_{1,1}$.
Thus the first identity is a cubic product of two beta periods,
not a new linear-period relation. Transcendence of each factor
separately would not decide transcendence of the product after
multiplication by $\pi$.

For $p=7$, write $B_{a,b}=B(a/7,b/7)$. The three optimal
identities are
\begin{align}
 \Gamma(1/7)^7
  &=\frac{\pi}{\sin(\pi/7)}
       B_{1,1}B_{2,1}B_{3,1}B_{4,1}B_{5,1},
       \label{p09:gpf:seventh-one}\\
 \Gamma(2/7)^7
  &=\frac{\pi}{7\sin(2\pi/7)}
       B_{2,2}B_{4,2}B_{6,2}B_{1,2}B_{3,2},
       \label{p09:gpf:seventh-two}\\
 \Gamma(3/7)^7
  &=\frac{2\pi}{49\sin(3\pi/7)}
       B_{3,3}B_{6,3}B_{2,3}B_{5,3}B_{1,3}.
       \label{p09:gpf:seventh-three}
\end{align}
Each uses five factors, and five is minimal. In all five identities
the equivalence
\[
 \Gamma(j/p)\in\overline{\mathbb Q}
 \quad\Longleftrightarrow\quad
 \pi\prod_{k=1}^{p-2}B(r_k/p,j/p)\in\overline{\mathbb Q}
\]
is exact, because the omitted coefficient is nonzero algebraic
and an algebraic $p$th power has an algebraic $p$th root.
This equivalence relocates the missing statement; it does not
prove it. For irrationality, even this final algebraicity
equivalence is too coarse: an algebraic irrational gamma value
would remain an unresolved rationality target.

\begin{center}
\begin{tabular}{llll}
\hline
Target & First power & Minimal beta cost & Outcome here\\
\hline
$\Gamma(1/5)$ & $5$ & $3$ & individual question unresolved\\
$\Gamma(2/5)$ & $5$ & $3$ & individual question unresolved\\
$\Gamma(1/7)$ & $7$ & $5$ & individual question unresolved\\
$\Gamma(2/7)$ & $7$ & $5$ & individual question unresolved\\
$\Gamma(3/7)$ & $7$ & $5$ & individual question unresolved\\
\hline
\end{tabular}
\end{center}

The exact formulas and vector calculations were checked with
rational arithmetic for every prime $5\leq p\leq97$ and every
$1\leq j\leq(p-1)/2$. These checks supplement the uniform proof;
they are not substitutes for it. For $p=5,7$, an exhaustive search
in the finite exponent word balls independently confirmed the
costs $3$ and $5$.

\subsection{What can be recycled next}

The elliptic quantile classifier from the preceding manuscript
uses a one-dimensional algebraic group: a linear relation between
elliptic logarithms can force an endomorphism relation and then
torsion. A higher-dimensional Fermat Jacobian supplies beta
periods, but normalizing or solving a quantile problem does not
turn the products in \eqref{p09:gpf:fifth-one}--\eqref{p09:gpf:seventh-three}
into linear forms. Taking logarithms is not a remedy: the resulting
logarithms are of beta values not known to be algebraic, so Baker's
theorem on logarithms of algebraic numbers cannot be invoked.

The useful new output is therefore an exact choice of the next
missing arithmetic statement. A direct attack on fifth gamma
values must go beyond the tested linear beta information; the
seventh case still escapes even the degree-two formal tests of
Theorem~\ref{p09:gpf:lowdegree}. The proof gives a reusable lower-bound
certificate for avoiding repeated attempts to extract one gamma
coordinate from inadequate multiplicative data. Historical
priority for this certificate formulation has not been established.


\section{Named targets and the remaining arithmetic step}
\label{p09:pfi:ledger}

The following compact ledger separates the fourteen named-constant
targets from the family theorems proved above.  The target is stated
explicitly: an irrationality result does not settle a transcendence
target.  ``Unresolved here'' describes this investigation and is not
a claim to have exhaustively classified every external preprint.
Kurepa's left-factorial conjecture is a fifteenth, related target,
rather than an additional constant.

\begingroup
\small
\setlength{\tabcolsep}{3pt}
\begin{longtable}{L{28mm}L{45mm}L{76mm}}
\toprule
Target & Transfer attempted & Exact remaining step\\
\midrule\endfirsthead
\toprule
Target & Transfer attempted & Exact remaining step\\
\midrule\endhead
$\gamma$: irrationality & Even Taylor schedules and actual
denominator clearance & Control the remaining sparse crossovers and
eliminate the auxiliary $\log n$ term; see
Theorem~\ref{p09:pg:gamma}.\\[3pt]
$\delta$: irrationality & Exact primitive Laguerre content,
prime-power trees & Establish a positive asymptotic margin after
division by the exact gcd; see \eqref{p09:pg:margin}.\\[3pt]
$G$: transcendence & Primitive full-rank determinant and basis
changes & A different kernel or a new uniform nonvanishing and
height estimate is required; rational basis changes preserve the
primitive polynomial.\\[3pt]
$\zeta(3)$: transcendence & Wilson trial functions and exact
norms & Exclude algebraic degrees beyond known irrationality;
the Wilson--Ap\'ery connection is prior work.\\[3pt]
$\zeta(5)$: transcendence & Exponential intervals with
$O(n)$ logarithmic denominator bound & Prove a genuinely stronger
arithmetic certificate; the bounds supplied here do not establish
transcendence.\\[3pt]
$\zeta(7)$: irrationality & The same explicit interval
construction & Improve the reduced-denominator constants or the
certificate enough to obtain a strict positive integerization
margin.\\[3pt]
$E_2$: transcendence & Full Lambert dilation and derivative
closure & Every proposed finite rational linear closure is
excluded by Theorem~\ref{p09:pfm:independence}; an arithmetic value
method outside this class is needed.\\[3pt]
$\sum_{n\ge1}F_n^{-1}$: transcendence & Remove the Lambert
component and retain the theta component & No nonzero finite
differential--dilation filter does this; the mixed value at the
algebraic point remains uncontrolled.\\[3pt]
$B_1$: irrationality & Algebraic recombinations of prime
cutoffs & Finite logarithmic cancellation is excluded by
Proposition~\ref{p09:pfm:mertens}; control of the infinite tail and
its interaction with $\gamma$ is absent.\\[3pt]
$\Gamma(1/5)$: transcendence & Optimal prime-grid beta
certificate & The exact formal cost is three factors for the
fifth power; the required actual period relation remains unproved.\\[3pt]
$\Gamma(2/5)$: transcendence & Its separate optimal beta
certificate & The same formal cost does not establish the
arithmetic nature of this gamma value.\\[3pt]
$\Gamma(1/7)$: transcendence & Optimal prime-grid beta
certificate & Five factors isolate the seventh power formally;
the needed numerical algebraic-independence input is missing.\\[3pt]
$\Gamma(2/7)$: transcendence & Its separate optimal beta
certificate & The explicit identity does not supply the missing
period-value transcendence step.\\[3pt]
$\Gamma(3/7)$: transcendence & Its separate optimal beta
certificate & The same limitation persists despite the proved
optimal formal depth.\\[3pt]
Kurepa conjecture & Transfer through shifted Laguerre
residue classes & The exact congruence bridge in
Proposition~\ref{p09:pg:kurepa} does not prove global nonvanishing.\\
\bottomrule
\end{longtable}
\endgroup

No new named-constant conjecture is closed by these attempts.
What changes is the content of the next required estimate.  A vague
denominator problem becomes an exact gcd and local residue problem;
a numerical entropy example becomes a minimum-alphabet theorem;
an inverse quantile construction becomes an exact degree algorithm;
and a failed additive transfer becomes a theorem excluding every
finite rational differential--dilation repair in the stated class.

\paragraph{What has and has not been verified.}
The manuscript includes proofs of the stated family classifications,
thresholds, and obstructions.  Separate reviews checked sensitive
steps, including CM twisting hypotheses, exact valuation identities,
the rational ternary entropy argument, and denominator clearance.
Small exact computations corroborate algebraic identities and
instances; they do not replace general proofs.  The assembled text
is a single TeX source with an embedded bibliography and no external
input files.  There is no Lean or comparable formal certification.
Historical novelty, completeness of the literature comparison, and
the appropriate journal level require further assessment.

% END SOURCE: KEIO_Arithmetic_Thresholds_and_Primitive_Forms.tex

\endgroup
\clearpage
\part{Sparse Certificates and Information Constants}
\label{merge:part10}
\begingroup

% BEGIN SOURCE: KEIO_Sparse_Certificates_and_Information_Constants.tex

\noindent\textbf{Editorial integration note.} 
The unresolved ternary entropy question is narrowed in Parts~\ref{merge:part11} and~\ref{merge:part12}. The latter excludes cubic probability fields. The three-jet functional result not yet obtained in this part appears with a precise derivative condition in Part~\ref{merge:part11}. The factorial-threshold cancellation count here is different from the final linear-threshold count.


\begin{quote}\small
\noindent\textit{Original part framing.}\par
\noindent\textbf{Scope and verification.}
This is a standalone successor to the earlier arithmetic-threshold
manuscript. Each use of a previous construction is restated, and the
additional conclusions are identified explicitly. The proofs use
classical external theorems, including Baker's logarithmic theorem and
orthogonal-polynomial identities. Independent proof reviews and exact
rational checks are not a formal verification of this manuscript.
The 2026 external sources report proofs of irrationality for
$\zeta(5)$ and Catalan's constant
\cite{p10:nz:Fauzan,p10:nz:Lai,p10:nz:OpenAICatalan}; their complete proofs and
Lean dependencies are not independently verified here. Our
corresponding targets are transcendence. Publication significance and historical priority
require further expert review.
\end{quote}
\section{The additional conclusions and their arithmetic limits}
\label{p10:next:introduction}

Rapid approximation, exact finite reconstruction, and functional
independence do not automatically give irrationality of a distinguished
constant. The preceding manuscript exposed several such gaps. This
successor asks which remaining assertions can actually be proved,
whether an unsuccessful transfer can be excluded throughout a larger
class, and whether the same obstruction yields arithmetic theorems in
another field.

\paragraph{From local content to global sparsity.}
The previous Laguerre content formula determined the exact common
divisor at a fixed index. Theorem~\ref{p10:ng:sparse} now proves
\[
 \#\{2\le n\le X:\log g_n\ge\varepsilon n\log n\}
                  =O_{A,\varepsilon}(\log X).
\]
This controls the sum of all local contributions without assuming
that every $p$-adic root is simple. It excludes bounded primitive
linear forms outside a logarithmically sparse set, but does not
eliminate an infinite exceptional subsequence. A different failed
route to Euler's constant is controlled by nonzero curvature:
Theorem~\ref{p10:nc:count} bounds the number of indices admitting a
bounded primitive even-Taylor crossover form by $O_B(X^{1/3})$.
The logarithmic term in those forms remains present.

\paragraph{From a transcendental entropy to an independent family.}
Theorem~\ref{p10:inf:independence} supplies the rational matrices
\[
 \rho_r=\frac12I_2+2^{-(r+3)/2}
                 \begin{pmatrix}1&1\\1&-1\end{pmatrix}
 \qquad(r\text{ an odd prime}).
\]
Their entropies and $1$ are linearly independent over
$\overline{\mathbb Q}$. This is stronger than proving each entropy
transcendental separately, but is not algebraic independence.
The channel theorem also identifies exactly how totally positive
algebraic inputs can lose only an algebraic amount of entropy.

\paragraph{A sharp finite-data obstruction.}
Theorem~\ref{p10:nmo:threshold} determines
\[
                         d(M)=M+1.
\]
Here $d(M)$ is the minimum common alphabet size for two positive
algebraic distributions with the same first $M$ power sums and
different arithmetic types of Shannon entropy: rational versus
transcendental. This is not the earlier polynomial-density degree
problem. Theorem~\ref{p10:entnext:four} gives a separate four-atom
distribution whose entries are all irrational algebraic numbers
and whose entropy is exactly one bit. The minimum for this latter
all-irrational problem remains either three or four.

\paragraph{Analytic tuning and nonlinear transfer.}
The Wilson shift theorem gives an exact exponential rate and shows
that equal degree and shift uniquely maximize it per total budget.
The Catalan construction produces strict rational intervals without
an auxiliary zeta value. These facts improve the construction but
do not supply a transcendence certificate. The Lambert results
instead address functions: they exclude nonlinear relations between
distinct differential--dilation jets, including mixed Lambert--Fibonacci
pairs, and exclude algebraic replacement by finite-level modular
functions. Appending the prime-harmonic function to such a pair gives
an algebraically independent triple over $\mathbb C(q)$
(Theorem~\ref{p10:nl:triple}). Evaluation at one algebraic point
is a further arithmetic step which is not established here.

\paragraph{Unrestricted gamma identities still have an exact cost.}
For a reduced rational argument $a/q$, Theorem~\ref{p10:gn:unrestricted}
determines the least number of beta factors in an integer beta word
representing $\Gamma(a/q)^m$, modulo algebraic factors and rational
powers of $\pi$, throughout the standard identity calculus:
\[
 \ell(m[a/q])=
 \begin{cases}|m|(q-2)/q,&q\mid2m,\\
               \infty,&q\nmid2m.
 \end{cases}
\]
The auxiliary arguments have no denominator restriction. This closes
a possible formal shortcut from the earlier prime-grid model. It
does not classify unknown numerical gamma relations outside the
standard identities.

The final ledger records what each named target gains and which
step is still missing. A new example, an obstruction theorem, an
external recent result, and a solution of a famous conjecture are
kept distinct throughout.


\section{Arithmetic information loss and independent entropy constants}
\label{p10:inf:section}

Throughout this section $H_2(p)=-\sum_i p_i\log p_i/\log2$.
All probability vectors have strictly positive entries; zero rows of a
channel are deleted. A rational channel is a matrix
$A=(A_{yi})\in\mathbb Q_{\geq0}^{m\times n}$ with
$\sum_y A_{yi}=1$. We study the signed entropy difference
\[
                 L_A(p)=H_2(p)-H_2(Ap).
\]
It is a conditional entropy when $A$ is a deterministic coarse graining,
but it need not be nonnegative for a general channel.

The prior entropy obstruction extends to a complete structural
classification for channels. Its additional input is the equality case
of relative-entropy data processing, proved below. We also obtain an
explicit infinite family of rational two-dimensional density matrices
whose entropies, together with $1$, are linearly independent over
$\overline{\mathbb Q}$. Linear independence is not algebraic
independence. Baker's logarithmic theorem is the classical arithmetic
input \cite{p07:rlog:BakerIII}; the arithmetic statements here are distinct
from the semialgebraic-level-set results of
\cite{p10:inf:BW,p10:inf:BM}. No exhaustive priority claim is made.

\subsection{Transport of logarithmic relations}

\begin{lemma}[Relations and multiplicative valuations]
\label{p10:inf:transport}
For positive algebraic numbers $a_1,\ldots,a_N$, every algebraic
linear relation among their real logarithms is an algebraic linear
combination of rational linear relations. Consequently:
\begin{enumerate}
\item If $\sigma(a_j)>0$ for all $j$ and
$\sum_j c_j\log a_j=0$, where $c_j$ are algebraic, then
$\sum_j c_j\log\sigma(a_j)=0$. The coefficients $c_j$ are
\emph{not} conjugated.
\item If $v$ is an additive valuation of a number field containing
the $a_j$, then $\sum_j c_j v(a_j)=0$.
\end{enumerate}
\end{lemma}
\begin{proof}
Choose a rational basis of the span of the logarithms. Baker's theorem
makes it independent over $\overline{\mathbb Q}$, so the two relation
spaces have the stated relation. A rational relation, after clearing
denominators, is equivalent to $\prod_j a_j^{k_j}=1$ with integers
$k_j$. Apply $\sigma$ or $v$ to this identity and then take the same
algebraic linear combination of the resulting relations. Positivity
justifies real logarithms in the first case.
\end{proof}

\subsection{A Galois criterion for exact information preservation}

Join input symbols $i,j$ by an edge when some row $y$ of $A$ has
$A_{yi}A_{yj}>0$. Let $\mathcal C(A)$ be the connected components
of this graph. Every nonzero output row belongs to exactly one
component; write $Y_C$ for those rows associated with $C$.

\begin{theorem}[Galois conjugates and channel sufficiency]
\label{p10:inf:sufficiency}
Let $p$ be a positive algebraic probability vector. If $L_A(p)$ is
algebraic and $q=\sigma(p)$ is a simultaneous positive real conjugate,
then
\[
           \frac{p_i}{q_i}=\frac{p_j}{q_j}
                 \qquad (i,j\in C,\ C\in\mathcal C(A)).
\]
Equivalently, the channel loses no relative entropy between $p$ and $q$:
\[
                         D(p\Vert q)=D(Ap\Vert Aq).
\]
Here $D$ uses natural logarithms. In particular, if the overlap graph
is connected and $p$ has a distinct positive conjugate, $L_A(p)$ is
transcendental.
\end{theorem}
\begin{proof}
Put $u=Ap$ and $v=Aq=\sigma(u)$. If $h=L_A(p)$ is algebraic,
Lemma~\ref{p10:inf:transport} applied to
\[
               -\sum_i p_i\log p_i+\sum_y u_y\log u_y
                                  -h\log2=0
\]
allows $p_i,u_y$ inside the logarithms to be replaced by $q_i,v_y$,
while retaining all coefficients. Subtraction yields
$D(p\Vert q)-D(u\Vert v)=0$.

For completeness, for each output $y$ define posterior vectors on
its positive support by
\[
 r_{i\mid y}=\frac{A_{yi}p_i}{u_y},\qquad
 s_{i\mid y}=\frac{A_{yi}q_i}{v_y}.
\]
A direct expansion gives
\[
 D(p\Vert q)-D(u\Vert v)=\sum_y u_y D(r_{\cdot\mid y}\Vert
 s_{\cdot\mid y})\geq0.
\]
The elementary logarithm inequality gives equality only when every
pair of posterior vectors agrees. This says
$p_i/q_i=u_y/v_y$ whenever $A_{yi}>0$, which is equivalent to
constancy on each connected component. For a connected graph
normalization forces the constant to be one.
\end{proof}

Call an algebraic number totally positive when all its conjugates
over $\mathbb Q$ are positive real numbers. The next theorem gives
a finite exact arithmetic test, not merely a sufficient obstruction.

\begin{theorem}[Complete classification for totally positive input]
\label{p10:inf:classification}
Suppose every $p_i$ is totally positive. For $C\in\mathcal C(A)$ put
\[
 P_C=\sum_{i\in C}p_i,\qquad r_i=\frac{p_i}{P_C}\ (i\in C),
 \qquad s_y=\sum_{i\in C}A_{yi}r_i\ (y\in Y_C).
\]
Then $L_A(p)$ is algebraic if and only if both of the following hold:
\begin{enumerate}
\item Every $r_i$ is rational.
\item For every odd prime $\ell$,
\[
 \sum_C P_C b_{C,\ell}=0,
 \qquad
 b_{C,\ell}=-\sum_{i\in C}r_i v_\ell(r_i)
                    +\sum_{y\in Y_C}s_y v_\ell(s_y).
\]
\end{enumerate}
Only finitely many primes occur. When these conditions hold,
\[
                         L_A(p)=\sum_C P_C b_{C,2}.
\]
For a connected overlap graph, algebraicity therefore forces $p$ to
be rational, and the algebraic entropy difference is rational.
\end{theorem}
\begin{proof}
Take a normal closure of the field generated by the $p_i$. Every
automorphism gives a positive conjugate. Theorem~\ref{p10:inf:sufficiency}
therefore implies $p_i/P_C=\sigma(p_i/P_C)$ for every automorphism,
so each $r_i$ is rational. The $s_y$ are then rational as well.

Both input and output have component masses $P_C$. Applying the
entropy chain rule separately to input and output cancels their
common mass entropy and gives
\[
 L_A(p)=\sum_C P_C\{H_2(r_C)-H_2(s_C)\}
       =\sum_\ell\left(\sum_C P_C b_{C,\ell}\right)
                         \frac{\log\ell}{\log2}.
\]
The logarithms of distinct primes are rationally independent by
unique factorization, hence linearly independent over $\overline{\mathbb Q}$
by Baker's theorem. Algebraicity of the last
display is therefore exactly cancellation of its odd-prime
coefficients. The displayed value follows, proving necessity and
sufficiency.
\end{proof}

\begin{corollary}[Conditional entropy]
\label{p10:inf:conditional}
For a positive joint law $p(x,y)$ whose entries are totally positive
algebraic numbers, algebraicity of $H_2(X\mid Y)$ forces every
conditional probability $p(x\mid y)$ to be rational. After this
structural test, the odd-prime test of
Theorem~\ref{p10:inf:classification} is necessary and sufficient.
\end{corollary}
\begin{proof}
Apply the theorem to the deterministic map $(x,y)\mapsto y$.
Its overlap components are exactly the fibres with fixed $y$.
\end{proof}

The conclusion concerns conditional probabilities, not the joint
probabilities themselves. For example, set
$a=(2+\sqrt2)/4$ and take the two columns of a positive $3\times2$
joint law to be
\[
           a(1/2,1/4,1/4),\qquad
           (1-a)(1/4,3/8,3/8).
\]
All entries are totally positive. This example illustrates the
structural test, but it does not pass the odd-prime test: its
conditional entropy has a nonzero $\log3$ coefficient and is
transcendental. To obtain an algebraic irrational conditional
entropy instead, use a first column with conditional distribution
$(1/2,1/4,1/8,1/8)$ and a second with $(1/4,1/4,1/4,1/4)$,
with masses $a,1-a$. The resulting strictly positive joint law has
\[
                 H_2(X\mid Y)=\frac74a+2(1-a)=2-\frac a4.
\]
Thus disconnected channels genuinely permit algebraic irrational
entropy differences, even with totally positive input.

\begin{corollary}[A cone of transcendental entropy differences]
\label{p10:inf:cone}
Let $p^{(j)}$ be finitely many positive algebraic laws, let $A_j$ be
rational channels, and let $c_j>0$ be real algebraic numbers.
Suppose a common field embedding makes all $q^{(j)}=\sigma(p^{(j)})$
positive. If
\[
                         \sum_j c_j L_{A_j}(p^{(j)})
\]
is algebraic, then every $A_j$ preserves relative entropy between
$p^{(j)}$ and $q^{(j)}$. Consequently, if just one pair has a strict
data-processing inequality, that weighted sum is transcendental.
\end{corollary}
\begin{proof}
Transport the combined logarithmic relation while retaining the
coefficients $c_j$. Subtraction gives a sum of nonnegative
data-processing gaps, weighted by the positive $c_j$, equal to zero.
Each gap must vanish.
\end{proof}

\subsection{Correlations with rational marginals}

\begin{theorem}[Marginally invisible algebraic correlations]
\label{p10:inf:mutual}
Let $p(x,y)$ be a positive algebraic joint law. If a distinct positive
conjugate $q=\sigma(p)$ has the same two marginals as $p$, then
$I_2(X;Y)$ is transcendental. More generally, for a positive law
$p(x,y,z)$, a distinct positive conjugate with the same $XZ$ and
$YZ$ marginals certifies transcendence of $I_2(X;Y\mid Z)$.
In particular, totally positive algebraic joint entries and rational
specified marginals have this property whenever the joint law has
an irrational entry.
\end{theorem}
\begin{proof}
The logarithmic formulas are
\[
 I_2(X;Y)=\frac1{\log2}\sum_{x,y}p(x,y)
                \log\frac{p(x,y)}{p_X(x)p_Y(y)}
\]
and
\[
 I_2(X;Y\mid Z)=\frac1{\log2}\sum_{x,y,z}p(x,y,z)
            \log\frac{p(x,y,z)p_Z(z)}{p_{XZ}(x,z)p_{YZ}(y,z)}.
\]
If the relevant value were algebraic, transporting bases and
subtracting would give $D(p\Vert q)=0$: all marginal logarithms
cancel by hypothesis. Thus $p=q$, a contradiction. For the final
claim, choose an automorphism moving an irrational entry; total
positivity makes the conjugate law positive and rational marginals
are fixed.
\end{proof}

\subsection{An infinite linearly independent family of quantum entropies}

\begin{theorem}[Independent entropy constants in dimension two]
\label{p10:inf:independence}
For each odd prime $r$, define the rational density matrix
\[
 \rho_r=\frac12 I_2+2^{-(r+3)/2}
                   \begin{pmatrix}1&1\\1&-1\end{pmatrix},
 \qquad h_r=-\operatorname{tr}(\rho_r\log_2\rho_r).
\]
Then $1$ and all the $h_r$ are linearly independent over
$\overline{\mathbb Q}$. In particular every nonzero finite algebraic
linear combination of the $h_r$ is transcendental.
\end{theorem}
\begin{proof}
The two eigenvalues are
\[
 \lambda_r^\pm=\frac12\pm\frac{\sqrt2}{2^{(r+3)/2}},
 \qquad
 \lambda_r^+\lambda_r^-=\frac{2^r-1}{2^{r+2}}.
\]
They are positive and lie in $K=\mathbb Q(\sqrt2)$.
For distinct primes $r,s$,
\[
          \gcd(2^r-1,2^s-1)=2^{\gcd(r,s)}-1=1.
\]
Fix a rational prime $q$ dividing $2^r-1$. It is odd. At a prime
ideal $\mathfrak q$ above $q$ where one of $\lambda_r^\pm$ has
positive valuation, the other has zero valuation: their sum is $1$,
and both are integral away from $2$. Existence follows by taking
the valuation of their product. For every other prime index $s$,
both $\lambda_s^\pm$ have zero $\mathfrak q$-valuation, by the
coprimality just proved. Also $v_{\mathfrak q}(2)=0$.

Suppose a finite relation
$b+\sum_r c_r h_r=0$ holds with algebraic coefficients. Multiplying
by $\log2$ gives
\[
 b\log2-\sum_r c_r\bigl(\lambda_r^+\log\lambda_r^+
                         +\lambda_r^-\log\lambda_r^-\bigr)=0.
\]
Apply the valuation consequence of Lemma~\ref{p10:inf:transport} at
the ideal chosen for index $r$. All terms disappear except
$-c_r\lambda_r^\pm v_{\mathfrak q}(\lambda_r^\pm)$.
Its two non-coefficient factors are nonzero, so $c_r=0$.
Repeat for every index, then obtain $b=0$.
\end{proof}

This proof needs neither an unproved independence conjecture nor
Zsigmondy's primitive-divisor theorem: the restriction to odd prime
indices supplies pairwise coprime norms by an elementary identity.
It does not prove algebraic independence of the entropy values.

\begin{corollary}[Higher-order correlations]
\label{p10:inf:parity}
For an odd prime $r$ put $t_r=\sqrt2/2^{(r+1)/2}$. For any $d\geq2$
define a law on $\{-1,1\}^d$ by
\[
        p_{r,d}(x_1,\ldots,x_d)
              =2^{-d}\left(1+t_r\prod_{i=1}^d x_i\right).
\]
Every proper marginal is uniform and rational, while every atom is
totally positive algebraic. Its total correlation and its conditional
mutual information between the first two coordinates given the
remaining coordinates both equal $1-h_r$. Consequently $1$ and
these correlation constants, as $r$ runs through odd primes, are
linearly independent over $\overline{\mathbb Q}$.
\end{corollary}
\begin{proof}
Summing over any missing coordinate cancels the parity term.
Each of the two parity classes has $2^{d-1}$ atoms and mass
$(1\pm t_r)/2=\lambda_r^\pm$, so the joint entropy is
$d-1+h_r$. Uniformity of proper marginals gives total correlation
$d-(d-1+h_r)=1-h_r$ and, using the entropy formula for conditional
mutual information, the same value for the stated conditional
mutual information. For $d=2$, the conditioning variable is empty.
The linear-independence assertion follows immediately from
Theorem~\ref{p10:inf:independence}.
\end{proof}

These are explicit arithmetic conclusions for information-theoretic
constants. They do not classify Euler's constant, Gompertz's constant,
or any gamma or zeta value. Their contribution is an exact channel
criterion and an infinite independent family, rather than a claimed
resolution of a classical named transcendence problem.


\section{A sharp limit on arithmetic inference from finitely many moments}
\label{p10:nmo:section}

For a positive probability vector $p=(p_1,\ldots,p_n)$, write
$s_k(p)=\sum_i p_i^k$. The normalization $s_1=1$ is included in the
moment count. For an integer $M\geq1$, let $d(M)$ be the least common
alphabet size of two positive algebraic probability vectors with
identical $s_1,\ldots,s_M$, one having rational Shannon entropy in
bits and the other having transcendental Shannon entropy in bits.

\begin{theorem}[Exact moment threshold]
\label{p10:nmo:threshold}
For every integer $M\geq1$,
\[
                              d(M)=M+1.
\]
The rational-entropy vector may be chosen dyadic. The other vector
may be chosen arbitrarily close to it, with its $M+1$ entries all
distinct, positive, and conjugate roots of one irreducible polynomial
over $\mathbb Q$. Thus its diagonal density matrix has rational
characteristic polynomial. The pair has identical integer-order
R\'enyi entropies of orders $2,\ldots,M$.
\end{theorem}
\begin{proof}
First suppose $n\leq M$. Newton's identities express every elementary
symmetric polynomial in $n$ variables as a rational polynomial in
$s_1,\ldots,s_n$. Equality of these power sums therefore determines
the monic polynomial $\prod_i(X-p_i)$, hence the unordered vector and
its Shannon entropy. This proves $d(M)\geq M+1$.

For the converse put $n=M+1$. Use the dyadic vector
\[
 p^{(0)}=(2^{-1},2^{-2},\ldots,2^{-(n-2)},
                         2^{-(n-1)},2^{-(n-1)})
\]
when $n\geq3$, and $p^{(0)}=(1/2,1/2)$ when $n=2$.
Its entries sum to one, and its entropy is rational because the
negative base-two logarithm of every entry is an integer. Set
\[
 f(X)=\prod_{i=1}^n(X-p_i^{(0)})\in\mathbb Z[1/2][X],
 \qquad a=2^{-(n-1)}.
\]
The root $a$ is double; the remaining roots, if present, are simple
and greater than $a$. Factor $f(X)=(X-a)^2g(X)$, where
\[
                         \operatorname{sign}g(a)=(-1)^{n-2}.
\]
Consequently, for sufficiently small nonzero $\tau$ with sign
$(-1)^{n-1}$, the polynomial $f(X)+\tau$ has $n$ distinct positive
real roots. To see this directly, the simple roots persist under a
small real perturbation, and near $a$ the equation is
$(X-a)^2=-\tau/g(X)$ with positive right-hand side. Its two roots
lie on opposite sides of $a$ and approach $a$ as $\tau\to0$.
They stay positive after shrinking the perturbation. Equivalently,
the same count follows by the intermediate value theorem on two
small intervals adjacent to $a$ and on disjoint intervals around
the simple roots. The degree then exhausts the possible roots.

Fix an odd rational prime $\ell$. Choose positive integers $b$ with
$\ell\nmid b$, tending to infinity, and put
\[
                         \tau=\frac{(-1)^{n-1}}{\ell b}.
\]
We next prove that $f+\tau$ is irreducible. Write
$f(X)=X^n+c_{n-1}X^{n-1}+\cdots+c_0$. All $c_j$ are integral at
$\ell$, since their denominators are powers of two. The reciprocal
polynomial
\[
 R(X)=\ell X^n\{f(1/X)+\tau\}
\]
has constant coefficient $\ell$, intermediate coefficients
$\ell c_{n-1},\ldots,\ell c_1$, and leading coefficient
$\ell c_0+(-1)^{n-1}/b$. Its leading coefficient is an
$\ell$-adic unit; every other coefficient is divisible by $\ell$;
and its constant coefficient is not divisible by $\ell^2$.
After division by the leading unit, Eisenstein's criterion over
$\mathbb Z_\ell$ applies. Hence $R$ is irreducible over
$\mathbb Q_\ell$, and therefore over $\mathbb Q$. Reciprocation
proves irreducibility of $f+\tau$ as well. Clearing denominators
prime to $\ell$ gives the same elementary Eisenstein argument over
the integers if preferred.

Let $p^{(b)}$ be the roots of $f+\tau$, ordered to converge to
$p^{(0)}$. The coefficient of $X^{n-1}$ is unchanged, so their sum
is one. In fact all coefficients except the constant coefficient
are unchanged. Newton's identities give
\[
                    s_k(p^{(b)})=s_k(p^{(0)})
                              \qquad(1\leq k\leq n-1=M).
\]
Each root has degree $n\geq2$, and all its conjugates are among
these positive roots. Thus the entries are totally positive and
irrational. The positive-conjugate entropy criterion
(the one-output case of Theorem~\ref{p10:inf:sufficiency}) proves that
$H_2(p^{(b)})$ is transcendental. Explicitly, the one-output channel
has output entropy zero and a connected overlap graph; an
automorphism moving one root provides the distinct positive
conjugate required by that theorem.

Finally, the integer-order R\'enyi entropy is
$\mathcal R_k(p)=(1-k)^{-1}\log_2s_k(p)$ for $k\geq2$.
The asserted equalities follow from the matching power sums.
Convergence of the roots supplies arbitrary closeness.
\end{proof}

\begin{example}[The first nontrivial spectral moment]
\label{p10:nmo:example}
For $M=2$ the dyadic vector $(1/2,1/4,1/4)$ has Shannon entropy
$3/2$ and second power sum $3/8$. The three roots of
\[
                 X^3-X^2+\frac5{16}X-\frac{23}{768}
\]
are positive, sum to one, and have the same second power sum $3/8$,
but their Shannon entropy is transcendental. This polynomial is
$(X-1/2)(X-1/4)^2+1/768$. Its sign is negative at $0$,
positive at $1/4$, negative at $5/12$, and positive at $1/2$,
so it has three positive roots in the intervening intervals.
The reciprocal Eisenstein argument uses $\ell=3$ and $b=256$.
\end{example}

\begin{corollary}[No uniform finite-moment arithmetic classifier]
\label{p10:nmo:nogo}
No fixed finite list of initial power sums can decide, for all
finite-dimensional algebraic probability vectors or spectra,
whether the Shannon entropy in bits is rational or transcendental.
For each $M$, even exact knowledge of $s_1,\ldots,s_M$ fails at
dimension $M+1$, and this dimension is optimal for such a pair.
\end{corollary}
\begin{proof}
The two vectors in Theorem~\ref{p10:nmo:threshold} produce exactly the
same supplied data and require different answers.
\end{proof}

This is a limitation of the specified observation data, not of all
algorithms. For a known dimension $n$, the first $n$ power sums do
determine the spectrum. The construction uses rational
characteristic polynomials, but does not assert that the displayed
diagonal matrices have rational entries. It does not resolve a
classical named irrationality conjecture, and historical novelty of
the sharp arithmetic formulation has not been established by an
exhaustive literature review.


\section{Four irrational probabilities with one bit of entropy}
\label{p10:entnext:section}

For a probability vector with positive entries, write
\[
 H_b(p)=-\sum_i p_i\log_b p_i,\qquad b>1.
\]
The preceding manuscript established that algebraic irrational entropy
in bits first occurs on three atoms. It also constructed five positive
irrational algebraic probabilities whose entropy is rational, without
claiming the minimum for that different problem. We reduce the latter
construction to four atoms. Its exact minimum is still either three or
four. The other results in this section classify the binary problem for
all algebraic logarithm bases, and give an exact critical-value reduction
for the remaining ternary problem. These statements use classical
Baker theory, rather than a new transcendence theorem.

\subsection{An exact four-atom certificate}

\begin{theorem}[Four irrational atoms, exactly one bit]
\label{p10:entnext:four}
Let $\varphi=(1+\sqrt5)/2$. The four positive algebraic irrational
numbers
\[
 p=\left(\frac{\varphi}{2},\frac{\varphi^{-5}}2,
                 \frac{\varphi^{-4}}2,\frac{\varphi^{-4}}2\right)
 =\left(\frac{1+\sqrt5}{4},\frac{5\sqrt5-11}{4},
                 \frac{7-3\sqrt5}{4},\frac{7-3\sqrt5}{4}\right)
\]
sum to one and satisfy $H_2(p)=1$.
\end{theorem}
\begin{proof}
The equality $\varphi^2=\varphi+1$ gives
\[
 \varphi^{-4}=5-3\varphi,\qquad
 \varphi^{-5}=5\varphi-8.
\]
Consequently
\[
 \varphi+\varphi^{-5}+2\varphi^{-4}=2,
 \qquad
 \varphi-5\varphi^{-5}-8\varphi^{-4}=0.
\]
Positivity follows either from the original expressions or from
$125>121$ and $49>45$. The nonzero $\sqrt5$ coefficients show that
every entry is irrational. For $m=(1,-5,-4,-4)$, the two displayed
identities say $\sum_i p_i=1$ and $\sum_i m_i p_i=0$. Substitution into
$\log_2p_i=-1+m_i\log_2\varphi$ gives
\[
                    H_2(p)=1-\log_2\varphi\sum_i m_i p_i=1.
\]
\end{proof}

A polynomial certificate explains the simultaneous cancellation. Put
\[
 F(t)=\frac{t+t^{-5}+2t^{-4}}2.
\]
Then
\begin{equation}\label{p10:entnext:factor}
 F(t)-1=\frac{(t^2-t-1)^2(t^2+1)}{2t^5}.
\end{equation}
Thus $F(\varphi)=1$ and $F'(\varphi)=0$. The four summands of $F$,
counting the repeated $t^{-4}/2$ separately, are exactly the four
probabilities. This is a verifiable algebraic identity, not a
high-precision numerical entropy calculation.

\begin{corollary}\label{p10:entnext:minimum}
Let $m_{\rm irr,rat}$ be the smallest number of positive algebraic
irrational probabilities with rational entropy in bits. Then
\[
                         3\le m_{\rm irr,rat}\le4.
\]
\end{corollary}
\begin{proof}
One positive probability equals one and hence is rational. The binary
case is excluded by Theorem~\ref{p10:entnext:binary} below, since the base
$2$ is not an algebraic unit. Theorem~\ref{p10:entnext:four} supplies the
upper bound.
\end{proof}

\subsection{The complete binary classification for an algebraic base}

We use Baker's theorem in the following standard form: a set of real
logarithms of positive algebraic numbers which is linearly independent
over $\mathbb Q$ remains linearly independent over
$\overline{\mathbb Q}$ \cite{p07:rlog:BakerIII}. Equivalently, every
algebraic-coefficient logarithmic relation is a linear combination of
rational-coefficient logarithmic relations. The rational-spectrum
entropy dichotomy is also recorded by Blakaj--Wolf
\cite[Theorem 2]{p10:entnext:BW}.

\begin{theorem}[Binary algebraic-base classification]
\label{p10:entnext:binary}
Let $b>1$ be algebraic and let $0<p<1$ be algebraic irrational. Then
$H_b(p,1-p)$ is algebraic if and only if there exist distinct positive
rational numbers $u,v$ such that
\[
                         p=b^{-u},\qquad 1-p=b^{-v}.
\]
In that case
\[
                  H_b(p,1-p)=up+v(1-p)
\]
is algebraic irrational, and $b$ is an algebraic unit.
\end{theorem}
\begin{proof}
Suppose $h=H_b(p,1-p)$ is algebraic. If $h$ were outside the rational
span of $1,p$, Baker's theorem applied after splitting the coefficients
with respect to the independent numbers $1,p,h$ would force
$\log b=0$. Thus $h=a+cp$ for $a,c\in\mathbb Q$. In the identity
\[
 p\log p+(1-p)\log(1-p)+(a+cp)\log b=0,
\]
splitting the independent algebraic coefficients $1,p$ gives
\[
 \log(1-p)+a\log b=0,
 \qquad \log p-\log(1-p)+c\log b=0.
\]
Hence $1-p=b^{-a}$ and $p=b^{-(a+c)}$. Positivity and $b>1$ imply
$a>0$ and $a+c>0$. Set $v=a$ and $u=a+c$. Equality $u=v$ would force
$p=1/2$, contrary to the hypothesis. The converse follows by direct
substitution, and the displayed affine expression in the irrational
number $p$ is irrational because $u-v\ne0$.

Choose a positive integer $d$ such that $m=du$ and $n=dv$ are distinct
positive integers, and put $\alpha=b^{-1/d}$. The identity
$\alpha^m+\alpha^n=1$ shows that $\alpha$ is a root of the monic
polynomial $X^{\max(m,n)}+X^{\min(m,n)}-1$. Its reciprocal is also
an algebraic integer, since the reciprocal polynomial is monic up to
sign. Thus $\alpha$, and therefore $b=\alpha^{-d}$, is an algebraic
unit.
\end{proof}

\begin{corollary}[A complete obstruction for nonunit bases]
\label{p10:entnext:nonunit}
If an algebraic $b>1$ is not an algebraic unit, every binary distribution
with irrational algebraic probabilities has transcendental entropy in
base $b$. This includes every rational base $b>1$.
\end{corollary}
\begin{proof}
The first assertion is the contrapositive of the theorem. A rational
algebraic unit is $1$ or $-1$, so no rational $b>1$ is a unit.
\end{proof}

For completeness, the rational-probability case has a different exact
criterion and must not be confused with Theorem~\ref{p10:entnext:binary}.
If $p=a/N$ with $0<a<N$, define the positive algebraic number
\[
 A=\left(\frac{N^N}{a^a(N-a)^{N-a}}\right)^{1/N}.
\]
Then $H_b(p,1-p)=\log A/\log b$ is algebraic if and only if
$A=b^r$ for some rational $r$; in that case the entropy equals $r$.
Indeed an algebraic irrational ratio would contradict the
Gelfond--Schneider theorem. The individual probabilities need not be
rational powers of $b$ in this rational-probability case.

\begin{example}[Golden and plastic bases]
For $b=\varphi$, the probabilities
$(\varphi^{-1},\varphi^{-2})$ have entropy
\[
                    H_\varphi=\frac{5-\sqrt5}{2}.
\]
Let $\rho>1$ be the real root of $X^3-X-1$. Then
\[
 \rho^{-2}+\rho^{-3}=1,
 \qquad \rho^{-1}+\rho^{-5}=1,
\]
the second identity following from $\rho^5=\rho^4+1$.
These produce two different binary distributions with algebraic
irrational entropy in base $\rho$. The irrationality of their
probabilities can be checked without a numerical approximation: if a
positive power of $\rho$ were rational, its nonreal conjugates would
have the same modulus as $\rho$, whereas their product of moduli is
$1/\rho$ and $\rho>1$. These examples show that the base-$2$ minimum
alphabet of three is sensitive to the choice of algebraic base.
\end{example}

\subsection{An exact reduction of the unresolved ternary case}

\begin{lemma}[Arithmetic dimension and logarithmic rank]
\label{p10:entnext:rank}
Let $p_1,\ldots,p_n>0$ be algebraic with $\sum_i p_i=1$, and let
$b>1$ be algebraic. Define
\[
 t=\dim_{\mathbb Q}\operatorname{span}_{\mathbb Q}\{p_1,\ldots,p_n\},
 \qquad
 r=\dim_{\mathbb Q}\operatorname{span}_{\mathbb Q}
        \{\log p_1,\ldots,\log p_n,\log b\}.
\]
If $h=H_b(p)$ is algebraic, then
\[
 h\in\operatorname{span}_{\mathbb Q}\{p_1,\ldots,p_n\},
 \qquad t+r\le n+1.
\]
\end{lemma}
\begin{proof}
If $h$ were outside the indicated rational span, split the entropy
identity with respect to a basis of that span together with $h$.
Baker's theorem would force $\log b=0$, a contradiction. Now choose a
rational basis $\theta_1,\ldots,\theta_t$ of the probability span and
expand every $p_i$ and $h$ in that basis. Splitting the entropy identity
produces $t$ rational logarithmic relations. Their coefficient rows
are independent because the submatrix formed by the $p_i$ coordinates
has row rank $t$. The rational relation space among the $n+1$
logarithms has dimension $n+1-r$, proving the inequality.
\end{proof}

\begin{theorem}[Three-monomial critical-value reduction]
\label{p10:entnext:ternary}
Suppose an algebraic ternary probability vector $p$ has at least one
irrational entry and $H_2(p)=h\in\mathbb Q$. Then there exist a
positive algebraic number $w$ and integers $m_1,m_2,m_3$, not all
zero, such that
\begin{equation}\label{p10:entnext:critical}
 p_i=2^{-h}w^{m_i},\qquad
 \sum_{i=1}^3w^{m_i}=2^h,\qquad
 \sum_{i=1}^3m_iw^{m_i}=0.
\end{equation}
Moreover $\log w$ and $\log2$ are rationally independent. Conversely,
any such algebraic data with $h\in\mathbb Q$ give a probability vector
of entropy $h$. The rational span of the probabilities and the rational
span of their logarithms together with $\log2$ both have dimension two
in the nonrational case.
\end{theorem}
\begin{proof}
In Lemma~\ref{p10:entnext:rank}, nonrationality gives $t\ge2$. If $r=1$,
each probability is a rational power of two. Such positive powers
summing to one must all have integral exponents: for a common exponent
denominator $d$, the independence of
$1,2^{1/d},\ldots,2^{(d-1)/d}$ follows from Eisenstein's criterion
for $X^d-2$, and every nonconstant basis coefficient in the sum is
nonnegative. Hence $r=1$ would make every probability rational.
Thus $r\ge2$, and $t+r\le4$ forces $t=r=2$.

Choose a positive algebraic $u$ such that $\log2,\log u$ form a
rational basis of the logarithm span. Write
$\log p_i=a_i\log2+b_i\log u$ with $a_i,b_i\in\mathbb Q$.
Baker's theorem gives
\[
                   \sum_i b_i p_i=0,
                   \qquad \sum_i a_i p_i=-h.
\]
The kernel of the rational map $(x_i)\mapsto\sum_i x_i p_i$ is
one-dimensional, and $b=(b_i)$ is a nonzero vector in it. As
$\sum_i(a_i+h)p_i=0$, there is a rational $\lambda$ with
$a_i+h=\lambda b_i$. Therefore
\[
                    p_i=2^{-h}(2^\lambda u)^{b_i}.
\]
Take a common positive denominator $d$ of the $b_i$, put
$w=(2^\lambda u)^{1/d}$, and set $m_i=db_i$. This gives
\eqref{p10:entnext:critical}. Rational independence is preserved by this
change of logarithm basis. The converse is obtained by substituting
$\log_2p_i=-h+m_i\log_2w$.
\end{proof}

The last theorem converts the missing lower bound into a precise
arithmetic problem about a stationary value of three Laurent
monomials. It does not prove that such data are impossible. In
particular, the four-atom construction proves an improved upper bound,
not the assertion $m_{\rm irr,rat}=4$. No claim that a previously named
irrationality conjecture is closed is made in this section.


% All labels and bibliography keys in this fragment begin with ng:.
\section{A global sparsity theorem for primitive Laguerre content}
\label{p10:ng:section}

The earlier analysis gave an exact formula for the common divisor of
the numerator and denominator of each Laguerre approximation to the
Euler--Gompertz constant.  It did not control that common divisor
globally as the index varied.  Here we prove a quantitative global
restriction: factorial-sized common divisors occur at only logarithmically
many indices.  The proof uses positive Jacobi matrices and the elementary
arithmetic of least common multiples, rather than a uniform hypothesis
about the simplicity of every local $p$-adic root.

\subsection{The fixed family and the arithmetic input}

Fix a positive integer $A$, and set
\[
 \delta_A=\int_0^\infty\frac{e^{-t}}{1+At}\,dt,
 \qquad Q_n=\sum_{j=0}^n j!\binom nj^2 A^j.
\]
Let $P_0=0$, $P_1=1$, $Q_0=1$, $Q_1=1+A$, and apply
\begin{equation}\label{p10:ng:rec}
 U_n=(1+(2n-1)A)U_{n-1}-(n-1)^2A^2U_{n-2}
 \qquad(n\ge2)
\end{equation}
to both sequences.  We write
\[
 g_n=\gcd(P_n,Q_n),\qquad
 q_n=Q_n/g_n,\qquad
 \ell_n=(Q_n\delta_A-P_n)/g_n.
\]
Thus $q_n$ is the actual reduced denominator and $\ell_n$ is the
corresponding primitive positive linear form.  These classical
approximants are reviewed in \cite{p10:ng:Lagarias}; recent continued-fraction
work on the same special function is discussed in \cite{p10:ng:MY}.

For self-containment, we record the two arithmetic facts needed below:
\begin{equation}\label{p10:ng:content}
 \gcd(Q_n,Q_{n-1})=1,
 \qquad
 g_n=\gcd\bigl(Q_n,((n-1)!)^2\bigr).
\end{equation}
Indeed, primes dividing $A$ never divide $Q_n$.  For a prime $p$ not
dividing $A$, Lucas's congruence gives $Q_{kp+r}\equiv Q_r\pmod p$
for $0\le r<p$.  Two consecutive zero terms would propagate backwards
under~\eqref{p10:ng:rec} to $Q_0\equiv0$, whereas $Q_0=1$.
The determinant identity
\[
 P_nQ_{n-1}-P_{n-1}Q_n=A^{2n-2}((n-1)!)^2
\]
then proves the second equality in~\eqref{p10:ng:content} by reduction
modulo $Q_n$.  These facts were established in the preceding manuscript;
the global estimates below are the additional conclusions.

\subsection{Separated indices are controlled by positive continuants}

For $1\le n<m$, let $C_{n,m}$ be the determinant of the symmetric
matrix whose rows and columns have indices $k=n+2,\ldots,m$, with
diagonal entries $1+(2k-1)A$ and off-diagonal entry $-kA$ between
$k$ and $k+1$.  The empty determinant, when $m=n+1$, is one.

\begin{lemma}[A uniform overlap bound]\label{p10:ng:overlap}
For every $1\le n<m$,
\begin{align}
 \gcd(Q_n,Q_m)&\mid C_{n,m},\label{p10:ng:divides}\\
 1\le C_{n,m}&\le ((2A+1)m)^{m-n-1}.\label{p10:ng:continuantbound}
\end{align}
\end{lemma}
\begin{proof}
Expansion of the determinant along its last row gives the same
three-term recurrence as~\eqref{p10:ng:rec}, starting with
$C_{n,n+1}=1$ and $C_{n,n+2}=1+(2n+3)A$.
Consequently
\[
                 Q_m\equiv C_{n,m}Q_{n+1}\pmod{Q_n}.
\]
Since $Q_n$ and $Q_{n+1}$ are coprime,
this congruence proves~\eqref{p10:ng:divides}.

The matrix defining $C_{n,m}$ is positive definite.  In an interior
row, its diagonal exceeds the sum of the absolute values of the two
off-diagonal entries by exactly one.  At boundary rows the excess is
larger.  Equivalently, its quadratic form is a sum of positive multiples
of squared adjacent differences, positive boundary terms, and the
Euclidean squared norm.  Its determinant is therefore positive.
Hadamard's inequality for positive definite matrices bounds the
determinant by the product of the diagonal entries.  Each is at most
$(2A+1)m$, proving~\eqref{p10:ng:continuantbound}.
\end{proof}

\begin{lemma}[Packing primitive contents]\label{p10:ng:packing}
If $1\le n_1<\cdots<n_h=N$, then
\begin{equation}\label{p10:ng:product}
 \prod_{i=1}^h g_{n_i}
 \le ((N-1)!)^2
       ((2A+1)N)^{\sum_{i<j}(n_j-n_i-1)}.
\end{equation}
\end{lemma}
\begin{proof}
For any positive integers $a_1,\ldots,a_h$,
\[
 \prod_i a_i\ \mid\
 \operatorname{lcm}(a_1,\ldots,a_h)
                      \prod_{i<j}\gcd(a_i,a_j).
\]
To verify this, order their valuations at a fixed prime as
$v_1\le\cdots\le v_h$.  The valuation on the right is
$v_h+\sum_{i=1}^{h-1}(h-i)v_i$, which is at least $\sum_i v_i$.
Apply this observation with $a_i=g_{n_i}$.
By~\eqref{p10:ng:content}, their least common multiple divides
$((N-1)!)^2$.  Each pairwise gcd divides $\gcd(Q_{n_i},Q_{n_j})$,
which is bounded by Lemma~\ref{p10:ng:overlap}.
\end{proof}

The use of positive continuants matters.  A formal polynomial
remainder would not suffice if it vanished at one of the selected
integer indices.  Positive definiteness excludes that possibility
uniformly for every gap and every positive integer $A$.

\subsection{Only logarithmically many factorial-sized cancellations}

\begin{theorem}[Global large-content sparsity]\label{p10:ng:sparse}
For every fixed positive integer $A$ and every $\varepsilon>0$,
\begin{equation}\label{p10:ng:count}
 \#\{2\le n\le X:\ \log g_n\ge\varepsilon n\log n\}
                    =O_{A,\varepsilon}(\log X).
\end{equation}
More precisely, there are constants $h\ge1$, $\eta>0$, and $N_0$
depending only on $A,\varepsilon$, such that every interval
$[N,(1+\eta)N]$ with $N\ge N_0$ contains at most $h-1$
members of this set.
\end{theorem}
\begin{proof}
It is enough to prove the assertion for $0<\varepsilon\le1$.
Choose an integer $h$ with $h\varepsilon>2$, and then choose $\eta>0$
so small that
\begin{equation}\label{p10:ng:choice}
          2(1+\eta)+\binom h2\eta<h\varepsilon.
\end{equation}
Suppose that $h$ exceptional indices $n_1<\cdots<n_h$ lie in
$[N,(1+\eta)N]$.  Their defining lower bounds give
\[
             \sum_{i=1}^h\log g_{n_i}
                         \ge h\varepsilon N\log N.
\]
On the other hand, Lemma~\ref{p10:ng:packing}, the bound
$n_j-n_i-1\le\eta N$, and Stirling's elementary upper bound yield
\[
 \sum_{i=1}^h\log g_{n_i}
 \le \left(2(1+\eta)+\binom h2\eta\right)N\log N
                                      +O_{A,h,\eta}(N).
\]
This contradicts~\eqref{p10:ng:choice} for all sufficiently large $N$.
Covering $[N_0,X]$ by $O_\eta(\log X)$ consecutive multiplicative
intervals proves~\eqref{p10:ng:count}.
\end{proof}

This is an unconditional bound for the sum of all local prime-power
contributions at almost every index, with an explicit order of
magnitude for the exceptional counting function.  It does not assert
the pointwise estimate $\log g_n=o(n\log n)$, because an infinite
exceptional subsequence is still permitted.

\subsection{Consequences for irrationality certificates}

The exact positive remainder is
\begin{equation}\label{p10:ng:remainder}
 R_n:=Q_n\delta_A-P_n
   =A^{2n}n!\int_0^\infty
          \frac{e^{-t}t^n}{(1+At)^{n+1}}\,dt>0.
\end{equation}
It follows either by Rodrigues's formula and integration by parts,
or directly from the Laguerre construction.  Only the coarse estimates
\begin{equation}\label{p10:ng:coarse}
       \log R_n=n\log n+O_A(n),\qquad
       \log Q_n=n\log n+O_A(n)
\end{equation}
are needed here.  For the first upper bound, use
$t^n/(1+At)^n\le A^{-n}$ in~\eqref{p10:ng:remainder}.
For its lower bound, restrict the integral to
$\sqrt n\le t\le\sqrt n+1$.  On that interval
\[
 \left(\frac{At}{1+At}\right)^n
       \ge \exp(-C_A\sqrt n),
 \qquad \frac{e^{-t}}{1+At}\ge
          \frac{\exp(-\sqrt n-1)}{1+A(\sqrt n+1)},
\]
and Stirling's formula completes the estimate.
For $Q_n$, its top summand gives $Q_n\ge n!A^n$, while
\[
 \frac{Q_n}{n!A^n}
  =L_n(-1/A)=\sum_{k=0}^n\binom nk\frac{A^{-k}}{k!}
  \le\sum_{k=0}^\infty\frac{(n/A)^k}{(k!)^2}
  \le e^{2\sqrt{n/A}}.
\]

\begin{corollary}[Primitive forms are large outside a logarithmic set]
\label{p10:ng:forms}
Fix $A\ge1$ and $0<\varepsilon<1$.  Apart from
$O_{A,\varepsilon}(\log X)$ indices $n\le X$ and a fixed finite set,
both $q_n$ and $\ell_n$ lie in
\[
 \left[\exp((1-\varepsilon)n\log n),\,
       \exp((1+\varepsilon)n\log n)\right].
\]
In particular, for every fixed $B>0$,
\begin{equation}\label{p10:ng:boundedcount}
          \#\{n\le X:\ 0<\ell_n\le B\}
                                    =O_{A,B}(\log X).
\end{equation}
\end{corollary}
\begin{proof}
Apply Theorem~\ref{p10:ng:sparse} with $\varepsilon/2$, and combine
$\log\ell_n=\log R_n-\log g_n$ and
$\log q_n=\log Q_n-\log g_n$ with~\eqref{p10:ng:coarse}.
The upper bounds follow from $g_n\ge1$.
\end{proof}

There is also a concrete lacunarity constraint on every possible
successful subsequence.  Set
\[
                    \lambda_*:=\frac{1+\sqrt{37}}6>1.
\]

\begin{theorem}[Lacunarity of bounded primitive forms]\label{p10:ng:lacunary}
Fix $A\ge1$ and $B>0$.  If $a<b<c$ are indices with
$\ell_a,\ell_b,\ell_c\le B$, then
\begin{equation}\label{p10:ng:triple}
                    3a+b\le3c+O_{A,B}(c/\log c).
\end{equation}
If $n_1<n_2<\cdots$ is an infinite sequence of such indices, then
for every $1<\lambda<\lambda_*$ there is $C>0$ such that
$n_j\ge C\lambda^j$.  Consequently
\[
 \limsup_{X\to\infty}
 \frac{\#\{n\le X:\ell_n\le B\}}{\log X}
             \le\frac1{\log\lambda_*}<6.028.
\]
The displayed constant is a proved upper bound, not a claim of optimality.
\end{theorem}
\begin{proof}
Boundedness of $\ell_n$ and~\eqref{p10:ng:coarse} imply
$\log g_n\ge n\log n-O_{A,B}(n)$.  Apply
Lemma~\ref{p10:ng:packing} to $a,b,c$.  The sum of the three
gap exponents is
$(b-a-1)+(c-a-1)+(c-b-1)=2(c-a)-3$.
Thus
\[
 a\log a+b\log b+c\log c
                   \le(4c-2a)\log c+O_{A,B}(c).
\]
For $0<u\le c$, $u\log(u/c)\ge-c/e$.  Replacing
the logarithms on the left by $\log c$ therefore changes the
bound by only $O(c)$, giving~\eqref{p10:ng:triple}.

For every $\tau>0$, the consecutive terms of the selected
subsequence eventually satisfy
\[
               (3+\tau)n_{j+2}\ge3n_j+n_{j+1}.
\]
The positive root $\lambda_\tau$ of
$(3+\tau)\lambda^2-\lambda-3=0$ tends to $\lambda_*$
as $\tau\downarrow0$.  Induction from two sufficiently late
initial terms gives $n_j\ge C_\tau\lambda_\tau^j$.
The asserted growth and counting bound follow.
\end{proof}

\subsection{What has and has not been closed}

For $A=1$, Theorems~\ref{p10:ng:sparse} and~\ref{p10:ng:lacunary}
close a definite global subproblem left by the preceding content
formula: large cancellation cannot occur at a positive-density set,
or even at more than logarithmically many indices up to $X$.
Any vanishing primitive-form proof of the irrationality of $\delta$
using this fixed classical family must select a geometrically sparse
subsequence.  The results neither exclude such a subsequence nor
construct one, and therefore do not prove $\delta$ irrational.

The earlier Kurepa equivalence is not resolved by a sparse-index
bound.  For any fixed prime $p$ with
$p\mid\sum_{j=0}^{p-1}j!$, congruence preservation makes $p$
divide the content along a whole residue class, but that contributes
only a fixed local valuation at typical indices.  This is entirely
compatible with Theorem~\ref{p10:ng:sparse}.  No conclusion about
the universal nonvanishing of the left factorial is drawn.

For Euler's constant $\gamma$, the separate Taylor-crossover route
still requires exponentially accurate cancellation and removal of
the remaining $\log n$ term.  The present theorem concerns the
Laguerre family for $\delta_A$; transferring it to the Taylor family
without a suitable common recurrence and content divisor would be
unjustified.



% All labels and bibliography keys in this fragment begin with nc:.
\section{Curvature turns a Taylor density-zero obstruction into a counting bound}
\label{p10:nc:section}

The even Taylor truncations considered in the preceding manuscript
could succeed only at an exponentially accurate crossover.  Their
possible bounded primitive forms were shown to have density zero.
Here a lattice-point argument gives a power bound for their number,
without an irrationality-measure estimate for the crossover slope.

For positive integers $n,M$, write
\[
 S_M(n)=\sum_{k=1}^M\frac{(-1)^{k-1}n^k}{k\,k!}
       =\frac{a_M(n)}{b_M(n)},
 \qquad \gcd(a_M(n),b_M(n))=1,
 \quad b_M(n)>0.
\]
For real $x>0$ and $u>0$, define
\begin{align*}
 E_u(x)&=\frac1{\Gamma(u+1)}\int_0^x\frac1t
                    \int_0^t e^{-s}(t-s)^u\,ds\,dt,\\
 T_x&=e^{-x}\int_0^\infty\frac{e^{-t}}{x+t}\,dt.
\end{align*}
Taylor's integral remainder gives, for even integers $M$,
\begin{equation}\label{p10:nc:identity}
                S_M(n)-\log n-\gamma=T_n-E_M(n).
\end{equation}
Let $\kappa>e$ be the unique solution of
$\kappa(\log\kappa-1)=1$.

\begin{theorem}[A quantitative bound for bounded crossover forms]
\label{p10:nc:count}
For every fixed $B>0$,
\begin{equation}\label{p10:nc:main}
 \#\left\{n\le X:\ \exists M\ge2\text{ even with }
       b_M(n)|S_M(n)-\log n-\gamma|\le B\right\}
                            =O_B(X^{1/3}).
\end{equation}
This statement neither proves that the set is finite nor removes
the $\log n$ term from the displayed forms.
\end{theorem}

\subsection{An exact smooth crossover with controlled curvature}

Set $F(x,c)=\log(E_{cx}(x)/T_x)$, with $c$ restricted to a compact
subinterval of $(0,\infty)$.  The following differentiated estimate
is the analytic input.

\begin{lemma}[A differentiated boundary expansion]\label{p10:nc:expansion}
Uniformly for $c$ in such a compact interval,
\begin{equation}\label{p10:nc:F}
 F(x,c)=x f(c)-\tfrac12\log x+g(c)+r(x,c),
\end{equation}
where
\[
 f(c)=1+c(1-\log c),\qquad
 g(c)=-\tfrac32\log c-\log(1+c)-\tfrac12\log(2\pi),
\]
and, for $0\le a,b\le2$,
\begin{equation}\label{p10:nc:derivatives}
           \partial_x^a\partial_c^b r(x,c)=O(x^{-1-a}).
\end{equation}
\end{lemma}
\begin{proof}
The change of variables $t=x-v$ gives
\[
 E_{cx}(x)=\frac{x^{cx-1}}{\Gamma(cx+1)}J(x,c),
 \qquad
 J(x,c)=\int_{\substack{v,s\ge0\\v+s\le x}}
       \frac{e^{-s}(1-(v+s)/x)^{cx}}{1-v/x}\,dv\,ds.
\]
The limiting integral is
\[
 J_0(c)=\int_0^\infty\int_0^\infty
               e^{-cv-(1+c)s}\,dv\,ds=\frac1{c(1+c)}.
\]
We verify, with the same range of derivatives,
\begin{equation}\label{p10:nc:J}
 \partial_x^a\partial_c^b(J(x,c)-J_0(c))=O(x^{-1-a}).
\end{equation}
On $v+s\le x/2$, put $y=v+s$ and use
\[
 cx\log(1-y/x)=-cy-\frac{cy^2}{2x}+O(y^3/x^2).
\]
Differentiation at fixed $c$ shows that the first and second
$x$ derivatives of the exponent beyond $-cy$ are bounded by
$C(y^2+y^3)x^{-2}$ and $C(y^2+y^4)x^{-3}$,
respectively, on this range.  The reciprocal factor $1/(1-v/x)$
has analogous bounds with powers of $v$.  Up to two $c$
derivatives only introduce additional polynomial factors in $y$.
The inequalities $\log(1-y/x)\le-y/x$ and
$c\ge c_0>0$ provide an integrable dominating function, a fixed
polynomial in $v+s$ times $e^{-c_0(v+s)}e^{-s}$.
The mean-value estimate for the exponential then proves
\eqref{p10:nc:J} on the inner range.

On the complementary range, write $z=1-y/x$.
The denominator is at most $z^{-1}$, while every derivative of
order under consideration introduces only finitely many additional
powers of $z^{-1}$, $x$, and $|\log z|$.
Split $z^{cx}=z^{cx/2}z^{cx/2}$.  The first factor is at most
$e^{-c_0x/4}$ when $y\ge x/2$; for all sufficiently large $x$,
the second absorbs all the stated singular powers and logarithms.
Thus this range, its differentiated integrals, and the corresponding
tails of $J_0$ are exponentially small.  More explicitly, extend the
integrand by zero to the fixed quadrant $v,s\ge0$.  For sufficiently
large $x$, its required derivatives vanish at $v+s=x$, since they
are bounded there by a fixed power of $x$ times
$z^{cx-C}(1+|\log z|)^C$.  Differentiation over the fixed quadrant
is therefore legitimate.  Contributions at the artificial splitting
boundary $v+s=x/2$ cancel, or can individually be absorbed in the
exponentially small error.  This proves~\eqref{p10:nc:J}.

Also
\[
 T_x=\frac{e^{-x}}x K(x),\qquad
 K(x)=\int_0^\infty\frac{e^{-t}}{1+t/x}\,dt=1+O(x^{-1}),
\]
where the first two differentiated remainders are
$O(x^{-2})$ and $O(x^{-3})$ by differentiation under the integral.
Finally, the usual differentiated Stirling expansion for
$\log\Gamma(cx+1)$ has an $O(x^{-1})$ remainder with the
derivative bounds in~\eqref{p10:nc:derivatives}; this also follows
directly by differentiating Binet's integral formula.
Combining these estimates, and using that $J_0(c)$ is bounded
away from zero on the chosen compact interval, proves the lemma.
\end{proof}

\begin{lemma}[The exact crossover curve]\label{p10:nc:curve}
For all sufficiently large real $x$, there is a unique $u(x)$ in
$(2x,4x)$ such that $E_{u(x)}(x)=T_x$.  It is smooth and satisfies
\begin{align}
 u(x)&=\kappa x-\frac{\log x}{2\log\kappa}+O(1),
                                                    \label{p10:nc:center}\\
 u''(x)&=\frac1{2\log\kappa\,x^2}
                    +O\left(\frac{(\log x)^2}{x^3}\right).
                                                    \label{p10:nc:curvature}
\end{align}
\end{lemma}
\begin{proof}
In Lemma~\ref{p10:nc:expansion}, $f(2)>0$, $f(4)<0$, and
$f'(c)=-\log c<0$ on $[2,4]$.  Thus $F$ has exactly one zero
there for all sufficiently large $x$, and its derivative with
respect to the real order $cx$ is bounded away from zero.
The implicit function theorem applies.

Write $u(x)=\kappa x+s(x)$.  Taylor expansion of~\eqref{p10:nc:F}
first yields $s(x)=O(\log x)$ and then the identity
\[
 0=f'(\kappa)s-\tfrac12\log x+g(\kappa)+\mathcal R(x,s),
 \qquad \mathcal R(x,s)=O((1+s^2)/x).
\]
For $|s|\le C\log x$, the differentiated estimates from
Lemma~\ref{p10:nc:expansion} imply
\begin{align*}
 \mathcal R_x&=O((1+s^2)/x^2),&
 \mathcal R_s&=O((1+|s|)/x),\\
 \mathcal R_{xx}&=O((1+s^2)/x^3),&
 \mathcal R_{xs}&=O((1+|s|)/x^2),&
 \mathcal R_{ss}&=O(x^{-1}).
\end{align*}
Implicit differentiation gives
\[
 s'(x)=\frac1{2f'(\kappa)x}
                 +O((\log x)^2/x^2),\qquad
 s''(x)=-\frac1{2f'(\kappa)x^2}
                 +O((\log x)^2/x^3).
\]
Since $f'(\kappa)=-\log\kappa$, these are the asserted formulas.
\end{proof}

\subsection{Bounded arithmetic forms must track the curve exponentially}

\begin{lemma}\label{p10:nc:tracking}
For every $B>0$ there are constants $c_B>0$ and $n_B$ such that,
whenever $n\ge n_B$ and an even $M\ge2$ satisfies the inequality
in~\eqref{p10:nc:main}, one has
\begin{equation}\label{p10:nc:trackingeq}
                         |M-u(n)|\le e^{-c_B n}.
\end{equation}
\end{lemma}
\begin{proof}
We give the arithmetic reduction explicitly.  For $M\le2n$,
restricting the defining integral for $E_M(n)$ to
$n-n/(M+1)\le t\le n$ and $0\le s\le1$ gives
\[
 E_M(n)\ge e^{-13}\frac{n^M}{(M+1)M!}\qquad(n\ge4).
\]
This tends to infinity uniformly for $2\le M\le2n$:
use $M!\le M^M$ when $M\le\sqrt n$, and Stirling's bound
with $en/M\ge e/2>1$ on the remaining range.
Since $T_n\to0$ and $b_M(n)\ge1$, bounded forms cannot have
such orders once $n$ is large.

For $M\ge2n$, every prime $M/2<p\le M$ with $p\nmid n$
occurs at exact valuation $-2$ in $S_M(n)$: its $k=p$ summand
has that valuation and every other summand has valuation at least
$-1$.  Therefore, with $\vartheta(y)=\sum_{p\le y}\log p$,
\[
 \log b_M(n)\ge2\{\vartheta(M)-\vartheta(M/2)\}-2\log n
                         =M+o(M)-2\log n.
\]
Jensen's inequality gives $T_n\ge e^{-n}/(n+1)$, so
$b_M(n)T_n\to\infty$ uniformly for $M\ge2n$.
It follows from~\eqref{p10:nc:identity} that
$|E_M(n)/T_n-1|\to0$ uniformly over bounded forms.

The real-order function $E_u(n)$ decreases for $u\ge2n$ when
$n$ is large: its logarithmic derivative is at most
$\log n-\psi(u+1)<0$.  Lemma~\ref{p10:nc:expansion} at $c=2,4$
thus confines $M/n$ to $[2,4]$, and the same lemma forces
$M/n\to\kappa$.  The denominator estimate now improves the
relative-error bound to
\[
 \left|\frac{E_M(n)}{T_n}-1\right|
                     \le\exp(-[\kappa-1+o(1)]n).
\]
The derivative of $\log(E_u(n)/T_n)$ with respect to $u$
is $-\log\kappa+o(1)$ near its zero, by
Lemma~\ref{p10:nc:expansion}.  The mean-value theorem proves
\eqref{p10:nc:trackingeq}, with any fixed $c_B<\kappa-1$ after
increasing the threshold $n_B$.
\end{proof}

\subsection{Integer triangles cannot follow nonzero curvature too often}

\begin{proof}[Proof of Theorem~\ref{p10:nc:count}]
On each interval $[N,2N]$, Lemma~\ref{p10:nc:curve} gives constants
$0<c_1<c_2$ independent of $N$ such that
\[
                 c_1N^{-2}\le u''(x)\le c_2N^{-2}.
\]
Choose a fixed small $\eta>0$ and set $H=\eta N^{2/3}$.
Suppose that three successful integers $n_1<n_2<n_3$ lie in
an interval of length at most $H$, and let $M_i$ be their
associated even orders.  By Lemma~\ref{p10:nc:tracking},
$|M_i-u(n_i)|\le e^{-c_BN}$.

The integer determinant
\[
 D=(n_2-n_1)(M_3-M_1)-(n_3-n_1)(M_2-M_1)
\]
differs from
\[
 D_u=(n_2-n_1)(n_3-n_1)(n_3-n_2)
                              u[n_1,n_2,n_3]
\]
by at most $4He^{-c_BN}$, where the divided difference satisfies
$u[n_1,n_2,n_3]=u''(\xi)/2$ for an intermediate point $\xi$.
Consequently $0<D_u\le c_2H^3/(2N^2)$.
Choose $\eta$ so that this upper bound is below $1/2$.
For large $N$, we then have $|D|<1$, hence $D=0$.
But the positive integer gaps give
$D_u\ge c_1N^{-2}$, whereas $4He^{-c_BN}=o(N^{-2})$;
this contradicts $D=0$.

There are therefore at most two successful integers in each
interval of length $H$.  Partitioning $[N,2N]$ gives
$O(N/H+1)=O(N^{1/3})$ successful integers.  Summation over
dyadic intervals proves~\eqref{p10:nc:main}.
\end{proof}

The curvature argument is independent of whether $\kappa$ has a
known finite irrationality measure.  It strengthens the previous
density-zero conclusion to an explicit counting estimate for this
specific failed approximation route.  It leaves open whether the
exceptional set is infinite and whether any of its forms actually
tends to zero; it does not settle the irrationality of $\gamma$.



\section{Shift optimization and positive Catalan certificates}
\label{p10:nz:section}

This section continues the Wilson trial construction, rather than
counting its earlier exponential-height estimate as a second result.
It determines the sharp exponential effect of changing the shift,
computes the leading constant in the transferred squared norms, and
removes an auxiliary-constant obstruction in the Catalan application.
None of these conclusions proves a new irrationality or transcendence
statement about the named constants.

The starting Stieltjes representation is due to Pr\'evost
\cite{p10:nz:Prevost}; Wilson remainder Pad\'e approximations were developed
further by Pr\'evost and Rivoal \cite{p10:nz:PR}.
Recent context is particularly important here: Fauzan's September 2026
preprint reports a proof of the irrationality of $\zeta(5)$ by normalized Hankel
determinants \cite{p10:nz:Fauzan}; Lai's October 2026 preprint uses the
determinant method to prove the linear independence of
$1,\zeta(3),\zeta(4)$ \cite{p10:nz:Lai}. Its revised version also records
OpenAI's reported Lean-formalized irrationality proof for Catalan's constant
\cite{p10:nz:OpenAICatalan}. Their complete proofs and Lean dependencies
are not independently verified here. These are external results, not discoveries
of this project. Our $\zeta(5)$ and Catalan targets below are
transcendence.

\subsection{The explicit family and its squared norms}
For integers $s\geq2$, let
\[
 d\nu_s(t)=w_s(y)\,dy,\qquad t=y^2,\qquad
 w_s(y)=\frac{2(2\pi)^{s-1}y^s}{(s-1)!}
                \sum_{\ell\geq1}\ell^{s-1}e^{-2\pi\ell y}.
\]
Its polynomial moments and Stieltjes transform are
\begin{align}
 m_{s,e}&=(-1)^eB_{2e+2}
                \frac{(2e+s)!}{(s-1)!(2e+2)!},\label{p10:nz:moments}\\
 S_s(z)&=\int_0^\infty\frac{d\nu_s(t)}{t+z}.\label{p10:nz:stieltjes}
\end{align}
For $a>0$ the Hurwitz formula reads
\begin{equation}\label{p10:nz:hurwitz}
 \zeta(s,a)=\frac1{2a^s}+\frac{a^{1-s}}{s-1}
                      +a^{1-s}S_s(a^2).
\end{equation}
Define
\begin{equation}\label{p10:nz:R}
 R_n(t)=\sum_{k=0}^n\frac{(-1)^k}{(k!)^2}
       \binom nk\binom{n+k+1}{k}\prod_{r=0}^{k-1}(t+r^2).
\end{equation}
The classical Wilson specialization $(0,1,1,1)$ gives
\begin{equation}\label{p10:nz:norm3}
 \int R_n^2\,d\nu_3=\frac{n+1}{4}.
\end{equation}
For the associated monic orthogonal polynomials, the recurrence
coefficients are
\[
 \beta_n=\frac{(n+1)^2(2n^2+4n+1)}{(2n+1)(2n+3)},\qquad
 \lambda_n=\frac{n^3(n+1)^3}{4(2n+1)^2}.
\]
Thus multiplication by $t$ in the orthonormal basis has diagonal
$\beta_n=n^2/2+O(n)$ and adjacent entries
$\sqrt{\lambda_n}=n^2/4+O(n)$.

\begin{theorem}[Exact leading transferred norm]\label{p10:nz:normleading}
For fixed odd $s\geq3$, put
\[
 \kappa_s=\frac{\pi^{s-3}}{2(s-1)!}
                         \binom{s-3}{(s-3)/2}.
\]
Then
\begin{equation}\label{p10:nz:normasymptotic}
 H_{s,n}:=\int R_n(t)^2\,d\nu_s(t)
       =\kappa_s n^{s-2}\bigl(1+O_s(n^{-1})\bigr).
\end{equation}
For $s=2$ one has the useful exact comparison
\begin{equation}\label{p10:nz:norm2}
                   0<H_{2,n}\leq\frac{n+1}{4}.
\end{equation}
\end{theorem}
\begin{proof}
Write $s=2k+3$. Directly from the defining densities,
\[
 \frac{w_s(y)}{w_3(y)}
    -\frac{2(2\pi)^{2k}}{(2k+2)!}y^{2k}
\]
is bounded on $(0,\infty)$. At infinity this follows by separating
the $\ell=1$ terms, whose remaining relative error decays
exponentially. At zero both densities tend to $1/\pi$; this limit
follows from the elementary Riemann-sum limit
$u^{m+1}\sum_{\ell\geq1}\ell^m e^{-u\ell}\to m!$.
Consequently the error in replacing $H_{s,n}$ by
$2(2\pi)^{2k}/(2k+2)!$ times $\int t^kR_n^2d\nu_3$ is $O_s(n)$.

Let $J$ denote the orthonormal Jacobi matrix. The identity
\[
 \int t^kR_n^2d\nu_3=\frac{n+1}{4}(J^k)_{nn}
\]
involves paths of length $k$ staying within $k$ of $n$. Each of the
finitely many coefficients on such paths, divided by $n^2$, tends
to $1/2$ on the diagonal and $1/4$ off the diagonal, with error
$O_k(n^{-1})$. Therefore
\[
 (J^k)_{nn}
 =n^{2k}[z^0]\left(\frac12+\frac{z+z^{-1}}4\right)^k
                 +O_k(n^{2k-1})
 =\frac1{4^k}\binom{2k}{k}n^{2k}+O_k(n^{2k-1}).
\]
The case $k=0$ is supplied exactly by \eqref{p10:nz:norm3}; for $k\geq1$
these formulas prove \eqref{p10:nz:normasymptotic}.
Finally geometric-series summation gives
\[
 \frac{w_2(y)}{w_3(y)}=\frac{\tanh(\pi y)}{\pi y}\leq1,
\]
which proves \eqref{p10:nz:norm2}.
\end{proof}

\subsection{The optimal shift at a fixed total budget}
For $\rho>0$ define
\[
 v=\sqrt{1+\rho^2},\qquad
 \Phi(\rho)=2\operatorname{arsinh}\rho
                    +2\rho\operatorname{arsinh}(1/\rho),
 \qquad
 C(\rho)=\frac{(v+\rho)^2}{4\pi^{3/2}\rho^{3/2}v^{1/2}}.
\]

\begin{theorem}[Uniform shift asymptotics]\label{p10:nz:shift}
Let $j,n$ be positive integers and suppose $\rho=j/n$ ranges over a
fixed compact subinterval of $(0,\infty)$. Uniformly in this range,
\begin{equation}\label{p10:nz:shiftasymptotic}
 R_n(-j^2)=C(\rho)n^{-3/2}e^{n\Phi(\rho)}
                                  (1+O(n^{-1})).
\end{equation}
For the Wilson trial intervals for $\zeta(s)$, $s\geq3$ odd,
the exact interval width is
\[
 W_{s,n,j}=\frac{H_{s,n}}{j^{s+1}R_n(-j^2)^2},
\]
and hence
\begin{equation}\label{p10:nz:widthexact}
 W_{s,n,j}=
 \frac{\kappa_s}{\rho^{s+1}C(\rho)^2}
 e^{-2n\Phi(\rho)}(1+O_s(n^{-1})).
\end{equation}
Among fixed positive shift ratios, the exponential decay rate per
total budget $n+j$ is uniquely maximized by $j/n=1$. The maximum is
$4\log(1+\sqrt2)$.
\end{theorem}
\begin{proof}
All nonzero terms in
\[
 R_n(-j^2)=\sum_{k=0}^{\min(n,j)}
 \binom nk\binom{n+k+1}{k}\binom jk\binom{j+k-1}{k}
\]
are positive. For $x=k/n$ in the interior of the summation range,
uniform Stirling expansion gives a term
$n^{-2}g_\rho(x)e^{nf_\rho(x)}(1+O(n^{-1}))$, where
\begin{align*}
 f_\rho(x)={}&(1+x)\log(1+x)-(1-x)\log(1-x)\\
             &+(\rho+x)\log(\rho+x)-(\rho-x)\log(\rho-x)
                       -4x\log x,\\
 g_\rho(x)={}&\frac{\rho(1+x)}{4\pi^2x^2}
       \sqrt{\frac{1+x}{(1-x)(\rho+x)(\rho-x)}}.
\end{align*}
The unique maximum satisfies
$(1-x^2)(\rho^2-x^2)=x^4$, so
\[
 x_*=\frac{\rho}{\sqrt{1+\rho^2}},\qquad
 f_\rho(x_*)=\Phi(\rho),\qquad
 f_\rho''(x_*)=-\frac{2(1+\rho^2)^{5/2}}{\rho^3}.
\]
For ratios in a compact set this saddle remains uniformly separated
from both endpoints. Strict concavity makes the complementary part
of the sum exponentially smaller. Taylor expansion on a shrinking
neighborhood of the saddle, followed by Gaussian summation, gives
\[
 R_n(-j^2)=n^{-3/2}
  g_\rho(x_*)\sqrt{\frac{2\pi}{|f_\rho''(x_*)|}}
  e^{n\Phi(\rho)}(1+O(n^{-1})).
\]
The Gaussian odd term cancels; the even remainder is $O(n^{-1})$.
The prefactor simplifies to $C(\rho)$, proving
\eqref{p10:nz:shiftasymptotic}. Combining it with
Theorem~\ref{p10:nz:normleading} proves \eqref{p10:nz:widthexact}.

The exponential rate per total budget is
$F(\rho)=2\Phi(\rho)/(1+\rho)$. Since
$\Phi'(\rho)=2\operatorname{arsinh}(1/\rho)$,
\[
 F'(\rho)=\frac{4}{(1+\rho)^2}
 \bigl(\operatorname{arsinh}(1/\rho)-\operatorname{arsinh}\rho\bigr).
\]
This is positive below $1$ and negative above $1$, proving the
strict optimality assertion.
\end{proof}

\begin{remark}[What parameter tuning cannot repair]\label{p10:nz:normalizer}
For integer shifts the previously proved common denominator is
\[
 Q_{s,n,j}=2(s-1)j^{s+1}L_j^s R_n(-j^2)^2 L_{2n}^4
                         (s-1)!L_{4n+s+1}^{s-1},
\]
where $L_m=\operatorname{lcm}(1,\ldots,m)$. The prime number theorem
and Theorem~\ref{p10:nz:shift} give, for $j/n\to\rho>0$,
\[
 \frac{\log Q_{s,n,j}}n
 =2\Phi(\rho)+4s+4+s\rho+o(1).
\]
The two occurrences of $2\Phi(\rho)$ cancel when this denominator
bound is compared with the interval width. No fixed positive shift
ratio makes this particular estimate certify irrationality.
This is a limitation of the common-denominator estimate, not a
lower bound on the reduced denominator and not an impossibility
theorem for the actual approximants.
\end{remark}

\subsection{An exact cancellation diagnostic}
For $z=j^2$, set $a=R_n(-z)$ and
$q=(a-R_n(t))/(t+z)$. Let $\mu_s$ denote integration against $\nu_s$.
The lower variational trial is exactly
\begin{equation}\label{p10:nz:defect}
 \mu_s\left(2\frac q a-(t+z)\frac{q^2}{a^2}\right)
       =\frac{\mu_s(q)}a+\frac{\mu_s(R_nq)}{a^2}.
\end{equation}
Indeed $(t+z)q=a-R_n$ proves the identity by multiplication.
For $s=3$, orthogonality makes the second term vanish, because
$\deg q<n$. This explains the loss of one power of $a$ in the
classical Gaussian Pad\'e case. For $s=5,7$ the exact rational number
$\mu_s(R_nq)$ is the obstruction to this same cancellation; one
cannot delete it on the basis of orthogonality for a different weight.
If an integer $M$ clears both moments in \eqref{p10:nz:defect}, the
remaining cancellation is determined exactly by
\[
 \gcd\bigl(aM\mu_s(q)+M\mu_s(R_nq),\;Ma^2\bigr).
\]
This identity separates orthogonality cancellation from genuine
arithmetic common factors.

\subsection{Strict Catalan intervals with linear logarithmic height}
Put $G=\sum_{v\geq0}(-1)^v/(2v+1)^2$. For any polynomial $p$ and $z>0$
define rational moment expressions
\begin{align*}
 \mathcal L_s(p;z)&=\int(2p-(t+z)p^2)\,d\nu_s,\\
 \mathcal U_s(p;z)&=\frac{m_{s,0}-\int(2tp-t(t+z)p^2)\,d\nu_s}{z}.
\end{align*}
For rational $c>0$ with $R_n(-c^2)\ne0$, put
\[
 p_{n,c}(t)=\frac{1-R_n(t)/R_n(-c^2)}{t+c^2},\quad
 L_{n,c}=\mathcal L_2(p_{n,c};c^2),\quad
 U_{n,c}=\mathcal U_2(p_{n,c};c^2).
\]
Completing squares proves
\begin{equation}\label{p10:nz:LU}
 L_{n,c}<S_2(c^2)<U_{n,c},\qquad
 U_{n,c}-L_{n,c}=\frac{H_{2,n}}{c^2R_n(-c^2)^2}.
\end{equation}

\begin{theorem}[A direct Catalan transfer]\label{p10:nz:catalan}
Let $n\geq1$, $a=n+1/4$, $b=n+3/4$, and
\[
 B_n=\sum_{v=0}^{n-1}\left(\frac1{(4v+1)^2}-\frac1{(4v+3)^2}\right)
       +\frac1{16}\left(\frac1{2a^2}+\frac1a
                              -\frac1{2b^2}-\frac1b\right).
\]
Then the rational numbers
\[
 \ell_n^G=B_n+\frac1{16}\left(\frac{L_{n,a}}a-\frac{U_{n,b}}b\right),
 \quad
 u_n^G=B_n+\frac1{16}\left(\frac{U_{n,a}}a-\frac{L_{n,b}}b\right)
\]
satisfy strict inequalities $\ell_n^G<G<u_n^G$. Moreover
\begin{align}
 u_n^G-\ell_n^G
 &=\frac{H_{2,n}}{16}
   \left(\frac1{a^3R_n(-a^2)^2}
                    +\frac1{b^3R_n(-b^2)^2}\right),\label{p10:nz:Gwidth}\\
 0<u_n^G-\ell_n^G
 &\leq \exp\bigl(-8\log(1+\sqrt2)n+O(\log n)\bigr),\label{p10:nz:Grate}\\
 \log\operatorname{den}(\ell_n^G),\quad
 \log\operatorname{den}(u_n^G)&=O(n).\label{p10:nz:Gheight}
\end{align}
No $\zeta(2)$ term or unproved determinant nonvanishing statement
occurs in these certificates.
\end{theorem}
\begin{proof}
The tail identity for the defining series and \eqref{p10:nz:hurwitz}
give
\[
 G=B_n+\frac1{16}\left(\frac{S_2(a^2)}a-\frac{S_2(b^2)}b\right).
\]
All summands defining $R_n(-a^2)$ and $R_n(-b^2)$ are positive:
the indices $k\leq n$ satisfy $k-1<a,b$. Thus the polynomials in
the statement are defined. Formula \eqref{p10:nz:LU} proves the strict
interval and its exact width. The saddle analysis in
Theorem~\ref{p10:nz:shift} also applies to real shifts $c>n$ with
$c=n+O(1)$: all terms remain positive, the same gamma-function
Stirling expansion holds, and the saddle remains interior.
Consequently $\log R_n(-c^2)=4\log(1+\sqrt2)n+O(\log n)$.
Equation \eqref{p10:nz:norm2} proves \eqref{p10:nz:Grate}.

Here are explicit arithmetic details for \eqref{p10:nz:Gheight}.
For integers $x$ and $k\geq0$,
\[
 2^{2k+v_2(k!)}\binom{x/4}{k}\in\mathbb Z.
\]
For each odd prime, the numerator
$\prod_{r=0}^{k-1}(x-4r)$ contains at least as many multiples of
each prime power as $k!$; this proves the assertion after taking
valuations. The same bound holds for
$\binom{x/4+k-1}{k}$. Since $v_2(k!)\leq k$, the polynomial
\[
 F_n(x)=2^{6n}R_n(-x^2/16)
\]
is integer-valued on $\mathbb Z$, of degree $2n$.
Write $k_1=4n+1$, $k_3=4n+3$, and $A_i=F_n(k_i)$, positive integers.
For $i=1,3$, the polynomial $A_i-F_n(x)$ vanishes at both
$x=k_i$ and $x=-k_i$. The elementary divided-difference property
of integer-valued polynomials implies
\[
 A_iL_{2n}^2p_{n,k_i/4}(-x^2/16)\in\operatorname{Int}(\mathbb Z).
\]
For completeness, if an integer-valued polynomial $f$ of degree
at most $d$ vanishes at an integer $c$, expansion in
$\binom{x-c}{r}$ gives
$L_df(x)/(x-c)\in\operatorname{Int}(\mathbb Z)$, since
$\binom{x-c}{r}/(x-c)=r^{-1}\binom{x-c-1}{r-1}$.
Applying this twice proves the displayed assertion.

Replacing $x$ by $4x$ preserves integer-valuedness. Therefore, with
$E=A_1^2A_3^2L_{2n}^4$, multiplying either moment polynomial
appearing in $L_{n,c}$ or $U_{n,c}$ by $16E$ makes its substitution
$t=-x^2$ integer-valued. The factor $16$ clears $c^2=k_i^2/16$.
For the Bernoulli functional $\mathcal B(x^r)=B_r$, one has
\[
 \mu_2(P)=\mathcal B\bigl(x^2P(-x^2)\bigr),\qquad
 L_{d+1}\mathcal B(f)\in\mathbb Z
 \quad(f\in\operatorname{Int}(\mathbb Z),\ \deg f\leq d).
\]
The last fact follows from
$\mathcal B(\binom{x}{r})=(-1)^r/(r+1)$, obtained by applying
$\mathcal B$ to $(1+z)^x$. All relevant degrees are at most $4n+2$.
An explicit, deliberately nonminimal, common denominator for the
two Catalan endpoints is consequently
\begin{equation}\label{p10:nz:Gden}
 Q_n^G=2^{12}(k_1k_3)^3 L_{4n+3}^3
                        A_1^2A_3^2L_{2n}^4.
\end{equation}
The two harmonic sums are cleared by $L_{4n+3}^2$; the remaining
powers of $a,b$ are covered by $2^{12}(k_1k_3)^3$.
Since $\log A_i=(6\log2+4\log(1+\sqrt2))n+O(\log n)$ and
$\log L_m=O(m)$, this proves \eqref{p10:nz:Gheight}.
\end{proof}

\paragraph{Arithmetic scope of the four attacks.}
For $\zeta(3)$, the exact defect \eqref{p10:nz:defect} vanishes, recovering
the classical orthogonality cancellation but providing no exclusion
of algebraic degrees greater than one. For $\zeta(5)$ and $\zeta(7)$,
the shift theorem proves that a change of the degree-to-shift ratio
does not improve the exponential rate at a fixed total budget; the
common-denominator estimate still has a strictly adverse margin.
For Catalan's constant the direct construction replaces the earlier
mixed-determinant nonvanishing problem by strict positive integrals,
but \eqref{p10:nz:Gden} does not even independently recover irrationality,
let alone establish transcendence. The remaining arithmetic question in each case concerns
the actual cancellation, not numerical convergence.

\paragraph{Exact finite checks, not asymptotic proof.}
The accompanying rational-arithmetic script constructs every
coefficient and Bernoulli moment with exact fractions. It verifies
\eqref{p10:nz:defect}, the $s=3$ orthogonality cancellation, and the exact
interval widths. For diagonal degree $n=50$, the ratios
$\log\operatorname{den}(\ell_{s,n,n})/n$ were approximately
$6.1096,11.3195,13.4737$ for $s=3,5,7$, while
$-\log W_{s,n,n}/n$ was approximately $7.0171,7.0064,7.0058$.
These data support no new irrationality claim. In particular,
finite data do not prove a lower asymptotic bound for denominators.
The Catalan script additionally checks the exact interval formula and
that the displayed common denominator clears both endpoints, for
$n=1,2,3,5,10,15,20$. It encloses the target using a separate Euler
transformation with a rigorous rational tail bound $2^{-400}$.


\section{Nonlinear independence beyond the Lambert obstruction}
\label{p10:nl:section}
The previous manuscript excluded rational linear differential--dilation
relations for the full divisor Lambert series.  We now retain the constant
beside its leading logarithm.  This supplies an algebraic independence
argument that the leading logarithm alone does not give.  The conclusions
concern functions, and do not establish transcendence of their values at
an algebraic point.

Put
\[
 L(q)=\sum_{n\ge1}\frac{q^n}{1-q^n}
     =\sum_{n\ge1}d(n)q^n,\qquad
 D=q\frac{d}{dq},\qquad
 F_{m,r}(q)=D^r[L(q^m)]\quad(m\ge1,\ r\ge0).
\]
In particular, differentiation is taken \emph{after} composition.
This convention accounts for the factor \(m^{-1}\), rather than
\(m^{-r-1}\), in the next formula.

\subsection{Two terms at a rational cusp}
\begin{lemma}[A logarithmic cusp with a power-saving error]
\label{p10:nl:cusp}
Let \(\xi\) be a primitive \(N\)-th root of unity and put
\(g=\gcd(N,m)\), \(N_m=N/g\).  For fixed \(N,m,r\), as \(t\downarrow0\),
\[
 F_{m,r}(\xi e^{-t})
 =\frac{r!}{mN_m\,t^{r+1}}
  \left(\log\frac1t+\gamma-2\log N_m-\log m+H_r\right)
   +O_{N,m,r}(t^{-r-1/2}),
 \tag{NL1}\label{p10:nl:cuspformula}
\]
where \(H_0=0\).
\end{lemma}
\begin{proof}
For a primitive \(N\)-th root \(\xi\), define
\[
 Z_\xi(s)=\sum_{u,v\ge1}\frac{\xi^{uv}}{(uv)^s}
 =N^{-2s}\sum_{a,b=1}^{N}\xi^{ab}
       \zeta(s,a/N)\zeta(s,b/N)\qquad(\Re s>1).
\]
Here \(\zeta(s,x)\) is the Hurwitz zeta function.  Its expansion at \(1\) is
\(\zeta(s,x)=(s-1)^{-1}-\psi(x)+O(s-1)\).
The elementary root sums
\[
 \sum_{a,b=1}^{N}\xi^{ab}=N,\qquad
 \sum_{a,b=1}^{N}\xi^{ab}\psi(a/N)=N\psi(1)=-N\gamma
\]
give
\[
 Z_\xi(s)=\frac{1/N}{(s-1)^2}
       +\frac{(2\gamma-2\log N)/N}{s-1}+O(1).
\]
Mellin inversion yields
\[
 L(\xi e^{-t})=\frac1{2\pi i}\int_{(c)}
                    \Gamma(s)Z_\xi(s)t^{-s}\,ds,\qquad c>1.
\]
After applying \((-\partial_t)^r\), replace \(\Gamma(s)\) by
\(\Gamma(s+r)\) and \(t^{-s}\) by \(t^{-s-r}\).
Move the contour to \(\Re s=1/2\).
The finite Hurwitz-zeta expression has polynomial growth on vertical
strips, while the gamma factor decays exponentially.  Thus the new
integral is \(O(t^{-r-1/2})\), and the only crossed pole is \(s=1\).
Since
\(\Gamma(r+1)=r!\) and \(\psi(r+1)=H_r-\gamma\), its residue is
\[
 \frac{r!}{Nt^{r+1}}
       \left(\log(1/t)+\gamma-2\log N+H_r\right).
\]
Finally apply this result to the primitive \(N_m\)-th root \(\xi^m\),
replace \(t\) by \(mt\), and multiply by \(m^r\).
\end{proof}

The expansion belongs to the classical theory of Lambert series at
roots of unity; much more extensive transseries are available in
\cite{p10:NLDorigoniKleinschmidt}.  The use made of its second term below
is the additional step in the present research program.

\subsection{Every pair of distinct jets is algebraically independent}
\begin{theorem}[Pairwise algebraic independence]
\label{p10:nl:pair}
If \((m,r)\ne(n,s)\), then
\[
 \boxed{F_{m,r}(q),\ F_{n,s}(q)
       \text{ are algebraically independent over }\mathbb C(q).}
\]
Consequently \(L(q)\) satisfies no nonzero first-order algebraic
differential equation over \(\mathbb C(q)\), and, for every \(m>1\),
there is no nonzero polynomial
\(P(q,X,Y)\in\mathbb C(q)[X,Y]\) with
\(P(q,L(q),L(q^m))=0\).
\end{theorem}
\begin{proof}
Suppose a nonzero polynomial \(P(q,X,Y)\) gives a relation.
Choose a prime \(N\) not dividing \(mn\) and a primitive \(N\)-th root
\(\xi\) at which all coefficients are finite and every coefficient
needed below is nonzero.  Such a choice is possible because a
nonzero rational function has finitely many zeros and poles.

First suppose \(r\ne s\). Give \(X,Y\) weights \(r+1,s+1\), and let
\(W\) be the maximum weight of a monomial occurring in \(P\).
Among its weight-\(W\) monomials choose one of maximum ordinary
degree.  This monomial is unique: the equations
\[
 (r+1)i+(s+1)j=W,\qquad i+j=d
\]
uniquely determine \((i,j)\).
Write \(u=\log(1/t)\).  Formula \eqref{p10:nl:cuspformula} shows that,
after multiplication by \(t^W\), the terms of weight \(W\) give a
polynomial in \(u\) whose leading coefficient is nonzero.
All smaller weights tend to zero; the power-saving remainder in
\eqref{p10:nl:cuspformula} remains \(o(1)\) after multiplication by
any fixed power of \(u\).  A nonzero polynomial in \(u\) cannot tend
to zero as \(u\to+\infty\), a contradiction.

Now suppose \(r=s\), so \(m\ne n\). Let \(P_d\) be the highest
ordinary homogeneous part of \(P\), chosen so that
\(P_d(\xi;X,Y)\ne0\). Set
\[
 A=\frac{r!}{Nm},\quad B=\frac{r!}{Nn},\quad
 C=\gamma-2\log N+H_r .
\]
Multiplication of the alleged relation by \(t^{(r+1)d}\) gives
\[
 P_d\bigl(\xi;A(u+C-\log m),B(u+C-\log n)\bigr)+o(1)=0.
 \tag{NL2}\label{p10:nl:affine}
\]
The displayed polynomial in \(u\) is nonzero.  Indeed, its two
linear arguments describe an affine line not passing through the
origin, because their normalized intercepts differ by
\(\log(n/m)\ne0\). A nonzero homogeneous polynomial in two variables
cannot vanish identically on such a line: its zero set is a finite
union of lines through the origin. This contradicts
\eqref{p10:nl:affine}.
\end{proof}

\begin{remark}[The scope of pairwise independence]
The theorem does not prove joint algebraic independence of three
or more jets, nor differential transcendence of \(L\).
It excludes algebraic differential equations involving one derivative,
but does not exclude equations of order two or higher.
Radial asymptotics have previously been used to prove algebraic
independence for Mahler functions \cite{p10:NLBrentCoonsZudilin};
our series is the full Lambert series, not a sparse Mahler series.
The related work of Pupyrev and B\'ezivin
\cite{p10:NLPupyrev,p10:NLBezivin} must also be distinguished:
B\'ezivin differentiates
\(\sum a^nq^n/(1-q^n)\) with respect to \(a\), whereas \(D\) here
differentiates the nome \(q\).
A comprehensive priority claim for Theorem~\ref{p10:nl:pair} is not made.
\end{remark}

\subsection{Algebraic exclusion of finite-level modular substitutions}
Let \(\mathcal K_{\mathrm{mod}}\) be the field of meromorphic functions
on the upper half-plane generated by \(q=e^{2\pi i\tau}\),
the quasimodular Eisenstein series \(E_2(\tau)\), and all
integer-weight meromorphic modular forms of finite congruence level
which are meromorphic at every cusp. Each individual element uses
only finitely many generators. We regard functions of \(q\) on the
upper half-plane by composition. This convention also accommodates
forms with a period greater than one.

\begin{lemma}[Logarithm-free leading monomials]
\label{p10:nl:modcusp}
At every rational cusp \(a/c\), a nonzero
\(h\in\mathcal K_{\mathrm{mod}}\), evaluated at
\(\tau=a/c+it/(2\pi)\), has an asymptotic
\[
 h(\tau)=C e^{-A/t}t^b(1+o(1)),
 \qquad C\ne0,\quad A\in\mathbb R,\quad b\in\mathbb Z.
 \tag{NL3}\label{p10:nl:modlead}
\]
The conclusion includes coefficients that are exponentially small
at the cusp.
\end{lemma}
\begin{proof}
Take a common congruence subgroup for the finitely many generators
in an expression for \(h\), and choose
\(\begin{pmatrix}a&b_0\\c&d\end{pmatrix}\in\mathrm{SL}_2(\mathbb Z)\).
The inverse-transformed variable is exactly
\[
 \tau'=-d/c+\frac{2\pi i}{c^2t}.
\]
A modular form of weight \(k\) becomes a constant times \(t^{-k}\)
times a convergent Laurent series in
\[
 Q=\exp(-\alpha/t),\qquad\alpha>0,
\]
after a common cusp width is chosen; phases are absorbed into
the coefficients. The transformation formula for \(E_2\) adds
only a multiple of \(t^{-1}\). The remaining generator
\(q=e^{2\pi ia/c}e^{-t}\) is analytic in \(t\).
These standard transformations are reviewed in \cite{p10:NLZagier}.

After denominators are cleared, the numerator and denominator of
\(h\) are convergent Laurent series in \(Q\), bounded below in
their \(Q\)-exponents, whose coefficients are meromorphic germs
in \(t\) with a common bound on pole order. They arise from finite
polynomials in the generators, so these bounds are uniform in the
\(Q\)-index. A nonzero such expression has a first nonzero
\(Q\)-coefficient. That coefficient has a first nonzero Laurent
term \(Ct^b\); the later \(Q\)-terms are exponentially smaller.
If all \(Q\)-coefficients vanished identically as germs, normal
convergence would make the original expression identically zero.
Taking the quotient of the two leading monomials proves
\eqref{p10:nl:modlead}.
\end{proof}

\begin{lemma}[A logarithm cannot be algebraic over a cusp field]
\label{p10:nl:logseparation}
Suppose every nonzero element of a coefficient field has
\eqref{p10:nl:modlead} at a fixed cusp. If
\[
 f(t)\sim C_0t^{-w}\log(1/t),\qquad C_0\ne0,
\]
then \(f\) is transcendental over that field.
\end{lemma}
\begin{proof}
In a nonzero sum \(\sum_j a_j f^j\), compare first the smallest
exponential exponent \(A_j\), then the smallest power \(b_j-wj\),
and finally the largest \(j\). The selected term uniquely
dominates the remaining terms: exponentials dominate powers,
and powers dominate powers of a logarithm.
\end{proof}

\begin{theorem}[Nonlinear modular obstruction for every rational filter]
\label{p10:nl:modular}
Let
\[
 F(q)=\sum_{m\in S}\sum_{r=0}^{R}
                   a_{m,r}(q)F_{m,r}(q),\qquad
 a_{m,r}\in\mathbb C(q),
\]
where \(S\) is a finite set of distinct positive integers and the
coefficients are not all zero. Then \(F\) is transcendental over
\(\mathcal K_{\mathrm{mod}}\).
\end{theorem}
\begin{proof}
Take the largest \(r=R_0\) with a nonzero coefficient and set
\[
 A_d(q)=\sum_{m\in S}\frac{\gcd(d,m)}{m}a_{m,R_0}(q),
                 \qquad d\in S.
\]
The gcd matrix on \(S\) is invertible. Indeed,
\[
 \gcd(d,m)=\sum_{e\mid d,\ e\mid m}\varphi(e)
\]
expresses it as a positive Gram matrix; the divisibility incidence
matrix has full row rank, since its columns indexed by \(S\) form
a triangular matrix with diagonal \(1\), after increasing ordering.
Hence some \(A_d\) is a nonzero rational function.

Choose a prime \(\ell\) not dividing any member of \(S\), put
\(N=d\ell\), and choose a primitive \(N\)-th root \(\xi\) avoiding
the finitely many relevant poles and zeros. Then
\(\gcd(N,m)=\gcd(d,m)\). Formula \eqref{p10:nl:cuspformula} implies
\[
 F(\xi e^{-t})\sim
       \frac{R_0! A_d(\xi)}{N}\,
                 t^{-R_0-1}\log(1/t).
\]
Apply Lemmas~\ref{p10:nl:modcusp} and \ref{p10:nl:logseparation}.
\end{proof}

This strengthens the earlier linear exclusion to all polynomial
relations between one such filter and finitely many modular
generators. It also excludes adjoining finitely many algebraic
roots of modular expressions, by transitivity of algebraic
extensions. It does not exclude identities involving integrals,
nonholomorphic completions, mock modular functions, or arbitrary
solutions of nonlinear differential equations.

\subsection{Fibonacci sums and a prime-harmonic transfer}
Define
\[
 \mathcal F(q)=\sum_{n\ge1}
          \frac{q^n}{1-(-1)^nq^{2n}},\qquad
 \vartheta(q)=\sum_{k\in\mathbb Z}q^{k^2}.
\]
Splitting even and odd \(n\) gives the exact identity
\[
 \mathcal F(q)=L(q^2)-L(q^4)
       +\frac{\vartheta(q)^2-\vartheta(-q)^2}{8}.
 \tag{NL4}\label{p10:nl:fib}
\]
For the odd part, the coefficient of \(q^j\) is zero for even \(j\)
and \(\sum_{d\mid j}\chi_4(d)\) for odd \(j\); the two-square
theta identity gives the last term in \eqref{p10:nl:fib}.
The theta squares are integer-weight modular forms of finite
level. Consequently Theorem~\ref{p10:nl:modular} shows that
\(\mathcal F\) is transcendental over \(\mathcal K_{\mathrm{mod}}\).
At \(q=e^{-t}\), the theta term is \(O(t^{-1})\), and
\[
 \mathcal F(e^{-t})\sim \frac1{4t}\log(1/t).
 \tag{NL5}\label{p10:nl:fibasym}
\]


\begin{lemma}[The Fibonacci cusp intercept]
\label{p10:nl:fibcusp}
Let \(N\) be an odd prime, \(\xi\) a primitive \(N\)-th root,
and \(N\nmid m\). Set
\(\mathcal F_{m,r}(q)=D^r[\mathcal F(q^m)]\) and
\(\epsilon_N=(-1)^{(N-1)/2}\). Then
\[
 \mathcal F_{m,r}(\xi e^{-t})
 =\frac{r!}{4Nm\,t^{r+1}}
 \left(\log(1/t)+\gamma-2\log N-\log m+H_r
                       +\epsilon_N\pi/2\right)
       +O_{N,m,r}(t^{-r-1/2}).
\]
\end{lemma}
\begin{proof}
Poisson summation on residue classes modulo \(N\) gives
\[
 \vartheta(\xi e^{-t})
 =\frac{\sqrt{\pi}}{N\sqrt t}
           \sum_{a=0}^{N-1}\xi^{a^2}
       +O(t^{-1/2}e^{-c_N/t}).
\]
The square of the quadratic Gauss sum is \(\epsilon_NN\),
independently of the primitive numerator in \(\xi\).
For \(\vartheta(-\xi e^{-t})\), use residue classes modulo \(2N\);
the zero-frequency Gauss sum vanishes, since the terms indexed
by \(a\) and \(a+N\) cancel. Consequently
\[
 \frac{\vartheta(\xi e^{-t})^2-\vartheta(-\xi e^{-t})^2}{8}
       =\frac{\epsilon_N\pi}{8Nt}
                    +O(t^{-1}e^{-c'_N/t}).
\]
The differentiated Poisson expansions have the same exponential
saving, with harmless additional powers of \(t^{-1}\).
Combine this with \eqref{p10:nl:fib} and Lemma~\ref{p10:nl:cusp}.
The constants involving \(\log 2\) cancel in
\(L(q^{2m})-L(q^{4m})\), giving the stated intercept.
\end{proof}

\begin{theorem}[Mixed Lambert--Fibonacci pairs]
\label{p10:nl:mixedpair}
Any two distinct members of the family
\[
 \{F_{m,r}:m\ge1,r\ge0\}
          \ \cup\ \{\mathcal F_{m,r}:m\ge1,r\ge0\}
\]
are algebraically independent over \(\mathbb C(q)\).
In particular, the full Fibonacci function satisfies no
first-order algebraic differential equation over \(\mathbb C(q)\).
\end{theorem}
\begin{proof}
For different derivative orders, repeat the unique
maximum-weight, maximum-degree argument of
Theorem~\ref{p10:nl:pair}.
For equal derivative orders, the normalized intercepts of the
two cusp expansions differ by a nonzero \(\log(m/n)\) for
distinct dilations in the same family.
For a mixed pair the difference has the form
\(\log(m/n)+\epsilon_N\pi/2\). At least one choice
\(\epsilon_N=1\) or \(-1\) makes this nonzero. Choose infinitely
many primes in the corresponding class modulo \(4\), avoiding
the finitely many coefficient zeros and poles.
The affine-line argument of Theorem~\ref{p10:nl:pair} applies.
\end{proof}

Now introduce the ordinary prime-harmonic generating function
\[
 \mathcal P(q)=\sum_{p\ {\rm prime}}\frac{q^p}{p}.
\]
This is not the prime zeta function \(\sum_p p^{-s}\).
The classical Mertens theorem gives
\[
 \mathcal P(e^{-t})=\log\log(1/t)+B_1+o(1).
 \tag{NL6}\label{p10:nl:primeasym}
\]
For completeness, let \(A(x)=\sum_{p\le x}1/p\). Partial summation gives
\(\mathcal P(e^{-t})=\int_{2t}^{\infty}e^{-u}A(u/t)\,du\).
Insert \(A(x)=\log\log x+B_1+o(1)\).
On \(u\in[t^{1/2},t^{-1/2}]\), the difference
\(\log\log(u/t)-\log\log(1/t)\) tends pointwise to zero and is
bounded by an integrable function after multiplication by \(e^{-u}\);
the lower and upper tails contribute \(o(1)\).
This proves \eqref{p10:nl:primeasym}. Quantitative versions of the
classical input are in \cite{p10:NLRosserSchoenfeld}.

\begin{theorem}[Two logarithmic scales over the modular field]
\label{p10:nl:primepair}
Both pairs
\[
 \boxed{\bigl(L(q),\mathcal P(q)\bigr),\qquad
        \bigl(\mathcal F(q),\mathcal P(q)\bigr)}
\]
are algebraically independent over \(\mathcal K_{\mathrm{mod}}\).
\end{theorem}
\begin{proof}
Write \(u=\log(1/t)\), and let \(f\) be either \(L\) or
\(\mathcal F\). Then
\(f(e^{-t})\sim Ct^{-1}u\), \(C\ne0\), while
\(\mathcal P(e^{-t})\sim\log u\).
In a nonzero polynomial
\(\sum_{i,j}a_{ij}f^i\mathcal P^j\), use
Lemma~\ref{p10:nl:modcusp} and compare, in order,
the smallest exponential exponent \(A_{ij}\), the smallest
power \(b_{ij}-i\), the largest \(i\), and the largest \(j\).
The selected term has a unique leading asymptotic
\[
 C_{ij}e^{-A_{ij}/t}t^{b_{ij}-i}u^i(\log u)^j.
\]
It cannot be canceled. This proves the theorem.
\end{proof}


\begin{theorem}[A third independent function from primes]
\label{p10:nl:triple}
Let \(f,g\) be any two distinct members of the family in
Theorem~\ref{p10:nl:mixedpair}. Then
\[
 \boxed{f(q),\ g(q),\ \mathcal P(q)
       \text{ are algebraically independent over }\mathbb C(q).}
\]
This assertion has coefficient field \(\mathbb C(q)\), rather
than \(\mathcal K_{\mathrm{mod}}\).
\end{theorem}
\begin{proof}
For a primitive root \(\xi\) of odd prime order \(N\), the Mertens
theorem in arithmetic progressions, followed by the same partial
summation as above, gives
\[
 \mathcal P(\xi e^{-t})
       =-\frac{1}{N-1}\log\log(1/t)+C_\xi+o(1).
 \tag{NL7}\label{p10:nl:primecusp}
\]
Indeed, each nonzero residue class has prime-harmonic sum
\((N-1)^{-1}\log\log x+C_a+o(1)\), while
\(\sum_{a=1}^{N-1}\xi^a=-1\); the prime \(N\) contributes only
a bounded term. The arithmetic-progression form of Mertens'
theorem is recorded, for example, in
\cite{p10:NLLanguascoZaccagnini}.

Suppose \(P(q,f,g,\mathcal P)=0\).
Choose a prime cusp as in Theorem~\ref{p10:nl:mixedpair}, avoiding
the finitely many nonzero rational coefficient zeros and poles.
Let \(u=\log(1/t)\).
If the derivative orders of \(f,g\) differ, take the largest
weighted degree in their variables and then the largest
ordinary degree in those variables. The selected pair of
exponents is unique. Its coefficient is a nonzero polynomial
in \(\mathcal P(\xi e^{-t})\), hence has a nonzero leading
power of \(\log u\), by \eqref{p10:nl:primecusp}. Other pairs have
smaller powers of \(u\), and lower weights have a power saving
in \(t\), so they cannot cancel this term.

For equal derivative orders, take the highest homogeneous part
in the two variables corresponding to \(f,g\), regarding its
coefficients as polynomials in a third indeterminate \(Z\).
Substitution of the affine cusp line is injective on homogeneous
polynomials in the first two variables, over \(\mathbb C[Z]\).
It therefore produces a nonzero polynomial \(Q(u,Z)\).
Taking \(Z=\mathcal P(\xi e^{-t})\), the largest power of \(u\)
in \(Q\), and then the largest power of \(Z\) in its coefficient,
has a unique nonzero leading term
\(c u^j(\log u)^k\).
The error from the two Lambert--Fibonacci cusp expansions is
\(O(t^{1/2})\) times fixed powers of \(u\) and \(\log u\),
hence tends to zero. This is again a contradiction.
\end{proof}

\begin{remark}[What did and did not transfer to the named constants]
The Erd\H{o}s--Borwein constant is \(L(1/2)\), and, if
\(\phi=(1+\sqrt5)/2\), the Fibonacci reciprocal sum is
\(\sqrt5\,\mathcal F(\phi^{-1})\).
The results above exclude complete families of functional
algebraic reductions, including nonlinear modular substitutions.
They supply no specialization theorem at \(1/2\) or
\(\phi^{-1}\). Bell--Smertnig's characterization
\cite{p10:NLBellSmertnig} already rules out the ordinary linear
Mahler framework for \(L\).
For \(B_1\), \eqref{p10:nl:primeasym} is a renormalized boundary limit.
Algebraic independence of the functions does not determine the
arithmetic nature of its constant term. Thus the transcendence
of the two reciprocal-sum constants and the irrationality of
\(B_1\) remain unresolved in this work.
\end{remark}



\section{Gamma isolation with unrestricted auxiliary denominators}
\label{p10:gn:section}

The preceding manuscript computed a beta-factor cost in a fixed prime
reflection model.  A natural repair of an unsuccessful isolation argument
is to introduce gamma values with new denominators and use Gauss
multiplication.  We show that this repair does not reduce that cost.
Moreover, we determine the complete cost functions for monomials in the
fifth and seventh gamma values.  These are theorems about the standard
identity calculus.  They are not proofs of the transcendence of the five
individual gamma values under investigation.

\subsection{The standard identity calculus}

All gamma and beta arguments in this section are positive rational
numbers.  Algebraic nonzero factors and rational powers of $\pi$ have
zero cost.  The identities allowed are recurrence, reflection, Gauss
multiplication, and their finite products, quotients, and rational powers.
Rational powers can be taken as positive real powers; any choices of
algebraic roots can be absorbed into the algebraic factor.

Here is a precise formal description.  Let $\mathcal D$ be the rational
vector space on symbols $[x]$, $x\in\mathbb Q/\mathbb Z$, modulo
\begin{align}
 [0]&=[1/2]=0, &[-x]&=-[x], \label{p10:gn:reflection}\\
 \sum_{k=0}^{d-1}[x+k/d]&=[dx] &&(d\geq1).
 \label{p10:gn:distribution}
\end{align}
A symbol $[x]$ represents $\Gamma(x)$ modulo algebraic factors and
rational powers of $\pi$.  Recurrence removes integer shifts,
reflection gives \eqref{p10:gn:reflection}, and Gauss multiplication gives
\eqref{p10:gn:distribution}.  When a representative is an integer, its
symbol is zero.  This also handles the limiting multiplication identity
at $x=0$.
For $x,y\in\mathbb Q/\mathbb Z$ write
\[
 b(x,y)=[x]+[y]-[x+y].
\]
It represents $B(x,y)=\Gamma(x)\Gamma(y)/\Gamma(x+y)$ after choosing
positive representatives.  The choices affect only algebraic factors.

For $v\in\mathcal D$ define
\begin{equation}\label{p10:gn:cost-definition}
 \ell(v)=\min\left\{\sum_i|n_i|:
 v=\sum_i n_i b(x_i,y_i),\quad n_i\in\mathbb Z\right\},
\end{equation}
with $\ell(v)=\infty$ when no such representation exists.  The
arguments $x_i,y_i$ have no denominator restriction.  The beta factors
in the final word have integer exponents; allowing arbitrary rational
exponents in that final word would define a different problem.
The equalities in \eqref{p10:gn:cost-definition} may nevertheless use
rational combinations of all standard gamma identities.

\begin{lemma}[Sawtooth certificates, including the endpoints]
\label{p10:gn:sawtooth}
Set
\[
 s(x)=\begin{cases}1-2\{x\},&x\notin\mathbb Z,\\0,&x\in\mathbb Z.
 \end{cases}
\]
Then
\begin{equation}\label{p10:gn:sawtooth-distribution}
 \sum_{k=0}^{d-1}s(x+k/d)=s(dx),
 \qquad s(-x)=-s(x),\qquad s(1/2)=0.
\end{equation}
Also $s(x)+s(y)-s(x+y)\in\{-1,0,1\}$.
Suppose a formal identity involves only finitely many standard relations,
and let $N$ be a common denominator of every argument appearing in that
proof.  For every integer $t$ coprime to $N$, substitution
$[x]\mapsto s(tx)$ respects every relation in the proof.
\end{lemma}

\begin{proof}
Put $u=\{dx\}$.  If $u\ne0$, the fractional parts of $x+k/d$ are
$(u+j)/d$, $0\leq j<d$, in some order.  Their sum is
$(d-1)/2+u$, so the sum of their sawtooth values is $1-2u$.
If $u=0$, the fractional parts are $0,1/d,\ldots,(d-1)/d$.
The zero term contributes zero and the remaining terms sum to zero.
This proves the distribution identity also at the discontinuities.
The oddness and the assertion about $1/2$ follow directly.
If neither $x$, $y$, nor $x+y$ is an integer, the beta defect is $1$
or $-1$ according as $\{x\}+\{y\}<1$ or $>1$; the omitted cases
have defect zero.

For a distribution relation with multiplier $d$, its arguments have
differences $1/d$, so $d\mid N$.  Since $(t,d)=1$, multiplication by
$t$ permutes the fractions $k/d$ modulo one.  Consequently
\[
 \sum_{k=0}^{d-1}s(tx+tk/d)
 =\sum_{k=0}^{d-1}s(tx+k/d)=s(tdx).
\]
The other relations were already checked.  A rational linear combination
of these relations is respected as well.
\end{proof}

A certificate need not extend to a real-valued functional on all of
$\mathcal D$ for one fixed $t$.  The proof above is deliberately finite:
for any proposed finite identity, a suitable unit $t$ is chosen after
all its auxiliary denominators have been included in $N$.
This avoids making an unjustified global invariance assertion.

\begin{theorem}[Exact isolation cost in the full standard calculus]
\label{p10:gn:unrestricted}
Let $x=a/q$ be reduced with $q\geq3$, and let $m\in\mathbb Z$.
Then
\begin{equation}\label{p10:gn:unrestricted-formula}
 \boxed{\ell(m[x])=
 \begin{cases}
 |m|(q-2)/q,&q\mid 2m,\\
 \infty,&q\nmid 2m.
 \end{cases}}
\end{equation}
This remains true even when arbitrarily large auxiliary denominators
and every standard gamma multiplication identity are permitted.
\end{theorem}

\begin{proof}
Consider any finite representation of $m[x]$ by an integer beta word,
including a finite proof from the standard relations.  Enlarge $N$ so
$q\mid N$.  Choose $t$ coprime to $N$ with $ta\equiv1\pmod q$.
Such a choice exists: impose the inverse of $a$ modulo $q$, and impose
$t\equiv1$ modulo each prime dividing $N$ but not $q$, then use the
Chinese remainder theorem.  Lemma~\ref{p10:gn:sawtooth} gives
\[
 m(1-2/q)=\sum_i n_i\epsilon_i,\qquad
 \epsilon_i\in\{-1,0,1\}.
\]
The right side is an integer.  Hence $q\mid 2m$, and its absolute
value is at most $\sum_i|n_i|$.  This proves both the obstruction and
the lower bound, including after saturation by rational powers of
standard identities.

If $q$ is odd, telescoping gives
\[
 \sum_{k=1}^{q-2}b(kx,x)=(q-1)[x]-[(q-1)x]=q[x].
\]
This is a word of length $q-2$, and $q\mid2m$ then means $q\mid m$.
If $q$ is even, $a$ is odd and $(q/2)x=1/2$ modulo one.  Thus
\[
 \sum_{k=1}^{q/2-1}b(kx,x)=(q/2)[x]-[(q/2)x]=(q/2)[x].
\]
This word has length $q/2-1$, and $q/2\mid m$.
Repetition and inversion attain the stated bound for every allowed $m$.
\end{proof}

For example, the least admissible powers at reduced denominators
$4,5,6,7,8,9,10$ are respectively $2,5,3,7,4,9,5$, with exact costs
$1,3,2,5,3,7,4$.  In particular, introducing denominators $10,15,35$,
or larger ones cannot shorten the three-factor certificate for a fifth
power or the five-factor certificate for a seventh power.

\subsection{Complete monomial cost at fifths and sevenths}

The single-coordinate theorem has a stronger lattice counterpart at the
two denominators relevant here.  Put
\[
 v_5(A,B)=A[1/5]+B[2/5],\qquad
 v_7(A,B,C)=A[1/7]+B[2/7]+C[3/7].
\]
All exponents below are integers and may be negative.

\begin{theorem}[The entire fifth-gamma lattice]
\label{p10:gn:fifths}
One has
\begin{equation}\label{p10:gn:fifth-formula}
 \ell(v_5(A,B))=
 \begin{cases}
 \bigl(|2A-B|+|A+2B|\bigr)/5,& B\equiv2A\pmod5,\\
 \infty,&B\not\equiv2A\pmod5.
 \end{cases}
\end{equation}
The finite value is also
$\max\{|3A+B|,|A-3B|\}/5$.
\end{theorem}

\begin{proof}
Applying the finite certificates with $t\equiv1,2\pmod5$ gives
respectively
\[
 (3A+B)/5,\qquad (A-3B)/5.
\]
Each must be an integer and its absolute value is a lower bound for
any word length, regardless of its auxiliary denominators.
Their integrality is equivalent to the displayed congruence.

Modulo \eqref{p10:gn:reflection}, the two beta symbols
$b(1/5,1/5)$ and $b(1/5,2/5)$ have vectors
$u=(2,-1)$ and $v=(1,2)$ in the ordered coordinates
$([1/5],[2/5])$.  Under the congruence hypothesis,
\[
 (A,B)=r u+s v,\qquad
 r=(2A-B)/5,\quad s=(A+2B)/5
\]
with $r,s\in\mathbb Z$.
The corresponding word has length $|r|+|s|$.  Its two certificate
values are $r+s$ and $r-s$, and
$\max\{|r+s|,|r-s|\}=|r|+|s|$.  Thus the lower bound is attained.
\end{proof}

\begin{theorem}[The entire seventh-gamma lattice]
\label{p10:gn:sevenths}
If $A+2B+3C\not\equiv0\pmod7$, then
$\ell(v_7(A,B,C))=\infty$.  Otherwise
\begin{equation}\label{p10:gn:seventh-formula}
 \boxed{\ell(v_7(A,B,C))=
 \frac1{7}\max\bigl\{|5A+3B+C|,\ |3A-B-5C|,\ |A-5B+3C|\bigr\}.}
\end{equation}
An optimal word can be computed by taking a median of four explicit
integers.
\end{theorem}

\begin{proof}
The certificate values for $t\equiv1,2,3\pmod7$ are precisely the
three linear forms divided by seven in \eqref{p10:gn:seventh-formula}.
Their integrality is equivalent to $A+2B+3C\equiv0\pmod7$.
For example, $5A+3B+C\equiv5(A+2B+3C)\pmod7$.
Their absolute values bound every word from below.

The four beta symbols with argument pairs
$(1,1),(1,2),(1,3),(2,2)$ divided by seven have coordinate vectors
\[
 u=(2,-1,0),\quad v=(1,1,-1),\quad
 w=(1,0,2),\quad z=(0,2,1),\qquad u+z=v+w.
\]
Set
\[
 r=\frac{2A-3B-C}{7},\qquad
 s=\frac{2A+4B-C}{7},\qquad
 h=\frac{A+2B+3C}{7}.
\]
These are integers under the congruence condition, and
$(A,B,C)=r u+s v+h w$.  For every integer $d$ the same vector is
\[
 (r+d)u+(s-d)v+(h-d)w+d z.
\]
Its word length is
\[
 |d+r|+|d-s|+|d-h|+|d|.
\]
Choose $d$ between the two middle order statistics of
$-r,s,h,0$.  Such an integer exists.  The minimum is the sum of the
two largest of these numbers minus the sum of the two smallest.
Equivalently it is
\[
 \max\{|r+s+h|,\ |r+s-h|,\ |r-s+h|\}.
\]
Indeed these are, up to sign, the three ways to subtract the sum of
one pair from the sum of its complementary pair.  Substitution gives
exactly \eqref{p10:gn:seventh-formula}, proving sharpness.
\end{proof}

These formulas solve an integer optimization problem for every exponent
vector, rather than testing only a bounded list of candidate identities.
The unrestricted-denominator lower bounds are essential: enumerating
beta words at denominators five and seven alone would not establish
this conclusion.

\subsection{What the formal theorem does and does not decide}

There is a natural map from $\mathcal D$ to the multiplicative classes
of actual gamma values modulo algebraic factors and rational powers of
$\pi$.  The standard identities show this map is well defined; its
injectivity is not known.  Rohrlich's conjectural description of
multiplicative gamma relations is exactly the missing type of input
\cite{p10:gn:waldschmidt}.  Under that conjecture, the cost formulas above
also govern actual numerical beta-product identities.  Without it,
an unexpected numerical relation may identify distinct formal classes.
In particular, a formal infinite cost for $\Gamma(1/5)$ does not prove
that this number is transcendental.  This is not a semantic distinction:
an algebraic gamma value would itself give a zero-cost numerical
identity that is invisible to the standard calculus.

For comparison there is a genuine, but classical, nonlinear arithmetic
subsystem inside the seventh-gamma problem.  Let
\[
 X=\Gamma(1/7),\quad Y=\Gamma(2/7),\quad Z=\Gamma(3/7),
 \qquad B_7=B(1/7,2/7)=XY/Z,
 \qquad Q_7=e^{\pi\sqrt7}.
\]
Reflection gives the exact identity
\begin{equation}\label{p10:gn:cs7}
 \operatorname{CS}(-7)=
 \frac{\sin(\pi/7)\sin(2\pi/7)}{\sin(3\pi/7)}\,
 \frac{B_7^2}{\pi}.
\end{equation}
The classical Chowla--Selberg and Nesterenko results imply that
$\pi$, $\operatorname{CS}(-7)$, and $Q_7$ are algebraically independent;
see the explicit formulation and recent quantitative refinements of
Cohen and Zudilin \cite{p10:gn:cohen-zudilin}.

\begin{proposition}[A completely classified classical CM subsystem]
\label{p10:gn:cm}
The numbers $\pi,B_7,Q_7$ are algebraically independent over
$\overline{\mathbb Q}$.  Consequently, for
$R\in\overline{\mathbb Q}(U,V,W)$, its evaluation
$R(\pi,B_7,Q_7)$ is algebraic if and only if $R$ is a constant rational
function.  Every nonzero denominator remains nonzero on evaluation.
\end{proposition}

\begin{proof}
By \eqref{p10:gn:cs7}, adjoining $B_7$ or adjoining
$\operatorname{CS}(-7)$ to $\overline{\mathbb Q}(\pi,Q_7)$ gives
fields that are algebraic over each other.  Their transcendence degree
is three, proving the first assertion.  If a nonconstant rational
function evaluated to an algebraic number $a$, its numerator minus
$a$ times its denominator would give a nonzero polynomial relation.
The same independence excludes a zero denominator.
\end{proof}

Proposition~\ref{p10:gn:cm} is stated to mark the usable arithmetic input,
not as a newly proved transcendence theorem.  It classifies nonlinear
expressions in a distinguished beta period.  The corresponding vector
$(1,1,-1)$ has exact beta cost one by Theorem~\ref{p10:gn:sevenths};
each isolated seventh power has cost five.  Knowing the arithmetic of
one CM beta period does not classify products of the other beta
periods in a five-factor isolating word.

For the five named individual targets
$\Gamma(1/5),\Gamma(2/5),\Gamma(1/7),\Gamma(2/7),\Gamma(3/7)$,
this section therefore establishes no irrationality or transcendence
closure.  Its new increment over the preceding manuscript is the
unrestricted-denominator theorem, its even and composite cases, and the
complete optimal monomial formulas at fifths and sevenths.
The primary 2026 discussion of gamma-value transcendence by Sheth
\cite{p10:gn:sheth} also makes clear that geometric transcendence does not
supply the missing individual-value result.


\section{Fourteen constant targets and one related conjecture}
\label{p10:next:ledger}

The named targets in this section are the arithmetic questions carried
forward from the previous manuscript. The word ``unresolved'' below
describes what the present arguments do not prove. It is not a claim
that no relevant external work exists. In particular, irrationality
of $\zeta(5)$ and Catalan's constant is reported in the cited 2026
external sources; neither is counted as a discovery of this manuscript,
and their complete proofs are not independently verified here.

\begingroup\small
\renewcommand{\arraystretch}{1.13}
\begin{longtable}{@{}L{0.25\textwidth}L{0.70\textwidth}@{}}
\toprule
Target & Additional result and remaining obstruction\\
\midrule
\endfirsthead
\toprule
Target & Additional result and remaining obstruction\\
\midrule
\endhead
$\gamma$: irrationality or transcendence &
Curvature bounds the indices admitting a bounded primitive even-Taylor
crossover form by $O_B(X^{1/3})$. An infinite exceptional set and
the logarithmic term in the form remain. No value conclusion.\\
$\delta$: irrationality &
Factorial-size primitive content occurs at only $O(\log X)$ indices
up to $X$. Bounded primitive forms must be sparse and lacunary.
An infinite exceptional certificate subsequence is not excluded.\\
$G$: transcendence &
Direct positive rational intervals avoid the auxiliary $\zeta(2)$
term and determinant nonvanishing gap. Their proved denominator bound
does not yield a transcendence criterion.\\
$\zeta(3)$: transcendence &
The exact cancellation diagnostic recovers the classical Gaussian
orthogonality cancellation. It does not exclude higher algebraic
degrees.\\
$\zeta(5)$: transcendence &
Exact norm and shift asymptotics; reduced-denominator experiments.
The reported external irrationality result is not strengthened to
transcendence here.\\
$\zeta(7)$: irrationality &
The degree-to-shift ratio is optimized exactly at a fixed total
budget. The common-denominator estimate retains a negative arithmetic
margin; a sufficient actual cancellation estimate is absent.\\
Erd\H{o}s--Borwein constant: transcendence &
Distinct Lambert jets are algebraically independent as functions,
and finite rational filters remain transcendental over the indicated
modular field. No specialization theorem establishes the value claim.\\
Reciprocal Fibonacci sum: transcendence &
A theta--Lambert decomposition transfers the nonlinear modular
obstruction. Functional independence still does not imply the
arithmetic nature of the specified sum.\\
Prime Mertens constant $B_1$: irrationality &
The prime-harmonic generating function has a distinct logarithmic
scale, giving functional exclusion and independence. Its constant
term is not shown irrational.\\
$\Gamma(1/5)$: transcendence &
Exact unrestricted beta-word cost, including the full fifth-gamma
monomial lattice. A formal isolation identity does not determine
the numerical value.\\
$\Gamma(2/5)$: transcendence &
The same complete fifth-gamma cost applies after every permitted
auxiliary denominator. No new value theorem.\\
$\Gamma(1/7)$: transcendence &
Exact unrestricted isolation and the complete seventh-gamma monomial
cost. Numerical algebraic relations remain beyond this calculus.\\
$\Gamma(2/7)$: transcendence &
The same seventh-gamma cost and its classical CM comparison apply.
No individual value conclusion.\\
$\Gamma(3/7)$: transcendence &
The same seventh-gamma cost closes the formal shortcut, without
proving transcendence of this individual gamma value.\\
Kurepa's left-factorial conjecture &
The previous local Laguerre-content connection can be compared with
the new global sparsity theorem. Sparse large content is compatible
with either outcome of the conjecture, so no resolution follows.\\
\bottomrule
\end{longtable}
\endgroup

\subsection*{What is newly closed inside the stated models}
The proved conclusions include a logarithmic counting bound for large
primitive cancellation; a complete totally positive channel-loss
classification; an infinite algebraically linearly independent family
of entropy constants; the exact finite-moment threshold $d(M)=M+1$;
the fixed-budget Wilson shift optimum; and the unrestricted standard
gamma-calculus cost. The explicit all-irrational four-atom entropy
construction and direct Catalan intervals are constructive additions.
The nonlinear Lambert conclusions strengthen the previous functional
obstruction and include independent triples with the prime-harmonic
function. These statements have their precise hypotheses and proofs
in the corresponding sections. They are not renamed solutions of the
named-constant questions in the table.

\subsection*{Verification and remaining research obligations}
Exact rational arithmetic checks accompany the Laguerre content,
Wilson and Catalan interval formulas, entropy identities, and finite
gamma costs. Such checks can detect a faulty formula but do not prove
an asymptotic assertion; those assertions are justified in the text.
The main remaining obligations for a named-constant breakthrough are
an arithmetic specialization theorem, an adequate primitive-height
bound, or elimination of a surviving exceptional subsequence. None
is silently substituted by convergence, functional independence,
or a finite numerical search.

% END SOURCE: KEIO_Sparse_Certificates_and_Information_Constants.tex

\endgroup
\clearpage
\part{Spectral Entropy and Constant Transfer}
\label{merge:part11}
\begingroup

% BEGIN SOURCE: KEIO_Spectral_Entropy_and_Arithmetic_Transfer.tex

\noindent\textbf{Editorial integration note.} 
The lower bound of degree three for a remaining ternary counterexample is superseded in Part~\ref{merge:part12} by degree at least four, or odd degree at least five for odd entropy denominator. The final part also distinguishes the invariant Catalan space from its convergent full-space repair and passes from formal Gamma signatures to a specified actual beta-period span.


\begin{quote}\small
\noindent\textit{Original part framing.}\par
\noindent\textbf{Scope.}
This standalone successor restates the constructions needed for its
proofs. It distinguishes totally real numbers (all conjugates real)
from totally positive numbers (all conjugates positive). Classical
external inputs, independent proof reviews, and exact finite checks
are identified; the manuscript is not formally verified.
The cited 2026 external sources report irrationality proofs for
$\zeta(5)$ and Catalan's constant
\cite{p11:zf:Fauzan,p11:zf:Lai,p10:nz:OpenAICatalan}. Their full proofs and formal
verification dependencies are not certified here. The corresponding
targets in this manuscript are transcendence, and those external
results are not counted as ours.
\end{quote}
\section{The new conclusions and their precise scope}
\label{p11:tr:introduction}

The starting point is an arithmetic mismatch. Rapid convergence,
positivity of a moment matrix, and algebraic independence of functions
can all be proved without deciding the arithmetic nature of a specified
constant. We treat the mismatch as a source of questions with exact
answers, and then return to the named targets to state what is still
missing. The claims below are additional conclusions of this research
program; none is a declaration of historical priority.

\paragraph{An exact arithmetic classifier for entropy powers.}
For a positive probability vector $p$ put
$E(p)=\prod_i p_i^{-p_i}$. Theorem~\ref{p11:ipe:classifier} proves, when
every $p_i$ is totally positive and $n\ge2$, that for nonzero real
algebraic $c$,
\[
 E(p)^c\in\overline{\mathbb Q}
 \quad\Longleftrightarrow\quad
 p_1,\ldots,p_n,c\in\mathbb Q.
\]
A symmetric relative-entropy identity forces every positive conjugate
vector to coincide with the original one. The result applies to any
density matrix with rational characteristic polynomial. It also makes
Laplacian integrality the exact criterion for algebraicity of the
effective rank of a finite graph. The final part of that section
strengthens the earlier linear independence of entropy constants to
an exclusion of algebraic-exponent products of their effective ranks.
It does not assert algebraic independence of those ranks.

\paragraph{A sharp threshold over all totally real fields.}
Theorem~\ref{p11:et:totallyreal} excludes rational entropy in bits for
every nonrational ternary probability vector whose entries are totally
real. Finite valuations first force a three-monomial generator to be
an algebraic unit with odd exponents. A count of real embeddings then
excludes every totally real probability field, without restricting
its degree. The explicit golden-ratio four-atom example gives the
matching upper bound:
\[
 \min\{\text{alphabet size with all entries irrational and totally real,
                     and rational entropy}\}=4.
\]
The statement concerns rational entropy; it does not classify the
entropy as algebraic or transcendental in every case. Without the
totally real hypothesis, the minimum remains three or four here.

\paragraph{A higher functional transfer, with a finite-order limit.}
Theorem~\ref{p11:tl:main} proves algebraic independence over $\mathbb C(q)$
of three distinct Lambert--Fibonacci differential--dilation jets,
provided their derivative orders are not all equal, together with the
prime-harmonic series. In particular the Lambert and Fibonacci
functions satisfy no algebraic differential equation of order at
most two. The proof varies prime cusp denominators to separate two
polynomial gradings. An explicit order-three relation in the leading
cusp model explains why that proof does not establish full
differential transcendence. These are function theorems; the named
values require a specialization argument that is not provided.

\paragraph{Positivity and the cost of higher algebraic degrees.}
Theorem~\ref{p11:zf:degree} converts a small primitive real-rooted
polynomial value and a single off-axis estimate into exclusion of a
specified algebraic degree. The exact Pell family in
Theorem~\ref{p11:zf:pell} reaches the quadratic boundary while retaining
positive Hankel matrices, simple roots, and primitive coefficients.
For Catalan's constant, Theorem~\ref{p11:zf:catalan} constructs a positive
Gram matrix $A_m+16G I_m$ with $A_m$ rational. Paired interpolation
cancels $\pi^2$ without losing determinant degree. Quantitative bounds
for primitive value and off-axis growth remain necessary.

\paragraph{Two broad arithmetic transfer classes can be measured exactly.}
The Laguerre parameter, unlike the index, has only simple local roots.
Theorem~\ref{p11:tg:law} therefore determines the complete distribution
of the primitive common divisor over parameter residues, with
independent truncated geometric valuations. Uniform exponential-tail
bounds follow. They describe parameter averages and do not settle
the fixed Euler--Gompertz parameter $A=1$.
For gamma values, Theorem~\ref{p11:gc:recognition} identifies exactly
which prime-denominator monomials can be reduced to elliptic
complex-multiplication periods using all standard gamma identities,
all auxiliary rational denominators, and rational powers. This is an
identity-calculus theorem; unknown numerical gamma relations remain
outside its conclusion.


\section{Algebraic powers of effective rank}
\label{p11:ipe:section}

For a positive probability vector $p=(p_i)_{i=1}^n$, define its
natural-log entropy and its effective number of outcomes by
\[
 S(p)=-\sum_i p_i\log p_i,
 \qquad E(p)=e^{S(p)}=\prod_i p_i^{-p_i}.
\]
Thus $E(p)=2^{H_2(p)}$. The quantity $E$ is also the effective rank
in the sense of Roy and Vetterli \cite{p11:ipe:RV}
of a density matrix when $p$ is its nonzero spectrum. It is not the
Gaussian-normalized entropy power used for continuous distributions.
The preceding manuscript proved transcendence of $H_2$ under a
rational-characteristic-polynomial hypothesis. Here we classify
\emph{all real algebraic powers of $E$} under the more general
hypothesis of totally positive algebraic probabilities.
The proof uses Baker's classical logarithmic theorem
\cite{p11:TLBaker}; historical priority of the resulting arithmetic
formulation has not been established by an exhaustive literature
review.

\subsection{A logarithmic coordinate lemma}

\begin{lemma}[When an algebraic linear combination is an algebraic logarithm]
\label{p11:ipe:coordinates}
Let $u_1,\ldots,u_t$ be real logarithms of positive algebraic numbers,
linearly independent over $\mathbb Q$. For algebraic real numbers
$b_1,\ldots,b_t$ one has
\[
       \exp\left(\sum_j b_j u_j\right)\in\overline{\mathbb Q}
                  \quad\Longleftrightarrow\quad b_j\in\mathbb Q
                        \text{ for every }j.
\]
Furthermore, algebraic linear relations among logarithms are the
algebraic span of their rational relations, and therefore remain
valid after conjugating their algebraic coefficients alone.
\end{lemma}
\begin{proof}
Rational coefficients give a product of rational powers of algebraic
numbers. Conversely, if the exponential is a positive algebraic
number $a$, consider the logarithms $u_1,\ldots,u_t,\log a$.
If $\log a$ were independent of the $u_j$ over $\mathbb Q$,
Baker's theorem would prohibit the displayed algebraic relation.
Hence $\log a=\sum_j r_j u_j$ with rational $r_j$.
Baker's theorem applied to the independent $u_j$ gives $b_j=r_j$.
The final assertion follows by choosing a rational basis of any
finite logarithmic span and comparing coefficients.
\end{proof}

\begin{theorem}[Positive-conjugate rigidity for algebraic powers]
\label{p11:ipe:positive}
Let $p=(p_1,\ldots,p_n)$ be a positive algebraic probability vector,
$n\geq2$, and let $c\ne0$ be real algebraic. Suppose $E(p)^c$ is
algebraic. For every field embedding $\sigma$ for which all
$q_i=\sigma(p_i)$ are positive, one has
\[
                         q_i=p_i\quad\text{for every }i,
                         \qquad \sigma(c)=c.
\]
Here the embedding is taken on a number field containing $p$ and $c$.
\end{theorem}
\begin{proof}
Put $A=E(p)^c>0$, which is algebraic by assumption. Then
\[
                      \log A+c\sum_i p_i\log p_i=0.
\]
Baker's rational-generation property allows conjugation of the
algebraic coefficients alone, retaining every logarithm. Writing
$d=\sigma(c)$, this gives
\[
                      \log A+d\sum_i q_i\log p_i=0.
\]
An extension of $\sigma$ to the algebraic closure may be used for
this coefficient operation; the base $A$ is not conjugated.
Subtract the two identities. Since $0<p_i<1$ and $q_i>0$, both
$\sum_i p_i\log p_i$ and $\sum_i q_i\log p_i$ are negative.
Thus
\[
             \frac dc=\frac{\sum_i p_i\log p_i}
                             {\sum_i q_i\log p_i}>0.
\]
In particular $d$ is real and has the same sign as $c$.

Set $r_i=cp_i-dq_i$. We have $\sum_i r_i\log p_i=0$.
Again by Baker's rational-generation property, this relation is an
algebraic linear combination of rational relations among the
$\log p_i$. Every rational relation, after clearing denominators,
is a multiplicative identity in the $p_i$. Conjugating its bases
by $\sigma$ and using positivity yields the same rational relation
among the $\log q_i$. Taking the original algebraic linear
combination, without conjugating the $r_i$, therefore gives
$\sum_i r_i\log q_i=0$. Subtraction now yields
\[
                cD(p\Vert q)+dD(q\Vert p)=0.
\]
Divide by $c$. Both divergences are nonnegative and $d/c>0$,
so each divergence is zero. Their strict equality condition gives
$p=q$. The earlier negative sums then give $d=c$.
\end{proof}

\begin{theorem}[Exact algebraic-power classifier]
\label{p11:ipe:classifier}
Let $p_1,\ldots,p_n>0$, $n\geq2$, sum to one, and suppose every
$p_i$ is totally positive over $\mathbb Q$. For every nonzero real
algebraic number $c$,
\[
 \boxed{\quad E(p)^c\in\overline{\mathbb Q}
       \quad\Longleftrightarrow\quad
       p_1,\ldots,p_n\in\mathbb Q\ \text{ and }\ c\in\mathbb Q.
       \quad}
\]
In particular, one irrational entry forces every nonzero real
algebraic power of the effective number of outcomes to be transcendental.
\end{theorem}
\begin{proof}
Take a normal field containing all entries and $c$. Total positivity
makes every conjugate vector positive. If $E(p)^c$ is algebraic,
Theorem~\ref{p11:ipe:positive} shows that every automorphism fixes each
$p_i$ and $c$, hence they are rational. Conversely, rational $p_i$
and rational $c$ give a finite product of rational powers of rational
numbers, which is algebraic.
\end{proof}

\begin{corollary}[Every algebraic logarithm base]
\label{p11:ipe:bases}
Under the irrational-entry hypothesis of
Theorem~\ref{p11:ipe:classifier}, the entropy $S(p)/\log b$ is
transcendental for every real algebraic base $b>1$.
\end{corollary}
\begin{proof}
If $a=S(p)/\log b$ were algebraic, then $a>0$ and
$E(p)^{1/a}=b$ would contradict Theorem~\ref{p11:ipe:classifier}.
\end{proof}

\begin{corollary}[Rational spectral data]
\label{p11:ipe:quantum}
Let $\rho$ be a density matrix with rational characteristic
polynomial and at least two positive eigenvalues. If any eigenvalue
is irrational, then
\[
                   \exp\bigl(c[-\operatorname{tr}(\rho\log\rho)]\bigr)
\]
is transcendental for every nonzero real algebraic $c$.
If all nonzero eigenvalues are rational, this expression is algebraic
exactly for rational $c$. This includes every rational-entry density
matrix. For a pure state the expression is always $1$ and is an
explicit exception.
\end{corollary}
\begin{proof}
Delete zero eigenvalues. Every conjugate of each remaining eigenvalue
is another positive root of the rational characteristic polynomial.
Thus the eigenvalues are totally positive, and
Theorem~\ref{p11:ipe:classifier} applies.
\end{proof}

\subsection{An exact criterion for graph constants}

For a finite simple graph $G$ with $m\geq1$ edges, let $L_G$ be its
combinatorial Laplacian and put
\[
       \rho_G=\frac{L_G}{2m},\qquad
       E_G=\exp\{-\operatorname{tr}(\rho_G\log\rho_G)\}.
\]
This graph density matrix and its entropy are classical
\cite{p11:ipe:BGS}. A graph is called Laplacian integral when every
eigenvalue of $L_G$ is an integer.

\begin{corollary}[Graph integrality is the exact arithmetic criterion]
\label{p11:ipe:graph}
Suppose $\operatorname{rank}L_G\geq2$. For every nonzero real
algebraic $c$,
\[
       E_G^c\in\overline{\mathbb Q}
       \quad\Longleftrightarrow\quad
       G\text{ is Laplacian integral and }c\in\mathbb Q.
\]
For the path $P_n$, $E_{P_n}^c$ is transcendental for every $n\geq4$
and every such $c$. For cycles $C_n$, $n\geq3$, one has
\[
       E_{C_n}^c\in\overline{\mathbb Q}
          \quad\Longleftrightarrow\quad
                  n\in\{3,4,6\}\text{ and }c\in\mathbb Q.
\]
The rank-one exception consists of one edge and any number of isolated
vertices; for these graphs $E_G^c=1$.
\end{corollary}
\begin{proof}
The matrix $\rho_G$ is rational and positive semidefinite, with trace
one. The eigenvalues of the integer matrix $L_G$ are algebraic
integers. They are rational exactly when they are integers, so the
first equivalence follows from Corollary~\ref{p11:ipe:quantum}.

The path eigenvalues are $2-2\cos(k\pi/n)$ for $0\leq k<n$.
For $n\geq4$, the first positive eigenvalue lies strictly between
zero and one, hence is not an integer. The cycle eigenvalues are
$2-2\cos(2\pi k/n)$, $0\leq k<n$. At $n=3,4,6$ all are
integers. For $n=5$, the eigenvalue $(5-\sqrt5)/2$ is irrational.
For $n\geq7$, the first positive eigenvalue lies strictly between
zero and one. These spectral formulas follow directly by applying
the Laplacian to the usual discrete cosine vectors for a path and
Fourier vectors for a cycle. Finally, Laplacian rank equals the
number of vertices minus the number of connected components, so
rank one means exactly one two-vertex component and isolated
vertices.
\end{proof}

The criterion propagates through standard graph operations. A
disjoint union, or a Cartesian product of nonempty graphs, is
Laplacian integral exactly when both factors are: the respective
spectra are the multiset union and all pairwise sums, and each
factor has a zero eigenvalue. Also a graph is Laplacian integral
exactly when its complement is, since on the orthogonal complement
of the all-ones vector the two Laplacians sum to $nI$.
Thus one nonintegral example yields whole families of canonical
transcendental effective-rank constants, subject only to the
explicit rank and nonempty-edge conditions above.

\begin{corollary}[Positive rational mixtures of spectral entropies]
\label{p11:ipe:mixture}
Let $p^{(1)},\ldots,p^{(k)}$ satisfy the hypotheses of
Theorem~\ref{p11:ipe:classifier}, and let $w_j>0$ be rational. If any
entry of any $p^{(j)}$ is irrational, then
\[
                  \exp\left(c\sum_{j=1}^k w_j S(p^{(j)})\right)
\]
is transcendental for every nonzero real algebraic $c$.
If all entries are rational, it is algebraic exactly when $c$ is
rational.
\end{corollary}
\begin{proof}
Suppose the displayed exponential is algebraic, and let $q^{(j)}$
be a simultaneous conjugate of the input vectors. Put $d=\sigma(c)$.
Coefficient-only conjugation gives the same logarithm as a weighted
sum of cross-entropies. Their positivity implies $d/c>0$.
Subtracting these two relations and then transporting their bases
as in Theorem~\ref{p11:ipe:positive} gives
\[
 \sum_j w_j\{cD(p^{(j)}\Vert q^{(j)})
                  +dD(q^{(j)}\Vert p^{(j)})\}=0.
\]
All summands have the same nonnegative sign after division by $c$,
so $q^{(j)}=p^{(j)}$ for every $j$, and then $d=c$.
Total positivity allows every automorphism, making all entries and
$c$ rational. Conversely rational data give an algebraic product.
\end{proof}

The total-positivity hypothesis cannot simply be dropped. Let $t=(\sqrt5-1)/2$ and $p=(t,t^2)$. Then
$t+t^2=1$ and
\[
                 E(p)=t^{-(2-t)},\qquad
                 E(p)^{1/(2-t)}=t^{-1}\in\overline{\mathbb Q}.
\]
Thus a positive algebraic vector with irrational entries can have
an algebraic nonzero algebraic power of its effective number of
outcomes. Its nonpositive conjugates prevent the Galois argument.

\subsection{Algebraic-exponent independence of an explicit family}

For each odd prime $r$, retain the preceding manuscript's rational
matrix
\[
 \rho_r=\frac12I_2+2^{-(r+3)/2}
                       \begin{pmatrix}1&1\\1&-1\end{pmatrix},
 \qquad h_r=H_2(\rho_r),\qquad E_r=2^{h_r}.
\]
The previously proved linear independence of $1,h_3,h_5,\ldots$
does not by itself decide algebraicity of products of the $E_r$.
The following strengthens it in that direction.

\begin{theorem}[No algebraic-exponent product is algebraic]
\label{p11:ipe:product}
For any finitely supported real algebraic coefficients $c_r$, not
all zero,
\[
                         \prod_r E_r^{c_r}
                         =2^{\sum_r c_rh_r}
\]
is transcendental. No positivity assumption on the $c_r$ is needed.
Consequently the $E_r$ are multiplicatively independent even modulo
positive algebraic numbers and with arbitrary real algebraic
exponents. This is not algebraic independence of the $E_r$.
\end{theorem}
\begin{proof}
The two eigenvalues are
\[
 \lambda_r^\pm=\frac12\pm\frac{\sqrt2}{2^{(r+3)/2}},\qquad
 \lambda_r^+\lambda_r^-=
                             \frac{2^r-1}{2^{r+2}}.
\]
For distinct odd primes $r,s$, the positive integers $2^r-1$ and
$2^s-1$ are coprime. Choose an odd prime $q\mid2^r-1$.
In $K=\mathbb Q(\sqrt2)$ there is a prime ideal
$\mathfrak q_+$ at which one eigenvalue, say $\lambda_r^+$,
has positive valuation and the other has zero valuation. The
sum of the eigenvalues is $1$, so they cannot both have positive
valuation. Their product proves existence, and their denominators
are powers of two. At the conjugate prime ideal $\mathfrak q_-$,
$\lambda_r^-$ has positive valuation and $\lambda_r^+$ has zero
valuation. These ideals are distinct. All eigenvalues with other
indices have valuation zero at both ideals.

Suppose $A=\prod_r E_r^{c_r}$ were algebraic. Then
\[
 \log A+\sum_r c_r\bigl(\lambda_r^+\log\lambda_r^+
                         +\lambda_r^-\log\lambda_r^-\bigr)=0.
\]
Every algebraic logarithmic relation is an algebraic combination
of rational ones. Each rational relation yields a multiplicative
identity and hence an identity of prime-ideal valuations.
Work in a number field containing $K$ and $A$, and extend each
of $\mathfrak q_+$ and $\mathfrak q_-$ to that field. Evaluating
the relation at the first extended valuation gives
$c_r\lambda_r^+\in\mathbb Q$; at the second it gives
$c_r\lambda_r^-\in\mathbb Q$. Valuations of algebraic numbers
are integers after normalization, and ramification only multiplies
the relevant nonzero valuations by positive integers.
If $c_r\ne0$, division gives
$\lambda_r^+/\lambda_r^-\in\mathbb Q$, which is false because
their sum is $1$ and each is irrational. Thus every $c_r=0$,
a contradiction.
\end{proof}

Together with the preceding manuscript's result, this provides
explicit constants for which every nonzero algebraic linear
combination $L=\sum c_rh_r$ is transcendental and, additionally,
$2^L$ is transcendental. No algebraic independence of $L$ and
$2^L$ is asserted. In particular these statements do not resolve
algebraic independence questions for general Gelfond--Schneider
numbers or for classical gamma and zeta constants.


\section{A unit obstruction and the exact totally-real entropy threshold}
\label{p11:et:section}

The preceding manuscript reduced the remaining three-atom problem to a
stationary value of three Laurent monomials. That reduction can be
strengthened substantially by using finite valuations. We prove that a
nonrational algebraic ternary vector with rational entropy would have
to arise from an algebraic unit with three odd exponents. This excludes
all such vectors in totally real number fields, of arbitrary degree.
Combined with the explicit four-atom example already obtained, the
least alphabet with totally real irrational algebraic probabilities
is exactly four. The corresponding problem over all
number fields remains open in this work.

Throughout, $H_2(p)=-\sum_i p_i\log_2p_i$, all probabilities are positive,
and a nonrational vector means that at least one entry is irrational.
An algebraic unit is an algebraic number whose value and reciprocal
are both algebraic integers.

\subsection{The earlier reduction, with primitive exponents}

\begin{lemma}[Three-monomial reduction]\label{p11:et:reduction}
Suppose $p=(p_1,p_2,p_3)$ is an algebraic, nonrational probability vector
and $h=H_2(p)\in\mathbb Q$. There are a positive algebraic number $w$ and
integers $m_1,m_2,m_3$ such that
\begin{equation}\label{p11:et:critical}
 p_i=2^{-h}w^{m_i},\qquad
 \sum_i w^{m_i}=2^h,\qquad \sum_i m_iw^{m_i}=0,
 \qquad \gcd(m_1,m_2,m_3)=1.
\end{equation}
The logarithms $\log w$ and $\log2$ are rationally independent. The
exponents are pairwise distinct, and their minimum is negative while
their maximum is positive.
\end{lemma}
\begin{proof}
We include the reduction to make the new obstruction self-contained.
Baker's logarithmic independence theorem \cite{p11:TLBaker} implies that
algebraic-coefficient relations among logarithms of positive algebraic
numbers are generated by rational-coefficient relations. Put
\[
 t=\dim_{\mathbb Q}\operatorname{span}_{\mathbb Q}\{p_1,p_2,p_3\},
 \quad
 r=\dim_{\mathbb Q}\operatorname{span}_{\mathbb Q}
       \{\log p_1,\log p_2,\log p_3,\log2\}.
\]
Splitting the entropy relation with respect to a rational basis of the
probabilities, which also spans $h$ because $\sum_i p_i=1$ and
$h\in\mathbb Q$, gives $t$ independent rational relations among the
four logarithms. Thus $t+r\le4$. Nonrationality gives $t\ge2$.
If $r=1$, all the $p_i$ are rational powers of two. Positive rational
powers of two summing to one must have integral exponents: clear a
common denominator $d$, use the independence of
$1,2^{1/d},\ldots,2^{(d-1)/d}$ supplied by Eisenstein's criterion for
$X^d-2$, and observe that all coefficients are nonnegative. Hence
$r=1$ would make all entries rational. Therefore $t=r=2$.

Choose a positive algebraic $u$ with $\log2,\log u$ a rational basis
of the logarithm span, and write
$\log p_i=a_i\log2+b_i\log u$ with rational $a_i,b_i$. Baker's theorem
then gives $\sum_i b_i p_i=0$ and $\sum_i a_i p_i=-h$.
The rational kernel of $(x_i)\mapsto\sum_i x_i p_i$ is one-dimensional;
as $b\ne0$, it follows that $a_i+h=\lambda b_i$ for one rational
$\lambda$. Clearing a common denominator in the $b_i$ gives
\eqref{p11:et:critical}, except possibly for primitivity. Replacing $w$ by
a positive integral power and dividing all $m_i$ by their common gcd
makes the exponents primitive.

If, for example, $m_1=m_2=a$ and $m_3=b\ne a$, stationarity and
normalization imply
$a(p_1+p_2)+bp_3=0$, while $p_1=p_2$. These two linear equations
make every probability rational. If all three exponents coincide,
stationarity makes that exponent zero, contrary to logarithmic
independence. Thus the exponents are distinct. Their weighted mean
zero, with positive weights, forces both a positive and a negative
exponent.
\end{proof}

\subsection{Finite valuations force a unit and odd exponents}

\begin{theorem}[Unit, parity, and norm restrictions]\label{p11:et:unit}
Under the hypotheses of Lemma~\ref{p11:et:reduction}, the normalized
generator $w$ is an algebraic unit, each $m_i$ is odd, and
\begin{equation}\label{p11:et:range}
                         1\le h<\log_2 3.
\end{equation}
Let $K=\mathbb Q(p_1,p_2,p_3)$ and $d=[K:\mathbb Q]$. Then
\begin{equation}\label{p11:et:norm}
              \left|N_{K/\mathbb Q}(p_i)\right|=2^{-hd}
              \quad(i=1,2,3),\qquad hd\in\mathbb Z.
\end{equation}
All three probabilities are irrational.
\end{theorem}
\begin{proof}
Put $C=2^h$ and $L=\mathbb Q(w)$. Since $h>0$ is rational, $C$ is
an algebraic integer: if $h=a/b$ in lowest terms, then $C$ satisfies
$X^b-2^a=0$. Also $C=\sum_iw^{m_i}\in L$.

Let $v$ be any finite additive valuation of $L$. If $v(w)>0$, the
term with the unique smallest exponent has strictly the smallest
valuation in $\sum_iw^{m_i}$, and that valuation is negative because
$\min_i m_i<0$. The ultrametric inequality with a unique minimum
therefore gives $v(C)<0$, contrary to integrality. If $v(w)<0$, the
unique largest exponent gives the same contradiction. Thus $v(w)=0$
at every finite place. Equivalently, $w$ is an algebraic unit.

Choose a prime of $L$ above $2$, reduce the two sum identities in
\eqref{p11:et:critical} in its residue field, and write $x$ for the
nonzero residue of $w$. Since $h>0$, the residue of $C$ is zero.
Thus
\[
             \sum_i x^{m_i}=0,
             \qquad \sum_{i:\ m_i\ {
m odd}}x^{m_i}=0.
\]
If exactly one exponent is odd, the second sum is a single nonzero
term. If exactly two are odd, subtracting the second sum from the
first gives a single nonzero term. Both are impossible. Primitivity
excludes three even exponents, so all three exponents are odd.
Consequently stationarity also gives
\[
                     C=\sum_i(1-m_i)w^{m_i}\in2\mathcal O_L,
\]
because each coefficient $1-m_i$ is even and every power of $w$ is
an algebraic integer. At a valuation normalized by $v(2)=1$, this
says $h=v(C)\ge1$. The strict entropy upper bound in
\eqref{p11:et:range} follows from $H_2(p)\le\log_2 3$, with equality only
for the rational uniform vector.

Primitivity supplies integers $c_i$ with $\sum_i c_i m_i=1$; hence
\[
                   w=\prod_i(Cp_i)^{c_i},\qquad L=K(C).
\]
Every conjugate of $C$ has absolute value $C$, while the absolute norm
of a unit is one. Taking norms of $p_i=C^{-1}w^{m_i}$ from $L$ to
$\mathbb Q$, and then using norm transitivity, proves
$|N_{K/\mathbb Q}(p_i)|=2^{-hd}$. This is rational, so $hd$ is an
integer, by uniqueness of prime factorization applied to a rational
power of two.

Finally, if $p_i$ were rational and positive, its norm identity would
give $p_i=2^{-h}$, so $w^{m_i}=1$. As $w>0$ and $w\ne1$, this forces
$m_i=0$, contradicting oddness.
\end{proof}

The lower bound $h\ge1$ is a new exclusion region for the remaining
ternary problem. It is not a statement that a nonrational ternary
vector of entropy exactly one bit exists.

\begin{proposition}[The generator has very few real conjugates]
\label{p11:et:signature}
Write $h=a/b$ in lowest positive terms. The field $L=\mathbb Q(w)$
from Theorem~\ref{p11:et:unit} has exactly one real embedding if $b$ is
odd and exactly two real embeddings if $b$ is even. In the latter
case the two real conjugates of $w$ are $w$ and $-w$.
In particular, $w$ is not totally real.
\end{proposition}
\begin{proof}
Set $F(x)=\sum_i x^{m_i}$. On the positive real axis,
\[
 \frac{d^2}{dy^2}F(e^y)=\sum_i m_i^2 e^{m_i y}>0.
\]
Thus $F$ has exactly one positive stationary point, namely $w$:
its first derivative in $y$ is strictly increasing and tends to
opposite infinities at the two ends, because the exponents have both
signs. Since all $m_i$ are odd, $F$ is odd on the nonzero real axis.
Every real conjugate of $w$ is a stationary point of $F$, and its
$F$-value is a real conjugate of $C$. Therefore it can only be $w$
or $-w$.

The number $C=2^{a/b}$ has degree $b$: Bezout's identity for $a,b$
shows $\mathbb Q(C)=\mathbb Q(2^{1/b})$, and $X^b-2$ is Eisenstein.
If $b$ is odd, $C$ has just one real conjugate, which is positive;
only $w$ is possible. If $b$ is even, $b$ divides $[L:\mathbb Q]$,
so that degree is even. Its real embeddings consequently have an even
count; as the count is positive and at most two, it is exactly two.

If $w$ were totally real, all its conjugates would be among $w,-w$.
Then $w^2$ would be rational, and also a positive rational algebraic
unit. This forces $w^2=1$, contradicting the rational independence
of $\log w,\log2$.
\end{proof}

\subsection{Closing the quadratic-field case}

\begin{theorem}[No nonrational quadratic ternary vector has rational entropy]
\label{p11:et:quadratic}
Let $K$ be a real quadratic number field and let
$p_1,p_2,p_3\in K$ be positive with sum one. If at least one $p_i$ is
irrational, then $H_2(p_1,p_2,p_3)\notin\mathbb Q$.
More generally, a nonrational ternary vector in any totally real field
cannot have entropy in $\tfrac12\mathbb Z$.
\end{theorem}
\begin{proof}
For the more general statement, suppose $h\in\tfrac12\mathbb Z$ and
let $K_0$ be the totally real field generated by the probabilities.
Theorem~\ref{p11:et:unit} gives a unit $w$ with
\[
                      \mathbb Q(w)=K_0(2^h)
                                    \subseteq K_0(\sqrt2).
\]
The compositum of totally real fields is totally real. This makes
$w$ totally real, contradicting Proposition~\ref{p11:et:signature}.

For a nonrational ternary vector in a real quadratic field, its
probability field has degree two. Formula~\eqref{p11:et:norm} gives
$2h\in\mathbb Z$, and \eqref{p11:et:range} leaves only
$h=1$ or $h=3/2$. Both have just been excluded.
\end{proof}


\subsection{Extension to every totally real number field}

The quadratic argument extends to arbitrary totally real fields. The
needed radical-degree fact has a short proof using the constant term
of a minimal polynomial; no unproved binomial irreducibility assumption
is needed.

\begin{lemma}[A rational radical over a totally real field]
\label{p11:et:radical}
Let $K$ be a totally real number field and let $C=2^{a/b}$, where
$a,b$ are coprime positive integers. Then
\[
 [K(C):K]=
 \begin{cases}
 b/2,& b\text{ is even and }\sqrt2\in K,\\
 b,&\text{otherwise}.
 \end{cases}
\]
In the first case the minimal polynomial is
$X^{b/2}-2^{(a-1)/2}\sqrt2$; in the second it is $X^b-2^a$.
\end{lemma}
\begin{proof}
Let $e=[K(C):K]$, and consider the minimal polynomial of $C$ over
$K$ at the given real embedding of $K$. Its $e$ roots belong to
$\{C\zeta:\zeta^b=1\}$. Its nonzero constant coefficient is real
and has absolute value $C^e$, so $C^e\in K$. A positive rational
power $2^{ae/b}$ has degree $b/\gcd(b,e)$ over $\mathbb Q$, by
Eisenstein and Bezout as in Proposition~\ref{p11:et:signature}. If this
degree exceeds two, some conjugates are nonreal. Since every element
of $K$ is totally real, we obtain $b/\gcd(b,e)\le2$.

As $1\le e\le b$, this forces $e=b$ when $b$ is odd, and
$e\in\{b/2,b\}$ when $b$ is even. The case $e=b/2$ puts
$C^e=2^{a/2}$ in $K$; as $a$ is odd, it forces $\sqrt2\in K$.
Conversely, if $\sqrt2\in K$, the displayed polynomial of degree
$b/2$ vanishes at $C$, so $e\le b/2$ and equality follows.
The degree-$b$ case has the other displayed minimal polynomial.
\end{proof}

\begin{theorem}[Irrational ternary entropy over totally real fields]
\label{p11:et:totallyreal}
Let $p_1,p_2,p_3$ be positive totally real algebraic numbers with sum
one. If at least one entry is irrational, then
\[
                           H_2(p_1,p_2,p_3)\notin\mathbb Q.
\]
\end{theorem}
\begin{proof}
Suppose instead that $h=a/b\in\mathbb Q$, in lowest positive terms.
The probability field $K=\mathbb Q(p_1,p_2,p_3)$ is totally real;
write $d=[K:\mathbb Q]>1$. The previous results give
$L=K(C)=\mathbb Q(w)$, with $C=2^{a/b}$, and a real-embedding count
$r_1(L)$ equal to one for odd $b$ and two for even $b$.

If $b$ is odd, Lemma~\ref{p11:et:radical} makes $X^b-2^a$ the minimal
polynomial over $K$. Under every real embedding of $K$ it has exactly
one real root. Counting extensions of embeddings therefore gives
$r_1(L)=d$, contradicting $r_1(L)=1$.

If $b$ is even and $\sqrt2\notin K$, the same minimal polynomial
has exactly two real roots under each embedding of $K$. Thus
$r_1(L)=2d$, contradicting $r_1(L)=2$.

It remains that $b$ is even and $\sqrt2\in K$. Put $n=b/2$.
The minimal polynomial is
\[
                  X^n-2^{(a-1)/2}\sqrt2.
\]
Exactly half of the $d$ real embeddings of $K$ send $\sqrt2$ to
$\sqrt2$, and half send it to $-\sqrt2$. If $n$ is odd, each
embedding has one real extension; if $n$ is even, the first half
have two real extensions and the second half have none. In both
cases $r_1(L)=d$. Hence $d=2$. The norm restriction $hd\in\mathbb Z$
forces $b\mid2$, so $b=2$ and $L=K(C)=K$ is totally real. This
contradicts the fact that $w$ is not totally real, proved in
Proposition~\ref{p11:et:signature}.
\end{proof}

\begin{corollary}[The exact minimum for totally real probabilities]
\label{p11:et:four}
Among probability vectors all of whose entries are positive totally
real algebraic irrational numbers, the least number of entries allowing
rational entropy in bits is exactly four. The minimum is already
attained within the quadratic field $\mathbb Q(\sqrt5)$.
\end{corollary}
\begin{proof}
One atom is necessarily rational. For two irrational algebraic
probabilities, rational entropy in base $2$ is impossible by Baker's
theorem: splitting the entropy relation with respect to $1,p$ forces
$p$ and $1-p$ to be rational powers of two, and positivity in the
Eisenstein basis used above makes them rational. Three atoms are
excluded by Theorem~\ref{p11:et:totallyreal}.

For the upper bound put $\varphi=(1+\sqrt5)/2$ and use the previously
constructed vector
\[
 \left(\frac{\varphi}{2},\frac{\varphi^{-5}}2,
                   \frac{\varphi^{-4}}2,\frac{\varphi^{-4}}2\right).
\]
All its entries are positive irrational elements of $\mathbb Q(\sqrt5)$.
The identities
\[
 \varphi+\varphi^{-5}+2\varphi^{-4}=2,
 \qquad \varphi-5\varphi^{-5}-8\varphi^{-4}=0
\]
show that they sum to one and have entropy exactly one bit.
\end{proof}

The corollary minimizes over all totally real probability vectors; it
does not assert that four atoms suffice in each prescribed number
field. The global minimum over unrestricted algebraic probabilities
remains either three or four. Any ternary counterexample would now
need a probability field that is not totally real, of degree at least
three, rational entropy in $[1,\log_2 3)$ with denominator dividing
that degree, and the unit and signature restrictions proved above.

\paragraph{Literature and scope.}
The logarithmic input is classical Baker theory. The rational-spectrum
rational-or-transcendental dichotomy appears explicitly in
Blakaj--Wolf \cite[Theorem 2]{p10:entnext:BW}. The finite-place unit obstruction,
parity restriction, and totally-real threshold above are additions
to this research program; the primary-source comparison does not
establish their historical priority. No arithmetic classification of
Euler's constant, the Gompertz constant, or an odd zeta value follows
from this entropy theorem.


\section{A third Lambert--Fibonacci jet and a fourth independent function}
\label{p11:tl:section}
The preceding two-jet argument used one rational cusp at a time.
We now vary the denominator of the cusp.  The residue \(1/N\)
separates ordinary polynomial degree from differential weight.
This proves a three-jet theorem, together with a fourth independent
function supplied by the prime-harmonic series.

Write
\[
 L(q)=\sum_{n\ge1}\frac{q^n}{1-q^n},\qquad
 \mathcal F(q)=\sum_{n\ge1}\frac{q^n}{1-(-1)^nq^{2n}},\qquad
 \mathcal P(q)=\sum_{p\ {\rm prime}}\frac{q^p}{p},
 \qquad D=q\frac{d}{dq}.
\]
For \(\eta\in\{0,1\}\), define
\[
 J_{m,r,\eta}(q)=
 \begin{cases}
 D^r[L(q^m)],&\eta=0,\\
 D^r[\mathcal F(q^m)],&\eta=1.
 \end{cases}
\]
Here \(m\ge1,\ r\ge0\), and differentiation follows composition.

\subsection{The fixed-cusp input}
For a prime \(N\equiv1\pmod4\), a primitive \(N\)-th root \(\xi\),
and \(N\nmid m\), the cusp expansions are
\[
 J_{m,r,\eta}(\xi e^{-t})
 =\frac{a_{m,r,\eta}}{N}t^{-w_r}
       \bigl(U+\gamma-2\log N+\beta_{m,r,\eta}\bigr)
       +O_{N,m,r}(t^{-w_r+1/2}),
 \tag{TL1}\label{p11:tl:expansion}
\]
where
\[
 U=\log(1/t),\quad
 w_r=r+1,\quad
 a_{m,r,\eta}=\frac{r!}{4^\eta m},\quad
 \beta_{m,r,\eta}=H_r-\log m+\eta\pi/2.
\]
We recall the derivation to make the use of the remainder explicit.
For \(\xi\) of order \(N\), the Mellin transform of \(L\) is
\[
 Z_\xi(s)=N^{-2s}\sum_{a,b=1}^{N}
       \xi^{ab}\zeta(s,a/N)\zeta(s,b/N)
 =\frac{1/N}{(s-1)^2}
  +\frac{(2\gamma-2\log N)/N}{s-1}+O(1).
\]
In the inverse Mellin integral for the \(r\)-th radial derivative,
the factor is \(\Gamma(s+r)t^{-s-r}\).
Shifting to \(\Re s=1/2\) gives the power-saving remainder in
\eqref{p11:tl:expansion}.  Its residue contains \(+H_r\).
These are standard Lambert cusp expansions; see
\cite{p10:NLDorigoniKleinschmidt} for a much more extensive treatment.

For the Fibonacci function use the exact identity
\[
 \mathcal F(q)=L(q^2)-L(q^4)
       +\frac{\vartheta(q)^2-\vartheta(-q)^2}{8},
 \qquad \vartheta(q)=\sum_{k\in\mathbb Z}q^{k^2}.
\]
At a primitive odd-prime cusp, Poisson summation and quadratic
Gauss sums give
\[
 \vartheta(\xi e^{-t})^2
     =\frac{(-1)^{(N-1)/2}\pi}{Nt}
                 +O(t^{-1}e^{-c_N/t}),\qquad
 \vartheta(-\xi e^{-t})^2=O(t^{-1}e^{-c_N/t}).
\]
The square of the Gauss sum is independent of the primitive
numerator. Differentiation preserves an exponential saving.
For \(N\equiv1\pmod4\), this gives the extra intercept \(\pi/2\)
in \eqref{p11:tl:expansion}; the two \(\log2\) contributions cancel.

The Mertens theorem in arithmetic progressions, followed by
partial summation, gives at the same cusp
\[
 \mathcal P(\xi e^{-t})
    =-\frac1{N-1}\log U+C_\xi+o(1).
 \tag{TL2}\label{p11:tl:prime}
\]
Indeed, every nonzero class modulo \(N\) contributes
\((N-1)^{-1}\log\log x+M(N,a)+o(1)\), and
\(\sum_{a=1}^{N-1}\xi^a=-1\); the prime \(N\) contributes
only a bounded term.  An explicit statement of the classical
input is \cite[Eq.~(1)]{p10:NLLanguascoZaccagnini}.

\begin{lemma}[Distinct cusp intercepts]
\label{p11:tl:distinct}
The real numbers \(\beta_{m,r,\eta}\) attached to distinct
triples \((m,r,\eta)\) are distinct.
\end{lemma}
\begin{proof}
An equality would imply
\[
 H_r-H_s=\log(m/n)-(\eta-\eta')\pi/2.
\]
If \(m/n\ne1\), the real logarithm \(\log(m/n)\) and
\(\log(-1)=i\pi\) are linearly independent over \(\mathbb Q\).
Baker's theorem~\cite{p11:TLBaker} implies that \(1,\log(m/n),i\pi\) are linearly
independent over \(\overline{\mathbb Q}\), contradicting the
displayed equality. If \(m/n=1\), irrationality of \(\pi\)
forces \(\eta=\eta'\); then strict increase of \(H_r\) forces
\(r=s\).  Thus the two triples coincide.
\end{proof}
The use of Baker's theorem here is only to distinguish the
known intercepts; it does not transfer to the values of \(L\)
or \(\mathcal F\).

\subsection{An injectivity lemma with two gradings}
\begin{lemma}
\label{p11:tl:injective}
Let \(w_1,w_2,w_3\) be positive integers, not all equal, and let
\(\beta_1,\beta_2,\beta_3\) be distinct complex numbers.
On the vector space of polynomials homogeneous of ordinary
degree \(d\) and weighted degree \(W\), the substitution
\[
 (X_1,X_2,X_3)\longmapsto
       (V+\beta_1,V+\beta_2,V+\beta_3)
\]
is injective. The assertion remains valid when the coefficients
belong to any purely transcendental extension of \(\mathbb C\).
\end{lemma}
\begin{proof}
The exponent vectors of the allowed monomials satisfy
\[
 e_1+e_2+e_3=d,\qquad
 w_1e_1+w_2e_2+w_3e_3=W.
\]
The integer kernel of these two independent equations has
rank one. Hence, after choosing one exponent vector \(e^0\),
every other vector is \(e^0+k\nu\), for a fixed nonzero integer
vector \(\nu\). Factoring out the image of the first monomial,
a proposed relation becomes a Laurent polynomial in
\[
 R(V)=\prod_{i=1}^{3}(V+\beta_i)^{\nu_i}.
\]
The rational function \(R\) is nonconstant, since its distinct
zeros and poles have orders \(\nu_i\), not all zero.
It is therefore transcendental over the coefficient field,
which does not contain \(V\). A nonzero Laurent polynomial
in \(R(V)\) cannot vanish.
\end{proof}

\begin{theorem}[Three jets, and four functions with primes]
\label{p11:tl:main}
Choose three distinct triples \((m_i,r_i,\eta_i)\),
\(1\le i\le3\), whose derivative orders \(r_i\) are not all equal.
Then
\[
 \boxed{J_{m_1,r_1,\eta_1}(q),\
        J_{m_2,r_2,\eta_2}(q),\
        J_{m_3,r_3,\eta_3}(q),\
        \mathcal P(q)
        \text{ are algebraically independent over }\mathbb C(q).}
\]
In particular, either of the triples
\[
 L,\ DL,\ D^2L
 \qquad\hbox{and}\qquad
 \mathcal F,\ D\mathcal F,\ D^2\mathcal F
\]
is algebraically independent over \(\mathbb C(q)\).
Thus neither \(L\) nor \(\mathcal F\) satisfies a nonzero
algebraic differential equation of order at most two over
\(\mathbb C(q)\).
\end{theorem}
\begin{proof}
Suppose
\[
 P(q,X_1,X_2,X_3,Z)\in\mathbb C[q,X_1,X_2,X_3,Z]
\]
is a nonzero polynomial giving an identity after the indicated
substitution; rational denominators have been cleared.
Let \(M\) bound its degree in \(q\).
Give \(X_i\) weight \(w_i=r_i+1\), and give \(Z\) weight zero.
Let \(W\) be the largest weight occurring, and decompose that
part by ordinary degree in the three \(X\)-variables:
\[
 P_W=\sum_d P_{W,d}.
\]
Set \(a_i=a_{m_i,r_i,\eta_i}\),
\(\beta_i=\beta_{m_i,r_i,\eta_i}\), and define
\[
 Q_d(q,V,Z)=P_{W,d}
        \bigl(q,a_1(V+\beta_1),a_2(V+\beta_2),
                    a_3(V+\beta_3),Z\bigr).
\]
Take any sufficiently large prime \(N\equiv1\pmod4\) not
dividing \(m_1m_2m_3\), and any primitive \(N\)-th root \(\xi\).
For this fixed \(N\), put \(V=U+\gamma-2\log N\).
The assumed identity and \eqref{p11:tl:expansion} imply
\[
 \sum_d N^{-d}
        Q_d\bigl(\xi,V,\mathcal P(\xi e^{-t})\bigr)
              +o(1)=0. \tag{TL3}\label{p11:tl:dominance}
\]
All error terms really are \(o(1)\): after multiplication by
\(t^W\), the jet error has a factor \(t^{1/2}\), lower weights
have a positive integer power of \(t\), and the remaining
factors grow only as fixed powers of \(U\) and \(\log U\).
Replacing the polynomial coefficients at \(q=\xi e^{-t}\)
by their values at \(\xi\) contributes only \(O(t)\) times
fixed powers of \(U\) and \(\log U\), again \(o(1)\).

The polynomial
\[
 Q_{N,\xi}(V,Z)=\sum_dN^{-d}Q_d(\xi,V,Z)
\]
must be zero. Otherwise, its highest power of \(V\), followed
by the highest power of \(Z\) in that coefficient, would give
a unique nonzero leading term
\(cV^j(\log V)^k\) after \eqref{p11:tl:prime} is inserted.
Such a term cannot satisfy \eqref{p11:tl:dominance}.
No coefficient extraction from the \(o(1)\) prime remainder
is needed; only its leading \(\log V\) growth is used.

Now use the \emph{exact root count}.  Every coefficient of
\(V^jZ^k\) in \(\sum_dN^{-d}Q_d(q,V,Z)\) is a polynomial in
\(q\) of degree at most \(M\). It vanishes at all \(N-1\)
primitive \(N\)-th roots. For \(N-1>M\), it vanishes identically.
Consequently
\[
 \sum_dN^{-d}Q_d(q,V,Z)=0
\]
as a polynomial identity for infinitely many primes \(N\).
Regarding this as a polynomial in \(1/N\), one obtains
\(Q_d(q,V,Z)=0\) for every \(d\).
By Lemmas~\ref{p11:tl:distinct} and \ref{p11:tl:injective}, over the
coefficient field \(\mathbb C(q,Z)\), each \(P_{W,d}\) must
vanish. This contradicts the definition of \(W\).
\end{proof}

\begin{remark}[A genuine increment and its scope]
The previous manuscript established independence of two jets
and, after adjoining \(\mathcal P\), three functions.
Theorem~\ref{p11:tl:main} increases these counts to three jets and
four functions whenever the derivative orders are not all equal.
It does not establish full differential transcendence or
independence of every finite collection of jets. It also says
nothing new about the transcendence of \(L(1/2)\), the
Fibonacci reciprocal sum, or the irrationality of \(B_1\).
A complete historical priority check for this three-jet
extension has not been completed.
\end{remark}

\subsection{An explicit obstruction at the next order}
The next proposition explains why the leading cusp family
alone cannot prove full differential transcendence.

\begin{proposition}[An order-three identity of the cusp model]
\label{p11:tl:barrier}
Define
\[
 \begin{split}
 \mathscr R(A,B,C,E)={}&A^2E^2-23ABCE-44B^2C^2\\
                   &+36B^3E+25AC^3.
 \end{split}
\]
It is homogeneous of ordinary degree \(4\) and weighted degree
\(10\) for weights \(1,2,3,4\), and
\[
 \mathscr R\bigl(a x v,\ a x^2(v+1),\
             a x^3(2v+3),\ a x^4(6v+11)\bigr)=0
 \tag{TL4}\label{p11:tl:modelrelation}
\]
identically in \(a,x,v\).
Nevertheless,
\[
 \mathscr R(L,DL,D^2L,D^3L)=-5q^4+O(q^5),
\]
so \eqref{p11:tl:modelrelation} is not a differential equation
for \(L\).
\end{proposition}
\begin{proof}
For the four arguments in \eqref{p11:tl:modelrelation}, elimination
of \(x\) uses
\[
 A x^2-3B x+C=0,\qquad
 4B x^2-5C x+E=0.
\]
The resultant is
\[
 (AE-4BC)^2-(12B^2-5AC)(5C^2-3BE),
\]
which expands to \(\mathscr R\).
The two homogeneity assertions are immediate.
Finally \(D^rL=q+O(q^2)\) for every \(r\), and the sum of the
five coefficients in \(\mathscr R\) is \(-5\).
\end{proof}
To exclude higher-order differential equations one must use
additional information beyond the two displayed cusp terms.
The proposition is a barrier to this proof mechanism, not a
barrier theorem for all possible methods.

\subsection{Why the stronger functional theorem does not decide a value}
\begin{proposition}[Specialization can destroy all finite-jet information]
\label{p11:tl:specialization}
Let \(\alpha\in\overline{\mathbb Q}\) lie inside the common
analytic domain of functions \(f_0,\ldots,f_s\) which are
algebraically independent over \(\mathbb C(q)\).
For arbitrary prescribed algebraic values \(c_0,\ldots,c_s\),
the functions
\[
 g_i(q)=(q-\alpha)f_i(q)+c_i
\]
are still algebraically independent over \(\mathbb C(q)\),
but \(g_i(\alpha)=c_i\).
If the original functions have algebraic Taylor coefficients,
so do the new ones.

More specifically, if \(f,Df,\ldots,D^Rf\) are independent,
then the first \(R+1\) jets of
\(g=(q-\alpha)^{R+1}f\) remain independent over
\(\mathbb C(q)\), while
\(g(\alpha)=Dg(\alpha)=\cdots=D^Rg(\alpha)=0\).
\end{proposition}
\begin{proof}
The first change of variables is invertible over
\(\mathbb C(q)\). In the second statement, Leibniz' rule gives
an invertible triangular linear transformation of the jets,
whose diagonal is \((q-\alpha)^{R+1}\).
The vanishing at \(\alpha\) follows from the zero multiplicity.
\end{proof}
This elementary proposition identifies the missing bridge:
functional independence, rational Taylor coefficients, and
even independent finite jets do not constitute a
specialization theorem. The arithmetic targets remain separate
from the new functional result.



\section{Positive determinant certificates and the algebraic-degree barrier}
\label{p11:zf:section}

The preceding research program supplied rapidly shrinking rational
intervals, but not a sufficient arithmetic margin. We now examine the
stronger polynomial certificates supplied by positive affine matrix
pencils. Two conclusions are obtained. A single off-axis determinant
controls the passage from irrationality to exclusion of higher algebraic
degrees. A separate interpolation construction removes the auxiliary
$\pi^2$ from a Catalan moment matrix while retaining positivity and full
degree. Neither conclusion is a new solution of a named-constant problem.

The external motivation is Fauzan's 2026 determinant construction and
Lai's subsequent independence argument \cite{p11:zf:Fauzan,p11:zf:Lai}.
For this investigation we read Fauzan's actual manuscript, including
Sections 2--4 and Appendix C, rather than relying only on its abstract.
Its complete proof and formal verification are not independently
certified here. The relative-matrix factorization used below already
appears in its Appendix C; our purpose is to extract an explicit
degree-sensitive criterion and test how sharp its obstruction is.

\subsection{A definite affine pencil has only real zeros}

\begin{lemma}[Spectrum and exact degree]\label{p11:zf:spectrum}
Let $A,B$ be rational symmetric $h\times h$ matrices, and suppose
$G=A+\theta B$ is positive definite for a real number $\theta$.
Put $P(X)=c\det(A+XB)$, where $c\ne0$ is rational. If
$C=G^{-1/2}BG^{-1/2}$ has nonzero eigenvalues
$\lambda_1,\ldots,\lambda_r$, then
\[
 P(X)=P(\theta)\prod_{j=1}^r(1+(X-\theta)\lambda_j).
\]
Consequently $\deg P=\operatorname{rank}B=r$, every zero of $P$ is
real, and the number of zeros below or above $\theta$, counted with
multiplicity, is the positive or negative inertia of $B$, respectively.
Moreover
\begin{equation}\label{p11:zf:budget}
 \log|P(\theta+i)|=\log|P(\theta)|
                    +\frac12\sum_{j=1}^r\log(1+\lambda_j^2).
\end{equation}
\end{lemma}
\begin{proof}
Factor $A+XB=G^{1/2}(I+(X-\theta)C)G^{1/2}$ and apply the spectral
theorem to the real symmetric matrix $C$. Congruence preserves inertia.
The zero corresponding to $\lambda_j$ is $\theta-1/\lambda_j$.
Taking absolute values at $X=\theta+i$ proves \eqref{p11:zf:budget}.
\end{proof}

In particular the off-axis quantity is an entire spectral budget,
not just the largest eigenvalue. One may also compute it through
positive matrices:
\[
 |\det(G+iB)|^2=\det G\,
                     \det(G+BG^{-1}B).
\]
Both matrices on the right are positive definite. This identity follows
from $\det(I+C^2)$ and does not require a choice of eigenvectors.

\begin{theorem}[A single off-axis test for algebraic degree]
\label{p11:zf:degree}
Let $\theta$ be real. Suppose $P_n\in\mathbb Z[X]$ have only real
zeros, $P_n(\theta)\ne0$, and $\deg P_n=o(T_n)$ for some
$T_n\to\infty$. Assume
\[
 \limsup_n T_n^{-1}\log|P_n(\theta)|\le-\tau<0,
 \qquad
 \limsup_n T_n^{-1}\log|P_n(\theta+i)|\le\eta.
\]
Then $\theta$ cannot be algebraic of degree $d$ whenever
\begin{equation}\label{p11:zf:threshold}
                         \tau>(d-1)\eta.
\end{equation}
In particular, $\eta\le0$ would imply transcendence.
\end{theorem}
\begin{proof}
For a real $u$ and arbitrary complex $z$, Cauchy--Schwarz gives
\[
 |z-\theta-u|\le
           \sqrt{1+|z-\theta|^2}\sqrt{1+u^2}.
\]
Factor a real-rooted polynomial $P$ into its linear factors and put
$u=r-\theta$ for each root $r$. Thus, if $h=\deg P$,
\begin{equation}\label{p11:zf:anchor}
 |P(z)|\le(1+|z-\theta|^2)^{h/2}|P(\theta+i)|.
\end{equation}
Suppose now that $\theta$ has degree $d$ with primitive minimal
polynomial $f$ of leading coefficient $a$, and conjugates
$\theta=\theta_1,\ldots,\theta_d$. Since $P_n(\theta)\ne0$,
$f$ and $P_n$ are coprime. Their resultant is a nonzero integer:
\[
 1\le\left|a^{\deg P_n}\prod_{j=1}^dP_n(\theta_j)\right|.
\]
Apply \eqref{p11:zf:anchor} to the $d-1$ other factors. The logarithm of
the right side is at most
\[
 \log|P_n(\theta)|+(d-1)\log|P_n(\theta+i)|+O_\theta(\deg P_n).
\]
After division by $T_n$, condition \eqref{p11:zf:threshold} makes its
upper limit negative, a contradiction.
\end{proof}

This is a specialized resultant criterion, not a replacement for
classical transcendence theory. Its practical advantage for a positive
pencil is that the missing arithmetic information is reduced to one
explicit off-axis determinant, rather than separate bounds at unknown
conjugates.

\subsection{The barrier is sharp even for positive Hankel pencils}

\begin{theorem}[An exact quadratic counterexample]
\label{p11:zf:pell}
Put $\lambda=1+\sqrt2$ and define positive integers $p_m,q_m$ by
\[
            p_m+q_m\sqrt2=\lambda^{2m+1}\qquad(m\ge0).
\]
For $n\ge1$, the primitive squarefree integer polynomial
\[
                P_n(X)=\prod_{m=n}^{2n-1}(q_mX-p_m)
\]
has degree $n$, all its zeros are real and less than $\sqrt2$, and
\begin{align}
 P_n(\sqrt2)&=\lambda^{-3n^2}>0,\label{p11:zf:pellsmall}\\
 \log|P_n(\sqrt2+i)|&=3n^2\log\lambda+O(n),\label{p11:zf:pellanchor}\\
 N_{\mathbb Q(\sqrt2)/\mathbb Q}(P_n(\sqrt2))&=(-1)^n.
 \label{p11:zf:pellnorm}
\end{align}
Up to a nonzero rational scalar, $P_n$ is the determinant of a
rational affine Hankel pencil which is positive definite at $\sqrt2$.
Thus positivity, full degree, simple real zeros, and quadratic
exponential decay do not by themselves exclude degree two.
\end{theorem}
\begin{proof}
Conjugation gives $p_m-q_m\sqrt2=-\lambda^{-(2m+1)}$, so
$p_m^2-2q_m^2=-1$, $\gcd(p_m,q_m)=1$, and
$q_m\sqrt2-p_m=\lambda^{-(2m+1)}$. The roots $p_m/q_m$ increase
strictly to $\sqrt2$. Each factor is primitive, so Gauss's lemma
makes their product primitive. Summing the exponents gives
\[
                  \sum_{m=n}^{2n-1}(2m+1)=3n^2,
\]
which proves \eqref{p11:zf:pellsmall}; the norm of each factor is $-1$.
Also $q_m=\lambda^{2m+1}(1+O(\lambda^{-4m-2}))/(2\sqrt2)$,
and
\[
 |q_m(\sqrt2+i)-p_m|=\sqrt{q_m^2+\lambda^{-4m-2}}.
\]
Multiplication proves \eqref{p11:zf:pellanchor}.

Choose distinct rational nodes $x_m=m$ and form
\[
 H_n(X)_{ab}=\sum_{m=n}^{2n-1}(q_mX-p_m)x_m^{a+b},
                     \qquad0\le a,b<n.
\]
Writing $V_{am}=x_m^a$ gives
$H_n(X)=V\operatorname{diag}(q_mX-p_m)V^T$. The Vandermonde
matrix is invertible. Thus $H_n(\sqrt2)$ is positive definite and
$\det H_n(X)=(\det V)^2P_n(X)$. Primitive normalization removes
the rational scalar. Finally \eqref{p11:zf:pellsmall} and
\eqref{p11:zf:pellanchor} give $\tau=\eta=3\log\lambda$ on the scale
$T_n=n^2$, exactly the boundary of \eqref{p11:zf:threshold} for $d=2$.
\end{proof}

\begin{proposition}[Amplification does not improve this ratio]
\label{p11:zf:amplification}
Suppose finitely many primitive real-rooted polynomial families
$P_{j,n}$ have limits
\[
 -T_n^{-1}\log|P_{j,n}(\theta)|\to\tau_j>0,\qquad
 T_n^{-1}\log|P_{j,n}(\theta+i)|\to\eta_j>0.
\]
For fixed nonnegative integers $a_j$, not all zero, their product
$R_n=\prod_jP_{j,n}^{a_j}$ is primitive, and its ratio of these two
rates is
\[
       \frac{\sum_j a_j\tau_j}{\sum_j a_j\eta_j}
                         \le\max_j\frac{\tau_j}{\eta_j}.
\]
At each fixed index, the same weighted-average identity holds for
arbitrary nonnegative powers. Consequently powers, direct sums of
determinant pencils, and Kronecker products do not improve this ratio
merely by amplification.
\end{proposition}
\begin{proof}
Gauss's lemma proves primitivity of the product; its two logarithms
add exactly. The ratio is a weighted average with positive weights
$a_j\eta_j$. A direct sum multiplies determinants. The determinant
of a Kronecker product is a product of powers of the original
determinants, so the same exact observation applies.
\end{proof}

\subsection{What the published quintic-zeta bounds currently certify}

For the normalization in Fauzan's Appendix C, put $\xi=\zeta(5)$.
The cited bounds have degree $37K/40$, value at most
$\exp(-79K^2/1600)$, and relative matrix norm at most
$\exp(6933K/500)$ for all sufficiently large admissible $K$
\cite[Appendix C]{p11:zf:Fauzan}. Lemma~\ref{p11:zf:spectrum} therefore gives
\[
 \tau=\frac{79}{1600},\qquad
 \eta=\frac{37}{40}\frac{6933}{500}-\frac{79}{1600}
           =\frac{511067}{40000},\qquad
 \frac\tau\eta=\frac{1975}{511067}<0.003865.
\]
These particular estimates certify degree one exclusion, but not
degree two exclusion. This is not a lower bound on the true
off-axis growth.

There is also a concrete content target. If the displayed integer
polynomial factors as $Q_K=c_KP_K$ with $P_K$ primitive and
\[
               \liminf K^{-2}\log c_K\ge\chi,
\]
then the same estimates give $\tau_P\ge\tau+\chi$ and
$\eta_P\le\eta-\chi$. Thus a sufficient additional content bound
for degree-$d$ exclusion is
\begin{equation}\label{p11:zf:contenttarget}
                  \chi>\frac{(d-1)\eta-\tau}{d}.
\end{equation}
For degree two this asks for $\chi>127273/20000=6.36365$.
No such content estimate is proved here. Passing to the primitive
polynomial is useful only when its quantitative common divisor is
actually controlled; a formal renormalization does not establish
\eqref{p11:zf:contenttarget}.

\subsection{A full-degree Catalan pencil with no auxiliary period}

Use the positive measure
\[
 d\nu_2(t)=4\pi y^2\sum_{\ell\ge1}\ell e^{-2\pi\ell y}\,dy,
 \qquad t=y^2,
\]
whose moments are $\mu_2(t^e)=(-1)^eB_{2e+2}$ and whose Stieltjes
transform satisfies
\[
 \mu_2((t+c^2)^{-1})=c\zeta(2,c)-1-\frac1{2c}.
\]
These are classical Hurwitz moment identities \cite{p11:zf:Prevost}.
Let $G=\sum_{k\ge0}(-1)^k/(2k+1)^2$, and set, for $j\ge0$,
\[
 a_j=j+\frac14,\qquad b_j=j+\frac34,
 \qquad K_j(t)=\frac1{a_j(t+a_j^2)}-
                                 \frac1{b_j(t+b_j^2)}.
\]
Notice that $K_j(t)>0$ on $[0,\infty)$, since
\[
 K_j(t)=\frac{(b_j-a_j)(t+a_j^2+a_jb_j+b_j^2)}
               {a_jb_j(t+a_j^2)(t+b_j^2)}.
\]

\begin{theorem}[Paired interpolation cancels the nuisance period]
\label{p11:zf:catalan}
For every integer $m\ge1$, let $f_0,\ldots,f_{m-1}\in\mathbb Q[t]$
be the unique polynomials of degree less than $2m$ satisfying
\[
 f_i(-a_j^2)=f_i(-b_j^2)=\delta_{ij}
                          \qquad(0\le i,j<m).
\]
Put $W_m=\sum_{j=0}^{m-1}K_j$. There is an explicitly computable
rational symmetric matrix $A_m$ such that
\begin{equation}\label{p11:zf:catalanmatrix}
 \left(\int W_m(t)f_i(t)f_j(t)d\nu_2(t)\right)_{i,j<m}
                              =A_m+16G I_m.
\end{equation}
The matrix on the left is positive definite. Hence
\[
                    D_m(X)=\det(A_m+16XI_m)
\]
has exact degree $m$, only real zeros, and every zero is strictly
less than $G$. The coefficient of $G$ has full rank $m$ even though
the coefficient of the auxiliary $\pi^2$ vanishes identically.
\end{theorem}
\begin{proof}
Interpolation at the $2m$ distinct rational nodes gives the $f_i$
uniquely. Explicitly $f_i$ is the sum of the two cardinal Lagrange
polynomials for its paired nodes. The classical identities
\[
 \zeta(2,1/4)=\pi^2+8G,\qquad
 \zeta(2,3/4)=\pi^2-8G
\]
follow by separating the odd reciprocal squares into their two
classes modulo four. The Hurwitz shift relation therefore gives
\[
 \mu_2((t+a_j^2)^{-1})=a_j\pi^2+8a_jG+r_j^+,
 \quad
 \mu_2((t+b_j^2)^{-1})=b_j\pi^2-8b_jG+r_j^-,
\]
with explicit rationals
\[
 r_j^+=-16a_j\sum_{k=0}^{j-1}(4k+1)^{-2}-1-(2a_j)^{-1},
 \quad
 r_j^-=-16b_j\sum_{k=0}^{j-1}(4k+3)^{-2}-1-(2b_j)^{-1}.
\]
For a polynomial $F$, division by $t+c^2$ leaves remainder
$F(-c^2)$; the quotient has rational moment under $\mu_2$.
Thus the coefficients of $\pi^2$ and $G$ in
$\mu_2(K_jf_if_k)$ are respectively
\[
 f_i(-a_j^2)f_k(-a_j^2)-f_i(-b_j^2)f_k(-b_j^2)=0,
\]
and
\[
 8\bigl(f_i(-a_j^2)f_k(-a_j^2)+
             f_i(-b_j^2)f_k(-b_j^2)\bigr)=16\delta_{ij}\delta_{kj}.
\]
Summation proves \eqref{p11:zf:catalanmatrix}, including rationality
and computability of $A_m$. The $f_i$ are linearly independent by
evaluation at the nodes. Since $W_m>0$ and $\nu_2$ has positive
density, their Gram matrix is positive definite. Lemma~\ref{p11:zf:spectrum}
with $B=16I_m$ proves every assertion about $D_m$.
\end{proof}

For $m=1$, the polynomial is $16X-88/9$, whose primitive multiple
is $18X-11$. For larger $m$ the construction is a rationally
computable algebraic lower certificate. We have not proved a useful
asymptotic smallness estimate for its primitive value at $G$, nor
the off-axis estimate needed in Theorem~\ref{p11:zf:degree}. The
construction closes the auxiliary-period and nonvanishing problems
for this family, not the transcendence problem.

\subsection{Six individual target outcomes}
\begin{enumerate}
\item $\zeta(3)$ transcendence: the positive-pencil test applies, but
the preceding orthogonal-polynomial cancellation does not supply
a sufficient off-axis budget. No new algebraic degree is excluded.
\item $\zeta(5)$ transcendence: the actual published coefficient and
relative-matrix bounds were inspected. Their exact degree ratio and
the sufficient new content target are displayed above. No higher
degree exclusion follows from those estimates.
\item $\zeta(7)$ irrationality: positive odd-zeta moment pencils have
real-rooted determinant polynomials, but no sufficient primitive-value
decay estimate is established.
\item $\zeta(9)$ irrationality: the same spectral criterion applies;
positivity alone cannot replace the missing local denominator estimate.
\item $\zeta(11)$ irrationality: likewise, no individual irrationality
conclusion follows. An existing disjunctive statement about several
odd zeta values is not an individual proof.
\item Catalan transcendence: Theorem~\ref{p11:zf:catalan} removes the
auxiliary period without loss of determinant degree. Primitive height
and off-axis growth remain to be controlled. The reported external
irrationality result is not counted as our result.
\end{enumerate}

Historical priority is not claimed for the specialized resultant
criterion, spectral facts, or interpolation formulation. The exact
Pell example demonstrates why the identified arithmetic barrier is
substantive even under all the displayed positivity hypotheses.


% All labels and bibliography keys in this fragment begin with tg:.
\section{An exact parameter law for primitive Gompertz cancellation}
\label{p11:tg:section}

The preceding index-sparsity theorem leaves open a geometrically
sparse sequence of unusually large common divisors for a fixed
Euler--Gompertz constant.  We now change the variable in the local
arithmetic analysis.  Roots as a function of the integer index need
not have proved simple derivatives.  In contrast, roots as a function
of the Laguerre parameter are always simple.  This distinction gives
an unconditional distribution law for the complete primitive content.
The resulting statement is about parameter averages; its limitation
at a prescribed parameter, especially $A=1$, is stated explicitly.

For positive integers $A,n$, use the classical family
\[
 \delta_A=\int_0^\infty\frac{e^{-t}}{1+At}\,dt,
 \qquad
 Q_n(A)=\sum_{j=0}^n j!\binom nj^2A^j
          =n!A^nL_n(-1/A).
\]
Let $P_0=0$, $P_1=1$, and give both $P_n,Q_n$ the recurrence
\[
 U_n=(1+(2n-1)A)U_{n-1}-(n-1)^2A^2U_{n-2}.
\]
The earlier exact-content calculation gives
\begin{equation}\label{p11:tg:content}
 g_n(A):=\gcd(P_n(A),Q_n(A))
       =\gcd(Q_n(A),D_n),\qquad D_n=((n-1)!)^2.
\end{equation}
This follows from adjacent-denominator coprimality and the determinant
$P_nQ_{n-1}-P_{n-1}Q_n=A^{2n-2}((n-1)!)^2$.
In particular $g_n(A)$ is a periodic function of $A$ modulo $D_n$.
The Laguerre approximations themselves are classical
\cite{p10:ng:Lagarias,p10:ng:MY}.

\subsection{Every parameter root is simple}

\begin{lemma}\label{p11:tg:simple}
Fix $n\ge1$ and a prime $p$, and put $r=n\bmod p$, $0\le r<p$.
As polynomials over $\mathbb F_p$,
\[
                    Q_n(A)=Q_r(A).
\]
The polynomial $Q_r(A)$ has degree $r$, constant term one, and
has no repeated root in an algebraic closure of $\mathbb F_p$.
If
\[
 \rho_p(n)=\#\{a\in\mathbb F_p:Q_n(a)=0\},
\]
then $0\le\rho_p(n)\le r\le p-1$, and for every $k\ge1$
there are exactly $\rho_p(n)$ roots of $Q_n(A)$ modulo $p^k$.
\end{lemma}
\begin{proof}
The terms with $j\ge p$ vanish modulo $p$ because of $j!$.
For $j<p$, Lucas's congruence gives $\binom nj\equiv\binom rj$,
which proves the polynomial identity.  Since $r!$ is nonzero modulo
$p$, its degree is exactly $r$.

For $r=0$ there are no roots.  For $1\le r<p$, the identity
$Q_r(A)=r!A^rL_r(-1/A)$ reduces simplicity to that of $L_r(x)$
at nonzero points.  The classical Laguerre differential equation
\cite[Table 18.8.1]{p11:tg:DLMF} is
\[
                 xL_r''+(1-x)L_r'+rL_r=0;
\]
it remains valid over $\mathbb F_p$, since the denominators in
$L_r$ are invertible.  Also $L_r(0)=1$.
If a nonzero root had multiplicity $m\ge2$, then
$2\le m\le r<p$.  The term of degree $m-2$ in the differential
equation would have nonzero coefficient
$x_0m(m-1)c$, where $c$ is the leading local coefficient.
The other terms have higher vanishing order, a contradiction.
The case $p=2$ is automatic since $r\le1$.

Thus every root modulo $p$ has a nonzero parameter derivative.
Taylor expansion shows that a root $a$ modulo $p^k$ has exactly
one lift $a+tp^k$ modulo $p^{k+1}$: the choice of $t\bmod p$
solves a linear equation with nonzero coefficient $Q_n'(a)$.
Induction proves the final assertion.
\end{proof}

For completeness, the same reduction proves the adjacent-denominator
coprimality used in \eqref{p11:tg:content}. If a prime $p$ divides both
$Q_n(A)$ and $Q_{n-1}(A)$, then $p\nmid A$ because each $Q_k(A)$
is congruent to one modulo $A$. Put $r=n\bmod p$. The case $r=0$
is impossible since $Q_n\equiv1\pmod p$. Otherwise the common root
is a root of $Q_r$ and $Q_{r-1}$ modulo $p$. If $r=1$, this
contradicts $Q_0=1$. If $r\ge2$, the recurrence, whose coefficient
$(r-1)^2A^2$ is nonzero modulo $p$, forces the root to be a root
of $Q_{r-2}$. Descent again reaches $Q_0=1$. Thus
$\gcd(Q_n(A),Q_{n-1}(A))=1$. The displayed adjacent determinant,
together with $\gcd(Q_n(A),A)=1$, now gives exactly
$\gcd(P_n(A),Q_n(A))=\gcd(Q_n(A),D_n)$.

The squarefreeness input is an elementary consequence of the
classical differential equation, rather than a new theorem about
Laguerre zeros.  Its use here eliminates the exceptional-root
hypothesis that remained in the index-variable analysis.

\subsection{Independent truncated geometric valuations}

For $p\mid D_n$, write $e_p=2v_p((n-1)!)$.
Let $A$ be uniform on any complete residue system modulo $D_n$,
and set $V_p=v_p(g_n(A))$.

\begin{theorem}[The exact parameter distribution]\label{p11:tg:law}
The random variables $V_p$, for $p\mid D_n$, are independent,
and their distributions are determined by
\begin{equation}\label{p11:tg:tail}
 \mathbb P(V_p\ge k)=\frac{\rho_p(n)}{p^k}
                      \qquad(1\le k\le e_p).
\end{equation}
In particular,
\begin{align}
 \mathbb E\log g_n(A)
 &=\sum_{p\mid D_n}
      \frac{\rho_p(n)\log p}{p-1}(1-p^{-e_p})
       \le\vartheta(n-1),\label{p11:tg:mean}\\
 \mathbb E[g_n(A)^s]
 &=\prod_{p\mid D_n}\left[
      1+\frac{\rho_p(n)(p^s-1)}p
                  \sum_{j=0}^{e_p-1}p^{(s-1)j}\right]
                      \qquad(s>0),\label{p11:tg:moment}
\end{align}
where $\vartheta(y)=\sum_{p\le y}\log p$.
\end{theorem}
\begin{proof}
The Chinese remainder theorem makes the residues of $A$ modulo the
distinct $p^{e_p}$ independent and uniform.  Lemma~\ref{p11:tg:simple}
counts exactly $\rho_p(n)p^{e_p-k}$ roots modulo $p^{e_p}$ of the
congruence $Q_n(A)\equiv0\pmod{p^k}$, proving~\eqref{p11:tg:tail}.
Sum the tail probabilities to get $\mathbb EV_p$, and use
\[
 p^{sV_p}=1+\sum_{k=1}^{e_p}
      (p^{sk}-p^{s(k-1)})\mathbf1_{\{V_p\ge k\}}
\]
to obtain~\eqref{p11:tg:moment}.  The inequality in~\eqref{p11:tg:mean}
uses $\rho_p(n)\le p-1$.
\end{proof}

\begin{corollary}[Exponential moment bound]\label{p11:tg:mgf}
For every fixed $0<s<1$,
\begin{equation}\label{p11:tg:mgfbound}
       \log\mathbb E[g_n(A)^s]
           \le s\vartheta(n-1)+O_s(n^s)=sn+o_s(n).
\end{equation}
Consequently, for every real $T$,
\begin{equation}\label{p11:tg:markov}
 \mathbb P\{\log g_n(A)\ge\vartheta(n-1)+T\}
                         \le\exp(-sT+O_s(n^s)).
\end{equation}
\end{corollary}
\begin{proof}
Extend the geometric sum in~\eqref{p11:tg:moment} to infinity and
use $\rho_p(n)\le p-1$.  The resulting factor is at most
\[
 1+\frac{(p-1)(p^s-1)}{p-p^s}
                 \le\frac{p^s}{1-p^{s-1}}.
\]
After taking logarithms, the additional error is at most
\[
 \sum_{p\le n-1}-\log(1-p^{s-1})
       \le C_s\sum_{j=2}^{n-1}j^{s-1}=O_s(n^s).
\]
The prime number theorem gives $\vartheta(n-1)=n+o(n)$.
Markov's inequality proves~\eqref{p11:tg:markov}.
\end{proof}

\subsection{Uniform exclusion for most finite parameter choices}

\begin{theorem}[A finite simultaneous parameter bound]\label{p11:tg:finite}
For every $\eta>0$, there are positive constants $C,c$ and $M_0$
such that, for all $N\ge M\ge M_0$, the fraction of residues
$A\bmod D_N$ for which at least one $n\in[M,N]$ satisfies
$\log g_n(A)>(1+\eta)n$ is at most $Ce^{-cM}$.
The same proportion among $1\le A\le H$ is at most
\begin{equation}\label{p11:tg:box}
                         Ce^{-cM}+\frac{D_N}{H}.
\end{equation}
\end{theorem}
\begin{proof}
Since $D_n\mid D_N$ for $n\le N$, the marginal distribution
at every such $n$ is the one in Theorem~\ref{p11:tg:law}.
Fix $0<s<1$.  Corollary~\ref{p11:tg:mgf} gives
\[
 \mathbb P\{\log g_n(A)>(1+\eta)n\}
                   \le e^{-s\eta n/2}
\]
for sufficiently large $n$.  Sum this geometric tail over
$M\le n\le N$; independence between different indices is not
needed.  Finally partition $[1,H]$ into complete blocks of length
$D_N$ and a remaining block of length less than $D_N$.
\end{proof}

This finite statement does not require a probabilistic interpretation
of the constants themselves.  It is an exact counting assertion
about ordinary positive integer parameters and all approximation
orders in a specified finite range.

We can also quantify the parameter rarity of bounded primitive
linear forms, uniformly over positive integer $A$.  Write
\[
 R_n(A)=Q_n(A)\delta_A-P_n(A),\qquad
 \ell_n(A)=R_n(A)/g_n(A).
\]
The positive remainder identity gives
\[
 R_n(A)=n!A^n\int_0^\infty
       \frac{e^{-t}}{1+At}\left(\frac{At}{1+At}\right)^n dt.
\]
Restricting to $\sqrt n\le t\le\sqrt n+1$, using $A\ge1$,
and using $\log(1+y)\le y$, we obtain the uniform lower bound
\begin{equation}\label{p11:tg:uniformR}
             R_n(A)\ge\frac{n!e^{-2\sqrt n-1}}{\sqrt n+2}.
\end{equation}
Thus for fixed $B>0$ and $0<s<1$, the fraction of
$A\in\{1,\ldots,D_N\}$ for which
$\ell_n(A)\le B$ for at least one $M\le n\le N$ is at most
\begin{equation}\label{p11:tg:bounded}
                      \exp(-sM\log M+O_{s,B}(M))
\end{equation}
as $M\to\infty$, uniformly in $N\ge M$.
Indeed, boundedness forces
$\log g_n(A)\ge n\log n-n-O_B(\sqrt n+\log n)$.
Apply~\eqref{p11:tg:mgfbound} and Markov's inequality at each $n$,
then sum the rapidly decreasing bounds.  The corresponding
proportion in $1\le A\le H$ has the additional error $D_N/H$.
Here periodic-block counting is applied to the necessary
gcd-threshold superset; the real-valued forms $\ell_n(A)$ themselves
are not asserted to be periodic in $A$.

\subsection{The almost-sure statement and its fixed-parameter limitation}

Equip the profinite integers $\widehat{\mathbb Z}$ with Haar
probability measure, and define $g_n(A)=\gcd(Q_n(A),D_n)$ through
the residue of the polynomial $Q_n(A)$ modulo $D_n$.
For every $\eta>0$, the exceptional probabilities
$\mathbb P\{\log g_n(A)>(1+\eta)n\}$ are summable by
Corollary~\ref{p11:tg:mgf}.  The first Borel--Cantelli lemma therefore
gives, after intersecting the full-measure sets for $\eta=1/j$,
$j=1,2,\ldots$,
\begin{equation}\label{p11:tg:haar}
       \limsup_{n\to\infty}\frac{\log g_n(A)}n\le1
            \quad\text{for Haar-almost every }A\in\widehat{\mathbb Z}.
\end{equation}
This is an all-large-orders conclusion for almost every profinite
parameter, stronger than an index-density statement for a fixed
parameter.

It cannot be specialized to $A=1$.  Every singleton, and hence the
whole embedded set of ordinary integers, has Haar measure zero in
$\widehat{\mathbb Z}$.  A full-measure theorem is logically
compatible with all prescribed integer parameters belonging to
its exceptional set.  Accordingly~\eqref{p11:tg:haar} does not close
the irrationality of $\delta$, any specified $\delta_A$, or
Euler's constant $\gamma$.  The new arithmetic result is the exact
parameter law and its quantitative finite simultaneous consequences;
the fixed-fibre arithmetic estimate remains an identifiable missing
step.

\subsection{Why curvature and a transcendental slope cannot finish the Euler attack}

The prior Euler--Taylor analysis gave an exact crossover curve with
positive curvature asymptotic to $a/x^2$ and a transcendental leading
slope.  It is tempting to infer that exponentially accurate integer
crossovers must eventually stop.  The following countermodel shows
that those two pieces of information alone cannot imply this.

\begin{proposition}[Exponentially close lattice points despite curvature]
\label{p11:tg:curvaturelimit}
Fix $a>0$, $b\in\mathbb R$, an arbitrary nonempty open interval
$I\subset\mathbb R$, and any function $\Psi:\mathbb N\to(0,\infty)$.
There are uncountably many transcendental $\alpha\in I$ for which
\begin{equation}\label{p11:tg:countermodel}
       |\alpha n-a\log n-b-2m|<e^{-\Psi(n)}
\end{equation}
holds for infinitely many pairs $(n,m)\in\mathbb N\times\mathbb Z$.
Every corresponding curve $f_\alpha(x)=\alpha x-a\log x-b$
has the exact strictly positive curvature $f_\alpha''(x)=a/x^2$.
\end{proposition}
\begin{proof}
For each fixed integer $n$, the centers
\[
                  \frac{2m+a\log n+b}{n}\qquad(m\in\mathbb Z)
\]
form a grid of spacing $2/n$.  Inside any nondegenerate compact
interval, a sufficiently large $n$ therefore gives at least two
well-separated centers in its interior.  Around each chosen center
take a closed interval of radius
\[
                   r_n=\frac{e^{-\Psi(n)}}{4n}.
\]
Choosing $n$ large enough makes these intervals disjoint, contained
in the parent interval, and of arbitrarily small diameter.
Every $\alpha$ in either child interval satisfies
\eqref{p11:tg:countermodel} at that $n$.

Start with a compact interval inside $I$.  At each level carry out
this construction for every current interval, choosing all new
indices larger than every previous index and all diameters below
$2^{-j}$ at level $j$.  The resulting binary Cantor construction
has uncountably many distinct limit points.  Every limit point
lies in infinitely many selected intervals and hence satisfies
\eqref{p11:tg:countermodel} infinitely often.  Removing the countable
set of algebraic numbers leaves uncountably many transcendental
limit points.  The derivative formula is immediate.
\end{proof}

This proposition does not replace the actual slope $\kappa$ by a
counterexample or assert that Euler's exact crossover has infinitely
many exceptional indices.  It establishes a precise limitation of
an attempted inference: positive $x^{-2}$ curvature, transcendence
of the slope, and even the exact leading linear--logarithmic shape
are compatible with arbitrarily accurate, infinitely sparse lattice
contacts.  A proof excluding all Euler crossovers must use additional
arithmetic information about the particular slope and the exact
integral defining the crossover, or a different approximation family.


\section{Exactly which prime-denominator gamma monomials enter the elliptic CM sector}
\label{p11:gc:section}

The preceding manuscript solved an unrestricted beta-word optimization
problem.  We now ask a different question: can one isolate the five
individual fifth and seventh gamma values by adjoining arbitrarily many
elliptic CM periods and using their Chowla--Selberg formulas?  We determine
the entire subspace accessible to this multiplicative strategy.  The
result permits auxiliary gamma values of arbitrary rational denominators,
all standard gamma identities, and rational powers.  The obstruction is a
Dirichlet-character component, rather than beta-word length.

\subsection{The precise identity calculus}

Let $\mathcal D$ be the rational vector space generated by $[x]$ for
$x\in\mathbb Q/\mathbb Z$, modulo
\begin{equation}\label{p11:gc:relations}
 [0]=[1/2]=0,\qquad [-x]=-[x],\qquad
 \sum_{k=0}^{m-1}[x+k/m]=[mx]\quad(m\ge1).
\end{equation}
The interpretation is the class of $\Gamma(x)$ modulo nonzero algebraic
factors and rational powers of $\pi$.  Recurrence, reflection and Gauss
multiplication give these relations.  Taking the quotient as a rational
vector space allows rational powers of the standard identities.
All asserted identities in this calculus have finite proofs.

For a negative fundamental discriminant $D$, put $d=|D|$ and let
$\chi_D$ be its primitive quadratic Dirichlet character.  Define
\begin{equation}\label{p11:gc:CMsymbol}
 c_D=\sum_{a=1}^{d-1}\chi_D(a)[a/d],\qquad
 \mathcal E=\operatorname{span}_{\mathbb Q}\{c_D:D<0
                  \text{ fundamental}\}\subset\mathcal D.
\end{equation}
Only finite linear combinations occur in this span.
This is precisely the multiplicative sector generated by the standard
Chowla--Selberg expressions for elliptic CM periods, after ignoring
algebraic factors and rational powers of $\pi$.  Indeed, if $h_D$ and
$w_D$ are the class number and number of roots of unity, a nonzero period
of a CM elliptic curve in the corresponding isogeny class is, up to a
nonzero algebraic factor,
\[
 \sqrt\pi\left(\prod_{a=1}^{d-1}
                 \Gamma(a/d)^{\chi_D(a)}\right)^{w_D/(4h_D)}.
\]
We use this classical formula only to identify the sector, not as a new
transcendence theorem; see Gross~\cite{p11:gc:gross}.
Periods of nonmaximal orders introduce no additional sector because
CM elliptic curves with the same imaginary quadratic field are isogenous
over $\overline{\mathbb Q}$.

For an odd prime $p$, put $h=(p-1)/2$ and
\[
 W_p=\operatorname{span}_{\mathbb Q}
                  \{[1/p],\ldots,[h/p]\}.
\]
The theorem below also proves these $h$ symbols are linearly independent
in the full standard calculus.

\begin{theorem}[Complete elliptic-CM recognition at prime denominators]
\label{p11:gc:recognition}
Let $\mathcal P$ be any finite set of odd primes.  An element
\[
 v=\sum_{p\in\mathcal P}\sum_{a=1}^{(p-1)/2} e_{p,a}[a/p],
 \qquad e_{p,a}\in\mathbb Q,
\]
belongs to $\mathcal E$ if and only if the following conditions hold
separately for every $p\in\mathcal P$:
\begin{enumerate}
 \item if $p\equiv1\pmod4$, then $e_{p,a}=0$ for every $a$;
 \item if $p\equiv3\pmod4$, there is $r_p\in\mathbb Q$ such that
       $e_{p,a}=r_p\left(\frac ap\right)$ for every $a$.
\end{enumerate}
Consequently
\[
 \dim_{\mathbb Q}\left(\left(\sum_{p\in\mathcal P}W_p\right)
                                      \cap\mathcal E\right)
 =\#\{p\in\mathcal P:p\equiv3\pmod4\}.
\]
All auxiliary rational denominators and all relations
\eqref{p11:gc:relations} are allowed in this assertion.
\end{theorem}

\subsection{Finite-proof character certificates}

Define the centered sawtooth
\[
 s(x)=\begin{cases}1-2\{x\},&x\notin\mathbb Z,\\
                  0,&x\in\mathbb Z.
       \end{cases}
\]
Its endpoint convention is essential.  Direct summation gives
\begin{equation}\label{p11:gc:sawtooth}
 s(-x)=-s(x),\quad s(1/2)=0,\quad
 \sum_{k=0}^{m-1}s(x+k/m)=s(mx).
\end{equation}
For completeness, if $u=\{mx\}\ne0$, the fractional parts on the
left are $(u+j)/m$, $0\le j<m$, in some order, and their sawtooth
sum is $1-2u$.  If $u=0$, their fractional parts are
$0,1/m,\ldots,(m-1)/m$ and the sum is zero.

Given any finite proof in $\mathcal D$, let $N$ contain every
denominator appearing in it.  For $t\in(\mathbb Z/N\mathbb Z)^\times$,
the substitution
\[
                         [x]\longmapsto s(tx)
\]
respects every relation in that proof.  For multiplication by $m$,
one has $m\mid N$, and $t$ permutes the shifts $k/m$ modulo one;
thus \eqref{p11:gc:sawtooth} applies.  This is a finite-proof certificate,
not an assertion that one fixed $t$ acts on every denominator globally.

For the CM symbol it gives
\begin{equation}\label{p11:gc:CMprofile}
 \sum_{a=1}^{d-1}\chi_D(a)s(ta/d)
       =K_D\chi_D(t),\qquad
 K_D=\sum_{a=1}^{d-1}\chi_D(a)s(a/d)=2L(0,\chi_D)\ne0.
\end{equation}
The first equality follows by the substitution $a\mapsto ta$ modulo
$d$.  The formula for $K_D$ follows from
$L(0,\chi)=-d^{-1}\sum_{a=1}^{d-1}a\chi(a)$.
For an odd primitive character, $L(0,\chi)$ is nonzero by the
Dirichlet functional equation and $L(1,\overline\chi)\ne0$.
These are classical Dirichlet $L$-function facts; see
\cite{p11:gc:dlmf}.
Thus every elliptic CM generator has a single quadratic-character
profile, independently of how large a common denominator is chosen.

\begin{lemma}[No invisible prime-denominator direction]
\label{p11:gc:injectivity}
For an odd prime $p$, the map
\[
 (e_1,\ldots,e_h)\longmapsto
 \Phi_e(t)=\sum_{a=1}^h e_a s(ta/p),
 \quad t\in(\mathbb Z/p\mathbb Z)^\times,
\]
is an isomorphism from $\mathbb C^h$ onto the space of odd functions
on $(\mathbb Z/p\mathbb Z)^\times$.
It preserves the character decomposition: the inverse image of the
line generated by a quadratic odd character is the same character
on the half-residue representatives.
\end{lemma}
\begin{proof}
Extend the coefficients oddly to the nonzero residues:
$E(a)=e_a$ for $1\le a\le h$, and $E(-a)=-e_a$.
Then
\[
 \Phi_e(t)=\frac12\sum_{a\in(\mathbb Z/p\mathbb Z)^\times}
                                    E(a)s(ta/p).
\]
For an odd Dirichlet character $\chi$ modulo $p$, its Fourier
coefficient equals
\[
 \sum_t\Phi_e(t)\overline{\chi(t)}
 =\frac12\left(\sum_a E(a)\chi(a)\right)
          \left(\sum_t s(t/p)\overline{\chi(t)}\right).
\]
The second factor is $2L(0,\overline\chi)\ne0$: every odd
character modulo the prime $p$ is primitive.
Fourier inversion therefore gives the claimed isomorphism on the odd
subspace.  For a quadratic character inversion does not change its
character line, proving the final assertion.
\end{proof}

\begin{proof}[Proof of Theorem~\ref{p11:gc:recognition}]
Suppose $v=\sum_D r_D c_D$ in $\mathcal D$, and choose $N$
containing all prime denominators, all $|D|$ on the right, and all
auxiliary denominators in a finite proof of this identity.  Its
profiles satisfy
\begin{equation}\label{p11:gc:profile-equality}
 \sum_{p\in\mathcal P}\Phi_{e_p}(t)
                   =\sum_D r_D K_D\chi_D(t)
 \qquad(t\in(\mathbb Z/N\mathbb Z)^\times).
\end{equation}
All characters are inflated to this common group.
The odd characters on the left have primitive conductor exactly $p$;
characters belonging to distinct primes have distinct conductors.
The characters on the right are primitive odd quadratic characters
of conductor $|D|$.  Distinct primitive characters remain distinct
when inflated to a common multiple of their conductors, and characters
of a finite abelian group are linearly independent.

There is no odd quadratic character of conductor $p$ when
$p\equiv1\pmod4$.  When $p\equiv3\pmod4$, the unique one is
$\chi_{-p}(a)=(a/p)$.  Equation~\eqref{p11:gc:profile-equality} and
Lemma~\ref{p11:gc:injectivity} therefore force precisely the coefficient
conditions in the theorem.  This argument also shows that a linear
combination of the prime-denominator symbols can vanish in
$\mathcal D$ only when every coefficient vanishes.

Conversely, for $p\equiv3\pmod4$, reflection and oddness give
\[
 c_{-p}=\sum_{a=1}^{p-1}\left(\frac ap\right)[a/p]
       =2\sum_{a=1}^{h}\left(\frac ap\right)[a/p].
\]
Each permitted direction thus belongs to $\mathcal E$ using only
reflection.  Summing these identities proves sufficiency and the
dimension formula.
\end{proof}

\subsection{The five targets and the exact remaining arithmetic gap}

For the particular joint attack, write
\[
 X_5=\Gamma(1/5),\quad Y_5=\Gamma(2/5),\qquad
 X_7=\Gamma(1/7),\quad Y_7=\Gamma(2/7),\quad Z_7=\Gamma(3/7).
\]
The theorem gives the complete criterion
\begin{equation}\label{p11:gc:five-seven}
 \begin{gathered}
 X_5^aY_5^bX_7^AY_7^BZ_7^C
 \text{ is standard-reducible to elliptic CM products}\\
 \Longleftrightarrow\quad a=b=0,\quad A=B=-C.
 \end{gathered}
\end{equation}
for rational exponents, with algebraic factors and rational powers of
$\pi$ ignored.  Thus mixing the fifth and seventh denominators does
not cancel their nonquadratic character components.  Nor can adding
CM elliptic curves of arbitrarily large discriminant cancel them.

The surviving seventh direction is
\[
                         B_7=X_7Y_7/Z_7=B(1/7,2/7).
\]
Its actual arithmetic is classical: reflection gives
\[
 \operatorname{CS}(-7)=
 \frac{\sin(\pi/7)\sin(2\pi/7)}{\sin(3\pi/7)}
                         \frac{B_7^2}{\pi},
\]
and the Chowla--Selberg--Nesterenko input implies algebraic
independence of $\pi,B_7,e^{\pi\sqrt7}$; see
Cohen--Zudilin~\cite{p11:gc:cohen-zudilin}.
This is the same classical CM direction used in the preceding
manuscript.  The new classification proves that it is the only one
available in the entire joint prime-denominator span by the stated
standard-calculus strategy.

More generally, for $p\equiv3\pmod4$ define
\[
 P_p=\prod_{a=1}^{(p-1)/2}\Gamma(a/p)^{\left(\frac ap\right)}.
\]
Reflection identifies a rational power of $P_p$, up to an algebraic
factor and a rational power of $\pi$, with
$\operatorname{CS}(-p)$.  Therefore $\pi$ and $P_p$ are
algebraically independent by the same classical CM theorem.
This last transcendence assertion is a consequence of known results,
not a newly settled conjecture.

The restriction to the standard calculus cannot be removed for arbitrary
gamma monomials with the results used here.  Rohrlich's conjecture would
make it legitimate to pass from a numerical multiplicative algebraic
relation to such a formal relation, and would then promote
\eqref{p11:gc:five-seven} to the corresponding assertion about actual
numerical CM-product identities.  We do not assume that conjecture.
W\"ustholz's theorem controls algebraic \emph{linear} relations between
periods; applying it to unrestricted products of beta periods is not
justified.  The tensor products that arise in those gamma-isolating
words are exactly where the attempted value transfer still stops.

For orientation, a known result of Vasil'ev and Grinspan already says
that at least two of $\pi,\Gamma(1/5),\Gamma(2/5)$ are algebraically
independent; see the discussion and original references in
Waldschmidt~\cite{p11:gc:waldschmidt}.  Our formal classification neither
improves that transcendence-degree bound nor selects which individual
gamma value is transcendental.  The five individual irrationality or
transcendence targets remain unclosed in this investigation.


\section{Sixteen named targets and the remaining arithmetic steps}
\label{p11:tr:ledger}

The entries below state outcomes of this manuscript. An unresolved
entry means that our arguments do not prove its target conclusion.
It does not erase external results or establish a complete survey of
the literature. In particular, the reported external irrationality
results for $G$ and $\zeta(5)$ are distinct from our transcendence
targets. The sixteen targets are counted as values, not as multiple
renamings of the same conditional criterion.

\begingroup\small
\renewcommand{\arraystretch}{1.12}
\begin{longtable}{@{}L{0.26\textwidth}L{0.69\textwidth}@{}}
\toprule
Named target & Additional conclusion and remaining step\\
\midrule
\endfirsthead
\toprule
Named target & Additional conclusion and remaining step\\
\midrule
\endhead
Euler's constant $\gamma$: irrationality &
An explicit transcendental-slope construction admits infinitely many
exponentially close lattice points despite curvature. Thus that
abstract curvature route cannot settle the fixed Euler target.
No irrationality conclusion follows.\\
Euler--Gompertz $\delta=\delta_1$: irrationality &
The exact parameter law and exponential tails control primitive
cancellation for most residues. They do not exclude exceptional
behavior at the prescribed parameter $A=1$.\\
Catalan $G$: transcendence &
The rational pencil $A_m+16XI_m$ has full degree and all zeros below
$G$; the $\pi^2$ coefficient vanishes exactly. Adequate primitive-value
and off-axis growth estimates are not proved.\\
$\zeta(3)$: transcendence &
The single off-axis degree criterion and sharp positive-Hankel
counterexample identify the required additional arithmetic margin.
No higher algebraic degree is newly excluded.\\
$\zeta(5)$: transcendence &
The inspected external determinant estimates have rate ratio
$1975/511067$. A sufficient quadratic-exclusion content rate is
$\chi>127273/20000$, but no such bound is established.\\
$\zeta(7)$: irrationality &
The positive-pencil criterion applies to the proposed certificate
framework; a sufficient primitive-value decay estimate is missing.\\
$\zeta(9)$: irrationality &
The same degree criterion requires arithmetic denominator control
that positivity and amplification do not supply.\\
$\zeta(11)$: irrationality &
No individual value conclusion follows. Independence or a disjunction
involving several odd zeta values does not prove this target.\\
Erd\H{o}s--Borwein $\sum_{n\ge1}(2^n-1)^{-1}$: transcendence &
The Lambert function has no algebraic differential equation of order
at most two and belongs to an independent four-function family.
Specialization at $q=1/2$ is not resolved.\\
Reciprocal Fibonacci sum $\sum_{n\ge1}F_n^{-1}$: transcendence &
The mixed Lambert--Fibonacci jet theorem strengthens functional
independence. The corresponding arithmetic specialization is
unproved.\\
Prime Mertens $B_1$: irrationality &
The prime-harmonic function adds a fourth independent function to
three jets. Its constant term is not thereby shown irrational.\\
$\Gamma(1/5)$: transcendence &
The fifth-denominator space has zero intersection with the entire
elliptic CM sector in the standard identity calculus. Numerical
transcendence requires information beyond that calculus.\\
$\Gamma(2/5)$: transcendence &
The same complete fifth-denominator obstruction holds with all
auxiliary denominators and all elliptic CM periods. No new individual
value theorem.\\
$\Gamma(1/7)$: transcendence &
The seventh-denominator CM intersection is exactly the Legendre-symbol
line. This coordinate is outside it; the formal obstruction does not
prove the value transcendental.\\
$\Gamma(2/7)$: transcendence &
The same exact one-dimensional CM recognition applies. Unknown
numerical multiplicative relations are not excluded.\\
$\Gamma(3/7)$: transcendence &
The same complete formal classification applies, without an
individual transcendence conclusion.\\
\bottomrule
\end{longtable}
\endgroup

The associated Kurepa left-factorial question is likewise not resolved.
The parameter-average law does not decide the special residue patterns
needed for that conjecture.

\subsection*{What is proved, and what is not counted as a closure}
The exact totally-real alphabet threshold, the totally-positive
algebraic-power classifier, the three-jet/four-function independence
theorem, the parameter distribution law, the gamma CM recognition
theorem, and the full-degree auxiliary-free Catalan pencil have
complete proofs in their stated settings. The determinant criterion
and its sharp Pell example clarify an algebraic-degree obstruction.
These are structural results. The number of newly resolved famous
constant conjectures is zero.

\subsection*{Verification and next obligations}
Independent proof reviews accompany the sections. Exact checks cover
paired Catalan interpolation and primitive determinants, Pell/Hankel
identities, Laguerre residue distributions, gamma character profiles,
and the Lambert cusp-model obstruction. Finite computations verify
formulas and examples; asymptotic conclusions rely on the proofs.
The principal remaining obligations are a suitable arithmetic
specialization theorem, a quantitative primitive determinant estimate,
or a fixed-parameter argument that goes beyond an average statement.
Neither functional independence nor rapid approximation is used as a
substitute for any of those obligations. Historical priority and
publication significance require further expert comparison.

% END SOURCE: KEIO_Spectral_Entropy_and_Arithmetic_Transfer.tex

\endgroup
\clearpage
\part{Quadratic Exceptions and Constant Transfer}
\label{merge:part12}
\begingroup

% BEGIN SOURCE: KEIO_Quadratic_Exceptions_and_Constant_Transfer.tex

\noindent\textbf{Editorial integration note.} 
This is the final source manuscript in the edition. Its exact entropy classifications, fixed-parameter cancellation bounds, Catalan repair, and actual beta-period calculation provide the latest conclusions in their respective branches. Its sixteen-target ledger states the remaining arithmetic gaps explicitly.


\begin{quote}\small
\noindent\textit{Original part framing.}\par
\noindent\textbf{Reading convention.}
This is a standalone mathematical successor to the earlier research
manuscripts, not a merger of their full contents. Definitions and
inputs needed for the present proofs are included. Totally positive
means that \emph{every} conjugate is positive; positivity of only the
displayed real embedding is weaker. Algebraic powers of positive
numbers always use the real logarithm. Classical external theorems,
proofs supplied here, finite checks, and unresolved transfers have
different roles. Independent proof reviews and exact computations do
not amount to formal verification.
\end{quote}
\section{Arithmetic classifications and the limits of transfer}
\label{p12:cf:introduction}

The main distinction is between a number whose conjugates satisfy
positivity constraints and a named constant approached by algebraic
numbers. The former setting allows logarithmic relations to be
transported and compared; in the latter, roots of approximating
polynomials are not known conjugates of the target. This distinction
explains both the rigidity results and the remaining barriers.

\paragraph{Relative entropy has a genuine quadratic exception.}
For $D(p\Vert r)=\sum_i p_i\log(p_i/r_i)$, with $r$ rational and
$p$ totally positive algebraic, Theorem~\ref{p12:icf:relative} proves that
an algebraic value of $\exp(cD(p\Vert r))$, for $c\ne0$ real
algebraic and $p\ne r$, forces
\[
             [\mathbb Q(p_1,\ldots,p_n):\mathbb Q]\le2.
\]
It gives an explicit multiplicative test for all quadratic exceptions.
The proof compares two independent transports of logarithmic
relations. Three distinct positive conjugate laws would require three
nonzero coefficients to have pairwise opposite signs. The resulting
degree barrier is sharp: binary laws have no nonrational exception,
whereas a concrete three-symbol law has one even for the uniform
reference. Thus the exact threshold in this specified problem is
three, not an unrestricted statement about all entropy problems.

\paragraph{Conditional laws and Markov sources have exact classifiers.}
For a totally positive algebraic joint law, an algebraic value of
$\exp(cH(X\mid Y))$ forces every conditional probability to be
rational. The remaining issue is an exact rational-line test on
prime-valuation coefficients (Theorem~\ref{p12:icf:conditional}).
Irrational marginals and irrational successful exponents can occur.
For a strictly positive stochastic matrix with totally positive
algebraic entries, the stationary constraints remove that freedom:
Theorem~\ref{p12:icf:markov} gives algebraicity precisely for rational
transition probabilities and rational exponent. These are arithmetic
classifications of explicitly specified entropy values.

\paragraph{The cubic ternary case closes inside a larger open problem.}
Theorem~\ref{p12:ef:cubic} excludes rational bit entropy for every positive
ternary algebraic probability vector whose probability field has
degree three, including cubic fields with a pair of complex
embeddings. A new restriction in the unit normal form says that the
greatest common divisor of exponent differences is a power of two.
For odd entropy denominator, the probability field equals the unit
generator field. The cubic case then reduces to two elementary unit
bounds. This does not exclude every ternary example over an arbitrary
number field: a remaining example would need degree at least four,
or odd degree at least five when the entropy denominator is odd.

\paragraph{Classical period theory upgrades a formal obstruction.}
Theorem~\ref{p12:thm:cf-fourteen} computes an actual
$\overline{\mathbb Q}$-linear basis of size fourteen for $1$ together
with the fifth and seventh beta-period spans. Its input is the
classical completeness theorem for beta values, not an assumption of
Rohrlich's conjecture for all Gamma monomials. Twelve displayed
transcendental ratios and an infinite family of exact
linear-fractional tests follow. These are explicit consequences of
classical transcendence theory; they are not twelve previously open
constant problems solved. No individual fifth or seventh Gamma
value is classified here.

\paragraph{Fixing the Gompertz parameter changes what can be averaged.}
For each fixed positive integer $A$, Theorem~\ref{p12:fg:roots} proves
$\rho_p(A)=O(p^{2/3})$ for the number of index residues at which a
prime divides the classical Laguerre denominator. Theorem~\ref{p12:fg:law}
then gives a limiting distribution and finite exponential moments
for a prime-support upper bound on normalized logarithmic content.
This applies in particular to $A=1$. A separate positive-continuant
packing argument gives the stronger content estimate
\[
 \frac1X\sum_{n\le X}\frac{\log g_n(A)}n
                   =O_A\bigl(X^{-1/3}\log(2X)\bigr),
\]
and only $O_{A,\varepsilon}((\log(2X))^4)$ indices up to $X$ have
$\log g_n(A)\ge\varepsilon n$ (Theorem~\ref{p12:fg:sublinear}).
On a single set of density one the reduced denominator and primitive
form both retain their first two factorial growth terms. These results
do not rule out every exceptional index.
Theorem~\ref{p12:fg:complexity} shows that a bounded primitive form would
require cancellation involving a growing number of primes.

\paragraph{Catalan convergence can be repaired without an arithmetic conclusion.}
A fixed quadratic polynomial invariant removes the auxiliary period
and produces positive Hankel pencils of fixed determinant degree.
Theorem~\ref{p12:zc:catalan} proves that their roots stay uniformly below
Catalan's constant. An explicit moment-null density provides a
quantitative gap. The full paired polynomial space contains an
additional sector omitted by that invariant construction.
Theorem~\ref{p12:zc:repair} proves that restoring the full space makes all
roots converge to $G$ from below at every fixed determinant degree.
The unproved step is a sufficiently strong comparison with primitive
integer coefficient height; analytic convergence alone supplies no
irrationality or transcendence certificate.

\paragraph{A Lambert optimization is a limit of specified estimates.}
Theorem~\ref{p12:lf:optimization} identifies the unique optimum of the
published common-shift denominator budget and proves
$\Lambda(x)>2/5$ uniformly. It therefore cannot reach the $1/3$
threshold for degree-three exclusion. Independent shifts in the
proved central cone do not improve this optimum. The result concerns
those sufficient clearing bounds, not a lower bound for the actual
primitive denominators and not a universal obstruction to other
Lambert constructions. The external October 2026 nonquadraticity
announcement is attributed separately and is not counted as a result
of this manuscript.

The sections below contain proofs and their assumptions. The final
ledger records what they imply, and fail to imply, for each named
constant. Historical priority remains a separate question from the
validity of the supplied arguments.


% BEGIN EMBEDDED information_results.tex

\section{Relative entropy: a degree-two barrier and its exact exceptions}
\label{p12:icf:section}

All logarithms in this section are real natural logarithms. For strictly
positive probability vectors write
\[
 D(p\Vert r)=\sum_i p_i\log\frac{p_i}{r_i},
 \qquad H(p)=-\sum_i p_i\log p_i.
\]
An algebraic number is \emph{totally positive} if every one of its
conjugates over $\mathbb Q$ is positive. This hypothesis is imposed on
the individual probabilities, not merely on their displayed values.
The conclusions below concern arithmetic values, rather than the
semialgebraicity of entropy level sets studied in
\cite{p10:inf:BW,p12:icf:BM}. They use classical logarithmic transcendence
\cite{p07:rlog:BakerIII}. Their historical priority has not been established
by an exhaustive literature review.

\subsection{Two independent transports of logarithmic relations}

\begin{lemma}[Baker transport and algebraic logarithms]
\label{p12:icf:transport}
Let $\beta_1,\ldots,\beta_N$ be positive algebraic numbers.
Every algebraic linear relation among the $\log\beta_i$ is an
algebraic linear combination of their rational relations.
Consequently the following operations are valid independently:
\begin{enumerate}
\item conjugate the algebraic coefficients of a relation, keeping
all logarithmic bases fixed;
\item conjugate its bases by an embedding under which all of them
remain positive, keeping all coefficients fixed.
\end{enumerate}
If $u_1,\ldots,u_s$ are rationally independent logarithms of positive
algebraic numbers and $b_j$ are real algebraic, then
\[
 \exp\Bigl(\sum_j b_j u_j\Bigr)\in\overline{\mathbb Q}
 \quad\Longleftrightarrow\quad b_j\in\mathbb Q\text{ for every }j.
\]
\end{lemma}
\begin{proof}
Choose a rational basis of the span of the logarithms. Baker's
linear independence theorem says that this basis remains independent
over $\overline{\mathbb Q}$. Comparing its coefficients proves the
rational-generation assertion and coefficient transport. A rational
relation, after clearing denominators, is a multiplicative identity
among the bases. Conjugation preserves this identity, and positivity
permits taking real logarithms; this proves base transport.
For the final assertion, if the exponential is algebraic $A>0$,
Baker's theorem forces $\log A$ into the rational span of the $u_j$.
Comparing coefficients gives $b_j\in\mathbb Q$. The converse is a
product of rational powers of algebraic numbers.
\end{proof}

\begin{theorem}[Exact rational-reference classification]
\label{p12:icf:relative}
Let $r$ be a strictly positive rational probability vector, let $p$
be a strictly positive probability vector with totally positive
algebraic entries, and put $K=\mathbb Q(p_1,\ldots,p_n)$.
For nonzero real algebraic $c$, the algebraicity of
$\exp(cD(p\Vert r))$ is classified as follows.
\begin{enumerate}
\item If $p=r$, the value is $1$ for every $c$.
\item If $p\ne r$ and $K=\mathbb Q$, it is algebraic exactly when
$c\in\mathbb Q$.
\item If $[K:\mathbb Q]\ge3$, it is transcendental for every such $c$.
\item Suppose $K=\mathbb Q(t)$, where $t=\sqrt d$ and $d>1$ is a
squarefree integer. Write $p_i=m_i+t v_i$ with $m_i,v_i\in\mathbb Q$.
The value is algebraic exactly when $c=a+bt\in K$, with
$a,b\in\mathbb Q$, and
\begin{equation}\label{p12:icf:quadratic-test}
       \prod_i\left(\frac{p_i}{r_i}\right)^{a v_i+b m_i}=1.
\end{equation}
All powers in this test are positive real powers with rational
exponents. Equivalently, clear a common denominator to obtain an
equality of algebraic numbers with integer exponents.
\end{enumerate}
In case (4), every successful nonzero exponent satisfies
$(a+bt)(a-bt)<0$.
\end{theorem}
\begin{proof}
The zero case follows from strict positivity of relative entropy.
Assume henceforth $p\ne r$, so $D(p\Vert r)>0$, and suppose
$A=\exp(cD(p\Vert r))$ is algebraic. Work in a finite normal field
$L$ containing $p$ and $c$. Put $\beta_i=p_i/r_i$.
For $\sigma\in\operatorname{Gal}(L/\mathbb Q)$ set
$p^\sigma=(\sigma p_i)_i$ and $c_\sigma=\sigma c$.
Coefficient transport in
\[
              \log A-\sum_i cp_i\log\beta_i=0
\]
gives, without changing the base $A$,
\[
        \sum_i c_\sigma p_i^\sigma\log\beta_i=\log A.
\]
Subtract the equations for $\sigma$ and $\tau$. Base transport
then gives the same resulting relation with bases $p_i^\sigma/r_i$
and, separately, with bases $p_i^\tau/r_i$. Subtracting these two
relations yields the key identity
\begin{equation}\label{p12:icf:pairwise}
 c_\sigma D(p^\sigma\Vert p^\tau)
       +c_\tau D(p^\tau\Vert p^\sigma)=0.
\end{equation}
Total positivity makes all these divergences well defined.

If $p^\sigma=p$, comparison with the identity embedding in the
fixed-base equation gives $c_\sigma=c$, because $D(p\Vert r)>0$.
Thus the stabilizer of $p$ fixes $c$, and $c\in K$. If
$p^\sigma\ne p$, equation~\eqref{p12:icf:pairwise} with the identity
embedding gives
\[
             \frac{c_\sigma}{c}
                =-\frac{D(p\Vert p^\sigma)}
                         {D(p^\sigma\Vert p)}<0.
\]
In particular all conjugates of $c$ are real. There cannot be three
distinct conjugate vectors: their three pairwise identities would
require three nonzero real numbers to have pairwise opposite signs.
The number of distinct conjugate vectors is $[K:\mathbb Q]$.
Therefore $[K:\mathbb Q]\le2$. If $K=\mathbb Q$, the already
proved inclusion $c\in K$ gives the necessary rationality of $c$;
its sufficiency is immediate from rational powers of rational numbers.

It remains to make the quadratic case exact. For $c=a+bt$,
\[
 cD(p\Vert r)=
 \sum_i(am_i+bdv_i)\log\beta_i
       +t\sum_i(av_i+bm_i)\log\beta_i.
\]
The exponential of the first sum is algebraic. The exponential of
the second sum before multiplication by $t$ is also algebraic,
since all its coefficients are rational. If that sum is nonzero,
Lemma~\ref{p12:icf:transport}, applied to its single nonzero logarithm,
shows that its exponential after multiplication by the irrational
algebraic $t$ is transcendental. Thus algebraicity is equivalent
to the second sum vanishing, precisely~\eqref{p12:icf:quadratic-test}.
Finally the two distinct conjugate vectors in
\eqref{p12:icf:pairwise} force $c$ and its quadratic conjugate to have
opposite signs.
\end{proof}

\begin{corollary}[Three positive conjugates suffice]
\label{p12:icf:local}
Let $p$ be a positive algebraic probability vector and $r$ a
positive rational probability vector. If $p$ has three distinct
conjugate vectors whose entries are all positive, then
$\exp(cD(p\Vert r))$ is transcendental for every nonzero real
algebraic $c$, even when other conjugate vectors are not positive.
\end{corollary}
\begin{proof}
The coefficient and base transports in the proof of
Theorem~\ref{p12:icf:relative} apply to any chosen positive conjugate
vectors. If there were an algebraic exponential, the pairwise
identity~\eqref{p12:icf:pairwise} would force three corresponding
nonzero conjugate coefficients to have pairwise negative real
ratios, which is impossible.
\end{proof}

\begin{example}[A cubic family with no exceptional reference]
\label{p12:icf:cubic}
For $j=1,2,3$ put
\[
                  p_j=\frac{3+2\cos(2\pi j/7)}8.
\]
These are the positive roots of
$512X^3-512X^2+152X-13$, and their sum is one.
The polynomial is irreducible modulo $3$ (its cubic reduction
has no root), hence irreducible over $\mathbb Q$. Thus
$[\mathbb Q(p_1,p_2,p_3):\mathbb Q]\ge3$, and every $p_j$ is
totally positive. For every strictly positive rational reference
$r$ and every nonzero real algebraic $c$,
\[
                    \prod_{j=1}^3(p_j/r_j)^{c p_j}
                    \quad\hbox{is transcendental}.
\]
The rational reference may vary freely; no quadratic escape is
possible for this one fixed cubic law.
\end{example}

\begin{corollary}[The exceptional exponent space is one rational line]
\label{p12:icf:line}
In the nonrational quadratic case of Theorem~\ref{p12:icf:relative},
the set
\[
       \{c\in\overline{\mathbb Q}\cap\mathbb R:
                       \exp(cD(p\Vert r))\in\overline{\mathbb Q}\}
\]
is either $\{0\}$ or $c_0\mathbb Q$ for one nonzero
$c_0\in K$ of negative norm. In particular a positive-norm
quadratic exponent can never be an exception.
\end{corollary}
\begin{proof}
The test~\eqref{p12:icf:quadratic-test} is the kernel over $\mathbb Q$
of the real linear expression
\[
  (a,b)\longmapsto
  a\sum_i v_i\log\beta_i+b\sum_i m_i\log\beta_i.
\]
Its two coefficients are not both zero, since otherwise
$D(p\Vert r)=0$. The kernel has rational dimension at most one.
The norm assertion follows from Theorem~\ref{p12:icf:relative}.
\end{proof}

\begin{example}[A genuine quadratic escape at a uniform reference]
\label{p12:icf:escape}
Let
\[
 t=\frac{\sqrt{17}}{36},\qquad
 p=\left(\frac14+t,\frac14-t,\frac12\right),\qquad
 r=\left(\frac13,\frac13,\frac13\right).
\]
All entries of $p$ are totally positive, and its multiset is
Galois invariant. Moreover $p_1p_2=4/81$, so
\[
 (3p_1)(3p_2)(3/2)^2=1,
 \qquad D(p\Vert r)=t\log\frac{p_1}{p_2}.
\]
It follows that
\[
       \boxed{\exp\bigl(D(p\Vert r)/t\bigr)
                         =\frac{9+\sqrt{17}}{9-\sqrt{17}}
                         \in\overline{\mathbb Q}.}
\]
Thus even a rational-characteristic-polynomial spectrum and a
uniform reference do not imply transcendence of every algebraic
power of the exponential relative entropy. The degree-two
exception in Theorem~\ref{p12:icf:relative} is necessary.
\end{example}

\begin{corollary}[Binary rigidity and the sharp three-symbol threshold]
\label{p12:icf:binary}
Let $p=(p_1,1-p_1)$ have totally positive algebraic entries and let
$r$ be a strictly positive rational binary law. If $p\ne r$, then
\[
 \exp(cD(p\Vert r))\in\overline{\mathbb Q}
 \quad\Longleftrightarrow\quad p_1\in\mathbb Q
                                  \text{ and }c\in\mathbb Q
\]
for nonzero real algebraic $c$. More generally, a nonrational
quadratic law has no exceptional exponent whenever its rational
reference lies on the affine line joining the law to its conjugate.
Consequently the minimum alphabet size allowing a nonrational
totally positive law, a rational reference, and a nonzero algebraic
exponent with algebraic exponential relative entropy is exactly $3$.
\end{corollary}
\begin{proof}
The degree barrier reduces the issue to the quadratic case
$p=m+t v$, $q=m-tv$, $c=a+bt$. A successful exponent has negative
norm, hence $b\ne0$ and $|a/b|<t$. Therefore
$z=m+(a/b)v$ is a strictly positive rational probability vector
on the open segment joining $p$ and $q$. Dividing
\eqref{p12:icf:quadratic-test} by $b$ gives
\[
                 \sum_i z_i\log p_i=\sum_i z_i\log r_i.
\]
Clear denominators in this rational logarithmic relation and
conjugate its bases. Since $z,r$ are rational, this also gives
$\sum_i z_i\log q_i=\sum_i z_i\log r_i$.
On the positive part of the affine line through $p,q$, the
function $x\mapsto\sum_i z_i\log x_i$ is strictly concave:
its second derivative along the nonzero direction $v$ is
$-\sum_i z_iv_i^2/x_i^2<0$. Three distinct points on this line
cannot have the same value. If $r$ lies on the line, $p,q,r$
would be three such points, a contradiction. Every binary
probability vector lies on the same affine line, proving binary
rigidity. One symbol has no nonrational law, and
Example~\ref{p12:icf:escape} supplies the three-symbol exception.
\end{proof}

\begin{corollary}[The Galois-averaged reference has no exceptions]
\label{p12:icf:average}
Let $p$ be a nonrational probability vector whose entries are
totally positive algebraic numbers. In a finite normal field
containing its entries, define
\[
              m_i=\frac1{|G|}\sum_{\sigma\in G}\sigma(p_i),
              \qquad G=\operatorname{Gal}(L/\mathbb Q).
\]
Then $m$ is a strictly positive rational probability vector and
\[
                 \exp(cD(p\Vert m))\text{ is transcendental}
\]
for every nonzero real algebraic $c$. If the multiset of entries
of $p$ is Galois invariant, one also has
\[
               D(p\Vert m)=H(m)-H(p)>0.
\]
Thus every nonzero real algebraic power of
$\exp(H(m)-H(p))$ is transcendental in this spectral case.
\end{corollary}
\begin{proof}
The average is fixed by $G$, positive, and normalized. If
$[\mathbb Q(p):\mathbb Q]\ge3$, use
Theorem~\ref{p12:icf:relative}. If this degree is two, $m$ is the
midpoint of the two conjugate vectors and
Corollary~\ref{p12:icf:binary} applies. Nonrationality excludes degree
one and ensures $p\ne m$.
For a Galois-invariant multiset, $m_i$ is constant on each
conjugacy orbit of entries. Conjugating $p$ permutes its entries
within these orbits, so $\sum_i p_i\log m_i$ is unchanged.
Averaging gives
$\sum_i p_i\log m_i=\sum_i m_i\log m_i$. Substitution into
the definition of relative entropy proves the entropy identity.
\end{proof}

\subsection{Conditional entropy and stationary Markov sources}

For a strictly positive joint law $p=(p_{xy})$ put
$P_y=\sum_xp_{xy}$ and $r_{x|y}=p_{xy}/P_y$. Let
\[
 F(p)=H(X\mid Y)=-\sum_{x,y}p_{xy}\log r_{x|y}.
\]
The $X$-alphabet is assumed to have at least two symbols.

\begin{theorem}[Exact conditional effective-number classifier]
\label{p12:icf:conditional}
Suppose all $p_{xy}$ are totally positive algebraic numbers.
For nonzero real algebraic $c$, the number $\exp(cF(p))$ is
algebraic if and only if both of the following hold:
\begin{enumerate}
\item Every conditional probability $r_{x|y}$ is rational.
\item For every rational prime $\ell$,
\begin{equation}\label{p12:icf:prime-test}
 cB_\ell\in\mathbb Q,\qquad
 B_\ell=-\sum_yP_y\sum_xr_{x|y}v_\ell(r_{x|y}).
\end{equation}
Here $v_\ell(\ell)=1$ and only finitely many $B_\ell$ are nonzero.
\end{enumerate}
If a conditional probability is irrational, every nonzero real
algebraic power of $\exp(F(p))$ is transcendental.
If all conditional probabilities are rational, the successful
exponents, together with zero, form either $\{0\}$ or one rational
line $\beta^{-1}\mathbb Q$: the latter occurs exactly when the
nonzero $B_\ell$ span a one-dimensional rational vector space
$\beta\mathbb Q$.
\end{theorem}
\begin{proof}
Suppose $A=\exp(cF(p))$ is algebraic and choose a normal field
containing $p,c$. For an automorphism $\sigma$ put
$q=\sigma p$, $Q_y=\sum_xq_{xy}$,
$s_{x|y}=q_{xy}/Q_y$, and $d=\sigma c$.
Coefficient transport gives
\[
 \log A=-c\sum_{xy}p_{xy}\log r_{x|y}
        =-d\sum_{xy}q_{xy}\log r_{x|y}.
\]
Both sums without the prefactors are strictly negative, since
$0<r_{x|y}<1$. Hence $d/c>0$ is real. The coefficient vector
$cp-dq$ annihilates the logarithms of $r$, and base transport
shows that it also annihilates the logarithms of $s$.
Subtracting the two relations yields
\[
 c\sum_yP_yD(r_{\cdot|y}\Vert s_{\cdot|y})
  +d\sum_yQ_yD(s_{\cdot|y}\Vert r_{\cdot|y})=0.
\]
All weights are positive after division by $c$. Thus every
conditional law equals its $\sigma$-conjugate. Since this holds
for all $\sigma$, each $r_{x|y}$ is rational.

For rational conditionals, prime factorization gives
$F(p)=\sum_\ell B_\ell\log\ell$ with the coefficients in
\eqref{p12:icf:prime-test}. Distinct prime logarithms are rationally
independent. Lemma~\ref{p12:icf:transport} proves both necessity and
sufficiency of that test. Finally, since $F(p)>0$, at least one
$B_\ell$ is nonzero. A common nonzero multiplier makes all the
$B_\ell$ rational precisely when their rational span has
dimension one; its successful multipliers are then exactly
$\beta^{-1}\mathbb Q$.
\end{proof}

\begin{example}[Irrational marginals and an irrational successful exponent]
\label{p12:icf:conditional-example}
Take two conditional laws on four symbols,
\[
 r^{(1)}=(1/2,1/4,1/8,1/8),\qquad
 r^{(2)}=(1/4,1/4,1/4,1/4),
\]
and marginal probabilities $P_1=a=(2+\sqrt2)/4$, $P_2=1-a$.
The joint entries $P_yr_x^{(y)}$ are totally positive algebraic,
and
\[
       F(p)=(2-a/4)\log2,
       \qquad \exp\left(\frac{F(p)}{2-a/4}\right)=2.
\]
The marginal and the successful exponent are irrational.
Consequently conditional entropy cannot inherit the stronger
claim that all entries and the exponent must be rational.
\end{example}

\begin{theorem}[Exact classifier for Markov entropy rates]
\label{p12:icf:markov}
Let $T$ be an $n\times n$ strictly positive stochastic matrix,
$n\ge2$, whose entries are totally positive algebraic numbers.
Let $\pi$ be its stationary probability vector and let
\[
             h(T)=-\sum_{i,j}\pi_iT_{ij}\log T_{ij}
\]
be the entropy rate of its stationary Markov chain. For every
nonzero real algebraic $c$,
\[
 \boxed{\quad \exp(c h(T))\in\overline{\mathbb Q}
       \quad\Longleftrightarrow\quad
       T_{ij}\in\mathbb Q\text{ for all }i,j
                   \ \text{and}\ c\in\mathbb Q.\quad}
\]
If an entry of $T$ is irrational, $h(T)/\log b$ is transcendental
for every algebraic real base $b>1$.
\end{theorem}
\begin{proof}
Strict positivity gives a unique stationary vector with positive
entries. Solving the stationary linear equations shows that
$\pi$ belongs to the number field generated by the entries of $T$.
Under every embedding, the conjugate matrix remains strictly
positive and stochastic, and the conjugate vector is its unique
normalized stationary vector. Thus each $\pi_i$ and each joint
probability $\pi_iT_{ij}$ is totally positive. Apply
Theorem~\ref{p12:icf:conditional} to consecutive states of the chain:
algebraicity forces every transition probability $T_{ij}$ to be
rational. Then the stationary equations have rational coefficients
and a unique solution, so $\pi$ is rational. Hence $h(T)$ is a
nonzero rational linear combination of prime logarithms. Its
positive exponential is algebraic, and
Lemma~\ref{p12:icf:transport} forces $c$ to be rational.
The converse follows from rational powers of rational numbers.
If $h(T)/\log b=a$ were algebraic, it would be positive and
$\exp(h(T)/a)=b$ would contradict the theorem.
\end{proof}

\begin{example}[An asymmetric algebraic Markov source]
For $a=(2+\sqrt2)/4$, put
\[
 T=\begin{pmatrix}a&1-a\\1/3&2/3\end{pmatrix},\qquad
 \pi=\left(\frac{1/3}{4/3-a},
                         \frac{1-a}{4/3-a}\right).
\]
Both conjugate transition matrices are strictly positive.
Every nonzero real algebraic power of $\exp(h(T))$ is
transcendental, although the rate involves an irrational stationary
mixture of two different row entropies.
\end{example}

\subsection{Scope of the transfer}

The degree barrier for relative entropy is not an unrestricted
relative-entropy rigidity claim: rational $p$ already makes
$\exp(D(p\Vert q))$ algebraic for every positive algebraic
reference $q$. The reference is rational in
Theorem~\ref{p12:icf:relative}, and its quadratic exceptions are real.
The conditional theorem concerns classical, strictly positive
joint laws. It does not assert a corresponding classification
for quantum conditional entropy or general hidden-Markov entropy
rates. None of these value classifications establishes algebraic
independence of the resulting transcendental constants or resolves
a previously named open conjecture.


% END EMBEDDED information_results.tex


% BEGIN EMBEDDED entropy_results.tex

\section{A two-primary obstruction and the cubic entropy case}
\label{p12:ef:section}

The previous manuscript excluded rational bit entropy for nonrational
ternary probabilities in totally real number fields. Here we exclude
\emph{every cubic probability field}, including fields with a complex
pair of embeddings. The new ingredient is an arithmetic restriction
on the differences of the three Laurent exponents: their greatest
common divisor is a power of two. This also identifies the probability
field with its unit-generator field whenever the entropy denominator
is odd. No unrestricted three-versus-four alphabet theorem is claimed.

Write $H_2(p)=-\sum_i p_i\log_2p_i$. Probabilities are strictly positive,
and ``nonrational vector'' means that at least one entry is irrational.
The probability field is $K=\mathbb Q(p_1,p_2,p_3)$.

\subsection{The inherited normal form}

We state and prove the previously obtained reduction to specify exactly
which inputs the new argument uses. Baker's logarithmic independence
theorem is the transcendence-theoretic input \cite{p07:rlog:BakerIII}.
The arithmetic of entropy values is also discussed in
\cite{p10:inf:BW}; the following field-degree question is separate
from the non-semialgebraicity questions treated there.

\begin{lemma}[Inherited unit normal form]\label{p12:ef:normal}
Suppose $p=(p_1,p_2,p_3)$ is a nonrational algebraic probability vector
and $h=H_2(p)\in\mathbb Q$. Write $h=a/b$ in lowest positive terms,
put $C=2^h$, and let $d=[K:\mathbb Q]$. There is a positive algebraic
unit $w\ne1$ and pairwise distinct odd integers $m_1,m_2,m_3$ such that
\begin{equation}\label{p12:ef:normalid}
 p_i=C^{-1}w^{m_i},\qquad
 \sum_iw^{m_i}=C,\qquad \sum_i m_iw^{m_i}=0,
 \qquad \gcd(m_1,m_2,m_3)=1.
\end{equation}
There are exponents of both signs. With $L=\mathbb Q(w)$ one has
\begin{equation}\label{p12:ef:inherited}
 L=K(C),\qquad 1\le h<\log_2 3,\qquad hd\in\mathbb Z,
 \qquad
 r_1(L)=\begin{cases}1&b\text{ odd},\\2&b\text{ even}.
 \end{cases}
\end{equation}
Here $r_1$ counts real embeddings. Every $p_i$ is irrational.
\end{lemma}
\begin{proof}
Let $t$ be the rational dimension of the span of the probabilities,
and let $r$ be that of $\log p_1,\log p_2,\log p_3,\log2$.
Splitting the entropy relation against a rational basis of the
probabilities, and applying Baker's theorem, gives $t$ independent
rational relations among the four logarithms. Thus $t+r\le4$.
Nonrationality gives $t\ge2$. If $r=1$, the probabilities are rational
powers of two. Clearing a common denominator of the exponents and
using the independence of $1,2^{1/q},\ldots,2^{(q-1)/q}$, by Eisenstein
for $X^q-2$, shows that a positive sum of such powers can equal one
only if all the powers have integral exponents. This contradicts
nonrationality. Hence $t=r=2$.

Write $\log p_i=\alpha_i\log2+\beta_i\log u$ in a rational basis,
with rational $\alpha_i,\beta_i$ and positive algebraic $u$.
Baker's theorem gives $\sum\beta_i p_i=0$ and
$\sum\alpha_i p_i=-h$. The rational kernel of
$(x_i)\mapsto\sum x_i p_i$ is one-dimensional, so
$\alpha_i+h=\lambda\beta_i$ for one rational $\lambda$.
Clearing denominators and making the exponent vector primitive gives
\eqref{p12:ef:normalid}, initially without the unit and parity assertions.
The logarithms of $w$ and $2$ are rationally independent.
Equal exponents would give two equal probabilities and, using
stationarity and normalization, make all three probabilities rational.
Thus the exponents are distinct. Their positive weighted mean zero
forces both signs.

The number $C$ is integral. At a finite valuation, $v(w)>0$ makes the
unique smallest exponent contribute a strictly negative, uniquely
minimal valuation to $\sum w^{m_i}$; $v(w)<0$ does the same with the
largest exponent. Both contradict integrality of $C$. Thus $w$ is a
unit. Reducing normalization and stationarity at a prime above two
shows that either all three exponents are odd or all three are even:
otherwise subtraction isolates one nonzero residue term. Primitivity
excludes the even case. Consequently
$C=\sum(1-m_i)w^{m_i}\in2\mathcal O_L$, so $h\ge1$.
The entropy bound is strict since the uniform vector is rational.

Bezout for the primitive exponents gives $w\in K(C)$, and the sum
identity gives $C\in\mathbb Q(w)$; hence $L=K(C)$. Taking absolute
norms in $p_i=C^{-1}w^{m_i}$ and descending from $L$ to $K$ gives
$|N_{K/\mathbb Q}(p_i)|=2^{-hd}$. Rationality of this number forces
$hd\in\mathbb Z$. A rational positive $p_i$ would then equal $C^{-1}$,
forcing $m_i=0$, contrary to oddness.

Finally $F(x)=\sum x^{m_i}$ has a unique positive stationary point,
because $y\mapsto F(e^y)$ is strictly convex and its derivative tends
to opposite infinities at the two ends. Oddness of the exponents
makes its only possible real stationary points $w,-w$.
Every real conjugate of $w$ must be one of these. The field
$\mathbb Q(C)=\mathbb Q(2^{1/b})$ has degree $b$ by Eisenstein and
Bezout. If $b$ is odd, its only real conjugate of $C$ is positive,
so only $w$ occurs. If $b$ is even, $[L:\mathbb Q]$ is even, and
the positive real-embedding count, at most two, must also be even.
This proves the signature assertion.
\end{proof}

\subsection{The new restriction on exponent differences}

\begin{theorem}[Two-primary exponent spacing]\label{p12:ef:spacing}
Under Lemma~\ref{p12:ef:normal}, set
$g=\gcd(m_2-m_1,m_3-m_1)>0$. Then
\begin{equation}\label{p12:ef:spacingid}
                    g=2^s\quad(s\ge1),\qquad
                    [L:K]\text{ is a power of }2.
\end{equation}
If the reduced denominator $b$ of $h$ is odd, then
\begin{equation}\label{p12:ef:odd}
                    L=K,\qquad r_1(K)=1,
                    \qquad [K:\mathbb Q]\text{ is odd}.
\end{equation}
\end{theorem}
\begin{proof}
All exponents are odd, so $g$ is even. Suppose an odd prime $\ell$
divides $g$. Then $m_i\equiv r\pmod\ell$ for a common residue $r$,
which is nonzero because the exponent vector is primitive. Write
$m_i=r+\ell n_i$ with integers $n_i$. Stationarity yields
\[
              0=rC+\ell\sum_i n_iw^{m_i}.
\]
The sum is an algebraic integer, while $rC$ is a unit at every prime
above $\ell$: indeed $C^b=2^a$ and $\ell\nmid2r$. Reduction at such
a prime is a contradiction. Therefore $g$ has no odd prime divisor.

The ratios $p_2/p_1=w^{m_2-m_1}$ and $p_3/p_1=w^{m_3-m_1}$ belong
to $K$. Bezout for these two differences gives $w^g\in K$.
The tower adjoining successively $w^{2^{s-1}},w^{2^{s-2}},\ldots,w$
has every step of degree at most two. Since $K\subseteq L=\mathbb Q(w)$,
this proves that $[L:K]$ is a power of two.

If $b$ is odd, Lemma~\ref{p12:ef:normal} gives $r_1(L)=1$. The identity
$[L:\mathbb Q]=r_1(L)+2r_2(L)$ makes $[L:\mathbb Q]$ odd.
Its factor $[L:K]$, being a power of two, must be one.
The remaining assertions follow.
\end{proof}

The exclusion concerns the gcd of the \emph{differences}, not merely
the gcd of the exponents. For example, every primitive triple of odd
exponents all congruent to $1$ modulo $3$ is excluded immediately,
without evaluating a stationary point.

\subsection{Two elementary cubic unit bounds}

\begin{lemma}[A small cubic unit exclusion]\label{p12:ef:cubicunit}
There is no algebraic unit $u$ of degree three with exactly one real
conjugate and $1<u<2^{1/3}$.
\end{lemma}
\begin{proof}
Write the other conjugates as $z,\bar z$. Positivity of $u$ and the
unit property give $u|z|^2=1$. The minimal polynomial therefore has
the form $X^3-sX^2+tX-1$, with $s,t\in\mathbb Z$, and
\[
 s=u+2\operatorname{Re}z,\qquad
 t=su-u^2+u^{-1}.
\]
Put $\theta=2^{1/3}$. Since $1<u<\theta<4/3$,
$-1<u-2/\sqrt u\le s\le u+2/\sqrt u<4$, whence
$s\in\{0,1,2,3\}$. For these four integers set
$f_s(x)=sx-x^2+x^{-1}$. For $x>1$,
\[
 f_s'(x)\le3-2x-x^{-2}
       =-\frac{(x-1)^2(2x+1)}{x^2}<0,
\]
so $f_s(u)<f_s(1)=s$. On the other hand,
\[
 f_s(u)-(s-1)
 =1-u^2+u^{-1}+s(u-1)
 >1-\theta^2+\theta^{-1}
 =1-\frac{\theta^2}{2}>0.
\]
Thus the integer $t$ lies strictly between $s-1$ and $s$, impossible.
\end{proof}

\begin{lemma}[A pure cubic unit exclusion]\label{p12:ef:purecubicunit}
In $E=\mathbb Q(\theta)$, $\theta=2^{1/3}$, there is no algebraic
unit $u$ satisfying $1<u<2\theta$.
\end{lemma}
\begin{proof}
We use the standard integral basis
$\mathcal O_E=\mathbb Z[\theta]$, with a verification recalled here
to keep the arithmetic input explicit; see also \cite{p12:ef:Conrad}.
The power-basis discriminant is $-108$, so its index can have only
the prime divisors two and three. Locally at two, $X^3-2$ is
Eisenstein, and $\theta$ is a uniformizer generating the full ring
of integers. Locally at three the generator $1+\theta$ satisfies
the Eisenstein polynomial
$X^3-3X^2+3X-3$. Thus neither prime divides the index, as claimed.

Write $u=A+B\theta+D\theta^2$ with integers $A,B,D$.
Let $\zeta=e^{2\pi i/3}$ and
$z=A+B\theta\zeta+D\theta^2\zeta^2$; the third conjugate is
$\bar z$. As $u$ is a positive unit, $u|z|^2=1$.
Fourier inversion gives
\[
 3A=u+z+\bar z,\quad
 3B\theta=u+\zeta^2z+\zeta\bar z,\quad
 3D\theta^2=u+\zeta z+\zeta^2\bar z.
\]
Each right side is real and belongs to
$[u-2/\sqrt u,u+2/\sqrt u]$.
The lower endpoint is greater than $-1$, so integrality gives
$A,B,D\ge0$. The function $x+2/\sqrt x$ is strictly increasing
for $x>1$. At $x=2\theta=\theta^4$ its value is $3\theta$.
Consequently
\[
              A<\theta<2,\qquad B<1,\qquad D<\theta^{-1}<1.
\]
It follows that $B=D=0$ and $A\in\{0,1\}$, contradicting $u>1$.
\end{proof}

These two lemmas are elementary unit bounds, not claims of new
results about the smallest Pisot number or the fundamental unit of
a pure cubic field. Their role here is to eliminate the two possible
rational entropy values in a cubic probability field.

\subsection{Closing the cubic probability-field case}

\begin{theorem}[No cubic ternary rational entropy]\label{p12:ef:cubic}
Let $p_1,p_2,p_3$ be positive algebraic numbers with sum one. If
$[\mathbb Q(p_1,p_2,p_3):\mathbb Q]=3$, then
\[
                         H_2(p_1,p_2,p_3)\notin\mathbb Q.
\]
No total-reality assumption is required.
\end{theorem}
\begin{proof}
Suppose $h$ is rational. The normal form gives $3h\in\mathbb Z$ and
$1\le h<\log_2 3<5/3$, so $h=1$ or $h=4/3$.
Both denominators are odd. Theorem~\ref{p12:ef:spacing} gives $L=K$,
of signature $(1,1)$. Replacing $w$ by $w^{-1}$ and all exponents
by their negatives if necessary, assume $w>1$.

First suppose $h=1$, so $C=2$. Two positive exponents would make
their two monomials alone exceed two. There is therefore one
positive odd exponent $a$, with the other two equal to $-b,-c$,
where $1\le b<c$ are odd. If $a=1$ and $b=1$, then
$w+w^{-1}+w^{-c}>2$. If $a=1$ and $b\ge3$, then $c\ge5$ and
\[
 F(\sqrt2)=\sqrt2+2^{-b/2}+2^{-c/2}
       \le\frac{11}{4\sqrt2}<2.
\]
But $w$ is the global minimum of $F$ on the positive real axis by
the strict convexity argument in Lemma~\ref{p12:ef:normal}; hence
$F(w)\le F(\sqrt2)<2$, another contradiction.
Thus $a\ge3$, and the normalization gives $w^a<2$, whence
$1<w<2^{1/3}$. This contradicts Lemma~\ref{p12:ef:cubicunit}.

Now suppose $h=4/3$. Since $C=2^{4/3}\in L=K$ and both
$\mathbb Q(C)$ and $K$ have degree three,
$K=\mathbb Q(2^{1/3})$. At least one exponent $a$ is positive,
so $w\le w^a<C=2\theta$. This contradicts
Lemma~\ref{p12:ef:purecubicunit}.
\end{proof}

\begin{corollary}[Degree barrier for the unresolved ternary case]
\label{p12:ef:barrier}
If a nonrational algebraic ternary probability vector has rational
bit entropy, its probability field has degree at least four.
If the reduced entropy denominator is odd, the probability field
has exactly one real embedding and odd degree at least five.
In particular, a one-bit example would require degree at least five.
\end{corollary}
\begin{proof}
Degree one is excluded by nonrationality and degree three by
Theorem~\ref{p12:ef:cubic}. For completeness, degree two is also
impossible: the real quadratic probability field $K$ is totally
real, the norm restriction leaves $h=1$ or $3/2$, and
$L=K(2^h)\subseteq K(\sqrt2)$ is totally real.
Yet Lemma~\ref{p12:ef:normal} permits at most the real conjugates
$w,-w$ for $w$. Total reality would force $w^2$ to be rational;
as a positive rational algebraic unit it equals one, a contradiction.
For odd denominator use Theorem~\ref{p12:ef:spacing} and exclude the
remaining odd degrees one and three.
\end{proof}

\paragraph{Scope of the advance.}
The cubic theorem rules out a family not covered by the earlier
totally-real theorem. It does not exclude non-totally-real fields
of degree four or higher, does not establish transcendence of these
entropy values, and does not classify algebraic irrational entropy.
The unrestricted least alphabet for entirely irrational algebraic
probabilities and rational bit entropy remains either three or four.
The exact value four for totally real probabilities belongs to the
previous manuscript and is not recounted as a new result here.


% END EMBEDDED entropy_results.tex


% BEGIN EMBEDDED gamma_results.tex

% Fragment: no preamble and no bibliography environment.
\section{Actual arithmetic separation of fifth and seventh beta periods}
\label{p12:sec:actual-beta-separation}

The previous formal Gamma calculus can be upgraded to a statement about
\emph{actual numbers} on a restricted, but infinite, family.  The upgrade
uses the classical Wolfart--W\"ustholz theorem; it does not establish the
Rohrlich conjecture for arbitrary Gamma monomials.  The resulting rank
calculation and linear-fractional classification below are explicit
specializations of classical period theory, not claimed solutions of
previously open individual Gamma-value problems.

\subsection{The precise classical input}

Write $u\sim v$ if $u/v\in\overline{\mathbb Q}^{\times}$.  We use the
following two parts of the beta-value theorem of Wolfart and W\"ustholz
\cite{p12:WW1985CF}.  The first is also stated as Theorem~5 in
\cite{p12:RivoalRoques2014CF}.
\begin{enumerate}
\item For positive rational $a_j,b_j$, algebraic linear dependence of
the numbers $B(a_j,b_j)$ implies that two of them have algebraic ratio.
Equivalently, representatives of distinct $\sim$-classes of positive
rational beta values are linearly independent over
$\overline{\mathbb Q}$.
\item An algebraic ratio of two rational beta values is accounted for
by the classical Gamma identities: recurrence, reflection, and Gauss
multiplication, with algebraic factors and rational powers admitted.
In particular, any invariant of these identities must assign the same
value to the two beta monomials.
\end{enumerate}
The second assertion is the beta-value case of the classical
Gauss--Legendre completeness statement in \cite{p12:WW1985CF}; the first
assertion alone would not justify our signature test.  This completeness
is known for these periods of Fermat Jacobians.  It is not known for
arbitrary Gamma monomials.  We also use Schneider's classical theorem
that $B(a,b)$ is transcendental when $a,b\in\mathbb Q\setminus\mathbb Z$
and $a+b\notin\mathbb Z_{\leq0}$; see Theorem~4 of
\cite{p12:RivoalRoques2014CF} and \cite{p12:Schneider1941CF}.

\subsection{A signature that survives unrestricted auxiliary denominators}

Put
\[
s(x)=\begin{cases}1-2\{x\},&x\notin\mathbb Z,\\0,&x\in\mathbb Z.
\end{cases}
\]
For $m\geq2$, $1\leq a,b<m$, and $t\in(\mathbb Z/m\mathbb Z)^\times$,
define
\[
\sigma_{m;a,b}(t)=s(ta/m)+s(tb/m)-s(t(a+b)/m).
\]
If $a+b=m$ the signature is identically zero.  Otherwise it takes its
values in $\{1,-1\}$, and is odd under $t\mapsto-t$.

Here is the verification that these are legitimate obstructions, even
when a proposed identity introduces other denominators.  The elementary
identity
\[
\sum_{j=0}^{h-1}s(x+j/h)=s(hx),\qquad h\geq1,
\]
holds at endpoints as well as away from them.  When $hx$ is not an
integer it follows by summing fractional parts; when $hx$ is an integer,
the nonintegral residues pair as $u,1-u$ and the integral residue has
value zero.  Also $s(x+1)=s(x)$ and $s(-x)=-s(x)$.
For a finite Gamma-identity derivation choose $L$ divisible by every
denominator and every multiplication index in the derivation.  For each
$t$ prime to $L$, assign $\Gamma(x)$ the additive value $s(tx)$ and assign
algebraic factors and powers of $\pi$ the value zero.  Recurrence and
reflection preserve this assignment.  In a multiplication identity,
$j\mapsto tj$ permutes the residues modulo its multiplication index,
so the displayed distribution identity proves preservation too.
Reduction of units modulo a divisor of $L$ is surjective.  Thus unequal
signatures cannot become equal by introducing auxiliary denominators.

\begin{lemma}[Actual separation from unequal signatures]
\label{p12:lem:cf-actual-signature}
Let $\beta=B(a/m,b/m)$ and $\beta'=B(c/n,d/n)$ with
$1\leq a,b<m$ and $1\leq c,d<n$.  Inflate the two signatures to the
units modulo $\operatorname{lcm}(m,n)$.  If the inflated signatures
differ, then $\beta/\beta'$ is transcendental.
\end{lemma}
\begin{proof}
An algebraic ratio would, by the classical beta completeness theorem,
admit a standard Gamma-identity derivation.  The finite-derivation
invariant just proved would then make the signatures equal.  This is a
contradiction.  Since the ratio is nonzero, failure to be algebraic is
precisely transcendence.
\end{proof}

\begin{theorem}[Separation of coprime conductors]
\label{p12:thm:cf-coprime-beta}
For $m\geq2$ let
\[
V_m=\operatorname{span}_{\overline{\mathbb Q}}
 \{B(a/m,b/m):1\leq a,b<m\}.
\]
If $m_1,\ldots,m_k\geq2$ are pairwise coprime, then the sum of the
quotient spaces $V_{m_i}/\overline{\mathbb Q}\pi$ is direct inside
$(V_{m_1}+\cdots+V_{m_k})/\overline{\mathbb Q}\pi$.  In particular,
\[
V_m\cap V_n=\overline{\mathbb Q}\pi\quad((m,n)=1),
\qquad
\dim\bigl(V_{m_1}+\cdots+V_{m_k}\bigr)
=1+\sum_{i=1}^k(\dim V_{m_i}-1).
\]
Moreover $1$ is linearly independent of this sum.
\end{theorem}
\begin{proof}
Each $V_m$ contains $\pi$ by reflection.  Its beta values with zero
signature are algebraic multiples of $\pi$.  A nonzero signature cannot
equal the zero signature.  A signature from conductor $m$ cannot equal
a nonzero signature from a coprime conductor $n$: by the Chinese
remainder theorem, such an equality would equate a function of an
arbitrary unit modulo $m$ with a function of an independent arbitrary
unit modulo $n$.  Both functions would be constant.  Their oddness
would force them to be zero, contrary to their values $\pm1$.
Choose representatives of the algebraic-ratio classes within each
conductor.  Lemma~\ref{p12:lem:cf-actual-signature} shows that the only class
shared by two conductors is the class of $\pi$.  The Wolfart--W\"ustholz
linear independence theorem now proves all stated assertions about the
sum and its dimension.  Finally include $B(1,1)=1$ among the beta
values.  Schneider's theorem says that none of the other classes is its
class; the same linear independence theorem includes $1$.
\end{proof}

\subsection{A complete fourteen-dimensional calculation}

Set
\[
X=\Gamma(1/5),\quad Y=\Gamma(2/5),\qquad
R=\Gamma(1/7),\quad S=\Gamma(2/7),\quad T=\Gamma(3/7),
\]
and define the actual positive numbers
\[
A=\frac{X^2}{Y},\quad B=\frac{XY^2}{\pi},\qquad
U=\frac{R^2}{S},\quad V=\frac{RS}{T},\quad
W=\frac{RT^2}{\pi},\quad Z=\frac{S^2T}{\pi}.
\]
Here $A=B(1/5,1/5)$ and $\sin(2\pi/5)B=B(1/5,2/5)$;
$U=B(1/7,1/7)$, $V=B(1/7,2/7)$,
$\sin(3\pi/7)W=B(1/7,3/7)$, and
$\sin(3\pi/7)Z=B(2/7,2/7)$.  These sine factors are nonzero algebraic
numbers.

\begin{theorem}[Exact mixed beta-period basis]
\label{p12:thm:cf-fourteen}
The following ordered list is linearly independent over
$\overline{\mathbb Q}$:
\[
\mathcal M=\left(1,\pi,A,B,\frac\pi A,\frac\pi B,
 U,V,W,Z,\frac\pi U,\frac\pi V,\frac\pi W,\frac\pi Z\right).
\]
It is a basis of $\overline{\mathbb Q}+V_5+V_7$.  Consequently
\[
\dim V_5=5,\qquad \dim V_7=9,\qquad
\dim(\overline{\mathbb Q}+V_5+V_7)=14.
\]
\end{theorem}
\begin{proof}
The signatures of $A,B$ on $t=1,2,3,4$ are respectively
\[
(1,1,-1,-1),\qquad(1,-1,1,-1).
\]
The signatures of $U,V,W,Z$ on $t=1,\ldots,6$ are respectively
\[
\begin{array}{c|rrrrrr}
U&1&1&1&-1&-1&-1\\
V&1&1&-1&1&-1&-1\\
W&1&-1&1&-1&1&-1\\
Z&1&-1&-1&1&1&-1
\end{array}
\]
The signature of $\pi/Q$ is the negative of that of $Q$.  These four
fifth signatures and eight seventh signatures are distinct and nonzero.
Each reciprocal is, up to an algebraic factor, a positive rational beta
value: for nonintegral $a,b,a+b$ with $0<a,b<1$,
\[
B(a,b)B(1-a,1-b)
=\frac{\pi\sin(\pi(a+b))}
 {(1-a-b)\sin(\pi a)\sin(\pi b)}.
\]
The expression remains positive when $a+b>1$ because both the numerator
sine and $1-a-b$ change sign.  The values $\pi$ and $1$ are represented
by $B(1/2,1/2)$ and $B(1,1)$ respectively.  Lemma~\ref{p12:lem:cf-actual-signature},
the coprime-conductor argument, and Schneider's theorem therefore
separate all fourteen algebraic-ratio classes.  Their linear
independence follows from Wolfart--W\"ustholz.

For completeness the spanning calculation uses only recurrence and
reflection.  To reduce $B(a/p,b/p)$, write a Gamma argument as $r/p$
modulo integers, replace $\Gamma(r/p)$ for $r>(p-1)/2$ by an algebraic
multiple of $\pi/\Gamma((p-r)/p)$, and collect exponents.  For $p=5$
the resulting classes are exactly
\[
\pi,\ A,\ B,\ \pi/A,\ \pi/B,
\]
and for $p=7$ they are exactly
\[
\pi,\ U,\ V,\ W,\ Z,\ \pi/U,\ \pi/V,\ \pi/W,\ \pi/Z.
\]
An exhaustive integer-exponent enumeration of the $16$ and $36$ pairs
is supplied with the manuscript; alternatively the displayed reduction
checks this finite assertion directly.  This proves spanning and all
dimensions.
\end{proof}

\subsection{An exact algebraicity test on an infinite explicit family}

For $c=(c_0,\ldots,c_{13})\in\overline{\mathbb Q}^{14}$ write
$L_c=\sum_j c_j\mathcal M_j$.
\begin{corollary}[Linear-fractional arithmetic classification]
\label{p12:cor:cf-projective}
For $d\ne0$, $L_d\ne0$, and
\[
\frac{L_c}{L_d}\in\overline{\mathbb Q}
\quad\Longleftrightarrow\quad
c=\alpha d\text{ for some }\alpha\in\overline{\mathbb Q}.
\]
In particular, $L_c$ is algebraic if and only if
$c_1=\cdots=c_{13}=0$.  Every ratio of two distinct entries of
$\mathcal M$ is transcendental.
\end{corollary}
\begin{proof}
The denominator cannot vanish by Theorem~\ref{p12:thm:cf-fourteen}.  If the
ratio is $\alpha\in\overline{\mathbb Q}$, then
$L_{c-\alpha d}=0$, so linear independence forces $c=\alpha d$.
The converse is immediate, and taking $d=(1,0,\ldots,0)$ proves the
algebraicity statement.
\end{proof}

This yields actual monomial exclusions that do not assert algebraic
independence of the five individual Gamma values.  Examples include
\[
\begin{gathered}
\frac{\pi X}{Y^3},\quad
\frac{X^4}{\pi Y^2},\quad
\frac{X^2Y^4}{\pi^3},\quad
\frac{RT}{S^2},\quad
\frac{\pi S}{T^3},\quad
\frac{\pi R^2}{S^3T},\\[1ex]
\frac{R^4}{\pi S^2},\quad
\frac{R^2S^2}{\pi T^2},\quad
\frac{R^2T^4}{\pi^3},\quad
\frac{S^4T^2}{\pi^3},\quad
\frac{X^2S}{YR^2},\quad
\frac{XY^2T}{\pi RS}.
\end{gathered}
\]
They are respectively $A/B$, $A/(\pi/A)$, $B/(\pi/B)$,
$U/V$, $V/W$, $U/Z$, $U/(\pi/U)$, $V/(\pi/V)$,
$W/(\pi/W)$, $Z/(\pi/Z)$, $A/U$, and $B/V$.
They are twelve explicit consequences of one classical period theorem
and the finite signature certificate, not twelve newly closed famous
conjectures.

For a polynomial version, replace $\pi$ with a formal variable $P$
in the fourteen Laurent monomials defining $\mathcal M$. Regard
$P,X,Y,R,S,T$ as independent variables. Evaluation at the
actual six constants is injective on their
$\overline{\mathbb Q}$-linear span.  Multiplication by one common Laurent
monomial turns this into an explicit finite-support polynomial family
having no nonzero vanishing member.  This is a complete statement for
that support, and extends to ratios of two such expressions by
Corollary~\ref{p12:cor:cf-projective}.

\subsection{The remaining boundary}

Linear independence of this fourteen-element list is compatible with
nonlinear relations among its entries.  Indeed the exact identity
\[
UZ=VW
\]
already holds.  The argument therefore cannot be promoted to algebraic
independence of the beta periods.  It also does not decide the
irrationality or transcendence of any one of $X,Y,R,S,T$ separately.
The new contribution to the present manuscript is an explicit actual
period-space calculation, its coprime-conductor separation proof, and
the resulting complete linear-fractional criterion.  The deep
transcendence input and its general beta-value completeness are
classical, and no claim of first historical discovery is made for
these specializations.


% END EMBEDDED gamma_results.tex


% BEGIN EMBEDDED gompertz_results.tex

% Fragment; all new labels use fg:.
\section{Fixed-parameter modular scarcity and the prime cost of cancellation}
\label{p12:fg:section}

The parameter law of the previous paper does not specialize to $A=1$.
Here the parameter remains fixed throughout.  We instead bound the
number of \emph{index residues} at which a prime divides the classical
Laguerre denominator.  No assertion about simplicity of these index
roots is required.  This yields a uniform bound on mean logarithmic
content at scale $n$, together with an obstruction to producing a small
primitive form using only a few primes.  It does not establish the
irrationality of the Euler--Gompertz constant.

Fix a positive integer $A$ and define
\[
 \delta_A=\int_0^\infty\frac{e^{-t}}{1+At}\,dt,\qquad
 Q_n=Q_n(A)=\sum_{j=0}^n j!\binom nj^2A^j,
\]
Let $P_0=0$, $P_1=1$, and give both sequences the recurrence
\[
 U_n=(1+(2n-1)A)U_{n-1}-(n-1)^2A^2U_{n-2}.
\]
For $n\ge1$, put
\[
 g_n=\gcd(Q_n,((n-1)!)^2),
 \quad \ell_n=\frac{Q_n\delta_A-P_n}{g_n}>0.
\]
These definitions have the following exact arithmetic and analytic
identities:
\begin{equation}\label{p12:fg:identities}
 g_n=\gcd(P_n,Q_n),\qquad
 Q_n\delta_A-P_n
   =n!A^n\int_0^\infty\frac{e^{-t}}{1+At}
                         \left(\frac{At}{1+At}\right)^n\,dt.
\end{equation}
In particular,
\begin{equation}\label{p12:fg:base}
 \log(Q_n\delta_A-P_n)=n\log n+O_A(n),\qquad
 v_p(g_n)\le\frac{2n}{p-1}\mathbf1_{\{p\mid Q_n\}}.
\end{equation}
The approximations themselves and their real and $p$-adic settings
are classical; see \cite{p12:fg:MY,p12:fg:MZ}.  The estimates below are consequences
of their recurrence, rather than new evaluations of $\delta_A$.

Here are short verifications to make the section self-contained.
The binomial formula implies $Q_{n+p}\equiv Q_n\pmod p$ for every
prime $p$, by deleting terms with $j\ge p$ and reducing the other
binomial coefficients. If a prime divided two consecutive $Q_n$,
it would not divide $A$; reducing indices modulo $p$ and using the
recurrence backwards would reach $Q_0=1$, a contradiction.
Thus $\gcd(Q_n,Q_{n-1})=1$. The recurrence also gives
\[
 P_nQ_{n-1}-P_{n-1}Q_n=A^{2n-2}((n-1)!)^2.
\]
Since $Q_n\equiv1\pmod A$, taking the gcd with $Q_n$ proves the
first identity in \eqref{p12:fg:identities}.

For the second, put $F_n(t)=n!A^nL_n(t)$, so
$F_n(-1/A)=Q_n$.  The polynomial quotient
$(Q_n-F_n(t))/(1+At)$ has integral $P_n$ against $e^{-t}dt$:
the Laguerre recurrence gives exactly the stated recurrence and
initial values after integration, since
$\int_0^\infty e^{-t}L_m(t)dt=0$ for $m\ge1$.
Rodrigues's identity
$n!e^{-t}L_n(t)=(d/dt)^n(e^{-t}t^n)$ and $n$ integrations by
parts now give the integral in \eqref{p12:fg:identities}; all boundary
terms vanish. Its upper bound is $n!A^n\delta_A$, while restricting
the integral to $[\sqrt n,\sqrt n+1]$ gives a lower bound
$n!A^n\exp(-O_A(\sqrt n+\log n))$. Stirling's estimate proves
the logarithmic assertion in \eqref{p12:fg:base}. The valuation bound
there follows directly from $g_n\mid((n-1)!)^2$ and Legendre's
formula.

\subsection{Short shifts constrain the index root set}

For a prime $p$, put
\[
 R_p(A)=\{r\in\mathbb F_p:Q_r(A)=0\pmod p\},\qquad
 \rho_p(A)=|R_p(A)|,
\]
where $r$ in $Q_r$ is represented by $0\le r<p$.
The binomial formula shows that $Q_{n+p}\equiv Q_n\pmod p$.
If $p\mid A$, then $R_p(A)$ is empty.

\begin{theorem}[Uniform scarcity of index roots]\label{p12:fg:roots}
For $p\nmid A$ and $1\le d<p$,
\begin{equation}\label{p12:fg:pair}
 |R_p(A)\cap(R_p(A)-d)|\le d-1.
\end{equation}
Consequently, for every integer $1\le h\le p$,
\begin{equation}\label{p12:fg:secondmoment}
 h\rho_p(A)^2\le p\rho_p(A)+\frac{p(h-1)(h-2)}3.
\end{equation}
In particular, with the absolute constant
$C_0=(1+\sqrt{7/3})/2$,
\begin{equation}\label{p12:fg:rootbound}
                         \rho_p(A)\le C_0p^{2/3}.
\end{equation}
\end{theorem}
\begin{proof}
Consecutive $Q_n$ cannot both vanish modulo $p$.  Indeed, after reducing
the indices modulo $p$, the recurrence
\[
 Q_k=(1+(2k-1)A)Q_{k-1}-(k-1)^2A^2Q_{k-2}
\]
allows descent to $Q_0=1$; an index congruent to zero already has value
one.  Define polynomials in $x$ by
\[
 C_1(x)=1,\quad C_2(x)=1+(2x+3)A,
\]
and, for $d\ge3$, by
\[
 C_d(x)=(1+(2x+2d-1)A)C_{d-1}(x)
                 -(x+d-1)^2A^2C_{d-2}(x).
\]
Reduction of the recurrence modulo $Q_n$ gives
$Q_{n+d}\equiv C_d(n)Q_{n+1}\pmod{Q_n}$.
The coefficient of $x^{d-1}$ in $C_d$ is $dA^{d-1}$: its recurrence
is $c_d=2Ac_{d-1}-A^2c_{d-2}$ with $c_1=1,c_2=2A$.
Thus $C_d$ has degree exactly $d-1$ over $\mathbb F_p$ when $d<p$.
If both $r$ and $r+d$ are index roots, then $Q_{r+1}\ne0$ and
$C_d(r)=0$.  This proves \eqref{p12:fg:pair}, including shifts that wrap
around the residue circle, by the periodicity already noted.

For each $x\in\mathbb F_p$ let
$t_x=\sum_{j=1}^h\mathbf1_{R_p(A)}(x+j)$.
Then $\sum_x t_x=h\rho_p(A)$ and
\begin{align*}
 \sum_x t_x^2
 &=h\rho_p(A)+2\sum_{d=1}^{h-1}(h-d)
                    |R_p(A)\cap(R_p(A)-d)|\\
 &\le h\rho_p(A)+\frac{h(h-1)(h-2)}3.
\end{align*}
Cauchy--Schwarz gives \eqref{p12:fg:secondmoment}.
Take $h=\lceil p^{1/3}\rceil$.  Since $p/h\le p^{2/3}$ and
$(h-1)(h-2)/h\le h-1\le p^{1/3}$, the resulting quadratic inequality is
\[
 \rho_p(A)^2\le p^{2/3}\rho_p(A)+\tfrac13p^{4/3}.
\]
Its positive root is \eqref{p12:fg:rootbound}.
\end{proof}

This argument concerns zeros in the discrete index, not zeros in the
Laguerre parameter.  In particular it neither asserts nor needs
squarefreeness of a polynomial interpolating the index sequence.
The method is an elementary short-shift counting argument.  Its
application here is proved in full; priority for this particular
application has not been established by the literature search.

\subsection{A limiting law for a prime-support upper bound}

Define
\[
 W_A(n)=2\sum_{\substack{p\le n\\p\mid Q_n(A)}}\frac{\log p}{p-1},
 \qquad b_p=\frac{2\log p}{p-1},\qquad
 a_p=\frac{\rho_p(A)}p.
\]
In particular $\log g_n/n\le W_A(n)$.
Let $B_p$ be independent Bernoulli random variables with
$\mathbb P(B_p=1)=a_p$.

\begin{theorem}[Fixed-parameter linear-scale control]\label{p12:fg:law}
The nonnegative random series $W_A=\sum_p b_p B_p$ converges almost
surely and in $L^1$.  As $N\to\infty$, the empirical distribution of
$W_A(n)$ for $1\le n\le N$ converges weakly to that of $W_A$, and its
first moment converges to
\begin{equation}\label{p12:fg:mean}
 \mu_A=2\sum_p\frac{\rho_p(A)\log p}{p(p-1)}<\infty.
\end{equation}
The constants $\mu_A$ are bounded by an absolute constant uniformly
in positive integer $A$.  Thus
\begin{equation}\label{p12:fg:gmean}
 \limsup_{N\to\infty}\frac1N\sum_{n=1}^N\frac{\log g_n(A)}n
                                  \le\mu_A.
\end{equation}
The moment generating function exists for every real $t$ and equals
\begin{equation}\label{p12:fg:mgf}
 M_A(t)=\mathbb E e^{tW_A}
       =\prod_p\left(1+a_p(e^{tb_p}-1)\right).
\end{equation}
There are absolute positive constants $c,L_0$ such that, for $L\ge L_0$,
\begin{equation}\label{p12:fg:tail}
 \limsup_{N\to\infty}\frac1N
    \#\{n\le N:\log g_n(A)\ge Ln\}
                        \le \exp(-c e^{L/2}).
\end{equation}
\end{theorem}
\begin{proof}
The root bound implies $a_p\le C_0p^{-1/3}$, and hence
$\sum_p a_pb_p<\infty$, even if the sum is enlarged to all integers
at least two.  Finite expectation proves convergence of the positive
random series and also convergence in $L^1$.

Fix $Y$.  By the Chinese remainder theorem, the divisibility events
$p\mid Q_n(A)$ for $p\le Y$ have independent limiting frequencies
$a_p$ as $n$ runs through integers.  This is a finite periodic statement.
The omission of terms $p>n$ changes this truncated assertion at only
finitely many indices.  To pass to all primes, use the elementary count
\[
 \#\{n\le N:p\mid Q_n(A)\}
       \le(N/p+1)\rho_p(A).
\]
For an absolute $C$, uniformly in $A$, this gives
\begin{align*}
 \frac1N\sum_{n\le N}\sum_{\substack{Y<p\le n\\p\mid Q_n}}b_p
 &\le C\sum_{p>Y}\frac{\log p}{p^{4/3}}
        +\frac C N\sum_{p\le N}\frac{\log p}{p^{1/3}}\\
 &\le C Y^{-1/3}\log(2Y)+C N^{-1/3}\log(2N).
\end{align*}
The final inequalities follow by comparison with sums over integers.
Truncation therefore proves convergence against bounded Lipschitz
functions, hence weak convergence, and also convergence of the first
moment.  This proves \eqref{p12:fg:mean} and \eqref{p12:fg:gmean}.

For fixed $t$, the tail of the logarithm of the product in
\eqref{p12:fg:mgf} is absolutely convergent because
$e^{tb_p}-1=O_t(b_p)$.  Independence and monotone convergence for
$t\ge0$ give the claimed identity; bounded convergence covers $t<0$.
We record a uniform estimate, for $t\ge2$:
\begin{equation}\label{p12:fg:mgfgrowth}
                  \log M_A(t)\le2t\log t+C t.
\end{equation}
For $p\le24t$, bound the logarithm of each factor by $tb_p$ and use
the classical elementary estimate
$\sum_{p\le x}\log p/(p-1)=\log x+O(1)$.
For $p>24t$, one has $tb_p\le(\log p)/6$; consequently
\[
 \log(1+a_p(e^{tb_p}-1))
 \le a_p tb_p e^{tb_p}
 \le C t\frac{\log p}{p^{7/6}}.
\]
Its tail sum is $O(t^{5/6}\log(2t))=O(t)$.
For completeness, the elementary prime-sum estimate follows from
$\log(m!)=\sum_{p\le m}v_p(m!)\log p$ and
$v_p(m!)=m/(p-1)+O(\log m/\log p)$, together with the elementary
Chebyshev bound $\pi(m)=O(m/\log m)$ and Stirling's estimate.

Weak convergence gives, for the closed set $[L,\infty)$, an upper
limit bounded by $\mathbb P(W_A\ge L)$.  Markov's inequality and
\eqref{p12:fg:mgfgrowth}, optimized at
$t=\exp((L-C-2)/2)$, bound this probability by
$\exp(-c e^{L/2})$ after enlarging $L_0$.
\end{proof}

The order of limits in \eqref{p12:fg:tail} matters: $L$ is fixed before
$N$ tends to infinity.  Substituting $L=\log N$ into this statement
is not justified.  In particular, the law does not eliminate the
sparse exceptional indices left by the earlier $O(\log N)$ theorem.
Unlike the preceding parameter law, however, all conclusions here
hold with $A=1$ fixed.

\subsection{Sublinear content with a quantitative exceptional set}

The support law gives more information about the distribution of prime
divisibility than about the content itself.  The latter admits a
stronger estimate by varying the number of indices in a continuant
packing argument.  This refinement also holds for every fixed $A$.

\begin{lemma}[Positive continuant packing]\label{p12:fg:packing}
If $1\le n_1<\cdots<n_h\le M$, then
\begin{equation}\label{p12:fg:packingbound}
 \sum_{i=1}^h\log g_{n_i}
 \le 2\log((M-1)!)+
       \log((2A+1)M)\sum_{i<j}(n_j-n_i-1).
\end{equation}
\end{lemma}
\begin{proof}
For $n<m$, the polynomial value $C_{m-n}(n)$ used in the proof of
Theorem~\ref{p12:fg:roots} is the determinant of a symmetric tridiagonal
matrix indexed by $k=n+2,\ldots,m$. Its diagonal entries are
$1+(2k-1)A$, with off-diagonal entry $-kA$ between $k$ and $k+1$.
The empty determinant is one. Expansion at the last row verifies
the defining recurrence. In an interior row the diagonal exceeds
the sum of the absolute off-diagonal entries by one; in boundary
rows it exceeds that sum by more. The matrix is therefore positive
definite. Hadamard's determinant inequality gives
\[
 1\le C_{m-n}(n)\le((2A+1)m)^{m-n-1}.
\]
The recurrence congruence and adjacent coprimality imply
$\gcd(Q_n,Q_m)\mid C_{m-n}(n)$.

For positive integers $a_1,\ldots,a_h$ one has
\[
 \prod_i a_i\mid\operatorname{lcm}(a_1,\ldots,a_h)
                         \prod_{i<j}\gcd(a_i,a_j).
\]
Indeed, ordering their valuations at a prime as
$v_1\le\cdots\le v_h$, the right side has valuation
$v_h+\sum_{i<h}(h-i)v_i\ge\sum_i v_i$.
Apply this to $a_i=g_{n_i}$. Their least common multiple divides
$((M-1)!)^2$, and each pairwise gcd divides the corresponding
positive continuant. Taking logarithms proves the claim.
\end{proof}

\begin{theorem}[Quantitative negligible content]\label{p12:fg:sublinear}
For each fixed positive integer $A$,
\begin{equation}\label{p12:fg:meanrate}
 \frac1X\sum_{n\le X}\frac{\log g_n(A)}n
                          =O_A(X^{-1/3}\log(2X)).
\end{equation}
For every fixed $\varepsilon>0$,
\begin{equation}\label{p12:fg:linearcount}
 \#\{n\le X:\log g_n(A)\ge\varepsilon n\}
                         =O_{A,\varepsilon}((\log(2X))^4).
\end{equation}
Consequently, outside a set of natural density zero,
\begin{align}
 \log(Q_n/g_n)&=n\log n+(\log A-1)n+o(n),\label{p12:fg:qsecond}\\
 \log\ell_n&=n\log n+(\log A-1)n+o(n).
                                                        \label{p12:fg:lsecond}
\end{align}
These are density statements; exceptional infinite subsequences
remain possible.
\end{theorem}
\begin{proof}
Work first in $[N,2N]$, with $N\ge2$.  An interval of span at most
$H$ containing $h$ selected indices satisfies, by
Lemma~\ref{p12:fg:packing},
\begin{equation}\label{p12:fg:localpack}
       \sum_i\log g_{n_i}\le C_A(N+h^2H)\log(2N)
\end{equation}
for a fixed constant $C_A\ge1$.
Partition $[N,2N]$ into intervals containing at most
$\lceil N^{1/3}\rceil$ consecutive integers and apply this inequality
to all indices in each interval. Each sum is $O_A(N\log(2N))$,
and there are $O(N^{2/3})$ intervals. Hence
\[
 \sum_{N\le n\le2N}\log g_n=O_A(N^{5/3}\log(2N)),\qquad
 \sum_{N\le n\le2N}\frac{\log g_n}n
                              =O_A(N^{2/3}\log(2N)).
\]
Dyadic summation proves \eqref{p12:fg:meanrate}.

For the sharper exceptional count, put $L=\log(2N)$,
$h=\lceil4C_A L/\varepsilon\rceil$, and
\[
                   H=\left\lfloor
                       \frac{\varepsilon N}{4C_AhL}
                     \right\rfloor.
\]
For sufficiently large $N$ these are positive integers, and
$H\asymp_{A,\varepsilon}N/L^2$.
If an interval of span at most $H$ contained $h$ exceptional indices,
their logarithmic contents would sum to at least $h\varepsilon N$.
But \eqref{p12:fg:localpack} is at most
\[
 C_A NL+C_Ah^2HL\le\tfrac14h\varepsilon N
                              +\tfrac14h\varepsilon N,
\]
a contradiction. Covering $[N,2N]$ by $O_{A,\varepsilon}(L^2)$
such intervals gives $O_{A,\varepsilon}(L^3)$ exceptional indices.
Summing over dyadic intervals proves \eqref{p12:fg:linearcount}.

For clarity, a single density-one set can support the stated $o(n)$.
For each integer $j\ge1$, let
$E_j=\{n:\log g_n\ge n/j\}$; each has density zero by
\eqref{p12:fg:linearcount}. Choose increasing cutoffs $T_j$ so large
that $|E_j\cap[1,X]|\le X/j$ for all $X\ge T_j$.
They may also be chosen so that $T_j\ge jT_{j-1}$.
The union of $E_j\cap[T_j,T_{j+1})$ has density zero: for
$T_j\le X<T_{j+1}$, its part below $T_j$ is at most
$T_{j-1}+T_j/(j-1)$, and its current part is at most $X/j$.
Outside this union $\log g_n=o(n)$.

Finally, the binomial formula and the exact remainder give
\[
 \log Q_n=\log(n!)+n\log A+O_A(\sqrt n),\qquad
 \log(Q_n\delta_A-P_n)
           =\log(n!)+n\log A+O_A(\sqrt n+\log n).
\]
For the first estimate use
$1\le Q_n/(n!A^n)\le\sum_{j\ge0}(n/A)^j/(j!)^2
\le e^{2\sqrt{n/A}}$; the second was proved after
\eqref{p12:fg:identities}. Subtracting $\log g_n=o(n)$ and applying
Stirling proves \eqref{p12:fg:qsecond}--\eqref{p12:fg:lsecond}.
\end{proof}

The new threshold in \eqref{p12:fg:linearcount} is linear in $n$.
The earlier logarithmic exceptional count concerned the larger
threshold $\varepsilon n\log n$.  These different thresholds must
not be conflated.  In particular, even the present stronger result
does not exclude all indices capable of giving an irrationality
certificate for $\delta_A$.

\subsection{A small primitive form requires many growing primes}

For $Y\ge2$, write
$g_{n,\le Y}=\prod_{p\le Y}p^{v_p(g_n)}$ and
$g_{n,>Y}=g_n/g_{n,\le Y}$.

\begin{theorem}[Necessary prime complexity]\label{p12:fg:complexity}
For every $Y\ge2$ and $n\ge2$,
\begin{equation}\label{p12:fg:budget}
                 \log g_{n,\le Y}\le2n\log Y+O(n),
\end{equation}
with an absolute implied constant.
Fix $B>0$ and $0<\theta<1/2$.  At every sufficiently large index
with $\ell_n\le B$,
\begin{align}
 \log g_{n,>n^\theta}
       &\ge(1-2\theta)n\log n-O_{A,B}(n),\label{p12:fg:largecontent}\\
 \#\{p>n^\theta:p\mid g_n\}
       &\ge\left(\frac{1-2\theta}{2}+o_{A,B,\theta}(1)\right)n^\theta.
                                                       \label{p12:fg:manyprimes}
\end{align}
More sharply, there are positive constants $c_{A,B},c'_{A,B}$ such
that any such sufficiently large index satisfies
\begin{equation}\label{p12:fg:sqrtprimes}
 \#\{p>c_{A,B}\sqrt n:p\mid g_n\}
                                  \ge c'_{A,B}\frac{\sqrt n}{\log n}.
\end{equation}
\end{theorem}
\begin{proof}
By \eqref{p12:fg:base} and the prime-sum estimate used above,
\[
 \log g_{n,\le Y}\le2n\sum_{p\le Y}\frac{\log p}{p-1}
                         \le2n\log Y+O(n).
\]
If $\ell_n\le B$, then \eqref{p12:fg:base} gives
$\log g_n\ge n\log n-O_{A,B}(n)$.
Subtracting \eqref{p12:fg:budget} at $Y=n^\theta$ proves
\eqref{p12:fg:largecontent}.  Every contributing prime is at most $n-1$,
and, when $p>n^\theta$, its contribution is at most
\[
 v_p(g_n)\log p\le\frac{2n\log n}{n^\theta-1}
                       =(2+o(1))n^{1-\theta}\log n.
\]
Division proves \eqref{p12:fg:manyprimes}.

Finally set $Y=c\sqrt n$.  The two preceding estimates give
$\log g_{n,>Y}\ge(-2\log c-C_{A,B})n$ for a fixed constant
$C_{A,B}$.  Choose $c>0$ small enough to make the parenthesis at
least one.  Each prime exceeding $Y$ contributes at most
$2n\log n/(Y-1)=O_{A,B}(\sqrt n\log n)$, so at least
$c'_{A,B}\sqrt n/\log n$ such primes are needed.
\end{proof}

Thus finitely many exceptional primes, no matter how their higher
valuations behave, cannot produce bounded primitive forms indefinitely.
The remaining arithmetic difficulty is simultaneous cancellation at
a growing collection of primes.  A bound on the number of roots
modulo each prime controls averages, but the Chinese remainder theorem
alone does not prevent one exceptionally small positive index from
meeting many of those congruences.  Neither the Kurepa conjecture nor
the irrationality of $\delta$ follows.  In particular, the known
congruence $Q_{p-1}(1)\equiv\sum_{j=0}^{p-1}j!\pmod p$ is a boundary
residue identity, not an equivalence between all the index-root
conditions used here and Kurepa's conjecture.


% END EMBEDDED gompertz_results.tex


% BEGIN EMBEDDED zeta_results.tex

\section{Polynomial period elimination creates a uniform Catalan barrier}
\label{p12:zc:section}

The previous paired interpolation construction removed $\pi^2$ from a
positive Catalan pencil. A tempting improvement is to replace the
changing interpolation space by powers of a single polynomial invariant,
thereby obtaining genuine Hankel matrices of arbitrarily large size while
the determinant degree remains fixed. We construct this improvement,
then prove that the restricted invariant space cannot produce arbitrarily
small primitive integer polynomial values. The obstruction is quantitative:
in the simplest case every determinant root stays an explicit positive
distance below Catalan's constant, independently of matrix size. We then
identify the omitted polynomial sector and prove that the full paired
spaces restore convergence at every fixed determinant degree. The
remaining primitive-height estimate is not supplied.

The signed density with all moments zero used below is classical in
Stieltjes moment theory; see Kleiber \cite[Section 4]{p12:zc:Kleiber} and
L\'opez-Garc\'ia \cite[Section 2]{p12:zc:Lopez}. Our application concerns the
specific period-eliminating pencil and its uniform arithmetic failure.
No historical priority claim is made for the general moment mechanism.

\subsection{A moment-null density with an exact negative-axis transform}

\begin{lemma}\label{p12:zc:null}
Let $p\ge4$ be an even integer, $c>0$, $\alpha=\pi/p$, and
\[
 g_c(u)=e^{-cu^{1/p}}\sin(c\tan\alpha\,u^{1/p})\qquad(u>0).
\]
All its nonnegative integer moments vanish. Moreover, for $s<0$,
\begin{equation}\label{p12:zc:transform}
 \int_0^\infty\frac{g_c(u)}{u-s}\,du
       =\pi\exp\{-c\sec\alpha\,(-s)^{1/p}\}>0.
\end{equation}
\end{lemma}
\begin{proof}
Put $z=c(1-i\tan\alpha)$. Substitution $u=y^p$ gives
\[
 \int_0^\infty u^k g_c(u)\,du
 =p\,\operatorname{Im}\frac{\Gamma(p(k+1))}{z^{p(k+1)}}=0,
\]
because $\arg z=-\pi/p$. All integrals are absolutely convergent.
For the transform, write $-s=v^p$ and integrate
\[
                  \frac{p y^{p-1}e^{-zy}}{y^p+v^p}
\]
over the sector between arguments zero and $\alpha$. The exponential
decays uniformly on its large circular arcs. The only pole on the
rotated boundary is $y=v e^{i\alpha}$, with positive real residue
$e^{-c\sec\alpha v}$. Rotate the boundary with a small indentation
at that pole. The principal-value integral on the ray is real:
its integrand, parametrized by $y=r e^{i\alpha}$, is
\[
                   \frac{p r^{p-1}e^{-c\sec\alpha r}}{r^p-v^p}.
\]
The half-residue contribution to the imaginary part of the original
integral is $\pi e^{-c\sec\alpha v}$. This proves
\eqref{p12:zc:transform}. Equivalently one may rotate to angles below
$\alpha$ and use $(x-i0)^{-1}=\operatorname{PV}(1/x)+i\pi\delta_0$;
the same sign and coefficient result.
\end{proof}

\begin{theorem}[Bounded evaluations and a fixed-degree gap]
\label{p12:zc:gap}
Let $w(u)>0$ be continuous on $(0,\infty)$, with
\[
 w(u)\asymp u^{-1/2}\quad(u\downarrow0),\qquad
 \log w(u)=-a u^{1/p}+O(\log u)\quad(u\to\infty),
\]
where $a>0$ and $p\ge4$ is even. Fix distinct negative points
$s_1,\ldots,s_k$ and positive weights $\omega_j$. Write
\[
 H_n=\left(\int_0^\infty u^{i+j}w(u)\,du\right)_{0\le i,j<n},
 \quad B_n=\left(\sum_{j=1}^k\omega_j s_j^{i+\ell}
                                    \right)_{0\le i,\ell<n}.
\]
For $c>a/2$, set
\[
 M_j=\frac{e^{2c\sec(\pi/p)(-s_j)^{1/p}}}{\pi^2}
       \int_0^\infty\frac{g_c(u)^2}{(u-s_j)^2w(u)}\,du,
 \qquad K=\sum_{j=1}^k\omega_jM_j.
\]
These constants are finite and positive, and $B_n\le K H_n$ for
every $n$. Consequently the polynomial
\[
                  Q_n(X)=\det(H_n+(X-\theta)B_n)
\]
has degree $k$ for $n\ge k$, all its zeros are real, and each zero
is at most $\theta-K^{-1}$. If $Q_n$ is a nonzero scalar multiple
of an integer polynomial $P_n$ of degree $k$, then
\begin{equation}\label{p12:zc:integergap}
                         |P_n(\theta)|\ge K^{-k}.
\end{equation}
This remains true for the primitive integer normalization.
\end{theorem}
\begin{proof}
Near zero, $g_c(u)=O(u^{1/p})$, so the integrand defining $M_j$
is $O(u^{2/p+1/2})$. At infinity it is bounded by a polynomial
factor times $e^{-(2c-a)u^{1/p}}$. Thus $M_j$ is finite.
For any polynomial $f$, the vanishing moments give
\[
 \int_0^\infty\frac{f(u)-f(s_j)}{u-s_j}g_c(u)\,du=0.
\]
Use \eqref{p12:zc:transform} and Cauchy--Schwarz to obtain
\[
                    |f(s_j)|^2\le M_j\int_0^\infty|f(u)|^2w(u)\,du.
\]
Summing with weights proves the matrix inequality. The positive
matrix $H_n^{-1/2}B_nH_n^{-1/2}$ has rank $k$ for $n\ge k$,
by the Vandermonde rank of the distinct $s_j$, and its nonzero
eigenvalues lie in $(0,K]$. The determinant factors accordingly,
with roots $\theta-1/\lambda_j\le\theta-1/K$.
Finally, the leading coefficient of a nonzero integer polynomial
has absolute value at least one. Multiplying its $k$ distances
from $\theta$ proves \eqref{p12:zc:integergap}.
\end{proof}

For $c>a$, the quotient $g_c/w$ is bounded, and is not identically
zero. Thus $w+\varepsilon g_c$ is positive for all sufficiently small
$|\varepsilon|$ and has exactly the same moments as $w$. This gives
an explicit proof of Stieltjes indeterminacy. The bounded-evaluation
argument is stronger for the present purpose because it provides a
uniform gap and requires only $c>a/2$.

\subsection{A genuine Hankel family for Catalan}

Use the classical positive Hurwitz moment measure
\cite{p12:zc:Prevost}
\[
 d\nu_2(t)=4\pi y^2\sum_{\ell\ge1}\ell e^{-2\pi\ell y}\,dy,
 \quad t=y^2,
\]
for which
\[
 \int t^e\,d\nu_2(t)=(-1)^eB_{2e+2},\qquad
 \int\frac{d\nu_2(t)}{t+c^2}
       =c\zeta(2,c)-1-\frac1{2c}.
\]
Let $G=\sum_{n\ge0}(-1)^n/(2n+1)^2$. Suppose that rational pairs
$a_j<b_j$, $1\le j\le k$, satisfy
\[
 a_j\in\tfrac14+\mathbb Z_{\ge0},\quad
 b_j\in\tfrac34+\mathbb Z_{\ge0},\quad a_j^2+b_j^2=C,
\]
and are distinct. Put
\[
 U(t)=t(t+C),\quad s_j=-a_j^2b_j^2<0,\quad
 W(t)=\sum_{j=1}^k\left(\frac1{a_j(t+a_j^2)}-
                                 \frac1{b_j(t+b_j^2)}\right).
\]
The $s_j$ are distinct, because their sum-of-squares and product
would otherwise determine the same unordered positive pair.
Each summand of $W$ is strictly positive on $[0,\infty)$.

\begin{theorem}[Full degree, followed by saturation]\label{p12:zc:catalan}
There are rational numbers $r_\ell$, computable by polynomial division
and Bernoulli numbers, such that for every $\ell\ge0$,
\[
              \int U(t)^\ell W(t)\,d\nu_2(t)
                    =r_\ell+16G\sum_{j=1}^k s_j^\ell.
\]
Thus the rational Hankel pencil
\[
 T_n(X)=\left(r_{i+\ell}+16X\sum_{j=1}^k s_j^{i+\ell}
                                             \right)_{0\le i,\ell<n}
\]
is positive definite at $X=G$, and $\det T_n$ has degree exactly
$k$ whenever $n\ge k$. For every fixed pole set there is an
$\varepsilon>0$, independent of $n$, such that every zero of
$\det T_n$ is at most $G-\varepsilon$. For its primitive integer
normalization $P_n$,
\[
                         |P_n(G)|\ge\varepsilon^k>0.
\]
In particular increasing the matrix size at a fixed pole set cannot
supply an irrationality or transcendence criterion based on values
tending to zero.
\end{theorem}
\begin{proof}
One has $U(-a_j^2)=U(-b_j^2)=s_j$. Polynomial division therefore
gives the same remainder $s_j^\ell$ at the paired poles. The two
Hurwitz transforms, divided by $a_j$ and $b_j$, have $\pi^2$
coefficients $1$ and $1$, and $G$ coefficients $8$ and $-8$.
Their difference has coefficient $16s_j^\ell$ in $G$ and zero
coefficient in $\pi^2$; all remaining terms are rational.

Push the positive measure $W(t)d\nu_2(t)$ forward by $u=U(t)$,
and call its density $w(u)$. Since $U'(t)=2t+C>0$, this density
is positive and continuous. At zero, $W(t)$ tends to a positive
constant and the $t$-density of $\nu_2$ is asymptotic to
$(2\pi\sqrt t)^{-1}$, so $w(u)\asymp u^{-1/2}$.
At infinity, $W(t)\asymp t^{-1}$ and
\[
                  w(u)\asymp u^{-3/4}e^{-2\pi u^{1/4}}.
\]
Theorem~\ref{p12:zc:gap} applies with $p=4$, $a=2\pi$,
$\theta=G$, and weights $\omega_j=16$. All assertions follow.
\end{proof}

\begin{proposition}[Arbitrarily many poles on a common circle]
\label{p12:zc:circles}
For every $k$ there is a pole set as above with at least $k$ pairs.
It can be chosen with $\log C=O((\log(k+1))^2)$.
\end{proposition}
\begin{proof}
For an integer $h\ge0$, let $N_j=10^{j+2}+1$ and form the
$2^h$ Gaussian integers
\[
             (1+3i)\prod_{j=1}^h(N_j+4\epsilon_j i),
                         \qquad\epsilon_j\in\{-1,1\}.
\]
Their coordinates $P,Q$ satisfy $P\equiv1$, $Q\equiv3\pmod4$.
All their arguments lie in $(\pi/4,\pi/2)$: the total perturbation
from $\arg(1+3i)$ is at most
$\sum_j4/10^{j+2}<1/100$.
They are distinct. Indeed, if $\theta_j=\arctan(4/N_j)$, then
$\theta_j>2/10^{j+2}>\sum_{l>j}\theta_l$, so a first differing
sign cannot be canceled by all later terms. Thus $0<P<Q$ and
$a=P/4$, $b=Q/4$ give $2^h$ distinct permissible pairs, all with
\[
                      C=\frac{10}{16}\prod_{j=1}^h(N_j^2+16).
\]
Taking $h=\lceil\log_2k\rceil$ proves both assertions.
\end{proof}

\subsection{An explicit gap for the smallest Catalan invariant}

\begin{theorem}\label{p12:zc:explicit}
Take $a=1/4$, $b=3/4$, $C=5/8$, and $U(t)=t(t+5/8)$.
For every $n\ge1$ the determinant polynomial in
Theorem~\ref{p12:zc:catalan} is linear. If its root is $x_n$, then
\begin{equation}\label{p12:zc:explicitgap}
 x_n\le G-\varepsilon_0,\qquad
 \varepsilon_0=\frac{39\pi^3(4-\pi)}{1600}e^{-2\sqrt6}>0.
\end{equation}
Its primitive integer normalization satisfies
$|P_n(G)|\ge\varepsilon_0$ for every $n$.
\end{theorem}
\begin{proof}
Here $s=-9/256$, and
\[
 W(t)=\frac{8}{3}\frac{t+13/16}{(t+1/16)(t+9/16)},
 \quad U(t)-s=(t+1/16)(t+9/16).
\]
Use Lemma~\ref{p12:zc:null} with $p=4$ and any $c>\pi$.
Write $E=\int_0^\infty g_c(u)^2/((u-s)^2w(u))\,du$.
On substituting $t=y^2$ and then $u=U(t)$, the geometric sum in
the measure gives exactly
\[
 E=\frac3{2\pi}\int_0^\infty
 \frac{e^{-2cU(y^2)^{1/4}}\sin^2(cU(y^2)^{1/4})
       (2y^2+5/8)^2\sinh^2(\pi y)}
      {(y^2+1/16)(y^2+9/16)(y^2+13/16)}\,dy.
\]
Since $U(y^2)^{1/4}\ge y$, $\sin^2\le1$, and
$\sinh^2(\pi y)\le e^{2\pi y}/4$, while
\[
 \frac{(2y^2+5/8)^2}
 {(y^2+1/16)(y^2+9/16)(y^2+13/16)}
 \le\frac{100}{9}\frac{16}{13}=\frac{1600}{117},
\]
we obtain $E\le100/(39\pi(c-\pi))$. The rational inequality
uses $(x+z)^2/(xz)=4+(x-z)^2/(xz)$ with
$x=y^2+1/16$ and $z=y^2+9/16$.
Also $(-s)^{1/4}=\sqrt3/4$, so
\[
 M_s\le\frac{100e^{c\sqrt6/2}}{39\pi^3(c-\pi)}.
\]
The gap in Theorem~\ref{p12:zc:gap}, with weight $16$, is at least
$1/(16M_s)$. Taking $c=4$ proves \eqref{p12:zc:explicitgap} and the
primitive integer bound.
\end{proof}

Exact rational computations for $n=1,2,5,10$ give respectively
\[
 x_n\approx0.6111111111,\quad0.6636009857,\quad
              0.6837781755,\quad0.6892645428.
\]
These are finite diagnostics, not the proof of saturation. The
uniform conclusion is \eqref{p12:zc:explicitgap}. It does not apply to
the previous moving-pole interpolation family, where the pole set
changes with the matrix dimension.

\subsection{The omitted sector repairs convergence, but not arithmetic height}

The saturation theorem concerns polynomials in $U(t)$ only. The full
space of paired polynomials is larger. For the common-circle pole set,
write $D(u)=\prod_{j=1}^k(u-s_j)$. Polynomial division by the monic
quadratic $t^2+Ct-u$ gives a unique decomposition
\[
                 f(t)=h(U(t))+(2t+C)r(U(t)).
\]
The involution $t\mapsto-C-t$ fixes the first summand and negates the
second. Consequently the paired constraints are equivalent to
$r(s_j)=0$ for every $j$, and hence
\begin{equation}\label{p12:zc:module}
 \{f:f(-a_j^2)=f(-b_j^2)\text{ for all }j\}
 =\mathbb Q[U]\ \oplus\ (2t+C)D(U)\mathbb Q[U].
\end{equation}
The first omitted polynomial has degree $2k+1$. Thus changing a basis
inside $\mathbb Q[U]$ cannot repair the saturation, whereas adjoining
the second summand changes the approximation problem itself.

\begin{lemma}[No bounded negative-point evaluation combination]
\label{p12:zc:unbounded}
For the measure $d\mu(t)=W(t)d\nu_2(t)$ above, no nonzero finite
linear combination of evaluations at distinct negative points is
bounded on the polynomial space with its $L^2(\mu)$ norm.
\end{lemma}
\begin{proof}
Suppose $\sum_j c_j f(-v_j^2)$ were bounded, with distinct $v_j>0$.
By extension to the Hilbert-space closure and the Riesz representation
theorem, there would be $h\in L^2(\mu)$ such that
\[
               \int f(t)h(t)\,d\mu(t)=\sum_jc_jf(-v_j^2)
\]
for every polynomial $f$. Complex coefficients, if present, can be
handled with the corresponding conjugate convention for $h$.
Symmetrize the measure under $t=y^2$ to an even measure $d\sigma(y)$
on $\mathbb R$. Its density has exponential tail $e^{-2\pi|y|}$
times a rational factor. Cauchy--Schwarz therefore makes
\[
                F(z)=\int_{\mathbb R} e^{izy}h(y^2)\,d\sigma(y)
\]
holomorphic on $|\operatorname{Im}z|<\pi$. Its Taylor coefficients
at zero are determined by the displayed moment identities and
symmetry, giving $F(z)=\sum_jc_j\cosh(v_jz)$ near zero and hence
throughout the strip. But $h(y^2)d\sigma(y)$ has finite total
variation, so $F$ is bounded on the real axis. A nonzero sum of
$\cosh(v_jz)$ with distinct positive $v_j$ is unbounded there,
as is seen from its largest exponent. This is a contradiction.
\end{proof}

\begin{theorem}[Full paired spaces recover every fixed determinant degree]
\label{p12:zc:repair}
Fix any $k$ distinct permissible quarter-integer pairs $a_j<b_j$;
the common-circle hypothesis is not required. Let $W$ be their
positive kernel sum and let
\[
 V_N=\{f\in\mathbb Q[t]:\deg f\le N,\
                    f(-a_j^2)=f(-b_j^2)\text{ for every }j\},
\]
assuming all $2k$ poles are distinct. For $N\ge2k-1$, choose any
rational basis of $V_N$ and form its Gram matrix in $L^2(W\nu_2)$.
It has the form $A_N+16G E_N^TE_N$, with $A_N$ rational, where
$E_N$ evaluates the basis at $-a_1^2,\ldots,-a_k^2$.
Then $\det(A_N+16X E_N^TE_N)$ has degree exactly $k$.
All its roots approach $G$ from below as $N\to\infty$.
After division by the leading coefficient these polynomials
converge coefficientwise to $(X-G)^k$.
\end{theorem}
\begin{proof}
The paired evaluation cancels the $\pi^2$ term and leaves coefficient
$16E_N^TE_N$ exactly as before. Paired Lagrange interpolation shows
that $E_N$ has rank $k$ once $N\ge2k-1$.
Let $H_N=A_N+16G E_N^TE_N>0$ and
$K_N=E_NH_N^{-1}E_N^T$. The matrices $K_N$ increase in the positive
semidefinite order as the polynomial spaces increase; their quadratic
forms are the squared norms of linear combinations of the evaluations.

We claim that the smallest eigenvalue of $K_N$ tends to infinity.
Otherwise take unit vectors $c_N$ along a subsequence with
$c_N^TK_Nc_N$ bounded, and pass to a limit $c\ne0$. Monotonicity
shows that $c^TK_Mc$ is bounded uniformly in $M$. Thus the functional
$L(f)=\sum_jc_jf(-a_j^2)$ is bounded on the union of the paired spaces.
Extend it to a bounded functional on $L^2(\mu)$. On all polynomials,
the difference from $\sum_jc_jf(-a_j^2)$ annihilates the kernel of
the surjective constraint map
\[
                  f\longmapsto(f(-a_j^2)-f(-b_j^2))_{j=1}^k.
\]
It is therefore a linear combination of these constraints. The
bounded extension would then be a nonzero linear combination of
negative-point evaluations on all polynomials: its coefficients at
the two members of each pair cannot all vanish unless $c=0$.
This contradicts Lemma~\ref{p12:zc:unbounded}.

The matrix determinant lemma now gives
\[
 \frac{\det(A_N+16X E_N^TE_N)}{\det H_N}
                         =\det(I_k+16(X-G)K_N).
\]
Its roots are $G-1/(16\lambda_j(K_N))$. All $k$ eigenvalues tend
to infinity, proving the root convergence and the final assertion.
\end{proof}

For the single pair $(1/4,3/4)$, the full spaces give roots
approximately $0.7047630881$, $0.8206648111$, and $0.8337339123$
at polynomial degree $3$, $10$, and $12$. The corresponding primitive
coefficient heights have logarithms approximately $26.98$, $136.72$,
and $177.88$. These finite values illustrate both the repaired
convergence and the unresolved arithmetic cost. The proof supplies
no rate relating primitive height to the small normalized value;
it therefore supplies no irrationality or transcendence certificate.

\subsection{What this does and does not transfer to other constants}

A rational positive reweighting of a Hurwitz moment measure retains
an exponential tail in $\sqrt t$. If a polynomial invariant $U$
of degree $d\ge2$, strictly increasing on $[0,\infty)$ with
$U(0)=0$ and $U'(0)>0$, replaces $t$, the pushed density instead
has exponential tail in $u^{1/(2d)}$. The same construction applies
with $p=2d$. Whenever the resulting coefficient of a target constant
is a fixed positive sum of evaluations at negative $u$-points,
Theorem~\ref{p12:zc:gap} prevents small primitive values at fixed rank.
This is a stated-condition obstruction, not a claim that all possible
zeta constructions have that form.

For $G$ transcendence we have proved a complete uniform obstruction
for the invariant Hankel route and an explicit primitive lower bound.
The full paired spaces repair analytic convergence at every fixed
determinant degree, but the primitive-height estimate remains open. For $\zeta(3)$ and $\zeta(5)$ transcendence, no higher
algebraic degree is excluded; the required off-axis or content bound
remains unproved. For $\zeta(7)$, $\zeta(9)$, and $\zeta(11)$
irrationality, no sufficient individual primitive-value estimate is
obtained. The entropy rigidity theorem cannot be applied to an
unknown constant merely because rational determinant polynomials
have algebraic roots: those roots are approximants, not conjugates
of the target. The number of newly resolved named-constant problems
in this section is zero.


% END EMBEDDED zeta_results.tex


% BEGIN EMBEDDED lambert_results.tex

\section{Lambert values: a sharp optimization limit for a published determinant budget}
\label{p12:lf:section}

The preceding manuscript established functional independence of three
Lambert--Fibonacci jets, with a fourth function supplied by prime sums.
That result did not specialize to an arithmetic value. Here we examine an
actual arithmetic determinant construction and optimize its shape. The
outcome is a precise limitation: the displayed local bounds cannot exclude
cubic values, even after optimizing the common shift. Separating the two
shifts also gives no improvement in the central parameter region.
This is a result about a particular certificate budget, not a theorem that
Lambert values can be cubic.

\subsection{The external arithmetic input and its status}
Write
\[
 F(q)=\sum_{n\ge1}\frac1{q^n-1}=L(1/q),\qquad q>1.
\]
Lambrinoudis's October 2026 manuscript \cite{p12:lf:Lambrinoudis} announces
that $1,F(q),F(q)^2$ are linearly independent over $\mathbb Q$ for
integer $q\ge2$. It supplies a general degree criterion with
\[
 \lambda=\frac{42963+428220/\pi^2}{202140}
        =0.4271829289\ldots,
 \qquad e\lambda\log M(q)<\log q.
\]
The source includes a claimed complete Lean formalization. We read the
actual mathematical source, including its local denominator argument;
we have not rebuilt that formalization or independently certified the
entire external theorem. The announcement is not counted as our result.
Its six-presentation refinement is stronger than the single-presentation
shape optimization below.

We retain the source's notation for its positive determinant. For pole
and moment shift $k$, matrix size $h$, the normalized analytic decay has
scale
\[
 \mathcal T_{h,k}=G_{h,k}
 +h\sum_{j=k}^{k+h-1}j
 +\frac{h(h+1)(2h+1)}6+\frac{kh(h+1)}2,
 \qquad
 G_{h,k}=k\left(hk+\frac{h(h-1)}2\right).
\]
The denominator budget comprises a pole at zero and cyclotomic poles.
All statements below optimize these explicit bounds; they do not assert
that the bounds equal the actual denominator orders.

\subsection{The complete common-shift budget}
Put $k/h\longrightarrow x\ge0$, and normalize all cubic quantities by
$h^3$. Define
\[
 S(x)=x^2+2x+\frac56,
 \qquad
 B(x)=\frac18+\frac{7/4+(9/2)x+3x^2}{\pi^2},
\]
and
\[
 A(x)=
 \begin{cases}
 \displaystyle\frac1{18}+\frac x3+\frac{2x^2}3+\frac{x^3}9,
       &0\le x\le1,\\[1ex]
 \displaystyle x^2+\frac16,&x\ge1.
 \end{cases}
 \qquad C(x)=A(x)+B(x),\quad \Lambda(x)=\frac{C(x)}{S(x)}.
\]

\begin{proposition}[Evaluation of the common-shift budget]
\label{p12:lf:budget}
The analytic decay scale is $S(x)$. The zero-place and cyclotomic
clearing estimates used in the source give the sufficient budgets
$A(x)$ and $B(x)$, respectively. In particular $x=1/20$ recovers the
source's displayed $h=60N,k=3N$ constants.
\end{proposition}
\begin{proof}
The formula for $S$ follows directly from $\mathcal T_{h,k}$.
For the zero place let $r\in[0,1]$ be the normalized column index and
$j\ge0$ the normalized selected node. The source's slotwise lower bound
has continuous cost
\[
 w_x(j,r)=(x+2-2r)j-\int_x^{x+1}\min(j,v)\,dv.
\]
The relaxed capacity bound is $j\ge r$ if $r<x$, and
$j\ge(x+r)/2$ otherwise. For $0\le x\le1$, the minimizing nodes are
\[
 j(r)=
 \begin{cases}
 r,&0\le r<x,\\
 (x+r)/2,&x\le r\le(x+2)/3,\\
 2r-1,&(x+2)/3\le r\le1.
 \end{cases}
\]
Indeed, before the pole window the cost is affine; inside the window
it is $j^2/2+(1-2r)j+x^2/2$. Its unconstrained minimizer is $2r-1$.
The expression after the window is nondecreasing. Direct integration
of these three pieces gives
\[
 x(x+1/2)-\int_0^1 w_x(j(r),r)\,dr
 =\frac1{18}+\frac x3+\frac{2x^2}3+\frac{x^3}9.
\]
For $x\ge1$ all capacity bounds are $j\ge r$, the minimum is $j=r$,
and the integral is $x/2-1/6$, giving the second expression for $A$.
Floors and endpoint rounding change the discrete sum by $O(h^2)$,
which does not affect its cubic limit.

For completeness, the cyclotomic envelope can be written using
$(u)_+=\max(u,0)$ as
\[
 \begin{split}
 E_h(n)&=3(h-n)_+-(h-2n)_+
             +\sum_{j\ge1}(h-(2j+1)n)_+,\\
 c_{h,k}(n)&=E_h(n)+2(h+k-n)_+-2(k+h/2-n)_+\\
           &\hspace{25mm}-2(h-n)_++2(h/2-n)_+.
 \end{split}
\]
The source's filtration bound is used when $n\le h/2$; its residue
rank bound gives this envelope for $n>h/2$. The totient asymptotic
and partial summation give
\[
 h^{-3}\sum_n\varphi(n)(Kh-mn)_+
        \longrightarrow\frac{K^3}{\pi^2m^2}.
\]
Applying it to the displayed tents, and using
$\sum_{j\ge1}(2j+1)^{-2}=\pi^2/8-1$, gives
\[
 B(x)=\frac18+
 \frac{7/4+2((1+x)^3-(1/2+x)^3-1+1/8)}{\pi^2},
\]
which simplifies to the stated formula.
\end{proof}

\begin{theorem}[The common-shift optimization and its cubic barrier]
\label{p12:lf:optimization}
The function $\Lambda$ has a unique global minimum. Its minimizer
$x_*$ is the unique positive root of
\[
 \frac{x^4}{9}+\frac{4x^3}{9}
 +\left(\frac{23}{18}+\frac3{2\pi^2}\right)x^2
 +\left(\frac34+\frac3{2\pi^2}\right)x
 +\frac1{4\pi^2}-\frac1{12}=0.
\]
Numerically,
\[
 x_*\simeq0.05873628,
 \qquad \Lambda(x_*)\simeq0.4271370864.
\]
Moreover the elementary uniform bound
\[
                         \boxed{\Lambda(x)>\frac25\quad(x\ge0)}
\]
holds. Hence this entire common-shift budget cannot certify degree-three
exclusion at an integer base, for which the criterion requires
$\Lambda<1/3$.
\end{theorem}
\begin{proof}
For $0\le x\le1$, the numerator of $\Lambda'(x)$ is exactly the
quartic displayed above. Its constant term is negative, all its other
coefficients are positive, and it is strictly increasing on $x\ge0$.
It has one positive zero, lying between zero and one. For $x\ge1$,
the numerator of $\Lambda'$ is
\[
 \left(2+\frac3{2\pi^2}\right)x^2
 +\left(\frac{13}{12}+\frac3{2\pi^2}\right)x
 +\frac1{4\pi^2}-\frac7{12},
\]
which is positive for $x\ge1$. The two formulas agree at $x=1$,
so the minimum is global and unique.

For the uniform inequality use $\pi^2<10$. If $0\le x\le1$,
\[
 C(x)-\frac25 S(x)
 >\frac1{45}-\frac{x}{60}+\frac{17x^2}{30}+\frac{x^3}{9}
 \ge\frac1{180}>0.
\]
If $x\ge1$, the same difference is greater than
\[
             \frac9{10}x^2-\frac7{20}x+\frac2{15}>0.
\]
The numerical values are diagnostics of the exact unique-root
characterization, not premises of the barrier proof.
\end{proof}

\subsection{Allowing independent moment and pole shifts}
It is natural to ask whether the common shift itself caused the loss.
Let the moment shift divided by $h$ tend to $a$, and the pole shift
 divided by $h$ tend to $b$. The determinant is still a positive Gram
determinant at the real Lambert value: only its positive scalar weight
changes. Its analytic scale becomes
\[
                        S(a,b)=ab+a+b+5/6.
\]
The augmented matrix in the external construction is symmetric in its
two indices. Transposition and the monomial factors in its Vandermonde
normalization therefore give the exact symmetry $Q_{a,b}=Q_{b,a}$
at the discrete integer-shift level.

\begin{proposition}[Asymmetric shifts do not improve the central cone]
\label{p12:lf:twoshift}
In the parameter region
\[
 a,b\ge0,\qquad 2a\le b+1,\qquad 2b\le a+1,
\]
put $s=a+b$ and $p=ab$. The same local bounds give
\[
 \begin{split}
 A(a,b)&=\frac1{18}+\frac s6+\frac{s^2}6-\frac{s^3}9+\frac{sp}2,\\
 B(a,b)&=\frac18+
 \frac{7/4+9s/4+9s^2/4+3s^3/4-(6+3s)p}{\pi^2}.
 \end{split}
\]
For fixed $s$, the ratio $(A+B)/(p+s+5/6)$ is minimized at
$a=b=s/2$. Thus independent shifts in this entire cone do not beat
the global minimum in Theorem~\ref{p12:lf:optimization}.
\end{proposition}
\begin{proof}
The zero-place cost is
\[
 w_{a,b}(j,r)=(a+2-2r)j-\int_b^{b+1}\min(j,v)\,dv.
\]
The minimizing nodes in the central cone are $r$ for $r<b$,
$(b+r)/2$ for $b\le r\le(2+2a-b)/3$, and
$2r-1-a+b$ thereafter. Integrating these expressions and subtracting
from the infinity normalization $a(b+1/2)$ gives $A(a,b)$.

For the cyclotomic bound, the residue matrix is supported on at most
$\tau_a(n)$ rows plus $\tau_b(n)$ columns, rather than twice one
common support size. Here the normalized support function is
$\tau_a(u)=\min(1,(1+a-u)_+)$. In the present cone $|a-b|\le1$,
and the cap appearing in the pole envelope is
\[
 \min(1,\tau_a(u)+\tau_b(u))
 =(1+a-u)_++(1+b-u)_+
               -2((1+a+b)/2-u)_+.
\]
The same tent summation as before gives
\[
 B(a,b)=\frac18+
      \frac{(1+a)^3+(1+b)^3-(1+a+b)^3/4}{\pi^2},
\]
which expands to the stated expression.

The cone implies $0\le s\le2$. For fixed $s$, the coefficient of $p$
in $A+B$ is
\[
                  \frac s2-\frac{6+3s}{\pi^2}<0,
\]
using $\pi^2<10$ and $s\le2$. The numerator is positive, whereas the
denominator increases with $p$. Thus the ratio strictly decreases
with $p$. Its largest allowed $p$ is $s^2/4$, attained at $a=b$;
that point remains inside the cone. This proves the assertion.
\end{proof}
No assertion is made here that asymmetric shifts outside this cone are
optimal or impossible to improve. Nor does the result cover all six
bordered presentations or better local denominator estimates.

\subsection{Arithmetic outcomes and the remaining step}
The external nonquadraticity announcement is an important new literature
input. Our contribution in this branch is the exact optimization and
its obstruction, not another proof of that announcement. The optimized
single-presentation rational-base cutoff would be approximately
$1-\Lambda(x_*)=0.5728629136$; the source already reports a stronger
six-presentation cutoff $0.573290$, so the numerical optimization is
not a new best irrationality region.

The transcendence of $L(1/2)$ is not proved. The full reciprocal Fibonacci
sum is not classified as transcendental, and $B_1$ is not classified as
irrational. Sparse Lambert series with exponents $d^h$ have arithmetic
specialization theorems \cite{p12:lf:BundschuhVaananen}, but they are different
functions from the full Lambert series. The modular component of the
Fibonacci decomposition likewise has classical transcendence results
\cite{p12:lf:Duverney}, which do not classify the full sum.

A cubic exclusion by this determinant route needs genuinely smaller local
clearing costs or a different certificate family; changing the common
shape, or introducing central-cone asymmetry, does not supply it.


% END EMBEDDED lambert_results.tex

\section{Sixteen named targets: individual outcomes}
\label{p12:cf:target-ledger}

This ledger records the conclusions of this manuscript, rather than
declaring the complete current status of every target in the
literature. A statement that a target is unclosed means that our
arguments do not prove its listed conclusion. External announcements
are not counted as our results. The twelve Gamma-ratio examples in
Section~\ref{p12:sec:actual-beta-separation} are not twelve additional
historically open problems.

For clarity, $\gamma=\lim_{n\to\infty}(\sum_{j=1}^n j^{-1}-\log n)$,
$\delta=\int_0^\infty e^{-t}/(1+t)\,dt$, and
$G=\sum_{n\ge0}(-1)^n/(2n+1)^2$.
Write $L(z)=\sum_{n\ge1}z^n/(1-z^n)$ for $|z|<1$,
let $F_1=F_2=1$ and $F_{n+2}=F_{n+1}+F_n$, and define
the prime Mertens constant by
$B_1=\lim_{x\to\infty}(\sum_{p\le x}p^{-1}-\log\log x)$.

\begingroup\small
\renewcommand{\arraystretch}{1.12}
\begin{longtable}{@{}L{0.255\textwidth}L{0.695\textwidth}@{}}
\toprule
Named target & Result of this follow-up and remaining step\\
\midrule\endfirsthead
\toprule
Named target & Result of this follow-up and remaining step\\
\midrule\endhead
\bottomrule\endfoot
Euler's constant $\gamma$: irrationality &
No new Euler-specific arithmetic certificate is obtained. Entropy
rigidity cannot be applied to an unknown target merely by approximating
it with algebraic numbers. The previous curvature obstruction is not
removed. This target remains unclosed here.\\
Euler--Gompertz $\delta$: irrationality &
The fixed-$A$ results include $A=1$. The mean of $\log g_n/n$ tends
to zero quantitatively, and only $O_\varepsilon((\log X)^4)$ indices
violate a fixed linear threshold. A bounded primitive form would
require many growing prime divisors. An exceptional infinite
subsequence is not excluded, so no irrationality conclusion follows.\\
Catalan $G$: transcendence &
The polynomial-invariant Hankel family has a uniform positive root
gap. Restoring the full paired polynomial space makes all roots
converge to $G$ at every fixed determinant degree. Primitive height
relative to the small normalized value remains uncontrolled.\\
$\zeta(3)$: transcendence &
The moment-space analysis explains how a choice of invariant can
destroy approximation. It gives no new exclusion of an algebraic
degree for $\zeta(3)$; sufficient primitive-value and off-axis bounds
are still absent.\\
$\zeta(5)$: transcendence &
No primitive-content gain sufficient for a transcendence certificate
is established. Statements about external irrationality results would
not settle this stronger target and are not counted here.\\
$\zeta(7)$: irrationality &
No individual integer linear forms with a sufficient primitive-value
estimate are constructed. The general moment obstruction is not an
irrationality proof for this value.\\
$\zeta(9)$: irrationality &
The denominator and content estimate needed for this individual value
remains missing. Positivity and fixed-degree root convergence do not
replace it.\\
$\zeta(11)$: irrationality &
No individual conclusion is obtained. A relation or a disjunction
involving several odd zeta values would not by itself establish this
target.\\
Erd\H{o}s--Borwein $L(1/2)$: transcendence &
The specified common-shift Lambert budget has a unique optimum and
never reaches the cubic-exclusion threshold. The separate external
nonquadraticity announcement is not our theorem and does not prove
transcendence.\\
Reciprocal Fibonacci sum $\sum_{n\ge1}F_n^{-1}$: transcendence &
The functional transfers and classical modular-component results do
not classify the full sum. The present Lambert budget optimization
provides no sufficient specialization theorem.\\
Prime Mertens $B_1$: irrationality &
No arithmetic classification of this constant is obtained. Functional
independence of a generating function is not a value theorem for its
limiting constant.\\
$\Gamma(1/5)$: transcendence &
Writing $X=\Gamma(1/5)$ and $Y=\Gamma(2/5)$, the actual period-space
calculation proves, for example, that $\pi X/Y^3$ is transcendental.
It does not isolate the arithmetic nature of $X$.\\
$\Gamma(2/5)$: transcendence &
The same exact span gives transcendence of $X^2Y^4/\pi^3$ and a
complete linear-fractional test on the stated support. No conclusion
for $Y$ alone follows.\\
$\Gamma(1/7)$: transcendence &
For $R=\Gamma(1/7)$, $S=\Gamma(2/7)$ and $T=\Gamma(3/7)$,
$R^4/(\pi S^2)$ is among the explicit classical-period consequences.
No individual conclusion for $R$ is proved.\\
$\Gamma(2/7)$: transcendence &
The ratio $\pi S/T^3$ is transcendental by the rank calculation.
That does not decide the arithmetic nature of $S$ separately.\\
$\Gamma(3/7)$: transcendence &
The ratio $R^2T^4/\pi^3$ is transcendental. The nonlinear identity
$UZ=VW$ also shows why the period basis is not an algebraically
independent family. No individual conclusion for $T$ follows.\\
\end{longtable}
\endgroup

The number of newly closed previously open named-constant conjectures
in this ledger is therefore \textbf{zero}. The related Kurepa
left-factorial conjecture is likewise not proved: its boundary
congruence does not encode all index-root conditions. This accounting
is compatible with proving exact classifications for the restricted
entropy and period-space questions developed above.

\section{Verification record and precise remaining problems}
\label{p12:cf:verification}

The mathematical claims were reviewed independently within the
research workflow. The checks verify exact finite identities and
sample cases; the proofs carry the universal quantifiers. In
particular, the Gamma certificate enumerates all $16$ fifth-denominator
and $36$ seventh-denominator beta pairs, checks their distinct
signatures on units modulo $35$, and verifies $9{,}300$ exact
distribution identities including endpoints. The twelve displayed
ratios are checked as integer exponent differences. The entropy
checks confirm the explicit examples and finite algebraic steps.
The Gompertz, Catalan and Lambert computations are reproducible
diagnostics of recurrences, moment identities and budget formulas.
They do not establish irrationality by numerical approximation.

The use of the full Wolfart--W\"ustholz beta completeness theorem is
explicit. Its primary publisher abstract and later precise research
statements were consulted; the paywalled full original article was
not inspected in this workflow. The weaker pair-class independence
statement alone would not justify the signature obstruction, so it
is not substituted for completeness. The external Lambert manuscript
and its mathematical source were inspected, while its announced Lean
verification was not rebuilt. No result is described here as formally
verified, and the literature search does not certify historical
priority of our specializations or formulations.

Four concrete problems emerge from the proved results. First, can a
nonrational ternary algebraic probability vector in a field of degree
at least four have rational bit entropy? The cubic exclusion and
two-primary spacing condition narrow that question but do not answer
it. Second, can the full paired Catalan spaces be equipped with a
primitive-height estimate strong enough to turn their repaired root
convergence into arithmetic information? Third, can simultaneous
Gompertz divisibility by a growing prime set be controlled at every
relevant index, beyond the limiting distribution and polylogarithmic
exceptional count proved here? Fourth,
can the Lambert clearing budget be lowered by a new local estimate or
a different presentation rather than by the shape changes already
optimized? None of these missing estimates is supplied by renaming
an approximating root as a conjugate of the target constant.

The complete relative-entropy, conditional-entropy and Markov
classifications stand independently of these open transfers. The
period calculation supplies a second form of exactness: within its
specified fourteen-dimensional support, algebraicity of every
linear-fractional expression is decided by proportionality of
coefficient vectors. Beyond those hypotheses and supports, further
arithmetic input is required.

% END SOURCE: KEIO_Quadratic_Exceptions_and_Constant_Transfer.tex

\endgroup
\clearpage
\addcontentsline{toc}{section}{Consolidated bibliography}
\renewcommand{\refname}{Consolidated bibliography}
\begin{thebibliography}{999}
\raggedright
\bibitem{p01:root:original}
Keio, \emph{Observation and Recovery in Arithmetic and Probability:
Prime Flows, Rare Events, and Infinite Mixtures}, research manuscript,
14 October 2026, 223 pp. AI-assisted working manuscript, not externally peer reviewed.
\bibitem{p01:bec17:Beukers}
F. Beukers, \emph{A refined version of the Siegel--Shidlovskii theorem},
Annals of Mathematics \textbf{163} (2006), 369--379,
Theorem~1.3 and Corollary~1.4,
\url{https://doi.org/10.4007/annals.2006.163.369},
\url{https://annals.math.princeton.edu/wp-content/uploads/annals-v163-n1-p08.pdf}.
\bibitem{p01:bec17:DLMF}
F. W. J. Olver et al., eds., \emph{NIST Digital Library of Mathematical Functions},
Sections 10.21(iii), 10.40(i), and 13.2--13.4.
\url{https://dlmf.nist.gov/10.21.iii};
\url{https://dlmf.nist.gov/10.40.i};
\url{https://dlmf.nist.gov/13.2};
\url{https://dlmf.nist.gov/13.3};
\url{https://dlmf.nist.gov/13.4}.
\bibitem{p01:bec17:Previous}
Keio, \emph{Observation and Recovery: Geometry, Mixtures, and Spectral Inversion},
unified sequel research manuscript, 16 October 2026, 447 pp.
Part X, ``Simultaneous gamma factors and reconstruction from boundary response'';
also the Bessel obstruction and spectral boundary sections.
AI-assisted working manuscript, not externally peer reviewed.
\bibitem{p01:ga17:EW}
E.~Ekstr\"om and H.~Wanntorp, Optimal stopping of a Brownian bridge,
\emph{Journal of Applied Probability} \textbf{46} (2009), 170--180.
doi:10.1239/jap/1238592123.
\url{https://doi.org/10.1239/jap/1238592123}.
\bibitem{p01:ga17:HL}
D.~Hobson and J.~Liu, Optimal stopping of a Brownian excursion and
an $\alpha$-dimensional Bessel bridge,
\emph{Journal of Applied Probability} (2026), First View, 1--18.
doi:10.1017/jpr.2026.10102.
\url{https://doi.org/10.1017/jpr.2026.10102}.
Preprint: \url{https://arxiv.org/abs/2504.19741}.
\bibitem{p01:ga17:FR}
S.~Fischler and T.~Rivoal, Relations between values of arithmetic
Gevrey series, and applications to values of the Gamma function,
\emph{Journal of Number Theory} \textbf{261} (2024), 36--54.
doi:10.1016/j.jnt.2024.02.016.
\url{https://arxiv.org/abs/2301.13518}.
\bibitem{p01:sc:GS}
M.~Karatarakis and F.~Wiedijk,
\emph{A formalization of the Gelfond--Schneider theorem},
arXiv:2603.24823v1 (2026), Theorem~3.1.
\href{https://arxiv.org/abs/2603.24823}{arXiv:2603.24823}.
Used solely as a primary proof source for the classical theorem.
\bibitem{p01:sc:LW}
S.~A.~Popescu,
\emph{A simple and self-contained proof for the Lindemann--Weierstrass
theorem}, arXiv:2306.14352v2 (2023).
\href{https://arxiv.org/abs/2306.14352}{arXiv:2306.14352}.
\bibitem{p01:sc:Baker}
A.~Baker,
\emph{Linear forms in the logarithms of algebraic numbers},
Mathematika \textbf{13} (1966), 204--216.
\href{https://doi.org/10.1112/S0025579300003971}
{doi:10.1112/S0025579300003971}.
\bibitem{p01:sc:fractal}
M.~L.~Lapidus, M.~van Frankenhuijsen, and E.~K.~Voskanian,
\emph{Quasiperiodic patterns of the complex dimensions of nonlattice
self-similar strings, via the LLL algorithm},
Mathematics \textbf{9} (2021), no.~6, article 591.
\href{https://doi.org/10.3390/math9060591}{doi:10.3390/math9060591};
\href{https://arxiv.org/abs/2009.03493}{arXiv:2009.03493}.
\bibitem{p01:sc:factorization}
N.~Lebowitz-Lockard,
\emph{On a theorem of Erd\H{o}s and Loxton},
arXiv:2504.03446 (2025).
\href{https://arxiv.org/abs/2504.03446}{arXiv:2504.03446}.
The introduction records the classical role of the real solution of
$\zeta(s)=2$ in ordered factorization.
\bibitem{p01:mc17:Ostrowski}
A.~Ostrowski, Sur les relations alg\'ebriques entre les int\'egrales
ind\'efinies, \emph{Acta Mathematica} \textbf{78} (1946), 315--318.
doi:10.1007/BF02421605.
\bibitem{p01:mc17:Srinivasan}
V.~R.~Srinivasan, \emph{Extensions by Antiderivatives, Exponentials of
Integrals and by Iterated Logarithms}, arXiv:0811.3004 (2008),
Theorem~2.5 and the discussion of the Kolchin--Ostrowski theorem.
\url{https://arxiv.org/abs/0811.3004}.
\bibitem{p01:mc17:Rivoal}
T.~Rivoal, On the arithmetic nature of the values of the gamma function,
Euler's constant, and Gompertz's constant,
\emph{Michigan Mathematical Journal} \textbf{61} (2012), 239--254.
doi:10.1307/mmj/1339011525.
\url{https://rivoal.perso.math.cnrs.fr/articles/gammater.pdf}.
\bibitem{p01:mc17:Powers}
M.~R.~Powers, \emph{On the Algebraic Independence of a Set of
Generalized Constants}, arXiv:2508.12123v3 (2025).
\url{https://arxiv.org/abs/2508.12123}.
This preprint supplies related probabilistic context; none of its new
claims is required for the proofs in this section.
\bibitem{p01:mc17:Waldschmidt}
M.~Waldschmidt, \emph{Schanuel's Conjecture: algebraic independence of
transcendental numbers}, Colloquium De Giorgi, 6 May 2014,
updated 20 May 2014, Sections~1--2.
\url{https://webusers.imj-prg.fr/~michel.waldschmidt/articles/pdf/ColloquiumDeGiorgiSchanuelConjecture2014.pdf}.
\bibitem{p01:ac17:VW}
W. Van Assche and T. Wolfs,
\emph{Rational approximation of Euler's constant using multiple orthogonal
polynomials}, Combinatorics and Number Theory \textbf{14} (2025), 141--162.
\href{https://doi.org/10.2140/cnt.2025.14.141}{doi:10.2140/cnt.2025.14.141};
\href{https://arxiv.org/abs/2404.09799}{arXiv:2404.09799}, version 3 (2025).
\bibitem{p01:ac17:Powers}
M. R. Powers, \emph{On the uniqueness of Ein(1) among linear combinations
of the Euler--Mascheroni and Euler--Gompertz constants}, preprint (2025),
\href{https://arxiv.org/abs/2510.00315}{arXiv:2510.00315}.
\bibitem{p01:ac17:Sun}
Z.-W. Sun, \emph{Catalan's constant is irrational}, preprint,
3 September 2026,
\href{https://arxiv.org/abs/2609.04176}{arXiv:2609.04176}.
This reference is a claim under discussion, not an input theorem here.
\bibitem{p01:ac17:Wachs}
D. Wachs, \emph{A note on a recent claimed proof of the irrationality
of Catalan's constant}, preprint, 16 September 2026,
\href{https://arxiv.org/abs/2609.22339}{arXiv:2609.22339}.
\bibitem{p01:ac17:BrownZudilin}
F. Brown and W. Zudilin, \emph{On cellular rational approximations to
$\zeta(5)$},
\href{https://arxiv.org/abs/2210.03391}{arXiv:2210.03391},
version 3, 29 January 2026.
\bibitem{p01:ci17:ABS}
B.~Adamczewski, J.~Bell, and D.~Smertnig,
\emph{A height gap theorem for coefficients of Mahler functions},
J. Eur. Math. Soc. \textbf{25} (2023), 2525--2571.
\url{https://doi.org/10.4171/JEMS/1244}.
Theorems~1.1 and~1.2; inspected publisher text:
\url{https://ems.press/content/serial-article-files/32888}.
\bibitem{p01:ci17:TT}
T.~Tao and J.~Ter\"av\"ainen,
\emph{Quantitative correlations and some problems on prime factors
of consecutive integers}, arXiv:2512.01739v2, 25 April 2026,
Theorem~1.3 and Section~1.3.
\url{https://arxiv.org/html/2512.01739v2}.
\bibitem{p01:ci17:Coons}
M.~Coons, \emph{On the rational approximation of the sum of the
reciprocals of the Fermat numbers}, Ramanujan J. \textbf{30} (2013),
39--52. Author preprint arXiv:1112.5072v1, Theorem~1.
\url{https://arxiv.org/html/1112.5072v1}.
\bibitem{p01:ci17:Adamczewski}
B.~Adamczewski, \emph{The many faces of the Kempner number},
J. Integer Sequences \textbf{16} (2013), Article~13.2.15.
\url{https://cs.uwaterloo.ca/journals/JIS/VOL16/Adamczewski/adam6.html}.
\url{https://arxiv.org/abs/1303.1685}.
\bibitem{p01:ci17:Yuen}
J.~Yuen, \emph{Kleene star of the primes is not regular in any base},
arXiv:2203.16088v2, 8 April 2022.
\url{https://arxiv.org/pdf/2203.16088}.
The introduction records the classical nonregularity theorem and
stronger finite-range state lower bounds.
\bibitem{p01:ci17:Erdos}
P.~Erd\H{o}s, \emph{On arithmetical properties of Lambert series},
J. Indian Math. Soc. (N.S.) \textbf{12} (1948), 63--66.
\url{https://users.renyi.hu/~p_erdos/1948-04.pdf}.
\bibitem{p01:ci17:Borwein}
P.~B.~Borwein, \emph{On the irrationality of certain series},
Math. Proc. Cambridge Philos. Soc. \textbf{112} (1992), 141--146.
\url{https://doi.org/10.1017/S030500410007081X}.
\url{https://www.cecm.sfu.ca/~pborwein/PAPERS/P59.pdf}.
\bibitem{p02:root:sequels}
Keio, \emph{Observation and Recovery: Geometry, Mixtures, and Spectral Inversion},
unified sequel research manuscript, 16 October 2026, 447 pp.
Part VIII, sharp large-order gamma-extraction separation; Part IX,
intermediate-rate factor region; Part X, simultaneous gamma factors.
AI-assisted working manuscript, not externally peer reviewed.
\bibitem{p02:root:previous}
Keio, \emph{Arithmetic of Observation Boundaries: Transcendence Transfers
and Diophantine Obstructions}, research manuscript, 17 October 2026,
41 pp. The immediate predecessor of the present paper.
AI-assisted working manuscript, not externally peer reviewed.
\bibitem{p02:sc2:GS}
M.~Karatarakis and F.~Wiedijk,
\emph{A formalization of the Gelfond--Schneider theorem},
arXiv:2603.24823v1 (2026), Theorem~3.1.
\url{https://arxiv.org/abs/2603.24823}.
Used as a primary proof account for the classical theorem only.
\bibitem{p02:sc2:Schanuel}
M.~Waldschmidt,
\emph{Schanuel's Conjecture: algebraic independence of transcendental
numbers}, Colloquium de Giorgi, 6 May 2014, updated 20 May 2014,
Sections~1--2.
\url{https://webusers.imj-prg.fr/~michel.waldschmidt/articles/pdf/ColloquiumDeGiorgiSchanuelConjecture2014.pdf}.
The conjecture is used only in the explicitly conditional corollary.
\bibitem{p02:ga18:HL}
D.~Hobson and J.~Liu, Optimal stopping of a Brownian excursion and
an $\alpha$-dimensional Bessel bridge,
\emph{Journal of Applied Probability} (2026), First View, 1--18.
doi:10.1017/jpr.2026.10102.
\url{https://doi.org/10.1017/jpr.2026.10102}.
\bibitem{p02:be18:Beukers}
F.~Beukers, A refined version of the Siegel--Shidlovskii theorem,
\emph{Annals of Mathematics} \textbf{163} (2006), 369--379,
Corollary~1.4.
\url{https://annals.math.princeton.edu/wp-content/uploads/annals-v163-n1-p08.pdf}.
\bibitem{p02:be18:Rivoal}
T.~Rivoal, On Galochkin's characterization of hypergeometric
$G$-functions, \emph{Moscow Journal of Combinatorics and Number Theory}
\textbf{11} (2022), no.~1, 11--19, Theorem~1.
\url{https://doi.org/10.2140/moscow.2022.11.11}.
Author's version: \url{https://rivoal.perso.math.cnrs.fr/articles/galochkin.pdf}.
\bibitem{p02:be18:MurtySaha}
M.~Ram Murty and E.~Saha, Transcendental values of the incomplete
gamma function and related questions,
\emph{Archiv der Mathematik} \textbf{105} (2015), no.~3, 271--283,
Corollary~1.
\url{https://doi.org/10.1007/s00013-015-0800-3}.
Author's copy: \url{https://mrammurty-ai.github.io/Ekata-Ram-IncompleteGamma.pdf}.
\bibitem{p02:be18:Arreche}
C.~E.~Arreche, A Galois-theoretic proof of the differential
transcendence of the incomplete Gamma function,
\emph{Journal of Algebra} \textbf{389} (2013), 119--127.
\url{https://doi.org/10.1016/j.jalgebra.2013.04.037}.
\url{https://arxiv.org/abs/1304.1917}.
\bibitem{p02:be18:DLMF}
F. W. J. Olver et al., eds., \emph{NIST Digital Library of Mathematical Functions},
Sections 10.25, 10.29, 10.40 and 13.2--13.6.
\url{https://dlmf.nist.gov/10.40}; \url{https://dlmf.nist.gov/13.3};
\url{https://dlmf.nist.gov/13.6}.
\bibitem{p02:ac18:VW}
T. Wolfs and W. Van Assche, \emph{Rational approximation of Euler's
constant using multiple orthogonal polynomials}, Combinatorics and
Number Theory \textbf{14} (2025), 141--162.
\href{https://doi.org/10.2140/cnt.2025.14.141}{doi:10.2140/cnt.2025.14.141};
\href{https://arxiv.org/abs/2404.09799}{arXiv:2404.09799}.
\bibitem{p02:ac18:Margulis}
G. A. Margulis, \emph{Indefinite quadratic forms and unipotent flows on
homogeneous spaces}, Banach Center Publications \textbf{23} (1989),
399--409. \href{https://doi.org/10.4064/-23-1-399-409}{doi:10.4064/-23-1-399-409}.
\bibitem{p02:ac18:Mixed}
S. Fischler and T. Rivoal, \emph{Relations between values of arithmetic
Gevrey series, and applications to values of the Gamma function},
Journal of Number Theory \textbf{261} (2024), 36--54.
\href{https://arxiv.org/abs/2301.13518}{arXiv:2301.13518}.
\bibitem{p02:ci18:Maynard2016}
J.~Maynard, \emph{Dense clusters of primes in subsets},
Compositio Math. \textbf{152} (2016), 1517--1554, Theorem~3.4.
\url{https://arxiv.org/abs/1405.2593}.
\bibitem{p02:ci18:AB2007}
B.~Adamczewski and Y.~Bugeaud,
\emph{On the complexity of algebraic numbers I. Expansions in integer bases},
Ann. of Math. \textbf{165} (2007), 547--565.
\url{https://doi.org/10.4007/annals.2007.165.547}.
\bibitem{p02:ci18:TaoTeravainen}
T.~Tao and J.~Ter\"av\"ainen,
\emph{Quantitative correlations and some problems on prime factors of
consecutive integers}, arXiv:2512.01739v2 (2026), Theorem~1.3.
\url{https://arxiv.org/html/2512.01739v2}.
\bibitem{p02:ci18:Nesterenko}
Yu.~V.~Nesterenko, \emph{Modular functions and transcendence questions},
Sbornik Math. \textbf{187} (1996), 1319--1348.
\url{https://doi.org/10.1070/SM1996v187n09ABEH000158}.
\bibitem{p02:ci18:Campbell}
J.~M.~Campbell, \emph{On the binary digits of the Erd\H{o}s--Borwein
constant}, arXiv:2605.24160v1 (2026).
\url{https://arxiv.org/html/2605.24160v1}.
\bibitem{p02:ci18:Disjunctivity}
\emph{Disjunctivity of the Erd\H{o}s--Borwein constant},
authorless public manuscript and conditional Lean development,
released September~7, 2026 in the CaptainSude repository.
\url{https://github.com/CaptainSude/erdos-borwein-disjunctivity}.
\bibitem{p03:root3:original}
Keio, \emph{Observation and Recovery in Arithmetic and Probability:
Prime Flows, Rare Events, and Infinite Mixtures}, integrated foundational
project manuscript, 223 pp. (2026). AI-assisted working manuscript.
\bibitem{p03:root3:sequels}
Keio, \emph{Observation and Recovery: Geometry, Mixtures, and Spectral
Inversion}, unified sequel project manuscript, 447 pp. (2026).
Parts VIII--X contain the gamma-extraction inputs cited here.
AI-assisted working manuscript.
\bibitem{p03:root3:previousI}
Keio, \emph{Arithmetic of Observation Boundaries: Transcendence Transfers
and Diophantine Obstructions}, project manuscript, 41 pp. (17 October 2026).
AI-assisted working manuscript; includes bridge-boundary classification
and endpoint-shape irrationality.
\bibitem{p03:root3:previousII}
Keio, \emph{Arithmetic of Optimal Values: Curvature, Special-Function
Elimination, and Positivity Barriers}, project manuscript, 34 pp.
(17 October 2026). AI-assisted working manuscript; immediate predecessor.
\bibitem{p03:root3:Catalan}
OpenAI, \emph{Catalan's constant is irrational}, OpenAI Math Release
preprint, 24 September 2026. The external claim is recorded for current
context; its full proof is not audited or invoked in this manuscript.
\url{https://github.com/openai/math/tree/main/preprints/Catalans-constant-is-irrational-September-24-2026}.
\bibitem{p03:sc3:Brownawell}
W.~Dale Brownawell,
\emph{On the Gelfond--Feldman measure of algebraic independence},
Compositio Mathematica \textbf{38} (1979), no.~3, 355--368,
Theorem~2 and the discussion on pp.~357--358.
\url{https://www.numdam.org/article/CM_1979__38_3_355_0.pdf}.
The arbitrary-base hypothesis is essential to the application here.
\bibitem{p03:sc3:GS}
M.~Karatarakis and F.~Wiedijk,
\emph{A formalization of the Gelfond--Schneider theorem},
arXiv:2603.24823v1 (2026), Theorem~3.1.
\url{https://arxiv.org/abs/2603.24823}.
Used as a primary proof account of the classical theorem only.
\bibitem{p03:ga19:Beukers}
F.~Beukers, A refined version of the Siegel--Shidlovskii theorem,
\emph{Annals of Mathematics} \textbf{163} (2006), 369--379.
\url{https://annals.math.princeton.edu/wp-content/uploads/annals-v163-n1-p08.pdf}.
\bibitem{p03:ga19:DLMF}
NIST Digital Library of Mathematical Functions,
\S13.7, Asymptotic expansions for large argument,
especially 13.7.2 and the separate remainder bounds in \S13.7(ii).
\url{https://dlmf.nist.gov/13.7}.
\bibitem{p03:ga19:FR}
S.~Fischler and T.~Rivoal, Zeros of $E$-functions and of exponential
polynomials defined over $\overline{\mathbb Q}$,
\emph{Rendiconti del Circolo Matematico di Palermo, Series 2}
\textbf{75} (2026), article 158.
doi:10.1007/s12215-026-01475-x.
Published version, Conjecture~1.4 and Theorem~1.5:
\url{https://www.imo.universite-paris-saclay.fr/~stephane.fischler/ezeros.pdf}.
\bibitem{p03:ga19:Dirichlet}
P.~G.~L.~Dirichlet, Beweis des Satzes, dass jede unbegrenzte
arithmetische Progression, deren erstes Glied und Differenz ganze
Zahlen ohne gemeinschaftlichen Factor sind, unendlich viele
Primzahlen enth\"alt,
\emph{Abhandlungen der K\"oniglich Preussischen Akademie der
Wissenschaften} (1837), 45--81.
English translation of the original memoir:
\url{https://arxiv.org/abs/0808.1408}.
\bibitem{p03:be19:DLMF}
F.~W.~J.~Olver et al., eds., \emph{NIST Digital Library of
Mathematical Functions}, Sections~10.25 and 10.29,
especially formulas~10.29.2 and 10.29.4.
\url{https://dlmf.nist.gov/10.29}.
\bibitem{p03:mi3:BPR}
S.~Basu, R.~Pollack, and M.-F.~Roy,
\emph{Algorithms in Real Algebraic Geometry}, second edition,
Springer, 2006. \url{https://doi.org/10.1007/3-540-33099-2}.
\bibitem{p03:mi3:Baker}
A.~Baker, \emph{Transcendental Number Theory}, Cambridge University
Press, 1975. \url{https://doi.org/10.1017/CBO9780511565977}.
\bibitem{p03:mi3:Dukkipati}
A.~Dukkipati, \emph{Maximum Entropy in the framework of Algebraic
Statistics: A First Step}, 2007.
\url{https://arxiv.org/abs/cs/0701117}.
See also the expanded report \emph{Algebraic Descriptions of Maximum
Entropy}, EURANDOM Report 2008-035,
\url{https://www.eurandom.tue.nl/reports/2008/035-report.pdf},
Propositions 2.1--2.3.
\bibitem{p03:mi3:toric}
B.~Sturmfels, S.~Telen, F.-X.~Vialard, and M.~von Renesse,
\emph{Toric Geometry of Entropic Regularization},
\url{https://arxiv.org/abs/2202.01571}, version 2, 2023.
\bibitem{p03:mi3:Blahut}
R.~E.~Blahut, Computation of channel capacity and rate-distortion
functions, \emph{IEEE Transactions on Information Theory}
\textbf{18} (1972), 460--473.
\url{https://doi.org/10.1109/TIT.1972.1054855}.
\bibitem{p03:mi3:Boche}
H.~Boche, R.~F.~Schaefer, and H.~V.~Poor,
\emph{Algorithmic Computability and Approximability of
Capacity-Achieving Input Distributions},
\url{https://arxiv.org/abs/2202.12617}, version 2, 2023;
especially Appendix B and the input representation in Section II.
\bibitem{p03:mi3:Liang}
X.-B.~Liang, An algebraic, analytic, and algorithmic investigation on
the capacity and capacity-achieving input probability distributions
of finite-input--finite-output discrete memoryless channels,
\emph{IEEE Transactions on Information Theory} \textbf{54} (2008),
1003--1023. \url{https://doi.org/10.1109/TIT.2007.915703}.
\bibitem{p03:mi3:IB}
N.~Tishby, F.~C.~Pereira, and W.~Bialek,
The information bottleneck method,
\emph{Proceedings of the 37th Annual Allerton Conference}, 1999,
368--377. \url{https://arxiv.org/abs/physics/0004057}.
\bibitem{p03:mi3:WynerZiv}
A.~D.~Wyner and J.~Ziv, A theorem on the entropy of certain binary
sequences and applications: Part I,
\emph{IEEE Transactions on Information Theory} \textbf{19} (1973),
769--772. \url{https://doi.org/10.1109/TIT.1973.1055107}.
\bibitem{p03:mi3:Hsu}
H.~Hsu, S.~Asoodeh, S.~Salamatian, and F.~P.~Calmon,
\emph{Generalizing Bottleneck Problems},
IEEE International Symposium on Information Theory, 2018.
\url{https://arxiv.org/abs/1802.05861}, version 3.
\bibitem{p03:mi3:BlakajWolf}
V.~Blakaj and M.~M.~Wolf, Transcendental properties of
entropy-constrained sets, \emph{Annales Henri Poincar\'e}
\textbf{24} (2023), 349--362.
\url{https://doi.org/10.1007/s00023-022-01227-4}.
\bibitem{p03:ac19:FRring}
S. Fischler and T. Rivoal, \emph{The ring of values of $E$-functions
is not a field}, preprint, arXiv:2610.16603, submitted 12 October 2026;
author-hosted version dated 14 October 2026.
\href{https://www.imo.universite-paris-saclay.fr/~fischler/Ecroix.pdf}{Author's PDF};
\href{https://arxiv.org/abs/2610.16603}{arXiv:2610.16603}.
\bibitem{p03:ac19:FRmixed}
S. Fischler and T. Rivoal, \emph{Relations between values of arithmetic
Gevrey series, and applications to values of the Gamma function},
Journal of Number Theory \textbf{261} (2024), 36--54.
\href{https://doi.org/10.1016/j.jnt.2024.02.016}{doi:10.1016/j.jnt.2024.02.016};
\href{https://www.imo.universite-paris-saclay.fr/~stephane.fischler/ssmixte.pdf}{Author's PDF}.
\bibitem{p03:ac19:Dusart}
P. Dusart, \emph{Estimates of Some Functions Over Primes without R.H.},
preprint (2010),
\href{https://arxiv.org/abs/1002.0442}{arXiv:1002.0442}.
\bibitem{p03:ci19:Trudgian2016}
T.~S.~Trudgian,
\emph{Updating the error term in the prime number theorem},
Ramanujan J. \textbf{39} (2016), 225--234.
\url{https://arxiv.org/abs/1401.2689}.
Theorem~2 gives the exponential error for the prime-counting function
used here; the proof only needs the resulting error of every fixed
inverse-logarithmic order.
\bibitem{p03:ci19:KanekoZagier1995}
M.~Kaneko and D.~Zagier,
\emph{A generalized Jacobi theta function and quasimodular forms},
in \emph{The Moduli Space of Curves}, Progress in Mathematics
\textbf{129}, Birkh\"auser, 1995, 165--172.
\url{https://www2.math.kyushu-u.ac.jp/~mkaneko/papers/genjacobi.pdf}.
See Section~1, especially the weight-two transformation law and
closure under differentiation.
\bibitem{p03:ci19:Pupyrev2005}
Yu.~A.~Pupyrev,
\emph{Linear and algebraic independence of $q$-zeta values},
Math. Notes \textbf{78} (2005), 563--568.
\url{https://doi.org/10.1007/s11006-005-0155-3}.
\bibitem{p03:ci19:Bezivin2008}
J.-P.~B\'ezivin,
\emph{Ind\'ependance lin\'eaire et alg\'ebrique de fonctions li\'ees
\`a la fonction $q$-dzeta},
Ann. Fac. Sci. Toulouse Math. (6) \textbf{17} (2008), 23--36.
\url{https://doi.org/10.5802/afst.1173}.
\bibitem{p04:root4:III}
Keio, \emph{Arithmetic Independence of Variational Constants:
Cubic Spectra, Fixed-Dimension Stopping, and Entropy Optima},
17 October 2026, 45 pp. AI-assisted project manuscript.
Standalone source: \texttt{KEIO\_Arithmetic\_Independence\_and\_Entropy\_Optima.tex}.
\bibitem{p04:ga20:JK}
A.~Jindal and S.~K.~Khanduja, An extension of Schur's irreducibility
result, \emph{Journal of Algebra} \textbf{664}, Part B (2025), 398--409.
doi:10.1016/j.jalgebra.2024.10.047.
Author version, Theorem~1.1 and the original Schur theorem in the
introduction: \url{https://arxiv.org/pdf/2305.04781}.
\bibitem{p04:ga20:LT}
R.~Levingston and D.~E.~Taylor, The theorem of Marggraff on primitive
permutation groups which contain a cycle,
\emph{Bulletin of the Australian Mathematical Society}
\textbf{15} (1976), 125--128.
doi:10.1017/S0004972700036832.
\url{https://doi.org/10.1017/S0004972700036832}.
\bibitem{p04:ga20:Jones}
G.~A.~Jones, Primitive permutation groups containing a cycle,
\emph{Bulletin of the Australian Mathematical Society}
\textbf{89} (2014), 159--165.
Section~3 on the Jordan--Marggraff theorem:
\url{https://arxiv.org/pdf/1209.5169}.
\bibitem{p04:ga20:FR}
S.~Fischler and T.~Rivoal, Zeros of $E$-functions and of exponential
polynomials defined over $\overline{\mathbb Q}$,
\emph{Rendiconti del Circolo Matematico di Palermo, Series 2}
\textbf{75} (2026), article 158.
doi:10.1007/s12215-026-01475-x.
Published version, definition of globally bounded $E$-functions
and Theorem~1.5:
\url{https://www.imo.universite-paris-saclay.fr/~stephane.fischler/ezeros.pdf}.
\bibitem{p04:mi4:Dubickas}
A.~Dubickas, Every Salem number is a difference of two Pisot numbers,
\emph{Proceedings of the Edinburgh Mathematical Society}
\textbf{66} (2023), 862--867.
\url{https://doi.org/10.1017/S0013091523000433}.
\bibitem{p04:mi4:Lind}
D.~A.~Lind, The entropies of topological Markov shifts and a related
class of algebraic integers, \emph{Ergodic Theory and Dynamical
Systems} \textbf{4} (1984), 283--300.
\url{https://doi.org/10.1017/S0143385700002443}.
\bibitem{p04:mi4:Poisson}
J.~A.~Adell, A.~Lekuona, and Y.~Yu,
Sharp bounds on the entropy of the Poisson law and related quantities,
\emph{IEEE Transactions on Information Theory}
\textbf{56} (2010), 2299--2306.
\url{https://doi.org/10.1109/TIT.2010.2044057}.
\url{https://arxiv.org/abs/1001.2897}.
\bibitem{p04:ac20:Lindemann}
F. Lindemann, \emph{Ueber die Zahl $\pi$}, Mathematische Annalen
\textbf{20} (1882), 213--225.
\href{https://doi.org/10.1007/BF01446522}{doi:10.1007/BF01446522}.
\bibitem{p04:sc4:Stanley}
R.~P.~Stanley,
\emph{Differentiably finite power series},
European Journal of Combinatorics \textbf{1} (1980), 175--188.
\url{https://math.mit.edu/~rstan/pubs/pubfiles/45.pdf}.
\bibitem{p04:sc4:RW}
D.~Roy and M.~Waldschmidt,
\emph{Approximation diophantienne et ind\'ependance alg\'ebrique
de logarithmes},
Annales scientifiques de l'\'Ecole Normale Sup\'erieure,
4e s\'erie, \textbf{30} (1997), no.~6, 753--796,
Corollary~1.2(a), pp.~756--757.
\url{https://www.numdam.org/article/ASENS_1997_4_30_6_753_0.pdf}.
\bibitem{p04:be20:DLMF}
F.~W.~J.~Olver et al., eds., \emph{NIST Digital Library of
Mathematical Functions}, Sections~10.29 and 10.39.
\url{https://dlmf.nist.gov/10.29};
\url{https://dlmf.nist.gov/10.39}.
\bibitem{p04:be20:Previous}
KEIO, \emph{Observation and Recovery: Geometry, Mixtures, and Spectral Inversion},
unified post-223-page manuscript, 16 October 2026, 447 pp.,
Part IX, theorem ``The unrestricted factor region'' and
lemma ``Uniform positivity at both ends.''
AI-assisted research manuscript, not externally peer reviewed.
\bibitem{p04:ci20:Trudgian2016}
T.~S.~Trudgian, \emph{Updating the error term in the prime number theorem},
Ramanujan J. \textbf{39} (2016), 225--234.
\url{https://arxiv.org/abs/1401.2689}. See Theorem~2.
\bibitem{p04:ci20:KanekoZagier1995}
M.~Kaneko and D.~Zagier,
\emph{A generalized Jacobi theta function and quasimodular forms},
in \emph{The Moduli Space of Curves}, Progress in Mathematics
\textbf{129}, Birkh\"auser, 1995, 165--172.
\url{https://www2.math.kyushu-u.ac.jp/~mkaneko/papers/genjacobi.pdf}.
\bibitem{p04:ci20:Zudilin2005}
W.~Zudilin, \emph{Approximations to $q$-logarithms and $q$-dilogarithms,
with applications to $q$-zeta values},
Zap. Nauchn. Sem. POMI \textbf{322} (2005), 107--124;
English translation, J. Math. Sci. (N.Y.) \textbf{137} (2006), 4673--4683.
\url{https://wain.mi-ras.ru/PS/qhyp.pdf}.
\bibitem{p04:ci20:TaoTeravainen2026}
T.~Tao and J.~Ter\"av\"ainen,
\emph{Quantitative correlations and some problems on prime factors
of consecutive integers}, preprint,
\url{https://arxiv.org/abs/2512.01739}, version~2, 25 April 2026.
\bibitem{p05:rv5:223}
Keio, \emph{Observation and Recovery in Arithmetic and Probability:
Prime Flows, Rare Events, and Infinite Mixtures}, 223-page foundational
project manuscript, 2026.
\bibitem{p05:rv5:447}
Keio, \emph{Observation and Recovery: Geometry, Mixtures, and Spectral
Inversion}, 447-page collected sequel manuscript, 2026. The chronological
ordering of its parts matters: some questions left open in earlier parts
are answered in later parts.
\bibitem{p05:rv5:IV}
Keio, \emph{Arithmetic Rigidity of Optimization and Special Functions:
Complete Planar Independence, Entropy Transitions, and Algebraic-Closure
Barriers}, 43-page preceding project manuscript, 17 October 2026.
\bibitem{p05:ga21:VS}
R.~A.~Vandermeulen and C.~D.~Scott,
On the identifiability of mixture models from grouped samples,
arXiv:1502.06644, version 2 (2022).
\url{https://arxiv.org/html/1502.06644v2}.
\bibitem{p05:ga21:BJR}
S.~Bonhomme, K.~Jochmans, and J.-M.~Robin,
Non-parametric estimation of finite mixtures from repeated measurements,
\emph{Journal of the Royal Statistical Society, Series B}
\textbf{78} (2016), 211--229.
\url{https://doi.org/10.1111/rssb.12110}.
\bibitem{p05:ga21:Simon}
T.~Simon, Comparing Fr\'echet and positive stable laws,
\emph{Electronic Journal of Probability} \textbf{19} (2014),
paper 16, 1--25.
\url{https://doi.org/10.1214/EJP.v19-3058};
\url{https://arxiv.org/html/1310.1888v2}.
\bibitem{p05:ga21:PX}
M.~R.~Powers and J.~Xu,
A criterion for extending continuous-mixture identifiability results,
\emph{Journal of Statistical Theory and Applications}
\textbf{24} (2025), 515--533.
\url{https://doi.org/10.1007/s44199-025-00124-z}.
\bibitem{p05:ga21:Previous}
Keio, \emph{Adaptive Recovery and Mixture Identification},
preceding project manuscript, 16 October 2026, sections
``Two joint replicates recover an arbitrary index--scale mixture''
and ``Contrasts and anchors recover correlated stable parameters.''
Internal branch source:
\texttt{research\_20261016\_round16/gaussian.tex};
integrated source:
\texttt{KEIO\_Adaptive\_Recovery\_and\_Mixture\_Identification.tex}.
\bibitem{p05:sv5:Lech}
C.~Lech, A note on recurring series,
\emph{Arkiv f\"or Matematik} \textbf{2} (1953), 417--421.
doi:10.1007/BF02590997.
\url{https://archive.ymsc.tsinghua.edu.cn/pacm_download/116/6892-11512_2007_Article_BF02590997.pdf}.
\bibitem{p05:sv5:BellLagarias}
J.~P.~Bell and J.~C.~Lagarias,
A Skolem--Mahler--Lech theorem for iterated automorphisms of $K$-algebras,
\emph{Canadian Journal of Mathematics} \textbf{67} (2015), 286--314.
doi:10.4153/CJM-2013-048-3.
Theorems~1.3 and~1.6.
\bibitem{p05:sv5:WuYang}
Y.~Wu and P.~Yang,
Optimal estimation of Gaussian mixtures via denoised method of moments,
\emph{Annals of Statistics} \textbf{48} (2020), 1981--2007.
doi:10.1214/19-AOS1873.
Full version: \url{https://arxiv.org/abs/1807.07237};
lemma and proposition numbers cited here refer to version~2.
\bibitem{p05:sv5:ExponentialSurvey}
J.~Heittokangas, K.~Ishizaki, K.~Tohge, and Z.-T.~Wen,
Value distribution of exponential polynomials and their role in the
theories of complex differential equations and oscillation theory,
\emph{Bulletin of the London Mathematical Society}
\textbf{55} (2023), 1--77.
doi:10.1112/blms.12719.
\url{https://arxiv.org/abs/2202.01574}.
\bibitem{p05:sv5:Skolem2026}
P.~Bacik, J.~Ouaknine, D.~Purser, and J.~Worrell,
On the $p$-adic Skolem problem,
in \emph{43rd International Symposium on Theoretical Aspects of
Computer Science}, LIPIcs \textbf{364} (2026), 8:1--8:20.
doi:10.4230/LIPIcs.STACS.2026.8.
\url{https://drops.dagstuhl.de/entities/document/10.4230/LIPIcs.STACS.2026.8}.
\bibitem{p05:ms5:Roy}
D.~Roy, Matrices whose coefficients are linear forms in logarithms,
\emph{Journal of Number Theory} \textbf{41} (1992), 22--47.
\url{https://doi.org/10.1016/0022-314X(92)90081-Y}.
\bibitem{p05:ms5:Diaz}
G.~Diaz, Produits et quotients de combinaisons lin\'eaires de logarithmes
de nombres alg\'ebriques: conjectures et r\'esultats partiels,
\emph{Journal de Th\'eorie des Nombres de Bordeaux}
\textbf{19} (2007), 373--391.
\url{https://doi.org/10.5802/jtnb.592}.
\bibitem{p05:ms5:Waldschmidt}
M.~Waldschmidt, \emph{The Four Exponentials Problem and the Schanuel
Conjecture}, author manuscript, May 26, 2022.
\url{https://webusers.imj-prg.fr/~michel.waldschmidt/articles/pdf/FourExponentialsSchanuel.pdf}.
\bibitem{p05:ms5:AGKS}
X.~Allamigeon, S.~Gaubert, R.~D.~Katz, and M.~Skomra,
Universal complexity bounds based on value iteration and application
to entropy games, in \emph{49th International Colloquium on Automata,
Languages, and Programming}, LIPIcs \textbf{229} (2022), 110:1--110:20.
\url{https://doi.org/10.4230/LIPIcs.ICALP.2022.110}.
\bibitem{p05:ms5:Boche}
H.~Boche, R.~F.~Schaefer, and H.~V.~Poor,
\emph{Algorithmic Computability and Approximability of
Capacity-Achieving Input Distributions}, 2022, version 2.
\url{https://arxiv.org/abs/2202.12617}.
\bibitem{p05:ar21:RivoalRoques}
T. Rivoal and J. Roques, \emph{$E$-functions of order 2 and units of
$E$-values}, preprint dated 6 November 2015, 16 pp.
\href{https://rivoal.perso.math.cnrs.fr/articles/unitse.pdf}{Author's PDF}.
\bibitem{p05:be21:Baricz}
\'{A}.~Baricz, D.~Prabhu K, S.~Singh, and A.~Vijesh V,
Infinitely divisible modified Bessel distributions,
\emph{Pacific Journal of Mathematics} \textbf{343} (2026), 261--313.
\url{https://doi.org/10.2140/pjm.2026.343.261}.
The version consulted for the external problem statements is
\url{https://arxiv.org/abs/2406.17721v2}, Section~4.
\bibitem{p05:be21:Capacity}
Keio, \emph{Observation and Recovery: Geometry, Mixtures, and Spectral Inversion},
unified post-223-page manuscript, October 2026, 447 pp.,
theorems ``The unrestricted factor region'' and
``The exact quartic cost of imposing ID,'' and lemma
``Uniform positivity at both ends.''
Original component: \texttt{research\_20261016\_round15/bessel.tex},
labels \texttt{be15:region}, \texttt{be15:quartic}, and
\texttt{be15:ends}.
AI-assisted research manuscript, not externally peer reviewed.
\bibitem{p05:be21:Endpoint}
Keio, ``Simultaneous gamma factors and reconstruction from
boundary response,'' in the unified post-223-page manuscript,
October 2026, lemma ``The scalar ingredient and its exact slope.''
Original component: \texttt{research\_20261016\_round16/bessel.tex},
label \texttt{be16:scalar}.
AI-assisted research manuscript, not externally peer reviewed.
\bibitem{p05:be21:Fractional}
Keio, ``Fractional extraction: exact contact geometry and
a sharp arithmetic rank barrier,'' in \emph{Arithmetic Rigidity of
Optimization and Special Functions}, 17 October 2026,
theorems ``Every real-shape contact is exactly double''
and ``Regularity and a universal slope inequality,'' and lemma
``Parameter derivatives at a contact.''
The labels in that manuscript are \texttt{be20:double},
\texttt{be20:frontier}, and \texttt{be20:response}.
AI-assisted research manuscript, not externally peer reviewed.
\bibitem{p05:in21:Tsukamoto}
M.~Tsukamoto, \emph{On an analogue of the Hurewicz theorem for mean
dimension}, preprint, 2022. \url{https://arxiv.org/abs/2206.03005}.
The introductory classical Hurewicz theorem is the input used here.
\bibitem{p05:in21:Kato}
H.~Kato, \emph{The nonexistence of expansive homeomorphisms on
hereditarily indecomposable compacta}, RIMS K\^oky\^uroku
\textbf{1578} (2008), 36--39, Theorem~2.3.
\url{https://www.kurims.kyoto-u.ac.jp/~kyodo/kokyuroku/contents/pdf/1578-08.pdf}.
\bibitem{p05:in21:Kechris}
A.~S.~Kechris, \emph{Classical Descriptive Set Theory}, Graduate Texts
in Mathematics \textbf{156}, Springer, 1995, Corollary~15.2.
\url{https://doi.org/10.1007/978-1-4612-4190-4}.
\bibitem{p05:in21:Stacks}
The Stacks Project Authors, \emph{Cohomology of Sheaves},
Section~20.13 (Leray spectral sequence) and Theorem~20.18.2,
Tag~09V6 (proper base change in topology).
\url{https://stacks.math.columbia.edu/tag/09V6}.
\bibitem{p05:in21:Lurie}
J.~Lurie, \emph{Higher Topos Theory}, Annals of Mathematics Studies
\textbf{170}, Princeton University Press, 2009; author's revised
version, 2012, Theorem~7.2.3.6 and Corollary~7.2.2.30.
\url{https://people.math.harvard.edu/~lurie/papers/highertopoi.pdf}.
\bibitem{p05:in21:KNU}
D.~Krieg, E.~Novak, and M.~Ullrich,
\emph{How many continuous measurements are needed to learn a vector?},
preprint, version~2, 22 September 2026.
\url{https://arxiv.org/abs/2412.06468}.
\bibitem{p05:in21:KishimotoMatsushita}
D.~Kishimoto and T.~Matsushita,
\emph{Van Kampen--Flores theorem and Stiefel--Whitney classes},
preprint, version~2, 2023.
\url{https://arxiv.org/abs/2304.08043}.
\bibitem{p05:in21:PriceWoodruff}
E.~Price and D.~P.~Woodruff, \emph{Lower bounds for adaptive sparse
recovery}, Proceedings of SODA 2013, 652--663.
\url{https://arxiv.org/abs/1205.3518}.
\bibitem{p06:lmod:BPR}
S.~Basu, R.~Pollack and M.-F.~Roy,
\emph{Algorithms in Real Algebraic Geometry}, second edition,
Springer, 2006. \url{https://doi.org/10.1007/3-540-33099-2}.
\bibitem{p06:lmod:Robinson}
J.~Robinson, Definability and decision problems
in arithmetic, \emph{Journal of Symbolic Logic} \textbf{14} (1949),
98--114. \url{https://doi.org/10.2307/2266510}.
Original text: \url{https://www.math.umd.edu/~laskow/713/Spring17/juliarobinson.pdf}.
\bibitem{p06:lmod:PillayTanovic}
A.~Pillay and P.~Tanovi\'c,
The number of countable models of first-order theories,
arXiv:2508.06854v2, 15 August 2026.
\url{https://arxiv.org/html/2508.06854v2}.
\bibitem{p06:lmod:Silver}
J.~H.~Silver, Counting the number of equivalence
classes of Borel and coanalytic equivalence relations,
\emph{Annals of Mathematical Logic} \textbf{18} (1980), 1--28.
\url{https://doi.org/10.1016/0003-4843(80)90002-9}.
\bibitem{p06:lmod:Krapp}
L.~S.~Krapp, Embedding the prime model of real
exponentiation into o-minimal exponential fields,
\emph{Bulletin of the London Mathematical Society} (2024).
\url{https://doi.org/10.1112/blms.12972}.
\bibitem{p06:lmod:ShelahProblems}
S.~Shelah, On what I do not understand
(and have something to say), model theory, arXiv:math/9910158v1,
1999, Problems 4.7 and 6.9.
\url{https://arxiv.org/pdf/math/9910158}.
\bibitem{p06:lmod:VaseyII}
S.~Vasey, Shelah's eventual categoricity
conjecture in universal classes. Part II,
\emph{Selecta Mathematica} \textbf{23} (2017), 1469--1506.
\url{https://doi.org/10.1007/s00029-016-0296-0};
\url{https://arxiv.org/abs/1602.02633}.
\bibitem{p06:lmod:OpenAIAEC}
OpenAI, \emph{Eventual categoricity for
abstract elementary classes}, primary preprint, 24 September 2026,
Theorem 1.2; proof claim, not independently certified in this report.
\url{https://github.com/openai/math/blob/main/preprints/Eventual-Categoricity-for-Abstract-Elementary-Classes-September-24-2026/paper.pdf}.
\bibitem{p06:lmod:OpenAIThreshold}
OpenAI, \emph{A CH obstruction to a
prescribed categoricity threshold}, primary preprint, 24 September 2026,
Theorem 1.1; proof claim, not independently certified in this report.
\url{https://github.com/openai/math/blob/main/preprints/A-CH-Obstruction-to-a-Prescribed-Categoricity-Threshold-September-24-2026/paper.pdf}.
\bibitem{p06:lmod:Espindola}
C.~Esp\'indola, A proof of Shelah's eventual
categoricity conjecture and an extension to accessible categories with
directed colimits, arXiv:1906.09169, 2019, revised 2022;
proof claim, not independently certified in this report.
\url{https://arxiv.org/abs/1906.09169}.
\bibitem{p06:lmod:JohnsonICM}
W.~Johnson, Henselianity and NIP fields,
\emph{Proceedings of the International Congress of Mathematicians
2026}, vol.~3, 97--112, Conjecture 4.2.
\url{https://doi.org/10.1137/25M1805412}.
\bibitem{p06:lmod:JohnsonDP}
W.~Johnson, Dp-finite fields VI: the
dp-finite Shelah conjecture, arXiv:2005.13989, 2020.
\url{https://arxiv.org/abs/2005.13989}.
\bibitem{p06:lmod:StableFields}
W.~Johnson, C.-M.~Tran, E.~Walsberg and
J.~Ye, The \'etale-open topology and the stable fields conjecture,
\emph{Journal of the European Mathematical Society} \textbf{26}
(2024), 4033--4070. \url{https://arxiv.org/abs/2009.02319}.
\bibitem{p06:lmod:BuildingTrees}
E.~Kaplan, A.~Matthews and E.~Walsberg,
Building trees in large fields, arXiv:2607.16942, 18 July 2026.
\url{https://arxiv.org/abs/2607.16942}.
\bibitem{p06:lmod:BaysKirby}
M.~Bays and J.~Kirby, Pseudo-exponential
maps, variants, and quasiminimality,
\emph{Algebra \& Number Theory} \textbf{12} (2018), 493--549.
\url{https://doi.org/10.2140/ant.2018.12.493};
\url{https://arxiv.org/abs/1512.04262}.
\bibitem{p06:lmod:ComplexPowers}
F.~Gallinaro and J.~Kirby,
Quasiminimality of complex powers,
\emph{Forum of Mathematics, Sigma} \textbf{12} (2024), e123.
\url{https://doi.org/10.1017/fms.2024.82};
\url{https://arxiv.org/abs/2304.06450}.
\bibitem{p06:lmod:Tent}
K.~Tent, From the Cherlin--Zilber Conjecture via
sharply 2-transitive groups to the Burnside problem,
arXiv:2606.18207, 16 June 2026.
\url{https://arxiv.org/abs/2606.18207}.
\bibitem{p06:lmod:HHL}
B.~Hart, E.~Hrushovski and M.~C.~Laskowski,
The uncountable spectra of countable theories,
\emph{Annals of Mathematics} \textbf{152} (2000), 207--257.
\url{https://annals.math.princeton.edu/articles/11619}.
\bibitem{p06:lmod:HrushovskiKueker}
E.~Hrushovski, Kueker's conjecture
for stable theories, \emph{Journal of Symbolic Logic} \textbf{54}
(1989), 207--220. \url{https://doi.org/10.2307/2275025}.
\bibitem{p06:lmod:TanovicKueker}
P.~Tanovi\'c, On Kueker's conjecture,
\emph{Journal of Symbolic Logic} \textbf{77} (2012), 1245--1256.
\url{https://doi.org/10.2178/jsl.7704100}.
\bibitem{p06:lmod:StableForking}
J.~Freitag and S.~Mutchnik,
A counterexample to the stable forking conjecture,
arXiv:2609.00436v1, 31 August 2026, Theorem 1.3 and Question 4.7;
counterexample claim, not independently fully checked in this report.
\url{https://arxiv.org/html/2609.00436v1}.
\bibitem{p06:lset:Gilson}
F. Gilson, \emph{A countable-support symmetric iteration separating PP from AC},
arXiv:2601.01855v8, 10 September 2026. Major correction withdrawing the final
construction; abstract and correction notice.
\url{https://arxiv.org/abs/2601.01855v8}.
\bibitem{p06:lset:Cohen}
P. J. Cohen, \emph{The independence of the continuum hypothesis},
Proceedings of the National Academy of Sciences USA 50 (1963), 1143--1148.
\url{https://doi.org/10.1073/pnas.50.6.1143}.
\bibitem{p06:lset:HanDoorn}
J. M. Han and F. van Doorn, \emph{A formal proof of the independence of the
continuum hypothesis}, arXiv:2102.02901 (2021).
\url{https://arxiv.org/abs/2102.02901}.
\bibitem{p06:lset:Woodin2026}
W. H. Woodin, \emph{The AD$^+$ duality program, the HOD conjecture, and the
Ultimate-$L$ conjecture}, Japanese Journal of Mathematics 21 (2026), 1--90.
\url{https://doi.org/10.1007/s11537-026-2541-4}.
\bibitem{p06:lset:ShelahPCF}
S. Shelah, \emph{Cardinal arithmetic for skeptics}, Bulletin of the American
Mathematical Society 26 (1992), 197--210, especially \S6.
\url{https://arxiv.org/abs/math/9201251}.
\bibitem{p06:lset:GitikPCF}
M. Gitik, \emph{PCF of a countable set can be uncountable}, author manuscript,
3 March 2026, revised version of \emph{Short extenders forcings II};
\S1, Lemma 7.12 and \S8, Question 3.
\url{https://www.math.tau.ac.il/~gitik/short%20extenders%20forcing%202-2026.pdf}.
\bibitem{p06:lset:LarsonOmega}
P. B. Larson, \emph{Three days of $\Omega$-logic}, lecture notes based on
November 2009 lectures; Definition 6.1, Theorems 6.4 and 9.2.
\url{https://paulblarson.github.io/omega_logic.pdf}.
\bibitem{p06:lset:ABL}
J. P. Aguilera, J. Bagaria and P. L\"ucke, \emph{Large cardinals, structural
reflection, and the HOD Conjecture}, arXiv:2411.11568v2 (2025);
\S1.2, Theorem B, Definition 3.13 and Corollary 3.15.
\url{https://arxiv.org/html/2411.11568v2}.
\bibitem{p06:lset:HKP}
J. D. Hamkins, G. Kirmayer and N. L. Perlmutter, \emph{Generalizations of
the Kunen inconsistency}, arXiv:1106.1951v2 (2012).
\url{https://arxiv.org/abs/1106.1951}.
\bibitem{p06:lset:McCallum}
R. McCallum, \emph{Inconsistency of Reinhardt cardinals with ZF},
arXiv:2601.10231v3, withdrawn 23 August 2026; withdrawal notice identifies
the continuity-of-Skolem-functions issue.
\url{https://arxiv.org/abs/2601.10231}.
\bibitem{p06:lset:Mildenberger}
H. Mildenberger, \emph{A solution to Roitman's problem}, arXiv:1404.7343,
withdrawn, final revision 20 February 2015.
\url{https://arxiv.org/abs/1404.7343}.
\bibitem{p06:lset:Pelayo}
J. de Jes\'us Pelayo-G\'omez, \emph{$\mathfrak d=\omega_1$ implies
$\mathfrak a=\omega_1$}, arXiv:2610.00639v1, 30 September 2026.
Theorems 1.1 and 1.4; Lemmas 3.1--3.3. Recent author proof claim, reviewed
locally here, not treated as a new result of this manuscript.
\url{https://arxiv.org/pdf/2610.00639v1}.
\bibitem{p06:lset:Trang}
N. Trang, \emph{PFA and guessing models}, Israel Journal of Mathematics 215
(2016); Theorems 0.2 and 0.4 of arXiv:1608.05726.
\url{https://arxiv.org/abs/1608.05726}.
\url{https://doi.org/10.1007/s11856-016-1390-x}.
\bibitem{p06:lset:VialeWeiss}
M. Viale and C. Wei\ss, \emph{On the consistency strength of the proper
forcing axiom}, Advances in Mathematics 228 (2011), 2672--2687;
Theorem 6.4 and Corollary 6.11 in arXiv:1012.2046.
\url{https://arxiv.org/html/1012.2046v1}.
\bibitem{p06:lset:ShelahBeth}
S. Shelah, \emph{Remarks in abstract model theory}, Annals of Pure and
Applied Logic 29 (1985), 255--288; Conclusion 6.11 and Claim 6.12.
\url{https://shelah.logic.at/files/95416/199.pdf}.
\bibitem{p06:lset:Vaananen2025}
J. V\"a\"an\"anen, \emph{Interpolation in model theory}, contemporary
author chapter, retrieved 16 October 2026; Proposition 14 and ensuing
discussion, and the Lindstr\"om theorem section.
\url{https://cibd.bitbucket.io/taci/chapter_07_modeltheory.pdf}.
\bibitem{p06:ldes:LutzSiskind}
P. Lutz and B. Siskind, \emph{Part 1 of Martin's Conjecture for
order-preserving and measure-preserving functions}, arXiv:2305.19646v3
(2024), Theorems 1.1, 4.8 and 4.13.
\url{https://arxiv.org/html/2305.19646v3}.
\bibitem{p06:ldes:Yu}
L. Yu, \emph{Order-preserving maps of the Turing degrees: regularity
and consistency strength}, arXiv:2610.14548v1 (10 October 2026),
Theorems 1.2 and 1.4; preprint.
\url{https://arxiv.org/html/2610.14548v1}.
\bibitem{p06:ldes:NaryshkinVaccaro}
P. Naryshkin and A. Vaccaro, \emph{Hyper-$u$-amenability and
hyperfiniteness of treeable equivalence relations}, Forum of
Mathematics, Sigma \textbf{14} (2026), e29, Theorem 2 and Corollary 3.
\url{https://doi.org/10.1017/fms.2026.10171}.
\bibitem{p06:ldes:ConleyMarks}
C. T. Conley and A. S. Marks, \emph{Hyperfiniteness and Borel
combinatorics}, arXiv:1611.02204, Introduction and Theorem 1.7.
\url{https://arxiv.org/html/1611.02204v1}.
\bibitem{p06:ldes:BeckerKechris}
H. Becker and A. S. Kechris, \emph{Borel actions of Polish groups},
arXiv:math/9304208 (1993).
\url{https://arxiv.org/abs/math/9304208}.
\bibitem{p06:ldes:Melleray}
J. Melleray, \emph{Some aspects of topological dynamics of Polish
groups (with an introduction to descriptive set theory)},
arXiv:2602.23799 (2026), Theorems 15.7 and 15.10 and discussion following
Theorem 15.10.
\url{https://arxiv.org/pdf/2602.23799}.
\bibitem{p06:ldes:Cameron}
P. J. Cameron, \emph{Around homogeneity}, arXiv:2602.22181v1
(2026), Section 8 and Question 8.1.
\url{https://arxiv.org/html/2602.22181v1}.
\bibitem{p06:ldes:BodorCameronSzabo}
B. Bodor, P. J. Cameron and C. Szab\'o, \emph{Infinitely many reducts
of homogeneous structures}, arXiv:1609.07694 (2016).
\url{https://arxiv.org/abs/1609.07694}.
\bibitem{p06:ldes:Leung}
S. Leung, \emph{Categoricity transfer for short AECs with amalgamation
over sets}, arXiv:2203.08956, Conjecture 1.2 and the main categoricity
transfer theorem.
\url{https://arxiv.org/pdf/2203.08956}.
\bibitem{p06:ldes:ChernikovShelah}
A. Chernikov and S. Shelah, \emph{On the number of Dedekind cuts and
two-cardinal models of dependent theories}, arXiv:1308.3099v3 (2014),
Section 4 and Theorem 4.5.
\url{https://arxiv.org/html/1308.3099v3}.
\bibitem{p06:ldes:DayMarks}
A. R. Day and A. S. Marks, \emph{On a question of Slaman and Steel},
Transactions of the American Mathematical Society, Series B
\textbf{13} (2026), article 11, Section 3.
\url{https://math.berkeley.edu/~marks/papers/mc_v17.pdf}.
\bibitem{p06:ldes:CoskeySchneider}
S. Coskey and S. Schneider, \emph{Cardinal characteristics and
countable Borel equivalence relations}, arXiv:1103.2312v4 (2016),
Definitions 1.2 and 3.3, Theorem 1.3 and Question 1.4.
\url{https://arxiv.org/html/1103.2312v4}.
\bibitem{p06:ldes:ArantKechrisLutz}
T. Arant, A. S. Kechris and P. Lutz, \emph{Borel graphable equivalence
relations}, Advances in Mathematics \textbf{487} (2026), 110765.
We use the theorem numbering of arXiv:2409.08624v2, especially
Question 3 and Theorem 65.
\url{https://arxiv.org/html/2409.08624v2}.
\url{https://doi.org/10.1016/j.aim.2025.110765}.
\bibitem{p06:ldes:FrischShinkoVidnyanszky}
J. Frisch, F. Shinko and Z. Vidny\'anszky,
\emph{Hyper-hyperfiniteness and complexity}, arXiv:2409.16445v2
(21 September 2026), Question 0.2 and Theorem 0.3; preprint.
\url{https://arxiv.org/html/2409.16445v2}.
\bibitem{p06:lcomp:SW}
T. A. Slaman and W. H. Woodin,
\emph{Definability in Degree Structures}, author manuscript (2005),
Theorem 7.5.3 and Conjecture 7.5.4.
\url{https://math.berkeley.edu/~slaman/talks/sw.pdf}.
\bibitem{p06:lcomp:SS}
T. A. Slaman and M. I. Soskova,
\emph{The $\Delta^0_2$ Turing degrees: automorphisms and definability},
author preprint, introduction and Corollary 4.2; accessed 17 October 2026.
\url{https://people.math.wisc.edu/~soskova/preprints/autdef.pdf}.
\bibitem{p06:lcomp:Merkle}
W. Merkle, J. S. Miller, A. Nies, J. Reimann, and F. Stephan,
\emph{Kolmogorov--Loveland randomness and stochasticity},
Annals of Pure and Applied Logic \textbf{138} (2006), 183--210.
\url{https://doi.org/10.1016/j.apal.2005.06.011}.
\bibitem{p06:lcomp:Petrovic}
T. Petrovi\'c,
\emph{How far is Kolmogorov--Loveland randomness from Martin--L\"of randomness?},
arXiv:2610.13585v1 (9 October 2026), Sections 2.6, 2.8 and 5,
Theorem 4.1.
\url{https://arxiv.org/html/2610.13585v1}.
\bibitem{p06:lcomp:Piazza}
M. Piazza, \emph{A finite algebraic target theorem for Kreisel's conjecture},
Archive for Mathematical Logic \textbf{65} (2026), 905--924,
published 29 September 2026, Sections 1 and 5.
\url{https://doi.org/10.1007/s00153-026-01029-z}.
\bibitem{p06:lcomp:Hrubes}
P. Hrube\v{s}, \emph{Theories very close to PA where Kreisel's conjecture is false},
Journal of Symbolic Logic \textbf{72} (2007), 123--137,
Introduction, Theorem 3, and Corollary 1.
\url{https://users.math.cas.cz/~hrubes/PDFs/KCMinus-JSL.pdf}.
\bibitem{p06:lcomp:SK}
P. G. Santos and R. Kahle,
\emph{Variants of Kreisel's conjecture on a new notion of provability},
Bulletin of Symbolic Logic \textbf{27} (2021), 337--350, Theorem 2.1.
\url{https://doi.org/10.1017/bsl.2021.68}.
\bibitem{p06:lcomp:HL}
K. Higuchi and P. Lutz,
\emph{A note on a conjecture of Sacks: it is harder to embed height three
partial orders than height two partial orders},
arXiv:2309.01876v2 (2023), Theorems 1.1, 1.2, and 1.7.
\url{https://arxiv.org/html/2309.01876v2}.
\bibitem{p06:lcomp:MG}
J. S. Miller and N. Greenberg,
\emph{An embedding into the Turing degrees},
author manuscript dated 8 November 2024, Theorem 1.1 and Proposition 1.2.
\url{https://homepages.ecs.vuw.ac.nz/~greenberg/Papers/88-sigma-ind.pdf}.
\bibitem{p06:lcomp:SkolemSurvey}
P. Bacik, T. Karimov, F. Luca, J. Nieuwveld, J. Ouaknine,
D. Purser, and J. Worrell,
\emph{A survey of the Skolem and Positivity Problems for linear recurrence sequences},
author manuscript (2026), 59 pp., Theorem 25 and Sections 5.3, 7, and 9.
\url{https://people.mpi-sws.org/~joel/publications/skolem_and_positivity_survey26.pdf}.
\bibitem{p06:lcomp:Manthe}
S. Manthe, \emph{The Borel monadic theory of order is decidable},
arXiv:2410.00887v2 (9 March 2026; first version 2024), Theorem 1.1.
\url{https://arxiv.org/html/2410.00887v2}.
\bibitem{p06:lcomp:ShelahOrder}
S. Shelah, \emph{The monadic theory of order},
Annals of Mathematics \textbf{102} (1975), 379--419;
author-posted version arXiv:2305.00968v1 (2023), Conjecture 0.3.
\url{https://arxiv.org/html/2305.00968v1}.
\bibitem{p06:lcomp:Neeman}
I. Neeman, \emph{Monadic theories of wellorders},
author manuscript dated 24 February 2008, Theorems 11--16 and
final discussion, pp. 9--14.
\url{https://www.math.ucla.edu/~ineeman/lmps.pdf/}.
\bibitem{p06:lcomp:AF}
S. Anscombe and A. Fehm,
\emph{Universal-existential theories of fields},
Proceedings of the Edinburgh Mathematical Society (2025),
doi:10.1017/S0013091525100771;
arXiv:2405.12771v4 (11 January 2026), Introduction and Theorems 1.1--1.2.
\url{https://arxiv.org/pdf/2405.12771}.
\bibitem{p06:lcomp:Gitin}
L. Gitin,
\emph{Undecidability of expansions of Laurent series fields by cyclic discrete subgroups},
arXiv:2408.13900v1 (2024), Theorem 9.
\url{https://arxiv.org/html/2408.13900v1}.
\bibitem{p06:lcomp:Wang}
W. Wang, \emph{On non-standard models of arithmetic with uncountable standard systems},
arXiv:2007.05885v1 (12 July 2020),
Theorems 1.2--1.3, Question 1.4, and Theorem 2.1.
\url{https://arxiv.org/html/2007.05885v1}.
\bibitem{p06:lcomp:Gitman}
V. Gitman, \emph{Scott's problem for proper Scott sets},
arXiv:0801.4364v1 (2008), Theorems 1.2--1.3 and 1.7, Corollary 1.4.
\url{https://arxiv.org/html/0801.4364v1}.
\bibitem{p06:lcomp:LNS}
S. Lempp, A. Nies, and T. A. Slaman,
\emph{The $\Pi_3$-theory of the computably enumerable Turing degrees is undecidable},
Transactions of the American Mathematical Society \textbf{350} (1998),
2719--2736, Introduction and main theorem.
\url{https://people.math.wisc.edu/~slempp/papers/pi3.pdf}.
\bibitem{p06:lcomp:Soskova}
M. I. Soskova, \emph{Fragments of the theory of the enumeration degrees},
Online Logic Seminar slides (23 April 2020), slides 2--4.
\url{https://people.math.wisc.edu/~soskova/talks/OnlineLogic2020.pdf}.
\bibitem{p06:lproof:FPS}
Y. Filmus, T. Pitassi and R. Santhanam,
\emph{Exponential Lower Bounds for $\mathrm{AC}^0$-Frege Imply Superpolynomial
Frege Lower Bounds}, ICALP 2011. Author manuscript, Theorem 2 and Corollary 1.
\url{https://yuvalfilmus.cs.technion.ac.il/Papers/FPS-ICALP.pdf}.
\bibitem{p06:lproof:REGR}
S. F. de Rezende, D. Engstr\"om, Y. Ghannane and K. Risse,
\emph{Superpolynomial Length Lower Bounds for Tree-Like Semantic Proof Systems
with Bounded Line Size}, preprint, 30 April 2026.
\url{https://arxiv.org/abs/2604.28172}.
\bibitem{p06:lproof:Freund}
A. Freund, \emph{Fra\"iss\'e's conjecture, partial impredicativity and
well-ordering principles, part I}, preprint, 19 June 2024,
Theorem 1.2 and introduction.
\url{https://arxiv.org/abs/2406.13485}.
\bibitem{p06:lproof:LLPY}
Q. Le Hou\'erou, L. Levy Patey and K. Yokoyama,
\emph{$\Pi^0_4$ conservation of Ramsey's theorem for pairs},
2024; corrected version 3, 6 May 2026,
Question 1.1 and Theorems 1.2, 1.4.
\url{https://arxiv.org/html/2404.18974v3}.
\bibitem{p06:lproof:HwangClaim}
H. Yu-Shiang Hwang,
\emph{$\Pi^1_1$ conservation of Ramsey's theorem for pairs over
$\mathsf B\Sigma^0_2$}, author-posted manuscript, 3 October 2026,
DOI 10.13140/RG.2.2.26676.41608.
Full-resolution claim not independently verified in this audit.
\url{https://www.researchgate.net/publication/415180723_P_conservation_of_Ramsey%27s_theorem_for_pairs_over_BS}.
\bibitem{p06:lproof:Csima}
B. F. Csima, D. D. Dzhafarov, D. R. Hirschfeldt, C. G. Jockusch Jr.,
R. Solomon and L. Brown Westrick,
\emph{The reverse mathematics of Hindman's theorem for sums of exactly two
elements}, preprint, version 3, 29 May 2018, introduction and Theorem 4.4.
\url{https://arxiv.org/abs/1804.09809}.
\bibitem{p06:lproof:LiuPatey}
L. Liu and L. Patey,
\emph{Hindman's theorem does not code $\varnothing^{(\omega)}$ in one
application}, preprint, 20 July 2026, Main Theorem 1.5 and following discussion.
\url{https://arxiv.org/html/2607.17666}.
\bibitem{p06:lproof:Liao}
Y. Liao, \emph{A recursive coloring function without $\Pi^0_3$ solutions for
Hindman's theorem}, Journal of Symbolic Logic \textbf{91}(2) (2026), 840--863,
DOI 10.1017/jsl.2024.49.  Also stated as Theorem 1.4 in
\cite{p06:lproof:LiuPatey}.
\url{https://doi.org/10.1017/jsl.2024.49}.
\bibitem{p06:lproof:Rathjen95}
M. Rathjen, \emph{Recent Advances in Ordinal Analysis:
$\Pi^1_2$--CA and Related Systems}, Bulletin of Symbolic Logic
\textbf{1}(4) (1995), 468--485.
\url{https://www.cambridge.org/core/journals/bulletin-of-symbolic-logic/article/recent-advances-in-ordinal-analysis-1-2-ca-and-related-systems/6ED09BFF442E1CB0BDF0ABD1292ED44E}.
\bibitem{p06:lproof:Rathjen05}
M. Rathjen, \emph{An ordinal analysis of parameter free
$\Pi^1_2$-comprehension}, Archive for Mathematical Logic \textbf{44} (2005),
263--362, DOI 10.1007/s00153-004-0232-4.
\url{https://link.springer.com/article/10.1007/s00153-004-0232-4}.
\bibitem{p06:lproof:Arai}
T. Arai, \emph{An ordinal analysis of $\Pi_N$-Collection},
preprint, version 4, 12 August 2025, Theorems 1.1--1.4.
\url{https://arxiv.org/abs/2311.12459v4}.
\bibitem{p06:lproof:Friedman99}
H. Friedman, \emph{Grand conjectures}, FOM posting, 16 April 1999,
Conjecture 1 (EFA) and Conjecture 2 (categorical methods).
\url{https://fomarchive.ugent.be/1999-April/003014.html}.
\bibitem{p06:lproof:Friedman10}
H. Friedman, author's answer to \emph{Status of Harvey Friedman's grand
conjecture}, MathOverflow, 2010.
\url{https://mathoverflow.net/questions/39452/status-of-harvey-friedmans-grand-conjecture/39559}.
\bibitem{p06:lproof:BGK}
M. H. S\o rensen, \emph{Problem 9: Is every weakly normalizing PTS also
strongly normalizing?}, TLCA list of open problems,
entry submitted 1993, revised 21 July 2014.
\url{https://tlca.di.unito.it/opltlca/opltlcasu16.html}.
\bibitem{p06:lproof:BHS}
G. Barthe, J. Hatcliff and M. H. S\o rensen,
\emph{Weak normalization implies strong normalization in a class of
non-dependent pure type systems}, Theoretical Computer Science
\textbf{269}(1--2) (2001), 317--361.
\url{https://www.sciencedirect.com/science/article/pii/S0304397501000123}.
\bibitem{p06:lproof:Poll}
E. Poll, \emph{Expansion postponement for normalising pure type systems},
Journal of Functional Programming \textbf{8}(1) (1998), 89--96,
Theorems 8 and 12, Corollary 13, DOI 10.1017/S095679689700292X.
\url{https://kar.kent.ac.uk/21695/1/expansion_postponement_for_normalising_poll.pdf}.
\bibitem{p06:lproof:KR}
M. Krombholz and M. Rathjen, \emph{Upper bounds on the graph minor theorem},
preprint, 30 June 2019; published in 2020,
Sections 1--3 and Theorems 5 and 12.
\url{https://arxiv.org/html/1907.00412v1}.
\bibitem{p06:lfin:css}
G. Cherlin, S. Shelah, and N. Shi,
\emph{Universal graphs with forbidden subgraphs and algebraic closure},
Advances in Applied Mathematics \textbf{22} (1999), 454--491.
\url{https://arxiv.org/abs/math/9809202}.
\bibitem{p06:lfin:tree}
G. Cherlin and S. Shelah,
\emph{Universal graphs with a forbidden subtree},
Journal of Combinatorial Theory, Series B \textbf{97} (2007), 293--333;
Theorem 1. \url{https://arxiv.org/abs/math/0512218}.
\bibitem{p06:lfin:shelah}
S. Shelah, \emph{Divide and conquer: dividing lines and universality},
Theoria \textbf{87} (2021), 259--348;
Definition 2.1 and Problems 2.2--2.3.
\url{https://shelah.logic.at/files/210561/theo.12289.pdf}.
\bibitem{p06:lfin:star2026}
OpenAI, \emph{Generalized Star Height at Most Three},
preprint dated 25 September 2026, Theorem 1.1, Section 1.1,
Lemma 10.1, and Proposition 10.2.
\url{https://github.com/openai/math/blob/main/preprints/Generalized-Star-Height-at-Most-Three-September-25-2026/article.pdf}.
\bibitem{p06:lfin:starlean}
OpenAI, \emph{Generalized star height at most three: Lean formalization scope},
2026, repository documentation and linked comparator statement;
accessed 17 October 2026. The comparator is a statement interface,
not by itself a proof artifact.
\url{https://github.com/openai/math/blob/main/lean/docs/134.md}.
\bibitem{p06:lfin:mt}
D. Macpherson and K. Tent, \emph{Omega-categorical pseudofinite groups},
Journal of Symbolic Logic, published online 15 January 2025,
Introduction, discussion after Theorem 1.4.
\url{https://doi.org/10.1017/jsl.2024.80}.
\bibitem{p06:lfin:ls}
S. Letzter and J. Sahasrabudhe,
\emph{On existentially complete triangle-free graphs},
2017, arXiv:1708.08817, Introduction and main nonexistence bound.
\url{https://arxiv.org/abs/1708.08817}.
\bibitem{p06:lfin:kk}
L. S. Krapp and S. Kuhlmann, \emph{Ordered transexponential fields},
Annals of Pure and Applied Logic \textbf{176} (2025), article 103541;
Question 2.3 and Theorem 3.6 in the arXiv version.
\url{https://arxiv.org/abs/2305.04607};
\url{https://doi.org/10.1016/j.apal.2024.103541}.
\bibitem{p06:lfin:palomar}
Palomar registry, \emph{A transexponential o-minimal structure},
indexed under MSC 03C64, accessed 17 October 2026.
Metadata only: source project and checked theorem were not retrieved
for this audit; no proof certification is inferred here.
\url{https://palomar-registry.org/subject?code=03C64&kind=msc}.
\bibitem{p06:lfin:podewski}
K. Krupi\'nski, P. Tanovi\'c, and F. O. Wagner,
\emph{Around Podewski's conjecture},
Fundamenta Mathematicae \textbf{222} (2013), 175--193;
Theorems 1 and 3, Lemma 2.1.
\url{https://arxiv.org/abs/1201.5709}.
\bibitem{p06:lfin:large}
W. Johnson, C.-M. Tran, E. Walsberg, and J. Ye,
\emph{Large implies henselian}, arXiv:2508.10886v2,
7 March 2026; Question 11.1 and Conjecture 11.3.
\url{https://arxiv.org/abs/2508.10886}.
\bibitem{p06:lfin:mell}
L. Stra\ss burger, \emph{On the decision problem for MELL},
Theoretical Computer Science \textbf{768} (2019), 91--98;
author report RR-9203 (2018), Theorem 4.4 and Section 5.
\url{https://www.lix.polytechnique.fr/~lutz/papers/RR-9203.pdf};
\url{https://doi.org/10.1016/j.tcs.2019.02.022}.
\bibitem{p06:lfin:bvass2026}
C. Bizi\`ere, J. Leroux, and G. Sutre,
\emph{Bridging the gap between plain VASS and branching VASS},
FOSSACS 2026; expanded version arXiv:2609.15869v1,
14 September 2026, Corollary 4.4, Remark 5.7, Section 6,
and Theorem 7.5.
\url{https://arxiv.org/abs/2609.15869};
\url{https://doi.org/10.1007/978-3-032-22730-0_4}.
\bibitem{p06:lfin:bvass2025}
C. Bizi\`ere, T. Hilaire, J. Leroux, and G. Sutre,
\emph{On the reachability problem for two-dimensional
branching VASS}, MFCS 2025, LIPIcs \textbf{345}, article 22.
\url{https://doi.org/10.4230/LIPIcs.MFCS.2025.22}.
\bibitem{p06:lfin:waldmann}
J. Waldmann, \emph{Normalization of S-terms is decidable},
RTA 1998, Lecture Notes in Computer Science \textbf{1379}, 138--150.
See also the author's thesis abstract, \emph{The Combinator S} (1998).
\url{https://www.imn.htwk-leipzig.de/~waldmann/pub/phdthesis/thesis.abs}.
\bibitem{p06:lfin:manzonetto}
G. Manzonetto, \emph{A lambda calculus satellite},
FSCD 2023, LIPIcs \textbf{260}, article 3,
Section 6.4, Problem 4.
\url{https://doi.org/10.4230/LIPIcs.FSCD.2023.3}.
\bibitem{p06:lfin:pcp}
T. Neary, \emph{Undecidability in binary tag systems and the Post
correspondence problem for five pairs of words}, STACS 2015,
LIPIcs \textbf{30}, 649--661; Definition 10 and Theorem 11.
\url{https://doi.org/10.4230/LIPIcs.STACS.2015.649}.
\bibitem{p07:rlog:BakerIII}
A.~Baker, Linear forms in the logarithms of
algebraic numbers (III), \emph{Mathematika} \textbf{14} (1967), 220--228.
\url{https://doi.org/10.1112/S0025579300003843}.
\bibitem{p07:rlog:BakerBook}
A.~Baker, \emph{Transcendental Number Theory},
Cambridge University Press, 1975, Chapter 2.
\url{https://doi.org/10.1017/CBO9780511565977}.
\bibitem{p07:rlog:Strzebonski}
A.~Strzebo\'nski, Real root isolation for
exp--log--arctan functions, \emph{Journal of Symbolic Computation}
\textbf{47} (2012), 282--314.
\url{https://doi.org/10.1016/j.jsc.2011.11.004}.
\bibitem{p07:rlog:BSP}
H.~Boche, R.~F.~Schaefer and H.~V.~Poor,
Algorithmic computability and approximability of capacity-achieving
input distributions, arXiv:2202.12617, 2022, revised 2023,
Appendix B, Theorem 4.
\url{https://arxiv.org/abs/2202.12617}; primary author/institution PDF:
\url{https://casa.rub.de/fileadmin/img/2_Forschung/2.1_Publikationen_PDFs/2022_Algorithmic_Computability_and_Approximability_of_Capacity-Achieving_Input_Distributions_Publication_ClusterofExcellence_CASA_Bochum.pdf}.
\bibitem{p07:rlog:Waldschmidt}
M.~Waldschmidt,
\emph{Introduction to transcendental numbers: exercises},
author's notes, updated 14 April 2021, p.~17.
\url{https://webusers.imj-prg.fr/~michel.waldschmidt/articles/pdf/TranscendenceIntroductionExercises.pdf}.
\bibitem{p07:rlog:BlakajWolf}
V.~Blakaj and M.~M.~Wolf,
Transcendental properties of entropy-constrained sets,
\emph{Annales Henri Poincar\'e} \textbf{24} (2023), 349--362,
published online 2022, Theorem 2.
\url{https://doi.org/10.1007/s00023-022-01227-4}.
\bibitem{p07:rlog:Yu}
R.~Yu, Tail criteria, no-go audits, and Ap\'ery-type
certificate obstructions for the irrationality of $e+\pi$,
arXiv:2606.17303, 15 June 2026.
\url{https://arxiv.org/abs/2606.17303}.
\bibitem{p07:rga:Sheth}
A.~Sheth, \emph{Real geometric transcendence for the Gamma function},
arXiv:2605.15117v1 (14 May 2026), Section 4.2, Conjecture 4.2
and its following paragraph.
\url{https://arxiv.org/pdf/2605.15117}.
\bibitem{p07:rga:Grinspan}
P.~Grinspan, \emph{Measures of simultaneous approximation for
quasi-periods of abelian varieties}, Journal of Number Theory
\textbf{94} (2002), 136--176.
\url{https://www.sciencedirect.com/science/article/pii/S0022314X01927337}.
\bibitem{p07:rga:Waldschmidt}
M.~Waldschmidt, \emph{An introduction to irrationality and transcendence
methods}, Arizona Winter School 2008, Lecture 3,
Section 3.3.2, pp. 90--91.
\url{https://webusers.imj-prg.fr/~michel.waldschmidt/articles/pdf/AWSLecture3.pdf}.
\bibitem{p07:rga:MPM}
V.~Mu\~noz and R.~P\'erez-Marco,
\emph{Stirling--Ramanujan constants are exponential periods},
Communications in Number Theory and Physics \textbf{18} (2024),
485--507, Section 1, Conjecture 1 and Corollary 3; distinction
between notions of period on pp. 491--492.
Also checked arXiv:2402.02660v4 (17 July 2026), Section 1.
\url{https://webusers.imj-prg.fr/~ricardo.perez-marco/publications/articles/Stirling-Ramanujan-constants.pdf}.
\bibitem{p07:rga:Pain}
J.-C.~Pain, \emph{Series representations for the logarithm of the
Glaisher--Kinkelin constant}, arXiv:2304.07629v1 (2023).
\url{https://arxiv.org/html/2304.07629v1}.
\bibitem{p07:rga:DLMF}
NIST Digital Library of Mathematical Functions,
\emph{Section 8.17: Incomplete beta functions},
Equations 8.17.1--8.17.4, accessed 17 October 2026.
\url{https://dlmf.nist.gov/8.17}.
\bibitem{p07:rga:BakerIII}
A.~Baker, \emph{Linear forms in the logarithms of algebraic numbers (III)},
Mathematika \textbf{14} (1967), 220--228.
\url{https://doi.org/10.1112/S0025579300003843}.
\bibitem{p07:rga:DhillonSaha}
S.~Dhillon and E.~Saha, \emph{On arithmetic nature of special values
of the incomplete beta function}, Acta Arithmetica \textbf{218}
(2025), 273--284, DOI 10.4064/aa240610-6-10.
Primary abstract and institutional bibliographic record checked;
full theorem statements not verified here.
\url{https://doi.org/10.4064/aa240610-6-10}.
\url{https://digitalcommons.isical.ac.in/journal-articles/5486/}.
\bibitem{p07:rga:SahaTalk}
E.~Saha, \emph{Arithmetic nature of the special values of the
incomplete beta function}, joint work with S.~Dhillon,
speed talk at Recurrence, Transcendence, and Diophantine Approximation,
14 July 2025, slides 1--4, compiled PDF pp.~219--223;
method and parameter ranges on slide 4/4, PDF p.~222.
\url{https://emresertoz.com/wp-content/uploads/2025/07/speed_talks.pdf}.
\bibitem{p07:rsc:DT}
D. Duverney and Y. Tachiya,
\emph{Refinement of the Chowla--Erd\H{o}s method and linear independence
of certain Lambert series}, author manuscript, 2019.
\url{https://danielduverney.fr/documents/theorie-des-nombres/DuverneyTachiya190522.pdf}.
\bibitem{p07:rsc:Campbell}
J. M. Campbell, \emph{On the binary digits of the Erd\H{o}s--Borwein
constant}, arXiv:2605.24160, 2026.
\url{https://arxiv.org/abs/2605.24160}.
\bibitem{p07:rsc:Waldschmidt}
M. Waldschmidt, \emph{Transcendental Number Theory: recent results
and open problems}, January 2022, slides 57--60.
\url{https://webusers.imj-prg.fr/~michel.waldschmidt/articles/pdf/TNTOpenPbs.pdf}.
\bibitem{p07:rsc:LZ}
A. Languasco and A. Zaccagnini,
\emph{Computing the Mertens and Meissel--Mertens constants for sums
over arithmetic progressions}, Experimental Mathematics
\textbf{19} (2010), 279--284; arXiv:0906.2132.
\url{https://arxiv.org/abs/0906.2132}.
\bibitem{p07:rsc:KL}
D. Keliher and E. S. Lee,
\emph{On the constants in Mertens' theorems for primes in arithmetic
progressions}, Integers \textbf{24A} (2024), A12.
\url{https://math.colgate.edu/~integers/a12Proc23/a12Proc23.pdf}.
\bibitem{p07:rsc:Adam}
B. Adamczewski, \emph{Mahler's method}, 2018 author version,
especially Theorem 4.1 and the rational--transcendental dichotomy
discussed immediately afterwards.
\url{https://adamczewski.perso.math.cnrs.fr/Mahler_selecta_new.pdf}.
\bibitem{p07:rsc:CDDM}
F. Chyzak, T. Dreyfus, P. Dumas and M. Mezzarobba,
\emph{Computing solutions of linear Mahler equations},
Mathematics of Computation \textbf{87} (2018), 2977--3021.
\url{https://arxiv.org/abs/1612.05518}.
\bibitem{p07:rsc:AZ}
C. E. Arreche and Y. Zhang,
\emph{Mahler discrete residues and summability for rational functions},
Proceedings of ISSAC 2022, 525--533.
\url{https://arxiv.org/abs/2202.09805}.
\bibitem{p07:rsc:Nguyen}
K. D. Nguyen,
\emph{Transcendental series of reciprocals of Fibonacci and Lucas
numbers}, arXiv:2009.02446, 2020, Theorems 1.2--1.3.
\url{https://arxiv.org/abs/2009.02446}.
\bibitem{p07:rsc:BS}
J. Bell and D. Smertnig,
\emph{Mahler series with multiplicative coefficient sequences},
arXiv:2603.23456, 2026.
\url{https://arxiv.org/abs/2603.23456}.
\bibitem{p07:rsc:Coons}
M. Coons, \emph{On the rational approximation of the sum of the
reciprocals of the Fermat numbers}, arXiv:1112.5072, 2011.
\url{https://arxiv.org/abs/1112.5072}.
\bibitem{p07:rgam:WV}
T. Wolfs and W. Van Assche,
\emph{Rational approximation of Euler's constant using multiple orthogonal polynomials},
Combinatorics and Number Theory \textbf{14} (2025), 141--162.
\href{https://doi.org/10.2140/cnt.2025.14.141}{doi:10.2140/cnt.2025.14.141};
\href{https://arxiv.org/html/2404.09799v3}{primary full text, Theorems 1.2--1.3}.
\bibitem{p07:rgam:MY}
N. Murabayashi and H. Yoshida,
\emph{Rational function approximations of the special function
$e^xE_1(x)$ and applications to irrationality of Euler--Gompertz constant
$\delta$}, preprint (2022).
\href{https://arxiv.org/abs/2210.06768}{arXiv:2210.06768}.
The proposed irrationality application retains an unproved sufficient
nonvanishing condition; it is not cited as an irrationality theorem.
\bibitem{p07:rgam:Powers}
M. R. Powers,
\emph{On the Uniqueness of Ein(1) among Linear Combinations of the
Euler--Mascheroni and Euler--Gompertz Constants}, preprint (2025).
\href{https://arxiv.org/abs/2510.00315}{arXiv:2510.00315}.
\bibitem{p07:rgam:FR}
S. Fischler and T. Rivoal,
\emph{The ring of values of $E$-functions is not a field},
preprint, arXiv:2610.16603, submitted 12 October 2026;
author PDF dated 14 October 2026, Proposition 3.1.
\href{https://arxiv.org/abs/2610.16603}{arXiv:2610.16603};
\href{https://www.imo.universite-paris-saclay.fr/~fischler/Ecroix.pdf}{author's PDF}.
The primary PDF retained from the earlier source audit was reread;
no arithmetic theorem about its coefficient constants is inferred.
\bibitem{p07:zrec:Fauzan}
A. Fauzan,
\emph{$\zeta(5)$ is irrational}, preprint, 17 September 2026,
Zenodo record 22826419.
\url{https://zenodo.org/records/22826419}.
\bibitem{p07:zrec:Romik}
D. Romik,
\emph{zeta5-irrationality: Lean formalization of the irrationality of
$\zeta(5)$}, public source repository, 2026; README,
\texttt{Zeta5/Axioms.lean}, \texttt{Zeta5/Basic.lean},
\texttt{Zeta5/Interface.lean}, \texttt{cert/Check.lean}, and
\texttt{docs/CERTIFICATION.md}, inspected 17 October 2026.
\url{https://github.com/danromik/zeta5-irrationality}.
The source interface was inspected; no complete kernel replay is claimed.
\bibitem{p07:zrec:Lai}
L. Lai,
\emph{The linear independence of $1$, $\zeta(3)$, $\zeta(4)$},
arXiv:2610.06838v2, 15 October 2026, especially Theorems 1.5 and 1.6.
\url{https://arxiv.org/html/2610.06838v2}.
\bibitem{p07:zrec:Sun}
Z.-W. Sun,
\emph{Catalan's constant is irrational},
arXiv:2609.04176v1, 3 September 2026.
\url{https://arxiv.org/html/2609.04176v1}.
This is the claim criticized in the next reference, not the later OpenAI proof.
\bibitem{p07:zrec:Wachs}
D. Wachs,
\emph{A note on a recent claimed proof of the irrationality of Catalan's
constant}, arXiv:2609.22339v1, 16 September 2026.
\url{https://arxiv.org/html/2609.22339v1}.
\bibitem{p07:zrec:OpenAI}
OpenAI,
\emph{Catalan's constant is irrational}, 24 September 2026,
Theorem 1.1 and Sections 1.2, 3--7; official paper and family 005 scope,
inspected 17 October 2026.
\url{https://github.com/openai/math/blob/main/preprints/Catalans-constant-is-irrational-September-24-2026/paper.pdf}.
Formalization scope:
\url{https://github.com/openai/math/blob/main/lean/docs/005.md}.
\bibitem{p07:zrec:Brown}
F. Brown,
\emph{Mellin transforms, transfinite diameter and rational approximations
of integrals}, arXiv:2604.20741v1, 22 April 2026, especially Criterion 1.3.
\url{https://arxiv.org/html/2604.20741v1}.
\bibitem{p07:zrec:DLMF}
NIST Digital Library of Mathematical Functions,
\emph{Classical orthogonal polynomials}, Section 18.3, Table 18.3.1,
normalizations for Jacobi, Legendre and Chebyshev polynomials;
accessed 17 October 2026.
\url{https://dlmf.nist.gov/18.3}.
\bibitem{p07:ric:BBC}
D.~H. Bailey, J.~M. Borwein and R.~E. Crandall,
On the Khintchine constant,
\emph{Mathematics of Computation} \textbf{66} (1997), 417--431.
\url{https://www.davidhbailey.com/dhbpapers/khinchine.pdf}.
\bibitem{p07:ric:AKS}
R.~L. Adler, M. Keane and M. Smorodinsky,
A construction of a normal number for the continued fraction transformation,
\emph{Journal of Number Theory} \textbf{13} (1981), 95--105.
\url{https://doi.org/10.1016/0022-314X(81)90031-7}.
\bibitem{p07:ric:Lanford}
O.~E. Lanford III,
A computer-assisted proof of the Feigenbaum conjectures,
\emph{Bulletin of the American Mathematical Society (N.S.)}
\textbf{6} (1982), 427--434.
\url{https://doi.org/10.1090/S0273-0979-1982-15008-X}.
\bibitem{p07:ric:Briggs}
K. Briggs, A precise calculation of the Feigenbaum constants,
\emph{Mathematics of Computation} \textbf{57} (1991), 435--439.
\url{https://doi.org/10.1090/S0025-5718-1991-1079009-6}.
\bibitem{p07:ric:KaramanosKotsireas}
K. Karamanos and I.~S. Kotsireas,
Statistical analysis of the first digits of the binary expansion of
Feigenbaum constants $\alpha$ and $\delta$,
\emph{Journal of the Franklin Institute} \textbf{342} (2005), 329--340.
\url{https://doi.org/10.1016/j.jfranklin.2004.11.004}.
\bibitem{p07:ric:Molteni}
A. Molteni,
An efficient method for the computation of the Feigenbaum constants to
high precision, preprint (2016), arXiv:1602.02357.
\url{https://arxiv.org/abs/1602.02357}.
\bibitem{p07:ric:HertlingSpandl}
P. Hertling and C. Spandl,
Computing a solution of Feigenbaum's functional equation in polynomial time,
\emph{Logical Methods in Computer Science} \textbf{10} (2014), no.~4,
article~7.
\url{https://arxiv.org/abs/1410.3277}.
\bibitem{p07:ric:Turing}
A.~M. Turing,
On computable numbers, with an application to the Entscheidungsproblem,
\emph{Proceedings of the London Mathematical Society (2)}
\textbf{42} (1936--1937), 230--265.
\url{https://doi.org/10.1112/plms/s2-42.1.230}.
\bibitem{p07:rq:Rstudent}
R Core Team, \emph{The Student t distribution}, R stats manual,
\url{https://stat.ethz.ch/R-manual/R-devel/RHOME/library/stats/html/TDist.html},
accessed 17 October 2026.
\bibitem{p08:cr:DhillonSaha}
S.~Dhillon and E.~Saha, \emph{On arithmetic nature of special values
of the incomplete beta function}, Acta Arithmetica \textbf{218}
(2025), 273--284, DOI 10.4064/aa240610-6-10.
Primary abstract and institutional record verified; exact overlap
with the inverse classification is not certified.
\url{https://doi.org/10.4064/aa240610-6-10}.
\url{https://digitalcommons.isical.ac.in/journal-articles/5486/}.
\bibitem{p08:cr:SahaTalk}
E.~Saha, \emph{Arithmetic nature of the special values of the
incomplete beta function}, joint work with S.~Dhillon,
speed talk at Recurrence, Transcendence, and Diophantine Approximation,
14 July 2025, slides 1--4, compiled PDF pp.~219--223;
method and parameter range on slide 4/4, compiled PDF p.~222.
\url{https://emresertoz.com/wp-content/uploads/2025/07/speed_talks.pdf}.
\bibitem{p08:cr:DLMF}
NIST Digital Library of Mathematical Functions,
\emph{Section 8.17: Incomplete beta functions},
Equations 8.17.1--8.17.4 and 8.17.20--8.17.21,
accessed 17 October 2026.
\url{https://dlmf.nist.gov/8.17}.
\bibitem{p08:cte:BlakajWolf}
V.~Blakaj and M.~M.~Wolf,
Transcendental properties of entropy-constrained sets,
\emph{Annales Henri Poincar\'e} \textbf{24} (2023), 349--362,
published online 27 August 2022, Theorem 2.
\url{https://doi.org/10.1007/s00023-022-01227-4}.
\bibitem{p08:cte:BSP}
H.~Boche, R.~F.~Schaefer and H.~V.~Poor,
Algorithmic computability and approximability of capacity-achieving
input distributions, arXiv:2202.12617, 2022, revised 2023,
Appendix B, Theorem 4. \url{https://arxiv.org/abs/2202.12617}.
\bibitem{p08:cte:Beukers}
F.~Beukers, A refined version of the
Siegel--Shidlovskii theorem, \emph{Annals of Mathematics}
\textbf{163} (2006), 369--379, Theorem 1.1.
\url{https://annals.math.princeton.edu/wp-content/uploads/annals-v163-n1-p08.pdf}.
\bibitem{p08:cte:Ostrowski}
A.~Ostrowski, Sur les relations alg\'ebriques
entre les int\'egrales ind\'efinies, \emph{Acta Mathematica}
\textbf{78} (1946), 315--318. \url{https://doi.org/10.1007/BF02421605}.
\bibitem{p08:cet:Wustholz}
G. W\"ustholz, Algebraische Punkte auf analytischen Untergruppen
algebraischer Gruppen, \emph{Ann. of Math.} (2) \textbf{129}
(1989), 501--517. English translation by W. Globke,
\url{https://wglobke.github.io/publications/Wuestholz1989.pdf}.
\bibitem{p08:cet:Chang}
C.-Y. Chang, Periods, logarithms and multiple zeta values,
\emph{Proceedings of the International Congress of Chinese
Mathematicians 2017}, International Press, 2020,
Theorem 1.2.1 (analytic subgroup theorem),
\url{https://www.math.nthu.edu.tw/~cychang/ICCM-2017.pdf}.
\bibitem{p08:rd:Schneider}
Th. Schneider, Arithmetische Untersuchungen elliptischer Integrale,
\emph{Mathematische Annalen} \textbf{113} (1937), 1--13.
\url{https://doi.org/10.1007/BF01571618}.
\bibitem{p08:rd:BertolinWaldschmidt}
C. Bertolin and M. Waldschmidt,
Variations on Schanuel's conjecture for elliptic and quasi-elliptic
functions I: the split case, \emph{Hardy--Ramanujan Journal}
\textbf{48} (2025), Part A, 10--29; arXiv:2504.14048,
Theorem~25(3).
\url{https://arxiv.org/html/2504.14048v1}.
\bibitem{p08:rd:Wustholz}
G. W\"ustholz,
Algebraische Punkte auf analytischen Untergruppen algebraischer Gruppen,
\emph{Annals of Mathematics (2)} \textbf{129} (1989), 501--517.
English translation of the original article:
\url{https://wglobke.github.io/publications/Wuestholz1989.pdf}.
\bibitem{p08:rd:Chang}
C.-Y. Chang, Periods, logarithms and multiple zeta values,
ICCM survey manuscript, Theorem~1.2.1.
\url{https://www.math.nthu.edu.tw/~cychang/ICCM-2017.pdf}.
\bibitem{p08:rd:Campbell}
J. Campbell, Automatic sequences and the Glaisher--Kinkelin constant,
\emph{Advances in Applied Mathematics} \textbf{158} (2024), 102721.
\url{https://doi.org/10.1016/j.aam.2024.102721}.
\bibitem{p08:rd:DLMF}
NIST Digital Library of Mathematical Functions, Section~25.4,
Reflection formulas, equation~25.4.2.
\url{https://dlmf.nist.gov/25.4.E2}.
\bibitem{p08:rd:BakerIII}
A. Baker, Linear forms in the logarithms of algebraic numbers (III),
\emph{Mathematika} \textbf{14} (1967), 220--228.
\url{https://doi.org/10.1112/S0025579300003843}.
\bibitem{p08:rd:PriorArc}
KEIO, \emph{Transcendence of Algebraic Arc Lengths and Inverse Periods:
Division Dichotomies, Beta Quantiles, and Arithmetic Transfers},
supplied manuscript, 12 October 2026; the elliptic-arc theorem and
algebraic-parameter beta theorem.
\bibitem{p08:ctm:Kubota}
K. K. Kubota, \emph{On the algebraic independence of holomorphic
solutions of certain functional equations and their values},
Mathematische Annalen \textbf{227} (1977), 9--50.
\url{https://doi.org/10.1007/BF01360961}.
\bibitem{p08:ctm:AF}
B. Adamczewski and C. Faverjon,
\emph{Mahler's method in several variables and finite automata},
Annals of Mathematics \textbf{204} (2026), 455--533.
\url{https://annals.math.princeton.edu/2026/204-2/p01}.
The theorem numbering quoted here is that of the author preprint
arXiv:2012.08283, Theorems 2.3 and 2.8, Corollary 2.5.
\url{https://arxiv.org/pdf/2012.08283}.
\bibitem{p08:ctm:AF18}
B. Adamczewski and C. Faverjon,
\emph{Mahler's method in several variables I: The theory of
regular singular systems}, author preprint, 2018,
Theorems 2.1 and 2.4.
\url{https://adamczewski.perso.math.cnrs.fr/RegularSingular1.pdf}.
\bibitem{p08:ctm:BV15}
P. Bundschuh and K. V\"a\"an\"anen,
\emph{Algebraic independence results on the generating Lambert
series of the powers of a fixed integer}, Hardy--Ramanujan Journal
\textbf{38} (2015), 36--44.
\url{https://hrj.episciences.org/1358/pdf}.
\bibitem{p08:ctm:BV14}
P. Bundschuh and K. V\"a\"an\"anen,
\emph{Algebraic independence of reciprocal sums of certain
Fibonacci-type numbers}, arXiv:1403.5510, 2014.
\url{https://arxiv.org/pdf/1403.5510}.
\bibitem{p08:ctm:Nguyen}
K. D. Nguyen, \emph{Transcendental series of reciprocals of Fibonacci
and Lucas numbers}, arXiv:2009.02446, 2020.
\url{https://arxiv.org/abs/2009.02446}.
\bibitem{p08:ctm:DT}
D. Duverney and Y. Tachiya,
\emph{Refinement of the Chowla--Erd\H{o}s method and linear independence
of certain Lambert series}, author manuscript, 2019, Remark 1.1.
\url{https://danielduverney.fr/documents/theorie-des-nombres/DuverneyTachiya190522.pdf}.
\bibitem{p08:ctm:Hirschhorn}
M. D. Hirschhorn and J. A. McGowan,
\emph{Algebraic consequences of Jacobi's two-square and four-square
theorems}, author manuscript, opening identity.
\url{https://web.maths.unsw.edu.au/~mikeh/webpapers/paper82.pdf}.
\bibitem{p08:ctm:Waldschmidt}
M. Waldschmidt, \emph{Transcendental Number Theory: recent results
and open problems}, January 2022, slides 58--59.
\url{https://webusers.imj-prg.fr/~michel.waldschmidt/articles/pdf/TNTOpenPbs.pdf}.
\bibitem{p08:rlap:WV}
T. Wolfs and W. Van Assche,
\emph{Rational approximation of Euler's constant using multiple orthogonal polynomials},
Combinatorics and Number Theory \textbf{14} (2025), 141--162.
\href{https://doi.org/10.2140/cnt.2025.14.141}{doi:10.2140/cnt.2025.14.141};
\href{https://arxiv.org/html/2404.09799v3}{primary full text}.
Theorems 1.2--1.3 concern improved approximants; the introduction also
records the classical Stieltjes--Laguerre Gompertz approximants.
\bibitem{p08:rlap:MY}
N. Murabayashi and H. Yoshida,
\emph{Rational function approximations of the special function
$e^xE_1(x)$ and applications to irrationality of Euler--Gompertz constant
$\delta$}, preprint, arXiv:2210.06768 (2022, subsequently revised).
\href{https://arxiv.org/abs/2210.06768}{arXiv:2210.06768}.
The final irrationality application requires an additional unproved
condition and is not used as a theorem of irrationality.
\bibitem{p08:dct:Fauzan}
A. Fauzan, \emph{$\zeta(5)$ is irrational}, 17 September 2026,
Zenodo 22826419, Sections 2.1--2.4 and Appendices C.2--C.3.
\url{https://zenodo.org/records/22826419}.
The primary PDF, including the functional, moment representation, relative
matrix estimate and coefficient-height bound, was inspected.
\bibitem{p08:dct:DLMF}
NIST Digital Library of Mathematical Functions,
Sections 25.6 and 25.11, equations 25.6.2 and 25.11.29
(Euler's even-zeta formula and Hermite's integral formula).
\url{https://dlmf.nist.gov/25.6#E2};
\url{https://dlmf.nist.gov/25.11#E29}.
\bibitem{p08:dct:OpenAI}
OpenAI, \emph{Catalan's constant is irrational}, 24 September 2026,
Section 2.1, especially equations (17) and (20).
\url{https://github.com/openai/math/blob/main/preprints/Catalans-constant-is-irrational-September-24-2026/paper.pdf}.
The local primary PDF and its exact moment recurrences were inspected.
\bibitem{p08:dct:Lai}
L. Lai, \emph{The linear independence of $1$, $\zeta(3)$, and $\zeta(4)$},
arXiv:2610.06838v2, 15 October 2026.
\url{https://arxiv.org/html/2610.06838v2}.
\bibitem{p08:dct:ZetaWindow}
gmDevi, \emph{At least one of $\zeta(7),\zeta(9),\ldots,\zeta(21)$ is
irrational}, Lean source repository and explicit theorem scope, 2026,
inspected 17 October 2026.
\url{https://github.com/gmDevi/zeta-7-21-lean}.
Its stated conclusion is a disjunction; no full Lean build was replayed here.
\bibitem{p08:dct:ZetaIssue}
huntrontrakkr, \emph{Anand's $\zeta(7)$ preprint: $I_{\mathrm{out}}$ in
Table 5/(122) does not follow from (59)}, primary calculation report,
3 October 2026; independent follow-up calculation dated 5 October 2026.
\url{https://github.com/gmDevi/zeta-7-21-lean/issues/1}.
The report concerns P. Anand, \emph{Zeta 7 is Irrational},
Zenodo 22920911, 23 September 2026,
\url{https://zenodo.org/records/22920911}.
\bibitem{p09:pfent:BW}
V.~Blakaj and M.~M.~Wolf, Transcendental properties
of entropy-constrained sets, \emph{Annales Henri Poincar\'e}
\textbf{24} (2023), 349--362, Theorem~2.
\url{https://doi.org/10.1007/s00023-022-01227-4}.
Author manuscript: \url{https://arxiv.org/abs/2111.10363}.
\bibitem{p09:ef:Chang}
C.-Y. Chang, Periods, logarithms and multiple zeta values,
\emph{Proceedings of the International Congress of Chinese
Mathematicians 2017}, International Press, 2020,
Theorem 1.2.1,
\url{https://www.math.nthu.edu.tw/~cychang/ICCM-2017.pdf}.
\bibitem{p09:ef:BCS}
A. Bourdon, P. L. Clark, and J. Stankewicz,
Torsion points on CM elliptic curves over real number fields,
\emph{Transactions of the American Mathematical Society}
\textbf{369} (2017), 8457--8496.
Preprint arXiv:1411.2742v3 (2015), Theorems 4.12 and 4.19,
\url{https://arxiv.org/html/1411.2742v3}.
\bibitem{p09:pg:Lagarias}
J.~C. Lagarias,
\emph{Euler's constant: Euler's work and modern developments},
Bull. Amer. Math. Soc. \textbf{50} (2013), 527--628;
\url{https://arxiv.org/abs/1303.1856}.
\bibitem{p09:pg:Bala}
P. Bala, congruence-preserving property recorded November 7, 2022,
in OEIS A002720, \emph{Number of partial permutations of an $n$-set};
\url{https://oeis.org/A002720/internal}.
\bibitem{p09:pg:Petojevic}
A. Petojevi\'c, S. Gordi\'c, M. Mandi\'c, and M. Gorjanac Ranitovi\'c,
\emph{New Equivalents of Kurepa's Hypothesis for Left Factorial},
Axioms \textbf{12} (2023), article 785;
\url{https://doi.org/10.3390/axioms12080785}.
\bibitem{p09:pg:Harvey}
D. Harvey,
\emph{The Accumulating Remainder Tree and Its Impact in Number Theory},
Proceedings of the International Congress of Mathematicians 2026,
Vol.~3, 414--433;
\url{https://doi.org/10.1137/25M1803139}.
\bibitem{p09:pz:Prevost2017}
M. Pr\'evost, \emph{Remainder Pad\'e approximants for the Hurwitz zeta
function}, arXiv:1709.05389 (2017); published in Results Math.\ \textbf{74}
(2019), Article 51. \url{https://arxiv.org/abs/1709.05389}.
\bibitem{p09:pz:PR2021}
M. Pr\'evost and T. Rivoal, \emph{Diagonal convergence of the remainder
Pad\'e approximants for the Hurwitz zeta function}, J. Number Theory
\textbf{222} (2021), 346--361.
\url{https://doi.org/10.1016/j.jnt.2020.10.019}.
\bibitem{p09:pz:DLMF}
NIST Digital Library of Mathematical Functions, Sections 18.25--18.26,
Wilson class orthogonal polynomials.
\url{https://dlmf.nist.gov/18.25}; \url{https://dlmf.nist.gov/18.26}.
\bibitem{p09:pz:Wilson1980}
J. A. Wilson, \emph{Some hypergeometric orthogonal polynomials}, SIAM J.
Math. Anal.\ \textbf{11} (1980), 690--701.
\url{https://doi.org/10.1137/0511064}.
\bibitem{p09:pfm:DLMF}
NIST Digital Library of Mathematical Functions,
Section 27.7, Lambert Series as Generating Functions,
\url{https://dlmf.nist.gov/27.7}.
\bibitem{p09:pfm:BS}
J. Bell and D. Smertnig,
\emph{Mahler series with multiplicative coefficient sequences},
preprint (2026), arXiv:2603.23456,
\url{https://arxiv.org/abs/2603.23456}.
\bibitem{p09:pfm:DT}
D. Duverney and Y. Tachiya,
\emph{Refinement of the Chowla--Erd\H{o}s method and linear
independence of certain Lambert series},
Forum Mathematicum \textbf{31} (2019), 1557--1566,
\url{https://doi.org/10.1515/forum-2018-0299}.
Author manuscript:
\url{https://danielduverney.fr/documents/theorie-des-nombres/DuverneyTachiya190522.pdf}.
\bibitem{p09:pfm:BakerI}
A. Baker,
\emph{Linear forms in the logarithms of algebraic numbers},
Mathematika \textbf{13} (1966), 204--216,
\url{https://doi.org/10.1112/S0025579300003971}.
\bibitem{p09:pfm:BakerIII}
A. Baker,
\emph{Linear forms in the logarithms of algebraic numbers (III)},
Mathematika \textbf{14} (1967), 220--228,
\url{https://doi.org/10.1112/S0025579300003843}.
\bibitem{p09:pfm:LZ}
A. Languasco and A. Zaccagnini,
\emph{Computing the Mertens and Meissel--Mertens constants for
sums over arithmetic progressions},
Experimental Mathematics \textbf{19} (2010), 279--284,
\url{https://doi.org/10.1080/10586458.2010.10390624}.
\bibitem{p09:gpf:WolfartWustholz}
J.~Wolfart and G.~W\"ustholz,
\emph{Der \"Uberlagerungsradius gewisser algebraischer Kurven und
die Werte der Betafunktion an rationalen Stellen},
Mathematische Annalen \textbf{273} (1985), 1--15.
\url{https://doi.org/10.1007/BF01455911}.
\bibitem{p09:gpf:ABP}
G.~W.~Anderson, W.~D.~Brownawell, and M.~A.~Papanikolas,
\emph{Determination of the algebraic relations among special
$\Gamma$-values in positive characteristic},
Annals of Mathematics \textbf{160} (2004), 237--313,
especially Sections 1.4.3--1.4.4 for the classical comparison.
\url{https://annals.math.princeton.edu/wp-content/uploads/annals-v160-n1-p06.pdf}.
\bibitem{p09:gpf:Waldschmidt}
M.~Waldschmidt, \emph{An introduction to irrationality and
transcendence methods}, Arizona Winter School 2008,
Lecture 3, Section 3.3.2, pp.~90--91.
\url{https://webusers.imj-prg.fr/~michel.waldschmidt/articles/pdf/AWSLecture3.pdf}.
\bibitem{p09:gpf:Sheth}
A.~Sheth, \emph{Real geometric transcendence for the Gamma function},
arXiv:2605.15117v1 (14 May 2026), discussion after Conjecture 4.2.
\url{https://arxiv.org/pdf/2605.15117}.
\bibitem{p09:gpf:DLMF}
NIST Digital Library of Mathematical Functions,
\emph{Sections 5.5: Functional relations, and 5.12: Beta function},
accessed 17 October 2026.
\url{https://dlmf.nist.gov/5.5}.
\url{https://dlmf.nist.gov/5.12}.
\bibitem{p10:inf:BW}
V.~Blakaj and M.~M.~Wolf, Transcendental properties
of entropy-constrained sets, \emph{Annales Henri Poincar\'e}
\textbf{24} (2023), 349--362.
\url{https://arxiv.org/abs/2111.10363}.
\bibitem{p10:inf:BM}
V.~Blakaj and C.~Manai, Transcendental properties
of entropy-constrained sets: Part II, \emph{J. Math. Phys.}
\textbf{65} (2024), 072203; arXiv:2302.09645.
\url{https://arxiv.org/abs/2302.09645}.
\bibitem{p10:entnext:BW}
V.~Blakaj and M.~M.~Wolf, Transcendental properties
of entropy-constrained sets, \emph{Annales Henri Poincar\'e}
\textbf{24} (2023), 349--362.
\url{https://doi.org/10.1007/s00023-022-01227-4}.
Author manuscript: \url{https://arxiv.org/html/2111.10363v2}.
\bibitem{p10:ng:Lagarias}
J. C. Lagarias, Euler's constant: Euler's work and modern developments,
\emph{Bull. Amer. Math. Soc.} \textbf{50} (2013), 527--628,
\href{https://arxiv.org/abs/1303.1856}{arXiv:1303.1856}.
\bibitem{p10:ng:MY}
N. Murabayashi and H. Yoshida,
Rational function approximations of the special function
$e^xE_1(x)$ and applications to irrationality of Euler--Gompertz
constant $\delta$, arXiv:2210.06768v2 (2024),
\url{https://arxiv.org/abs/2210.06768}.
\bibitem{p10:nz:Prevost}
M. Pr\'evost, \emph{Remainder Pad\'e approximants for the Hurwitz zeta
function}, Results Math. 74 (2019), article 51;
\url{https://arxiv.org/abs/1709.05389}.
\bibitem{p10:nz:PR}
M. Pr\'evost and T. Rivoal, \emph{Diagonal convergence of the remainder
Pad\'e approximants for the Hurwitz zeta function}, J. Number Theory
222 (2021), 346--361. \url{https://doi.org/10.1016/j.jnt.2020.10.019}.
\bibitem{p10:nz:Fauzan}
A. Fauzan, \emph{$\zeta(5)$ is irrational}, preprint, September 17, 2026.
\url{https://doi.org/10.5281/zenodo.22826419}.
\bibitem{p10:nz:Lai}
L. Lai, \emph{The linear independence of $1$, $\zeta(3)$, and $\zeta(4)$},
preprint, October 2026, revised version 2.
\url{https://arxiv.org/abs/2610.06838v2}.
\bibitem{p10:nz:OpenAICatalan}
OpenAI, \emph{Catalan's constant is irrational}, OpenAI Math Release,
September 24, 2026, with accompanying Lean 4 formalization.
\url{https://github.com/openai/math}.
\bibitem{p10:NLDorigoniKleinschmidt}
D. Dorigoni and A. Kleinschmidt,
Resurgent expansion of Lambert series and iterated Eisenstein integrals,
\emph{Commun. Number Theory Phys.} \textbf{15} (2021), 1--57.
\url{https://arxiv.org/abs/2001.11035}.
\bibitem{p10:NLBrentCoonsZudilin}
R. P. Brent, M. Coons and W. Zudilin,
Algebraic independence of Mahler functions via radial asymptotics,
\emph{Int. Math. Res. Not.} \textbf{2016} (2016), 571--603.
\url{https://arxiv.org/abs/1412.7906}.
\bibitem{p10:NLPupyrev}
Yu. A. Pupyrev,
Linear and algebraic independence of \(q\)-zeta values,
\emph{Math. Notes} \textbf{78} (2005), 563--568.
\url{https://doi.org/10.1007/s11006-005-0155-3}.
\bibitem{p10:NLBezivin}
J.-P. B\'ezivin,
Ind\'ependance lin\'eaire et alg\'ebrique de fonctions li\'ees
\`a la fonction \(q\)-dzeta,
\emph{Ann. Fac. Sci. Toulouse Math.} (6) \textbf{17} (2008), 23--36.
\url{https://doi.org/10.5802/afst.1173}.
\bibitem{p10:NLZagier}
D. Zagier, Elliptic modular forms and their applications,
in \emph{The 1-2-3 of Modular Forms}, Springer, 2008, pp. 1--103.
\url{https://people.mpim-bonn.mpg.de/zagier/files-restricted/doi/10.1007/978-3-540-74119-0/fulltext.pdf}.
\bibitem{p10:NLRosserSchoenfeld}
J. B. Rosser and L. Schoenfeld,
Approximate formulas for some functions of prime numbers,
\emph{Illinois J. Math.} \textbf{6} (1962), 64--94.
\url{https://doi.org/10.1215/ijm/1255631807}.
\bibitem{p10:NLBellSmertnig}
J. Bell and D. Smertnig,
Mahler series with multiplicative coefficient sequences,
preprint, 2026.
\url{https://arxiv.org/abs/2603.23456}.
\bibitem{p10:NLLanguascoZaccagnini}
A. Languasco and A. Zaccagnini,
Computing the Mertens and Meissel--Mertens constants for sums
over arithmetic progressions,
\emph{Experimental Mathematics} \textbf{19} (2010), 279--284.
\url{https://arxiv.org/abs/0906.2132}.
\bibitem{p10:gn:waldschmidt}
M.~Waldschmidt, \emph{Arizona Winter School lectures, Lecture 5:
Gamma values and periods}, 2008, especially Conjecture 5.25 and the
universal odd distribution formulation of Rohrlich's conjecture.
\url{https://webusers.imj-prg.fr/~michel.waldschmidt/articles/pdf/AWSLecture5.pdf}.
\bibitem{p10:gn:cohen-zudilin}
H.~Cohen and W.~Zudilin, \emph{Continued Fractions and Irrationality
Measures for Chowla--Selberg Gamma Quotients}, arXiv:2510.00215v2, 2025.
\url{https://arxiv.org/abs/2510.00215}.
\bibitem{p10:gn:sheth}
A.~Sheth, \emph{Real geometric transcendence for the Gamma function},
arXiv:2605.15117, 2026.
\url{https://arxiv.org/abs/2605.15117}.
\bibitem{p11:ipe:RV}
O.~Roy and M.~Vetterli, The effective rank: a
measure of effective dimensionality, \emph{Proc. EUSIPCO} (2007).
\url{https://www.eurasip.org/Proceedings/Eusipco/Eusipco2007/Papers/a5p-h05.pdf}.
\bibitem{p11:ipe:BGS}
S.~L.~Braunstein, S.~Ghosh and S.~Severini,
The Laplacian of a graph as a density matrix: a basic combinatorial
approach to separability of mixed states,
\emph{Annals of Combinatorics} \textbf{10} (2006), no.~3.
\url{https://arxiv.org/abs/quant-ph/0406165}.
\bibitem{p11:TLBaker}
A. Baker,
Linear forms in the logarithms of algebraic numbers, I and III,
\emph{Mathematika} \textbf{13} (1966), 204--216;
\textbf{14} (1967), 220--228.
\url{https://doi.org/10.1112/S0025579300003971}.
\bibitem{p11:zf:Fauzan}
A. Fauzan, \emph{$\zeta(5)$ is irrational}, preprint, September 17, 2026,
32 pp., especially Sections 2--4 and Appendix C.
\url{https://doi.org/10.5281/zenodo.22826419}.
\bibitem{p11:zf:Lai}
L. Lai, \emph{The linear independence of $1$, $\zeta(3)$, and $\zeta(4)$},
preprint, October 2026, version 2.
\url{https://arxiv.org/abs/2610.06838v2}.
\bibitem{p11:zf:Prevost}
M. Pr\'evost, \emph{Remainder Pad\'e approximants for the Hurwitz zeta
function}, Results Math. 74 (2019), article 51.
\url{https://arxiv.org/abs/1709.05389}.
\bibitem{p11:tg:DLMF}
NIST Digital Library of Mathematical Functions, Section 18.8,
Table 18.8.1, classical orthogonal-polynomial differential equations,
\url{https://dlmf.nist.gov/18.8}.
\bibitem{p11:gc:gross}
B.~H.~Gross, \emph{On the periods of abelian varieties}, 2018,
especially Sections 3--4.
\url{https://people.math.harvard.edu/~gross/preprints/cs.pdf}.
\bibitem{p11:gc:dlmf}
F.~W.~J.~Olver et al., eds., \emph{NIST Digital Library of Mathematical
Functions}, Section 25.15, Dirichlet $L$-functions.
\url{https://dlmf.nist.gov/25.15}.
\bibitem{p11:gc:cohen-zudilin}
H.~Cohen and W.~Zudilin, \emph{Continued Fractions and Irrationality
Measures for Chowla--Selberg Gamma Quotients}, arXiv:2510.00215, 2025.
\url{https://arxiv.org/abs/2510.00215}.
\bibitem{p11:gc:waldschmidt}
M.~Waldschmidt, \emph{Elliptic functions and transcendence}, survey,
2006, Sections 4.3 and 5.2.
\url{https://webusers.imj-prg.fr/~michel.waldschmidt/articles/pdf/SurveyTrdceEllipt2005.pdf}.
\bibitem{p12:icf:BM}
V.~Blakaj and C.~Manai, Transcendental properties
of entropy-constrained sets II, \emph{Journal of Mathematical Physics}
\textbf{65} (2024), 072203.
\url{https://arxiv.org/abs/2302.09645}.
\bibitem{p12:ef:Conrad}
K.~Conrad, \emph{The splitting field of $X^3-2$ over $\mathbb Q$},
expository notes, pp.~1--3,
\url{https://kconrad.math.uconn.edu/blurbs/gradnumthy/Qw2.pdf}
(accessed 17 October 2026).
\bibitem{p12:WW1985CF}
J.~Wolfart and G.~W\"ustholz,
\emph{Der \"Uberlagerungsradius gewisser algebraischer Kurven und die
Werte der Betafunktion an rationalen Stellen},
Math. Ann. \textbf{273} (1985), 1--15.
\url{https://doi.org/10.1007/BF01455911}.
\bibitem{p12:RivoalRoques2014CF}
T.~Rivoal and J.~Roques,
\emph{Hadamard products of algebraic functions},
J. Number Theory \textbf{145} (2014), 579--603.
Author manuscript, Theorems~4--5 and Section~3:
\url{https://rivoal.perso.math.cnrs.fr/articles/algcoef.pdf}.
\bibitem{p12:Schneider1941CF}
Th.~Schneider,
\emph{Zur Theorie der Abelschen Funktionen und Integrale},
J. Reine Angew. Math. \textbf{183} (1941), 110--128.
\bibitem{p12:fg:MY}
N. Murabayashi and H. Yoshida,
\emph{Rational function approximations of the special function
$e^xE_1(x)$ and applications to irrationality of Euler--Gompertz
constant $\delta$}, arXiv:2210.06768v2 (2024),
\url{https://arxiv.org/abs/2210.06768}.
\bibitem{p12:fg:MZ}
T. Matala-aho and W. Zudilin,
\emph{Euler's factorial series and global relations},
Journal of Number Theory \textbf{186} (2018), 202--210,
\url{https://arxiv.org/abs/1703.02633}.
\bibitem{p12:zc:Kleiber}
C. Kleiber, \emph{The generalized lognormal distribution and the
Stieltjes moment problem}, author manuscript, Section 4.
\url{https://arxiv.org/abs/1301.1277}.
\bibitem{p12:zc:Lopez}
M. L\'opez-Garc\'ia, \emph{Operators on the vanishing moments subspace
and Stieltjes classes for M-indeterminate probability distributions}, Arab Journal of Mathematical Sciences
\textbf{28} (2022), 229--242.
\url{https://doi.org/10.1108/AJMS-04-2021-0083}.
\bibitem{p12:zc:Prevost}
M. Pr\'evost, \emph{Remainder Pad\'e approximants for the Hurwitz zeta
function}, Results Math. \textbf{74} (2019), article 51.
\url{https://arxiv.org/abs/1709.05389}.
\bibitem{p12:lf:Lambrinoudis}
C. Lambrinoudis,
\emph{Nonquadraticity of Lambert Series Values},
self-published manuscript and Lean repository, October 2026.
\url{https://github.com/clambro/lambert-nonquadraticity}.
\bibitem{p12:lf:BundschuhVaananen}
P. Bundschuh and K. V\"a\"an\"anen,
Algebraic independence results on the generating Lambert series of the
powers of a fixed integer,
\emph{Hardy--Ramanujan Journal} \textbf{38} (2015).
\url{https://doi.org/10.46298/hrj.2015.1358}.
\bibitem{p12:lf:Duverney}
D. Duverney, K. Nishioka, K. Nishioka and I. Shiokawa,
Transcendence of Rogers--Ramanujan continued fraction and reciprocal
sums of Fibonacci numbers,
\emph{Proc. Japan Acad. Ser. A Math. Sci.} \textbf{73} (1997), 140--142.
\url{https://doi.org/10.3792/pjaa.73.140}.
\end{thebibliography}
\end{document}
